% Options for packages loaded elsewhere
\PassOptionsToPackage{unicode}{hyperref}
\PassOptionsToPackage{hyphens}{url}
\PassOptionsToPackage{dvipsnames,svgnames,x11names}{xcolor}
\documentclass[
  11pt,
  a4paper,
  openright]{book}
\usepackage{xcolor}
\usepackage[inner=1in,outer=0.75in,top=0.85in,bottom=1in]{geometry}
\usepackage{amsmath,amssymb}
\setcounter{secnumdepth}{5}
\usepackage{iftex}
\ifPDFTeX
  \usepackage[T1]{fontenc}
  \usepackage[utf8]{inputenc}
  \usepackage{textcomp} % provide euro and other symbols
\else % if luatex or xetex
  \usepackage{unicode-math} % this also loads fontspec
  \defaultfontfeatures{Scale=MatchLowercase}
  \defaultfontfeatures[\rmfamily]{Ligatures=TeX,Scale=1}
\fi
\usepackage{lmodern}
\ifPDFTeX\else
  % xetex/luatex font selection
\fi
% Use upquote if available, for straight quotes in verbatim environments
\IfFileExists{upquote.sty}{\usepackage{upquote}}{}
\IfFileExists{microtype.sty}{% use microtype if available
  \usepackage[]{microtype}
  \UseMicrotypeSet[protrusion]{basicmath} % disable protrusion for tt fonts
}{}
\makeatletter
\@ifundefined{KOMAClassName}{% if non-KOMA class
  \IfFileExists{parskip.sty}{%
    \usepackage{parskip}
  }{% else
    \setlength{\parindent}{0pt}
    \setlength{\parskip}{6pt plus 2pt minus 1pt}}
}{% if KOMA class
  \KOMAoptions{parskip=half}}
\makeatother
\usepackage{color}
\usepackage{fancyvrb}
\newcommand{\VerbBar}{|}
\newcommand{\VERB}{\Verb[commandchars=\\\{\}]}
\DefineVerbatimEnvironment{Highlighting}{Verbatim}{commandchars=\\\{\}}
% Add ',fontsize=\small' for more characters per line
\usepackage{framed}
\definecolor{shadecolor}{RGB}{35,38,41}
\newenvironment{Shaded}{\begin{snugshade}}{\end{snugshade}}
\newcommand{\AlertTok}[1]{\textcolor[rgb]{0.58,0.85,0.30}{\textbf{\colorbox[rgb]{0.30,0.12,0.14}{#1}}}}
\newcommand{\AnnotationTok}[1]{\textcolor[rgb]{0.25,0.50,0.35}{#1}}
\newcommand{\AttributeTok}[1]{\textcolor[rgb]{0.16,0.50,0.73}{#1}}
\newcommand{\BaseNTok}[1]{\textcolor[rgb]{0.96,0.45,0.00}{#1}}
\newcommand{\BuiltInTok}[1]{\textcolor[rgb]{0.50,0.55,0.55}{#1}}
\newcommand{\CharTok}[1]{\textcolor[rgb]{0.24,0.68,0.91}{#1}}
\newcommand{\CommentTok}[1]{\textcolor[rgb]{0.48,0.49,0.49}{#1}}
\newcommand{\CommentVarTok}[1]{\textcolor[rgb]{0.50,0.55,0.55}{#1}}
\newcommand{\ConstantTok}[1]{\textcolor[rgb]{0.15,0.68,0.68}{\textbf{#1}}}
\newcommand{\ControlFlowTok}[1]{\textcolor[rgb]{0.99,0.74,0.29}{\textbf{#1}}}
\newcommand{\DataTypeTok}[1]{\textcolor[rgb]{0.16,0.50,0.73}{#1}}
\newcommand{\DecValTok}[1]{\textcolor[rgb]{0.96,0.45,0.00}{#1}}
\newcommand{\DocumentationTok}[1]{\textcolor[rgb]{0.64,0.20,0.25}{#1}}
\newcommand{\ErrorTok}[1]{\textcolor[rgb]{0.85,0.27,0.33}{\underline{#1}}}
\newcommand{\ExtensionTok}[1]{\textcolor[rgb]{0.00,0.60,1.00}{\textbf{#1}}}
\newcommand{\FloatTok}[1]{\textcolor[rgb]{0.96,0.45,0.00}{#1}}
\newcommand{\FunctionTok}[1]{\textcolor[rgb]{0.56,0.27,0.68}{#1}}
\newcommand{\ImportTok}[1]{\textcolor[rgb]{0.15,0.68,0.38}{#1}}
\newcommand{\InformationTok}[1]{\textcolor[rgb]{0.77,0.36,0.00}{#1}}
\newcommand{\KeywordTok}[1]{\textcolor[rgb]{0.81,0.81,0.76}{\textbf{#1}}}
\newcommand{\NormalTok}[1]{\textcolor[rgb]{0.81,0.81,0.76}{#1}}
\newcommand{\OperatorTok}[1]{\textcolor[rgb]{0.81,0.81,0.76}{#1}}
\newcommand{\OtherTok}[1]{\textcolor[rgb]{0.15,0.68,0.38}{#1}}
\newcommand{\PreprocessorTok}[1]{\textcolor[rgb]{0.15,0.68,0.38}{#1}}
\newcommand{\RegionMarkerTok}[1]{\textcolor[rgb]{0.16,0.50,0.73}{\colorbox[rgb]{0.08,0.19,0.26}{#1}}}
\newcommand{\SpecialCharTok}[1]{\textcolor[rgb]{0.24,0.68,0.91}{#1}}
\newcommand{\SpecialStringTok}[1]{\textcolor[rgb]{0.85,0.27,0.33}{#1}}
\newcommand{\StringTok}[1]{\textcolor[rgb]{0.96,0.31,0.31}{#1}}
\newcommand{\VariableTok}[1]{\textcolor[rgb]{0.15,0.68,0.68}{#1}}
\newcommand{\VerbatimStringTok}[1]{\textcolor[rgb]{0.85,0.27,0.33}{#1}}
\newcommand{\WarningTok}[1]{\textcolor[rgb]{0.85,0.27,0.33}{#1}}
\usepackage{longtable,booktabs,array}
\newcounter{none} % for unnumbered tables
\usepackage{calc} % for calculating minipage widths
% Correct order of tables after \paragraph or \subparagraph
\usepackage{etoolbox}
\makeatletter
\patchcmd\longtable{\par}{\if@noskipsec\mbox{}\fi\par}{}{}
\makeatother
% Allow footnotes in longtable head/foot
\IfFileExists{footnotehyper.sty}{\usepackage{footnotehyper}}{\usepackage{footnote}}
\makesavenoteenv{longtable}
\usepackage{graphicx}
\makeatletter
\newsavebox\pandoc@box
\newcommand*\pandocbounded[1]{% scales image to fit in text height/width
  \sbox\pandoc@box{#1}%
  \Gscale@div\@tempa{\textheight}{\dimexpr\ht\pandoc@box+\dp\pandoc@box\relax}%
  \Gscale@div\@tempb{\linewidth}{\wd\pandoc@box}%
  \ifdim\@tempb\p@<\@tempa\p@\let\@tempa\@tempb\fi% select the smaller of both
  \ifdim\@tempa\p@<\p@\scalebox{\@tempa}{\usebox\pandoc@box}%
  \else\usebox{\pandoc@box}%
  \fi%
}
% Set default figure placement to htbp
\def\fps@figure{htbp}
\makeatother
\setlength{\emergencystretch}{3em} % prevent overfull lines
\providecommand{\tightlist}{%
  \setlength{\itemsep}{0pt}\setlength{\parskip}{0pt}}
% ═══════════════════════════════════════════════════════════════════════════════
% PDF Digital Signatures — Technical Reference
% LaTeX preamble for professional book typesetting
% ═══════════════════════════════════════════════════════════════════════════════

% ── Page geometry ────────────────────────────────────────────────────────────
% Geometry is set by pandoc via -V geometry="..." to avoid option clash.
% Do not load \usepackage{geometry} here.

% ── Typography ───────────────────────────────────────────────────────────────
\usepackage{fontspec}
\setmainfont{Noto Serif}[
    BoldFont={Noto Serif Bold},
    ItalicFont={Noto Serif Italic},
    BoldItalicFont={Noto Serif Bold Italic}
]
\setsansfont{Noto Sans}[
    BoldFont={Noto Sans Bold},
    ItalicFont={Noto Sans Italic}
]
\setmonofont{DejaVu Sans Mono}[
    Scale=0.85
]

\usepackage{microtype}          % Microtypographic improvements
\usepackage{setspace}
\setstretch{1.08}               % Slightly increased line spacing for readability

% ── Colors ───────────────────────────────────────────────────────────────────
\usepackage{xcolor}
\definecolor{accent}{HTML}{0f3460}
\definecolor{linkcolor}{HTML}{1a73e8}
\definecolor{codebackground}{HTML}{1e1e2e}
\definecolor{codetext}{HTML}{cdd6f4}
\definecolor{notecolor}{HTML}{0f3460}
\definecolor{warningcolor}{HTML}{e65100}
\definecolor{cautioncolor}{HTML}{b71c1c}
\definecolor{tipcolor}{HTML}{2e7d32}
\definecolor{importantcolor}{HTML}{6a1b9a}
\definecolor{sidebarcolor}{HTML}{f5f5f5}

% ── Headers and footers ─────────────────────────────────────────────────────
\usepackage{fancyhdr}
\pagestyle{fancy}
\fancyhf{}
\fancyhead[LE]{\small\textit{\leftmark}}
\fancyhead[RO]{\small\textit{\rightmark}}
\fancyfoot[C]{\small\thepage}
\renewcommand{\headrulewidth}{0.4pt}
\renewcommand{\footrulewidth}{0pt}

% ── Chapter and section styling ──────────────────────────────────────────────
\usepackage{titlesec}

\titleformat{\chapter}[display]
    {\normalfont\Huge\bfseries\sffamily\color{accent}}
    {\chaptertitlename\ \thechapter}
    {20pt}
    {\Huge}
\titlespacing*{\chapter}{0pt}{-30pt}{30pt}

\titleformat{\section}
    {\normalfont\Large\bfseries\sffamily\color{accent}}
    {\thesection}{1em}{}
\titleformat{\subsection}
    {\normalfont\large\bfseries\sffamily}
    {\thesubsection}{1em}{}
\titleformat{\subsubsection}
    {\normalfont\normalsize\bfseries\sffamily}
    {\thesubsubsection}{1em}{}

% ── Code listings ────────────────────────────────────────────────────────────
\usepackage{listings}
\lstdefinestyle{default}{
    backgroundcolor=\color{codebackground},
    basicstyle=\small\ttfamily\color{codetext},
    breakatwhitespace=false,
    breaklines=true,
    captionpos=b,
    keepspaces=true,
    showspaces=false,
    showstringspaces=false,
    showtabs=false,
    tabsize=4,
    frame=single,
    rulecolor=\color{codebackground},
    numbers=left,
    numberstyle=\tiny\color{gray},
    numbersep=8pt,
    xleftmargin=1.5em,
    framexleftmargin=1.5em,
    aboveskip=1em,
    belowskip=0.8em,
    % Language-specific keyword colors
    keywordstyle=\color[HTML]{89b4fa}\bfseries,
    stringstyle=\color[HTML]{a6e3a1},
    commentstyle=\color[HTML]{6c7086}\itshape,
}
\lstset{style=default}

% Java language definition
\lstdefinelanguage{Java21}{
    morekeywords={abstract,assert,boolean,break,byte,case,catch,char,class,
        const,continue,default,do,double,else,enum,extends,final,finally,
        float,for,goto,if,implements,import,instanceof,int,interface,long,
        native,new,package,private,protected,public,return,short,static,
        strictfp,super,switch,synchronized,this,throw,throws,transient,
        try,var,void,volatile,while,yield,record,sealed,permits,when},
    sensitive=true,
    morecomment=[l]{//},
    morecomment=[s]{/*}{*/},
    morestring=[b]",
    morestring=[b]',
}

% ASN.1 definition
\lstdefinelanguage{ASN1}{
    morekeywords={SEQUENCE,SET,CHOICE,OPTIONAL,IMPLICIT,EXPLICIT,
        OF,INTEGER,BOOLEAN,OCTET,STRING,BIT,NULL,OBJECT,IDENTIFIER,
        ANY,DEFINED,BY},
    sensitive=true,
    morecomment=[l]{--},
    morestring=[b]",
}

% ── Admonitions (NOTE, WARNING, etc.) ────────────────────────────────────────
\usepackage{tcolorbox}
\tcbuselibrary{breakable,skins}

\newtcolorbox{notebox}{
    colback=notecolor!5, colframe=notecolor!70,
    fonttitle=\bfseries\sffamily, title=Note,
    breakable, left=4pt, right=4pt, top=4pt, bottom=4pt,
    boxrule=0.5pt, arc=2pt
}
\newtcolorbox{warningbox}{
    colback=warningcolor!5, colframe=warningcolor!70,
    fonttitle=\bfseries\sffamily, title=Warning,
    breakable, left=4pt, right=4pt, top=4pt, bottom=4pt,
    boxrule=0.5pt, arc=2pt
}
\newtcolorbox{cautionbox}{
    colback=cautioncolor!5, colframe=cautioncolor!70,
    fonttitle=\bfseries\sffamily, title=Caution,
    breakable, left=4pt, right=4pt, top=4pt, bottom=4pt,
    boxrule=0.5pt, arc=2pt
}
\newtcolorbox{importantbox}{
    colback=importantcolor!5, colframe=importantcolor!70,
    fonttitle=\bfseries\sffamily, title=Important,
    breakable, left=4pt, right=4pt, top=4pt, bottom=4pt,
    boxrule=0.5pt, arc=2pt
}
\newtcolorbox{tipbox}{
    colback=tipcolor!5, colframe=tipcolor!70,
    fonttitle=\bfseries\sffamily, title=Tip,
    breakable, left=4pt, right=4pt, top=4pt, bottom=4pt,
    boxrule=0.5pt, arc=2pt
}
\newtcolorbox{sidebarbox}{
    colback=sidebarcolor, colframe=gray!40,
    fonttitle=\bfseries\sffamily,
    breakable, left=6pt, right=6pt, top=6pt, bottom=6pt,
    boxrule=0.3pt, arc=3pt
}

% ── Tables ───────────────────────────────────────────────────────────────────
\usepackage{booktabs}
\usepackage{longtable}
\usepackage{array}
\renewcommand{\arraystretch}{1.3}   % More vertical breathing room in tables

% ── Graphics & TikZ ──────────────────────────────────────────────────────────
\usepackage{graphicx}
\usepackage{float}
\usepackage{caption}

% TikZ — all diagrams are rendered natively in LaTeX
\usepackage{tikz}
\usetikzlibrary{
    arrows.meta,        % Arrow tips
    positioning,        % Relative node placement
    calc,               % Coordinate calculations
    decorations.pathreplacing, % Braces
    shapes.geometric,   % Diamond, ellipse, etc.
    shapes.multipart,   % Split nodes (for ASN.1 structures)
    fit,                % Fit node around others
    backgrounds,        % Background layers
    chains,             % Chain positioning
    matrix              % Matrix of nodes
}

% forest — for tree diagrams (AcroForm field tree, ASN.1, certificate chain)
\usepackage[edges]{forest}
\captionsetup{
    font={small,sf},
    labelfont={bf,color=accent},
    format=hang,
    margin=1cm
}

% ── Table of Contents ───────────────────────────────────────────────────────
% Widen the number fields so deep section numbers (e.g. 12.2.11) don't
% overlap with the title text.
\usepackage{tocloft}
\setlength{\cftsecnumwidth}{3.5em}      % section (e.g. 12.2)
\setlength{\cftsubsecnumwidth}{4.5em}   % subsection (e.g. 12.2.11)
\setlength{\cftsubsubsecnumwidth}{5.5em}% subsubsection (e.g. 12.2.11.1)

% ── Cross-references and links ───────────────────────────────────────────────
\usepackage{hyperref}
\hypersetup{
    colorlinks=true,
    linkcolor=accent,
    citecolor=accent,
    urlcolor=linkcolor,
    bookmarks=true,
    bookmarksnumbered=true,
    pdfauthor={PDFalyzer Studio Contributors},
    pdftitle={PDF Digital Signatures — A Technical Reference},
    pdfsubject={PDF, Digital Signatures, eIDAS, PAdES, CMS, X.509},
    pdfkeywords={PDF, digital signatures, PAdES, CAdES, eIDAS, X.509, CMS, PKCS7, PDFBox, BouncyCastle}
}

% ── Index ────────────────────────────────────────────────────────────────────
\usepackage{makeidx}
\makeindex

% ── Bibliography ─────────────────────────────────────────────────────────────
\usepackage[
    backend=biber,
    style=numeric,
    sorting=nyt,
    maxnames=3,
    minnames=1
]{biblatex}

% ── Misc ─────────────────────────────────────────────────────────────────────
\usepackage{enumitem}               % Better list formatting
\setlist{noitemsep,topsep=3pt}      % Compact lists

\usepackage{csquotes}               % Smart quotes
\usepackage{xurl}                   % Better URL line breaking

% ── Override pandoc syntax highlighting to plain monospace ─────────────────
% Pandoc generates \begin{Shaded}\begin{Highlighting}[]\end{Highlighting}\end{Shaded}
% with token commands like \NormalTok, \KeywordTok, etc.
% We redefine everything to plain \ttfamily output — no colors, no background box.
\usepackage{fancyvrb}
\DefineVerbatimEnvironment{Highlighting}{Verbatim}{
    commandchars=\\\{\},
    fontsize=\small,
    xleftmargin=1.5em,
    breaklines=true
}
\renewenvironment{Shaded}{}{}
% Redefine all pandoc token commands to plain text
\renewcommand{\AlertTok}[1]{#1}
\renewcommand{\AnnotationTok}[1]{#1}
\renewcommand{\AttributeTok}[1]{#1}
\renewcommand{\BaseNTok}[1]{#1}
\renewcommand{\BuiltInTok}[1]{#1}
\renewcommand{\CharTok}[1]{#1}
\renewcommand{\CommentTok}[1]{#1}
\renewcommand{\CommentVarTok}[1]{#1}
\renewcommand{\ConstantTok}[1]{#1}
\renewcommand{\ControlFlowTok}[1]{#1}
\renewcommand{\DataTypeTok}[1]{#1}
\renewcommand{\DecValTok}[1]{#1}
\renewcommand{\DocumentationTok}[1]{#1}
\renewcommand{\ErrorTok}[1]{#1}
\renewcommand{\ExtensionTok}[1]{#1}
\renewcommand{\FloatTok}[1]{#1}
\renewcommand{\FunctionTok}[1]{#1}
\renewcommand{\ImportTok}[1]{#1}
\renewcommand{\InformationTok}[1]{#1}
\renewcommand{\KeywordTok}[1]{#1}
\renewcommand{\NormalTok}[1]{#1}
\renewcommand{\OperatorTok}[1]{#1}
\renewcommand{\OtherTok}[1]{#1}
\renewcommand{\PreprocessorTok}[1]{#1}
\renewcommand{\RegionMarkerTok}[1]{#1}
\renewcommand{\SpecialCharTok}[1]{#1}
\renewcommand{\SpecialStringTok}[1]{#1}
\renewcommand{\StringTok}[1]{#1}
\renewcommand{\VariableTok}[1]{#1}
\renewcommand{\VerbatimStringTok}[1]{#1}
\renewcommand{\WarningTok}[1]{#1}
\usepackage{bookmark}
\IfFileExists{xurl.sty}{\usepackage{xurl}}{} % add URL line breaks if available
\urlstyle{same}
\hypersetup{
  pdftitle={PDF Digital Signatures --- A Technical Reference},
  pdfauthor={PDFalyzer Studio Contributors},
  colorlinks=true,
  linkcolor={accent},
  filecolor={Maroon},
  citecolor={Blue},
  urlcolor={linkcolor},
  pdfcreator={LaTeX via pandoc}}

\title{PDF Digital Signatures --- A Technical Reference}
\author{PDFalyzer Studio Contributors}
\date{2026-03-06}

\begin{document}
\frontmatter
\maketitle

{
\setcounter{tocdepth}{2}
\tableofcontents
}
\mainmatter
\section{Preface}\label{_preface}

PDF digital signatures are one of those technologies that seem
straightforward until you actually have to implement them. You read the
spec, write some code, run a test --- and the signature validates in
your tool but fails in Adobe Acrobat. Or it works in Acrobat but fails
the EU's validation service. Or it works everywhere on Monday and breaks
on Friday because a timestamp authority renewed its certificate.

If you've been there, this book is for you.

This is a technical reference for engineers, developers, and security
researchers who need to understand PDF digital signatures at the deepest
level --- from the raw bytes in the file to the legal framework that
gives those bytes legal standing across 27 EU member states. Every claim
is backed by precise references to ISO 32000-2:2020, the ETSI PAdES
standards, relevant RFCs, and the eIDAS regulation. Code examples use
Apache PDFBox 3.0.0 and BouncyCastle 1.80 on Java 21, because that's
what most of us actually use in production.

\subsection{Who This Book Is For}\label{_who_this_book_is_for}

You are a software engineer, a security researcher, or a technically
curious person who needs to \emph{understand} PDF signatures, not just
use them. Maybe you're building a signing service. Maybe you're
debugging why a customer's signature shows "invalid" in one viewer but
"valid" in another. Maybe you're the person who got handed a government
compliance checklist that says "PAdES-B-LTA" and nobody in your
organization knows what that means.

We assume you know how to program, that you've encountered public-key
cryptography before (even if the details are fuzzy), and that you're
comfortable reading data structures. We do \emph{not} assume you know
anything about PDF internals, CMS, or the European regulatory landscape.
Parts I and II build that foundation from scratch.

\subsection{What This Book Is Not}\label{_what_this_book_is_not}

This is not a tutorial. There are no "Hello World" signing examples in
Chapter 1. If you need to sign a PDF in the next 20 minutes, copy the
code from Chapter 28 and come back later to understand why it works.

This is not a product manual. We use PDFBox and BouncyCastle because
they're open-source, widely deployed, and well-documented. The concepts
apply equally to iText, OpenPDF, or any other library --- the PDF
specification doesn't care what language you write in.

This is not a legal guide. We explain the eIDAS regulation and its
technical requirements in detail, but we are engineers, not lawyers.
When legal precision matters --- and with qualified electronic
signatures, it always does --- consult someone with a law degree.

\subsection{How This Book Is
Organized}\label{_how_this_book_is_organized}

The book is structured in six parts, each building on the previous:

\textbf{Part I: Cryptographic Foundations} (Chapters 1--6) starts from
first principles. If you can explain the difference between RSA and
ECDSA, between a certificate and a key pair, between CMS and PKCS\#7 ---
skip ahead. If any of those distinctions are hazy, start here. We cover
just enough cryptography to make the rest of the book make sense, with a
focus on the specific algorithms and structures you'll encounter in PDF
signatures.

\textbf{Part II: PDF Internals} (Chapters 7--11) dissects the PDF file
format itself. You'll learn how objects, cross-reference tables, and
incremental updates work --- because you cannot understand PDF
signatures without understanding incremental saves. You'll also learn
about AcroForm fields, annotations, and coordinate systems, because
signature fields are form fields, signature appearances are annotations,
and the coordinate system \emph{will} trip you up.

\textbf{Part III: PDF Digital Signatures} (Chapters 12--18) is the heart
of the book. This is where we take apart the \texttt{/Sig} dictionary
key by key, trace the byte range byte by byte, and walk through the CMS
structure field by field. We cover approval and certification
signatures, visual appearances, validation logic, PAdES profiles, and
long-term validation. If you only read one part of this book, read this
one.

\textbf{Part IV: eIDAS \& Qualified Signatures} (Chapters 19--23) enters
the regulatory world. The European eIDAS regulation created a legal
framework where a digital signature can have the same legal standing as
a handwritten one --- but only if it meets very specific technical
requirements. We explain what those requirements are, how the EU Trusted
List works, what a Qualified Trust Service Provider does, and how remote
signing via the Cloud Signature Consortium API fits into the picture.

\textbf{Part V: Security, Incidents \& Outlook} (Chapters 24--27) covers
what can go wrong --- and what has gone wrong. The Ruhr University
Bochum research team published devastating attacks against PDF signature
validation in 2019--2021. We analyze every attack technically, explain
why it worked, and show how to defend against it. We also cover
certificate authority breaches, industry oddities, and the emerging EUDI
wallet.

\textbf{Part VI: Code \& Implementation} (Chapters 28--31) is where
theory meets production. Complete, tested Java code for signing, CMS
construction, verification, and analysis --- not snippets, but full
classes you can drop into a project. We also provide test cases and
reference PDFs for every scenario discussed in the book.

\subsection{Conventions Used in This
Book}\label{_conventions_used_in_this_book}

PDF dictionary keys are written in monospace with a leading slash:
\texttt{/SubFilter}, \texttt{/ByteRange}, \texttt{/Contents}. When we
show a full dictionary, it's presented as a code listing with callout
annotations explaining each key.

Spec references follow this pattern: "ISO 32000-2:2020 §12.8.1" means
ISO 32000-2:2020, section 12.8.1. RFC references: "RFC 5652 ---
Cryptographic Message Syntax (CMS) §5" means RFC 5652, section 5. We use
these abbreviations consistently throughout.

Code examples are Java 21 using PDFBox 3.0.0 and BouncyCastle 1.80.
Every code listing compiles and runs --- the source is in the book's
repository alongside test cases that verify the examples.

Admonition blocks carry specific meaning:

Additional context, background, or alternative approaches.

Practical advice from experience --- the kind of thing a senior
colleague would tell you over coffee.

Something you must understand to avoid a specific class of bugs.

A common mistake or misunderstanding that leads to broken signatures.

A security-relevant issue that could lead to signature bypasses or
compliance failures.

\subsection{A Note on the PDF
Specification}\label{_a_note_on_the_pdf_specification}

The PDF specification has a complicated history. Adobe published it as a
proprietary format in 1993, released it as ISO 32000-1 in 2008, and then
ISO 32000-2:2020 superseded it in 2020 with significant changes to the
digital signature sections.

Throughout this book, we reference ISO 32000-2:2020 as the authoritative
source. Where the behavior of older PDF versions matters (and for
digital signatures, it often does --- plenty of documents in the wild
still use PDF 1.7 features), we note the differences explicitly.

The PAdES standards (ETSI EN 319 142) add an additional layer of
requirements on top of ISO 32000-2. They don't replace the PDF spec ---
they constrain it. A PAdES-compliant signature is always a valid ISO
32000-2 signature, but not every valid ISO 32000-2 signature is
PAdES-compliant. This distinction matters enormously in the EU, and
increasingly, everywhere else.

\subsection{Three Decades of PDF: A
Timeline}\label{_three_decades_of_pdf_a_timeline}

PDF was born as a page description language for printing. Digital
signatures were an afterthought --- bolted on six years later in PDF
1.3, using a cryptographic framework (PKCS\#7) that predated PDF itself.
Understanding why PDF signatures work the way they do --- and why they
break the way they break --- requires understanding the history that
shaped them.

The following table traces every major milestone in the evolution of PDF
and, specifically, the digital signature infrastructure built on top of
it. Some of these are format changes. Some are regulatory events. Some
are security incidents that forced the entire industry to rethink
assumptions it had held for years. All of them left marks on the code
you will write or debug.

\begin{longtable}[]{@{}
  >{\raggedright\arraybackslash}p{(\linewidth - 4\tabcolsep) * \real{0.1250}}
  >{\raggedright\arraybackslash}p{(\linewidth - 4\tabcolsep) * \real{0.2500}}
  >{\raggedright\arraybackslash}p{(\linewidth - 4\tabcolsep) * \real{0.6250}}@{}}
\caption{The evolution of PDF and PDF digital signatures,
1993--2026}\tabularnewline
\toprule\noalign{}
\begin{minipage}[b]{\linewidth}\raggedright
Year
\end{minipage} & \begin{minipage}[b]{\linewidth}\raggedright
Event
\end{minipage} & \begin{minipage}[b]{\linewidth}\raggedright
What happened and why it matters
\end{minipage} \\
\midrule\noalign{}
\endfirsthead
\toprule\noalign{}
\begin{minipage}[b]{\linewidth}\raggedright
Year
\end{minipage} & \begin{minipage}[b]{\linewidth}\raggedright
Event
\end{minipage} & \begin{minipage}[b]{\linewidth}\raggedright
What happened and why it matters
\end{minipage} \\
\midrule\noalign{}
\endhead
\bottomrule\noalign{}
\endlastfoot
1993 & \textbf{PDF 1.0} released (Acrobat 1.0) & Adobe published the
Portable Document Format as a proprietary specification for reliable
cross-platform document rendering. No forms, no interactivity, no
signatures --- just pages, text, images, and fonts. The file structure
established here (header, body, cross-reference table, trailer) survives
unchanged to this day. Every PDF you sign in 2026 still starts with
\texttt{\%PDF-1.x} or \texttt{\%PDF-2.0} and ends with
\texttt{\%\%EOF}. \\
1996 & \textbf{PDF 1.2} (Acrobat 3.0) & Introduced \textbf{AcroForm} ---
interactive form fields that users could fill in. This matters for
signatures because signature fields \emph{are} form fields. The
\texttt{/AcroForm} dictionary, the field tree, widget annotations ---
all of the infrastructure that Chapter 9 dissects --- was born here. PDF
1.2 also added Unicode support and Flate compression, both of which are
ubiquitous in modern PDFs. \\
1999 & \textbf{PDF 1.3} --- digital signatures introduced (Acrobat 4.0)
& The moment PDF gained the ability to carry cryptographic signatures.
Adobe defined the \texttt{/Sig} dictionary type, the \texttt{/ByteRange}
array, the \texttt{/Contents} hex string for the signature value, and
the \texttt{/SubFilter} key with its initial value
\texttt{adbe.pkcs7.detached}. The design decision to embed the CMS
signature \emph{inside} the PDF file (rather than as a detached sidecar)
and to define the signed region as a byte range with a gap for the
signature value itself --- these were architectural choices made in 1999
that we still live with. Every quirk in Chapter 12 traces back to this
release. \\
2001 & \textbf{PDF 1.4} (Acrobat 5.0) & Added \textbf{transparency} (the
imaging model's biggest change ever), 128-bit RC4 encryption, and
\textbf{tagged PDF} for accessibility. For signatures: PDF 1.4
formalized the \texttt{adbe.pkcs7.sha1} SubFilter (where the SHA-1 hash
of the entire document is embedded as the CMS content, rather than using
a detached signature). This SubFilter is now deprecated but still
appears in legacy documents. Acrobat 5 also introduced
\textbf{certification signatures} (\texttt{/DocMDP} transform) --- the
distinction between "I approve this document" and "I certify this is the
authoritative version" starts here. \\
2003 & \textbf{PDF 1.5} (Acrobat 6.0) & Introduced \textbf{object
streams} and \textbf{cross-reference streams} --- a complete reworking
of how PDF files store their internal structure. Instead of the classic
ASCII cross-reference table, PDF 1.5 allows binary cross-reference data
compressed inside stream objects. This matters for signatures because
the byte range must cover the \emph{entire} file structure, and
cross-reference streams change where objects physically live in the
file. PDF 1.5 also added \textbf{XFA forms} (Adobe's XML-based form
format), which would become a persistent source of interoperability
problems for signatures on XFA documents --- a problem so annoying that
ISO 32000-2:2020 deprecated XFA entirely. \\
2005 & \textbf{PDF 1.6} (Acrobat 7.0) & \textbf{AES-128 encryption}
replaced the weak RC4-based encryption from earlier versions. Embedded
file streams (attachments) and 3D content were added. For signatures:
PDF 1.6 refined the interaction between encryption and signing --- a
signed, encrypted PDF requires careful handling to ensure the signature
covers the encrypted content correctly. The \texttt{/NeedAppearances}
flag behavior was also clarified, which matters because a missing
appearance stream on a signature field causes some viewers to show a
blank signature block. \\
2006 & \textbf{PDF 1.7} (Acrobat 8.0) & The last version of PDF that
Adobe defined unilaterally before handing the specification to ISO. PDF
1.7 solidified the \textbf{AES-256 encryption} scheme and finalized the
AcroForm/XFA coexistence model. The vast majority of signed PDFs in the
wild today are PDF 1.7 documents --- not because of any technical
limitation of PDF 2.0, but because template authors, form generators,
and enterprise software have not yet migrated. When you debug a
signature failure in production, you are almost certainly looking at a
PDF 1.7 file. \\
2008 & \textbf{ISO 32000-1:2008} --- PDF becomes an ISO standard & Adobe
donated the PDF 1.7 specification to ISO, making it an open
international standard. The technical content was essentially identical
to Adobe's PDF Reference 1.7, but the governance changed: future changes
would go through the ISO committee process, not Adobe's product team.
For signatures, this was the moment PDF digital signatures gained the
credibility of an ISO-backed specification --- a prerequisite for
government and regulated-industry adoption. \\
2009 & \textbf{PAdES} --- ETSI TS 102 778 (Parts 1--5) & The European
Telecommunications Standards Institute published the first \textbf{PDF
Advanced Electronic Signatures} standards. PAdES defined profiles for
how CMS signatures should be embedded in PDF to meet European legal
requirements: which signed attributes are mandatory (ESS
signing-certificate-v2), which SubFilter values to use
(\texttt{ETSI.CAdES.detached}), and how to embed timestamps, revocation
data, and long-term validation information. PAdES didn't replace ISO
32000 --- it constrained it. A PAdES-compliant signature is always a
valid ISO 32000 signature, but the reverse is not true. This distinction
--- and the interoperability problems it causes when a PDF viewer
validates against ISO 32000 rules but not PAdES rules --- is a recurring
theme in this book (Chapters 17 and 26). \\
2011 & \textbf{DigiNotar breach} --- CA infrastructure crisis & A Dutch
certificate authority, DigiNotar, was completely compromised by an
attacker who issued over 500 fraudulent certificates --- including
certificates for google.com, mozilla.org, and intelligence agencies. The
certificates were used in man-in-the-middle attacks against Iranian
dissidents. Every major browser vendor revoked trust in DigiNotar within
weeks; the company went bankrupt within months. For PDF signatures:
DigiNotar was a wake-up call about the fragility of the certificate
trust model that all digital signatures depend on. If a CA can be
compromised, any signature issued by that CA's certificates becomes
suspect. This incident accelerated the adoption of Certificate
Transparency logs and tighter CA audit requirements --- both of which
affected the trust infrastructure behind PDF signatures. Chapter 25
covers this and other CA breaches in detail. \\
2012 & \textbf{Adobe Approved Trust List (AATL)} formalized & Adobe
formalized its \textbf{AATL} program --- a curated list of root
certificates that Acrobat trusts for signature validation. If your
signing certificate chains up to an AATL root, Acrobat shows a green
checkmark. If it doesn't --- even if the certificate is perfectly valid,
issued by a reputable CA, and chains to a trust anchor in your OS trust
store --- Acrobat shows a warning. This created a de facto monopoly on
"valid" PDF signatures: to get a green checkmark in the world's dominant
PDF viewer, you need an AATL-listed certificate, which costs
significantly more than a standard code-signing certificate. The AATL
remains one of the most contentious aspects of the PDF signature
ecosystem (Chapter 26). \\
2014 & \textbf{eIDAS Regulation} (EU 910/2014) enacted & The European
Union passed the \textbf{Electronic Identification, Authentication and
Trust Services} regulation, creating a legal framework where digital
signatures can have the same legal standing as handwritten signatures
--- but only if they meet strict technical requirements. eIDAS defined
three tiers: \textbf{electronic signatures} (any electronic indication
of signing intent), \textbf{advanced electronic signatures} (uniquely
linked to the signatory, capable of identifying them, under their sole
control, linked to the data so any change is detectable), and
\textbf{qualified electronic signatures} (advanced + created by a
qualified device + based on a qualified certificate issued by a
qualified trust service provider). Only QES has automatic legal
equivalence to a handwritten signature across all EU member states.
PAdES became the primary technical mechanism for creating
eIDAS-compliant signatures in PDF. Chapters 19--23 cover the entire
eIDAS ecosystem. \\
2016 & \textbf{eIDAS fully effective} (July 1, 2016) & The eIDAS
regulation became fully effective across all EU member states, with the
\textbf{EU Trusted List (EUTL)} operational. The EUTL is a
machine-readable, cryptographically signed list of every Qualified Trust
Service Provider in the EU --- the authoritative source for determining
whether a signing certificate qualifies for the highest legal tier. For
developers: if your application needs to validate whether a PDF
signature is a qualified electronic signature, you must check the
signing certificate against the EUTL, not just against a generic trust
store. The EUTL XML schema, its signature validation, and the LOT (List
of Trusted Lists) hierarchy are covered in Chapter 20. \\
2017 & \textbf{ISO 32000-2:2017} --- PDF 2.0, first edition & The first
major revision of the PDF specification under ISO governance. PDF 2.0
deprecated XFA forms, removed several legacy features, and made
significant changes to the digital signature sections. Notably: the
\texttt{adbe.pkcs7.sha1} SubFilter was deprecated (SHA-1 is
cryptographically broken for collision resistance since the SHAttered
attack of February 2017), and the specification added the
\texttt{ETSI.RFC3161} SubFilter for document-level timestamps. PDF 2.0
also clarified the byte range coverage rules and tightened the
requirements around incremental updates after signing --- changes that
were partly motivated by early reports of the signature validation
problems that Bochum would systematically expose two years later. \\
2019 & \textbf{"1 Trillion Dollar Refund"} --- Bochum attacks (CCS 2019)
& The Ruhr University Bochum team published their landmark paper at ACM
CCS 2019, systematically breaking PDF signature validation in every
major viewer. They identified three attack classes: \textbf{Universal
Signature Forgery (USF)} --- crafting a signed PDF with a garbage or
empty signature value that viewers accepted as valid;
\textbf{Incremental Saving Attack (ISA)} --- appending unsigned changes
to a signed PDF that viewers failed to flag; and \textbf{Signature
Wrapping Attack (SWA)} --- relocating the signature's byte range to
point at attacker-controlled data. Of 22 desktop viewers tested, all
were vulnerable to at least one attack. The paper title referenced the
fact that signed PDFs are used for tax refund claims, insurance payouts,
and financial contracts --- if signatures could be forged, the financial
exposure was enormous. This research is the subject of Chapter 24. \\
2020 & \textbf{ISO 32000-2:2020} --- PDF 2.0, second edition & The
second edition of PDF 2.0 incorporated corrections and clarifications,
including tightened language around signature validation motivated by
the Bochum attacks. In the same year, Adobe issued
\textbf{CVE-2020-9592} and \textbf{CVE-2020-9596} and patched Acrobat
against USF, ISA, and SWA. Foxit, Nitro, LibreOffice, Okular, and every
other major viewer released their own patches. The PAdES standards were
also updated (ETSI EN 319 142-1 V1.1.1) to align with the second edition
of ISO 32000-2. \\
2021 & \textbf{Shadow Attacks} (NDSS 2021) & The Bochum team published a
follow-up: \textbf{Shadow Attacks}. Unlike the 2019 attacks (which
modified a PDF \emph{after} signing), Shadow Attacks prepared malicious
content \emph{before} the document was signed. The attacker created a
PDF with hidden content (invisible text, overlapping images, replaced
form fields) that the signer could not see during signing --- but that
became visible when the recipient opened the document. The signature
validated correctly because the signed bytes were genuinely signed by
the legitimate signer. The attack exploited the gap between what the
signer \emph{saw} and what the bytes \emph{contained} --- a fundamental
problem in any visual-document signing model. Three variants were
demonstrated: \textbf{Hide} (hiding content behind overlapping
elements), \textbf{Replace} (swapping visible content via interactive
form features), and \textbf{Hide-and-Replace} (combining both). Again,
all major viewers were vulnerable. \\
2022 & \textbf{ETSI EN 319 142-1 V1.2.1} --- PAdES update & ETSI updated
the PAdES standard to align with ISO 32000-2:2020 and to incorporate
lessons from the Bochum attacks. The update clarified requirements
around signature scope validation, incremental update handling, and the
relationship between the byte range and the visible document content. It
also formalized the PAdES-B-LTA profile's requirements for archival
timestamps in the DSS (Document Security Store) --- the mechanism that
enables long-term validation decades after the signing certificate has
expired. \\
2024 & \textbf{eIDAS 2.0} (EU 2024/1183) enacted & The EU adopted the
revised eIDAS regulation, which introduced the \textbf{European Digital
Identity (EUDI) Wallet} --- a smartphone-based identity wallet that
every EU member state must offer to its citizens by 2026. The EUDI
wallet will carry government-issued identity attestations (PID ---
Person Identification Data) and qualified electronic attestations of
attributes (QEAA), enabling citizens to sign documents with qualified
electronic signatures from their phone. This represents the most
ambitious attempt to make QES mainstream --- moving it from specialized
HSM-backed signing services to consumer smartphones. The Architecture
Reference Framework (ARF) defines the technical specifications. Chapter
27 covers the EUDI wallet and its impact on PDF signing. \\
2025--26 & \textbf{EUDI wallet pilots \& deployment} & EU member states
are conducting large-scale pilot programs for the EUDI wallet across
sectors: banking (identity verification), travel (digital travel
credentials), education (diplomas), healthcare (prescriptions), and
government services. The pilots test the wallet's signing capabilities,
including PAdES-B-LTA signature creation from mobile devices, remote
QSCD integration, and cross-border signature validation. The Cloud
Signature Consortium (CSC) API --- the standard interface for remote
signing --- is being adapted for wallet integration. The outcome of
these pilots will determine whether QES becomes a mainstream consumer
technology or remains a specialized enterprise tool. \\
\end{longtable}

The table omits minor point releases and focuses on the events that
changed what developers need to know. If you are working with PDF
signatures today, every row in that table has left its fingerprint on
the code you are writing --- whether it's the byte range gap from 1999,
the PAdES constraints from 2009, the trust model questions raised by
DigiNotar in 2011, or the validation rewrites forced by Bochum in
2019--2021.

\subsection{Acknowledgments}\label{_acknowledgments}

This book would not exist without the open-source communities that make
PDF signature work accessible to everyone. The \textbf{Apache PDFBox}
team --- particularly Andreas Lehmkühler, Tilman Hausherr, and the many
contributors who maintain the signing infrastructure --- have built the
most widely used open-source PDF library in the Java ecosystem. The
\textbf{Legion of the Bouncy Castle} (yes, that's their actual name)
gave us BouncyCastle, which provides the CMS and X.509 implementations
that most Java signing code depends on, whether it knows it or not.

Special acknowledgment goes to the security researchers at \textbf{Ruhr
University Bochum} --- Vladislav Mladenov, Christian Mainka, Karsten
Meyer zu Selhausen, Martin Grothe, and Jörg Schwenk --- whose systematic
analysis of PDF signature validation between 2019 and 2021 exposed
fundamental flaws in how every major PDF viewer handled signatures.
Their papers --- \emph{"1 Trillion Dollar Refund: How To Spoof PDF
Signatures"} (CCS 2019), \emph{"Shadow Attacks: Hiding and Replacing
Content in Signed PDFs"} (NDSS 2021), and several follow-ups ---
demonstrated that attackers could forge signatures, hide content changes
behind valid signatures, and bypass validation in Adobe Acrobat, Foxit
Reader, and dozens of other viewers. The industry patched, but the
research fundamentally changed how we think about PDF signature
security. Their work is cited throughout this book, and Chapter 24 is
dedicated to a technical analysis of their attacks.

What they found was this: PDF viewers did not actually verify what they
claimed to verify. The Bochum team identified three major attack
classes. \textbf{Universal Signature Forgery (USF)} showed that some
viewers only checked whether a signature \emph{existed} --- not whether
its cryptographic value was correct. An attacker could craft a PDF with
an empty or garbage \texttt{/Contents} value, and viewers like Adobe
Acrobat DC (before the patch) would display it as "signed and valid."
\textbf{Incremental Saving Attack (ISA)} exploited the incremental
update mechanism that PDF signatures depend on: an attacker could append
a new revision to a signed PDF --- changing text, adding pages,
replacing content entirely --- and because viewers did not properly
check whether the appended changes fell outside the signed byte range,
the signature still showed as valid. \textbf{Signature Wrapping Attack
(SWA)} relocated the signed content within the file so that the byte
range pointed at attacker-controlled data while the original signature
remained structurally present. The follow-up \textbf{Shadow Attacks}
paper (NDSS 2021) went further: an attacker could prepare a PDF
\emph{before} it was signed, embedding hidden content that would only
become visible after signing --- the signer would see one document, but
the recipient would see another, and the signature would validate on
both views.

The team tested 22 desktop PDF viewers and 7 online validation services.
Every single one was vulnerable to at least one attack class. Adobe
issued CVE-2020-9592 and CVE-2020-9596 and patched Acrobat in May 2020.
Foxit, Nitro, LibreOffice, and every other major viewer followed with
their own patches. The research forced a fundamental rethinking of PDF
signature validation --- in particular, how viewers handle incremental
updates after signing, how byte range coverage is verified, and how the
relationship between the CMS signature and the actual document bytes is
validated. Before Bochum, most viewers treated signature validation as a
checkbox. After Bochum, it became a security-critical code path that
required adversarial thinking.

The \textbf{European Telecommunications Standards Institute (ETSI)}
deserves credit for the PAdES standards that brought PDF signatures into
the regulated world. Whatever one may think about the complexity of ETSI
EN 319 142, it solved the real problem of making PDF signatures legally
interoperable across borders.

And the PDF specification itself is the product of hundreds of engineers
at Adobe, ISO, and ETSI who somehow managed to build a digital signature
framework into a page description language originally designed for
printing. That it works as well as it does is a testament to their
effort. That it has the quirks it has is a testament to the fact that
three decades of backward compatibility leaves marks.

\chapter{Cryptographic Foundations}\label{_cryptographic_foundations}

\chapter{Asymmetric Cryptography Primer}\label{ch01}

This chapter builds the cryptographic foundation that everything else in
this book depends on. We start with the key distribution problem that
asymmetric cryptography was invented to solve, work through RSA and
elliptic curve algorithms at a conceptual level, and end with practical
Java code that generates the key pairs you will use to sign PDFs.

If you can already explain why RSA-2048 and ECDSA with P-256 provide
roughly equivalent security, and you know why BouncyCastle exists
alongside the JDK's built-in providers, skip to Chapter 2. Everyone
else: this is where your journey into PDF signatures begins.

\section{The Key Distribution
Problem}\label{_the_key_distribution_problem}

Before we talk about RSA, elliptic curves, or any of the algorithms that
make PDF signatures possible, we need to understand the problem they
were invented to solve.

Symmetric cryptography --- where the same key encrypts and
decrypts --- has been around for millennia. Caesar ciphers, Enigma
machines, AES-256: they all share a single fundamental weakness. Both
parties must possess the same secret key. If Alice wants to send an
encrypted message to Bob, she and Bob must first agree on a shared
secret. In person, this is trivial --- write the key on a piece of paper
and hand it over. Over a network, it is a nightmare.

This is the \textbf{key distribution problem}: how do you securely share
a secret with someone you have never met, over a channel that an
adversary controls?

For most of human history, the answer was "you don't --- you use a
trusted courier." Governments ran key distribution networks. Military
organizations employed dedicated personnel whose sole job was
transporting codebooks. Banks used armored vehicles. None of this scales
to the internet, where you need to establish secure communications with
millions of strangers, instantly, without any prior arrangement.

In 1976, Whitfield Diffie and Martin Hellman published "New Directions
in Cryptography," describing a system where two parties could agree on a
shared secret over a public channel without ever having communicated
before. A year later, Ronald Rivest, Adi Shamir, and Leonard Adleman
published the RSA algorithm --- the first practical public-key
cryptosystem. (Clifford Cocks at GCHQ had independently discovered an
equivalent system in 1973, but it was classified until 1997.
Cryptography has a habit of being invented multiple times.)

The insight behind all of these systems is the same:
\textbf{mathematical one-way functions}. Operations that are easy to
perform in one direction but computationally infeasible to reverse.
Multiplying two large primes takes microseconds; factoring their product
takes longer than the age of the universe. This asymmetry --- easy
forward, hard backward --- is the foundation of everything that follows.

\section{Key Pairs: The Fundamental
Asymmetry}\label{_key_pairs_the_fundamental_asymmetry}

An asymmetric key pair consists of two mathematically related keys:

\begin{itemize}
\item
  A \textbf{private key} --- kept secret, never shared, never
  transmitted
\item
  A \textbf{public key} --- derived from the private key, freely
  distributable
\end{itemize}

The relationship between them has two critical properties:

\begin{enumerate}
\def\labelenumi{\arabic{enumi}.}
\item
  \textbf{It is computationally infeasible to derive the private key
  from the public key.} Given a public key, there is no known efficient
  algorithm to compute the corresponding private key. The security of
  the entire system rests on this.
\item
  \textbf{Operations performed with one key can only be reversed with
  the other.} Data encrypted with the public key can only be decrypted
  with the private key. Data signed with the private key can only be
  verified with the public key.
\end{enumerate}

This second property is what makes digital signatures possible. When you
sign a PDF, your private key produces a cryptographic value that is
mathematically bound to both your key and the document content. Anyone
who has your public key --- which is embedded in your certificate, which
is embedded in the PDF itself --- can verify that the signature was
produced by someone holding the corresponding private key.

They cannot forge a new signature. They cannot modify the document
without invalidating the signature. They cannot extract your private key
from the signature or the public key.

This is the fundamental asymmetry that gives asymmetric cryptography its
name, and it is the reason PDF digital signatures work at all.

The terms "encrypt" and "decrypt" are technically imprecise when applied
to digital signatures. Signing is not encryption in the classical
sense --- you are not hiding data; you are producing a proof. The
distinction matters in formal cryptographic contexts and in legal
frameworks like eIDAS (Regulation (EU) No 910/2014 (eIDAS)).

Throughout this book, we use "sign" and "verify" for signature
operations, reserving "encrypt" and "decrypt" for actual confidentiality
operations. Many older references (and some library APIs, including
parts of the JCA) conflate the two. Do not let the terminology confuse
you.

\section{RSA: The Workhorse}\label{_rsa_the_workhorse}

RSA, named after Rivest, Shamir, and Adleman, remains the most widely
deployed asymmetric algorithm in the world. Every PDF viewer supports
it. Every certificate authority issues RSA certificates. Every signing
library implements it. If you encounter a signed PDF in production,
there is an overwhelming probability that it was signed with RSA.

\subsection{The Mathematical Basis}\label{_the_mathematical_basis}

The security of RSA rests on the \textbf{integer factorization problem}:
given the product of two large primes, recover the original primes.
Multiplication is trivial; factoring is believed to be computationally
infeasible for sufficiently large numbers. Here is how this translates
into a key pair.

\begin{enumerate}
\def\labelenumi{\arabic{enumi}.}
\item
  \textbf{Choose two large primes, \emph{p} and \emph{q}.} These primes
  are generated independently using a cryptographically secure random
  number generator. For RSA-2048, each prime is approximately 1024 bits
  long. Finding primes of this size is itself nontrivial --- it requires
  probabilistic primality tests (Miller-Rabin is the standard), and
  several candidates may be tested before a probable prime is found.
  This is why RSA key generation is noticeably slower than ECDSA key
  generation.
\item
  \textbf{Compute the modulus \emph{n = p * q}.} This product \emph{n}
  appears in both the public and private keys. For RSA-2048, \emph{n} is
  a 2048-bit number. Computing \emph{n} from \emph{p} and \emph{q} takes
  microseconds. Recovering \emph{p} and \emph{q} from \emph{n} is the
  hard problem.
\item
  \textbf{Compute the totient.} Euler's totient function for \emph{n} is
  \emph{phi(n) = (p - 1)(q - 1)}. This value is used to derive the
  private exponent and is then discarded. It must never be stored or
  revealed --- knowing \emph{phi(n)} is equivalent to knowing the
  factorization of \emph{n}.
\item
  \textbf{Choose the public exponent \emph{e}.} In practice, \emph{e} is
  almost always 65537 (0x10001). This value was chosen because it is
  prime, it has a low Hamming weight (only two bits set), and it makes
  modular exponentiation during verification fast. You will see 65537 in
  virtually every RSA certificate you examine. Some ancient systems used
  \emph{e} = 3, which is faster but has known vulnerabilities when
  combined with naive padding.
\item
  \textbf{Compute the private exponent \emph{d}.} \emph{d} is the
  modular inverse of \emph{e} modulo \emph{phi(n)}: \emph{d * e = 1 mod
  phi(n)}. This is computed using the extended Euclidean algorithm.
  Knowing \emph{d} requires knowing \emph{phi(n)}, which requires
  knowing \emph{p} and \emph{q}.
\item
  \textbf{The public key} is the pair (\emph{n}, \emph{e}).
\item
  \textbf{The private key} is the pair (\emph{n}, \emph{d}). In
  practice, the private key also stores \emph{p}, \emph{q}, and several
  precomputed values (the CRT --- Chinese Remainder
  Theorem --- parameters) that make signing four times faster.
\end{enumerate}

\subsection{Why Large Primes Matter}\label{_why_large_primes_matter}

The security of RSA depends on the inability of an attacker to factor
\emph{n} and recover \emph{p} and \emph{q}. The difficulty of factoring
grows super-polynomially with the size of \emph{n}. The best known
classical factoring algorithm is the General Number Field Sieve (GNFS),
whose complexity is roughly:

\begin{verbatim}
exp( (1.923 + o(1)) * (ln n)^(1/3) * (ln ln n)^(2/3) )
\end{verbatim}

This means that doubling the key size does \emph{not} merely double the
work --- it increases it by many orders of magnitude. Factoring a
768-bit RSA modulus was achieved in 2009 by a team of researchers using
hundreds of machines over two years. Factoring a 1024-bit modulus is
estimated to be within the reach of a well-funded nation-state.
Factoring a 2048-bit modulus is estimated to require more computation
than humanity can currently muster.

\subsection{RSA Key Sizes vs Security
Levels}\label{_rsa_key_sizes_vs_security_levels}

The size of an RSA key is the bit length of the modulus \emph{n}. Larger
modulus means harder factoring means more security --- but also slower
operations.

\begin{longtable}[]{@{}
  >{\raggedright\arraybackslash}p{(\linewidth - 8\tabcolsep) * \real{0.1111}}
  >{\raggedright\arraybackslash}p{(\linewidth - 8\tabcolsep) * \real{0.1111}}
  >{\raggedright\arraybackslash}p{(\linewidth - 8\tabcolsep) * \real{0.2222}}
  >{\raggedright\arraybackslash}p{(\linewidth - 8\tabcolsep) * \real{0.2222}}
  >{\raggedright\arraybackslash}p{(\linewidth - 8\tabcolsep) * \real{0.3333}}@{}}
\caption{RSA key sizes and their practical
characteristics}\tabularnewline
\toprule\noalign{}
\begin{minipage}[b]{\linewidth}\raggedright
Key Size
\end{minipage} & \begin{minipage}[b]{\linewidth}\raggedright
Security Level
\end{minipage} & \begin{minipage}[b]{\linewidth}\raggedright
Sign Time (relative)
\end{minipage} & \begin{minipage}[b]{\linewidth}\raggedright
Signature Size
\end{minipage} & \begin{minipage}[b]{\linewidth}\raggedright
Status
\end{minipage} \\
\midrule\noalign{}
\endfirsthead
\toprule\noalign{}
\begin{minipage}[b]{\linewidth}\raggedright
Key Size
\end{minipage} & \begin{minipage}[b]{\linewidth}\raggedright
Security Level
\end{minipage} & \begin{minipage}[b]{\linewidth}\raggedright
Sign Time (relative)
\end{minipage} & \begin{minipage}[b]{\linewidth}\raggedright
Signature Size
\end{minipage} & \begin{minipage}[b]{\linewidth}\raggedright
Status
\end{minipage} \\
\midrule\noalign{}
\endhead
\bottomrule\noalign{}
\endlastfoot
1024-bit & \textasciitilde80-bit & 1x (baseline) & 128 bytes &
\textbf{Broken.} NIST disallowed after 2013. Do not use for any purpose.
You will still encounter these in legacy PDFs --- they were common
before 2010. \\
2048-bit & \textasciitilde112-bit & \textasciitilde4x & 256 bytes &
\textbf{Current minimum.} NIST SP 800-57 minimum through 2030. This is
what most CAs issue today. Adobe AATL requires at least 2048-bit RSA. \\
3072-bit & \textasciitilde128-bit & \textasciitilde9x & 384 bytes &
\textbf{Recommended for new deployments.} Provides 128-bit
security --- the level NIST considers adequate through at least 2030 and
likely well beyond. \\
4096-bit & \textasciitilde140-bit & \textasciitilde16x & 512 bytes &
\textbf{Conservative choice.} Noticeable performance impact for
high-volume signing. Some government requirements mandate this. The
additional security margin over 3072 is modest relative to the
performance cost. \\
\end{longtable}

RSA-1024 is not merely "weak" --- it is \textbf{factored}. In 2010, a
team factored a 768-bit RSA modulus. The computational cost of factoring
1024-bit moduli is within reach of well-funded adversaries. If you
encounter a signed PDF with an RSA-1024 certificate, the signature
provides no meaningful security guarantee, even if the viewer shows it
as "valid." Many viewers do not check key size during validation.

\subsection{Performance Impact of Key
Size}\label{_performance_impact_of_key_size}

RSA has an unusual asymmetry in its operation costs:
\textbf{verification is fast; signing is slow}.

This is because the public exponent \emph{e} is small (65537), making
the modular exponentiation during verification much cheaper than the
exponentiation with the large private exponent \emph{d} during signing.
For PDF signatures, this works in our favor: documents are signed once
but verified many times.

However, RSA's performance degrades sharply as key sizes increase. The
signing operation scales roughly as O(k\textsuperscript{3}) where
\emph{k} is the key size in bits. Doubling the key size makes signing
approximately eight times slower. For a desktop application signing a
single document, this is irrelevant. For a server processing thousands
of signing requests per second, the difference between 2048-bit and
4096-bit RSA is significant.

\subsection{Why RSA-1024 Persists in the
Wild}\label{_why_rsa_1024_persists_in_the_wild}

If RSA-1024 is broken, why do you still encounter it in signed PDFs? The
answer is a combination of institutional inertia, cost, and the peculiar
nature of long-lived documents.

Government agencies issued RSA-1024 certificates throughout the 2000s
for document signing infrastructure. Many of those signed documents
still exist, still circulate, and still need to be validated. Some
organizations operate signing infrastructure on hardware security
modules (HSMs) that cannot be easily re-keyed --- the cost of replacing
an HSM fleet across dozens of offices can be enormous. Other systems
operate in air-gapped environments where updating cryptographic
infrastructure requires physical site visits.

The practical consequence: your validation code must \emph{handle}
RSA-1024 gracefully, even if it correctly flags it as insecure. Refusing
to parse a legacy document because the key size is too small is
unhelpful. Displaying it as "valid" without qualification is
irresponsible. The right approach is what PDFalyzer does: show the
signature, show the key size, and let the user draw the appropriate
conclusion.

\textbf{How to check the key size of a certificate in a signed PDF:}

In PDFalyzer, open the Signatures tab and look at the "Public Key"
section of any signature. The key size is displayed directly.

From the command line, extract the certificate and use OpenSSL:

\begin{Shaded}
\begin{Highlighting}[]
\ExtensionTok{openssl}\NormalTok{ x509 }\AttributeTok{{-}in}\NormalTok{ cert.pem }\AttributeTok{{-}text} \AttributeTok{{-}noout} \KeywordTok{|} \FunctionTok{grep} \StringTok{"Public{-}Key"}
\end{Highlighting}
\end{Shaded}

In Java with BouncyCastle, it is one line:

\begin{Shaded}
\begin{Highlighting}[]
\DataTypeTok{int}\NormalTok{ keySize }\OperatorTok{=} \OperatorTok{((}\BuiltInTok{RSAPublicKey}\OperatorTok{)}\NormalTok{ cert}\OperatorTok{.}\FunctionTok{getPublicKey}\OperatorTok{()).}\FunctionTok{getModulus}\OperatorTok{().}\FunctionTok{bitLength}\OperatorTok{();}
\end{Highlighting}
\end{Shaded}

If you see 1024 or less, the certificate was issued before modern
standards were enforced, and the signature's cryptographic value is nil.

\section{Elliptic Curve Cryptography: Smaller Keys, Same
Security}\label{_elliptic_curve_cryptography_smaller_keys_same_security}

In the mid-1980s, Neal Koblitz and Victor Miller independently proposed
using elliptic curves over finite fields as the basis for public-key
cryptography. The core idea: the \textbf{elliptic curve discrete
logarithm problem (ECDLP)} is believed to be harder than integer
factorization for equivalent key sizes. This means we can achieve the
same security level with dramatically smaller keys.

\subsection{Elliptic Curve Basics}\label{_elliptic_curve_basics}

An elliptic curve, for cryptographic purposes, is the set of points
\emph{(x, y)} satisfying an equation of the form
\emph{y\textsuperscript{2} = x\textsuperscript{3} + ax + b} over a
finite field. The "finite field" part is important: we are not working
with real numbers and smooth curves. We are working with integers modulo
a large prime \emph{p}, and the curve is a discrete set of points.

The cryptographic operation on these points is \textbf{point addition}:
given two points on the curve, there is a geometric rule for computing a
third point. Repeated point addition (scalar multiplication) gives us
the one-way function: given a point \emph{G} and a scalar \emph{d},
computing \emph{Q = d * G} is efficient. Recovering \emph{d} from
\emph{Q} and \emph{G} (the ECDLP) is computationally infeasible.

You do not need to understand the geometry to use ECDSA effectively.
What matters is the practical consequence: the ECDLP is harder than
integer factorization for the same key size, which is why a 256-bit
elliptic curve key provides the same security as a 3072-bit RSA key.

\subsection{The NIST Curves: P-256, P-384,
P-521}\label{_the_nist_curves_p_256_p_384_p_521}

The standard curves for ECDSA, defined by NIST in FIPS 186-4:

\begin{longtable}[]{@{}
  >{\raggedright\arraybackslash}p{(\linewidth - 6\tabcolsep) * \real{0.1667}}
  >{\raggedright\arraybackslash}p{(\linewidth - 6\tabcolsep) * \real{0.1667}}
  >{\raggedright\arraybackslash}p{(\linewidth - 6\tabcolsep) * \real{0.1667}}
  >{\raggedright\arraybackslash}p{(\linewidth - 6\tabcolsep) * \real{0.5000}}@{}}
\caption{NIST standard curves for ECDSA}\tabularnewline
\toprule\noalign{}
\begin{minipage}[b]{\linewidth}\raggedright
Curve
\end{minipage} & \begin{minipage}[b]{\linewidth}\raggedright
Security Level
\end{minipage} & \begin{minipage}[b]{\linewidth}\raggedright
Signature Size
\end{minipage} & \begin{minipage}[b]{\linewidth}\raggedright
Notes
\end{minipage} \\
\midrule\noalign{}
\endfirsthead
\toprule\noalign{}
\begin{minipage}[b]{\linewidth}\raggedright
Curve
\end{minipage} & \begin{minipage}[b]{\linewidth}\raggedright
Security Level
\end{minipage} & \begin{minipage}[b]{\linewidth}\raggedright
Signature Size
\end{minipage} & \begin{minipage}[b]{\linewidth}\raggedright
Notes
\end{minipage} \\
\midrule\noalign{}
\endhead
\bottomrule\noalign{}
\endlastfoot
P-256 (secp256r1) & 128-bit & 64 bytes & The workhorse. Universally
supported. Adobe Acrobat, PDFBox, iText, BouncyCastle --- everyone
supports P-256. If you are choosing an ECDSA curve for PDF signing, this
is the default choice. \\
P-384 (secp384r1) & 192-bit & 96 bytes & Required by certain government
profiles (e.g., NSA Suite B, CNSS Policy 15). Broader support than P-521
but less than P-256. \\
P-521 (secp521r1) & 256-bit & 132 bytes & Maximum security. The name is
not a typo --- it is 521 bits, not 512. Rarely used in practice for PDF
signatures because the security margin over P-384 is excessive for most
threat models, and some tools have spotty support. \\
\end{longtable}

The curve names are a source of endless confusion. "P-256" is the NIST
name. "secp256r1" is the SECG (Standards for Efficient Cryptography
Group) name. "prime256v1" is the OpenSSL name. They are all the same
curve. You will encounter all three names in certificate metadata,
library APIs, and documentation. When PDFBox reports a signature
algorithm, it uses the JCA name; when BouncyCastle constructs the CMS
structure, it uses OIDs; when you read the certificate in OpenSSL, it
says "prime256v1." They are all P-256.

\subsection{Why ECDSA Signatures Are
Smaller}\label{_why_ecdsa_signatures_are_smaller}

The practical consequence of the ECDLP's difficulty is striking.
RSA-3072 provides the same security as ECC-256, but the RSA key is
twelve times larger. This size difference cascades through every part of
the system: smaller keys mean smaller certificates, smaller CMS
structures in the PDF's \texttt{/Contents} field, faster network
transmission of signed documents, and less storage overhead.

\begin{longtable}[]{@{}
  >{\raggedright\arraybackslash}p{(\linewidth - 6\tabcolsep) * \real{0.2000}}
  >{\raggedright\arraybackslash}p{(\linewidth - 6\tabcolsep) * \real{0.2000}}
  >{\raggedright\arraybackslash}p{(\linewidth - 6\tabcolsep) * \real{0.2000}}
  >{\raggedright\arraybackslash}p{(\linewidth - 6\tabcolsep) * \real{0.4000}}@{}}
\caption{Equivalent security levels across algorithms (NIST SP 800-57
Part 1, Rev. 5)}\tabularnewline
\toprule\noalign{}
\begin{minipage}[b]{\linewidth}\raggedright
Security Level
\end{minipage} & \begin{minipage}[b]{\linewidth}\raggedright
RSA Key Size
\end{minipage} & \begin{minipage}[b]{\linewidth}\raggedright
ECC Key Size
\end{minipage} & \begin{minipage}[b]{\linewidth}\raggedright
Ratio
\end{minipage} \\
\midrule\noalign{}
\endfirsthead
\toprule\noalign{}
\begin{minipage}[b]{\linewidth}\raggedright
Security Level
\end{minipage} & \begin{minipage}[b]{\linewidth}\raggedright
RSA Key Size
\end{minipage} & \begin{minipage}[b]{\linewidth}\raggedright
ECC Key Size
\end{minipage} & \begin{minipage}[b]{\linewidth}\raggedright
Ratio
\end{minipage} \\
\midrule\noalign{}
\endhead
\bottomrule\noalign{}
\endlastfoot
112-bit & 2048-bit & 224-bit & 9:1 \\
128-bit & 3072-bit & 256-bit & 12:1 \\
192-bit & 7680-bit & 384-bit & 20:1 \\
256-bit & 15360-bit & 521-bit & 30:1 \\
\end{longtable}

For PDF signatures specifically, the \texttt{/Contents} hex string must
be large enough to hold the entire CMS \texttt{SignedData} structure,
including the signer's certificate chain. With RSA-4096 certificates and
a chain of three certificates, this blob can exceed 10 KB. With ECDSA
P-256, the same chain might be 3 KB. Chapter 13 covers why the size of
\texttt{/Contents} matters --- it determines the byte range gap, which
must be estimated before signing.

\subsection{ECDSA Adoption in PDF
Signatures}\label{_ecdsa_adoption_in_pdf_signatures}

ECDSA adoption in PDF has been slower than in TLS, where P-256 became
the dominant algorithm years ago. The reasons are institutional rather
than technical.

Certificate authorities have been issuing RSA signing certificates for
two decades, and existing workflows are built around RSA. Many
enterprise signing solutions --- particularly those using HSMs
provisioned in the 2010s --- have RSA keys burned into hardware. The
Adobe AATL program includes ECDSA certificates from several providers,
but the majority of AATL roots remain RSA.

That said, new deployments increasingly choose ECDSA. The eIDAS 2.0
Architecture Reference Framework (Regulation (EU) 2024/1183 (eIDAS 2.0))
lists ECDSA as a recommended algorithm for the EUDI wallet. If you are
starting a new signing project and you have control over the certificate
procurement process, ECDSA with P-256 is the technically superior
choice.

The random nonce \emph{k} in ECDSA is a notorious source of catastrophic
failures. If the same \emph{k} is used for two different messages, the
private key can be algebraically recovered. This is not a theoretical
concern --- it has happened in production. Sony's PlayStation 3
code-signing key was extracted in 2010 because they used a constant
value for \emph{k}. The Bitcoin Android wallet vulnerability in 2013 was
caused by a weak random number generator producing predictable \emph{k}
values. Modern implementations use deterministic nonce generation (RFC
6979) to eliminate this class of bugs entirely. BouncyCastle 1.80
supports RFC 6979 deterministic signatures.

\section{EdDSA and Ed25519: The Future
Direction}\label{_eddsa_and_ed25519_the_future_direction}

Edwards-curve Digital Signature Algorithm (EdDSA) is a modern signature
scheme designed by Daniel J. Bernstein and colleagues, published as RFC
8032 in 2017. It addresses several practical shortcomings of ECDSA while
providing equivalent or better security.

\subsection{Bernstein's Work and Design
Philosophy}\label{_bernsteins_work_and_design_philosophy}

Daniel J. Bernstein --- "djb" to everyone in the cryptographic
community --- designed Curve25519 in 2005 and the Ed25519 signature
scheme in 2011 with a very specific philosophy: eliminate every source
of implementation error.

The NIST curves (P-256, P-384) have been criticized for their opaque
parameter generation --- NIST never fully explained how the "random"
seed values were chosen, which fueled suspicion (especially after the
Dual\_EC\_DRBG backdoor scandal in 2013, where the NSA was found to have
inserted a backdoor into a NIST-standardized random number generator).
Bernstein designed Curve25519 with fully transparent, verifiable
parameters --- every constant has a documented justification.

More importantly, Ed25519 was designed to be \emph{implementable}
without the subtle pitfalls that plague ECDSA:

\begin{enumerate}
\def\labelenumi{\arabic{enumi}.}
\item
  \textbf{Deterministic signatures.} EdDSA does not use a random nonce.
  The nonce is derived deterministically from the private key and the
  message. This eliminates the entire class of nonce-reuse
  vulnerabilities. There is no Sony PS3 scenario with EdDSA --- it is
  structurally impossible.
\item
  \textbf{Constant-time by design.} The algorithm was designed from the
  ground up to run in constant time. There are no data-dependent
  branches or memory accesses that could leak information about the
  private key through timing, power analysis, or cache behavior.
\item
  \textbf{No parameter negotiation.} You use Ed25519 (128-bit security)
  or Ed448 (224-bit security). That is it. No curves to choose, no hash
  functions to configure, no padding schemes to get wrong.
\end{enumerate}

\subsection{Ed25519 and Ed448}\label{_ed25519_and_ed448}

EdDSA defines two concrete instances:

\begin{itemize}
\item
  \textbf{Ed25519}: Based on Curve25519 (a twisted Edwards curve).
  Provides \textasciitilde128-bit security. Signatures are 64 bytes,
  public keys are 32 bytes. This is the one everyone uses.
\item
  \textbf{Ed448}: Based on Curve448 (the "Goldilocks" curve, named by
  Mike Hamburg because 448 = 224 * 2, a "just right" security level).
  Provides \textasciitilde224-bit security. Signatures are 114 bytes,
  public keys are 57 bytes.
\end{itemize}

\subsection{EdDSA and PDF: Adoption
Status}\label{_eddsa_and_pdf_adoption_status}

Here is where the good news ends, at least for PDF signatures.

As of 2026, \textbf{EdDSA support in the PDF signature ecosystem is
minimal}. ISO 32000-2:2020 does not explicitly mention EdDSA. The PAdES
standards (ETSI EN 319 142-1 V1.1.1 (2016-04)) reference ECDSA but not
EdDSA. Adobe Acrobat does not support EdDSA signatures --- a PDF signed
with Ed25519 will show as "unknown algorithm" or "signature invalid" in
Acrobat.

\textbf{Do not use EdDSA for PDF signatures in production today.}

While the algorithm itself is excellent, the ecosystem does not support
it. A signature that cannot be validated by the dominant PDF viewer is,
for all practical purposes, useless.

That said, EdDSA adoption is accelerating in the broader PKI ecosystem
(TLS, SSH, code signing), and it is a matter of when --- not if --- PDF
viewers add support. The eIDAS 2.0 Architecture Reference Framework
(Regulation (EU) 2024/1183 (eIDAS 2.0)) includes EdDSA in its list of
acceptable algorithms for the EUDI wallet. When EUDI wallet-based PDF
signing becomes mainstream (Chapter 27), EdDSA support in PDF viewers
will follow.

\section{Algorithm Comparison}\label{_algorithm_comparison}

\begin{longtable}[]{@{}
  >{\raggedright\arraybackslash}p{(\linewidth - 10\tabcolsep) * \real{0.2500}}
  >{\raggedright\arraybackslash}p{(\linewidth - 10\tabcolsep) * \real{0.1250}}
  >{\raggedright\arraybackslash}p{(\linewidth - 10\tabcolsep) * \real{0.1250}}
  >{\raggedright\arraybackslash}p{(\linewidth - 10\tabcolsep) * \real{0.1250}}
  >{\raggedright\arraybackslash}p{(\linewidth - 10\tabcolsep) * \real{0.1250}}
  >{\raggedright\arraybackslash}p{(\linewidth - 10\tabcolsep) * \real{0.2500}}@{}}
\caption{Algorithm comparison for PDF digital signatures (as of
2026)}\tabularnewline
\toprule\noalign{}
\begin{minipage}[b]{\linewidth}\raggedright
Property
\end{minipage} & \begin{minipage}[b]{\linewidth}\raggedright
RSA-3072
\end{minipage} & \begin{minipage}[b]{\linewidth}\raggedright
ECDSA P-256
\end{minipage} & \begin{minipage}[b]{\linewidth}\raggedright
ECDSA P-384
\end{minipage} & \begin{minipage}[b]{\linewidth}\raggedright
Ed25519
\end{minipage} & \begin{minipage}[b]{\linewidth}\raggedright
Notes
\end{minipage} \\
\midrule\noalign{}
\endfirsthead
\toprule\noalign{}
\begin{minipage}[b]{\linewidth}\raggedright
Property
\end{minipage} & \begin{minipage}[b]{\linewidth}\raggedright
RSA-3072
\end{minipage} & \begin{minipage}[b]{\linewidth}\raggedright
ECDSA P-256
\end{minipage} & \begin{minipage}[b]{\linewidth}\raggedright
ECDSA P-384
\end{minipage} & \begin{minipage}[b]{\linewidth}\raggedright
Ed25519
\end{minipage} & \begin{minipage}[b]{\linewidth}\raggedright
Notes
\end{minipage} \\
\midrule\noalign{}
\endhead
\bottomrule\noalign{}
\endlastfoot
Security level & 128-bit & 128-bit & 192-bit & \textasciitilde128-bit &
All except P-384 provide equivalent security. \\
Public key size & \textasciitilde400 bytes & \textasciitilde90 bytes &
\textasciitilde120 bytes & 32 bytes & Encoded sizes
(SubjectPublicKeyInfo). \\
Signature size & 384 bytes & 64 bytes & 96 bytes & 64 bytes & Raw
signature value before CMS packaging. \\
Keygen speed & Slow (probabilistic prime search) & Fast & Fast & Fast &
RSA-3072 keygen: 200-1500ms. EC: under 5ms. \\
Sign speed & Moderate & Fast & Fast & Fast & RSA scales poorly with key
size. \\
Verify speed & Fast & Moderate & Moderate & Fast & RSA verification is
cheap (small \emph{e}). \\
Adobe Acrobat & Full support & Full support & Full support & No support
& Acrobat drives compatibility requirements. \\
PDF spec support & ISO 32000-2:2020 & ISO 32000-2:2020 & ISO
32000-2:2020 & Not mentioned & EdDSA OIDs are not in the PDF spec. \\
PAdES support & Yes & Yes & Yes & No & PAdES references ECDSA, not
EdDSA. \\
\end{longtable}

\begin{longtable}[]{@{}
  >{\raggedright\arraybackslash}p{(\linewidth - 8\tabcolsep) * \real{0.3333}}
  >{\raggedright\arraybackslash}p{(\linewidth - 8\tabcolsep) * \real{0.1667}}
  >{\raggedright\arraybackslash}p{(\linewidth - 8\tabcolsep) * \real{0.1667}}
  >{\raggedright\arraybackslash}p{(\linewidth - 8\tabcolsep) * \real{0.1667}}
  >{\raggedright\arraybackslash}p{(\linewidth - 8\tabcolsep) * \real{0.1667}}@{}}
\caption{Algorithm support across PDF viewers and libraries (as of early
2026)}\tabularnewline
\toprule\noalign{}
\begin{minipage}[b]{\linewidth}\raggedright
Viewer / Library
\end{minipage} & \begin{minipage}[b]{\linewidth}\raggedright
RSA
\end{minipage} & \begin{minipage}[b]{\linewidth}\raggedright
ECDSA (P-256)
\end{minipage} & \begin{minipage}[b]{\linewidth}\raggedright
ECDSA (P-384/P-521)
\end{minipage} & \begin{minipage}[b]{\linewidth}\raggedright
EdDSA
\end{minipage} \\
\midrule\noalign{}
\endfirsthead
\toprule\noalign{}
\begin{minipage}[b]{\linewidth}\raggedright
Viewer / Library
\end{minipage} & \begin{minipage}[b]{\linewidth}\raggedright
RSA
\end{minipage} & \begin{minipage}[b]{\linewidth}\raggedright
ECDSA (P-256)
\end{minipage} & \begin{minipage}[b]{\linewidth}\raggedright
ECDSA (P-384/P-521)
\end{minipage} & \begin{minipage}[b]{\linewidth}\raggedright
EdDSA
\end{minipage} \\
\midrule\noalign{}
\endhead
\bottomrule\noalign{}
\endlastfoot
Adobe Acrobat DC & Yes & Yes & Yes (P-384); Partial (P-521) & No \\
Foxit Reader & Yes & Yes & Yes (P-384); No (P-521) & No \\
PDFBox 3.0.0 & Yes & Yes & Yes & Partial (via BouncyCastle provider) \\
iText 8.x & Yes & Yes & Yes & No \\
BouncyCastle 1.80 & Yes & Yes & Yes & Yes (Ed25519, Ed448) \\
LibreOffice Draw & Yes & Yes & Partial & No \\
\end{longtable}

\section{Connecting to PDF: OIDs and
SignerInfo}\label{_connecting_to_pdf_oids_and_signerinfo}

When a PDF signature is created, the signing algorithm is encoded as an
ASN.1 Object Identifier (OID) inside the CMS \texttt{SignedData}
structure (RFC 5652 --- Cryptographic Message Syntax (CMS)).
Specifically, it appears in the
\texttt{SignerInfo.digestEncryptionAlgorithm} field (the CMS name for
the signature algorithm --- "digestEncryptionAlgorithm" is a historical
misnomer that has caused confusion for two decades, but we are stuck
with it).

\begin{longtable}[]{@{}
  >{\raggedright\arraybackslash}p{(\linewidth - 4\tabcolsep) * \real{0.3750}}
  >{\raggedright\arraybackslash}p{(\linewidth - 4\tabcolsep) * \real{0.3750}}
  >{\raggedright\arraybackslash}p{(\linewidth - 4\tabcolsep) * \real{0.2500}}@{}}
\caption{Common algorithm OIDs you will encounter in PDF
signatures}\tabularnewline
\toprule\noalign{}
\begin{minipage}[b]{\linewidth}\raggedright
Algorithm
\end{minipage} & \begin{minipage}[b]{\linewidth}\raggedright
OID
\end{minipage} & \begin{minipage}[b]{\linewidth}\raggedright
Where You See It
\end{minipage} \\
\midrule\noalign{}
\endfirsthead
\toprule\noalign{}
\begin{minipage}[b]{\linewidth}\raggedright
Algorithm
\end{minipage} & \begin{minipage}[b]{\linewidth}\raggedright
OID
\end{minipage} & \begin{minipage}[b]{\linewidth}\raggedright
Where You See It
\end{minipage} \\
\midrule\noalign{}
\endhead
\bottomrule\noalign{}
\endlastfoot
RSA (PKCS\#1 v1.5) & \texttt{1.2.840.113549.1.1.1} & Legacy signatures.
The hash algorithm is specified separately in
\texttt{digestAlgorithm}. \\
SHA-256 with RSA & \texttt{1.2.840.113549.1.1.11} & Most common in
modern PDF signatures. \\
SHA-384 with RSA & \texttt{1.2.840.113549.1.1.12} & Higher-security
RSA. \\
SHA-512 with RSA & \texttt{1.2.840.113549.1.1.13} & Government and
defense contexts. \\
RSASSA-PSS & \texttt{1.2.840.113549.1.1.10} & Probabilistic padding.
Hash specified in parameters. \\
ECDSA with SHA-256 & \texttt{1.2.840.10045.4.3.2} & Growing adoption. \\
ECDSA with SHA-384 & \texttt{1.2.840.10045.4.3.3} & Common with P-384
curves. \\
Ed25519 & \texttt{1.3.101.112} & Not yet seen in production PDF
signatures. \\
\end{longtable}

When you encounter the OID \texttt{1.2.840.113549.1.1.1} (plain RSA) in
the \texttt{digestEncryptionAlgorithm} field, do not panic. This does
not mean the signature used no hash. It means the CMS was constructed in
the older style where the signature algorithm OID says "RSA" and the
hash algorithm is determined solely by the \texttt{digestAlgorithm}
field in the \texttt{SignerInfo}. Both styles --- the combined OID like
\texttt{1.2.840.113549.1.1.11} (SHA-256 with RSA) and the split
approach --- are valid per RFC 5652 --- Cryptographic Message Syntax
(CMS). PDFBox and BouncyCastle handle both transparently.

\section{Java's Cryptographic Provider Model and
BouncyCastle}\label{_javas_cryptographic_provider_model_and_bouncycastle}

Before we write any code, we need to understand how Java organizes its
cryptographic implementations, because this directly affects how you
generate keys, create signatures, and interact with BouncyCastle.

\subsection{JCA, JCE, and the Provider
Architecture}\label{_jca_jce_and_the_provider_architecture}

The \textbf{Java Cryptography Architecture (JCA)} defines abstract
interfaces --- \texttt{KeyPairGenerator}, \texttt{Signature},
\texttt{MessageDigest}, \texttt{Cipher} --- and a \textbf{provider
model} where concrete implementations are plugged in at runtime. The
\textbf{Java Cryptography Extension (JCE)} extends JCA with additional
capabilities. In modern Java (since JDK 9), the distinction is mostly
historical.

\begin{Shaded}
\begin{Highlighting}[]
\CommentTok{// This does NOT specify which implementation to use}
\BuiltInTok{KeyPairGenerator}\NormalTok{ kpg }\OperatorTok{=} \BuiltInTok{KeyPairGenerator}\OperatorTok{.}\FunctionTok{getInstance}\OperatorTok{(}\StringTok{"RSA"}\OperatorTok{);}

\CommentTok{// The JCA searches registered providers in priority order}
\CommentTok{// and returns the first one that supports "RSA"}
\end{Highlighting}
\end{Shaded}

The JDK ships with built-in providers: \texttt{SUN},
\texttt{SunRsaSign}, \texttt{SunEC}, \texttt{SunJCE}. These cover the
basics --- RSA, ECDSA with NIST curves, AES, SHA-2 hashes. For many
applications, the built-in providers are sufficient.

For PDF signatures, they are not.

\subsection{Why BouncyCastle Exists}\label{_why_bouncycastle_exists}

The \textbf{Legion of the Bouncy Castle} (yes, that is the
organization's legal name, registered in Australia) provides the most
comprehensive open-source cryptographic library for Java. BouncyCastle
is not an alternative to the JCA --- it is a \textbf{JCA provider} that
plugs into the same framework.

We need BouncyCastle for PDF signatures for several reasons:

\begin{enumerate}
\def\labelenumi{\arabic{enumi}.}
\item
  \textbf{CMS/PKCS\#7 implementation.} The JDK does not include a CMS
  library. BouncyCastle's \texttt{bcpkix} module provides
  \texttt{CMSSignedDataGenerator}, \texttt{CMSSignedData}, and
  everything else needed to construct and parse the CMS
  \texttt{SignedData} structures that PDF signatures require (RFC 5652
  --- Cryptographic Message Syntax (CMS)).
\item
  \textbf{Algorithm breadth.} BouncyCastle supports algorithms the JDK
  does not, including EdDSA (Ed25519, Ed448), certain elliptic curves,
  and legacy algorithms that appear in older signed PDFs.
\item
  \textbf{ASN.1 library.} BouncyCastle includes a complete ASN.1 parsing
  and generation library --- essential for working with X.509
  certificates, CMS structures, and OCSP responses at the byte level.
\item
  \textbf{Certificate generation.} The \texttt{bcpkix} module includes
  \texttt{X509v3CertificateBuilder} for creating X.509 certificates
  programmatically --- something you cannot do with the JDK alone.
\end{enumerate}

\subsection{Provider Registration}\label{_provider_registration}

\begin{Shaded}
\begin{Highlighting}[]
\KeywordTok{import} \ImportTok{java}\OperatorTok{.}\ImportTok{security}\OperatorTok{.}\ImportTok{Security}\OperatorTok{;}
\KeywordTok{import} \ImportTok{org}\OperatorTok{.}\ImportTok{bouncycastle}\OperatorTok{.}\ImportTok{jce}\OperatorTok{.}\ImportTok{provider}\OperatorTok{.}\ImportTok{BouncyCastleProvider}\OperatorTok{;}

\CommentTok{// Register BouncyCastle as a JCA provider}
\CommentTok{// This must happen before any cryptographic operations}
\BuiltInTok{Security}\OperatorTok{.}\FunctionTok{addProvider}\OperatorTok{(}\KeywordTok{new} \FunctionTok{BouncyCastleProvider}\OperatorTok{());}
\end{Highlighting}
\end{Shaded}

In a Spring Boot application, register the BouncyCastle provider in a
\texttt{@PostConstruct} method of a \texttt{@Configuration} class. This
ensures it is available before any bean that needs cryptographic
operations is initialized. PDFalyzer does this in
\texttt{BouncyCastleConfig.java}. If you forget this step, you will get
\texttt{NoSuchAlgorithmException} or \texttt{NoSuchProviderException} at
runtime --- and the error message will not tell you that a missing
provider registration is the cause.

\section{Practical Code: Generating Key
Pairs}\label{_practical_code_generating_key_pairs}

\subsection{RSA Key Pair Generation}\label{_rsa_key_pair_generation}

\begin{Shaded}
\begin{Highlighting}[]
\KeywordTok{import} \ImportTok{java}\OperatorTok{.}\ImportTok{security}\OperatorTok{.}\ImportTok{KeyPair}\OperatorTok{;}
\KeywordTok{import} \ImportTok{java}\OperatorTok{.}\ImportTok{security}\OperatorTok{.}\ImportTok{KeyPairGenerator}\OperatorTok{;}
\KeywordTok{import} \ImportTok{java}\OperatorTok{.}\ImportTok{security}\OperatorTok{.}\ImportTok{Security}\OperatorTok{;}
\KeywordTok{import} \ImportTok{java}\OperatorTok{.}\ImportTok{security}\OperatorTok{.}\ImportTok{interfaces}\OperatorTok{.}\ImportTok{RSAPrivateKey}\OperatorTok{;}
\KeywordTok{import} \ImportTok{java}\OperatorTok{.}\ImportTok{security}\OperatorTok{.}\ImportTok{interfaces}\OperatorTok{.}\ImportTok{RSAPublicKey}\OperatorTok{;}

\KeywordTok{import} \ImportTok{org}\OperatorTok{.}\ImportTok{bouncycastle}\OperatorTok{.}\ImportTok{jce}\OperatorTok{.}\ImportTok{provider}\OperatorTok{.}\ImportTok{BouncyCastleProvider}\OperatorTok{;}

\KeywordTok{public} \KeywordTok{class}\NormalTok{ RsaKeyPairExample }\OperatorTok{\{}

    \KeywordTok{public} \DataTypeTok{static} \BuiltInTok{KeyPair} \FunctionTok{generateRsaKeyPair}\OperatorTok{(}\DataTypeTok{int}\NormalTok{ keySize}\OperatorTok{)} \KeywordTok{throws} \BuiltInTok{Exception} \OperatorTok{\{}
        \BuiltInTok{Security}\OperatorTok{.}\FunctionTok{addProvider}\OperatorTok{(}\KeywordTok{new} \FunctionTok{BouncyCastleProvider}\OperatorTok{());}

        \BuiltInTok{KeyPairGenerator}\NormalTok{ kpg }\OperatorTok{=} \BuiltInTok{KeyPairGenerator}\OperatorTok{.}\FunctionTok{getInstance}\OperatorTok{(}\StringTok{"RSA"}\OperatorTok{,} \StringTok{"BC"}\OperatorTok{);}
\NormalTok{        kpg}\OperatorTok{.}\FunctionTok{initialize}\OperatorTok{(}\NormalTok{keySize}\OperatorTok{);}
        \ControlFlowTok{return}\NormalTok{ kpg}\OperatorTok{.}\FunctionTok{generateKeyPair}\OperatorTok{();}
    \OperatorTok{\}}

    \KeywordTok{public} \DataTypeTok{static} \DataTypeTok{void} \FunctionTok{main}\OperatorTok{(}\BuiltInTok{String}\OperatorTok{[]}\NormalTok{ args}\OperatorTok{)} \KeywordTok{throws} \BuiltInTok{Exception} \OperatorTok{\{}
        \BuiltInTok{KeyPair}\NormalTok{ kp }\OperatorTok{=} \FunctionTok{generateRsaKeyPair}\OperatorTok{(}\DecValTok{3072}\OperatorTok{);}

        \BuiltInTok{RSAPublicKey}\NormalTok{ pub }\OperatorTok{=} \OperatorTok{(}\BuiltInTok{RSAPublicKey}\OperatorTok{)}\NormalTok{ kp}\OperatorTok{.}\FunctionTok{getPublic}\OperatorTok{();}
        \BuiltInTok{RSAPrivateKey}\NormalTok{ priv }\OperatorTok{=} \OperatorTok{(}\BuiltInTok{RSAPrivateKey}\OperatorTok{)}\NormalTok{ kp}\OperatorTok{.}\FunctionTok{getPrivate}\OperatorTok{();}

        \BuiltInTok{System}\OperatorTok{.}\FunctionTok{out}\OperatorTok{.}\FunctionTok{println}\OperatorTok{(}\StringTok{"Algorithm:    "} \OperatorTok{+}\NormalTok{ pub}\OperatorTok{.}\FunctionTok{getAlgorithm}\OperatorTok{());}
        \BuiltInTok{System}\OperatorTok{.}\FunctionTok{out}\OperatorTok{.}\FunctionTok{println}\OperatorTok{(}\StringTok{"Key size:     "} \OperatorTok{+}\NormalTok{ pub}\OperatorTok{.}\FunctionTok{getModulus}\OperatorTok{().}\FunctionTok{bitLength}\OperatorTok{()} \OperatorTok{+} \StringTok{" bits"}\OperatorTok{);}
        \BuiltInTok{System}\OperatorTok{.}\FunctionTok{out}\OperatorTok{.}\FunctionTok{println}\OperatorTok{(}\StringTok{"Pub exponent: "} \OperatorTok{+}\NormalTok{ pub}\OperatorTok{.}\FunctionTok{getPublicExponent}\OperatorTok{());}
        \BuiltInTok{System}\OperatorTok{.}\FunctionTok{out}\OperatorTok{.}\FunctionTok{println}\OperatorTok{(}\StringTok{"Pub format:   "} \OperatorTok{+}\NormalTok{ pub}\OperatorTok{.}\FunctionTok{getFormat}\OperatorTok{());}  \CommentTok{// X.509}
        \BuiltInTok{System}\OperatorTok{.}\FunctionTok{out}\OperatorTok{.}\FunctionTok{println}\OperatorTok{(}\StringTok{"Priv format:  "} \OperatorTok{+}\NormalTok{ priv}\OperatorTok{.}\FunctionTok{getFormat}\OperatorTok{());} \CommentTok{// PKCS\#8}
    \OperatorTok{\}}
\OperatorTok{\}}
\end{Highlighting}
\end{Shaded}

The output will show a 3072-bit key with public exponent 65537, the
public key in X.509 (SubjectPublicKeyInfo) format, and the private key
in PKCS\#8 format. These format names --- X.509 for public keys, PKCS\#8
for private keys --- are the standard encodings that appear in PEM
files. Chapter 5 covers key material formats in detail.

\subsection{ECDSA Key Pair Generation}\label{_ecdsa_key_pair_generation}

\begin{Shaded}
\begin{Highlighting}[]
\KeywordTok{import} \ImportTok{java}\OperatorTok{.}\ImportTok{security}\OperatorTok{.}\ImportTok{KeyPair}\OperatorTok{;}
\KeywordTok{import} \ImportTok{java}\OperatorTok{.}\ImportTok{security}\OperatorTok{.}\ImportTok{KeyPairGenerator}\OperatorTok{;}
\KeywordTok{import} \ImportTok{java}\OperatorTok{.}\ImportTok{security}\OperatorTok{.}\ImportTok{Security}\OperatorTok{;}
\KeywordTok{import} \ImportTok{java}\OperatorTok{.}\ImportTok{security}\OperatorTok{.}\ImportTok{interfaces}\OperatorTok{.}\ImportTok{ECPrivateKey}\OperatorTok{;}
\KeywordTok{import} \ImportTok{java}\OperatorTok{.}\ImportTok{security}\OperatorTok{.}\ImportTok{interfaces}\OperatorTok{.}\ImportTok{ECPublicKey}\OperatorTok{;}
\KeywordTok{import} \ImportTok{java}\OperatorTok{.}\ImportTok{security}\OperatorTok{.}\ImportTok{spec}\OperatorTok{.}\ImportTok{ECGenParameterSpec}\OperatorTok{;}

\KeywordTok{import} \ImportTok{org}\OperatorTok{.}\ImportTok{bouncycastle}\OperatorTok{.}\ImportTok{jce}\OperatorTok{.}\ImportTok{provider}\OperatorTok{.}\ImportTok{BouncyCastleProvider}\OperatorTok{;}

\KeywordTok{public} \KeywordTok{class}\NormalTok{ EcdsaKeyPairExample }\OperatorTok{\{}

    \KeywordTok{public} \DataTypeTok{static} \BuiltInTok{KeyPair} \FunctionTok{generateEcdsaKeyPair}\OperatorTok{(}\BuiltInTok{String}\NormalTok{ curveName}\OperatorTok{)} \KeywordTok{throws} \BuiltInTok{Exception} \OperatorTok{\{}
        \BuiltInTok{Security}\OperatorTok{.}\FunctionTok{addProvider}\OperatorTok{(}\KeywordTok{new} \FunctionTok{BouncyCastleProvider}\OperatorTok{());}

        \BuiltInTok{KeyPairGenerator}\NormalTok{ kpg }\OperatorTok{=} \BuiltInTok{KeyPairGenerator}\OperatorTok{.}\FunctionTok{getInstance}\OperatorTok{(}\StringTok{"EC"}\OperatorTok{,} \StringTok{"BC"}\OperatorTok{);}
\NormalTok{        kpg}\OperatorTok{.}\FunctionTok{initialize}\OperatorTok{(}\KeywordTok{new} \BuiltInTok{ECGenParameterSpec}\OperatorTok{(}\NormalTok{curveName}\OperatorTok{));}
        \ControlFlowTok{return}\NormalTok{ kpg}\OperatorTok{.}\FunctionTok{generateKeyPair}\OperatorTok{();}
    \OperatorTok{\}}

    \KeywordTok{public} \DataTypeTok{static} \DataTypeTok{void} \FunctionTok{main}\OperatorTok{(}\BuiltInTok{String}\OperatorTok{[]}\NormalTok{ args}\OperatorTok{)} \KeywordTok{throws} \BuiltInTok{Exception} \OperatorTok{\{}
        \CommentTok{// P{-}256: use the SECG name, not the NIST name}
        \BuiltInTok{KeyPair}\NormalTok{ kp }\OperatorTok{=} \FunctionTok{generateEcdsaKeyPair}\OperatorTok{(}\StringTok{"secp256r1"}\OperatorTok{);}

        \BuiltInTok{ECPublicKey}\NormalTok{ pub }\OperatorTok{=} \OperatorTok{(}\BuiltInTok{ECPublicKey}\OperatorTok{)}\NormalTok{ kp}\OperatorTok{.}\FunctionTok{getPublic}\OperatorTok{();}
        \BuiltInTok{ECPrivateKey}\NormalTok{ priv }\OperatorTok{=} \OperatorTok{(}\BuiltInTok{ECPrivateKey}\OperatorTok{)}\NormalTok{ kp}\OperatorTok{.}\FunctionTok{getPrivate}\OperatorTok{();}

        \BuiltInTok{System}\OperatorTok{.}\FunctionTok{out}\OperatorTok{.}\FunctionTok{println}\OperatorTok{(}\StringTok{"Algorithm:   "} \OperatorTok{+}\NormalTok{ pub}\OperatorTok{.}\FunctionTok{getAlgorithm}\OperatorTok{());}
        \BuiltInTok{System}\OperatorTok{.}\FunctionTok{out}\OperatorTok{.}\FunctionTok{println}\OperatorTok{(}\StringTok{"Curve order: "} \OperatorTok{+}\NormalTok{ pub}\OperatorTok{.}\FunctionTok{getParams}\OperatorTok{().}\FunctionTok{getOrder}\OperatorTok{().}\FunctionTok{bitLength}\OperatorTok{()} \OperatorTok{+} \StringTok{" bits"}\OperatorTok{);}
        \BuiltInTok{System}\OperatorTok{.}\FunctionTok{out}\OperatorTok{.}\FunctionTok{println}\OperatorTok{(}\StringTok{"Pub format:  "} \OperatorTok{+}\NormalTok{ pub}\OperatorTok{.}\FunctionTok{getFormat}\OperatorTok{());}  \CommentTok{// X.509}
        \BuiltInTok{System}\OperatorTok{.}\FunctionTok{out}\OperatorTok{.}\FunctionTok{println}\OperatorTok{(}\StringTok{"Priv format: "} \OperatorTok{+}\NormalTok{ priv}\OperatorTok{.}\FunctionTok{getFormat}\OperatorTok{());} \CommentTok{// PKCS\#8}

        \CommentTok{// The public key is a point on the curve (x, y)}
        \BuiltInTok{System}\OperatorTok{.}\FunctionTok{out}\OperatorTok{.}\FunctionTok{println}\OperatorTok{(}\StringTok{"Pub X:       "} \OperatorTok{+}\NormalTok{ pub}\OperatorTok{.}\FunctionTok{getW}\OperatorTok{().}\FunctionTok{getAffineX}\OperatorTok{().}\FunctionTok{bitLength}\OperatorTok{()} \OperatorTok{+} \StringTok{" bits"}\OperatorTok{);}
        \BuiltInTok{System}\OperatorTok{.}\FunctionTok{out}\OperatorTok{.}\FunctionTok{println}\OperatorTok{(}\StringTok{"Pub Y:       "} \OperatorTok{+}\NormalTok{ pub}\OperatorTok{.}\FunctionTok{getW}\OperatorTok{().}\FunctionTok{getAffineY}\OperatorTok{().}\FunctionTok{bitLength}\OperatorTok{()} \OperatorTok{+} \StringTok{" bits"}\OperatorTok{);}
    \OperatorTok{\}}
\OperatorTok{\}}
\end{Highlighting}
\end{Shaded}

Note the curve name: \texttt{"secp256r1"}, not \texttt{"P-256"}. The JCA
uses SECG names in \texttt{ECGenParameterSpec}. BouncyCastle accepts
both, but for portability, use the SECG name.

\subsection{Performance Comparison}\label{_performance_comparison}

\begin{Shaded}
\begin{Highlighting}[]
\KeywordTok{import} \ImportTok{java}\OperatorTok{.}\ImportTok{security}\OperatorTok{.}\ImportTok{KeyPair}\OperatorTok{;}
\KeywordTok{import} \ImportTok{java}\OperatorTok{.}\ImportTok{security}\OperatorTok{.}\ImportTok{KeyPairGenerator}\OperatorTok{;}
\KeywordTok{import} \ImportTok{java}\OperatorTok{.}\ImportTok{security}\OperatorTok{.}\ImportTok{Security}\OperatorTok{;}
\KeywordTok{import} \ImportTok{java}\OperatorTok{.}\ImportTok{security}\OperatorTok{.}\ImportTok{spec}\OperatorTok{.}\ImportTok{ECGenParameterSpec}\OperatorTok{;}

\KeywordTok{import} \ImportTok{org}\OperatorTok{.}\ImportTok{bouncycastle}\OperatorTok{.}\ImportTok{jce}\OperatorTok{.}\ImportTok{provider}\OperatorTok{.}\ImportTok{BouncyCastleProvider}\OperatorTok{;}

\KeywordTok{public} \KeywordTok{class}\NormalTok{ KeyPairComparisonExample }\OperatorTok{\{}

    \KeywordTok{public} \DataTypeTok{static} \DataTypeTok{void} \FunctionTok{main}\OperatorTok{(}\BuiltInTok{String}\OperatorTok{[]}\NormalTok{ args}\OperatorTok{)} \KeywordTok{throws} \BuiltInTok{Exception} \OperatorTok{\{}
        \BuiltInTok{Security}\OperatorTok{.}\FunctionTok{addProvider}\OperatorTok{(}\KeywordTok{new} \FunctionTok{BouncyCastleProvider}\OperatorTok{());}

        \CommentTok{// RSA{-}3072: \textasciitilde{}128{-}bit security}
        \DataTypeTok{long}\NormalTok{ start }\OperatorTok{=} \BuiltInTok{System}\OperatorTok{.}\FunctionTok{nanoTime}\OperatorTok{();}
        \BuiltInTok{KeyPairGenerator}\NormalTok{ rsaKpg }\OperatorTok{=} \BuiltInTok{KeyPairGenerator}\OperatorTok{.}\FunctionTok{getInstance}\OperatorTok{(}\StringTok{"RSA"}\OperatorTok{,} \StringTok{"BC"}\OperatorTok{);}
\NormalTok{        rsaKpg}\OperatorTok{.}\FunctionTok{initialize}\OperatorTok{(}\DecValTok{3072}\OperatorTok{);}
        \BuiltInTok{KeyPair}\NormalTok{ rsaKp }\OperatorTok{=}\NormalTok{ rsaKpg}\OperatorTok{.}\FunctionTok{generateKeyPair}\OperatorTok{();}
        \DataTypeTok{long}\NormalTok{ rsaTime }\OperatorTok{=} \BuiltInTok{System}\OperatorTok{.}\FunctionTok{nanoTime}\OperatorTok{()} \OperatorTok{{-}}\NormalTok{ start}\OperatorTok{;}

        \CommentTok{// ECDSA P{-}256: \textasciitilde{}128{-}bit security}
\NormalTok{        start }\OperatorTok{=} \BuiltInTok{System}\OperatorTok{.}\FunctionTok{nanoTime}\OperatorTok{();}
        \BuiltInTok{KeyPairGenerator}\NormalTok{ ecKpg }\OperatorTok{=} \BuiltInTok{KeyPairGenerator}\OperatorTok{.}\FunctionTok{getInstance}\OperatorTok{(}\StringTok{"EC"}\OperatorTok{,} \StringTok{"BC"}\OperatorTok{);}
\NormalTok{        ecKpg}\OperatorTok{.}\FunctionTok{initialize}\OperatorTok{(}\KeywordTok{new} \BuiltInTok{ECGenParameterSpec}\OperatorTok{(}\StringTok{"secp256r1"}\OperatorTok{));}
        \BuiltInTok{KeyPair}\NormalTok{ ecKp }\OperatorTok{=}\NormalTok{ ecKpg}\OperatorTok{.}\FunctionTok{generateKeyPair}\OperatorTok{();}
        \DataTypeTok{long}\NormalTok{ ecTime }\OperatorTok{=} \BuiltInTok{System}\OperatorTok{.}\FunctionTok{nanoTime}\OperatorTok{()} \OperatorTok{{-}}\NormalTok{ start}\OperatorTok{;}

        \BuiltInTok{System}\OperatorTok{.}\FunctionTok{out}\OperatorTok{.}\FunctionTok{println}\OperatorTok{(}\StringTok{"RSA{-}3072 keygen:    "} \OperatorTok{+}\NormalTok{ rsaTime }\OperatorTok{/} \DecValTok{1\_000\_000} \OperatorTok{+} \StringTok{" ms"}\OperatorTok{);}
        \BuiltInTok{System}\OperatorTok{.}\FunctionTok{out}\OperatorTok{.}\FunctionTok{println}\OperatorTok{(}\StringTok{"ECDSA P{-}256 keygen: "} \OperatorTok{+}\NormalTok{ ecTime }\OperatorTok{/} \DecValTok{1\_000\_000} \OperatorTok{+} \StringTok{" ms"}\OperatorTok{);}
        \BuiltInTok{System}\OperatorTok{.}\FunctionTok{out}\OperatorTok{.}\FunctionTok{println}\OperatorTok{(}\StringTok{"RSA pub key size:   "} \OperatorTok{+}\NormalTok{ rsaKp}\OperatorTok{.}\FunctionTok{getPublic}\OperatorTok{().}\FunctionTok{getEncoded}\OperatorTok{().}\FunctionTok{length} \OperatorTok{+} \StringTok{" bytes"}\OperatorTok{);}
        \BuiltInTok{System}\OperatorTok{.}\FunctionTok{out}\OperatorTok{.}\FunctionTok{println}\OperatorTok{(}\StringTok{"EC pub key size:    "} \OperatorTok{+}\NormalTok{ ecKp}\OperatorTok{.}\FunctionTok{getPublic}\OperatorTok{().}\FunctionTok{getEncoded}\OperatorTok{().}\FunctionTok{length} \OperatorTok{+} \StringTok{" bytes"}\OperatorTok{);}
        \BuiltInTok{System}\OperatorTok{.}\FunctionTok{out}\OperatorTok{.}\FunctionTok{println}\OperatorTok{(}\StringTok{"RSA priv key size:  "} \OperatorTok{+}\NormalTok{ rsaKp}\OperatorTok{.}\FunctionTok{getPrivate}\OperatorTok{().}\FunctionTok{getEncoded}\OperatorTok{().}\FunctionTok{length} \OperatorTok{+} \StringTok{" bytes"}\OperatorTok{);}
        \BuiltInTok{System}\OperatorTok{.}\FunctionTok{out}\OperatorTok{.}\FunctionTok{println}\OperatorTok{(}\StringTok{"EC priv key size:   "} \OperatorTok{+}\NormalTok{ ecKp}\OperatorTok{.}\FunctionTok{getPrivate}\OperatorTok{().}\FunctionTok{getEncoded}\OperatorTok{().}\FunctionTok{length} \OperatorTok{+} \StringTok{" bytes"}\OperatorTok{);}
    \OperatorTok{\}}
\OperatorTok{\}}
\end{Highlighting}
\end{Shaded}

On a typical machine, RSA-3072 key generation takes 200-1500 ms (it
varies because finding large primes is probabilistic), while ECDSA P-256
key generation completes in under 5 ms. The encoded key sizes will show
the RSA public key at roughly 400 bytes versus 90 bytes for the EC
public key.

These numbers matter in practice. If you are building a self-signed
certificate generator for testing (as we do in Chapter 29), RSA key
generation is the bottleneck. If you are deploying a signing service
that provisions per-user keys, the hundred-fold speed advantage of ECDSA
key generation is worth considering.

\section{Choosing an Algorithm for PDF
Signatures}\label{_choosing_an_algorithm_for_pdf_signatures}

Given all of the above, what should you actually use?

\textbf{For maximum compatibility: RSA-2048 or RSA-3072 with SHA-256.}
Every PDF viewer, every library, every validation service supports this
combination. If your signed PDFs need to be validated by the widest
possible audience --- and they usually do --- RSA with SHA-256 is the
safe choice. Use 3072-bit keys for new deployments; 2048-bit is
acceptable for existing systems but provides only 112-bit security.

\textbf{For smaller signatures and modern infrastructure: ECDSA with
P-256 and SHA-256.} If you control the validation environment (e.g., an
internal document management system) or if you are targeting modern
viewers (Adobe Acrobat 2020+, PDFBox 3.x, iText 8.x), ECDSA P-256 gives
you 128-bit security with much smaller key material. The CMS blob in
\texttt{/Contents} will be noticeably smaller, which simplifies byte
range estimation.

\textbf{For future-proofing: watch EdDSA, but do not deploy it yet.}
EdDSA is the superior algorithm by every technical measure, but the PDF
ecosystem has not caught up. When Adobe Acrobat adds EdDSA
support --- likely driven by EUDI wallet adoption --- it will become the
preferred choice.

\textbf{The algorithm choice affects your certificate, not your code.}

The PDF signing code in PDFBox does not care whether you give it an RSA
key or an ECDSA key. The
\texttt{SignatureInterface.sign(InputStream\ content)} method receives
the document bytes and returns a CMS blob. Whether that blob contains an
RSA signature or an ECDSA signature depends entirely on what key
material you provide.

This means you can switch algorithms by changing your certificate and
keystore --- the signing code remains identical. We demonstrate this in
Chapter 29.

\section{Summary}\label{_summary}

Asymmetric cryptography provides two operations that make digital
signatures possible: signing with a private key, and verifying with the
corresponding public key. The security of the system rests on
mathematical problems --- integer factorization for RSA, the elliptic
curve discrete logarithm for ECDSA and EdDSA --- that are easy to set up
but computationally infeasible to reverse.

RSA remains the dominant algorithm for PDF signatures due to universal
support across all viewers and libraries. ECDSA with P-256 offers
equivalent security at 128 bits with dramatically smaller keys and
faster operations. EdDSA (Ed25519, Ed448) is technically superior but
lacks ecosystem support for PDF signatures as of 2026.

Java's JCA provider model allows BouncyCastle to extend the JDK's
cryptographic capabilities, and BouncyCastle is essential for PDF
signatures because it provides the CMS, ASN.1, and certificate
generation capabilities that the JDK lacks.

In the next chapter, we examine hash functions and the mechanics of the
signing operation itself --- how the document's bytes are reduced to a
fixed-length digest, how that digest is signed, and what happens when
the hash algorithm is compromised (as SHA-1 was, with devastating
consequences for the PDF ecosystem).

\chapter{Hash Functions \& Digital Signing}\label{ch02}

This chapter covers the two primitives that every digital signature
depends on: cryptographic hash functions and the sign/verify operation.
We go deep on the SHA family, explain why MD5 and SHA-1 are still
haunting PDF signatures decades after they were broken, and walk through
the exact sequence of bytes that turns a document into a verifiable
commitment.

If you already know what a hash function does, skip to
\hyperref[ch02-sha-family]{The SHA Family} for the forensic details of
which algorithms you will actually encounter in the wild. If you want
code, jump to \hyperref[ch02-java-hashing]{Java: Computing Hashes}.

\section{What Hash Functions Do}\label{_what_hash_functions_do}

A cryptographic hash function takes an input of arbitrary length and
produces a fixed-size output --- the \emph{digest} or \emph{hash value}.
That is it. One operation, no keys involved, purely deterministic. Feed
in the same bytes, get the same digest, every time, on every machine.

But the properties that make this useful for digital signatures are
specific and non-negotiable:

\begin{longtable}[]{@{}
  >{\raggedright\arraybackslash}p{(\linewidth - 2\tabcolsep) * \real{0.2857}}
  >{\raggedright\arraybackslash}p{(\linewidth - 2\tabcolsep) * \real{0.7143}}@{}}
\caption{Properties of a cryptographic hash function}\tabularnewline
\toprule\noalign{}
\begin{minipage}[b]{\linewidth}\raggedright
Property
\end{minipage} & \begin{minipage}[b]{\linewidth}\raggedright
What It Means
\end{minipage} \\
\midrule\noalign{}
\endfirsthead
\toprule\noalign{}
\begin{minipage}[b]{\linewidth}\raggedright
Property
\end{minipage} & \begin{minipage}[b]{\linewidth}\raggedright
What It Means
\end{minipage} \\
\midrule\noalign{}
\endhead
\bottomrule\noalign{}
\endlastfoot
\textbf{Deterministic} & Same input always produces the same output. No
randomness, no state. \texttt{SHA-256("hello")} is the same value on
your laptop, on a server in Frankfurt, and on a Mars rover running Java
21. Always. \\
\textbf{Fixed-size output} & Regardless of whether the input is 1 byte
or 1 terabyte, the output is always the same length. SHA-256 produces 32
bytes (256 bits). SHA-512 produces 64 bytes (512 bits). This fixed-size
property is what makes hash functions useful as input to signing
algorithms. \\
\textbf{Avalanche effect} & Change a single bit of the input, and on
average 50\% of the output bits change. The output looks completely
unrelated to the original. There is no way to look at two digests and
determine whether their inputs were similar. \\
\textbf{Pre-image resistance} & Given a hash value \texttt{h}, it is
computationally infeasible to find any input \texttt{m} such that
\texttt{hash(m)\ =\ h}. You cannot reverse a hash function. This is not
encryption --- there is no key, no decryption, no way back. \\
\textbf{Second pre-image resistance} & Given an input \texttt{m1}, it is
computationally infeasible to find a different input \texttt{m2} such
that \texttt{hash(m1)\ =\ hash(m2)}. You cannot find a substitute
document that produces the same hash. \\
\textbf{Collision resistance} & It is computationally infeasible to find
\emph{any} two distinct inputs \texttt{m1} and \texttt{m2} such that
\texttt{hash(m1)\ =\ hash(m2)}. This is stronger than second pre-image
resistance --- the attacker gets to choose both inputs. \\
\end{longtable}

The distinction between second pre-image resistance and collision
resistance matters enormously for digital signatures. An attacker who
breaks collision resistance can prepare two documents with the same hash
\emph{before} either is signed --- get one signed, then substitute the
other. An attacker who breaks second pre-image resistance can find a
substitute for an \emph{already signed} document. Collision attacks are
easier (birthday paradox --- they require roughly 2\^{}(n/2) work for an
n-bit hash), which is why we retire algorithms at the collision
resistance boundary, not the pre-image boundary.

\section{Why We Hash Before Signing}\label{_why_we_hash_before_signing}

You might wonder: why not just sign the document directly?

Three reasons, and they all matter:

\begin{enumerate}
\def\labelenumi{\arabic{enumi}.}
\item
  \textbf{Performance.} RSA signing is slow. ECDSA signing is faster but
  still orders of magnitude slower than hashing. A 10 MB PDF would take
  an unacceptable time to sign directly with RSA-2048. Hashing reduces
  the input to 32 bytes (for SHA-256), and signing 32 bytes is fast
  regardless of the original document size.
\item
  \textbf{Fixed-size input.} RSA can only operate on data smaller than
  its modulus (typically 256 bytes for RSA-2048). You cannot feed a 10
  MB document directly into RSA. Hashing solves this by producing a
  fixed-size value that fits within the algorithm's constraints.
\item
  \textbf{Security proof structure.} The security of digital signature
  schemes is typically proven in two parts: the hash function provides
  collision resistance, and the signing algorithm provides
  unforgeability. Separating the two allows each component to be
  analyzed (and replaced) independently. When SHA-1 was broken, we could
  switch to SHA-256 without changing the RSA or ECDSA implementation.
\end{enumerate}

The process is always the same:

\begin{verbatim}
Sign:   digest = hash(data)  -->  signature = sign(digest, privateKey)
Verify: digest = hash(data)  -->  result    = verify(digest, signature, publicKey)
\end{verbatim}

The signer computes the hash of the data, then signs the hash with their
private key. The verifier computes the hash of the same data
independently, then checks the signature against that hash using the
signer's public key. If the hashes match and the signature is valid, the
data has not been modified since signing.

The verifier computes the hash independently. The hash is not sent
alongside the signature --- it is recomputed from the original data. If
an attacker modifies even one byte of the data, the verifier's
independently computed hash will differ from the hash that was signed,
and verification will fail. This is the fundamental integrity guarantee
of digital signatures.

\section{The SHA Family}\label{ch02-sha-family}

The Secure Hash Algorithm family, designed by the NSA and published by
NIST, is what you will encounter in virtually every PDF signature
created after 2010.

\subsection{SHA-1: The Algorithm That Refused to
Die}\label{_sha_1_the_algorithm_that_refused_to_die}

SHA-1 produces a 160-bit (20-byte) digest. It was published in 1995 as
FIPS PUB 180-1 and became the default hash algorithm for essentially
everything: TLS certificates, code signing, Git commits, and yes, PDF
signatures.

The theoretical weaknesses appeared early. In 2005, Xiaoyun Wang, Yiqun
Lisa Yin, and Hongbo Yu published an attack reducing the collision
resistance of SHA-1 from the expected 2\^{}80 operations to
approximately 2\^{}69 --- still impractical, but a clear signal that the
algorithm was weakening. NIST recommended transitioning away from SHA-1
by 2010 for digital signatures.

The practical break came on February 23, 2017.

\subsubsection{The SHAttered Attack}\label{_the_shattered_attack}

A team from CWI Amsterdam (Marc Stevens) and Google (Elie Bursztein,
Pierre Karpman) published the first practical SHA-1 collision. They
produced two distinct PDF files --- \texttt{shattered-1.pdf} and
\texttt{shattered-2.pdf} --- that had identical SHA-1 hashes but
displayed completely different content. One showed a blue background
with the text "SHAttered"; the other showed a red background with
different text.

The numbers:

\begin{itemize}
\item
  The attack required approximately 2\^{}63.1 SHA-1 computations
\item
  That translated to roughly 6,500 CPU-years and 110 GPU-years of
  computation
\item
  Google estimated the cost at approximately \$110,000 using cloud
  computing resources
\item
  The collision was found using a technique called an
  \emph{identical-prefix collision attack}, where both files share a
  common prefix and the collision is achieved by appending carefully
  crafted "near-collision blocks"
\end{itemize}

\begin{longtable}[]{@{}
  >{\raggedright\arraybackslash}p{(\linewidth - 2\tabcolsep) * \real{0.2000}}
  >{\raggedright\arraybackslash}p{(\linewidth - 2\tabcolsep) * \real{0.8000}}@{}}
\caption{How the SHAttered collision was constructed}\tabularnewline
\toprule\noalign{}
\begin{minipage}[b]{\linewidth}\raggedright
Component
\end{minipage} & \begin{minipage}[b]{\linewidth}\raggedright
Description
\end{minipage} \\
\midrule\noalign{}
\endfirsthead
\toprule\noalign{}
\begin{minipage}[b]{\linewidth}\raggedright
Component
\end{minipage} & \begin{minipage}[b]{\linewidth}\raggedright
Description
\end{minipage} \\
\midrule\noalign{}
\endhead
\bottomrule\noalign{}
\endlastfoot
Common prefix & Both PDFs share an identical header and initial
structure --- valid PDF objects, page definitions, and the beginning of
a JPEG image stream. \\
Collision blocks & 128 bytes of carefully computed data inserted into
each file at the same offset. These blocks differ between the two files
but are constructed so that the SHA-1 internal state converges to the
same value after processing them. \\
Divergent suffix & After the collision blocks, the remainder of each
file can be \emph{anything}. Since the SHA-1 state is identical at that
point, any shared suffix produces the same final hash. The attackers
used this to embed different JPEG image data, producing visually
distinct documents. \\
\end{longtable}

The SHAttered attack has direct consequences for PDF signatures. If a
PDF signature uses SHA-1 as its digest algorithm (SubFilter
\texttt{adbe.pkcs7.sha1} or a CMS SignedData with SHA-1 as the digest
algorithm), an attacker who controls the document content \emph{before}
signing can prepare two versions with the same SHA-1 hash. Get one
version signed, and the other is also "signed" --- the signature
validates on both. This is not a theoretical concern: the SHAttered team
demonstrated it with actual PDF files.

\subsubsection{The Cost of \$110,000}\label{_the_cost_of_110000}

It is worth dwelling on the attack cost because it illustrates how
cryptographic security erodes over time. When SHAttered was published in
2017, \$110,000 in cloud compute seemed significant. By 2020, the cost
had dropped to an estimated \$45,000 due to GPU price improvements. A
2020 paper by Ga\{eacute\}tan Leurent and Thomas Peyrin ("SHA-1 is a
Shambles") demonstrated a \emph{chosen-prefix} collision attack (more
powerful than identical-prefix) for an estimated \$45,000.

By 2026, a SHA-1 collision can likely be generated for under \$20,000.
This is well within the budget of a motivated individual, let alone a
criminal organization or nation-state. The cost trajectory is clear: it
only gets cheaper.

\subsubsection{Consequences for PDF and the Deprecation
Timeline}\label{_consequences_for_pdf_and_the_deprecation_timeline}

The SHAttered attack triggered a cascade of changes across the industry:

\begin{itemize}
\item
  ISO 32000-2:2020 deprecated the \texttt{adbe.pkcs7.sha1} SubFilter
\item
  Adobe Acrobat began showing warnings for SHA-1 signatures (initially a
  yellow warning, later escalated to an error on newer Acrobat versions)
\item
  The PAdES standards (ETSI EN 319 142-1 V1.1.1 (2016-04)) require
  SHA-256 or stronger for all profiles
\item
  The CA/Browser Forum prohibited SHA-1 in TLS certificates effective
  January 1, 2017 --- a deadline that was already set before SHAttered,
  but the attack validated the urgency
\item
  Git, which used SHA-1 for commit hashes, began its multi-year
  migration to SHA-256 (still ongoing as of 2026)
\item
  Major browsers dropped SHA-1 certificate support: Chrome (January
  2017), Firefox (January 2017), Edge (May 2017)
\end{itemize}

If you encounter a PDF with a SHA-1 based signature, do not immediately
panic. A SHA-1 signature is not \emph{forged} --- it was created
legitimately at a time when SHA-1 was considered acceptable. The risk is
that someone \emph{could} have prepared a collision before signing. For
documents signed before 2017 by trusted parties, the practical risk is
low. For documents signed after 2017 with SHA-1, ask pointed questions
about why.

\subsection{SHA-256, SHA-384, SHA-512: The Modern
Workhorses}\label{_sha_256_sha_384_sha_512_the_modern_workhorses}

These are the algorithms you should be using. All three are members of
the SHA-2 family, published by NIST in 2001 as FIPS PUB 180-2.

\begin{longtable}[]{@{}
  >{\raggedright\arraybackslash}p{(\linewidth - 8\tabcolsep) * \real{0.1250}}
  >{\raggedright\arraybackslash}p{(\linewidth - 8\tabcolsep) * \real{0.1250}}
  >{\raggedright\arraybackslash}p{(\linewidth - 8\tabcolsep) * \real{0.1250}}
  >{\raggedright\arraybackslash}p{(\linewidth - 8\tabcolsep) * \real{0.2500}}
  >{\raggedright\arraybackslash}p{(\linewidth - 8\tabcolsep) * \real{0.3750}}@{}}
\caption{SHA-2 family members used in PDF signatures}\tabularnewline
\toprule\noalign{}
\begin{minipage}[b]{\linewidth}\raggedright
Algorithm
\end{minipage} & \begin{minipage}[b]{\linewidth}\raggedright
Output (bits)
\end{minipage} & \begin{minipage}[b]{\linewidth}\raggedright
Output (bytes)
\end{minipage} & \begin{minipage}[b]{\linewidth}\raggedright
Block size
\end{minipage} & \begin{minipage}[b]{\linewidth}\raggedright
Usage in PDF signatures
\end{minipage} \\
\midrule\noalign{}
\endfirsthead
\toprule\noalign{}
\begin{minipage}[b]{\linewidth}\raggedright
Algorithm
\end{minipage} & \begin{minipage}[b]{\linewidth}\raggedright
Output (bits)
\end{minipage} & \begin{minipage}[b]{\linewidth}\raggedright
Output (bytes)
\end{minipage} & \begin{minipage}[b]{\linewidth}\raggedright
Block size
\end{minipage} & \begin{minipage}[b]{\linewidth}\raggedright
Usage in PDF signatures
\end{minipage} \\
\midrule\noalign{}
\endhead
\bottomrule\noalign{}
\endlastfoot
SHA-256 & 256 & 32 & 512 bits (64 bytes) & The default for virtually all
modern PDF signatures. PAdES-B-B requires SHA-256 as the minimum. If you
are signing a PDF today, use SHA-256 unless you have a specific reason
for a stronger algorithm. \\
SHA-384 & 384 & 48 & 1024 bits (128 bytes) & Used when organizational
policy requires stronger-than-256-bit hashing. Common in government and
defense contexts. SHA-384 is actually a truncated SHA-512 with different
initial values --- it uses the 64-bit word size of SHA-512, which makes
it faster than SHA-256 on 64-bit platforms. \\
SHA-512 & 512 & 64 & 1024 bits (128 bytes) & Maximum strength in the
SHA-2 family. Rare in PDF signatures because SHA-256 is considered
sufficient, and SHA-512 produces larger digests that increase CMS
signature size. Used in high-security environments. \\
\end{longtable}

SHA-384's performance characteristic is surprising and worth knowing.
Because it operates on 64-bit words (like SHA-512) rather than 32-bit
words (like SHA-256), SHA-384 is actually \emph{faster} than SHA-256 on
modern 64-bit processors. If you need to hash large documents and your
platform is 64-bit (it is), SHA-384 gives you more security bits at
equal or better performance. This is one of those facts that sounds
wrong but is measurably true.

\subsubsection{SHA-3: The Insurance
Policy}\label{_sha_3_the_insurance_policy}

SHA-3 (Keccak), published as FIPS 202 in 2015, is a completely different
design from SHA-2 --- it uses a sponge construction rather than the
Merkle-Damgard construction of SHA-1 and SHA-2. SHA-3 was selected by
NIST specifically because it has \emph{no} structural similarity to
SHA-2, meaning any attack on SHA-2 would not apply.

You will not encounter SHA-3 in PDF signatures today. No PDF
specification references it, no PAdES profile requires it, and no
mainstream PDF library defaults to it. It exists as insurance: if SHA-2
is ever broken, SHA-3 is ready to replace it without a multi-year
standardization process.

BouncyCastle 1.80 supports SHA-3 for signing, so the migration path
exists when needed.

\section{MD5: The Ghost in the
Machine}\label{_md5_the_ghost_in_the_machine}

MD5 produces a 128-bit (16-byte) digest. It was designed by Ronald
Rivest (the "R" in RSA) in 1991 and published as RFC 1321. It was once
the standard hash function for everything.

MD5 is \emph{broken}. Not theoretically, not in some academic
sense --- practically, completely, irrecoverably broken for collision
resistance.

\subsection{The Wang-Yu Attacks (2004)}\label{_the_wang_yu_attacks_2004}

In August 2004, Xiaoyun Wang and Hongbo Yu (the same Wang who later
weakened SHA-1) presented the first practical MD5 collision at CRYPTO
2004. They could find two distinct inputs with the same MD5 hash in less
than an hour on a standard PC. The collision required approximately
2\^{}39 MD5 operations --- billions of times easier than the expected
2\^{}64.

By 2006, Vlastimil Klima had reduced the attack to \emph{one minute} on
a single notebook computer. By 2008, Alexander Sotirov and Marc Stevens
demonstrated a rogue CA attack using MD5 collisions, creating a
fraudulent intermediate CA certificate that chained to a legitimate
root. This was not an academic exercise --- they obtained a real
certificate from a real CA (RapidSSL, which used MD5 for certificate
signing) and used MD5 collisions to create a rogue CA certificate that
any browser would trust.

Today, MD5 collisions can be generated in \emph{seconds} on commodity
hardware. The algorithm is fully broken for any security-sensitive
purpose.

\subsection{Why MD5 Still Appears in Legacy
PDFs}\label{_why_md5_still_appears_in_legacy_pdfs}

Despite being broken for collision resistance since 2004, you will
encounter MD5 in:

\begin{itemize}
\item
  Old PDF signatures created before 2005 using
  \texttt{adbe.pkcs7.detached} with MD5 as the digest algorithm
\item
  Legacy government and enterprise PDFs that were signed by systems
  deployed in the early 2000s
\item
  Some PKCS\#12 keystores that use MD5-based key derivation for the
  encryption of the private key (this is a different use of MD5 --- key
  derivation, not signature hashing --- and while not ideal, the risk
  profile is different from using MD5 for signing)
\item
  The \texttt{/ID} array in PDF trailers, which traditionally contained
  MD5 hashes of file metadata (this is not security-sensitive)
\end{itemize}

If you encounter a PDF signed with MD5 as the digest algorithm, the
signature provides \emph{no} collision resistance. An attacker who
controlled the document before signing could have prepared a collision.
For documents signed by trusted parties before 2004, the practical risk
is limited (the attack was not known yet). For anything signed after
2004, the signature's integrity guarantee is meaningless against a
determined attacker.

\section{The Signing Process Step by
Step}\label{_the_signing_process_step_by_step}

Now we connect the hash function to the asymmetric cryptography from
Chapter 1 and trace the exact sequence of operations that produces a PDF
digital signature.

\subsection{Step 1: Hash the Document}\label{_step_1_hash_the_document}

The signing library reads the bytes of the PDF specified by the
\texttt{/ByteRange} array. This is \emph{not} the entire file --- it is
everything except the \texttt{/Contents} hex string, which is where the
signature will be placed. The gap exists because you cannot include the
signature value in the data being signed (you would need to know the
signature before computing it --- a circular dependency).

The library feeds these bytes through the chosen hash algorithm:

\begin{verbatim}
SHA-256( byteRange[0..offset1] || byteRange[offset2..end] ) = digest (32 bytes)
\end{verbatim}

This 32-byte digest is the \emph{fingerprint} of the document at the
moment of signing.

\subsection{Step 2: Pad the Hash}\label{_step_2_pad_the_hash}

Before the hash can be signed with the private key, it must be padded to
the correct size for the signing algorithm. For RSA, there are two
padding schemes:

\textbf{PKCS\#1 v1.5 padding} (the traditional scheme): The hash is
prefixed with a DER-encoded \texttt{DigestInfo} structure that
identifies the hash algorithm, then padded with \texttt{0xFF} bytes to
fill the modulus size. The structure looks like:

\begin{Shaded}
\begin{Highlighting}[]
\NormalTok{DigestInfo ::= }\DataTypeTok{SEQUENCE}\NormalTok{ \{}
\NormalTok{    digestAlgorithm  AlgorithmIdentifier,  }\CommentTok{{-}{-} OID for SHA{-}256}
\NormalTok{    digest           }\DataTypeTok{OCTET} \DataTypeTok{STRING}           \CommentTok{{-}{-} the 32{-}byte hash}
\NormalTok{\}}
\end{Highlighting}
\end{Shaded}

The padded message is:
\texttt{0x00\ 0x01\ {[}0xFF\ padding{]}\ 0x00\ {[}DigestInfo\ DER\ bytes{]}}.

\textbf{RSASSA-PSS} (Probabilistic Signature Scheme): A more modern
padding scheme that includes random salt, making signatures
non-deterministic. PSS provides a tighter security proof than PKCS\#1
v1.5, but it is less commonly used in PDF signatures because the entire
PDF ecosystem was built around PKCS\#1 v1.5. Some CAs issue PSS
certificates, and you may encounter the OID
\texttt{1.2.840.113549.1.1.10} (RSASSA-PSS) in the
\texttt{SignerInfo.signatureAlgorithm} field.

For ECDSA, there is no padding step --- the hash is used directly as
input to the ECDSA signing equation.

The choice of padding scheme is not just a technical detail. PKCS\#1
v1.5 has known theoretical weaknesses (Bleichenbacher's 1998 attack on
encryption; related issues apply to signatures in certain implementation
scenarios). RSASSA-PSS is provably secure under standard assumptions. In
practice, PKCS\#1 v1.5 remains safe for signatures (as opposed to
encryption, where the Bleichenbacher attack is devastating), but
security-conscious deployments should prefer PSS. The ETSI EN 319 142-1
V1.1.1 (2016-04) standards allow both.

\subsection{Step 3: Encrypt with Private
Key}\label{_step_3_encrypt_with_private_key}

The padded hash (for RSA) or raw hash (for ECDSA) is processed with the
private key:

For RSA: \emph{s = paddedHash\textsuperscript{d} mod n}, where \emph{d}
is the private exponent and \emph{n} is the modulus. The result \emph{s}
is the signature value --- a number the same size as the modulus (256
bytes for RSA-2048, 384 bytes for RSA-3072).

For ECDSA: The algorithm generates a nonce \emph{k} (either random or
deterministic per RFC 6979), computes a curve point \emph{(x1, y1) = k *
G}, and derives the signature pair \emph{(r, s)} from \emph{x1}, the
hash, and the private key \emph{d}. The signature is the pair \emph{(r,
s)}, each approximately the same size as the curve order (32 bytes each
for P-256, totaling 64 bytes).

\subsection{Step 4: Package as CMS
SignedData}\label{_step_4_package_as_cms_signeddata}

The raw signature value is not written directly into the PDF. It is
wrapped in a CMS \texttt{SignedData} structure (RFC 5652 ---
Cryptographic Message Syntax (CMS) Section 5) that includes:

\begin{itemize}
\item
  The signature value itself
\item
  The signer's certificate (and often the full certificate chain)
\item
  The \texttt{SignerInfo} containing the digest algorithm, signature
  algorithm, and signed attributes
\item
  The \texttt{messageDigest} signed attribute --- the hash from Step 1
\item
  The \texttt{signingTime} attribute (optional but common)
\item
  The \texttt{contentType} attribute (always \texttt{id-data} for PDF)
\end{itemize}

This CMS blob is DER-encoded, then hex-encoded (each byte becomes two
hex characters), and written into the \texttt{/Contents} entry of the
signature dictionary.

\section{The Verification Process Step by
Step}\label{_the_verification_process_step_by_step}

Verification is the mirror image of signing, with one crucial addition:
the verifier must independently recompute the hash and compare it
against what the signer claims to have signed.

\subsection{Step 1: Extract the CMS
SignedData}\label{_step_1_extract_the_cms_signeddata}

The validator reads the \texttt{/Contents} hex string from the signature
dictionary and decodes it from hex to DER to obtain the CMS
\texttt{SignedData} structure.

\subsection{Step 2: Re-Hash the Byte
Range}\label{_step_2_re_hash_the_byte_range}

The validator reads the \texttt{/ByteRange} array, extracts the
specified byte ranges from the PDF file, and computes the hash using the
algorithm specified in the \texttt{SignerInfo.digestAlgorithm} field.

This is the critical integrity check: if anyone has modified the PDF
bytes after signing, this hash will differ from the original.

\subsection{Step 3: Compare Hashes}\label{_step_3_compare_hashes}

The validator extracts the \texttt{messageDigest} attribute from the CMS
\texttt{SignerInfo} signed attributes. This attribute contains the hash
that the signer computed over the byte range content. The validator
compares this with the independently computed hash from Step 2.

If they do not match, the document has been modified and validation
fails immediately. There is no need to check the cryptographic
signature.

\subsection{Step 4: Verify the Cryptographic
Signature}\label{_step_4_verify_the_cryptographic_signature}

If the hashes match, the validator checks that the signature over the
signed attributes (which include the \texttt{messageDigest} attribute)
is cryptographically valid. For RSA: compute \emph{m\textquotesingle{} =
s\textsuperscript{e} mod n} and verify that the result matches the
padded hash of the signed attributes. For ECDSA: verify the \emph{(r,
s)} pair against the hash using the public key point \emph{Q}.

If this check passes, the signature is \textbf{cryptographically
valid} --- it was produced by someone holding the private key that
corresponds to the public key in the certificate.

Cryptographic validity is necessary but not sufficient. A signature can
be cryptographically valid and still be \emph{untrustworthy}:

\begin{itemize}
\item
  The certificate might be expired.
\item
  The certificate might have been revoked.
\item
  The certificate might not chain to a trusted root.
\item
  The signed byte range might not cover the entire document (the Bochum
  attacks exploited exactly this --- see Chapter 24).
\item
  The certificate might lack the \texttt{digitalSignature} or
  \texttt{nonRepudiation} key usage extension (see Chapter 3).
\end{itemize}

Signature \textbf{validation} --- as opposed to mere cryptographic
verification --- requires checking all of these conditions. Chapter 16
covers the full validation process.

\section{DigestMethod in PDF
Signatures}\label{_digestmethod_in_pdf_signatures}

The hash algorithm used for a PDF signature is determined by the
\texttt{/SubFilter} value and the CMS
\texttt{SignerInfo.digestAlgorithm} field.

\begin{longtable}[]{@{}
  >{\raggedright\arraybackslash}p{(\linewidth - 4\tabcolsep) * \real{0.2857}}
  >{\raggedright\arraybackslash}p{(\linewidth - 4\tabcolsep) * \real{0.2857}}
  >{\raggedright\arraybackslash}p{(\linewidth - 4\tabcolsep) * \real{0.4286}}@{}}
\caption{How the hash algorithm is specified}\tabularnewline
\toprule\noalign{}
\begin{minipage}[b]{\linewidth}\raggedright
SubFilter
\end{minipage} & \begin{minipage}[b]{\linewidth}\raggedright
Hash Determination
\end{minipage} & \begin{minipage}[b]{\linewidth}\raggedright
Notes
\end{minipage} \\
\midrule\noalign{}
\endfirsthead
\toprule\noalign{}
\begin{minipage}[b]{\linewidth}\raggedright
SubFilter
\end{minipage} & \begin{minipage}[b]{\linewidth}\raggedright
Hash Determination
\end{minipage} & \begin{minipage}[b]{\linewidth}\raggedright
Notes
\end{minipage} \\
\midrule\noalign{}
\endhead
\bottomrule\noalign{}
\endlastfoot
\texttt{adbe.pkcs7.detached} & CMS \texttt{digestAlgorithm} field & The
hash algorithm is specified inside the CMS structure. Can be any
algorithm the signer chooses (SHA-256, SHA-384, SHA-512, or legacy
SHA-1/MD5). \\
\texttt{adbe.pkcs7.sha1} & Always SHA-1 & Deprecated by ISO
32000-2:2020. The SHA-1 hash of the entire byte range is embedded as the
CMS content. Do not create new signatures with this SubFilter. \\
\texttt{ETSI.CAdES.detached} & CMS \texttt{digestAlgorithm} field &
PAdES-compliant signatures. Must use SHA-256 or stronger per ETSI EN 319
142-1 V1.1.1 (2016-04). This is the recommended SubFilter for new
signatures. \\
\texttt{ETSI.RFC3161} & CMS \texttt{digestAlgorithm} field &
Document-level timestamp tokens (RFC 3161 --- Internet X.509 PKI
Time-Stamp Protocol (TSP)). The hash algorithm is specified in the
timestamp request. \\
\end{longtable}

\section{Hash Algorithm OIDs}\label{_hash_algorithm_oids}

When a hash algorithm appears inside a CMS SignedData structure, it is
identified by an ASN.1 Object Identifier (OID). You need to know these
OIDs because they appear in the raw bytes of every PDF signature, and
when something goes wrong, you will be reading hex dumps.

\begin{longtable}[]{@{}
  >{\raggedright\arraybackslash}p{(\linewidth - 4\tabcolsep) * \real{0.2857}}
  >{\raggedright\arraybackslash}p{(\linewidth - 4\tabcolsep) * \real{0.4286}}
  >{\raggedright\arraybackslash}p{(\linewidth - 4\tabcolsep) * \real{0.2857}}@{}}
\caption{Hash algorithm OIDs used in CMS (RFC 5652 --- Cryptographic
Message Syntax (CMS) and NIST)}\tabularnewline
\toprule\noalign{}
\begin{minipage}[b]{\linewidth}\raggedright
Algorithm
\end{minipage} & \begin{minipage}[b]{\linewidth}\raggedright
OID
\end{minipage} & \begin{minipage}[b]{\linewidth}\raggedright
Notes
\end{minipage} \\
\midrule\noalign{}
\endfirsthead
\toprule\noalign{}
\begin{minipage}[b]{\linewidth}\raggedright
Algorithm
\end{minipage} & \begin{minipage}[b]{\linewidth}\raggedright
OID
\end{minipage} & \begin{minipage}[b]{\linewidth}\raggedright
Notes
\end{minipage} \\
\midrule\noalign{}
\endhead
\bottomrule\noalign{}
\endlastfoot
MD5 & \texttt{1.2.840.113549.2.5} & Legacy only. Do not use for new
signatures. \\
SHA-1 & \texttt{1.3.14.3.2.26} & Deprecated. Still accepted by some
validators for backward compatibility. \\
SHA-224 & \texttt{2.16.840.1.101.3.4.2.4} & Rarely used. Exists for
completeness. \\
SHA-256 & \texttt{2.16.840.1.101.3.4.2.1} & The default for modern PDF
signatures. PAdES minimum. \\
SHA-384 & \texttt{2.16.840.1.101.3.4.2.2} & Used in higher-security
contexts. \\
SHA-512 & \texttt{2.16.840.1.101.3.4.2.3} & Maximum SHA-2 strength. \\
SHA3-256 & \texttt{2.16.840.1.101.3.4.2.8} & Future use. Not currently
seen in PDF signatures. \\
SHA3-384 & \texttt{2.16.840.1.101.3.4.2.9} & Future use. \\
SHA3-512 & \texttt{2.16.840.1.101.3.4.2.10} & Future use. \\
\end{longtable}

The OID \texttt{2.16.840.1.101.3.4.2.1} (SHA-256) is the one you will
see most often. The prefix \texttt{2.16.840.1.101.3.4} is the NIST
algorithm arc. The \texttt{.2} subtree is hash algorithms, and
\texttt{.1} is SHA-256 specifically. Memorize this OID --- you will
recognize it instantly in ASN.1 dumps and hex views.

The OID structure \texttt{1.2.840.113549} is the RSA Security arc (the
company, now part of EMC/Dell). MD5 (\texttt{…\hspace{0pt}2.5}) and
various PKCS algorithms live under this prefix. The OID
\texttt{1.3.14.3.2.26} for SHA-1 is under the OIW (Open Information
Warehouse) arc --- a historical artifact from the 1990s standardization
process.

\section{Signature Algorithm OIDs}\label{_signature_algorithm_oids}

The signing operation itself also has OIDs, which combine the hash
algorithm with the public-key algorithm:

\begin{longtable}[]{@{}
  >{\raggedright\arraybackslash}p{(\linewidth - 4\tabcolsep) * \real{0.3750}}
  >{\raggedright\arraybackslash}p{(\linewidth - 4\tabcolsep) * \real{0.3750}}
  >{\raggedright\arraybackslash}p{(\linewidth - 4\tabcolsep) * \real{0.2500}}@{}}
\caption{Common signature algorithm OIDs in PDF
signatures}\tabularnewline
\toprule\noalign{}
\begin{minipage}[b]{\linewidth}\raggedright
Algorithm Combination
\end{minipage} & \begin{minipage}[b]{\linewidth}\raggedright
OID
\end{minipage} & \begin{minipage}[b]{\linewidth}\raggedright
Notes
\end{minipage} \\
\midrule\noalign{}
\endfirsthead
\toprule\noalign{}
\begin{minipage}[b]{\linewidth}\raggedright
Algorithm Combination
\end{minipage} & \begin{minipage}[b]{\linewidth}\raggedright
OID
\end{minipage} & \begin{minipage}[b]{\linewidth}\raggedright
Notes
\end{minipage} \\
\midrule\noalign{}
\endhead
\bottomrule\noalign{}
\endlastfoot
SHA-256 with RSA & \texttt{1.2.840.113549.1.1.11} & Most common
combination in current PDF signatures. \\
SHA-384 with RSA & \texttt{1.2.840.113549.1.1.12} & Higher-security
RSA. \\
SHA-512 with RSA & \texttt{1.2.840.113549.1.1.13} & Maximum SHA-2 +
RSA. \\
SHA-256 with ECDSA & \texttt{1.2.840.10045.4.3.2} & Growing in usage.
Smaller signatures than RSA. \\
SHA-384 with ECDSA & \texttt{1.2.840.10045.4.3.3} & Common with P-384
curves. \\
SHA-1 with RSA & \texttt{1.2.840.113549.1.1.5} & Deprecated. Legacy
signatures only. \\
RSASSA-PSS & \texttt{1.2.840.113549.1.1.10} & Probabilistic RSA. The
hash algorithm is specified in the parameters, not in the OID itself. \\
\end{longtable}

When parsing CMS SignedData from a PDF signature, you may encounter a
mismatch between the \texttt{digestAlgorithm} in the SignerInfo (which
specifies the hash algorithm) and the hash algorithm implied by the
\texttt{signatureAlgorithm} OID. For example, the
\texttt{digestAlgorithm} might say SHA-256, but the
\texttt{signatureAlgorithm} OID might be \texttt{1.2.840.113549.1.1.1}
(RSA encryption, with no hash specified). This is technically allowed by
RFC 5652 --- Cryptographic Message Syntax (CMS) (the signature algorithm
can be "RSA" and the hash is determined by the \texttt{digestAlgorithm}
field), but it confuses many parsers. Always use the
\texttt{digestAlgorithm} field as the authoritative source for which
hash was used.

\section{When Hash and Signature Algorithms
Collide}\label{_when_hash_and_signature_algorithms_collide}

An interesting question arises: what happens when you pair a hash
algorithm with a signature algorithm designed for a different security
level? For example, SHA-256 (128-bit security) with ECDSA P-384 (192-bit
security)?

The answer is straightforward: \textbf{the overall security level is the
minimum of the two components}. SHA-256 with P-384 provides 128-bit
security, not 192-bit. The hash function is the bottleneck.

Standard pairings:

\begin{itemize}
\item
  \textbf{SHA-256 with RSA-3072 or ECDSA P-256}: 128-bit security across
  the board. The default for most use cases.
\item
  \textbf{SHA-384 with RSA-7680 or ECDSA P-384}: 192-bit security.
  Government/defense contexts.
\item
  \textbf{SHA-512 with ECDSA P-521}: 256-bit security. Rarely needed.
\end{itemize}

Mismatched pairings are not \emph{wrong} --- they work correctly --- but
they waste resources. Using SHA-512 with RSA-2048 gives you 112-bit
security (limited by RSA) while computing an unnecessarily expensive
hash. Using SHA-256 with P-384 gives you 128-bit security (limited by
SHA-256) while using a curve that provides more security than the hash
can deliver.

In practice, almost everyone uses SHA-256 with either RSA-2048/3072 or
ECDSA P-256. The matched security levels make the overall system
coherent, and SHA-256 is fast, universally supported, and has no known
weaknesses. Do not overthink this unless you have specific compliance
requirements that mandate higher security levels.

\section{Java: Computing Hashes}\label{ch02-java-hashing}

The \texttt{java.security.MessageDigest} class is Java's built-in hash
computation API. For basic hashing, you do not need BouncyCastle --- the
JDK provides SHA-256 and SHA-512 out of the box.

\begin{Shaded}
\begin{Highlighting}[]
\KeywordTok{import} \ImportTok{java}\OperatorTok{.}\ImportTok{security}\OperatorTok{.}\ImportTok{MessageDigest}\OperatorTok{;}
\KeywordTok{import} \ImportTok{java}\OperatorTok{.}\ImportTok{security}\OperatorTok{.}\ImportTok{NoSuchAlgorithmException}\OperatorTok{;}

\KeywordTok{public} \KeywordTok{class}\NormalTok{ HashExample }\OperatorTok{\{}

    \CommentTok{/**}
     \CommentTok{*}\NormalTok{ Compute the SHA}\CommentTok{{-}256}\NormalTok{ hash of the given data}\CommentTok{.}
     \CommentTok{*}\NormalTok{ Returns the raw }\CommentTok{32{-}}\NormalTok{byte digest}\CommentTok{.}
     \CommentTok{*/}
    \KeywordTok{public} \DataTypeTok{static} \DataTypeTok{byte}\OperatorTok{[]} \FunctionTok{sha256}\OperatorTok{(}\DataTypeTok{byte}\OperatorTok{[]}\NormalTok{ data}\OperatorTok{)} \KeywordTok{throws} \BuiltInTok{NoSuchAlgorithmException} \OperatorTok{\{}
        \BuiltInTok{MessageDigest}\NormalTok{ md }\OperatorTok{=} \BuiltInTok{MessageDigest}\OperatorTok{.}\FunctionTok{getInstance}\OperatorTok{(}\StringTok{"SHA{-}256"}\OperatorTok{);}
        \ControlFlowTok{return}\NormalTok{ md}\OperatorTok{.}\FunctionTok{digest}\OperatorTok{(}\NormalTok{data}\OperatorTok{);}
    \OperatorTok{\}}

    \CommentTok{/**}
     \CommentTok{*}\NormalTok{ Compute a hash using a specified algorithm}\CommentTok{.}
     \CommentTok{*}\NormalTok{ Algorithm names}\CommentTok{:} \CommentTok{"}\NormalTok{SHA}\CommentTok{{-}256",} \CommentTok{"}\NormalTok{SHA}\CommentTok{{-}384",} \CommentTok{"}\NormalTok{SHA}\CommentTok{{-}512",} \CommentTok{"}\NormalTok{SHA}\CommentTok{{-}1",} \CommentTok{"}\NormalTok{MD5}\CommentTok{".}
     \CommentTok{*/}
    \KeywordTok{public} \DataTypeTok{static} \DataTypeTok{byte}\OperatorTok{[]} \FunctionTok{hash}\OperatorTok{(}\DataTypeTok{byte}\OperatorTok{[]}\NormalTok{ data}\OperatorTok{,} \BuiltInTok{String}\NormalTok{ algorithm}\OperatorTok{)}
            \KeywordTok{throws} \BuiltInTok{NoSuchAlgorithmException} \OperatorTok{\{}
        \BuiltInTok{MessageDigest}\NormalTok{ md }\OperatorTok{=} \BuiltInTok{MessageDigest}\OperatorTok{.}\FunctionTok{getInstance}\OperatorTok{(}\NormalTok{algorithm}\OperatorTok{);}
        \ControlFlowTok{return}\NormalTok{ md}\OperatorTok{.}\FunctionTok{digest}\OperatorTok{(}\NormalTok{data}\OperatorTok{);}
    \OperatorTok{\}}

    \CommentTok{/**}
     \CommentTok{*}\NormalTok{ Demonstrate the avalanche effect}\CommentTok{:}\NormalTok{ one}\CommentTok{{-}}\NormalTok{bit change produces}
     \CommentTok{*}\NormalTok{ a completely different hash}\CommentTok{.}
     \CommentTok{*/}
    \KeywordTok{public} \DataTypeTok{static} \DataTypeTok{void} \FunctionTok{demonstrateAvalanche}\OperatorTok{()} \KeywordTok{throws} \BuiltInTok{NoSuchAlgorithmException} \OperatorTok{\{}
        \DataTypeTok{byte}\OperatorTok{[]}\NormalTok{ original }\OperatorTok{=} \StringTok{"The quick brown fox jumps over the lazy dog"}\OperatorTok{.}\FunctionTok{getBytes}\OperatorTok{();}
        \DataTypeTok{byte}\OperatorTok{[]}\NormalTok{ modified }\OperatorTok{=} \StringTok{"The quick brown fox jumps over the lazy dog."}\OperatorTok{.}\FunctionTok{getBytes}\OperatorTok{();} 

        \DataTypeTok{byte}\OperatorTok{[]}\NormalTok{ hash1 }\OperatorTok{=} \FunctionTok{sha256}\OperatorTok{(}\NormalTok{original}\OperatorTok{);}
        \DataTypeTok{byte}\OperatorTok{[]}\NormalTok{ hash2 }\OperatorTok{=} \FunctionTok{sha256}\OperatorTok{(}\NormalTok{modified}\OperatorTok{);}

        \CommentTok{// hash1 and hash2 will be completely different despite}
        \CommentTok{// the inputs differing by a single period character}
    \OperatorTok{\}}
\OperatorTok{\}}
\end{Highlighting}
\end{Shaded}

\begin{itemize}
\item
  The only difference is a trailing period. The SHA-256 outputs differ
  in approximately 50\% of their bits.
\end{itemize}

For PDF signing, you typically do not call \texttt{MessageDigest}
directly. Instead, the signing library (PDFBox + BouncyCastle) computes
the hash as part of the CMS signing operation. But knowing how to
compute hashes independently is essential for debugging --- when a
signature fails validation, one of the first things to check is whether
the hash of the byte range matches the \texttt{messageDigest} attribute
in the CMS SignerInfo.

\section{Java: Signing with
BouncyCastle}\label{_java_signing_with_bouncycastle}

Here is the complete signing flow using BouncyCastle 1.80. This code
signs a byte array and produces a CMS SignedData structure --- exactly
what ends up in a PDF signature's \texttt{/Contents} value.

\begin{Shaded}
\begin{Highlighting}[]
\KeywordTok{import} \ImportTok{java}\OperatorTok{.}\ImportTok{security}\OperatorTok{.}\ImportTok{KeyStore}\OperatorTok{;}
\KeywordTok{import} \ImportTok{java}\OperatorTok{.}\ImportTok{security}\OperatorTok{.}\ImportTok{PrivateKey}\OperatorTok{;}
\KeywordTok{import} \ImportTok{java}\OperatorTok{.}\ImportTok{security}\OperatorTok{.}\ImportTok{Security}\OperatorTok{;}
\KeywordTok{import} \ImportTok{java}\OperatorTok{.}\ImportTok{security}\OperatorTok{.}\ImportTok{cert}\OperatorTok{.}\ImportTok{Certificate}\OperatorTok{;}
\KeywordTok{import} \ImportTok{java}\OperatorTok{.}\ImportTok{security}\OperatorTok{.}\ImportTok{cert}\OperatorTok{.}\ImportTok{X509Certificate}\OperatorTok{;}
\KeywordTok{import} \ImportTok{java}\OperatorTok{.}\ImportTok{util}\OperatorTok{.}\ImportTok{ArrayList}\OperatorTok{;}
\KeywordTok{import} \ImportTok{java}\OperatorTok{.}\ImportTok{util}\OperatorTok{.}\ImportTok{List}\OperatorTok{;}

\KeywordTok{import} \ImportTok{org}\OperatorTok{.}\ImportTok{bouncycastle}\OperatorTok{.}\ImportTok{cert}\OperatorTok{.}\ImportTok{jcajce}\OperatorTok{.}\ImportTok{JcaCertStore}\OperatorTok{;}
\KeywordTok{import} \ImportTok{org}\OperatorTok{.}\ImportTok{bouncycastle}\OperatorTok{.}\ImportTok{cms}\OperatorTok{.}\ImportTok{CMSProcessableByteArray}\OperatorTok{;}
\KeywordTok{import} \ImportTok{org}\OperatorTok{.}\ImportTok{bouncycastle}\OperatorTok{.}\ImportTok{cms}\OperatorTok{.}\ImportTok{CMSSignedData}\OperatorTok{;}
\KeywordTok{import} \ImportTok{org}\OperatorTok{.}\ImportTok{bouncycastle}\OperatorTok{.}\ImportTok{cms}\OperatorTok{.}\ImportTok{CMSSignedDataGenerator}\OperatorTok{;}
\KeywordTok{import} \ImportTok{org}\OperatorTok{.}\ImportTok{bouncycastle}\OperatorTok{.}\ImportTok{cms}\OperatorTok{.}\ImportTok{jcajce}\OperatorTok{.}\ImportTok{JcaSignerInfoGeneratorBuilder}\OperatorTok{;}
\KeywordTok{import} \ImportTok{org}\OperatorTok{.}\ImportTok{bouncycastle}\OperatorTok{.}\ImportTok{jce}\OperatorTok{.}\ImportTok{provider}\OperatorTok{.}\ImportTok{BouncyCastleProvider}\OperatorTok{;}
\KeywordTok{import} \ImportTok{org}\OperatorTok{.}\ImportTok{bouncycastle}\OperatorTok{.}\ImportTok{operator}\OperatorTok{.}\ImportTok{ContentSigner}\OperatorTok{;}
\KeywordTok{import} \ImportTok{org}\OperatorTok{.}\ImportTok{bouncycastle}\OperatorTok{.}\ImportTok{operator}\OperatorTok{.}\ImportTok{jcajce}\OperatorTok{.}\ImportTok{JcaContentSignerBuilder}\OperatorTok{;}
\KeywordTok{import} \ImportTok{org}\OperatorTok{.}\ImportTok{bouncycastle}\OperatorTok{.}\ImportTok{operator}\OperatorTok{.}\ImportTok{jcajce}\OperatorTok{.}\ImportTok{JcaDigestCalculatorProviderBuilder}\OperatorTok{;}

\KeywordTok{public} \KeywordTok{class}\NormalTok{ SigningExample }\OperatorTok{\{}

    \DataTypeTok{static} \OperatorTok{\{}
        \BuiltInTok{Security}\OperatorTok{.}\FunctionTok{addProvider}\OperatorTok{(}\KeywordTok{new} \FunctionTok{BouncyCastleProvider}\OperatorTok{());}
    \OperatorTok{\}}

    \CommentTok{/**}
     \CommentTok{*}\NormalTok{ Sign data and return DER}\CommentTok{{-}}\NormalTok{encoded CMS SignedData}\CommentTok{.}
     \CommentTok{*}
     \CommentTok{*} \AnnotationTok{@param}\CommentTok{ }\AnnotationTok{dataToSign}\CommentTok{     }\NormalTok{the raw bytes to sign }\CommentTok{(}\NormalTok{byte range content}\CommentTok{)}
     \CommentTok{*} \AnnotationTok{@param}\CommentTok{ }\AnnotationTok{keyStore}\CommentTok{       }\NormalTok{loaded PKCS}\CommentTok{\#12}\NormalTok{ key store}
     \CommentTok{*} \AnnotationTok{@param}\CommentTok{ }\AnnotationTok{alias}\CommentTok{          }\NormalTok{key entry alias}
     \CommentTok{*} \AnnotationTok{@param}\CommentTok{ }\AnnotationTok{password}\CommentTok{       }\NormalTok{key password}
     \CommentTok{*} \AnnotationTok{@param}\CommentTok{ }\AnnotationTok{hashAlgorithm}\CommentTok{  "}\NormalTok{SHA}\CommentTok{{-}256",} \CommentTok{"}\NormalTok{SHA}\CommentTok{{-}384",}\NormalTok{ or }\CommentTok{"}\NormalTok{SHA}\CommentTok{{-}512"}
     \CommentTok{*} \AnnotationTok{@return}\NormalTok{               DER}\CommentTok{{-}}\NormalTok{encoded CMS SignedData}
     \CommentTok{*/}
    \KeywordTok{public} \DataTypeTok{static} \DataTypeTok{byte}\OperatorTok{[]} \FunctionTok{signData}\OperatorTok{(}\DataTypeTok{byte}\OperatorTok{[]}\NormalTok{ dataToSign}\OperatorTok{,} \BuiltInTok{KeyStore}\NormalTok{ keyStore}\OperatorTok{,}
                                   \BuiltInTok{String}\NormalTok{ alias}\OperatorTok{,} \DataTypeTok{char}\OperatorTok{[]}\NormalTok{ password}\OperatorTok{,}
                                   \BuiltInTok{String}\NormalTok{ hashAlgorithm}\OperatorTok{)} \KeywordTok{throws} \BuiltInTok{Exception} \OperatorTok{\{}

        \BuiltInTok{PrivateKey}\NormalTok{ privateKey }\OperatorTok{=} \OperatorTok{(}\BuiltInTok{PrivateKey}\OperatorTok{)}\NormalTok{ keyStore}\OperatorTok{.}\FunctionTok{getKey}\OperatorTok{(}\NormalTok{alias}\OperatorTok{,}\NormalTok{ password}\OperatorTok{);}    
        \BuiltInTok{Certificate}\OperatorTok{[]}\NormalTok{ chain }\OperatorTok{=}\NormalTok{ keyStore}\OperatorTok{.}\FunctionTok{getCertificateChain}\OperatorTok{(}\NormalTok{alias}\OperatorTok{);}
        \BuiltInTok{X509Certificate}\NormalTok{ signingCert }\OperatorTok{=} \OperatorTok{(}\BuiltInTok{X509Certificate}\OperatorTok{)}\NormalTok{ chain}\OperatorTok{[}\DecValTok{0}\OperatorTok{];}

        \BuiltInTok{List}\OperatorTok{\textless{}}\BuiltInTok{Certificate}\OperatorTok{\textgreater{}}\NormalTok{ certList }\OperatorTok{=} \KeywordTok{new} \BuiltInTok{ArrayList}\OperatorTok{\textless{}\textgreater{}();}                           
        \ControlFlowTok{for} \OperatorTok{(}\BuiltInTok{Certificate}\NormalTok{ cert }\OperatorTok{:}\NormalTok{ chain}\OperatorTok{)} \OperatorTok{\{}
\NormalTok{            certList}\OperatorTok{.}\FunctionTok{add}\OperatorTok{(}\NormalTok{cert}\OperatorTok{);}
        \OperatorTok{\}}

        \BuiltInTok{String}\NormalTok{ signatureAlgorithm }\OperatorTok{=}\NormalTok{ hashAlgorithm }\OperatorTok{+} \StringTok{"withRSA"}\OperatorTok{;}                   

\NormalTok{        ContentSigner signer }\OperatorTok{=} \KeywordTok{new} \FunctionTok{JcaContentSignerBuilder}\OperatorTok{(}\NormalTok{signatureAlgorithm}\OperatorTok{)}
                \OperatorTok{.}\FunctionTok{setProvider}\OperatorTok{(}\StringTok{"BC"}\OperatorTok{)}
                \OperatorTok{.}\FunctionTok{build}\OperatorTok{(}\NormalTok{privateKey}\OperatorTok{);}

\NormalTok{        CMSSignedDataGenerator generator }\OperatorTok{=} \KeywordTok{new} \FunctionTok{CMSSignedDataGenerator}\OperatorTok{();}
\NormalTok{        generator}\OperatorTok{.}\FunctionTok{addSignerInfoGenerator}\OperatorTok{(}
                \KeywordTok{new} \FunctionTok{JcaSignerInfoGeneratorBuilder}\OperatorTok{(}
                        \KeywordTok{new} \FunctionTok{JcaDigestCalculatorProviderBuilder}\OperatorTok{()}
                                \OperatorTok{.}\FunctionTok{setProvider}\OperatorTok{(}\StringTok{"BC"}\OperatorTok{)}
                                \OperatorTok{.}\FunctionTok{build}\OperatorTok{()}
                \OperatorTok{).}\FunctionTok{build}\OperatorTok{(}\NormalTok{signer}\OperatorTok{,}\NormalTok{ signingCert}\OperatorTok{)}
        \OperatorTok{);}
\NormalTok{        generator}\OperatorTok{.}\FunctionTok{addCertificates}\OperatorTok{(}\KeywordTok{new} \FunctionTok{JcaCertStore}\OperatorTok{(}\NormalTok{certList}\OperatorTok{));}                    

\NormalTok{        CMSProcessableByteArray content }\OperatorTok{=} \KeywordTok{new} \FunctionTok{CMSProcessableByteArray}\OperatorTok{(}\NormalTok{dataToSign}\OperatorTok{);}
\NormalTok{        CMSSignedData signedData }\OperatorTok{=}\NormalTok{ generator}\OperatorTok{.}\FunctionTok{generate}\OperatorTok{(}\NormalTok{content}\OperatorTok{,} \KeywordTok{false}\OperatorTok{);}            

        \ControlFlowTok{return}\NormalTok{ signedData}\OperatorTok{.}\FunctionTok{getEncoded}\OperatorTok{();}                                           
    \OperatorTok{\}}
\OperatorTok{\}}
\end{Highlighting}
\end{Shaded}

\begin{itemize}
\item
  Extract the private key from the key store. In production, use an HSM
  (see Chapter 5).
\item
  Include the full certificate chain so validators can build the path
  without fetching intermediates.
\item
  For ECDSA keys, use \texttt{hashAlgorithm\ +\ "withECDSA"} instead.
  The algorithm name must match the key type.
\item
  Embedding certificates in the CMS is critical --- without them, the
  validator has no chain to verify.
\item
  The \texttt{false} parameter means \emph{detached} signature: the
  content is not included in the CMS structure.
\item
  The returned bytes are DER-encoded and will be hex-encoded to fill
  \texttt{/Contents}.
\end{itemize}

The \texttt{false} parameter in
\texttt{generator.generate(content,\ false)} is the single most
important boolean in PDF signing. It means "detached signature" --- the
signed content is \emph{not} embedded in the CMS structure. PDF requires
detached signatures because the content (the byte range) lives in the
PDF file itself. If you pass \texttt{true}, the CMS will include the
content, and some validators will reject it because the
\texttt{/SubFilter} says \texttt{adbe.pkcs7.detached} but the CMS
contains encapsulated content.

\section{Debugging Hash Mismatches}\label{_debugging_hash_mismatches}

When a PDF signature fails validation, the most common root cause is a
hash mismatch: the hash of the byte range content does not match the
\texttt{messageDigest} signed attribute in the CMS SignerInfo.

Here is how to debug this systematically:

\begin{Shaded}
\begin{Highlighting}[]
\KeywordTok{import} \ImportTok{java}\OperatorTok{.}\ImportTok{security}\OperatorTok{.}\ImportTok{MessageDigest}\OperatorTok{;}

\KeywordTok{import} \ImportTok{org}\OperatorTok{.}\ImportTok{bouncycastle}\OperatorTok{.}\ImportTok{asn1}\OperatorTok{.}\ImportTok{ASN1OctetString}\OperatorTok{;}
\KeywordTok{import} \ImportTok{org}\OperatorTok{.}\ImportTok{bouncycastle}\OperatorTok{.}\ImportTok{asn1}\OperatorTok{.}\ImportTok{cms}\OperatorTok{.}\ImportTok{Attribute}\OperatorTok{;}
\KeywordTok{import} \ImportTok{org}\OperatorTok{.}\ImportTok{bouncycastle}\OperatorTok{.}\ImportTok{asn1}\OperatorTok{.}\ImportTok{cms}\OperatorTok{.}\ImportTok{CMSAttributes}\OperatorTok{;}
\KeywordTok{import} \ImportTok{org}\OperatorTok{.}\ImportTok{bouncycastle}\OperatorTok{.}\ImportTok{cms}\OperatorTok{.}\ImportTok{CMSSignedData}\OperatorTok{;}
\KeywordTok{import} \ImportTok{org}\OperatorTok{.}\ImportTok{bouncycastle}\OperatorTok{.}\ImportTok{cms}\OperatorTok{.}\ImportTok{SignerInformation}\OperatorTok{;}
\KeywordTok{import} \ImportTok{org}\OperatorTok{.}\ImportTok{bouncycastle}\OperatorTok{.}\ImportTok{util}\OperatorTok{.}\ImportTok{encoders}\OperatorTok{.}\ImportTok{Hex}\OperatorTok{;}

\KeywordTok{public} \KeywordTok{class}\NormalTok{ HashDebugExample }\OperatorTok{\{}

    \CommentTok{/**}
     \CommentTok{*}\NormalTok{ Compare the hash in the CMS messageDigest attribute against}
     \CommentTok{*}\NormalTok{ a freshly computed hash of the byte range content}\CommentTok{.}
     \CommentTok{*/}
    \KeywordTok{public} \DataTypeTok{static} \DataTypeTok{void} \FunctionTok{debugHashMismatch}\OperatorTok{(}\DataTypeTok{byte}\OperatorTok{[]}\NormalTok{ cmsBytes}\OperatorTok{,} \DataTypeTok{byte}\OperatorTok{[]}\NormalTok{ byteRangeContent}\OperatorTok{)}
            \KeywordTok{throws} \BuiltInTok{Exception} \OperatorTok{\{}

\NormalTok{        CMSSignedData signedData }\OperatorTok{=} \KeywordTok{new} \FunctionTok{CMSSignedData}\OperatorTok{(}\NormalTok{cmsBytes}\OperatorTok{);}

        \ControlFlowTok{for} \OperatorTok{(}\NormalTok{SignerInformation signerInfo }\OperatorTok{:}\NormalTok{ signedData}\OperatorTok{.}\FunctionTok{getSignerInfos}\OperatorTok{().}\FunctionTok{getSigners}\OperatorTok{())} \OperatorTok{\{}

            \CommentTok{// 1. Get the digest algorithm from the SignerInfo}
            \BuiltInTok{String}\NormalTok{ digestAlgOid }\OperatorTok{=}\NormalTok{ signerInfo}\OperatorTok{.}\FunctionTok{getDigestAlgOID}\OperatorTok{();}
            \BuiltInTok{String}\NormalTok{ digestAlgName }\OperatorTok{=} \FunctionTok{mapOidToName}\OperatorTok{(}\NormalTok{digestAlgOid}\OperatorTok{);}

            \CommentTok{// 2. Compute the hash of the byte range content independently}
            \BuiltInTok{MessageDigest}\NormalTok{ md }\OperatorTok{=} \BuiltInTok{MessageDigest}\OperatorTok{.}\FunctionTok{getInstance}\OperatorTok{(}\NormalTok{digestAlgName}\OperatorTok{);}
            \DataTypeTok{byte}\OperatorTok{[]}\NormalTok{ computedHash }\OperatorTok{=}\NormalTok{ md}\OperatorTok{.}\FunctionTok{digest}\OperatorTok{(}\NormalTok{byteRangeContent}\OperatorTok{);}

            \CommentTok{// 3. Extract the messageDigest attribute from signed attributes}
            \BuiltInTok{Attribute}\NormalTok{ mdAttr }\OperatorTok{=}\NormalTok{ signerInfo}\OperatorTok{.}\FunctionTok{getSignedAttributes}\OperatorTok{()}
                    \OperatorTok{.}\FunctionTok{get}\OperatorTok{(}\NormalTok{CMSAttributes}\OperatorTok{.}\FunctionTok{messageDigest}\OperatorTok{);}
            \DataTypeTok{byte}\OperatorTok{[]}\NormalTok{ embeddedHash }\OperatorTok{=}\NormalTok{ ASN1OctetString}\OperatorTok{.}\FunctionTok{getInstance}\OperatorTok{(}
\NormalTok{                    mdAttr}\OperatorTok{.}\FunctionTok{getAttrValues}\OperatorTok{().}\FunctionTok{getObjectAt}\OperatorTok{(}\DecValTok{0}\OperatorTok{)}
            \OperatorTok{).}\FunctionTok{getOctets}\OperatorTok{();}

            \CommentTok{// 4. Compare}
            \DataTypeTok{boolean}\NormalTok{ match }\OperatorTok{=} \BuiltInTok{MessageDigest}\OperatorTok{.}\FunctionTok{isEqual}\OperatorTok{(}\NormalTok{computedHash}\OperatorTok{,}\NormalTok{ embeddedHash}\OperatorTok{);}
            \CommentTok{// If match is false: the byte range was extracted incorrectly,}
            \CommentTok{// or the PDF was modified after signing}
        \OperatorTok{\}}
    \OperatorTok{\}}

    \KeywordTok{private} \DataTypeTok{static} \BuiltInTok{String} \FunctionTok{mapOidToName}\OperatorTok{(}\BuiltInTok{String}\NormalTok{ oid}\OperatorTok{)} \OperatorTok{\{}
        \ControlFlowTok{return} \ControlFlowTok{switch} \OperatorTok{(}\NormalTok{oid}\OperatorTok{)} \OperatorTok{\{}
            \ControlFlowTok{case} \StringTok{"2.16.840.1.101.3.4.2.1"} \OperatorTok{{-}\textgreater{}} \StringTok{"SHA{-}256"}\OperatorTok{;}
            \ControlFlowTok{case} \StringTok{"2.16.840.1.101.3.4.2.2"} \OperatorTok{{-}\textgreater{}} \StringTok{"SHA{-}384"}\OperatorTok{;}
            \ControlFlowTok{case} \StringTok{"2.16.840.1.101.3.4.2.3"} \OperatorTok{{-}\textgreater{}} \StringTok{"SHA{-}512"}\OperatorTok{;}
            \ControlFlowTok{case} \StringTok{"1.3.14.3.2.26"}           \OperatorTok{{-}\textgreater{}} \StringTok{"SHA{-}1"}\OperatorTok{;}
            \ControlFlowTok{case} \StringTok{"1.2.840.113549.2.5"}      \OperatorTok{{-}\textgreater{}} \StringTok{"MD5"}\OperatorTok{;}
            \KeywordTok{default} \OperatorTok{{-}\textgreater{}} \ControlFlowTok{throw} \KeywordTok{new} \BuiltInTok{IllegalArgumentException}\OperatorTok{(}\StringTok{"Unknown OID: "} \OperatorTok{+}\NormalTok{ oid}\OperatorTok{);}
        \OperatorTok{\};}
    \OperatorTok{\}}
\OperatorTok{\}}
\end{Highlighting}
\end{Shaded}

When debugging a hash mismatch, the most common causes are: (1) the byte
range was extracted incorrectly --- off by one byte, or not accounting
for the gap correctly; (2) the PDF was modified after signing (perhaps
by a tool that did a full save instead of incremental save); (3) the
byte range array values do not add up to cover the entire file (there
are bytes before or after the byte range that are not part of the signed
content but should be --- see Chapter 13 for the gory details).

\section{Summary}\label{_summary_2}

Digital signatures are built on two primitives: hash functions provide
integrity (any modification is detectable), and asymmetric signing
provides authentication (only the private key holder could have created
the signature) and non-repudiation (the signer cannot deny having
signed).

For PDF signatures:

\begin{itemize}
\item
  \textbf{Use SHA-256} as the digest algorithm. It is the PAdES minimum,
  universally supported, and has no known weaknesses.
\item
  \textbf{SHA-1 is deprecated.} The SHAttered attack proved practical
  collision attacks are feasible. Do not create new SHA-1 signatures.
\item
  \textbf{MD5 is broken.} Collisions can be generated in seconds. Legacy
  only.
\item
  \textbf{The hash is always computed over the byte range}, not the
  entire PDF file. The byte range excludes the \texttt{/Contents} value
  (which contains the signature itself --- you cannot sign your own
  signature).
\item
  \textbf{CMS detached signatures} are the standard: the content is not
  embedded in the CMS structure.
\item
  \textbf{PKCS\#1 v1.5 padding} is dominant in PDF signatures, but
  \textbf{RSASSA-PSS} is gaining adoption for new deployments.
\end{itemize}

The next chapter introduces X.509 certificates --- the mechanism that
binds a public key to an identity, turning a raw cryptographic operation
into a trust relationship.

\section{X.509 Certificates}\label{ch03}

\begin{quote}
A digital signature proves that someone with a specific private key
signed some data. But \emph{who} holds that private key? A public key is
just a number --- a very large number, but a number with no name
attached.

X.509 certificates solve this problem. A certificate is a signed
statement that binds a public key to an identity: "This public key
belongs to Alice Smith at Acme Corp, and I, DigiCert, vouch for it."
This chapter takes apart the X.509 certificate format, field by field,
extension by extension, with the level of detail you need to debug
real-world certificate problems in PDF signatures.
\end{quote}

\subsection{What a Certificate Is}\label{_what_a_certificate_is}

At its core, an X.509 certificate is a data structure containing three
things:

\begin{enumerate}
\def\labelenumi{\arabic{enumi}.}
\item
  A \textbf{subject} --- the identity (person, organization, or device)
  that the certificate represents
\item
  A \textbf{public key} --- the subject's public key, which can be used
  to verify signatures created with the corresponding private key
\item
  A \textbf{digital signature} from a \textbf{Certificate Authority
  (CA)} --- a trusted third party that vouches for the binding between
  the subject and the public key
\end{enumerate}

The CA's signature is what gives the certificate its authority. Without
it, you have a self-signed statement: "I am Alice, and here is my public
key, trust me." With it, you have a verifiable claim backed by an entity
that (presumably) verified Alice's identity before signing.

The X.509 standard is defined in ITU-T Recommendation X.509 (originally
from 1988, extensively revised since) and profiled for Internet use in
RFC 5280 --- Internet X.509 PKI Certificate and CRL Profile. When we
talk about X.509 certificates in the context of PDF signatures, we mean
the RFC 5280 --- Internet X.509 PKI Certificate and CRL Profile
profile --- the Internet PKI profile that every browser, every operating
system, and every PDF viewer actually implements.

\subsection{ASN.1 and DER Encoding}\label{_asn_1_and_der_encoding}

Before we look at the certificate structure, we need to talk about how
certificates are encoded, because you \emph{will} encounter raw
certificate bytes when debugging PDF signatures, and you need to
understand what you're looking at.

Certificates are defined using \textbf{ASN.1} (Abstract Syntax Notation
One), a schema language for describing data structures. ASN.1 is to
certificates what JSON Schema is to JSON --- it defines the structure,
types, and constraints. But unlike JSON, ASN.1 data is encoded in a
compact binary format.

The encoding used for certificates is \textbf{DER} (Distinguished
Encoding Rules), a subset of BER (Basic Encoding Rules) that produces
exactly one encoding for any given value. This uniqueness property is
critical: because a certificate's signature covers its DER encoding,
there can be only one valid encoding --- otherwise the signature
wouldn't verify.

\subsubsection{Tag-Length-Value}\label{_tag_length_value}

DER uses a Tag-Length-Value (TLV) structure. Every data element consists
of:

\begin{itemize}
\item
  \textbf{Tag} (1 or more bytes): identifies the type (SEQUENCE,
  INTEGER, BIT STRING, OID, etc.)
\item
  \textbf{Length} (1 or more bytes): the number of bytes in the value
\item
  \textbf{Value} (variable): the actual data
\end{itemize}

\begin{longtable}[]{@{}
  >{\raggedright\arraybackslash}p{(\linewidth - 4\tabcolsep) * \real{0.1667}}
  >{\raggedright\arraybackslash}p{(\linewidth - 4\tabcolsep) * \real{0.3333}}
  >{\raggedright\arraybackslash}p{(\linewidth - 4\tabcolsep) * \real{0.5000}}@{}}
\caption{Common ASN.1 tags in certificates}\tabularnewline
\toprule\noalign{}
\begin{minipage}[b]{\linewidth}\raggedright
Tag (hex)
\end{minipage} & \begin{minipage}[b]{\linewidth}\raggedright
ASN.1 Type
\end{minipage} & \begin{minipage}[b]{\linewidth}\raggedright
Where You'll See It
\end{minipage} \\
\midrule\noalign{}
\endfirsthead
\toprule\noalign{}
\begin{minipage}[b]{\linewidth}\raggedright
Tag (hex)
\end{minipage} & \begin{minipage}[b]{\linewidth}\raggedright
ASN.1 Type
\end{minipage} & \begin{minipage}[b]{\linewidth}\raggedright
Where You'll See It
\end{minipage} \\
\midrule\noalign{}
\endhead
\bottomrule\noalign{}
\endlastfoot
\texttt{30} & SEQUENCE (constructed) & The outermost wrapper of the
certificate, and of most sub-structures (tbsCertificate, extensions,
etc.) \\
\texttt{31} & SET (constructed) & Relative Distinguished Name components
in the issuer and subject fields \\
\texttt{02} & INTEGER & Serial number, RSA modulus and exponent, version
number \\
\texttt{03} & BIT STRING & The signature value, and the subject public
key value \\
\texttt{04} & OCTET STRING & Extension values, key identifiers \\
\texttt{05} & NULL & Algorithm parameters (RSA with SHA-256 has NULL
parameters) \\
\texttt{06} & OBJECT IDENTIFIER & Algorithm OIDs, extension OIDs, policy
OIDs \\
\texttt{0C} & UTF8String & Subject and issuer name components \\
\texttt{13} & PrintableString & Country codes and other
restricted-character fields \\
\texttt{17} & UTCTime & Validity dates (for dates before 2050) \\
\texttt{18} & GeneralizedTime & Validity dates (for dates from 2050
onward, per RFC 5280 --- Internet X.509 PKI Certificate and CRL Profile
§4.1.2.5) \\
\texttt{A0} & Context-specific {[}0{]} & Version number (explicit tag),
extensions container \\
\texttt{A3} & Context-specific {[}3{]} & Extensions container in
tbsCertificate \\
\end{longtable}

When you see \texttt{30\ 82\ 03\ A5} at the start of a binary blob,
you're looking at a SEQUENCE with a two-byte length (the \texttt{82}
means "the next 2 bytes are the length"), specifically \texttt{03\ A5} =
933 bytes. This is almost certainly a certificate or a CMS structure.
The \texttt{82} prefix on the length byte is the tell: \texttt{81} means
one-byte length follows, \texttt{82} means two-byte length follows.
Certificate structures are almost always in the \texttt{82} range
(258-65535 bytes).

\subsubsection{PEM vs. DER}\label{_pem_vs_der}

You will encounter certificates in two physical formats:

\begin{itemize}
\item
  \textbf{DER}: raw binary. File extensions: \texttt{.der},
  \texttt{.cer}, \texttt{.crt} (ambiguous)
\item
  \textbf{PEM}: Base64-encoded DER wrapped in
  \texttt{-\/-\/-\/-\/-BEGIN\ CERTIFICATE-\/-\/-\/-\/-} /
  \texttt{-\/-\/-\/-\/-END\ CERTIFICATE-\/-\/-\/-\/-} headers. File
  extensions: \texttt{.pem}, \texttt{.crt} (ambiguous)
\end{itemize}

\begin{verbatim}
-----BEGIN CERTIFICATE-----
MIIFYjCCBEqgAwIBAgIQd70NbNs2+RrqIQ/E8FjTDTANBgkqhkiG9w0BAQsFAMDQ
... (Base64-encoded DER bytes) ...
-----END CERTIFICATE-----
\end{verbatim}

The content between the headers is the DER encoding, Base64-encoded with
line breaks every 64 characters. Nothing more, nothing less. To get the
DER bytes, strip the headers and Base64-decode.

The \texttt{.crt} extension is ambiguous --- it can be either DER or
PEM. If you're loading a certificate and get a parsing error, try the
other format. Java's \texttt{CertificateFactory} handles both, but
BouncyCastle's \texttt{PEMParser} expects PEM, and raw
\texttt{X509CertificateHolder(byte{[}{]})} expects DER.

\subsection{Certificate Structure}\label{_certificate_structure}

An X.509v3 certificate has three top-level fields, defined in RFC 5280
--- Internet X.509 PKI Certificate and CRL Profile Section 4.1:

\begin{verbatim}
Certificate ::= SEQUENCE {
    tbsCertificate       TBSCertificate,         
    signatureAlgorithm   AlgorithmIdentifier,     
    signatureValue       BIT STRING               
}
\end{verbatim}

\begin{itemize}
\item
  The "to be signed" certificate --- all the data that the CA signs
\item
  The algorithm the CA used to create the signature (must match the
  algorithm in tbsCertificate)
\item
  The CA's digital signature over the DER encoding of tbsCertificate
\end{itemize}

The CA signs \texttt{tbsCertificate} --- the entire inner structure,
DER-encoded. The \texttt{signatureAlgorithm} and \texttt{signatureValue}
are outside the signed region (they can't be inside it --- the signature
can't sign itself, just as a PDF's \texttt{/Contents} can't be part of
its own byte range).

\subsubsection{tbsCertificate}\label{_tbscertificate}

This is where all the interesting data lives:

\begin{verbatim}
TBSCertificate ::= SEQUENCE {
    version          [0] EXPLICIT INTEGER DEFAULT v1,  
    serialNumber         INTEGER,                       
    signature            AlgorithmIdentifier,           
    issuer               Name,                          
    validity             Validity,                      
    subject              Name,                          
    subjectPublicKeyInfo SubjectPublicKeyInfo,          
    issuerUniqueID   [1] IMPLICIT BIT STRING OPTIONAL,  
    subjectUniqueID  [2] IMPLICIT BIT STRING OPTIONAL,  
    extensions       [3] EXPLICIT Extensions OPTIONAL    
}
\end{verbatim}

\begin{itemize}
\item
  Version: 0 = v1, 1 = v2, 2 = v3. All modern certificates are v3. If
  you see v1, it's either very old or self-signed for testing.
\item
  Serial number: unique within the issuing CA. Must be a positive
  integer, at most 20 bytes. CAs use random serial numbers to prevent
  hash-based attacks.
\item
  Signature algorithm: must match the outer \texttt{signatureAlgorithm}
  field. This redundancy is specified by RFC 5280 --- Internet X.509 PKI
  Certificate and CRL Profile to prevent algorithm substitution attacks.
\item
  Issuer: the Distinguished Name (DN) of the CA that signed this
  certificate.
\item
  Validity: \texttt{notBefore} and \texttt{notAfter} timestamps defining
  the certificate's validity period.
\item
  Subject: the Distinguished Name of the entity this certificate
  represents.
\item
  Subject public key: the algorithm and the actual public key bytes.
\item
  Issuer/subject unique IDs: deprecated, never used in practice. Ignore
  them.
\item
  Extensions: the v3 extensions that make certificates actually useful.
  See \hyperref[ch03-extensions]{Certificate Extensions} below.
\end{itemize}

\paragraph{Distinguished Names}\label{_distinguished_names}

The issuer and subject fields use X.500 Distinguished Names --- a
hierarchical naming scheme from the 1980s that looks like this:

\begin{verbatim}
CN=Alice Smith, O=Acme Corporation, L=San Francisco, ST=California, C=US
\end{verbatim}

Each component is a Relative Distinguished Name (RDN) identified by an
OID:

\begin{longtable}[]{@{}
  >{\raggedright\arraybackslash}p{(\linewidth - 4\tabcolsep) * \real{0.1667}}
  >{\raggedright\arraybackslash}p{(\linewidth - 4\tabcolsep) * \real{0.3333}}
  >{\raggedright\arraybackslash}p{(\linewidth - 4\tabcolsep) * \real{0.5000}}@{}}
\caption{Common RDN OIDs}\tabularnewline
\toprule\noalign{}
\begin{minipage}[b]{\linewidth}\raggedright
Abbreviation
\end{minipage} & \begin{minipage}[b]{\linewidth}\raggedright
OID
\end{minipage} & \begin{minipage}[b]{\linewidth}\raggedright
Meaning
\end{minipage} \\
\midrule\noalign{}
\endfirsthead
\toprule\noalign{}
\begin{minipage}[b]{\linewidth}\raggedright
Abbreviation
\end{minipage} & \begin{minipage}[b]{\linewidth}\raggedright
OID
\end{minipage} & \begin{minipage}[b]{\linewidth}\raggedright
Meaning
\end{minipage} \\
\midrule\noalign{}
\endhead
\bottomrule\noalign{}
\endlastfoot
CN & \texttt{2.5.4.3} & Common Name --- the certificate holder's name \\
O & \texttt{2.5.4.10} & Organization \\
OU & \texttt{2.5.4.11} & Organizational Unit (deprecated in TLS
certificates since 2022, but still used in document signing
certificates) \\
L & \texttt{2.5.4.7} & Locality (city) \\
ST & \texttt{2.5.4.8} & State or Province \\
C & \texttt{2.5.4.6} & Country (two-letter ISO 3166-1 code) \\
serialNumber & \texttt{2.5.4.5} & Subject serial number (not the
certificate serial number --- this is a personal or organizational
identifier, like a tax ID or VAT number) \\
E / emailAddress & \texttt{1.2.840.113549.1.9.1} & Email address
(legacy; modern certificates use the Subject Alternative Name extension
instead) \\
\end{longtable}

The DN ordering in X.509 certificates is \emph{most specific first} (CN
before O before C), which is the opposite of LDAP directory convention
(C before O before CN). This catches people off guard when comparing DNs
between different tools. Java's
\texttt{X509Certificate.getSubjectX500Principal().getName()} returns RFC
2253 format (CN first), while some tools show LDAP order (C first). When
comparing issuer DNs for chain building, always compare the DER
encoding, not the string representation.

\paragraph{Validity Period}\label{_validity_period}

\begin{verbatim}
Validity ::= SEQUENCE {
    notBefore   Time,    -- certificate becomes valid
    notAfter    Time     -- certificate expires
}
\end{verbatim}

RFC 5280 --- Internet X.509 PKI Certificate and CRL Profile §4.1.2.5
specifies that dates through 2049 use UTCTime (\texttt{YYMMDDHHMMSSZ})
and dates in 2050 and beyond use GeneralizedTime
(\texttt{YYYYMMDDHHMMSSZ}). This means every certificate issued today
uses UTCTime for both \texttt{notBefore} and
\texttt{notAfter} --- unless the certificate is valid until 2050 or
later (some root CA certificates are).

For PDF signatures, the validity period has a nuanced interpretation. A
signature made while the certificate was valid remains meaningful
\emph{after} the certificate expires --- the certificate was valid
\emph{at the time of signing}. This is why timestamps are critical: a
timestamp proves that the signature existed at a specific point in time,
allowing validation even after certificate expiration. Without a
timestamp, a validator must decide whether to check the certificate
against the current time (strict) or the claimed signing time (lenient).
Different validators make different choices, which is one of the many
reasons the same signature shows "valid" in one viewer and "invalid" in
another.

\subsection{Certificate Extensions}\label{ch03-extensions}

Extensions are what make X.509v3 certificates useful for real-world
applications. Without extensions, a certificate is just a name bound to
a public key --- no information about what the key can be used for,
whether the certificate can issue other certificates, or how to check if
it's been revoked.

Each extension has three components:

\begin{itemize}
\item
  \textbf{OID}: identifies the extension type
\item
  \textbf{Critical flag}: if true, a system that doesn't understand this
  extension \emph{must} reject the certificate
\item
  \textbf{Value}: the extension-specific data (an OCTET STRING
  containing DER-encoded content)
\end{itemize}

\subsubsection{KeyUsage}\label{_keyusage}

OID: \texttt{2.5.29.15} \textbar{} Defined in: RFC 5280 --- Internet
X.509 PKI Certificate and CRL Profile §4.2.1.3 \textbar{} Typically:
\textbf{Critical}

The KeyUsage extension is a bit string that specifies what the public
key in this certificate may be used for.

\begin{longtable}[]{@{}
  >{\raggedright\arraybackslash}p{(\linewidth - 4\tabcolsep) * \real{0.1667}}
  >{\raggedright\arraybackslash}p{(\linewidth - 4\tabcolsep) * \real{0.1667}}
  >{\raggedright\arraybackslash}p{(\linewidth - 4\tabcolsep) * \real{0.6667}}@{}}
\caption{KeyUsage bits relevant to PDF signing}\tabularnewline
\toprule\noalign{}
\begin{minipage}[b]{\linewidth}\raggedright
Bit
\end{minipage} & \begin{minipage}[b]{\linewidth}\raggedright
Name
\end{minipage} & \begin{minipage}[b]{\linewidth}\raggedright
Relevance to PDF Signing
\end{minipage} \\
\midrule\noalign{}
\endfirsthead
\toprule\noalign{}
\begin{minipage}[b]{\linewidth}\raggedright
Bit
\end{minipage} & \begin{minipage}[b]{\linewidth}\raggedright
Name
\end{minipage} & \begin{minipage}[b]{\linewidth}\raggedright
Relevance to PDF Signing
\end{minipage} \\
\midrule\noalign{}
\endhead
\bottomrule\noalign{}
\endlastfoot
0 & \texttt{digitalSignature} & \textbf{Required for PDF signing.}
Indicates the key can be used to verify digital signatures (other than
certificate and CRL signatures). This is the bit you need for signing
documents. \\
1 & \texttt{nonRepudiation} (contentCommitment) & Indicates the key is
used for signatures where the signer cannot later deny having signed.
Some certificate profiles for document signing set this bit instead of
or in addition to \texttt{digitalSignature}. The RFC 5280 --- Internet
X.509 PKI Certificate and CRL Profile renamed this to
\texttt{contentCommitment} but the old name persists in most tools. \\
5 & \texttt{keyCertSign} & Indicates the key is used to verify
signatures on certificates. Set on CA certificates. If this bit is set
on an end-entity certificate, something is wrong. \\
6 & \texttt{cRLSign} & Indicates the key is used to verify signatures on
CRLs. Set on CA certificates that issue CRLs. \\
\end{longtable}

Some document signing certificates have \emph{only}
\texttt{nonRepudiation} (bit 1) set, without \texttt{digitalSignature}
(bit 0). This is technically valid per RFC 5280 --- Internet X.509 PKI
Certificate and CRL Profile, but it causes problems: some validators
check for \texttt{digitalSignature} and reject certificates that only
have \texttt{nonRepudiation}. The safest approach for CAs issuing
document signing certificates is to set both bits. If you're validating
signatures, check for \emph{either} bit 0 or bit 1.

\subsubsection{ExtendedKeyUsage}\label{_extendedkeyusage}

OID: \texttt{2.5.29.37} \textbar{} Defined in: RFC 5280 --- Internet
X.509 PKI Certificate and CRL Profile §4.2.1.12 \textbar{} Typically:
\textbf{Non-critical}

While KeyUsage is a coarse-grained bitmask, ExtendedKeyUsage specifies
more precise purposes using OIDs.

\begin{longtable}[]{@{}
  >{\raggedright\arraybackslash}p{(\linewidth - 4\tabcolsep) * \real{0.3333}}
  >{\raggedright\arraybackslash}p{(\linewidth - 4\tabcolsep) * \real{0.2222}}
  >{\raggedright\arraybackslash}p{(\linewidth - 4\tabcolsep) * \real{0.4444}}@{}}
\caption{ExtendedKeyUsage OIDs relevant to PDF signing}\tabularnewline
\toprule\noalign{}
\begin{minipage}[b]{\linewidth}\raggedright
OID
\end{minipage} & \begin{minipage}[b]{\linewidth}\raggedright
Name
\end{minipage} & \begin{minipage}[b]{\linewidth}\raggedright
Meaning
\end{minipage} \\
\midrule\noalign{}
\endfirsthead
\toprule\noalign{}
\begin{minipage}[b]{\linewidth}\raggedright
OID
\end{minipage} & \begin{minipage}[b]{\linewidth}\raggedright
Name
\end{minipage} & \begin{minipage}[b]{\linewidth}\raggedright
Meaning
\end{minipage} \\
\midrule\noalign{}
\endhead
\bottomrule\noalign{}
\endlastfoot
\texttt{1.3.6.1.5.5.7.3.4} & \texttt{id-kp-emailProtection} &
Historically used for document signing certificates (S/MIME). Many older
PDF signing certificates have this EKU. \\
\texttt{1.3.6.1.4.1.311.10.3.12} & \texttt{id-kp-documentSigning}
(Microsoft) & Microsoft's document signing EKU. Common on certificates
from CAs that target Windows environments. \\
\texttt{1.2.840.113583.1.1.5} & Adobe Authentic Documents Trust &
Adobe-specific EKU for AATL certificates. If a certificate has this EKU
and chains to an AATL root, Acrobat shows maximum trust. \\
\texttt{1.3.6.1.5.5.7.3.9} & \texttt{id-kp-OCSPSigning} & For OCSP
responder certificates only. Not for document signing. \\
\end{longtable}

When a PDF signature shows as "untrusted" in Adobe Acrobat but validates
fine in other viewers, the first thing to check is the ExtendedKeyUsage.
Acrobat is stricter about EKU than most other validators. A certificate
without an explicit document signing EKU may be rejected by Acrobat even
if it has \texttt{digitalSignature} in KeyUsage.

\subsubsection{BasicConstraints}\label{_basicconstraints}

OID: \texttt{2.5.29.19} \textbar{} Defined in: RFC 5280 --- Internet
X.509 PKI Certificate and CRL Profile §4.2.1.9 \textbar{} Typically:
\textbf{Critical}

\begin{verbatim}
BasicConstraints ::= SEQUENCE {
    cA                BOOLEAN DEFAULT FALSE,
    pathLenConstraint INTEGER (0..MAX) OPTIONAL
}
\end{verbatim}

This extension distinguishes CA certificates from end-entity
certificates:

\begin{itemize}
\item
  \texttt{cA:\ TRUE} --- this certificate can issue other certificates
\item
  \texttt{cA:\ FALSE} (or absent) --- this is an end-entity certificate,
  it cannot issue other certificates
\item
  \texttt{pathLenConstraint} --- if present, limits how many
  intermediate CAs can follow this one in a chain
\end{itemize}

If you encounter a PDF signed by a certificate with
\texttt{BasicConstraints:\ cA=TRUE}, something is deeply wrong. CA
certificates should never be used for document signing --- their private
keys should be locked in HSMs and used only to sign other certificates.
A document signed by a CA certificate suggests either a misconfiguration
or a test setup that accidentally used the CA key for signing.

\subsubsection{SubjectKeyIdentifier and
AuthorityKeyIdentifier}\label{_subjectkeyidentifier_and_authoritykeyidentifier}

\textbf{SubjectKeyIdentifier} (OID: \texttt{2.5.29.14}): a hash of the
certificate's public key, used as a compact identifier. Typically the
SHA-1 hash of the BIT STRING value of the subjectPublicKey (not the
entire SubjectPublicKeyInfo structure --- just the key bytes).

\textbf{AuthorityKeyIdentifier} (OID: \texttt{2.5.29.35}): identifies
the CA certificate that signed this one. Contains the CA's
SubjectKeyIdentifier (keyIdentifier), and optionally the CA's issuer
name and serial number.

These two extensions are the primary mechanism for chain building:
follow the AuthorityKeyIdentifier from a certificate to find the CA
certificate whose SubjectKeyIdentifier matches.

The SubjectKeyIdentifier is typically a 20-byte SHA-1 hash. Yes,
SHA-1 --- even in 2026. This is not a security concern because the SKI
is used as an identifier, not for integrity verification. There is no
attack value in finding a collision for a SubjectKeyIdentifier.

\subsubsection{CRLDistributionPoints}\label{_crldistributionpoints}

OID: \texttt{2.5.29.31} \textbar{} Defined in: RFC 5280 --- Internet
X.509 PKI Certificate and CRL Profile §4.2.1.13

Contains one or more URLs where CRLs (Certificate Revocation Lists) can
be downloaded to check if this certificate has been revoked.

\begin{verbatim}
CRLDistributionPoints ::= SEQUENCE SIZE (1..MAX) OF DistributionPoint

DistributionPoint ::= SEQUENCE {
    distributionPoint [0] DistributionPointName OPTIONAL,
    reasons           [1] ReasonFlags OPTIONAL,
    cRLIssuer         [2] GeneralNames OPTIONAL
}
\end{verbatim}

In practice, this is almost always a single HTTP URL:

\begin{verbatim}
http://crl.example.com/issuing-ca.crl
\end{verbatim}

\subsubsection{AuthorityInfoAccess}\label{_authorityinfoaccess}

OID: \texttt{1.3.6.1.5.5.7.1.1} \textbar{} Defined in: RFC 5280 ---
Internet X.509 PKI Certificate and CRL Profile §4.2.2.1

This extension is critical for chain building and revocation checking.
It contains two types of access descriptions:

\begin{longtable}[]{@{}
  >{\raggedright\arraybackslash}p{(\linewidth - 4\tabcolsep) * \real{0.2222}}
  >{\raggedright\arraybackslash}p{(\linewidth - 4\tabcolsep) * \real{0.3333}}
  >{\raggedright\arraybackslash}p{(\linewidth - 4\tabcolsep) * \real{0.4444}}@{}}
\caption{AuthorityInfoAccess methods}\tabularnewline
\toprule\noalign{}
\begin{minipage}[b]{\linewidth}\raggedright
Method
\end{minipage} & \begin{minipage}[b]{\linewidth}\raggedright
OID
\end{minipage} & \begin{minipage}[b]{\linewidth}\raggedright
Purpose
\end{minipage} \\
\midrule\noalign{}
\endfirsthead
\toprule\noalign{}
\begin{minipage}[b]{\linewidth}\raggedright
Method
\end{minipage} & \begin{minipage}[b]{\linewidth}\raggedright
OID
\end{minipage} & \begin{minipage}[b]{\linewidth}\raggedright
Purpose
\end{minipage} \\
\midrule\noalign{}
\endhead
\bottomrule\noalign{}
\endlastfoot
\textbf{OCSP} & \texttt{1.3.6.1.5.5.7.48.1} & URL of the OCSP responder
for checking certificate revocation in real time. The validator sends a
request with the certificate's serial number and gets back a signed
response saying "good," "revoked," or "unknown." \\
\textbf{CA Issuers} & \texttt{1.3.6.1.5.5.7.48.2} & URL where the
issuing CA's certificate can be downloaded. This is how validators fetch
missing intermediate certificates during chain building. If a PDF
signature doesn't embed the full chain, the validator can use this URL
to find the intermediates. \\
\end{longtable}

The "CA Issuers" access method is the safety net for incomplete
certificate chains. When a signing tool embeds only the end-entity
certificate in the CMS (not the full chain), the validator must fetch
the intermediates. It uses the CA Issuers URL to download the next
certificate up the chain, then follows \emph{that} certificate's CA
Issuers URL, and so on until it reaches a trusted root. This
works --- when you have network access. In offline validation scenarios,
the chain must be complete. Always embed the full chain in your CMS
signatures.

\subsection{Certificate Profiles}\label{_certificate_profiles}

Not all certificates are created equal. The purpose of a certificate is
determined by its extensions, particularly KeyUsage and
ExtendedKeyUsage.

\subsubsection{End-Entity vs. CA
Certificates}\label{_end_entity_vs_ca_certificates}

\begin{longtable}[]{@{}
  >{\raggedright\arraybackslash}p{(\linewidth - 4\tabcolsep) * \real{0.2500}}
  >{\raggedright\arraybackslash}p{(\linewidth - 4\tabcolsep) * \real{0.3750}}
  >{\raggedright\arraybackslash}p{(\linewidth - 4\tabcolsep) * \real{0.3750}}@{}}
\caption{End-entity vs. CA certificate comparison}\tabularnewline
\toprule\noalign{}
\begin{minipage}[b]{\linewidth}\raggedright
Property
\end{minipage} & \begin{minipage}[b]{\linewidth}\raggedright
End-Entity Certificate
\end{minipage} & \begin{minipage}[b]{\linewidth}\raggedright
CA Certificate
\end{minipage} \\
\midrule\noalign{}
\endfirsthead
\toprule\noalign{}
\begin{minipage}[b]{\linewidth}\raggedright
Property
\end{minipage} & \begin{minipage}[b]{\linewidth}\raggedright
End-Entity Certificate
\end{minipage} & \begin{minipage}[b]{\linewidth}\raggedright
CA Certificate
\end{minipage} \\
\midrule\noalign{}
\endhead
\bottomrule\noalign{}
\endlastfoot
BasicConstraints & \texttt{cA:\ FALSE} or absent & \texttt{cA:\ TRUE},
optionally with pathLenConstraint \\
KeyUsage & \texttt{digitalSignature}, \texttt{nonRepudiation},
\texttt{keyEncipherment}, \texttt{dataEncipherment} (varies) &
\texttt{keyCertSign}, \texttt{cRLSign} \\
Purpose & Signs documents, authenticates TLS sessions, signs email &
Signs other certificates and CRLs \\
Private key location & Software (PKCS\#12), HSM, smart card & HSM
(always, for production CAs) \\
Validity period & 1-3 years typically & 10-30 years for roots, 5-10
years for intermediates \\
\end{longtable}

\subsubsection{Document Signing
Certificates}\label{_document_signing_certificates}

A certificate profile specifically for document signing typically has:

\begin{itemize}
\item
  KeyUsage: \texttt{digitalSignature} and/or \texttt{nonRepudiation}
\item
  ExtendedKeyUsage: document signing OID (varies by CA)
\item
  Subject: identifies the signer (person or organization)
\item
  Validity: 1-3 years
\item
  CRL and OCSP: both present for revocation checking
\end{itemize}

\subsubsection{How Different CAs Structure Their
Certificates}\label{_how_different_cas_structure_their_certificates}

Different CAs make different choices, and these choices affect
interoperability:

\begin{longtable}[]{@{}
  >{\raggedright\arraybackslash}p{(\linewidth - 6\tabcolsep) * \real{0.1818}}
  >{\raggedright\arraybackslash}p{(\linewidth - 6\tabcolsep) * \real{0.2727}}
  >{\raggedright\arraybackslash}p{(\linewidth - 6\tabcolsep) * \real{0.2727}}
  >{\raggedright\arraybackslash}p{(\linewidth - 6\tabcolsep) * \real{0.2727}}@{}}
\caption{CA certificate profile comparison}\tabularnewline
\toprule\noalign{}
\begin{minipage}[b]{\linewidth}\raggedright
Property
\end{minipage} & \begin{minipage}[b]{\linewidth}\raggedright
DigiCert
\end{minipage} & \begin{minipage}[b]{\linewidth}\raggedright
GlobalSign
\end{minipage} & \begin{minipage}[b]{\linewidth}\raggedright
National CAs (e.g., Austrian A-Trust)
\end{minipage} \\
\midrule\noalign{}
\endfirsthead
\toprule\noalign{}
\begin{minipage}[b]{\linewidth}\raggedright
Property
\end{minipage} & \begin{minipage}[b]{\linewidth}\raggedright
DigiCert
\end{minipage} & \begin{minipage}[b]{\linewidth}\raggedright
GlobalSign
\end{minipage} & \begin{minipage}[b]{\linewidth}\raggedright
National CAs (e.g., Austrian A-Trust)
\end{minipage} \\
\midrule\noalign{}
\endhead
\bottomrule\noalign{}
\endlastfoot
Key algorithm & RSA 2048/4096 or ECC P-256/P-384 & RSA 2048/4096 or ECC
P-256/P-384 & RSA 2048/4096 (ECC adoption varies) \\
Hash algorithm & SHA-256 or SHA-384 & SHA-256 or SHA-384 & SHA-256 \\
Subject DN & CN = person/org name, O = organization & CN = person/org
name, O = organization & CN = person name, serialNumber = national ID, C
= country code \\
EKU & Adobe document signing OID + email protection & Adobe document
signing OID & Varies widely --- some use EU-specific OIDs for qualified
certificates \\
AATL membership & Yes (since 2012) & Yes & Varies --- many EU national
CAs are on AATL \\
eIDAS qualified & Some certificate types & Some certificate types &
Typically yes (this is their primary purpose) \\
OCSP & Yes, with OCSP stapling & Yes & Yes (required by eIDAS) \\
Certificate Transparency & Yes (required by Chrome/Apple) & Yes &
Sometimes (not required for non-TLS certificates) \\
\end{longtable}

When debugging why a PDF signature is trusted in one viewer but not
another, the CA's AATL status is usually the answer. Adobe Acrobat
trusts AATL certificates. Other viewers (macOS Preview, Okular,
Firefox's built-in PDF viewer) use the OS trust store, which is
different. A DigiCert document signing certificate will be trusted in
Acrobat (AATL member), probably trusted in Windows (DigiCert roots are
in the Windows trust store), and possibly untrusted in Firefox (which
uses its own trust store). Same certificate, same signature, three
different trust decisions.

\subsection{Common Certificate Problems in the
Wild}\label{_common_certificate_problems_in_the_wild}

Years of debugging PDF signatures have revealed a pattern of recurring
certificate problems. Here are the ones you will encounter most often:

\subsubsection{Missing Extensions}\label{_missing_extensions}

\textbf{Symptom:} Signature validates cryptographically but shows
"untrusted" or "cannot determine if the signer is who they claim to be."

\textbf{Cause:} The signing certificate is missing critical
extensions --- typically ExtendedKeyUsage or SubjectKeyIdentifier. Some
legacy CAs (particularly older European government CAs) issued
certificates with minimal extensions that predate the current RFC 5280
--- Internet X.509 PKI Certificate and CRL Profile profile requirements.

\textbf{Fix:} The certificate needs to be reissued with proper
extensions. There is no workaround --- you cannot add extensions to an
existing certificate without re-signing it (which means the CA must do
it).

\subsubsection{Wrong Key Usage}\label{_wrong_key_usage}

\textbf{Symptom:} Signature fails validation with "key usage does not
permit digital signature" or similar.

\textbf{Cause:} The certificate has KeyUsage set to
\texttt{keyEncipherment} or \texttt{dataEncipherment} but not
\texttt{digitalSignature} or \texttt{nonRepudiation}. This happens when
someone uses a TLS server certificate (designed for key exchange) to
sign a PDF document.

\textbf{Fix:} Use a certificate with the correct KeyUsage for document
signing. A TLS certificate is not a signing certificate.

\subsubsection{Expired Certificates}\label{_expired_certificates}

\textbf{Symptom:} "Certificate has expired" warning, even though the
document was signed years ago when the certificate was valid.

\textbf{Cause:} The validator is checking the certificate against the
\emph{current} time instead of the \emph{signing} time. Without a
trusted timestamp, the validator has no proof of when the signature was
actually created.

\textbf{Fix:} Include a timestamp in the signature (see Chapter 18 on
LTV). A timestamp from a trusted TSA proves the signature existed at a
specific time, allowing validators to check the certificate's validity
against that time rather than the current time.

\subsubsection{Self-Signed
Certificates}\label{_self_signed_certificates}

\textbf{Symptom:} "The signer's identity is unknown" or "certificate is
not issued by a trusted CA."

\textbf{Cause:} The signing certificate is self-signed (the issuer and
subject are the same, and it signed itself). Self-signed certificates
have no external trust anchor.

\textbf{Fix:} For production use, obtain a certificate from a recognized
CA. For internal/testing use, you can add the self-signed certificate to
the viewer's trust store manually --- but this defeats the purpose of
the CA trust model.

Self-signed certificates are perfectly valid for testing and
development. PDFalyzer's self-signed certificate generation feature is
designed exactly for this purpose. The cryptographic integrity of the
signature is identical regardless of whether the certificate is
self-signed or CA-issued --- the difference is purely about trust.

\subsubsection{Incomplete Certificate
Chain}\label{_incomplete_certificate_chain}

\textbf{Symptom:} "Unable to build certificate chain" or "issuer
certificate not found."

\textbf{Cause:} The CMS signature embeds only the end-entity
certificate, not the intermediate CA certificates. The validator cannot
build a path from the end-entity to a trusted root.

\textbf{Fix:} Include the full certificate chain (end-entity + all
intermediates) in the CMS \texttt{certificates} field. Do NOT include
the root CA --- the root must already be in the validator's trust store.

This is one of the most common PDF signature problems. Always embed the
full chain (minus the root) in your CMS signatures. The signing code in
Chapter 2 already does this --- note how it includes \texttt{chain}
(from the PKCS\#12 key store) in the \texttt{CMSSignedDataGenerator}. If
your PKCS\#12 file doesn't contain the intermediates, you need to add
them before signing.

\subsection{Java: Loading and Inspecting
Certificates}\label{_java_loading_and_inspecting_certificates}

Here is how to load certificates and extract their fields using
BouncyCastle 1.80.

\begin{Shaded}
\begin{Highlighting}[]
\KeywordTok{import} \ImportTok{java}\OperatorTok{.}\ImportTok{io}\OperatorTok{.}\ImportTok{FileInputStream}\OperatorTok{;}
\KeywordTok{import} \ImportTok{java}\OperatorTok{.}\ImportTok{security}\OperatorTok{.}\ImportTok{cert}\OperatorTok{.}\ImportTok{CertificateFactory}\OperatorTok{;}
\KeywordTok{import} \ImportTok{java}\OperatorTok{.}\ImportTok{security}\OperatorTok{.}\ImportTok{cert}\OperatorTok{.}\ImportTok{X509Certificate}\OperatorTok{;}
\KeywordTok{import} \ImportTok{java}\OperatorTok{.}\ImportTok{util}\OperatorTok{.}\ImportTok{Date}\OperatorTok{;}
\KeywordTok{import} \ImportTok{java}\OperatorTok{.}\ImportTok{util}\OperatorTok{.}\ImportTok{List}\OperatorTok{;}

\KeywordTok{import} \ImportTok{org}\OperatorTok{.}\ImportTok{bouncycastle}\OperatorTok{.}\ImportTok{asn1}\OperatorTok{.}\ImportTok{ASN1ObjectIdentifier}\OperatorTok{;}
\KeywordTok{import} \ImportTok{org}\OperatorTok{.}\ImportTok{bouncycastle}\OperatorTok{.}\ImportTok{asn1}\OperatorTok{.}\ImportTok{x509}\OperatorTok{.}\ImportTok{AuthorityInformationAccess}\OperatorTok{;}
\KeywordTok{import} \ImportTok{org}\OperatorTok{.}\ImportTok{bouncycastle}\OperatorTok{.}\ImportTok{asn1}\OperatorTok{.}\ImportTok{x509}\OperatorTok{.}\ImportTok{BasicConstraints}\OperatorTok{;}
\KeywordTok{import} \ImportTok{org}\OperatorTok{.}\ImportTok{bouncycastle}\OperatorTok{.}\ImportTok{asn1}\OperatorTok{.}\ImportTok{x509}\OperatorTok{.}\ImportTok{CRLDistPoint}\OperatorTok{;}
\KeywordTok{import} \ImportTok{org}\OperatorTok{.}\ImportTok{bouncycastle}\OperatorTok{.}\ImportTok{asn1}\OperatorTok{.}\ImportTok{x509}\OperatorTok{.}\ImportTok{DistributionPoint}\OperatorTok{;}
\KeywordTok{import} \ImportTok{org}\OperatorTok{.}\ImportTok{bouncycastle}\OperatorTok{.}\ImportTok{asn1}\OperatorTok{.}\ImportTok{x509}\OperatorTok{.}\ImportTok{Extension}\OperatorTok{;}
\KeywordTok{import} \ImportTok{org}\OperatorTok{.}\ImportTok{bouncycastle}\OperatorTok{.}\ImportTok{asn1}\OperatorTok{.}\ImportTok{x509}\OperatorTok{.}\ImportTok{GeneralName}\OperatorTok{;}
\KeywordTok{import} \ImportTok{org}\OperatorTok{.}\ImportTok{bouncycastle}\OperatorTok{.}\ImportTok{asn1}\OperatorTok{.}\ImportTok{x509}\OperatorTok{.}\ImportTok{KeyUsage}\OperatorTok{;}
\KeywordTok{import} \ImportTok{org}\OperatorTok{.}\ImportTok{bouncycastle}\OperatorTok{.}\ImportTok{cert}\OperatorTok{.}\ImportTok{X509CertificateHolder}\OperatorTok{;}
\KeywordTok{import} \ImportTok{org}\OperatorTok{.}\ImportTok{bouncycastle}\OperatorTok{.}\ImportTok{cert}\OperatorTok{.}\ImportTok{jcajce}\OperatorTok{.}\ImportTok{JcaX509CertificateHolder}\OperatorTok{;}

\KeywordTok{public} \KeywordTok{class}\NormalTok{ CertificateInspector }\OperatorTok{\{}

    \CommentTok{/**}
     \CommentTok{*}\NormalTok{ Load a certificate from a DER or PEM file and inspect its fields}\CommentTok{.}
     \CommentTok{*/}
    \KeywordTok{public} \DataTypeTok{static} \DataTypeTok{void} \FunctionTok{inspectCertificate}\OperatorTok{(}\BuiltInTok{String}\NormalTok{ filePath}\OperatorTok{)} \KeywordTok{throws} \BuiltInTok{Exception} \OperatorTok{\{}
        \CommentTok{// Load the certificate using JCA}
        \BuiltInTok{CertificateFactory}\NormalTok{ cf }\OperatorTok{=} \BuiltInTok{CertificateFactory}\OperatorTok{.}\FunctionTok{getInstance}\OperatorTok{(}\StringTok{"X.509"}\OperatorTok{);}
        \BuiltInTok{X509Certificate}\NormalTok{ javaCert}\OperatorTok{;}
        \ControlFlowTok{try} \OperatorTok{(}\BuiltInTok{FileInputStream}\NormalTok{ fis }\OperatorTok{=} \KeywordTok{new} \BuiltInTok{FileInputStream}\OperatorTok{(}\NormalTok{filePath}\OperatorTok{))} \OperatorTok{\{}
\NormalTok{            javaCert }\OperatorTok{=} \OperatorTok{(}\BuiltInTok{X509Certificate}\OperatorTok{)}\NormalTok{ cf}\OperatorTok{.}\FunctionTok{generateCertificate}\OperatorTok{(}\NormalTok{fis}\OperatorTok{);}
        \OperatorTok{\}}

        \CommentTok{// Wrap in BouncyCastle holder for extension access}
\NormalTok{        X509CertificateHolder holder }\OperatorTok{=} \KeywordTok{new} \FunctionTok{JcaX509CertificateHolder}\OperatorTok{(}\NormalTok{javaCert}\OperatorTok{);}

        \CommentTok{// {-}{-}{-} Basic fields {-}{-}{-}}
        \BuiltInTok{String}\NormalTok{ subject }\OperatorTok{=}\NormalTok{ javaCert}\OperatorTok{.}\FunctionTok{getSubjectX500Principal}\OperatorTok{().}\FunctionTok{getName}\OperatorTok{();}      
        \BuiltInTok{String}\NormalTok{ issuer }\OperatorTok{=}\NormalTok{ javaCert}\OperatorTok{.}\FunctionTok{getIssuerX500Principal}\OperatorTok{().}\FunctionTok{getName}\OperatorTok{();}
        \BuiltInTok{String}\NormalTok{ serial }\OperatorTok{=}\NormalTok{ javaCert}\OperatorTok{.}\FunctionTok{getSerialNumber}\OperatorTok{().}\FunctionTok{toString}\OperatorTok{(}\DecValTok{16}\OperatorTok{);}            
        \DataTypeTok{int}\NormalTok{ version }\OperatorTok{=}\NormalTok{ javaCert}\OperatorTok{.}\FunctionTok{getVersion}\OperatorTok{();}
        \BuiltInTok{Date}\NormalTok{ notBefore }\OperatorTok{=}\NormalTok{ javaCert}\OperatorTok{.}\FunctionTok{getNotBefore}\OperatorTok{();}
        \BuiltInTok{Date}\NormalTok{ notAfter }\OperatorTok{=}\NormalTok{ javaCert}\OperatorTok{.}\FunctionTok{getNotAfter}\OperatorTok{();}
        \BuiltInTok{String}\NormalTok{ sigAlg }\OperatorTok{=}\NormalTok{ javaCert}\OperatorTok{.}\FunctionTok{getSigAlgName}\OperatorTok{();}
        \BuiltInTok{String}\NormalTok{ pubKeyAlg }\OperatorTok{=}\NormalTok{ javaCert}\OperatorTok{.}\FunctionTok{getPublicKey}\OperatorTok{().}\FunctionTok{getAlgorithm}\OperatorTok{();}
        \DataTypeTok{int}\NormalTok{ keySize }\OperatorTok{=} \FunctionTok{getKeySize}\OperatorTok{(}\NormalTok{javaCert}\OperatorTok{);}

        \CommentTok{// {-}{-}{-} KeyUsage {-}{-}{-}}
        \DataTypeTok{boolean}\OperatorTok{[]}\NormalTok{ keyUsageBits }\OperatorTok{=}\NormalTok{ javaCert}\OperatorTok{.}\FunctionTok{getKeyUsage}\OperatorTok{();}                    
        \ControlFlowTok{if} \OperatorTok{(}\NormalTok{keyUsageBits }\OperatorTok{!=} \KeywordTok{null}\OperatorTok{)} \OperatorTok{\{}
            \DataTypeTok{boolean}\NormalTok{ digitalSignature }\OperatorTok{=}\NormalTok{ keyUsageBits}\OperatorTok{[}\DecValTok{0}\OperatorTok{];}
            \DataTypeTok{boolean}\NormalTok{ nonRepudiation }\OperatorTok{=}\NormalTok{ keyUsageBits}\OperatorTok{[}\DecValTok{1}\OperatorTok{];}
            \DataTypeTok{boolean}\NormalTok{ keyEncipherment }\OperatorTok{=}\NormalTok{ keyUsageBits}\OperatorTok{[}\DecValTok{2}\OperatorTok{];}
            \DataTypeTok{boolean}\NormalTok{ keyCertSign }\OperatorTok{=}\NormalTok{ keyUsageBits}\OperatorTok{[}\DecValTok{5}\OperatorTok{];}
        \OperatorTok{\}}

        \CommentTok{// {-}{-}{-} BasicConstraints {-}{-}{-}}
\NormalTok{        Extension bcExt }\OperatorTok{=}\NormalTok{ holder}\OperatorTok{.}\FunctionTok{getExtension}\OperatorTok{(}\NormalTok{Extension}\OperatorTok{.}\FunctionTok{basicConstraints}\OperatorTok{);}
        \ControlFlowTok{if} \OperatorTok{(}\NormalTok{bcExt }\OperatorTok{!=} \KeywordTok{null}\OperatorTok{)} \OperatorTok{\{}
\NormalTok{            BasicConstraints bc }\OperatorTok{=}\NormalTok{ BasicConstraints}\OperatorTok{.}\FunctionTok{getInstance}\OperatorTok{(}
\NormalTok{                    bcExt}\OperatorTok{.}\FunctionTok{getParsedValue}\OperatorTok{());}
            \DataTypeTok{boolean}\NormalTok{ isCA }\OperatorTok{=}\NormalTok{ bc}\OperatorTok{.}\FunctionTok{isCA}\OperatorTok{();}
            \DataTypeTok{int}\NormalTok{ pathLen }\OperatorTok{=}\NormalTok{ bc}\OperatorTok{.}\FunctionTok{getPathLenConstraint}\OperatorTok{()} \OperatorTok{!=} \KeywordTok{null}
                    \OperatorTok{?}\NormalTok{ bc}\OperatorTok{.}\FunctionTok{getPathLenConstraint}\OperatorTok{().}\FunctionTok{intValue}\OperatorTok{()} \OperatorTok{:} \OperatorTok{{-}}\DecValTok{1}\OperatorTok{;}
        \OperatorTok{\}}

        \CommentTok{// {-}{-}{-} ExtendedKeyUsage {-}{-}{-}}
        \BuiltInTok{List}\OperatorTok{\textless{}}\BuiltInTok{String}\OperatorTok{\textgreater{}}\NormalTok{ ekuOids }\OperatorTok{=}\NormalTok{ javaCert}\OperatorTok{.}\FunctionTok{getExtendedKeyUsage}\OperatorTok{();}              
        \ControlFlowTok{if} \OperatorTok{(}\NormalTok{ekuOids }\OperatorTok{!=} \KeywordTok{null}\OperatorTok{)} \OperatorTok{\{}
            \ControlFlowTok{for} \OperatorTok{(}\BuiltInTok{String}\NormalTok{ oid }\OperatorTok{:}\NormalTok{ ekuOids}\OperatorTok{)} \OperatorTok{\{}
                \CommentTok{// Map OID to human{-}readable name}
                \BuiltInTok{String}\NormalTok{ name }\OperatorTok{=} \FunctionTok{mapEkuOid}\OperatorTok{(}\NormalTok{oid}\OperatorTok{);}
            \OperatorTok{\}}
        \OperatorTok{\}}

        \CommentTok{// {-}{-}{-} AuthorityInfoAccess (OCSP + CA Issuers) {-}{-}{-}}
\NormalTok{        Extension aiaExt }\OperatorTok{=}\NormalTok{ holder}\OperatorTok{.}\FunctionTok{getExtension}\OperatorTok{(}\NormalTok{Extension}\OperatorTok{.}\FunctionTok{authorityInfoAccess}\OperatorTok{);}
        \ControlFlowTok{if} \OperatorTok{(}\NormalTok{aiaExt }\OperatorTok{!=} \KeywordTok{null}\OperatorTok{)} \OperatorTok{\{}
\NormalTok{            AuthorityInformationAccess aia }\OperatorTok{=}\NormalTok{ AuthorityInformationAccess}
                    \OperatorTok{.}\FunctionTok{getInstance}\OperatorTok{(}\NormalTok{aiaExt}\OperatorTok{.}\FunctionTok{getParsedValue}\OperatorTok{());}

            \CommentTok{// OCSP responder URLs}
            \DataTypeTok{var}\NormalTok{ ocspAccessDescriptions }\OperatorTok{=}\NormalTok{ aia}\OperatorTok{.}\FunctionTok{getAccessDescriptions}\OperatorTok{();}
            \ControlFlowTok{for} \OperatorTok{(}\DataTypeTok{var}\NormalTok{ ad }\OperatorTok{:}\NormalTok{ ocspAccessDescriptions}\OperatorTok{)} \OperatorTok{\{}
                \ControlFlowTok{if} \OperatorTok{(}\NormalTok{ad}\OperatorTok{.}\FunctionTok{getAccessMethod}\OperatorTok{().}\FunctionTok{equals}\OperatorTok{(}
                        \KeywordTok{new} \FunctionTok{ASN1ObjectIdentifier}\OperatorTok{(}\StringTok{"1.3.6.1.5.5.7.48.1"}\OperatorTok{)))} \OperatorTok{\{}
                    \BuiltInTok{String}\NormalTok{ ocspUrl }\OperatorTok{=}\NormalTok{ ad}\OperatorTok{.}\FunctionTok{getAccessLocation}\OperatorTok{()}
                            \OperatorTok{.}\FunctionTok{getName}\OperatorTok{().}\FunctionTok{toString}\OperatorTok{();}                          
                \OperatorTok{\}}
            \OperatorTok{\}}
        \OperatorTok{\}}

        \CommentTok{// {-}{-}{-} CRL Distribution Points {-}{-}{-}}
\NormalTok{        Extension crlExt }\OperatorTok{=}\NormalTok{ holder}\OperatorTok{.}\FunctionTok{getExtension}\OperatorTok{(}\NormalTok{Extension}\OperatorTok{.}\FunctionTok{cRLDistributionPoints}\OperatorTok{);}
        \ControlFlowTok{if} \OperatorTok{(}\NormalTok{crlExt }\OperatorTok{!=} \KeywordTok{null}\OperatorTok{)} \OperatorTok{\{}
\NormalTok{            CRLDistPoint crlDistPoint }\OperatorTok{=}\NormalTok{ CRLDistPoint}\OperatorTok{.}\FunctionTok{getInstance}\OperatorTok{(}
\NormalTok{                    crlExt}\OperatorTok{.}\FunctionTok{getParsedValue}\OperatorTok{());}
            \ControlFlowTok{for} \OperatorTok{(}\NormalTok{DistributionPoint dp }\OperatorTok{:}\NormalTok{ crlDistPoint}\OperatorTok{.}\FunctionTok{getDistributionPoints}\OperatorTok{())} \OperatorTok{\{}
                \CommentTok{// Extract URL from the distribution point}
                \DataTypeTok{var}\NormalTok{ dpName }\OperatorTok{=}\NormalTok{ dp}\OperatorTok{.}\FunctionTok{getDistributionPoint}\OperatorTok{();}
                \ControlFlowTok{if} \OperatorTok{(}\NormalTok{dpName }\OperatorTok{!=} \KeywordTok{null}\OperatorTok{)} \OperatorTok{\{}
                    \DataTypeTok{var}\NormalTok{ names }\OperatorTok{=}\NormalTok{ org}\OperatorTok{.}\FunctionTok{bouncycastle}\OperatorTok{.}\FunctionTok{asn1}\OperatorTok{.}\FunctionTok{x509}\OperatorTok{.}\FunctionTok{GeneralNames}
                            \OperatorTok{.}\FunctionTok{getInstance}\OperatorTok{(}\NormalTok{dpName}\OperatorTok{.}\FunctionTok{getName}\OperatorTok{());}
                    \ControlFlowTok{for} \OperatorTok{(}\NormalTok{GeneralName name }\OperatorTok{:}\NormalTok{ names}\OperatorTok{.}\FunctionTok{getNames}\OperatorTok{())} \OperatorTok{\{}
                        \ControlFlowTok{if} \OperatorTok{(}\NormalTok{name}\OperatorTok{.}\FunctionTok{getTagNo}\OperatorTok{()} \OperatorTok{==}\NormalTok{ GeneralName}\OperatorTok{.}\FunctionTok{uniformResourceIdentifier}\OperatorTok{)} \OperatorTok{\{}
                            \BuiltInTok{String}\NormalTok{ crlUrl }\OperatorTok{=}\NormalTok{ name}\OperatorTok{.}\FunctionTok{getName}\OperatorTok{().}\FunctionTok{toString}\OperatorTok{();}
                        \OperatorTok{\}}
                    \OperatorTok{\}}
                \OperatorTok{\}}
            \OperatorTok{\}}
        \OperatorTok{\}}
    \OperatorTok{\}}

    \KeywordTok{private} \DataTypeTok{static} \DataTypeTok{int} \FunctionTok{getKeySize}\OperatorTok{(}\BuiltInTok{X509Certificate}\NormalTok{ cert}\OperatorTok{)} \OperatorTok{\{}
        \BuiltInTok{String}\NormalTok{ algorithm }\OperatorTok{=}\NormalTok{ cert}\OperatorTok{.}\FunctionTok{getPublicKey}\OperatorTok{().}\FunctionTok{getAlgorithm}\OperatorTok{();}
        \ControlFlowTok{if} \OperatorTok{(}\StringTok{"RSA"}\OperatorTok{.}\FunctionTok{equals}\OperatorTok{(}\NormalTok{algorithm}\OperatorTok{))} \OperatorTok{\{}
            \DataTypeTok{var}\NormalTok{ rsaKey }\OperatorTok{=} \OperatorTok{(}\NormalTok{java}\OperatorTok{.}\FunctionTok{security}\OperatorTok{.}\FunctionTok{interfaces}\OperatorTok{.}\FunctionTok{RSAPublicKey}\OperatorTok{)}\NormalTok{ cert}\OperatorTok{.}\FunctionTok{getPublicKey}\OperatorTok{();}
            \ControlFlowTok{return}\NormalTok{ rsaKey}\OperatorTok{.}\FunctionTok{getModulus}\OperatorTok{().}\FunctionTok{bitLength}\OperatorTok{();}
        \OperatorTok{\}} \ControlFlowTok{else} \ControlFlowTok{if} \OperatorTok{(}\StringTok{"EC"}\OperatorTok{.}\FunctionTok{equals}\OperatorTok{(}\NormalTok{algorithm}\OperatorTok{))} \OperatorTok{\{}
            \DataTypeTok{var}\NormalTok{ ecKey }\OperatorTok{=} \OperatorTok{(}\NormalTok{java}\OperatorTok{.}\FunctionTok{security}\OperatorTok{.}\FunctionTok{interfaces}\OperatorTok{.}\FunctionTok{ECPublicKey}\OperatorTok{)}\NormalTok{ cert}\OperatorTok{.}\FunctionTok{getPublicKey}\OperatorTok{();}
            \ControlFlowTok{return}\NormalTok{ ecKey}\OperatorTok{.}\FunctionTok{getParams}\OperatorTok{().}\FunctionTok{getOrder}\OperatorTok{().}\FunctionTok{bitLength}\OperatorTok{();}
        \OperatorTok{\}}
        \ControlFlowTok{return} \OperatorTok{{-}}\DecValTok{1}\OperatorTok{;}
    \OperatorTok{\}}

    \KeywordTok{private} \DataTypeTok{static} \BuiltInTok{String} \FunctionTok{mapEkuOid}\OperatorTok{(}\BuiltInTok{String}\NormalTok{ oid}\OperatorTok{)} \OperatorTok{\{}
        \ControlFlowTok{return} \ControlFlowTok{switch} \OperatorTok{(}\NormalTok{oid}\OperatorTok{)} \OperatorTok{\{}
            \ControlFlowTok{case} \StringTok{"1.3.6.1.5.5.7.3.4"} \OperatorTok{{-}\textgreater{}} \StringTok{"emailProtection"}\OperatorTok{;}
            \ControlFlowTok{case} \StringTok{"1.3.6.1.4.1.311.10.3.12"} \OperatorTok{{-}\textgreater{}} \StringTok{"documentSigning (Microsoft)"}\OperatorTok{;}
            \ControlFlowTok{case} \StringTok{"1.2.840.113583.1.1.5"} \OperatorTok{{-}\textgreater{}} \StringTok{"adobeAuthenticDocumentsTrust"}\OperatorTok{;}
            \ControlFlowTok{case} \StringTok{"1.3.6.1.5.5.7.3.1"} \OperatorTok{{-}\textgreater{}} \StringTok{"serverAuth (TLS)"}\OperatorTok{;}
            \ControlFlowTok{case} \StringTok{"1.3.6.1.5.5.7.3.9"} \OperatorTok{{-}\textgreater{}} \StringTok{"OCSPSigning"}\OperatorTok{;}
            \KeywordTok{default} \OperatorTok{{-}\textgreater{}} \StringTok{"unknown ("} \OperatorTok{+}\NormalTok{ oid }\OperatorTok{+} \StringTok{")"}\OperatorTok{;}
        \OperatorTok{\};}
    \OperatorTok{\}}
\OperatorTok{\}}
\end{Highlighting}
\end{Shaded}

\begin{itemize}
\item
  Returns the subject DN in RFC 2253 format (CN first). For display, you
  might want to parse this into individual components.
\item
  Serial numbers are typically displayed in hexadecimal. Some CAs use
  very large random serial numbers (up to 20 bytes per RFC 5280 ---
  Internet X.509 PKI Certificate and CRL Profile).
\item
  \texttt{getKeyUsage()} returns a boolean array where index 0 =
  digitalSignature, 1 = nonRepudiation, \ldots\hspace{0pt}, 5 =
  keyCertSign, 6 = cRLSign. Returns null if the extension is absent.
\item
  Java's \texttt{getExtendedKeyUsage()} returns OIDs as strings. You
  need to map them to human-readable names yourself.
\item
  The OCSP URL is what you send revocation check requests to. This URL
  must be accessible from the validator's network.
\end{itemize}

\subsection{Dissecting a Real Signing
Certificate}\label{_dissecting_a_real_signing_certificate}

Let's walk through a real document signing certificate, field by field,
to see how the theory maps to practice.

\begin{verbatim}
Certificate:
    Data:
        Version: 3 (0x2)
        Serial Number:
            0a:55:14:f1:28:1d:45:87:21:f1:27:3c:e5:b6:4f:2a
        Signature Algorithm: sha256WithRSAEncryption
        Issuer: C=US, O=DigiCert Inc, CN=DigiCert SHA2 Assured ID CA
        Validity
            Not Before: Jan 15 00:00:00 2024 GMT
            Not After : Jan 14 23:59:59 2027 GMT
        Subject: C=US, ST=California, L=San Francisco,
                 O=Acme Corporation, CN=Acme Document Signing
        Subject Public Key Info:
            Public Key Algorithm: rsaEncryption
                RSA Public-Key: (2048 bit)
\end{verbatim}

Let's annotate each field:

\begin{longtable}[]{@{}
  >{\raggedright\arraybackslash}p{(\linewidth - 2\tabcolsep) * \real{0.2857}}
  >{\raggedright\arraybackslash}p{(\linewidth - 2\tabcolsep) * \real{0.7143}}@{}}
\caption{Certificate field analysis}\tabularnewline
\toprule\noalign{}
\begin{minipage}[b]{\linewidth}\raggedright
Field
\end{minipage} & \begin{minipage}[b]{\linewidth}\raggedright
Analysis
\end{minipage} \\
\midrule\noalign{}
\endfirsthead
\toprule\noalign{}
\begin{minipage}[b]{\linewidth}\raggedright
Field
\end{minipage} & \begin{minipage}[b]{\linewidth}\raggedright
Analysis
\end{minipage} \\
\midrule\noalign{}
\endhead
\bottomrule\noalign{}
\endlastfoot
Version 3 & Expected. All modern certificates are v3 to support
extensions. \\
Serial Number & 16 bytes (128 bits) of random data. DigiCert uses large
random serials to prevent hash-based attacks on the serial number. \\
Signature Algorithm & \texttt{sha256WithRSAEncryption} --- SHA-256 hash
with RSA signature. This is the algorithm the CA (DigiCert) used to sign
this certificate. \\
Issuer & \texttt{DigiCert\ SHA2\ Assured\ ID\ CA} --- this is an
intermediate CA certificate. The "SHA2" in the name indicates it's a
post-SHA-1 intermediate. You'll need to find \emph{this} certificate to
build the chain. \\
Validity & 3-year validity (Jan 2024 to Jan 2027). Typical for document
signing certificates. The signer needs to renew before Jan 14, 2027, or
new signatures will be created with an expired certificate. \\
Subject & Identifies the organization (Acme Corporation) and the purpose
(Document Signing). The CN often contains the signer's name or the
certificate's purpose rather than a hostname. \\
Public Key & RSA 2048-bit. The minimum for document signing. Some
organizations require 4096-bit, but 2048 is widely accepted and
faster. \\
\end{longtable}

The extensions would look like:

\begin{verbatim}
X509v3 extensions:
    X509v3 Key Usage: critical
        Digital Signature
    X509v3 Extended Key Usage:
        1.2.840.113583.1.1.5
    X509v3 Basic Constraints: critical
        CA:FALSE
    X509v3 Subject Key Identifier:
        A1:B2:C3:D4:E5:F6:...
    X509v3 Authority Key Identifier:
        keyid:F8:7E:34:22:...
    Authority Information Access:
        OCSP - URI:http://ocsp.digicert.com
        CA Issuers - URI:http://cacerts.digicert.com/DigiCertSHA2AssuredIDCA.crt
    X509v3 CRL Distribution Points:
        Full Name:
          URI:http://crl3.digicert.com/sha2-assured-id-ca.crl
          URI:http://crl4.digicert.com/sha2-assured-id-ca.crl
\end{verbatim}

\begin{longtable}[]{@{}
  >{\raggedright\arraybackslash}p{(\linewidth - 2\tabcolsep) * \real{0.2857}}
  >{\raggedright\arraybackslash}p{(\linewidth - 2\tabcolsep) * \real{0.7143}}@{}}
\caption{Extension analysis}\tabularnewline
\toprule\noalign{}
\begin{minipage}[b]{\linewidth}\raggedright
Extension
\end{minipage} & \begin{minipage}[b]{\linewidth}\raggedright
Analysis
\end{minipage} \\
\midrule\noalign{}
\endfirsthead
\toprule\noalign{}
\begin{minipage}[b]{\linewidth}\raggedright
Extension
\end{minipage} & \begin{minipage}[b]{\linewidth}\raggedright
Analysis
\end{minipage} \\
\midrule\noalign{}
\endhead
\bottomrule\noalign{}
\endlastfoot
KeyUsage: digitalSignature & Correct for document signing. This
certificate can create digital signatures. The \texttt{critical} flag
means validators must enforce this. \\
ExtendedKeyUsage: 1.2.840.113583.1.1.5 & This is the Adobe Authentic
Documents Trust OID. It means this certificate is intended for PDF
document signing and is likely on the AATL. \\
BasicConstraints: CA:FALSE & Correct. This is an end-entity certificate.
It cannot issue other certificates. \\
SubjectKeyIdentifier & A hash of the public key. Used for chain building
(other certificates reference this to identify this certificate as their
issuer). \\
AuthorityKeyIdentifier & Points to the issuing CA's
SubjectKeyIdentifier. Follow this to find
\texttt{DigiCert\ SHA2\ Assured\ ID\ CA}. \\
AIA: OCSP & The OCSP responder URL. Send a request here to check if this
certificate has been revoked. \\
AIA: CA Issuers & The URL to download the issuing CA's certificate. If
the CMS signature doesn't include the intermediate, a validator can
fetch it here. \\
CRL Distribution Points & Two CRL URLs (redundancy). Download the CRL to
check for revocation without OCSP. \\
\end{longtable}

DigiCert provides two CRL URLs (crl3 and crl4) for redundancy. If one
CDN is down, the other should still work. This is a best practice --- if
your CA provides only one CRL URL and it becomes unreachable, validators
that use CRL checking will fail to validate signatures. OCSP is the
preferred revocation checking method (real-time, smaller responses), but
CRL is the fallback.

\subsection{Summary}\label{_summary_3}

X.509 certificates are the bridge between raw cryptography and trust.
They bind a public key to an identity, and the CA's signature on that
binding is what makes the trust model work --- or fail, when CAs are
compromised (DigiNotar), when extensions are missing, or when the chain
is incomplete.

For PDF signatures:

\begin{itemize}
\item
  \textbf{Every signed PDF embeds one or more X.509 certificates} in the
  CMS SignedData structure.
\item
  \textbf{The signing certificate must have \texttt{digitalSignature} or
  \texttt{nonRepudiation}} in KeyUsage.
\item
  \textbf{The full chain} (end-entity + intermediates, not the root)
  should be embedded in the CMS.
\item
  \textbf{AIA extensions} provide fallback URLs for chain building and
  OCSP revocation checking.
\item
  \textbf{Trust is viewer-dependent}: Acrobat uses AATL, browsers use
  their own trust stores, OS-level trust stores vary. The same
  certificate can be trusted in one context and untrusted in another.
\end{itemize}

The next chapter follows these certificates up the chain --- from the
end-entity signing certificate, through intermediates, to the root CA
that anchors the entire trust hierarchy.

\chapter{Certificate Chains \& Trust}\label{ch04}

\section{Why Chains Exist}\label{_why_chains_exist}

In Chapter 3 we established that an X.509 certificate binds a public key
to an identity. A certificate authority (CA) vouches for that binding by
signing the certificate with its own private key. But who vouches for
the CA?

This is the trust delegation problem. A root CA cannot practically issue
certificates to every person and organization on the planet. The
operational burden would be staggering: vetting applicants, maintaining
infrastructure, handling revocations, passing audits --- all for
millions of subscribers. Even if it could, concentrating all certificate
issuance in a single entity would create a catastrophic single point of
failure. If that one root CA were compromised, every certificate it had
ever issued would become suspect.

The solution is \textbf{hierarchical delegation}. A root CA issues
certificates not (primarily) to end users, but to \textbf{intermediate
CAs}. Those intermediate CAs, in turn, issue certificates to other
intermediates or to end-entity subscribers. The root CA's private key
stays in a vault --- literally, in most cases, inside an HSM in a
physically secured facility that is powered on only for scheduled
signing ceremonies.

This creates a chain:

\begin{enumerate}
\def\labelenumi{\arabic{enumi}.}
\item
  The \textbf{root CA} self-signs its own certificate. It is the trust
  anchor --- the starting point that must be accepted on faith (or, more
  precisely, accepted because it appears in a trust store you have
  chosen to rely on).
\item
  One or more \textbf{intermediate CAs} hold certificates signed by the
  root CA (or by another intermediate above them). They perform the
  day-to-day work of issuing certificates.
\item
  The \textbf{end-entity} (sometimes called the "leaf") certificate is
  the one that belongs to the actual signer --- the person or service
  that signs PDFs. This certificate is signed by an intermediate CA.
\end{enumerate}

The term "chain" is somewhat misleading. In practice, a certificate's
trust path can be a \textbf{directed acyclic graph} rather than a simple
linear chain. An intermediate CA might hold multiple
certificates --- one signed by Root A, another signed by Root
B --- enabling cross-certification between different PKI hierarchies.
However, for PDF signatures, you almost always deal with simple linear
chains of depth two or three: root → intermediate → end-entity, or root
→ intermediate → intermediate → end-entity. The rare cases where
cross-certification matters are in government PKIs and the EU Trusted
List hierarchy, which we cover later in this chapter.

\section{Chain Structure in Detail}\label{_chain_structure_in_detail}

A typical PDF signing chain has three certificates:

\begin{longtable}[]{@{}
  >{\raggedright\arraybackslash}p{(\linewidth - 4\tabcolsep) * \real{0.1429}}
  >{\raggedright\arraybackslash}p{(\linewidth - 4\tabcolsep) * \real{0.2857}}
  >{\raggedright\arraybackslash}p{(\linewidth - 4\tabcolsep) * \real{0.5714}}@{}}
\caption{Certificate chain for a PDF signing certificate}\tabularnewline
\toprule\noalign{}
\begin{minipage}[b]{\linewidth}\raggedright
Position
\end{minipage} & \begin{minipage}[b]{\linewidth}\raggedright
Role
\end{minipage} & \begin{minipage}[b]{\linewidth}\raggedright
Characteristics
\end{minipage} \\
\midrule\noalign{}
\endfirsthead
\toprule\noalign{}
\begin{minipage}[b]{\linewidth}\raggedright
Position
\end{minipage} & \begin{minipage}[b]{\linewidth}\raggedright
Role
\end{minipage} & \begin{minipage}[b]{\linewidth}\raggedright
Characteristics
\end{minipage} \\
\midrule\noalign{}
\endhead
\bottomrule\noalign{}
\endlastfoot
1 (leaf) & End-entity (signer) & Contains the signer's public key. Has
\texttt{keyUsage} with \texttt{digitalSignature} or
\texttt{nonRepudiation} set. Has \texttt{extendedKeyUsage} with document
signing or email protection OIDs. Issued by the intermediate CA. Short
validity period (1-3 years). \\
2 & Intermediate CA & Contains the intermediate CA's public key. Has
\texttt{basicConstraints} with \texttt{cA=TRUE} and a
\texttt{pathLenConstraint}. Has \texttt{keyUsage} with
\texttt{keyCertSign} and \texttt{cRLSign}. Issued by the root CA. Medium
validity period (5-10 years). \\
3 (anchor) & Root CA & Self-signed: the issuer and subject are
identical. Has \texttt{basicConstraints} with \texttt{cA=TRUE}. Present
in trust stores. Long validity period (20-30 years). Never transmitted
in the CMS structure of a PDF signature (validators must already have it
in their trust store). \\
\end{longtable}

The CMS \texttt{SignedData} structure in a PDF signature's
\texttt{/Contents} field typically includes certificates at positions 1
and 2 --- the end-entity and intermediate(s). The root CA certificate is
\textbf{not} included. This is by design: the validator must already
trust the root, and including it in the signature would just waste
space. The validator checks the chain by building a path from the
end-entity up through the intermediates until it reaches a certificate
that matches one in its trust store.

Some signing implementations include the root certificate in the CMS
\texttt{certificates} field. This is not wrong per se --- RFC 5652 ---
Cryptographic Message Syntax (CMS) allows it --- but it serves no
purpose. A validator will not trust a root certificate just because it
appears inside a signature. Trust comes from the trust store, not from
the document. If anything, including the root wastes space in the
\texttt{/Contents} hex string, which matters when you are estimating the
size of the content gap (Chapter 13). PDFBox does not include the root
by default. iText does. Adobe Acrobat tolerates both approaches.

A chain with a \textbf{missing intermediate} is the single most common
cause of "certificate chain not trusted" errors in PDF signatures. The
end-entity certificate is present, the root is in the trust store, but
the intermediate that links them is nowhere to be found. The CMS
structure does not contain it, and the validator cannot download it (or
is configured not to). Always include the full chain (minus the root) in
your CMS \texttt{certificates} field.

\section{Building the Chain}\label{_building_the_chain}

Given a bag of certificates --- perhaps from a CMS \texttt{SignedData}
structure, perhaps from a PKCS\#12 file, perhaps assembled from multiple
sources --- how do you construct the actual chain?

The algorithm relies on two X.509 extensions that exist specifically for
this purpose:

\subsection{AuthorityKeyIdentifier and
SubjectKeyIdentifier}\label{_authoritykeyidentifier_and_subjectkeyidentifier}

Every well-formed certificate contains:

\begin{itemize}
\item
  \textbf{SubjectKeyIdentifier (SKI)}: an extension on the certificate
  itself, identifying the public key it contains. Typically a SHA-1 hash
  of the public key's bit string. Defined in RFC 5280 --- Internet X.509
  PKI Certificate and CRL Profile section 4.2.1.2.
\item
  \textbf{AuthorityKeyIdentifier (AKI)}: an extension identifying the
  certificate that signed this one. Contains the SKI of the issuer's
  certificate. Defined in RFC 5280 --- Internet X.509 PKI Certificate
  and CRL Profile section 4.2.1.1.
\end{itemize}

The chain-building algorithm is conceptually simple:

\begin{enumerate}
\def\labelenumi{\arabic{enumi}.}
\item
  Start with the end-entity certificate (the signer's cert from the CMS
  \texttt{SignerInfo}).
\item
  Read its AKI extension. This gives you the SKI of the issuer.
\item
  Search the available certificates for one whose SKI matches that AKI.
\item
  When you find it, that is the next link in the chain. Read \emph{its}
  AKI.
\item
  Repeat until you find a certificate where the subject equals the
  issuer (self-signed) or one that matches a trust anchor.
\end{enumerate}

\begin{Shaded}
\begin{Highlighting}[]
\KeywordTok{import} \ImportTok{java}\OperatorTok{.}\ImportTok{security}\OperatorTok{.}\ImportTok{cert}\OperatorTok{.}\ImportTok{X509Certificate}\OperatorTok{;}
\KeywordTok{import} \ImportTok{java}\OperatorTok{.}\ImportTok{util}\OperatorTok{.}\ImportTok{ArrayList}\OperatorTok{;}
\KeywordTok{import} \ImportTok{java}\OperatorTok{.}\ImportTok{util}\OperatorTok{.}\ImportTok{Arrays}\OperatorTok{;}
\KeywordTok{import} \ImportTok{java}\OperatorTok{.}\ImportTok{util}\OperatorTok{.}\ImportTok{Collection}\OperatorTok{;}
\KeywordTok{import} \ImportTok{java}\OperatorTok{.}\ImportTok{util}\OperatorTok{.}\ImportTok{List}\OperatorTok{;}

\KeywordTok{import} \ImportTok{org}\OperatorTok{.}\ImportTok{bouncycastle}\OperatorTok{.}\ImportTok{asn1}\OperatorTok{.}\ImportTok{x509}\OperatorTok{.}\ImportTok{AuthorityKeyIdentifier}\OperatorTok{;}
\KeywordTok{import} \ImportTok{org}\OperatorTok{.}\ImportTok{bouncycastle}\OperatorTok{.}\ImportTok{asn1}\OperatorTok{.}\ImportTok{x509}\OperatorTok{.}\ImportTok{Extension}\OperatorTok{;}
\KeywordTok{import} \ImportTok{org}\OperatorTok{.}\ImportTok{bouncycastle}\OperatorTok{.}\ImportTok{asn1}\OperatorTok{.}\ImportTok{x509}\OperatorTok{.}\ImportTok{SubjectKeyIdentifier}\OperatorTok{;}
\KeywordTok{import} \ImportTok{org}\OperatorTok{.}\ImportTok{bouncycastle}\OperatorTok{.}\ImportTok{cert}\OperatorTok{.}\ImportTok{X509CertificateHolder}\OperatorTok{;}
\KeywordTok{import} \ImportTok{org}\OperatorTok{.}\ImportTok{bouncycastle}\OperatorTok{.}\ImportTok{cert}\OperatorTok{.}\ImportTok{jcajce}\OperatorTok{.}\ImportTok{JcaX509CertificateHolder}\OperatorTok{;}

\CommentTok{/**}
 \CommentTok{*}\NormalTok{ Build a certificate chain from an unordered collection of certificates}\CommentTok{.}
 \CommentTok{*}\NormalTok{ Matches AuthorityKeyIdentifier }\CommentTok{{-}\textgreater{}}\NormalTok{ SubjectKeyIdentifier}\CommentTok{.}
 \CommentTok{*/}
\KeywordTok{public} \KeywordTok{class}\NormalTok{ ChainBuilder }\OperatorTok{\{}

    \KeywordTok{public} \DataTypeTok{static} \BuiltInTok{List}\OperatorTok{\textless{}}\BuiltInTok{X509Certificate}\OperatorTok{\textgreater{}} \FunctionTok{buildChain}\OperatorTok{(}
            \BuiltInTok{X509Certificate}\NormalTok{ endEntity}\OperatorTok{,}
            \BuiltInTok{Collection}\OperatorTok{\textless{}}\BuiltInTok{X509Certificate}\OperatorTok{\textgreater{}}\NormalTok{ candidates}\OperatorTok{)} \KeywordTok{throws} \BuiltInTok{Exception} \OperatorTok{\{}

        \BuiltInTok{List}\OperatorTok{\textless{}}\BuiltInTok{X509Certificate}\OperatorTok{\textgreater{}}\NormalTok{ chain }\OperatorTok{=} \KeywordTok{new} \BuiltInTok{ArrayList}\OperatorTok{\textless{}\textgreater{}();}
\NormalTok{        chain}\OperatorTok{.}\FunctionTok{add}\OperatorTok{(}\NormalTok{endEntity}\OperatorTok{);}

        \BuiltInTok{X509Certificate}\NormalTok{ current }\OperatorTok{=}\NormalTok{ endEntity}\OperatorTok{;}
        \DataTypeTok{int}\NormalTok{ maxDepth }\OperatorTok{=} \DecValTok{10}\OperatorTok{;} \CommentTok{// prevent infinite loops from malformed chains}

        \ControlFlowTok{for} \OperatorTok{(}\DataTypeTok{int}\NormalTok{ i }\OperatorTok{=} \DecValTok{0}\OperatorTok{;}\NormalTok{ i }\OperatorTok{\textless{}}\NormalTok{ maxDepth}\OperatorTok{;}\NormalTok{ i}\OperatorTok{++)} \OperatorTok{\{}
            \CommentTok{// Self{-}signed? We\textquotesingle{}ve reached a root.}
            \ControlFlowTok{if} \OperatorTok{(}\NormalTok{current}\OperatorTok{.}\FunctionTok{getSubjectX500Principal}\OperatorTok{()}
                    \OperatorTok{.}\FunctionTok{equals}\OperatorTok{(}\NormalTok{current}\OperatorTok{.}\FunctionTok{getIssuerX500Principal}\OperatorTok{()))} \OperatorTok{\{}
                \ControlFlowTok{break}\OperatorTok{;}
            \OperatorTok{\}}

            \DataTypeTok{byte}\OperatorTok{[]}\NormalTok{ akiBytes }\OperatorTok{=} \FunctionTok{getAuthorityKeyIdentifier}\OperatorTok{(}\NormalTok{current}\OperatorTok{);}
            \ControlFlowTok{if} \OperatorTok{(}\NormalTok{akiBytes }\OperatorTok{==} \KeywordTok{null}\OperatorTok{)} \OperatorTok{\{}
                \ControlFlowTok{break}\OperatorTok{;} \CommentTok{// no AKI {-}{-} cannot chain further}
            \OperatorTok{\}}

            \BuiltInTok{X509Certificate}\NormalTok{ issuer }\OperatorTok{=} \FunctionTok{findBySki}\OperatorTok{(}\NormalTok{candidates}\OperatorTok{,}\NormalTok{ akiBytes}\OperatorTok{);}
            \ControlFlowTok{if} \OperatorTok{(}\NormalTok{issuer }\OperatorTok{==} \KeywordTok{null}\OperatorTok{)} \OperatorTok{\{}
                \ControlFlowTok{break}\OperatorTok{;} \CommentTok{// issuer not in our candidate set}
            \OperatorTok{\}}

\NormalTok{            chain}\OperatorTok{.}\FunctionTok{add}\OperatorTok{(}\NormalTok{issuer}\OperatorTok{);}
\NormalTok{            current }\OperatorTok{=}\NormalTok{ issuer}\OperatorTok{;}
        \OperatorTok{\}}

        \ControlFlowTok{return}\NormalTok{ chain}\OperatorTok{;}
    \OperatorTok{\}}

    \KeywordTok{private} \DataTypeTok{static} \DataTypeTok{byte}\OperatorTok{[]} \FunctionTok{getAuthorityKeyIdentifier}\OperatorTok{(}
            \BuiltInTok{X509Certificate}\NormalTok{ cert}\OperatorTok{)} \KeywordTok{throws} \BuiltInTok{Exception} \OperatorTok{\{}

\NormalTok{        X509CertificateHolder holder }\OperatorTok{=} \KeywordTok{new} \FunctionTok{JcaX509CertificateHolder}\OperatorTok{(}\NormalTok{cert}\OperatorTok{);}
\NormalTok{        Extension akiExt }\OperatorTok{=}\NormalTok{ holder}\OperatorTok{.}\FunctionTok{getExtension}\OperatorTok{(}
\NormalTok{                Extension}\OperatorTok{.}\FunctionTok{authorityKeyIdentifier}\OperatorTok{);}
        \ControlFlowTok{if} \OperatorTok{(}\NormalTok{akiExt }\OperatorTok{==} \KeywordTok{null}\OperatorTok{)} \OperatorTok{\{}
            \ControlFlowTok{return} \KeywordTok{null}\OperatorTok{;}
        \OperatorTok{\}}
\NormalTok{        AuthorityKeyIdentifier aki }\OperatorTok{=}
\NormalTok{                AuthorityKeyIdentifier}\OperatorTok{.}\FunctionTok{getInstance}\OperatorTok{(}\NormalTok{akiExt}\OperatorTok{.}\FunctionTok{getParsedValue}\OperatorTok{());}
        \ControlFlowTok{return}\NormalTok{ aki}\OperatorTok{.}\FunctionTok{getKeyIdentifier}\OperatorTok{();}
    \OperatorTok{\}}

    \KeywordTok{private} \DataTypeTok{static} \BuiltInTok{X509Certificate} \FunctionTok{findBySki}\OperatorTok{(}
            \BuiltInTok{Collection}\OperatorTok{\textless{}}\BuiltInTok{X509Certificate}\OperatorTok{\textgreater{}}\NormalTok{ certs}\OperatorTok{,}
            \DataTypeTok{byte}\OperatorTok{[]}\NormalTok{ targetSki}\OperatorTok{)} \KeywordTok{throws} \BuiltInTok{Exception} \OperatorTok{\{}

        \ControlFlowTok{for} \OperatorTok{(}\BuiltInTok{X509Certificate}\NormalTok{ candidate }\OperatorTok{:}\NormalTok{ certs}\OperatorTok{)} \OperatorTok{\{}
\NormalTok{            X509CertificateHolder holder }\OperatorTok{=}
                    \KeywordTok{new} \FunctionTok{JcaX509CertificateHolder}\OperatorTok{(}\NormalTok{candidate}\OperatorTok{);}
\NormalTok{            Extension skiExt }\OperatorTok{=}\NormalTok{ holder}\OperatorTok{.}\FunctionTok{getExtension}\OperatorTok{(}
\NormalTok{                    Extension}\OperatorTok{.}\FunctionTok{subjectKeyIdentifier}\OperatorTok{);}
            \ControlFlowTok{if} \OperatorTok{(}\NormalTok{skiExt }\OperatorTok{==} \KeywordTok{null}\OperatorTok{)} \OperatorTok{\{}
                \ControlFlowTok{continue}\OperatorTok{;}
            \OperatorTok{\}}
\NormalTok{            SubjectKeyIdentifier ski }\OperatorTok{=}
\NormalTok{                    SubjectKeyIdentifier}\OperatorTok{.}\FunctionTok{getInstance}\OperatorTok{(}\NormalTok{skiExt}\OperatorTok{.}\FunctionTok{getParsedValue}\OperatorTok{());}
            \ControlFlowTok{if} \OperatorTok{(}\BuiltInTok{Arrays}\OperatorTok{.}\FunctionTok{equals}\OperatorTok{(}\NormalTok{ski}\OperatorTok{.}\FunctionTok{getKeyIdentifier}\OperatorTok{(),}\NormalTok{ targetSki}\OperatorTok{))} \OperatorTok{\{}
                \ControlFlowTok{return}\NormalTok{ candidate}\OperatorTok{;}
            \OperatorTok{\}}
        \OperatorTok{\}}
        \ControlFlowTok{return} \KeywordTok{null}\OperatorTok{;}
    \OperatorTok{\}}
\OperatorTok{\}}
\end{Highlighting}
\end{Shaded}

In practice, chain building can fail for several reasons. The AKI
extension might be absent (legal but rare in modern certificates). The
SKI might not match because the issuing CA generated the SKI using a
non-standard method (SHA-1 of the public key is the norm, but RFC 5280
--- Internet X.509 PKI Certificate and CRL Profile does not mandate it).
Or the intermediate certificate might simply not be available. When
debugging chain failures, always dump both the AKI of the child and the
SKI of the expected parent, and compare them byte by byte.

\subsection{Fallback: Issuer Distinguished Name
Matching}\label{_fallback_issuer_distinguished_name_matching}

If the AKI/SKI matching fails --- or if the certificates lack these
extensions (which happens with very old CAs) --- the fallback is
matching the certificate's \textbf{Issuer DN} against candidate
certificates\textquotesingle{} \textbf{Subject DN}. This is less
reliable because distinguished name comparison is notoriously tricky:
leading/trailing whitespace, different string encodings (PrintableString
vs UTF8String), case sensitivity of certain components --- all of these
can cause a mismatch even when the certificates are part of the same
chain.

RFC 5280 --- Internet X.509 PKI Certificate and CRL Profile section 7.1
defines the comparison rules for distinguished names, and they are
considerably more complex than string equality. BouncyCastle's
\texttt{X500Name.equals()} handles this correctly; Java's
\texttt{X500Principal.equals()} usually does too, but edge cases exist
with certain encoding combinations.

\section{Path Validation: RFC 5280 Section
6}\label{_path_validation_rfc_5280_section_6}

Building the chain gives you an ordered list of certificates.
\textbf{Path validation} verifies that this chain is actually
valid --- that every link is cryptographically sound and that the
certificates meet all the required constraints.

RFC 5280 --- Internet X.509 PKI Certificate and CRL Profile section 6
defines the algorithm in six phases. The full specification runs to
several pages of pseudocode; here we distill the essential checks that
matter for PDF signature validation.

\subsection{The Six Steps}\label{_the_six_steps}

\begin{longtable}[]{@{}
  >{\raggedright\arraybackslash}p{(\linewidth - 4\tabcolsep) * \real{0.1250}}
  >{\raggedright\arraybackslash}p{(\linewidth - 4\tabcolsep) * \real{0.2500}}
  >{\raggedright\arraybackslash}p{(\linewidth - 4\tabcolsep) * \real{0.6250}}@{}}
\caption{Certificate path validation (RFC 5280 --- Internet X.509 PKI
Certificate and CRL Profile section 6)}\tabularnewline
\toprule\noalign{}
\begin{minipage}[b]{\linewidth}\raggedright
Step
\end{minipage} & \begin{minipage}[b]{\linewidth}\raggedright
Name
\end{minipage} & \begin{minipage}[b]{\linewidth}\raggedright
What It Checks
\end{minipage} \\
\midrule\noalign{}
\endfirsthead
\toprule\noalign{}
\begin{minipage}[b]{\linewidth}\raggedright
Step
\end{minipage} & \begin{minipage}[b]{\linewidth}\raggedright
Name
\end{minipage} & \begin{minipage}[b]{\linewidth}\raggedright
What It Checks
\end{minipage} \\
\midrule\noalign{}
\endhead
\bottomrule\noalign{}
\endlastfoot
1 & Initialization & Set the initial trust anchor (root CA certificate),
the expected issuer name, the expected public key, and policy
constraints. These are the "inputs" to the validation algorithm. \\
2 & Basic processing & For each certificate in the chain (from root to
leaf): verify the signature on the certificate using the issuer's public
key, check that the certificate is within its validity period (notBefore
⇐ current time ⇐ notAfter), check that the certificate has not been
revoked. \\
3 & Name constraints & Check that the subject name and any subject
alternative names comply with any name constraints imposed by
certificates higher in the chain. A CA can restrict its subordinates to
only issue certificates for certain domains or email addresses. \\
4 & Policy processing & Process the certificate policies extension and
any policy constraints. This determines which policies are valid for the
entire chain. Policy processing is the most complex part of path
validation and is rarely used in PDF signature contexts --- but some
government PKIs rely on it. \\
5 & Basic constraints & Verify that every CA certificate in the chain
has \texttt{basicConstraints} with \texttt{cA=TRUE}, and that the
\texttt{pathLenConstraint} is not exceeded. A \texttt{pathLenConstraint}
of 0 means the CA can only issue end-entity certificates, not
sub-CAs. \\
6 & Wrap-up & After processing all certificates, check the final
certificate (the end-entity) against the expected usage. For PDF
signing, this means checking \texttt{keyUsage} for
\texttt{digitalSignature} or \texttt{nonRepudiation}, and optionally
checking \texttt{extendedKeyUsage}. \\
\end{longtable}

\begin{Shaded}
\begin{Highlighting}[]
\KeywordTok{import} \ImportTok{java}\OperatorTok{.}\ImportTok{security}\OperatorTok{.}\ImportTok{cert}\OperatorTok{.}\ImportTok{CertPathValidator}\OperatorTok{;}
\KeywordTok{import} \ImportTok{java}\OperatorTok{.}\ImportTok{security}\OperatorTok{.}\ImportTok{cert}\OperatorTok{.}\ImportTok{CertificateFactory}\OperatorTok{;}
\KeywordTok{import} \ImportTok{java}\OperatorTok{.}\ImportTok{security}\OperatorTok{.}\ImportTok{cert}\OperatorTok{.}\ImportTok{CertPath}\OperatorTok{;}
\KeywordTok{import} \ImportTok{java}\OperatorTok{.}\ImportTok{security}\OperatorTok{.}\ImportTok{cert}\OperatorTok{.}\ImportTok{PKIXParameters}\OperatorTok{;}
\KeywordTok{import} \ImportTok{java}\OperatorTok{.}\ImportTok{security}\OperatorTok{.}\ImportTok{cert}\OperatorTok{.}\ImportTok{PKIXCertPathValidatorResult}\OperatorTok{;}
\KeywordTok{import} \ImportTok{java}\OperatorTok{.}\ImportTok{security}\OperatorTok{.}\ImportTok{cert}\OperatorTok{.}\ImportTok{TrustAnchor}\OperatorTok{;}
\KeywordTok{import} \ImportTok{java}\OperatorTok{.}\ImportTok{security}\OperatorTok{.}\ImportTok{cert}\OperatorTok{.}\ImportTok{X509Certificate}\OperatorTok{;}
\KeywordTok{import} \ImportTok{java}\OperatorTok{.}\ImportTok{util}\OperatorTok{.}\ImportTok{Collections}\OperatorTok{;}
\KeywordTok{import} \ImportTok{java}\OperatorTok{.}\ImportTok{util}\OperatorTok{.}\ImportTok{List}\OperatorTok{;}

\CommentTok{/**}
 \CommentTok{*}\NormalTok{ Validate a certificate chain using the JCA PKIX path validator}\CommentTok{.}
 \CommentTok{*}\NormalTok{ This performs the full RFC }\CommentTok{5280}\NormalTok{ section }\CommentTok{6}\NormalTok{ algorithm}\CommentTok{.}
 \CommentTok{*/}
\KeywordTok{public} \KeywordTok{class}\NormalTok{ PathValidationExample }\OperatorTok{\{}

    \KeywordTok{public} \DataTypeTok{static} \BuiltInTok{PKIXCertPathValidatorResult} \FunctionTok{validateChain}\OperatorTok{(}
            \BuiltInTok{List}\OperatorTok{\textless{}}\BuiltInTok{X509Certificate}\OperatorTok{\textgreater{}}\NormalTok{ chain}\OperatorTok{,}
            \BuiltInTok{X509Certificate}\NormalTok{ trustAnchor}\OperatorTok{)} \KeywordTok{throws} \BuiltInTok{Exception} \OperatorTok{\{}

        \CommentTok{// The chain list should NOT include the trust anchor}
        \BuiltInTok{CertificateFactory}\NormalTok{ cf }\OperatorTok{=} \BuiltInTok{CertificateFactory}\OperatorTok{.}\FunctionTok{getInstance}\OperatorTok{(}\StringTok{"X.509"}\OperatorTok{);}
        \BuiltInTok{CertPath}\NormalTok{ certPath }\OperatorTok{=}\NormalTok{ cf}\OperatorTok{.}\FunctionTok{generateCertPath}\OperatorTok{(}\NormalTok{chain}\OperatorTok{);}

        \BuiltInTok{TrustAnchor}\NormalTok{ anchor }\OperatorTok{=} \KeywordTok{new} \BuiltInTok{TrustAnchor}\OperatorTok{(}\NormalTok{trustAnchor}\OperatorTok{,} \KeywordTok{null}\OperatorTok{);}
        \BuiltInTok{PKIXParameters}\NormalTok{ params }\OperatorTok{=} \KeywordTok{new} \BuiltInTok{PKIXParameters}\OperatorTok{(}
                \BuiltInTok{Collections}\OperatorTok{.}\FunctionTok{singleton}\OperatorTok{(}\NormalTok{anchor}\OperatorTok{));}

        \CommentTok{// Disable revocation checking here {-}{-} we handle it separately}
\NormalTok{        params}\OperatorTok{.}\FunctionTok{setRevocationEnabled}\OperatorTok{(}\KeywordTok{false}\OperatorTok{);}

        \BuiltInTok{CertPathValidator}\NormalTok{ validator }\OperatorTok{=}
                \BuiltInTok{CertPathValidator}\OperatorTok{.}\FunctionTok{getInstance}\OperatorTok{(}\StringTok{"PKIX"}\OperatorTok{);}

        \ControlFlowTok{return} \OperatorTok{(}\BuiltInTok{PKIXCertPathValidatorResult}\OperatorTok{)}
\NormalTok{                validator}\OperatorTok{.}\FunctionTok{validate}\OperatorTok{(}\NormalTok{certPath}\OperatorTok{,}\NormalTok{ params}\OperatorTok{);}
    \OperatorTok{\}}
\OperatorTok{\}}
\end{Highlighting}
\end{Shaded}

The JCA \texttt{CertPathValidator} expects the chain in a specific
order: \textbf{end-entity first, root last (or root omitted)}. If you
pass the chain in the wrong order, validation will fail with a confusing
error message about "unable to find valid certification path to
requested target." This is the single most common mistake when using the
JCA path validation API.

\subsection{Validity Period Checking and Signing
Time}\label{_validity_period_checking_and_signing_time}

There is a subtlety in validity period checking that matters enormously
for PDF signatures. A certificate might be expired \emph{now} but was
valid at the time the document was signed. Should the signature be
considered valid?

The answer depends on whether the signature includes a trusted
timestamp. If the CMS \texttt{SignedData} contains a timestamp token in
the unsigned attributes (see Chapter 18), and that timestamp was
produced while the certificate was still valid, then the signature
remains valid even after the certificate expires. This is the foundation
of \textbf{long-term validation (LTV)} --- the mechanism that allows a
signed PDF to remain verifiable for decades.

Without a timestamp, most validators check the certificate against the
current time. A signature made with a certificate that has since expired
will show warnings or errors in Adobe Acrobat, even if the signature was
perfectly valid when it was created.

This is why timestamps matter so much for PDF signatures. A timestamp
cryptographically proves \emph{when} the signature was created, allowing
validators to check the certificate against the signing time rather than
the current time. If you are building a signing service, always include
a timestamp. Chapter 18 covers timestamps and LTV in exhaustive detail.

\section{Revocation Checking}\label{_revocation_checking}

A certificate may be revoked before it expires. The private key might
have been compromised, the subscriber might have left the organization,
or the CA might have discovered that the certificate was issued
incorrectly. Revocation checking is the mechanism by which validators
discover whether a certificate has been revoked.

There are two mechanisms: \textbf{Certificate Revocation Lists (CRLs)}
and the \textbf{Online Certificate Status Protocol (OCSP)}. Both have
strengths and weaknesses, and a robust validator should support both.

\subsection{CRL Deep Dive}\label{_crl_deep_dive}

A Certificate Revocation List is a signed data structure published by a
CA that lists all certificates the CA has revoked. CRLs are defined in
RFC 5280 --- Internet X.509 PKI Certificate and CRL Profile section 5.

\subsubsection{CRL Structure}\label{_crl_structure}

A CRL contains:

\begin{longtable}[]{@{}
  >{\raggedright\arraybackslash}p{(\linewidth - 2\tabcolsep) * \real{0.3333}}
  >{\raggedright\arraybackslash}p{(\linewidth - 2\tabcolsep) * \real{0.6667}}@{}}
\caption{CRL structure (simplified from RFC 5280 --- Internet X.509 PKI
Certificate and CRL Profile section 5.1)}\tabularnewline
\toprule\noalign{}
\begin{minipage}[b]{\linewidth}\raggedright
Field
\end{minipage} & \begin{minipage}[b]{\linewidth}\raggedright
Description
\end{minipage} \\
\midrule\noalign{}
\endfirsthead
\toprule\noalign{}
\begin{minipage}[b]{\linewidth}\raggedright
Field
\end{minipage} & \begin{minipage}[b]{\linewidth}\raggedright
Description
\end{minipage} \\
\midrule\noalign{}
\endhead
\bottomrule\noalign{}
\endlastfoot
\texttt{tbsCertList} & The "to-be-signed" portion, containing the issuer
name, the \texttt{thisUpdate} timestamp (when the CRL was issued), the
\texttt{nextUpdate} timestamp (when the next CRL will be issued), and
the list of revoked certificates. \\
\texttt{signatureAlgorithm} & The algorithm used to sign the CRL (e.g.,
SHA-256 with RSA). \\
\texttt{signatureValue} & The CA's signature over the
\texttt{tbsCertList}, proving the CRL is authentic and unmodified. \\
\end{longtable}

Each entry in the \texttt{revokedCertificates} list contains:

\begin{itemize}
\item
  The \textbf{serial number} of the revoked certificate
\item
  The \textbf{revocation date}
\item
  Optional \textbf{CRL entry extensions}, including the
  \texttt{reasonCode} (keyCompromise, cACompromise, affiliationChanged,
  superseded, cessationOfOperation, certificateHold, removeFromCRL,
  privilegeWithdrawn, aACompromise)
\end{itemize}

\subsubsection{CRL Distribution Points}\label{_crl_distribution_points}

How does a validator find the CRL? The certificate itself contains a
\textbf{CRL Distribution Points} extension (RFC 5280 --- Internet X.509
PKI Certificate and CRL Profile section 4.2.1.13) with one or more URLs
where the CRL can be downloaded. Typically these are HTTP URLs, though
LDAP URLs appear in older enterprise PKIs.

\begin{Shaded}
\begin{Highlighting}[]
\NormalTok{CRLDistributionPoints ::= }\DataTypeTok{SEQUENCE} \KeywordTok{SIZE}\NormalTok{ (}\DecValTok{1}\NormalTok{..MAX) }\DataTypeTok{OF}\NormalTok{ DistributionPoint}

\NormalTok{DistributionPoint ::= }\DataTypeTok{SEQUENCE}\NormalTok{ \{}
\NormalTok{    distributionPoint  [}\DecValTok{0}\NormalTok{] DistributionPointName }\KeywordTok{OPTIONAL}\NormalTok{,}
\NormalTok{    reasons            [}\DecValTok{1}\NormalTok{] ReasonFlags }\KeywordTok{OPTIONAL}\NormalTok{,}
\NormalTok{    cRLIssuer          [}\DecValTok{2}\NormalTok{] GeneralNames }\KeywordTok{OPTIONAL}
\NormalTok{\}}
\end{Highlighting}
\end{Shaded}

\subsubsection{The Size Problem}\label{_the_size_problem}

CRLs have a fundamental scalability problem: they grow monotonically.
Every revoked certificate is listed until the certificate's natural
expiration date passes. A busy CA that revokes hundreds or thousands of
certificates per year will produce CRLs that are megabytes in size.

Let's Encrypt, for example, issues millions of certificates. A full CRL
from Let's Encrypt would be enormous. This is why Let's Encrypt does not
publish CRLs for individual certificates (they rely on OCSP and, more
recently, CRL sharding).

For PDF signatures, the CRL size matters because long-term validation
(LTV) requires embedding revocation data \emph{inside the PDF}. A 2 MB
CRL embedded in the Document Security Store (DSS) inflates the PDF by 2
MB. If the chain has three certificates and each CRL is large, the
overhead becomes significant.

\subsubsection{Delta CRLs}\label{_delta_crls}

Delta CRLs (RFC 5280 --- Internet X.509 PKI Certificate and CRL Profile
section 5.2.4) were designed to mitigate the size problem. A delta CRL
contains only the changes since the last "base" CRL was published. In
theory, this allows validators to download a small delta rather than a
large full CRL.

In practice, delta CRLs are rarely used. Most CAs do not publish them.
Most validators do not request them. The overhead of tracking base/delta
CRL relationships adds complexity that the industry has largely chosen
to avoid, preferring OCSP for real-time status and full CRLs for
archival purposes.

\subsection{OCSP Deep Dive}\label{_ocsp_deep_dive}

The \textbf{Online Certificate Status Protocol} (RFC 6960 --- Online
Certificate Status Protocol (OCSP)) provides real-time, per-certificate
revocation status. Instead of downloading a list of all revoked
certificates, the validator sends a request asking "is this specific
certificate revoked?" and gets back a signed response.

\subsubsection{OCSP Request Structure}\label{_ocsp_request_structure}

An OCSP request identifies the certificate by:

\begin{itemize}
\item
  The hash of the issuer's distinguished name
\item
  The hash of the issuer's public key
\item
  The certificate's serial number
\item
  An optional nonce to prevent replay attacks
\end{itemize}

\begin{Shaded}
\begin{Highlighting}[]
\NormalTok{OCSPRequest ::= }\DataTypeTok{SEQUENCE}\NormalTok{ \{}
\NormalTok{    tbsRequest     TBSRequest,}
\NormalTok{    optionalSignature [}\DecValTok{0}\NormalTok{] EXPLICIT Signature }\KeywordTok{OPTIONAL}
\NormalTok{\}}

\NormalTok{TBSRequest ::= }\DataTypeTok{SEQUENCE}\NormalTok{ \{}
\NormalTok{    version        [}\DecValTok{0}\NormalTok{] EXPLICIT Version }\KeywordTok{DEFAULT}\NormalTok{ v1,}
\NormalTok{    requestorName  [}\DecValTok{1}\NormalTok{] EXPLICIT GeneralName }\KeywordTok{OPTIONAL}\NormalTok{,}
\NormalTok{    requestList    }\DataTypeTok{SEQUENCE} \DataTypeTok{OF}\NormalTok{ Request,}
\NormalTok{    requestExtensions [}\DecValTok{2}\NormalTok{] EXPLICIT Extensions }\KeywordTok{OPTIONAL}
\NormalTok{\}}

\NormalTok{Request ::= }\DataTypeTok{SEQUENCE}\NormalTok{ \{}
\NormalTok{    reqCert        CertID,}
\NormalTok{    singleRequestExtensions [}\DecValTok{0}\NormalTok{] EXPLICIT Extensions }\KeywordTok{OPTIONAL}
\NormalTok{\}}

\NormalTok{CertID ::= }\DataTypeTok{SEQUENCE}\NormalTok{ \{}
\NormalTok{    hashAlgorithm  AlgorithmIdentifier,}
\NormalTok{    issuerNameHash }\DataTypeTok{OCTET} \DataTypeTok{STRING}\NormalTok{,}
\NormalTok{    issuerKeyHash  }\DataTypeTok{OCTET} \DataTypeTok{STRING}\NormalTok{,}
\NormalTok{    serialNumber   CertificateSerialNumber}
\NormalTok{\}}
\end{Highlighting}
\end{Shaded}

\subsubsection{OCSP Response Structure}\label{_ocsp_response_structure}

The responder returns one of: \textbf{good} (not revoked),
\textbf{revoked} (with the revocation time and optional reason), or
\textbf{unknown} (the responder does not know about this certificate).

The response is signed, either by the CA itself or by a specially
designated \textbf{OCSP responder certificate} --- a certificate issued
by the CA with the \texttt{id-pkix-ocsp-nocheck} extension, indicating
it is authorized to sign OCSP responses on behalf of the CA.

\subsubsection{The Nonce Problem}\label{_the_nonce_problem}

OCSP supports a nonce extension: the requester includes a random value,
and the responder echoes it back. This prevents an attacker from
replaying an old "good" response for a certificate that has since been
revoked.

In practice, many OCSP responders \textbf{ignore the nonce}. They
pre-compute responses for all active certificates and serve them from
cache, which is incompatible with per-request nonces. Adobe's OCSP
responder, for example, does not support nonces. The security trade-off
is real: without nonces, a replay window exists equal to the response's
\texttt{nextUpdate} time. For most PDF signature scenarios, this is
acceptable --- the replay window is typically hours, not days.

When embedding OCSP responses in a PDF for LTV, nonce-less responses are
actually preferable. A response with a nonce is bound to a specific
request and cannot be reused. A cached, nonce-less response can be
embedded in the PDF and validated later by anyone, which is exactly what
LTV requires. This is a case where the "less secure" option (no nonce)
is the correct choice for the specific use case.

\subsubsection{OCSP Stapling}\label{_ocsp_stapling}

\textbf{OCSP stapling} is a mechanism where the entity presenting a
certificate also presents a recent OCSP response for that certificate,
pre-fetched from the OCSP responder. In TLS, the server staples the OCSP
response to the TLS handshake.

For PDF signatures, the equivalent is embedding the OCSP response in the
CMS \texttt{SignedData} structure or in the PDF's Document Security
Store (DSS). This eliminates the need for the validator to contact the
OCSP responder at all --- the revocation proof travels with the
document, just like the certificate chain.

\subsection{CRL vs OCSP: When to Use
Which}\label{_crl_vs_ocsp_when_to_use_which}

\begin{longtable}[]{@{}
  >{\raggedright\arraybackslash}p{(\linewidth - 4\tabcolsep) * \real{0.2500}}
  >{\raggedright\arraybackslash}p{(\linewidth - 4\tabcolsep) * \real{0.3750}}
  >{\raggedright\arraybackslash}p{(\linewidth - 4\tabcolsep) * \real{0.3750}}@{}}
\caption{CRL vs OCSP comparison for PDF signatures}\tabularnewline
\toprule\noalign{}
\begin{minipage}[b]{\linewidth}\raggedright
Aspect
\end{minipage} & \begin{minipage}[b]{\linewidth}\raggedright
CRL
\end{minipage} & \begin{minipage}[b]{\linewidth}\raggedright
OCSP
\end{minipage} \\
\midrule\noalign{}
\endfirsthead
\toprule\noalign{}
\begin{minipage}[b]{\linewidth}\raggedright
Aspect
\end{minipage} & \begin{minipage}[b]{\linewidth}\raggedright
CRL
\end{minipage} & \begin{minipage}[b]{\linewidth}\raggedright
OCSP
\end{minipage} \\
\midrule\noalign{}
\endhead
\bottomrule\noalign{}
\endlastfoot
Granularity & All revocations for a CA & Single certificate status \\
Freshness & Updated periodically (hours to days) & Real-time or
near-real-time (minutes) \\
Size & Can be megabytes (proportional to revocations) & Typically 1-3 KB
per response \\
Availability & Can be cached locally; works offline once downloaded &
Requires network access (unless stapled/embedded) \\
LTV suitability & Good: can be embedded in DSS & Good: can be embedded
in DSS (prefer nonce-less) \\
Privacy & Good: the CA does not learn which certificates you check &
Poor: the OCSP responder learns exactly which certificates you are
validating (which documents you are opening) \\
\end{longtable}

For PDF signing services, use OCSP for real-time validation during
signing, and embed both the OCSP response and relevant CRLs in the DSS
for long-term validation. Belt and suspenders. Adobe Acrobat's LTV
process follows exactly this approach.

\begin{Shaded}
\begin{Highlighting}[]
\KeywordTok{import} \ImportTok{java}\OperatorTok{.}\ImportTok{io}\OperatorTok{.}\ImportTok{InputStream}\OperatorTok{;}
\KeywordTok{import} \ImportTok{java}\OperatorTok{.}\ImportTok{io}\OperatorTok{.}\ImportTok{OutputStream}\OperatorTok{;}
\KeywordTok{import} \ImportTok{java}\OperatorTok{.}\ImportTok{math}\OperatorTok{.}\ImportTok{BigInteger}\OperatorTok{;}
\KeywordTok{import} \ImportTok{java}\OperatorTok{.}\ImportTok{net}\OperatorTok{.}\ImportTok{HttpURLConnection}\OperatorTok{;}
\KeywordTok{import} \ImportTok{java}\OperatorTok{.}\ImportTok{net}\OperatorTok{.}\ImportTok{URI}\OperatorTok{;}
\KeywordTok{import} \ImportTok{java}\OperatorTok{.}\ImportTok{security}\OperatorTok{.}\ImportTok{cert}\OperatorTok{.}\ImportTok{X509Certificate}\OperatorTok{;}

\KeywordTok{import} \ImportTok{org}\OperatorTok{.}\ImportTok{bouncycastle}\OperatorTok{.}\ImportTok{asn1}\OperatorTok{.}\ImportTok{DEROctetString}\OperatorTok{;}
\KeywordTok{import} \ImportTok{org}\OperatorTok{.}\ImportTok{bouncycastle}\OperatorTok{.}\ImportTok{asn1}\OperatorTok{.}\ImportTok{ocsp}\OperatorTok{.}\ImportTok{OCSPObjectIdentifiers}\OperatorTok{;}
\KeywordTok{import} \ImportTok{org}\OperatorTok{.}\ImportTok{bouncycastle}\OperatorTok{.}\ImportTok{asn1}\OperatorTok{.}\ImportTok{x509}\OperatorTok{.}\ImportTok{Extension}\OperatorTok{;}
\KeywordTok{import} \ImportTok{org}\OperatorTok{.}\ImportTok{bouncycastle}\OperatorTok{.}\ImportTok{asn1}\OperatorTok{.}\ImportTok{x509}\OperatorTok{.}\ImportTok{Extensions}\OperatorTok{;}
\KeywordTok{import} \ImportTok{org}\OperatorTok{.}\ImportTok{bouncycastle}\OperatorTok{.}\ImportTok{cert}\OperatorTok{.}\ImportTok{jcajce}\OperatorTok{.}\ImportTok{JcaX509CertificateHolder}\OperatorTok{;}
\KeywordTok{import} \ImportTok{org}\OperatorTok{.}\ImportTok{bouncycastle}\OperatorTok{.}\ImportTok{cert}\OperatorTok{.}\ImportTok{ocsp}\OperatorTok{.}\ImportTok{BasicOCSPResp}\OperatorTok{;}
\KeywordTok{import} \ImportTok{org}\OperatorTok{.}\ImportTok{bouncycastle}\OperatorTok{.}\ImportTok{cert}\OperatorTok{.}\ImportTok{ocsp}\OperatorTok{.}\ImportTok{CertificateID}\OperatorTok{;}
\KeywordTok{import} \ImportTok{org}\OperatorTok{.}\ImportTok{bouncycastle}\OperatorTok{.}\ImportTok{cert}\OperatorTok{.}\ImportTok{ocsp}\OperatorTok{.}\ImportTok{OCSPReq}\OperatorTok{;}
\KeywordTok{import} \ImportTok{org}\OperatorTok{.}\ImportTok{bouncycastle}\OperatorTok{.}\ImportTok{cert}\OperatorTok{.}\ImportTok{ocsp}\OperatorTok{.}\ImportTok{OCSPReqBuilder}\OperatorTok{;}
\KeywordTok{import} \ImportTok{org}\OperatorTok{.}\ImportTok{bouncycastle}\OperatorTok{.}\ImportTok{cert}\OperatorTok{.}\ImportTok{ocsp}\OperatorTok{.}\ImportTok{OCSPResp}\OperatorTok{;}
\KeywordTok{import} \ImportTok{org}\OperatorTok{.}\ImportTok{bouncycastle}\OperatorTok{.}\ImportTok{cert}\OperatorTok{.}\ImportTok{ocsp}\OperatorTok{.}\ImportTok{SingleResp}\OperatorTok{;}
\KeywordTok{import} \ImportTok{org}\OperatorTok{.}\ImportTok{bouncycastle}\OperatorTok{.}\ImportTok{operator}\OperatorTok{.}\ImportTok{jcajce}\OperatorTok{.}\ImportTok{JcaDigestCalculatorProviderBuilder}\OperatorTok{;}

\CommentTok{/**}
 \CommentTok{*}\NormalTok{ Send an OCSP request and check the response}\CommentTok{.}
 \CommentTok{*}\NormalTok{ Demonstrates the full OCSP round}\CommentTok{{-}}\NormalTok{trip with BouncyCastle}\CommentTok{.}
 \CommentTok{*/}
\KeywordTok{public} \KeywordTok{class}\NormalTok{ OcspCheckExample }\OperatorTok{\{}

    \KeywordTok{public} \DataTypeTok{static}\NormalTok{ SingleResp }\FunctionTok{checkOcsp}\OperatorTok{(}
            \BuiltInTok{X509Certificate}\NormalTok{ cert}\OperatorTok{,}
            \BuiltInTok{X509Certificate}\NormalTok{ issuer}\OperatorTok{,}
            \BuiltInTok{String}\NormalTok{ ocspUrl}\OperatorTok{)} \KeywordTok{throws} \BuiltInTok{Exception} \OperatorTok{\{}

        \CommentTok{// Build the CertificateID that identifies the cert to check}
\NormalTok{        CertificateID certId }\OperatorTok{=} \KeywordTok{new} \FunctionTok{CertificateID}\OperatorTok{(}
                \KeywordTok{new} \FunctionTok{JcaDigestCalculatorProviderBuilder}\OperatorTok{()}
                        \OperatorTok{.}\FunctionTok{build}\OperatorTok{()}
                        \OperatorTok{.}\FunctionTok{get}\OperatorTok{(}\NormalTok{CertificateID}\OperatorTok{.}\FunctionTok{HASH\_SHA1}\OperatorTok{),}
                \KeywordTok{new} \FunctionTok{JcaX509CertificateHolder}\OperatorTok{(}\NormalTok{issuer}\OperatorTok{),}
\NormalTok{                cert}\OperatorTok{.}\FunctionTok{getSerialNumber}\OperatorTok{());}

        \CommentTok{// Build the OCSP request}
\NormalTok{        OCSPReqBuilder reqBuilder }\OperatorTok{=} \KeywordTok{new} \FunctionTok{OCSPReqBuilder}\OperatorTok{();}
\NormalTok{        reqBuilder}\OperatorTok{.}\FunctionTok{addRequest}\OperatorTok{(}\NormalTok{certId}\OperatorTok{);}

        \CommentTok{// Add a nonce (optional {-}{-} many responders ignore it)}
        \BuiltInTok{BigInteger}\NormalTok{ nonce }\OperatorTok{=} \BuiltInTok{BigInteger}\OperatorTok{.}\FunctionTok{valueOf}\OperatorTok{(}\BuiltInTok{System}\OperatorTok{.}\FunctionTok{currentTimeMillis}\OperatorTok{());}
\NormalTok{        Extension nonceExt }\OperatorTok{=} \KeywordTok{new} \FunctionTok{Extension}\OperatorTok{(}
\NormalTok{                OCSPObjectIdentifiers}\OperatorTok{.}\FunctionTok{id\_pkix\_ocsp\_nonce}\OperatorTok{,}
                \KeywordTok{false}\OperatorTok{,}
                \KeywordTok{new} \FunctionTok{DEROctetString}\OperatorTok{(}\NormalTok{nonce}\OperatorTok{.}\FunctionTok{toByteArray}\OperatorTok{()));}
\NormalTok{        reqBuilder}\OperatorTok{.}\FunctionTok{setRequestExtensions}\OperatorTok{(}
                \KeywordTok{new} \FunctionTok{Extensions}\OperatorTok{(}\KeywordTok{new}\NormalTok{ Extension}\OperatorTok{[]\{}\NormalTok{ nonceExt }\OperatorTok{\}));}

\NormalTok{        OCSPReq ocspReq }\OperatorTok{=}\NormalTok{ reqBuilder}\OperatorTok{.}\FunctionTok{build}\OperatorTok{();}

        \CommentTok{// Send the request via HTTP POST}
        \BuiltInTok{URI}\NormalTok{ uri }\OperatorTok{=} \KeywordTok{new} \BuiltInTok{URI}\OperatorTok{(}\NormalTok{ocspUrl}\OperatorTok{);}
        \BuiltInTok{HttpURLConnection}\NormalTok{ conn }\OperatorTok{=}
                \OperatorTok{(}\BuiltInTok{HttpURLConnection}\OperatorTok{)}\NormalTok{ uri}\OperatorTok{.}\FunctionTok{toURL}\OperatorTok{().}\FunctionTok{openConnection}\OperatorTok{();}
\NormalTok{        conn}\OperatorTok{.}\FunctionTok{setRequestMethod}\OperatorTok{(}\StringTok{"POST"}\OperatorTok{);}
\NormalTok{        conn}\OperatorTok{.}\FunctionTok{setRequestProperty}\OperatorTok{(}\StringTok{"Content{-}Type"}\OperatorTok{,}
                \StringTok{"application/ocsp{-}request"}\OperatorTok{);}
\NormalTok{        conn}\OperatorTok{.}\FunctionTok{setDoOutput}\OperatorTok{(}\KeywordTok{true}\OperatorTok{);}

        \ControlFlowTok{try} \OperatorTok{(}\BuiltInTok{OutputStream}\NormalTok{ os }\OperatorTok{=}\NormalTok{ conn}\OperatorTok{.}\FunctionTok{getOutputStream}\OperatorTok{())} \OperatorTok{\{}
\NormalTok{            os}\OperatorTok{.}\FunctionTok{write}\OperatorTok{(}\NormalTok{ocspReq}\OperatorTok{.}\FunctionTok{getEncoded}\OperatorTok{());}
        \OperatorTok{\}}

        \CommentTok{// Parse the response}
\NormalTok{        OCSPResp ocspResp}\OperatorTok{;}
        \ControlFlowTok{try} \OperatorTok{(}\BuiltInTok{InputStream}\NormalTok{ is }\OperatorTok{=}\NormalTok{ conn}\OperatorTok{.}\FunctionTok{getInputStream}\OperatorTok{())} \OperatorTok{\{}
\NormalTok{            ocspResp }\OperatorTok{=} \KeywordTok{new} \FunctionTok{OCSPResp}\OperatorTok{(}\NormalTok{is}\OperatorTok{);}
        \OperatorTok{\}}

        \ControlFlowTok{if} \OperatorTok{(}\NormalTok{ocspResp}\OperatorTok{.}\FunctionTok{getStatus}\OperatorTok{()} \OperatorTok{!=}\NormalTok{ OCSPResp}\OperatorTok{.}\FunctionTok{SUCCESSFUL}\OperatorTok{)} \OperatorTok{\{}
            \ControlFlowTok{throw} \KeywordTok{new} \BuiltInTok{RuntimeException}\OperatorTok{(}
                    \StringTok{"OCSP response status: "} \OperatorTok{+}\NormalTok{ ocspResp}\OperatorTok{.}\FunctionTok{getStatus}\OperatorTok{());}
        \OperatorTok{\}}

\NormalTok{        BasicOCSPResp basicResp }\OperatorTok{=}
                \OperatorTok{(}\NormalTok{BasicOCSPResp}\OperatorTok{)}\NormalTok{ ocspResp}\OperatorTok{.}\FunctionTok{getResponseObject}\OperatorTok{();}

        \CommentTok{// Find the response for our certificate}
        \ControlFlowTok{for} \OperatorTok{(}\NormalTok{SingleResp singleResp }\OperatorTok{:}\NormalTok{ basicResp}\OperatorTok{.}\FunctionTok{getResponses}\OperatorTok{())} \OperatorTok{\{}
            \ControlFlowTok{if} \OperatorTok{(}\NormalTok{singleResp}\OperatorTok{.}\FunctionTok{getCertID}\OperatorTok{().}\FunctionTok{equals}\OperatorTok{(}\NormalTok{certId}\OperatorTok{))} \OperatorTok{\{}
                \ControlFlowTok{return}\NormalTok{ singleResp}\OperatorTok{;}
            \OperatorTok{\}}
        \OperatorTok{\}}

        \ControlFlowTok{return} \KeywordTok{null}\OperatorTok{;} \CommentTok{// no matching response found}
    \OperatorTok{\}}
\OperatorTok{\}}
\end{Highlighting}
\end{Shaded}

\section{Trust Stores}\label{_trust_stores}

A trust store is a collection of root CA certificates that a system has
decided to trust. Without a trust store, there is no anchor for chain
validation --- every chain ends in a self-signed certificate, and if
that certificate is not in a trust store, the chain is untrusted
regardless of how cryptographically correct it is.

Different ecosystems maintain different trust stores, and understanding
which one applies to your PDF signatures is critical.

\subsection{Adobe Approved Trust List
(AATL)}\label{_adobe_approved_trust_list_aatl}

The \textbf{Adobe Approved Trust List} is a curated set of root
certificates that Adobe Acrobat and Adobe Reader trust for signature
validation. If your signing certificate chains to an AATL root, Acrobat
shows a green checkmark with "Signed and all signatures are valid." If
it does not --- even if the certificate is perfectly valid, issued by a
reputable CA, and trusted by your operating system --- Acrobat shows a
yellow warning triangle.

This is the single most important trust store for PDF signatures in
practice, because Adobe Acrobat has approximately 75\% of the worldwide
PDF viewer market share.

\subsubsection{How AATL Works}\label{_how_aatl_works}

Adobe publishes the AATL as a signed file that Acrobat downloads
automatically (the URL is hardcoded). The list is refreshed every 30
days. It contains approximately 250 root certificates from CAs around
the world.

To get on the AATL, a CA must:

\begin{enumerate}
\def\labelenumi{\arabic{enumi}.}
\item
  Pass a WebTrust or ETSI audit for PDF signing
\item
  Meet Adobe's technical requirements (minimum RSA-2048, SHA-256 or
  stronger)
\item
  Sign an agreement with Adobe
\item
  Pay an annual fee (reportedly in the range of \$10,000-\$50,000 USD,
  though Adobe does not publish the exact figure)
\end{enumerate}

\subsubsection{The Monopoly Problem}\label{_the_monopoly_problem}

The AATL creates a de facto monopoly on "valid" PDF signatures in the
world's dominant viewer. If you want your signed PDFs to show a green
checkmark in Acrobat, you need a certificate from an AATL-member CA.
Those certificates cost significantly more than standard TLS or
code-signing certificates --- typically \$200-\$500 per year for a basic
document signing certificate, and considerably more for organizational
certificates.

This pricing creates a barrier to entry. An individual developer who
wants to sign PDFs "properly" (by Acrobat's standards) must pay for an
AATL-rooted certificate, even though the cryptographic security of a
self-signed certificate is identical. The trust model imposes a
financial cost on trust itself.

This is not unique to Adobe --- the TLS ecosystem has a similar dynamic
with browser trust stores. But in TLS, Let's Encrypt disrupted the
pricing model by offering free certificates. There is no Let's Encrypt
for PDF signing certificates. The AATL's requirements (particularly the
audit requirements) make a free document signing CA impractical.

\subsection{EU Trusted List (EUTL)}\label{_eu_trusted_list_eutl}

The \textbf{European Union Trusted List} is a fundamentally different
approach to trust. Instead of a single vendor curating a list, the EUTL
is a hierarchy of machine-readable, cryptographically signed XML
documents maintained by the 27 EU member states under the authority of
the eIDAS regulation (Regulation (EU) No 910/2014 (eIDAS)).

\subsubsection{Structure: The LOT
Hierarchy}\label{_structure_the_lot_hierarchy}

The EUTL follows a \textbf{List of Trusted Lists (LOT)} architecture:

\begin{enumerate}
\def\labelenumi{\arabic{enumi}.}
\item
  The \textbf{European Commission} publishes the LOT --- a signed XML
  document listing the URLs of every member state's national trusted
  list.
\item
  Each \textbf{member state} publishes its own trusted list --- a signed
  XML document listing every Qualified Trust Service Provider (QTSP) in
  that country, along with the services they offer and the certificates
  associated with those services.
\item
  Each QTSP entry includes the \textbf{service type} (e.g., issuing
  qualified certificates for electronic signatures), the \textbf{service
  status} (granted, withdrawn, etc.), and the relevant certificates.
\end{enumerate}

The LOT is defined by ETSI TS 119 612 V2.1.1 (2015-07) and is publicly
accessible at a well-known URL maintained by the European Commission.

\subsubsection{How EUTL Differs from
AATL}\label{_how_eutl_differs_from_aatl}

\begin{longtable}[]{@{}
  >{\raggedright\arraybackslash}p{(\linewidth - 4\tabcolsep) * \real{0.2500}}
  >{\raggedright\arraybackslash}p{(\linewidth - 4\tabcolsep) * \real{0.3750}}
  >{\raggedright\arraybackslash}p{(\linewidth - 4\tabcolsep) * \real{0.3750}}@{}}
\caption{AATL vs EUTL comparison}\tabularnewline
\toprule\noalign{}
\begin{minipage}[b]{\linewidth}\raggedright
Aspect
\end{minipage} & \begin{minipage}[b]{\linewidth}\raggedright
AATL
\end{minipage} & \begin{minipage}[b]{\linewidth}\raggedright
EUTL
\end{minipage} \\
\midrule\noalign{}
\endfirsthead
\toprule\noalign{}
\begin{minipage}[b]{\linewidth}\raggedright
Aspect
\end{minipage} & \begin{minipage}[b]{\linewidth}\raggedright
AATL
\end{minipage} & \begin{minipage}[b]{\linewidth}\raggedright
EUTL
\end{minipage} \\
\midrule\noalign{}
\endhead
\bottomrule\noalign{}
\endlastfoot
Governance & Adobe (private company) & European Commission + member
states (public authority) \\
Legal basis & Adobe's terms of service & eIDAS regulation (Regulation
(EU) No 910/2014 (eIDAS)), legally binding across 27 member states \\
Format & Proprietary signed file & Machine-readable XML, signed per ETSI
TS 119 612 V2.1.1 (2015-07) \\
Scope & Any CA willing to pay and pass audit & Only Qualified Trust
Service Providers recognized under eIDAS \\
Cost to CAs & Annual fee to Adobe & No fee to EU; the QTSP status is
granted by the national supervisory body \\
Viewer support & Adobe Acrobat (dominant), some others & Adobe Acrobat
(imports EUTL), EU DSS validation framework, national validators \\
Legal effect & None (Acrobat showing "valid" has no legal meaning) &
Qualified electronic signatures have the legal equivalent of handwritten
signatures across the EU \\
\end{longtable}

Adobe Acrobat can import the EUTL in addition to the AATL. When
configured to use both, Acrobat validates against the union of AATL and
EUTL trust anchors. In European deployments, this is the standard
configuration.

A signature that validates against the EUTL as a qualified electronic
signature has legal standing that an AATL-only signature does not. The
EUTL is not just a technical trust store --- it is a legal instrument.
If you are operating in the EU, the EUTL is the authoritative trust
store for determining whether a signature meets the highest legal tier.
Chapter 19 covers the eIDAS legal framework in detail.

\subsection{Operating System Trust
Stores}\label{_operating_system_trust_stores}

Every major operating system maintains its own root CA trust store:

\begin{itemize}
\item
  \textbf{Windows}: The Microsoft Root Certificate Program. Certificates
  are distributed via Windows Update. The store is managed through the
  \texttt{certmgr.msc} snap-in or via \texttt{certutil}. Programs using
  the Windows CryptoAPI (Schannel, .NET) use this store automatically.
\item
  \textbf{macOS}: The Apple Root Certificate Program. Certificates are
  in the System Keychain. Managed through Keychain Access or
  \texttt{security} command-line tool. Programs using Security.framework
  (Safari, macOS apps) use this store.
\item
  \textbf{Linux}: Varies by distribution. Debian/Ubuntu use
  \texttt{/etc/ssl/certs/} and the \texttt{ca-certificates} package. Red
  Hat/Fedora use \texttt{/etc/pki/tls/certs/} and the \texttt{ca-trust}
  system. Updated via package manager.
\end{itemize}

For PDF signatures, OS trust stores are relevant primarily when using
system-level PDF viewers (Preview on macOS, Evince on Linux) or when a
Java application is configured to use the OS trust store instead of
Java's cacerts. Adobe Acrobat on Windows and macOS does \textbf{not} use
the OS trust store for signature validation by default --- it uses the
AATL (and optionally the EUTL).

\subsection{Java cacerts}\label{_java_cacerts}

Java ships with its own trust store: a PKCS\#12 file (or JKS in older
JDKs) located at \texttt{\$JAVA\_HOME/lib/security/cacerts}. The default
password is the string \texttt{changeit} (and almost nobody changes it).

The Java cacerts store is maintained by the JDK vendor (Oracle, Eclipse
Adoptium, Amazon Corretto, etc.) and typically mirrors the major
browser/OS trust stores with some additions and removals.

\begin{Shaded}
\begin{Highlighting}[]
\KeywordTok{import} \ImportTok{java}\OperatorTok{.}\ImportTok{io}\OperatorTok{.}\ImportTok{FileInputStream}\OperatorTok{;}
\KeywordTok{import} \ImportTok{java}\OperatorTok{.}\ImportTok{security}\OperatorTok{.}\ImportTok{KeyStore}\OperatorTok{;}
\KeywordTok{import} \ImportTok{java}\OperatorTok{.}\ImportTok{security}\OperatorTok{.}\ImportTok{cert}\OperatorTok{.}\ImportTok{X509Certificate}\OperatorTok{;}
\KeywordTok{import} \ImportTok{java}\OperatorTok{.}\ImportTok{util}\OperatorTok{.}\ImportTok{Enumeration}\OperatorTok{;}

\CommentTok{/**}
 \CommentTok{*}\NormalTok{ List all trusted root certificates in the Java cacerts store}\CommentTok{.}
 \CommentTok{*/}
\KeywordTok{public} \KeywordTok{class}\NormalTok{ CacertsExample }\OperatorTok{\{}

    \KeywordTok{public} \DataTypeTok{static} \DataTypeTok{void} \FunctionTok{listTrustedRoots}\OperatorTok{()} \KeywordTok{throws} \BuiltInTok{Exception} \OperatorTok{\{}
        \BuiltInTok{String}\NormalTok{ cacertsPath }\OperatorTok{=} \BuiltInTok{System}\OperatorTok{.}\FunctionTok{getProperty}\OperatorTok{(}\StringTok{"java.home"}\OperatorTok{)}
                \OperatorTok{+} \StringTok{"/lib/security/cacerts"}\OperatorTok{;}

        \BuiltInTok{KeyStore}\NormalTok{ ks }\OperatorTok{=} \BuiltInTok{KeyStore}\OperatorTok{.}\FunctionTok{getInstance}\OperatorTok{(}\StringTok{"PKCS12"}\OperatorTok{);}
        \ControlFlowTok{try} \OperatorTok{(}\BuiltInTok{FileInputStream}\NormalTok{ fis }\OperatorTok{=} \KeywordTok{new} \BuiltInTok{FileInputStream}\OperatorTok{(}\NormalTok{cacertsPath}\OperatorTok{))} \OperatorTok{\{}
\NormalTok{            ks}\OperatorTok{.}\FunctionTok{load}\OperatorTok{(}\NormalTok{fis}\OperatorTok{,} \StringTok{"changeit"}\OperatorTok{.}\FunctionTok{toCharArray}\OperatorTok{());}
        \OperatorTok{\}}

        \BuiltInTok{Enumeration}\OperatorTok{\textless{}}\BuiltInTok{String}\OperatorTok{\textgreater{}}\NormalTok{ aliases }\OperatorTok{=}\NormalTok{ ks}\OperatorTok{.}\FunctionTok{aliases}\OperatorTok{();}
        \DataTypeTok{int}\NormalTok{ count }\OperatorTok{=} \DecValTok{0}\OperatorTok{;}
        \ControlFlowTok{while} \OperatorTok{(}\NormalTok{aliases}\OperatorTok{.}\FunctionTok{hasMoreElements}\OperatorTok{())} \OperatorTok{\{}
            \BuiltInTok{String}\NormalTok{ alias }\OperatorTok{=}\NormalTok{ aliases}\OperatorTok{.}\FunctionTok{nextElement}\OperatorTok{();}
            \ControlFlowTok{if} \OperatorTok{(}\NormalTok{ks}\OperatorTok{.}\FunctionTok{isCertificateEntry}\OperatorTok{(}\NormalTok{alias}\OperatorTok{))} \OperatorTok{\{}
                \BuiltInTok{X509Certificate}\NormalTok{ cert }\OperatorTok{=} \OperatorTok{(}\BuiltInTok{X509Certificate}\OperatorTok{)}
\NormalTok{                        ks}\OperatorTok{.}\FunctionTok{getCertificate}\OperatorTok{(}\NormalTok{alias}\OperatorTok{);}
\NormalTok{                count}\OperatorTok{++;}
                \CommentTok{// cert.getSubjectX500Principal().getName()}
                \CommentTok{// gives the root CA\textquotesingle{}s distinguished name}
            \OperatorTok{\}}
        \OperatorTok{\}}
        \CommentTok{// Typical count: 100{-}150 root certificates}
    \OperatorTok{\}}
\OperatorTok{\}}
\end{Highlighting}
\end{Shaded}

To add a custom root CA to the Java trust store (for testing or internal
PKIs):

\begin{Shaded}
\begin{Highlighting}[]
\ExtensionTok{keytool} \AttributeTok{{-}importcert} \DataTypeTok{\textbackslash{}}
    \AttributeTok{{-}alias} \StringTok{"my{-}internal{-}ca"} \DataTypeTok{\textbackslash{}}
    \AttributeTok{{-}file}\NormalTok{ my{-}root{-}ca.pem }\DataTypeTok{\textbackslash{}}
    \AttributeTok{{-}keystore} \VariableTok{$JAVA\_HOME}\NormalTok{/lib/security/cacerts }\DataTypeTok{\textbackslash{}}
    \AttributeTok{{-}storepass}\NormalTok{ changeit}
\end{Highlighting}
\end{Shaded}

Modifying the system cacerts file affects all Java applications on that
machine. For production deployments, prefer creating a separate trust
store and configuring your application to use it via
\texttt{javax.net.ssl.trustStore} system property or programmatic
\texttt{TrustManagerFactory} configuration. Do not modify the global
cacerts unless you are administering the entire machine.

\subsection{Custom Trust Stores}\label{_custom_trust_stores}

Sometimes none of the standard trust stores are appropriate:

\begin{itemize}
\item
  \textbf{Internal PKIs}: Your organization runs its own CA hierarchy
  for internal document signing. The root CA is not in AATL, EUTL, or
  any OS trust store, because it was never intended to be publicly
  trusted.
\item
  \textbf{Restricted trust}: You want to validate PDF signatures but
  only trust a specific subset of CAs --- perhaps only EUTL-listed
  QTSPs, not the entire AATL.
\item
  \textbf{Air-gapped environments}: Systems with no network access
  cannot download CRL/OCSP responses or refresh trust stores. All trust
  material must be pre-loaded.
\end{itemize}

In Java, creating a custom trust store is straightforward:

\begin{Shaded}
\begin{Highlighting}[]
\KeywordTok{import} \ImportTok{java}\OperatorTok{.}\ImportTok{io}\OperatorTok{.}\ImportTok{FileOutputStream}\OperatorTok{;}
\KeywordTok{import} \ImportTok{java}\OperatorTok{.}\ImportTok{security}\OperatorTok{.}\ImportTok{KeyStore}\OperatorTok{;}
\KeywordTok{import} \ImportTok{java}\OperatorTok{.}\ImportTok{security}\OperatorTok{.}\ImportTok{cert}\OperatorTok{.}\ImportTok{CertificateFactory}\OperatorTok{;}
\KeywordTok{import} \ImportTok{java}\OperatorTok{.}\ImportTok{security}\OperatorTok{.}\ImportTok{cert}\OperatorTok{.}\ImportTok{X509Certificate}\OperatorTok{;}
\KeywordTok{import} \ImportTok{java}\OperatorTok{.}\ImportTok{io}\OperatorTok{.}\ImportTok{FileInputStream}\OperatorTok{;}

\CommentTok{/**}
 \CommentTok{*}\NormalTok{ Create a custom trust store containing only your organization}\CommentTok{\textquotesingle{}}\NormalTok{s}
 \CommentTok{*}\NormalTok{ root CA certificate}\CommentTok{. }\NormalTok{Use this for internal PKI validation}\CommentTok{.}
 \CommentTok{*/}
\KeywordTok{public} \KeywordTok{class}\NormalTok{ CustomTrustStoreExample }\OperatorTok{\{}

    \KeywordTok{public} \DataTypeTok{static} \DataTypeTok{void} \FunctionTok{createCustomTrustStore}\OperatorTok{(}
            \BuiltInTok{String}\NormalTok{ rootCertPath}\OperatorTok{,}
            \BuiltInTok{String}\NormalTok{ trustStorePath}\OperatorTok{,}
            \BuiltInTok{String}\NormalTok{ password}\OperatorTok{)} \KeywordTok{throws} \BuiltInTok{Exception} \OperatorTok{\{}

        \BuiltInTok{CertificateFactory}\NormalTok{ cf }\OperatorTok{=} \BuiltInTok{CertificateFactory}\OperatorTok{.}\FunctionTok{getInstance}\OperatorTok{(}\StringTok{"X.509"}\OperatorTok{);}
        \BuiltInTok{X509Certificate}\NormalTok{ rootCert}\OperatorTok{;}
        \ControlFlowTok{try} \OperatorTok{(}\BuiltInTok{FileInputStream}\NormalTok{ fis }\OperatorTok{=} \KeywordTok{new} \BuiltInTok{FileInputStream}\OperatorTok{(}\NormalTok{rootCertPath}\OperatorTok{))} \OperatorTok{\{}
\NormalTok{            rootCert }\OperatorTok{=} \OperatorTok{(}\BuiltInTok{X509Certificate}\OperatorTok{)}\NormalTok{ cf}\OperatorTok{.}\FunctionTok{generateCertificate}\OperatorTok{(}\NormalTok{fis}\OperatorTok{);}
        \OperatorTok{\}}

        \BuiltInTok{KeyStore}\NormalTok{ trustStore }\OperatorTok{=} \BuiltInTok{KeyStore}\OperatorTok{.}\FunctionTok{getInstance}\OperatorTok{(}\StringTok{"PKCS12"}\OperatorTok{);}
\NormalTok{        trustStore}\OperatorTok{.}\FunctionTok{load}\OperatorTok{(}\KeywordTok{null}\OperatorTok{,} \KeywordTok{null}\OperatorTok{);} \CommentTok{// initialize empty store}
\NormalTok{        trustStore}\OperatorTok{.}\FunctionTok{setCertificateEntry}\OperatorTok{(}\StringTok{"internal{-}root"}\OperatorTok{,}\NormalTok{ rootCert}\OperatorTok{);}

        \ControlFlowTok{try} \OperatorTok{(}\BuiltInTok{FileOutputStream}\NormalTok{ fos }\OperatorTok{=}
                \KeywordTok{new} \BuiltInTok{FileOutputStream}\OperatorTok{(}\NormalTok{trustStorePath}\OperatorTok{))} \OperatorTok{\{}
\NormalTok{            trustStore}\OperatorTok{.}\FunctionTok{store}\OperatorTok{(}\NormalTok{fos}\OperatorTok{,}\NormalTok{ password}\OperatorTok{.}\FunctionTok{toCharArray}\OperatorTok{());}
        \OperatorTok{\}}
    \OperatorTok{\}}
\OperatorTok{\}}
\end{Highlighting}
\end{Shaded}

\section{Trust Store Comparison}\label{_trust_store_comparison}

\begin{longtable}[]{@{}
  >{\raggedright\arraybackslash}p{(\linewidth - 8\tabcolsep) * \real{0.2000}}
  >{\raggedright\arraybackslash}p{(\linewidth - 8\tabcolsep) * \real{0.2000}}
  >{\raggedright\arraybackslash}p{(\linewidth - 8\tabcolsep) * \real{0.2000}}
  >{\raggedright\arraybackslash}p{(\linewidth - 8\tabcolsep) * \real{0.2000}}
  >{\raggedright\arraybackslash}p{(\linewidth - 8\tabcolsep) * \real{0.2000}}@{}}
\caption{Trust store comparison for PDF signature
validation}\tabularnewline
\toprule\noalign{}
\begin{minipage}[b]{\linewidth}\raggedright
Characteristic
\end{minipage} & \begin{minipage}[b]{\linewidth}\raggedright
AATL
\end{minipage} & \begin{minipage}[b]{\linewidth}\raggedright
EUTL
\end{minipage} & \begin{minipage}[b]{\linewidth}\raggedright
OS Store
\end{minipage} & \begin{minipage}[b]{\linewidth}\raggedright
Java cacerts
\end{minipage} \\
\midrule\noalign{}
\endfirsthead
\toprule\noalign{}
\begin{minipage}[b]{\linewidth}\raggedright
Characteristic
\end{minipage} & \begin{minipage}[b]{\linewidth}\raggedright
AATL
\end{minipage} & \begin{minipage}[b]{\linewidth}\raggedright
EUTL
\end{minipage} & \begin{minipage}[b]{\linewidth}\raggedright
OS Store
\end{minipage} & \begin{minipage}[b]{\linewidth}\raggedright
Java cacerts
\end{minipage} \\
\midrule\noalign{}
\endhead
\bottomrule\noalign{}
\endlastfoot
Primary consumer & Adobe Acrobat/Reader & EU DSS, Acrobat (configured) &
OS-native apps, browsers & Java applications \\
Approximate root count & \textasciitilde250 & Varies by member state;
\textasciitilde100 QTSPs total & \textasciitilde150-200 &
\textasciitilde100-150 \\
Update mechanism & Auto-download (30-day cycle) & LOT refresh per ETSI
TS 119 612 V2.1.1 (2015-07) & OS update (Windows Update, etc.) & JDK
update \\
Legal standing & None & Qualified signatures per eIDAS & None & None \\
PDF viewer support & Adobe Acrobat (primary) & Adobe Acrobat (opt-in),
EU validators & Evince, Preview, system viewers & PDFBox-based apps
(depends on config) \\
Cost for CA inclusion & Annual fee + audit & National supervisory body
approval & Vendor program requirements & JDK vendor decision \\
\end{longtable}

\section{Debugging Chain Failures}\label{_debugging_chain_failures}

When a PDF signature shows "certificate chain not trusted," here is the
systematic approach:

\subsection{Step 1: Extract the
Certificates}\label{_step_1_extract_the_certificates}

Open the PDF in PDFalyzer or use PDFBox to extract the CMS
\texttt{SignedData} from \texttt{/Contents}. List every certificate in
the CMS \texttt{certificates} field. Note the subject DN, issuer DN,
serial number, SKI, and AKI of each.

\subsection{Step 2: Build the Chain
Manually}\label{_step_2_build_the_chain_manually}

Starting from the signing certificate, match AKI to SKI. Can you build a
complete chain from the signer to a self-signed root? If a link is
missing, that is your problem --- the CMS is missing an intermediate.

\subsection{Step 3: Check the Trust
Anchor}\label{_step_3_check_the_trust_anchor}

Is the root certificate in the relevant trust store? If you are
validating in Acrobat, check the AATL. If you are validating in your
Java application, check cacerts or your custom trust store. The root
being "well-known" is not enough --- it must be in the specific trust
store the validator is using.

\subsection{Step 4: Check Validity
Periods}\label{_step_4_check_validity_periods}

Is every certificate in the chain currently within its validity period?
Or, if the signature has a timestamp, was every certificate valid at the
time the timestamp was created?

\subsection{Step 5: Check Revocation}\label{_step_5_check_revocation}

Is every certificate in the chain un-revoked? Can the validator reach
the CRL distribution points or OCSP responders listed in the
certificates? Firewall rules, proxy settings, and DNS resolution
failures are common culprits in enterprise environments.

\subsection{Step 6: Check Key Usage}\label{_step_6_check_key_usage}

Does the signing certificate have the correct \texttt{keyUsage} bits
set? Some validators (including Adobe Acrobat) require
\texttt{nonRepudiation} (also called \texttt{contentCommitment}) for
document signatures. Others accept \texttt{digitalSignature}. If neither
is set, the certificate is not suitable for signing regardless of chain
validity.

When you encounter a chain failure in production, resist the urge to
immediately add the root certificate to your trust store. First,
understand \emph{why} the chain is failing. Is it a missing
intermediate? A revoked certificate? A validator misconfiguration?
Adding a root to a trust store should be a deliberate decision, not a
debugging shortcut. Every root you add is a root you trust to issue
certificates for any purpose --- including, potentially, certificates
that could impersonate other entities.

\section{Cross-Certification and Bridge
CAs}\label{_cross_certification_and_bridge_cas}

We mentioned earlier that trust paths can be graphs, not just chains.
\textbf{Cross-certification} is the mechanism that creates these more
complex topologies.

When two independent PKI hierarchies need to trust each other's
certificates, they can cross-certify: Root A issues a certificate for
Root B's public key (or an intermediate in B's hierarchy), and vice
versa. This creates a trust bridge between the two hierarchies without
requiring end-entity certificates to be re-issued.

The US Federal PKI uses a \textbf{Bridge CA} architecture: the Federal
Bridge CA cross-certifies with individual agency CAs, enabling
certificates issued by the Department of Defense to be validated by
systems that trust the Department of State's root, and vice versa.

For PDF signatures, cross-certification is relevant when a signing
certificate chains to a root that is not directly in the validator's
trust store, but \emph{is} cross-certified with a root that is. The
chain-building algorithm must consider cross-certificates as additional
links, potentially finding a path through an unexpected route.

This adds complexity to chain building, and not all validators handle it
correctly. Adobe Acrobat supports cross-certification for AATL and EUTL
roots. Java's \texttt{CertPathValidator} supports it if you provide the
cross-certificates as additional material via
\texttt{PKIXBuilderParameters.addCertStore()}. PDFBox relies on
BouncyCastle for chain building and validation, and BouncyCastle handles
cross-certification correctly in most cases.

If you are debugging a chain failure and the certificates seem correct
but the path does not reach a trust anchor, check for
cross-certificates. The missing link might not be a traditional
intermediate CA but a cross-certificate that bridges two PKI
hierarchies. Cross-certificates are sometimes stored in LDAP directories
rather than embedded in the CMS structure, making them easy to overlook.

\section{Summary}\label{_summary_4}

Certificate chains exist because no single CA can practically serve
everyone. The hierarchical delegation model --- root CAs issuing to
intermediate CAs, intermediates issuing to end-entities --- distributes
the operational burden while maintaining a chain of trust back to a
known anchor.

Chain building matches AuthorityKeyIdentifier to SubjectKeyIdentifier,
walking from the end-entity up through intermediates to the root. Path
validation (RFC 5280 --- Internet X.509 PKI Certificate and CRL Profile
section 6) verifies that every link in the chain is cryptographically
sound, temporally valid, and compliant with the constraints imposed by
certificates higher in the chain.

Revocation checking --- via CRLs for bulk status and OCSP for real-time,
per-certificate status --- ensures that compromised or withdrawn
certificates are detected. For PDF signatures, embedding revocation data
in the document (via the DSS or CMS unsigned attributes) enables
long-term validation without network access.

The trust store anchoring the chain determines whether a signature is
considered "trusted." In the PDF world, the AATL dominates because Adobe
Acrobat dominates. In the EU, the EUTL provides legally binding trust
anchors for qualified electronic signatures. OS trust stores and Java's
cacerts serve broader ecosystems but are not the primary validation
paths for most PDF viewers.

When chain validation fails, the debugging process is methodical:
extract certificates, build the chain manually, check the trust anchor,
verify validity periods, check revocation status, and verify key usage.
Most chain failures in production are caused by missing intermediates,
not by actual trust problems.

The next chapter examines the formats in which keys and certificates are
stored and transported --- PKCS\#12, PEM, DER, JKS, and
PKCS\#8 --- because loading the wrong format is the quickest way to turn
a working certificate chain into an inscrutable exception stack trace.

\chapter{Key Material Formats}\label{ch05}

\section{PKCS\#12 / PFX: The Universal
Container}\label{_pkcs12_pfx_the_universal_container}

\textbf{PKCS\#12} is the format you will use most often for PDF signing.
It bundles a private key, the corresponding certificate, and (usually)
the full certificate chain into a single, password-protected file. The
standard file extensions are \texttt{.p12} and \texttt{.pfx}.

\subsection{History}\label{_history}

The lineage is tangled, as most things in cryptography are.

Microsoft created \textbf{PFX} (Personal Information Exchange) in the
mid-1990s as a container format for exporting keys and certificates from
Internet Explorer and IIS. RSA Laboratories, working on their PKCS
(Public-Key Cryptography Standards) series, developed \textbf{PKCS\#12}
as a platform-neutral alternative. The two formats converged over time:
PKCS\#12 v1.0 (1999) incorporated Microsoft's PFX design, and
\textbf{RFC 7292} (2014) formalized the standard.

Today, \texttt{.p12} and \texttt{.pfx} are effectively interchangeable
extensions for the same format. Some tools (notably Windows) prefer
\texttt{.pfx}; others (Java, OpenSSL) prefer \texttt{.p12}. The binary
content is identical.

You will occasionally encounter ancient \texttt{.pfx} files from the
late 1990s that use a pre-standard PFX format incompatible with
PKCS\#12. These are rare in 2026 but not extinct. OpenSSL will refuse to
parse them, producing an error about "unknown bag type." If you
encounter one, the only reliable conversion path is through a legacy
Windows system or Internet Explorer's certificate export.

\subsection{Internal Structure}\label{_internal_structure}

A PKCS\#12 file is a DER-encoded ASN.1 structure containing:

\begin{longtable}[]{@{}
  >{\raggedright\arraybackslash}p{(\linewidth - 2\tabcolsep) * \real{0.2857}}
  >{\raggedright\arraybackslash}p{(\linewidth - 2\tabcolsep) * \real{0.7143}}@{}}
\caption{PKCS\#12 internal structure}\tabularnewline
\toprule\noalign{}
\begin{minipage}[b]{\linewidth}\raggedright
Component
\end{minipage} & \begin{minipage}[b]{\linewidth}\raggedright
Description
\end{minipage} \\
\midrule\noalign{}
\endfirsthead
\toprule\noalign{}
\begin{minipage}[b]{\linewidth}\raggedright
Component
\end{minipage} & \begin{minipage}[b]{\linewidth}\raggedright
Description
\end{minipage} \\
\midrule\noalign{}
\endhead
\bottomrule\noalign{}
\endlastfoot
PFX PDU & The outermost wrapper. Contains a version number (always 3),
the authenticated safe, and an optional MAC (Message Authentication
Code) for integrity verification. \\
AuthenticatedSafe & A SEQUENCE of ContentInfo values, each containing
one or more SafeBag objects. The ContentInfo can be unencrypted
(\texttt{data}), encrypted with a password-based cipher
(\texttt{encryptedData}), or encrypted with a public key
(\texttt{envelopedData} --- rarely used). \\
SafeBag & The actual content container. Types include \texttt{keyBag}
(unencrypted private key), \texttt{pkcs8ShroudedKeyBag} (encrypted
private key --- the most common), \texttt{certBag} (certificate),
\texttt{crlBag} (CRL), and \texttt{secretBag} (arbitrary secret). \\
MAC & An HMAC computed over the AuthenticatedSafe using a key derived
from the password. This verifies that the password is correct and that
the file has not been tampered with. The MAC uses PBKDF2 or a similar
key derivation function. \\
\end{longtable}

The typical PKCS\#12 file for PDF signing contains:

\begin{enumerate}
\def\labelenumi{\arabic{enumi}.}
\item
  One \texttt{pkcs8ShroudedKeyBag} --- the private key, encrypted with a
  password-derived key
\item
  One or more \texttt{certBag} entries --- the signer's certificate and
  the intermediate CA certificate(s)
\end{enumerate}

The private key encryption within the PKCS\#12 uses a password-based
encryption scheme. Historically, this was \textbf{PBE-SHA1-3DES} (also
known as \texttt{pbeWithSHAAnd3-KeyTripleDES-CBC}), which is the default
in many older tools and in Java's PKCS\#12 implementation prior to JDK
17.

PBE-SHA1-3DES is weak by modern standards. SHA-1 for key derivation and
Triple DES for encryption are both deprecated. OpenSSL 3.x defaults to
AES-256-CBC with PBKDF2-HMAC-SHA256 for new PKCS\#12 files, which is
much stronger. However, some tools --- including older versions of Adobe
Acrobat --- cannot read PKCS\#12 files encrypted with the newer
algorithms. If you generate a PKCS\#12 with OpenSSL 3.x and Acrobat
cannot open it, try re-exporting with the legacy flag:
\texttt{openssl\ pkcs12\ -export\ -legacy\ -in\ cert.pem\ -inkey\ key.pem\ -out\ keystore.pfx}.

\subsection{Loading PKCS\#12 in Java}\label{_loading_pkcs12_in_java}

\begin{Shaded}
\begin{Highlighting}[]
\KeywordTok{import} \ImportTok{java}\OperatorTok{.}\ImportTok{io}\OperatorTok{.}\ImportTok{FileInputStream}\OperatorTok{;}
\KeywordTok{import} \ImportTok{java}\OperatorTok{.}\ImportTok{security}\OperatorTok{.}\ImportTok{KeyStore}\OperatorTok{;}
\KeywordTok{import} \ImportTok{java}\OperatorTok{.}\ImportTok{security}\OperatorTok{.}\ImportTok{PrivateKey}\OperatorTok{;}
\KeywordTok{import} \ImportTok{java}\OperatorTok{.}\ImportTok{security}\OperatorTok{.}\ImportTok{cert}\OperatorTok{.}\ImportTok{Certificate}\OperatorTok{;}
\KeywordTok{import} \ImportTok{java}\OperatorTok{.}\ImportTok{security}\OperatorTok{.}\ImportTok{cert}\OperatorTok{.}\ImportTok{X509Certificate}\OperatorTok{;}
\KeywordTok{import} \ImportTok{java}\OperatorTok{.}\ImportTok{util}\OperatorTok{.}\ImportTok{Enumeration}\OperatorTok{;}

\CommentTok{/**}
 \CommentTok{*}\NormalTok{ Load a PKCS}\CommentTok{\#12}\NormalTok{ keystore and extract the signing key and chain}\CommentTok{.}
 \CommentTok{*/}
\KeywordTok{public} \KeywordTok{class}\NormalTok{ Pkcs12LoadExample }\OperatorTok{\{}

    \KeywordTok{public} \DataTypeTok{static} \DataTypeTok{void} \FunctionTok{loadAndInspect}\OperatorTok{(}
            \BuiltInTok{String}\NormalTok{ p12Path}\OperatorTok{,} \BuiltInTok{String}\NormalTok{ password}\OperatorTok{)} \KeywordTok{throws} \BuiltInTok{Exception} \OperatorTok{\{}

        \BuiltInTok{KeyStore}\NormalTok{ ks }\OperatorTok{=} \BuiltInTok{KeyStore}\OperatorTok{.}\FunctionTok{getInstance}\OperatorTok{(}\StringTok{"PKCS12"}\OperatorTok{);}
        \ControlFlowTok{try} \OperatorTok{(}\BuiltInTok{FileInputStream}\NormalTok{ fis }\OperatorTok{=} \KeywordTok{new} \BuiltInTok{FileInputStream}\OperatorTok{(}\NormalTok{p12Path}\OperatorTok{))} \OperatorTok{\{}
\NormalTok{            ks}\OperatorTok{.}\FunctionTok{load}\OperatorTok{(}\NormalTok{fis}\OperatorTok{,}\NormalTok{ password}\OperatorTok{.}\FunctionTok{toCharArray}\OperatorTok{());}
        \OperatorTok{\}}

        \CommentTok{// Enumerate all entries}
        \BuiltInTok{Enumeration}\OperatorTok{\textless{}}\BuiltInTok{String}\OperatorTok{\textgreater{}}\NormalTok{ aliases }\OperatorTok{=}\NormalTok{ ks}\OperatorTok{.}\FunctionTok{aliases}\OperatorTok{();}
        \ControlFlowTok{while} \OperatorTok{(}\NormalTok{aliases}\OperatorTok{.}\FunctionTok{hasMoreElements}\OperatorTok{())} \OperatorTok{\{}
            \BuiltInTok{String}\NormalTok{ alias }\OperatorTok{=}\NormalTok{ aliases}\OperatorTok{.}\FunctionTok{nextElement}\OperatorTok{();}

            \ControlFlowTok{if} \OperatorTok{(}\NormalTok{ks}\OperatorTok{.}\FunctionTok{isKeyEntry}\OperatorTok{(}\NormalTok{alias}\OperatorTok{))} \OperatorTok{\{}
                \BuiltInTok{PrivateKey}\NormalTok{ key }\OperatorTok{=} \OperatorTok{(}\BuiltInTok{PrivateKey}\OperatorTok{)}\NormalTok{ ks}\OperatorTok{.}\FunctionTok{getKey}\OperatorTok{(}
\NormalTok{                        alias}\OperatorTok{,}\NormalTok{ password}\OperatorTok{.}\FunctionTok{toCharArray}\OperatorTok{());}    
                \BuiltInTok{Certificate}\OperatorTok{[]}\NormalTok{ chain }\OperatorTok{=}\NormalTok{ ks}\OperatorTok{.}\FunctionTok{getCertificateChain}\OperatorTok{(}\NormalTok{alias}\OperatorTok{);}

                \CommentTok{// key.getAlgorithm() {-}\textgreater{} "RSA" or "EC"}
                \CommentTok{// key.getFormat() {-}\textgreater{} "PKCS\#8"}
                \CommentTok{// chain[0] is the end{-}entity certificate}
                \CommentTok{// chain[chain.length{-}1] is closest to the root}

                \BuiltInTok{X509Certificate}\NormalTok{ signingCert }\OperatorTok{=}
                        \OperatorTok{(}\BuiltInTok{X509Certificate}\OperatorTok{)}\NormalTok{ chain}\OperatorTok{[}\DecValTok{0}\OperatorTok{];}
                \CommentTok{// signingCert.getSubjectX500Principal() {-}\textgreater{} signer DN}
                \CommentTok{// signingCert.getNotAfter() {-}\textgreater{} expiration date}
            \OperatorTok{\}}
        \OperatorTok{\}}
    \OperatorTok{\}}
\OperatorTok{\}}
\end{Highlighting}
\end{Shaded}

\begin{itemize}
\item
  The password for the key entry is typically the same as the store
  password for PKCS\#12 files, but the API allows them to be different.
  Some tools (notably \texttt{keytool}) set the key password equal to
  the store password when creating PKCS\#12 files.
\end{itemize}

When a PKCS\#12 file fails to load with
\texttt{IOException:\ keystore\ password\ was\ incorrect}, the error is
misleading. It can mean (a) the password is actually wrong, (b) the file
is corrupted, (c) the file uses an encryption algorithm not supported by
your JDK, or (d) the file is not actually PKCS\#12 format at all
(perhaps it is PEM-encoded and was given a \texttt{.p12} extension by
mistake). Always check the file format first:
\texttt{openssl\ pkcs12\ -info\ -in\ file.p12\ -nokeys\ -passin\ pass:test\ 2\textgreater{}\&1\ \textbar{}\ head\ -5}
will quickly tell you whether the file is valid PKCS\#12 and whether the
password works.

\section{PEM: The Text Format}\label{_pem_the_text_format}

\textbf{PEM} (Privacy-Enhanced Mail) is the most human-friendly key and
certificate format --- and, paradoxically, has nothing to do with email
privacy anymore.

\subsection{History}\label{_history_2}

PEM was defined in RFCs 1421-1424 (1993) as a framework for email
encryption. The framework itself was complex, poorly adopted, and
quickly superseded by S/MIME and PGP. But the \textbf{encoding
format} --- Base64-encoded DER data bracketed by header and footer
lines --- became the de facto standard for exchanging cryptographic
objects in text form.

The format is simple:

\begin{Shaded}
\begin{Highlighting}[]
\NormalTok{{-}{-}{-}{-}{-}BEGIN CERTIFICATE{-}{-}{-}{-}{-}}
\NormalTok{MIIDXTCCAkWgAwIBAgIJAJC1HiIAZAiUMA0GCSqGSIb3Qw0teleBAUAMEUxCzAJ}
\NormalTok{BgNVBAYTAkFVMRMwEQYDVQQIDApTb21lLVN0YXRlMSEwHwYDVQQKDBhJbnRlcm5l}
\NormalTok{... (Base64{-}encoded DER data, typically 64 or 76 characters per line) ...}
\NormalTok{dGhlIGNlcnRpZmljYXRl}
\NormalTok{{-}{-}{-}{-}{-}END CERTIFICATE{-}{-}{-}{-}{-}}
\end{Highlighting}
\end{Shaded}

The header line identifies the content type. Common headers include:

\begin{longtable}[]{@{}
  >{\raggedright\arraybackslash}p{(\linewidth - 2\tabcolsep) * \real{0.4286}}
  >{\raggedright\arraybackslash}p{(\linewidth - 2\tabcolsep) * \real{0.5714}}@{}}
\caption{PEM header types}\tabularnewline
\toprule\noalign{}
\begin{minipage}[b]{\linewidth}\raggedright
Header
\end{minipage} & \begin{minipage}[b]{\linewidth}\raggedright
Content
\end{minipage} \\
\midrule\noalign{}
\endfirsthead
\toprule\noalign{}
\begin{minipage}[b]{\linewidth}\raggedright
Header
\end{minipage} & \begin{minipage}[b]{\linewidth}\raggedright
Content
\end{minipage} \\
\midrule\noalign{}
\endhead
\bottomrule\noalign{}
\endlastfoot
\texttt{-\/-\/-\/-\/-BEGIN\ CERTIFICATE-\/-\/-\/-\/-} & X.509
certificate (DER-encoded) \\
\texttt{-\/-\/-\/-\/-BEGIN\ PRIVATE\ KEY-\/-\/-\/-\/-} & PKCS\#8
unencrypted private key \\
\texttt{-\/-\/-\/-\/-BEGIN\ ENCRYPTED\ PRIVATE\ KEY-\/-\/-\/-\/-} &
PKCS\#8 encrypted private key \\
\texttt{-\/-\/-\/-\/-BEGIN\ RSA\ PRIVATE\ KEY-\/-\/-\/-\/-} & PKCS\#1
RSA-specific private key (legacy OpenSSL format) \\
\texttt{-\/-\/-\/-\/-BEGIN\ PUBLIC\ KEY-\/-\/-\/-\/-} &
SubjectPublicKeyInfo (X.509) public key \\
\texttt{-\/-\/-\/-\/-BEGIN\ CERTIFICATE\ REQUEST-\/-\/-\/-\/-} &
PKCS\#10 certificate signing request (CSR) \\
\texttt{-\/-\/-\/-\/-BEGIN\ X509\ CRL-\/-\/-\/-\/-} & Certificate
Revocation List \\
\end{longtable}

\subsection{Multiple Objects in One
File}\label{_multiple_objects_in_one_file}

A PEM file can contain multiple objects --- this is one of its most
useful features. A common pattern for PDF signing is to concatenate the
certificate chain into a single file:

\begin{Shaded}
\begin{Highlighting}[]
\NormalTok{{-}{-}{-}{-}{-}BEGIN CERTIFICATE{-}{-}{-}{-}{-}}
\NormalTok{(end{-}entity certificate)}
\NormalTok{{-}{-}{-}{-}{-}END CERTIFICATE{-}{-}{-}{-}{-}}
\NormalTok{{-}{-}{-}{-}{-}BEGIN CERTIFICATE{-}{-}{-}{-}{-}}
\NormalTok{(intermediate CA certificate)}
\NormalTok{{-}{-}{-}{-}{-}END CERTIFICATE{-}{-}{-}{-}{-}}
\NormalTok{{-}{-}{-}{-}{-}BEGIN CERTIFICATE{-}{-}{-}{-}{-}}
\NormalTok{(root CA certificate {-}{-} optional)}
\NormalTok{{-}{-}{-}{-}{-}END CERTIFICATE{-}{-}{-}{-}{-}}
\end{Highlighting}
\end{Shaded}

Java's \texttt{CertificateFactory.generateCertificates(InputStream)}
(note the plural) reads all certificates from such a file.
BouncyCastle's \texttt{PEMParser} reads any PEM objects sequentially.

\subsection{Loading PEM with
BouncyCastle}\label{_loading_pem_with_bouncycastle}

\begin{Shaded}
\begin{Highlighting}[]
\KeywordTok{import} \ImportTok{java}\OperatorTok{.}\ImportTok{io}\OperatorTok{.}\ImportTok{FileReader}\OperatorTok{;}
\KeywordTok{import} \ImportTok{java}\OperatorTok{.}\ImportTok{security}\OperatorTok{.}\ImportTok{KeyPair}\OperatorTok{;}
\KeywordTok{import} \ImportTok{java}\OperatorTok{.}\ImportTok{security}\OperatorTok{.}\ImportTok{PrivateKey}\OperatorTok{;}
\KeywordTok{import} \ImportTok{java}\OperatorTok{.}\ImportTok{security}\OperatorTok{.}\ImportTok{cert}\OperatorTok{.}\ImportTok{X509Certificate}\OperatorTok{;}

\KeywordTok{import} \ImportTok{org}\OperatorTok{.}\ImportTok{bouncycastle}\OperatorTok{.}\ImportTok{cert}\OperatorTok{.}\ImportTok{X509CertificateHolder}\OperatorTok{;}
\KeywordTok{import} \ImportTok{org}\OperatorTok{.}\ImportTok{bouncycastle}\OperatorTok{.}\ImportTok{cert}\OperatorTok{.}\ImportTok{jcajce}\OperatorTok{.}\ImportTok{JcaX509CertificateConverter}\OperatorTok{;}
\KeywordTok{import} \ImportTok{org}\OperatorTok{.}\ImportTok{bouncycastle}\OperatorTok{.}\ImportTok{openssl}\OperatorTok{.}\ImportTok{PEMDecryptorProvider}\OperatorTok{;}
\KeywordTok{import} \ImportTok{org}\OperatorTok{.}\ImportTok{bouncycastle}\OperatorTok{.}\ImportTok{openssl}\OperatorTok{.}\ImportTok{PEMEncryptedKeyPair}\OperatorTok{;}
\KeywordTok{import} \ImportTok{org}\OperatorTok{.}\ImportTok{bouncycastle}\OperatorTok{.}\ImportTok{openssl}\OperatorTok{.}\ImportTok{PEMKeyPair}\OperatorTok{;}
\KeywordTok{import} \ImportTok{org}\OperatorTok{.}\ImportTok{bouncycastle}\OperatorTok{.}\ImportTok{openssl}\OperatorTok{.}\ImportTok{PEMParser}\OperatorTok{;}
\KeywordTok{import} \ImportTok{org}\OperatorTok{.}\ImportTok{bouncycastle}\OperatorTok{.}\ImportTok{openssl}\OperatorTok{.}\ImportTok{jcajce}\OperatorTok{.}\ImportTok{JcaPEMKeyConverter}\OperatorTok{;}
\KeywordTok{import} \ImportTok{org}\OperatorTok{.}\ImportTok{bouncycastle}\OperatorTok{.}\ImportTok{openssl}\OperatorTok{.}\ImportTok{jcajce}\OperatorTok{.}\ImportTok{JcePEMDecryptorProviderBuilder}\OperatorTok{;}

\CommentTok{/**}
 \CommentTok{*}\NormalTok{ Load certificates and private keys from PEM files using BouncyCastle}\CommentTok{.}
 \CommentTok{*/}
\KeywordTok{public} \KeywordTok{class}\NormalTok{ PemLoadExample }\OperatorTok{\{}

    \CommentTok{/**}
     \CommentTok{*}\NormalTok{ Load a private key from a PEM file}\CommentTok{.}
     \CommentTok{*}\NormalTok{ Handles both encrypted and unencrypted keys}\CommentTok{.}
     \CommentTok{*/}
    \KeywordTok{public} \DataTypeTok{static} \BuiltInTok{PrivateKey} \FunctionTok{loadPrivateKey}\OperatorTok{(}
            \BuiltInTok{String}\NormalTok{ pemPath}\OperatorTok{,} \BuiltInTok{String}\NormalTok{ password}\OperatorTok{)} \KeywordTok{throws} \BuiltInTok{Exception} \OperatorTok{\{}

        \ControlFlowTok{try} \OperatorTok{(}\NormalTok{PEMParser parser }\OperatorTok{=} \KeywordTok{new} \FunctionTok{PEMParser}\OperatorTok{(}\KeywordTok{new} \BuiltInTok{FileReader}\OperatorTok{(}\NormalTok{pemPath}\OperatorTok{)))} \OperatorTok{\{}
            \BuiltInTok{Object}\NormalTok{ obj }\OperatorTok{=}\NormalTok{ parser}\OperatorTok{.}\FunctionTok{readObject}\OperatorTok{();}

\NormalTok{            JcaPEMKeyConverter converter }\OperatorTok{=} \KeywordTok{new} \FunctionTok{JcaPEMKeyConverter}\OperatorTok{()}
                    \OperatorTok{.}\FunctionTok{setProvider}\OperatorTok{(}\StringTok{"BC"}\OperatorTok{);}

            \ControlFlowTok{if} \OperatorTok{(}\NormalTok{obj }\KeywordTok{instanceof}\NormalTok{ PEMEncryptedKeyPair}\OperatorTok{)} \OperatorTok{\{}
\NormalTok{                PEMDecryptorProvider decryptor }\OperatorTok{=}
                        \KeywordTok{new} \FunctionTok{JcePEMDecryptorProviderBuilder}\OperatorTok{()}
                                \OperatorTok{.}\FunctionTok{build}\OperatorTok{(}\NormalTok{password}\OperatorTok{.}\FunctionTok{toCharArray}\OperatorTok{());}
\NormalTok{                PEMKeyPair decrypted }\OperatorTok{=}
                        \OperatorTok{((}\NormalTok{PEMEncryptedKeyPair}\OperatorTok{)}\NormalTok{ obj}\OperatorTok{).}\FunctionTok{decryptKeyPair}\OperatorTok{(}\NormalTok{decryptor}\OperatorTok{);}
                \ControlFlowTok{return}\NormalTok{ converter}\OperatorTok{.}\FunctionTok{getKeyPair}\OperatorTok{(}\NormalTok{decrypted}\OperatorTok{).}\FunctionTok{getPrivate}\OperatorTok{();}
            \OperatorTok{\}}

            \ControlFlowTok{if} \OperatorTok{(}\NormalTok{obj }\KeywordTok{instanceof}\NormalTok{ PEMKeyPair}\OperatorTok{)} \OperatorTok{\{}
                \ControlFlowTok{return}\NormalTok{ converter}\OperatorTok{.}\FunctionTok{getKeyPair}\OperatorTok{((}\NormalTok{PEMKeyPair}\OperatorTok{)}\NormalTok{ obj}\OperatorTok{).}\FunctionTok{getPrivate}\OperatorTok{();}
            \OperatorTok{\}}

            \ControlFlowTok{if} \OperatorTok{(}\NormalTok{obj }\KeywordTok{instanceof}\NormalTok{ org}\OperatorTok{.}\FunctionTok{bouncycastle}\OperatorTok{.}\FunctionTok{asn1}\OperatorTok{.}\FunctionTok{pkcs}\OperatorTok{.}\FunctionTok{PrivateKeyInfo}\OperatorTok{)} \OperatorTok{\{}
                \ControlFlowTok{return}\NormalTok{ converter}\OperatorTok{.}\FunctionTok{getPrivateKey}\OperatorTok{(}
                        \OperatorTok{(}\NormalTok{org}\OperatorTok{.}\FunctionTok{bouncycastle}\OperatorTok{.}\FunctionTok{asn1}\OperatorTok{.}\FunctionTok{pkcs}\OperatorTok{.}\FunctionTok{PrivateKeyInfo}\OperatorTok{)}\NormalTok{ obj}\OperatorTok{);}
            \OperatorTok{\}}

            \ControlFlowTok{throw} \KeywordTok{new} \BuiltInTok{IllegalArgumentException}\OperatorTok{(}
                    \StringTok{"Unexpected PEM object: "} \OperatorTok{+}\NormalTok{ obj}\OperatorTok{.}\FunctionTok{getClass}\OperatorTok{().}\FunctionTok{getName}\OperatorTok{());}
        \OperatorTok{\}}
    \OperatorTok{\}}

    \CommentTok{/**}
     \CommentTok{*}\NormalTok{ Load an X}\CommentTok{.509}\NormalTok{ certificate from a PEM file}\CommentTok{.}
     \CommentTok{*/}
    \KeywordTok{public} \DataTypeTok{static} \BuiltInTok{X509Certificate} \FunctionTok{loadCertificate}\OperatorTok{(}
            \BuiltInTok{String}\NormalTok{ pemPath}\OperatorTok{)} \KeywordTok{throws} \BuiltInTok{Exception} \OperatorTok{\{}

        \ControlFlowTok{try} \OperatorTok{(}\NormalTok{PEMParser parser }\OperatorTok{=} \KeywordTok{new} \FunctionTok{PEMParser}\OperatorTok{(}\KeywordTok{new} \BuiltInTok{FileReader}\OperatorTok{(}\NormalTok{pemPath}\OperatorTok{)))} \OperatorTok{\{}
            \BuiltInTok{Object}\NormalTok{ obj }\OperatorTok{=}\NormalTok{ parser}\OperatorTok{.}\FunctionTok{readObject}\OperatorTok{();}

            \ControlFlowTok{if} \OperatorTok{(}\NormalTok{obj }\KeywordTok{instanceof}\NormalTok{ X509CertificateHolder}\OperatorTok{)} \OperatorTok{\{}
                \ControlFlowTok{return} \KeywordTok{new} \FunctionTok{JcaX509CertificateConverter}\OperatorTok{()}
                        \OperatorTok{.}\FunctionTok{setProvider}\OperatorTok{(}\StringTok{"BC"}\OperatorTok{)}
                        \OperatorTok{.}\FunctionTok{getCertificate}\OperatorTok{((}\NormalTok{X509CertificateHolder}\OperatorTok{)}\NormalTok{ obj}\OperatorTok{);}
            \OperatorTok{\}}

            \ControlFlowTok{throw} \KeywordTok{new} \BuiltInTok{IllegalArgumentException}\OperatorTok{(}
                    \StringTok{"Not a certificate: "} \OperatorTok{+}\NormalTok{ obj}\OperatorTok{.}\FunctionTok{getClass}\OperatorTok{().}\FunctionTok{getName}\OperatorTok{());}
        \OperatorTok{\}}
    \OperatorTok{\}}
\OperatorTok{\}}
\end{Highlighting}
\end{Shaded}

The \texttt{PEMParser} class reads objects sequentially. If your PEM
file contains a private key followed by certificates, you must call
\texttt{parser.readObject()} once for the key and then again for each
certificate. A common mistake is to read only the first object and close
the parser, missing the certificates entirely.

\section{DER: The Binary Encoding}\label{_der_the_binary_encoding}

\textbf{DER} (Distinguished Encoding Rules) is the canonical binary
encoding for ASN.1 data structures. Every cryptographic object you
encounter --- certificates, private keys, CMS structures, OCSP
responses --- has a DER encoding. PEM is just Base64-encoded DER with
headers.

\subsection{Relationship to PEM}\label{_relationship_to_pem}

The conversion between DER and PEM is trivial:

\begin{itemize}
\item
  \textbf{DER to PEM}: Base64-encode the DER bytes, split into
  64-character lines, wrap with the appropriate header/footer.
\item
  \textbf{PEM to DER}: Strip the header/footer lines, concatenate the
  remaining lines, Base64-decode.
\end{itemize}

DER files have no standard extension --- you will see \texttt{.der},
\texttt{.cer}, \texttt{.crt}, and even \texttt{.pem} used for
DER-encoded data (incorrectly, but it happens). The reliable way to
distinguish DER from PEM is to check whether the file starts with
\texttt{-\/-\/-\/-\/-BEGIN}: if yes, PEM; if no, DER (or something else
entirely).

\subsection{When DER Matters}\label{_when_der_matters}

For PDF signatures, DER encoding is everywhere:

\begin{itemize}
\item
  The CMS \texttt{SignedData} in the \texttt{/Contents} field is
  DER-encoded (then hex-encoded)
\item
  Each certificate in the CMS is DER-encoded
\item
  OCSP responses embedded in the DSS are DER-encoded
\item
  CRLs embedded in the DSS are DER-encoded
\end{itemize}

When you parse these structures with BouncyCastle, the library handles
the DER encoding/decoding transparently. You only deal with DER directly
when debugging at the byte level --- dumping the hex content of
\texttt{/Contents} and feeding it to \texttt{openssl\ asn1parse}, for
example.

\begin{Shaded}
\begin{Highlighting}[]
\CommentTok{\# Dump the ASN.1 structure of a DER{-}encoded CMS blob}
\ExtensionTok{openssl}\NormalTok{ asn1parse }\AttributeTok{{-}inform}\NormalTok{ DER }\AttributeTok{{-}in}\NormalTok{ signature.der}

\CommentTok{\# Convert DER certificate to PEM}
\ExtensionTok{openssl}\NormalTok{ x509 }\AttributeTok{{-}inform}\NormalTok{ DER }\AttributeTok{{-}in}\NormalTok{ cert.der }\AttributeTok{{-}outform}\NormalTok{ PEM }\AttributeTok{{-}out}\NormalTok{ cert.pem}

\CommentTok{\# Convert PEM certificate to DER}
\ExtensionTok{openssl}\NormalTok{ x509 }\AttributeTok{{-}inform}\NormalTok{ PEM }\AttributeTok{{-}in}\NormalTok{ cert.pem }\AttributeTok{{-}outform}\NormalTok{ DER }\AttributeTok{{-}out}\NormalTok{ cert.der}
\end{Highlighting}
\end{Shaded}

\section{JKS: Java KeyStore (Legacy)}\label{_jks_java_keystore_legacy}

\textbf{JKS} (Java KeyStore) was Java's original proprietary keystore
format, introduced in JDK 1.2 (1998). It served the Java ecosystem for
nearly two decades and is still encountered in legacy enterprise
systems.

\subsection{Why JKS Is Being Replaced}\label{_why_jks_is_being_replaced}

JKS has several problems:

\begin{enumerate}
\def\labelenumi{\arabic{enumi}.}
\item
  \textbf{Proprietary format}: Only Java tools can read it. OpenSSL
  cannot parse JKS. Non-Java applications cannot use JKS keystores.
\item
  \textbf{Weak encryption}: The private key protection in JKS uses a
  custom, non-standard algorithm based on password-derived key material
  and a proprietary obfuscation scheme. It is not PBE by any recognized
  standard.
\item
  \textbf{Case-insensitive aliases}: JKS converts all aliases to
  lowercase, which can cause subtle bugs when migrating to PKCS\#12
  (which is case-sensitive).
\item
  \textbf{No standard for interoperability}: Migrating keys between JKS
  and non-Java systems always requires format conversion.
\end{enumerate}

Since JDK 9, the default keystore type in Java is \textbf{PKCS12}, not
JKS. The \texttt{keytool} utility creates PKCS\#12 keystores by default.
If you are starting a new project, always use PKCS\#12.

To check which keystore type you are dealing with:
\texttt{keytool\ -list\ -keystore\ mystore.jks} will prompt for a
password and display entries. If you see "Keystore type: jks" in the
output, it is a JKS file. If you see "Keystore type: PKCS12", it is
PKCS\#12. To convert JKS to PKCS\#12:
\texttt{keytool\ -importkeystore\ -srckeystore\ old.jks\ -destkeystore\ new.p12\ -deststoretype\ PKCS12}.

\section{PKCS\#8: Private Key Format}\label{_pkcs8_private_key_format}

\textbf{PKCS\#8} (RFC 5958, originally RSA PKCS\#8) is the standard
format for private keys, independent of the key algorithm. When you
extract a private key from a PKCS\#12 file in Java
(\texttt{keyStore.getKey()}), the key's \texttt{getFormat()} method
returns \texttt{"PKCS\#8"} --- that is the encoding format of the key
bytes.

\subsection{Unencrypted vs Encrypted}\label{_unencrypted_vs_encrypted}

PKCS\#8 defines two structures:

\begin{itemize}
\item
  \textbf{PrivateKeyInfo} (unencrypted): Contains the algorithm
  identifier and the raw private key data. Stored in PEM files with the
  header \texttt{-\/-\/-\/-\/-BEGIN\ PRIVATE\ KEY-\/-\/-\/-\/-}.
\item
  \textbf{EncryptedPrivateKeyInfo}: The PrivateKeyInfo wrapped in a
  password-based encryption layer. Stored in PEM files with the header
  \texttt{-\/-\/-\/-\/-BEGIN\ ENCRYPTED\ PRIVATE\ KEY-\/-\/-\/-\/-}.
\end{itemize}

\begin{Shaded}
\begin{Highlighting}[]
\NormalTok{PrivateKeyInfo ::= }\DataTypeTok{SEQUENCE}\NormalTok{ \{}
\NormalTok{    version                Version,}
\NormalTok{    privateKeyAlgorithm    AlgorithmIdentifier,}
\NormalTok{    privateKey             }\DataTypeTok{OCTET} \DataTypeTok{STRING}
\NormalTok{\}}

\NormalTok{EncryptedPrivateKeyInfo ::= }\DataTypeTok{SEQUENCE}\NormalTok{ \{}
\NormalTok{    encryptionAlgorithm    AlgorithmIdentifier,}
\NormalTok{    encryptedData          }\DataTypeTok{OCTET} \DataTypeTok{STRING}
\NormalTok{\}}
\end{Highlighting}
\end{Shaded}

The \texttt{privateKey} octet string contains the algorithm-specific key
data --- an RSA private key (PKCS\#1 structure) for RSA, or an EC
private key structure for ECDSA. PKCS\#8 wraps these algorithm-specific
formats in a uniform container.

PKCS\#8 is the format that the JCA uses internally. When you call
\texttt{privateKey.getEncoded()}, the returned bytes are PKCS\#8
\texttt{PrivateKeyInfo}. When you call \texttt{publicKey.getEncoded()},
the returned bytes are \texttt{SubjectPublicKeyInfo} (the X.509 public
key format, defined in RFC 5280 --- Internet X.509 PKI Certificate and
CRL Profile). These format names --- PKCS\#8 for private keys, X.509 for
public keys --- appear in Java's \texttt{Key.getFormat()} and in PEM
headers.

\section{PKCS\#1: RSA-Specific Key
Format}\label{_pkcs1_rsa_specific_key_format}

\textbf{PKCS\#1} (RFC 8017) defines the RSA-specific private key
structure:

\begin{Shaded}
\begin{Highlighting}[]
\NormalTok{RSAPrivateKey ::= }\DataTypeTok{SEQUENCE}\NormalTok{ \{}
\NormalTok{    version           Version,}
\NormalTok{    modulus           }\DataTypeTok{INTEGER}\NormalTok{,  }\CommentTok{{-}{-} n}
\NormalTok{    publicExponent    }\DataTypeTok{INTEGER}\NormalTok{,  }\CommentTok{{-}{-} e}
\NormalTok{    privateExponent   }\DataTypeTok{INTEGER}\NormalTok{,  }\CommentTok{{-}{-} d}
\NormalTok{    prime1            }\DataTypeTok{INTEGER}\NormalTok{,  }\CommentTok{{-}{-} p}
\NormalTok{    prime2            }\DataTypeTok{INTEGER}\NormalTok{,  }\CommentTok{{-}{-} q}
\NormalTok{    exponent1         }\DataTypeTok{INTEGER}\NormalTok{,  }\CommentTok{{-}{-} d mod (p{-}1)}
\NormalTok{    exponent2         }\DataTypeTok{INTEGER}\NormalTok{,  }\CommentTok{{-}{-} d mod (q{-}1)}
\NormalTok{    coefficient       }\DataTypeTok{INTEGER}   \CommentTok{{-}{-} (inverse of q) mod p}
\NormalTok{\}}
\end{Highlighting}
\end{Shaded}

This is the format stored inside a PKCS\#8 \texttt{PrivateKeyInfo} when
the algorithm is RSA. It is also the format used in PEM files with the
header
\texttt{-\/-\/-\/-\/-BEGIN\ RSA\ PRIVATE\ KEY-\/-\/-\/-\/-} --- OpenSSL's
"traditional" format.

The distinction between PKCS\#1 and PKCS\#8 confuses many developers:

\begin{itemize}
\item
  \texttt{-\/-\/-\/-\/-BEGIN\ RSA\ PRIVATE\ KEY-\/-\/-\/-\/-} is PKCS\#1
  (RSA-specific, no algorithm identifier wrapper)
\item
  \texttt{-\/-\/-\/-\/-BEGIN\ PRIVATE\ KEY-\/-\/-\/-\/-} is PKCS\#8
  (algorithm-agnostic, contains an algorithm identifier)
\end{itemize}

Both contain the same RSA key data. PKCS\#8 simply wraps the PKCS\#1
structure with an additional layer that identifies the algorithm. If you
try to load a PKCS\#1 PEM file using code that expects PKCS\#8, you will
get a parsing error --- and vice versa.

OpenSSL's default output format changed between versions. OpenSSL 1.x
used PKCS\#1
(\texttt{-\/-\/-\/-\/-BEGIN\ RSA\ PRIVATE\ KEY-\/-\/-\/-\/-}) by default
for RSA keys. OpenSSL 3.x uses PKCS\#8
(\texttt{-\/-\/-\/-\/-BEGIN\ PRIVATE\ KEY-\/-\/-\/-\/-}) by default. If
you are following a tutorial written for an older OpenSSL version, the
key files it produces may be in a different format than your current
OpenSSL generates. Always check the PEM header.

\section{Format Conversion
Reference}\label{_format_conversion_reference}

This matrix shows how to convert between formats using \texttt{openssl}
and \texttt{keytool}. Every PDF signing project eventually needs at
least one of these conversions.

\begin{longtable}[]{@{}
  >{\raggedright\arraybackslash}p{(\linewidth - 4\tabcolsep) * \real{0.2222}}
  >{\raggedright\arraybackslash}p{(\linewidth - 4\tabcolsep) * \real{0.2222}}
  >{\raggedright\arraybackslash}p{(\linewidth - 4\tabcolsep) * \real{0.5556}}@{}}
\caption{Key and certificate format conversion commands}\tabularnewline
\toprule\noalign{}
\begin{minipage}[b]{\linewidth}\raggedright
From
\end{minipage} & \begin{minipage}[b]{\linewidth}\raggedright
To
\end{minipage} & \begin{minipage}[b]{\linewidth}\raggedright
Command
\end{minipage} \\
\midrule\noalign{}
\endfirsthead
\toprule\noalign{}
\begin{minipage}[b]{\linewidth}\raggedright
From
\end{minipage} & \begin{minipage}[b]{\linewidth}\raggedright
To
\end{minipage} & \begin{minipage}[b]{\linewidth}\raggedright
Command
\end{minipage} \\
\midrule\noalign{}
\endhead
\bottomrule\noalign{}
\endlastfoot
PKCS\#12 & PEM (key + cert) &
\texttt{openssl\ pkcs12\ -in\ store.p12\ -out\ bundle.pem\ -nodes} \\
PKCS\#12 & Separate PEM files &
\texttt{openssl\ pkcs12\ -in\ store.p12\ -nocerts\ -out\ key.pem} (key)
then \texttt{openssl\ pkcs12\ -in\ store.p12\ -nokeys\ -out\ cert.pem}
(cert) \\
PEM (key + cert) & PKCS\#12 &
\texttt{openssl\ pkcs12\ -export\ -in\ cert.pem\ -inkey\ key.pem\ -out\ store.p12\ -name\ "alias"} \\
PEM certificate & DER certificate &
\texttt{openssl\ x509\ -in\ cert.pem\ -outform\ DER\ -out\ cert.der} \\
DER certificate & PEM certificate &
\texttt{openssl\ x509\ -inform\ DER\ -in\ cert.der\ -outform\ PEM\ -out\ cert.pem} \\
PKCS\#1 key & PKCS\#8 key &
\texttt{openssl\ pkcs8\ -topk8\ -in\ rsa.pem\ -out\ pkcs8.pem\ -nocrypt} \\
PKCS\#8 key & PKCS\#1 key &
\texttt{openssl\ rsa\ -in\ pkcs8.pem\ -out\ rsa.pem\ -traditional} \\
JKS & PKCS\#12 &
\texttt{keytool\ -importkeystore\ -srckeystore\ old.jks\ -destkeystore\ new.p12\ -deststoretype\ PKCS12} \\
PKCS\#12 & JKS &
\texttt{keytool\ -importkeystore\ -srckeystore\ store.p12\ -srcstoretype\ PKCS12\ -destkeystore\ new.jks\ -deststoretype\ JKS} \\
PEM certificate & JKS (trust entry) &
\texttt{keytool\ -importcert\ -file\ cert.pem\ -keystore\ trust.jks\ -alias\ "mycert"} \\
\end{longtable}

The \texttt{-nodes} flag in the OpenSSL PKCS\#12 export command means
"no DES" --- the private key will not be encrypted in the output PEM
file. This is convenient for testing but dangerous in production. Never
leave unencrypted private keys on disk in production environments. Use
\texttt{-aes256} instead of \texttt{-nodes} to encrypt the output
private key with AES-256.

\section{HSM Integration: When Keys Cannot
Leave}\label{_hsm_integration_when_keys_cannot_leave}

A \textbf{Hardware Security Module (HSM)} is a dedicated hardware device
that generates, stores, and uses cryptographic keys internally. The
private key \textbf{never leaves the HSM} --- all signing operations
happen inside the device. This provides the strongest possible key
protection: even if the server is completely compromised, the attacker
cannot extract the private key.

\subsection{PKCS\#11}\label{_pkcs11}

The standard interface for HSM access is \textbf{PKCS\#11} (also called
Cryptoki), defined by OASIS (originally by RSA Laboratories). PKCS\#11
defines a C API for cryptographic operations: key generation, signing,
encryption, etc. The HSM vendor provides a shared library (\texttt{.so}
on Linux, \texttt{.dll} on Windows) that implements the PKCS\#11
interface.

Java accesses PKCS\#11 via the \textbf{SunPKCS11 provider} --- a JCA
provider that bridges the Java cryptographic API to a PKCS\#11 library.

\begin{Shaded}
\begin{Highlighting}[]
\KeywordTok{import} \ImportTok{java}\OperatorTok{.}\ImportTok{security}\OperatorTok{.}\ImportTok{KeyStore}\OperatorTok{;}
\KeywordTok{import} \ImportTok{java}\OperatorTok{.}\ImportTok{security}\OperatorTok{.}\ImportTok{Provider}\OperatorTok{;}
\KeywordTok{import} \ImportTok{java}\OperatorTok{.}\ImportTok{security}\OperatorTok{.}\ImportTok{Security}\OperatorTok{;}

\CommentTok{/**}
 \CommentTok{*}\NormalTok{ Load keys from an HSM via PKCS}\CommentTok{\#11.}
 \CommentTok{*}\NormalTok{ The private key never leaves the HSM }\CommentTok{{-}{-}}\NormalTok{ signing happens on}\CommentTok{{-}}\NormalTok{device}\CommentTok{.}
 \CommentTok{*/}
\KeywordTok{public} \KeywordTok{class}\NormalTok{ Pkcs11Example }\OperatorTok{\{}

    \KeywordTok{public} \DataTypeTok{static} \BuiltInTok{KeyStore} \FunctionTok{loadFromHsm}\OperatorTok{(}
            \BuiltInTok{String}\NormalTok{ pkcs11ConfigPath}\OperatorTok{,} \BuiltInTok{String}\NormalTok{ pin}\OperatorTok{)} \KeywordTok{throws} \BuiltInTok{Exception} \OperatorTok{\{}

        \CommentTok{// Configure the SunPKCS11 provider}
        \BuiltInTok{Provider}\NormalTok{ p }\OperatorTok{=} \BuiltInTok{Security}\OperatorTok{.}\FunctionTok{getProvider}\OperatorTok{(}\StringTok{"SunPKCS11"}\OperatorTok{);}
\NormalTok{        p }\OperatorTok{=}\NormalTok{ p}\OperatorTok{.}\FunctionTok{configure}\OperatorTok{(}\NormalTok{pkcs11ConfigPath}\OperatorTok{);}      
        \BuiltInTok{Security}\OperatorTok{.}\FunctionTok{addProvider}\OperatorTok{(}\NormalTok{p}\OperatorTok{);}

        \CommentTok{// Load the "keystore" {-}{-} this actually enumerates}
        \CommentTok{// objects on the HSM token}
        \BuiltInTok{KeyStore}\NormalTok{ ks }\OperatorTok{=} \BuiltInTok{KeyStore}\OperatorTok{.}\FunctionTok{getInstance}\OperatorTok{(}\StringTok{"PKCS11"}\OperatorTok{,}\NormalTok{ p}\OperatorTok{);}
\NormalTok{        ks}\OperatorTok{.}\FunctionTok{load}\OperatorTok{(}\KeywordTok{null}\OperatorTok{,}\NormalTok{ pin}\OperatorTok{.}\FunctionTok{toCharArray}\OperatorTok{());}        

        \ControlFlowTok{return}\NormalTok{ ks}\OperatorTok{;}
        \CommentTok{// ks.getKey() returns a key handle, not the actual key bytes}
        \CommentTok{// Signing with this key executes on the HSM}
    \OperatorTok{\}}
\OperatorTok{\}}
\end{Highlighting}
\end{Shaded}

\begin{itemize}
\item
  The config file specifies the PKCS\#11 library path, slot number, and
  other HSM-specific parameters.
\item
  The PIN authenticates to the HSM token. Unlike PKCS\#12 passwords, HSM
  PINs are typically short numeric codes enforced by the hardware.
\end{itemize}

\subsection{Why This Matters for PDF
Signing}\label{_why_this_matters_for_pdf_signing}

When you use PDFBox's \texttt{SignatureInterface.sign()} with an
HSM-backed key, the signing operation is transparently delegated to the
HSM. The \texttt{PrivateKey} object returned by \texttt{ks.getKey()} is
a \textbf{handle}, not the actual key material. The JCA provider
intercepts the signing call, forwards it to the HSM via PKCS\#11, and
returns the result. Your signing code does not change --- the same
\texttt{CMSSignedDataGenerator} code from Chapter 2 works with HSM keys
and software keys identically.

This is the elegance of the JCA provider model: the signing code is
agnostic to key storage. Whether the key lives in a PKCS\#12 file, a JKS
keystore, or an HSM, the API is the same.

For eIDAS qualified electronic signatures, the private key \emph{must}
be stored in a \textbf{Qualified Signature Creation Device
(QSCD)} --- an HSM or smart card that has been certified to Common
Criteria EAL 4+ or equivalent. Software keystores (PKCS\#12, JKS) are
not QSCDs. If you are building a signing service that needs to produce
qualified electronic signatures, HSM integration is not optional --- it
is a legal requirement under Regulation (EU) No 910/2014 (eIDAS) Article
29.

\section{Why Format Conversion Goes
Wrong}\label{_why_format_conversion_goes_wrong}

Format conversions are a perennial source of frustration. Here are the
most common failure modes:

\subsection{Password Encoding
Mismatch}\label{_password_encoding_mismatch}

PKCS\#12 passwords can be encoded as either \textbf{UTF-8} or
\textbf{UTF-16BE} (BMPString in ASN.1 terminology) for the key
derivation function. The PKCS\#12 standard specifies BMPString, but some
implementations use UTF-8.

If you create a PKCS\#12 with a password containing non-ASCII characters
(umlauts, accented characters, CJK characters) using OpenSSL, and then
try to open it in Java, the key derivation may produce a different key
because the two implementations encode the password differently.

The fix: use only ASCII characters in PKCS\#12 passwords. This is not a
theoretical recommendation --- it is a practical necessity in
cross-platform environments.

\subsection{OpenSSL Legacy Format vs
PKCS\#8}\label{_openssl_legacy_format_vs_pkcs8}

As mentioned above, OpenSSL's default private key format changed between
major versions. A script that worked with OpenSSL 1.1.1 may produce keys
that your PKCS\#8-expecting Java code cannot parse when run with OpenSSL
3.x, or vice versa.

Always check the PEM header and convert as needed.

\subsection{Certificate Chain Order in
PKCS\#12}\label{_certificate_chain_order_in_pkcs12}

The PKCS\#12 specification does not mandate a specific order for
certificates in the file. Some tools store the end-entity certificate
first, followed by intermediates. Others store them in reverse order.
Java's \texttt{KeyStore.getCertificateChain()} returns them in order
(end-entity first), regardless of how they are stored in the
file --- but only if the certificates are associated with the correct
key alias.

If the chain is missing intermediates (because they were not included
when the PKCS\#12 was created), \texttt{getCertificateChain()} returns
only the end-entity certificate. This produces the classic "chain not
trusted" error at signing or validation time.

When creating a PKCS\#12 for PDF signing, always include the full chain
(minus the root). With OpenSSL:
\texttt{openssl\ pkcs12\ -export\ -in\ signer.pem\ -inkey\ signer-key.pem\ -certfile\ intermediate-chain.pem\ -out\ signing.p12}.
The \texttt{-certfile} flag adds the intermediate certificates. Without
it, only the end-entity certificate is included.

\subsection{DER vs PEM Confusion}\label{_der_vs_pem_confusion}

Renaming a file's extension does not change its format. A file named
\texttt{cert.pem} that is actually DER-encoded will cause a Base64
decode error. A file named \texttt{cert.der} that is actually
PEM-encoded will cause an ASN.1 parse error.

To reliably detect the format:

\begin{Shaded}
\begin{Highlighting}[]
\CommentTok{\# If the file starts with "{-}{-}{-}{-}{-}}\RegionMarkerTok{BEGIN}\CommentTok{", it is PEM}
\FunctionTok{head} \AttributeTok{{-}1}\NormalTok{ cert.pem}

\CommentTok{\# Try parsing as DER; if it fails, it is probably PEM (or corrupted)}
\ExtensionTok{openssl}\NormalTok{ x509 }\AttributeTok{{-}inform}\NormalTok{ DER }\AttributeTok{{-}in}\NormalTok{ cert.der }\AttributeTok{{-}noout} \DecValTok{2}\OperatorTok{\textgreater{}\&}\DecValTok{1}

\CommentTok{\# The file command on Linux/macOS can often identify the format}
\FunctionTok{file}\NormalTok{ cert.pem}
\end{Highlighting}
\end{Shaded}

\section{Summary}\label{_summary_5}

Key material comes in several formats, each with its own history,
strengths, and failure modes:

\begin{itemize}
\item
  \textbf{PKCS\#12} is the universal container for private keys and
  certificate chains. Use it for PDF signing keystores.
\item
  \textbf{PEM} is the human-readable, Base64-encoded format for
  individual certificates, keys, and other cryptographic objects. Use it
  for interchange with OpenSSL and when you need to inspect contents
  visually.
\item
  \textbf{DER} is the binary ASN.1 encoding underlying everything. You
  encounter it inside PDF signatures (the CMS blob in
  \texttt{/Contents}), in OCSP responses, and in CRLs.
\item
  \textbf{JKS} is Java's legacy keystore format. Migrate to PKCS\#12.
\item
  \textbf{PKCS\#8} is the standard private key format
  (algorithm-agnostic). Java uses it internally.
\item
  \textbf{PKCS\#1} is the RSA-specific key format. Watch for the PEM
  header to distinguish it from PKCS\#8.
\end{itemize}

For HSM-backed signing (mandatory for eIDAS qualified signatures), keys
are accessed via PKCS\#11 and never leave the hardware. The JCA provider
model makes this transparent to your signing code.

Format conversion problems --- password encoding, legacy vs modern
defaults, missing chain certificates, DER/PEM confusion --- are among
the most common issues in PDF signing deployments. The conversion
reference table in this chapter and a systematic approach to checking
file formats will save you hours of debugging.

In the next chapter, we examine the CMS SignedData structure --- the
envelope that carries the cryptographic signature, the certificate
chain, and the signed attributes from the signer to the validator, all
encoded inside the PDF's \texttt{/Contents} hex string.

\chapter{CMS/PKCS\#7 Signed Data}\label{ch06}

Every PDF digital signature carries its cryptographic payload inside a
CMS \texttt{SignedData} structure --- a binary envelope defined by RFC
5652 --- Cryptographic Message Syntax (CMS) that bundles the signature
value, the signer's certificate chain, the digest algorithm, and a set
of signed metadata attributes into a single DER-encoded blob. This
chapter dissects that envelope layer by layer, from the outermost
\texttt{ContentInfo} wrapper down to the individual signed attributes,
with annotated ASN.1 dumps and BouncyCastle 1.80 code on Java 21. If you
understand CMS \texttt{SignedData}, you understand the heart of what
lives inside the \texttt{/Contents} hex string of every signed PDF.

If you have parsed CMS structures before, skip to
\hyperref[ch06-signerinfo-deep-dive]{SignerInfo: The Core of the
Signature} for the signed attributes encoding trap that has caused more
validation failures than any other single issue. If you want code, jump
to \hyperref[ch06-java-parsing]{Java: Parsing CMS SignedData}.

\section{From PKCS\#7 to CMS: A
History}\label{_from_pkcs7_to_cms_a_history}

The Cryptographic Message Syntax did not spring into existence fully
formed. It evolved over sixteen years through three distinct phases,
each driven by different organizations and different needs.
Understanding this lineage matters because the older formats are still
embedded in millions of signed PDFs, and the naming confusion between
"PKCS\#7" and "CMS" persists in specifications, code comments, and error
messages to this day.

\subsection{PKCS\#7: RSA Laboratories,
1993}\label{_pkcs7_rsa_laboratories_1993}

RSA Laboratories published \textbf{PKCS\#7} (Public-Key Cryptography
Standard \#7) in November 1993 as a general-purpose format for signed
and encrypted data. The full title was "Cryptographic Message Syntax
Standard," which is ironic because it was not yet the CMS we know
today --- it was a proprietary RSA specification, not an IETF standard.

PKCS\#7 version 1.5 defined the core structure that survives almost
unchanged:

\begin{itemize}
\item
  A \texttt{ContentInfo} wrapper with an OID identifying the content
  type
\item
  A \texttt{SignedData} type containing digest algorithms, signer
  information, certificates, and CRLs
\item
  A \texttt{SignerInfo} type containing the actual signature and
  per-signer metadata
\item
  Support for both \emph{attached} (encapsulated) and \emph{detached}
  signatures
\end{itemize}

The design was influenced by the PEM (Privacy-Enhanced Mail) framework
from RFCs 1421-1424 and by the emerging S/MIME work at RSA. The goal was
a format that could carry signed data for email, file signing, code
signing, and any other context where a digital signature needed to
travel alongside or separately from its content.

The "PKCS\#7" name persists everywhere. The PDF specification uses
\texttt{adbe.pkcs7.detached} and \texttt{adbe.pkcs7.sha1} as SubFilter
values. Java's \texttt{KeyStore.getInstance("PKCS7")} exists. OpenSSL
has \texttt{pkcs7} commands. In all of these contexts, the actual format
in use is CMS as defined by RFC 5652 --- Cryptographic Message Syntax
(CMS), not the original 1993 PKCS\#7. The names are legacy; the
structures have evolved.

\subsection{S/MIME and RFC 2315, 1998}\label{_smime_and_rfc_2315_1998}

The S/MIME (Secure/Multipurpose Internet Mail Extensions) working group
at the IETF took PKCS\#7 and brought it into the standards process.
\textbf{RFC 2315} (March 1998) republished PKCS\#7 v1.5 essentially
verbatim as an informational RFC, giving it a stable, citable reference
in the IETF document series.

But the S/MIME group needed enhancements that RSA's original design did
not provide: better algorithm agility, support for attribute
certificates, and a cleaner separation between the content being signed
and the signature metadata. This work produced the first true CMS
specification.

\subsection{CMS: RFC 2630 (1999) and RFC 3852
(2004)}\label{_cms_rfc_2630_1999_and_rfc_3852_2004}

\textbf{RFC 2630} (June 1999) was the first IETF-standards-track CMS
specification. It renamed the format from "PKCS\#7" to "Cryptographic
Message Syntax" and introduced several changes:

\begin{itemize}
\item
  The \texttt{SignedData} version number advanced from 1 to 3 (with
  version selection rules based on which optional fields are present)
\item
  \texttt{EncapsulatedContentInfo} replaced the old \texttt{ContentInfo}
  for the inner content, allowing non-data content types
\item
  The \texttt{SignerIdentifier} was formalized with two choices:
  \texttt{issuerAndSerialNumber} (the classic) and
  \texttt{subjectKeyIdentifier} (new)
\item
  Attribute certificates were added to the \texttt{certificates} field
  as an option alongside X.509 certificates
\item
  The \texttt{digestAlgorithms} SET was required to be consistent with
  the algorithms actually used in \texttt{SignerInfo} entries
\end{itemize}

\textbf{RFC 3852} (July 2004) was a revision that clarified ambiguities
and tightened requirements, particularly around the encoding of signed
attributes.

\subsection{RFC 5652: The Current Standard
(2009)}\label{_rfc_5652_the_current_standard_2009}

\textbf{RFC 5652 --- Cryptographic Message Syntax (CMS)} (September
2009) is the current and authoritative CMS specification. It obsoletes
RFC 3852 and is the normative reference used by ISO 32000-2:2020, the
PAdES standards (ETSI EN 319 142-1 V1.1.1 (2016-04)), and BouncyCastle.

The changes from RFC 3852 were modest --- mostly editorial
clarifications and alignment with updated ASN.1 modules. The core
structures are stable and have not changed since 1999. When we refer to
"CMS" in this book, we mean RFC 5652 --- Cryptographic Message Syntax
(CMS) unless otherwise noted.

\begin{longtable}[]{@{}
  >{\raggedright\arraybackslash}p{(\linewidth - 4\tabcolsep) * \real{0.1429}}
  >{\raggedright\arraybackslash}p{(\linewidth - 4\tabcolsep) * \real{0.2857}}
  >{\raggedright\arraybackslash}p{(\linewidth - 4\tabcolsep) * \real{0.5714}}@{}}
\caption{CMS evolution timeline}\tabularnewline
\toprule\noalign{}
\begin{minipage}[b]{\linewidth}\raggedright
Year
\end{minipage} & \begin{minipage}[b]{\linewidth}\raggedright
Document
\end{minipage} & \begin{minipage}[b]{\linewidth}\raggedright
Significance
\end{minipage} \\
\midrule\noalign{}
\endfirsthead
\toprule\noalign{}
\begin{minipage}[b]{\linewidth}\raggedright
Year
\end{minipage} & \begin{minipage}[b]{\linewidth}\raggedright
Document
\end{minipage} & \begin{minipage}[b]{\linewidth}\raggedright
Significance
\end{minipage} \\
\midrule\noalign{}
\endhead
\bottomrule\noalign{}
\endlastfoot
1993 & PKCS\#7 v1.5 & RSA Laboratories\textquotesingle{} original
specification. Defined SignedData, EnvelopedData, and the ContentInfo
wrapper. Still referenced by name in PDF SubFilter values. \\
1998 & RFC 2315 & PKCS\#7 republished as an informational RFC. No
technical changes, but gave the format an IETF-citable reference. \\
1999 & RFC 2630 & First standards-track CMS. Introduced version 3
SignedData, SubjectKeyIdentifier as an alternative SignerIdentifier, and
formalized attribute handling. \\
2004 & RFC 3852 & Clarified signed attributes encoding requirements.
Tightened consistency rules for digestAlgorithms SET. \\
2009 & RFC 5652 --- Cryptographic Message Syntax (CMS) & Current
standard. Obsoletes RFC 3852. The normative reference for all modern PDF
signature work. \\
\end{longtable}

\section{Why PDF Uses CMS}\label{_why_pdf_uses_cms}

When Adobe added digital signatures to PDF 1.3 in 1999, they needed a
container format for the cryptographic signature. They chose PKCS\#7
(which was becoming CMS at exactly the same time) for pragmatic reasons:

\begin{enumerate}
\def\labelenumi{\arabic{enumi}.}
\item
  \textbf{It was the only real option.} In 1999, the choices for
  cryptographic message formats were PKCS\#7/CMS and PGP. PKCS\#7/CMS
  was the enterprise and government standard; PGP was associated with
  personal email encryption. For a document format targeting business
  use, PKCS\#7/CMS was the obvious choice.
\item
  \textbf{It bundles everything needed for validation.} A CMS
  \texttt{SignedData} structure can carry the signature value, the
  signer's certificate, intermediate CA certificates, CRLs, and
  timestamp tokens --- everything a validator needs without making
  network requests. This self-contained property is critical for PDF,
  where a signed document might be validated years after signing, on an
  air-gapped system with no network access.
\item
  \textbf{It supports detached signatures natively.} PDF needs the
  signed content (the byte range) to live in the PDF file, not inside
  the CMS structure. CMS's detached signature mode was designed for
  exactly this use case.
\end{enumerate}

The practical consequence is that the \texttt{/Contents} value in every
PDF signature dictionary is a hex-encoded, DER-encoded CMS
\texttt{SignedData} structure. When ISO 32000-2:2020 says the
\texttt{/Contents} field contains "the DER-encoded PKCS\#7 binary data
object," it means a CMS \texttt{SignedData} as defined by RFC 5652 ---
Cryptographic Message Syntax (CMS). The terms are used interchangeably
in the spec, which is historically understandable but terminologically
sloppy.

\section{The ContentInfo Wrapper}\label{_the_contentinfo_wrapper}

Every CMS structure is wrapped in a \texttt{ContentInfo} --- the
outermost ASN.1 envelope that identifies what type of CMS object is
inside.

\begin{Shaded}
\begin{Highlighting}[]
\NormalTok{ContentInfo ::= }\DataTypeTok{SEQUENCE}\NormalTok{ \{}
\NormalTok{    contentType     ContentType,     }\CommentTok{{-}{-} OID identifying the CMS type}
\NormalTok{    content    [}\DecValTok{0}\NormalTok{]  EXPLICIT ANY DEFINED BY contentType}
\NormalTok{\}}

\NormalTok{ContentType ::= OBJECT IDENTIFIER}
\end{Highlighting}
\end{Shaded}

For a PDF signature, the \texttt{contentType} is always:

\begin{verbatim}
1.2.840.113549.1.7.2   -- id-signedData
\end{verbatim}

This OID lives under the RSA Security arc (\texttt{1.2.840.113549}) in
the PKCS\#7 namespace (\texttt{.1.7}), with \texttt{.2} identifying
\texttt{SignedData} specifically. The \texttt{content} field contains
the \texttt{SignedData} structure itself, tagged with EXPLICIT
\texttt{{[}0{]}}.

\begin{longtable}[]{@{}
  >{\raggedright\arraybackslash}p{(\linewidth - 4\tabcolsep) * \real{0.3333}}
  >{\raggedright\arraybackslash}p{(\linewidth - 4\tabcolsep) * \real{0.3333}}
  >{\raggedright\arraybackslash}p{(\linewidth - 4\tabcolsep) * \real{0.3333}}@{}}
\caption{CMS ContentType OIDs}\tabularnewline
\toprule\noalign{}
\begin{minipage}[b]{\linewidth}\raggedright
OID
\end{minipage} & \begin{minipage}[b]{\linewidth}\raggedright
Name
\end{minipage} & \begin{minipage}[b]{\linewidth}\raggedright
Usage in PDF
\end{minipage} \\
\midrule\noalign{}
\endfirsthead
\toprule\noalign{}
\begin{minipage}[b]{\linewidth}\raggedright
OID
\end{minipage} & \begin{minipage}[b]{\linewidth}\raggedright
Name
\end{minipage} & \begin{minipage}[b]{\linewidth}\raggedright
Usage in PDF
\end{minipage} \\
\midrule\noalign{}
\endhead
\bottomrule\noalign{}
\endlastfoot
\texttt{1.2.840.113549.1.7.1} & \texttt{id-data} & Used as the
\texttt{eContentType} inside SignedData to indicate that the signed
content is arbitrary data (the PDF byte range) \\
\texttt{1.2.840.113549.1.7.2} & \texttt{id-signedData} & The outermost
ContentInfo type for all PDF signatures \\
\texttt{1.2.840.113549.1.7.3} & \texttt{id-envelopedData} & Not used in
PDF signatures (used for encryption) \\
\texttt{1.2.840.113549.1.9.16.1.4} & \texttt{id-ct-TSTInfo} & Used as
the \texttt{eContentType} for RFC 3161 timestamp tokens embedded in PDF
signatures \\
\end{longtable}

The OID \texttt{1.2.840.113549.1.7.2} is the first thing a CMS parser
encounters when reading a PDF signature. If this OID is missing or
wrong, the entire structure is unparsable. When you are debugging a
broken signature by examining raw hex, look for the DER encoding of this
OID near the beginning of the \texttt{/Contents} value:
\texttt{06\ 09\ 2A\ 86\ 48\ 86\ F7\ 0D\ 01\ 07\ 02}. Those eleven bytes
(tag \texttt{06}, length \texttt{09}, then the nine OID bytes) are the
signature of a CMS \texttt{SignedData} structure.

\section{The SignedData Structure}\label{_the_signeddata_structure}

Inside the \texttt{ContentInfo} wrapper, the \texttt{SignedData}
structure contains everything related to the signature:

\begin{Shaded}
\begin{Highlighting}[]
\NormalTok{SignedData ::= }\DataTypeTok{SEQUENCE}\NormalTok{ \{}
\NormalTok{    version            CMSVersion,                          }
\NormalTok{    digestAlgorithms   DigestAlgorithmIdentifiers,          }
\NormalTok{    encapContentInfo   EncapsulatedContentInfo,             }
\NormalTok{    certificates  [}\DecValTok{0}\NormalTok{]  IMPLICIT CertificateSet }\KeywordTok{OPTIONAL}\NormalTok{,    }
\NormalTok{    crls          [}\DecValTok{1}\NormalTok{]  IMPLICIT RevocationInfoChoices }\KeywordTok{OPTIONAL}\NormalTok{, }
\NormalTok{    signerInfos        SignerInfos                          }
\NormalTok{\}}
\end{Highlighting}
\end{Shaded}

\begin{itemize}
\item
  Version number: 1, 3, 4, or 5, depending on which optional fields are
  present. For PDF signatures with X.509 certificates and
  \texttt{issuerAndSerialNumber} in SignerInfo, this is typically
  \textbf{1}.
\item
  SET OF digest algorithm identifiers. Must include every algorithm used
  by any SignerInfo in this structure. For a single-signer PDF, this is
  a set with one element (e.g., SHA-256).
\item
  The content that was signed. For PDF detached signatures, the
  \texttt{eContent} field is \textbf{absent} --- the content lives in
  the PDF file itself.
\item
  The certificate collection. Typically contains the signer's end-entity
  certificate plus all intermediate CA certificates needed to build the
  chain to a trust anchor.
\item
  CRLs or OCSP responses for revocation checking. Rarely populated in
  basic PDF signatures; more common in PAdES-B-LT and PAdES-B-LTA
  profiles where revocation data is embedded for long-term validation.
\item
  SET OF SignerInfo structures --- one per signer. PDF signatures almost
  always have exactly one SignerInfo.
\end{itemize}

\subsection{Version Number Rules}\label{_version_number_rules}

The version field is not arbitrary --- RFC 5652 --- Cryptographic
Message Syntax (CMS) Section 5.1 specifies strict rules:

\begin{longtable}[]{@{}
  >{\raggedright\arraybackslash}p{(\linewidth - 2\tabcolsep) * \real{0.2000}}
  >{\raggedright\arraybackslash}p{(\linewidth - 2\tabcolsep) * \real{0.8000}}@{}}
\caption{SignedData version selection (RFC 5652 --- Cryptographic
Message Syntax (CMS) Section 5.1)}\tabularnewline
\toprule\noalign{}
\begin{minipage}[b]{\linewidth}\raggedright
Version
\end{minipage} & \begin{minipage}[b]{\linewidth}\raggedright
Condition
\end{minipage} \\
\midrule\noalign{}
\endfirsthead
\toprule\noalign{}
\begin{minipage}[b]{\linewidth}\raggedright
Version
\end{minipage} & \begin{minipage}[b]{\linewidth}\raggedright
Condition
\end{minipage} \\
\midrule\noalign{}
\endhead
\bottomrule\noalign{}
\endlastfoot
1 & No attribute certificates in \texttt{certificates}, no
\texttt{OtherRevocationInfoFormat} in \texttt{crls}, no version-3
\texttt{SignerInfo}, and the \texttt{encapContentInfo.eContentType} is
\texttt{id-data} \\
3 & The \texttt{encapContentInfo.eContentType} is other than
\texttt{id-data} (e.g., \texttt{id-ct-TSTInfo} for a timestamp token),
OR the structure contains \texttt{SignerInfo} version 3 \\
4 & Attribute certificates with version 2 are present in
\texttt{certificates} \\
5 & \texttt{OtherRevocationInfoFormat} is present in \texttt{crls} \\
\end{longtable}

For the vast majority of PDF signatures, the version is \textbf{1}.
Timestamp tokens embedded as unsigned attributes use version \textbf{3}
because their inner \texttt{eContentType} is \texttt{id-ct-TSTInfo}
rather than \texttt{id-data}.

Some older signing libraries (notably early versions of iText before
version 5.3) hardcoded the version to 1 regardless of the content type.
This produced technically non-compliant CMS structures that happened to
validate in most tools but failed strict validators like the EU's DSS
library. If you encounter a CMS blob with version 1 but a
non-\texttt{id-data} content type, that is a bug in the signing library,
not a feature.

\subsection{The digestAlgorithms SET}\label{_the_digestalgorithms_set}

This is a SET OF \texttt{AlgorithmIdentifier}, declared at the top level
of \texttt{SignedData} to allow a parser to identify all digest
algorithms in use before reading any \texttt{SignerInfo} entries. The
intent is optimization: a validator reading a large detached signature
can start computing digests in parallel before it reaches the
\texttt{SignerInfo} structures.

For PDF signatures with a single signer using SHA-256:

\begin{Shaded}
\begin{Highlighting}[]
\NormalTok{digestAlgorithms:}
  \DataTypeTok{SET}\NormalTok{ (}\DecValTok{1}\NormalTok{ element)}
    \DataTypeTok{SEQUENCE}
\NormalTok{      OID: }\DecValTok{2}\NormalTok{.}\DecValTok{16}\NormalTok{.}\DecValTok{840}\NormalTok{.}\DecValTok{1}\NormalTok{.}\DecValTok{101}\NormalTok{.}\DecValTok{3}\NormalTok{.}\DecValTok{4}\NormalTok{.}\DecValTok{2}\NormalTok{.}\DecValTok{1}\NormalTok{  (SHA}\DecValTok{{-}256}\NormalTok{)}
      \DataTypeTok{NULL}\NormalTok{                           (parameters, absent for SHA}\DecValTok{{-}2}\NormalTok{)}
\end{Highlighting}
\end{Shaded}

The \texttt{NULL} parameters field for SHA-2 algorithms is a persistent
source of interoperability issues. RFC 5652 --- Cryptographic Message
Syntax (CMS) says the parameters "SHOULD be omitted" for SHA-256 and
related algorithms, but some older implementations include an explicit
\texttt{NULL} (tag \texttt{05\ 00}), and some validators reject
structures where it is absent. BouncyCastle 1.80 includes the explicit
\texttt{NULL} by default, which is the safe choice for maximum
compatibility.

\subsection{EncapsulatedContentInfo: Detached vs
Encapsulated}\label{_encapsulatedcontentinfo_detached_vs_encapsulated}

\begin{Shaded}
\begin{Highlighting}[]
\NormalTok{EncapsulatedContentInfo ::= }\DataTypeTok{SEQUENCE}\NormalTok{ \{}
\NormalTok{    eContentType    ContentType,}
\NormalTok{    eContent   [}\DecValTok{0}\NormalTok{]  EXPLICIT }\DataTypeTok{OCTET} \DataTypeTok{STRING} \KeywordTok{OPTIONAL}
\NormalTok{\}}
\end{Highlighting}
\end{Shaded}

This is where detached and encapsulated signatures diverge:

\begin{itemize}
\item
  \textbf{Detached} (PDF standard): \texttt{eContent} is
  \textbf{absent}. The signed content lives outside the CMS
  structure --- in the PDF file, specified by the \texttt{/ByteRange}
  array. The \texttt{eContentType} is \texttt{id-data}
  (\texttt{1.2.840.113549.1.7.1}).
\item
  \textbf{Encapsulated}: \texttt{eContent} is \textbf{present} and
  contains the signed data. This mode is used in S/MIME email signatures
  where the message body is embedded in the CMS structure.
\end{itemize}

For the SubFilter \texttt{adbe.pkcs7.sha1} (deprecated), the CMS uses a
hybrid approach: the \texttt{eContent} field contains the SHA-1 hash of
the byte range (not the byte range itself), making it technically an
encapsulated signature over a hash rather than a detached signature over
the content. This design was misguided and is one of the reasons
\texttt{adbe.pkcs7.sha1} was deprecated.

If you are building a PDF signing implementation and you accidentally
pass \texttt{true} for the encapsulate parameter in
\texttt{CMSSignedDataGenerator.generate()}, the entire byte range
content will be embedded inside the CMS structure, inside the
\texttt{/Contents} hex string, inside the PDF. The result is a signature
that is dramatically larger than necessary (potentially megabytes),
wastes space, and will be rejected by validators that expect a detached
signature when the SubFilter is \texttt{adbe.pkcs7.detached} or
\texttt{ETSI.CAdES.detached}.

\subsection{The certificates Field}\label{_the_certificates_field}

The \texttt{certificates} field is an IMPLICIT SET that can contain:

\begin{itemize}
\item
  X.509 certificates (\texttt{Certificate} type from RFC 5280 ---
  Internet X.509 PKI Certificate and CRL Profile)
\item
  Attribute certificates (v1 or v2) --- rarely used in PDF signatures
\item
  Other certificate formats --- essentially never used
\end{itemize}

For PDF signatures, this field should contain at a minimum the signer's
end-entity certificate and all intermediate CA certificates needed to
chain to a trust anchor. The root CA certificate is traditionally
\emph{not} included because the validator must already trust
it --- including it adds bytes without adding trust.

The ordering of certificates in this SET is \textbf{undefined} by RFC
5652 --- Cryptographic Message Syntax (CMS). A validator must be
prepared to receive certificates in any order and build the chain by
matching issuer/subject Distinguished Names and Authority/Subject Key
Identifiers.

\section{SignerInfo: The Core of the
Signature}\label{ch06-signerinfo-deep-dive}

Each \texttt{SignerInfo} structure represents one signer's contribution
to the \texttt{SignedData}:

\begin{Shaded}
\begin{Highlighting}[]
\NormalTok{SignerInfo ::= }\DataTypeTok{SEQUENCE}\NormalTok{ \{}
\NormalTok{    version               CMSVersion,                    }
\NormalTok{    sid                   SignerIdentifier,               }
\NormalTok{    digestAlgorithm       DigestAlgorithmIdentifier,      }
\NormalTok{    signedAttrs      [}\DecValTok{0}\NormalTok{]  IMPLICIT SignedAttributes }\KeywordTok{OPTIONAL}\NormalTok{, }
\NormalTok{    signatureAlgorithm    SignatureAlgorithmIdentifier,   }
\NormalTok{    signature             SignatureValue,                 }
\NormalTok{    unsignedAttrs    [}\DecValTok{1}\NormalTok{]  IMPLICIT UnsignedAttributes }\KeywordTok{OPTIONAL} 
\NormalTok{\}}
\end{Highlighting}
\end{Shaded}

\begin{itemize}
\item
  Version 1 if \texttt{sid} is \texttt{issuerAndSerialNumber}, version 3
  if \texttt{sid} is \texttt{subjectKeyIdentifier}.
\item
  Identifies which certificate belongs to this signer.
\item
  The hash algorithm used to compute the message digest (e.g., SHA-256).
\item
  Attributes that are included in the signature computation --- critical
  for PDF signatures.
\item
  The signature algorithm (e.g., SHA-256 with RSA, or ECDSA with
  SHA-256).
\item
  The raw signature bytes.
\item
  Attributes added after signing --- used for timestamps,
  counter-signatures.
\end{itemize}

\subsection{SignerIdentifier: Finding the Signer's
Certificate}\label{_signeridentifier_finding_the_signers_certificate}

\begin{Shaded}
\begin{Highlighting}[]
\NormalTok{SignerIdentifier ::= }\DataTypeTok{CHOICE}\NormalTok{ \{}
\NormalTok{    issuerAndSerialNumber    IssuerAndSerialNumber,}
\NormalTok{    subjectKeyIdentifier [}\DecValTok{0}\NormalTok{] SubjectKeyIdentifier}
\NormalTok{\}}
\end{Highlighting}
\end{Shaded}

The \texttt{issuerAndSerialNumber} form is dominant in PDF signatures.
It contains the issuer DN and serial number of the signer's
certificate --- enough to uniquely identify the certificate within the
\texttt{certificates} SET.

The \texttt{subjectKeyIdentifier} form uses the Subject Key Identifier
extension from the certificate. It is more compact but less commonly
used because older validators do not always support it.

\subsection{Signed Attributes: What Gets
Signed}\label{_signed_attributes_what_gets_signed}

When signed attributes are present (and for PDF signatures, they
\textbf{must} be present per ISO 32000-2:2020), the signature is
computed over the DER encoding of the signed attributes --- \emph{not}
over the content itself. The content's digest is included as a signed
attribute, creating an indirect binding:

\begin{verbatim}
content  --(hash)-->  messageDigest attribute
                            |
signed attributes  --(sign)-->  signature value
\end{verbatim}

The mandatory signed attributes for PDF signatures are:

\begin{longtable}[]{@{}
  >{\raggedright\arraybackslash}p{(\linewidth - 4\tabcolsep) * \real{0.2222}}
  >{\raggedright\arraybackslash}p{(\linewidth - 4\tabcolsep) * \real{0.3333}}
  >{\raggedright\arraybackslash}p{(\linewidth - 4\tabcolsep) * \real{0.4444}}@{}}
\caption{Required signed attributes in PDF CMS
signatures}\tabularnewline
\toprule\noalign{}
\begin{minipage}[b]{\linewidth}\raggedright
Attribute
\end{minipage} & \begin{minipage}[b]{\linewidth}\raggedright
OID
\end{minipage} & \begin{minipage}[b]{\linewidth}\raggedright
Description
\end{minipage} \\
\midrule\noalign{}
\endfirsthead
\toprule\noalign{}
\begin{minipage}[b]{\linewidth}\raggedright
Attribute
\end{minipage} & \begin{minipage}[b]{\linewidth}\raggedright
OID
\end{minipage} & \begin{minipage}[b]{\linewidth}\raggedright
Description
\end{minipage} \\
\midrule\noalign{}
\endhead
\bottomrule\noalign{}
\endlastfoot
\texttt{contentType} & \texttt{1.2.840.113549.1.9.3} & Must be
\texttt{id-data} (\texttt{1.2.840.113549.1.7.1}). Identifies the type of
content that was signed. Required by RFC 5652 --- Cryptographic Message
Syntax (CMS) whenever signed attributes are present. \\
\texttt{messageDigest} & \texttt{1.2.840.113549.1.9.4} & The hash of the
content (the PDF byte range). This is the attribute that binds the CMS
signature to the PDF document bytes. The validator recomputes this hash
and compares it against the value in this attribute. \\
\texttt{signingTime} & \texttt{1.2.840.113549.1.9.5} & The time the
signer claims to have signed. This is self-asserted and not
cryptographically verified --- use a timestamp authority for trusted
time (see \texttt{unsignedAttrs}). \\
\texttt{ESS-signing-certificate-v2} &
\texttt{1.2.840.113549.1.9.16.2.47} & Required by PAdES (ETSI EN 319
122-1 V1.1.1 (2016-04)). Contains a hash of the signer's certificate,
binding the signature to a specific certificate and preventing
certificate substitution attacks. Without this attribute, an attacker
could replace the signer's certificate with a different certificate that
has the same public key. \\
\end{longtable}

The \texttt{signingTime} attribute is optional per RFC 5652 ---
Cryptographic Message Syntax (CMS) but effectively mandatory in
practice. Every PDF viewer displays the signing time, and a signature
without one looks suspicious to users. However, the signing time is
\emph{not} cryptographically trustworthy --- the signer controls the
clock. For legally binding timestamps, use an RFC 3161 timestamp token
in the unsigned attributes (Chapter 18).

\subsection{The ESS-signing-certificate-v2
Attribute}\label{_the_ess_signing_certificate_v2_attribute}

This attribute deserves special attention because it is the dividing
line between a basic CMS signature and a CAdES-compliant
signature --- and therefore between a signature that merely works and
one that meets European regulatory requirements.

The structure contains:

\begin{Shaded}
\begin{Highlighting}[]
\NormalTok{SigningCertificateV2 ::= }\DataTypeTok{SEQUENCE}\NormalTok{ \{}
\NormalTok{    certs        }\DataTypeTok{SEQUENCE} \DataTypeTok{OF}\NormalTok{ ESSCertIDv2,}
\NormalTok{    policies     }\DataTypeTok{SEQUENCE} \DataTypeTok{OF}\NormalTok{ PolicyInformation }\KeywordTok{OPTIONAL}
\NormalTok{\}}

\NormalTok{ESSCertIDv2 ::= }\DataTypeTok{SEQUENCE}\NormalTok{ \{}
\NormalTok{    hashAlgorithm  AlgorithmIdentifier }\KeywordTok{DEFAULT}\NormalTok{ \{sha}\DecValTok{{-}256}\NormalTok{\},}
\NormalTok{    certHash       Hash,              }\CommentTok{{-}{-} hash of the signer\textquotesingle{}s certificate}
\NormalTok{    issuerSerial   IssuerSerial }\KeywordTok{OPTIONAL}
\NormalTok{\}}
\end{Highlighting}
\end{Shaded}

The \texttt{certHash} is a hash of the signer's DER-encoded certificate.
A validator computes this hash independently and compares it against the
attribute value, ensuring that the certificate in the CMS structure is
the one the signer intended to use.

Without this attribute, a subtle attack is possible: if two certificates
share the same public key (which can happen if a certificate is
re-issued with different extensions or validity periods), an attacker
could swap the certificate in the CMS structure without invalidating the
signature. The ESS-signing-certificate-v2 attribute prevents this by
binding the signature to the exact certificate, not just the public key.

PAdES baseline profiles (PAdES-B-B and above) require
ESS-signing-certificate-v2. A PDF signature that uses the SubFilter
\texttt{ETSI.CAdES.detached} but lacks this attribute is non-compliant
with ETSI EN 319 142-1 V1.1.1 (2016-04). The EU's DSS validation library
will flag it as a format failure. If you are implementing PDF signing
and need European compliance, this attribute is not optional.

\section{The Signed Attributes DER Encoding
Trap}\label{_the_signed_attributes_der_encoding_trap}

This section describes the single most common cause of CMS signature
validation failures. It has tripped up every signing library author at
least once, and it will trip you up too if you are building or debugging
a CMS implementation at the byte level.

\subsection{The Problem}\label{_the_problem}

RFC 5652 --- Cryptographic Message Syntax (CMS) Section 5.4 specifies
how signed attributes are encoded for signature computation:

\begin{quote}
The result of the message digest calculation process depends on whether
the signedAttrs field is present. When the field is present,
{[}\ldots\hspace{0pt}{]} the complete DER encoding of the SignedAttrs
value contained in the signedAttrs field is {[}used{]}.
{[}\ldots\hspace{0pt}{]} the IMPLICIT {[}0{]} tag in the signedAttrs is
not used for the DER encoding, rather an EXPLICIT SET OF tag is used.

---  RFC 5652 --- Cryptographic Message Syntax (CMS) Section 5.4
\end{quote}

Here is what this means in practice:

\begin{enumerate}
\def\labelenumi{\arabic{enumi}.}
\item
  In the \texttt{SignerInfo} structure, the \texttt{signedAttrs} field
  is tagged \texttt{{[}0{]}\ IMPLICIT}, meaning the ASN.1 tag byte is
  \texttt{A0} (context class, constructed, tag 0).
\item
  But when computing the signature over the signed attributes, the tag
  must be replaced with \texttt{31} (UNIVERSAL SET OF) --- the tag that
  would be used if the attributes were a standalone SET OF, not an
  implicitly tagged field.
\item
  The rest of the encoding (length bytes and content) remains identical.
\end{enumerate}

This means the bytes that are signed are \emph{not} the bytes that
appear in the CMS structure. You must re-tag the signed attributes
before computing or verifying the signature.

\subsection{Why This Exists}\label{_why_this_exists}

The reason is rooted in ASN.1 encoding theory. The \texttt{signedAttrs}
field uses IMPLICIT tagging to save one byte in the encoding (implicit
tagging replaces the original tag; explicit tagging adds a wrapper). But
the signature must be computed over the \emph{canonical} encoding of the
attributes as a SET OF, not over their implicitly tagged representation
within the SignerInfo. The two encodings differ only in the first byte
(\texttt{A0} vs \texttt{31}), but that single byte difference is enough
to invalidate the signature.

If you are computing a CMS signature manually (perhaps because you are
building a remote signing integration where the hash is computed locally
and the signature is applied remotely), you \textbf{must} re-tag the
signed attributes before hashing them. Use tag \texttt{31} (SET OF), not
\texttt{A0} (IMPLICIT {[}0{]}). BouncyCastle handles this correctly in
its \texttt{CMSSignedDataGenerator}, but if you are working at a lower
level --- constructing CMS structures from raw ASN.1 objects, or
implementing a signature interface for an HSM --- this is the byte that
will break you.

\begin{Shaded}
\begin{Highlighting}[]
\KeywordTok{import} \ImportTok{org}\OperatorTok{.}\ImportTok{bouncycastle}\OperatorTok{.}\ImportTok{asn1}\OperatorTok{.}\ImportTok{ASN1EncodableVector}\OperatorTok{;}
\KeywordTok{import} \ImportTok{org}\OperatorTok{.}\ImportTok{bouncycastle}\OperatorTok{.}\ImportTok{asn1}\OperatorTok{.}\ImportTok{ASN1Set}\OperatorTok{;}
\KeywordTok{import} \ImportTok{org}\OperatorTok{.}\ImportTok{bouncycastle}\OperatorTok{.}\ImportTok{asn1}\OperatorTok{.}\ImportTok{BERTags}\OperatorTok{;}
\KeywordTok{import} \ImportTok{org}\OperatorTok{.}\ImportTok{bouncycastle}\OperatorTok{.}\ImportTok{asn1}\OperatorTok{.}\ImportTok{DERSet}\OperatorTok{;}

\CommentTok{/**}
 \CommentTok{*}\NormalTok{ Demonstrate the signed attributes re}\CommentTok{{-}}\NormalTok{tagging requirement}\CommentTok{.}
 \CommentTok{*}\NormalTok{ In practice}\CommentTok{,}\NormalTok{ BouncyCastle handles this internally }\CommentTok{{-}{-}}\NormalTok{ this}
 \CommentTok{*}\NormalTok{ code is for understanding}\CommentTok{,}\NormalTok{ not production use}\CommentTok{.}
 \CommentTok{*/}
\KeywordTok{public} \KeywordTok{class}\NormalTok{ SignedAttrsEncodingExample }\OperatorTok{\{}

    \CommentTok{/**}
     \CommentTok{*}\NormalTok{ Re}\CommentTok{{-}}\NormalTok{encode signed attributes with SET OF tag for signing}\CommentTok{.}
     \CommentTok{*}\NormalTok{ The input is the IMPLICIT }\CommentTok{[0]{-}}\NormalTok{tagged bytes from SignerInfo}\CommentTok{;}
     \CommentTok{*}\NormalTok{ the output replaces tag }\CommentTok{0xA0}\NormalTok{ with }\CommentTok{0x31.}
     \CommentTok{*/}
    \KeywordTok{public} \DataTypeTok{static} \DataTypeTok{byte}\OperatorTok{[]} \FunctionTok{retagForSigning}\OperatorTok{(}\DataTypeTok{byte}\OperatorTok{[]}\NormalTok{ implicitTaggedAttrs}\OperatorTok{)} \OperatorTok{\{}
        \DataTypeTok{byte}\OperatorTok{[]}\NormalTok{ retagged }\OperatorTok{=}\NormalTok{ implicitTaggedAttrs}\OperatorTok{.}\FunctionTok{clone}\OperatorTok{();}
        \CommentTok{// Replace IMPLICIT [0] tag (0xA0) with SET OF tag (0x31)}
        \ControlFlowTok{if} \OperatorTok{(}\NormalTok{retagged}\OperatorTok{[}\DecValTok{0}\OperatorTok{]} \OperatorTok{==} \OperatorTok{(}\DataTypeTok{byte}\OperatorTok{)} \BaseNTok{0xA0}\OperatorTok{)} \OperatorTok{\{}
\NormalTok{            retagged}\OperatorTok{[}\DecValTok{0}\OperatorTok{]} \OperatorTok{=} \OperatorTok{(}\DataTypeTok{byte}\OperatorTok{)} \OperatorTok{(}\NormalTok{BERTags}\OperatorTok{.}\FunctionTok{SET} \OperatorTok{|}\NormalTok{ BERTags}\OperatorTok{.}\FunctionTok{CONSTRUCTED}\OperatorTok{);} \CommentTok{// 0x31}
        \OperatorTok{\}}
        \ControlFlowTok{return}\NormalTok{ retagged}\OperatorTok{;}
    \OperatorTok{\}}
\OperatorTok{\}}
\end{Highlighting}
\end{Shaded}

When debugging a signature that fails cryptographic verification but the
\texttt{messageDigest} attribute matches the content hash, the
re-tagging issue is the first thing to check. Extract the signed
attributes bytes from the CMS structure, verify that the first byte is
\texttt{31} (not \texttt{A0}) when fed into the signature verification,
and compare the hash of those bytes against the decrypted signature
value.

\section{Unsigned Attributes: Post-Signature
Additions}\label{_unsigned_attributes_post_signature_additions}

Unsigned attributes are added to the \texttt{SignerInfo} \emph{after}
the signature has been computed. They are not covered by the signature
and can be added or modified without invalidating it. This is by
design --- the primary use case is adding information that is not
available at signing time.

The most important unsigned attributes for PDF signatures:

\begin{longtable}[]{@{}
  >{\raggedright\arraybackslash}p{(\linewidth - 4\tabcolsep) * \real{0.2222}}
  >{\raggedright\arraybackslash}p{(\linewidth - 4\tabcolsep) * \real{0.3333}}
  >{\raggedright\arraybackslash}p{(\linewidth - 4\tabcolsep) * \real{0.4444}}@{}}
\caption{Common unsigned attributes in PDF CMS
signatures}\tabularnewline
\toprule\noalign{}
\begin{minipage}[b]{\linewidth}\raggedright
Attribute
\end{minipage} & \begin{minipage}[b]{\linewidth}\raggedright
OID
\end{minipage} & \begin{minipage}[b]{\linewidth}\raggedright
Description
\end{minipage} \\
\midrule\noalign{}
\endfirsthead
\toprule\noalign{}
\begin{minipage}[b]{\linewidth}\raggedright
Attribute
\end{minipage} & \begin{minipage}[b]{\linewidth}\raggedright
OID
\end{minipage} & \begin{minipage}[b]{\linewidth}\raggedright
Description
\end{minipage} \\
\midrule\noalign{}
\endhead
\bottomrule\noalign{}
\endlastfoot
\texttt{signature-time-stamp} & \texttt{1.2.840.113549.1.9.16.2.14} & An
RFC 3161 timestamp token over the signature value. Provides a trusted,
third-party assertion of the time the signature existed. This is the
mechanism for timestamping described in Chapter 18. \\
\texttt{counterSignature} & \texttt{1.2.840.113549.1.9.6} & A
counter-signature --- another signer's signature over the original
signature value. Rarely used in PDF signatures directly; timestamps are
the modern replacement. \\
\end{longtable}

The timestamp token in the \texttt{signature-time-stamp} attribute is
itself a CMS \texttt{SignedData} structure, with its
\texttt{eContentType} set to \texttt{id-ct-TSTInfo}
(\texttt{1.2.840.113549.1.9.16.1.4}). This creates a nested
CMS-within-CMS structure: the outer \texttt{SignedData} is the PDF
signature, and the inner \texttt{SignedData} (inside an unsigned
attribute of the outer) is the timestamp.

\section{CAdES: CMS Advanced Electronic
Signatures}\label{_cades_cms_advanced_electronic_signatures}

The basic CMS \texttt{SignedData} structure defined by RFC 5652 ---
Cryptographic Message Syntax (CMS) is cryptographically complete but
lacks features needed for European regulatory compliance. \textbf{CAdES}
(CMS Advanced Electronic Signatures), defined in ETSI EN 319 122-1
V1.1.1 (2016-04), extends CMS with additional signed attributes and
profile rules.

The key additions CAdES makes to CMS for PDF signing:

\begin{enumerate}
\def\labelenumi{\arabic{enumi}.}
\item
  \textbf{ESS-signing-certificate-v2} (discussed above) --- mandatory in
  all CAdES profiles, binds the signature to a specific certificate.
\item
  \textbf{Signature policy identifier} --- optional attribute that
  references a specific signature policy (a set of rules defining the
  meaning and validation requirements of the signature). Used in some
  government contexts.
\item
  \textbf{Commitment type indication} --- optional attribute declaring
  what the signer intends by signing (proof of origin, proof of receipt,
  proof of approval, etc.).
\item
  \textbf{Signer location} --- optional attribute with the signer's
  claimed location at signing time.
\end{enumerate}

The relationship between CMS, CAdES, and PAdES is hierarchical:

\begin{verbatim}
CMS (RFC 5652)
  |
  +-- CAdES (ETSI EN 319 122) = CMS + mandatory signed attributes
        |
        +-- PAdES (ETSI EN 319 142) = CAdES embedded in PDF per ISO 32000-2
\end{verbatim}

A PAdES signature is a CAdES signature is a CMS signature. Each layer
adds constraints but does not remove any. A basic
\texttt{adbe.pkcs7.detached} signature is CMS but not necessarily CAdES.
An \texttt{ETSI.CAdES.detached} signature is CMS, CAdES, and PAdES.

\section{Old PKCS\#7 vs Modern CMS: What
Changed}\label{_old_pkcs7_vs_modern_cms_what_changed}

If you work with PDF signatures long enough, you will encounter CMS
structures produced by older libraries that do not quite match what RFC
5652 --- Cryptographic Message Syntax (CMS) specifies. The most common
issues trace back to libraries that implemented the original PKCS\#7
v1.5 rather than CMS, or that implemented an intermediate version with
bugs that were never fixed because "it worked in Acrobat."

\begin{longtable}[]{@{}
  >{\raggedright\arraybackslash}p{(\linewidth - 4\tabcolsep) * \real{0.2500}}
  >{\raggedright\arraybackslash}p{(\linewidth - 4\tabcolsep) * \real{0.3750}}
  >{\raggedright\arraybackslash}p{(\linewidth - 4\tabcolsep) * \real{0.3750}}@{}}
\caption{Differences between legacy PKCS\#7 and modern CMS in the
wild}\tabularnewline
\toprule\noalign{}
\begin{minipage}[b]{\linewidth}\raggedright
Issue
\end{minipage} & \begin{minipage}[b]{\linewidth}\raggedright
Legacy PKCS\#7 Behavior
\end{minipage} & \begin{minipage}[b]{\linewidth}\raggedright
Modern CMS (RFC 5652 --- Cryptographic Message Syntax (CMS))
\end{minipage} \\
\midrule\noalign{}
\endfirsthead
\toprule\noalign{}
\begin{minipage}[b]{\linewidth}\raggedright
Issue
\end{minipage} & \begin{minipage}[b]{\linewidth}\raggedright
Legacy PKCS\#7 Behavior
\end{minipage} & \begin{minipage}[b]{\linewidth}\raggedright
Modern CMS (RFC 5652 --- Cryptographic Message Syntax (CMS))
\end{minipage} \\
\midrule\noalign{}
\endhead
\bottomrule\noalign{}
\endlastfoot
SignedData version & Always 1, regardless of content type & Version
depends on which optional fields are present (1, 3, 4, or 5) \\
digestAlgorithms consistency & May list algorithms not actually used by
any SignerInfo & Must be consistent with SignerInfo entries \\
Signed attributes encoding & Some libraries used BER (not DER) encoding
for the SET OF, causing non-deterministic byte ordering & Must use DER,
which enforces a specific ordering of SET elements \\
ESS-signing-certificate & Absent (not part of original PKCS\#7) &
Required by CAdES and PAdES \\
Certificate ordering & Some libraries assumed specific ordering & No
ordering defined; validators must handle any order \\
NULL parameters in AlgorithmIdentifier & Inconsistent --- some included,
some omitted & SHA-2 algorithms SHOULD omit parameters, but explicit
NULL is tolerated \\
\end{longtable}

Early versions of iText (before 5.3.x) and some older .NET signing
libraries produced CMS structures with BER-encoded signed attributes.
BER allows multiple valid encodings for the same data --- specifically,
SET OF elements may appear in any order. DER requires a specific
ordering (by encoded value). If the signing library used one ordering
and the validator expects another, the hash of the signed attributes
will not match and the signature will fail cryptographic verification.
The signature is not actually broken --- it was computed over
BER-encoded bytes, and if you verify against the same bytes, it passes.
But strict DER-only validators will reject it. This is the most common
"it works in Acrobat but fails in DSS" scenario.

\section{Java: Parsing CMS SignedData}\label{ch06-java-parsing}

Here is how to parse a CMS \texttt{SignedData} structure extracted from
a PDF signature's \texttt{/Contents} field, using BouncyCastle 1.80:

\begin{Shaded}
\begin{Highlighting}[]
\KeywordTok{import} \ImportTok{java}\OperatorTok{.}\ImportTok{security}\OperatorTok{.}\ImportTok{cert}\OperatorTok{.}\ImportTok{X509Certificate}\OperatorTok{;}
\KeywordTok{import} \ImportTok{java}\OperatorTok{.}\ImportTok{util}\OperatorTok{.}\ImportTok{Collection}\OperatorTok{;}

\KeywordTok{import} \ImportTok{org}\OperatorTok{.}\ImportTok{bouncycastle}\OperatorTok{.}\ImportTok{asn1}\OperatorTok{.}\ImportTok{ASN1OctetString}\OperatorTok{;}
\KeywordTok{import} \ImportTok{org}\OperatorTok{.}\ImportTok{bouncycastle}\OperatorTok{.}\ImportTok{asn1}\OperatorTok{.}\ImportTok{ASN1UTCTime}\OperatorTok{;}
\KeywordTok{import} \ImportTok{org}\OperatorTok{.}\ImportTok{bouncycastle}\OperatorTok{.}\ImportTok{asn1}\OperatorTok{.}\ImportTok{cms}\OperatorTok{.}\ImportTok{Attribute}\OperatorTok{;}
\KeywordTok{import} \ImportTok{org}\OperatorTok{.}\ImportTok{bouncycastle}\OperatorTok{.}\ImportTok{asn1}\OperatorTok{.}\ImportTok{cms}\OperatorTok{.}\ImportTok{CMSAttributes}\OperatorTok{;}
\KeywordTok{import} \ImportTok{org}\OperatorTok{.}\ImportTok{bouncycastle}\OperatorTok{.}\ImportTok{cert}\OperatorTok{.}\ImportTok{X509CertificateHolder}\OperatorTok{;}
\KeywordTok{import} \ImportTok{org}\OperatorTok{.}\ImportTok{bouncycastle}\OperatorTok{.}\ImportTok{cert}\OperatorTok{.}\ImportTok{jcajce}\OperatorTok{.}\ImportTok{JcaX509CertificateConverter}\OperatorTok{;}
\KeywordTok{import} \ImportTok{org}\OperatorTok{.}\ImportTok{bouncycastle}\OperatorTok{.}\ImportTok{cms}\OperatorTok{.}\ImportTok{CMSSignedData}\OperatorTok{;}
\KeywordTok{import} \ImportTok{org}\OperatorTok{.}\ImportTok{bouncycastle}\OperatorTok{.}\ImportTok{cms}\OperatorTok{.}\ImportTok{SignerInformation}\OperatorTok{;}
\KeywordTok{import} \ImportTok{org}\OperatorTok{.}\ImportTok{bouncycastle}\OperatorTok{.}\ImportTok{cms}\OperatorTok{.}\ImportTok{SignerInformationStore}\OperatorTok{;}
\KeywordTok{import} \ImportTok{org}\OperatorTok{.}\ImportTok{bouncycastle}\OperatorTok{.}\ImportTok{util}\OperatorTok{.}\ImportTok{Store}\OperatorTok{;}
\KeywordTok{import} \ImportTok{org}\OperatorTok{.}\ImportTok{bouncycastle}\OperatorTok{.}\ImportTok{util}\OperatorTok{.}\ImportTok{encoders}\OperatorTok{.}\ImportTok{Hex}\OperatorTok{;}

\CommentTok{/**}
 \CommentTok{*}\NormalTok{ Parse and inspect a CMS SignedData structure from a PDF signature}\CommentTok{.}
 \CommentTok{*}\NormalTok{ The input is the raw DER bytes from the }\CommentTok{/}\NormalTok{Contents hex string}\CommentTok{.}
 \CommentTok{*/}
\KeywordTok{public} \KeywordTok{class}\NormalTok{ CmsParsingExample }\OperatorTok{\{}

    \KeywordTok{public} \DataTypeTok{static} \DataTypeTok{void} \FunctionTok{inspectCmsSignature}\OperatorTok{(}\DataTypeTok{byte}\OperatorTok{[]}\NormalTok{ cmsBytes}\OperatorTok{)} \KeywordTok{throws} \BuiltInTok{Exception} \OperatorTok{\{}
\NormalTok{        CMSSignedData signedData }\OperatorTok{=} \KeywordTok{new} \FunctionTok{CMSSignedData}\OperatorTok{(}\NormalTok{cmsBytes}\OperatorTok{);}               

        \CommentTok{// {-}{-} SignedData version {-}{-}}
        \DataTypeTok{int}\NormalTok{ version }\OperatorTok{=}\NormalTok{ signedData}\OperatorTok{.}\FunctionTok{getVersion}\OperatorTok{();}
        \CommentTok{// version 1 = standard PDF signature}
        \CommentTok{// version 3 = timestamp token or SubjectKeyIdentifier used}

        \CommentTok{// {-}{-} Certificates embedded in the CMS {-}{-}}
\NormalTok{        Store}\OperatorTok{\textless{}}\NormalTok{X509CertificateHolder}\OperatorTok{\textgreater{}}\NormalTok{ certStore }\OperatorTok{=}\NormalTok{ signedData}\OperatorTok{.}\FunctionTok{getCertificates}\OperatorTok{();}
\NormalTok{        JcaX509CertificateConverter converter }\OperatorTok{=}
                \KeywordTok{new} \FunctionTok{JcaX509CertificateConverter}\OperatorTok{().}\FunctionTok{setProvider}\OperatorTok{(}\StringTok{"BC"}\OperatorTok{);}

        \CommentTok{// {-}{-} Iterate over SignerInfos (usually exactly one) {-}{-}}
\NormalTok{        SignerInformationStore signerInfoStore }\OperatorTok{=}\NormalTok{ signedData}\OperatorTok{.}\FunctionTok{getSignerInfos}\OperatorTok{();}

        \ControlFlowTok{for} \OperatorTok{(}\NormalTok{SignerInformation signerInfo }\OperatorTok{:}\NormalTok{ signerInfoStore}\OperatorTok{.}\FunctionTok{getSigners}\OperatorTok{())} \OperatorTok{\{}    

            \CommentTok{// Digest algorithm (e.g., "2.16.840.1.101.3.4.2.1" = SHA{-}256)}
            \BuiltInTok{String}\NormalTok{ digestOid }\OperatorTok{=}\NormalTok{ signerInfo}\OperatorTok{.}\FunctionTok{getDigestAlgOID}\OperatorTok{();}

            \CommentTok{// Signature algorithm (e.g., "1.2.840.113549.1.1.11" = SHA256withRSA)}
            \BuiltInTok{String}\NormalTok{ sigAlgOid }\OperatorTok{=}\NormalTok{ signerInfo}\OperatorTok{.}\FunctionTok{getEncryptionAlgOID}\OperatorTok{();}

            \CommentTok{// {-}{-} Signed attributes {-}{-}}
            \ControlFlowTok{if} \OperatorTok{(}\NormalTok{signerInfo}\OperatorTok{.}\FunctionTok{getSignedAttributes}\OperatorTok{()} \OperatorTok{!=} \KeywordTok{null}\OperatorTok{)} \OperatorTok{\{}

                \CommentTok{// messageDigest: the hash of the PDF byte range}
                \BuiltInTok{Attribute}\NormalTok{ mdAttr }\OperatorTok{=}\NormalTok{ signerInfo}\OperatorTok{.}\FunctionTok{getSignedAttributes}\OperatorTok{()}
                        \OperatorTok{.}\FunctionTok{get}\OperatorTok{(}\NormalTok{CMSAttributes}\OperatorTok{.}\FunctionTok{messageDigest}\OperatorTok{);}
                \ControlFlowTok{if} \OperatorTok{(}\NormalTok{mdAttr }\OperatorTok{!=} \KeywordTok{null}\OperatorTok{)} \OperatorTok{\{}
                    \DataTypeTok{byte}\OperatorTok{[]}\NormalTok{ digest }\OperatorTok{=}\NormalTok{ ASN1OctetString}\OperatorTok{.}\FunctionTok{getInstance}\OperatorTok{(}
\NormalTok{                            mdAttr}\OperatorTok{.}\FunctionTok{getAttrValues}\OperatorTok{().}\FunctionTok{getObjectAt}\OperatorTok{(}\DecValTok{0}\OperatorTok{)}
                    \OperatorTok{).}\FunctionTok{getOctets}\OperatorTok{();}
                    \CommentTok{// Hex.toHexString(digest) for human{-}readable output}
                \OperatorTok{\}}

                \CommentTok{// signingTime: the claimed signing time}
                \BuiltInTok{Attribute}\NormalTok{ stAttr }\OperatorTok{=}\NormalTok{ signerInfo}\OperatorTok{.}\FunctionTok{getSignedAttributes}\OperatorTok{()}
                        \OperatorTok{.}\FunctionTok{get}\OperatorTok{(}\NormalTok{CMSAttributes}\OperatorTok{.}\FunctionTok{signingTime}\OperatorTok{);}
                \ControlFlowTok{if} \OperatorTok{(}\NormalTok{stAttr }\OperatorTok{!=} \KeywordTok{null}\OperatorTok{)} \OperatorTok{\{}
\NormalTok{                    ASN1UTCTime time }\OperatorTok{=}\NormalTok{ ASN1UTCTime}\OperatorTok{.}\FunctionTok{getInstance}\OperatorTok{(}
\NormalTok{                            stAttr}\OperatorTok{.}\FunctionTok{getAttrValues}\OperatorTok{().}\FunctionTok{getObjectAt}\OperatorTok{(}\DecValTok{0}\OperatorTok{));}
                    \CommentTok{// time.getDate() returns a java.util.Date}
                \OperatorTok{\}}
            \OperatorTok{\}}

            \CommentTok{// {-}{-} Find the signer\textquotesingle{}s certificate {-}{-}}
            \BuiltInTok{Collection}\OperatorTok{\textless{}}\NormalTok{X509CertificateHolder}\OperatorTok{\textgreater{}}\NormalTok{ matches }\OperatorTok{=}
\NormalTok{                    certStore}\OperatorTok{.}\FunctionTok{getMatches}\OperatorTok{(}\NormalTok{signerInfo}\OperatorTok{.}\FunctionTok{getSID}\OperatorTok{());}                 
            \ControlFlowTok{for} \OperatorTok{(}\NormalTok{X509CertificateHolder holder }\OperatorTok{:}\NormalTok{ matches}\OperatorTok{)} \OperatorTok{\{}
                \BuiltInTok{X509Certificate}\NormalTok{ cert }\OperatorTok{=}\NormalTok{ converter}\OperatorTok{.}\FunctionTok{getCertificate}\OperatorTok{(}\NormalTok{holder}\OperatorTok{);}
                \CommentTok{// cert.getSubjectX500Principal() = signer identity}
                \CommentTok{// cert.getNotAfter() = certificate expiration}
            \OperatorTok{\}}

            \CommentTok{// {-}{-} Raw signature bytes {-}{-}}
            \DataTypeTok{byte}\OperatorTok{[]}\NormalTok{ signatureBytes }\OperatorTok{=}\NormalTok{ signerInfo}\OperatorTok{.}\FunctionTok{getSignature}\OperatorTok{();}                 
            \CommentTok{// signatureBytes.length = 256 for RSA{-}2048, 384 for RSA{-}3072, etc.}
        \OperatorTok{\}}
    \OperatorTok{\}}
\OperatorTok{\}}
\end{Highlighting}
\end{Shaded}

\begin{itemize}
\item
  The \texttt{CMSSignedData} constructor parses the DER bytes into an
  in-memory object graph.
\item
  PDF signatures have exactly one SignerInfo. Multiple signers use
  multiple signature dictionaries, not multiple SignerInfos in one CMS
  structure.
\item
  The \texttt{getSID()} method returns a selector that matches against
  the \texttt{SignerIdentifier} --- either
  \texttt{issuerAndSerialNumber} or \texttt{subjectKeyIdentifier}.
\item
  The raw signature value. For RSA, this is the encrypted hash. For
  ECDSA, this is the DER-encoded \texttt{(r,\ s)} pair.
\end{itemize}

\section{Java: Building CMS
SignedData}\label{_java_building_cms_signeddata}

Creating a CMS \texttt{SignedData} for PDF signing is the reverse of
parsing. The following example uses \texttt{CMSSignedDataGenerator} with
explicit control over the signed attributes:

\begin{Shaded}
\begin{Highlighting}[]
\KeywordTok{import} \ImportTok{java}\OperatorTok{.}\ImportTok{security}\OperatorTok{.}\ImportTok{PrivateKey}\OperatorTok{;}
\KeywordTok{import} \ImportTok{java}\OperatorTok{.}\ImportTok{security}\OperatorTok{.}\ImportTok{Security}\OperatorTok{;}
\KeywordTok{import} \ImportTok{java}\OperatorTok{.}\ImportTok{security}\OperatorTok{.}\ImportTok{cert}\OperatorTok{.}\ImportTok{Certificate}\OperatorTok{;}
\KeywordTok{import} \ImportTok{java}\OperatorTok{.}\ImportTok{security}\OperatorTok{.}\ImportTok{cert}\OperatorTok{.}\ImportTok{X509Certificate}\OperatorTok{;}
\KeywordTok{import} \ImportTok{java}\OperatorTok{.}\ImportTok{util}\OperatorTok{.}\ImportTok{ArrayList}\OperatorTok{;}
\KeywordTok{import} \ImportTok{java}\OperatorTok{.}\ImportTok{util}\OperatorTok{.}\ImportTok{List}\OperatorTok{;}

\KeywordTok{import} \ImportTok{org}\OperatorTok{.}\ImportTok{bouncycastle}\OperatorTok{.}\ImportTok{cert}\OperatorTok{.}\ImportTok{jcajce}\OperatorTok{.}\ImportTok{JcaCertStore}\OperatorTok{;}
\KeywordTok{import} \ImportTok{org}\OperatorTok{.}\ImportTok{bouncycastle}\OperatorTok{.}\ImportTok{cms}\OperatorTok{.}\ImportTok{CMSProcessableByteArray}\OperatorTok{;}
\KeywordTok{import} \ImportTok{org}\OperatorTok{.}\ImportTok{bouncycastle}\OperatorTok{.}\ImportTok{cms}\OperatorTok{.}\ImportTok{CMSSignedData}\OperatorTok{;}
\KeywordTok{import} \ImportTok{org}\OperatorTok{.}\ImportTok{bouncycastle}\OperatorTok{.}\ImportTok{cms}\OperatorTok{.}\ImportTok{CMSSignedDataGenerator}\OperatorTok{;}
\KeywordTok{import} \ImportTok{org}\OperatorTok{.}\ImportTok{bouncycastle}\OperatorTok{.}\ImportTok{cms}\OperatorTok{.}\ImportTok{jcajce}\OperatorTok{.}\ImportTok{JcaSignerInfoGeneratorBuilder}\OperatorTok{;}
\KeywordTok{import} \ImportTok{org}\OperatorTok{.}\ImportTok{bouncycastle}\OperatorTok{.}\ImportTok{jce}\OperatorTok{.}\ImportTok{provider}\OperatorTok{.}\ImportTok{BouncyCastleProvider}\OperatorTok{;}
\KeywordTok{import} \ImportTok{org}\OperatorTok{.}\ImportTok{bouncycastle}\OperatorTok{.}\ImportTok{operator}\OperatorTok{.}\ImportTok{ContentSigner}\OperatorTok{;}
\KeywordTok{import} \ImportTok{org}\OperatorTok{.}\ImportTok{bouncycastle}\OperatorTok{.}\ImportTok{operator}\OperatorTok{.}\ImportTok{jcajce}\OperatorTok{.}\ImportTok{JcaContentSignerBuilder}\OperatorTok{;}
\KeywordTok{import} \ImportTok{org}\OperatorTok{.}\ImportTok{bouncycastle}\OperatorTok{.}\ImportTok{operator}\OperatorTok{.}\ImportTok{jcajce}\OperatorTok{.}\ImportTok{JcaDigestCalculatorProviderBuilder}\OperatorTok{;}

\CommentTok{/**}
 \CommentTok{*}\NormalTok{ Build a detached CMS SignedData suitable for embedding}
 \CommentTok{*}\NormalTok{ in a PDF signature}\CommentTok{\textquotesingle{}}\NormalTok{s }\CommentTok{/}\NormalTok{Contents field}\CommentTok{.}
 \CommentTok{*/}
\KeywordTok{public} \KeywordTok{class}\NormalTok{ CmsBuildingExample }\OperatorTok{\{}

    \DataTypeTok{static} \OperatorTok{\{}
        \BuiltInTok{Security}\OperatorTok{.}\FunctionTok{addProvider}\OperatorTok{(}\KeywordTok{new} \FunctionTok{BouncyCastleProvider}\OperatorTok{());}
    \OperatorTok{\}}

    \CommentTok{/**}
     \CommentTok{*}\NormalTok{ Create a detached CMS SignedData over the given byte range content}\CommentTok{.}
     \CommentTok{*}
     \CommentTok{*} \AnnotationTok{@param}\CommentTok{ }\AnnotationTok{byteRangeContent}\CommentTok{  }\NormalTok{the PDF bytes specified by }\CommentTok{/}\NormalTok{ByteRange}
     \CommentTok{*} \AnnotationTok{@param}\CommentTok{ }\AnnotationTok{privateKey}\CommentTok{        }\NormalTok{the signer}\CommentTok{\textquotesingle{}}\NormalTok{s private key}
     \CommentTok{*} \AnnotationTok{@param}\CommentTok{ }\AnnotationTok{chain}\CommentTok{             }\NormalTok{the certificate chain }\CommentTok{(}\NormalTok{end}\CommentTok{{-}}\NormalTok{entity first}\CommentTok{)}
     \CommentTok{*} \AnnotationTok{@return}\NormalTok{ DER}\CommentTok{{-}}\NormalTok{encoded CMS SignedData bytes}
     \CommentTok{*/}
    \KeywordTok{public} \DataTypeTok{static} \DataTypeTok{byte}\OperatorTok{[]} \FunctionTok{createDetachedSignature}\OperatorTok{(}
            \DataTypeTok{byte}\OperatorTok{[]}\NormalTok{ byteRangeContent}\OperatorTok{,}
            \BuiltInTok{PrivateKey}\NormalTok{ privateKey}\OperatorTok{,}
            \BuiltInTok{Certificate}\OperatorTok{[]}\NormalTok{ chain}\OperatorTok{)} \KeywordTok{throws} \BuiltInTok{Exception} \OperatorTok{\{}

        \BuiltInTok{X509Certificate}\NormalTok{ signingCert }\OperatorTok{=} \OperatorTok{(}\BuiltInTok{X509Certificate}\OperatorTok{)}\NormalTok{ chain}\OperatorTok{[}\DecValTok{0}\OperatorTok{];}

        \CommentTok{// Determine signature algorithm from key type}
        \BuiltInTok{String}\NormalTok{ sigAlg }\OperatorTok{=} \ControlFlowTok{switch} \OperatorTok{(}\NormalTok{privateKey}\OperatorTok{.}\FunctionTok{getAlgorithm}\OperatorTok{())} \OperatorTok{\{}
            \ControlFlowTok{case} \StringTok{"RSA"} \OperatorTok{{-}\textgreater{}} \StringTok{"SHA256withRSA"}\OperatorTok{;}
            \ControlFlowTok{case} \StringTok{"EC"}  \OperatorTok{{-}\textgreater{}} \StringTok{"SHA256withECDSA"}\OperatorTok{;}
            \KeywordTok{default} \OperatorTok{{-}\textgreater{}} \ControlFlowTok{throw} \KeywordTok{new} \BuiltInTok{IllegalArgumentException}\OperatorTok{(}
                    \StringTok{"Unsupported key type: "} \OperatorTok{+}\NormalTok{ privateKey}\OperatorTok{.}\FunctionTok{getAlgorithm}\OperatorTok{());}
        \OperatorTok{\};}

\NormalTok{        ContentSigner signer }\OperatorTok{=} \KeywordTok{new} \FunctionTok{JcaContentSignerBuilder}\OperatorTok{(}\NormalTok{sigAlg}\OperatorTok{)}
                \OperatorTok{.}\FunctionTok{setProvider}\OperatorTok{(}\StringTok{"BC"}\OperatorTok{)}
                \OperatorTok{.}\FunctionTok{build}\OperatorTok{(}\NormalTok{privateKey}\OperatorTok{);}

\NormalTok{        CMSSignedDataGenerator generator }\OperatorTok{=} \KeywordTok{new} \FunctionTok{CMSSignedDataGenerator}\OperatorTok{();}

\NormalTok{        generator}\OperatorTok{.}\FunctionTok{addSignerInfoGenerator}\OperatorTok{(}
                \KeywordTok{new} \FunctionTok{JcaSignerInfoGeneratorBuilder}\OperatorTok{(}
                        \KeywordTok{new} \FunctionTok{JcaDigestCalculatorProviderBuilder}\OperatorTok{()}
                                \OperatorTok{.}\FunctionTok{setProvider}\OperatorTok{(}\StringTok{"BC"}\OperatorTok{)}
                                \OperatorTok{.}\FunctionTok{build}\OperatorTok{()}
                \OperatorTok{).}\FunctionTok{build}\OperatorTok{(}\NormalTok{signer}\OperatorTok{,}\NormalTok{ signingCert}\OperatorTok{)}                                  
        \OperatorTok{);}

        \CommentTok{// Include the full certificate chain}
        \BuiltInTok{List}\OperatorTok{\textless{}}\BuiltInTok{Certificate}\OperatorTok{\textgreater{}}\NormalTok{ certList }\OperatorTok{=} \KeywordTok{new} \BuiltInTok{ArrayList}\OperatorTok{\textless{}\textgreater{}();}
        \ControlFlowTok{for} \OperatorTok{(}\BuiltInTok{Certificate}\NormalTok{ cert }\OperatorTok{:}\NormalTok{ chain}\OperatorTok{)} \OperatorTok{\{}
\NormalTok{            certList}\OperatorTok{.}\FunctionTok{add}\OperatorTok{(}\NormalTok{cert}\OperatorTok{);}
        \OperatorTok{\}}
\NormalTok{        generator}\OperatorTok{.}\FunctionTok{addCertificates}\OperatorTok{(}\KeywordTok{new} \FunctionTok{JcaCertStore}\OperatorTok{(}\NormalTok{certList}\OperatorTok{));}

        \CommentTok{// Generate DETACHED signature (false = do not encapsulate content)}
\NormalTok{        CMSProcessableByteArray content }\OperatorTok{=}
                \KeywordTok{new} \FunctionTok{CMSProcessableByteArray}\OperatorTok{(}\NormalTok{byteRangeContent}\OperatorTok{);}
\NormalTok{        CMSSignedData signedData }\OperatorTok{=}\NormalTok{ generator}\OperatorTok{.}\FunctionTok{generate}\OperatorTok{(}\NormalTok{content}\OperatorTok{,} \KeywordTok{false}\OperatorTok{);}        

        \ControlFlowTok{return}\NormalTok{ signedData}\OperatorTok{.}\FunctionTok{getEncoded}\OperatorTok{();}
    \OperatorTok{\}}
\OperatorTok{\}}
\end{Highlighting}
\end{Shaded}

\begin{itemize}
\item
  \texttt{JcaSignerInfoGeneratorBuilder.build()} automatically adds the
  standard signed attributes (\texttt{contentType},
  \texttt{messageDigest}, \texttt{signingTime}). For CAdES compliance,
  you need to add ESS-signing-certificate-v2 explicitly using a
  \texttt{CMSAttributeTableGenerator}.
\item
  The \texttt{false} parameter is critical: detached signature. The byte
  range content is hashed but not embedded in the CMS structure.
\end{itemize}

\section{Annotated ASN.1 Dump of a Real PDF
Signature}\label{_annotated_asn_1_dump_of_a_real_pdf_signature}

The following is the output of \texttt{openssl\ asn1parse} run against
the DER-decoded \texttt{/Contents} of an actual PDF signature. Each line
is annotated with what it represents in the CMS structure.

\begin{Shaded}
\begin{Highlighting}[]
    \DecValTok{0}\NormalTok{:d=}\DecValTok{0}\NormalTok{  hl=}\DecValTok{4}\NormalTok{ l=}\DecValTok{2417}\NormalTok{ cons: }\DataTypeTok{SEQUENCE}              \CommentTok{{-}{-} ContentInfo}
    \DecValTok{4}\NormalTok{:d=}\DecValTok{1}\NormalTok{  hl=}\DecValTok{2}\NormalTok{ l=   }\DecValTok{9}\NormalTok{ prim: OBJECT   :pkcs7{-}signedData}
                                                     \CommentTok{{-}{-} OID 1.2.840.113549.1.7.2}
   \DecValTok{15}\NormalTok{:d=}\DecValTok{1}\NormalTok{  hl=}\DecValTok{4}\NormalTok{ l=}\DecValTok{2402}\NormalTok{ cons: cont [ }\DecValTok{0}\NormalTok{ ]             }\CommentTok{{-}{-} EXPLICIT [0] content}
   \DecValTok{19}\NormalTok{:d=}\DecValTok{2}\NormalTok{  hl=}\DecValTok{4}\NormalTok{ l=}\DecValTok{2398}\NormalTok{ cons: }\DataTypeTok{SEQUENCE}                \CommentTok{{-}{-} SignedData}
   \DecValTok{23}\NormalTok{:d=}\DecValTok{3}\NormalTok{  hl=}\DecValTok{2}\NormalTok{ l=   }\DecValTok{1}\NormalTok{ prim: }\DataTypeTok{INTEGER}\NormalTok{  :}\DecValTok{01}            \CommentTok{{-}{-} version = 1}
   \DecValTok{26}\NormalTok{:d=}\DecValTok{3}\NormalTok{  hl=}\DecValTok{2}\NormalTok{ l=  }\DecValTok{15}\NormalTok{ cons: }\DataTypeTok{SET}                     \CommentTok{{-}{-} digestAlgorithms}
   \DecValTok{28}\NormalTok{:d=}\DecValTok{4}\NormalTok{  hl=}\DecValTok{2}\NormalTok{ l=  }\DecValTok{13}\NormalTok{ cons: }\DataTypeTok{SEQUENCE}
   \DecValTok{30}\NormalTok{:d=}\DecValTok{5}\NormalTok{  hl=}\DecValTok{2}\NormalTok{ l=   }\DecValTok{9}\NormalTok{ prim: OBJECT   :sha256        }\CommentTok{{-}{-} OID 2.16.840.1.101.3.4.2.1}
   \DecValTok{41}\NormalTok{:d=}\DecValTok{5}\NormalTok{  hl=}\DecValTok{2}\NormalTok{ l=   }\DecValTok{0}\NormalTok{ prim: }\DataTypeTok{NULL}                    \CommentTok{{-}{-} parameters (explicit NULL)}
   \DecValTok{43}\NormalTok{:d=}\DecValTok{3}\NormalTok{  hl=}\DecValTok{2}\NormalTok{ l=  }\DecValTok{11}\NormalTok{ cons: }\DataTypeTok{SEQUENCE}                \CommentTok{{-}{-} encapContentInfo}
   \DecValTok{45}\NormalTok{:d=}\DecValTok{4}\NormalTok{  hl=}\DecValTok{2}\NormalTok{ l=   }\DecValTok{9}\NormalTok{ prim: OBJECT   :pkcs7{-}data    }\CommentTok{{-}{-} eContentType = id{-}data}
                                                     \CommentTok{{-}{-} OID 1.2.840.113549.1.7.1}
                                                     \CommentTok{{-}{-} (eContent ABSENT = detached)}
   \DecValTok{56}\NormalTok{:d=}\DecValTok{3}\NormalTok{  hl=}\DecValTok{4}\NormalTok{ l=}\DecValTok{1735}\NormalTok{ cons: cont [ }\DecValTok{0}\NormalTok{ ]             }\CommentTok{{-}{-} certificates [0] IMPLICIT}
   \DecValTok{60}\NormalTok{:d=}\DecValTok{4}\NormalTok{  hl=}\DecValTok{4}\NormalTok{ l= }\DecValTok{872}\NormalTok{ cons: }\DataTypeTok{SEQUENCE}                \CommentTok{{-}{-} Certificate \#1 (end{-}entity)}
  \DecValTok{936}\NormalTok{:d=}\DecValTok{4}\NormalTok{  hl=}\DecValTok{4}\NormalTok{ l= }\DecValTok{855}\NormalTok{ cons: }\DataTypeTok{SEQUENCE}                \CommentTok{{-}{-} Certificate \#2 (intermediate CA)}
 \DecValTok{1795}\NormalTok{:d=}\DecValTok{3}\NormalTok{  hl=}\DecValTok{4}\NormalTok{ l= }\DecValTok{624}\NormalTok{ cons: }\DataTypeTok{SET}                     \CommentTok{{-}{-} signerInfos}
 \DecValTok{1799}\NormalTok{:d=}\DecValTok{4}\NormalTok{  hl=}\DecValTok{4}\NormalTok{ l= }\DecValTok{620}\NormalTok{ cons: }\DataTypeTok{SEQUENCE}                \CommentTok{{-}{-} SignerInfo}
 \DecValTok{1803}\NormalTok{:d=}\DecValTok{5}\NormalTok{  hl=}\DecValTok{2}\NormalTok{ l=   }\DecValTok{1}\NormalTok{ prim: }\DataTypeTok{INTEGER}\NormalTok{  :}\DecValTok{01}            \CommentTok{{-}{-} version = 1}
 \DecValTok{1806}\NormalTok{:d=}\DecValTok{5}\NormalTok{  hl=}\DecValTok{3}\NormalTok{ l= }\DecValTok{141}\NormalTok{ cons: }\DataTypeTok{SEQUENCE}                \CommentTok{{-}{-} sid = issuerAndSerialNumber}
 \DecValTok{1809}\NormalTok{:d=}\DecValTok{6}\NormalTok{  hl=}\DecValTok{3}\NormalTok{ l= }\DecValTok{118}\NormalTok{ cons: }\DataTypeTok{SEQUENCE}                \CommentTok{{-}{-}   issuer DN}
 \DecValTok{1930}\NormalTok{:d=}\DecValTok{6}\NormalTok{  hl=}\DecValTok{2}\NormalTok{ l=  }\DecValTok{19}\NormalTok{ prim: }\DataTypeTok{INTEGER}\NormalTok{  :...           }\CommentTok{{-}{-}   serialNumber}
 \DecValTok{1951}\NormalTok{:d=}\DecValTok{5}\NormalTok{  hl=}\DecValTok{2}\NormalTok{ l=  }\DecValTok{13}\NormalTok{ cons: }\DataTypeTok{SEQUENCE}                \CommentTok{{-}{-} digestAlgorithm}
 \DecValTok{1953}\NormalTok{:d=}\DecValTok{6}\NormalTok{  hl=}\DecValTok{2}\NormalTok{ l=   }\DecValTok{9}\NormalTok{ prim: OBJECT   :sha256        }\CommentTok{{-}{-}   SHA{-}256}
 \DecValTok{1964}\NormalTok{:d=}\DecValTok{6}\NormalTok{  hl=}\DecValTok{2}\NormalTok{ l=   }\DecValTok{0}\NormalTok{ prim: }\DataTypeTok{NULL}
 \DecValTok{1966}\NormalTok{:d=}\DecValTok{5}\NormalTok{  hl=}\DecValTok{3}\NormalTok{ l= }\DecValTok{189}\NormalTok{ cons: cont [ }\DecValTok{0}\NormalTok{ ]             }\CommentTok{{-}{-} signedAttrs [0] IMPLICIT}
 \DecValTok{1969}\NormalTok{:d=}\DecValTok{6}\NormalTok{  hl=}\DecValTok{2}\NormalTok{ l=  }\DecValTok{24}\NormalTok{ cons: }\DataTypeTok{SEQUENCE}                \CommentTok{{-}{-}   contentType attribute}
 \DecValTok{1995}\NormalTok{:d=}\DecValTok{6}\NormalTok{  hl=}\DecValTok{2}\NormalTok{ l=  }\DecValTok{47}\NormalTok{ cons: }\DataTypeTok{SEQUENCE}                \CommentTok{{-}{-}   messageDigest attribute}
 \DecValTok{2044}\NormalTok{:d=}\DecValTok{6}\NormalTok{  hl=}\DecValTok{2}\NormalTok{ l=  }\DecValTok{28}\NormalTok{ cons: }\DataTypeTok{SEQUENCE}                \CommentTok{{-}{-}   signingTime attribute}
 \DecValTok{2074}\NormalTok{:d=}\DecValTok{6}\NormalTok{  hl=}\DecValTok{2}\NormalTok{ l=  }\DecValTok{79}\NormalTok{ cons: }\DataTypeTok{SEQUENCE}                \CommentTok{{-}{-}   ESS{-}signing{-}certificate{-}v2}
 \DecValTok{2155}\NormalTok{:d=}\DecValTok{5}\NormalTok{  hl=}\DecValTok{2}\NormalTok{ l=  }\DecValTok{13}\NormalTok{ cons: }\DataTypeTok{SEQUENCE}                \CommentTok{{-}{-} signatureAlgorithm}
 \DecValTok{2157}\NormalTok{:d=}\DecValTok{6}\NormalTok{  hl=}\DecValTok{2}\NormalTok{ l=   }\DecValTok{9}\NormalTok{ prim: OBJECT   :sha256WithRSAEncryption}
                                                     \CommentTok{{-}{-} OID 1.2.840.113549.1.1.11}
 \DecValTok{2168}\NormalTok{:d=}\DecValTok{6}\NormalTok{  hl=}\DecValTok{2}\NormalTok{ l=   }\DecValTok{0}\NormalTok{ prim: }\DataTypeTok{NULL}
 \DecValTok{2170}\NormalTok{:d=}\DecValTok{5}\NormalTok{  hl=}\DecValTok{4}\NormalTok{ l= }\DecValTok{256}\NormalTok{ prim: }\DataTypeTok{OCTET} \DataTypeTok{STRING}            \CommentTok{{-}{-} signature (256 bytes = RSA{-}2048)}
 \DecValTok{2430}\NormalTok{:d=}\DecValTok{5}\NormalTok{  hl=}\DecValTok{2}\NormalTok{ l=   }\DecValTok{7}\NormalTok{ cons: cont [ }\DecValTok{1}\NormalTok{ ]             }\CommentTok{{-}{-} unsignedAttrs [1] IMPLICIT}
                                                     \CommentTok{{-}{-} (timestamp token would be here)}
\end{Highlighting}
\end{Shaded}

The offsets in this dump are from the start of the DER-encoded CMS blob,
not from the start of the PDF file. To get the absolute file offset, add
the starting offset of the \texttt{/Contents} hex string (as specified
by \texttt{/ByteRange}) plus the hex-to-binary decoding offset. When
debugging with a hex editor, remember that the \texttt{/Contents} value
is hex-encoded: each byte in the ASN.1 dump corresponds to two hex
characters in the PDF file.

\section{Summary}\label{_summary_6}

CMS \texttt{SignedData} is the envelope that carries every PDF digital
signature. Its lineage stretches from RSA's PKCS\#7 in 1993 through the
S/MIME working group to the current RFC 5652 --- Cryptographic Message
Syntax (CMS) standard. The core design has been stable for over
twenty-five years, which is both its strength (mature, well-understood,
universally supported) and its burden (legacy quirks from 1993 are still
with us).

The key points for PDF signature work:

\begin{itemize}
\item
  \textbf{ContentInfo wraps everything.} The OID
  \texttt{1.2.840.113549.1.7.2} identifies the CMS as
  \texttt{SignedData}. Learn to recognize its DER encoding
  (\texttt{06\ 09\ 2A\ 86\ 48\ 86\ F7\ 0D\ 01\ 07\ 02}) in hex dumps.
\item
  \textbf{PDF signatures are detached.} The \texttt{eContent} in
  \texttt{EncapsulatedContentInfo} is absent. The signed content lives
  in the PDF file, addressed by \texttt{/ByteRange}.
\item
  \textbf{Signed attributes are signed, not the content.} The signature
  covers the DER-encoded signed attributes, which include the
  \texttt{messageDigest} (hash of the content). This two-level
  indirection is essential to understand.
\item
  \textbf{The DER re-tagging trap is real.} The signed attributes must
  be re-tagged from \texttt{A0} to \texttt{31} before hashing for
  signature computation. BouncyCastle handles this, but any lower-level
  implementation must get this right.
\item
  \textbf{ESS-signing-certificate-v2 is the CAdES requirement.} Without
  it, the signature is CMS-valid but not CAdES-valid and therefore not
  PAdES-compliant.
\item
  \textbf{Unsigned attributes carry timestamps.} The
  \texttt{signature-time-stamp} attribute contains a nested CMS
  structure --- a timestamp token from an RFC 3161 TSA.
\item
  \textbf{Legacy PKCS\#7 differences still cause interop failures.} BER
  vs DER encoding of signed attributes is the most common culprit.
\end{itemize}

The next part of this book turns to PDF internals --- the file
structure, incremental updates, and form fields that the CMS
\texttt{SignedData} is embedded within. Understanding the container is
meaningless without understanding the document it lives in.

\chapter{PDF Internals}\label{_pdf_internals}

\chapter{PDF File Structure}\label{ch07}

Every PDF file on earth --- whether it is a single-page receipt from
your corner bakery or a 4,000-page securities filing with twelve nested
signatures --- is built from the same four structural components,
arranged in the same order, governed by the same rules that John Warnock
and Chuck Geschke laid down at Adobe in 1993. If you want to understand
why a signature's \texttt{/ByteRange} array contains raw file offsets,
why cross-reference tables use fixed-width 20-byte entries, or why some
PDFs have three \texttt{\%\%EOF} markers, you have to understand this
structure at the byte level.

This chapter dismantles a PDF file from the first byte to the last. We
will trace headers, decode every object type the specification defines,
read cross-reference tables in both the classic ASCII format and the
newer compressed stream format, examine the trailer dictionary, walk
through a minimal PDF byte by byte, and then connect all of it back to
the one thing we care about most: digital signatures.

\section{The Four Sections}\label{_the_four_sections}

ISO 32000-2:2020 Section 7.5 defines the layout. Every conforming PDF
file consists of these four parts, in this order:

\begin{enumerate}
\def\labelenumi{\arabic{enumi}.}
\item
  \textbf{Header} --- a single line declaring the PDF version, plus an
  optional binary comment
\item
  \textbf{Body} --- the collection of indirect objects that define
  everything in the document: pages, fonts, images, form fields,
  annotations, and signatures
\item
  \textbf{Cross-reference table} --- an index that maps object numbers
  to byte offsets in the file, enabling random access without linear
  scanning
\item
  \textbf{Trailer} --- a dictionary that tells the parser where to find
  the cross-reference table, which object is the document root, and (for
  incremental updates) where to find the previous cross-reference
  section
\end{enumerate}

This four-part structure has been stable since PDF 1.0. PDF 1.5
introduced an alternative binary representation for the cross-reference
data (cross-reference streams), and PDF 2.0 deprecated some legacy
metadata keys, but the fundamental anatomy has not changed in over three
decades.

\subsection{The Header}\label{_the_header}

A PDF file begins with a version identifier on the first line:

\begin{verbatim}
%PDF-1.7
\end{verbatim}

That is nine bytes: the five characters \texttt{\%PDF-}, the version
number, and a line ending. The \texttt{\%} character is the PDF comment
delimiter, which means the header line is technically a comment --- a
design choice that lets non-PDF-aware tools skip it without error.

The version in the header is the \emph{file} version. ISO 32000-2:2020
Section 7.5.2 specifies that a \texttt{/Version} key in the document
catalog may override this value. In practice, many generators write
\texttt{\%PDF-1.7} in the header and then set \texttt{/Version\ /2.0} in
the catalog when they need PDF 2.0 features. This dual-version mechanism
exists for backward compatibility: an older parser that does not
understand the catalog \texttt{/Version} will still see a version number
it can process (or at least attempt to process) in the header.

\subsubsection{The Binary Comment
Trick}\label{_the_binary_comment_trick}

Immediately after the header line, you will almost always find a second
comment line containing four or more bytes with values above 127 (0x80):

\begin{verbatim}
%âãÏÓ
\end{verbatim}

This is not decoration. Its purpose is purely practical: it forces
file-type detection heuristics --- in web servers, email gateways, FTP
clients, operating systems --- to classify the file as binary rather
than ASCII text. Without this line, a transfer in ASCII mode (remember
the FTP \texttt{ASCII} vs \texttt{BINARY} setting?) could mangle the
binary stream data inside the PDF. ISO 32000-2:2020 Section 7.5.2
recommends but does not require this binary hint.

You will occasionally encounter PDFs without the binary
comment --- usually from minimalist generators or test harnesses. The
file is still valid, but it is more vulnerable to transfer corruption.

When debugging PDF files at the byte level, the header is your first
sanity check. If \texttt{\%PDF-} is not the first five bytes, something
has happened upstream: a Unicode BOM was prepended, an HTML wrapper was
added, a proxy rewrote the response body, or the file was opened as text
and line endings were converted. We have encountered every one of these
in production systems.

\subsection{The Body}\label{_the_body}

The body contains all the PDF objects that constitute the document. Page
trees, content streams, font dictionaries, image data, form fields,
metadata, and --- critically for us --- signature dictionaries and their
CMS payloads.

Objects in the body are written as \textbf{indirect objects}, each
wrapped in a numbered envelope:

\begin{verbatim}
7 0 obj
<< /Type /Page /MediaBox [0 0 612 792] /Contents 12 0 R >>
endobj
\end{verbatim}

The prefix \texttt{7\ 0\ obj} means "object number 7, generation number
0." The keyword \texttt{endobj} marks the end of the object. Between
them is the object's value --- here, a dictionary describing a page.

Objects may appear in any order. There is no requirement that object 1
precedes object 2, or that page objects come before font objects. The
cross-reference table handles the mapping from logical object number to
physical byte offset. This is a fundamental design principle:
\textbf{the physical ordering of objects in the file has no bearing on
the logical structure of the document}.

Different PDF generators produce wildly different orderings. Adobe
Acrobat tends to cluster related objects together. Apache PDFBox writes
objects in the order they were allocated in memory. Linearized PDFs (the
"fast web view" format described in ISO 32000-2:2020 Section 7.6)
deliberately reorder objects so that the first page can be displayed
before the entire file is downloaded. None of this changes the logical
content --- but it changes the byte offsets, which in turn changes the
byte range that a signature covers.

Some tools assume objects appear in numerical order. They do not. If you
are writing a PDF parser or debugging one, never assume object N appears
before object N+1 in the file. Always consult the cross-reference table.

\subsection{The Cross-Reference Table}\label{_the_cross_reference_table}

After the body comes the cross-reference table (commonly called the
"xref table") --- an index that records the byte offset where each
object starts in the file.

\begin{verbatim}
xref
0 8
0000000000 65535 f
0000000015 00000 n
0000000078 00000 n
0000000178 00000 n
0000000331 00000 n
0000000409 00000 n
0000000553 00000 n
0000000715 00000 n
\end{verbatim}

The keyword \texttt{xref} marks the beginning. The line \texttt{0\ 8}
means "this section starts at object number 0 and contains 8 entries."
Each subsequent line is exactly 20 bytes (including the line terminator)
and contains three fields:

\begin{itemize}
\item
  \textbf{Byte offset} (10 digits, zero-padded) --- the position in the
  file where the object starts
\item
  \textbf{Generation number} (5 digits, zero-padded) --- used for object
  versioning; almost always \texttt{00000}
\item
  \textbf{In-use flag} --- \texttt{n} means the object is in use;
  \texttt{f} means it is free (deleted)
\end{itemize}

The very first entry (object 0) is always free, with generation number
65535. It serves as the head of a linked list of free objects. Each free
entry's "byte offset" field actually contains the object number of the
\emph{next} free entry, forming a chain. This mechanism allows PDF
writers to reuse deleted object numbers --- a feature that mattered when
PDFs were small and object numbers precious, and that is now mostly
vestigial.

The 20-byte-per-entry format is non-negotiable. This rigidity is
deliberate: it enables random access to any entry via arithmetic
(\texttt{entry\_offset\ =\ xref\_start\ +\ header\_bytes\ +\ object\_number\ *\ 20})
without scanning the entire table. This random-access property is what
allows a viewer to display a single page from a 10,000-page document
without loading the whole file into memory.

The generation number increments when an object is deleted and its
number is later reused, preventing stale references (\texttt{7\ 0\ R})
from resolving to the wrong content. In practice, nearly every object
you encounter in the wild has generation 0. The mechanism exists, the
specification defines it carefully, and production PDFs almost never
exercise it.

\subsection{The Trailer}\label{_the_trailer}

Immediately after the cross-reference table comes the trailer:

\begin{verbatim}
trailer
<<
  /Size 8
  /Root 1 0 R
  /Info 6 0 R
  /ID [<a9f2c74b...> <a9f2c74b...>]
>>
startxref
4892
%%EOF
\end{verbatim}

The trailer dictionary is the most important dictionary in the file,
because \emph{everything} starts here. A PDF parser begins at the
\textbf{end} of the file. It searches backward for \texttt{\%\%EOF},
then backward for \texttt{startxref}, reads the byte offset (4892), and
jumps to that offset to find the cross-reference table. From the xref
table, it reads the trailer dictionary. From the trailer dictionary's
\texttt{/Root} key, it finds the document catalog. From the catalog, it
can navigate to every object in the document.

\textbf{Parsing a PDF starts from the end, not the beginning.} This
backward-search strategy is fundamental to how incremental updates work
(Chapter 8), and it is the reason a PDF can be extended by appending
data without rewriting the beginning of the file.

\begin{longtable}[]{@{}
  >{\raggedright\arraybackslash}p{(\linewidth - 4\tabcolsep) * \real{0.2000}}
  >{\raggedright\arraybackslash}p{(\linewidth - 4\tabcolsep) * \real{0.2000}}
  >{\raggedright\arraybackslash}p{(\linewidth - 4\tabcolsep) * \real{0.6000}}@{}}
\caption{Trailer Dictionary Keys --- ISO 32000-2:2020 Section 7.5.5,
Table 15}\tabularnewline
\toprule\noalign{}
\begin{minipage}[b]{\linewidth}\raggedright
Key
\end{minipage} & \begin{minipage}[b]{\linewidth}\raggedright
Type
\end{minipage} & \begin{minipage}[b]{\linewidth}\raggedright
Description
\end{minipage} \\
\midrule\noalign{}
\endfirsthead
\toprule\noalign{}
\begin{minipage}[b]{\linewidth}\raggedright
Key
\end{minipage} & \begin{minipage}[b]{\linewidth}\raggedright
Type
\end{minipage} & \begin{minipage}[b]{\linewidth}\raggedright
Description
\end{minipage} \\
\midrule\noalign{}
\endhead
\bottomrule\noalign{}
\endlastfoot
\texttt{/Size} & integer & Total number of entries in the
cross-reference table (including object 0). One greater than the highest
object number used. \\
\texttt{/Root} & indirect reference & The document catalog --- root of
the entire logical document tree. \\
\texttt{/Info} & indirect reference & Document information dictionary.
Optional and deprecated in PDF 2.0 (replaced by XMP metadata), but still
universally present in real-world files. \\
\texttt{/ID} & array of two byte strings & File identifiers. The first
element is assigned at creation and never changes; the second is updated
on each save. Signatures depend on \texttt{/ID} --- PAdES profiles
strongly recommend it. \\
\texttt{/Prev} & integer & Byte offset of the \emph{previous}
cross-reference section. This is how incremental updates chain together
(Chapter 8). \\
\texttt{/Encrypt} & dictionary or reference & Encryption dictionary. The
signature value itself (\texttt{/Contents}) is never encrypted even in
encrypted PDFs, but the surrounding content may be. \\
\end{longtable}

The \texttt{/ID} array is not cosmetic. Some signing tools use the first
element as input to seed value dictionaries. Some validators check
\texttt{/ID} consistency across incremental revisions. If a tool strips
or regenerates \texttt{/ID} during processing, existing signatures may
fail validation. We have seen this happen with document management
systems that "normalize" PDFs on upload.

\section{PDF Object Types}\label{_pdf_object_types}

Every value in a PDF file is one of exactly ten types, defined in ISO
32000-2:2020 Section 7.3. Understanding them is necessary for reading
raw PDF structure, debugging signatures at the COS level, and working
with PDFBox's object model.

\subsection{Boolean, Integer, Real,
Null}\label{_boolean_integer_real_null}

\textbf{Boolean} --- the keywords \texttt{true} or \texttt{false}. Used
in flags like \texttt{/NeedAppearances\ true}.

\textbf{Integer} --- a signed whole number with no decimal point:
\texttt{42}, \texttt{-7}, \texttt{0}. The four values in a
\texttt{/ByteRange} array are integers. Integers in PDF have no fixed
bit width; implementations typically use 64-bit storage.

\textbf{Real} --- a signed number with a decimal point: \texttt{3.14},
\texttt{-0.001}, \texttt{72.0}. Used in coordinate values,
transformation matrices, and font metrics.

\textbf{Null} --- the keyword \texttt{null}. A dictionary key mapped to
\texttt{null} is equivalent to the key being absent. Null references are
resolved as "object not found."

\subsection{String (Literal and
Hexadecimal)}\label{_string_literal_and_hexadecimal}

PDF has two string representations:

\textbf{Literal string} --- parentheses-delimited:
\texttt{(Hello\ World)}. Backslash escapes handle special characters:
\texttt{\textbackslash{}(}, \texttt{\textbackslash{})},
\texttt{\textbackslash{}\textbackslash{}}, \texttt{\textbackslash{}n},
\texttt{\textbackslash{}r}, \texttt{\textbackslash{}t}, and octal
escapes like \texttt{\textbackslash{}101} (the letter A). Unbalanced
parentheses must be escaped; balanced ones need not be.

\textbf{Hexadecimal string} --- angle-bracket-delimited hex pairs:
\texttt{\textless{}48656C6C6F\textgreater{}}. Each pair of hex digits
represents one byte. The \texttt{/Contents} value in a signature
dictionary is always a hex string because it contains raw binary CMS
data. The hex string is right-padded with zeros to fill the space
reserved during signing.

\subsection{Name}\label{_name}

A slash-prefixed identifier: \texttt{/Type}, \texttt{/Page},
\texttt{/Sig}, \texttt{/ByteRange}. Names are atoms, not
strings --- they serve as dictionary keys and enumeration values. Names
containing characters outside the printable ASCII range use a
\texttt{\#xx} hex escape: \texttt{/F\#23oo} represents the name
\texttt{F\#oo}.

\subsection{Array}\label{_array}

An ordered collection enclosed in square brackets:
\texttt{{[}0\ 14000\ 30408\ 14592{]}}. Arrays can contain any mix of
types, including other arrays and dictionaries. The \texttt{/ByteRange}
value is an array of four integers.

\subsection{Dictionary}\label{_dictionary}

Key-value pairs enclosed in double angle brackets:

\begin{verbatim}
<< /Type /Sig /Filter /Adobe.PPKLite /SubFilter /adbe.pkcs7.detached >>
\end{verbatim}

Keys are always names. Values can be any type. Dictionaries are the
workhorse of PDF --- every page, every font, every annotation, every
signature is a dictionary or contains one.

\subsection{Stream}\label{_stream}

A dictionary followed by raw binary data, delimited by \texttt{stream}
and \texttt{endstream}:

\begin{verbatim}
8 0 obj
<< /Type /XObject /Subtype /Image /Length 12345 /Filter /FlateDecode >>
stream
...binary data (12345 bytes, possibly compressed)...
endstream
endobj
\end{verbatim}

Streams are the only mechanism for storing bulk data: page drawing
instructions, images, embedded font programs, cross-reference data, and
CMap tables. The \texttt{/Length} key must be exact --- if it is wrong,
the parser reads too few or too many bytes, corrupting everything that
follows.

A wrong \texttt{/Length} is one of the most common forms of PDF
corruption from buggy generators. When debugging a corrupt PDF, always
compare the \texttt{/Length} value against the actual byte count between
\texttt{stream} and \texttt{endstream}. QPDF's \texttt{-\/-check} flag
catches this automatically.

\subsection{Indirect Object and Indirect
Reference}\label{_indirect_object_and_indirect_reference}

An \textbf{indirect object} wraps any value with an object number and
generation number: \texttt{12\ 0\ obj\ …\hspace{0pt}\ endobj}. An
\textbf{indirect reference} is a pointer to such an object:
\texttt{12\ 0\ R}. The cross-reference table resolves these: given
\texttt{12\ 0\ R}, the parser looks up object 12, generation 0, in the
xref table, seeks to the recorded byte offset, and reads the object.

Indirect references are the glue that holds a PDF together. The catalog
references the page tree, the page tree references individual pages,
pages reference content streams and annotations, annotations reference
appearance streams --- all through indirect references.

\section{Cross-Reference Streams}\label{ch07-xref-streams}

PDF 1.5 introduced \textbf{cross-reference streams} --- an alternative
to the classic ASCII xref table that stores the same data in a
compressed binary stream object.

\subsection{Why They Exist}\label{_why_they_exist}

The classic xref table has two problems. First, the 20-bytes-per-entry
format is wasteful: a PDF with 100,000 objects needs a 2 MB xref table
of pure ASCII text. Cross-reference streams compress this data 5-10x
using standard stream filters (typically FlateDecode). Second, the
classic format has no entry type for objects stored inside object
streams (covered next); cross-reference streams introduce a Type 2 entry
for exactly this purpose.

\subsection{Structure}\label{_structure}

\begin{verbatim}
42 0 obj
<<
  /Type /XRef
  /Size 100
  /W [1 3 1]
  /Index [0 100]
  /Root 1 0 R
  /Info 5 0 R
  /ID [<abc123> <abc123>]
  /Filter /FlateDecode
  /Length 487
>>
stream
...compressed binary data...
endstream
endobj
\end{verbatim}

The stream dictionary serves double duty: it contains both the stream
parameters (\texttt{/Filter}, \texttt{/Length}) and the trailer
parameters (\texttt{/Root}, \texttt{/Info}, \texttt{/ID},
\texttt{/Prev}, \texttt{/Size}). There is no separate \texttt{trailer}
keyword when cross-reference streams are used.

The \texttt{/W} array specifies field widths in bytes for each entry.
\texttt{/W\ {[}1\ 3\ 1{]}} means: 1 byte for the type field, 3 bytes for
field 2, 1 byte for field 3 --- 5 bytes total per entry.

\begin{longtable}[]{@{}
  >{\raggedright\arraybackslash}p{(\linewidth - 4\tabcolsep) * \real{0.1111}}
  >{\raggedright\arraybackslash}p{(\linewidth - 4\tabcolsep) * \real{0.4444}}
  >{\raggedright\arraybackslash}p{(\linewidth - 4\tabcolsep) * \real{0.4444}}@{}}
\caption{Cross-Reference Stream Entry Types --- ISO 32000-2:2020 Section
7.5.8, Table 18}\tabularnewline
\toprule\noalign{}
\begin{minipage}[b]{\linewidth}\raggedright
Type
\end{minipage} & \begin{minipage}[b]{\linewidth}\raggedright
Field 2 Meaning
\end{minipage} & \begin{minipage}[b]{\linewidth}\raggedright
Field 3 Meaning
\end{minipage} \\
\midrule\noalign{}
\endfirsthead
\toprule\noalign{}
\begin{minipage}[b]{\linewidth}\raggedright
Type
\end{minipage} & \begin{minipage}[b]{\linewidth}\raggedright
Field 2 Meaning
\end{minipage} & \begin{minipage}[b]{\linewidth}\raggedright
Field 3 Meaning
\end{minipage} \\
\midrule\noalign{}
\endhead
\bottomrule\noalign{}
\endlastfoot
0 & Next free object number & Generation number of the free entry \\
1 & Byte offset of the object in the file & Generation number \\
2 & Object number of the containing object stream & Index of this object
within that object stream \\
\end{longtable}

Type 2 is the key innovation. It says "this object is not a standalone
indirect object in the file body; it is packed inside object stream N,
at position M." Classic xref tables cannot express this.

The \texttt{/Index} array specifies which object numbers are covered.
\texttt{/Index\ {[}0\ 100{]}} means entries for objects 0 through 99. A
cross-reference stream can have multiple subsections:
\texttt{/Index\ {[}0\ 50\ 200\ 10{]}} covers objects 0-49 and 200-209.

Some PDF libraries produce files that start with a classic xref table
but switch to cross-reference streams in incremental updates (or vice
versa). This is legal per the specification, but it trips up parsers
that expect consistency. PDFBox handles the transition transparently;
not all tools do.

\section{Object Streams}\label{_object_streams}

PDF 1.5 also introduced \textbf{object streams}: multiple indirect
objects packed into a single compressed stream, eliminating the repeated
\texttt{N\ 0\ obj}\ldots\hspace{0pt}\texttt{endobj} framing overhead.

\begin{verbatim}
10 0 obj
<< /Type /ObjStm /N 5 /First 32 /Filter /FlateDecode /Length 418 >>
stream
...compressed data containing 5 objects...
endstream
endobj
\end{verbatim}

The \texttt{/N} key gives the count of packed objects; \texttt{/First}
gives the byte offset within the decompressed data where the first
object's value begins (the preceding bytes contain a table of object
numbers and offsets).

Not everything can live in an object stream. ISO 32000-2:2020 Section
7.5.7 excludes: stream objects (they have their own
\texttt{stream}/\texttt{endstream} framing), the encryption dictionary,
objects with non-zero generation numbers, and the cross-reference stream
itself. Signature dictionaries \emph{can} technically live in object
streams (they are dictionaries, not streams), but signing tools almost
never place them there because it would complicate byte-range
calculation.

If you need to find a signature dictionary in a raw PDF, do not search
for \texttt{/Type\ /Sig} in the file bytes --- it might be compressed
inside an object stream. Use \texttt{qpdf\ -\/-show-object=N} where N is
the object number. QPDF decompresses object streams transparently.

\section{Hex Dump Walkthrough: A Minimal Valid
PDF}\label{_hex_dump_walkthrough_a_minimal_valid_pdf}

Let us trace through a complete, minimal PDF file, annotating every
section:

\begin{verbatim}
%PDF-1.7                                          
%âãÏÓ                                             

1 0 obj                                            
<< /Type /Catalog /Pages 2 0 R >>
endobj

2 0 obj                                            
<< /Type /Pages /Kids [3 0 R] /Count 1 >>
endobj

3 0 obj                                            
<< /Type /Page /Parent 2 0 R /MediaBox [0 0 612 792]
   /Contents 4 0 R
   /Resources << /Font << /F1 5 0 R >> >> >>
endobj

4 0 obj                                            
<< /Length 44 >>
stream
BT /F1 24 Tf 100 700 Td (Hello) Tj ET
endstream
endobj

5 0 obj                                            
<< /Type /Font /Subtype /Type1 /BaseFont /Helvetica >>
endobj

xref                                               
0 6
0000000000 65535 f
0000000015 00000 n
0000000066 00000 n
0000000121 00000 n
0000000286 00000 n
0000000384 00000 n

trailer                                            
<< /Size 6 /Root 1 0 R >>
startxref                                          
453
%%EOF                                              
\end{verbatim}

\begin{itemize}
\item
  Header: PDF version 1.7
\item
  Binary hint: four bytes above 0x80, forcing binary classification
\item
  Catalog: root of the document tree; \texttt{/Pages} points to the page
  tree root
\item
  Pages node: page tree root with one child page
\item
  Page: defines the media size (US Letter), references a content stream
  and font resources
\item
  Content stream: drawing instructions --- Begin Text, set font F1 at
  24pt, move to (100,700), show "Hello", End Text
\item
  Font: a reference to standard Helvetica (not embedded --- viewers
  provide a substitute)
\item
  Cross-reference table: six entries mapping object numbers to byte
  offsets
\item
  Trailer: \texttt{/Size\ 6} (six xref entries), \texttt{/Root} points
  to the catalog
\item
  \texttt{startxref}: byte offset of the \texttt{xref} keyword
\item
  \texttt{\%\%EOF}: end-of-file marker
\end{itemize}

A parser reads this \emph{backward}: find \texttt{\%\%EOF}, find
\texttt{startxref}, jump to byte 453, parse the xref table, read the
trailer, follow \texttt{/Root} to object 1, and navigate from there.

This minimal PDF omits a \texttt{/Info} dictionary, \texttt{/ID} array,
and font embedding. Production PDFs have dozens to hundreds of objects,
but the structural skeleton is identical.

\section{How Different Generators Produce Different
Structures}\label{_how_different_generators_produce_different_structures}

Not all PDFs are created equal at the byte level, and the differences
matter when you are computing byte ranges for signatures.

\textbf{Adobe Acrobat} produces linearized files by default, with
objects reordered for progressive display and a hint table at the
beginning. It uses cross-reference streams in modern versions.

\textbf{Apache PDFBox} (3.0.0) writes objects in memory-allocation
order, uses classic xref tables by default, and does not linearize. Its
incremental save mode (\texttt{saveIncremental()}) produces clean
append-only updates.

\textbf{iText / OpenPDF} frequently uses cross-reference streams even
for PDF 1.7 files, and its object ordering varies by version.

\textbf{LibreOffice} produces classic xref tables with object ordering
that appears semi-random (it reflects internal processing order).

\textbf{Microsoft Word} (Save as PDF) produces well-structured files
with cross-reference streams and object streams, heavily compressed.

\textbf{QPDF} can transform between any of these representations:
linearize, de-linearize, convert xref tables to streams or vice versa,
unpack object streams, and normalize whitespace. It is the Swiss Army
knife of PDF structural debugging.

When comparing two versions of a signed PDF at the structural level:

\begin{Shaded}
\begin{Highlighting}[]
\CommentTok{\# Show all xref entries and their origins}
\ExtensionTok{qpdf} \AttributeTok{{-}{-}show{-}xref}\NormalTok{ document.pdf}

\CommentTok{\# Convert to JSON for semantic comparison}
\ExtensionTok{qpdf} \AttributeTok{{-}{-}json}\NormalTok{ document{-}v1.pdf }\OperatorTok{\textgreater{}}\NormalTok{ /tmp/v1.json}
\ExtensionTok{qpdf} \AttributeTok{{-}{-}json}\NormalTok{ document{-}v2.pdf }\OperatorTok{\textgreater{}}\NormalTok{ /tmp/v2.json}
\FunctionTok{diff}\NormalTok{ /tmp/v1.json /tmp/v2.json}

\CommentTok{\# Decompress everything for human{-}readable byte{-}level comparison}
\ExtensionTok{qpdf} \AttributeTok{{-}{-}qdf} \AttributeTok{{-}{-}object{-}streams}\OperatorTok{=}\NormalTok{disable input.pdf readable.pdf}
\end{Highlighting}
\end{Shaded}

\section{Why File Structure Matters for
Signatures}\label{_why_file_structure_matters_for_signatures}

We have spent this chapter on raw PDF anatomy for a specific reason:
\textbf{PDF digital signatures operate at the byte level}, not the
logical-object level.

The \texttt{/ByteRange} array in a signature dictionary specifies
exactly which bytes in the physical file are hashed and signed.
Verification reads those exact bytes, computes a cryptographic hash, and
compares it against the hash embedded in the CMS \texttt{SignedData}
structure. If a single byte differs --- because an object was
renumbered, a cross-reference entry was reformatted, whitespace was
normalized, or a stream was recompressed --- the hash will not match,
and the signature is invalid.

This has concrete consequences:

\textbf{Re-saving breaks signatures.} Loading a signed PDF into memory
and writing it back out --- even without changing any logical
content --- rewrites the file from scratch: new object ordering, new
byte offsets, new xref entries. The bytes change. The signature dies.
Signed PDFs must be saved incrementally (Chapter 8).

\textbf{Linearization affects byte layout.} Linearizing a PDF after
signing rearranges objects for progressive display. The bytes move. The
signature breaks. Always linearize \emph{before} signing, never after.

\textbf{Cross-reference format affects layout.} Whether a file uses
classic xref tables or xref streams changes the physical bytes in the
file. A signature computed against one format will not validate if the
format changes, even when the logical content is identical.

\textbf{Object streams affect findability.} If a signature dictionary is
packed inside an object stream, byte-level tools that search for
\texttt{/Type\ /Sig} in the raw file will miss it. Always use proper PDF
parsing, never raw byte scanning, when searching for signature-related
objects.

\section{PDFBox COS Layer: Mapping PDF Types to
Java}\label{_pdfbox_cos_layer_mapping_pdf_types_to_java}

PDFBox's COS (Carousel Object System) layer provides a one-to-one
mapping from PDF types to Java classes:

\begin{longtable}[]{@{}
  >{\raggedright\arraybackslash}p{(\linewidth - 4\tabcolsep) * \real{0.2222}}
  >{\raggedright\arraybackslash}p{(\linewidth - 4\tabcolsep) * \real{0.3333}}
  >{\raggedright\arraybackslash}p{(\linewidth - 4\tabcolsep) * \real{0.4444}}@{}}
\caption{PDFBox COS Type Mapping}\tabularnewline
\toprule\noalign{}
\begin{minipage}[b]{\linewidth}\raggedright
PDF Type
\end{minipage} & \begin{minipage}[b]{\linewidth}\raggedright
PDFBox Class
\end{minipage} & \begin{minipage}[b]{\linewidth}\raggedright
Notes
\end{minipage} \\
\midrule\noalign{}
\endfirsthead
\toprule\noalign{}
\begin{minipage}[b]{\linewidth}\raggedright
PDF Type
\end{minipage} & \begin{minipage}[b]{\linewidth}\raggedright
PDFBox Class
\end{minipage} & \begin{minipage}[b]{\linewidth}\raggedright
Notes
\end{minipage} \\
\midrule\noalign{}
\endhead
\bottomrule\noalign{}
\endlastfoot
Boolean & \texttt{COSBoolean} & Singleton
\texttt{TRUE}/\texttt{FALSE} \\
Integer & \texttt{COSInteger} & Wraps \texttt{long}; common values
cached \\
Real & \texttt{COSFloat} & Wraps \texttt{float} \\
String & \texttt{COSString} & Both literal and hex;
\texttt{setForceHexForm(true)} for hex output \\
Name & \texttt{COSName} & Interned: \texttt{COSName.TYPE},
\texttt{COSName.SIG}, custom via
\texttt{COSName.getPDFName("ByteRange")} \\
Null & \texttt{COSNull} & Singleton \texttt{COSNull.NULL} \\
Array & \texttt{COSArray} & Ordered list of \texttt{COSBase} \\
Dictionary & \texttt{COSDictionary} & \texttt{getItem()} returns raw
value; \texttt{getDictionaryObject()} dereferences indirect refs \\
Stream & \texttt{COSStream} & Extends \texttt{COSDictionary};
\texttt{createInputStream()} for decoded data \\
Indirect ref & \texttt{COSObject} & Wraps \texttt{COSBase} with object
number and generation \\
\end{longtable}

The distinction between \texttt{getItem()} and
\texttt{getDictionaryObject()} is critical when working with signatures.
\texttt{getItem()} returns the raw value --- if the entry is an indirect
reference, you get the \texttt{COSObject} wrapper (which carries the
object number). \texttt{getDictionaryObject()} follows the reference and
returns the target dictionary. Use \texttt{getDictionaryObject()} when
you need the content; use \texttt{getItem()} when you need the object
number (e.g., to correlate with xref byte offsets).

Here is how to traverse the file structure from trailer to signature
dictionaries using PDFBox 3.0.0:

\begin{Shaded}
\begin{Highlighting}[]
\KeywordTok{import} \ImportTok{org}\OperatorTok{.}\ImportTok{apache}\OperatorTok{.}\ImportTok{pdfbox}\OperatorTok{.}\ImportTok{Loader}\OperatorTok{;}
\KeywordTok{import} \ImportTok{org}\OperatorTok{.}\ImportTok{apache}\OperatorTok{.}\ImportTok{pdfbox}\OperatorTok{.}\ImportTok{cos}\OperatorTok{.}\ImportTok{COSArray}\OperatorTok{;}
\KeywordTok{import} \ImportTok{org}\OperatorTok{.}\ImportTok{apache}\OperatorTok{.}\ImportTok{pdfbox}\OperatorTok{.}\ImportTok{cos}\OperatorTok{.}\ImportTok{COSDictionary}\OperatorTok{;}
\KeywordTok{import} \ImportTok{org}\OperatorTok{.}\ImportTok{apache}\OperatorTok{.}\ImportTok{pdfbox}\OperatorTok{.}\ImportTok{cos}\OperatorTok{.}\ImportTok{COSName}\OperatorTok{;}
\KeywordTok{import} \ImportTok{org}\OperatorTok{.}\ImportTok{apache}\OperatorTok{.}\ImportTok{pdfbox}\OperatorTok{.}\ImportTok{cos}\OperatorTok{.}\ImportTok{COSString}\OperatorTok{;}
\KeywordTok{import} \ImportTok{org}\OperatorTok{.}\ImportTok{apache}\OperatorTok{.}\ImportTok{pdfbox}\OperatorTok{.}\ImportTok{pdmodel}\OperatorTok{.}\ImportTok{PDDocument}\OperatorTok{;}
\KeywordTok{import} \ImportTok{org}\OperatorTok{.}\ImportTok{apache}\OperatorTok{.}\ImportTok{pdfbox}\OperatorTok{.}\ImportTok{pdmodel}\OperatorTok{.}\ImportTok{interactive}\OperatorTok{.}\ImportTok{form}\OperatorTok{.}\ImportTok{PDAcroForm}\OperatorTok{;}
\KeywordTok{import} \ImportTok{org}\OperatorTok{.}\ImportTok{apache}\OperatorTok{.}\ImportTok{pdfbox}\OperatorTok{.}\ImportTok{pdmodel}\OperatorTok{.}\ImportTok{interactive}\OperatorTok{.}\ImportTok{form}\OperatorTok{.}\ImportTok{PDField}\OperatorTok{;}
\KeywordTok{import} \ImportTok{org}\OperatorTok{.}\ImportTok{apache}\OperatorTok{.}\ImportTok{pdfbox}\OperatorTok{.}\ImportTok{pdmodel}\OperatorTok{.}\ImportTok{interactive}\OperatorTok{.}\ImportTok{form}\OperatorTok{.}\ImportTok{PDSignatureField}\OperatorTok{;}

\ControlFlowTok{try} \OperatorTok{(}\NormalTok{PDDocument doc }\OperatorTok{=}\NormalTok{ Loader}\OperatorTok{.}\FunctionTok{loadPDF}\OperatorTok{(}\NormalTok{pdfBytes}\OperatorTok{))} \OperatorTok{\{}
    \CommentTok{// Inspect the trailer}
\NormalTok{    COSDictionary trailer }\OperatorTok{=}\NormalTok{ doc}\OperatorTok{.}\FunctionTok{getDocument}\OperatorTok{().}\FunctionTok{getTrailer}\OperatorTok{();}
\NormalTok{    COSArray id }\OperatorTok{=}\NormalTok{ trailer}\OperatorTok{.}\FunctionTok{getCOSArray}\OperatorTok{(}\NormalTok{COSName}\OperatorTok{.}\FunctionTok{ID}\OperatorTok{);}
\NormalTok{    log}\OperatorTok{.}\FunctionTok{info}\OperatorTok{(}\StringTok{"File ID present: \{\}"}\OperatorTok{,}\NormalTok{ id }\OperatorTok{!=} \KeywordTok{null}\OperatorTok{);}
\NormalTok{    log}\OperatorTok{.}\FunctionTok{info}\OperatorTok{(}\StringTok{"Trailer /Size: \{\}"}\OperatorTok{,}\NormalTok{ trailer}\OperatorTok{.}\FunctionTok{getInt}\OperatorTok{(}\NormalTok{COSName}\OperatorTok{.}\FunctionTok{SIZE}\OperatorTok{));}

    \CommentTok{// Navigate to signatures through the AcroForm}
\NormalTok{    PDAcroForm acroForm }\OperatorTok{=}\NormalTok{ doc}\OperatorTok{.}\FunctionTok{getDocumentCatalog}\OperatorTok{().}\FunctionTok{getAcroForm}\OperatorTok{();}
    \ControlFlowTok{if} \OperatorTok{(}\NormalTok{acroForm }\OperatorTok{==} \KeywordTok{null}\OperatorTok{)} \OperatorTok{\{}
\NormalTok{        log}\OperatorTok{.}\FunctionTok{info}\OperatorTok{(}\StringTok{"No AcroForm {-}{-} no form fields or signatures"}\OperatorTok{);}
        \ControlFlowTok{return}\OperatorTok{;}
    \OperatorTok{\}}

    \ControlFlowTok{for} \OperatorTok{(}\NormalTok{PDField field }\OperatorTok{:}\NormalTok{ acroForm}\OperatorTok{.}\FunctionTok{getFieldTree}\OperatorTok{())} \OperatorTok{\{}
        \ControlFlowTok{if} \OperatorTok{(}\NormalTok{field }\KeywordTok{instanceof}\NormalTok{ PDSignatureField sigField}\OperatorTok{)} \OperatorTok{\{}
\NormalTok{            COSDictionary sigDict }\OperatorTok{=}\NormalTok{ sigField}\OperatorTok{.}\FunctionTok{getCOSObject}\OperatorTok{();}
\NormalTok{            COSDictionary sigValue }\OperatorTok{=}\NormalTok{ sigDict}\OperatorTok{.}\FunctionTok{getCOSDictionary}\OperatorTok{(}\NormalTok{COSName}\OperatorTok{.}\FunctionTok{V}\OperatorTok{);}

            \ControlFlowTok{if} \OperatorTok{(}\NormalTok{sigValue }\OperatorTok{!=} \KeywordTok{null}\OperatorTok{)} \OperatorTok{\{}
\NormalTok{                COSName subFilter }\OperatorTok{=}\NormalTok{ sigValue}\OperatorTok{.}\FunctionTok{getCOSName}\OperatorTok{(}\NormalTok{COSName}\OperatorTok{.}\FunctionTok{SUB\_FILTER}\OperatorTok{);}
\NormalTok{                COSArray byteRange }\OperatorTok{=}\NormalTok{ sigValue}\OperatorTok{.}\FunctionTok{getCOSArray}\OperatorTok{(}
\NormalTok{                    COSName}\OperatorTok{.}\FunctionTok{getPDFName}\OperatorTok{(}\StringTok{"ByteRange"}\OperatorTok{));}
\NormalTok{                COSString contents }\OperatorTok{=} \OperatorTok{(}\NormalTok{COSString}\OperatorTok{)}\NormalTok{ sigValue}\OperatorTok{.}\FunctionTok{getDictionaryObject}\OperatorTok{(}
\NormalTok{                    COSName}\OperatorTok{.}\FunctionTok{CONTENTS}\OperatorTok{);}

\NormalTok{                log}\OperatorTok{.}\FunctionTok{info}\OperatorTok{(}\StringTok{"Signature \textquotesingle{}\{\}\textquotesingle{}: SubFilter=\{\}, ByteRange=\{\}"}\OperatorTok{,}
\NormalTok{                    field}\OperatorTok{.}\FunctionTok{getFullyQualifiedName}\OperatorTok{(),}\NormalTok{ subFilter}\OperatorTok{,}\NormalTok{ byteRange}\OperatorTok{);}
\NormalTok{                log}\OperatorTok{.}\FunctionTok{info}\OperatorTok{(}\StringTok{"  /Contents hex length: \{\} bytes"}\OperatorTok{,}
\NormalTok{                    contents}\OperatorTok{.}\FunctionTok{getBytes}\OperatorTok{().}\FunctionTok{length}\OperatorTok{);}
            \OperatorTok{\}}
        \OperatorTok{\}}
    \OperatorTok{\}}
\OperatorTok{\}}
\end{Highlighting}
\end{Shaded}

\section{Common Structural Problems and
Debugging}\label{_common_structural_problems_and_debugging}

\subsection{Incorrect /Length Values}\label{_incorrect_length_values}

A stream's \texttt{/Length} does not match its actual byte count. This
corrupts parsing downstream and can shift xref offsets, making
byte-range verification fail. QPDF catches this:

\begin{Shaded}
\begin{Highlighting}[]
\ExtensionTok{qpdf} \AttributeTok{{-}{-}check}\NormalTok{ suspect.pdf}
\end{Highlighting}
\end{Shaded}

\subsection{Missing /ID Array}\label{_missing_id_array}

Some generators omit \texttt{/ID} entirely. While technically optional,
its absence can trip up validators that check it as part of integrity
verification. Use the PDFBox code above to detect missing \texttt{/ID}
programmatically.

\subsection{Multiple \%\%EOF Markers}\label{_multiple_eof_markers}

A PDF with incremental updates (Chapter 8) will have multiple
\texttt{\%\%EOF} markers. This is normal for signed documents. But some
web-based PDF processors strip everything after the first
\texttt{\%\%EOF}, destroying all incremental updates and all signatures.

\subsection{Trailing Bytes After Final
\%\%EOF}\label{_trailing_bytes_after_final_eof}

Bytes after the last \texttt{\%\%EOF} are not part of the PDF. For
signed documents, the \texttt{/ByteRange} of the outermost signature
must cover every byte up to the end of the file. Trailing bytes outside
the byte range create a validation gap that the Bochum ISA attacks
(Chapter 24) exploited.

\subsection{Cross-Reference Table
Corruption}\label{_cross_reference_table_corruption}

Byte offsets in the xref table do not match actual object positions.
This happens when a tool modifies the body without rebuilding the xref,
or when binary transfer corrupts bytes.

QPDF can repair broken xref tables, but repair involves a full rewrite,
which invalidates all signatures. Use \texttt{-\/-check} for diagnosis
without modification. For signed documents, repair is not an
option --- you need the original uncorrupted file.

\section{What Comes Next}\label{_what_comes_next}

We have now dissected the complete physical structure of a PDF file:
header, body, cross-reference table (in both classic ASCII and
compressed stream formats), object streams, trailer, and
\texttt{\%\%EOF}. We have mapped every PDF object type to its PDFBox COS
class, walked through a minimal file byte by byte, and connected the
physical structure to the requirements of digital signatures.

Chapter 8 builds directly on this foundation with \textbf{incremental
updates} --- the append-only mechanism that makes PDF signatures
possible in the first place. If this chapter was about the structure of
a single snapshot, Chapter 8 is about how those snapshots stack on top
of each other, forming the layered history that signatures depend on to
survive.

\chapter{Incremental Updates}\label{ch08}

\section{What Incremental Updates
Are}\label{_what_incremental_updates_are}

A PDF file, at its simplest, has four parts: a header
(\texttt{\%PDF-1.x} or \texttt{\%PDF-2.0}), a body of numbered objects,
a cross-reference table (or cross-reference stream) that maps object
numbers to byte offsets, and a trailer that points to the root object
and the cross-reference table. The file ends with the \texttt{\%\%EOF}
marker.

Chapter 7 covered this structure in detail. What Chapter 7 did not cover
is what happens when you \emph{modify} a PDF.

There are exactly two ways to save changes to a PDF file:

\begin{enumerate}
\def\labelenumi{\arabic{enumi}.}
\item
  \textbf{Full rewrite} --- rebuild the entire file from scratch,
  writing a new body, a new cross-reference table, and a new trailer.
  The old file is replaced entirely.
\item
  \textbf{Incremental update} (also called \textbf{incremental save} or
  \textbf{append mode}) --- leave the original file untouched and
  \emph{append} new data to the end. The new data contains only the
  objects that changed (or were added), a new cross-reference section
  mapping those objects, and a new trailer. The file now has two
  \texttt{\%\%EOF} markers: the original one (buried in the middle of
  the file) and the new one at the end.
\end{enumerate}

ISO 32000-2:2020 Section 7.5.6 defines incremental updates:

\begin{quote}
The updated PDF file shall contain the original data of the PDF file
plus appended update data. Updates may be appended to a file multiple
times. The update data added each time a file is updated shall contain
only those objects that have been added or modified, along with a
cross-reference section and trailer for those objects.

---  ISO 32000-2:2020 Section 7.5.6
\end{quote}

The key phrase is "the original data of the PDF file" --- the original
bytes are preserved, byte for byte, untouched. This is not a diff or a
delta. It is raw appending.

\section{Why Incremental Updates
Exist}\label{_why_incremental_updates_exist}

Incremental updates solve three distinct problems, and understanding all
three explains why the mechanism is designed the way it is.

\subsection{Problem 1: Performance on Large
Documents}\label{_problem_1_performance_on_large_documents}

A 500-page PDF with embedded fonts and high-resolution images might be
200 MB. If a user fills in a single form field --- changing perhaps 50
bytes of actual content --- a full rewrite would require reading,
processing, and writing all 200 MB. An incremental update appends a few
kilobytes.

This was especially important in the 1990s when the format was designed.
Disk I/O was slow, memory was scarce, and rewriting a large file for a
minor change was unacceptable. Even today, incremental saves are
noticeably faster for large documents.

\subsection{Problem 2: Preserving Previous
State}\label{_problem_2_preserving_previous_state}

When a PDF is incrementally updated, the previous revision is still
there --- it is the bytes before the appended data. This means you can
"undo" an incremental update by truncating the file at the previous
\texttt{\%\%EOF} boundary, recovering the exact document as it existed
before the change.

This property is not a side effect. It is a design goal. ISO
32000-2:2020 Section 7.5.6 explicitly states that conforming readers
should be able to retrieve previous revisions of a document by following
the chain of \texttt{/Prev} pointers in the trailers.

For workflows where audit trails matter --- legal documents, contracts,
government forms --- this built-in version history is valuable. You do
not need an external version control system. The PDF file \emph{is} the
version history.

\subsection{Problem 3: Digital
Signatures}\label{_problem_3_digital_signatures}

This is the reason we are here, and it deserves its own section
(\hyperref[ch08-why-signatures-need-incremental]{Why Incremental Updates
Are Critical for Signatures}). The short version: a digital signature
covers a specific byte range. If you rewrite the file, the bytes change,
and the signature becomes invalid. Incremental updates allow
modifications \emph{after} signing without disturbing the signed bytes.

Problem 3 is not optional. If you are building any tool that edits PDFs
which may contain signatures --- even if your tool has nothing to do
with signing --- you must support incremental saves. A full rewrite of a
signed PDF destroys every signature in it. There is no workaround, no
clever trick, no way to reassemble the pieces. The signed bytes are
gone.

\section{Anatomy of an Incremental
Update}\label{_anatomy_of_an_incremental_update}

Let us walk through exactly what gets appended during an incremental
update. Suppose we start with a simple PDF containing three objects: a
catalog (object 1), a page (object 2), and a page content stream (object
3).

\subsection{The Original File}\label{_the_original_file}

\begin{verbatim}
%PDF-1.7
1 0 obj                         
<< /Type /Catalog /Pages 2 0 R >>
endobj

2 0 obj
<< /Type /Page /Contents 3 0 R /MediaBox [0 0 612 792] >>
endobj

3 0 obj
<< /Length 44 >>
stream
BT /F1 12 Tf 100 700 Td (Hello) Tj ET
endstream
endobj

xref                            
0 4
0000000000 65535 f
0000000009 00000 n
0000000066 00000 n
0000000148 00000 n

trailer                         
<< /Size 4 /Root 1 0 R >>

startxref                       
252
%%EOF                           
\end{verbatim}

\begin{itemize}
\item
  Three objects in the body
\item
  Cross-reference table mapping object numbers to byte offsets
\item
  Trailer with \texttt{/Size} (total object count) and \texttt{/Root}
  (catalog reference)
\item
  \texttt{startxref} gives the byte offset of the xref table
\item
  End-of-file marker
\end{itemize}

Nothing surprising here. This is the standard PDF structure from Chapter
7. Let us say this file is 350 bytes long.

\subsection{Appending an Incremental
Update}\label{_appending_an_incremental_update}

Now suppose we want to change the text on the page from "Hello" to
"Hello World". We need a new version of object 3. Instead of rewriting
the entire file, we append:

\begin{verbatim}
3 0 obj                         
<< /Length 50 >>
stream
BT /F1 12 Tf 100 700 Td (Hello World) Tj ET
endstream
endobj

xref                            
0 1
0000000000 65535 f
3 1                             
0000000351 00000 n

trailer                         
<< /Size 4 /Root 1 0 R /Prev 252 >>

startxref                       
460
%%EOF                           
\end{verbatim}

\begin{itemize}
\item
  The new version of object 3, with updated content
\item
  A new cross-reference section --- but only for changed objects
\item
  The xref entry for object 3 now points to byte offset 351 (the start
  of the new object 3)
\item
  The new trailer includes a \texttt{/Prev} key pointing to the
  \emph{previous} xref table (at byte 252)
\item
  \texttt{startxref} points to the new xref section
\item
  A second \texttt{\%\%EOF} marker
\end{itemize}

The file now contains \emph{two} copies of object 3: the original at
byte offset 148 and the new one at byte offset 351. The original
cross-reference table (at byte 252) still points to the old version. The
new cross-reference section (at byte 460) points to the new version.

A conforming reader starts from the end of the file, finds the last
\texttt{startxref}, reads the new xref section, and sees that object 3
is at byte 351. For objects 1 and 2, the new xref section has no entries
--- so the reader follows the \texttt{/Prev} pointer to the previous
xref table and finds them there.

The result: the reader sees the updated document. But the old bytes are
still present, untouched, in the first 350 bytes of the file.

\subsection{\texorpdfstring{The \texttt{/Prev}
Chain}{The /Prev Chain}}\label{_the_prev_chain}

The \texttt{/Prev} key in each trailer forms a linked list of
cross-reference sections, ordered from newest to oldest. A reader
reconstructing the full object map starts with the newest xref section,
then walks the \texttt{/Prev} chain, adding entries for any object
numbers it has not yet seen. The first entry wins --- newer revisions
override older ones.

For a PDF that has been incrementally updated three times, the chain
looks like this:

\begin{verbatim}
%%EOF  (revision 3 — current)
  startxref -> xref_3
  trailer_3: /Prev -> xref_2

%%EOF  (revision 2)
  startxref -> xref_2
  trailer_2: /Prev -> xref_1

%%EOF  (revision 1 — original)
  startxref -> xref_1
  trailer_1: (no /Prev — this is the first revision)
\end{verbatim}

Each revision's trailer may also update the \texttt{/Root} reference if
the catalog object was modified. In practice, the \texttt{/Root} often
stays the same --- the catalog object number does not change, but the
catalog \emph{object} might be replaced with a new version that
references new page trees or updated AcroForm dictionaries.

\textbf{Counting \texttt{\%\%EOF} markers tells you the number of
revisions.} A freshly created PDF has one \texttt{\%\%EOF}. After one
incremental save, it has two. After \emph{n} incremental saves, it has
\emph{n} + 1\ldots\hspace{0pt} unless some tool has done a full rewrite
at some point, collapsing all revisions back to one. A signed PDF with
two signatures typically has three \texttt{\%\%EOF} markers: the
original document, the first signature, and the second signature.

\section{How Objects Are Overridden}\label{_how_objects_are_overridden}

The override mechanism is simple: same object number, new xref entry.

When a conforming reader encounters object \texttt{3\ 0} in both the new
xref section and a previous one, it uses the entry from the newest
section. The old object still exists in the file --- it occupies bytes,
it can be found by parsing --- but it is logically superseded.

This has important implications:

\begin{itemize}
\item
  \textbf{No object is ever deleted.} You can add an xref entry marking
  an object as "free" (generation number incremented, next free object
  pointer updated), but the bytes of the old object remain in the file.
  Deleting content from a PDF does not remove it from the file --- it
  merely hides it behind a newer xref section.
\item
  \textbf{Object numbers are stable.} Object 3 is always object 3,
  whether you are looking at revision 1 or revision 5. The
  \emph{content} of object 3 may differ across revisions, but the number
  does not change. This is critical for signatures because the
  \texttt{/ByteRange} and cross-reference tables use byte offsets and
  object numbers as their addressing scheme.
\item
  \textbf{Generation numbers exist but are rarely used.} The spec allows
  incrementing the generation number when an object is freed and reused,
  but in practice, most tools reuse generation 0 for everything. Do not
  write code that assumes generation numbers will be non-zero --- you
  will almost never see them in the wild.
\end{itemize}

Some tools claim to "optimize" PDFs by removing unused objects and
rebuilding the xref table. This is a full rewrite. If the PDF contains
signatures, every signature is destroyed. Tools that do this include
some versions of Ghostscript's \texttt{pdfwrite} device, QPDF's
\texttt{-\/-linearize} option (when not combined with
\texttt{-\/-preserve-unreferenced-objects}), and various online "PDF
optimizer" services. If you are processing signed PDFs in an automated
pipeline, test your tools carefully.

\section{Cross-Reference Streams and Incremental
Updates}\label{_cross_reference_streams_and_incremental_updates}

PDF 1.5 introduced cross-reference streams as an alternative to the
traditional ASCII cross-reference table. Instead of a plaintext table,
the cross-reference data is stored as a compressed binary stream object.
The trailer dictionary is merged into the stream dictionary --- there is
no separate \texttt{trailer} keyword.

Incremental updates work identically with cross-reference streams. The
appended data contains a new cross-reference stream object with entries
for modified objects, and the stream dictionary includes a
\texttt{/Prev} key pointing to the previous cross-reference stream's
byte offset.

\begin{verbatim}
15 0 obj                                    
<< /Type /XRef
   /Size 16
   /W [1 3 1]                              
   /Root 1 0 R
   /Prev 8432                              
   /Filter /FlateDecode
   /Length 87 >>
stream
...compressed binary xref data...
endstream
endobj

startxref
9150                                        
%%EOF
\end{verbatim}

\begin{itemize}
\item
  The cross-reference stream is itself a PDF object (object 15 in this
  case)
\item
  \texttt{/W} defines the width of each field in the binary data: 1 byte
  for type, 3 bytes for offset, 1 byte for generation
\item
  \texttt{/Prev} points to the previous cross-reference stream, just as
  in the classic trailer
\item
  \texttt{startxref} points to this stream object
\end{itemize}

The \texttt{/Prev} chain works exactly the same way. A reader walks the
chain, building the complete object map from newest to oldest, using the
first entry encountered for each object number.

A PDF file can mix classic xref tables and cross-reference streams
across revisions. The original file might use a classic table, and an
incremental update might append a cross-reference stream (or vice
versa). A conforming reader must handle both. PDFBox handles this
transparently --- \texttt{Loader.loadPDF()} follows the \texttt{/Prev}
chain regardless of whether each link is a table or a stream.

\section{Why Incremental Updates Are Critical for
Signatures}\label{ch08-why-signatures-need-incremental}

This is the most important section in this chapter. Everything else ---
the performance benefits, the revision history, the object override
mechanics --- is useful context. But for digital signatures, incremental
updates are not a convenience. They are a \emph{requirement}.

\subsection{The Fundamental
Constraint}\label{_the_fundamental_constraint}

A PDF digital signature covers a specific byte range (see Chapter 13 for
the complete deep dive). The \texttt{/ByteRange} array in the signature
dictionary specifies exactly which bytes in the file are hashed and
signed. For the signature to remain valid, those bytes must not change.
Not a single one.

If you do a full rewrite of a signed PDF --- reorganizing objects,
recomputing byte offsets, rebuilding the cross-reference table --- the
bytes at the signed offsets will be different. The hash will not match.
The signature is dead.

An incremental update avoids this problem by \emph{not touching the
original bytes}. The signed region (which covers everything from the
start of the file through the \texttt{/Contents} gap and out to the
\texttt{\%\%EOF} of the revision that was signed) remains intact. New
content is appended \emph{after} the signed region.

\subsection{Signing Is an Incremental
Update}\label{_signing_is_an_incremental_update}

When PDFBox (or any compliant library) signs a PDF, it performs an
incremental update. The signing operation appends:

\begin{itemize}
\item
  A new or updated signature field object (with the \texttt{/V}
  reference pointing to the signature dictionary)
\item
  The signature dictionary object itself (containing
  \texttt{/ByteRange}, \texttt{/Contents}, \texttt{/Filter},
  \texttt{/SubFilter}, and the rest)
\item
  An updated AcroForm dictionary (with the new field added to
  \texttt{/Fields} and \texttt{/SigFlags} set)
\item
  A new cross-reference section mapping these new/updated objects
\item
  A new trailer with \texttt{/Prev} pointing to the previous revision
\item
  A new \texttt{\%\%EOF}
\end{itemize}

The signature's byte range covers everything from byte 0 to the end of
this new \texttt{\%\%EOF} --- \emph{including} the bytes of all previous
revisions. The signature literally covers the entire file, minus the gap
where the CMS blob sits.

\subsection{Multiple Signatures}\label{_multiple_signatures}

Now consider what happens when a second person signs the same PDF. The
second signature is another incremental update, appended after the
first. The second signature's byte range starts at byte 0 and extends to
the end of \emph{its} \texttt{\%\%EOF} --- which means it covers the
first signature's bytes too.

The first signature's byte range covers bytes up to the first
signature's \texttt{\%\%EOF}. The second signature's appended data falls
\emph{after} that range. As far as the first signature is concerned,
those bytes do not exist --- they are outside the range it covers, and
that is fine.

This is how multiple signatures coexist:

\begin{verbatim}
|  Original document          |  Revision 1 (%%EOF)
|  + First signature append   |  Revision 2 (%%EOF)
|  + Second signature append  |  Revision 3 (%%EOF)

First signature's byte range:  [0 .. end of revision 2]   (minus its /Contents gap)
Second signature's byte range: [0 .. end of revision 3]   (minus its /Contents gap)
\end{verbatim}

The first signature validates against revisions 1 and 2. The second
signature validates against revisions 1, 2, and 3. Neither signature
invalidates the other.

This only works because each signature is appended incrementally. If the
second signer's tool did a full rewrite instead of an incremental save,
the first signature's bytes would be destroyed, and both signatures
would be invalid. This is why \texttt{/SigFlags} bit 2 (AppendOnly)
exists --- it tells editors "you must use incremental save, or you will
break things."

\subsection{Form Filling After
Signing}\label{_form_filling_after_signing}

The same principle applies to form filling after a certification
signature with \texttt{/DocMDP} permission level 2 (which allows form
filling and signing but no other changes). When a user fills in a form
field, the PDF viewer should save the change as an incremental update.
The filled field value is a new object appended to the file, and the
certification signature remains valid because its byte range is
untouched.

If the viewer does a full save instead, the certification signature
breaks. The viewer has just violated the signer's explicit instruction
(via \texttt{/DocMDP}) that the document should only be modified
incrementally.

\section{What Goes Wrong: Flattening and Full
Rewrites}\label{_what_goes_wrong_flattening_and_full_rewrites}

Now that we understand why incremental updates matter, let us catalog
the ways they get destroyed.

\subsection{The Flattening Problem}\label{_the_flattening_problem}

"Flattening" a PDF typically means merging all incremental updates into
a single revision --- rewriting the file with one body, one
cross-reference table, one trailer, and one \texttt{\%\%EOF}. Some tools
call this "optimizing," "linearizing," "compressing," or just "saving."

The result is a smaller, cleaner file. It is also a file in which every
digital signature is invalid.

Common tools that flatten by default or by option:

\begin{itemize}
\item
  \textbf{Ghostscript} (\texttt{gs\ -sDEVICE=pdfwrite}) --- always does
  a full rewrite; there is no incremental mode
\item
  \textbf{QPDF} with \texttt{-\/-linearize} --- reorganizes the file for
  progressive loading, destroying incremental history
\item
  \textbf{Some online PDF tools} --- "compress PDF" or "optimize PDF"
  services typically do a full rewrite
\item
  \textbf{LibreOffice} (depending on version and configuration) ---
  older versions did not always use incremental save, even when
  \texttt{/SigFlags} bit 2 was set
\item
  \textbf{Adobe Acrobat} "Save As" vs "Save" --- "Save" uses incremental
  update; "Save As" does a full rewrite (destroying signatures unless
  the user re-signs)
\end{itemize}

If you are building an automated PDF processing pipeline and signed PDFs
will pass through it, you must audit every tool in the pipeline to
ensure it uses incremental saves (or does not modify the file at all). A
single non-incremental save anywhere in the pipeline will invalidate
every signature. This is the number one cause of "signatures were valid
when I sent the PDF but invalid when the recipient opened it."

\subsection{The "Save" vs "Save As"
Trap}\label{_the_save_vs_save_as_trap}

Adobe Acrobat's behavior is the one most users encounter:

\begin{itemize}
\item
  \textbf{File \textgreater{} Save} performs an incremental update.
  Signatures are preserved.
\item
  \textbf{File \textgreater{} Save As} writes a new file from scratch.
  Every signature in the document is destroyed.
\end{itemize}

Acrobat shows a warning dialog when you attempt Save As on a signed
document, but many users click through warnings without reading them.
Other PDF editors may not even warn.

When training users who work with signed PDFs, the single most important
instruction is: "Never use Save As. Always use Save." This one piece of
advice prevents more signature breakage than any amount of technical
configuration.

\section{Extracting a Specific
Revision}\label{_extracting_a_specific_revision}

Because incremental updates are purely additive --- new bytes appended
after the previous \texttt{\%\%EOF} --- extracting a previous revision
is trivial: truncate the file at the appropriate \texttt{\%\%EOF}
boundary.

If a PDF has three \texttt{\%\%EOF} markers (three revisions), you can
extract:

\begin{itemize}
\item
  \textbf{Revision 1}: all bytes from the start of the file up to and
  including the first \texttt{\%\%EOF}
\item
  \textbf{Revision 2}: all bytes from the start up to and including the
  second \texttt{\%\%EOF}
\item
  \textbf{Revision 3}: the entire file (up to and including the third
  \texttt{\%\%EOF})
\end{itemize}

This is not a diff extraction --- each "revision" is a complete,
self-contained PDF file that a reader can open and render. Revision 1 is
the original document. Revision 2 is the original plus the first set of
changes. And so on.

\subsection{Extracting with PDFBox}\label{_extracting_with_pdfbox}

PDFBox does not provide a built-in "extract revision N" method, but the
approach is straightforward: find the \emph{N}-th \texttt{\%\%EOF}
marker and copy bytes from 0 to that position.

\begin{Shaded}
\begin{Highlighting}[]
\KeywordTok{import} \ImportTok{java}\OperatorTok{.}\ImportTok{io}\OperatorTok{.}\ImportTok{ByteArrayOutputStream}\OperatorTok{;}
\KeywordTok{import} \ImportTok{java}\OperatorTok{.}\ImportTok{io}\OperatorTok{.}\ImportTok{IOException}\OperatorTok{;}
\KeywordTok{import} \ImportTok{java}\OperatorTok{.}\ImportTok{nio}\OperatorTok{.}\ImportTok{charset}\OperatorTok{.}\ImportTok{StandardCharsets}\OperatorTok{;}
\KeywordTok{import} \ImportTok{java}\OperatorTok{.}\ImportTok{util}\OperatorTok{.}\ImportTok{ArrayList}\OperatorTok{;}
\KeywordTok{import} \ImportTok{java}\OperatorTok{.}\ImportTok{util}\OperatorTok{.}\ImportTok{List}\OperatorTok{;}

\CommentTok{/**}
 \CommentTok{*}\NormalTok{ Extracts individual revisions from a PDF with incremental updates}\CommentTok{.}
 \CommentTok{*}\NormalTok{ Each revision is a complete}\CommentTok{,}\NormalTok{ valid PDF file}\CommentTok{.}
 \CommentTok{*/}
\KeywordTok{public} \KeywordTok{class}\NormalTok{ RevisionExtractor }\OperatorTok{\{}

    \KeywordTok{private} \DataTypeTok{static} \DataTypeTok{final} \DataTypeTok{byte}\OperatorTok{[]}\NormalTok{ EOF\_MARKER }\OperatorTok{=} \StringTok{"\%\%EOF"}\OperatorTok{.}\FunctionTok{getBytes}\OperatorTok{(}\NormalTok{StandardCharsets}\OperatorTok{.}\FunctionTok{US\_ASCII}\OperatorTok{);}

    \CommentTok{/**}
     \CommentTok{*}\NormalTok{ Find byte offsets of all }\CommentTok{\%\%}\NormalTok{EOF markers in the PDF}\CommentTok{.}
     \CommentTok{*/}
    \KeywordTok{public} \DataTypeTok{static} \BuiltInTok{List}\OperatorTok{\textless{}}\BuiltInTok{Integer}\OperatorTok{\textgreater{}} \FunctionTok{findEofOffsets}\OperatorTok{(}\DataTypeTok{byte}\OperatorTok{[]}\NormalTok{ pdfBytes}\OperatorTok{)} \OperatorTok{\{}
        \BuiltInTok{List}\OperatorTok{\textless{}}\BuiltInTok{Integer}\OperatorTok{\textgreater{}}\NormalTok{ offsets }\OperatorTok{=} \KeywordTok{new} \BuiltInTok{ArrayList}\OperatorTok{\textless{}\textgreater{}();}
        \ControlFlowTok{for} \OperatorTok{(}\DataTypeTok{int}\NormalTok{ i }\OperatorTok{=} \DecValTok{0}\OperatorTok{;}\NormalTok{ i }\OperatorTok{\textless{}=}\NormalTok{ pdfBytes}\OperatorTok{.}\FunctionTok{length} \OperatorTok{{-}}\NormalTok{ EOF\_MARKER}\OperatorTok{.}\FunctionTok{length}\OperatorTok{;}\NormalTok{ i}\OperatorTok{++)} \OperatorTok{\{}
            \DataTypeTok{boolean}\NormalTok{ match }\OperatorTok{=} \KeywordTok{true}\OperatorTok{;}
            \ControlFlowTok{for} \OperatorTok{(}\DataTypeTok{int}\NormalTok{ j }\OperatorTok{=} \DecValTok{0}\OperatorTok{;}\NormalTok{ j }\OperatorTok{\textless{}}\NormalTok{ EOF\_MARKER}\OperatorTok{.}\FunctionTok{length}\OperatorTok{;}\NormalTok{ j}\OperatorTok{++)} \OperatorTok{\{}
                \ControlFlowTok{if} \OperatorTok{(}\NormalTok{pdfBytes}\OperatorTok{[}\NormalTok{i }\OperatorTok{+}\NormalTok{ j}\OperatorTok{]} \OperatorTok{!=}\NormalTok{ EOF\_MARKER}\OperatorTok{[}\NormalTok{j}\OperatorTok{])} \OperatorTok{\{}
\NormalTok{                    match }\OperatorTok{=} \KeywordTok{false}\OperatorTok{;}
                    \ControlFlowTok{break}\OperatorTok{;}
                \OperatorTok{\}}
            \OperatorTok{\}}
            \ControlFlowTok{if} \OperatorTok{(}\NormalTok{match}\OperatorTok{)} \OperatorTok{\{}
                \CommentTok{// Include the newline after \%\%EOF if present}
                \DataTypeTok{int}\NormalTok{ end }\OperatorTok{=}\NormalTok{ i }\OperatorTok{+}\NormalTok{ EOF\_MARKER}\OperatorTok{.}\FunctionTok{length}\OperatorTok{;}
                \ControlFlowTok{if} \OperatorTok{(}\NormalTok{end }\OperatorTok{\textless{}}\NormalTok{ pdfBytes}\OperatorTok{.}\FunctionTok{length} \OperatorTok{\&\&}\NormalTok{ pdfBytes}\OperatorTok{[}\NormalTok{end}\OperatorTok{]} \OperatorTok{==} \CharTok{\textquotesingle{}\textbackslash{}r\textquotesingle{}}\OperatorTok{)}\NormalTok{ end}\OperatorTok{++;}
                \ControlFlowTok{if} \OperatorTok{(}\NormalTok{end }\OperatorTok{\textless{}}\NormalTok{ pdfBytes}\OperatorTok{.}\FunctionTok{length} \OperatorTok{\&\&}\NormalTok{ pdfBytes}\OperatorTok{[}\NormalTok{end}\OperatorTok{]} \OperatorTok{==} \CharTok{\textquotesingle{}\textbackslash{}n\textquotesingle{}}\OperatorTok{)}\NormalTok{ end}\OperatorTok{++;}
\NormalTok{                offsets}\OperatorTok{.}\FunctionTok{add}\OperatorTok{(}\NormalTok{end}\OperatorTok{);}
            \OperatorTok{\}}
        \OperatorTok{\}}
        \ControlFlowTok{return}\NormalTok{ offsets}\OperatorTok{;}
    \OperatorTok{\}}

    \CommentTok{/**}
     \CommentTok{*}\NormalTok{ Extract revision N }\CommentTok{(1{-}}\NormalTok{based}\CommentTok{)}\NormalTok{ as a complete PDF byte array}\CommentTok{.}
     \CommentTok{*/}
    \KeywordTok{public} \DataTypeTok{static} \DataTypeTok{byte}\OperatorTok{[]} \FunctionTok{extractRevision}\OperatorTok{(}\DataTypeTok{byte}\OperatorTok{[]}\NormalTok{ pdfBytes}\OperatorTok{,} \DataTypeTok{int}\NormalTok{ revision}\OperatorTok{)} \OperatorTok{\{}
        \BuiltInTok{List}\OperatorTok{\textless{}}\BuiltInTok{Integer}\OperatorTok{\textgreater{}}\NormalTok{ eofOffsets }\OperatorTok{=} \FunctionTok{findEofOffsets}\OperatorTok{(}\NormalTok{pdfBytes}\OperatorTok{);}
        \ControlFlowTok{if} \OperatorTok{(}\NormalTok{revision }\OperatorTok{\textless{}} \DecValTok{1} \OperatorTok{||}\NormalTok{ revision }\OperatorTok{\textgreater{}}\NormalTok{ eofOffsets}\OperatorTok{.}\FunctionTok{size}\OperatorTok{())} \OperatorTok{\{}
            \ControlFlowTok{throw} \KeywordTok{new} \BuiltInTok{IllegalArgumentException}\OperatorTok{(}
                \StringTok{"Revision "} \OperatorTok{+}\NormalTok{ revision }\OperatorTok{+} \StringTok{" out of range. "}
                \OperatorTok{+} \StringTok{"File has "} \OperatorTok{+}\NormalTok{ eofOffsets}\OperatorTok{.}\FunctionTok{size}\OperatorTok{()} \OperatorTok{+} \StringTok{" revisions."}\OperatorTok{);}
        \OperatorTok{\}}
        \DataTypeTok{int}\NormalTok{ endOffset }\OperatorTok{=}\NormalTok{ eofOffsets}\OperatorTok{.}\FunctionTok{get}\OperatorTok{(}\NormalTok{revision }\OperatorTok{{-}} \DecValTok{1}\OperatorTok{);}
        \DataTypeTok{byte}\OperatorTok{[]}\NormalTok{ result }\OperatorTok{=} \KeywordTok{new} \DataTypeTok{byte}\OperatorTok{[}\NormalTok{endOffset}\OperatorTok{];}
        \BuiltInTok{System}\OperatorTok{.}\FunctionTok{arraycopy}\OperatorTok{(}\NormalTok{pdfBytes}\OperatorTok{,} \DecValTok{0}\OperatorTok{,}\NormalTok{ result}\OperatorTok{,} \DecValTok{0}\OperatorTok{,}\NormalTok{ endOffset}\OperatorTok{);}
        \ControlFlowTok{return}\NormalTok{ result}\OperatorTok{;}
    \OperatorTok{\}}
\OperatorTok{\}}
\end{Highlighting}
\end{Shaded}

This approach works for classic incremental updates. It does \emph{not}
handle linearized PDFs correctly (see
\hyperref[ch08-linearized]{Linearized PDFs and Incremental Updates}),
because linearized files have a \texttt{\%\%EOF} at the end of the first
part that is not a revision boundary. For production code, you should
also handle the case where \texttt{\%\%EOF} appears inside a stream
object (rare but technically possible in pathological files).

\subsection{Extracting with QPDF}\label{_extracting_with_qpdf}

QPDF provides a more robust approach from the command line:

\begin{Shaded}
\begin{Highlighting}[]
\CommentTok{\# Show all cross{-}reference sections and their offsets}
\ExtensionTok{qpdf} \AttributeTok{{-}{-}show{-}xref}\NormalTok{ input.pdf}

\CommentTok{\# Extract a specific page range (not exactly revision extraction,}
\CommentTok{\# but useful for isolating content)}
\ExtensionTok{qpdf}\NormalTok{ input.pdf }\AttributeTok{{-}{-}pages}\NormalTok{ . 1{-}5 }\AttributeTok{{-}{-}}\NormalTok{ output.pdf}
\end{Highlighting}
\end{Shaded}

QPDF does not have a built-in "extract revision N" command, but
\texttt{-\/-show-xref} gives you all the information you need to
determine the \texttt{\%\%EOF} boundaries manually.

\section{How PDFBox Handles Incremental
Saves}\label{_how_pdfbox_handles_incremental_saves}

PDFBox 3.0.0 provides two distinct save methods, and choosing the wrong
one is a common source of signature breakage.

\subsection{\texorpdfstring{\texttt{doc.save()} --- Full
Rewrite}{doc.save() --- Full Rewrite}}\label{_doc_save_full_rewrite}

The standard \texttt{save()} method writes the entire document from
scratch:

\begin{Shaded}
\begin{Highlighting}[]
\ControlFlowTok{try} \OperatorTok{(}\NormalTok{PDDocument doc }\OperatorTok{=}\NormalTok{ Loader}\OperatorTok{.}\FunctionTok{loadPDF}\OperatorTok{(}\NormalTok{inputBytes}\OperatorTok{))} \OperatorTok{\{}
    \CommentTok{// Make changes...}
\NormalTok{    doc}\OperatorTok{.}\FunctionTok{save}\OperatorTok{(}\NormalTok{outputStream}\OperatorTok{);}  \CommentTok{// Full rewrite — destroys signatures!}
\OperatorTok{\}}
\end{Highlighting}
\end{Shaded}

This produces a clean, single-revision PDF. All incremental update
history is lost. All signatures are invalid.

Use \texttt{save()} only when:

\begin{itemize}
\item
  The document has no signatures and you do not need to preserve
  revision history
\item
  You are creating a new document from scratch
\item
  You explicitly want to flatten the file (and accept the consequences)
\end{itemize}

\subsection{\texorpdfstring{\texttt{doc.saveIncremental()} --- Append
Mode}{doc.saveIncremental() --- Append Mode}}\label{_doc_saveincremental_append_mode}

The \texttt{saveIncremental()} method appends only the changed objects:

\begin{Shaded}
\begin{Highlighting}[]
\ControlFlowTok{try} \OperatorTok{(}\NormalTok{PDDocument doc }\OperatorTok{=}\NormalTok{ Loader}\OperatorTok{.}\FunctionTok{loadPDF}\OperatorTok{(}\NormalTok{inputBytes}\OperatorTok{))} \OperatorTok{\{}
    \CommentTok{// Make changes...}
\NormalTok{    doc}\OperatorTok{.}\FunctionTok{saveIncremental}\OperatorTok{(}\NormalTok{outputStream}\OperatorTok{);}  \CommentTok{// Appends — preserves signatures}
\OperatorTok{\}}
\end{Highlighting}
\end{Shaded}

This is what you must use when the document contains signatures. PDFBox
tracks which objects have been modified (via the COSUpdateInfo
mechanism) and writes only those objects, plus a new cross-reference
section and trailer, to the output.

\texttt{saveIncremental()} requires that the document was loaded from a
stream or byte array --- PDFBox needs the original bytes to copy them to
the output before appending. If you loaded the document from a
\texttt{File} and then try to write incrementally to a different output
stream, the original bytes are read from the file. If you try to
overwrite the same file, you will corrupt it. Always load from bytes and
write to a new output.

\subsection{\texorpdfstring{The \texttt{saveIncremental()} and Signing
Workflow}{The saveIncremental() and Signing Workflow}}\label{_the_saveincremental_and_signing_workflow}

When signing a PDF, you do not call \texttt{saveIncremental()} directly.
Instead, you use the signing API:

\begin{Shaded}
\begin{Highlighting}[]
\ControlFlowTok{try} \OperatorTok{(}\NormalTok{PDDocument doc }\OperatorTok{=}\NormalTok{ Loader}\OperatorTok{.}\FunctionTok{loadPDF}\OperatorTok{(}\NormalTok{inputBytes}\OperatorTok{))} \OperatorTok{\{}
\NormalTok{    PDSignature signature }\OperatorTok{=} \KeywordTok{new} \FunctionTok{PDSignature}\OperatorTok{();}
\NormalTok{    signature}\OperatorTok{.}\FunctionTok{setFilter}\OperatorTok{(}\NormalTok{PDSignature}\OperatorTok{.}\FunctionTok{FILTER\_ADOBE\_PPKLITE}\OperatorTok{);}
\NormalTok{    signature}\OperatorTok{.}\FunctionTok{setSubFilter}\OperatorTok{(}\NormalTok{PDSignature}\OperatorTok{.}\FunctionTok{SUBFILTER\_ADBE\_PKCS7\_DETACHED}\OperatorTok{);}
\NormalTok{    signature}\OperatorTok{.}\FunctionTok{setSignDate}\OperatorTok{(}\BuiltInTok{Calendar}\OperatorTok{.}\FunctionTok{getInstance}\OperatorTok{());}

\NormalTok{    SignatureOptions options }\OperatorTok{=} \KeywordTok{new} \FunctionTok{SignatureOptions}\OperatorTok{();}
\NormalTok{    options}\OperatorTok{.}\FunctionTok{setPreferredSignatureSize}\OperatorTok{(}\DecValTok{16384}\OperatorTok{);}  \CommentTok{// Reserve 16KB for CMS}

\NormalTok{    doc}\OperatorTok{.}\FunctionTok{addSignature}\OperatorTok{(}\NormalTok{signature}\OperatorTok{,}\NormalTok{ signingInterface}\OperatorTok{,}\NormalTok{ options}\OperatorTok{);}

    \BuiltInTok{ByteArrayOutputStream}\NormalTok{ baos }\OperatorTok{=} \KeywordTok{new} \BuiltInTok{ByteArrayOutputStream}\OperatorTok{();}
\NormalTok{    doc}\OperatorTok{.}\FunctionTok{saveIncremental}\OperatorTok{(}\NormalTok{baos}\OperatorTok{);}  
    \DataTypeTok{byte}\OperatorTok{[]}\NormalTok{ signedPdf }\OperatorTok{=}\NormalTok{ baos}\OperatorTok{.}\FunctionTok{toByteArray}\OperatorTok{();}
\OperatorTok{\}}
\end{Highlighting}
\end{Shaded}

\begin{itemize}
\item
  \texttt{saveIncremental()} here does three things: writes the original
  bytes, appends the new objects (including the signature dictionary
  with a placeholder \texttt{/Contents}), computes the hash over the
  byte ranges, calls your \texttt{SignatureInterface.sign()} method to
  get the CMS bytes, and writes them into the placeholder.
\end{itemize}

The signing operation is \emph{always} an incremental save. PDFBox
enforces this --- there is no \texttt{doc.save()} variant that supports
signing, because signing with a full rewrite is logically impossible
(the signed byte range would cover a file that no longer exists).

\section{\texorpdfstring{Multiple \texttt{\%\%EOF} Markers: What They
Mean}{Multiple \%\%EOF Markers: What They Mean}}\label{_multiple_eof_markers_what_they_mean}

Every \texttt{\%\%EOF} marker in a PDF represents the end of a revision.
Counting them tells you the document's history:

\begin{itemize}
\item
  \textbf{1 \texttt{\%\%EOF}}: A single-revision PDF. Never been
  incrementally updated. Either freshly created, or previously
  incrementally updated but later flattened by a full rewrite.
\item
  \textbf{2 \texttt{\%\%EOF}}: One incremental update after the
  original. Common for a single signature, a single form fill, or a
  single annotation addition.
\item
  \textbf{3 \texttt{\%\%EOF}}: Two incremental updates. Common for a
  document with two signatures, or one signature plus a subsequent form
  fill.
\item
  \textbf{N \texttt{\%\%EOF}}: N-1 incremental updates after the
  original. Documents that have been through complex workflows (multiple
  signers, multiple form fill rounds, annotation rounds) can have many
  revisions.
\end{itemize}

A quick way to count revisions in a PDF is:
\texttt{grep\ -c\ \textquotesingle{}\%\%EOF\textquotesingle{}\ document.pdf}.
This works for most files, but be aware that \texttt{\%\%EOF} could
theoretically appear inside a stream object (unlikely in practice). For
reliable counting, parse the \texttt{/Prev} chain instead.

\subsection{What Revision Counts Tell You About a Signed
PDF}\label{_what_revision_counts_tell_you_about_a_signed_pdf}

For signature validation, the revision count carries diagnostic
information:

\begin{itemize}
\item
  A signed PDF with exactly 2 \texttt{\%\%EOF} markers is the simplest
  case: original document + one signature. The signature covers
  everything in revision 1 and 2.
\item
  A signed PDF with 3 \texttt{\%\%EOF} markers might mean: original +
  signature + post-signature modification (form fill, annotation, second
  signature). The first signature covers revisions 1-2. Whatever is in
  revision 3 was added after signing --- whether it is permissible
  depends on the \texttt{/TransformMethod} and \texttt{/DocMDP}
  settings.
\item
  A signed PDF with only 1 \texttt{\%\%EOF} is suspicious. Either the
  signature was created during the initial save (some tools do this), or
  the file was flattened after signing (which means the signature is
  almost certainly invalid, even if the bytes were carefully
  reconstructed --- see the discussion in Chapter 24 about the
  Incremental Saving Attack).
\end{itemize}

\section{Tracing the Revision
History}\label{_tracing_the_revision_history}

To reconstruct the complete revision history of a PDF, you walk the
\texttt{/Prev} chain from the newest trailer back to the oldest:

\begin{Shaded}
\begin{Highlighting}[]
\KeywordTok{import} \ImportTok{org}\OperatorTok{.}\ImportTok{apache}\OperatorTok{.}\ImportTok{pdfbox}\OperatorTok{.}\ImportTok{Loader}\OperatorTok{;}
\KeywordTok{import} \ImportTok{org}\OperatorTok{.}\ImportTok{apache}\OperatorTok{.}\ImportTok{pdfbox}\OperatorTok{.}\ImportTok{cos}\OperatorTok{.}\ImportTok{COSDictionary}\OperatorTok{;}
\KeywordTok{import} \ImportTok{org}\OperatorTok{.}\ImportTok{apache}\OperatorTok{.}\ImportTok{pdfbox}\OperatorTok{.}\ImportTok{cos}\OperatorTok{.}\ImportTok{COSName}\OperatorTok{;}
\KeywordTok{import} \ImportTok{org}\OperatorTok{.}\ImportTok{apache}\OperatorTok{.}\ImportTok{pdfbox}\OperatorTok{.}\ImportTok{pdmodel}\OperatorTok{.}\ImportTok{PDDocument}\OperatorTok{;}

\KeywordTok{import} \ImportTok{java}\OperatorTok{.}\ImportTok{util}\OperatorTok{.}\ImportTok{ArrayList}\OperatorTok{;}
\KeywordTok{import} \ImportTok{java}\OperatorTok{.}\ImportTok{util}\OperatorTok{.}\ImportTok{List}\OperatorTok{;}

\CommentTok{/**}
 \CommentTok{*}\NormalTok{ Traces the }\CommentTok{/}\NormalTok{Prev chain to enumerate all cross}\CommentTok{{-}}\NormalTok{reference sections}\CommentTok{.}
 \CommentTok{*/}
\KeywordTok{public} \KeywordTok{class}\NormalTok{ RevisionTracer }\OperatorTok{\{}

    \KeywordTok{public} \DataTypeTok{static} \BuiltInTok{List}\OperatorTok{\textless{}}\BuiltInTok{Long}\OperatorTok{\textgreater{}} \FunctionTok{traceXrefOffsets}\OperatorTok{(}\DataTypeTok{byte}\OperatorTok{[]}\NormalTok{ pdfBytes}\OperatorTok{)} \KeywordTok{throws} \BuiltInTok{Exception} \OperatorTok{\{}
        \BuiltInTok{List}\OperatorTok{\textless{}}\BuiltInTok{Long}\OperatorTok{\textgreater{}}\NormalTok{ offsets }\OperatorTok{=} \KeywordTok{new} \BuiltInTok{ArrayList}\OperatorTok{\textless{}\textgreater{}();}
        \ControlFlowTok{try} \OperatorTok{(}\NormalTok{PDDocument doc }\OperatorTok{=}\NormalTok{ Loader}\OperatorTok{.}\FunctionTok{loadPDF}\OperatorTok{(}\NormalTok{pdfBytes}\OperatorTok{))} \OperatorTok{\{}
\NormalTok{            COSDictionary trailer }\OperatorTok{=}\NormalTok{ doc}\OperatorTok{.}\FunctionTok{getDocument}\OperatorTok{().}\FunctionTok{getTrailer}\OperatorTok{();}
            \CommentTok{// The last startxref value}
            \DataTypeTok{long}\NormalTok{ currentOffset }\OperatorTok{=}\NormalTok{ doc}\OperatorTok{.}\FunctionTok{getDocument}\OperatorTok{().}\FunctionTok{getStartXref}\OperatorTok{();}
\NormalTok{            offsets}\OperatorTok{.}\FunctionTok{add}\OperatorTok{(}\NormalTok{currentOffset}\OperatorTok{);}

            \CommentTok{// Walk the /Prev chain}
            \ControlFlowTok{while} \OperatorTok{(}\NormalTok{trailer}\OperatorTok{.}\FunctionTok{containsKey}\OperatorTok{(}\NormalTok{COSName}\OperatorTok{.}\FunctionTok{PREV}\OperatorTok{))} \OperatorTok{\{}
                \DataTypeTok{long}\NormalTok{ prev }\OperatorTok{=}\NormalTok{ trailer}\OperatorTok{.}\FunctionTok{getLong}\OperatorTok{(}\NormalTok{COSName}\OperatorTok{.}\FunctionTok{PREV}\OperatorTok{);}
\NormalTok{                offsets}\OperatorTok{.}\FunctionTok{add}\OperatorTok{(}\NormalTok{prev}\OperatorTok{);}
                \CommentTok{// Note: to get the previous trailer, we would need to}
                \CommentTok{// parse the xref section at the /Prev offset.}
                \CommentTok{// PDFBox\textquotesingle{}s internal parser does this, but the API}
                \CommentTok{// does not expose individual revision trailers directly.}
                \ControlFlowTok{break}\OperatorTok{;} \CommentTok{// Simplified — full implementation needs xref parsing}
            \OperatorTok{\}}
        \OperatorTok{\}}
        \ControlFlowTok{return}\NormalTok{ offsets}\OperatorTok{;}
    \OperatorTok{\}}
\OperatorTok{\}}
\end{Highlighting}
\end{Shaded}

PDFBox's public API does not expose a convenient "iterate over all
revisions" method. The internal \texttt{COSParser} follows the
\texttt{/Prev} chain during loading, but the individual revision
trailers are not directly accessible after parsing. For production
revision analysis, consider using QPDF's \texttt{-\/-show-xref} output
(see \hyperref[ch08-debugging]{Debugging Incremental Updates}) or
writing a custom parser that reads the raw bytes and follows the
\texttt{/Prev} pointers manually.

\section{Linearized PDFs and Incremental Updates}\label{ch08-linearized}

Linearized PDFs (sometimes called "optimized for fast web view") have a
special structure designed for progressive loading over HTTP. A
linearized PDF reorganizes its objects so that the first page can be
displayed before the entire file has been downloaded.

The linearization adds a \textbf{linearization dictionary} (object 1,
typically) and restructures the file into two parts separated by an
internal \texttt{\%\%EOF} --- but this internal \texttt{\%\%EOF} is
\emph{not} a revision boundary. It is part of the linearization
structure.

This creates a compatibility problem with incremental updates:

\begin{itemize}
\item
  A linearized PDF already has an unusual structure with cross-reference
  data at both the beginning and end of the file.
\item
  An incremental update appends data after the final \texttt{\%\%EOF},
  which is fine for non-linearized PDFs but technically invalidates the
  linearization (because the linearization dictionary's \texttt{/L} key
  records the total file length, which is now wrong).
\item
  Some viewers re-linearize on save, which is a full rewrite, which
  destroys signatures.
\end{itemize}

ISO 32000-2:2020 Section 7.5.6 states:

\begin{quote}
If an incremental update is made to a linearized PDF file, the result is
a valid non-linearized PDF file.
\end{quote}

In other words, an incremental update on a linearized PDF destroys the
linearization but produces a valid PDF. This is acceptable --- the file
simply loses its "fast web view" optimization.

If your workflow requires both linearization and digital signatures, you
must linearize \emph{before} signing, not after. Linearizing after
signing is a full rewrite that destroys signatures. Signing after
linearizing produces a valid signed PDF that is no longer linearized.
There is no way to have a linearized, signed PDF that remains linearized
after signing --- the incremental update required by signing inherently
breaks the linearization structure.

\section{The Incremental Saving Attack}\label{ch08-isa}

The incremental update mechanism that makes PDF signatures possible is
also the mechanism that attackers have exploited to \emph{bypass} those
signatures.

The \textbf{Incremental Saving Attack (ISA)}, described by the Ruhr
University Bochum team in their CCS 2019 paper, works as follows:

\begin{enumerate}
\def\labelenumi{\arabic{enumi}.}
\item
  The attacker starts with a legitimately signed PDF
\item
  The attacker appends a new incremental update to the file --- adding,
  modifying, or replacing objects (pages, text, images, form field
  values)
\item
  The attacker does \emph{not} touch the signed bytes; the incremental
  update is appended after the signed region
\item
  The attacker presents the modified file to a PDF viewer
\item
  The viewer reads the file from the end, uses the newest
  cross-reference section, and renders the attacker's content
\item
  The viewer validates the signature by checking the original byte range
  --- which is untouched --- and reports "Signature is valid"
\end{enumerate}

The signature \emph{is} valid --- over the bytes it covers. But the
viewer is displaying content from \emph{outside} the signed byte range.
The user sees a green checkmark and attacker-controlled content.

The defense against ISA is conceptually simple but difficult to
implement correctly: a viewer must detect that content has been appended
after the signed revision and warn the user. Specifically, the viewer
must check whether the byte range of the last signature covers the
entire file. If the file extends beyond the end of the signed byte
range, there are unsigned changes.

Whether those changes are permissible (e.g., a subsequent signature, an
allowed form fill under DocMDP level 2) or malicious requires deeper
analysis of what the appended objects actually do --- which is where the
complexity lies.

Chapter 24 covers the ISA and its variants in complete technical detail,
including the specific techniques attackers used to bypass even the
"unsigned content detected" warnings that some viewers implemented as an
initial defense.

\section{Debugging Incremental Updates}\label{ch08-debugging}

When something goes wrong with a signed PDF --- the signature shows as
invalid, a form field value disappeared, a viewer crashes --- the first
thing to check is the incremental update structure.

\subsection{QPDF: Your Best Friend}\label{_qpdf_your_best_friend}

QPDF is an indispensable tool for PDF debugging. For incremental update
analysis, the key commands are:

\begin{Shaded}
\begin{Highlighting}[]
\CommentTok{\# Show all cross{-}reference entries with their types and offsets}
\ExtensionTok{qpdf} \AttributeTok{{-}{-}show{-}xref}\NormalTok{ signed{-}document.pdf}

\CommentTok{\# Show the structure of the file (objects, streams, references)}
\ExtensionTok{qpdf} \AttributeTok{{-}{-}show{-}object}\OperatorTok{=}\NormalTok{trailer signed{-}document.pdf}

\CommentTok{\# Check the file for structural errors}
\ExtensionTok{qpdf} \AttributeTok{{-}{-}check}\NormalTok{ signed{-}document.pdf}

\CommentTok{\# Show a specific object}
\ExtensionTok{qpdf} \AttributeTok{{-}{-}show{-}object}\OperatorTok{=}\NormalTok{42,0 signed{-}document.pdf}
\end{Highlighting}
\end{Shaded}

The \texttt{-\/-show-xref} output is particularly valuable. It shows
every object number, its byte offset, and which cross-reference section
it belongs to. From this output, you can determine:

\begin{itemize}
\item
  How many cross-reference sections exist (i.e., how many revisions)
\item
  Which objects were added or modified in each revision
\item
  Whether the xref entries are consistent (offsets point to actual
  objects)
\end{itemize}

\subsection{Comparing Revisions}\label{_comparing_revisions}

To understand what changed between revisions, extract each revision (by
truncating at \texttt{\%\%EOF} boundaries as described earlier) and
compare them:

\begin{Shaded}
\begin{Highlighting}[]
\CommentTok{\# Extract revision 1 (original document)}
\CommentTok{\# Suppose the first \%\%EOF is at byte 45230}
\FunctionTok{head} \AttributeTok{{-}c}\NormalTok{ 45230 signed{-}document.pdf }\OperatorTok{\textgreater{}}\NormalTok{ revision1.pdf}

\CommentTok{\# Extract revision 2 (after first signature)}
\CommentTok{\# Suppose the second \%\%EOF is at byte 67890}
\FunctionTok{head} \AttributeTok{{-}c}\NormalTok{ 67890 signed{-}document.pdf }\OperatorTok{\textgreater{}}\NormalTok{ revision2.pdf}

\CommentTok{\# Compare the objects in each revision}
\ExtensionTok{qpdf} \AttributeTok{{-}{-}show{-}xref}\NormalTok{ revision1.pdf }\OperatorTok{\textgreater{}}\NormalTok{ xref1.txt}
\ExtensionTok{qpdf} \AttributeTok{{-}{-}show{-}xref}\NormalTok{ revision2.pdf }\OperatorTok{\textgreater{}}\NormalTok{ xref2.txt}
\FunctionTok{diff}\NormalTok{ xref1.txt xref2.txt}
\end{Highlighting}
\end{Shaded}

The diff will show you exactly which objects were added or modified in
the incremental update. For a signature, you should see the signature
dictionary object, the updated AcroForm, and possibly an updated catalog
or page annotation list.

\subsection{Hex-Level Inspection}\label{_hex_level_inspection}

For the most stubborn problems, there is no substitute for looking at
the raw bytes. Open the PDF in a hex editor and search for
\texttt{\%\%EOF}. Each occurrence marks a revision boundary. The bytes
between two consecutive \texttt{\%\%EOF} markers (plus the preceding
\texttt{startxref} and trailer) constitute one incremental update.

\textbf{A debugging checklist for incremental update problems:}

\begin{enumerate}
\def\labelenumi{\arabic{enumi}.}
\item
  Count \texttt{\%\%EOF} markers --- does the number match the expected
  number of revisions?
\item
  Check each trailer's \texttt{/Prev} value --- does the chain form a
  consistent linked list?
\item
  Verify that \texttt{startxref} in each revision points to a valid xref
  section
\item
  For signed PDFs: does the byte range of each signature end at a
  \texttt{\%\%EOF} boundary?
\item
  Are there any bytes after the last \texttt{\%\%EOF}? (There should not
  be --- trailing garbage can confuse some parsers.)
\item
  Open each extracted revision independently --- does it render
  correctly?
\item
  Compare the object sets across revisions --- do the changes make sense
  for the expected workflow?
\end{enumerate}

\section{Real-World Scenarios}\label{_real_world_scenarios}

The theory of incremental updates is clean. The practice is messy. Here
are the scenarios you will encounter most often.

\subsection{Scenario 1: Single
Signature}\label{_scenario_1_single_signature}

The simplest case. A document is created, then signed. The file has two
revisions (two \texttt{\%\%EOF} markers). The signature's byte range
covers the entire file. No complications.

\subsection{Scenario 2: Two Approval
Signatures}\label{_scenario_2_two_approval_signatures}

Alice signs first, then Bob. The file has three revisions: original,
Alice's signature, Bob's signature. Alice's byte range covers revisions
1-2. Bob's byte range covers revisions 1-3. When validating Alice's
signature, a viewer should check the byte range of Alice's signature and
warn if changes between revision 2 and 3 go beyond what Alice's
\texttt{/TransformMethod} permits (typically, adding more signatures is
allowed).

\subsection{Scenario 3: Certification +
Approval}\label{_scenario_3_certification_approval}

The document author applies a certification signature with
\texttt{/DocMDP} permission level 2 (form fill and sign allowed). Later,
a user fills form fields and applies an approval signature. The file has
three or four revisions (original, certification, form fill, approval
--- or the form fill and approval might be in the same incremental
update). The certification signature permits the form fill and the
additional signature. A viewer must verify that the changes in the later
revisions are within the scope of the DocMDP permissions.

\subsection{Scenario 4: Annotations After
Signing}\label{_scenario_4_annotations_after_signing}

A user adds a sticky note annotation to a signed PDF. If the editor
saves incrementally, a new revision is appended with the annotation
object and an updated page annotations array. The original signature's
byte range is untouched. Whether the annotation is \emph{allowed}
depends on the signature's \texttt{/TransformMethod}:

\begin{itemize}
\item
  DocMDP level 1 (no changes allowed): the annotation violates the
  permissions. A conforming viewer should warn that the document has
  been modified in a way that exceeds the signer's permissions.
\item
  DocMDP level 2 (form fill and sign only): annotations are not form
  filling, so this is a violation.
\item
  DocMDP level 3 (annotations, form filling, and signing allowed): the
  annotation is permitted.
\end{itemize}

The incremental update itself is structurally fine in all three cases.
The question is whether the \emph{content} of the update is semantically
permissible. This distinction --- structural validity vs. semantic
permissibility --- is at the heart of the Incremental Saving Attack
(Chapter 24).

\subsection{Scenario 5: The Flattened
Disaster}\label{_scenario_5_the_flattened_disaster}

A signed PDF passes through an automated document management system that
"optimizes" PDFs for storage. The optimization tool does a full rewrite.
The output PDF has one \texttt{\%\%EOF}, no \texttt{/Prev} chain, and
the signature dictionary's byte range references offsets that no longer
correspond to the correct content. The signature is permanently invalid.
The signed bytes are gone.

There is no recovery. The original file must be retrieved from backup
(if it exists) or the document must be re-signed.

This scenario is \emph{extremely} common in enterprise environments.
Document management systems, email gateways, PDF/A conversion tools, and
print spoolers all potentially rewrite PDFs. If your organization
processes signed PDFs, audit your entire document pipeline for
non-incremental saves. Every tool that touches a signed PDF must either
leave it completely unmodified or use incremental saves exclusively.

\subsection{Scenario 6: The Double-Linearization
Trap}\label{_scenario_6_the_double_linearization_trap}

A document is created with linearization enabled (for fast web view),
signed, and then processed by a tool that re-linearizes it. The
re-linearization is a full rewrite. The signature is destroyed.

The fix: disable linearization in any tool that processes signed PDFs,
or ensure that no post-signing processing step involves linearization.

\section{Summary}\label{_summary_7}

Incremental updates are the mechanism that makes PDF digital signatures
possible. Without them, signing a PDF would require that the document
never be modified again --- no additional signatures, no form filling,
no annotations, nothing. With them, a signed PDF can evolve through a
controlled sequence of revisions, each one adding new content without
disturbing the bytes that previous signatures cover.

The rules are simple:

\begin{itemize}
\item
  A signed PDF must be saved incrementally. Always.
\item
  The byte range of each signature covers a specific set of revisions.
\item
  New content is appended after the signed byte range.
\item
  The \texttt{/Prev} chain in the trailers links all revisions together.
\item
  Extracting a specific revision means truncating at the appropriate
  \texttt{\%\%EOF} boundary.
\item
  A full rewrite destroys all incremental history and all signatures.
\item
  Linearization and incremental updates are incompatible after the first
  incremental save.
\end{itemize}

The next chapter (Chapter 9) examines AcroForm fields --- the form
infrastructure that signature fields are built on. Chapter 12 will bring
incremental updates and signature architecture together, and Chapter 13
will dissect the byte range in complete detail. If the byte range is
where signatures live or die, incremental updates are the ground they
stand on.

\chapter{AcroForm \& Interactive Forms}\label{ch09}

\section{The AcroForm Dictionary}\label{_the_acroform_dictionary}

Every PDF that contains interactive form fields has an
\texttt{/AcroForm} entry in the document catalog. This is the root of
all form infrastructure in the document: the field list, the signature
flags, the default appearance string, the default resources. If
\texttt{/AcroForm} is missing from the catalog, the document has no
interactive forms --- period. Some tools create signature fields and
forget to create the AcroForm entry, which leads to signatures that
exist as orphaned objects in the file but are invisible to every
conforming viewer.

The AcroForm dictionary lives at:

\begin{verbatim}
trailer -> /Root (Catalog) -> /AcroForm
\end{verbatim}

Here is a representative AcroForm dictionary with all keys that matter
for signature work:

\begin{verbatim}
<<
  /Fields [10 0 R 11 0 R 12 0 R]    % 
  /SigFlags 3                         % 
  /NeedAppearances false               % 
  /DA (/Helv 0 Tf 0 g)                % 
  /DR <<                               % 
    /Font << /Helv 8 0 R /ZaDb 9 0 R >>
  >>
>>
\end{verbatim}

\begin{itemize}
\item
  \texttt{/Fields} --- array of references to the top-level field
  dictionaries. Every field in the document must be reachable from this
  array (directly or through parent-child relationships).
\item
  \texttt{/SigFlags} --- a bit field that controls signature-related
  behavior. Value 3 means both bit 1 (SignaturesExist) and bit 2
  (AppendOnly) are set.
\item
  \texttt{/NeedAppearances} --- tells the viewer whether to regenerate
  appearance streams before rendering.
\item
  \texttt{/DA} --- default appearance string, used for fields that do
  not have their own \texttt{/DA}.
\item
  \texttt{/DR} --- default resources dictionary, providing fonts and
  other resources for appearance generation.
\end{itemize}

Let's examine each of these in the depth they deserve.

\subsection{/Fields --- The Root of the Field
Tree}\label{_fieldsthe_root_of_the_field_tree}

The \texttt{/Fields} array contains indirect references to the
document's top-level field dictionaries. "Top-level" means fields that
have no parent --- they are the roots of potentially nested field trees.
A flat form with ten independent fields will have ten entries in
\texttt{/Fields}. A form with a hierarchical namespace (like
\texttt{address.city}, \texttt{address.zip}) will have a single
\texttt{address} parent in \texttt{/Fields}, with \texttt{city} and
\texttt{zip} as children.

A field that is not reachable from \texttt{/Fields} does not exist as
far as a conforming reader is concerned. We have seen PDFs where a
signing tool created a signature field object, set its \texttt{/V} to
point at a signature dictionary, added a widget annotation to a
page --- and forgot to add the field to \texttt{/Fields}. The signature
data is physically present in the file. The CMS bytes are valid. But no
viewer will find it, because the AcroForm field tree does not reference
it.

If you are debugging a "missing signature" problem, check
\texttt{/Fields} first.

\subsection{/SigFlags --- The Two Bits That Protect
Everything}\label{_sigflagsthe_two_bits_that_protect_everything}

We cover \texttt{/SigFlags} in depth here because getting it wrong
destroys signatures --- silently, irreversibly, and usually long after
the signing ceremony is over.

\begin{longtable}[]{@{}
  >{\raggedright\arraybackslash}p{(\linewidth - 4\tabcolsep) * \real{0.2500}}
  >{\raggedright\arraybackslash}p{(\linewidth - 4\tabcolsep) * \real{0.1250}}
  >{\raggedright\arraybackslash}p{(\linewidth - 4\tabcolsep) * \real{0.6250}}@{}}
\caption{SigFlags bit definitions --- ISO 32000-2:2020 S12.7.2, Table
225}\tabularnewline
\toprule\noalign{}
\begin{minipage}[b]{\linewidth}\raggedright
Name
\end{minipage} & \begin{minipage}[b]{\linewidth}\raggedright
Bit
\end{minipage} & \begin{minipage}[b]{\linewidth}\raggedright
What it does
\end{minipage} \\
\midrule\noalign{}
\endfirsthead
\toprule\noalign{}
\begin{minipage}[b]{\linewidth}\raggedright
Name
\end{minipage} & \begin{minipage}[b]{\linewidth}\raggedright
Bit
\end{minipage} & \begin{minipage}[b]{\linewidth}\raggedright
What it does
\end{minipage} \\
\midrule\noalign{}
\endhead
\bottomrule\noalign{}
\endlastfoot
\texttt{SignaturesExist} & 1 & Declares that the document contains at
least one signature field with a value. Conforming readers use this as a
signal to check for signatures and display validation status. Without
it, some viewers will not even look for signatures. \\
\texttt{AppendOnly} & 2 & Mandates that any modifications to the
document \textbf{shall} use incremental updates only. This is the single
most important protection for existing signatures. If a viewer or editor
does a full save (rewriting the entire file from scratch), every
signature in the document becomes invalid --- the bytes have changed,
the byte ranges no longer match, and the cryptographic verification will
fail. \\
\end{longtable}

The only correct value for \texttt{/SigFlags} in a signed document is
\textbf{3} (both bits set). A value of 1 (SignaturesExist without
AppendOnly) is technically allowed by the spec but is a bug in
practice --- it tells viewers that signatures exist but does not protect
them from being destroyed by a full save.

Many PDF generation tools --- including some versions of LibreOffice,
older versions of iText, and various government form generators --- set
\texttt{/SigFlags} to 1 or omit it entirely. If you receive a signed PDF
with \texttt{/SigFlags} missing or set to 1, the signatures are at risk.
Any editor that performs a full save will destroy them.

AppendOnly matters because incremental updates are what make
multi-signature workflows possible --- each signer appends their changes
without moving the bytes that previous signers\textquotesingle{} byte
ranges reference. Without AppendOnly, an editor may do a full rewrite,
destroying every signature.

We have also seen the converse: \texttt{/SigFlags} is 0 but signed
fields exist with valid CMS data. Acrobat shows "No Signatures" because
bit 1 was never set, so it never looked.

The AppendOnly flag is advisory, not enforced by the file format. A
malicious or poorly written tool can ignore it. The flag protects
against \emph{accidental} destruction by well-behaved editors.
Intentional tampering is what the cryptographic signature itself guards
against.

\subsection{/NeedAppearances --- The Flag That Breaks
Signatures}\label{_needappearancesthe_flag_that_breaks_signatures}

The \texttt{/NeedAppearances} flag (boolean, default \texttt{false})
tells the viewer to regenerate appearance streams before rendering. For
ordinary fields this is a convenience --- skip generating appearances
and let the viewer handle it. For signature fields, it is a disaster.

When a viewer regenerates appearances, it modifies the document bytes.
If the document is signed, the signature becomes invalid. It does not
matter that the viewer was "just" fixing up visuals --- the bytes
changed.

If \texttt{/NeedAppearances} is \texttt{true} in a signed document,
either: (1) the viewer regenerates appearances and the signature fails,
or (2) the viewer skips regeneration and some fields show blank
rectangles. \textbf{Never set \texttt{/NeedAppearances} to \texttt{true}
in a document that will be signed.} Always generate complete appearance
streams for every field before signing.

\subsection{/DA --- Default Appearance
String}\label{_dadefault_appearance_string}

The \texttt{/DA} (Default Appearance) string is a PDF content stream
fragment that specifies the default font, size, and color for variable
text in form fields. It typically looks like:

\begin{verbatim}
/Helv 0 Tf 0 g
\end{verbatim}

This means: use the font aliased as \texttt{/Helv}, at size 0
(auto-size), in black (\texttt{0\ g}). The font name \texttt{/Helv} must
be defined in the \texttt{/DR} (Default Resources) dictionary.

Why does this matter for signatures? Because if a signing tool creates a
visible signature appearance that contains text --- the signer's name,
the date, the reason --- and the font referenced in the appearance
stream is not available in \texttt{/DR} \emph{and} not embedded in the
appearance's own \texttt{/Resources}, some viewers will fall back to
\texttt{/DA} and render the text in Helvetica or not at all.

Always embed the font resources directly in the signature field's
appearance stream \texttt{/Resources} dictionary. Do not rely on
\texttt{/DR}. The appearance stream is a self-contained Form XObject and
should carry everything it needs to render correctly, regardless of what
the parent AcroForm provides.

\subsection{/DR --- Default Resources}\label{_drdefault_resources}

The \texttt{/DR} dictionary provides fonts, color spaces, and other
resources available to all form fields. The two fonts you will encounter
in virtually every AcroForm document are \texttt{/Helv} (Helvetica, for
text) and \texttt{/ZaDb} (ZapfDingbats, for checkbox ticks and radio
button bullets).

The \texttt{/DR} fonts are \emph{not} the same as page-level fonts. A
page's \texttt{/Resources} dictionary is independent. If a signature
appearance stream references a font from the page's resources but not
from \texttt{/DR} or the appearance's own resources, rendering may fail
in viewers that strictly follow the spec.

\section{Field Tree Hierarchy}\label{_field_tree_hierarchy}

\subsection{Parent-Child Relationships and Fully Qualified
Names}\label{_parent_child_relationships_and_fully_qualified_names}

The AcroForm field tree is not flat. Fields can be nested using
\texttt{/Parent} and \texttt{/Kids} keys. A parent field
\texttt{\textless{}\textless{}\ /T\ (address)\ /Kids\ {[}20\ 0\ R\ 21\ 0\ R{]}\ \textgreater{}\textgreater{}}
with a child field
\texttt{\textless{}\textless{}\ /T\ (city)\ /Parent\ 19\ 0\ R\ /FT\ /Tx\ \textgreater{}\textgreater{}}
produces the fully qualified name \texttt{address.city}.

The fully qualified name is constructed by concatenating \texttt{/T}
(partial name) values from root to leaf, separated by periods. This is
the name you use in PDFBox's \texttt{acroForm.getField("address.city")},
in JavaScript, and in FDF data files.

\textbf{Debugging field-not-found errors:} If
\texttt{acroForm.getField("Signature1")} returns null, check whether the
field is nested. It might be \texttt{form.Signature1} or
\texttt{page1.Signature1}. Use \texttt{acroForm.getFieldTree()} to
iterate every field in the hierarchy and inspect the fully qualified
names.

\subsection{Terminal vs. Non-Terminal
Fields}\label{_terminal_vs_non_terminal_fields}

A \textbf{terminal field} is a field with no \texttt{/Kids}, or whose
\texttt{/Kids} are all widget annotations (not other fields). Terminal
fields have a \texttt{/FT} (field type) key --- \texttt{/Tx},
\texttt{/Btn}, \texttt{/Ch}, or \texttt{/Sig}.

A \textbf{non-terminal field} is a pure grouping node. It has
\texttt{/Kids} that are other fields, and it does \emph{not} have
\texttt{/FT}. Its only purpose is to provide a namespace prefix.

Some PDF generators create non-terminal fields that also carry
\texttt{/FT}. This is technically a spec violation (ISO 32000-2:2020
S12.7.4.1), but it happens. PDFBox handles it gracefully in most cases,
but other libraries may choke. If you are building a field tree, keep
non-terminal nodes clean --- no \texttt{/FT}, no \texttt{/V}, no
\texttt{/Rect}.

\section{Field Types}\label{_field_types}

ISO 32000-2:2020 S12.7.4 defines four field types. All of them can
appear in the same AcroForm, side by side, and a single document
frequently contains all four --- especially government and legal forms
that combine text input, checkboxes, dropdown selections, and signature
fields.

\subsection{/Tx --- Text Fields}\label{_txtext_fields}

Text fields hold user-entered text (value in \texttt{/V} as a text
string). Key properties: \texttt{/MaxLen} (maximum character count),
\texttt{/Q} (text alignment: 0=left, 1=center, 2=right). Text fields
interact with signatures in one important way: if a text field is filled
\emph{after} a document is signed, the fill operation is an incremental
update. Whether this invalidates the signature depends on the
\texttt{/Lock} and \texttt{/FieldMDP} settings (covered in Chapter 14).

\subsection{/Btn --- Button Fields}\label{_btnbutton_fields}

Button fields come in three flavors, distinguished by the \texttt{/Ff}
(field flags) bits: \textbf{pushbuttons} (bit 17 set, no persistent
value), \textbf{checkboxes} (bit 17 clear, bit 16 clear, two-state
toggles with values like \texttt{/Yes} or \texttt{/Off}), and
\textbf{radio buttons} (bit 17 clear, bit 16 set, mutually exclusive
options). For signatures, button fields matter primarily in the context
of form flattening and field locking.

\subsection{/Ch --- Choice Fields}\label{_chchoice_fields}

Choice fields present a list of options: \textbf{combo boxes} (bit 18
set, a dropdown) and \textbf{list boxes} (bit 18 clear, a scrollable
list). Options are defined in the \texttt{/Opt} array, where each entry
can be a simple text string or a two-element array
\texttt{{[}export\_value,\ display\_text{]}}.

\subsection{/Sig --- Signature Fields}\label{_sigsignature_fields}

Signature fields are where form infrastructure meets cryptographic
infrastructure. A \texttt{/Sig} field is, structurally, just another
form field in the AcroForm tree. But its \texttt{/V} (value) is not a
text string or a name --- it is a \textbf{signature dictionary}
(\texttt{/Type\ /Sig}), which is a complex structure containing the byte
range, the CMS data, the signing filter, timestamps, and more.

A signature field before signing:

\begin{verbatim}
<<
  /FT /Sig
  /T (Signature1)
  /V null                    % or /V absent entirely
  /Rect [100 600 300 650]
  /Subtype /Widget
  /P 5 0 R
>>
\end{verbatim}

The same field after signing:

\begin{verbatim}
<<
  /FT /Sig
  /T (Signature1)
  /V 42 0 R                  % reference to the /Sig dictionary
  /Rect [100 600 300 650]
  /Subtype /Widget
  /P 5 0 R
  /AP << /N 43 0 R >>        % appearance stream added by the signing tool
>>
\end{verbatim}

Object 42 is the signature dictionary (covered in depth in Chapter 12).
Object 43 is the visual appearance of the signature.

The signature field is \emph{not} the signature. The field is the
\textbf{container}. The signature dictionary (the \texttt{/V} value) is
the \textbf{content}. A document can have signature fields with no
value --- these are placeholder fields, pre-positioned by a template
author, waiting for a signer to fill them in.

This distinction matters because:

\begin{itemize}
\item
  A document with an empty signature field still has an entry in
  \texttt{/Fields}.
\item
  The \texttt{/SigFlags} should \emph{not} have bit 1 set until at least
  one field actually has a value.
\item
  Validation logic should check \texttt{/V}, not just the presence of
  \texttt{/FT\ /Sig} fields.
\end{itemize}

What makes \texttt{/Sig} fields unique: (1) the value is a complex
dictionary, not a simple string; (2) signing computes a hash over the
entire file, not just the field; (3) the \texttt{/SV} (Seed Value) and
\texttt{/Lock} dictionaries allow the field to constrain which
algorithms, certificates, and post-signing modifications are acceptable;
and (4) the widget's \texttt{/Rect} determines visibility --- a
zero-size rect means an invisible signature.

\section{Widget Annotations}\label{_widget_annotations}

\subsection{The Connection Between Fields and
Pages}\label{_the_connection_between_fields_and_pages}

A form field is an abstract data object (name, type, value). A widget
annotation is its visual representation on a specific page (coordinates,
size, appearance).

In the most common case, the field and widget are \textbf{merged into a
single PDF object} --- the same dictionary carries both \texttt{/FT}
(field type) and \texttt{/Subtype\ /Widget}, both \texttt{/T} (field
name) and \texttt{/Rect} (position). The spec (ISO 32000-2:2020
S12.5.6.19) also allows separating them: a field's \texttt{/Kids} array
can contain widget annotations, each with \texttt{/Parent} pointing back
to the field. This allows a single field to have widgets on multiple
pages --- unusual, but legal.

When debugging, check whether the field and widget are merged or
separate. \texttt{PDSignatureField.getWidgets()} returns the list of
widgets. An empty list means no visual representation --- fine for
invisible signatures, a bug if you expected a visible one.

\subsection{Page Association}\label{_page_association}

A widget must be associated with a page via two links: the widget's
\texttt{/P} key referencing the page dictionary, and the page's
\texttt{/Annots} array containing a reference to the widget. Both should
be present. If the widget is missing from \texttt{/Annots}, it will not
render --- even if \texttt{/P} points to the correct page.

Some signing tools add the widget to the field tree but forget to add it
to the page's \texttt{/Annots} array. The signature is valid
cryptographically but invisible on every page. This is one of the most
common "invisible signature" bugs that is not actually an invisible
signature --- it is a missing annotation reference.

\section{Appearance Streams}\label{_appearance_streams}

\subsection{Normal, Rollover, and Down}\label{_normal_rollover_and_down}

The \texttt{/AP} dictionary can define three appearance streams:
\texttt{/N} (Normal --- the only one that matters for signatures),
\texttt{/R} (Rollover --- mouse hover, rarely used), and \texttt{/D}
(Down --- click state, irrelevant for signatures).

Each can be a single Form XObject stream or a \textbf{subdictionary}
mapping state names to streams. Checkboxes use a subdictionary
(\texttt{\textless{}\textless{}\ /Yes\ 60\ 0\ R\ /Off\ 61\ 0\ R\ \textgreater{}\textgreater{}}).
Signature fields always use a single stream --- signatures do not have
states.

\subsection{Why Appearance Streams Matter for
Signatures}\label{_why_appearance_streams_matter_for_signatures}

If the appearance stream is missing, the viewer either renders a blank
rectangle, tries to generate one on the fly (dangerous if the document
is signed --- see \texttt{/NeedAppearances} above), or shows a generic
placeholder. For signed fields, the only correct approach is to have the
signing tool generate a complete appearance stream before signing.

The appearance stream is included in the signed byte range. Modifying it
after signing invalidates the signature. However, the appearance is
\emph{not} semantically bound to the CMS content --- the appearance
could say "Signed by Alice" while the certificate belongs to Bob. The
viewer must verify both independently.

\subsection{Generating Appearances with
PDFBox}\label{_generating_appearances_with_pdfbox}

PDFBox provides the \texttt{PDAcroForm.refreshAppearances()} method (and
its per-field equivalent) to generate appearance streams for all fields.
This is the equivalent of setting \texttt{/NeedAppearances} to
\texttt{true} and having the viewer regenerate --- except you do it
programmatically, before signing, so the appearances are part of the
signed bytes.

\begin{Shaded}
\begin{Highlighting}[]
\ControlFlowTok{try} \OperatorTok{(}\NormalTok{PDDocument doc }\OperatorTok{=}\NormalTok{ Loader}\OperatorTok{.}\FunctionTok{loadPDF}\OperatorTok{(}\NormalTok{pdfBytes}\OperatorTok{))} \OperatorTok{\{}
\NormalTok{    PDAcroForm acroForm }\OperatorTok{=}\NormalTok{ doc}\OperatorTok{.}\FunctionTok{getDocumentCatalog}\OperatorTok{().}\FunctionTok{getAcroForm}\OperatorTok{();}
    \ControlFlowTok{if} \OperatorTok{(}\NormalTok{acroForm }\OperatorTok{!=} \KeywordTok{null}\OperatorTok{)} \OperatorTok{\{}
\NormalTok{        acroForm}\OperatorTok{.}\FunctionTok{setNeedAppearances}\OperatorTok{(}\KeywordTok{false}\OperatorTok{);}  \CommentTok{// clear the flag}
\NormalTok{        acroForm}\OperatorTok{.}\FunctionTok{refreshAppearances}\OperatorTok{();}        \CommentTok{// generate all appearances}
    \OperatorTok{\}}
\OperatorTok{\}}
\end{Highlighting}
\end{Shaded}

\texttt{refreshAppearances()} will overwrite existing appearance
streams. If you have custom signature appearances (hand-drawn images,
company logos), call this \emph{before} setting up your custom
appearances, or call it only on non-signature fields. For signature
appearances, most signing workflows build the appearance stream manually
using a \texttt{PDFormXObject} --- Chapter 15 covers this in detail.

\section{XFA Forms: The Deprecated
Complication}\label{_xfa_forms_the_deprecated_complication}

\subsection{What XFA Was and Why It Was
Deprecated}\label{_what_xfa_was_and_why_it_was_deprecated}

XFA (XML Forms Architecture) was Adobe's attempt to replace the limited
AcroForm system with a powerful XML-based form description language. An
XFA form embeds XML data inside the PDF file (in the \texttt{/XFA} key
of the AcroForm dictionary), and a viewer that supports XFA renders the
form from XML rather than from the AcroForm field tree. XFA could do
things AcroForm could not: dynamic layouts, complex validation rules,
repeating sections, and data binding to external schemas.

ISO 32000-2:2020 deprecated XFA for three reasons: (1) only Adobe
implemented it fully --- interoperability was abysmal; (2) XFA and
AcroForm coexist awkwardly in the same file, with the two
representations frequently disagreeing about field values; and (3) XFA's
scripting capabilities have been exploited in multiple CVEs.

\subsection{Why XFA Still Causes
Problems}\label{_why_xfa_still_causes_problems}

Deprecation does not mean disappearance. Millions of XFA-based PDF forms
are still in circulation --- government tax forms, insurance
applications, healthcare records. When these documents are signed, the
signature covers the entire file, including the XFA data.

The problems:

\begin{itemize}
\item
  \textbf{PDFBox does not support XFA rendering.} If you load an XFA
  document in PDFBox, you get the AcroForm fallback. The field values
  you see may differ from what an XFA-aware viewer shows. If you sign
  based on the AcroForm representation, the signature covers bytes that
  an XFA viewer interprets differently.
\item
  \textbf{Field names may not match.} XFA and AcroForm can use different
  naming conventions. A signature field named
  \texttt{form1{[}0{]}.\#subform{[}0{]}.SignatureField1{[}0{]}} in XFA
  might be \texttt{SignatureField1} in the AcroForm fallback.
\item
  \textbf{Validation ambiguity.} If the XFA representation says one
  thing and the signed AcroForm says another, which one did the signer
  actually see and intend to sign?
\end{itemize}

If you encounter an XFA document in a signing workflow, the safest
approach is to flatten the XFA to pure AcroForm (removing the
\texttt{/XFA} key and ensuring all field values are represented in the
AcroForm structure) \emph{before} any signatures are applied. PDFBox's
\texttt{PDAcroForm.flatten()} can help, but test thoroughly --- XFA
flattening is lossy by nature.

\section{Form Flattening and
Signatures}\label{_form_flattening_and_signatures}

Form flattening converts interactive fields into static page content:
field values are "burned in" to page content streams, widget annotations
are removed, and the AcroForm structure is deleted. The result is a
static document that looks like the filled-in form but cannot be edited.

Flattening destroys every signature in the document. The signature
fields themselves are removed (they are form fields), the byte ranges no
longer point at the original bytes, and the AcroForm structure
referencing the signatures is gone.

\textbf{Never flatten a signed document.} This sounds obvious, but it
happens constantly in automated workflows --- sign, then flatten "for
archival" or "to prevent editing." The result is a static PDF with no
valid signatures. If you need to prevent editing after signing, use
field locking (\texttt{/Lock}) and certification signatures
(\texttt{/DocMDP}) --- not flattening.

\section{PDFBox AcroForm API}\label{_pdfbox_acroform_api}

\subsection{PDAcroForm: The Entry
Point}\label{_pdacroform_the_entry_point}

In PDFBox 3.0.0, the AcroForm is accessed through the document catalog:

\begin{Shaded}
\begin{Highlighting}[]
\ControlFlowTok{try} \OperatorTok{(}\NormalTok{PDDocument doc }\OperatorTok{=}\NormalTok{ Loader}\OperatorTok{.}\FunctionTok{loadPDF}\OperatorTok{(}\NormalTok{pdfBytes}\OperatorTok{))} \OperatorTok{\{}
\NormalTok{    PDAcroForm acroForm }\OperatorTok{=}\NormalTok{ doc}\OperatorTok{.}\FunctionTok{getDocumentCatalog}\OperatorTok{().}\FunctionTok{getAcroForm}\OperatorTok{();}
    \ControlFlowTok{if} \OperatorTok{(}\NormalTok{acroForm }\OperatorTok{==} \KeywordTok{null}\OperatorTok{)} \OperatorTok{\{}
        \ControlFlowTok{return}\OperatorTok{;}  \CommentTok{// No interactive forms in this document}
    \OperatorTok{\}}
    \ControlFlowTok{for} \OperatorTok{(}\NormalTok{PDField field }\OperatorTok{:}\NormalTok{ acroForm}\OperatorTok{.}\FunctionTok{getFields}\OperatorTok{())} \OperatorTok{\{}
\NormalTok{        log}\OperatorTok{.}\FunctionTok{debug}\OperatorTok{(}\StringTok{"Field: \{\} (type: \{\}, fqn: \{\})"}\OperatorTok{,}
\NormalTok{            field}\OperatorTok{.}\FunctionTok{getPartialName}\OperatorTok{(),}\NormalTok{ field}\OperatorTok{.}\FunctionTok{getFieldType}\OperatorTok{(),}
\NormalTok{            field}\OperatorTok{.}\FunctionTok{getFullyQualifiedName}\OperatorTok{());}
    \OperatorTok{\}}
\OperatorTok{\}}
\end{Highlighting}
\end{Shaded}

\texttt{acroForm.getFields()} returns only top-level fields. To traverse
the full tree including children, use \texttt{acroForm.getFieldTree()},
which returns a flat iterable over every field in the hierarchy.

\subsection{Field Tree Traversal}\label{_field_tree_traversal}

\texttt{acroForm.getFieldTree()} returns a flat iterable over every
field in the hierarchy --- almost always what you want. For finding
signature fields specifically:

\begin{Shaded}
\begin{Highlighting}[]
\ControlFlowTok{for} \OperatorTok{(}\NormalTok{PDField field }\OperatorTok{:}\NormalTok{ acroForm}\OperatorTok{.}\FunctionTok{getFieldTree}\OperatorTok{())} \OperatorTok{\{}
    \ControlFlowTok{if} \OperatorTok{(}\NormalTok{field }\KeywordTok{instanceof}\NormalTok{ PDSignatureField sigField}\OperatorTok{)} \OperatorTok{\{}
\NormalTok{        log}\OperatorTok{.}\FunctionTok{info}\OperatorTok{(}\StringTok{"Signature field: \{\} (signed: \{\})"}\OperatorTok{,}
\NormalTok{            sigField}\OperatorTok{.}\FunctionTok{getFullyQualifiedName}\OperatorTok{(),}
\NormalTok{            sigField}\OperatorTok{.}\FunctionTok{getSignature}\OperatorTok{()} \OperatorTok{!=} \KeywordTok{null}\OperatorTok{);}
    \OperatorTok{\}}
\OperatorTok{\}}
\end{Highlighting}
\end{Shaded}

If you need tree depth or explicit parent-child relationships, use
\texttt{PDNonTerminalField.getChildren()} and recurse manually.

\subsection{PDSignatureField and
PDSignature}\label{_pdsignaturefield_and_pdsignature}

\texttt{PDSignatureField} is the PDFBox class representing a
\texttt{/Sig} field. It provides access to the signature value (the
\texttt{/V} dictionary) through \texttt{getSignature()}, which returns a
\texttt{PDSignature} or \texttt{null} if the field is unsigned.

\begin{Shaded}
\begin{Highlighting}[]
\NormalTok{PDSignatureField sigField }\OperatorTok{=} \OperatorTok{(}\NormalTok{PDSignatureField}\OperatorTok{)}\NormalTok{ acroForm}\OperatorTok{.}\FunctionTok{getField}\OperatorTok{(}\StringTok{"Signature1"}\OperatorTok{);}
\ControlFlowTok{if} \OperatorTok{(}\NormalTok{sigField }\OperatorTok{==} \KeywordTok{null}\OperatorTok{)} \OperatorTok{\{} \ControlFlowTok{return}\OperatorTok{;} \OperatorTok{\}}  \CommentTok{// field not found {-}{-} check fully qualified name}

\NormalTok{PDSignature sig }\OperatorTok{=}\NormalTok{ sigField}\OperatorTok{.}\FunctionTok{getSignature}\OperatorTok{();}
\ControlFlowTok{if} \OperatorTok{(}\NormalTok{sig }\OperatorTok{==} \KeywordTok{null}\OperatorTok{)} \OperatorTok{\{} \ControlFlowTok{return}\OperatorTok{;} \OperatorTok{\}}       \CommentTok{// field exists but is unsigned}

\NormalTok{log}\OperatorTok{.}\FunctionTok{info}\OperatorTok{(}\StringTok{"Filter: \{\}"}\OperatorTok{,}\NormalTok{ sig}\OperatorTok{.}\FunctionTok{getFilter}\OperatorTok{());}       \CommentTok{// e.g. "Adobe.PPKLite"}
\NormalTok{log}\OperatorTok{.}\FunctionTok{info}\OperatorTok{(}\StringTok{"SubFilter: \{\}"}\OperatorTok{,}\NormalTok{ sig}\OperatorTok{.}\FunctionTok{getSubFilter}\OperatorTok{());} \CommentTok{// e.g. "adbe.pkcs7.detached"}
\NormalTok{log}\OperatorTok{.}\FunctionTok{info}\OperatorTok{(}\StringTok{"Name: \{\}"}\OperatorTok{,}\NormalTok{ sig}\OperatorTok{.}\FunctionTok{getName}\OperatorTok{());}
\NormalTok{log}\OperatorTok{.}\FunctionTok{info}\OperatorTok{(}\StringTok{"Sign date: \{\}"}\OperatorTok{,}\NormalTok{ sig}\OperatorTok{.}\FunctionTok{getSignDate}\OperatorTok{());}
\end{Highlighting}
\end{Shaded}

\texttt{PDAcroForm.getField(String\ name)} uses the \textbf{fully
qualified name}. If the field is nested (e.g.,
\texttt{signatures.approval}), you must pass the full path. Passing just
\texttt{approval} will return \texttt{null}. This is one of the most
common sources of "field not found" bugs.

\subsection{Creating Signature Fields
Programmatically}\label{_creating_signature_fields_programmatically}

Creating a signature field involves three steps: create the field,
create the widget annotation, and wire them into the AcroForm and page
structure.

\begin{Shaded}
\begin{Highlighting}[]
\CommentTok{// Ensure the AcroForm exists}
\NormalTok{PDAcroForm acroForm }\OperatorTok{=}\NormalTok{ doc}\OperatorTok{.}\FunctionTok{getDocumentCatalog}\OperatorTok{().}\FunctionTok{getAcroForm}\OperatorTok{();}
\ControlFlowTok{if} \OperatorTok{(}\NormalTok{acroForm }\OperatorTok{==} \KeywordTok{null}\OperatorTok{)} \OperatorTok{\{}
\NormalTok{    acroForm }\OperatorTok{=} \KeywordTok{new} \FunctionTok{PDAcroForm}\OperatorTok{(}\NormalTok{doc}\OperatorTok{);}
\NormalTok{    doc}\OperatorTok{.}\FunctionTok{getDocumentCatalog}\OperatorTok{().}\FunctionTok{setAcroForm}\OperatorTok{(}\NormalTok{acroForm}\OperatorTok{);}
\OperatorTok{\}}

\CommentTok{// Create the signature field and configure its widget}
\NormalTok{PDSignatureField sigField }\OperatorTok{=} \KeywordTok{new} \FunctionTok{PDSignatureField}\OperatorTok{(}\NormalTok{acroForm}\OperatorTok{);}
\NormalTok{sigField}\OperatorTok{.}\FunctionTok{setPartialName}\OperatorTok{(}\StringTok{"Signature1"}\OperatorTok{);}
\NormalTok{PDAnnotationWidget widget }\OperatorTok{=}\NormalTok{ sigField}\OperatorTok{.}\FunctionTok{getWidgets}\OperatorTok{().}\FunctionTok{get}\OperatorTok{(}\DecValTok{0}\OperatorTok{);}
\NormalTok{PDPage targetPage }\OperatorTok{=}\NormalTok{ doc}\OperatorTok{.}\FunctionTok{getPage}\OperatorTok{(}\DecValTok{0}\OperatorTok{);}
\NormalTok{widget}\OperatorTok{.}\FunctionTok{setRectangle}\OperatorTok{(}\KeywordTok{new} \FunctionTok{PDRectangle}\OperatorTok{(}\DecValTok{100}\OperatorTok{,} \DecValTok{600}\OperatorTok{,} \DecValTok{200}\OperatorTok{,} \DecValTok{50}\OperatorTok{));}
\NormalTok{widget}\OperatorTok{.}\FunctionTok{setPage}\OperatorTok{(}\NormalTok{targetPage}\OperatorTok{);}

\CommentTok{// Wire into AcroForm and page structure}
\NormalTok{targetPage}\OperatorTok{.}\FunctionTok{getAnnotations}\OperatorTok{().}\FunctionTok{add}\OperatorTok{(}\NormalTok{widget}\OperatorTok{);}
\BuiltInTok{List}\OperatorTok{\textless{}}\NormalTok{PDField}\OperatorTok{\textgreater{}}\NormalTok{ fields }\OperatorTok{=} \KeywordTok{new} \BuiltInTok{ArrayList}\OperatorTok{\textless{}\textgreater{}(}\NormalTok{acroForm}\OperatorTok{.}\FunctionTok{getFields}\OperatorTok{());}
\NormalTok{fields}\OperatorTok{.}\FunctionTok{add}\OperatorTok{(}\NormalTok{sigField}\OperatorTok{);}
\NormalTok{acroForm}\OperatorTok{.}\FunctionTok{setFields}\OperatorTok{(}\NormalTok{fields}\OperatorTok{);}
\NormalTok{acroForm}\OperatorTok{.}\FunctionTok{getCOSObject}\OperatorTok{().}\FunctionTok{setInt}\OperatorTok{(}\NormalTok{COSName}\OperatorTok{.}\FunctionTok{SIG\_FLAGS}\OperatorTok{,} \DecValTok{3}\OperatorTok{);}
\end{Highlighting}
\end{Shaded}

\textbf{Common mistakes when creating signature fields:}

\begin{enumerate}
\def\labelenumi{\arabic{enumi}.}
\item
  \textbf{Forgetting to add the widget to the page's \texttt{/Annots}
  array.} The field exists in the AcroForm but is invisible.
\item
  \textbf{Forgetting to add the field to \texttt{/Fields}.} The widget
  is on the page but no viewer considers it a form field.
\item
  \textbf{Not setting \texttt{/SigFlags}.} The field is there, it gets
  signed, but editors may do a full save and destroy it.
\item
  \textbf{Using \texttt{new\ PDRectangle(x,\ y,\ w,\ h)} vs. raw array.}
  PDFBox's \texttt{PDRectangle} constructor takes (x, y, width, height),
  but the raw PDF \texttt{/Rect} array is
  \texttt{{[}llx\ lly\ urx\ ury{]}}. PDFBox handles the conversion, but
  if you're debugging raw COS dictionaries, remember the coordinate
  difference.
\end{enumerate}

\section{Common Form/Signature Interaction
Bugs}\label{_common_formsignature_interaction_bugs}

Years of debugging have revealed a predictable set of bugs that occur at
the intersection of AcroForm infrastructure and signature functionality.
These are not edge cases --- they are the everyday reality of working
with signed PDF forms.

\subsection{Missing Appearance
Streams}\label{_missing_appearance_streams}

\textbf{Symptom:} Signature validates, but the rectangle on the page is
blank. \textbf{Cause:} No \texttt{/AP} entry on the widget, or
\texttt{/NeedAppearances} is \texttt{true} and the viewer destroyed the
original appearance during regeneration. \textbf{Fix:} Always generate a
complete \texttt{PDAppearanceStream} during signing and set it as the
normal appearance on the widget (see Chapter 15 for details).

\subsection{Wrong Field Hierarchy}\label{_wrong_field_hierarchy}

\textbf{Symptom:} PDFBox finds the field but Acrobat does not, or vice
versa. \textbf{Cause:} Inconsistent parent-child chain --- a child
references a parent that does not list it in \texttt{/Kids}, or the
fully qualified name computed by concatenating \texttt{/T} values does
not match expectations. \textbf{Fix:} Validate the tree bidirectionally:
every child's \texttt{/Parent} must list it in \texttt{/Kids}, and vice
versa.

\subsection{SigFlags Not Set or Wrong}\label{_sigflags_not_set_or_wrong}

\textbf{Symptom:} Signatures are valid, but a subsequent edit does a
full save and destroys them all. \textbf{Cause:} \texttt{/SigFlags} is
0, 1, or absent. \textbf{Fix:} Set \texttt{/SigFlags} to 3 when creating
the first signature field --- before signing, not after.

\subsection{Orphaned Signature Fields}\label{_orphaned_signature_fields}

\textbf{Symptom:} Signature panel shows a signature with no visual
element on any page; or shows nothing despite signature data existing in
the file. \textbf{Cause:} Widget is missing from a page's
\texttt{/Annots} (variant 1), or the field is not referenced from
\texttt{/Fields} (variant 2). \textbf{Fix:} Verify both paths: field →
AcroForm \texttt{/Fields}, and widget → page \texttt{/Annots}.

\subsection{Acrobat vs. PDFBox Rendering
Differences}\label{_acrobat_vs_pdfbox_rendering_differences}

\textbf{Symptom:} A form that looks perfect in Acrobat renders
differently in PDFBox (or vice versa).

\textbf{Common causes:} Acrobat aggressively substitutes missing fonts
(PDFBox does not --- it uses only what is embedded or available on the
system). Acrobat regenerates appearances when \texttt{/NeedAppearances}
is \texttt{true} (PDFBox requires an explicit
\texttt{refreshAppearances()} call). Acrobat is lenient about malformed
COS dictionaries (a field with \texttt{/FT} on a non-terminal node may
work in Acrobat but cause unexpected behavior in PDFBox). And Acrobat
interprets \texttt{/DA} strings flexibly --- font size 0 means
"auto-size," missing color operators are silently defaulted --- while
PDFBox parses more strictly.

When debugging cross-viewer differences, compare the raw COS structure.
In PDFBox: \texttt{field.getCOSObject().toString()}. In Acrobat:
Preflight tool or JavaScript console
(\texttt{this.getField("name").valueAsString}). The differences are
almost always at the COS level.

\subsection{Real-World Scenario: The Government Tax
Form}\label{_real_world_scenario_the_government_tax_form}

A government tax authority distributes a PDF form designed in Adobe
LiveCycle Designer. The form uses XFA internally with an AcroForm
fallback. A citizen fills it in using Chrome's built-in PDF viewer
(which uses the AcroForm fallback because it does not support XFA). The
citizen's tax advisor opens the filled form in Acrobat (which
\emph{does} support XFA), sees different field values because the XFA
data was never updated, and signs the document.

The signature is technically valid --- it covers the bytes that were
present at signing time. But the signed content includes XFA data that
disagrees with the AcroForm data. There is no clean fix short of
removing XFA from the document entirely. This is precisely why ISO
32000-2:2020 deprecated XFA, and why the deprecation matters for
real-world signature workflows.

\chapter{PDF Coordinate System}\label{_pdf_coordinate_system}

The PDF coordinate system is the single most common source of signature
placement bugs. It is conceptually simple --- an origin at the
bottom-left corner, units in points --- but it collides head-on with how
every screen coordinate system works, how page rotation transforms
visual layout without changing coordinates, and how annotation
rectangles encode position in a way that rewards precision and punishes
guessing.

If you have ever placed a signature at what you thought was the bottom
of a page and found it at the top, or created a perfectly valid
signature that appeared as a mirror image on a rotated page, or spent an
afternoon trying to figure out why your annotation rectangle produced a
zero-height invisible box --- this chapter is for you.

We will cover the coordinate system from first principles, work through
every page box that ISO 32000-2:2020 defines, explain how page rotation
interacts with coordinates, show how the Current Transformation Matrix
works, and then tie it all back to signature placement with practical
PDFBox code and debugging techniques.

\section{User Space: The Default Coordinate
System}\label{_user_space_the_default_coordinate_system}

Every PDF page has a \textbf{user space} coordinate system defined in
ISO 32000-2:2020 Section 8.3.2. The rules are few but non-negotiable:

\begin{itemize}
\item
  The \textbf{origin} is at the \textbf{bottom-left} corner of the page.
\item
  The \textbf{x-axis} increases to the \textbf{right}.
\item
  The \textbf{y-axis} increases \textbf{upward}.
\item
  One unit equals \textbf{1/72 of an inch} --- which is exactly one
  \textbf{typographic point}.
\end{itemize}

That last property is not a coincidence. PDF was born as a page
description language for printing, and the typographic point has been
the standard unit of measure in the printing industry since the 18th
century. When the specification says a page is 612 units wide and 792
units tall, that's 8.5 by 11 inches --- US Letter size. A4 paper (the
ISO 216 standard used nearly everywhere outside the US) is 595.28 by
841.89 points (210 by 297 millimetres).

The 1/72-inch definition is a convenient approximation. The traditional
printer's point (the Didot point, the Pica point) varies slightly
depending on which historical tradition you follow. PDF sidesteps this
ambiguity by defining one user space unit as \emph{exactly} 1/72 inch,
period. ISO 32000-2:2020 Section 8.3.2.3 makes this explicit.

The bottom-left origin is where most coordinate confusion begins. Every
windowing system you have used --- HTML, Java Swing, JavaFX, Android,
iOS, Win32 --- places the origin at the \textbf{top-left} corner with
the y-axis increasing \textbf{downward}. We are trained to think of y=0
as the top. In PDF, y=0 is the bottom.

This means that when you want to place a signature 50 points from the
top of a US Letter page, the y-coordinate is not 50 --- it is 792 minus
50 minus the height of the signature rectangle. We will return to this
conversion formula repeatedly throughout this chapter because it is the
root cause of the most common signature positioning bugs.

\section{Page Boxes: The Five
Rectangles}\label{_page_boxes_the_five_rectangles}

A PDF page does not have a single "size." It has up to five nested
rectangles, each serving a different purpose. ISO 32000-2:2020 Section
14.11.2 defines them all. Understanding which box matters for signature
placement --- and which boxes most documents actually define ---
prevents an entire category of positioning errors.

\subsection{MediaBox (Required)}\label{_mediabox_required}

The \textbf{MediaBox} defines the physical medium --- the full extent of
the page as it would be printed or rendered. It is the only required
page box; every other box defaults to the MediaBox if absent.

For a US Letter page:

\begin{verbatim}
/MediaBox [0 0 612 792]
\end{verbatim}

The array is \texttt{{[}llx\ lly\ urx\ ury{]}} --- lower-left x,
lower-left y, upper-right x, upper-right y. The MediaBox almost always
starts at the origin \texttt{{[}0\ 0\ …\hspace{0pt}{]}}, but it does not
have to. A MediaBox of \texttt{{[}100\ 100\ 712\ 892{]}} is perfectly
legal and defines the same 612-by-792 point area, just shifted. Most
tools normalize to zero, but when debugging coordinates, never assume.

\subsection{CropBox}\label{_cropbox}

The \textbf{CropBox} defines the region to which the page contents are
clipped when displayed or printed. Think of it as the visible window
into the MediaBox. If absent, the CropBox defaults to the MediaBox.

A common use case: a document designed for A4 paper with crop marks. The
MediaBox includes the crop marks (larger than the final page), and the
CropBox clips to the actual content area.

For signature placement, the CropBox is the relevant boundary in most
viewers. Adobe Acrobat positions annotations relative to the CropBox,
not the MediaBox, when they differ.

If you compute signature coordinates against the MediaBox but the
document has a smaller CropBox, your signature may appear to be shifted
or partially off-screen. Always check \texttt{page.getCropBox()} in
PDFBox and use its dimensions as the reference frame.

\subsection{BleedBox}\label{_bleedbox}

The \textbf{BleedBox} defines the region to which the page contents are
clipped when output in a production (printing) environment. It is
slightly larger than the TrimBox to allow for bleed --- the area where
ink extends beyond the final cut edge to avoid white margins from
imprecise cutting.

The BleedBox is relevant for professional printing workflows and has no
meaningful impact on digital signature placement. We mention it for
completeness.

\subsection{TrimBox}\label{_trimbox}

The \textbf{TrimBox} defines the intended final dimensions of the page
after trimming. In a print-production workflow, the TrimBox is the
"real" page boundary --- the finished product.

Like the BleedBox, the TrimBox is irrelevant for signature placement in
digital workflows. However, if you are building a signing service that
handles print-ready PDFs, be aware that the TrimBox may be significantly
smaller than the MediaBox, and your users\textquotesingle{} mental model
of "the page" aligns with the TrimBox, not the MediaBox.

\subsection{ArtBox}\label{_artbox}

The \textbf{ArtBox} defines the extent of the page's "meaningful
content" as intended by the author. It is rarely used in practice and
has no impact on signature placement.

\subsection{Which Box Matters for
Signatures?}\label{_which_box_matters_for_signatures}

For digital signature placement, the decision tree is simple:

\begin{enumerate}
\def\labelenumi{\arabic{enumi}.}
\item
  \textbf{Use the CropBox} as your reference frame. In PDFBox,
  \texttt{page.getCropBox()} returns the CropBox if defined, otherwise
  it falls back to the MediaBox. This is the correct default.
\item
  If you need to verify that your annotation rectangle falls within the
  visible area, check against the CropBox.
\item
  Never assume the CropBox starts at \texttt{{[}0\ 0\ …\hspace{0pt}{]}}.
  A CropBox of \texttt{{[}50\ 50\ 562\ 742{]}} shifts the visible
  origin.
\end{enumerate}

PDFBox's \texttt{PDPage.getCropBox()} handles the fallback-to-MediaBox
logic automatically. You almost never need to call
\texttt{getMediaBox()} directly for positioning purposes. The one
exception is when you need to know the physical page dimensions for
coordinate translation on pages with a non-zero CropBox origin.

\section{Annotation Rectangles}\label{_annotation_rectangles}

Every annotation --- including signature widget annotations --- has a
\texttt{/Rect} entry that specifies its position on the page. The
\texttt{/Rect} is an array of four numbers:

\begin{verbatim}
/Rect [x1 y1 x2 y2]
\end{verbatim}

Where \texttt{{[}x1\ y1{]}} is the \textbf{lower-left} corner and
\texttt{{[}x2\ y2{]}} is the \textbf{upper-right} corner, both in user
space coordinates relative to the page's coordinate system.

For a signature that should appear 100 points from the left edge, 50
points from the bottom edge, and is 200 points wide by 60 points tall:

\begin{verbatim}
/Rect [100 50 300 110]
\end{verbatim}

That gives us: lower-left at (100, 50), upper-right at (100+200, 50+60)
= (300, 110).

The convention is lower-left to upper-right. If you accidentally swap
these --- providing upper-left to lower-right, for example --- some
viewers will silently normalize the rectangle, while others will render
a zero-area or inverted annotation. PDFBox's \texttt{PDRectangle}
constructor takes \texttt{(x,\ y,\ width,\ height)} where
\texttt{(x,\ y)} is the lower-left corner, so it handles this correctly
if you feed it the right values.

\section{Why Signature Positioning Goes
Wrong}\label{_why_signature_positioning_goes_wrong}

Here is the bug that every developer writes exactly once. We want a
signature at the visual top-right corner of a US Letter page. In screen
coordinates (top-left origin), that would be approximately
\texttt{(400,\ 30)} for a rectangle starting near the top-right. So we
write:

\begin{Shaded}
\begin{Highlighting}[]
\NormalTok{PDRectangle rect }\OperatorTok{=} \KeywordTok{new} \FunctionTok{PDRectangle}\OperatorTok{(}\DecValTok{400}\OperatorTok{,} \DecValTok{30}\OperatorTok{,} \DecValTok{200}\OperatorTok{,} \DecValTok{50}\OperatorTok{);}
\end{Highlighting}
\end{Shaded}

And the signature appears at the \textbf{bottom-right} of the page.
Because y=30 in PDF user space is 30 points from the \emph{bottom}, not
30 points from the top.

The correct conversion is:

\begin{Shaded}
\begin{Highlighting}[]
\DataTypeTok{float}\NormalTok{ pageHeight }\OperatorTok{=}\NormalTok{ page}\OperatorTok{.}\FunctionTok{getCropBox}\OperatorTok{().}\FunctionTok{getHeight}\OperatorTok{();}
\DataTypeTok{float}\NormalTok{ desiredTopOffset }\OperatorTok{=} \DecValTok{30}\OperatorTok{;}    \CommentTok{// 30 points from the visual top}
\DataTypeTok{float}\NormalTok{ signatureHeight }\OperatorTok{=} \DecValTok{50}\OperatorTok{;}

\DataTypeTok{float}\NormalTok{ pdfY }\OperatorTok{=}\NormalTok{ pageHeight }\OperatorTok{{-}}\NormalTok{ desiredTopOffset }\OperatorTok{{-}}\NormalTok{ signatureHeight}\OperatorTok{;}
\NormalTok{PDRectangle rect }\OperatorTok{=} \KeywordTok{new} \FunctionTok{PDRectangle}\OperatorTok{(}\DecValTok{400}\OperatorTok{,}\NormalTok{ pdfY}\OperatorTok{,} \DecValTok{200}\OperatorTok{,}\NormalTok{ signatureHeight}\OperatorTok{);}
\end{Highlighting}
\end{Shaded}

The formula is:
\texttt{pdfY\ =\ pageHeight\ -\ topOffset\ -\ rectangleHeight}.

This is the single most important formula in this chapter. Memorize it,
put it in a utility method, and never compute it inline.

If your signing UI lets users click a position on a rendered page, the
click coordinates are in screen space (origin top-left). You must invert
the y-axis before creating the PDF annotation rectangle. Failing to do
so produces signatures that appear to be vertically "mirrored" --- a
signature placed at the top appears at the bottom, and vice versa.

\section{Page Rotation}\label{_page_rotation}

The \texttt{/Rotate} key in a page dictionary specifies a rotation angle
that viewers should apply when displaying the page. Legal values are 0,
90, 180, and 270 (degrees, clockwise).

Here is the critical insight that trips up even experienced developers:
\textbf{page rotation does not change the coordinate system}. It changes
how the viewer \emph{displays} the coordinates, but the coordinates
themselves --- and the annotation rectangles --- remain in the original,
unrotated space.

Consider a US Letter page with \texttt{/Rotate\ 90}. The viewer displays
it in landscape orientation: 792 points wide and 612 points tall. But
the MediaBox is still \texttt{{[}0\ 0\ 612\ 792{]}} and the coordinate
system still has its origin at the bottom-left corner of the
\emph{unrotated} page.

This means that if you compute an annotation rectangle for a
90-degree-rotated page using the visual (rotated) dimensions, your
coordinates will be wrong. The annotation will appear in a completely
unexpected location because the viewer applies the rotation \emph{after}
positioning the annotation in the unrotated coordinate space.

\subsection{Rotation and Signature
Placement}\label{_rotation_and_signature_placement}

To place a signature at a specific visual position on a rotated page, we
must think in the \emph{unrotated} coordinate space and account for the
rotation.

For a page rotated 90 degrees clockwise:

\begin{itemize}
\item
  The visual width (what the user sees) corresponds to the unrotated
  height (792 for US Letter).
\item
  The visual height corresponds to the unrotated width (612 for US
  Letter).
\item
  A point at visual position (vx, vy) --- where vy is measured from the
  visual top --- maps to unrotated coordinates (vy\_from\_bottom, vx)
  approximately, depending on the rotation direction.
\end{itemize}

The general approach in PDFBox:

\begin{Shaded}
\begin{Highlighting}[]
\NormalTok{PDRectangle cropBox }\OperatorTok{=}\NormalTok{ page}\OperatorTok{.}\FunctionTok{getCropBox}\OperatorTok{();}
\DataTypeTok{int}\NormalTok{ rotation }\OperatorTok{=}\NormalTok{ page}\OperatorTok{.}\FunctionTok{getRotation}\OperatorTok{();}

\DataTypeTok{float}\NormalTok{ pageWidth}\OperatorTok{,}\NormalTok{ pageHeight}\OperatorTok{;}
\ControlFlowTok{if} \OperatorTok{(}\NormalTok{rotation }\OperatorTok{==} \DecValTok{90} \OperatorTok{||}\NormalTok{ rotation }\OperatorTok{==} \DecValTok{270}\OperatorTok{)} \OperatorTok{\{}
    \CommentTok{// Visual dimensions are swapped}
\NormalTok{    pageWidth }\OperatorTok{=}\NormalTok{ cropBox}\OperatorTok{.}\FunctionTok{getHeight}\OperatorTok{();}
\NormalTok{    pageHeight }\OperatorTok{=}\NormalTok{ cropBox}\OperatorTok{.}\FunctionTok{getWidth}\OperatorTok{();}
\OperatorTok{\}} \ControlFlowTok{else} \OperatorTok{\{}
\NormalTok{    pageWidth }\OperatorTok{=}\NormalTok{ cropBox}\OperatorTok{.}\FunctionTok{getWidth}\OperatorTok{();}
\NormalTok{    pageHeight }\OperatorTok{=}\NormalTok{ cropBox}\OperatorTok{.}\FunctionTok{getHeight}\OperatorTok{();}
\OperatorTok{\}}

\CommentTok{// Now compute visual position}
\DataTypeTok{float}\NormalTok{ visualX }\OperatorTok{=} \DecValTok{400}\OperatorTok{;}                          \CommentTok{// from left of visual page}
\DataTypeTok{float}\NormalTok{ visualTopOffset }\OperatorTok{=} \DecValTok{30}\OperatorTok{;}                   \CommentTok{// from top of visual page}
\DataTypeTok{float}\NormalTok{ sigWidth }\OperatorTok{=} \DecValTok{200}\OperatorTok{;}
\DataTypeTok{float}\NormalTok{ sigHeight }\OperatorTok{=} \DecValTok{50}\OperatorTok{;}

\CommentTok{// Convert to PDF coordinates depending on rotation}
\DataTypeTok{float}\NormalTok{ pdfX}\OperatorTok{,}\NormalTok{ pdfY}\OperatorTok{;}
\ControlFlowTok{switch} \OperatorTok{(}\NormalTok{rotation}\OperatorTok{)} \OperatorTok{\{}
    \ControlFlowTok{case} \DecValTok{0}\OperatorTok{:}
\NormalTok{        pdfX }\OperatorTok{=}\NormalTok{ cropBox}\OperatorTok{.}\FunctionTok{getLowerLeftX}\OperatorTok{()} \OperatorTok{+}\NormalTok{ visualX}\OperatorTok{;}
\NormalTok{        pdfY }\OperatorTok{=}\NormalTok{ cropBox}\OperatorTok{.}\FunctionTok{getUpperRightY}\OperatorTok{()} \OperatorTok{{-}}\NormalTok{ visualTopOffset }\OperatorTok{{-}}\NormalTok{ sigHeight}\OperatorTok{;}
        \ControlFlowTok{break}\OperatorTok{;}
    \ControlFlowTok{case} \DecValTok{90}\OperatorTok{:}
\NormalTok{        pdfX }\OperatorTok{=}\NormalTok{ cropBox}\OperatorTok{.}\FunctionTok{getLowerLeftX}\OperatorTok{()} \OperatorTok{+}\NormalTok{ visualTopOffset}\OperatorTok{;}
\NormalTok{        pdfY }\OperatorTok{=}\NormalTok{ cropBox}\OperatorTok{.}\FunctionTok{getLowerLeftY}\OperatorTok{()} \OperatorTok{+}\NormalTok{ visualX}\OperatorTok{;}
        \CommentTok{// Width and height swap in the annotation rect}
        \DataTypeTok{float}\NormalTok{ temp }\OperatorTok{=}\NormalTok{ sigWidth}\OperatorTok{;}
\NormalTok{        sigWidth }\OperatorTok{=}\NormalTok{ sigHeight}\OperatorTok{;}
\NormalTok{        sigHeight }\OperatorTok{=}\NormalTok{ temp}\OperatorTok{;}
        \ControlFlowTok{break}\OperatorTok{;}
    \ControlFlowTok{case} \DecValTok{180}\OperatorTok{:}
\NormalTok{        pdfX }\OperatorTok{=}\NormalTok{ cropBox}\OperatorTok{.}\FunctionTok{getUpperRightX}\OperatorTok{()} \OperatorTok{{-}}\NormalTok{ visualX }\OperatorTok{{-}}\NormalTok{ sigWidth}\OperatorTok{;}
\NormalTok{        pdfY }\OperatorTok{=}\NormalTok{ cropBox}\OperatorTok{.}\FunctionTok{getLowerLeftY}\OperatorTok{()} \OperatorTok{+}\NormalTok{ visualTopOffset}\OperatorTok{;}
        \ControlFlowTok{break}\OperatorTok{;}
    \ControlFlowTok{case} \DecValTok{270}\OperatorTok{:}
\NormalTok{        pdfX }\OperatorTok{=}\NormalTok{ cropBox}\OperatorTok{.}\FunctionTok{getUpperRightX}\OperatorTok{()} \OperatorTok{{-}}\NormalTok{ visualTopOffset }\OperatorTok{{-}}\NormalTok{ sigHeight}\OperatorTok{;}
\NormalTok{        pdfY }\OperatorTok{=}\NormalTok{ cropBox}\OperatorTok{.}\FunctionTok{getUpperRightY}\OperatorTok{()} \OperatorTok{{-}}\NormalTok{ visualX }\OperatorTok{{-}}\NormalTok{ sigWidth}\OperatorTok{;}
\NormalTok{        temp }\OperatorTok{=}\NormalTok{ sigWidth}\OperatorTok{;}
\NormalTok{        sigWidth }\OperatorTok{=}\NormalTok{ sigHeight}\OperatorTok{;}
\NormalTok{        sigHeight }\OperatorTok{=}\NormalTok{ temp}\OperatorTok{;}
        \ControlFlowTok{break}\OperatorTok{;}
    \KeywordTok{default}\OperatorTok{:}
        \ControlFlowTok{throw} \KeywordTok{new} \BuiltInTok{IllegalArgumentException}\OperatorTok{(}\StringTok{"Invalid rotation: "} \OperatorTok{+}\NormalTok{ rotation}\OperatorTok{);}
\OperatorTok{\}}

\NormalTok{PDRectangle signatureRect }\OperatorTok{=} \KeywordTok{new} \FunctionTok{PDRectangle}\OperatorTok{(}\NormalTok{pdfX}\OperatorTok{,}\NormalTok{ pdfY}\OperatorTok{,}\NormalTok{ sigWidth}\OperatorTok{,}\NormalTok{ sigHeight}\OperatorTok{);}
\end{Highlighting}
\end{Shaded}

Wrap this coordinate transformation in a dedicated utility method and
unit-test it for all four rotation values. This is the kind of code
where an off-by-one error in a single case produces a bug that only
manifests for rotated pages --- which may be 5\% of your input
documents, ensuring the bug survives initial testing and surfaces in
production.

\section{Transformation Matrices and the
CTM}\label{_transformation_matrices_and_the_ctm}

ISO 32000-2:2020 Section 8.3.3 defines the \textbf{Current
Transformation Matrix (CTM)} --- a 3x3 affine transformation matrix that
maps user space coordinates to device space coordinates. Every content
stream operates under a CTM that controls translation, scaling,
rotation, and skewing.

The CTM is represented as six numbers \texttt{{[}a\ b\ c\ d\ e\ f{]}},
corresponding to the matrix:

\begin{verbatim}
| a  b  0 |
| c  d  0 |
| e  f  1 |
\end{verbatim}

For signature appearances, the CTM matters because the appearance stream
(the Form XObject that draws the visual signature) has its own
coordinate system. The \texttt{/BBox} of the appearance stream defines
the drawing area, and the viewer maps this BBox into the annotation
rectangle using a transformation derived from the rectangle dimensions
and the BBox dimensions.

When the appearance stream's BBox does not match the annotation
rectangle's aspect ratio, the viewer scales the appearance to fit ---
potentially distorting the signature image. If the BBox is
\texttt{{[}0\ 0\ 200\ 50{]}} but the annotation rectangle is 200 by 100
points, the appearance will be stretched vertically by a factor of 2.

When creating signature appearances in PDFBox, always ensure the
appearance stream's BBox matches the dimensions of the annotation
rectangle. If you use \texttt{PDRectangle} for both and create them from
the same width and height values, this happens automatically. The bug
occurs when one is constructed independently of the other, or when the
annotation rectangle is modified after the appearance is created.

Content streams use operators like \texttt{cm} (concatenate matrix) to
modify the CTM:

\begin{verbatim}
1 0 0 1 100 200 cm    % translate by (100, 200)
0.5 0 0 0.5 0 0 cm    % scale by 50%
\end{verbatim}

For signature placement, you rarely need to manipulate the CTM directly.
PDFBox handles the coordinate mapping between the appearance stream and
the annotation rectangle. But when debugging a signature that appears at
the wrong size or position, checking the CTM in the appearance stream's
content is a valuable diagnostic step.

\section{Signature Appearance Rectangles and Page
Coordinates}\label{_signature_appearance_rectangles_and_page_coordinates}

The annotation rectangle (\texttt{/Rect}) positions the signature on the
page. The appearance stream's BBox defines the drawing canvas. The
relationship between them determines how the visual signature is
rendered.

Here is the complete chain:

\begin{enumerate}
\def\labelenumi{\arabic{enumi}.}
\item
  The page has a coordinate system defined by its MediaBox (or CropBox).
\item
  The annotation's \texttt{/Rect} specifies where on the page the
  annotation sits, in page coordinates.
\item
  The appearance stream has a \texttt{/BBox} that defines its internal
  drawing area.
\item
  The viewer computes a transformation from the BBox to the Rect ---
  scaling and translating the appearance stream content to fill the
  annotation rectangle.
\item
  If the page has a \texttt{/Rotate} value, the viewer rotates the
  entire visual output \emph{after} all of the above.
\end{enumerate}

This chain explains why a signature appearance that looks correct on one
page can appear distorted on another: if the target page has different
dimensions or a different rotation, and you reuse the same appearance
stream without adjusting the annotation rectangle, the scaling
transformation changes.

\section{PDFBox Coordinate Handling}\label{_pdfbox_coordinate_handling}

PDFBox provides \texttt{PDRectangle} as the primary class for working
with rectangles in PDF user space. The key methods:

\begin{Shaded}
\begin{Highlighting}[]
\NormalTok{PDRectangle mediaBox }\OperatorTok{=}\NormalTok{ page}\OperatorTok{.}\FunctionTok{getMediaBox}\OperatorTok{();}   \CommentTok{// Physical page extent}
\NormalTok{PDRectangle cropBox }\OperatorTok{=}\NormalTok{ page}\OperatorTok{.}\FunctionTok{getCropBox}\OperatorTok{();}     \CommentTok{// Visible page extent (falls back to MediaBox)}
\DataTypeTok{int}\NormalTok{ rotation }\OperatorTok{=}\NormalTok{ page}\OperatorTok{.}\FunctionTok{getRotation}\OperatorTok{();}           \CommentTok{// 0, 90, 180, or 270}

\CommentTok{// PDRectangle stores lower{-}left corner + width/height}
\DataTypeTok{float}\NormalTok{ llx }\OperatorTok{=}\NormalTok{ cropBox}\OperatorTok{.}\FunctionTok{getLowerLeftX}\OperatorTok{();}
\DataTypeTok{float}\NormalTok{ lly }\OperatorTok{=}\NormalTok{ cropBox}\OperatorTok{.}\FunctionTok{getLowerLeftY}\OperatorTok{();}
\DataTypeTok{float}\NormalTok{ urx }\OperatorTok{=}\NormalTok{ cropBox}\OperatorTok{.}\FunctionTok{getUpperRightX}\OperatorTok{();}
\DataTypeTok{float}\NormalTok{ ury }\OperatorTok{=}\NormalTok{ cropBox}\OperatorTok{.}\FunctionTok{getUpperRightY}\OperatorTok{();}
\DataTypeTok{float}\NormalTok{ width }\OperatorTok{=}\NormalTok{ cropBox}\OperatorTok{.}\FunctionTok{getWidth}\OperatorTok{();}
\DataTypeTok{float}\NormalTok{ height }\OperatorTok{=}\NormalTok{ cropBox}\OperatorTok{.}\FunctionTok{getHeight}\OperatorTok{();}
\end{Highlighting}
\end{Shaded}

\texttt{PDRectangle} predefined constants include
\texttt{PDRectangle.A4} (595.28 x 841.89) and
\texttt{PDRectangle.LETTER} (612 x 792). These are useful for creating
new pages but should never be used as the reference for positioning on
an existing page --- always query the actual page's CropBox.

\subsection{Practical Example: Placing a Visible
Signature}\label{_practical_example_placing_a_visible_signature}

The following complete example places a visible signature rectangle at a
specific visual position on a page, correctly handling rotation:

\begin{Shaded}
\begin{Highlighting}[]
\KeywordTok{import} \ImportTok{org}\OperatorTok{.}\ImportTok{apache}\OperatorTok{.}\ImportTok{pdfbox}\OperatorTok{.}\ImportTok{Loader}\OperatorTok{;}
\KeywordTok{import} \ImportTok{org}\OperatorTok{.}\ImportTok{apache}\OperatorTok{.}\ImportTok{pdfbox}\OperatorTok{.}\ImportTok{pdmodel}\OperatorTok{.}\ImportTok{PDDocument}\OperatorTok{;}
\KeywordTok{import} \ImportTok{org}\OperatorTok{.}\ImportTok{apache}\OperatorTok{.}\ImportTok{pdfbox}\OperatorTok{.}\ImportTok{pdmodel}\OperatorTok{.}\ImportTok{PDPage}\OperatorTok{;}
\KeywordTok{import} \ImportTok{org}\OperatorTok{.}\ImportTok{apache}\OperatorTok{.}\ImportTok{pdfbox}\OperatorTok{.}\ImportTok{pdmodel}\OperatorTok{.}\ImportTok{common}\OperatorTok{.}\ImportTok{PDRectangle}\OperatorTok{;}

\KeywordTok{public} \KeywordTok{class}\NormalTok{ SignaturePositionCalculator }\OperatorTok{\{}

    \CommentTok{/**}
     \CommentTok{*}\NormalTok{ Computes the PDF annotation rectangle for a signature at a given}
     \CommentTok{*}\NormalTok{ visual position}\CommentTok{,}\NormalTok{ accounting for page rotation and CropBox offset}\CommentTok{.}
     \CommentTok{*}
     \CommentTok{*} \AnnotationTok{@param}\CommentTok{ }\AnnotationTok{page}\CommentTok{         }\NormalTok{the target page}
     \CommentTok{*} \AnnotationTok{@param}\CommentTok{ }\AnnotationTok{visualX}\CommentTok{      }\NormalTok{x offset from the left of the visual page }\CommentTok{(}\NormalTok{points}\CommentTok{)}
     \CommentTok{*} \AnnotationTok{@param}\CommentTok{ }\AnnotationTok{visualY}\CommentTok{      }\NormalTok{y offset from the top of the visual page }\CommentTok{(}\NormalTok{points}\CommentTok{)}
     \CommentTok{*} \AnnotationTok{@param}\CommentTok{ }\AnnotationTok{visualWidth}\CommentTok{  }\NormalTok{desired visual width of the signature }\CommentTok{(}\NormalTok{points}\CommentTok{)}
     \CommentTok{*} \AnnotationTok{@param}\CommentTok{ }\AnnotationTok{visualHeight}\CommentTok{ }\NormalTok{desired visual height of the signature }\CommentTok{(}\NormalTok{points}\CommentTok{)}
     \CommentTok{*} \AnnotationTok{@return}\NormalTok{ the annotation rectangle in PDF user space coordinates}
     \CommentTok{*/}
    \KeywordTok{public} \DataTypeTok{static}\NormalTok{ PDRectangle }\FunctionTok{computeSignatureRect}\OperatorTok{(}
\NormalTok{            PDPage page}\OperatorTok{,}
            \DataTypeTok{float}\NormalTok{ visualX}\OperatorTok{,} \DataTypeTok{float}\NormalTok{ visualY}\OperatorTok{,}
            \DataTypeTok{float}\NormalTok{ visualWidth}\OperatorTok{,} \DataTypeTok{float}\NormalTok{ visualHeight}\OperatorTok{)} \OperatorTok{\{}

\NormalTok{        PDRectangle crop }\OperatorTok{=}\NormalTok{ page}\OperatorTok{.}\FunctionTok{getCropBox}\OperatorTok{();}
        \DataTypeTok{int}\NormalTok{ rotation }\OperatorTok{=}\NormalTok{ page}\OperatorTok{.}\FunctionTok{getRotation}\OperatorTok{();}

        \DataTypeTok{float}\NormalTok{ pdfX}\OperatorTok{,}\NormalTok{ pdfY}\OperatorTok{,}\NormalTok{ pdfW}\OperatorTok{,}\NormalTok{ pdfH}\OperatorTok{;}

        \ControlFlowTok{switch} \OperatorTok{(}\NormalTok{rotation}\OperatorTok{)} \OperatorTok{\{}
            \ControlFlowTok{case} \DecValTok{0}\OperatorTok{:}
\NormalTok{                pdfX }\OperatorTok{=}\NormalTok{ crop}\OperatorTok{.}\FunctionTok{getLowerLeftX}\OperatorTok{()} \OperatorTok{+}\NormalTok{ visualX}\OperatorTok{;}
\NormalTok{                pdfY }\OperatorTok{=}\NormalTok{ crop}\OperatorTok{.}\FunctionTok{getUpperRightY}\OperatorTok{()} \OperatorTok{{-}}\NormalTok{ visualY }\OperatorTok{{-}}\NormalTok{ visualHeight}\OperatorTok{;}
\NormalTok{                pdfW }\OperatorTok{=}\NormalTok{ visualWidth}\OperatorTok{;}
\NormalTok{                pdfH }\OperatorTok{=}\NormalTok{ visualHeight}\OperatorTok{;}
                \ControlFlowTok{break}\OperatorTok{;}
            \ControlFlowTok{case} \DecValTok{90}\OperatorTok{:}
\NormalTok{                pdfX }\OperatorTok{=}\NormalTok{ crop}\OperatorTok{.}\FunctionTok{getLowerLeftX}\OperatorTok{()} \OperatorTok{+}\NormalTok{ visualY}\OperatorTok{;}
\NormalTok{                pdfY }\OperatorTok{=}\NormalTok{ crop}\OperatorTok{.}\FunctionTok{getLowerLeftY}\OperatorTok{()} \OperatorTok{+}\NormalTok{ visualX}\OperatorTok{;}
\NormalTok{                pdfW }\OperatorTok{=}\NormalTok{ visualHeight}\OperatorTok{;}
\NormalTok{                pdfH }\OperatorTok{=}\NormalTok{ visualWidth}\OperatorTok{;}
                \ControlFlowTok{break}\OperatorTok{;}
            \ControlFlowTok{case} \DecValTok{180}\OperatorTok{:}
\NormalTok{                pdfX }\OperatorTok{=}\NormalTok{ crop}\OperatorTok{.}\FunctionTok{getUpperRightX}\OperatorTok{()} \OperatorTok{{-}}\NormalTok{ visualX }\OperatorTok{{-}}\NormalTok{ visualWidth}\OperatorTok{;}
\NormalTok{                pdfY }\OperatorTok{=}\NormalTok{ crop}\OperatorTok{.}\FunctionTok{getLowerLeftY}\OperatorTok{()} \OperatorTok{+}\NormalTok{ visualY}\OperatorTok{;}
\NormalTok{                pdfW }\OperatorTok{=}\NormalTok{ visualWidth}\OperatorTok{;}
\NormalTok{                pdfH }\OperatorTok{=}\NormalTok{ visualHeight}\OperatorTok{;}
                \ControlFlowTok{break}\OperatorTok{;}
            \ControlFlowTok{case} \DecValTok{270}\OperatorTok{:}
\NormalTok{                pdfX }\OperatorTok{=}\NormalTok{ crop}\OperatorTok{.}\FunctionTok{getUpperRightX}\OperatorTok{()} \OperatorTok{{-}}\NormalTok{ visualY }\OperatorTok{{-}}\NormalTok{ visualHeight}\OperatorTok{;}
\NormalTok{                pdfY }\OperatorTok{=}\NormalTok{ crop}\OperatorTok{.}\FunctionTok{getUpperRightY}\OperatorTok{()} \OperatorTok{{-}}\NormalTok{ visualX }\OperatorTok{{-}}\NormalTok{ visualWidth}\OperatorTok{;}
\NormalTok{                pdfW }\OperatorTok{=}\NormalTok{ visualHeight}\OperatorTok{;}
\NormalTok{                pdfH }\OperatorTok{=}\NormalTok{ visualWidth}\OperatorTok{;}
                \ControlFlowTok{break}\OperatorTok{;}
            \KeywordTok{default}\OperatorTok{:}
                \ControlFlowTok{throw} \KeywordTok{new} \BuiltInTok{IllegalArgumentException}\OperatorTok{(}
                    \StringTok{"Unsupported rotation: "} \OperatorTok{+}\NormalTok{ rotation}\OperatorTok{);}
        \OperatorTok{\}}

        \ControlFlowTok{return} \KeywordTok{new} \FunctionTok{PDRectangle}\OperatorTok{(}\NormalTok{pdfX}\OperatorTok{,}\NormalTok{ pdfY}\OperatorTok{,}\NormalTok{ pdfW}\OperatorTok{,}\NormalTok{ pdfH}\OperatorTok{);}
    \OperatorTok{\}}
\OperatorTok{\}}
\end{Highlighting}
\end{Shaded}

Call this with visual coordinates --- "200 points from the left, 50
points from the top, 180 wide, 40 tall" --- and it returns the correct
PDF rectangle regardless of rotation.

\section{Common Coordinate Bugs}\label{_common_coordinate_bugs}

Over years of building PDF signing tools, we have seen the same bugs
recur. Here is a catalog.

\subsection{Flipped Signatures (Y-axis
Inversion)}\label{_flipped_signatures_y_axis_inversion}

\textbf{Symptom}: The signature appears at the vertical opposite of
where it should be. A signature placed near the top of the page appears
near the bottom.

\textbf{Cause}: Using screen-space y-coordinates (origin top-left)
directly as PDF y-coordinates (origin bottom-left) without inverting.

\textbf{Fix}: Apply the formula
\texttt{pdfY\ =\ cropBox.getUpperRightY()\ -\ screenY\ -\ signatureHeight}.

\subsection{Signatures on the Wrong Page Position (Rotated
Pages)}\label{_signatures_on_the_wrong_page_position_rotated_pages}

\textbf{Symptom}: The signature appears in a completely unexpected
location --- not just vertically flipped, but seemingly random --- and
only on certain documents.

\textbf{Cause}: The page has a non-zero \texttt{/Rotate} value, and the
coordinate calculation does not account for it. Many developers test
exclusively with non-rotated pages and never encounter this until a
customer submits a rotated document.

\textbf{Fix}: Always check \texttt{page.getRotation()} and apply the
rotation-aware coordinate transformation shown above.

\subsection{Zero-Area Annotation (Invisible Signature Becomes Truly
Invisible)}\label{_zero_area_annotation_invisible_signature_becomes_truly_invisible}

\textbf{Symptom}: The signature field exists (it shows in the signature
panel), but nothing is visible on the page.

\textbf{Cause}: The annotation rectangle has zero width or height,
typically because the lower-left and upper-right corners are identical,
or because a coordinate calculation produced a negative dimension that
the viewer clamped to zero.

\textbf{Fix}: Always validate that the computed rectangle has positive
width and height before creating the annotation. Log the rectangle
values during development:

\begin{Shaded}
\begin{Highlighting}[]
\NormalTok{log}\OperatorTok{.}\FunctionTok{debug}\OperatorTok{(}\StringTok{"Signature rect: x=\{\}, y=\{\}, w=\{\}, h=\{\}"}\OperatorTok{,}
\NormalTok{    rect}\OperatorTok{.}\FunctionTok{getLowerLeftX}\OperatorTok{(),}\NormalTok{ rect}\OperatorTok{.}\FunctionTok{getLowerLeftY}\OperatorTok{(),}
\NormalTok{    rect}\OperatorTok{.}\FunctionTok{getWidth}\OperatorTok{(),}\NormalTok{ rect}\OperatorTok{.}\FunctionTok{getHeight}\OperatorTok{());}
\end{Highlighting}
\end{Shaded}

\subsection{Signature Appears
Off-Page}\label{_signature_appears_off_page}

\textbf{Symptom}: The signature field exists, the signature validates,
but nothing is visible. Unlike the zero-area case, the rectangle has
positive dimensions --- it's just outside the visible CropBox.

\textbf{Cause}: Coordinates were computed against the MediaBox, but the
CropBox is smaller. Or the CropBox has a non-zero origin that was not
accounted for.

\textbf{Fix}: Always compute against \texttt{page.getCropBox()} and
include the CropBox's lower-left offset in your calculations.

\subsection{Distorted Signature
Appearance}\label{_distorted_signature_appearance}

\textbf{Symptom}: The signature image or text appears stretched or
squashed.

\textbf{Cause}: The appearance stream's BBox has a different aspect
ratio than the annotation rectangle. The viewer scales the BBox to fit
the Rect, distorting the content.

\textbf{Fix}: Ensure the appearance stream's BBox dimensions match the
annotation rectangle's width and height.

\section{Debugging Coordinates with QPDF and
RUPS}\label{_debugging_coordinates_with_qpdf_and_rups}

When a signature appears in the wrong place, you need to inspect the raw
PDF objects. Two tools are indispensable.

\subsection{QPDF}\label{_qpdf}

QPDF (\url{https://qpdf.sourceforge.io/}) is a command-line tool that
can dump every object in a PDF in a human-readable JSON format:

\begin{Shaded}
\begin{Highlighting}[]
\ExtensionTok{qpdf} \AttributeTok{{-}{-}json}\NormalTok{ input.pdf }\OperatorTok{\textgreater{}}\NormalTok{ /tmp/structure.json}
\end{Highlighting}
\end{Shaded}

Then search for the annotation rectangle:

\begin{Shaded}
\begin{Highlighting}[]
\FunctionTok{grep} \AttributeTok{{-}A2} \StringTok{\textquotesingle{}"Rect"\textquotesingle{}}\NormalTok{ /tmp/structure.json}
\end{Highlighting}
\end{Shaded}

You will see the raw \texttt{/Rect} array values. Compare them against
the page's \texttt{/MediaBox} and \texttt{/CropBox} to verify they fall
within the visible area.

QPDF can also show page-level dictionaries:

\begin{Shaded}
\begin{Highlighting}[]
\ExtensionTok{qpdf} \AttributeTok{{-}{-}show{-}object}\OperatorTok{=}\NormalTok{3 input.pdf   }\CommentTok{\# show object 3 (a page, for example)}
\end{Highlighting}
\end{Shaded}

\subsection{RUPS (Reading and Updating PDF
Syntax)}\label{_rups_reading_and_updating_pdf_syntax}

RUPS is a Java Swing application from the iText team that provides a
tree view of every PDF object. It is invaluable for visual debugging
because you can navigate the object tree, inspect annotation rectangles,
and see the page boxes side by side.

To inspect a signature's coordinates in RUPS:

\begin{enumerate}
\def\labelenumi{\arabic{enumi}.}
\item
  Open the PDF in RUPS.
\item
  Navigate to \textbf{Pages} → \textbf{Page N} → \textbf{Annots} (the
  annotations array).
\item
  Find the signature widget annotation.
\item
  Check the \texttt{/Rect} value --- this is the annotation rectangle in
  page coordinates.
\item
  Go back to the page and check \texttt{/MediaBox}, \texttt{/CropBox},
  and \texttt{/Rotate}.
\item
  Verify that the Rect falls within the CropBox and that you've
  accounted for the rotation.
\end{enumerate}

RUPS also lets you inspect the appearance stream's BBox. Navigate into
the annotation's \texttt{/AP} → \texttt{/N} and check the \texttt{/BBox}
entry. If the BBox dimensions don't match the Rect dimensions, you've
found your distortion bug.

\subsection{PDFBox Programmatic
Inspection}\label{_pdfbox_programmatic_inspection}

You can also inspect coordinates programmatically, which is useful in
automated tests:

\begin{Shaded}
\begin{Highlighting}[]
\ControlFlowTok{try} \OperatorTok{(}\NormalTok{PDDocument doc }\OperatorTok{=}\NormalTok{ Loader}\OperatorTok{.}\FunctionTok{loadPDF}\OperatorTok{(}\KeywordTok{new} \BuiltInTok{File}\OperatorTok{(}\StringTok{"signed.pdf"}\OperatorTok{)))} \OperatorTok{\{}
    \ControlFlowTok{for} \OperatorTok{(}\NormalTok{PDPage page }\OperatorTok{:}\NormalTok{ doc}\OperatorTok{.}\FunctionTok{getPages}\OperatorTok{())} \OperatorTok{\{}
\NormalTok{        log}\OperatorTok{.}\FunctionTok{info}\OperatorTok{(}\StringTok{"Page MediaBox: \{\}"}\OperatorTok{,}\NormalTok{ page}\OperatorTok{.}\FunctionTok{getMediaBox}\OperatorTok{());}
\NormalTok{        log}\OperatorTok{.}\FunctionTok{info}\OperatorTok{(}\StringTok{"Page CropBox:  \{\}"}\OperatorTok{,}\NormalTok{ page}\OperatorTok{.}\FunctionTok{getCropBox}\OperatorTok{());}
\NormalTok{        log}\OperatorTok{.}\FunctionTok{info}\OperatorTok{(}\StringTok{"Page Rotation: \{\}"}\OperatorTok{,}\NormalTok{ page}\OperatorTok{.}\FunctionTok{getRotation}\OperatorTok{());}

        \ControlFlowTok{for} \OperatorTok{(}\DataTypeTok{var}\NormalTok{ annot }\OperatorTok{:}\NormalTok{ page}\OperatorTok{.}\FunctionTok{getAnnotations}\OperatorTok{())} \OperatorTok{\{}
            \ControlFlowTok{if} \OperatorTok{(}\StringTok{"Sig"}\OperatorTok{.}\FunctionTok{equals}\OperatorTok{(}\NormalTok{annot}\OperatorTok{.}\FunctionTok{getSubtype}\OperatorTok{()))} \OperatorTok{\{}
\NormalTok{                PDRectangle rect }\OperatorTok{=}\NormalTok{ annot}\OperatorTok{.}\FunctionTok{getRectangle}\OperatorTok{();}
\NormalTok{                log}\OperatorTok{.}\FunctionTok{info}\OperatorTok{(}\StringTok{"Sig annotation Rect: [\{\}, \{\}, \{\}, \{\}]"}\OperatorTok{,}
\NormalTok{                    rect}\OperatorTok{.}\FunctionTok{getLowerLeftX}\OperatorTok{(),}\NormalTok{ rect}\OperatorTok{.}\FunctionTok{getLowerLeftY}\OperatorTok{(),}
\NormalTok{                    rect}\OperatorTok{.}\FunctionTok{getUpperRightX}\OperatorTok{(),}\NormalTok{ rect}\OperatorTok{.}\FunctionTok{getUpperRightY}\OperatorTok{());}
            \OperatorTok{\}}
        \OperatorTok{\}}
    \OperatorTok{\}}
\OperatorTok{\}}
\end{Highlighting}
\end{Shaded}

When writing integration tests for signature placement, always test with
at least four documents: one for each rotation value (0, 90, 180, 270).
Also test with a document that has a CropBox smaller than its MediaBox.
These five test cases catch the vast majority of coordinate bugs before
they reach production.

\section{Summary}\label{_summary_8}

The PDF coordinate system is a solved problem --- the rules are clear
and have not changed since PDF 1.0. The difficulty lies entirely in the
mismatch between PDF's bottom-left origin and the top-left origin of
every UI framework, compounded by page rotation's effect on visual
layout.

For signature placement, remember these principles:

\begin{itemize}
\item
  Always use the CropBox (via \texttt{page.getCropBox()}) as your
  reference frame.
\item
  Always invert the y-axis when converting from screen coordinates:
  \texttt{pdfY\ =\ cropBoxTop\ -\ screenY\ -\ height}.
\item
  Always check \texttt{page.getRotation()} and transform coordinates
  accordingly.
\item
  Always ensure the appearance stream BBox matches the annotation
  rectangle dimensions.
\item
  Always test with rotated pages and non-standard CropBoxes.
\end{itemize}

These are not complex algorithms. They are simple arithmetic applied
consistently. The bugs come from skipping a step, not from the math
being hard.

\chapter{Font Embedding in PDF}\label{_font_embedding_in_pdf}

Fonts are the invisible infrastructure of every PDF document. When they
work, nobody notices. When they break, the document becomes
unreadable --- and if that document contains a digital signature, the
signature appearance becomes a row of blank rectangles or missing-glyph
squares that make your legally binding signature look like a rendering
error.

This chapter explains how fonts work inside PDF files, how they get
embedded (or not), what goes wrong when they don't, and how all of this
matters specifically for digital signatures. We cover the font types
defined by ISO 32000-2:2020, the encoding models, the mechanics of
embedding and subsetting, and the practical reality of getting text to
render correctly in signature appearance streams across different
viewers.

If you have ever seen a signature block that shows \texttt{.notdef}
squares where a signer's name should be, or a PDF where the signature
text vanishes entirely on one viewer but renders fine on another, this
chapter explains why.

\section{Why Fonts Matter for Digital
Signatures}\label{_why_fonts_matter_for_digital_signatures}

A PDF digital signature has two independent representations: the
\textbf{cryptographic} representation (the CMS \texttt{SignedData}
structure, the byte range, the certificate chain) and the
\textbf{visual} representation (the appearance stream that shows the
signer's name, date, and reason on the page). Chapter 15 covers how
appearance streams are structured; this chapter covers the font
infrastructure that makes the \emph{text} inside those appearances
render correctly.

The visual appearance of a signature is a Form XObject --- a miniature
PDF page with its own resource dictionary, its own content stream, and
its own font references. When that content stream contains a \texttt{Tj}
or \texttt{TJ} operator to draw text, the PDF viewer must resolve the
font reference, find the glyph outlines, and render them. If any step in
that chain fails --- the font reference points to a font that isn't
embedded, the encoding doesn't map the character code to the right
glyph, the font descriptor is missing, the font file stream is
corrupt --- the text doesn't render.

This matters more for signatures than for regular document text, for
three reasons:

\begin{enumerate}
\def\labelenumi{\arabic{enumi}.}
\item
  \textbf{Signature appearances are generated programmatically.} When a
  human creates a PDF in a word processor, the word processor embeds the
  fonts it uses. When your signing code creates a signature appearance,
  \emph{your code} is responsible for font handling. If you get it
  wrong, nobody catches the error until a user reports that the
  signature looks broken.
\item
  \textbf{Signature appearances must survive across viewers.} A
  signature created on a server using PDFBox must render correctly in
  Adobe Acrobat, Foxit Reader, browser-based viewers using pdf.js,
  LibreOffice, and Preview.app. Each viewer has its own font fallback
  behavior, its own interpretation of encoding, and its own handling of
  missing fonts. The only way to guarantee consistent rendering is to
  embed the font.
\item
  \textbf{Signature text often contains names with non-ASCII
  characters.} Signer names like "Muller" are fine. Signer names like
  "Muller" with an umlaut, "Gonzalez" with an accent, or any East Asian,
  Arabic, or Cyrillic name will fail if the font and encoding aren't set
  up correctly. In a global signing system, "works for ASCII names" is
  not good enough.
\end{enumerate}

\section{Font Types in PDF}\label{_font_types_in_pdf}

ISO 32000-2:2020 Section 9.6 defines several font types. The
distinctions matter because each type has different embedding rules,
different encoding models, and different failure modes.

\subsection{Type 1 Fonts (PostScript
Era)}\label{_type_1_fonts_postscript_era}

Type 1 fonts are the original PostScript outline fonts, defined by Adobe
in 1984. They use cubic Bezier curves for glyph outlines and a custom
PostScript-based encoding scheme.

In PDF, a Type 1 font dictionary has \texttt{/Subtype\ /Type1} and
references a font program (the glyph outlines) via the
\texttt{/FontDescriptor} dictionary's \texttt{/FontFile} entry. The font
program is a stream containing the PostScript font data in PFB (Printer
Font Binary) format.

Type 1 fonts are limited to a maximum of 256 glyphs per font
instance --- a consequence of the single-byte encoding model inherited
from PostScript. This makes them unsuitable for CJK scripts, complex
scripts like Arabic or Devanagari, or any language that needs more than
256 distinct glyphs.

Adobe officially ended support for Type 1 fonts in January 2023,
removing the ability to author content with Type 1 fonts in their
Creative Cloud applications. Existing PDFs with embedded Type 1 font
programs still render correctly (the outlines are in the file), but this
signals the industry's long-term direction: TrueType and OpenType are
the future.

\subsection{The Standard 14 Fonts}\label{_the_standard_14_fonts}

PDF defines fourteen fonts that conforming readers are expected to have
available without embedding. These are the \textbf{Standard 14 fonts},
specified in ISO 32000-2:2020 Section 9.6.2.2:

\begin{longtable}[]{@{}
  >{\raggedright\arraybackslash}p{(\linewidth - 4\tabcolsep) * \real{0.2500}}
  >{\raggedright\arraybackslash}p{(\linewidth - 4\tabcolsep) * \real{0.3750}}
  >{\raggedright\arraybackslash}p{(\linewidth - 4\tabcolsep) * \real{0.3750}}@{}}
\caption{The Standard 14 font families}\tabularnewline
\toprule\noalign{}
\begin{minipage}[b]{\linewidth}\raggedright
Family
\end{minipage} & \begin{minipage}[b]{\linewidth}\raggedright
Regular / Bold / Italic / BoldItalic
\end{minipage} & \begin{minipage}[b]{\linewidth}\raggedright
Notes
\end{minipage} \\
\midrule\noalign{}
\endfirsthead
\toprule\noalign{}
\begin{minipage}[b]{\linewidth}\raggedright
Family
\end{minipage} & \begin{minipage}[b]{\linewidth}\raggedright
Regular / Bold / Italic / BoldItalic
\end{minipage} & \begin{minipage}[b]{\linewidth}\raggedright
Notes
\end{minipage} \\
\midrule\noalign{}
\endhead
\bottomrule\noalign{}
\endlastfoot
Helvetica & Helvetica, Helvetica-Bold, Helvetica-Oblique,
Helvetica-BoldOblique & Sans-serif. The default choice for most
programmatic PDF generation. \\
Times & Times-Roman, Times-Bold, Times-Italic, Times-BoldItalic & Serif.
Common in academic and legal documents. \\
Courier & Courier, Courier-Bold, Courier-Oblique, Courier-BoldOblique &
Monospace. Used for code listings and form data. \\
Symbol & Symbol & Mathematical and Greek symbols. \\
ZapfDingbats & ZapfDingbats & Decorative symbols (checkmarks, arrows,
etc.). \\
\end{longtable}

The critical word is "expected." ISO 32000-2:2020 Section 9.6.2.2 states
that conforming readers \emph{should} provide these fonts, but PDF 2.0
actually recommends embedding even the Standard 14 fonts to ensure
consistent rendering. The practical reality is that most viewers do
provide metrics-compatible substitutes for these fonts, but the
substitutes are not identical --- a document typeset with Helvetica on
one viewer may render with Arial or Liberation Sans on another, causing
slight differences in character widths, line breaks, and text
positioning.

For signature appearances, these differences are usually cosmetic and
acceptable. But if your signature appearance calculates text positions
based on precise Helvetica metrics (e.g., centering text in the
signature rectangle), a substitute font with slightly different widths
will misalign the text.

For signature appearances, the Standard 14 fonts are "good enough" for
ASCII text. They are available everywhere, they don't require embedding,
and the rendering differences between substitute fonts are typically
invisible at the font sizes used in signature blocks (8-12 points).

For non-ASCII text --- accented characters, umlauts, CJK
characters --- the Standard 14 fonts are \emph{not} sufficient, because
their built-in encoding (WinAnsiEncoding) cannot represent characters
outside the Windows-1252 code page. If your signers have names with
characters outside that range, you must embed a TrueType or OpenType
font with the appropriate glyph coverage.

\subsubsection{PDFBox 3.x and the Standard 14 API
Change}\label{_pdfbox_3_x_and_the_standard_14_api_change}

If you've migrated from PDFBox 2.x to 3.x, you've likely encountered
this:

\begin{Shaded}
\begin{Highlighting}[]
\CommentTok{// PDFBox 2.x {-}{-} this no longer compiles in 3.x}
\NormalTok{PDFont font }\OperatorTok{=}\NormalTok{ PDType1Font}\OperatorTok{.}\FunctionTok{HELVETICA}\OperatorTok{;}
\end{Highlighting}
\end{Shaded}

In PDFBox 2.x, the Standard 14 fonts were available as static constants
on \texttt{PDType1Font}. PDFBox 3.x removed these constants and replaced
them with an enum-based API:

\begin{Shaded}
\begin{Highlighting}[]
\CommentTok{// PDFBox 3.x {-}{-} the correct way}
\NormalTok{PDFont font }\OperatorTok{=} \KeywordTok{new} \FunctionTok{PDType1Font}\OperatorTok{(}\NormalTok{Standard14Fonts}\OperatorTok{.}\FunctionTok{FontName}\OperatorTok{.}\FunctionTok{HELVETICA}\OperatorTok{);}

\CommentTok{// Other examples:}
\NormalTok{PDFont bold }\OperatorTok{=} \KeywordTok{new} \FunctionTok{PDType1Font}\OperatorTok{(}\NormalTok{Standard14Fonts}\OperatorTok{.}\FunctionTok{FontName}\OperatorTok{.}\FunctionTok{HELVETICA\_BOLD}\OperatorTok{);}
\NormalTok{PDFont courier }\OperatorTok{=} \KeywordTok{new} \FunctionTok{PDType1Font}\OperatorTok{(}\NormalTok{Standard14Fonts}\OperatorTok{.}\FunctionTok{FontName}\OperatorTok{.}\FunctionTok{COURIER}\OperatorTok{);}
\NormalTok{PDFont times }\OperatorTok{=} \KeywordTok{new} \FunctionTok{PDType1Font}\OperatorTok{(}\NormalTok{Standard14Fonts}\OperatorTok{.}\FunctionTok{FontName}\OperatorTok{.}\FunctionTok{TIMES\_ROMAN}\OperatorTok{);}
\end{Highlighting}
\end{Shaded}

The reason for this change is that the old static constants were shared
mutable objects --- every caller that obtained
\texttt{PDType1Font.HELVETICA} got the same instance, which could cause
subtle threading and state pollution issues when multiple documents were
processed concurrently. The new constructor creates a fresh instance
each time, eliminating the shared-state problem.

If you search the internet for PDFBox font examples, the majority of
results still show the 2.x API (\texttt{PDType1Font.HELVETICA}). These
examples will not compile with PDFBox 3.0.0. Always verify that code
examples match the major version of PDFBox you are using.

\subsection{TrueType Fonts}\label{_truetype_fonts}

TrueType fonts, developed jointly by Apple and Microsoft in the late
1980s as a competitor to Adobe's Type 1, use quadratic Bezier curves for
glyph outlines. They are the most common font format in modern computing
and the recommended choice for PDF font embedding.

In PDF, a TrueType font dictionary has \texttt{/Subtype\ /TrueType}. The
font program is embedded via the \texttt{/FontDescriptor} dictionary's
\texttt{/FontFile2} entry, which contains the raw TrueType font data (a
\texttt{.ttf} file, essentially).

TrueType fonts in PDF can be either \textbf{simple} (single-byte
encoding, max 256 glyphs per instance) or embedded as part of a
\textbf{composite} font (Type 0 + CIDFont, described below) for full
Unicode coverage. When used as a simple font, TrueType has the same
256-glyph limitation as Type 1. The composite font mechanism removes
this limitation.

\begin{Shaded}
\begin{Highlighting}[]
\CommentTok{// Loading and embedding a TrueType font in PDFBox 3.x}
\NormalTok{PDFont font }\OperatorTok{=}\NormalTok{ PDType0Font}\OperatorTok{.}\FunctionTok{load}\OperatorTok{(}\NormalTok{document}\OperatorTok{,} \KeywordTok{new} \BuiltInTok{File}\OperatorTok{(}\StringTok{"/path/to/DejaVuSans.ttf"}\OperatorTok{));}

\CommentTok{// The font is now embedded in the document.}
\CommentTok{// It will be automatically subsetted when the document is saved,}
\CommentTok{// including only the glyphs actually used.}
\end{Highlighting}
\end{Shaded}

\texttt{PDType0Font.load()} is the PDFBox 3.x method for loading
TrueType and OpenType fonts. Despite the name suggesting Type 0, this
method creates the full composite font structure (Type 0 wrapper +
CIDFont descendant) needed for Unicode text support. The font file is
embedded in the document, and PDFBox automatically subsets it on save.

\subsection{Type 0 / CIDFont (Composite
Fonts)}\label{_type_0_cidfont_composite_fonts}

Type 0 fonts --- also called composite fonts --- are the mechanism PDF
uses to support large character sets (CJK scripts, Arabic, Devanagari,
and any script requiring more than 256 glyphs).

A Type 0 font consists of two linked dictionaries:

\begin{enumerate}
\def\labelenumi{\arabic{enumi}.}
\item
  \textbf{The Type 0 font dictionary} (\texttt{/Subtype\ /Type0}): the
  "wrapper" that the content stream references. It contains a
  \texttt{/DescendantFonts} array pointing to the CIDFont dictionary and
  a \texttt{/Encoding} entry specifying the CMap that maps character
  codes to CIDs (Character IDs).
\item
  \textbf{The CIDFont dictionary} (\texttt{/Subtype\ /CIDFontType0} or
  \texttt{/CIDFontType2}): the actual font data. CIDFontType0 wraps a
  PostScript CID-keyed font; CIDFontType2 wraps a TrueType font. In
  practice, CIDFontType2 (TrueType-based) is overwhelmingly more common.
\end{enumerate}

\begin{verbatim}
% Type 0 composite font structure in a PDF
10 0 obj
<< /Type /Font
   /Subtype /Type0
   /BaseFont /ABCDEF+DejaVuSans         
   /Encoding /Identity-H                 
   /DescendantFonts [11 0 R]            
   /ToUnicode 12 0 R                     
>>
endobj

11 0 obj
<< /Type /Font
   /Subtype /CIDFontType2               
   /BaseFont /ABCDEF+DejaVuSans
   /CIDSystemInfo << /Registry (Adobe) /Ordering (Identity) /Supplement 0 >>
   /FontDescriptor 13 0 R
   /W [3 [278] 36 [722] 68 [556 500]]  
>>
endobj
\end{verbatim}

\begin{itemize}
\item
  The \texttt{/BaseFont} name includes a subset prefix
  (ABCDEF+) --- we'll explain this below
\item
  \texttt{/Identity-H} means character codes map directly to CIDs
  (horizontal writing mode)
\item
  The descendant CIDFont dictionary
\item
  The \texttt{/ToUnicode} CMap enables text extraction --- critical for
  accessibility and search
\item
  CIDFontType2: the actual glyph data comes from an embedded TrueType
  font
\item
  The \texttt{/W} (Widths) array defines per-glyph widths for precise
  text positioning
\end{itemize}

When PDFBox's \texttt{PDType0Font.load()} embeds a TrueType font, it
creates exactly this structure. You don't need to build it manually, but
understanding it helps when debugging font issues in existing PDFs.

\section{Font Embedding vs.
Referencing}\label{_font_embedding_vs_referencing}

There are two ways a PDF can use a font: \textbf{embed} it (include the
font program inside the PDF file) or \textbf{reference} it (include only
the font name and metrics, relying on the viewer to provide the actual
glyphs).

\subsection{Referenced Fonts}\label{_referenced_fonts}

A referenced font includes the font dictionary and metrics but
\emph{not} the font program. The viewer must find a matching font on the
system or use a fallback. The Standard 14 fonts are the canonical
example --- they're designed to be used without embedding because every
viewer is expected to have compatible substitutes.

Referenced non-standard fonts are a recipe for rendering failures. If
your PDF references "CompanyBrandFont-Regular" without embedding it, the
viewer will substitute whatever it can find --- usually a generic serif
or sans-serif font --- and the text will look wrong. Character widths
will differ, glyphs may be missing, and text positioning will be off.

For signature appearances, \textbf{never reference non-standard fonts}.
The signing server has the font file; the recipient's PDF viewer
probably doesn't.

\subsection{Embedded Fonts}\label{_embedded_fonts}

An embedded font includes the complete font program (or a subset of it)
as a stream inside the PDF. The viewer uses the embedded data to render
glyphs, regardless of what fonts are installed on the system.

Embedding guarantees consistent rendering across all viewers and all
platforms. It increases the PDF file size --- a full TrueType font can
be several megabytes --- but subsetting (described below) reduces the
impact dramatically.

The font program is stored as a stream referenced by the
\texttt{/FontDescriptor} dictionary:

\begin{longtable}[]{@{}
  >{\raggedright\arraybackslash}p{(\linewidth - 4\tabcolsep) * \real{0.2500}}
  >{\raggedright\arraybackslash}p{(\linewidth - 4\tabcolsep) * \real{0.2500}}
  >{\raggedright\arraybackslash}p{(\linewidth - 4\tabcolsep) * \real{0.5000}}@{}}
\caption{Font embedding streams by font type --- ISO 32000-2:2020 Table
123}\tabularnewline
\toprule\noalign{}
\begin{minipage}[b]{\linewidth}\raggedright
Font Type
\end{minipage} & \begin{minipage}[b]{\linewidth}\raggedright
FontDescriptor Key
\end{minipage} & \begin{minipage}[b]{\linewidth}\raggedright
Stream Contents
\end{minipage} \\
\midrule\noalign{}
\endfirsthead
\toprule\noalign{}
\begin{minipage}[b]{\linewidth}\raggedright
Font Type
\end{minipage} & \begin{minipage}[b]{\linewidth}\raggedright
FontDescriptor Key
\end{minipage} & \begin{minipage}[b]{\linewidth}\raggedright
Stream Contents
\end{minipage} \\
\midrule\noalign{}
\endhead
\bottomrule\noalign{}
\endlastfoot
Type 1 & \texttt{/FontFile} & PFB (Printer Font Binary) data --- the
Type 1 font program \\
TrueType & \texttt{/FontFile2} & Complete TrueType font file
(\texttt{.ttf}) \\
CFF / OpenType & \texttt{/FontFile3} & CFF (Compact Font Format) data,
or OpenType font program. The stream's \texttt{/Subtype} distinguishes:
\texttt{/CIDFontType0C} for CID-keyed CFF, \texttt{/Type1C} for simple
CFF, \texttt{/OpenType} for OpenType. \\
\end{longtable}

The distinction between \texttt{/FontFile}, \texttt{/FontFile2}, and
\texttt{/FontFile3} is not just metadata --- the stream contents are in
fundamentally different binary formats. A viewer that tries to parse a
CFF stream as TrueType (or vice versa) will get corrupted glyph outlines
or crash. Getting the key wrong is a common source of font rendering
failures in programmatically generated PDFs.

\section{Font Subsetting}\label{_font_subsetting}

Embedding a full font file --- especially a CJK font with tens of
thousands of glyphs --- adds significant bulk to the PDF. \textbf{Font
subsetting} solves this by embedding only the glyphs actually used in
the document.

A typical signature appearance uses fewer than 100 distinct glyphs (the
signer's name, a date, a reason string). Subsetting a 2 MB TrueType font
down to just those glyphs produces a subset of perhaps 10-30 KB.

\subsection{The Subset Prefix
Convention}\label{_the_subset_prefix_convention}

ISO 32000-2:2020 Section 9.6.4 defines a naming convention for subsetted
fonts: the \texttt{/BaseFont} name is prefixed with six uppercase
letters followed by a plus sign.

\begin{verbatim}
/BaseFont /ABCDEF+DejaVuSans
\end{verbatim}

The prefix "ABCDEF+" indicates that this is a subset --- it does not
contain all glyphs from the original DejaVuSans font. The six letters
are arbitrary (PDFBox generates them randomly); their only purpose is to
make the subset's \texttt{/BaseFont} name unique within the document,
preventing different subsets of the same font from being confused with
each other.

A font \emph{without} the subset prefix:

\begin{verbatim}
/BaseFont /DejaVuSans
\end{verbatim}

This indicates the full font is embedded (or the font is referenced
without embedding).

When debugging font issues, the subset prefix is your first clue. If you
see \texttt{/ABCDEF+Helvetica}, the font is subsetted. If a glyph is
missing from the rendered output, it may not have been included in the
subset --- meaning the glyph was not used (or not \emph{known} to be
used) when the subset was created. This is a common bug when signature
text is modified after the font was subsetted.

\subsection{How PDFBox Handles
Subsetting}\label{_how_pdfbox_handles_subsetting}

PDFBox 3.0.0 performs automatic font subsetting when saving a document.
When you load a TrueType font with \texttt{PDType0Font.load()}, PDFBox
tracks which glyphs are used during text drawing operations and includes
only those glyphs in the final embedded subset.

\begin{Shaded}
\begin{Highlighting}[]
\CommentTok{// PDFBox tracks glyph usage automatically}
\NormalTok{PDFont font }\OperatorTok{=}\NormalTok{ PDType0Font}\OperatorTok{.}\FunctionTok{load}\OperatorTok{(}\NormalTok{document}\OperatorTok{,} \KeywordTok{new} \BuiltInTok{File}\OperatorTok{(}\StringTok{"DejaVuSans.ttf"}\OperatorTok{));}

\CommentTok{// When we write text, PDFBox notes which glyphs are needed}
\ControlFlowTok{try} \OperatorTok{(}\BuiltInTok{OutputStream}\NormalTok{ os }\OperatorTok{=}\NormalTok{ formXObject}\OperatorTok{.}\FunctionTok{getStream}\OperatorTok{().}\FunctionTok{createOutputStream}\OperatorTok{())} \OperatorTok{\{}
    \CommentTok{// ... write content stream using this font ...}
\OperatorTok{\}}

\CommentTok{// On save, PDFBox subsets the font to include only used glyphs}
\NormalTok{document}\OperatorTok{.}\FunctionTok{save}\OperatorTok{(}\NormalTok{outputStream}\OperatorTok{);}
\end{Highlighting}
\end{Shaded}

There is one catch. If you write raw content stream operators directly
(as we must for signature appearance Form XObjects in PDFBox 3.x, since
\texttt{PDPageContentStream} does not support Form XObjects), PDFBox may
not automatically track glyph usage. In this case, you must encode the
text through the font to ensure the glyph tracking is triggered:

\begin{Shaded}
\begin{Highlighting}[]
\CommentTok{// Ensure glyph tracking when writing raw content streams}
\DataTypeTok{byte}\OperatorTok{[]}\NormalTok{ encodedText }\OperatorTok{=}\NormalTok{ font}\OperatorTok{.}\FunctionTok{encode}\OperatorTok{(}\StringTok{"Digitally signed by: Muller"}\OperatorTok{);}
\CommentTok{// Use encodedText in your content stream, not the raw string}
\end{Highlighting}
\end{Shaded}

If you bypass PDFBox's text encoding and write raw character codes
directly into the content stream, the subset may not include all
required glyphs. The result: some characters render as \texttt{.notdef}
(the missing glyph symbol, typically a small rectangle) in viewers. This
is one of the most common causes of broken signature appearances.

\section{Encoding}\label{_encoding}

Encoding is the mapping between \textbf{character codes} (the byte
values in the content stream) and \textbf{glyphs} (the visual shapes the
viewer renders). Getting encoding right is essential for correct text
rendering, and getting it wrong is the most common cause of garbled or
missing text in PDF documents.

\subsection{Simple Font Encoding
(Single-Byte)}\label{_simple_font_encoding_single_byte}

For simple fonts (Type 1 and simple TrueType), each character code is a
single byte (0-255). The encoding maps each byte value to a glyph name,
and the font program maps glyph names to outlines.

PDF defines two built-in encodings:

\textbf{WinAnsiEncoding} (ISO 32000-2:2020 Table D.1): Corresponds
closely to Windows-1252. Covers Western European characters: basic
Latin, accented characters (a, e, u, n, etc.), and common symbols. This
is the default for most PDF generators and the encoding used by the
Standard 14 fonts.

\textbf{MacRomanEncoding} (ISO 32000-2:2020 Table D.2): The classic
Macintosh character set. Rarely used in modern PDFs, but you'll
encounter it in documents created by older Mac applications.

A font dictionary can also specify a \textbf{custom encoding} via the
\texttt{/Encoding} key, using a \texttt{/Differences} array that
overrides individual code-to-glyph mappings:

\begin{verbatim}
/Encoding <<
  /Type /Encoding
  /BaseEncoding /WinAnsiEncoding
  /Differences [
    128 /Euro             
    160 /space /exclamdown /cent /sterling
  ]
>>
\end{verbatim}

\begin{itemize}
\item
  Maps byte value 128 to the Euro sign glyph, overriding the
  WinAnsiEncoding mapping
\end{itemize}

\subsection{The 256-Glyph Wall}\label{_the_256_glyph_wall}

Single-byte encoding means a simple font can address at most 256 glyphs.
This is fine for English and most Western European languages, but it's a
hard wall for:

\begin{itemize}
\item
  \textbf{CJK languages}: Chinese alone requires thousands of glyphs
\item
  \textbf{Arabic and Hebrew}: need contextual glyph shaping (the same
  character looks different depending on position in a word)
\item
  \textbf{Devanagari, Thai, Georgian}: scripts with large glyph
  inventories
\item
  \textbf{Mixed-script text}: a signature appearance showing a name in
  Cyrillic with a date in Latin
\end{itemize}

For any of these, you need a composite font (Type 0 + CIDFont) with a
multi-byte encoding.

\subsection{Composite Font Encoding}\label{_composite_font_encoding}

Composite fonts use a \textbf{CMap} (Character Map) to map multi-byte
character codes to CIDs. The most common CMap in modern PDFs is
\texttt{/Identity-H}, which maps character codes directly to CIDs as a
simple identity function --- code 0 maps to CID 0, code 1 to CID 1, and
so on.

With \texttt{/Identity-H}, the character codes in the content stream are
raw CID values encoded as two-byte big-endian integers. This gives
access to the full 65,536 CID range, more than enough for any single
font.

The downside of \texttt{/Identity-H} is that the character codes have no
inherent semantic meaning --- they're just glyph indices. Without
additional mapping information, a viewer cannot extract the text content
or determine what Unicode characters the glyphs represent. That's where
the \texttt{/ToUnicode} CMap comes in.

\section{The ToUnicode CMap}\label{_the_tounicode_cmap}

The \texttt{/ToUnicode} entry in a font dictionary maps character codes
to Unicode values. It is not used for rendering --- the viewer uses the
font's glyph outlines for that. Instead, it is used for \textbf{text
extraction}: copy-paste, search, and accessibility.

For signature appearances, the \texttt{/ToUnicode} CMap matters because:

\begin{enumerate}
\def\labelenumi{\arabic{enumi}.}
\item
  \textbf{Accessibility tools} read the \texttt{/ToUnicode} CMap to
  extract text from the appearance stream. A signature appearance
  without \texttt{/ToUnicode} mappings will appear as empty or garbled
  text to screen readers.
\item
  \textbf{PDF validators} (including veraPDF) check for the presence of
  \texttt{/ToUnicode} CMaps as a compliance requirement for PDF/A. If
  your signed PDFs must be PDF/A-compliant, the fonts in your signature
  appearance must have proper \texttt{/ToUnicode} entries.
\item
  \textbf{Text search} in the document depends on \texttt{/ToUnicode}.
  Without it, searching for the signer's name will not find the text in
  the signature appearance.
\end{enumerate}

A \texttt{/ToUnicode} CMap is a PDF stream containing a CMap program in
PostScript-like syntax:

\begin{verbatim}
/CIDInit /ProcSet findresource begin
12 dict begin
begincmap
/CIDSystemInfo
  << /Registry (Adobe) /Ordering (UCS) /Supplement 0 >> def
/CMapName /Adobe-Identity-UCS def
/CMapType 2 def
1 begincodespacerange
<0000> <FFFF>
endcodespacerange
5 beginbfchar
<0024> <0044>        % CID 36 -> U+0044 (D)
<004C> <0069>        % CID 76 -> U+0069 (i)
<004E> <0067>        % CID 78 -> U+0067 (g)
endbfchar
endcmap
CMapName currentdict /CMap defineresource pop
end
end
\end{verbatim}

PDFBox generates the \texttt{/ToUnicode} CMap automatically when you use
\texttt{PDType0Font.load()} and write text through the font's encoding
methods. You do not need to create it manually --- but understanding the
structure helps when you're debugging a PDF where text extraction
returns garbled results.

In PDFBox 3.x, the method \texttt{PDFont.getToUnicodeCMap()} is
\emph{not} publicly accessible. To inspect the \texttt{/ToUnicode} CMap
of a font, you must access the COS dictionary directly:
\texttt{font.getCOSObject().getDictionaryObject(COSName.TO\_UNICODE)}.
This is a common stumbling block when writing font inspection or
diagnostic tools.

\section{Font Descriptors}\label{_font_descriptors}

Every embedded font (and many referenced fonts) has a \textbf{font
descriptor} dictionary (\texttt{/Type\ /FontDescriptor}) that provides
metrics and metadata about the font. The font descriptor is
essential --- it tells the viewer how to render the font even before it
looks at the glyph outlines.

\begin{longtable}[]{@{}
  >{\raggedright\arraybackslash}p{(\linewidth - 4\tabcolsep) * \real{0.2500}}
  >{\raggedright\arraybackslash}p{(\linewidth - 4\tabcolsep) * \real{0.1250}}
  >{\raggedright\arraybackslash}p{(\linewidth - 4\tabcolsep) * \real{0.6250}}@{}}
\caption{Key entries in the font descriptor --- ISO 32000-2:2020 Table
122}\tabularnewline
\toprule\noalign{}
\begin{minipage}[b]{\linewidth}\raggedright
Key
\end{minipage} & \begin{minipage}[b]{\linewidth}\raggedright
Type
\end{minipage} & \begin{minipage}[b]{\linewidth}\raggedright
Purpose
\end{minipage} \\
\midrule\noalign{}
\endfirsthead
\toprule\noalign{}
\begin{minipage}[b]{\linewidth}\raggedright
Key
\end{minipage} & \begin{minipage}[b]{\linewidth}\raggedright
Type
\end{minipage} & \begin{minipage}[b]{\linewidth}\raggedright
Purpose
\end{minipage} \\
\midrule\noalign{}
\endhead
\bottomrule\noalign{}
\endlastfoot
\texttt{/FontName} & name & The PostScript name of the font (e.g.,
\texttt{/DejaVuSans}) \\
\texttt{/Flags} & integer & Bit field describing the font: FixedPitch
(bit 1), Serif (bit 2), Symbolic (bit 3), Script (bit 4), Nonsymbolic
(bit 6), Italic (bit 7), AllCap (bit 17), SmallCap (bit 18), ForceBold
(bit 19) \\
\texttt{/FontBBox} & array & The bounding box of all glyphs:
\texttt{{[}llx\ lly\ urx\ ury{]}} in glyph coordinate space \\
\texttt{/ItalicAngle} & number & Angle of the italic slant (0 for
upright fonts) \\
\texttt{/Ascent} & number & The maximum height above the baseline \\
\texttt{/Descent} & number & The maximum depth below the baseline
(typically negative) \\
\texttt{/CapHeight} & number & The height of capital letters \\
\texttt{/StemV} & number & The dominant stem width (vertical
strokes) --- affects rendering weight \\
\texttt{/FontFile} & stream & Type 1 font program (PFB format) \\
\texttt{/FontFile2} & stream & TrueType font program \\
\texttt{/FontFile3} & stream & CFF or OpenType font program (with
\texttt{/Subtype} to distinguish) \\
\end{longtable}

A font descriptor with all metrics but \emph{no} font file stream
describes a \textbf{referenced} font --- the viewer must find a
substitute. If the descriptor includes a font file stream, the font is
\textbf{embedded}. Checking whether \texttt{/FontFile},
\texttt{/FontFile2}, or \texttt{/FontFile3} exists is the definitive way
to determine if a font is embedded.

\section{PDFBox Font Handling in
3.x}\label{_pdfbox_font_handling_in_3_x}

PDFBox 3.0.0 made several breaking changes to its font API compared to
2.x. Understanding these changes is essential for writing correct font
handling code.

\subsection{Loading Fonts}\label{_loading_fonts}

\begin{Shaded}
\begin{Highlighting}[]
\CommentTok{// Standard 14 font (no embedding needed)}
\NormalTok{PDFont helvetica }\OperatorTok{=} \KeywordTok{new} \FunctionTok{PDType1Font}\OperatorTok{(}\NormalTok{Standard14Fonts}\OperatorTok{.}\FunctionTok{FontName}\OperatorTok{.}\FunctionTok{HELVETICA}\OperatorTok{);}

\CommentTok{// Embed a TrueType font from a file}
\NormalTok{PDFont dejaVu }\OperatorTok{=}\NormalTok{ PDType0Font}\OperatorTok{.}\FunctionTok{load}\OperatorTok{(}\NormalTok{document}\OperatorTok{,} \KeywordTok{new} \BuiltInTok{File}\OperatorTok{(}\StringTok{"DejaVuSans.ttf"}\OperatorTok{));}

\CommentTok{// Embed a TrueType font from an InputStream}
\ControlFlowTok{try} \OperatorTok{(}\BuiltInTok{InputStream}\NormalTok{ is }\OperatorTok{=} \FunctionTok{getClass}\OperatorTok{().}\FunctionTok{getResourceAsStream}\OperatorTok{(}\StringTok{"/fonts/NotoSans{-}Regular.ttf"}\OperatorTok{))} \OperatorTok{\{}
\NormalTok{    PDFont notoSans }\OperatorTok{=}\NormalTok{ PDType0Font}\OperatorTok{.}\FunctionTok{load}\OperatorTok{(}\NormalTok{document}\OperatorTok{,}\NormalTok{ is}\OperatorTok{);}
\OperatorTok{\}}

\CommentTok{// Embed a TrueType font WITHOUT subsetting}
\CommentTok{// (include all glyphs {-}{-} larger file, but no risk of missing glyphs)}
\NormalTok{PDFont fullFont }\OperatorTok{=}\NormalTok{ PDType0Font}\OperatorTok{.}\FunctionTok{load}\OperatorTok{(}\NormalTok{document}\OperatorTok{,} \KeywordTok{new} \BuiltInTok{File}\OperatorTok{(}\StringTok{"DejaVuSans.ttf"}\OperatorTok{),} \KeywordTok{false}\OperatorTok{);}
\end{Highlighting}
\end{Shaded}

The \texttt{false} parameter to \texttt{PDType0Font.load()} disables
subsetting. This is rarely needed, but can be useful during debugging to
rule out subsetting as the cause of missing glyphs.

\subsection{The getEncoding() Removal}\label{_the_getencoding_removal}

PDFBox 2.x provided a \texttt{PDFont.getEncoding()} method that returned
the font's encoding object. PDFBox 3.x \textbf{removed} this method.

If you need to inspect a font's encoding in PDFBox 3.0.0, you must
access the COS dictionary directly:

\begin{Shaded}
\begin{Highlighting}[]
\CommentTok{// PDFBox 2.x {-}{-} no longer works}
\CommentTok{// Encoding encoding = font.getEncoding();}

\CommentTok{// PDFBox 3.x {-}{-} access the COS dictionary}
\NormalTok{COSBase encodingObj }\OperatorTok{=}\NormalTok{ font}\OperatorTok{.}\FunctionTok{getCOSObject}\OperatorTok{().}\FunctionTok{getDictionaryObject}\OperatorTok{(}\NormalTok{COSName}\OperatorTok{.}\FunctionTok{ENCODING}\OperatorTok{);}

\ControlFlowTok{if} \OperatorTok{(}\NormalTok{encodingObj }\KeywordTok{instanceof}\NormalTok{ COSName encodingName}\OperatorTok{)} \OperatorTok{\{}
    \CommentTok{// Named encoding: /WinAnsiEncoding, /MacRomanEncoding, etc.}
\NormalTok{    log}\OperatorTok{.}\FunctionTok{info}\OperatorTok{(}\StringTok{"Font uses named encoding: \{\}"}\OperatorTok{,}\NormalTok{ encodingName}\OperatorTok{.}\FunctionTok{getName}\OperatorTok{());}
\OperatorTok{\}} \ControlFlowTok{else} \ControlFlowTok{if} \OperatorTok{(}\NormalTok{encodingObj }\KeywordTok{instanceof}\NormalTok{ COSDictionary encodingDict}\OperatorTok{)} \OperatorTok{\{}
    \CommentTok{// Custom encoding with /Differences array}
\NormalTok{    COSArray differences }\OperatorTok{=}\NormalTok{ encodingDict}\OperatorTok{.}\FunctionTok{getDictionaryObject}\OperatorTok{(}
\NormalTok{            COSName}\OperatorTok{.}\FunctionTok{getPDFName}\OperatorTok{(}\StringTok{"Differences"}\OperatorTok{),}\NormalTok{ COSArray}\OperatorTok{.}\FunctionTok{class}\OperatorTok{);}
    \ControlFlowTok{if} \OperatorTok{(}\NormalTok{differences }\OperatorTok{!=} \KeywordTok{null}\OperatorTok{)} \OperatorTok{\{}
\NormalTok{        log}\OperatorTok{.}\FunctionTok{info}\OperatorTok{(}\StringTok{"Font has custom encoding with \{\} difference entries"}\OperatorTok{,}
\NormalTok{                differences}\OperatorTok{.}\FunctionTok{size}\OperatorTok{());}
    \OperatorTok{\}}
\OperatorTok{\}}
\end{Highlighting}
\end{Shaded}

\subsection{Inspecting Font Embedding
Programmatically}\label{_inspecting_font_embedding_programmatically}

A common diagnostic task is checking whether a font in a PDF is
embedded. Here's a complete example:

\begin{Shaded}
\begin{Highlighting}[]
\KeywordTok{private} \DataTypeTok{boolean} \FunctionTok{isFontEmbedded}\OperatorTok{(}\NormalTok{PDFont font}\OperatorTok{)} \OperatorTok{\{}
\NormalTok{    PDFontDescriptor descriptor }\OperatorTok{=}\NormalTok{ font}\OperatorTok{.}\FunctionTok{getFontDescriptor}\OperatorTok{();}
    \ControlFlowTok{if} \OperatorTok{(}\NormalTok{descriptor }\OperatorTok{==} \KeywordTok{null}\OperatorTok{)} \OperatorTok{\{}
        \CommentTok{// Standard 14 fonts may not have a descriptor}
        \ControlFlowTok{return} \KeywordTok{false}\OperatorTok{;}
    \OperatorTok{\}}
    \ControlFlowTok{return}\NormalTok{ descriptor}\OperatorTok{.}\FunctionTok{getFontFile}\OperatorTok{()} \OperatorTok{!=} \KeywordTok{null}
        \OperatorTok{||}\NormalTok{ descriptor}\OperatorTok{.}\FunctionTok{getFontFile2}\OperatorTok{()} \OperatorTok{!=} \KeywordTok{null}
        \OperatorTok{||}\NormalTok{ descriptor}\OperatorTok{.}\FunctionTok{getFontFile3}\OperatorTok{()} \OperatorTok{!=} \KeywordTok{null}\OperatorTok{;}
\OperatorTok{\}}
\end{Highlighting}
\end{Shaded}

\section{Why Signature Text Shows as Squares or
Blanks}\label{_why_signature_text_shows_as_squares_or_blanks}

This section addresses the most common font-related failures in
signature appearances, what causes them, and how to fix them.

\subsection{Missing Font}\label{_missing_font}

\textbf{Symptom}: The signature block is empty, or shows substituted
text in a different font than intended.

\textbf{Cause}: The signature appearance references a font that is not
embedded and not available on the viewer's system.

\textbf{Fix}: Always embed fonts used in signature appearances. For
basic ASCII text, the Standard 14 fonts are a safe choice because
viewers provide substitutes. For anything else, embed a TrueType font
using \texttt{PDType0Font.load()}.

\subsection{Wrong Encoding}\label{_wrong_encoding}

\textbf{Symptom}: Some characters render correctly, but accented
characters (e, u, n) show as wrong characters or \texttt{.notdef}
squares.

\textbf{Cause}: The text contains characters outside the font's encoding
range. WinAnsiEncoding covers codes 32-255 corresponding to
Windows-1252, which includes most Western European accented characters.
But characters outside this range --- em-dash at certain positions,
certain ligatures, or any non-Latin script --- will fail.

\textbf{Fix}: Use a composite font (Type 0) with \texttt{/Identity-H}
encoding and an embedded TrueType font. This provides full Unicode
coverage. In PDFBox: use \texttt{PDType0Font.load()} instead of the
Standard 14 fonts.

\subsection{Unembedded Font with Non-Standard
Glyphs}\label{_unembedded_font_with_non_standard_glyphs}

\textbf{Symptom}: Text renders as squares or blank space on some viewers
but works on others.

\textbf{Cause}: The font is referenced but not embedded. Viewers that
happen to have a compatible font installed render the text correctly;
viewers without the font substitute a generic font that may not have the
required glyphs.

\textbf{Fix}: Embed the font.

\subsection{Subset Missing Glyphs}\label{_subset_missing_glyphs}

\textbf{Symptom}: Most text renders correctly, but specific characters
are missing.

\textbf{Cause}: The font was subsetted, but the subset does not include
all required glyphs. This typically happens when text is written to the
content stream \emph{after} the font subsetting decision was made, or
when glyph usage tracking was bypassed.

\textbf{Fix}: Ensure all text is written through PDFBox's font encoding
methods \emph{before} saving. If the problem persists, try loading the
font with subsetting disabled
(\texttt{PDType0Font.load(document,\ file,\ false)}) to confirm
subsetting is the cause.

When writing raw content stream operators for a Form XObject (which
PDFBox 3.x requires for signature appearances), make sure to encode text
through the font \emph{before} writing the encoded bytes to the stream.
Skipping this step means PDFBox won't know which glyphs are needed and
may produce a subset that's missing characters.

\begin{Shaded}
\begin{Highlighting}[]
\CommentTok{// WRONG: writing raw text bypasses glyph tracking}
\NormalTok{os}\OperatorTok{.}\FunctionTok{write}\OperatorTok{(}\StringTok{"(John Muller) Tj"}\OperatorTok{.}\FunctionTok{getBytes}\OperatorTok{(}\NormalTok{StandardCharsets}\OperatorTok{.}\FunctionTok{US\_ASCII}\OperatorTok{));}

\CommentTok{// RIGHT: encode through the font first}
\DataTypeTok{byte}\OperatorTok{[]}\NormalTok{ encoded }\OperatorTok{=}\NormalTok{ font}\OperatorTok{.}\FunctionTok{encode}\OperatorTok{(}\StringTok{"John Muller"}\OperatorTok{);}
\CommentTok{// Then write the encoded bytes as a hex string in the content stream}
\NormalTok{os}\OperatorTok{.}\FunctionTok{write}\OperatorTok{((}\StringTok{"\textless{}"} \OperatorTok{+} \FunctionTok{bytesToHex}\OperatorTok{(}\NormalTok{encoded}\OperatorTok{)} \OperatorTok{+} \StringTok{"\textgreater{} Tj"}\OperatorTok{).}\FunctionTok{getBytes}\OperatorTok{(}\NormalTok{StandardCharsets}\OperatorTok{.}\FunctionTok{US\_ASCII}\OperatorTok{));}
\end{Highlighting}
\end{Shaded}

\section{Font Licensing and the fsType
Flag}\label{_font_licensing_and_the_fstype_flag}

TrueType and OpenType fonts contain a licensing flag called
\textbf{fsType} in the OS/2 table. This flag indicates what the font
vendor permits in terms of embedding and editing:

\begin{longtable}[]{@{}
  >{\raggedright\arraybackslash}p{(\linewidth - 4\tabcolsep) * \real{0.1429}}
  >{\raggedright\arraybackslash}p{(\linewidth - 4\tabcolsep) * \real{0.2857}}
  >{\raggedright\arraybackslash}p{(\linewidth - 4\tabcolsep) * \real{0.5714}}@{}}
\caption{fsType embedding permission bits}\tabularnewline
\toprule\noalign{}
\begin{minipage}[b]{\linewidth}\raggedright
Bits
\end{minipage} & \begin{minipage}[b]{\linewidth}\raggedright
Name
\end{minipage} & \begin{minipage}[b]{\linewidth}\raggedright
Meaning
\end{minipage} \\
\midrule\noalign{}
\endfirsthead
\toprule\noalign{}
\begin{minipage}[b]{\linewidth}\raggedright
Bits
\end{minipage} & \begin{minipage}[b]{\linewidth}\raggedright
Name
\end{minipage} & \begin{minipage}[b]{\linewidth}\raggedright
Meaning
\end{minipage} \\
\midrule\noalign{}
\endhead
\bottomrule\noalign{}
\endlastfoot
0x0000 & Installable & No restrictions. The font may be embedded and
used freely. \\
0x0002 & Restricted & The font must not be embedded in documents.
Period. \\
0x0004 & Preview \& Print & The font may be embedded in read-only
documents. The recipient can view and print but should not edit. \\
0x0008 & Editable & The font may be embedded, and the recipient may edit
the document. \\
\end{longtable}

Most free and open-source fonts (DejaVu, Noto, Liberation) have fsType
0x0000 (Installable). Many commercial fonts (fonts you pay for) have
restrictive fsType values.

PDF generators typically respect the fsType flag --- PDFBox will emit a
warning (but not throw an exception) when embedding a font with
restrictive fsType. Some tools refuse to embed the font entirely.

For signature appearances:

\begin{itemize}
\item
  \textbf{Use fonts with fsType 0x0000} (Installable) to avoid any
  licensing complications.
\item
  The \textbf{Google Noto font family} (\texttt{Noto\ Sans},
  \texttt{Noto\ Serif}, \texttt{Noto\ Sans\ CJK}) covers virtually every
  script in Unicode and is Apache-2.0 licensed with fsType 0x0000. It is
  the best general-purpose choice for internationalized signature text.
\item
  \textbf{DejaVu Sans} covers Latin, Greek, Cyrillic, Armenian,
  Georgian, and many symbols. Free, well-tested, and widely deployed.
\item
  \textbf{Liberation Sans/Serif/Mono} are metrics-compatible with
  Arial/Times New Roman/Courier New. Licensed under SIL Open Font
  License.
\end{itemize}

Even if your PDF generator embeds a font with restrictive fsType, the
PDF viewer will render it. The fsType flag is a licensing statement, not
a technical enforcement mechanism. But distributing PDFs with improperly
embedded commercial fonts exposes your organization to licensing claims.
For server-side signing operations that produce thousands of PDFs, use
only fonts whose license explicitly permits embedding.

\section{Practical Advice for Signature
Appearances}\label{_practical_advice_for_signature_appearances}

Drawing from everything in this chapter, here is the decision framework
for choosing and configuring fonts in signature appearance streams:

\subsection{For ASCII-Only Signatures}\label{_for_ascii_only_signatures}

If your signer names, dates, and reasons are guaranteed to be ASCII only
(a-z, A-Z, 0-9, basic punctuation):

\begin{Shaded}
\begin{Highlighting}[]
\CommentTok{// Simple, reliable, zero{-}cost approach}
\NormalTok{PDFont font }\OperatorTok{=} \KeywordTok{new} \FunctionTok{PDType1Font}\OperatorTok{(}\NormalTok{Standard14Fonts}\OperatorTok{.}\FunctionTok{FontName}\OperatorTok{.}\FunctionTok{HELVETICA}\OperatorTok{);}
\end{Highlighting}
\end{Shaded}

This works everywhere. No embedding needed, no licensing concerns, no
subsetting complications. The text will render in Helvetica or a
metrics-compatible substitute on every viewer.

\subsection{For International
Signatures}\label{_for_international_signatures}

If signer names may contain non-ASCII characters (accented Latin,
Cyrillic, CJK, Arabic, or anything beyond Windows-1252):

\begin{Shaded}
\begin{Highlighting}[]
\CommentTok{// Embed a Unicode{-}capable font}
\ControlFlowTok{try} \OperatorTok{(}\BuiltInTok{InputStream}\NormalTok{ is }\OperatorTok{=} \FunctionTok{getClass}\OperatorTok{().}\FunctionTok{getResourceAsStream}\OperatorTok{(}\StringTok{"/fonts/NotoSans{-}Regular.ttf"}\OperatorTok{))} \OperatorTok{\{}
\NormalTok{    PDFont font }\OperatorTok{=}\NormalTok{ PDType0Font}\OperatorTok{.}\FunctionTok{load}\OperatorTok{(}\NormalTok{document}\OperatorTok{,}\NormalTok{ is}\OperatorTok{);}
    \CommentTok{// Now you can write any Unicode character the font covers}
\OperatorTok{\}}
\end{Highlighting}
\end{Shaded}

Bundle the font file with your application (in
\texttt{src/main/resources/fonts/} for a Maven project). Do not rely on
system-installed fonts --- the signing server may not have the same
fonts as your development machine.

\subsection{Font Size Guidelines}\label{_font_size_guidelines}

Signature appearance text is typically rendered at small sizes (8-12
points) in a rectangle that's often only 150-300 points wide. Keep these
guidelines in mind:

\begin{itemize}
\item
  \textbf{10pt} is a good default for signer name and date
\item
  \textbf{8pt} works for secondary information (reason, location,
  certificate details)
\item
  \textbf{Never go below 7pt} --- some viewers render small text poorly,
  and the text becomes unreadable when printed
\item
  \textbf{Calculate text width} before rendering to ensure it fits the
  signature rectangle; truncation or overflow beyond the bounding box
  results in clipped text that viewers may or may not display
\end{itemize}

\section{Debugging Font Issues}\label{_debugging_font_issues}

When a signature appearance renders incorrectly, the following
diagnostic steps will identify the root cause.

\subsection{Step 1: Identify the Font}\label{_step_1_identify_the_font}

Open the PDF in a tool that shows the raw object structure (PDFalyzer,
QPDF, or \texttt{PDFDebugger} from PDFBox's tools module). Navigate to
the signature field's widget annotation, then to \texttt{/AP\ /N} (the
normal appearance stream), then to \texttt{/Resources\ /Font}. Note the
font name, subtype, and whether a font descriptor with a font file
stream exists.

\subsection{Step 2: Check Embedding}\label{_step_2_check_embedding}

Look for \texttt{/FontFile}, \texttt{/FontFile2}, or \texttt{/FontFile3}
in the font descriptor. If none exists, the font is not
embedded --- that's likely the problem.

\subsection{Step 3: Check Encoding}\label{_step_3_check_encoding}

For simple fonts, check the \texttt{/Encoding} entry. Is it
\texttt{/WinAnsiEncoding}? If the problematic characters are outside
Windows-1252, the encoding is the cause.

For composite fonts, check the \texttt{/Encoding} entry on the Type 0
dictionary. It should be \texttt{/Identity-H} for a properly structured
modern PDF.

\subsection{Step 4: Check the ToUnicode
CMap}\label{_step_4_check_the_tounicode_cmap}

If the text \emph{renders} correctly but copy-paste or search returns
garbled results, the \texttt{/ToUnicode} CMap is missing or incorrect.

\subsection{Step 5: veraPDF
Validation}\label{_step_5_verapdf_validation}

veraPDF provides font-specific validation rules that catch a wide range
of font issues:

\begin{Shaded}
\begin{Highlighting}[]
\CommentTok{\# Run veraPDF validation focusing on font issues}
\ExtensionTok{verapdf} \AttributeTok{{-}{-}format}\NormalTok{ text }\AttributeTok{{-}{-}flavour}\NormalTok{ 2b input.pdf }\DecValTok{2}\OperatorTok{\textgreater{}\&}\DecValTok{1} \KeywordTok{|} \FunctionTok{grep} \AttributeTok{{-}i}\NormalTok{ font}
\end{Highlighting}
\end{Shaded}

veraPDF checks for embedded fonts (required for PDF/A), valid
\texttt{/ToUnicode} mappings, correct font descriptor entries, and
consistent glyph widths. Its validation messages are often the fastest
way to identify exactly what's wrong with a font configuration.

\subsection{Step 6: PDFBox Font
Inspection}\label{_step_6_pdfbox_font_inspection}

For programmatic inspection, PDFBox can enumerate all fonts in a page or
appearance stream:

\begin{Shaded}
\begin{Highlighting}[]
\KeywordTok{private} \DataTypeTok{void} \FunctionTok{inspectAppearanceFonts}\OperatorTok{(}\NormalTok{PDAnnotationWidget widget}\OperatorTok{)} \OperatorTok{\{}
\NormalTok{    PDAppearanceStream appearance }\OperatorTok{=}\NormalTok{ widget}\OperatorTok{.}\FunctionTok{getAppearance}\OperatorTok{().}\FunctionTok{getNormalAppearance}\OperatorTok{();}
    \ControlFlowTok{if} \OperatorTok{(}\NormalTok{appearance }\OperatorTok{==} \KeywordTok{null}\OperatorTok{)} \OperatorTok{\{}
\NormalTok{        log}\OperatorTok{.}\FunctionTok{warn}\OperatorTok{(}\StringTok{"Signature widget has no normal appearance stream"}\OperatorTok{);}
        \ControlFlowTok{return}\OperatorTok{;}
    \OperatorTok{\}}

\NormalTok{    PDResources resources }\OperatorTok{=}\NormalTok{ appearance}\OperatorTok{.}\FunctionTok{getResources}\OperatorTok{();}
    \ControlFlowTok{if} \OperatorTok{(}\NormalTok{resources }\OperatorTok{==} \KeywordTok{null}\OperatorTok{)} \OperatorTok{\{}
\NormalTok{        log}\OperatorTok{.}\FunctionTok{warn}\OperatorTok{(}\StringTok{"Appearance stream has no resources"}\OperatorTok{);}
        \ControlFlowTok{return}\OperatorTok{;}
    \OperatorTok{\}}

    \ControlFlowTok{for} \OperatorTok{(}\NormalTok{COSName fontName }\OperatorTok{:}\NormalTok{ resources}\OperatorTok{.}\FunctionTok{getFontNames}\OperatorTok{())} \OperatorTok{\{}
\NormalTok{        PDFont font }\OperatorTok{=}\NormalTok{ resources}\OperatorTok{.}\FunctionTok{getFont}\OperatorTok{(}\NormalTok{fontName}\OperatorTok{);}
\NormalTok{        log}\OperatorTok{.}\FunctionTok{info}\OperatorTok{(}\StringTok{"Font /\{\}: type=\{\}, baseFont=\{\}, embedded=\{\}"}\OperatorTok{,}
\NormalTok{                fontName}\OperatorTok{.}\FunctionTok{getName}\OperatorTok{(),}
\NormalTok{                font}\OperatorTok{.}\FunctionTok{getClass}\OperatorTok{().}\FunctionTok{getSimpleName}\OperatorTok{(),}
\NormalTok{                font}\OperatorTok{.}\FunctionTok{getName}\OperatorTok{(),}
                \FunctionTok{isFontEmbedded}\OperatorTok{(}\NormalTok{font}\OperatorTok{));}

        \CommentTok{// Check for ToUnicode CMap}
\NormalTok{        COSBase toUnicode }\OperatorTok{=}\NormalTok{ font}\OperatorTok{.}\FunctionTok{getCOSObject}\OperatorTok{()}
                \OperatorTok{.}\FunctionTok{getDictionaryObject}\OperatorTok{(}\NormalTok{COSName}\OperatorTok{.}\FunctionTok{TO\_UNICODE}\OperatorTok{);}
\NormalTok{        log}\OperatorTok{.}\FunctionTok{info}\OperatorTok{(}\StringTok{"  ToUnicode CMap: \{\}"}\OperatorTok{,}
\NormalTok{                toUnicode }\OperatorTok{!=} \KeywordTok{null} \OperatorTok{?} \StringTok{"present"} \OperatorTok{:} \StringTok{"MISSING"}\OperatorTok{);}
    \OperatorTok{\}}
\OperatorTok{\}}
\end{Highlighting}
\end{Shaded}

When debugging font issues in production, always capture the raw bytes
of the appearance stream and the font dictionary's COS structure. The
visual symptom (squares, blanks, wrong characters) tells you
\emph{something} is wrong; the COS structure tells you \emph{what}. A
tool that shows both the rendered appearance and the underlying object
tree --- the way PDFalyzer does --- saves enormous debugging time
compared to flipping between a viewer and a hex editor.

\section{Summary}\label{_summary_9}

Fonts in PDF are a deceptively deep topic. For digital signatures
specifically, the key takeaways are:

\begin{itemize}
\item
  \textbf{Embed your fonts.} The Standard 14 fonts are safe for ASCII
  text; everything else must be embedded.
\item
  \textbf{Use composite fonts (Type 0) for non-ASCII text.} Simple fonts
  are limited to 256 glyphs.
\item
  \textbf{Encode text through the font API} to ensure glyph tracking for
  subsetting works correctly.
\item
  \textbf{Use freely licensed fonts} (Noto Sans, DejaVu, Liberation) to
  avoid fsType complications.
\item
  \textbf{Test in multiple viewers} --- what works in Adobe Acrobat may
  fail in pdf.js or LibreOffice.
\item
  \textbf{When debugging, look at the COS structure}, not just the
  rendered output. The font dictionary, font descriptor, encoding, and
  \texttt{/ToUnicode} CMap tell the full story.
\end{itemize}

The next chapter --- Chapter 12 --- moves from the document
infrastructure we've built up across Part II into the core subject of
this book: the architecture of PDF digital signatures themselves.

\chapter{PDF Digital Signatures}\label{_pdf_digital_signatures}

\section{Signature Architecture}\label{ch12}

\begin{quote}
This chapter dissects the internal architecture of PDF digital
signatures as specified in ISO 32000-2:2020 §12.8. We examine every
dictionary, every key, and every byte --- from the AcroForm field tree
down to the raw CMS \texttt{SignedData} blob embedded in
\texttt{/Contents}.

But more importantly, we talk about what actually matters when you're
staring at a broken signature at 2 AM, wondering why Adobe says
"invalid" while your code says "fine."
\end{quote}

\subsection{Overview: What a PDF Signature Actually
Is}\label{_overview_what_a_pdf_signature_actually_is}

Here's the thing most people get wrong about PDF signatures: they think
of them as something \emph{attached to} a document, like a wax seal on
an envelope. They're not.

A PDF digital signature is a set of \textbf{PDF objects embedded within
the document itself}. The signature \emph{is part of the document}. This
seemingly minor architectural choice has profound consequences --- it's
why PDF signatures are simultaneously powerful and maddening, why they
survive copy operations but break on re-saves, and why you can have a
perfectly valid signature on a document that's been modified after
signing (yes, really).

The four pieces of a PDF signature:

\begin{enumerate}
\def\labelenumi{\arabic{enumi}.}
\item
  A \textbf{signature field} (\texttt{PDSignatureField}) in the AcroForm
  field tree --- the form field that "holds" the signature
\item
  A \textbf{signature dictionary} (\texttt{/Type\ /Sig}) --- the actual
  cryptographic payload and metadata
\item
  A \textbf{widget annotation} --- the visual rectangle on a page (may
  be zero-size for invisible signatures)
\item
  An \textbf{appearance stream} (\texttt{/AP}) --- what the viewer
  renders in the signature rectangle
\end{enumerate}

The signature field is a form field like any other (\texttt{/FT\ /Sig}).
The signature dictionary is the \emph{value} of that field
(\texttt{/V}). An unsigned field has \texttt{/V} absent or null --- the
field exists but has no signature value yet.

Think of it like an empty checkbox on a paper form. The box is there,
waiting to be checked. A document author can place signature fields in a
PDF template, and signers fill them in later --- each one independently,
at different times, possibly on different continents.

\textbf{When you encounter a PDF with signature problems, start here.}
Open the document in a COS-level viewer (like PDFalyzer, or even a hex
editor) and check:

\begin{enumerate}
\def\labelenumi{\arabic{enumi}.}
\item
  Does the AcroForm exist? (Some PDFs have signature fields floating
  without an AcroForm --- instant headache.)
\item
  Does \texttt{/SigFlags} include bit 2 (AppendOnly)? If not, any editor
  can destroy all signatures with a simple save.
\item
  Is the \texttt{/V} entry present on the signature field? If absent,
  the field is unsigned.
\item
  Does the widget actually have a page reference? Some tools create
  orphaned widgets.
\end{enumerate}

This five-minute checklist saves hours of debugging.

\subsection{Signature Fields in the AcroForm
Tree}\label{ch12-field-tree}

\begin{quote}
A signature field (field type Sig) shall represent a digital signature
that may or may not have a signature value.

---  ISO 32000-2:2020 §12.7.4.5
\end{quote}

Signature fields live in the document's interactive form structure:

\subsubsection{The SigFlags Trap}\label{_the_sigflags_trap}

The \texttt{/SigFlags} entry is one of those things that's easy to
overlook and devastating to get wrong.

\begin{longtable}[]{@{}
  >{\raggedright\arraybackslash}p{(\linewidth - 8\tabcolsep) * \real{0.1818}}
  >{\raggedright\arraybackslash}p{(\linewidth - 8\tabcolsep) * \real{0.0909}}
  >{\raggedright\arraybackslash}p{(\linewidth - 8\tabcolsep) * \real{0.0909}}
  >{\raggedright\arraybackslash}p{(\linewidth - 8\tabcolsep) * \real{0.1818}}
  >{\raggedright\arraybackslash}p{(\linewidth - 8\tabcolsep) * \real{0.4545}}@{}}
\caption{SigFlags --- ISO 32000-2:2020 §12.7.2, Table
225}\tabularnewline
\toprule\noalign{}
\begin{minipage}[b]{\linewidth}\raggedright
Name
\end{minipage} & \begin{minipage}[b]{\linewidth}\raggedright
Bit
\end{minipage} & \begin{minipage}[b]{\linewidth}\raggedright
Req
\end{minipage} & \begin{minipage}[b]{\linewidth}\raggedright
Spec Ref
\end{minipage} & \begin{minipage}[b]{\linewidth}\raggedright
Effect
\end{minipage} \\
\midrule\noalign{}
\endfirsthead
\toprule\noalign{}
\begin{minipage}[b]{\linewidth}\raggedright
Name
\end{minipage} & \begin{minipage}[b]{\linewidth}\raggedright
Bit
\end{minipage} & \begin{minipage}[b]{\linewidth}\raggedright
Req
\end{minipage} & \begin{minipage}[b]{\linewidth}\raggedright
Spec Ref
\end{minipage} & \begin{minipage}[b]{\linewidth}\raggedright
Effect
\end{minipage} \\
\midrule\noalign{}
\endhead
\bottomrule\noalign{}
\endlastfoot
\texttt{SignaturesExist} & 1 & Required & §12.7.2 & Tells the viewer
"this document contains at least one signature field." Conforming
readers \textbf{shall not} allow saving in a way that invalidates
existing signatures. \\
\texttt{AppendOnly} & 2 & Required & §12.7.4.5 & The document
\textbf{shall} be saved using incremental updates only. This is the
mechanism that protects existing signatures from being destroyed by a
full rewrite. \\
\end{longtable}

Setting \texttt{/SigFlags} to 3 (both bits) is mandatory when a document
contains signatures. If bit 2 is missing, a naive PDF editor may do a
full save and destroy all signatures. Many real-world PDFs get this
wrong.

Here's what happens in practice: someone creates a PDF with LibreOffice,
adds a signature field, signs it --- and LibreOffice doesn't set bit 2.
The document gets emailed to a colleague who opens it in some random PDF
editor to add an annotation, hits save, and the editor does a full
rewrite. Every signature is now invalid. The colleague has no idea what
happened. The signer gets an angry phone call.

\textbf{The fix is always the same:} if you're building a tool that
creates signature fields, set \texttt{/SigFlags} to 3 the moment any
field has a value. And if you're building a tool that edits PDFs, check
for \texttt{/SigFlags} bit 2 before saving --- if it's set, you
\emph{must} use incremental save. If your tool doesn't support
incremental save\ldots\hspace{0pt} well, you need to support it now.
Welcome to PDF.

\subsubsection{Field Dictionary
Structure}\label{_field_dictionary_structure}

A signature field's COS dictionary looks like this. Every key matters,
and I'll explain why:

Here is the complete COS dictionary of a typical signature field. In
most PDFs, the field and its widget annotation are merged into a single
object --- the spec (ISO 32000-2:2020 §12.5.6.19) allows splitting them,
but few tools bother:

\begin{verbatim}
<<
  /Type /Annot
  /Subtype /Widget
  /FT /Sig
  /T (Signature1)
  /V 42 0 R
  /Rect [100 500 300 550]
  /AP << /N 43 0 R >>
  /P 7 0 R
  /Lock <<
    /Type /SigFieldLock
    /Action /All
  >>
  /SV <<
    /Type /SV
    /Filter /Adobe.PPKLite
    /SubFilter [/adbe.pkcs7.detached /ETSI.CAdES.detached]
    /DigestMethod [/SHA256 /SHA384 /SHA512]
    /Reasons [(Approved)(Reviewed)]
    /MDP << /P 2 >>
  >>
>>
\end{verbatim}

The keys fall into three groups: annotation/widget keys (how it looks
and where it lives), signature field keys (what it is and what it points
to), and sub-dictionaries that constrain signing behavior.

\begin{longtable}[]{@{}
  >{\raggedright\arraybackslash}p{(\linewidth - 8\tabcolsep) * \real{0.1667}}
  >{\raggedright\arraybackslash}p{(\linewidth - 8\tabcolsep) * \real{0.1667}}
  >{\raggedright\arraybackslash}p{(\linewidth - 8\tabcolsep) * \real{0.0833}}
  >{\raggedright\arraybackslash}p{(\linewidth - 8\tabcolsep) * \real{0.1667}}
  >{\raggedright\arraybackslash}p{(\linewidth - 8\tabcolsep) * \real{0.4167}}@{}}
\caption{Annotation/Widget Keys --- ISO 32000-2:2020
§12.5.6.19}\tabularnewline
\toprule\noalign{}
\begin{minipage}[b]{\linewidth}\raggedright
Key
\end{minipage} & \begin{minipage}[b]{\linewidth}\raggedright
Type
\end{minipage} & \begin{minipage}[b]{\linewidth}\raggedright
Req
\end{minipage} & \begin{minipage}[b]{\linewidth}\raggedright
Spec Ref
\end{minipage} & \begin{minipage}[b]{\linewidth}\raggedright
Description
\end{minipage} \\
\midrule\noalign{}
\endfirsthead
\toprule\noalign{}
\begin{minipage}[b]{\linewidth}\raggedright
Key
\end{minipage} & \begin{minipage}[b]{\linewidth}\raggedright
Type
\end{minipage} & \begin{minipage}[b]{\linewidth}\raggedright
Req
\end{minipage} & \begin{minipage}[b]{\linewidth}\raggedright
Spec Ref
\end{minipage} & \begin{minipage}[b]{\linewidth}\raggedright
Description
\end{minipage} \\
\midrule\noalign{}
\endhead
\bottomrule\noalign{}
\endlastfoot
\texttt{/Type} & name & Optional & §12.5.6 & Always \texttt{/Annot} for
the merged field+widget. \\
\texttt{/Subtype} & name & Required & §12.5.6 & Always \texttt{/Widget}
--- this is a widget annotation. \\
\texttt{/Rect} & rectangle & Required & §12.5.6 & Widget rectangle in
page coordinates: \texttt{{[}llx\ lly\ urx\ ury{]}}. If
\texttt{{[}0\ 0\ 0\ 0{]}}, the signature is invisible. \\
\texttt{/P} & reference & Optional & §12.5.6 & Page reference --- which
page this widget is on. Some poorly generated PDFs omit this, and then
viewers have to guess. Don't be that tool. \\
\texttt{/AP} & dictionary & Optional & §12.5.6.19 & Appearance stream
dictionary. \texttt{/N} is the normal appearance --- what the viewer
draws. Without this, many viewers show a generic placeholder. Adobe is
particularly opinionated about appearances. \\
\end{longtable}

\begin{longtable}[]{@{}
  >{\raggedright\arraybackslash}p{(\linewidth - 6\tabcolsep) * \real{0.2500}}
  >{\raggedright\arraybackslash}p{(\linewidth - 6\tabcolsep) * \real{0.2500}}
  >{\raggedright\arraybackslash}p{(\linewidth - 6\tabcolsep) * \real{0.1250}}
  >{\raggedright\arraybackslash}p{(\linewidth - 6\tabcolsep) * \real{0.3750}}@{}}
\caption{Signature Field Keys --- ISO 32000-2:2020
§12.7.5.5}\tabularnewline
\toprule\noalign{}
\begin{minipage}[b]{\linewidth}\raggedright
Key
\end{minipage} & \begin{minipage}[b]{\linewidth}\raggedright
Type
\end{minipage} & \begin{minipage}[b]{\linewidth}\raggedright
Req
\end{minipage} & \begin{minipage}[b]{\linewidth}\raggedright
Spec Ref
\end{minipage} \\
\midrule\noalign{}
\endfirsthead
\toprule\noalign{}
\begin{minipage}[b]{\linewidth}\raggedright
Key
\end{minipage} & \begin{minipage}[b]{\linewidth}\raggedright
Type
\end{minipage} & \begin{minipage}[b]{\linewidth}\raggedright
Req
\end{minipage} & \begin{minipage}[b]{\linewidth}\raggedright
Spec Ref
\end{minipage} \\
\midrule\noalign{}
\endhead
\bottomrule\noalign{}
\endlastfoot
\texttt{/FT} & name & Required & §12.7.4.1 \\
\multicolumn{4}{@{}>{\raggedright\arraybackslash}p{(\linewidth - 6\tabcolsep) * \real{1.0000} + 6\tabcolsep}@{}}{%
Field type --- \texttt{/Sig} marks this as a signature field.
Non-negotiable. If you see \texttt{/FT\ /Tx} on something that's
supposed to be a signature, someone made a mistake.} \\
\texttt{/T} & text string & Required & §12.7.4.1 \\
\multicolumn{4}{@{}>{\raggedright\arraybackslash}p{(\linewidth - 6\tabcolsep) * \real{1.0000} + 6\tabcolsep}@{}}{%
Partial field name --- used to identify the field programmatically. This
is how your code finds "the signature field called Signature1."} \\
\texttt{/V} & reference & Optional & §12.7.5.5 \\
\multicolumn{4}{@{}>{\raggedright\arraybackslash}p{(\linewidth - 6\tabcolsep) * \real{1.0000} + 6\tabcolsep}@{}}{%
Reference to the \textbf{signature dictionary} (\texttt{/Type\ /Sig})
--- the cryptographic payload. If absent, the field is unsigned. When
present, this is the key entry that connects the form field to the
actual signature. The signature dictionary is detailed in
\hyperref[ch12-sig-dict]{The Signature Dictionary ()}.} \\
\end{longtable}

The remaining two keys --- \texttt{/Lock} and \texttt{/SV} --- are
sub-dictionaries that deserve their own tables.

\begin{longtable}[]{@{}
  >{\raggedright\arraybackslash}p{(\linewidth - 6\tabcolsep) * \real{0.2500}}
  >{\raggedright\arraybackslash}p{(\linewidth - 6\tabcolsep) * \real{0.2500}}
  >{\raggedright\arraybackslash}p{(\linewidth - 6\tabcolsep) * \real{0.1250}}
  >{\raggedright\arraybackslash}p{(\linewidth - 6\tabcolsep) * \real{0.3750}}@{}}
\caption{Lock Dictionary (\protect\texttt{/Type\ /SigFieldLock}) --- ISO
32000-2:2020 §12.7.5.5, Table 237}\tabularnewline
\toprule\noalign{}
\begin{minipage}[b]{\linewidth}\raggedright
Key
\end{minipage} & \begin{minipage}[b]{\linewidth}\raggedright
Type
\end{minipage} & \begin{minipage}[b]{\linewidth}\raggedright
Req
\end{minipage} & \begin{minipage}[b]{\linewidth}\raggedright
Spec Ref
\end{minipage} \\
\midrule\noalign{}
\endfirsthead
\toprule\noalign{}
\begin{minipage}[b]{\linewidth}\raggedright
Key
\end{minipage} & \begin{minipage}[b]{\linewidth}\raggedright
Type
\end{minipage} & \begin{minipage}[b]{\linewidth}\raggedright
Req
\end{minipage} & \begin{minipage}[b]{\linewidth}\raggedright
Spec Ref
\end{minipage} \\
\midrule\noalign{}
\endhead
\bottomrule\noalign{}
\endlastfoot
\texttt{/Type} & name & Optional & §12.7.5.5 \\
\multicolumn{4}{@{}>{\raggedright\arraybackslash}p{(\linewidth - 6\tabcolsep) * \real{1.0000} + 6\tabcolsep}@{}}{%
Always \texttt{/SigFieldLock}.} \\
\texttt{/Action} & name & Required & §12.7.5.5, Table 237 \\
\multicolumn{4}{@{}>{\raggedright\arraybackslash}p{(\linewidth - 6\tabcolsep) * \real{1.0000} + 6\tabcolsep}@{}}{%
What to lock after signing: \texttt{/All} (lock everything),
\texttt{/Include} (lock fields listed in \texttt{/Fields}), or
\texttt{/Exclude} (lock everything \emph{except} fields listed in
\texttt{/Fields}). This is how approval signatures can lock specific
fields --- see \hyperref[ch14]{???}.} \\
\texttt{/Fields} & array & Conditional & §12.7.5.5, Table 237 \\
\multicolumn{4}{@{}>{\raggedright\arraybackslash}p{(\linewidth - 6\tabcolsep) * \real{1.0000} + 6\tabcolsep}@{}}{%
Array of field names to include or exclude from locking. Required when
\texttt{/Action} is \texttt{/Include} or \texttt{/Exclude}.} \\
\end{longtable}

\subsubsection{\texorpdfstring{Seed Value Dictionary (\texttt{/SV}): The
Unsung
Hero}{Seed Value Dictionary (/SV): The Unsung Hero}}\label{_seed_value_dictionary_sv_the_unsung_hero}

\begin{quote}
The seed value dictionary provides constraint information that shall be
used at the time the signature is applied.

---  ISO 32000-2:2020 §12.7.5.5
\end{quote}

The seed value dictionary is one of the most powerful and underused
features in the entire PDF signature specification. Most people don't
even know it exists, which is a shame, because it solves a problem that
plagues every organization deploying PDF signatures at scale:
\textbf{how do you ensure signers use the right algorithm, the right
format, and the right handler?}

\begin{longtable}[]{@{}
  >{\raggedright\arraybackslash}p{(\linewidth - 6\tabcolsep) * \real{0.2500}}
  >{\raggedright\arraybackslash}p{(\linewidth - 6\tabcolsep) * \real{0.2500}}
  >{\raggedright\arraybackslash}p{(\linewidth - 6\tabcolsep) * \real{0.1250}}
  >{\raggedright\arraybackslash}p{(\linewidth - 6\tabcolsep) * \real{0.3750}}@{}}
\caption{Seed Value Dictionary (\protect\texttt{/Type\ /SV}) --- ISO
32000-2:2020 §12.7.5.5, Table 238}\tabularnewline
\toprule\noalign{}
\begin{minipage}[b]{\linewidth}\raggedright
Key
\end{minipage} & \begin{minipage}[b]{\linewidth}\raggedright
Type
\end{minipage} & \begin{minipage}[b]{\linewidth}\raggedright
Req
\end{minipage} & \begin{minipage}[b]{\linewidth}\raggedright
Spec Ref
\end{minipage} \\
\midrule\noalign{}
\endfirsthead
\toprule\noalign{}
\begin{minipage}[b]{\linewidth}\raggedright
Key
\end{minipage} & \begin{minipage}[b]{\linewidth}\raggedright
Type
\end{minipage} & \begin{minipage}[b]{\linewidth}\raggedright
Req
\end{minipage} & \begin{minipage}[b]{\linewidth}\raggedright
Spec Ref
\end{minipage} \\
\midrule\noalign{}
\endhead
\bottomrule\noalign{}
\endlastfoot
\texttt{/Filter} & name & Optional & §12.7.5.5, Table 238 \\
\multicolumn{4}{@{}>{\raggedright\arraybackslash}p{(\linewidth - 6\tabcolsep) * \real{1.0000} + 6\tabcolsep}@{}}{%
Required signature handler (e.g., \texttt{/Adobe.PPKLite}). Use this to
ensure your documents work with specific software.} \\
\texttt{/SubFilter} & array of names & Optional & §12.7.5.5, Table
238 \\
\multicolumn{4}{@{}>{\raggedright\arraybackslash}p{(\linewidth - 6\tabcolsep) * \real{1.0000} + 6\tabcolsep}@{}}{%
Acceptable signature formats. Set this to
\texttt{{[}/ETSI.CAdES.detached{]}} if you need PAdES compliance. The
signer's software \textbf{must} use one of these formats.} \\
\texttt{/DigestMethod} & array of names & Optional & §12.7.5.5, Table
238 \\
\multicolumn{4}{@{}>{\raggedright\arraybackslash}p{(\linewidth - 6\tabcolsep) * \real{1.0000} + 6\tabcolsep}@{}}{%
Acceptable hash algorithms. Set this to
\texttt{{[}/SHA256\ /SHA384\ /SHA512{]}} --- and for the love of all
that is holy, \textbf{do not include \texttt{/SHA1}}. It's 2026. SHA-1
has been broken since 2017.} \\
\texttt{/Reasons} & array of strings & Optional & §12.7.5.5, Table
238 \\
\multicolumn{4}{@{}>{\raggedright\arraybackslash}p{(\linewidth - 6\tabcolsep) * \real{1.0000} + 6\tabcolsep}@{}}{%
Pre-defined reason strings the signer must choose from. Great for audit
trails --- instead of "asdf" in the reason field, signers pick from
"Approved", "Reviewed", "Authorized".} \\
\texttt{/LegalAttestation} & array of strings & Optional & §12.7.5.5,
Table 238 \\
\multicolumn{4}{@{}>{\raggedright\arraybackslash}p{(\linewidth - 6\tabcolsep) * \real{1.0000} + 6\tabcolsep}@{}}{%
Legal attestation strings --- rarely used, but available if you need
legal boilerplate attached to the signature.} \\
\texttt{/Ff} & integer & Optional & §12.7.5.5, Table 238 \\
\multicolumn{4}{@{}>{\raggedright\arraybackslash}p{(\linewidth - 6\tabcolsep) * \real{1.0000} + 6\tabcolsep}@{}}{%
Flags indicating which seed values are \textbf{required} (vs.
recommended). If bit 1 is set on a constraint, it's mandatory --- the
signer's software must enforce it, not just suggest it.} \\
\end{longtable}

The \texttt{/MDP} and \texttt{/TimeStamp} entries are sub-dictionaries
with their own structure:

\begin{longtable}[]{@{}
  >{\raggedright\arraybackslash}p{(\linewidth - 8\tabcolsep) * \real{0.1667}}
  >{\raggedright\arraybackslash}p{(\linewidth - 8\tabcolsep) * \real{0.1667}}
  >{\raggedright\arraybackslash}p{(\linewidth - 8\tabcolsep) * \real{0.0833}}
  >{\raggedright\arraybackslash}p{(\linewidth - 8\tabcolsep) * \real{0.1667}}
  >{\raggedright\arraybackslash}p{(\linewidth - 8\tabcolsep) * \real{0.4167}}@{}}
\caption{Seed Value MDP Dictionary (\protect\texttt{/MDP}) --- ISO
32000-2:2020 §12.7.5.5}\tabularnewline
\toprule\noalign{}
\begin{minipage}[b]{\linewidth}\raggedright
Key
\end{minipage} & \begin{minipage}[b]{\linewidth}\raggedright
Type
\end{minipage} & \begin{minipage}[b]{\linewidth}\raggedright
Req
\end{minipage} & \begin{minipage}[b]{\linewidth}\raggedright
Spec Ref
\end{minipage} & \begin{minipage}[b]{\linewidth}\raggedright
Description
\end{minipage} \\
\midrule\noalign{}
\endfirsthead
\toprule\noalign{}
\begin{minipage}[b]{\linewidth}\raggedright
Key
\end{minipage} & \begin{minipage}[b]{\linewidth}\raggedright
Type
\end{minipage} & \begin{minipage}[b]{\linewidth}\raggedright
Req
\end{minipage} & \begin{minipage}[b]{\linewidth}\raggedright
Spec Ref
\end{minipage} & \begin{minipage}[b]{\linewidth}\raggedright
Description
\end{minipage} \\
\midrule\noalign{}
\endhead
\bottomrule\noalign{}
\endlastfoot
\texttt{/P} & integer & Required & §12.7.5.5 & The DocMDP permission
level this signature must enforce: 1 (no changes), 2 (form fills +
signatures), 3 (form fills + signatures + annotations). This is how a
template author says "whoever signs this field must also lock the
document." \\
\end{longtable}

\begin{longtable}[]{@{}
  >{\raggedright\arraybackslash}p{(\linewidth - 8\tabcolsep) * \real{0.1667}}
  >{\raggedright\arraybackslash}p{(\linewidth - 8\tabcolsep) * \real{0.1667}}
  >{\raggedright\arraybackslash}p{(\linewidth - 8\tabcolsep) * \real{0.0833}}
  >{\raggedright\arraybackslash}p{(\linewidth - 8\tabcolsep) * \real{0.1667}}
  >{\raggedright\arraybackslash}p{(\linewidth - 8\tabcolsep) * \real{0.4167}}@{}}
\caption{Seed Value TimeStamp Dictionary (\protect\texttt{/TimeStamp})
--- ISO 32000-2:2020 §12.7.5.5}\tabularnewline
\toprule\noalign{}
\begin{minipage}[b]{\linewidth}\raggedright
Key
\end{minipage} & \begin{minipage}[b]{\linewidth}\raggedright
Type
\end{minipage} & \begin{minipage}[b]{\linewidth}\raggedright
Req
\end{minipage} & \begin{minipage}[b]{\linewidth}\raggedright
Spec Ref
\end{minipage} & \begin{minipage}[b]{\linewidth}\raggedright
Description
\end{minipage} \\
\midrule\noalign{}
\endfirsthead
\toprule\noalign{}
\begin{minipage}[b]{\linewidth}\raggedright
Key
\end{minipage} & \begin{minipage}[b]{\linewidth}\raggedright
Type
\end{minipage} & \begin{minipage}[b]{\linewidth}\raggedright
Req
\end{minipage} & \begin{minipage}[b]{\linewidth}\raggedright
Spec Ref
\end{minipage} & \begin{minipage}[b]{\linewidth}\raggedright
Description
\end{minipage} \\
\midrule\noalign{}
\endhead
\bottomrule\noalign{}
\endlastfoot
\texttt{/URL} & ASCII string & Required & §12.7.5.5 & URL of the RFC
3161 timestamp server. \\
\texttt{/Ff} & integer & Optional & §12.7.5.5 & If bit 1 is set,
timestamping is required (not just recommended). If you need PAdES-B-T,
set this. The signer's software will contact the TSA automatically. \\
\end{longtable}

\textbf{Practical advice for template authors:}

If you're creating PDF templates that others will sign, use seed values
aggressively. Don't trust signers to pick the right options --- they
won't. A financial services company once sent us PDFs where half the
signatures used SHA-1 because some signers had old software and nobody
constrained the algorithm. Those signatures were technically valid but
would fail compliance audits.

Set \texttt{/DigestMethod} to \texttt{{[}/SHA256{]}}, set
\texttt{/SubFilter} to \texttt{{[}/ETSI.CAdES.detached{]}}, and set the
\texttt{/Ff} flags to make them mandatory. Your future self will thank
you.

\subsection{\texorpdfstring{The Signature Dictionary
(\texttt{/Type\ /Sig})}{The Signature Dictionary (/Type /Sig)}}\label{ch12-sig-dict}

The signature dictionary is the heart of a PDF digital signature. It's
referenced from the field's \texttt{/V} entry, and it's where the real
action is. Everything else --- the field, the widget, the appearance ---
is just scaffolding. This dictionary is where the cryptographic proof
lives.

Here is the complete structure of a typical signature dictionary ---
every key you'll encounter in production:

\begin{verbatim}
<<
  /Type /Sig
  /Filter /Adobe.PPKLite
  /SubFilter /ETSI.CAdES.detached
  /ByteRange [0 15000 15500 34500]
  /Contents <308206A7...00000000>
  /M (D:20260306120000+00'00')
  /Name (John Doe)
  /Location (Vienna, Austria)
  /Reason (Approved)
  /ContactInfo (john@example.com)
  /Reference [
    <<
      /Type /SigRef
      /TransformMethod /DocMDP
      /TransformParams <<
        /Type /TransformParams
        /P 2
        /V /1.2
      >>
      /DigestMethod /SHA256
    >>
  ]
  /Prop_Build <<
    /App << /Name /PDFalyzer >>
    /Filter << /Name /Adobe.PPKLite >>
  >>
>>
\end{verbatim}

The table below dissects each key. The core cryptographic keys come
first, then metadata, then sub-dictionaries.

\subsubsection{Core Signature Keys}\label{_core_signature_keys}

These five keys define the cryptographic signature itself. Get any of
them wrong and the signature is broken.

\begin{longtable}[]{@{}
  >{\raggedright\arraybackslash}p{(\linewidth - 6\tabcolsep) * \real{0.2857}}
  >{\raggedright\arraybackslash}p{(\linewidth - 6\tabcolsep) * \real{0.1429}}
  >{\raggedright\arraybackslash}p{(\linewidth - 6\tabcolsep) * \real{0.1429}}
  >{\raggedright\arraybackslash}p{(\linewidth - 6\tabcolsep) * \real{0.4286}}@{}}
\caption{Core keys of the Signature Dictionary
(\protect\texttt{/Type\ /Sig}) --- ISO 32000-2:2020 §12.8.1, Table
255}\tabularnewline
\toprule\noalign{}
\begin{minipage}[b]{\linewidth}\raggedright
Key
\end{minipage} & \begin{minipage}[b]{\linewidth}\raggedright
Type
\end{minipage} & \begin{minipage}[b]{\linewidth}\raggedright
Req
\end{minipage} & \begin{minipage}[b]{\linewidth}\raggedright
Spec Ref
\end{minipage} \\
\midrule\noalign{}
\endfirsthead
\toprule\noalign{}
\begin{minipage}[b]{\linewidth}\raggedright
Key
\end{minipage} & \begin{minipage}[b]{\linewidth}\raggedright
Type
\end{minipage} & \begin{minipage}[b]{\linewidth}\raggedright
Req
\end{minipage} & \begin{minipage}[b]{\linewidth}\raggedright
Spec Ref
\end{minipage} \\
\midrule\noalign{}
\endhead
\bottomrule\noalign{}
\endlastfoot
\texttt{/Type} & name & Optional & §12.8.1 \\
\multicolumn{4}{@{}>{\raggedright\arraybackslash}p{(\linewidth - 6\tabcolsep) * \real{1.0000} + 6\tabcolsep}@{}}{%
Always \texttt{/Sig} when present. Identifies this dictionary as a
signature value. Optional per the spec, but always include it ---
defensive parsers rely on it.} \\
\texttt{/Filter} & name & Required & §12.8.1, Table 255 \\
\multicolumn{4}{@{}>{\raggedright\arraybackslash}p{(\linewidth - 6\tabcolsep) * \real{1.0000} + 6\tabcolsep}@{}}{%
The signature handler name. In theory, this identifies which software
can process the signature. In practice, it's always
\texttt{/Adobe.PPKLite}. Always. Even when the signing software has
nothing to do with Adobe. The name stuck because Adobe defined the
original signature spec in 1999 and everyone just\ldots\hspace{0pt} went
with it.} \\
\texttt{/SubFilter} & name & Required & §12.8.1, Table 255 \\
\multicolumn{4}{@{}>{\raggedright\arraybackslash}p{(\linewidth - 6\tabcolsep) * \real{1.0000} + 6\tabcolsep}@{}}{%
The actual encoding format of the signature value. \emph{This} is the
key that determines how \texttt{/Contents} is structured --- see
\hyperref[ch12-subfilter]{SubFilter Values: A Brief History of Getting
It Right} for the full history and current recommendation.} \\
\texttt{/ByteRange} & array & Required & §12.8.1, Table 255 \\
\multicolumn{4}{@{}>{\raggedright\arraybackslash}p{(\linewidth - 6\tabcolsep) * \real{1.0000} + 6\tabcolsep}@{}}{%
Four integers \texttt{{[}offset1\ length1\ offset2\ length2{]}} defining
which bytes of the file are signed. The gap between the two ranges is
where \texttt{/Contents} lives. The most important array in the entire
dictionary --- a deep dive is in \hyperref[ch13]{???}.} \\
\texttt{/Contents} & string & Required & §12.8.1, Table 255 \\
\multicolumn{4}{@{}>{\raggedright\arraybackslash}p{(\linewidth - 6\tabcolsep) * \real{1.0000} + 6\tabcolsep}@{}}{%
Hex-encoded CMS \texttt{SignedData} blob (RFC 5652 --- Cryptographic
Message Syntax (CMS) §5), padded with zeros. This is the actual
cryptographic signature. Everything else in this dictionary is metadata.
This is the payload. See \hyperref[ch12-contents]{The Value: CMS
SignedData}.} \\
\end{longtable}

\subsubsection{Metadata Keys}\label{_metadata_keys}

These keys carry informational data about the signing event. They are
part of the signed bytes and cannot be changed after signing --- but
they are set by whoever controls the signing process, \textbf{not
verified against any external authority}.

\begin{longtable}[]{@{}
  >{\raggedright\arraybackslash}p{(\linewidth - 8\tabcolsep) * \real{0.1667}}
  >{\raggedright\arraybackslash}p{(\linewidth - 8\tabcolsep) * \real{0.1667}}
  >{\raggedright\arraybackslash}p{(\linewidth - 8\tabcolsep) * \real{0.0833}}
  >{\raggedright\arraybackslash}p{(\linewidth - 8\tabcolsep) * \real{0.1667}}
  >{\raggedright\arraybackslash}p{(\linewidth - 8\tabcolsep) * \real{0.4167}}@{}}
\caption{Metadata keys of the Signature Dictionary --- ISO 32000-2:2020
§12.8.1, Table 255}\tabularnewline
\toprule\noalign{}
\begin{minipage}[b]{\linewidth}\raggedright
Key
\end{minipage} & \begin{minipage}[b]{\linewidth}\raggedright
Type
\end{minipage} & \begin{minipage}[b]{\linewidth}\raggedright
Req
\end{minipage} & \begin{minipage}[b]{\linewidth}\raggedright
Spec Ref
\end{minipage} & \begin{minipage}[b]{\linewidth}\raggedright
Description
\end{minipage} \\
\midrule\noalign{}
\endfirsthead
\toprule\noalign{}
\begin{minipage}[b]{\linewidth}\raggedright
Key
\end{minipage} & \begin{minipage}[b]{\linewidth}\raggedright
Type
\end{minipage} & \begin{minipage}[b]{\linewidth}\raggedright
Req
\end{minipage} & \begin{minipage}[b]{\linewidth}\raggedright
Spec Ref
\end{minipage} & \begin{minipage}[b]{\linewidth}\raggedright
Description
\end{minipage} \\
\midrule\noalign{}
\endhead
\bottomrule\noalign{}
\endlastfoot
\texttt{/M} & date string & Optional & §12.8.1 & Signing time in PDF
date format. \textbf{Do not trust this.} It's the signer's local clock.
Use the CMS timestamp attribute for authoritative time. \\
\texttt{/Name} & text string & Optional & §12.8.1 & Signer's name.
\textbf{Do not trust this either.} It's free text. The authoritative
name is in the X.509 certificate inside \texttt{/Contents}. \\
\texttt{/Location} & text string & Optional & §12.8.1 & Claimed signing
location. Informational, unverified. \\
\texttt{/Reason} & text string & Optional & §12.8.1 & Claimed signing
reason (e.g., "Approved", "Reviewed"). Informational, unverified. \\
\texttt{/ContactInfo} & text string & Optional & §12.8.1 & Contact info.
Informational, optional, you see the pattern. \\
\end{longtable}

The \texttt{/M}, \texttt{/Name}, \texttt{/Location}, \texttt{/Reason},
and \texttt{/ContactInfo} fields are \textbf{unauthenticated metadata}.
Yes, they're part of the signed bytes. Yes, they can't be changed after
signing. But they're set \emph{before} signing, by whoever controls the
signing process. If a malicious signer puts "Barack Obama" in the
\texttt{/Name} field and signs with a self-signed certificate, the
signature will be structurally valid. \textbf{The only trustworthy
identity information is the X.509 certificate in \texttt{/Contents}.}
Never display \texttt{/Name} to users as if it were verified identity.
This mistake is shockingly common in PDF viewer implementations.

\subsubsection{\texorpdfstring{Signature Reference Dictionary
(\texttt{/Reference})}{Signature Reference Dictionary (/Reference)}}\label{_signature_reference_dictionary_reference}

The \texttt{/Reference} entry is an array of \textbf{signature reference
dictionaries} (\texttt{/Type\ /SigRef}) that define transform methods
--- rules about what can and cannot change after signing. This is the
mechanism behind certification signatures and field-level locking
(covered in depth in \hyperref[ch14]{???}).

\begin{longtable}[]{@{}
  >{\raggedright\arraybackslash}p{(\linewidth - 6\tabcolsep) * \real{0.2500}}
  >{\raggedright\arraybackslash}p{(\linewidth - 6\tabcolsep) * \real{0.2500}}
  >{\raggedright\arraybackslash}p{(\linewidth - 6\tabcolsep) * \real{0.1250}}
  >{\raggedright\arraybackslash}p{(\linewidth - 6\tabcolsep) * \real{0.3750}}@{}}
\caption{Signature Reference Dictionary
(\protect\texttt{/Type\ /SigRef}) --- ISO 32000-2:2020 §12.8.1, Table
256}\tabularnewline
\toprule\noalign{}
\begin{minipage}[b]{\linewidth}\raggedright
Key
\end{minipage} & \begin{minipage}[b]{\linewidth}\raggedright
Type
\end{minipage} & \begin{minipage}[b]{\linewidth}\raggedright
Req
\end{minipage} & \begin{minipage}[b]{\linewidth}\raggedright
Spec Ref
\end{minipage} \\
\midrule\noalign{}
\endfirsthead
\toprule\noalign{}
\begin{minipage}[b]{\linewidth}\raggedright
Key
\end{minipage} & \begin{minipage}[b]{\linewidth}\raggedright
Type
\end{minipage} & \begin{minipage}[b]{\linewidth}\raggedright
Req
\end{minipage} & \begin{minipage}[b]{\linewidth}\raggedright
Spec Ref
\end{minipage} \\
\midrule\noalign{}
\endhead
\bottomrule\noalign{}
\endlastfoot
\texttt{/Type} & name & Optional & §12.8.1, Table 256 \\
\multicolumn{4}{@{}>{\raggedright\arraybackslash}p{(\linewidth - 6\tabcolsep) * \real{1.0000} + 6\tabcolsep}@{}}{%
Always \texttt{/SigRef}.} \\
\texttt{/TransformMethod} & name & Required & §12.8.1, Table 256 \\
\multicolumn{4}{@{}>{\raggedright\arraybackslash}p{(\linewidth - 6\tabcolsep) * \real{1.0000} + 6\tabcolsep}@{}}{%
The transform method: \texttt{/DocMDP} (document modification detection
and prevention), \texttt{/FieldMDP} (field-level locking), or
\texttt{/UR} (usage rights --- deprecated in PDF 2.0). DocMDP is by far
the most common. See \hyperref[ch14]{???}.} \\
\texttt{/TransformParams} & dictionary & Optional & §12.8.1, Table
256 \\
\multicolumn{4}{@{}>{\raggedright\arraybackslash}p{(\linewidth - 6\tabcolsep) * \real{1.0000} + 6\tabcolsep}@{}}{%
Parameters for the transform method. For DocMDP, this contains
\texttt{/P} (permission level: 1 = no changes, 2 = form fills +
signatures, 3 = form fills + signatures + annotations) and \texttt{/V}
(version).} \\
\texttt{/DigestMethod} & name & Required & §12.8.1, Table 256 \\
\multicolumn{4}{@{}>{\raggedright\arraybackslash}p{(\linewidth - 6\tabcolsep) * \real{1.0000} + 6\tabcolsep}@{}}{%
Hash algorithm for this reference (e.g., \texttt{/SHA256}). Must match
the algorithm used in the CMS structure.} \\
\end{longtable}

\subsubsection{\texorpdfstring{Build Properties Dictionary
(\texttt{/Prop\_Build})}{Build Properties Dictionary (/Prop\_Build)}}\label{_build_properties_dictionary_prop_build}

The \texttt{/Prop\_Build} dictionary identifies the software that
created the signature. It's optional and purely informational --- but
surprisingly useful for debugging.

\begin{longtable}[]{@{}
  >{\raggedright\arraybackslash}p{(\linewidth - 8\tabcolsep) * \real{0.1667}}
  >{\raggedright\arraybackslash}p{(\linewidth - 8\tabcolsep) * \real{0.1667}}
  >{\raggedright\arraybackslash}p{(\linewidth - 8\tabcolsep) * \real{0.0833}}
  >{\raggedright\arraybackslash}p{(\linewidth - 8\tabcolsep) * \real{0.1667}}
  >{\raggedright\arraybackslash}p{(\linewidth - 8\tabcolsep) * \real{0.4167}}@{}}
\caption{Build Properties Dictionary --- ISO 32000-2:2020 §12.8.1, Table
258}\tabularnewline
\toprule\noalign{}
\begin{minipage}[b]{\linewidth}\raggedright
Key
\end{minipage} & \begin{minipage}[b]{\linewidth}\raggedright
Type
\end{minipage} & \begin{minipage}[b]{\linewidth}\raggedright
Req
\end{minipage} & \begin{minipage}[b]{\linewidth}\raggedright
Spec Ref
\end{minipage} & \begin{minipage}[b]{\linewidth}\raggedright
Description
\end{minipage} \\
\midrule\noalign{}
\endfirsthead
\toprule\noalign{}
\begin{minipage}[b]{\linewidth}\raggedright
Key
\end{minipage} & \begin{minipage}[b]{\linewidth}\raggedright
Type
\end{minipage} & \begin{minipage}[b]{\linewidth}\raggedright
Req
\end{minipage} & \begin{minipage}[b]{\linewidth}\raggedright
Spec Ref
\end{minipage} & \begin{minipage}[b]{\linewidth}\raggedright
Description
\end{minipage} \\
\midrule\noalign{}
\endhead
\bottomrule\noalign{}
\endlastfoot
\texttt{/App} & dictionary & Optional & §12.8.1 & Application that
created the signature. Contains \texttt{/Name} (e.g.,
\texttt{/PDFalyzer}, \texttt{/AdobeAcrobatDC}). \\
\texttt{/Filter} & dictionary & Optional & §12.8.1 & Signature handler
info. Contains \texttt{/Name} (e.g., \texttt{/Adobe.PPKLite}). \\
\texttt{/PubSec} & dictionary & Optional & §12.8.1 & Public security
handler info. Rarely populated. \\
\end{longtable}

\textbf{The \texttt{/Prop\_Build} dictionary is your friend when
debugging.}

When a signature doesn't validate, the first thing I check is
\texttt{/Prop\_Build}. It tells you which software created the
signature. If it says \texttt{/iText} with a version from 2015, you know
you might be dealing with an old CMS structure. If it says
\texttt{/Adobe\ Acrobat} and the signature is invalid, something is
genuinely wrong because Adobe's own tool should produce signatures that
Adobe validates. If \texttt{/Prop\_Build} is missing entirely, the
signing software was lazy --- which tells you something about the
quality of the rest of the signature, too.

\subsubsection{SubFilter Values: A Brief History of Getting It
Right}\label{ch12-subfilter}

The \texttt{/SubFilter} value determines how the \texttt{/Contents}
bytes are structured. Getting this right is the difference between
"signature works everywhere" and "signature works in Adobe but nowhere
else" (or vice versa).

The history of SubFilter values is essentially the history of PDF
signatures trying --- and sometimes failing --- to keep up with
cryptographic best practices:

The first two SubFilter values are deprecated --- they use SHA-1, which
has been broken for collision resistance since the SHAttered attack of
February 2017. You will encounter them in legacy documents, but you
should never use them for new signatures.

\begin{longtable}[]{@{}
  >{\raggedright\arraybackslash}p{(\linewidth - 6\tabcolsep) * \real{0.2500}}
  >{\raggedright\arraybackslash}p{(\linewidth - 6\tabcolsep) * \real{0.2500}}
  >{\raggedright\arraybackslash}p{(\linewidth - 6\tabcolsep) * \real{0.1250}}
  >{\raggedright\arraybackslash}p{(\linewidth - 6\tabcolsep) * \real{0.3750}}@{}}
\caption{Deprecated SubFilter Values (do not use for new
signatures)}\tabularnewline
\toprule\noalign{}
\begin{minipage}[b]{\linewidth}\raggedright
SubFilter
\end{minipage} & \begin{minipage}[b]{\linewidth}\raggedright
Standard
\end{minipage} & \begin{minipage}[b]{\linewidth}\raggedright
Req
\end{minipage} & \begin{minipage}[b]{\linewidth}\raggedright
Spec Ref
\end{minipage} \\
\midrule\noalign{}
\endfirsthead
\toprule\noalign{}
\begin{minipage}[b]{\linewidth}\raggedright
SubFilter
\end{minipage} & \begin{minipage}[b]{\linewidth}\raggedright
Standard
\end{minipage} & \begin{minipage}[b]{\linewidth}\raggedright
Req
\end{minipage} & \begin{minipage}[b]{\linewidth}\raggedright
Spec Ref
\end{minipage} \\
\midrule\noalign{}
\endhead
\bottomrule\noalign{}
\endlastfoot
\texttt{/adbe.x509.rsa\_sha1} & PDF 1.3+ & Deprecated & ISO 32000-2:2020
§12.8.3.2 \\
\multicolumn{4}{@{}>{\raggedright\arraybackslash}p{(\linewidth - 6\tabcolsep) * \real{1.0000} + 6\tabcolsep}@{}}{%
\textbf{Do not use.} Raw RSA signature with the certificate in a
separate \texttt{/Cert} key instead of inside \texttt{/Contents}. No CMS
wrapper. SHA-1 only. This was PDF's first attempt at digital signatures
in 1999. It made sense at the time. It doesn't anymore.} \\
\texttt{/adbe.pkcs7.sha1} & PDF 1.3+ & Deprecated & ISO 32000-2:2020
§12.8.3.4 \\
\multicolumn{4}{@{}>{\raggedright\arraybackslash}p{(\linewidth - 6\tabcolsep) * \real{1.0000} + 6\tabcolsep}@{}}{%
\textbf{Do not use.} CMS \texttt{SignedData} where the SHA-1 digest of
the byte range is \emph{encapsulated} in the CMS. This is the worst of
both worlds: SHA-1 (broken) and encapsulated content (unnecessary
complexity). Nobody misses this format.} \\
\end{longtable}

The following three values are the current standard. For new signatures,
use \texttt{/ETSI.CAdES.detached} unless you have a specific reason not
to.

\begin{longtable}[]{@{}
  >{\raggedright\arraybackslash}p{(\linewidth - 6\tabcolsep) * \real{0.2500}}
  >{\raggedright\arraybackslash}p{(\linewidth - 6\tabcolsep) * \real{0.2500}}
  >{\raggedright\arraybackslash}p{(\linewidth - 6\tabcolsep) * \real{0.1250}}
  >{\raggedright\arraybackslash}p{(\linewidth - 6\tabcolsep) * \real{0.3750}}@{}}
\caption{Current SubFilter Values}\tabularnewline
\toprule\noalign{}
\begin{minipage}[b]{\linewidth}\raggedright
SubFilter
\end{minipage} & \begin{minipage}[b]{\linewidth}\raggedright
Standard
\end{minipage} & \begin{minipage}[b]{\linewidth}\raggedright
Req
\end{minipage} & \begin{minipage}[b]{\linewidth}\raggedright
Spec Ref
\end{minipage} \\
\midrule\noalign{}
\endfirsthead
\toprule\noalign{}
\begin{minipage}[b]{\linewidth}\raggedright
SubFilter
\end{minipage} & \begin{minipage}[b]{\linewidth}\raggedright
Standard
\end{minipage} & \begin{minipage}[b]{\linewidth}\raggedright
Req
\end{minipage} & \begin{minipage}[b]{\linewidth}\raggedright
Spec Ref
\end{minipage} \\
\midrule\noalign{}
\endhead
\bottomrule\noalign{}
\endlastfoot
\texttt{/adbe.pkcs7.detached} & PDF 1.3+ & Current & ISO 32000-2:2020
§12.8.3.3 \\
\multicolumn{4}{@{}>{\raggedright\arraybackslash}p{(\linewidth - 6\tabcolsep) * \real{1.0000} + 6\tabcolsep}@{}}{%
CMS \texttt{SignedData} with detached content. The workhorse format ---
widely supported since PDF 1.3. The signed bytes are the PDF byte range,
not embedded in CMS. If you need maximum compatibility with older
software, this is your format.} \\
\texttt{/ETSI.CAdES.detached} & PDF 2.0 / PAdES & Modern & ETSI EN 319
142-1 V1.1.1 (2016-04) §5.1 \\
\multicolumn{4}{@{}>{\raggedright\arraybackslash}p{(\linewidth - 6\tabcolsep) * \real{1.0000} + 6\tabcolsep}@{}}{%
CMS \texttt{SignedData} with mandatory CAdES signed attributes,
specifically the \texttt{signing-certificate-v2} attribute. Required for
PAdES-B-B and higher. \textbf{This is the modern standard. Use this.}
The CMS structure is identical to \texttt{/adbe.pkcs7.detached} at the
container level --- the difference is one signed attribute
(\textasciitilde50 bytes).} \\
\texttt{/ETSI.RFC3161} & PDF 2.0 / PAdES & Modern & ETSI EN 319 142-1
V1.1.1 (2016-04) §5.4 \\
\multicolumn{4}{@{}>{\raggedright\arraybackslash}p{(\linewidth - 6\tabcolsep) * \real{1.0000} + 6\tabcolsep}@{}}{%
RFC 3161 timestamp token --- not a signature, but a document-level
timestamp. Used for PAdES-B-T and long-term archival. You'll encounter
these when dealing with LTV (Long-Term Validation).} \\
\end{longtable}

Here's a trap that catches people all the time: Adobe Acrobat and most
viewers accept \emph{all} SubFilter values for backward compatibility.
Your signature will "work" in Acrobat regardless of which SubFilter you
use. But \textbf{only} \texttt{/ETSI.CAdES.detached} is valid for
PAdES-compliant signatures under eIDAS. If you need legal standing in
the EU --- and increasingly, that's everyone --- this is the only
acceptable value. Don't let "it works in Acrobat" lull you into using
\texttt{/adbe.pkcs7.detached} for new signatures. Acrobat is not the
arbiter of legal validity; the eIDAS regulation is.

\textbf{The compatibility question everyone asks:}

"Can I use \texttt{/ETSI.CAdES.detached} if my users have old PDF
readers?"

Yes. The CMS structure is identical to \texttt{/adbe.pkcs7.detached} at
the container level. Old readers that don't understand the SubFilter
name will typically fall back to parsing the \texttt{/Contents} as
generic CMS/PKCS\#7 --- and it will work. The difference is one signed
attribute (\texttt{signing-certificate-v2}), which old readers simply
ignore.

The only scenario where it breaks: truly ancient software (pre-2010)
that strictly checks the SubFilter name against a hardcoded whitelist
and rejects unknown values. If your users are running PDF readers from
2008, they have bigger problems than SubFilter compatibility.

\subsection{\texorpdfstring{The \texttt{/Contents} Value: CMS
SignedData}{The /Contents Value: CMS SignedData}}\label{ch12-contents}

The \texttt{/Contents} entry is a \textbf{hex-encoded byte string}
containing a DER-encoded CMS \texttt{SignedData} structure (RFC 5652 ---
Cryptographic Message Syntax (CMS) §5).

This is where the actual cryptographic signature lives. Everything else
in the signature dictionary is metadata. This is the payload.

\subsubsection{Why Hex-Encoded?}\label{_why_hex_encoded}

PDF strings can be literal \texttt{(…\hspace{0pt})} or hex
\texttt{\textless{}…\hspace{0pt}\textgreater{}}. The hex form is used
because CMS blobs contain arbitrary binary data (DER-encoded ASN.1), and
literal strings would need extensive escaping for bytes like
\texttt{\textbackslash{}(}, \texttt{\textbackslash{})},
\texttt{\textbackslash{}\textbackslash{}}. Hex encoding doubles the size
but is safe and unambiguous.

A typical \texttt{/Contents} value looks like
\texttt{\textless{}308206A7…\hspace{0pt}several\ thousand\ hex\ characters…\hspace{0pt}0000000000\textgreater{}}.
The leading bytes (\texttt{3082…\hspace{0pt}}) are the DER-encoded ASN.1
SEQUENCE tag, and the trailing zeros are padding. Which brings us to one
of the more elegant hacks in the PDF specification:

\subsubsection{Pre-allocation: The Chicken-and-Egg
Problem}\label{_pre_allocation_the_chicken_and_egg_problem}

Here's the fundamental problem with PDF signatures:

\begin{enumerate}
\def\labelenumi{\arabic{enumi}.}
\item
  The signature covers the file's bytes (including the signature
  dictionary itself)
\item
  But the signature dictionary must contain the CMS blob
\item
  But the CMS blob can only be computed \emph{after} you know which
  bytes to sign
\item
  But you can't know which bytes to sign until the dictionary is written
\end{enumerate}

The solution is a trick that every PDF signing implementation uses:
\textbf{pre-allocate} a large enough \texttt{/Contents} value filled
with zeros, write the file, \emph{then} overwrite the zeros with the
actual CMS bytes:

The pre-allocated size must be estimated accurately, and this is where
things go wrong in the real world:

\begin{itemize}
\item
  \textbf{Too small} --- the CMS blob doesn't fit --- signing fails.
  You'll get an exception from PDFBox and no output file. This is the
  most common "my signing suddenly broke" bug, usually triggered by
  adding a certificate chain or a timestamp that makes the CMS blob
  bigger than expected.
\item
  \textbf{Too large} --- wasted bytes, but the signature is valid. A few
  KB of zeros at the end of the file never hurt anyone.
\end{itemize}

\textbf{Always err on the side of too large.} Here are sensible
defaults:

{\def\LTcaptype{none} % do not increment counter
\begin{longtable}[]{@{}
  >{\raggedright\arraybackslash}p{(\linewidth - 4\tabcolsep) * \real{0.3750}}
  >{\raggedright\arraybackslash}p{(\linewidth - 4\tabcolsep) * \real{0.2500}}
  >{\raggedright\arraybackslash}p{(\linewidth - 4\tabcolsep) * \real{0.3750}}@{}}
\toprule\noalign{}
\begin{minipage}[b]{\linewidth}\raggedright
Scenario
\end{minipage} & \begin{minipage}[b]{\linewidth}\raggedright
Hex Size
\end{minipage} & \begin{minipage}[b]{\linewidth}\raggedright
Bytes
\end{minipage} \\
\midrule\noalign{}
\endhead
\bottomrule\noalign{}
\endlastfoot
Basic signature (no chain, no timestamp) & \texttt{0x2500} (PDFBox
default) & \textasciitilde9 KB \\
Signature with certificate chain & \texttt{0x4000} & \textasciitilde16
KB \\
PAdES-B-T (with RFC 3161 timestamp) & \texttt{0x6000} &
\textasciitilde24 KB \\
PAdES-B-LT (embedded OCSP + CRLs) & \texttt{0x18000} & \textasciitilde96
KB \\
PAdES-B-LTA (archival, large CRL) & \texttt{0x28000} &
\textasciitilde160 KB \\
\end{longtable}
}

I've seen production systems fail because someone used the PDFBox
default (\texttt{0x2500}) for PAdES-B-T signatures. The first 500
signatures worked fine because the TSA responses were small. Then the
TSA renewed its certificate with a longer chain, the responses grew by 2
KB, and every signing operation started failing on a Friday afternoon.
Set it to \texttt{0x6000} minimum and forget about it.

\subsubsection{CMS SignedData: The Cryptographic
Core}\label{_cms_signeddata_the_cryptographic_core}

The CMS blob inside \texttt{/Contents} follows RFC 5652 ---
Cryptographic Message Syntax (CMS) §5. If you've worked with S/MIME
email signatures or code signing, this structure will look familiar ---
it's the same \texttt{SignedData} format used across the entire PKI
ecosystem:

\begin{Shaded}
\begin{Highlighting}[]
\NormalTok{SignedData ::= }\DataTypeTok{SEQUENCE}\NormalTok{ \{}
\NormalTok{    version            CMSVersion,              }
\NormalTok{    digestAlgorithms   DigestAlgorithmIdentifiers,  }
\NormalTok{    encapContentInfo   EncapsulatedContentInfo,     }
\NormalTok{    certificates   [}\DecValTok{0}\NormalTok{] IMPLICIT CertificateSet }\KeywordTok{OPTIONAL}\NormalTok{, }
\NormalTok{    crls           [}\DecValTok{1}\NormalTok{] IMPLICIT RevocationInfoChoices }\KeywordTok{OPTIONAL}\NormalTok{, }
\NormalTok{    signerInfos        SignerInfos              }
\NormalTok{\}}
\end{Highlighting}
\end{Shaded}

\begin{itemize}
\item
  Version: 1 for basic CMS, 3 if attribute certificates are present
\item
  Set of digest algorithms used (typically just \texttt{\{SHA-256\}})
\item
  Content type: \texttt{id-data} (1.2.840.113549.1.7.1) with
  \textbf{absent} content --- this is what "detached" means. The signed
  data (the PDF byte range) is not inside the CMS structure. It's
  outside, in the PDF file itself.
\item
  Certificates --- the signer's certificate and optionally the full
  chain (intermediates, possibly root). \textbf{Always include the full
  chain.} I know it makes the blob bigger. I know storage is a concern.
  Do it anyway. A signature without a chain is a signature that can only
  be validated by someone who already has all your certificates. That's
  almost nobody.
\item
  CRLs --- rarely embedded here in the modern PAdES world. In PAdES-B-LT
  and later, revocation data goes in the DSS dictionary instead.
\item
  Signer information --- the actual signature value and signed
  attributes. This is where the cryptography lives.
\end{itemize}

\paragraph{SignerInfo: Where the Cryptography
Lives}\label{_signerinfo_where_the_cryptography_lives}

\begin{Shaded}
\begin{Highlighting}[]
\NormalTok{SignerInfo ::= }\DataTypeTok{SEQUENCE}\NormalTok{ \{}
\NormalTok{    version                  CMSVersion,}
\NormalTok{    sid                      SignerIdentifier,    }
\NormalTok{    digestAlgorithm          DigestAlgorithmIdentifier, }
\NormalTok{    signedAttrs          [}\DecValTok{0}\NormalTok{] IMPLICIT SignedAttributes }\KeywordTok{OPTIONAL}\NormalTok{, }
\NormalTok{    signatureAlgorithm       SignatureAlgorithmIdentifier, }
\NormalTok{    signature                SignatureValue,      }
\NormalTok{    unsignedAttrs        [}\DecValTok{1}\NormalTok{] IMPLICIT UnsignedAttributes }\KeywordTok{OPTIONAL} 
\NormalTok{\}}
\end{Highlighting}
\end{Shaded}

\begin{itemize}
\item
  Identifies the signer's certificate --- either by issuer DN + serial
  number, or by subject key identifier. Both work. Issuer+serial is more
  traditional; subject key identifier is more robust when certificates
  are reissued.
\item
  Hash algorithm used to digest the PDF byte range. SHA-256 is the
  minimum acceptable. SHA-384 or SHA-512 if you're paranoid or work in
  defense.
\item
  \textbf{Signed attributes} --- authenticated metadata. See the
  critical attributes table below. These are \emph{signed}, meaning any
  modification invalidates the signature.
\item
  Signature algorithm (e.g., SHA256withRSA, SHA256withECDSA). Note: the
  \emph{signature} is computed over the DER-encoded signed attributes,
  not directly over the PDF bytes. This is a CMS subtlety that trips up
  many implementers.
\item
  The actual cryptographic signature bytes.
\item
  \textbf{Unsigned attributes} --- unauthenticated extensions.
  Timestamps go here.
\end{itemize}

\paragraph{The Signed Attributes That
Matter}\label{_the_signed_attributes_that_matter}

The first two attributes are required by CMS whenever signed attributes
are present. The remaining attributes add PAdES compliance and
defense-in-depth.

\begin{longtable}[]{@{}
  >{\raggedright\arraybackslash}p{(\linewidth - 6\tabcolsep) * \real{0.3333}}
  >{\raggedright\arraybackslash}p{(\linewidth - 6\tabcolsep) * \real{0.2222}}
  >{\raggedright\arraybackslash}p{(\linewidth - 6\tabcolsep) * \real{0.1111}}
  >{\raggedright\arraybackslash}p{(\linewidth - 6\tabcolsep) * \real{0.3333}}@{}}
\caption{Required CMS Signed Attributes --- RFC 5652 --- Cryptographic
Message Syntax (CMS) §11}\tabularnewline
\toprule\noalign{}
\begin{minipage}[b]{\linewidth}\raggedright
Attribute
\end{minipage} & \begin{minipage}[b]{\linewidth}\raggedright
OID
\end{minipage} & \begin{minipage}[b]{\linewidth}\raggedright
Req
\end{minipage} & \begin{minipage}[b]{\linewidth}\raggedright
Spec Ref
\end{minipage} \\
\midrule\noalign{}
\endfirsthead
\toprule\noalign{}
\begin{minipage}[b]{\linewidth}\raggedright
Attribute
\end{minipage} & \begin{minipage}[b]{\linewidth}\raggedright
OID
\end{minipage} & \begin{minipage}[b]{\linewidth}\raggedright
Req
\end{minipage} & \begin{minipage}[b]{\linewidth}\raggedright
Spec Ref
\end{minipage} \\
\midrule\noalign{}
\endhead
\bottomrule\noalign{}
\endlastfoot
\texttt{content-type} & 1.2.840.113549.1.9.3 & Required & RFC 5652 ---
Cryptographic Message Syntax (CMS) §11.1 \\
\multicolumn{4}{@{}>{\raggedright\arraybackslash}p{(\linewidth - 6\tabcolsep) * \real{1.0000} + 6\tabcolsep}@{}}{%
Must be \texttt{id-data} (1.2.840.113549.1.7.1). Required by CMS when
signed attributes are present. Forgetting this is a classic
implementation bug --- the signature appears valid in some tools but
fails in others.} \\
\texttt{message-digest} & 1.2.840.113549.1.9.4 & Required & RFC 5652 ---
Cryptographic Message Syntax (CMS) §11.2 \\
\multicolumn{4}{@{}>{\raggedright\arraybackslash}p{(\linewidth - 6\tabcolsep) * \real{1.0000} + 6\tabcolsep}@{}}{%
SHA-256 hash of the PDF byte range. This is what's actually verified
during validation. If this hash doesn't match the file bytes, the
signature is invalid. Full stop.} \\
\end{longtable}

\begin{longtable}[]{@{}
  >{\raggedright\arraybackslash}p{(\linewidth - 6\tabcolsep) * \real{0.3333}}
  >{\raggedright\arraybackslash}p{(\linewidth - 6\tabcolsep) * \real{0.2222}}
  >{\raggedright\arraybackslash}p{(\linewidth - 6\tabcolsep) * \real{0.1111}}
  >{\raggedright\arraybackslash}p{(\linewidth - 6\tabcolsep) * \real{0.3333}}@{}}
\caption{PAdES/CAdES Signed Attributes}\tabularnewline
\toprule\noalign{}
\begin{minipage}[b]{\linewidth}\raggedright
Attribute
\end{minipage} & \begin{minipage}[b]{\linewidth}\raggedright
OID
\end{minipage} & \begin{minipage}[b]{\linewidth}\raggedright
Req
\end{minipage} & \begin{minipage}[b]{\linewidth}\raggedright
Spec Ref
\end{minipage} \\
\midrule\noalign{}
\endfirsthead
\toprule\noalign{}
\begin{minipage}[b]{\linewidth}\raggedright
Attribute
\end{minipage} & \begin{minipage}[b]{\linewidth}\raggedright
OID
\end{minipage} & \begin{minipage}[b]{\linewidth}\raggedright
Req
\end{minipage} & \begin{minipage}[b]{\linewidth}\raggedright
Spec Ref
\end{minipage} \\
\midrule\noalign{}
\endhead
\bottomrule\noalign{}
\endlastfoot
\texttt{signing-certificate-v2} & 1.2.840.113549.1.9.16.2.47 & Required
(PAdES) & RFC 5035 §3 \\
\multicolumn{4}{@{}>{\raggedright\arraybackslash}p{(\linewidth - 6\tabcolsep) * \real{1.0000} + 6\tabcolsep}@{}}{%
Hash of the signer's certificate (SHA-256). \textbf{Required for PAdES /
CAdES.} This single attribute prevents certificate substitution attacks
--- an attacker can't swap in a different certificate without breaking
the signature. It's \textasciitilde50 bytes. There is no reason not to
include it.} \\
\texttt{signing-time} & 1.2.840.113549.1.9.5 & Optional & RFC 5652 ---
Cryptographic Message Syntax (CMS) §11.3 \\
\multicolumn{4}{@{}>{\raggedright\arraybackslash}p{(\linewidth - 6\tabcolsep) * \real{1.0000} + 6\tabcolsep}@{}}{%
Claimed signing time. Not authoritative --- it's whatever the signer's
clock said. For authoritative time, you need a timestamp (see unsigned
attributes). Optional, but almost everyone includes it.} \\
\texttt{algorithm-protection} & 1.2.840.113549.1.9.52 & Recommended &
RFC 6211 \\
\multicolumn{4}{@{}>{\raggedright\arraybackslash}p{(\linewidth - 6\tabcolsep) * \real{1.0000} + 6\tabcolsep}@{}}{%
Protects against algorithm substitution attacks by binding the digest
and signature algorithms into the signed attributes. Good hygiene to
include.} \\
\end{longtable}

\textbf{The \texttt{signing-certificate-v2} attribute is what makes the
difference between \texttt{/adbe.pkcs7.detached} and
\texttt{/ETSI.CAdES.detached}.}

Without it: plain CMS (PKCS\#7) --- vulnerable to certificate
substitution attacks. With it: CAdES --- the signer's certificate is
bound to the signature by hash.

This single attribute is the gateway to PAdES compliance. It's about 50
bytes in the CMS structure. There is no reason not to include it, even
if you're using \texttt{/adbe.pkcs7.detached} as your SubFilter.

Certificate substitution attacks are not theoretical. In 2019,
researchers at Ruhr University Bochum demonstrated practical attacks
against multiple PDF viewers that exploited the absence of certificate
binding. The fix was trivially simple: check
\texttt{signing-certificate-v2}. Several major viewers had to be
patched.

\paragraph{Unsigned Attributes: The Afterthought That Became
Essential}\label{_unsigned_attributes_the_afterthought_that_became_essential}

\begin{longtable}[]{@{}
  >{\raggedright\arraybackslash}p{(\linewidth - 6\tabcolsep) * \real{0.3333}}
  >{\raggedright\arraybackslash}p{(\linewidth - 6\tabcolsep) * \real{0.2222}}
  >{\raggedright\arraybackslash}p{(\linewidth - 6\tabcolsep) * \real{0.1111}}
  >{\raggedright\arraybackslash}p{(\linewidth - 6\tabcolsep) * \real{0.3333}}@{}}
\caption{Unsigned Attributes --- RFC 5652 --- Cryptographic Message
Syntax (CMS) §11}\tabularnewline
\toprule\noalign{}
\begin{minipage}[b]{\linewidth}\raggedright
Attribute
\end{minipage} & \begin{minipage}[b]{\linewidth}\raggedright
OID
\end{minipage} & \begin{minipage}[b]{\linewidth}\raggedright
Req
\end{minipage} & \begin{minipage}[b]{\linewidth}\raggedright
Spec Ref
\end{minipage} \\
\midrule\noalign{}
\endfirsthead
\toprule\noalign{}
\begin{minipage}[b]{\linewidth}\raggedright
Attribute
\end{minipage} & \begin{minipage}[b]{\linewidth}\raggedright
OID
\end{minipage} & \begin{minipage}[b]{\linewidth}\raggedright
Req
\end{minipage} & \begin{minipage}[b]{\linewidth}\raggedright
Spec Ref
\end{minipage} \\
\midrule\noalign{}
\endhead
\bottomrule\noalign{}
\endlastfoot
\texttt{signature-time-stamp} & 1.2.840.113549.1.9.16.2.14 & Required
(PAdES-B-T) & RFC 3161 --- Internet X.509 PKI Time-Stamp Protocol (TSP),
ETSI EN 319 142-1 V1.1.1 (2016-04) §5.3 \\
\multicolumn{4}{@{}>{\raggedright\arraybackslash}p{(\linewidth - 6\tabcolsep) * \real{1.0000} + 6\tabcolsep}@{}}{%
RFC 3161 timestamp token over the signature value. Provides
authoritative proof of signing time from a trusted TSA. This is the only
way to prove \emph{when} a signature was made. The timestamp token has
its own signature from the TSA, making it self-authenticating --- which
is why it's safe as an unsigned attribute.} \\
\texttt{certificate-values} & 1.2.840.113549.1.9.16.2.23 & Optional &
RFC 5126 §5.6.1 \\
\multicolumn{4}{@{}>{\raggedright\arraybackslash}p{(\linewidth - 6\tabcolsep) * \real{1.0000} + 6\tabcolsep}@{}}{%
Full certificate chain values. Used in CAdES-X-Long. In the PAdES world,
these typically go in the DSS dictionary instead (see
\hyperref[ch18]{???}).} \\
\texttt{revocation-values} & 1.2.840.113549.1.9.16.2.24 & Optional & RFC
5126 §5.6.1 \\
\multicolumn{4}{@{}>{\raggedright\arraybackslash}p{(\linewidth - 6\tabcolsep) * \real{1.0000} + 6\tabcolsep}@{}}{%
OCSP responses and CRLs. Used in CAdES-X-Long. In PAdES, these go in the
DSS dictionary instead.} \\
\end{longtable}

The timestamp attribute deserves special attention. Unsigned attributes
are, by definition, not protected by the signature --- they can be added
after signing without invalidating anything. This is \emph{by design}. A
timestamp proves that the signature existed at a certain time, and it's
intentionally added after the signature is created (because you need the
signature value to timestamp it).

This is also why you should never put anything security-critical in
unsigned attributes that isn't self-authenticating (like a timestamp
token, which has its own signature from the TSA).

\subsection{The Signing Process: Step by
Step}\label{ch12-signing-process}

The process of creating a PDF digital signature is intricate because the
signature must cover bytes that include the signature dictionary itself
--- a chicken-and-egg problem solved by the \texttt{/Contents}
pre-allocation we discussed above.

Here's the full sequence, and understanding it is essential for
debugging:

\subsubsection{Step-by-Step Detail}\label{_step_by_step_detail}

{\def\LTcaptype{none} % do not increment counter
\begin{longtable}[]{@{}
  >{\raggedright\arraybackslash}p{(\linewidth - 2\tabcolsep) * \real{0.0909}}
  >{\raggedright\arraybackslash}p{(\linewidth - 2\tabcolsep) * \real{0.9091}}@{}}
\toprule\noalign{}
\begin{minipage}[b]{\linewidth}\raggedright
Step
\end{minipage} & \begin{minipage}[b]{\linewidth}\raggedright
Detail
\end{minipage} \\
\midrule\noalign{}
\endhead
\bottomrule\noalign{}
\endlastfoot
\textbf{1} & \textbf{Create or locate the signature field.} If the
document already has an unsigned field (\texttt{/FT\ /Sig} with no
\texttt{/V}), use it. Otherwise, create a new field in the AcroForm
\texttt{/Fields} array with a widget annotation on the desired page.
\emph{Practical note:} if the document was prepared by someone else (a
template), always prefer using existing unsigned fields. Creating new
fields in a document that already has signatures is technically possible
via incremental save, but it makes some validators nervous. \\
\textbf{2} & \textbf{Build the signature dictionary} with placeholder
values. \texttt{/ByteRange} is set to \texttt{{[}0\ 0\ 0\ 0{]}} (will be
overwritten). \texttt{/Contents} is filled with \texttt{00} bytes at the
estimated maximum size. All other keys (\texttt{/Filter},
\texttt{/SubFilter}, \texttt{/M}, \texttt{/Name}, \texttt{/Reason},
\texttt{/Reference}) are set to their final values \emph{now} --- they
become part of the signed bytes and can't change later. \\
\textbf{3} & \textbf{Write the PDF using incremental save.} The new
objects (signature dict, updated field, updated xref) are appended after
the existing \texttt{\%\%EOF}. This is critical --- a full rewrite would
invalidate any existing signatures. \emph{This is non-negotiable.} Even
if the document has no prior signatures, use incremental save. It's the
only safe habit. \\
\textbf{4} & \textbf{Compute actual byte offsets.} After writing, the
exact byte positions of the \texttt{/Contents} hex string are known.
Update \texttt{/ByteRange} to
\texttt{{[}0,\ contentsStart,\ contentsEnd,\ fileLength\ -\ contentsEnd{]}}.
The two segments cover everything \emph{except} the \texttt{/Contents}
value itself. \\
\textbf{5} & \textbf{Read the signed bytes.} Concatenate
\texttt{file{[}0\ ..\ contentsStart-1{]}} and
\texttt{file{[}contentsEnd\ ..\ fileLength-1{]}}. These are the exact
bytes the signature covers --- everything except the \texttt{/Contents}
hex value. \emph{Debug tip:} if you suspect byte range issues, dump
these bytes to a file and hash them manually. Compare with the
\texttt{message-digest} attribute in the CMS. \\
\textbf{6} & \textbf{Hash the signed bytes.} Compute SHA-256 (or
SHA-384/512) over the concatenated byte ranges. This hash becomes the
\texttt{message-digest} signed attribute in the CMS structure. \\
\textbf{7--8} & \textbf{Sign the hash.} The private key (RSA or ECDSA)
signs the DER-encoded signed attributes (which include the hash). This
may happen locally or via a remote HSM/signing service. In the remote
case, only the hash leaves your system --- the private key never does.
More on remote signing in \hyperref[ch22]{???}. \\
\textbf{9} & \textbf{Build the CMS \texttt{SignedData} structure.}
Include: digest algorithm, signer certificate, certificate chain, signed
attributes (content-type, message-digest, signing-certificate-v2),
signature value. DER-encode the result. \\
\textbf{10} & \textbf{Write CMS bytes into \texttt{/Contents}.}
Hex-encode the DER bytes and write them at the exact byte offset of the
\texttt{/Contents} value in the file. Remaining space is left as
\texttt{00} padding. No other bytes in the file are modified. \\
\end{longtable}
}

Step 10 is a \textbf{raw byte write} --- not a PDF object update. The
file is opened as a random-access binary stream, the cursor is
positioned at the \texttt{/Contents} hex string's start offset, and the
hex chars are overwritten in place. If the CMS blob is larger than the
pre-allocated space, signing fails. If even a single byte is written at
the wrong offset, the signature is invalid and the PDF may be corrupted.
This is why you don't implement PDF signing from scratch --- you use
PDFBox or a similar library that handles the byte-level bookkeeping for
you.

\subsection{Incremental Save: The Foundation of
Multi-Signature}\label{ch12-incremental}

Every signature operation uses \textbf{incremental save} (ISO
32000-2:2020 §7.5.6). The original file bytes are never modified --- new
objects are appended after the existing \texttt{\%\%EOF}.

This is the mechanism that makes multiple independent signatures
possible. Without incremental save, adding a second signature would
modify the bytes covered by the first signature, invalidating it. With
incremental save, the first signature's bytes remain untouched --- the
second signature's data is simply appended at the end.

Each signature's \texttt{/ByteRange} covers everything up to and
including its own incremental update --- but \textbf{not} the updates
that come after:

\begin{itemize}
\item
  Signature 1's byte range:
  \texttt{{[}0,\ …\hspace{0pt},\ …\hspace{0pt},\ M{]}} --- covers
  original + update \#1
\item
  Signature 2's byte range:
  \texttt{{[}0,\ …\hspace{0pt},\ …\hspace{0pt},\ P{]}} --- covers
  original + update \#1 + update \#2
\end{itemize}

This is why Signature 1 remains valid after Signature 2 is added:
Signature 2's bytes are \textbf{appended after} the range that Signature
1 covers. From Signature 1's perspective, those bytes are "trailing" ---
present in the file but outside the signed range.

\textbf{"The file has been modified since signing" --- don't panic.}

This is the single most common false alarm in PDF signature validation.
A validator detects trailing bytes after a signature's byte range and
flags it as "modified." But in the vast majority of cases, those
trailing bytes are simply a later incremental update --- another
signature, a form fill, or an annotation.

\textbf{What to check:}

\begin{enumerate}
\def\labelenumi{\arabic{enumi}.}
\item
  Is the trailing data a valid incremental update? (Look for xref table
  + trailer + \%\%EOF.)
\item
  Does it correspond to an expected operation? (Another signature, a
  form field being filled.)
\item
  Did the byte range change? (The original signature's bytes at their
  original offsets should be byte-for-byte identical.)
\end{enumerate}

If the answers are yes, yes, and no --- the signature is fine. The
"modification" is just the normal operation of PDF's append-only
architecture.

\textbf{When to actually worry:}

\begin{itemize}
\item
  Trailing bytes that aren't a valid incremental update --- raw data
  appended to the file
\item
  An incremental update that modifies objects covered by an existing
  signature in unexpected ways (e.g., page content changes)
\item
  Multiple \%\%EOF markers where the trailing section isn't a proper PDF
  update
\end{itemize}

\subsection{Widget Annotations and Page Placement}\label{ch12-widget}

Every signature field has an associated widget annotation (ISO
32000-2:2020 §12.5.6.19) that defines its visual presence on a page.

\subsubsection{\texorpdfstring{The \texttt{/Rect} Coordinate
System}{The /Rect Coordinate System}}\label{_the_rect_coordinate_system}

The widget's \texttt{/Rect} array specifies
\texttt{{[}llx,\ lly,\ urx,\ ury{]}} in \textbf{default user space}
coordinates:

PDF coordinates have \textbf{origin at lower-left} with Y increasing
upward. This is the opposite of screen coordinates (origin top-left, Y
increasing downward). This will bite you. It bites everyone. You will
write the conversion wrong the first time. Here it is, so you can copy
it:

\begin{Shaded}
\begin{Highlighting}[]
\CommentTok{// viewport = page.getViewport(\{ scale: scale \})}
\KeywordTok{var}\NormalTok{ left }\OperatorTok{=}\NormalTok{ bbox[}\DecValTok{0}\NormalTok{] }\OperatorTok{*}\NormalTok{ scale}\OperatorTok{;}                              
\KeywordTok{var}\NormalTok{ top  }\OperatorTok{=}\NormalTok{ viewport}\OperatorTok{.}\AttributeTok{height} \OperatorTok{{-}}\NormalTok{ (bbox[}\DecValTok{1}\NormalTok{] }\OperatorTok{+}\NormalTok{ bbox[}\DecValTok{3}\NormalTok{]) }\OperatorTok{*}\NormalTok{ scale}\OperatorTok{;} 
\KeywordTok{var}\NormalTok{ width }\OperatorTok{=}\NormalTok{ bbox[}\DecValTok{2}\NormalTok{] }\OperatorTok{*}\NormalTok{ scale}\OperatorTok{;}                             
\KeywordTok{var}\NormalTok{ height }\OperatorTok{=}\NormalTok{ bbox[}\DecValTok{3}\NormalTok{] }\OperatorTok{*}\NormalTok{ scale}\OperatorTok{;}                            
\end{Highlighting}
\end{Shaded}

\begin{itemize}
\item
  X translates directly (both coordinate systems have X increasing
  rightward)
\item
  Y must be flipped: PDF Y=0 is bottom, screen Y=0 is top. Note the
  \texttt{bbox{[}1{]}\ +\ bbox{[}3{]}} --- that's lower-left-Y plus
  height, giving the top edge in PDF space, which maps to the top edge
  in screen space.
\item
  Width scales linearly
\item
  Height scales linearly
\end{itemize}

Where \texttt{bbox} is
\texttt{{[}lowerLeftX,\ lowerLeftY,\ width,\ height{]}} as extracted by
PDFBox's \texttt{PDRectangle}. Note that PDFBox gives you
\texttt{{[}x,\ y,\ width,\ height{]}} while the raw PDF \texttt{/Rect}
is \texttt{{[}llx,\ lly,\ urx,\ ury{]}} --- PDFBox does the conversion
for you. Don't convert twice.

\subsubsection{Invisible Signatures}\label{_invisible_signatures}

A signature can be \textbf{invisible} --- having no visual appearance on
the page:

\begin{itemize}
\item
  \textbf{Zero-size rect:} \texttt{/Rect\ {[}0\ 0\ 0\ 0{]}} --- the
  widget exists but has no area
\item
  \textbf{No widget at all:} The field has no annotation entry (valid
  per spec, but some viewers complain)
\end{itemize}

Invisible signatures are common for:

\begin{itemize}
\item
  Certification signatures --- the document author certifies the
  document without cluttering the layout
\item
  Machine-generated signatures --- automated workflows where nobody's
  looking at the PDF
\item
  Counter-signatures --- additional approval signatures in complex
  chains where visual space is limited
\end{itemize}

\textbf{Should your signature be visible or invisible?}

For approval signatures intended for human signers: make them visible.
Users expect to \emph{see} their signature on the page. A visible
signature provides a visual anchor --- "I signed here" --- and helps in
audit scenarios.

For certification signatures and automated workflows: invisible is
usually fine. The certification status is shown in the viewer's
signature panel, not on the page.

For machine-to-machine: always invisible. No human is looking at the
page. Adding a visual appearance just wastes space and processing time.

One exception: if you're in a regulated industry where auditors review
printed copies of signed PDFs (yes, this still happens in 2026), visible
signatures are essential because auditors can't check digital signatures
on paper. In that case, include the signer's name, the signing time, and
ideally a note like "Digitally signed --- verify electronically."

\subsection{\texorpdfstring{AcroForm \texttt{/SigFlags}: Document-Level
Signature
Metadata}{AcroForm /SigFlags: Document-Level Signature Metadata}}\label{ch12-pdfsigflags}

\begin{quote}
If the document contains signature fields, the AcroForm dictionary shall
include a SigFlags entry.

---  ISO 32000-2:2020 §12.7.2, Table 225
\end{quote}

\begin{longtable}[]{@{}
  >{\raggedright\arraybackslash}p{(\linewidth - 6\tabcolsep) * \real{0.2857}}
  >{\raggedright\arraybackslash}p{(\linewidth - 6\tabcolsep) * \real{0.1429}}
  >{\raggedright\arraybackslash}p{(\linewidth - 6\tabcolsep) * \real{0.1429}}
  >{\raggedright\arraybackslash}p{(\linewidth - 6\tabcolsep) * \real{0.4286}}@{}}
\caption{SigFlags Bit Values --- ISO 32000-2:2020 §12.7.2, Table
225}\tabularnewline
\toprule\noalign{}
\begin{minipage}[b]{\linewidth}\raggedright
Name
\end{minipage} & \begin{minipage}[b]{\linewidth}\raggedright
Bit
\end{minipage} & \begin{minipage}[b]{\linewidth}\raggedright
Req
\end{minipage} & \begin{minipage}[b]{\linewidth}\raggedright
Spec Ref
\end{minipage} \\
\midrule\noalign{}
\endfirsthead
\toprule\noalign{}
\begin{minipage}[b]{\linewidth}\raggedright
Name
\end{minipage} & \begin{minipage}[b]{\linewidth}\raggedright
Bit
\end{minipage} & \begin{minipage}[b]{\linewidth}\raggedright
Req
\end{minipage} & \begin{minipage}[b]{\linewidth}\raggedright
Spec Ref
\end{minipage} \\
\midrule\noalign{}
\endhead
\bottomrule\noalign{}
\endlastfoot
\texttt{SignaturesExist} & 1 (value 1) & Required & §12.7.2, Table
225 \\
\multicolumn{4}{@{}>{\raggedright\arraybackslash}p{(\linewidth - 6\tabcolsep) * \real{1.0000} + 6\tabcolsep}@{}}{%
At least one signature field has a value. Readers should present
signature-specific UI (e.g., the blue bar at the top of Adobe Acrobat).
Must be set whenever any signature field contains a \texttt{/V}
entry.} \\
\texttt{AppendOnly} & 2 (value 2) & Required & §12.7.2, Table 225 \\
\multicolumn{4}{@{}>{\raggedright\arraybackslash}p{(\linewidth - 6\tabcolsep) * \real{1.0000} + 6\tabcolsep}@{}}{%
The document \textbf{shall} only be saved incrementally. A conforming
reader that allows editing \textbf{must not} do a full save --- it must
append. This is the single most important flag for signature
preservation. Without it, a naive editor may do a full rewrite and
destroy every signature.} \\
\end{longtable}

In practice, both bits should be set (\texttt{/SigFlags\ 3}) whenever
the document contains any signed signature field. This was already
discussed above, but it bears repeating: if you get this wrong,
signatures will be destroyed by innocent save operations, and the blame
will be on your tool.

\subsection{PDFBox API Mapping}\label{ch12-pdfbox-api}

For those of you who skipped here --- I don't blame you, this is the
section with the code --- here's how the spec concepts map to PDFBox 3.x
classes:

{\def\LTcaptype{none} % do not increment counter
\begin{longtable}[]{@{}
  >{\raggedright\arraybackslash}p{(\linewidth - 4\tabcolsep) * \real{0.3000}}
  >{\raggedright\arraybackslash}p{(\linewidth - 4\tabcolsep) * \real{0.4000}}
  >{\raggedright\arraybackslash}p{(\linewidth - 4\tabcolsep) * \real{0.3000}}@{}}
\toprule\noalign{}
\begin{minipage}[b]{\linewidth}\raggedright
PDF Concept
\end{minipage} & \begin{minipage}[b]{\linewidth}\raggedright
PDFBox Class
\end{minipage} & \begin{minipage}[b]{\linewidth}\raggedright
Key Methods
\end{minipage} \\
\midrule\noalign{}
\endhead
\bottomrule\noalign{}
\endlastfoot
Signature field & \texttt{PDSignatureField} & \texttt{getSignature()},
\texttt{getWidgets()}, \texttt{getPartialName()} \\
Signature dictionary & \texttt{PDSignature} & \texttt{getByteRange()},
\texttt{getContents(byte{[}{]})}, \texttt{getSubFilter()},
\texttt{getName()}, \texttt{getSignDate()} \\
AcroForm & \texttt{PDAcroForm} & \texttt{getFields()},
\texttt{getFieldTree()}, \texttt{setSignaturesExist()},
\texttt{setAppendOnly()} \\
Widget annotation & \texttt{PDAnnotationWidget} &
\texttt{getRectangle()}, \texttt{getPage()},
\texttt{getNormalAppearanceStream()} \\
Appearance stream & \texttt{PDFormXObject} & \texttt{getStream()},
\texttt{getBBox()}, \texttt{getResources()} \\
Adding a signature & \texttt{PDDocument} &
\texttt{addSignature(PDSignature,\ SignatureInterface,\ SignatureOptions)} \\
Incremental save & \texttt{PDDocument} &
\texttt{saveIncremental(OutputStream)} \\
\end{longtable}
}

\subsubsection{Minimal Signing Code (Key
Lines)}\label{_minimal_signing_code_key_lines}

The essential PDFBox signing code in \textasciitilde20 lines. This is
the skeleton --- the full, production-ready version is in Appendix A,
but these are the lines that matter:

\begin{Shaded}
\begin{Highlighting}[]
\NormalTok{PDDocument doc }\OperatorTok{=}\NormalTok{ Loader}\OperatorTok{.}\FunctionTok{loadPDF}\OperatorTok{(}\NormalTok{pdfBytes}\OperatorTok{);}                 
\NormalTok{PDAcroForm acroForm }\OperatorTok{=}\NormalTok{ doc}\OperatorTok{.}\FunctionTok{getDocumentCatalog}\OperatorTok{().}\FunctionTok{getAcroForm}\OperatorTok{();}

\NormalTok{PDSignature signature }\OperatorTok{=} \KeywordTok{new} \FunctionTok{PDSignature}\OperatorTok{();}
\NormalTok{signature}\OperatorTok{.}\FunctionTok{setFilter}\OperatorTok{(}\NormalTok{PDSignature}\OperatorTok{.}\FunctionTok{FILTER\_ADOBE\_PPKLITE}\OperatorTok{);}         
\NormalTok{signature}\OperatorTok{.}\FunctionTok{setSubFilter}\OperatorTok{(}\NormalTok{PDSignature}\OperatorTok{.}\FunctionTok{SUBFILTER\_ETSI\_CADES\_DETACHED}\OperatorTok{);} 
\NormalTok{signature}\OperatorTok{.}\FunctionTok{setName}\OperatorTok{(}\StringTok{"John Doe"}\OperatorTok{);}
\NormalTok{signature}\OperatorTok{.}\FunctionTok{setSignDate}\OperatorTok{(}\BuiltInTok{Calendar}\OperatorTok{.}\FunctionTok{getInstance}\OperatorTok{());}

\NormalTok{SignatureOptions options }\OperatorTok{=} \KeywordTok{new} \FunctionTok{SignatureOptions}\OperatorTok{();}
\NormalTok{options}\OperatorTok{.}\FunctionTok{setPreferredSignatureSize}\OperatorTok{(}\BaseNTok{0x6000}\OperatorTok{);}                     
\NormalTok{options}\OperatorTok{.}\FunctionTok{setVisualSignature}\OperatorTok{(}\NormalTok{appearanceStream}\OperatorTok{);}                  

\NormalTok{doc}\OperatorTok{.}\FunctionTok{addSignature}\OperatorTok{(}\NormalTok{signature}\OperatorTok{,}\NormalTok{ signingInterface}\OperatorTok{,}\NormalTok{ options}\OperatorTok{);}        
\NormalTok{doc}\OperatorTok{.}\FunctionTok{saveIncremental}\OperatorTok{(}\NormalTok{outputStream}\OperatorTok{);}                             
\end{Highlighting}
\end{Shaded}

\begin{itemize}
\item
  Always load from bytes. Don't create a new \texttt{PDDocument()} and
  try to \texttt{saveIncremental()} --- PDFBox 3.x will throw
  \texttt{IllegalStateException:\ document\ was\ not\ loaded\ from\ a\ file\ or\ a\ stream}.
  Yes, this error message has confused every single person who's ever
  encountered it.
\item
  Filter is always \texttt{Adobe.PPKLite} for interoperability. Don't be
  creative here.
\item
  SubFilter: use \texttt{ETSI.CAdES.detached} for PAdES compliance. See
  the earlier section for why.
\item
  Pre-allocate 24KB for the CMS blob. \textbf{Not} the PDFBox default of
  9KB. See the pre-allocation table above.
\item
  Visual appearance --- set this to \texttt{null} for invisible
  signatures. Or provide an \texttt{InputStream} of a pre-built
  appearance.
\item
  Register the signature and the \texttt{SignatureInterface}. The
  interface's \texttt{sign()} method will be called during
  \texttt{saveIncremental()}.
\item
  \textbf{Incremental save.} There is no other option. Do not call
  \texttt{save()} on a document you're signing.
\end{itemize}

\subsubsection{\texorpdfstring{The \texttt{SignatureInterface}: Your
Extension
Point}{The SignatureInterface: Your Extension Point}}\label{_the_signatureinterface_your_extension_point}

\begin{Shaded}
\begin{Highlighting}[]
\KeywordTok{public} \KeywordTok{interface}\NormalTok{ SignatureInterface }\OperatorTok{\{}
    \DataTypeTok{byte}\OperatorTok{[]} \FunctionTok{sign}\OperatorTok{(}\BuiltInTok{InputStream}\NormalTok{ content}\OperatorTok{)} \KeywordTok{throws} \BuiltInTok{IOException}\OperatorTok{;} 
\OperatorTok{\}}
\end{Highlighting}
\end{Shaded}

\begin{itemize}
\item
  Called by PDFBox during \texttt{saveIncremental()}. The
  \texttt{InputStream} provides the concatenated byte ranges --- this is
  the data you need to hash and sign. Your implementation builds the CMS
  \texttt{SignedData} blob and returns the DER-encoded bytes.
\end{itemize}

This is the \textbf{only thing you need to implement}. PDFBox handles
the PDF structure, the byte range calculation, the incremental save, the
widget creation, the appearance stream embedding. You handle the
cryptography.

\textbf{Keep your \texttt{SignatureInterface} implementation stateless
and testable.}

Don't put PDF logic in there. Don't reference the \texttt{PDDocument}.
Your \texttt{sign()} method receives an \texttt{InputStream} of bytes
and returns CMS bytes. That's it. This makes it trivially testable ---
pass in any byte stream, check that the output is valid CMS. It also
makes it easy to swap between local signing (BouncyCastle) and remote
signing (HSM, cloud) without touching any PDF code.

\subsection{Common Mistakes and How to Avoid
Them}\label{ch12-common-mistakes}

After years of working with PDF signatures, the same problems come up
over and over. Here's the hit list:

{\def\LTcaptype{none} % do not increment counter
\begin{longtable}[]{@{}
  >{\raggedright\arraybackslash}p{(\linewidth - 4\tabcolsep) * \real{0.1000}}
  >{\raggedright\arraybackslash}p{(\linewidth - 4\tabcolsep) * \real{0.4000}}
  >{\raggedright\arraybackslash}p{(\linewidth - 4\tabcolsep) * \real{0.5000}}@{}}
\toprule\noalign{}
\begin{minipage}[b]{\linewidth}\raggedright
\#
\end{minipage} & \begin{minipage}[b]{\linewidth}\raggedright
Mistake
\end{minipage} & \begin{minipage}[b]{\linewidth}\raggedright
Fix
\end{minipage} \\
\midrule\noalign{}
\endhead
\bottomrule\noalign{}
\endlastfoot
1 & \textbf{Trusting \texttt{/Name} for signer identity.} Many systems
display the \texttt{/Name} field from the signature dictionary as if
it's the verified signer name. It's not. It's free-text set by whoever
controls the signing process. & Always extract the signer's identity
from the X.509 certificate in the CMS structure. Display the
certificate's Subject CN (or the full Subject DN for enterprise
contexts). \\
2 & \textbf{Using \texttt{PDDocument.save()} instead of
\texttt{saveIncremental()}.} Full save rewrites the entire file,
destroying all existing signatures. & Always use
\texttt{saveIncremental()}. Make it a code review checklist item: "Does
this code path call \texttt{save()} on a signed document? Reject." \\
3 & \textbf{Pre-allocating too little space for \texttt{/Contents}.}
Works fine in testing with simple certificates, breaks in production
when chains get longer or timestamps are added. & Use \texttt{0x6000}
(24KB) minimum. For PAdES-B-LT, use \texttt{0x18000} (96KB). Test with
your actual TSA and certificate chain, not with test certificates. \\
4 & \textbf{Forgetting \texttt{signing-certificate-v2} in the CMS
attributes.} Technically valid as PKCS\#7, but not PAdES-compliant and
vulnerable to certificate substitution. & Use BouncyCastle's
\texttt{JcaSignerInfoGeneratorBuilder} with the correct
\texttt{DigestCalculatorProvider}. It adds the attribute automatically
when configured correctly. See the code in Appendix A. \\
5 & \textbf{Not including the certificate chain.} A signature with just
the signer's certificate can't be validated by anyone who doesn't
already have the issuer certificates. & Always include the full chain
(signer + intermediates). You can omit the root CA --- validators
typically have their own trust stores. But include everything else. \\
6 & \textbf{Creating a fresh \texttt{PDDocument} and calling
\texttt{saveIncremental()}.} PDFBox 3.x requires the document to have
been loaded from bytes. & Create a blank PDF, save it to bytes, load it
back, then sign. See the test code in this chapter. \\
7 & \textbf{Setting \texttt{/SigFlags} to 1 instead of 3.} Bit 1 alone
means "signatures exist" but doesn't enforce append-only saving. &
Always set \texttt{/SigFlags} to 3. No exceptions. \\
8 & \textbf{Modifying the PDF between adding the signature and calling
\texttt{saveIncremental()}.} Any change to the document graph between
\texttt{addSignature()} and \texttt{saveIncremental()} may change byte
offsets, breaking the signature. & Set up everything before calling
\texttt{addSignature()}. Once you've added the signature, the very next
call should be \texttt{saveIncremental()}. Nothing in between. \\
\end{longtable}
}

\subsection{Key Points}\label{ch12-key-points}

\begin{itemize}
\item
  A PDF signature = signature field + signature dictionary + widget +
  appearance stream --- all embedded in the document
\item
  The signature dictionary's \texttt{/Contents} holds a CMS
  \texttt{SignedData} blob (hex-encoded, zero-padded)
\item
  \texttt{/ByteRange} defines exactly which file bytes are signed ---
  everything except \texttt{/Contents} itself
\item
  \texttt{/SubFilter} determines the signature format: use
  \texttt{/ETSI.CAdES.detached} for PAdES compliance --- there's no good
  reason not to
\item
  The \texttt{signing-certificate-v2} attribute is what distinguishes
  CAdES from plain PKCS\#7 --- and what prevents certificate
  substitution attacks
\item
  All signatures use incremental save --- original bytes are never
  modified, ever
\item
  \texttt{/SigFlags} must be set to 3 (SignaturesExist + AppendOnly) to
  protect existing signatures --- getting this wrong destroys other
  people's signatures
\item
  Seed value dictionaries (\texttt{/SV}) can constrain how a field is
  signed --- use them in templates to prevent signer mistakes
\item
  Widget \texttt{/Rect} uses PDF coordinates (origin lower-left, Y up)
  --- you will get the conversion wrong the first time, and that's OK
\item
  PDFBox's \texttt{SignatureInterface.sign(InputStream)} is the single
  extension point for all signing logic --- keep it clean, stateless,
  and testable
\item
  When in doubt, pre-allocate more space for \texttt{/Contents} ---
  wasted bytes are cheap, failed signatures are expensive
\item
  "The file has been modified since signing" is usually not a problem
  --- it's usually just a later incremental update
\end{itemize}

\subsection{Specification References}\label{ch12-spec-refs}

\chapter{Byte Range Deep Dive}\label{_byte_range_deep_dive}

The byte range is where PDF signatures live or die. Everything in
Chapter 12 --- the \texttt{/Sig} dictionary, the \texttt{/SubFilter},
the certificate chain, the CMS structure --- all of it is meaningless if
the byte range is wrong. A single off-by-one error in the
\texttt{/ByteRange} array, and the signature is cryptographically
invalid. A subtle gap in coverage, and an attacker can inject content
that the signature appears to endorse.

This chapter takes the byte range apart byte by byte. We will trace
exactly which bytes are signed, where the gap is, why the gap exists,
what lives in the gap, and what goes wrong when any of these assumptions
are violated. We will look at hex dumps of real PDFs, visualize the byte
range geometry, walk through multi-signature scenarios, and examine the
exact attacks that exploit byte range weaknesses.

If Chapter 12 taught you \emph{what} a byte range is, this chapter
teaches you \emph{how} it actually works in the raw file --- and how to
verify that it works correctly.

\section{The Byte Range Model}\label{_the_byte_range_model}

\subsection{What the Spec Says}\label{_what_the_spec_says}

ISO 32000-2:2020 §12.8.1 defines the \texttt{/ByteRange} array as:

\begin{quote}
An array of pairs of integers (starting byte offset, length in bytes)
that shall describe the exact byte range for the digest computation.
Multiple discontiguous byte ranges shall be used to describe a digest
that does not include the signature value itself.
\end{quote}

In practice, every PDF signature uses exactly \textbf{one pair of byte
ranges} --- two segments --- with a gap between them. The gap contains
the \texttt{/Contents} value (the hex-encoded CMS signature). The two
segments cover everything \emph{else} in the file.

The \texttt{/ByteRange} array always has exactly four integers:

\begin{verbatim}
/ByteRange [offset1 length1 offset2 length2]
\end{verbatim}

Where:

\begin{itemize}
\item
  \texttt{offset1} = 0 (always --- the signed region starts at the first
  byte of the file)
\item
  \texttt{length1} = number of bytes from the start of the file to the
  opening \texttt{\textless{}} of the \texttt{/Contents} hex string
\item
  \texttt{offset2} = byte offset immediately after the closing
  \texttt{\textgreater{}} of the \texttt{/Contents} hex string
\item
  \texttt{length2} = number of bytes from \texttt{offset2} to the end of
  the file
\end{itemize}

The math is simple:

\begin{itemize}
\item
  \textbf{Signed bytes} = \texttt{{[}0,\ length1)} +
  \texttt{{[}offset2,\ offset2\ +\ length2)}
\item
  \textbf{Unsigned gap} = \texttt{{[}length1,\ offset2)} --- this
  contains only the \texttt{/Contents} hex string value
\item
  \textbf{Total file size} = \texttt{offset2\ +\ length2}
\end{itemize}

\subsection{Visualizing the Byte
Range}\label{_visualizing_the_byte_range}

To make this concrete, let's look at a real signed PDF. Suppose we have
a signed file of 45,000 bytes with a 32KB (32,768 hex characters =
16,384 bytes) reserved \texttt{/Contents} value:

\begin{verbatim}
File layout (not to scale):

Byte 0                                                     Byte 44,999
|================================================================|
|  PDF header, objects, pages,    |  /Contents  |  xref, trailer |
|  /Sig dictionary (up to the    |  hex string |  %%EOF         |
|  opening < of /Contents)       |  (the gap)  |                |
|================================|=============|================|
|<---- length1 = 14,000 -------->|<-- gap  --->|<-- length2 -->|
|                                 |  = 16,408   |   = 14,592    |
|                                 |             |                |
|  SIGNED (segment 1)            |  NOT SIGNED |  SIGNED (seg 2)|
|================================|=============|================|

/ByteRange [0 14000 30408 14592]
\end{verbatim}

The gap is slightly larger than 16,384 bytes because it includes the
\texttt{\textless{}} and \texttt{\textgreater{}} delimiters that bracket
the hex string, plus the `/Contents ` key prefix. The exact gap size
depends on how PDFBox (or whatever library produced the file) serializes
the signature dictionary.

\subsection{Why the Gap Exists: The Chicken-and-Egg
Problem}\label{_why_the_gap_exists_the_chicken_and_egg_problem}

The gap exists because of a fundamental bootstrapping problem:

\begin{enumerate}
\def\labelenumi{\arabic{enumi}.}
\item
  To \textbf{compute the signature}, you need to hash the bytes of the
  file
\item
  To \textbf{write the file}, you need to include the signature (inside
  \texttt{/Contents})
\item
  But the signature is \emph{part of} the file --- so how can you hash a
  file that doesn't exist yet?
\end{enumerate}

The solution: \textbf{exclude the signature value from the hash}. The
\texttt{/Contents} hex string is the one part of the file that is
\emph{not} covered by the signature. Everything else --- the PDF header,
all objects, the signature dictionary itself (including
\texttt{/ByteRange}, \texttt{/Filter}, \texttt{/SubFilter}, and all
other keys), the cross-reference table, and the trailer --- is signed.

This means the signing process works like this:

\begin{enumerate}
\def\labelenumi{\arabic{enumi}.}
\item
  Write the entire PDF file, including the \texttt{/Sig} dictionary, but
  leave the \texttt{/Contents} value as a placeholder (all zeros)
\item
  Record the byte offsets of where the placeholder starts and ends
\item
  Compute the hash over everything \emph{except} the placeholder:
  \texttt{hash\ =\ SHA256(segment1\ \textbar{}\textbar{}\ segment2)}
\item
  Construct the CMS \texttt{SignedData} using this hash
\item
  Hex-encode the CMS bytes and write them into the placeholder space
\item
  If the CMS is shorter than the reserved space (it almost always is),
  pad the remainder with zeros
\end{enumerate}

This is why PDFBox's \texttt{SignatureInterface.sign(InputStream)}
receives an \texttt{InputStream} --- that stream contains the
concatenation of segment 1 and segment 2 (the bytes \emph{outside} the
gap), ready to be hashed.

The \texttt{/ByteRange} array itself is part of the signed region. This
is critical for security: if an attacker could modify the
\texttt{/ByteRange} after signing, they could redirect the signature to
cover different content. The byte range is committed before the hash is
computed, and any change to it invalidates the signature.

\section{The Gap in Detail}\label{_the_gap_in_detail}

\subsection{What's Actually in the
Gap}\label{_whats_actually_in_the_gap}

Let's look at a hex dump of the gap region. In a properly signed PDF,
the gap contains exactly one thing: the hex-encoded CMS
\texttt{SignedData}, padded with zeros:

\begin{verbatim}
Byte 14000:  3C                                           -- opening '<'
Byte 14001:  30 38 30 36 30 39 ...                        -- CMS bytes (hex-encoded)
             ...
             ... (several KB of hex characters)
             ...
Byte 28500:  30 30 30 30 30 30 30 30 30 30 30 30 ...     -- zero padding ("000000...")
             ...
Byte 30407:  3E                                           -- closing '>'
\end{verbatim}

The hex encoding doubles the size: a 4,096-byte CMS signature becomes
8,192 hex characters. The reserved space is typically 16KB (32,768 hex
characters = 16,384 bytes of actual CMS capacity), which is enough for
most PAdES-B-B signatures.

\textbf{PAdES-B-LTA signatures are much larger} than PAdES-B-B. A B-LTA
signature includes:

\begin{itemize}
\item
  The CMS \texttt{SignedData} itself (\textasciitilde2--4 KB)
\item
  The signing certificate and full chain (\textasciitilde2--5 KB)
\item
  The timestamp token from the TSA (\textasciitilde2--4 KB)
\item
  OCSP responses for each certificate in the chain (\textasciitilde1--3
  KB each)
\item
  CRL data if OCSP is not available (\textasciitilde5--50 KB)
\end{itemize}

A B-LTA signature can easily reach 20--30 KB, which requires 40--60 KB
of hex-encoded space. If you're producing B-LTA signatures, reserve at
least \textbf{64 KB} for \texttt{/Contents}:

\begin{Shaded}
\begin{Highlighting}[]
\CommentTok{// PDFBox: set signature content size (in bytes, before hex encoding)}
\NormalTok{SignatureOptions options }\OperatorTok{=} \KeywordTok{new} \FunctionTok{SignatureOptions}\OperatorTok{();}
\NormalTok{options}\OperatorTok{.}\FunctionTok{setPreferredSignatureSize}\OperatorTok{(}\DecValTok{65536}\OperatorTok{);} \CommentTok{// 64 KB}
\NormalTok{doc}\OperatorTok{.}\FunctionTok{addSignature}\OperatorTok{(}\NormalTok{signature}\OperatorTok{,}\NormalTok{ signingInterface}\OperatorTok{,}\NormalTok{ options}\OperatorTok{);}
\end{Highlighting}
\end{Shaded}

If the CMS exceeds the reserved space, PDFBox throws an
\texttt{IOException} --- the signing fails silently in the sense that no
file is produced. This is the \#1 cause of "signing worked in
development but fails in production" --- development certificates have
short chains and no timestamps, production certificates have full
chains, OCSP responses, and timestamps that exceed the default buffer.

\subsection{The Zero-Padding Question}\label{_the_zero_padding_question}

The zero padding at the end of the \texttt{/Contents} hex string is not
part of the CMS. It's just filler --- the hex string \texttt{00} decodes
to a null byte, and CMS parsers stop reading at the end of the
DER-encoded ASN.1 structure, ignoring trailing zeros.

Some validators are stricter than others about this padding:

\begin{itemize}
\item
  \textbf{Adobe Acrobat}: Ignores trailing zero padding --- works fine
\item
  \textbf{LibreOffice}: Ignores trailing zeros --- works fine
\item
  \textbf{Some custom validators}: May compute the hash of the entire
  \texttt{/Contents} hex string (including padding) and fail because the
  hash doesn't match what the signer computed over just the CMS bytes.
  This is a bug in the validator, not in the signed PDF.
\end{itemize}

When writing validation code, always parse the \texttt{/Contents} hex
string as DER and stop at the end of the ASN.1 structure. Do not hash
the raw hex string --- decode it to binary, find the end of the DER
sequence (the total length indicated in the ASN.1 header), and ignore
everything after that.

\section{Incremental Saves and the Byte
Range}\label{_incremental_saves_and_the_byte_range}

\subsection{How Signing Creates an Incremental
Update}\label{_how_signing_creates_an_incremental_update}

When PDFBox signs a PDF, it performs an \textbf{incremental save} --- it
appends new data to the end of the existing file rather than rewriting
the entire file. This is not optional for signing; it is how signing
must work.

An incrementally saved PDF looks like this:

\begin{figure}
\centering
\pandocbounded{\includegraphics[keepaspectratio,alt={byterange incremental}]{diagrams/pdf/byterange-incremental.pdf}}
\caption{PDF file before and after signing --- incremental update
appends new objects without modifying the original file}
\end{figure}

The "new objects" include the signature dictionary (or modifications to
an existing signature field), any new annotation objects for the visual
appearance, and updated AcroForm entries. The "new xref" is a
cross-reference section that records the locations of these new objects,
and the "new trailer" points to the new xref.

The entire appended region --- including the new \texttt{\%\%EOF} --- is
covered by the byte range of the signature. The byte range starts at
byte 0 (the original file's header) and extends to the end of the
appended data, with only the \texttt{/Contents} gap excluded.

This means:

\begin{itemize}
\item
  The signature covers the \textbf{original file} in its entirety
\item
  The signature covers the \textbf{new objects} added during signing
  (the signature dictionary, appearances, etc.)
\item
  The signature covers the \textbf{new cross-reference table} and
  \textbf{new trailer}
\item
  The only thing not covered is the CMS signature value itself
\end{itemize}

\subsection{What Happens After Signing: The Trailing Bytes
Problem}\label{_what_happens_after_signing_the_trailing_bytes_problem}

Here is one of the most security-critical aspects of byte range
handling: \textbf{what happens to bytes that appear after the signed
region?}

Consider this scenario:

\begin{figure}
\centering
\pandocbounded{\includegraphics[keepaspectratio,alt={byterange single}]{diagrams/pdf/byterange-single.pdf}}
\caption{Byte range coverage --- the signature covers two segments with
a gap for \protect\texttt{/Contents}}
\end{figure}

If \texttt{offset2\ +\ length2} equals the total file size, the
signature covers everything --- good.

But what if someone appends data \emph{after} the signed region?

\begin{figure}
\centering
\pandocbounded{\includegraphics[keepaspectratio,alt={byterange trailing}]{diagrams/pdf/byterange-trailing.pdf}}
\caption{Trailing bytes after the signed region --- the Incremental
Saving Attack surface}
\end{figure}

The byte range still validates --- the hash over the signed segments
hasn't changed. But the file now contains unsigned data at the end. If
this unsigned data is a valid incremental update (with its own xref,
trailer, and \texttt{\%\%EOF}), a PDF viewer might load it and display
the \emph{attacker's} content while showing the original signature as
valid.

This is exactly the \textbf{Incremental Saving Attack (ISA)} that the
Bochum team demonstrated in 2019 (Chapter 24). The attack works because
PDF viewers before the patches would happily load incremental updates
appended after the signed byte range, display the new content, and still
report the signature as valid --- because the signed bytes hadn't been
modified.

\textbf{The trailing bytes check is the most critical validation step.}

When verifying a PDF signature, you \textbf{must} verify that:

\begin{verbatim}
offset2 + length2 == actual file size
\end{verbatim}

If the file is larger than the signed region, something was appended
after signing. This doesn't necessarily mean the file was tampered with
--- a second legitimate signature would also append data --- but it
means you must validate the second revision's signature independently.

If there is unsigned trailing data that is \emph{not} a valid subsequent
signature revision, the document has been modified after signing and the
signature should be treated as covering only the original content.

Chapter 24 provides the full attack taxonomy and defense code.

\subsection{Multi-Signature Scenarios}\label{_multi_signature_scenarios}

Enterprise PDF workflows often involve multiple signatures --- a
document might be signed by the author, countersigned by a manager, and
certified by a compliance officer. Each signature creates its own
incremental update.

The byte range geometry for a multiply-signed PDF looks like this:

\begin{figure}
\centering
\pandocbounded{\includegraphics[keepaspectratio,alt={byterange multi}]{diagrams/pdf/byterange-multi.pdf}}
\caption{Multi-signature byte range nesting --- each subsequent
signature covers everything except its own gap}
\end{figure}

Each signature covers everything before its own gap (including previous
signatures\textquotesingle{} gaps --- the hex bytes of
\texttt{/Contents1} and \texttt{/Contents2} are part of the signed data
for signatures 2 and 3). This nesting means:

\begin{itemize}
\item
  \textbf{Signature 1} covers the original PDF + the objects and xref
  added during the first signing
\item
  \textbf{Signature 2} covers everything that signature 1 covers PLUS
  the objects added during the second signing. This means signature 2
  implicitly endorses signature 1's existence (though not necessarily
  its validity)
\item
  \textbf{Signature 3} covers everything, providing the broadest
  coverage
\end{itemize}

Each signature is a snapshot of the document at the time it was applied.
Signature 1 saw the document in state S1. Signature 2 saw it in state S2
(which includes signature 1). If you need to know what the document
looked like when signature 1 was applied, you must extract revision 1 by
truncating the file at the end of signature 1's byte range ---
everything after that was added later.

PDFBox provides this capability:

\begin{Shaded}
\begin{Highlighting}[]
\CommentTok{// Extract the revision that a specific signature covers}
\NormalTok{PDDocument doc }\OperatorTok{=}\NormalTok{ Loader}\OperatorTok{.}\FunctionTok{loadPDF}\OperatorTok{(}\NormalTok{signedFile}\OperatorTok{);}
\BuiltInTok{List}\OperatorTok{\textless{}}\NormalTok{PDSignature}\OperatorTok{\textgreater{}}\NormalTok{ signatures }\OperatorTok{=}\NormalTok{ doc}\OperatorTok{.}\FunctionTok{getSignatureDictionaries}\OperatorTok{();}
\ControlFlowTok{for} \OperatorTok{(}\NormalTok{PDSignature sig }\OperatorTok{:}\NormalTok{ signatures}\OperatorTok{)} \OperatorTok{\{}
    \DataTypeTok{int}\OperatorTok{[]}\NormalTok{ byteRange }\OperatorTok{=}\NormalTok{ sig}\OperatorTok{.}\FunctionTok{getByteRange}\OperatorTok{();}
    \DataTypeTok{int}\NormalTok{ revisionEnd }\OperatorTok{=}\NormalTok{ byteRange}\OperatorTok{[}\DecValTok{2}\OperatorTok{]} \OperatorTok{+}\NormalTok{ byteRange}\OperatorTok{[}\DecValTok{3}\OperatorTok{];}

    \CommentTok{// Extract this revision}
    \DataTypeTok{byte}\OperatorTok{[]}\NormalTok{ revision }\OperatorTok{=} \KeywordTok{new} \DataTypeTok{byte}\OperatorTok{[}\NormalTok{revisionEnd}\OperatorTok{];}
    \ControlFlowTok{try} \OperatorTok{(}\BuiltInTok{InputStream}\NormalTok{ is }\OperatorTok{=} \KeywordTok{new} \BuiltInTok{FileInputStream}\OperatorTok{(}\NormalTok{signedFile}\OperatorTok{))} \OperatorTok{\{}
\NormalTok{        is}\OperatorTok{.}\FunctionTok{readNBytes}\OperatorTok{(}\NormalTok{revision}\OperatorTok{,} \DecValTok{0}\OperatorTok{,}\NormalTok{ revisionEnd}\OperatorTok{);}
    \OperatorTok{\}}
    \CommentTok{// \textquotesingle{}revision\textquotesingle{} is now the PDF as it existed when this signature was applied}
\OperatorTok{\}}
\end{Highlighting}
\end{Shaded}

\subsection{Permission-Based Restrictions: What Can Change After
Signing?}\label{_permission_based_restrictions_what_can_change_after_signing}

Not all changes after signing are attacks. The PDF specification defines
three \textbf{modification detection and prevention (MDP)} permission
levels that control what changes are allowed after a certification
signature:

\begin{longtable}[]{@{}
  >{\raggedright\arraybackslash}p{(\linewidth - 6\tabcolsep) * \real{0.1250}}
  >{\raggedright\arraybackslash}p{(\linewidth - 6\tabcolsep) * \real{0.2500}}
  >{\raggedright\arraybackslash}p{(\linewidth - 6\tabcolsep) * \real{0.1250}}
  >{\raggedright\arraybackslash}p{(\linewidth - 6\tabcolsep) * \real{0.5000}}@{}}
\caption{DocMDP Permission Levels (ISO 32000-2:2020
§12.8.2.2.2)}\tabularnewline
\toprule\noalign{}
\begin{minipage}[b]{\linewidth}\raggedright
Level
\end{minipage} & \begin{minipage}[b]{\linewidth}\raggedright
Name
\end{minipage} & \begin{minipage}[b]{\linewidth}\raggedright
Req
\end{minipage} & \begin{minipage}[b]{\linewidth}\raggedright
Description
\end{minipage} \\
\midrule\noalign{}
\endfirsthead
\toprule\noalign{}
\begin{minipage}[b]{\linewidth}\raggedright
Level
\end{minipage} & \begin{minipage}[b]{\linewidth}\raggedright
Name
\end{minipage} & \begin{minipage}[b]{\linewidth}\raggedright
Req
\end{minipage} & \begin{minipage}[b]{\linewidth}\raggedright
Description
\end{minipage} \\
\midrule\noalign{}
\endhead
\bottomrule\noalign{}
\endlastfoot
P=1 & No changes & Optional & No changes whatsoever are allowed. Any
modification --- including filling form fields or adding annotations ---
invalidates the signature. This is the most restrictive level. \\
P=2 & Form filling + signing & Default & Users may fill form fields and
apply additional signatures. Other changes (adding pages, modifying
text, changing form field definitions) invalidate the signature. This is
the most common level for government forms. \\
P=3 & Form filling + signing + annotations & Optional & Users may fill
form fields, sign, and add annotations (comments, highlights, stamps).
This is the most permissive level while still detecting unauthorized
modifications. \\
\end{longtable}

When validating a signature with MDP restrictions, you must check not
only whether the byte range is intact, but also whether the
\emph{content} of subsequent incremental updates falls within the
allowed permissions. This is where validation gets complex --- you need
to compare the PDF objects in the signed revision against the current
revision and determine whether the differences are permitted.

Adobe Acrobat does this analysis automatically and shows warnings like
"The document has been altered since the signature was applied, but the
changes are permitted." Building this analysis yourself requires
understanding of the PDF object model at a level covered in Chapters 7
and 8.

\section{Size Estimation: How Much Space to
Reserve}\label{_size_estimation_how_much_space_to_reserve}

\subsection{The Problem}\label{_the_problem_2}

Before you can sign a PDF, you must write it to disk with a placeholder
for \texttt{/Contents}. The placeholder must be large enough to hold the
final CMS signature. If it's too small, the signing fails. If it's too
large, you waste space (and the file contains kilobytes of zero
padding).

The challenge: the exact CMS size depends on:

\begin{itemize}
\item
  \textbf{The signing certificate} --- larger key sizes (RSA-4096 vs
  RSA-2048) produce larger signatures
\item
  \textbf{The certificate chain} --- more intermediates = more bytes
\item
  \textbf{The signed attributes} --- PAdES requires
  ESS-signing-certificate-v2, content-type, message-digest; more
  attributes = more bytes
\item
  \textbf{Timestamps} --- a TSA response is itself a CMS
  \texttt{SignedData}, typically 2--4 KB
\item
  \textbf{OCSP responses} --- each certificate in the chain may have a
  1--3 KB OCSP response
\item
  \textbf{CRL data} --- if a certificate's revocation information is
  provided via CRL instead of OCSP, CRLs can be 5--500 KB depending on
  the CA's revocation volume
\end{itemize}

You won't know the exact sizes until you've fetched the timestamp and
revocation data --- but you need to commit the space \emph{before}
computing the hash.

\subsection{Estimation Strategies}\label{_estimation_strategies}

\textbf{Strategy 1: Fixed generous buffer (simple, recommended)}

Reserve a fixed buffer that's large enough for worst-case scenarios:

\begin{Shaded}
\begin{Highlighting}[]
\CommentTok{// Signature size estimates by PAdES profile}
\DataTypeTok{int}\NormalTok{ estimatedSize}\OperatorTok{;}
\ControlFlowTok{switch} \OperatorTok{(}\NormalTok{profile}\OperatorTok{)} \OperatorTok{\{}
    \ControlFlowTok{case}\NormalTok{ B\_B}\OperatorTok{:}\NormalTok{   estimatedSize }\OperatorTok{=} \DecValTok{16384}\OperatorTok{;}  \ControlFlowTok{break}\OperatorTok{;} \CommentTok{// 16 KB — basic CMS, no timestamp}
    \ControlFlowTok{case}\NormalTok{ B\_T}\OperatorTok{:}\NormalTok{   estimatedSize }\OperatorTok{=} \DecValTok{32768}\OperatorTok{;}  \ControlFlowTok{break}\OperatorTok{;} \CommentTok{// 32 KB — CMS + timestamp token}
    \ControlFlowTok{case}\NormalTok{ B\_LT}\OperatorTok{:}\NormalTok{  estimatedSize }\OperatorTok{=} \DecValTok{65536}\OperatorTok{;}  \ControlFlowTok{break}\OperatorTok{;} \CommentTok{// 64 KB — CMS + timestamp + OCSP/CRL}
    \ControlFlowTok{case}\NormalTok{ B\_LTA}\OperatorTok{:}\NormalTok{ estimatedSize }\OperatorTok{=} \DecValTok{131072}\OperatorTok{;} \ControlFlowTok{break}\OperatorTok{;} \CommentTok{// 128 KB — CMS + timestamp + OCSP/CRL + archival TS}
\OperatorTok{\}}
\NormalTok{SignatureOptions options }\OperatorTok{=} \KeywordTok{new} \FunctionTok{SignatureOptions}\OperatorTok{();}
\NormalTok{options}\OperatorTok{.}\FunctionTok{setPreferredSignatureSize}\OperatorTok{(}\NormalTok{estimatedSize}\OperatorTok{);}
\end{Highlighting}
\end{Shaded}

This wastes some space (a B-B signature typically needs only 4--8 KB of
the 16 KB reserved), but the waste is negligible compared to the content
of a typical PDF. A 100 KB PDF with 16 KB of zero padding in
\texttt{/Contents} is 116 KB --- nobody cares.

\textbf{Strategy 2: Pre-compute exact size (complex, rarely necessary)}

If you're generating millions of signed PDFs per day and storage costs
matter, you can:

\begin{enumerate}
\def\labelenumi{\arabic{enumi}.}
\item
  Fetch the timestamp and revocation data \emph{before} creating the
  placeholder
\item
  Compute the CMS size exactly
\item
  Reserve only the exact space needed (plus a small safety margin)
\end{enumerate}

This is significantly more complex because you need to construct a "dry
run" CMS to measure it, and the timestamp/OCSP data might change between
the dry run and the actual signing. In practice, this optimization is
almost never worth the complexity.

\textbf{Just use 64 KB (65536 bytes) as your default.} It's enough for
PAdES-B-LT in the vast majority of cases, it wastes a trivial amount of
space, and it eliminates the entire class of "signature too large for
reserved space" failures. Only reduce it if you have a specific reason
(embedded systems with extreme storage constraints) and only increase it
if you're producing B-LTA signatures with CRLs from large CAs.

\subsection{What Happens When the Estimate Is
Wrong}\label{_what_happens_when_the_estimate_is_wrong}

If the CMS signature is \textbf{smaller} than the reserved space: no
problem. The remaining space is filled with zero bytes. All conforming
validators ignore the padding.

If the CMS signature is \textbf{larger} than the reserved space:
\textbf{the signing fails}. PDFBox throws an \texttt{IOException} with a
message like:

\begin{verbatim}
IOException: Can't write signature, not enough space
\end{verbatim}

When this happens:

\begin{enumerate}
\def\labelenumi{\arabic{enumi}.}
\item
  The output file is corrupted (partial write) --- discard it
\item
  Increase \texttt{SignatureOptions.setPreferredSignatureSize()} and
  retry
\item
  To debug: log the actual CMS size from your
  \texttt{SignatureInterface.sign()} method and compare it to the
  reserved size
\end{enumerate}

\begin{Shaded}
\begin{Highlighting}[]
\AttributeTok{@Override}
\KeywordTok{public} \DataTypeTok{byte}\OperatorTok{[]} \FunctionTok{sign}\OperatorTok{(}\BuiltInTok{InputStream}\NormalTok{ content}\OperatorTok{)} \KeywordTok{throws} \BuiltInTok{IOException} \OperatorTok{\{}
    \DataTypeTok{byte}\OperatorTok{[]}\NormalTok{ cmsBytes }\OperatorTok{=} \FunctionTok{createCMSSignature}\OperatorTok{(}\NormalTok{content}\OperatorTok{);}
\NormalTok{    log}\OperatorTok{.}\FunctionTok{info}\OperatorTok{(}\StringTok{"CMS signature size: \{\} bytes (reserved: \{\} bytes)"}\OperatorTok{,}
\NormalTok{             cmsBytes}\OperatorTok{.}\FunctionTok{length}\OperatorTok{,}\NormalTok{ reservedSize}\OperatorTok{);}
    \ControlFlowTok{return}\NormalTok{ cmsBytes}\OperatorTok{;}
\OperatorTok{\}}
\end{Highlighting}
\end{Shaded}

\section{Byte Range Verification: How to Check
It}\label{_byte_range_verification_how_to_check_it}

\subsection{The Verification
Algorithm}\label{_the_verification_algorithm}

When validating a signature, the byte range verification follows these
steps:

\begin{enumerate}
\def\labelenumi{\arabic{enumi}.}
\item
  \textbf{Read the \texttt{/ByteRange} array} from the signature
  dictionary: \texttt{{[}offset1,\ length1,\ offset2,\ length2{]}}
\item
  \textbf{Verify structural sanity}:

  \begin{itemize}
  \item
    \texttt{offset1} must be 0
  \item
    \texttt{length1} must be positive
  \item
    \texttt{offset2} must equal \texttt{length1} + gap size (where gap
    size = the length of the \texttt{/Contents} hex string including
    delimiters)
  \item
    \texttt{length2} must be positive
  \item
    \texttt{offset2\ +\ length2} should equal the file size (or the size
    of the revision, for multi-signature documents)
  \end{itemize}
\item
  \textbf{Verify the gap contains only \texttt{/Contents}}: the bytes in
  \texttt{{[}length1,\ offset2)} should be a hex string containing the
  CMS signature value
\item
  \textbf{Extract the signed bytes}: concatenate
  \texttt{{[}0,\ length1)} and
  \texttt{{[}offset2,\ offset2\ +\ length2)}
\item
  \textbf{Compute the hash}: \texttt{hash\ =\ SHA-256(signed\_bytes)}
  (or whatever digest algorithm the CMS specifies)
\item
  \textbf{Verify the CMS}: parse the CMS \texttt{SignedData} from
  \texttt{/Contents}, extract the signed \texttt{messageDigest}
  attribute, and compare it to the hash you just computed. If they
  match, and the CMS cryptographic verification passes, the byte range
  is intact.
\end{enumerate}

\subsection{Implementation in Java}\label{_implementation_in_java}

\begin{Shaded}
\begin{Highlighting}[]
\CommentTok{/**}
 \CommentTok{*}\NormalTok{ Verifies the byte range integrity of a PDF signature}\CommentTok{.}
 \CommentTok{*}\NormalTok{ Returns true if the byte range covers the expected regions}
 \CommentTok{*}\NormalTok{ and the hash matches the CMS digest}\CommentTok{.}
 \CommentTok{*/}
\KeywordTok{public} \DataTypeTok{boolean} \FunctionTok{verifyByteRange}\OperatorTok{(}\DataTypeTok{byte}\OperatorTok{[]}\NormalTok{ pdfBytes}\OperatorTok{,}\NormalTok{ PDSignature signature}\OperatorTok{)}
        \KeywordTok{throws} \BuiltInTok{IOException}\OperatorTok{,}\NormalTok{ CMSException }\OperatorTok{\{}
    \DataTypeTok{int}\OperatorTok{[]}\NormalTok{ byteRange }\OperatorTok{=}\NormalTok{ signature}\OperatorTok{.}\FunctionTok{getByteRange}\OperatorTok{();}

    \CommentTok{// Step 1: Structural checks}
    \ControlFlowTok{if} \OperatorTok{(}\NormalTok{byteRange}\OperatorTok{.}\FunctionTok{length} \OperatorTok{!=} \DecValTok{4}\OperatorTok{)} \OperatorTok{\{}
\NormalTok{        log}\OperatorTok{.}\FunctionTok{error}\OperatorTok{(}\StringTok{"ByteRange must have exactly 4 elements, got \{\}"}\OperatorTok{,}\NormalTok{ byteRange}\OperatorTok{.}\FunctionTok{length}\OperatorTok{);}
        \ControlFlowTok{return} \KeywordTok{false}\OperatorTok{;}
    \OperatorTok{\}}

    \DataTypeTok{int}\NormalTok{ offset1 }\OperatorTok{=}\NormalTok{ byteRange}\OperatorTok{[}\DecValTok{0}\OperatorTok{];}
    \DataTypeTok{int}\NormalTok{ length1 }\OperatorTok{=}\NormalTok{ byteRange}\OperatorTok{[}\DecValTok{1}\OperatorTok{];}
    \DataTypeTok{int}\NormalTok{ offset2 }\OperatorTok{=}\NormalTok{ byteRange}\OperatorTok{[}\DecValTok{2}\OperatorTok{];}
    \DataTypeTok{int}\NormalTok{ length2 }\OperatorTok{=}\NormalTok{ byteRange}\OperatorTok{[}\DecValTok{3}\OperatorTok{];}

    \ControlFlowTok{if} \OperatorTok{(}\NormalTok{offset1 }\OperatorTok{!=} \DecValTok{0}\OperatorTok{)} \OperatorTok{\{}
\NormalTok{        log}\OperatorTok{.}\FunctionTok{error}\OperatorTok{(}\StringTok{"ByteRange offset1 must be 0, got \{\}"}\OperatorTok{,}\NormalTok{ offset1}\OperatorTok{);}
        \ControlFlowTok{return} \KeywordTok{false}\OperatorTok{;}
    \OperatorTok{\}}

    \CommentTok{// Step 2: Check coverage}
    \DataTypeTok{int}\NormalTok{ signedEnd }\OperatorTok{=}\NormalTok{ offset2 }\OperatorTok{+}\NormalTok{ length2}\OperatorTok{;}
    \ControlFlowTok{if} \OperatorTok{(}\NormalTok{signedEnd }\OperatorTok{!=}\NormalTok{ pdfBytes}\OperatorTok{.}\FunctionTok{length}\OperatorTok{)} \OperatorTok{\{}
\NormalTok{        log}\OperatorTok{.}\FunctionTok{warn}\OperatorTok{(}\StringTok{"ByteRange ends at \{\} but file size is \{\} — "} \OperatorTok{+}
                 \StringTok{"\{\} trailing bytes are unsigned"}\OperatorTok{,}
\NormalTok{                 signedEnd}\OperatorTok{,}\NormalTok{ pdfBytes}\OperatorTok{.}\FunctionTok{length}\OperatorTok{,}\NormalTok{ pdfBytes}\OperatorTok{.}\FunctionTok{length} \OperatorTok{{-}}\NormalTok{ signedEnd}\OperatorTok{);}
        \CommentTok{// This is a }\AlertTok{WARNING}\CommentTok{, not necessarily an error —}
        \CommentTok{// subsequent signatures would extend the file}
    \OperatorTok{\}}

    \CommentTok{// Step 3: Extract and hash signed bytes}
    \BuiltInTok{MessageDigest}\NormalTok{ digest }\OperatorTok{=} \BuiltInTok{MessageDigest}\OperatorTok{.}\FunctionTok{getInstance}\OperatorTok{(}\StringTok{"SHA{-}256"}\OperatorTok{);}
\NormalTok{    digest}\OperatorTok{.}\FunctionTok{update}\OperatorTok{(}\NormalTok{pdfBytes}\OperatorTok{,}\NormalTok{ offset1}\OperatorTok{,}\NormalTok{ length1}\OperatorTok{);}       \CommentTok{// segment 1}
\NormalTok{    digest}\OperatorTok{.}\FunctionTok{update}\OperatorTok{(}\NormalTok{pdfBytes}\OperatorTok{,}\NormalTok{ offset2}\OperatorTok{,}\NormalTok{ length2}\OperatorTok{);}        \CommentTok{// segment 2}
    \DataTypeTok{byte}\OperatorTok{[]}\NormalTok{ computedHash }\OperatorTok{=}\NormalTok{ digest}\OperatorTok{.}\FunctionTok{digest}\OperatorTok{();}

    \CommentTok{// Step 4: Parse CMS and compare}
    \DataTypeTok{byte}\OperatorTok{[]}\NormalTok{ cmsBytes }\OperatorTok{=}\NormalTok{ signature}\OperatorTok{.}\FunctionTok{getContents}\OperatorTok{();}
\NormalTok{    CMSSignedData cmsData }\OperatorTok{=} \KeywordTok{new} \FunctionTok{CMSSignedData}\OperatorTok{(}\NormalTok{cmsBytes}\OperatorTok{);}
\NormalTok{    SignerInformation signer }\OperatorTok{=}\NormalTok{ cmsData}\OperatorTok{.}\FunctionTok{getSignerInfos}\OperatorTok{()}
            \OperatorTok{.}\FunctionTok{getSigners}\OperatorTok{().}\FunctionTok{iterator}\OperatorTok{().}\FunctionTok{next}\OperatorTok{();}

    \CommentTok{// Extract the messageDigest from signed attributes}
    \BuiltInTok{Attribute}\NormalTok{ digestAttr }\OperatorTok{=}\NormalTok{ signer}\OperatorTok{.}\FunctionTok{getSignedAttributes}\OperatorTok{()}
            \OperatorTok{.}\FunctionTok{get}\OperatorTok{(}\NormalTok{CMSAttributes}\OperatorTok{.}\FunctionTok{messageDigest}\OperatorTok{);}
    \DataTypeTok{byte}\OperatorTok{[]}\NormalTok{ cmsDigest }\OperatorTok{=}\NormalTok{ ASN1OctetString}\OperatorTok{.}\FunctionTok{getInstance}\OperatorTok{(}
\NormalTok{            digestAttr}\OperatorTok{.}\FunctionTok{getAttrValues}\OperatorTok{().}\FunctionTok{getObjectAt}\OperatorTok{(}\DecValTok{0}\OperatorTok{)).}\FunctionTok{getOctets}\OperatorTok{();}

    \DataTypeTok{boolean}\NormalTok{ hashMatch }\OperatorTok{=} \BuiltInTok{MessageDigest}\OperatorTok{.}\FunctionTok{isEqual}\OperatorTok{(}\NormalTok{computedHash}\OperatorTok{,}\NormalTok{ cmsDigest}\OperatorTok{);}
    \ControlFlowTok{if} \OperatorTok{(!}\NormalTok{hashMatch}\OperatorTok{)} \OperatorTok{\{}
\NormalTok{        log}\OperatorTok{.}\FunctionTok{error}\OperatorTok{(}\StringTok{"Hash mismatch: computed=\{\} vs CMS=\{\}"}\OperatorTok{,}
\NormalTok{                  Hex}\OperatorTok{.}\FunctionTok{toHexString}\OperatorTok{(}\NormalTok{computedHash}\OperatorTok{),}\NormalTok{ Hex}\OperatorTok{.}\FunctionTok{toHexString}\OperatorTok{(}\NormalTok{cmsDigest}\OperatorTok{));}
    \OperatorTok{\}}
    \ControlFlowTok{return}\NormalTok{ hashMatch}\OperatorTok{;}
\OperatorTok{\}}
\end{Highlighting}
\end{Shaded}

The code above verifies that the \textbf{hash} matches --- it does not
verify the \textbf{CMS signature} (the RSA/ECDSA computation over the
signed attributes). Full verification requires both:

\begin{enumerate}
\def\labelenumi{\arabic{enumi}.}
\item
  Byte range hash matches CMS \texttt{messageDigest} attribute (shown
  above)
\item
  CMS signature over \texttt{signedAttrs} verifies with the signer's
  certificate
\end{enumerate}

Chapter 31 provides complete verification code that covers both steps
plus certificate chain validation.

\section{Hex Dump Walkthrough: A Real Signed
PDF}\label{_hex_dump_walkthrough_a_real_signed_pdf}

Let's trace the byte range through an actual signed PDF, byte by byte.
This is the kind of analysis you do when debugging --- you open a hex
editor (or use \texttt{xxd}, \texttt{hexdump}, or QPDF) and look at the
raw bytes.

\subsection{Finding the Signature
Dictionary}\label{_finding_the_signature_dictionary}

First, locate the signature dictionary in the file. With QPDF:

\begin{Shaded}
\begin{Highlighting}[]
\ExtensionTok{qpdf} \AttributeTok{{-}{-}show{-}object}\OperatorTok{=}\NormalTok{15 signed{-}document.pdf}
\end{Highlighting}
\end{Shaded}

(Where object 15 is the signature dictionary --- find it by tracing from
the AcroForm → Fields → {[}signature field{]} → /V.)

Output:

\begin{verbatim}
15 0 obj
<<
  /ByteRange [0 14231 47007 1849]
  /ContactInfo ()
  /Contents <3082...0000>
  /Filter /Adobe.PPKLite
  /Location ()
  /M (D:20260308120000+01'00')
  /Name (Test Signer)
  /Reason ()
  /SubFilter /adbe.pkcs7.detached
  /Type /Sig
>>
endobj
\end{verbatim}

Now we know:

\begin{itemize}
\item
  \textbf{Segment 1}: bytes \texttt{{[}0,\ 14231)} --- 14,231 bytes from
  file start
\item
  \textbf{Gap}: bytes \texttt{{[}14231,\ 47007)} --- 32,776 bytes
  containing \texttt{/Contents}
\item
  \textbf{Segment 2}: bytes \texttt{{[}47007,\ 48856)} --- 1,849 bytes
  at file end
\item
  \textbf{Total file size}: 47,007 + 1,849 = 48,856 bytes
\end{itemize}

\subsection{Examining the Gap
Boundaries}\label{_examining_the_gap_boundaries}

Let's look at the bytes right around the gap:

\begin{Shaded}
\begin{Highlighting}[]
\ExtensionTok{xxd} \AttributeTok{{-}s}\NormalTok{ 14220 }\AttributeTok{{-}l}\NormalTok{ 30 signed{-}document.pdf}
\end{Highlighting}
\end{Shaded}

Output:

\begin{verbatim}
0000378c: 2f43 6f6e 7465 6e74 7320 3c30 3830 3230  /Contents <08020
0000379c: 3036 3032 3031 3030 3330 3832 3032 3032  060201003082020
\end{verbatim}

We can see \texttt{/Contents\ \textless{}} at offset
\textasciitilde14,223. The \texttt{\textless{}} delimiter is at the byte
just before the hex string begins. The \texttt{/ByteRange} says segment
1 ends at offset 14,231 --- which is the \texttt{\textless{}} character.
Actually, the exact offset of \texttt{\textless{}} depends on how PDFBox
serializes the dictionary.

Now the end of the gap:

\begin{Shaded}
\begin{Highlighting}[]
\ExtensionTok{xxd} \AttributeTok{{-}s}\NormalTok{ 47000 }\AttributeTok{{-}l}\NormalTok{ 20 signed{-}document.pdf}
\end{Highlighting}
\end{Shaded}

Output:

\begin{verbatim}
0000b798: 3030 3030 3030 3030 3e0a 2f46 696c 7465  00000000>./Filte
\end{verbatim}

The \texttt{\textgreater{}} delimiter is at byte 47,007 (approximately),
followed by a newline and the next dictionary key. Segment 2 starts
immediately after the \texttt{\textgreater{}}.

\subsection{Verifying No Trailing
Bytes}\label{_verifying_no_trailing_bytes}

\begin{Shaded}
\begin{Highlighting}[]
\FunctionTok{wc} \AttributeTok{{-}c}\NormalTok{ signed{-}document.pdf}
\CommentTok{\# Output: 48856 signed{-}document.pdf}
\end{Highlighting}
\end{Shaded}

Byte range says: offset2 + length2 = 47,007 + 1,849 = 48,856. File size
= 48,856. They match --- no trailing unsigned bytes. The signature
covers the entire file (minus the gap).

\section{Common Byte Range Problems}\label{_common_byte_range_problems}

\subsection{Problem 1: Trailing Bytes After Signed
Region}\label{_problem_1_trailing_bytes_after_signed_region}

\textbf{Symptom}: Signature validates cryptographically, but the file
contains content not covered by the signature.

\textbf{Cause}: Something was appended to the file after signing ---
possibly a legitimate incremental update (second signature, form fill),
or possibly an Incremental Saving Attack.

\textbf{Detection}:

\begin{Shaded}
\begin{Highlighting}[]
\DataTypeTok{int}\NormalTok{ signedEnd }\OperatorTok{=}\NormalTok{ byteRange}\OperatorTok{[}\DecValTok{2}\OperatorTok{]} \OperatorTok{+}\NormalTok{ byteRange}\OperatorTok{[}\DecValTok{3}\OperatorTok{];}
\DataTypeTok{int}\NormalTok{ fileSize }\OperatorTok{=}\NormalTok{ pdfBytes}\OperatorTok{.}\FunctionTok{length}\OperatorTok{;}
\ControlFlowTok{if} \OperatorTok{(}\NormalTok{signedEnd }\OperatorTok{\textless{}}\NormalTok{ fileSize}\OperatorTok{)} \OperatorTok{\{}
    \DataTypeTok{int}\NormalTok{ trailingBytes }\OperatorTok{=}\NormalTok{ fileSize }\OperatorTok{{-}}\NormalTok{ signedEnd}\OperatorTok{;}
\NormalTok{    log}\OperatorTok{.}\FunctionTok{warn}\OperatorTok{(}\StringTok{"\{\} bytes of unsigned data after signed region"}\OperatorTok{,}\NormalTok{ trailingBytes}\OperatorTok{);}
    \CommentTok{// Analyze the trailing data to determine if it\textquotesingle{}s a legitimate revision}
\OperatorTok{\}}
\end{Highlighting}
\end{Shaded}

\textbf{Resolution}: Parse the trailing data. If it contains a valid
cross-reference section and trailer pointing to a subsequent signature,
it's likely a legitimate multi-signature document. If it contains
modified objects (changed page content, replaced fonts, new pages), it
may be an attack. Chapter 24 provides detailed analysis code.

\subsection{\texorpdfstring{Problem 2: \texttt{/Contents} Size
Exceeded}{Problem 2: /Contents Size Exceeded}}\label{_problem_2_contents_size_exceeded}

\textbf{Symptom}:
\texttt{IOException:\ Can’t\ write\ signature,\ not\ enough\ space}
during signing.

\textbf{Cause}: The CMS \texttt{SignedData} (after hex encoding) exceeds
the reserved \texttt{/Contents} space.

\textbf{Fix}: Increase
\texttt{SignatureOptions.setPreferredSignatureSize()}. See the size
estimation section above for recommended values by PAdES profile.

\subsection{Problem 3: ByteRange Doesn't Start at
Zero}\label{_problem_3_byterange_doesnt_start_at_zero}

\textbf{Symptom}: Validator reports "invalid byte range" or "signature
does not cover document header."

\textbf{Cause}: Extremely rare --- would indicate a buggy signing
library. The spec requires that the signed region start at byte 0 to
cover the PDF header.

\textbf{Resolution}: This is unfixable after the fact --- the PDF must
be re-signed.

\subsection{\texorpdfstring{Problem 4: Gap Contains More Than Just
\texttt{/Contents}}{Problem 4: Gap Contains More Than Just /Contents}}\label{_problem_4_gap_contains_more_than_just_contents}

\textbf{Symptom}: Validator reports "unexpected data in signature gap"
or hash mismatch.

\textbf{Cause}: The bytes between the two signed segments should contain
\emph{only} the \texttt{/Contents} hex string. If any other data exists
in this gap --- for example, if a tool inserted extra whitespace or a
stray dictionary key --- the byte range geometry is broken.

\textbf{Detection}: Parse the gap region and verify it matches the
pattern \texttt{\textless{}{[}0-9A-Fa-f{]}+\textgreater{}} with optional
zero padding.

\subsection{Problem 5: Hash Algorithm
Mismatch}\label{_problem_5_hash_algorithm_mismatch}

\textbf{Symptom}: Computed hash doesn't match CMS \texttt{messageDigest}
attribute, even though the byte range looks correct.

\textbf{Cause}: You're hashing with a different algorithm than the one
specified in the CMS. The CMS \texttt{SignerInfo.digestAlgorithm} tells
you which algorithm was used --- it might be SHA-256, SHA-384, SHA-512,
or (for legacy documents) SHA-1.

\textbf{Fix}: Always read the digest algorithm from the CMS before
computing the hash:

\begin{Shaded}
\begin{Highlighting}[]
\NormalTok{SignerInformation signer }\OperatorTok{=}\NormalTok{ cmsData}\OperatorTok{.}\FunctionTok{getSignerInfos}\OperatorTok{().}\FunctionTok{getSigners}\OperatorTok{().}\FunctionTok{iterator}\OperatorTok{().}\FunctionTok{next}\OperatorTok{();}
\BuiltInTok{String}\NormalTok{ digestAlgOid }\OperatorTok{=}\NormalTok{ signer}\OperatorTok{.}\FunctionTok{getDigestAlgOID}\OperatorTok{();}
\CommentTok{// Map OID to algorithm name}
\BuiltInTok{String}\NormalTok{ algName }\OperatorTok{=} \ControlFlowTok{switch} \OperatorTok{(}\NormalTok{digestAlgOid}\OperatorTok{)} \OperatorTok{\{}
    \ControlFlowTok{case} \StringTok{"2.16.840.1.101.3.4.2.1"} \OperatorTok{{-}\textgreater{}} \StringTok{"SHA{-}256"}\OperatorTok{;}
    \ControlFlowTok{case} \StringTok{"2.16.840.1.101.3.4.2.2"} \OperatorTok{{-}\textgreater{}} \StringTok{"SHA{-}384"}\OperatorTok{;}
    \ControlFlowTok{case} \StringTok{"2.16.840.1.101.3.4.2.3"} \OperatorTok{{-}\textgreater{}} \StringTok{"SHA{-}512"}\OperatorTok{;}
    \ControlFlowTok{case} \StringTok{"1.3.14.3.2.26"}          \OperatorTok{{-}\textgreater{}} \StringTok{"SHA{-}1"}\OperatorTok{;}   \CommentTok{// deprecated}
    \KeywordTok{default} \OperatorTok{{-}\textgreater{}} \ControlFlowTok{throw} \KeywordTok{new} \BuiltInTok{IllegalArgumentException}\OperatorTok{(}\StringTok{"Unknown digest OID: "} \OperatorTok{+}\NormalTok{ digestAlgOid}\OperatorTok{);}
\OperatorTok{\};}
\BuiltInTok{MessageDigest}\NormalTok{ digest }\OperatorTok{=} \BuiltInTok{MessageDigest}\OperatorTok{.}\FunctionTok{getInstance}\OperatorTok{(}\NormalTok{algName}\OperatorTok{);}
\end{Highlighting}
\end{Shaded}

\section{Key Takeaways}\label{_key_takeaways}

\begin{itemize}
\item
  The byte range divides a signed PDF into exactly \textbf{two signed
  segments} with \textbf{one unsigned gap} between them
\item
  The gap contains only the hex-encoded CMS signature value (the
  \texttt{/Contents} string) --- nothing else
\item
  \texttt{offset2\ +\ length2} must equal the file size; any trailing
  bytes are unsigned and potentially malicious
\item
  Each signature in a multi-signature document covers all previous
  signatures, creating a nested coverage model
\item
  Always reserve generous space for \texttt{/Contents} --- 64 KB is a
  safe default for most PAdES profiles
\item
  The \texttt{/ByteRange} array itself is signed --- it cannot be
  modified after signing without detection
\item
  Byte range verification is the \textbf{first step} of signature
  validation --- if it fails, there is no need to check the CMS
\item
  The Incremental Saving Attack (ISA) exploits the trailing bytes
  problem; the trailing bytes check is your primary defense
\end{itemize}

\chapter{Approval vs. Certification
Signatures}\label{_approval_vs_certification_signatures}

PDF digital signatures come in two fundamentally different flavors, and
confusing them is one of the most common mistakes in signature
implementation. An \textbf{approval signature} says "I have reviewed and
agree with this document." A \textbf{certification signature} says "I am
the author of this document, and I am defining what changes are allowed
after I sign it."

The difference is not just semantic --- it is structural. A
certification signature carries a \texttt{/DocMDP} (Modification
Detection and Prevention) transform that defines permission levels,
restricts what subsequent signers and form fillers can do, and triggers
different validation behavior in viewers. Getting this wrong means your
signatures either restrict too much (form fields that can't be filled
after signing) or too little (no protection against unauthorized
modifications).

This chapter explains both signature types, the DocMDP permission model,
Lock and FieldMDP dictionaries, and the interaction rules that govern
what happens when multiple signatures coexist.

\section{The Two Signature Types}\label{_the_two_signature_types}

\subsection{Approval Signatures}\label{_approval_signatures}

An \textbf{approval signature} is what most people think of when they
hear "digital signature." It means: "I, the signer, have reviewed the
document and give my approval."

Technically, an approval signature is a \texttt{/Sig} dictionary with:

\begin{itemize}
\item
  \texttt{/Filter}: \texttt{/Adobe.PPKLite}
\item
  \texttt{/SubFilter}: \texttt{adbe.pkcs7.detached} or
  \texttt{ETSI.CAdES.detached}
\item
  \texttt{/ByteRange}: covers the document up to this point
\item
  No \texttt{/TransformMethod} in any associated signature reference
  dictionary
\end{itemize}

Approval signatures have these properties:

\begin{itemize}
\item
  \textbf{Multiple allowed}: A document can have any number of approval
  signatures. Each one adds an incremental update.
\item
  \textbf{No permission restrictions}: An approval signature does not
  limit what subsequent users can do with the document. After you apply
  an approval signature, someone else can add another signature, fill
  form fields, add annotations, or even modify the document (though
  modifications will cause the signature to show as "modified since
  signing").
\item
  \textbf{Order matters}: Each approval signature covers everything up
  to the point it was applied. Signature 2 covers signature 1. If the
  document is modified between signature 1 and signature 2, signature 1
  will show as "document modified" while signature 2 will show as
  "valid" (because it was applied after the modification).
\end{itemize}

\subsection{Certification Signatures}\label{_certification_signatures}

A \textbf{certification signature} (also called an "author signature" or
"MDP signature") is structurally different. It carries a
\texttt{/TransformMethod} of \texttt{/DocMDP} in its signature reference
dictionary, which embeds a set of permission rules into the signature.

Technically, a certification signature includes:

\begin{itemize}
\item
  Everything an approval signature has
\item
  A \texttt{/Reference} array in the signature dictionary containing a
  \texttt{/SigRef} (Signature Reference) dictionary
\item
  The \texttt{/SigRef} dictionary specifies
  \texttt{/TransformMethod\ /DocMDP}
\item
  A \texttt{/TransformParams} dictionary with a \texttt{/P} key (the
  permission level)
\end{itemize}

And at the document level:

\begin{itemize}
\item
  The document catalog's \texttt{/Perms} dictionary includes a
  \texttt{/DocMDP} entry pointing to the certifying signature
\end{itemize}

Certification signatures have these properties:

\begin{itemize}
\item
  \textbf{Only one per document}: A document can have at most one
  certification signature. It is always the first signature applied. If
  you try to add a certification signature to a document that already
  has one, the behavior is undefined --- some viewers will reject it,
  others will ignore the previous one.
\item
  \textbf{Defines permissions}: The \texttt{/P} value (1, 2, or 3)
  controls what changes are allowed after certification. This is the
  defining feature that makes certification signatures different from
  approval signatures.
\item
  \textbf{First in the chain}: The certification signature is always the
  foundational signature. All subsequent signatures are approval
  signatures that build on the certified baseline.
\end{itemize}

\begin{verbatim}
<< /Type /Sig
   /Filter /Adobe.PPKLite
   /SubFilter /adbe.pkcs7.detached
   /ByteRange [0 14231 47007 1849]
   /Contents <3082...>
   /Reference [
     << /Type /SigRef
        /TransformMethod /DocMDP
        /TransformParams << /Type /TransformParams /P 2 /V /1.2 >>
     >>
   ]
   /Name (Document Author)
   /M (D:20260308120000+01'00')
>>
\end{verbatim}

The document catalog connects this signature to the document-level
permission system:

\begin{verbatim}
<< /Type /Catalog
   /AcroForm << ... >>
   /Perms << /DocMDP 15 0 R >>    % Reference to the certifying signature object
   /Pages 3 0 R
>>
\end{verbatim}

\section{The DocMDP Permission
Model}\label{_the_docmdp_permission_model}

\subsection{Permission Level 1: No Changes
Allowed}\label{_permission_level_1_no_changes_allowed}

\texttt{/P\ 1} is the strictest level. After certification with P=1,
\textbf{no changes whatsoever are allowed} to the document. This means:

\begin{itemize}
\item
  No form fields can be filled
\item
  No additional signatures can be added
\item
  No annotations can be created
\item
  No pages can be added or removed
\item
  Any incremental update invalidates the certification
\end{itemize}

This level is appropriate for documents that are complete and final ---
archived records, published regulations, completed certificates. In
practice, it is rarely used because most signing workflows require at
least form filling or countersigning.

P=1 is so restrictive that it prevents the document from being
countersigned. If your workflow requires multiple signatures (author
signs, then reviewer countersigns), \textbf{do not use P=1}. Use P=2
instead --- it allows additional signatures while preventing content
modifications.

\subsection{Permission Level 2: Form Filling and Signing
Only}\label{_permission_level_2_form_filling_and_signing_only}

\texttt{/P\ 2} is the \textbf{default and most common} permission level.
After certification with P=2:

\begin{itemize}
\item
  \textbf{Allowed}: Filling form fields, applying additional (approval)
  signatures, adding signature appearances
\item
  \textbf{Not allowed}: Modifying page content, adding pages, removing
  pages, changing form field definitions (adding/removing fields),
  modifying JavaScript actions
\item
  Any change beyond form filling and signing invalidates the
  certification
\end{itemize}

P=2 is the standard for government forms, contracts, and regulated
documents. The document author sets up the form (fields, layout,
instructions), certifies it, and then the form recipient fills in the
fields and applies their own approval signature --- without being able
to alter the form structure or content.

This is the level you should use unless you have a specific reason to
choose P=1 or P=3.

\subsection{Permission Level 3: Form Filling, Signing, and
Annotations}\label{_permission_level_3_form_filling_signing_and_annotations}

\texttt{/P\ 3} is the most permissive certification level. After
certification with P=3:

\begin{itemize}
\item
  \textbf{Allowed}: Everything P=2 allows, plus adding annotations
  (comments, highlights, stamps, sticky notes, text markup)
\item
  \textbf{Not allowed}: Modifying page content, adding or removing
  pages, changing form field definitions
\end{itemize}

P=3 is appropriate for review workflows where multiple people need to
comment on a certified document without altering its content. Think of a
contract that needs review comments from legal, compliance, and
management before final signatures.

\subsection{What "Invalidates the Certification" Actually
Means}\label{_what_invalidates_the_certification_actually_means}

When we say a change "invalidates the certification," we mean:

\begin{itemize}
\item
  The \textbf{byte range} of the certification signature remains intact
  --- the hash still matches
\item
  The \textbf{CMS signature} still verifies cryptographically
\item
  But a \textbf{conforming viewer} detects that the modifications in
  subsequent revisions exceed the allowed permissions and displays a
  warning
\end{itemize}

In Adobe Acrobat, this shows as: "The document has been altered or
corrupted since the Signature was applied." In LibreOffice, it shows as:
"Document has been modified since it was last signed."

The exact wording varies by viewer, and --- this is important ---
\textbf{not all viewers enforce MDP permissions correctly}. Some viewers
only check whether the byte range is intact and don't analyze the
incremental updates for permission compliance. This is one of the
weaknesses that the Bochum team exploited in the Incremental Saving
Attack (Chapter 24).

\section{Lock Dictionaries: Field-Level
Protection}\label{_lock_dictionaries_field_level_protection}

\subsection{What Lock Dictionaries Do}\label{_what_lock_dictionaries_do}

Beyond document-level MDP permissions, PDF supports \textbf{field-level
locking} via the \texttt{/Lock} dictionary on a signature field. A Lock
dictionary specifies which form fields should be locked (made read-only)
once this particular signature is applied.

The Lock dictionary lives in the signature field dictionary (not the
signature value dictionary), and it applies when the field is signed:

\begin{verbatim}
<< /Type /Annot
   /Subtype /Widget
   /FT /Sig
   /T (SignatureField1)
   /V 15 0 R               % Reference to the /Sig dictionary
   /Lock <<
     /Type /SigFieldLock
     /Action /Include       % or /Exclude or /All
     /Fields [(FieldA) (FieldB)]
   >>
>>
\end{verbatim}

\subsection{Lock Actions}\label{_lock_actions}

\begin{longtable}[]{@{}
  >{\raggedright\arraybackslash}p{(\linewidth - 6\tabcolsep) * \real{0.2222}}
  >{\raggedright\arraybackslash}p{(\linewidth - 6\tabcolsep) * \real{0.1111}}
  >{\raggedright\arraybackslash}p{(\linewidth - 6\tabcolsep) * \real{0.1111}}
  >{\raggedright\arraybackslash}p{(\linewidth - 6\tabcolsep) * \real{0.5556}}@{}}
\caption{Signature Field Lock Actions (ISO 32000-2:2020
§12.8.2.4)}\tabularnewline
\toprule\noalign{}
\begin{minipage}[b]{\linewidth}\raggedright
Action
\end{minipage} & \begin{minipage}[b]{\linewidth}\raggedright
Type
\end{minipage} & \begin{minipage}[b]{\linewidth}\raggedright
Req
\end{minipage} & \begin{minipage}[b]{\linewidth}\raggedright
Description
\end{minipage} \\
\midrule\noalign{}
\endfirsthead
\toprule\noalign{}
\begin{minipage}[b]{\linewidth}\raggedright
Action
\end{minipage} & \begin{minipage}[b]{\linewidth}\raggedright
Type
\end{minipage} & \begin{minipage}[b]{\linewidth}\raggedright
Req
\end{minipage} & \begin{minipage}[b]{\linewidth}\raggedright
Description
\end{minipage} \\
\midrule\noalign{}
\endhead
\bottomrule\noalign{}
\endlastfoot
\texttt{/All} & name & --- & Lock all fields in the document. No
\texttt{/Fields} array is needed. After this signature is applied, every
form field in the document becomes read-only. This is the simplest but
most restrictive lock. \\
\texttt{/Include} & name & --- & Lock only the fields listed in the
\texttt{/Fields} array. The listed fields become read-only; all other
fields remain fillable. Use this when you want to protect specific
fields (e.g., "Price" and "Total") while allowing others (e.g.,
"Customer Signature") to be filled. \\
\texttt{/Exclude} & name & --- & Lock all fields \textbf{except} those
listed in the \texttt{/Fields} array. The listed fields remain fillable;
everything else is locked. Use this when most fields should be locked
but a few specific fields must remain editable (e.g., a "Date Received"
field that an administrator fills after the signer has signed). \\
\end{longtable}

\subsection{\texorpdfstring{The \texttt{/Fields}
Array}{The /Fields Array}}\label{_the_fields_array}

The \texttt{/Fields} array in the Lock dictionary contains the
\textbf{fully qualified field names} of the fields to include or
exclude. In a hierarchical field tree, the fully qualified name is the
concatenation of parent and child names separated by periods:

\begin{verbatim}
/Fields [(PersonalInfo.FirstName) (PersonalInfo.LastName) (Signature2)]
\end{verbatim}

Field name matching is case-sensitive and must be exact. A typo in the
Lock dictionary's \texttt{/Fields} array means the field won't be locked
--- and there will be no error message. The signature applies
successfully, the Lock dictionary is present, but the misspelled field
remains editable.

\textbf{Always test your Lock configurations} by:

\begin{enumerate}
\def\labelenumi{\arabic{enumi}.}
\item
  Signing the document with the Lock dictionary
\item
  Opening it in Adobe Acrobat
\item
  Attempting to edit the fields that should be locked
\item
  Verifying that Acrobat prevents editing and shows the fields as
  read-only
\end{enumerate}

\section{FieldMDP: Fine-Grained Modification
Detection}\label{_fieldmdp_fine_grained_modification_detection}

\subsection{FieldMDP vs. DocMDP}\label{_fieldmdp_vs_docmdp}

While DocMDP controls what categories of changes are allowed (form
filling, annotations, etc.), \textbf{FieldMDP} controls which specific
fields can be modified. It is a more targeted mechanism.

FieldMDP uses the same \texttt{/Reference} array mechanism as DocMDP but
with \texttt{/TransformMethod\ /FieldMDP}:

\begin{verbatim}
<< /Type /SigRef
   /TransformMethod /FieldMDP
   /TransformParams <<
     /Type /TransformParams
     /Action /Include
     /Fields [(ProtectedField1) (ProtectedField2)]
     /V /2.0
   >>
>>
\end{verbatim}

The \texttt{/Action} and \texttt{/Fields} keys work identically to Lock
dictionary actions --- \texttt{/Include}, \texttt{/Exclude}, and
\texttt{/All} with the same semantics.

The difference from Lock dictionaries:

\begin{itemize}
\item
  \textbf{Lock dictionary}: Controls whether a field is
  \textbf{editable} in the viewer UI (read-only vs. writable). It's a
  presentation-level restriction --- the viewer \emph{should} prevent
  the user from typing in a locked field.
\item
  \textbf{FieldMDP}: Controls whether modifications to specific fields
  are \textbf{detected} during validation. If a FieldMDP-protected field
  is changed after signing, the signature shows as modified. It's a
  validation-level mechanism.
\end{itemize}

In practice, Lock and FieldMDP are often used together: Lock prevents
casual editing in a UI, and FieldMDP ensures that any bypass of the Lock
(via direct file manipulation) is detected during validation.

\section{Interaction Rules: When Multiple Signatures
Coexist}\label{_interaction_rules_when_multiple_signatures_coexist}

\subsection{The Certification + Approval
Pattern}\label{_the_certification_approval_pattern}

The most common multi-signature pattern is:

\begin{enumerate}
\def\labelenumi{\arabic{enumi}.}
\item
  \textbf{Author creates} the document with form fields
\item
  \textbf{Author certifies} with P=2 (form filling and signing allowed)
  and optionally locks specific fields
\item
  \textbf{Recipient fills} form fields
\item
  \textbf{Recipient applies} an approval signature (which adds an
  incremental update with the filled values and the new signature)
\end{enumerate}

This works because:

\begin{itemize}
\item
  The certification (P=2) allows form filling and additional signatures
\item
  The approval signature is applied to a document that already has a
  certification
\item
  The certification signature remains valid because the changes (form
  fill + new signature) are within the P=2 permissions
\end{itemize}

\subsection{What Happens When Permissions
Conflict}\label{_what_happens_when_permissions_conflict}

Things get complicated when:

\textbf{Scenario 1: Approval signature before certification}

Technically, any signature can be applied to any document. If someone
applies an approval signature first and then a certification signature
second, the certification signature covers (signs over) the approval
signature. But the certification's MDP permissions only apply to changes
\emph{after} the certification --- they don't retroactively validate the
approval signature's changes.

In practice, this scenario is unusual and may confuse validators.
\textbf{Always apply the certification signature first.}

\textbf{Scenario 2: Lock dictionary conflicts with MDP}

If the certification says P=2 (form filling allowed) but a Lock
dictionary locks all fields --- the Lock takes precedence for the locked
fields. The MDP permission defines the ceiling (what's \emph{possible}),
and Lock dictionaries further restrict within that ceiling.

\textbf{Scenario 3: Multiple Lock dictionaries from different
signatures}

Each signature field can have its own Lock dictionary. When multiple
signatures with Lock dictionaries coexist:

\begin{itemize}
\item
  Locks are cumulative --- each new signature can lock additional fields
\item
  Locks are permanent --- once a field is locked by any signature, no
  subsequent signature can unlock it
\item
  A field is locked if \emph{any} Lock dictionary in the chain locks it
\end{itemize}

This means the first signer can lock fields A and B, the second signer
can lock fields C and D, and after both signatures are applied, fields
A, B, C, and D are all locked.

\begin{verbatim}
// Signature 1 (Author): Lock fields A, B
SignatureField1.Lock = { Action: Include, Fields: [A, B] }

// After Sig1: A=locked, B=locked, C=editable, D=editable

// Signature 2 (Reviewer): Lock fields C, D
SignatureField2.Lock = { Action: Include, Fields: [C, D] }

// After Sig2: A=locked, B=locked, C=locked, D=locked
\end{verbatim}

\section{Implementation: Certification Signatures with
PDFBox}\label{_implementation_certification_signatures_with_pdfbox}

\subsection{Creating a Certification
Signature}\label{_creating_a_certification_signature}

\begin{Shaded}
\begin{Highlighting}[]
\NormalTok{PDDocument doc }\OperatorTok{=}\NormalTok{ Loader}\OperatorTok{.}\FunctionTok{loadPDF}\OperatorTok{(}\NormalTok{inputFile}\OperatorTok{);}

\CommentTok{// Create the signature}
\NormalTok{PDSignature signature }\OperatorTok{=} \KeywordTok{new} \FunctionTok{PDSignature}\OperatorTok{();}
\NormalTok{signature}\OperatorTok{.}\FunctionTok{setFilter}\OperatorTok{(}\NormalTok{PDSignature}\OperatorTok{.}\FunctionTok{FILTER\_ADOBE\_PPKLITE}\OperatorTok{);}
\NormalTok{signature}\OperatorTok{.}\FunctionTok{setSubFilter}\OperatorTok{(}\NormalTok{PDSignature}\OperatorTok{.}\FunctionTok{SUBFILTER\_ADBE\_PKCS7\_DETACHED}\OperatorTok{);}
\NormalTok{signature}\OperatorTok{.}\FunctionTok{setName}\OperatorTok{(}\StringTok{"Document Author"}\OperatorTok{);}
\NormalTok{signature}\OperatorTok{.}\FunctionTok{setSignDate}\OperatorTok{(}\BuiltInTok{Calendar}\OperatorTok{.}\FunctionTok{getInstance}\OperatorTok{());}

\CommentTok{// Set certification (MDP) permissions — P=2 is recommended}
\NormalTok{doc}\OperatorTok{.}\FunctionTok{addSignature}\OperatorTok{(}\NormalTok{signature}\OperatorTok{,}\NormalTok{ signingInterface}\OperatorTok{);}

\CommentTok{// Set the access permissions}
\CommentTok{// 2 = form filling and signing allowed}
\NormalTok{AccessPermission accessPermission }\OperatorTok{=} \KeywordTok{new} \FunctionTok{AccessPermission}\OperatorTok{();}
\NormalTok{accessPermission}\OperatorTok{.}\FunctionTok{setDocMDP}\OperatorTok{(}\DecValTok{2}\OperatorTok{);}
\NormalTok{doc}\OperatorTok{.}\FunctionTok{setDocumentPermissions}\OperatorTok{(}\NormalTok{signature}\OperatorTok{,}\NormalTok{ accessPermission}\OperatorTok{);}

\NormalTok{doc}\OperatorTok{.}\FunctionTok{saveIncremental}\OperatorTok{(}\NormalTok{outputStream}\OperatorTok{);}
\NormalTok{doc}\OperatorTok{.}\FunctionTok{close}\OperatorTok{();}
\end{Highlighting}
\end{Shaded}

PDFBox's API for certification signatures evolved between versions. In
PDFBox 2.x, you used \texttt{doc.addSignature()} with a
\texttt{SignatureOptions} and separately configured the MDP via the
catalog's \texttt{/Perms} dictionary. In PDFBox 3.x, the
\texttt{AccessPermission} class provides a cleaner API, but the
underlying mechanism is the same --- writing the \texttt{/DocMDP}
reference into the catalog.

\subsection{Adding Lock Dictionaries}\label{_adding_lock_dictionaries}

Lock dictionaries are configured on the signature field before signing:

\begin{Shaded}
\begin{Highlighting}[]
\CommentTok{// Get or create the signature field}
\NormalTok{PDSignatureField sigField }\OperatorTok{=} \FunctionTok{findOrCreateSignatureField}\OperatorTok{(}\NormalTok{doc}\OperatorTok{,} \StringTok{"AuthorSignature"}\OperatorTok{);}

\CommentTok{// Create lock dictionary}
\NormalTok{COSDictionary lockDict }\OperatorTok{=} \KeywordTok{new} \FunctionTok{COSDictionary}\OperatorTok{();}
\NormalTok{lockDict}\OperatorTok{.}\FunctionTok{setName}\OperatorTok{(}\NormalTok{COSName}\OperatorTok{.}\FunctionTok{TYPE}\OperatorTok{,} \StringTok{"SigFieldLock"}\OperatorTok{);}
\NormalTok{lockDict}\OperatorTok{.}\FunctionTok{setName}\OperatorTok{(}\NormalTok{COSName}\OperatorTok{.}\FunctionTok{getPDFName}\OperatorTok{(}\StringTok{"Action"}\OperatorTok{),} \StringTok{"Include"}\OperatorTok{);}

\NormalTok{COSArray fieldsToLock }\OperatorTok{=} \KeywordTok{new} \FunctionTok{COSArray}\OperatorTok{();}
\NormalTok{fieldsToLock}\OperatorTok{.}\FunctionTok{add}\OperatorTok{(}\KeywordTok{new} \FunctionTok{COSString}\OperatorTok{(}\StringTok{"PriceField"}\OperatorTok{));}
\NormalTok{fieldsToLock}\OperatorTok{.}\FunctionTok{add}\OperatorTok{(}\KeywordTok{new} \FunctionTok{COSString}\OperatorTok{(}\StringTok{"TotalField"}\OperatorTok{));}
\NormalTok{fieldsToLock}\OperatorTok{.}\FunctionTok{add}\OperatorTok{(}\KeywordTok{new} \FunctionTok{COSString}\OperatorTok{(}\StringTok{"ContractTerms"}\OperatorTok{));}
\NormalTok{lockDict}\OperatorTok{.}\FunctionTok{setItem}\OperatorTok{(}\NormalTok{COSName}\OperatorTok{.}\FunctionTok{getPDFName}\OperatorTok{(}\StringTok{"Fields"}\OperatorTok{),}\NormalTok{ fieldsToLock}\OperatorTok{);}

\CommentTok{// Attach lock dictionary to the signature field}
\NormalTok{sigField}\OperatorTok{.}\FunctionTok{getCOSObject}\OperatorTok{().}\FunctionTok{setItem}\OperatorTok{(}\NormalTok{COSName}\OperatorTok{.}\FunctionTok{getPDFName}\OperatorTok{(}\StringTok{"Lock"}\OperatorTok{),}\NormalTok{ lockDict}\OperatorTok{);}

\CommentTok{// Now sign the document}
\NormalTok{doc}\OperatorTok{.}\FunctionTok{addSignature}\OperatorTok{(}\NormalTok{signature}\OperatorTok{,}\NormalTok{ signingInterface}\OperatorTok{);}
\NormalTok{doc}\OperatorTok{.}\FunctionTok{saveIncremental}\OperatorTok{(}\NormalTok{outputStream}\OperatorTok{);}
\end{Highlighting}
\end{Shaded}

\section{Seed Value Dictionaries: Constraining How a Signature Field Is
Signed}\label{_seed_value_dictionaries_constraining_how_a_signature_field_is_signed}

\subsection{Purpose of Seed Values}\label{_purpose_of_seed_values}

A \textbf{seed value dictionary} (\texttt{/SV}) on a signature field
constrains \emph{how} the field can be signed. It does not sign anything
itself --- it provides instructions to the signing tool about what
certificate to use, what algorithm to use, which SubFilter is required,
and other signing parameters.

This is the mechanism that government form creators use to ensure that
form fillers sign with the correct certificate type, use the correct
algorithm, and produce signatures that will pass validation.

\subsection{Key Seed Value Entries}\label{_key_seed_value_entries}

\begin{longtable}[]{@{}
  >{\raggedright\arraybackslash}p{(\linewidth - 6\tabcolsep) * \real{0.2857}}
  >{\raggedright\arraybackslash}p{(\linewidth - 6\tabcolsep) * \real{0.2857}}
  >{\raggedright\arraybackslash}p{(\linewidth - 6\tabcolsep) * \real{0.1429}}
  >{\raggedright\arraybackslash}p{(\linewidth - 6\tabcolsep) * \real{0.2857}}@{}}
\caption{Seed Value Dictionary Keys (ISO 32000-2:2020
§12.8.2.3)}\tabularnewline
\toprule\noalign{}
\begin{minipage}[b]{\linewidth}\raggedright
Key
\end{minipage} & \begin{minipage}[b]{\linewidth}\raggedright
Type
\end{minipage} & \begin{minipage}[b]{\linewidth}\raggedright
Req
\end{minipage} & \begin{minipage}[b]{\linewidth}\raggedright
Spec Ref
\end{minipage} \\
\midrule\noalign{}
\endfirsthead
\toprule\noalign{}
\begin{minipage}[b]{\linewidth}\raggedright
Key
\end{minipage} & \begin{minipage}[b]{\linewidth}\raggedright
Type
\end{minipage} & \begin{minipage}[b]{\linewidth}\raggedright
Req
\end{minipage} & \begin{minipage}[b]{\linewidth}\raggedright
Spec Ref
\end{minipage} \\
\midrule\noalign{}
\endhead
\bottomrule\noalign{}
\endlastfoot
\texttt{/Filter} & name & Optional & §12.8.2.3, Table 237 \\
\multicolumn{4}{@{}>{\raggedright\arraybackslash}p{(\linewidth - 6\tabcolsep) * \real{1.0000} + 6\tabcolsep}@{}}{%
The signature handler to use. In practice, always
\texttt{/Adobe.PPKLite}. If set, the signing tool should use this
handler --- if a different handler is selected, the signing tool should
warn the user.} \\
\texttt{/SubFilter} & array & Optional & §12.8.2.3, Table 237 \\
\multicolumn{4}{@{}>{\raggedright\arraybackslash}p{(\linewidth - 6\tabcolsep) * \real{1.0000} + 6\tabcolsep}@{}}{%
Array of acceptable SubFilter values. The signing tool should use one of
these SubFilter values. Example:
\texttt{{[}/adbe.pkcs7.detached\ /ETSI.CAdES.detached{]}} means either
detached format is acceptable.} \\
\texttt{/DigestMethod} & array & Optional & §12.8.2.3, Table 237 \\
\multicolumn{4}{@{}>{\raggedright\arraybackslash}p{(\linewidth - 6\tabcolsep) * \real{1.0000} + 6\tabcolsep}@{}}{%
Array of acceptable digest algorithms. Example:
\texttt{{[}/SHA256\ /SHA384\ /SHA512{]}}. If set, the signing tool must
use one of these algorithms. \texttt{/SHA1} is deprecated and should not
be included.} \\
\texttt{/Reasons} & array & Optional & §12.8.2.3, Table 237 \\
\multicolumn{4}{@{}>{\raggedright\arraybackslash}p{(\linewidth - 6\tabcolsep) * \real{1.0000} + 6\tabcolsep}@{}}{%
Array of text strings presenting acceptable signing reasons. The signing
tool should let the user choose from this list. Example:
\texttt{{[}(I\ approve\ this\ document)\ (I\ have\ reviewed\ this\ document)\ (I\ am\ the\ authorized\ signatory){]}}.} \\
\texttt{/MDP} & dictionary & Optional & §12.8.2.3, Table 237 \\
\multicolumn{4}{@{}>{\raggedright\arraybackslash}p{(\linewidth - 6\tabcolsep) * \real{1.0000} + 6\tabcolsep}@{}}{%
The MDP parameters dictionary, specifying the default DocMDP permission
level for signatures applied to this field.} \\
\texttt{/TimeStamp} & dictionary & Optional & §12.8.2.3, Table 237 \\
\multicolumn{4}{@{}>{\raggedright\arraybackslash}p{(\linewidth - 6\tabcolsep) * \real{1.0000} + 6\tabcolsep}@{}}{%
Timestamp seed value dictionary specifying a timestamp authority URL and
whether a timestamp is required. If the \texttt{/Ff} flag is set to 1
(required), the signing tool must include a timestamp from the specified
TSA.} \\
\texttt{/Ff} & integer (flags) & Optional & §12.8.2.3, Table 237 \\
\multicolumn{4}{@{}>{\raggedright\arraybackslash}p{(\linewidth - 6\tabcolsep) * \real{1.0000} + 6\tabcolsep}@{}}{%
Flags indicating which seed values are mandatory (bit 1 set) vs.
recommended (bit 1 clear). If a seed value's corresponding Ff bit is
set, the signing tool must respect the constraint --- failure to do so
should prevent the signature from being applied.} \\
\texttt{/LegalAttestation} & array & Optional & §12.8.2.3 \\
\multicolumn{4}{@{}>{\raggedright\arraybackslash}p{(\linewidth - 6\tabcolsep) * \real{1.0000} + 6\tabcolsep}@{}}{%
Array of legal attestation strings. The signer may be required to accept
one of these attestations as part of the signing ceremony.} \\
\end{longtable}

\subsection{Timestamp Seed Values}\label{_timestamp_seed_values}

The \texttt{/TimeStamp} sub-dictionary in seed values is particularly
important for PAdES compliance:

\begin{verbatim}
/SV <<
  /TimeStamp <<
    /URL (http://tsa.example.com/timestamp)
    /Ff 1              % 1 = timestamp required, 0 = optional
  >>
  /SubFilter [/ETSI.CAdES.detached]
  /DigestMethod [/SHA256 /SHA384 /SHA512]
  /Ff 1                % All seed values are mandatory
>>
\end{verbatim}

This seed value dictionary says: "When signing this field, you must use
CAdES detached format, SHA-256 or better, and you must include a
timestamp from \texttt{tsa.example.com}. All of these requirements are
mandatory (Ff=1)."

A conforming signing tool (Adobe Acrobat, for example) will enforce
these constraints automatically. PDFBox-based signing code must read and
respect these seed values manually --- the library does not enforce them
automatically.

\textbf{Always read the seed value dictionary before signing.} If you're
building a signing service that signs arbitrary PDFs, the seed values
may impose constraints that your default configuration doesn't meet.
Ignoring them produces a signature that appears valid cryptographically
but violates the form author's intent --- and may fail validation in
systems that check seed value compliance.

\begin{Shaded}
\begin{Highlighting}[]
\NormalTok{PDSignatureField sigField }\OperatorTok{=} \KeywordTok{...}\OperatorTok{;} \CommentTok{// the field you\textquotesingle{}re about to sign}
\NormalTok{COSDictionary sv }\OperatorTok{=}\NormalTok{ sigField}\OperatorTok{.}\FunctionTok{getCOSObject}\OperatorTok{()}
    \OperatorTok{.}\FunctionTok{getCOSDictionary}\OperatorTok{(}\NormalTok{COSName}\OperatorTok{.}\FunctionTok{getPDFName}\OperatorTok{(}\StringTok{"SV"}\OperatorTok{));}
\ControlFlowTok{if} \OperatorTok{(}\NormalTok{sv }\OperatorTok{!=} \KeywordTok{null}\OperatorTok{)} \OperatorTok{\{}
    \CommentTok{// Check SubFilter requirements}
\NormalTok{    COSArray subFilters }\OperatorTok{=}\NormalTok{ sv}\OperatorTok{.}\FunctionTok{getCOSArray}\OperatorTok{(}\NormalTok{COSName}\OperatorTok{.}\FunctionTok{getPDFName}\OperatorTok{(}\StringTok{"SubFilter"}\OperatorTok{));}
    \CommentTok{// Check DigestMethod requirements}
\NormalTok{    COSArray digestMethods }\OperatorTok{=}\NormalTok{ sv}\OperatorTok{.}\FunctionTok{getCOSArray}\OperatorTok{(}\NormalTok{COSName}\OperatorTok{.}\FunctionTok{getPDFName}\OperatorTok{(}\StringTok{"DigestMethod"}\OperatorTok{));}
    \CommentTok{// Check Timestamp requirements}
\NormalTok{    COSDictionary ts }\OperatorTok{=}\NormalTok{ sv}\OperatorTok{.}\FunctionTok{getCOSDictionary}\OperatorTok{(}\NormalTok{COSName}\OperatorTok{.}\FunctionTok{getPDFName}\OperatorTok{(}\StringTok{"TimeStamp"}\OperatorTok{));}
    \ControlFlowTok{if} \OperatorTok{(}\NormalTok{ts }\OperatorTok{!=} \KeywordTok{null}\OperatorTok{)} \OperatorTok{\{}
        \BuiltInTok{String}\NormalTok{ tsaUrl }\OperatorTok{=}\NormalTok{ ts}\OperatorTok{.}\FunctionTok{getString}\OperatorTok{(}\NormalTok{COSName}\OperatorTok{.}\FunctionTok{getPDFName}\OperatorTok{(}\StringTok{"URL"}\OperatorTok{));}
        \DataTypeTok{int}\NormalTok{ tsFlags }\OperatorTok{=}\NormalTok{ ts}\OperatorTok{.}\FunctionTok{getInt}\OperatorTok{(}\NormalTok{COSName}\OperatorTok{.}\FunctionTok{getPDFName}\OperatorTok{(}\StringTok{"Ff"}\OperatorTok{),} \DecValTok{0}\OperatorTok{);}
        \DataTypeTok{boolean}\NormalTok{ timestampRequired }\OperatorTok{=} \OperatorTok{(}\NormalTok{tsFlags }\OperatorTok{\&} \DecValTok{1}\OperatorTok{)} \OperatorTok{!=} \DecValTok{0}\OperatorTok{;}
    \OperatorTok{\}}
\OperatorTok{\}}
\end{Highlighting}
\end{Shaded}

\section{Real-World Patterns and
Mistakes}\label{_real_world_patterns_and_mistakes}

\subsection{Mistake 1: Using Certification When Approval Is
Needed}\label{_mistake_1_using_certification_when_approval_is_needed}

\textbf{Scenario}: A developer creates a multi-signature workflow where
three people sign sequentially. They make the first signature a
certification signature with P=1 (no changes allowed).

\textbf{Result}: The second and third signatures fail or invalidate the
certification, because P=1 prohibits additional signatures.

\textbf{Fix}: Use P=2 for the first (certification) signature if
additional signatures are needed. Or use approval signatures for all
three if no MDP restrictions are needed.

\subsection{Mistake 2: Not Setting MDP When It's
Needed}\label{_mistake_2_not_setting_mdp_when_its_needed}

\textbf{Scenario}: A government form is meant to be filled and signed by
the citizen, with no other modifications allowed. The form creator signs
it as a plain approval signature without MDP.

\textbf{Result}: The form has no modification restrictions. Anyone can
modify the form structure, add fields, remove fields, or alter content
--- and the original signature only indicates "document has been
modified" without specifying what was allowed.

\textbf{Fix}: Use a certification signature with P=2 to allow form
filling while preventing structural changes.

\subsection{Mistake 3: Lock Dictionary with Wrong Field
Names}\label{_mistake_3_lock_dictionary_with_wrong_field_names}

\textbf{Scenario}: The signature is configured to lock fields "Price",
"Total", and "Terms". But the actual field names in the PDF are
"price\_field", "total\_amount", and "contract\_terms" (because the form
designer used different naming).

\textbf{Result}: The Lock dictionary is present and syntactically
correct, but no fields are actually locked --- the names don't match.
There is no error message or warning.

\textbf{Fix}: Enumerate all field names before configuring the Lock
dictionary. In PDFBox:

\begin{Shaded}
\begin{Highlighting}[]
\NormalTok{PDAcroForm acroForm }\OperatorTok{=}\NormalTok{ doc}\OperatorTok{.}\FunctionTok{getDocumentCatalog}\OperatorTok{().}\FunctionTok{getAcroForm}\OperatorTok{();}
\ControlFlowTok{for} \OperatorTok{(}\NormalTok{PDField field }\OperatorTok{:}\NormalTok{ acroForm}\OperatorTok{.}\FunctionTok{getFieldTree}\OperatorTok{())} \OperatorTok{\{}
\NormalTok{    log}\OperatorTok{.}\FunctionTok{info}\OperatorTok{(}\StringTok{"Field: \{\} (type: \{\})"}\OperatorTok{,}
\NormalTok{             field}\OperatorTok{.}\FunctionTok{getFullyQualifiedName}\OperatorTok{(),}\NormalTok{ field}\OperatorTok{.}\FunctionTok{getFieldType}\OperatorTok{());}
\OperatorTok{\}}
\end{Highlighting}
\end{Shaded}

\subsection{Mistake 4: Assuming All Viewers Enforce
MDP}\label{_mistake_4_assuming_all_viewers_enforce_mdp}

\textbf{Scenario}: A developer relies on MDP P=2 to prevent content
modifications, assuming all viewers will refuse to let users edit
non-form content after certification.

\textbf{Result}: Adobe Acrobat enforces MDP correctly. Foxit enforces it
mostly correctly. LibreOffice enforces it partially. PDFBox (when used
as a library) does not enforce MDP at all --- it allows any modification
programmatically. Browser-based PDF viewers (Chrome, Firefox) often
ignore MDP entirely because they don't support editing.

\textbf{Lesson}: MDP is a \emph{validation} mechanism, not an
\emph{access control} mechanism. It tells you whether changes were
permitted, not whether they were prevented. A determined user with a hex
editor or a programmatic tool can always modify a PDF --- MDP ensures
you can \emph{detect} that the modification occurred.

\section{Key Takeaways}\label{_key_takeaways_2}

\begin{itemize}
\item
  \textbf{Approval signatures}: "I approve this document." Multiple
  allowed, no restrictions on subsequent changes, each covers all
  previous content.
\item
  \textbf{Certification signatures}: "I authored this document and
  define what can change." Only one per document, must be first, carries
  DocMDP permissions (P=1/2/3).
\item
  \textbf{P=2 is the default and most common} permission level ---
  allows form filling and additional signatures while preventing content
  modifications.
\item
  \textbf{Lock dictionaries} control which fields become read-only after
  a specific signature is applied. Cumulative and permanent.
\item
  \textbf{FieldMDP} provides validation-level detection of changes to
  specific fields --- complements Lock dictionaries.
\item
  \textbf{Seed value dictionaries} constrain how a signature field must
  be signed --- SubFilter, digest algorithm, timestamp requirements.
\item
  MDP is a \textbf{detection mechanism}, not a prevention mechanism. Not
  all viewers enforce it. Always verify compliance programmatically.
\item
  \textbf{Always apply certification signatures first}, then approval
  signatures. The certification establishes the permission framework.
\end{itemize}

\chapter{Visual Appearance of
Signatures}\label{_visual_appearance_of_signatures}

A digital signature has two faces. The cryptographic face --- the CMS
\texttt{SignedData}, the byte range, the certificate chain --- is what
makes the signature \emph{valid}. The visual face --- what the user sees
in the PDF viewer --- is what makes the signature \emph{trustworthy to a
human reader}. These two faces are entirely independent, and
understanding their independence is critical for both security and
usability.

The visual appearance of a PDF signature is an \textbf{annotation
appearance stream} --- a miniature PDF drawing that shows the signer's
name, date, reason, and optionally a handwritten signature image, a
company logo, or other graphical elements. This appearance stream is
part of the signed bytes (it lives inside the signature field's widget
annotation), but it is \emph{not} cryptographically bound to the CMS in
any meaningful way. A valid CMS signature does not guarantee that the
visual appearance accurately represents the signer --- it only
guarantees that the bytes haven't changed.

This distinction is exactly what Shadow Attacks (Chapter 24) exploit: an
attacker can create a PDF where the visual appearance shows one thing to
the signer and something different to the recipient, with the signature
validating correctly in both cases.

This chapter explains how signature appearances work, how to create them
with PDFBox, what the layer model means, and how different viewers
handle appearance rendering.

\section{Appearance Streams: The
Basics}\label{_appearance_streams_the_basics}

\subsection{What an Appearance Stream
Is}\label{_what_an_appearance_stream_is}

Every annotation in a PDF --- not just signatures --- can have an
\textbf{appearance stream}. An appearance stream is a Form XObject (a
\texttt{/Type\ /XObject\ /Subtype\ /Form} stream) that contains the
drawing instructions for how the annotation looks.

For a signature widget annotation, the appearance stream is stored under
the \texttt{/AP} (Appearance) key:

\begin{verbatim}
<< /Type /Annot
   /Subtype /Widget
   /FT /Sig
   /T (SignatureField1)
   /V 15 0 R                    % the /Sig dictionary
   /Rect [100 600 300 650]      % position on the page (x1 y1 x2 y2)
   /AP <<
     /N 20 0 R                  % Normal appearance stream
   >>
>>
\end{verbatim}

The \texttt{/N} key holds the \textbf{normal} appearance --- the one
shown when the annotation is not being interacted with. Annotations can
also have \texttt{/R} (rollover --- shown on mouse hover) and
\texttt{/D} (down --- shown when clicked), but signature appearances
only use \texttt{/N}.

The appearance stream (object 20 in this example) is a Form XObject:

\begin{verbatim}
20 0 obj
<< /Type /XObject
   /Subtype /Form
   /BBox [0 0 200 50]          % bounding box in the stream's coordinate space
   /Resources << /Font << /F1 21 0 R >> >>
   /Length 142
>>
stream
BT
/F1 10 Tf
0 35 Td
(Digitally signed by: Test Signer) Tj
0 -15 Td
(Date: 2026-03-08 12:00:00 UTC) Tj
0 -15 Td
(Reason: I approve this document) Tj
ET
endstream
endobj
\end{verbatim}

This is a complete miniature PDF page: it has a bounding box, font
resources, and drawing instructions using the standard PDF content
stream operators (\texttt{BT}/\texttt{ET} for text blocks, \texttt{Tf}
for font selection, \texttt{Td} for positioning, \texttt{Tj} for showing
text).

\subsection{Visible vs. Invisible
Signatures}\label{_visible_vs_invisible_signatures}

Not all signatures need a visual appearance. A signature can be:

\textbf{Visible}: The signature field has a non-zero \texttt{/Rect}
(annotation rectangle) and an appearance stream. The user sees a
signature block on a specific page.

\textbf{Invisible}: The signature field either has a zero-size
\texttt{/Rect} (\texttt{{[}0\ 0\ 0\ 0{]}}), or has no appearance stream,
or is not associated with any page. The signature exists
cryptographically but doesn't appear on any page.

Invisible signatures are common in automated signing workflows where the
PDF is being signed programmatically (e.g., a server signing all
outgoing invoices) and there's no need for a visual indicator on the
document. The signature is still listed in the viewer's signature panel
--- the user can still inspect it, verify the certificate, and check the
validation status.

Even for invisible signatures, some viewers (notably Adobe Acrobat)
display a blue banner at the top of the document saying "Signed and all
signatures are valid" or similar. This banner is not part of the PDF
content --- it's the viewer's UI.

\section{The Layer Model: n0, n2, and
n4}\label{_the_layer_model_n0_n2_and_n4}

\subsection{Adobe's Legacy Layer
System}\label{_adobes_legacy_layer_system}

Adobe Acrobat uses a layer-based approach to compose signature
appearances. This system was designed by Adobe for their own
implementation and is not formally part of the ISO 32000 specification
--- but because Adobe Acrobat is the dominant viewer, understanding it
matters.

The layers are:

\begin{itemize}
\item
  \textbf{n0} (Layer 0 --- Background): Contains the background graphic
  and a static validity indicator. Adobe Acrobat uses this to show the
  blue "valid" checkmark or the red "invalid" X. This layer is generated
  and updated by Acrobat itself --- it is not typically set by the
  signing application.
\item
  \textbf{n2} (Layer 2 --- Content): Contains the signer's information
  --- name, date, reason, location, and optionally a handwritten
  signature image or company logo. \textbf{This is the layer you create}
  when building a visual signature appearance.
\item
  \textbf{n4} (Layer 4 --- Text overlay): An optional text overlay,
  rarely used.
\end{itemize}

In practice, most signing applications (including PDFBox-based
implementations) create a single appearance stream that serves as the
entire appearance --- they don't use the n0/n2/n4 layer decomposition.
Adobe Acrobat, when it encounters a signature without layer-based
appearances, creates its own composite appearance using the provided
stream as the content layer and adding its own validity indicators.

The layer model is an implementation detail of Adobe Acrobat, not a spec
requirement. When you create a signature appearance with PDFBox, you
create a single Form XObject for the \texttt{/AP\ /N} entry. Acrobat
will overlay its own validity graphics on top. Other viewers
(LibreOffice, Foxit, browser-based viewers) each handle the appearance
differently --- some show your appearance as-is, some add their own
decorations.

\subsection{What This Means for Your
Implementation}\label{_what_this_means_for_your_implementation}

When building a signature appearance with PDFBox:

\begin{enumerate}
\def\labelenumi{\arabic{enumi}.}
\item
  \textbf{Create a single Form XObject} with the signer's information
\item
  \textbf{Don't worry about n0/n2/n4 layers} --- the single appearance
  works across all viewers
\item
  \textbf{Set reasonable fonts and sizes} --- the appearance will be
  rendered in the space defined by the signature field's \texttt{/Rect}
\item
  \textbf{Include key information}: signer's name, signing date/time,
  reason (if applicable)
\item
  \textbf{Test in multiple viewers} --- Acrobat, Foxit, LibreOffice, and
  a browser-based viewer --- to ensure the appearance renders correctly
  in each
\end{enumerate}

\section{Creating Signature Appearances with
PDFBox}\label{_creating_signature_appearances_with_pdfbox}

\subsection{Basic Text Appearance}\label{_basic_text_appearance}

The simplest appearance shows the signer's name and date:

\begin{Shaded}
\begin{Highlighting}[]
\KeywordTok{private}\NormalTok{ PDFormXObject }\FunctionTok{createSignatureAppearance}\OperatorTok{(}\NormalTok{PDDocument doc}\OperatorTok{,}
\NormalTok{        PDRectangle rect}\OperatorTok{,} \BuiltInTok{String}\NormalTok{ signerName}\OperatorTok{,} \BuiltInTok{Calendar}\NormalTok{ signDate}\OperatorTok{)}
        \KeywordTok{throws} \BuiltInTok{IOException} \OperatorTok{\{}

\NormalTok{    PDFormXObject form }\OperatorTok{=} \KeywordTok{new} \FunctionTok{PDFormXObject}\OperatorTok{(}\NormalTok{doc}\OperatorTok{.}\FunctionTok{getDocument}\OperatorTok{().}\FunctionTok{createCOSStream}\OperatorTok{());}
\NormalTok{    form}\OperatorTok{.}\FunctionTok{setResources}\OperatorTok{(}\KeywordTok{new} \FunctionTok{PDResources}\OperatorTok{());}
\NormalTok{    form}\OperatorTok{.}\FunctionTok{setBBox}\OperatorTok{(}\KeywordTok{new} \FunctionTok{PDRectangle}\OperatorTok{(}\NormalTok{rect}\OperatorTok{.}\FunctionTok{getWidth}\OperatorTok{(),}\NormalTok{ rect}\OperatorTok{.}\FunctionTok{getHeight}\OperatorTok{()));}

\NormalTok{    PDFont font }\OperatorTok{=} \KeywordTok{new} \FunctionTok{PDType1Font}\OperatorTok{(}\NormalTok{Standard14Fonts}\OperatorTok{.}\FunctionTok{FontName}\OperatorTok{.}\FunctionTok{HELVETICA}\OperatorTok{);}
\NormalTok{    form}\OperatorTok{.}\FunctionTok{getResources}\OperatorTok{().}\FunctionTok{put}\OperatorTok{(}\NormalTok{COSName}\OperatorTok{.}\FunctionTok{getPDFName}\OperatorTok{(}\StringTok{"F1"}\OperatorTok{),}\NormalTok{ font}\OperatorTok{);}

    \DataTypeTok{float}\NormalTok{ fontSize }\OperatorTok{=} \DecValTok{8}\OperatorTok{;}
    \BuiltInTok{String}\NormalTok{ dateStr }\OperatorTok{=} \KeywordTok{new} \BuiltInTok{SimpleDateFormat}\OperatorTok{(}\StringTok{"yyyy{-}MM{-}dd HH:mm:ss z"}\OperatorTok{)}
            \OperatorTok{.}\FunctionTok{format}\OperatorTok{(}\NormalTok{signDate}\OperatorTok{.}\FunctionTok{getTime}\OperatorTok{());}

    \CommentTok{// Build the content stream manually}
    \ControlFlowTok{try} \OperatorTok{(}\BuiltInTok{OutputStream}\NormalTok{ os }\OperatorTok{=}\NormalTok{ form}\OperatorTok{.}\FunctionTok{getStream}\OperatorTok{().}\FunctionTok{createOutputStream}\OperatorTok{())} \OperatorTok{\{}
        \BuiltInTok{String}\NormalTok{ content }\OperatorTok{=} \BuiltInTok{String}\OperatorTok{.}\FunctionTok{join}\OperatorTok{(}\StringTok{"}\SpecialCharTok{\textbackslash{}n}\StringTok{"}\OperatorTok{,}
            \StringTok{"BT"}\OperatorTok{,}
            \StringTok{"/F1 "} \OperatorTok{+}\NormalTok{ fontSize }\OperatorTok{+} \StringTok{" Tf"}\OperatorTok{,}
            \StringTok{"0 "} \OperatorTok{+} \OperatorTok{(}\NormalTok{rect}\OperatorTok{.}\FunctionTok{getHeight}\OperatorTok{()} \OperatorTok{{-}}\NormalTok{ fontSize }\OperatorTok{{-}} \DecValTok{2}\OperatorTok{)} \OperatorTok{+} \StringTok{" Td"}\OperatorTok{,}
            \StringTok{"("} \OperatorTok{+} \FunctionTok{escapeForPdf}\OperatorTok{(}\NormalTok{signerName}\OperatorTok{)} \OperatorTok{+} \StringTok{") Tj"}\OperatorTok{,}
            \StringTok{"0 {-}"} \OperatorTok{+} \OperatorTok{(}\NormalTok{fontSize }\OperatorTok{+} \DecValTok{2}\OperatorTok{)} \OperatorTok{+} \StringTok{" Td"}\OperatorTok{,}
            \StringTok{"(Date: "} \OperatorTok{+} \FunctionTok{escapeForPdf}\OperatorTok{(}\NormalTok{dateStr}\OperatorTok{)} \OperatorTok{+} \StringTok{") Tj"}\OperatorTok{,}
            \StringTok{"ET"}
        \OperatorTok{);}
\NormalTok{        os}\OperatorTok{.}\FunctionTok{write}\OperatorTok{(}\NormalTok{content}\OperatorTok{.}\FunctionTok{getBytes}\OperatorTok{(}\NormalTok{StandardCharsets}\OperatorTok{.}\FunctionTok{US\_ASCII}\OperatorTok{));}
    \OperatorTok{\}}

    \ControlFlowTok{return}\NormalTok{ form}\OperatorTok{;}
\OperatorTok{\}}

\KeywordTok{private} \BuiltInTok{String} \FunctionTok{escapeForPdf}\OperatorTok{(}\BuiltInTok{String}\NormalTok{ text}\OperatorTok{)} \OperatorTok{\{}
    \ControlFlowTok{return}\NormalTok{ text}\OperatorTok{.}\FunctionTok{replace}\OperatorTok{(}\StringTok{"}\SpecialCharTok{\textbackslash{}\textbackslash{}}\StringTok{"}\OperatorTok{,} \StringTok{"}\SpecialCharTok{\textbackslash{}\textbackslash{}\textbackslash{}\textbackslash{}}\StringTok{"}\OperatorTok{)}
               \OperatorTok{.}\FunctionTok{replace}\OperatorTok{(}\StringTok{"("}\OperatorTok{,} \StringTok{"}\SpecialCharTok{\textbackslash{}\textbackslash{}}\StringTok{("}\OperatorTok{)}
               \OperatorTok{.}\FunctionTok{replace}\OperatorTok{(}\StringTok{")"}\OperatorTok{,} \StringTok{"}\SpecialCharTok{\textbackslash{}\textbackslash{}}\StringTok{)"}\OperatorTok{);}
\OperatorTok{\}}
\end{Highlighting}
\end{Shaded}

In PDFBox 3.x, you \textbf{cannot} use \texttt{PDPageContentStream} to
write to a Form XObject. \texttt{PDPageContentStream} is strictly for
page content. For Form XObjects, you must write raw PDF operators to the
stream's output stream, as shown above.

This is one of the breaking changes from PDFBox 2.x, where
\texttt{PDPageContentStream} could be used with Form XObjects via a
constructor overload that accepted a \texttt{PDFormXObject} parameter.
That constructor no longer exists in 3.x.

\subsection{Appearance with Signature
Image}\label{_appearance_with_signature_image}

Many signing workflows include a handwritten signature image --- either
captured via a canvas on a web page or tablet, or uploaded as a PNG/JPEG
file. Adding an image to the appearance requires embedding it as a PDF
image XObject:

\begin{Shaded}
\begin{Highlighting}[]
\KeywordTok{private}\NormalTok{ PDFormXObject }\FunctionTok{createAppearanceWithImage}\OperatorTok{(}\NormalTok{PDDocument doc}\OperatorTok{,}
\NormalTok{        PDRectangle rect}\OperatorTok{,} \BuiltInTok{String}\NormalTok{ signerName}\OperatorTok{,} \BuiltInTok{Calendar}\NormalTok{ signDate}\OperatorTok{,}
        \DataTypeTok{byte}\OperatorTok{[]}\NormalTok{ signatureImagePng}\OperatorTok{)} \KeywordTok{throws} \BuiltInTok{IOException} \OperatorTok{\{}

\NormalTok{    PDFormXObject form }\OperatorTok{=} \KeywordTok{new} \FunctionTok{PDFormXObject}\OperatorTok{(}\NormalTok{doc}\OperatorTok{.}\FunctionTok{getDocument}\OperatorTok{().}\FunctionTok{createCOSStream}\OperatorTok{());}
\NormalTok{    form}\OperatorTok{.}\FunctionTok{setResources}\OperatorTok{(}\KeywordTok{new} \FunctionTok{PDResources}\OperatorTok{());}
\NormalTok{    form}\OperatorTok{.}\FunctionTok{setBBox}\OperatorTok{(}\KeywordTok{new} \FunctionTok{PDRectangle}\OperatorTok{(}\NormalTok{rect}\OperatorTok{.}\FunctionTok{getWidth}\OperatorTok{(),}\NormalTok{ rect}\OperatorTok{.}\FunctionTok{getHeight}\OperatorTok{()));}

    \CommentTok{// Add font}
\NormalTok{    PDFont font }\OperatorTok{=} \KeywordTok{new} \FunctionTok{PDType1Font}\OperatorTok{(}\NormalTok{Standard14Fonts}\OperatorTok{.}\FunctionTok{FontName}\OperatorTok{.}\FunctionTok{HELVETICA}\OperatorTok{);}
\NormalTok{    form}\OperatorTok{.}\FunctionTok{getResources}\OperatorTok{().}\FunctionTok{put}\OperatorTok{(}\NormalTok{COSName}\OperatorTok{.}\FunctionTok{getPDFName}\OperatorTok{(}\StringTok{"F1"}\OperatorTok{),}\NormalTok{ font}\OperatorTok{);}

    \CommentTok{// Add signature image}
\NormalTok{    PDImageXObject image }\OperatorTok{=}\NormalTok{ PDImageXObject}\OperatorTok{.}\FunctionTok{createFromByteArray}\OperatorTok{(}
\NormalTok{            doc}\OperatorTok{,}\NormalTok{ signatureImagePng}\OperatorTok{,} \StringTok{"signature"}\OperatorTok{);}
\NormalTok{    form}\OperatorTok{.}\FunctionTok{getResources}\OperatorTok{().}\FunctionTok{put}\OperatorTok{(}\NormalTok{COSName}\OperatorTok{.}\FunctionTok{getPDFName}\OperatorTok{(}\StringTok{"Img0"}\OperatorTok{),}\NormalTok{ image}\OperatorTok{);}

    \CommentTok{// Layout: image on the left, text on the right}
    \DataTypeTok{float}\NormalTok{ imgWidth }\OperatorTok{=}\NormalTok{ rect}\OperatorTok{.}\FunctionTok{getWidth}\OperatorTok{()} \OperatorTok{*} \FloatTok{0.4f}\OperatorTok{;}
    \DataTypeTok{float}\NormalTok{ imgHeight }\OperatorTok{=}\NormalTok{ rect}\OperatorTok{.}\FunctionTok{getHeight}\OperatorTok{()} \OperatorTok{{-}} \DecValTok{4}\OperatorTok{;}
    \DataTypeTok{float}\NormalTok{ textX }\OperatorTok{=}\NormalTok{ imgWidth }\OperatorTok{+} \DecValTok{5}\OperatorTok{;}
    \DataTypeTok{float}\NormalTok{ fontSize }\OperatorTok{=} \DecValTok{7}\OperatorTok{;}
    \BuiltInTok{String}\NormalTok{ dateStr }\OperatorTok{=} \KeywordTok{new} \BuiltInTok{SimpleDateFormat}\OperatorTok{(}\StringTok{"yyyy{-}MM{-}dd HH:mm:ss z"}\OperatorTok{)}
            \OperatorTok{.}\FunctionTok{format}\OperatorTok{(}\NormalTok{signDate}\OperatorTok{.}\FunctionTok{getTime}\OperatorTok{());}

    \ControlFlowTok{try} \OperatorTok{(}\BuiltInTok{OutputStream}\NormalTok{ os }\OperatorTok{=}\NormalTok{ form}\OperatorTok{.}\FunctionTok{getStream}\OperatorTok{().}\FunctionTok{createOutputStream}\OperatorTok{())} \OperatorTok{\{}
        \BuiltInTok{StringBuilder}\NormalTok{ sb }\OperatorTok{=} \KeywordTok{new} \BuiltInTok{StringBuilder}\OperatorTok{();}
        \CommentTok{// Draw image}
\NormalTok{        sb}\OperatorTok{.}\FunctionTok{append}\OperatorTok{(}\StringTok{"q}\SpecialCharTok{\textbackslash{}n}\StringTok{"}\OperatorTok{);}  \CommentTok{// save graphics state}
\NormalTok{        sb}\OperatorTok{.}\FunctionTok{append}\OperatorTok{(}\NormalTok{imgWidth}\OperatorTok{).}\FunctionTok{append}\OperatorTok{(}\StringTok{" 0 0 "}\OperatorTok{).}\FunctionTok{append}\OperatorTok{(}\NormalTok{imgHeight}\OperatorTok{)}
          \OperatorTok{.}\FunctionTok{append}\OperatorTok{(}\StringTok{" 2 2 cm}\SpecialCharTok{\textbackslash{}n}\StringTok{"}\OperatorTok{);}  \CommentTok{// transform matrix}
\NormalTok{        sb}\OperatorTok{.}\FunctionTok{append}\OperatorTok{(}\StringTok{"/Img0 Do}\SpecialCharTok{\textbackslash{}n}\StringTok{"}\OperatorTok{);}  \CommentTok{// draw image}
\NormalTok{        sb}\OperatorTok{.}\FunctionTok{append}\OperatorTok{(}\StringTok{"Q}\SpecialCharTok{\textbackslash{}n}\StringTok{"}\OperatorTok{);}  \CommentTok{// restore graphics state}
        \CommentTok{// Draw text}
\NormalTok{        sb}\OperatorTok{.}\FunctionTok{append}\OperatorTok{(}\StringTok{"BT}\SpecialCharTok{\textbackslash{}n}\StringTok{"}\OperatorTok{);}
\NormalTok{        sb}\OperatorTok{.}\FunctionTok{append}\OperatorTok{(}\StringTok{"/F1 "}\OperatorTok{).}\FunctionTok{append}\OperatorTok{(}\NormalTok{fontSize}\OperatorTok{).}\FunctionTok{append}\OperatorTok{(}\StringTok{" Tf}\SpecialCharTok{\textbackslash{}n}\StringTok{"}\OperatorTok{);}
\NormalTok{        sb}\OperatorTok{.}\FunctionTok{append}\OperatorTok{(}\NormalTok{textX}\OperatorTok{).}\FunctionTok{append}\OperatorTok{(}\StringTok{" "}\OperatorTok{)}
          \OperatorTok{.}\FunctionTok{append}\OperatorTok{(}\NormalTok{rect}\OperatorTok{.}\FunctionTok{getHeight}\OperatorTok{()} \OperatorTok{{-}}\NormalTok{ fontSize }\OperatorTok{{-}} \DecValTok{2}\OperatorTok{).}\FunctionTok{append}\OperatorTok{(}\StringTok{" Td}\SpecialCharTok{\textbackslash{}n}\StringTok{"}\OperatorTok{);}
\NormalTok{        sb}\OperatorTok{.}\FunctionTok{append}\OperatorTok{(}\StringTok{"(Digitally signed by:) Tj}\SpecialCharTok{\textbackslash{}n}\StringTok{"}\OperatorTok{);}
\NormalTok{        sb}\OperatorTok{.}\FunctionTok{append}\OperatorTok{(}\StringTok{"0 {-}"}\OperatorTok{).}\FunctionTok{append}\OperatorTok{(}\NormalTok{fontSize }\OperatorTok{+} \DecValTok{2}\OperatorTok{).}\FunctionTok{append}\OperatorTok{(}\StringTok{" Td}\SpecialCharTok{\textbackslash{}n}\StringTok{"}\OperatorTok{);}
\NormalTok{        sb}\OperatorTok{.}\FunctionTok{append}\OperatorTok{(}\StringTok{"("}\OperatorTok{).}\FunctionTok{append}\OperatorTok{(}\FunctionTok{escapeForPdf}\OperatorTok{(}\NormalTok{signerName}\OperatorTok{)).}\FunctionTok{append}\OperatorTok{(}\StringTok{") Tj}\SpecialCharTok{\textbackslash{}n}\StringTok{"}\OperatorTok{);}
\NormalTok{        sb}\OperatorTok{.}\FunctionTok{append}\OperatorTok{(}\StringTok{"0 {-}"}\OperatorTok{).}\FunctionTok{append}\OperatorTok{(}\NormalTok{fontSize }\OperatorTok{+} \DecValTok{2}\OperatorTok{).}\FunctionTok{append}\OperatorTok{(}\StringTok{" Td}\SpecialCharTok{\textbackslash{}n}\StringTok{"}\OperatorTok{);}
\NormalTok{        sb}\OperatorTok{.}\FunctionTok{append}\OperatorTok{(}\StringTok{"(Date: "}\OperatorTok{).}\FunctionTok{append}\OperatorTok{(}\FunctionTok{escapeForPdf}\OperatorTok{(}\NormalTok{dateStr}\OperatorTok{)).}\FunctionTok{append}\OperatorTok{(}\StringTok{") Tj}\SpecialCharTok{\textbackslash{}n}\StringTok{"}\OperatorTok{);}
\NormalTok{        sb}\OperatorTok{.}\FunctionTok{append}\OperatorTok{(}\StringTok{"ET}\SpecialCharTok{\textbackslash{}n}\StringTok{"}\OperatorTok{);}
\NormalTok{        os}\OperatorTok{.}\FunctionTok{write}\OperatorTok{(}\NormalTok{sb}\OperatorTok{.}\FunctionTok{toString}\OperatorTok{().}\FunctionTok{getBytes}\OperatorTok{(}\NormalTok{StandardCharsets}\OperatorTok{.}\FunctionTok{US\_ASCII}\OperatorTok{));}
    \OperatorTok{\}}

    \ControlFlowTok{return}\NormalTok{ form}\OperatorTok{;}
\OperatorTok{\}}
\end{Highlighting}
\end{Shaded}

\subsection{Attaching the Appearance to the
Signature}\label{_attaching_the_appearance_to_the_signature}

After creating the appearance Form XObject, attach it to the signature
field's widget annotation:

\begin{Shaded}
\begin{Highlighting}[]
\CommentTok{// Create signature field and widget annotation}
\NormalTok{PDSignature signature }\OperatorTok{=} \KeywordTok{new} \FunctionTok{PDSignature}\OperatorTok{();}
\CommentTok{// ... set signature properties ...}

\CommentTok{// Set the visual rectangle on page 1}
\NormalTok{PDPage page }\OperatorTok{=}\NormalTok{ doc}\OperatorTok{.}\FunctionTok{getPage}\OperatorTok{(}\DecValTok{0}\OperatorTok{);}
\NormalTok{PDRectangle rect }\OperatorTok{=} \KeywordTok{new} \FunctionTok{PDRectangle}\OperatorTok{(}\DecValTok{100}\OperatorTok{,} \DecValTok{600}\OperatorTok{,} \DecValTok{200}\OperatorTok{,} \DecValTok{50}\OperatorTok{);}

\NormalTok{SignatureOptions options }\OperatorTok{=} \KeywordTok{new} \FunctionTok{SignatureOptions}\OperatorTok{();}
\NormalTok{options}\OperatorTok{.}\FunctionTok{setPreferredSignatureSize}\OperatorTok{(}\DecValTok{65536}\OperatorTok{);}
\NormalTok{options}\OperatorTok{.}\FunctionTok{setPage}\OperatorTok{(}\DecValTok{0}\OperatorTok{);}  \CommentTok{// page index (0{-}based)}
\NormalTok{options}\OperatorTok{.}\FunctionTok{setVisualSignature}\OperatorTok{(}
    \FunctionTok{createSignatureAppearance}\OperatorTok{(}\NormalTok{doc}\OperatorTok{,}\NormalTok{ rect}\OperatorTok{,} \StringTok{"Test Signer"}\OperatorTok{,} \BuiltInTok{Calendar}\OperatorTok{.}\FunctionTok{getInstance}\OperatorTok{()));}

\NormalTok{doc}\OperatorTok{.}\FunctionTok{addSignature}\OperatorTok{(}\NormalTok{signature}\OperatorTok{,}\NormalTok{ signingInterface}\OperatorTok{,}\NormalTok{ options}\OperatorTok{);}
\NormalTok{doc}\OperatorTok{.}\FunctionTok{saveIncremental}\OperatorTok{(}\NormalTok{outputStream}\OperatorTok{);}
\end{Highlighting}
\end{Shaded}

\section{Appearance Rendering Across
Viewers}\label{_appearance_rendering_across_viewers}

\subsection{How Different Viewers Handle
Appearances}\label{_how_different_viewers_handle_appearances}

One of the most frustrating aspects of PDF signature appearances is that
different viewers render them differently --- even when the appearance
stream is identical.

\begin{longtable}[]{@{}
  >{\raggedright\arraybackslash}p{(\linewidth - 2\tabcolsep) * \real{0.2857}}
  >{\raggedright\arraybackslash}p{(\linewidth - 2\tabcolsep) * \real{0.7143}}@{}}
\caption{Signature appearance rendering behavior across major
viewers}\tabularnewline
\toprule\noalign{}
\begin{minipage}[b]{\linewidth}\raggedright
Viewer
\end{minipage} & \begin{minipage}[b]{\linewidth}\raggedright
Behavior
\end{minipage} \\
\midrule\noalign{}
\endfirsthead
\toprule\noalign{}
\begin{minipage}[b]{\linewidth}\raggedright
Viewer
\end{minipage} & \begin{minipage}[b]{\linewidth}\raggedright
Behavior
\end{minipage} \\
\midrule\noalign{}
\endhead
\bottomrule\noalign{}
\endlastfoot
\textbf{Adobe Acrobat} & Overlays its own validity indicator (blue
checkmark / red X) on layer n0. Your appearance stream is rendered on
layer n2. Acrobat may resize your content to fit its layer model. If
\texttt{/NeedAppearances} is true in the AcroForm, Acrobat may
regenerate appearances from scratch. \\
\textbf{Foxit Reader} & Renders the appearance stream mostly as
provided. Adds a small validity icon in the corner. Generally the most
faithful rendering of custom appearances. \\
\textbf{LibreOffice Draw} & Renders the appearance stream as-is. Does
not overlay validity indicators on the appearance itself --- shows
validation status in a separate panel. Good for testing that your raw
appearance looks correct. \\
\textbf{Chrome (built-in viewer)} & Renders appearance streams but does
not validate signatures. No validity indicators. The appearance is what
you see --- useful for checking visual layout. \\
\textbf{Firefox (pdf.js)} & Renders appearance streams. Does not
validate signatures (though experimental validation support has been
discussed). Like Chrome, useful for layout testing. \\
\textbf{Okular (KDE)} & Renders appearances and provides signature
validation. May show validity status differently from other viewers. \\
\end{longtable}

\textbf{Always test your signature appearances in at least three
viewers} --- Adobe Acrobat, a non-Adobe desktop viewer (Foxit or
LibreOffice), and a browser. Appearances that look perfect in Acrobat
may render differently in other viewers due to:

\begin{itemize}
\item
  Different font fallback behavior (if the font isn't embedded)
\item
  Different handling of the coordinate system and bounding box
\item
  Different overlay/decoration behavior
\item
  Different treatment of transparency and blending
\end{itemize}

\subsection{\texorpdfstring{The \texttt{/NeedAppearances}
Problem}{The /NeedAppearances Problem}}\label{_the_needappearances_problem}

The AcroForm dictionary has a boolean flag \texttt{/NeedAppearances}
that tells viewers whether to regenerate form field appearances from
scratch. If this flag is \texttt{true}, some viewers (notably Adobe
Acrobat) will ignore the existing appearance streams and generate their
own based on the field's value.

For signature fields, this is usually not what you want. If
\texttt{/NeedAppearances} is true and Acrobat regenerates the
signature's appearance, your carefully crafted visual (with the signer's
image, custom layout, etc.) gets replaced with Acrobat's default
rendering.

\textbf{Always set \texttt{/NeedAppearances} to \texttt{false}} (or
remove the key entirely --- the default is \texttt{false}) before
signing:

\begin{Shaded}
\begin{Highlighting}[]
\NormalTok{PDAcroForm acroForm }\OperatorTok{=}\NormalTok{ doc}\OperatorTok{.}\FunctionTok{getDocumentCatalog}\OperatorTok{().}\FunctionTok{getAcroForm}\OperatorTok{();}
\ControlFlowTok{if} \OperatorTok{(}\NormalTok{acroForm }\OperatorTok{!=} \KeywordTok{null}\OperatorTok{)} \OperatorTok{\{}
\NormalTok{    acroForm}\OperatorTok{.}\FunctionTok{getCOSObject}\OperatorTok{().}\FunctionTok{setBoolean}\OperatorTok{(}\NormalTok{COSName}\OperatorTok{.}\FunctionTok{NEED\_APPEARANCES}\OperatorTok{,} \KeywordTok{false}\OperatorTok{);}
\OperatorTok{\}}
\end{Highlighting}
\end{Shaded}

\section{Positioning: Getting the Rectangle
Right}\label{_positioning_getting_the_rectangle_right}

\subsection{The Coordinate System}\label{_the_coordinate_system}

PDF uses a \textbf{bottom-left origin} coordinate system where (0, 0) is
at the lower-left corner of the page. The y-axis goes \emph{up}, not
down. This is the opposite of most screen coordinate systems and is a
common source of positioning errors.

The signature field's \texttt{/Rect} array defines the annotation
rectangle as \texttt{{[}x1\ y1\ x2\ y2{]}}, where:

\begin{itemize}
\item
  \texttt{(x1,\ y1)} is the lower-left corner
\item
  \texttt{(x2,\ y2)} is the upper-right corner
\end{itemize}

So a signature field at the bottom-right of an A4 page (595 x 842
points) might have:

\begin{verbatim}
/Rect [395 50 575 100]    % 180pt wide, 50pt tall, near bottom-right
\end{verbatim}

A common mistake: specifying the rectangle with screen coordinates
(origin at top-left, y going down). If you compute
\texttt{y\ =\ pageHeight\ -\ desiredY}, make sure you do this for
\emph{both} y1 and y2, and that y1 \textless{} y2.

\textbf{Debug tip}: If your signature appearance shows up in the wrong
position, check:

\begin{enumerate}
\def\labelenumi{\arabic{enumi}.}
\item
  Is the page rotated? (A \texttt{/Rotate} key on the page dictionary
  changes the visible coordinate system)
\item
  Is the \texttt{MediaBox} the expected size? (Custom page sizes shift
  everything)
\item
  Did you flip y1 and y2? (The lower-left corner must have the smaller y
  value)
\end{enumerate}

\subsection{Handling Page Rotation}\label{_handling_page_rotation}

PDF pages can have a \texttt{/Rotate} key that specifies a rotation
angle (0, 90, 180, or 270 degrees). This rotation affects how the page
is \emph{displayed} but does not change the coordinate system of the
page content or annotations.

This means that on a 90-degree rotated page (landscape orientation
achieved via rotation), the annotation rectangle coordinates are still
in the unrotated coordinate space. To place a signature at the visual
bottom-right of a rotated page, you need to account for the rotation.

\begin{Shaded}
\begin{Highlighting}[]
\NormalTok{PDPage page }\OperatorTok{=}\NormalTok{ doc}\OperatorTok{.}\FunctionTok{getPage}\OperatorTok{(}\NormalTok{pageIndex}\OperatorTok{);}
\DataTypeTok{int}\NormalTok{ rotation }\OperatorTok{=}\NormalTok{ page}\OperatorTok{.}\FunctionTok{getRotation}\OperatorTok{();}
\NormalTok{PDRectangle mediaBox }\OperatorTok{=}\NormalTok{ page}\OperatorTok{.}\FunctionTok{getMediaBox}\OperatorTok{();}

\DataTypeTok{float}\NormalTok{ sigWidth }\OperatorTok{=} \DecValTok{200}\OperatorTok{;}
\DataTypeTok{float}\NormalTok{ sigHeight }\OperatorTok{=} \DecValTok{50}\OperatorTok{;}
\DataTypeTok{float}\NormalTok{ margin }\OperatorTok{=} \DecValTok{30}\OperatorTok{;}

\NormalTok{PDRectangle sigRect}\OperatorTok{;}
\ControlFlowTok{switch} \OperatorTok{(}\NormalTok{rotation}\OperatorTok{)} \OperatorTok{\{}
    \ControlFlowTok{case} \DecValTok{0}\OperatorTok{:}   \CommentTok{// Normal portrait}
\NormalTok{        sigRect }\OperatorTok{=} \KeywordTok{new} \FunctionTok{PDRectangle}\OperatorTok{(}
\NormalTok{            mediaBox}\OperatorTok{.}\FunctionTok{getWidth}\OperatorTok{()} \OperatorTok{{-}}\NormalTok{ sigWidth }\OperatorTok{{-}}\NormalTok{ margin}\OperatorTok{,}
\NormalTok{            margin}\OperatorTok{,}
\NormalTok{            sigWidth}\OperatorTok{,}\NormalTok{ sigHeight}\OperatorTok{);}
        \ControlFlowTok{break}\OperatorTok{;}
    \ControlFlowTok{case} \DecValTok{90}\OperatorTok{:}  \CommentTok{// Landscape (rotated clockwise)}
\NormalTok{        sigRect }\OperatorTok{=} \KeywordTok{new} \FunctionTok{PDRectangle}\OperatorTok{(}
\NormalTok{            margin}\OperatorTok{,}
\NormalTok{            margin}\OperatorTok{,}
\NormalTok{            sigHeight}\OperatorTok{,}\NormalTok{ sigWidth}\OperatorTok{);}  \CommentTok{// swap width/height}
        \ControlFlowTok{break}\OperatorTok{;}
    \ControlFlowTok{case} \DecValTok{180}\OperatorTok{:} \CommentTok{// Upside{-}down portrait}
\NormalTok{        sigRect }\OperatorTok{=} \KeywordTok{new} \FunctionTok{PDRectangle}\OperatorTok{(}
\NormalTok{            margin}\OperatorTok{,}
\NormalTok{            mediaBox}\OperatorTok{.}\FunctionTok{getHeight}\OperatorTok{()} \OperatorTok{{-}}\NormalTok{ sigHeight }\OperatorTok{{-}}\NormalTok{ margin}\OperatorTok{,}
\NormalTok{            sigWidth}\OperatorTok{,}\NormalTok{ sigHeight}\OperatorTok{);}
        \ControlFlowTok{break}\OperatorTok{;}
    \ControlFlowTok{case} \DecValTok{270}\OperatorTok{:} \CommentTok{// Landscape (rotated counter{-}clockwise)}
\NormalTok{        sigRect }\OperatorTok{=} \KeywordTok{new} \FunctionTok{PDRectangle}\OperatorTok{(}
\NormalTok{            mediaBox}\OperatorTok{.}\FunctionTok{getWidth}\OperatorTok{()} \OperatorTok{{-}}\NormalTok{ sigHeight }\OperatorTok{{-}}\NormalTok{ margin}\OperatorTok{,}
\NormalTok{            mediaBox}\OperatorTok{.}\FunctionTok{getHeight}\OperatorTok{()} \OperatorTok{{-}}\NormalTok{ sigWidth }\OperatorTok{{-}}\NormalTok{ margin}\OperatorTok{,}
\NormalTok{            sigHeight}\OperatorTok{,}\NormalTok{ sigWidth}\OperatorTok{);}
        \ControlFlowTok{break}\OperatorTok{;}
    \KeywordTok{default}\OperatorTok{:}
\NormalTok{        sigRect }\OperatorTok{=} \KeywordTok{new} \FunctionTok{PDRectangle}\OperatorTok{(}\NormalTok{margin}\OperatorTok{,}\NormalTok{ margin}\OperatorTok{,}\NormalTok{ sigWidth}\OperatorTok{,}\NormalTok{ sigHeight}\OperatorTok{);}
\OperatorTok{\}}
\end{Highlighting}
\end{Shaded}

\section{Security Implications of Visual
Appearances}\label{_security_implications_of_visual_appearances}

\subsection{The Appearance/Content
Disconnect}\label{_the_appearancecontent_disconnect}

The visual appearance of a signature is stored in the PDF's object
structure, which is covered by the byte range. But the appearance is not
cryptographically bound to the \emph{identity} of the signer. The
appearance can show anything --- "Signed by Alice" while the CMS is
signed by Bob's certificate.

A conforming viewer should display the certificate information from the
CMS (not from the appearance) in its signature properties panel. But
what the user \emph{sees} on the page is the appearance --- and most
users trust the visual appearance more than the signature properties
panel.

This is a deliberate design trade-off in PDF: the signer (or the signing
application) controls what the appearance looks like, and the viewer
overlays validity information. The validity indicators (checkmarks,
status banners) come from the viewer's own analysis of the CMS --- they
don't come from the appearance stream.

\subsection{Shadow Attacks and Visual
Appearances}\label{_shadow_attacks_and_visual_appearances}

Shadow Attacks (Chapter 24) exploit the gap between what a signer sees
and what the document contains. In the \textbf{Replace} variant, the
attacker creates a form field with a specific visual appearance, then
changes the field's value (and thus its appearance) after the signer
certifies the document. Because the certification allows form filling
(P=2 or P=3), the change is permitted --- but the visual content of the
document has changed dramatically.

The defense is the same as for all Shadow Attacks: validate not just the
cryptographic integrity but also the \emph{semantic} integrity --- what
changed between revisions, and were those changes consistent with the
displayed content?

\textbf{For maximum security}: When validating a signed PDF, don't just
check the signature's validity. Extract the revision that was signed (by
truncating the file at the byte range end --- see Chapter 13) and
compare its visual rendering to the current document. If they differ,
the document has been modified since signing --- even if the
modifications were technically "permitted" by the MDP level.

\section{Common Appearance Problems}\label{_common_appearance_problems}

\subsection{Problem: Blank Signature
Rectangle}\label{_problem_blank_signature_rectangle}

\textbf{Symptom}: The signature field shows as an empty box in the
viewer.

\textbf{Causes}:

\begin{itemize}
\item
  No appearance stream set (the \texttt{/AP} dictionary is missing)
\item
  \texttt{/NeedAppearances} is \texttt{true}, and the viewer regenerated
  an empty appearance
\item
  The Form XObject has an empty stream (no drawing instructions)
\item
  The bounding box (\texttt{/BBox}) is {[}0 0 0 0{]} or has zero
  dimensions
\end{itemize}

\subsection{Problem: Text Appears Clipped or
Off-Screen}\label{_problem_text_appears_clipped_or_off_screen}

\textbf{Symptom}: Only part of the signature text is visible, or the
text appears outside the signature rectangle.

\textbf{Causes}:

\begin{itemize}
\item
  The text positioning (\texttt{Td} operators) exceeds the Form
  XObject's \texttt{/BBox}
\item
  The font size is too large for the available space
\item
  The \texttt{/BBox} doesn't match the annotation's \texttt{/Rect}
  dimensions
\end{itemize}

\textbf{Fix}: Ensure the Form XObject's \texttt{/BBox} width and height
match the annotation \texttt{/Rect} dimensions, and that all drawing
instructions fall within the \texttt{/BBox}.

\subsection{Problem: Font Not
Rendering}\label{_problem_font_not_rendering}

\textbf{Symptom}: Text appears as squares, dots, or is completely
invisible.

\textbf{Causes}:

\begin{itemize}
\item
  The font resource is not declared in the Form XObject's
  \texttt{/Resources}
\item
  The font name in the content stream (\texttt{/F1}) doesn't match the
  resource key
\item
  Non-ASCII characters used with a non-embedded Standard 14 font
  (Standard 14 fonts only support WinAnsiEncoding)
\end{itemize}

\textbf{Fix}: For text beyond basic Latin characters (accented letters,
CJK, etc.), embed a TrueType or Type0 font. Standard 14 fonts
(Helvetica, Times, Courier, etc.) only support a limited character set.

\section{Key Takeaways}\label{_key_takeaways_3}

\begin{itemize}
\item
  Signature appearance is an \textbf{annotation appearance stream} --- a
  Form XObject with drawing instructions
\item
  The appearance is \textbf{visually independent} from the cryptographic
  signature --- they can show different things
\item
  \textbf{Invisible signatures} have zero-size rectangles or no
  appearance stream --- they exist cryptographically but don't show on
  the page
\item
  The \textbf{n0/n2/n4 layer model} is Adobe's implementation detail,
  not a spec requirement --- create a single Form XObject for maximum
  compatibility
\item
  \textbf{PDFBox 3.x} requires writing raw PDF operators to Form XObject
  streams --- \texttt{PDPageContentStream} only works for pages
\item
  \textbf{Always set \texttt{/NeedAppearances} to \texttt{false}} before
  signing to prevent viewers from regenerating your custom appearance
\item
  \textbf{Test in multiple viewers} --- Acrobat, Foxit, LibreOffice, and
  browsers all render appearances differently
\item
  The \textbf{coordinate system} is bottom-left origin with y going up
  --- account for page rotation when positioning
\item
  Visual appearances are a \textbf{security surface} --- Shadow Attacks
  exploit the disconnect between appearance and content
\end{itemize}

\chapter{Signature Validation}\label{_signature_validation}

Signing a PDF is the easy part. Validating a signed PDF --- correctly,
securely, and in a way that matches what other validators produce --- is
where the real complexity lives. A signature can fail validation for
dozens of reasons, and the error messages from different viewers are
often vague or contradictory. Adobe says "valid." LibreOffice says
"signature cannot be verified." Your custom validator throws a
\texttt{CMSException}. Which one is right?

This chapter walks through the complete validation process, from the
first byte range check to the final certificate chain verification. We
cover what each validation step checks, what can go wrong at each step,
how different viewers handle each step, and how to debug validation
failures systematically.

\section{The Five Steps of Signature
Validation}\label{_the_five_steps_of_signature_validation}

Signature validation is not a single check --- it is a pipeline of five
independent verifications, each of which must pass. A failure at any
step means the signature is invalid (or at least questionable),
regardless of whether the other steps pass.

\subsection{Step 1: Structural Integrity (Byte
Range)}\label{_step_1_structural_integrity_byte_range}

\textbf{Question answered}: "Does the signature cover the expected bytes
of the file?"

\textbf{What to check}:

\begin{enumerate}
\def\labelenumi{\arabic{enumi}.}
\item
  The \texttt{/ByteRange} array has exactly 4 elements:
  \texttt{{[}offset1\ length1\ offset2\ length2{]}}
\item
  \texttt{offset1} is 0 (signature covers from the start of the file)
\item
  \texttt{offset2\ +\ length2} equals the file size (or revision size
  for multi-signature documents)
\item
  The gap \texttt{{[}length1,\ offset2)} contains only the
  \texttt{/Contents} hex string
\item
  No trailing unsigned bytes exist after the signed region (unless they
  are a subsequent legitimate revision)
\end{enumerate}

\textbf{Common failures}:

\begin{itemize}
\item
  Trailing bytes after the signed region --- possible Incremental Saving
  Attack (Chapter 24)
\item
  ByteRange doesn't start at 0 --- broken signing implementation
\item
  Gap contains unexpected data --- corrupted or manually edited PDF
\end{itemize}

Chapter 13 covers this step in exhaustive detail.

\subsection{Step 2: Hash Verification (Digest
Match)}\label{_step_2_hash_verification_digest_match}

\textbf{Question answered}: "Does the hash of the signed bytes match
what the CMS claims was signed?"

\textbf{What to check}:

\begin{enumerate}
\def\labelenumi{\arabic{enumi}.}
\item
  Read the digest algorithm from the CMS
  \texttt{SignerInfo.digestAlgorithm}
\item
  Compute the hash over the byte range segments:
  \texttt{hash\ =\ algo(segment1\ \textbar{}\textbar{}\ segment2)}
\item
  Extract the \texttt{messageDigest} attribute from the CMS
  \texttt{signedAttrs}
\item
  Compare the computed hash to the CMS-provided hash
\end{enumerate}

\textbf{Common failures}:

\begin{itemize}
\item
  Hash algorithm mismatch --- you computed SHA-256 but the CMS used
  SHA-384
\item
  Byte range computation error --- you read the wrong bytes
\item
  Modified file --- someone changed the file after signing
\end{itemize}

\begin{Shaded}
\begin{Highlighting}[]
\CommentTok{// Step 2: Compute and compare hash}
\BuiltInTok{String}\NormalTok{ digestAlgOid }\OperatorTok{=}\NormalTok{ signerInfo}\OperatorTok{.}\FunctionTok{getDigestAlgOID}\OperatorTok{();}
\BuiltInTok{String}\NormalTok{ algName }\OperatorTok{=} \FunctionTok{mapOidToAlgorithm}\OperatorTok{(}\NormalTok{digestAlgOid}\OperatorTok{);}
\BuiltInTok{MessageDigest}\NormalTok{ digest }\OperatorTok{=} \BuiltInTok{MessageDigest}\OperatorTok{.}\FunctionTok{getInstance}\OperatorTok{(}\NormalTok{algName}\OperatorTok{);}

\DataTypeTok{int}\OperatorTok{[]}\NormalTok{ byteRange }\OperatorTok{=}\NormalTok{ signature}\OperatorTok{.}\FunctionTok{getByteRange}\OperatorTok{();}
\NormalTok{digest}\OperatorTok{.}\FunctionTok{update}\OperatorTok{(}\NormalTok{pdfBytes}\OperatorTok{,}\NormalTok{ byteRange}\OperatorTok{[}\DecValTok{0}\OperatorTok{],}\NormalTok{ byteRange}\OperatorTok{[}\DecValTok{1}\OperatorTok{]);}
\NormalTok{digest}\OperatorTok{.}\FunctionTok{update}\OperatorTok{(}\NormalTok{pdfBytes}\OperatorTok{,}\NormalTok{ byteRange}\OperatorTok{[}\DecValTok{2}\OperatorTok{],}\NormalTok{ byteRange}\OperatorTok{[}\DecValTok{3}\OperatorTok{]);}
\DataTypeTok{byte}\OperatorTok{[]}\NormalTok{ computedHash }\OperatorTok{=}\NormalTok{ digest}\OperatorTok{.}\FunctionTok{digest}\OperatorTok{();}

\BuiltInTok{Attribute}\NormalTok{ digestAttr }\OperatorTok{=}\NormalTok{ signerInfo}\OperatorTok{.}\FunctionTok{getSignedAttributes}\OperatorTok{()}
        \OperatorTok{.}\FunctionTok{get}\OperatorTok{(}\NormalTok{CMSAttributes}\OperatorTok{.}\FunctionTok{messageDigest}\OperatorTok{);}
\DataTypeTok{byte}\OperatorTok{[]}\NormalTok{ claimedHash }\OperatorTok{=}\NormalTok{ ASN1OctetString}\OperatorTok{.}\FunctionTok{getInstance}\OperatorTok{(}
\NormalTok{        digestAttr}\OperatorTok{.}\FunctionTok{getAttrValues}\OperatorTok{().}\FunctionTok{getObjectAt}\OperatorTok{(}\DecValTok{0}\OperatorTok{)).}\FunctionTok{getOctets}\OperatorTok{();}

\ControlFlowTok{if} \OperatorTok{(!}\BuiltInTok{MessageDigest}\OperatorTok{.}\FunctionTok{isEqual}\OperatorTok{(}\NormalTok{computedHash}\OperatorTok{,}\NormalTok{ claimedHash}\OperatorTok{))} \OperatorTok{\{}
\NormalTok{    log}\OperatorTok{.}\FunctionTok{error}\OperatorTok{(}\StringTok{"Hash mismatch — document was modified after signing"}\OperatorTok{);}
    \ControlFlowTok{return}\NormalTok{ ValidationResult}\OperatorTok{.}\FunctionTok{HASH\_MISMATCH}\OperatorTok{;}
\OperatorTok{\}}
\end{Highlighting}
\end{Shaded}

\subsection{Step 3: Cryptographic Verification (CMS
Signature)}\label{_step_3_cryptographic_verification_cms_signature}

\textbf{Question answered}: "Was the CMS \texttt{SignedData} actually
created by the claimed signer's private key?"

\textbf{What to check}:

\begin{enumerate}
\def\labelenumi{\arabic{enumi}.}
\item
  Extract the signer's certificate from the CMS
  \texttt{SignedData.certificates} collection
\item
  Verify that the \texttt{SignerInfo.signatureAlgorithm} matches the
  certificate's key type (RSA cert with RSA signature, ECDSA cert with
  ECDSA signature)
\item
  Verify the signature over the \texttt{signedAttrs} using the signer's
  public key
\item
  Verify that the signed attributes include the required fields:
  \texttt{contentType}, \texttt{messageDigest}, and (for PAdES)
  \texttt{ESS-signing-certificate-v2}
\end{enumerate}

\textbf{Common failures}:

\begin{itemize}
\item
  Wrong certificate used for verification --- the CMS contains multiple
  certificates, and you picked an intermediate instead of the signer
\item
  Algorithm mismatch --- rare, but can happen with misconfigured
  BouncyCastle providers
\item
  Missing \texttt{ESS-signing-certificate-v2} --- not a cryptographic
  failure, but a PAdES compliance failure
\end{itemize}

\begin{Shaded}
\begin{Highlighting}[]
\CommentTok{// Step 3: CMS cryptographic verification}
\NormalTok{CMSSignedData cmsData }\OperatorTok{=} \KeywordTok{new} \FunctionTok{CMSSignedData}\OperatorTok{(}
        \KeywordTok{new} \FunctionTok{CMSProcessableByteArray}\OperatorTok{(}\NormalTok{signedBytes}\OperatorTok{),}\NormalTok{ signature}\OperatorTok{.}\FunctionTok{getContents}\OperatorTok{());}

\NormalTok{SignerInformationStore signerStore }\OperatorTok{=}\NormalTok{ cmsData}\OperatorTok{.}\FunctionTok{getSignerInfos}\OperatorTok{();}
\NormalTok{SignerInformation signerInfo }\OperatorTok{=}\NormalTok{ signerStore}\OperatorTok{.}\FunctionTok{getSigners}\OperatorTok{().}\FunctionTok{iterator}\OperatorTok{().}\FunctionTok{next}\OperatorTok{();}

\CommentTok{// Find the signer\textquotesingle{}s certificate}
\NormalTok{X509CertificateHolder signerCert }\OperatorTok{=} \FunctionTok{findSignerCertificate}\OperatorTok{(}
\NormalTok{        cmsData}\OperatorTok{.}\FunctionTok{getCertificates}\OperatorTok{(),}\NormalTok{ signerInfo}\OperatorTok{.}\FunctionTok{getSID}\OperatorTok{());}

\CommentTok{// Verify the CMS signature}
\DataTypeTok{boolean}\NormalTok{ cmsValid }\OperatorTok{=}\NormalTok{ signerInfo}\OperatorTok{.}\FunctionTok{verify}\OperatorTok{(}
        \KeywordTok{new} \FunctionTok{JcaSimpleSignerInfoVerifierBuilder}\OperatorTok{()}
                \OperatorTok{.}\FunctionTok{setProvider}\OperatorTok{(}\StringTok{"BC"}\OperatorTok{)}
                \OperatorTok{.}\FunctionTok{build}\OperatorTok{(}\NormalTok{signerCert}\OperatorTok{));}

\ControlFlowTok{if} \OperatorTok{(!}\NormalTok{cmsValid}\OperatorTok{)} \OperatorTok{\{}
\NormalTok{    log}\OperatorTok{.}\FunctionTok{error}\OperatorTok{(}\StringTok{"CMS signature verification failed — "} \OperatorTok{+}
              \StringTok{"signature was not created by the claimed certificate"}\OperatorTok{);}
    \ControlFlowTok{return}\NormalTok{ ValidationResult}\OperatorTok{.}\FunctionTok{CMS\_INVALID}\OperatorTok{;}
\OperatorTok{\}}
\end{Highlighting}
\end{Shaded}

\subsection{Step 4: Certificate Chain
Validation}\label{_step_4_certificate_chain_validation}

\textbf{Question answered}: "Is the signer's certificate trustworthy?
Was it valid at the time of signing?"

\textbf{What to check}:

\begin{enumerate}
\def\labelenumi{\arabic{enumi}.}
\item
  Build the certificate chain from the signer's certificate up to a
  trusted root
\item
  Verify each certificate's signature (each cert is signed by its
  issuer)
\item
  Check validity periods --- was each certificate valid at the signing
  time?
\item
  Check revocation status --- was any certificate revoked at or before
  the signing time?
\item
  Check key usage --- does the signer's certificate allow digital
  signatures (\texttt{digitalSignature} or \texttt{nonRepudiation} key
  usage)?
\item
  Check basic constraints --- are intermediate certificates properly
  marked as CA certificates?
\end{enumerate}

\textbf{Common failures}:

\begin{itemize}
\item
  \textbf{Self-signed certificate}: The signer's certificate is not
  issued by a recognized CA. This is technically valid (the CMS
  signature will verify) but the certificate is not trusted. Most
  viewers show "signature is valid but the identity of the signer could
  not be verified."
\item
  \textbf{Expired certificate}: The certificate was valid when the
  signature was created, but has since expired. With timestamps, this is
  not a problem (the timestamp proves the signature was created before
  expiration). Without timestamps, viewers may warn about an expired
  certificate.
\item
  \textbf{Revoked certificate}: A certificate in the chain was revoked
  (compromised, or the CA issued it in error). Revocation checking
  requires OCSP or CRL lookup --- which may fail if the revocation
  servers are unreachable.
\item
  \textbf{Missing intermediate}: The CMS doesn't include all
  intermediate certificates, and the validator can't build the chain.
  This is the most common real-world certificate chain failure.
\item
  \textbf{Not in trust store}: The root certificate is not in the
  validator's trust store. Different validators have different trust
  stores --- Adobe uses AATL, operating systems use their own, and EU
  validators use the EUTL.
\end{itemize}

\subsection{Step 5: Trust Evaluation}\label{_step_5_trust_evaluation}

\textbf{Question answered}: "Do I trust the entity that signed this
document?"

This is the most subjective step, and it's where validators diverge the
most. Trust evaluation is not a purely technical question --- it depends
on context:

\textbf{Adobe Acrobat trust model}: Acrobat trusts certificates that
chain to a root in the \textbf{Adobe Approved Trust List (AATL)}. If the
certificate chains to an AATL root, Acrobat shows a green checkmark. If
not --- even if the certificate is perfectly valid, issued by a
reputable CA, and chains to a root in the operating system's trust store
--- Acrobat shows a yellow warning: "At least one signature has
problems."

This is one of the most frustrating aspects of PDF signature validation
for developers. Your signature is cryptographically valid, the
certificate is issued by a recognized CA, the certificate is not expired
or revoked --- but Acrobat shows a warning because the CA is not on
Adobe's list. Getting on the AATL requires the CA to apply to Adobe and
pay fees.

\textbf{EU validation model}: EU validators (like the DSS demonstration
webapp, or validators used by EU government services) trust certificates
that chain to a root in the \textbf{EU Trusted List (EUTL)}. The EUTL is
the authoritative registry of Qualified Trust Service Providers. A
signature is a "qualified electronic signature" --- with the legal
weight of a handwritten signature under eIDAS --- only if the
certificate chains to an EUTL-listed QTSP.

\textbf{Operating system trust model}: Some applications (like
LibreOffice) use the operating system's certificate store for trust
evaluation. This means the set of trusted roots varies by OS and by the
user's system configuration.

\textbf{Custom trust model}: Your application can define its own trust
store --- a specific set of root certificates that you trust for your
use case. This is common in enterprise environments where the signing
certificates are issued by an internal CA.

\textbf{When debugging trust failures, ask three questions:}

\begin{enumerate}
\def\labelenumi{\arabic{enumi}.}
\item
  What trust store is the validator using? (AATL, EUTL, OS store,
  custom)
\item
  Is the signer's root certificate in that trust store?
\item
  Are all intermediates available? (Either embedded in the CMS or
  fetchable via the certificate's Authority Information Access
  extension)
\end{enumerate}

The answer to question 1 determines everything else. A signature can be
"valid" in one validator and "untrusted" in another, purely because of
trust store differences.

\section{Validation Differences Across
Viewers}\label{_validation_differences_across_viewers}

One of the most confusing aspects of PDF signature validation is that
different viewers produce different results for the same PDF. This is
not a bug --- it's a consequence of the five-step validation model,
where each viewer makes different choices at each step.

\begin{longtable}[]{@{}
  >{\raggedright\arraybackslash}p{(\linewidth - 10\tabcolsep) * \real{0.2857}}
  >{\raggedright\arraybackslash}p{(\linewidth - 10\tabcolsep) * \real{0.1429}}
  >{\raggedright\arraybackslash}p{(\linewidth - 10\tabcolsep) * \real{0.1429}}
  >{\raggedright\arraybackslash}p{(\linewidth - 10\tabcolsep) * \real{0.1429}}
  >{\raggedright\arraybackslash}p{(\linewidth - 10\tabcolsep) * \real{0.1429}}
  >{\raggedright\arraybackslash}p{(\linewidth - 10\tabcolsep) * \real{0.1429}}@{}}
\caption{Validation behavior comparison across major PDF
viewers}\tabularnewline
\toprule\noalign{}
\begin{minipage}[b]{\linewidth}\raggedright
Check
\end{minipage} & \begin{minipage}[b]{\linewidth}\raggedright
Acrobat
\end{minipage} & \begin{minipage}[b]{\linewidth}\raggedright
Foxit
\end{minipage} & \begin{minipage}[b]{\linewidth}\raggedright
LibreOffice
\end{minipage} & \begin{minipage}[b]{\linewidth}\raggedright
PDFBox
\end{minipage} & \begin{minipage}[b]{\linewidth}\raggedright
EU DSS
\end{minipage} \\
\midrule\noalign{}
\endfirsthead
\toprule\noalign{}
\begin{minipage}[b]{\linewidth}\raggedright
Check
\end{minipage} & \begin{minipage}[b]{\linewidth}\raggedright
Acrobat
\end{minipage} & \begin{minipage}[b]{\linewidth}\raggedright
Foxit
\end{minipage} & \begin{minipage}[b]{\linewidth}\raggedright
LibreOffice
\end{minipage} & \begin{minipage}[b]{\linewidth}\raggedright
PDFBox
\end{minipage} & \begin{minipage}[b]{\linewidth}\raggedright
EU DSS
\end{minipage} \\
\midrule\noalign{}
\endhead
\bottomrule\noalign{}
\endlastfoot
Byte range trailing bytes & Warns & Warns & Partial & Manual & Warns \\
Hash verification & Yes & Yes & Yes & Manual & Yes \\
CMS signature & Yes & Yes & Yes & Manual & Yes \\
Certificate chain & AATL & Mixed & OS store & Manual & EUTL \\
Revocation (OCSP) & Yes & Yes & Limited & Manual & Yes \\
Revocation (CRL) & Yes & Yes & Yes & Manual & Yes \\
MDP enforcement & Full & Mostly & Partial & None & Full \\
Timestamp verification & Yes & Yes & Partial & Manual & Yes \\
PAdES compliance & Partial & Partial & No & Manual & Full \\
\end{longtable}

"Manual" for PDFBox means the library provides the primitives (signature
extraction, byte range access, CMS parsing) but does not perform the
validation automatically --- you must write the validation logic
yourself. This is by design: PDFBox is a PDF library, not a validation
service.

EU DSS provides the most complete automated validation, including
ETSI-compliant validation reports.

\section{Debugging Validation
Failures}\label{_debugging_validation_failures}

\subsection{The Systematic Approach}\label{_the_systematic_approach}

When a signature fails validation, resist the urge to guess. Instead,
walk through the five steps systematically:

\textbf{Step 1: Is the byte range intact?}

\begin{Shaded}
\begin{Highlighting}[]
\CommentTok{\# Check byte range with QPDF}
\ExtensionTok{qpdf} \AttributeTok{{-}{-}show{-}object}\OperatorTok{=\textless{}}\NormalTok{sig{-}obj{-}number}\OperatorTok{\textgreater{}}\NormalTok{ document.pdf}
\CommentTok{\# Look at /ByteRange — does offset2 + length2 == file size?}
\FunctionTok{wc} \AttributeTok{{-}c}\NormalTok{ document.pdf}
\end{Highlighting}
\end{Shaded}

\textbf{Step 2: Does the hash match?}

\begin{Shaded}
\begin{Highlighting}[]
\CommentTok{// Compute the hash yourself and compare to the CMS messageDigest}
\CommentTok{// (Code in Chapter 13)}
\end{Highlighting}
\end{Shaded}

\textbf{Step 3: Does the CMS verify?}

\begin{Shaded}
\begin{Highlighting}[]
\CommentTok{\# Extract /Contents, decode hex, verify with OpenSSL}
\ExtensionTok{openssl}\NormalTok{ cms }\AttributeTok{{-}inform}\NormalTok{ DER }\AttributeTok{{-}in}\NormalTok{ signature.der }\AttributeTok{{-}verify} \DataTypeTok{\textbackslash{}}
    \AttributeTok{{-}content}\NormalTok{ signed{-}bytes.bin }\AttributeTok{{-}noverify} \AttributeTok{{-}out}\NormalTok{ /dev/null}
\CommentTok{\# {-}noverify skips certificate verification — tests only the CMS math}
\end{Highlighting}
\end{Shaded}

\textbf{Step 4: Is the certificate chain valid?}

\begin{Shaded}
\begin{Highlighting}[]
\CommentTok{\# Extract the signer\textquotesingle{}s certificate and verify the chain}
\ExtensionTok{openssl}\NormalTok{ cms }\AttributeTok{{-}inform}\NormalTok{ DER }\AttributeTok{{-}in}\NormalTok{ signature.der }\AttributeTok{{-}noout} \AttributeTok{{-}cmsout} \AttributeTok{{-}print} \KeywordTok{|} \DataTypeTok{\textbackslash{}}
    \FunctionTok{grep} \AttributeTok{{-}A}\NormalTok{ 5 }\StringTok{"cert:"}
\ExtensionTok{openssl}\NormalTok{ verify }\AttributeTok{{-}CAfile}\NormalTok{ trusted{-}roots.pem signer{-}cert.pem}
\end{Highlighting}
\end{Shaded}

\textbf{Step 5: Is the certificate trusted?}

Check the signer's root against the relevant trust store (AATL list,
EUTL, OS store).

\subsection{Common Error Messages and Their
Meaning}\label{_common_error_messages_and_their_meaning}

\begin{longtable}[]{@{}
  >{\raggedright\arraybackslash}p{(\linewidth - 4\tabcolsep) * \real{0.3333}}
  >{\raggedright\arraybackslash}p{(\linewidth - 4\tabcolsep) * \real{0.2222}}
  >{\raggedright\arraybackslash}p{(\linewidth - 4\tabcolsep) * \real{0.4444}}@{}}
\caption{PDF signature error messages decoded}\tabularnewline
\toprule\noalign{}
\begin{minipage}[b]{\linewidth}\raggedright
Error Message (Viewer)
\end{minipage} & \begin{minipage}[b]{\linewidth}\raggedright
Step
\end{minipage} & \begin{minipage}[b]{\linewidth}\raggedright
What It Actually Means
\end{minipage} \\
\midrule\noalign{}
\endfirsthead
\toprule\noalign{}
\begin{minipage}[b]{\linewidth}\raggedright
Error Message (Viewer)
\end{minipage} & \begin{minipage}[b]{\linewidth}\raggedright
Step
\end{minipage} & \begin{minipage}[b]{\linewidth}\raggedright
What It Actually Means
\end{minipage} \\
\midrule\noalign{}
\endhead
\bottomrule\noalign{}
\endlastfoot
"Signature is INVALID" (Acrobat) & 2 or 3 & Hash mismatch or CMS
verification failure. The document was modified after signing, or the
CMS is malformed. \\
"At least one signature has problems" (Acrobat) & 5 & The certificate
does not chain to an AATL root. The signature may be cryptographically
perfect --- Acrobat just doesn't trust the CA. \\
"Document has been altered or corrupted since the Signature was applied"
(Acrobat) & 1 or MDP & Either the byte range doesn't cover the full file
(trailing bytes), or modifications after signing exceed the MDP
permissions. \\
"Signature is valid, but the identity of the signer could not be
verified" (Acrobat) & 4 or 5 & CMS is valid, but the certificate chain
is incomplete or the root is not trusted. \\
"The signature cannot be verified" (LibreOffice) & 3, 4, or 5 & Generic
error --- LibreOffice doesn't distinguish well between CMS failure and
trust failure. Check the signature details panel. \\
\texttt{CMSException:\ message-digest\ attribute\ value\ does\ not\ match\ calculated\ value}
(BouncyCastle) & 2 & Hash mismatch. The bytes you passed to the verifier
don't match what was signed. \\
\texttt{CertPathValidatorException:\ validity\ check\ failed} (Java) & 4
& A certificate in the chain has expired. \\
\texttt{CertPathValidatorException:\ no\ valid\ CRL\ found} (Java) & 4 &
Revocation checking failed because the CRL distribution point is
unreachable or the CRL is not cached. \\
\end{longtable}

\section{Timestamp Verification}\label{_timestamp_verification}

\subsection{Why Timestamps Matter for
Validation}\label{_why_timestamps_matter_for_validation}

A digital signature proves that a specific certificate holder signed
specific bytes. But it doesn't prove \emph{when} the signing happened.
Without a timestamp, you can only say "the signature was created
sometime before now."

This matters because certificates expire and can be revoked. If a
certificate was valid on March 1 but expired on March 31, and you
validate the signature on April 15:

\begin{itemize}
\item
  \textbf{Without timestamp}: The certificate is expired. Some
  validators will flag the signature as questionable.
\item
  \textbf{With timestamp}: The timestamp proves the signature was
  created on March 1, when the certificate was valid. The signature
  remains valid despite the certificate's expiration.
\end{itemize}

\subsection{Verifying the Timestamp}\label{_verifying_the_timestamp}

A signature timestamp is embedded as an unsigned attribute
(\texttt{id-aa-signatureTimeStampToken}, OID 1.2.840.113549.1.9.16.2.14)
in the CMS \texttt{SignerInfo}:

\begin{Shaded}
\begin{Highlighting}[]
\CommentTok{// Extract and verify the signature timestamp}
\NormalTok{AttributeTable unsignedAttrs }\OperatorTok{=}\NormalTok{ signerInfo}\OperatorTok{.}\FunctionTok{getUnsignedAttributes}\OperatorTok{();}
\ControlFlowTok{if} \OperatorTok{(}\NormalTok{unsignedAttrs }\OperatorTok{!=} \KeywordTok{null}\OperatorTok{)} \OperatorTok{\{}
    \BuiltInTok{Attribute}\NormalTok{ tsAttr }\OperatorTok{=}\NormalTok{ unsignedAttrs}\OperatorTok{.}\FunctionTok{get}\OperatorTok{(}
\NormalTok{            PKCSObjectIdentifiers}\OperatorTok{.}\FunctionTok{id\_aa\_signatureTimeStampToken}\OperatorTok{);}
    \ControlFlowTok{if} \OperatorTok{(}\NormalTok{tsAttr }\OperatorTok{!=} \KeywordTok{null}\OperatorTok{)} \OperatorTok{\{}
\NormalTok{        CMSSignedData tsToken }\OperatorTok{=} \KeywordTok{new} \FunctionTok{CMSSignedData}\OperatorTok{(}
\NormalTok{                tsAttr}\OperatorTok{.}\FunctionTok{getAttrValues}\OperatorTok{().}\FunctionTok{getObjectAt}\OperatorTok{(}\DecValTok{0}\OperatorTok{).}\FunctionTok{toASN1Primitive}\OperatorTok{().}\FunctionTok{getEncoded}\OperatorTok{());}

        \CommentTok{// Verify the timestamp\textquotesingle{}s CMS signature}
\NormalTok{        SignerInformation tsSigner }\OperatorTok{=}\NormalTok{ tsToken}\OperatorTok{.}\FunctionTok{getSignerInfos}\OperatorTok{()}
                \OperatorTok{.}\FunctionTok{getSigners}\OperatorTok{().}\FunctionTok{iterator}\OperatorTok{().}\FunctionTok{next}\OperatorTok{();}
\NormalTok{        X509CertificateHolder tsCert }\OperatorTok{=} \FunctionTok{findSignerCertificate}\OperatorTok{(}
\NormalTok{                tsToken}\OperatorTok{.}\FunctionTok{getCertificates}\OperatorTok{(),}\NormalTok{ tsSigner}\OperatorTok{.}\FunctionTok{getSID}\OperatorTok{());}

        \DataTypeTok{boolean}\NormalTok{ tsValid }\OperatorTok{=}\NormalTok{ tsSigner}\OperatorTok{.}\FunctionTok{verify}\OperatorTok{(}
                \KeywordTok{new} \FunctionTok{JcaSimpleSignerInfoVerifierBuilder}\OperatorTok{()}
                        \OperatorTok{.}\FunctionTok{setProvider}\OperatorTok{(}\StringTok{"BC"}\OperatorTok{).}\FunctionTok{build}\OperatorTok{(}\NormalTok{tsCert}\OperatorTok{));}

        \ControlFlowTok{if} \OperatorTok{(}\NormalTok{tsValid}\OperatorTok{)} \OperatorTok{\{}
            \CommentTok{// Extract the timestamp time}
\NormalTok{            TimeStampToken tst }\OperatorTok{=} \KeywordTok{new} \FunctionTok{TimeStampToken}\OperatorTok{(}\NormalTok{tsToken}\OperatorTok{);}
            \BuiltInTok{Date}\NormalTok{ signingTime }\OperatorTok{=}\NormalTok{ tst}\OperatorTok{.}\FunctionTok{getTimeStampInfo}\OperatorTok{().}\FunctionTok{getGenTime}\OperatorTok{();}
\NormalTok{            log}\OperatorTok{.}\FunctionTok{info}\OperatorTok{(}\StringTok{"Signature was timestamped at: \{\}"}\OperatorTok{,}\NormalTok{ signingTime}\OperatorTok{);}

            \CommentTok{// Verify that the timestamp covers the signature value}
            \DataTypeTok{byte}\OperatorTok{[]}\NormalTok{ tsDigest }\OperatorTok{=}\NormalTok{ tst}\OperatorTok{.}\FunctionTok{getTimeStampInfo}\OperatorTok{()}
                    \OperatorTok{.}\FunctionTok{getMessageImprintDigest}\OperatorTok{();}
            \CommentTok{// This digest should match the hash of the signature value}
        \OperatorTok{\}}
    \OperatorTok{\}}
\OperatorTok{\}}
\end{Highlighting}
\end{Shaded}

\subsection{Timestamp Chain of Trust}\label{_timestamp_chain_of_trust}

The timestamp itself is a CMS signature --- signed by the TSA's
certificate. This certificate has its own chain and trust requirements:

\begin{enumerate}
\def\labelenumi{\arabic{enumi}.}
\item
  The TSA certificate must have the \texttt{id-kp-timeStamping} extended
  key usage (OID 1.3.6.1.5.5.7.3.8)
\item
  The TSA certificate chain must be valid at the time of verification
\item
  The TSA's root must be in your trust store
\end{enumerate}

If the TSA's certificate has also expired by the time you validate
(years later), you need \textbf{long-term validation (LTV)} --- which is
the subject of Chapter 18.

\section{Building a Complete Validation
Pipeline}\label{_building_a_complete_validation_pipeline}

\subsection{Architecture}\label{_architecture}

A production validation pipeline typically follows this structure:

\begin{verbatim}
Input: PDF file
  │
  ├─→ Step 1: Extract all signatures
  │     └─→ For each signature:
  │           ├─→ Step 2: Byte range check
  │           ├─→ Step 3: Hash verification
  │           ├─→ Step 4: CMS verification
  │           ├─→ Step 5: Certificate chain
  │           ├─→ Step 6: Revocation check
  │           ├─→ Step 7: Timestamp verification (if present)
  │           ├─→ Step 8: MDP compliance (if certification sig)
  │           └─→ Step 9: Trust evaluation
  │
  └─→ Output: Validation report
        ├─→ Per-signature results
        ├─→ Overall document status
        └─→ Detailed diagnostics
\end{verbatim}

Chapter 31 provides a complete Java implementation of this pipeline.

\subsection{When to Use EU DSS
Instead}\label{_when_to_use_eu_dss_instead}

If you need to produce \textbf{ETSI-compliant validation reports} (ETSI
TS 119 102-2) --- for example, because you're building a validation
service for a government or regulated industry --- use the EU DSS
framework (Chapter 23) instead of building the pipeline yourself. EU DSS
handles all validation steps, produces standards-compliant reports, and
has been tested against the full range of PAdES profiles and edge cases.

Building your own pipeline is appropriate when:

\begin{itemize}
\item
  You need custom trust evaluation logic
\item
  You need to validate signatures in a specific context (e.g., your own
  application's signing workflow)
\item
  You want to understand every step of the validation process (which is
  the point of this book)
\item
  You're building a specialized tool (forensic analysis, security
  research)
\end{itemize}

\section{Key Takeaways}\label{_key_takeaways_4}

\begin{itemize}
\item
  Signature validation is a \textbf{five-step pipeline}: byte range →
  hash → CMS → certificate chain → trust
\item
  Each step is independent --- a failure at any step means the signature
  is invalid or questionable
\item
  \textbf{Different viewers produce different results} because they use
  different trust stores and enforce different rules
\item
  Adobe Acrobat uses \textbf{AATL} for trust; EU validators use
  \textbf{EUTL}; LibreOffice uses the \textbf{OS trust store}
\item
  The most common real-world failure is \textbf{trust evaluation} ---
  the signature is cryptographically perfect but the root CA isn't in
  the validator's trust store
\item
  \textbf{Timestamps} prove when a signature was created --- essential
  for handling expired certificates and building long-term validation
\item
  \textbf{Debug systematically}: walk through the five steps in order,
  using QPDF, OpenSSL, and BouncyCastle to isolate the failure
\item
  \textbf{EU DSS} provides the most complete automated validation with
  ETSI-compliant reports
\item
  PDFBox provides the \textbf{primitives} for validation but does not
  automate the pipeline --- you build it yourself
\end{itemize}

\chapter{PAdES Profiles (ETSI EN 319
142)}\label{_pades_profiles_etsi_en_319_142}

ISO 32000-2 defines \emph{how} a digital signature is embedded in a PDF.
PAdES defines \emph{what} that signature must contain to be legally
meaningful in Europe --- and, increasingly, everywhere else.

PAdES --- PDF Advanced Electronic Signatures --- is a set of standards
published by the \textbf{European Telecommunications Standards Institute
(ETSI)} that constrain PDF signatures to meet specific compliance
requirements. A PAdES signature is always a valid ISO 32000-2 signature,
but the reverse is not true. PAdES adds mandatory signed attributes,
requires specific SubFilter values, mandates timestamp handling, and
defines four progressively more complete signature profiles.

Understanding PAdES is not optional if your signatures need to be
validated by European government services, if your customers require
eIDAS-compliant signatures, or if you want your signatures to remain
verifiable for years or decades. Even if you're not operating in the EU,
PAdES profiles have become the de facto standard for "properly
structured" PDF signatures worldwide --- Adobe Acrobat increasingly
favors PAdES-structured signatures in its validation logic.

This chapter explains each PAdES profile, what it requires, how the
profiles relate to each other, and how to produce each one.

\section{The Four PAdES Baseline
Profiles}\label{_the_four_pades_baseline_profiles}

PAdES defines four \textbf{baseline profiles}, each building on the
previous one. They are formally named PAdES-B-B, PAdES-B-T, PAdES-B-LT,
and PAdES-B-LTA. The naming follows the pattern:
PAdES-B(aseline)-\textless level\textgreater.

\begin{longtable}[]{@{}
  >{\raggedright\arraybackslash}p{(\linewidth - 4\tabcolsep) * \real{0.2000}}
  >{\raggedright\arraybackslash}p{(\linewidth - 4\tabcolsep) * \real{0.2000}}
  >{\raggedright\arraybackslash}p{(\linewidth - 4\tabcolsep) * \real{0.6000}}@{}}
\caption{PAdES Baseline Profile Summary}\tabularnewline
\toprule\noalign{}
\begin{minipage}[b]{\linewidth}\raggedright
Profile
\end{minipage} & \begin{minipage}[b]{\linewidth}\raggedright
Full Name
\end{minipage} & \begin{minipage}[b]{\linewidth}\raggedright
What It Adds
\end{minipage} \\
\midrule\noalign{}
\endfirsthead
\toprule\noalign{}
\begin{minipage}[b]{\linewidth}\raggedright
Profile
\end{minipage} & \begin{minipage}[b]{\linewidth}\raggedright
Full Name
\end{minipage} & \begin{minipage}[b]{\linewidth}\raggedright
What It Adds
\end{minipage} \\
\midrule\noalign{}
\endhead
\bottomrule\noalign{}
\endlastfoot
\textbf{B-B} & Baseline --- Basic & The minimum PAdES signature. CMS
\texttt{SignedData} with mandatory signed attributes including
\texttt{ESS-signing-certificate-v2}. The signature itself --- no
timestamps, no revocation data. \\
\textbf{B-T} & Baseline --- Timestamp & B-B + a \textbf{signature
timestamp} from a trusted TSA. Proves when the signature was created.
The timestamp is an unsigned CMS attribute
(\texttt{id-aa-signatureTimeStampToken}). \\
\textbf{B-LT} & Baseline --- Long-Term & B-T + embedded
\textbf{revocation data} (OCSP responses and/or CRLs) and the full
certificate chain in the PDF's \textbf{DSS} (Document Security Store)
dictionary. Everything needed to validate the signature offline, without
contacting any external servers. \\
\textbf{B-LTA} & Baseline --- Long-Term with Archive timestamp & B-LT +
a \textbf{document timestamp} signature that covers the entire file
(including the DSS dictionary). Enables re-timestamping decades later,
creating an indefinite chain of trust that outlives any individual
certificate. \\
\end{longtable}

\subsection{The Progression: Why Four
Levels?}\label{_the_progression_why_four_levels}

Each profile addresses a specific failure mode:

\begin{itemize}
\item
  \textbf{B-B} fails when: the signing certificate expires (validators
  may warn), or when you need to prove \emph{when} the signature was
  created (there's no timestamp)
\item
  \textbf{B-T} fails when: the OCSP responder or CRL distribution point
  goes offline years later, and you can't check revocation status
\item
  \textbf{B-LT} fails when: the timestamp's certificate, or the signing
  certificate's root CA, expires or is removed from trust stores decades
  later
\item
  \textbf{B-LTA} handles all of these by providing a renewable chain of
  timestamps
\end{itemize}

In practice:

\begin{itemize}
\item
  \textbf{B-B} is sufficient for short-lived documents that will be
  validated within the signing certificate's validity period (typically
  1-3 years)
\item
  \textbf{B-T} is the minimum for any document with legal significance
  --- you need to prove when it was signed
\item
  \textbf{B-LT} is required for long-term archival (10+ years) where you
  can't assume external services will remain available
\item
  \textbf{B-LTA} is the gold standard for documents that must be
  verifiable indefinitely --- government records, notarized documents,
  regulatory filings
\end{itemize}

\section{PAdES-B-B: The Baseline
Signature}\label{_pades_b_b_the_baseline_signature}

\subsection{Requirements}\label{_requirements}

A PAdES-B-B signature must include:

\begin{enumerate}
\def\labelenumi{\arabic{enumi}.}
\item
  \textbf{SubFilter}: Must be \texttt{ETSI.CAdES.detached} (not
  \texttt{adbe.pkcs7.detached})
\item
  \textbf{Mandatory signed attributes in CMS}:

  \begin{itemize}
  \item
    \texttt{content-type} (OID 1.2.840.113549.1.9.3) --- must be
    \texttt{id-data}
  \item
    \texttt{message-digest} (OID 1.2.840.113549.1.9.4) --- the hash of
    the signed bytes
  \item
    \texttt{ESS-signing-certificate-v2} (OID 1.2.840.113549.1.9.16.2.47)
    --- identifies the signing certificate by its hash, preventing
    certificate substitution attacks
  \end{itemize}
\item
  \textbf{Signing certificate embedded in CMS}: The signer's certificate
  (and ideally the full chain up to but not including the root) must be
  embedded in the CMS \texttt{SignedData.certificates} collection
\item
  \textbf{Digest algorithm}: Must be SHA-256 or stronger (SHA-384,
  SHA-512). SHA-1 is \textbf{not} PAdES-compliant.
\end{enumerate}

\subsection{The ESS-signing-certificate-v2
Attribute}\label{_the_ess_signing_certificate_v2_attribute_2}

This attribute is the key difference between a basic ISO 32000-2
signature and a PAdES-B-B signature. It binds the CMS signature to a
specific certificate by including a hash of the certificate in the
signed attributes.

Why this matters: without \texttt{ESS-signing-certificate-v2}, an
attacker could theoretically substitute a different certificate into the
CMS \texttt{certificates} collection --- one with the same public key
but different identity information. The CMS signature would still verify
(same key), but the identity would be wrong. The ESS attribute prevents
this by cryptographically committing to a specific certificate.

\begin{Shaded}
\begin{Highlighting}[]
\CommentTok{// Creating the ESS{-}signing{-}certificate{-}v2 attribute with BouncyCastle}
\BuiltInTok{MessageDigest}\NormalTok{ certDigest }\OperatorTok{=} \BuiltInTok{MessageDigest}\OperatorTok{.}\FunctionTok{getInstance}\OperatorTok{(}\StringTok{"SHA{-}256"}\OperatorTok{);}
\DataTypeTok{byte}\OperatorTok{[]}\NormalTok{ certHash }\OperatorTok{=}\NormalTok{ certDigest}\OperatorTok{.}\FunctionTok{digest}\OperatorTok{(}\NormalTok{signingCertificate}\OperatorTok{.}\FunctionTok{getEncoded}\OperatorTok{());}

\NormalTok{ESSCertIDv2 essCertId }\OperatorTok{=} \KeywordTok{new} \FunctionTok{ESSCertIDv2}\OperatorTok{(}
        \KeywordTok{new} \FunctionTok{AlgorithmIdentifier}\OperatorTok{(}\NormalTok{NISTObjectIdentifiers}\OperatorTok{.}\FunctionTok{id\_sha256}\OperatorTok{),}
\NormalTok{        certHash}\OperatorTok{);}

\NormalTok{SigningCertificateV2 signingCertV2 }\OperatorTok{=} \KeywordTok{new} \FunctionTok{SigningCertificateV2}\OperatorTok{(}
        \KeywordTok{new}\NormalTok{ ESSCertIDv2}\OperatorTok{[]} \OperatorTok{\{}\NormalTok{ essCertId }\OperatorTok{\});}

\CommentTok{// Add to signed attributes}
\NormalTok{CMSAttributeTableGenerator signedAttrGen }\OperatorTok{=} \KeywordTok{new} \FunctionTok{DefaultSignedAttributeTableGenerator}\OperatorTok{(}
        \KeywordTok{new} \FunctionTok{AttributeTable}\OperatorTok{(}\KeywordTok{new} \BuiltInTok{Hashtable}\OperatorTok{\textless{}\textgreater{}()} \OperatorTok{\{\{}
            \FunctionTok{put}\OperatorTok{(}\NormalTok{PKCSObjectIdentifiers}\OperatorTok{.}\FunctionTok{id\_aa\_signingCertificateV2}\OperatorTok{,}
                \KeywordTok{new} \BuiltInTok{Attribute}\OperatorTok{(}\NormalTok{PKCSObjectIdentifiers}\OperatorTok{.}\FunctionTok{id\_aa\_signingCertificateV2}\OperatorTok{,}
                    \KeywordTok{new} \FunctionTok{DERSet}\OperatorTok{(}\NormalTok{signingCertV2}\OperatorTok{)));}
        \OperatorTok{\}\}));}
\end{Highlighting}
\end{Shaded}

A common mistake: using \texttt{ESS-signing-certificate} (v1, OID
1.2.840.113549.1.9.16.2.12) instead of
\texttt{ESS-signing-certificate-v2} (OID 1.2.840.113549.1.9.16.2.47).
The v1 attribute uses SHA-1 for the certificate hash, which is
deprecated. PAdES requires v2, which uses SHA-256 by default.

iText 2.x (and old versions of some other libraries) produced v1
attributes. If you encounter a PDF with \texttt{ESS-signing-certificate}
(v1), it's likely from a legacy signing implementation.

\subsection{\texorpdfstring{SubFilter: \texttt{ETSI.CAdES.detached} vs.
\texttt{adbe.pkcs7.detached}}{SubFilter: ETSI.CAdES.detached vs. adbe.pkcs7.detached}}\label{_subfilter_etsi_cades_detached_vs_adbe_pkcs7_detached}

PAdES requires the SubFilter value \texttt{ETSI.CAdES.detached}. This is
structurally identical to \texttt{adbe.pkcs7.detached} --- both produce
a detached CMS signature embedded in the \texttt{/Contents} hex string.
The difference is semantic: \texttt{ETSI.CAdES.detached} signals that
the signature follows the CAdES (CMS Advanced Electronic Signatures)
profile, which includes the ESS-signing-certificate-v2 attribute and
other CAdES requirements.

In practice, many validators accept both SubFilter values for PAdES
validation. Adobe Acrobat treats \texttt{ETSI.CAdES.detached} and
\texttt{adbe.pkcs7.detached} similarly. But for strict PAdES compliance
--- especially when dealing with EU government validators --- use
\texttt{ETSI.CAdES.detached}.

\section{PAdES-B-T: Adding a
Timestamp}\label{_pades_b_t_adding_a_timestamp}

\subsection{What B-T Adds}\label{_what_b_t_adds}

PAdES-B-T is PAdES-B-B plus a \textbf{signature timestamp} --- a
timestamp token from a Trusted Timestamp Authority (TSA) that proves
when the signature was created.

The timestamp is embedded as an \textbf{unsigned attribute}
(\texttt{id-aa-signatureTimeStampToken}, OID 1.2.840.113549.1.9.16.2.14)
in the CMS \texttt{SignerInfo}. It's unsigned because it's applied
\emph{after} the signature --- you sign first, then request a timestamp
over the signature value.

\subsection{How Timestamping Works}\label{_how_timestamping_works}

The timestamp workflow:

\begin{enumerate}
\def\labelenumi{\arabic{enumi}.}
\item
  \textbf{Sign the document} --- produce the CMS \texttt{SignedData}
  with all signed attributes (B-B level)
\item
  \textbf{Hash the signature value} --- compute SHA-256 over the CMS
  \texttt{SignerInfo.signature} bytes
\item
  \textbf{Send a timestamp request} --- send the hash to a TSA using RFC
  3161 (HTTP POST to the TSA's URL)
\item
  \textbf{Receive the timestamp response} --- the TSA returns a
  \texttt{TimeStampToken} (which is itself a CMS \texttt{SignedData}
  signed by the TSA's certificate)
\item
  \textbf{Embed the timestamp} --- add the \texttt{TimeStampToken} as an
  unsigned attribute in the CMS
\end{enumerate}

\begin{Shaded}
\begin{Highlighting}[]
\CommentTok{// Requesting a timestamp from a TSA}
\NormalTok{TimeStampRequestGenerator tsReqGen }\OperatorTok{=} \KeywordTok{new} \FunctionTok{TimeStampRequestGenerator}\OperatorTok{();}
\NormalTok{tsReqGen}\OperatorTok{.}\FunctionTok{setCertReq}\OperatorTok{(}\KeywordTok{true}\OperatorTok{);}  \CommentTok{// Request TSA\textquotesingle{}s certificate in the response}

\DataTypeTok{byte}\OperatorTok{[]}\NormalTok{ signatureHash }\OperatorTok{=} \BuiltInTok{MessageDigest}\OperatorTok{.}\FunctionTok{getInstance}\OperatorTok{(}\StringTok{"SHA{-}256"}\OperatorTok{)}
        \OperatorTok{.}\FunctionTok{digest}\OperatorTok{(}\NormalTok{signerInfo}\OperatorTok{.}\FunctionTok{getSignature}\OperatorTok{());}

\NormalTok{TimeStampRequest tsReq }\OperatorTok{=}\NormalTok{ tsReqGen}\OperatorTok{.}\FunctionTok{generate}\OperatorTok{(}
\NormalTok{        TSPAlgorithms}\OperatorTok{.}\FunctionTok{SHA256}\OperatorTok{,}\NormalTok{ signatureHash}\OperatorTok{);}

\CommentTok{// Send to TSA via HTTP}
\DataTypeTok{byte}\OperatorTok{[]}\NormalTok{ tsReqBytes }\OperatorTok{=}\NormalTok{ tsReq}\OperatorTok{.}\FunctionTok{getEncoded}\OperatorTok{();}
\BuiltInTok{HttpURLConnection}\NormalTok{ conn }\OperatorTok{=} \OperatorTok{(}\BuiltInTok{HttpURLConnection}\OperatorTok{)}
        \KeywordTok{new} \BuiltInTok{URL}\OperatorTok{(}\NormalTok{tsaUrl}\OperatorTok{).}\FunctionTok{openConnection}\OperatorTok{();}
\NormalTok{conn}\OperatorTok{.}\FunctionTok{setDoOutput}\OperatorTok{(}\KeywordTok{true}\OperatorTok{);}
\NormalTok{conn}\OperatorTok{.}\FunctionTok{setRequestMethod}\OperatorTok{(}\StringTok{"POST"}\OperatorTok{);}
\NormalTok{conn}\OperatorTok{.}\FunctionTok{setRequestProperty}\OperatorTok{(}\StringTok{"Content{-}Type"}\OperatorTok{,} \StringTok{"application/timestamp{-}query"}\OperatorTok{);}
\NormalTok{conn}\OperatorTok{.}\FunctionTok{getOutputStream}\OperatorTok{().}\FunctionTok{write}\OperatorTok{(}\NormalTok{tsReqBytes}\OperatorTok{);}

\NormalTok{TimeStampResponse tsResp }\OperatorTok{=} \KeywordTok{new} \FunctionTok{TimeStampResponse}\OperatorTok{(}\NormalTok{conn}\OperatorTok{.}\FunctionTok{getInputStream}\OperatorTok{());}
\NormalTok{tsResp}\OperatorTok{.}\FunctionTok{validate}\OperatorTok{(}\NormalTok{tsReq}\OperatorTok{);}  \CommentTok{// Verify the response matches our request}

\NormalTok{TimeStampToken tsToken }\OperatorTok{=}\NormalTok{ tsResp}\OperatorTok{.}\FunctionTok{getTimeStampToken}\OperatorTok{();}
\end{Highlighting}
\end{Shaded}

\subsection{Choosing a TSA}\label{_choosing_a_tsa}

Choosing a TSA is a non-trivial decision. The TSA's certificate must be
trusted by the validators that will check your signatures. Key
considerations:

\begin{itemize}
\item
  \textbf{For Adobe Acrobat}: The TSA's root must be in the AATL, or the
  timestamp will show a warning
\item
  \textbf{For EU compliance}: The TSA should be a Qualified Trust
  Service Provider (QTSP) listed in the EUTL
\item
  \textbf{For availability}: The TSA must be reachable when you sign ---
  if the TSA is down, your signing fails (or you produce a B-B signature
  without a timestamp)
\item
  \textbf{For cost}: Some TSAs are free (FreeTSA, DigiCert's free tier),
  others charge per timestamp or per certificate
\end{itemize}

Chapter 18 covers timestamps in depth, including TSA selection, failover
strategies, and the Document Security Store.

\section{PAdES-B-LT: Long-Term Data}\label{_pades_b_lt_long_term_data}

\subsection{What B-LT Adds}\label{_what_b_lt_adds}

PAdES-B-LT is PAdES-B-T plus \textbf{long-term validation data} embedded
directly in the PDF. This data ensures that the signature can be
validated years from now, even if the signing certificate's CA has gone
out of business, the OCSP responder has been decommissioned, and the CRL
distribution points return 404.

The long-term data is stored in the \textbf{DSS} (Document Security
Store) dictionary, which is added to the document catalog:

\begin{verbatim}
<< /Type /Catalog
   /AcroForm << ... >>
   /DSS <<
     /Certs [30 0 R 31 0 R 32 0 R]       % Certificate chain as DER-encoded streams
     /OCSPs [33 0 R]                       % OCSP responses as DER-encoded streams
     /CRLs [34 0 R]                        % CRLs as DER-encoded streams
     /VRI <<
       /A1B2C3D4... <<                     % Hash-keyed VRI (Validation Related Info)
         /Cert [30 0 R 31 0 R]
         /OCSP [33 0 R]
         /TU (D:20260308120000Z)           % Time of validation
       >>
     >>
   >>
>>
\end{verbatim}

\subsection{The DSS Dictionary}\label{_the_dss_dictionary}

\begin{longtable}[]{@{}
  >{\raggedright\arraybackslash}p{(\linewidth - 6\tabcolsep) * \real{0.2857}}
  >{\raggedright\arraybackslash}p{(\linewidth - 6\tabcolsep) * \real{0.2857}}
  >{\raggedright\arraybackslash}p{(\linewidth - 6\tabcolsep) * \real{0.1429}}
  >{\raggedright\arraybackslash}p{(\linewidth - 6\tabcolsep) * \real{0.2857}}@{}}
\caption{DSS Dictionary Keys (ISO 32000-2:2020
§12.8.4.3)}\tabularnewline
\toprule\noalign{}
\begin{minipage}[b]{\linewidth}\raggedright
Key
\end{minipage} & \begin{minipage}[b]{\linewidth}\raggedright
Type
\end{minipage} & \begin{minipage}[b]{\linewidth}\raggedright
Req
\end{minipage} & \begin{minipage}[b]{\linewidth}\raggedright
Spec Ref
\end{minipage} \\
\midrule\noalign{}
\endfirsthead
\toprule\noalign{}
\begin{minipage}[b]{\linewidth}\raggedright
Key
\end{minipage} & \begin{minipage}[b]{\linewidth}\raggedright
Type
\end{minipage} & \begin{minipage}[b]{\linewidth}\raggedright
Req
\end{minipage} & \begin{minipage}[b]{\linewidth}\raggedright
Spec Ref
\end{minipage} \\
\midrule\noalign{}
\endhead
\bottomrule\noalign{}
\endlastfoot
\texttt{/Certs} & array of streams & Optional & §12.8.4.3 \\
\multicolumn{4}{@{}>{\raggedright\arraybackslash}p{(\linewidth - 6\tabcolsep) * \real{1.0000} + 6\tabcolsep}@{}}{%
DER-encoded certificates --- the full chain for every signature and
timestamp in the document. Include the signer's certificate, all
intermediates, and optionally the root.} \\
\texttt{/OCSPs} & array of streams & Optional & §12.8.4.3 \\
\multicolumn{4}{@{}>{\raggedright\arraybackslash}p{(\linewidth - 6\tabcolsep) * \real{1.0000} + 6\tabcolsep}@{}}{%
DER-encoded OCSP responses for certificate revocation checking. Each
response covers a specific certificate at a specific point in time.} \\
\texttt{/CRLs} & array of streams & Optional & §12.8.4.3 \\
\multicolumn{4}{@{}>{\raggedright\arraybackslash}p{(\linewidth - 6\tabcolsep) * \real{1.0000} + 6\tabcolsep}@{}}{%
DER-encoded Certificate Revocation Lists. Used when OCSP is not
available for a certificate.} \\
\texttt{/VRI} & dictionary & Optional & §12.8.4.3 \\
\multicolumn{4}{@{}>{\raggedright\arraybackslash}p{(\linewidth - 6\tabcolsep) * \real{1.0000} + 6\tabcolsep}@{}}{%
Validation Related Information --- a dictionary keyed by uppercase hex
SHA-1 hashes of signature values, mapping each signature to its specific
validation data.} \\
\end{longtable}

\subsection{Why B-LT Matters}\label{_why_b_lt_matters}

Without B-LT, validating a signature requires contacting external
servers:

\begin{itemize}
\item
  OCSP responder to check revocation status
\item
  CRL distribution point to download the revocation list
\item
  CA's AIA (Authority Information Access) URL to fetch intermediate
  certificates
\end{itemize}

These servers may be unreachable years from now. CAs shut down, merge,
change URLs, or decommission infrastructure. With B-LT, everything
needed for validation is embedded in the PDF itself --- the document is
self-contained.

\section{PAdES-B-LTA: Archival
Timestamps}\label{_pades_b_lta_archival_timestamps}

\subsection{What B-LTA Adds}\label{_what_b_lta_adds}

PAdES-B-LTA is PAdES-B-LT plus a \textbf{document timestamp} --- a
timestamp signature that covers the entire PDF file, including the DSS
dictionary.

This timestamp is not part of the original signature's CMS --- it's a
separate signature dictionary with \texttt{/SubFilter\ /ETSI.RFC3161}:

\begin{verbatim}
<< /Type /Sig
   /Filter /Adobe.PPKLite
   /SubFilter /ETSI.RFC3161
   /ByteRange [0 52000 84776 2100]
   /Contents <3082...>
>>
\end{verbatim}

This timestamp signature covers the entire file (including the original
signature, the DSS, and everything else) and proves that all of this
data existed at the time the timestamp was created.

\subsection{The Archival Chain}\label{_the_archival_chain}

The power of B-LTA is that it creates a \textbf{renewable chain of
trust}:

\begin{enumerate}
\def\labelenumi{\arabic{enumi}.}
\item
  The original signature (signed by the signer's certificate, valid for
  1-3 years)
\item
  The signature timestamp (signed by the TSA's certificate, valid for
  10-20 years)
\item
  The archival timestamp (signed by the TSA's certificate, covers
  everything)
\item
  Years later: add another archival timestamp (covers the previous
  archival timestamp)
\item
  Repeat indefinitely
\end{enumerate}

Each archival timestamp extends the provable lifetime of the signature.
Even after the original signing certificate has expired, been revoked,
and the CA has gone out of business --- the chain of archival timestamps
proves that the signature was valid at the time it was created.

This is the mechanism that national archives use to preserve the legal
validity of digitally signed documents for decades.

\subsection{Re-Timestamping}\label{_re_timestamping}

To extend a B-LTA signature's validity:

\begin{enumerate}
\def\labelenumi{\arabic{enumi}.}
\item
  Fetch fresh OCSP responses / CRLs for the archival timestamp's
  certificate chain
\item
  Add them to the DSS dictionary
\item
  Apply a new archival timestamp (document timestamp) over the entire
  file
\end{enumerate}

This can be done without the original signer's involvement --- it's a
maintenance operation performed by the document custodian.

\textbf{Automated re-timestamping} should be a scheduled process in any
document management system that stores B-LTA signed documents. Before
the current archival timestamp's TSA certificate expires (check the
certificate's \texttt{notAfter} date), apply a new archival timestamp.
This keeps the chain of trust alive.

\section{Profile Comparison Matrix}\label{_profile_comparison_matrix}

\begin{longtable}[]{@{}
  >{\raggedright\arraybackslash}p{(\linewidth - 8\tabcolsep) * \real{0.4286}}
  >{\raggedright\arraybackslash}p{(\linewidth - 8\tabcolsep) * \real{0.1429}}
  >{\raggedright\arraybackslash}p{(\linewidth - 8\tabcolsep) * \real{0.1429}}
  >{\raggedright\arraybackslash}p{(\linewidth - 8\tabcolsep) * \real{0.1429}}
  >{\raggedright\arraybackslash}p{(\linewidth - 8\tabcolsep) * \real{0.1429}}@{}}
\caption{PAdES profile feature comparison}\tabularnewline
\toprule\noalign{}
\begin{minipage}[b]{\linewidth}\raggedright
Feature
\end{minipage} & \begin{minipage}[b]{\linewidth}\raggedright
B-B
\end{minipage} & \begin{minipage}[b]{\linewidth}\raggedright
B-T
\end{minipage} & \begin{minipage}[b]{\linewidth}\raggedright
B-LT
\end{minipage} & \begin{minipage}[b]{\linewidth}\raggedright
B-LTA
\end{minipage} \\
\midrule\noalign{}
\endfirsthead
\toprule\noalign{}
\begin{minipage}[b]{\linewidth}\raggedright
Feature
\end{minipage} & \begin{minipage}[b]{\linewidth}\raggedright
B-B
\end{minipage} & \begin{minipage}[b]{\linewidth}\raggedright
B-T
\end{minipage} & \begin{minipage}[b]{\linewidth}\raggedright
B-LT
\end{minipage} & \begin{minipage}[b]{\linewidth}\raggedright
B-LTA
\end{minipage} \\
\midrule\noalign{}
\endhead
\bottomrule\noalign{}
\endlastfoot
CMS with ESS-signing-certificate-v2 & Yes & Yes & Yes & Yes \\
Signature timestamp & No & Yes & Yes & Yes \\
Embedded revocation data (DSS) & No & No & Yes & Yes \\
Embedded certificate chain (DSS) & No & No & Yes & Yes \\
Document timestamp (archival) & No & No & No & Yes \\
Offline validation & No & Partial & Yes & Yes \\
Validity beyond cert expiry & No & Yes (with TS) & Yes & Yes
(indefinite) \\
EU/eIDAS suitability & Minimal & Short-term & Good & \textbf{Full
compliance} \\
Typical use case & Internal docs & Contracts & Regulated docs &
Archival \\
\end{longtable}

\section{Migrating Between Profiles}\label{_migrating_between_profiles}

One of the elegant properties of PAdES profiles is that you can
\textbf{upgrade} a signature's profile without re-signing:

\begin{itemize}
\item
  \textbf{B-B → B-T}: Add a timestamp (unsigned attribute to the CMS).
  Requires the original file and the CMS.
\item
  \textbf{B-T → B-LT}: Fetch revocation data and certificates, create
  the DSS dictionary, add as incremental update.
\item
  \textbf{B-LT → B-LTA}: Add a document timestamp covering the entire
  file.
\end{itemize}

This progression means you can sign with B-B today (fast, no TSA
dependency) and upgrade to B-LTA later as a batch operation. EU DSS
provides methods for exactly this:

\begin{Shaded}
\begin{Highlighting}[]
\CommentTok{// EU DSS: Extend a B{-}B signature to B{-}LTA}
\NormalTok{PAdESService service }\OperatorTok{=} \KeywordTok{new} \FunctionTok{PAdESService}\OperatorTok{(}\NormalTok{certificateVerifier}\OperatorTok{);}
\NormalTok{service}\OperatorTok{.}\FunctionTok{setTspSource}\OperatorTok{(}\NormalTok{tspSource}\OperatorTok{);}

\NormalTok{DSSDocument extendedDoc }\OperatorTok{=}\NormalTok{ service}\OperatorTok{.}\FunctionTok{extendDocument}\OperatorTok{(}
\NormalTok{    signedDocument}\OperatorTok{,}
    \KeywordTok{new} \FunctionTok{PAdESSignatureParameters}\OperatorTok{()} \OperatorTok{\{\{}
        \FunctionTok{setSignatureLevel}\OperatorTok{(}\NormalTok{SignatureLevel}\OperatorTok{.}\FunctionTok{PAdES\_BASELINE\_LTA}\OperatorTok{);}
    \OperatorTok{\}\});}
\end{Highlighting}
\end{Shaded}

\section{Key Takeaways}\label{_key_takeaways_5}

\begin{itemize}
\item
  PAdES constrains ISO 32000-2 signatures for legal and regulatory
  compliance --- a PAdES signature is always a valid PDF signature, but
  not vice versa
\item
  \textbf{Four profiles}: B-B (basic), B-T (timestamped), B-LT
  (long-term data), B-LTA (archival)
\item
  \textbf{ESS-signing-certificate-v2} is the mandatory attribute that
  makes a signature PAdES-compliant --- without it, you have a valid ISO
  32000-2 signature but not PAdES
\item
  \textbf{SubFilter \texttt{ETSI.CAdES.detached}} is required for strict
  PAdES compliance
\item
  The \textbf{DSS dictionary} (B-LT and above) embeds certificates, OCSP
  responses, and CRLs directly in the PDF for offline validation
\item
  \textbf{B-LTA} provides indefinite validity through renewable archival
  timestamps
\item
  Profiles can be \textbf{upgraded} without re-signing --- B-B → B-T →
  B-LT → B-LTA as incremental updates
\item
  For EU/eIDAS compliance, \textbf{B-LTA is the target} --- anything
  less may not meet long-term archival requirements
\end{itemize}

\chapter{Timestamps \& Long-Term
Validation}\label{_timestamps_long_term_validation}

A digital signature without a timestamp is like a contract without a
date. You can prove \emph{who} signed and \emph{what} they signed, but
not \emph{when}. And in the world of digital signatures, "when" is often
the most important question --- because certificates expire, keys get
compromised, algorithms get broken, and the difference between "signed
before the certificate expired" and "signed after the certificate
expired" can be the difference between a legally binding document and an
invalid one.

This chapter covers everything about timestamps and long-term validation
(LTV): how RFC 3161 timestamps work, how they're embedded in PDF
signatures, how the Document Security Store (DSS) enables offline
validation, how archival timestamps create indefinite validity, and how
to build a complete LTV infrastructure.

\section{RFC 3161: The Timestamp
Protocol}\label{_rfc_3161_the_timestamp_protocol}

\subsection{What a Timestamp Proves}\label{_what_a_timestamp_proves}

An RFC 3161 timestamp is a cryptographic proof that a specific piece of
data existed at a specific point in time. The proof comes from a
\textbf{Trusted Timestamp Authority (TSA)} --- a server operated by a
trusted third party that:

\begin{enumerate}
\def\labelenumi{\arabic{enumi}.}
\item
  Receives a hash of the data to be timestamped
\item
  Combines the hash with the current time
\item
  Signs the combination with the TSA's private key
\item
  Returns the signed timestamp token
\end{enumerate}

The timestamp token is a CMS \texttt{SignedData} structure (yes, a
timestamp is itself a CMS signature) containing a \texttt{TSTInfo}
(Timestamp Token Info) structure with:

\begin{itemize}
\item
  The hash algorithm and hash value of the timestamped data
\item
  The timestamp time (UTC)
\item
  The TSA's name
\item
  A serial number (unique per timestamp)
\item
  An accuracy indication (how precise the time is)
\end{itemize}

\subsection{The Timestamp Request/Response
Flow}\label{_the_timestamp_requestresponse_flow}

\begin{figure}
\centering
\pandocbounded{\includegraphics[keepaspectratio,alt={seq tsa timestamp}]{diagrams/pdf/seq-tsa-timestamp.pdf}}
\caption{RFC 3161 timestamp request/response flow}
\end{figure}

The protocol uses HTTP POST with content type
\texttt{application/timestamp-query} for the request and
\texttt{application/timestamp-reply} for the response. The request and
response are DER-encoded ASN.1 structures defined in RFC 3161.

\subsection{Implementation with
BouncyCastle}\label{_implementation_with_bouncycastle}

\begin{Shaded}
\begin{Highlighting}[]
\CommentTok{/**}
 \CommentTok{*}\NormalTok{ Requests a timestamp from a TSA for the given data}\CommentTok{.}
 \CommentTok{*}\NormalTok{ Returns the TimeStampToken that can be embedded as an unsigned CMS attribute}\CommentTok{.}
 \CommentTok{*/}
\KeywordTok{public}\NormalTok{ TimeStampToken }\FunctionTok{requestTimestamp}\OperatorTok{(}\DataTypeTok{byte}\OperatorTok{[]}\NormalTok{ dataToTimestamp}\OperatorTok{,} \BuiltInTok{String}\NormalTok{ tsaUrl}\OperatorTok{)}
        \KeywordTok{throws} \BuiltInTok{Exception} \OperatorTok{\{}

    \CommentTok{// 1. Hash the data}
    \BuiltInTok{MessageDigest}\NormalTok{ digest }\OperatorTok{=} \BuiltInTok{MessageDigest}\OperatorTok{.}\FunctionTok{getInstance}\OperatorTok{(}\StringTok{"SHA{-}256"}\OperatorTok{);}
    \DataTypeTok{byte}\OperatorTok{[]}\NormalTok{ hash }\OperatorTok{=}\NormalTok{ digest}\OperatorTok{.}\FunctionTok{digest}\OperatorTok{(}\NormalTok{dataToTimestamp}\OperatorTok{);}

    \CommentTok{// 2. Create the timestamp request}
\NormalTok{    TimeStampRequestGenerator reqGen }\OperatorTok{=} \KeywordTok{new} \FunctionTok{TimeStampRequestGenerator}\OperatorTok{();}
\NormalTok{    reqGen}\OperatorTok{.}\FunctionTok{setCertReq}\OperatorTok{(}\KeywordTok{true}\OperatorTok{);}  \CommentTok{// Include TSA\textquotesingle{}s cert in the response}
\NormalTok{    TimeStampRequest tsReq }\OperatorTok{=}\NormalTok{ reqGen}\OperatorTok{.}\FunctionTok{generate}\OperatorTok{(}\NormalTok{TSPAlgorithms}\OperatorTok{.}\FunctionTok{SHA256}\OperatorTok{,}\NormalTok{ hash}\OperatorTok{);}

    \CommentTok{// 3. Send to TSA via HTTP}
    \DataTypeTok{byte}\OperatorTok{[]}\NormalTok{ reqBytes }\OperatorTok{=}\NormalTok{ tsReq}\OperatorTok{.}\FunctionTok{getEncoded}\OperatorTok{();}
    \BuiltInTok{HttpURLConnection}\NormalTok{ conn }\OperatorTok{=} \OperatorTok{(}\BuiltInTok{HttpURLConnection}\OperatorTok{)}
            \KeywordTok{new} \BuiltInTok{URL}\OperatorTok{(}\NormalTok{tsaUrl}\OperatorTok{).}\FunctionTok{openConnection}\OperatorTok{();}
\NormalTok{    conn}\OperatorTok{.}\FunctionTok{setDoOutput}\OperatorTok{(}\KeywordTok{true}\OperatorTok{);}
\NormalTok{    conn}\OperatorTok{.}\FunctionTok{setRequestMethod}\OperatorTok{(}\StringTok{"POST"}\OperatorTok{);}
\NormalTok{    conn}\OperatorTok{.}\FunctionTok{setRequestProperty}\OperatorTok{(}\StringTok{"Content{-}Type"}\OperatorTok{,} \StringTok{"application/timestamp{-}query"}\OperatorTok{);}
\NormalTok{    conn}\OperatorTok{.}\FunctionTok{setConnectTimeout}\OperatorTok{(}\DecValTok{10000}\OperatorTok{);}
\NormalTok{    conn}\OperatorTok{.}\FunctionTok{setReadTimeout}\OperatorTok{(}\DecValTok{30000}\OperatorTok{);}

    \ControlFlowTok{try} \OperatorTok{(}\BuiltInTok{OutputStream}\NormalTok{ out }\OperatorTok{=}\NormalTok{ conn}\OperatorTok{.}\FunctionTok{getOutputStream}\OperatorTok{())} \OperatorTok{\{}
\NormalTok{        out}\OperatorTok{.}\FunctionTok{write}\OperatorTok{(}\NormalTok{reqBytes}\OperatorTok{);}
    \OperatorTok{\}}

    \CommentTok{// 4. Read the response}
    \DataTypeTok{byte}\OperatorTok{[]}\NormalTok{ respBytes}\OperatorTok{;}
    \ControlFlowTok{try} \OperatorTok{(}\BuiltInTok{InputStream}\NormalTok{ in }\OperatorTok{=}\NormalTok{ conn}\OperatorTok{.}\FunctionTok{getInputStream}\OperatorTok{())} \OperatorTok{\{}
\NormalTok{        respBytes }\OperatorTok{=}\NormalTok{ in}\OperatorTok{.}\FunctionTok{readAllBytes}\OperatorTok{();}
    \OperatorTok{\}}

    \CommentTok{// 5. Parse and validate the response}
\NormalTok{    TimeStampResponse tsResp }\OperatorTok{=} \KeywordTok{new} \FunctionTok{TimeStampResponse}\OperatorTok{(}\NormalTok{respBytes}\OperatorTok{);}
\NormalTok{    tsResp}\OperatorTok{.}\FunctionTok{validate}\OperatorTok{(}\NormalTok{tsReq}\OperatorTok{);}  \CommentTok{// Verify it matches our request}

    \ControlFlowTok{if} \OperatorTok{(}\NormalTok{tsResp}\OperatorTok{.}\FunctionTok{getStatus}\OperatorTok{()} \OperatorTok{!=}\NormalTok{ PKIStatus}\OperatorTok{.}\FunctionTok{GRANTED}
            \OperatorTok{\&\&}\NormalTok{ tsResp}\OperatorTok{.}\FunctionTok{getStatus}\OperatorTok{()} \OperatorTok{!=}\NormalTok{ PKIStatus}\OperatorTok{.}\FunctionTok{GRANTED\_WITH\_MODS}\OperatorTok{)} \OperatorTok{\{}
        \ControlFlowTok{throw} \KeywordTok{new} \FunctionTok{TSPException}\OperatorTok{(}\StringTok{"TSA rejected request: status="} \OperatorTok{+}
\NormalTok{                tsResp}\OperatorTok{.}\FunctionTok{getStatus}\OperatorTok{()} \OperatorTok{+} \StringTok{", statusString="} \OperatorTok{+}
\NormalTok{                tsResp}\OperatorTok{.}\FunctionTok{getStatusString}\OperatorTok{());}
    \OperatorTok{\}}

\NormalTok{    TimeStampToken tsToken }\OperatorTok{=}\NormalTok{ tsResp}\OperatorTok{.}\FunctionTok{getTimeStampToken}\OperatorTok{();}
\NormalTok{    log}\OperatorTok{.}\FunctionTok{info}\OperatorTok{(}\StringTok{"Timestamp obtained: time=\{\}, serial=\{\}, TSA=\{\}"}\OperatorTok{,}
\NormalTok{            tsToken}\OperatorTok{.}\FunctionTok{getTimeStampInfo}\OperatorTok{().}\FunctionTok{getGenTime}\OperatorTok{(),}
\NormalTok{            tsToken}\OperatorTok{.}\FunctionTok{getTimeStampInfo}\OperatorTok{().}\FunctionTok{getSerialNumber}\OperatorTok{(),}
\NormalTok{            tsToken}\OperatorTok{.}\FunctionTok{getTimeStampInfo}\OperatorTok{().}\FunctionTok{getTsa}\OperatorTok{());}

    \ControlFlowTok{return}\NormalTok{ tsToken}\OperatorTok{;}
\OperatorTok{\}}
\end{Highlighting}
\end{Shaded}

\subsection{TSA Selection Criteria}\label{_tsa_selection_criteria}

Choosing the right TSA is critical. A timestamp is only as trustworthy
as the TSA that produced it.

\begin{longtable}[]{@{}
  >{\raggedright\arraybackslash}p{(\linewidth - 2\tabcolsep) * \real{0.3333}}
  >{\raggedright\arraybackslash}p{(\linewidth - 2\tabcolsep) * \real{0.6667}}@{}}
\caption{TSA selection criteria for different use cases}\tabularnewline
\toprule\noalign{}
\begin{minipage}[b]{\linewidth}\raggedright
Criterion
\end{minipage} & \begin{minipage}[b]{\linewidth}\raggedright
Considerations
\end{minipage} \\
\midrule\noalign{}
\endfirsthead
\toprule\noalign{}
\begin{minipage}[b]{\linewidth}\raggedright
Criterion
\end{minipage} & \begin{minipage}[b]{\linewidth}\raggedright
Considerations
\end{minipage} \\
\midrule\noalign{}
\endhead
\bottomrule\noalign{}
\endlastfoot
\textbf{Trust store compatibility} & For Adobe Acrobat: TSA's root must
be in AATL. For EU compliance: TSA must be a QTSP in the EUTL. For
general use: TSA's root should be widely trusted. \\
\textbf{Availability \& SLA} & TSAs can go down. For production signing,
use a TSA with a published SLA. Consider a fallback TSA if the primary
is unreachable. \\
\textbf{Cost} & Free TSAs exist (FreeTSA.org, some DigiCert endpoints)
but may not be suitable for production. Commercial TSAs charge per
timestamp or per annual subscription. Costs range from free to several
thousand euros per year. \\
\textbf{Accuracy} & TSAs report their time accuracy (e.g., ±1 second).
For most use cases, sub-second accuracy is not critical. \\
\textbf{Geographic location} & Some regulations require the TSA to be
operated within specific jurisdictions. EU regulations may require a TSA
operated by an EU-based QTSP. \\
\textbf{Response time} & Network latency + TSA processing time. Free
TSAs may be slower. For high-volume signing (thousands per minute), TSA
response time matters. \\
\textbf{Certificate lifetime} & The TSA's signing certificate has an
expiry date. Longer-lived TSA certificates mean your timestamps are
verifiable for longer before re-timestamping is needed. \\
\end{longtable}

\textbf{Never hardcode a single TSA URL} without a fallback strategy. If
the TSA is unreachable when your application needs to sign, the signing
fails. At minimum, implement:

\begin{enumerate}
\def\labelenumi{\arabic{enumi}.}
\item
  A configurable TSA URL (not hardcoded)
\item
  A reasonable timeout (10-30 seconds)
\item
  A retry with backoff (the TSA may be temporarily overloaded)
\item
  Optionally, a fallback TSA URL
\item
  A graceful degradation path: if no TSA is available, produce a B-B
  signature and log a warning for later B-T upgrade
\end{enumerate}

\section{TSA Health Probing}\label{_tsa_health_probing}

\subsection{Why You Need Probing}\label{_why_you_need_probing}

Before your application signs a document, you need to know which TSAs
are actually reachable. A signing operation that fails at the
timestamping step --- after the user has already authenticated, after
the CSC API has consumed a SAD token, after the CMS signature has been
computed --- is one of the most frustrating failure modes in digital
signing. You can avoid it by probing TSAs in the background.

\subsection{Building a TSA Probe}\label{_building_a_tsa_probe}

A TSA probe sends a minimal RFC 3161 timestamp request --- using a dummy
hash, with \texttt{certReq=false} to minimize response size --- and
measures whether the server responds correctly:

\begin{Shaded}
\begin{Highlighting}[]
\CommentTok{/**}
 \CommentTok{*}\NormalTok{ Probes a TSA server by sending a minimal timestamp request}\CommentTok{.}
 \CommentTok{*}\NormalTok{ Returns the response time in milliseconds}\CommentTok{,}\NormalTok{ or }\CommentTok{{-}1}\NormalTok{ if the probe fails}\CommentTok{.}
 \CommentTok{*/}
\KeywordTok{public} \DataTypeTok{long} \FunctionTok{probeTsa}\OperatorTok{(}\BuiltInTok{String}\NormalTok{ tsaUrl}\OperatorTok{,} \DataTypeTok{int}\NormalTok{ timeoutMs}\OperatorTok{)} \OperatorTok{\{}
    \ControlFlowTok{try} \OperatorTok{\{}
        \CommentTok{// Hash of dummy data — we don\textquotesingle{}t care about the content}
        \DataTypeTok{byte}\OperatorTok{[]}\NormalTok{ dummyHash }\OperatorTok{=} \BuiltInTok{MessageDigest}\OperatorTok{.}\FunctionTok{getInstance}\OperatorTok{(}\StringTok{"SHA{-}256"}\OperatorTok{)}
                \OperatorTok{.}\FunctionTok{digest}\OperatorTok{(}\StringTok{"tsa{-}probe"}\OperatorTok{.}\FunctionTok{getBytes}\OperatorTok{(}\NormalTok{StandardCharsets}\OperatorTok{.}\FunctionTok{UTF\_8}\OperatorTok{));}

\NormalTok{        TimeStampRequestGenerator gen }\OperatorTok{=} \KeywordTok{new} \FunctionTok{TimeStampRequestGenerator}\OperatorTok{();}
\NormalTok{        gen}\OperatorTok{.}\FunctionTok{setCertReq}\OperatorTok{(}\KeywordTok{false}\OperatorTok{);}  \CommentTok{// Minimize response size}
\NormalTok{        TimeStampRequest req }\OperatorTok{=}\NormalTok{ gen}\OperatorTok{.}\FunctionTok{generate}\OperatorTok{(}\NormalTok{TSPAlgorithms}\OperatorTok{.}\FunctionTok{SHA256}\OperatorTok{,}\NormalTok{ dummyHash}\OperatorTok{);}

        \DataTypeTok{long}\NormalTok{ start }\OperatorTok{=} \BuiltInTok{System}\OperatorTok{.}\FunctionTok{nanoTime}\OperatorTok{();}

        \BuiltInTok{HttpURLConnection}\NormalTok{ conn }\OperatorTok{=} \OperatorTok{(}\BuiltInTok{HttpURLConnection}\OperatorTok{)}
                \KeywordTok{new} \BuiltInTok{URL}\OperatorTok{(}\NormalTok{tsaUrl}\OperatorTok{).}\FunctionTok{openConnection}\OperatorTok{();}
\NormalTok{        conn}\OperatorTok{.}\FunctionTok{setDoOutput}\OperatorTok{(}\KeywordTok{true}\OperatorTok{);}
\NormalTok{        conn}\OperatorTok{.}\FunctionTok{setRequestMethod}\OperatorTok{(}\StringTok{"POST"}\OperatorTok{);}
\NormalTok{        conn}\OperatorTok{.}\FunctionTok{setRequestProperty}\OperatorTok{(}\StringTok{"Content{-}Type"}\OperatorTok{,}
                \StringTok{"application/timestamp{-}query"}\OperatorTok{);}
\NormalTok{        conn}\OperatorTok{.}\FunctionTok{setConnectTimeout}\OperatorTok{(}\NormalTok{timeoutMs}\OperatorTok{);}
\NormalTok{        conn}\OperatorTok{.}\FunctionTok{setReadTimeout}\OperatorTok{(}\NormalTok{timeoutMs}\OperatorTok{);}
\NormalTok{        conn}\OperatorTok{.}\FunctionTok{getOutputStream}\OperatorTok{().}\FunctionTok{write}\OperatorTok{(}\NormalTok{req}\OperatorTok{.}\FunctionTok{getEncoded}\OperatorTok{());}

        \CommentTok{// Validate Content{-}Type — some misconfigured servers}
        \CommentTok{// return an HTML error page with HTTP 200}
        \BuiltInTok{String}\NormalTok{ contentType }\OperatorTok{=}\NormalTok{ conn}\OperatorTok{.}\FunctionTok{getContentType}\OperatorTok{();}
        \ControlFlowTok{if} \OperatorTok{(}\NormalTok{contentType }\OperatorTok{==} \KeywordTok{null}
                \OperatorTok{||} \OperatorTok{!}\NormalTok{contentType}\OperatorTok{.}\FunctionTok{contains}\OperatorTok{(}\StringTok{"timestamp{-}reply"}\OperatorTok{))} \OperatorTok{\{}
\NormalTok{            log}\OperatorTok{.}\FunctionTok{warn}\OperatorTok{(}\StringTok{"TSA \{\} returned unexpected Content{-}Type: \{\}"}\OperatorTok{,}
\NormalTok{                    tsaUrl}\OperatorTok{,}\NormalTok{ contentType}\OperatorTok{);}
            \ControlFlowTok{return} \OperatorTok{{-}}\DecValTok{1}\OperatorTok{;}
        \OperatorTok{\}}

        \DataTypeTok{byte}\OperatorTok{[]}\NormalTok{ respBytes }\OperatorTok{=}\NormalTok{ conn}\OperatorTok{.}\FunctionTok{getInputStream}\OperatorTok{().}\FunctionTok{readAllBytes}\OperatorTok{();}
\NormalTok{        TimeStampResponse resp }\OperatorTok{=} \KeywordTok{new} \FunctionTok{TimeStampResponse}\OperatorTok{(}\NormalTok{respBytes}\OperatorTok{);}

        \DataTypeTok{long}\NormalTok{ elapsed }\OperatorTok{=} \OperatorTok{(}\BuiltInTok{System}\OperatorTok{.}\FunctionTok{nanoTime}\OperatorTok{()} \OperatorTok{{-}}\NormalTok{ start}\OperatorTok{)} \OperatorTok{/} \DecValTok{1\_000\_000}\OperatorTok{;}

        \ControlFlowTok{if} \OperatorTok{(}\NormalTok{resp}\OperatorTok{.}\FunctionTok{getStatus}\OperatorTok{()} \OperatorTok{==}\NormalTok{ PKIStatus}\OperatorTok{.}\FunctionTok{GRANTED}
                \OperatorTok{||}\NormalTok{ resp}\OperatorTok{.}\FunctionTok{getStatus}\OperatorTok{()} \OperatorTok{==}\NormalTok{ PKIStatus}\OperatorTok{.}\FunctionTok{GRANTED\_WITH\_MODS}\OperatorTok{)} \OperatorTok{\{}
            \ControlFlowTok{return}\NormalTok{ elapsed}\OperatorTok{;}
        \OperatorTok{\}}
        \ControlFlowTok{return} \OperatorTok{{-}}\DecValTok{1}\OperatorTok{;}
    \OperatorTok{\}} \ControlFlowTok{catch} \OperatorTok{(}\BuiltInTok{Exception}\NormalTok{ e}\OperatorTok{)} \OperatorTok{\{}
        \ControlFlowTok{return} \OperatorTok{{-}}\DecValTok{1}\OperatorTok{;}
    \OperatorTok{\}}
\OperatorTok{\}}
\end{Highlighting}
\end{Shaded}

\textbf{Content-Type validation is essential.} Some TSA endpoints sit
behind reverse proxies or WAFs that return an HTTP 200 response with an
HTML error page when the backend is down. Without checking
\texttt{Content-Type:\ application/timestamp-reply}, your application
might accept a garbage response and only fail later when trying to parse
it as ASN.1. Always validate the response \texttt{Content-Type} before
parsing.

\subsection{Probe Strategy}\label{_probe_strategy}

For production applications, run probes on a schedule:

\begin{enumerate}
\def\labelenumi{\arabic{enumi}.}
\item
  \textbf{At startup}: Probe all configured TSAs concurrently using a
  thread pool. Mark each as "online" or "offline" with its measured
  response time.
\item
  \textbf{Periodically}: Re-probe every 5-15 minutes. TSA availability
  can change --- maintenance windows, network issues, certificate
  rotations.
\item
  \textbf{On failure}: If a signing operation fails at the timestamping
  step, immediately probe all TSAs and update their status. The failure
  may be transient (timeout) or permanent (certificate expired, endpoint
  removed).
\end{enumerate}

Present the probe results to the user: a green/red status indicator with
response time helps them choose the fastest available TSA for their
signing operation.

\section{Caching TSA and OCSP
Responses}\label{_caching_tsa_and_ocsp_responses}

\subsection{Why Caching Matters}\label{_why_caching_matters}

Timestamp and OCSP requests are network calls to external servers. In a
high-volume signing scenario --- batch signing hundreds of documents ---
making a fresh TSA request and OCSP lookup for every single signature is
wasteful and slow. More importantly, it creates an availability
dependency: if the TSA or OCSP responder goes down mid-batch, the entire
batch fails.

\subsection{TSA Response Caching}\label{_tsa_response_caching}

TSA responses are \textbf{not cacheable} in the traditional sense ---
each timestamp token contains a unique serial number and a specific
timestamp, making every response unique by design. However, you can
optimize for batch scenarios:

\begin{enumerate}
\def\labelenumi{\arabic{enumi}.}
\item
  \textbf{Connection pooling}: Reuse HTTP connections to the TSA server.
  TLS handshakes are expensive; a keep-alive connection eliminates that
  overhead for subsequent requests.
\item
  \textbf{Pre-warming}: Establish the connection before the signing flow
  begins. The first request to a TSA often takes 2-5x longer than
  subsequent requests due to DNS resolution, TLS negotiation, and
  server-side session setup.
\item
  \textbf{Batch authorization}: Some CSC providers and TSA operators
  support batch requests --- signing multiple hashes or timestamping
  multiple values in a single API call. If your TSA supports it, prefer
  batch requests over individual calls.
\end{enumerate}

\textbf{Never reuse a timestamp token across multiple signatures.} Each
timestamp proves existence at a specific instant for a specific hash.
Sharing a timestamp token between two signatures is technically invalid
(the hash won't match) and a compliance violation.

\subsection{OCSP Response Caching}\label{_ocsp_response_caching}

Unlike timestamps, OCSP responses \textbf{are cacheable} --- and you
should cache them aggressively. An OCSP response contains a
\texttt{thisUpdate} and \texttt{nextUpdate} field that define a validity
window. Within that window, the response is authoritative: the CA is
asserting that the certificate's revocation status has not changed since
\texttt{thisUpdate} and will not change before \texttt{nextUpdate}.

\begin{Shaded}
\begin{Highlighting}[]
\CommentTok{/**}
 \CommentTok{*}\NormalTok{ OCSP response cache}\CommentTok{.}
 \CommentTok{*}\NormalTok{ Caches responses by certificate serial number}\CommentTok{,}\NormalTok{ respects nextUpdate}\CommentTok{.}
 \CommentTok{*/}
\KeywordTok{private} \DataTypeTok{final} \BuiltInTok{Map}\OperatorTok{\textless{}}\BuiltInTok{BigInteger}\OperatorTok{,}\NormalTok{ CachedOcspResponse}\OperatorTok{\textgreater{}}\NormalTok{ ocspCache}
        \OperatorTok{=} \KeywordTok{new} \BuiltInTok{ConcurrentHashMap}\OperatorTok{\textless{}\textgreater{}();}

\KeywordTok{public} \DataTypeTok{byte}\OperatorTok{[]} \FunctionTok{getOcspResponse}\OperatorTok{(}\BuiltInTok{X509Certificate}\NormalTok{ cert}\OperatorTok{,}
        \BuiltInTok{X509Certificate}\NormalTok{ issuer}\OperatorTok{)} \KeywordTok{throws} \BuiltInTok{Exception} \OperatorTok{\{}
    \BuiltInTok{BigInteger}\NormalTok{ serial }\OperatorTok{=}\NormalTok{ cert}\OperatorTok{.}\FunctionTok{getSerialNumber}\OperatorTok{();}
\NormalTok{    CachedOcspResponse cached }\OperatorTok{=}\NormalTok{ ocspCache}\OperatorTok{.}\FunctionTok{get}\OperatorTok{(}\NormalTok{serial}\OperatorTok{);}

    \ControlFlowTok{if} \OperatorTok{(}\NormalTok{cached }\OperatorTok{!=} \KeywordTok{null} \OperatorTok{\&\&}\NormalTok{ cached}\OperatorTok{.}\FunctionTok{isValid}\OperatorTok{())} \OperatorTok{\{}
\NormalTok{        log}\OperatorTok{.}\FunctionTok{debug}\OperatorTok{(}\StringTok{"OCSP cache hit for serial=\{\}"}\OperatorTok{,}\NormalTok{ serial}\OperatorTok{);}
        \ControlFlowTok{return}\NormalTok{ cached}\OperatorTok{.}\FunctionTok{responseBytes}\OperatorTok{();}
    \OperatorTok{\}}

    \CommentTok{// Cache miss or expired — fetch fresh response}
    \DataTypeTok{byte}\OperatorTok{[]}\NormalTok{ respBytes }\OperatorTok{=} \FunctionTok{fetchOCSPResponse}\OperatorTok{(}\NormalTok{cert}\OperatorTok{,}\NormalTok{ issuer}\OperatorTok{);}
    \ControlFlowTok{if} \OperatorTok{(}\NormalTok{respBytes }\OperatorTok{!=} \KeywordTok{null}\OperatorTok{)} \OperatorTok{\{}
\NormalTok{        OCSPResp ocspResp }\OperatorTok{=} \KeywordTok{new} \FunctionTok{OCSPResp}\OperatorTok{(}\NormalTok{respBytes}\OperatorTok{);}
\NormalTok{        BasicOCSPResp basic }\OperatorTok{=} \OperatorTok{(}\NormalTok{BasicOCSPResp}\OperatorTok{)}
\NormalTok{                ocspResp}\OperatorTok{.}\FunctionTok{getResponseObject}\OperatorTok{();}
        \BuiltInTok{Date}\NormalTok{ nextUpdate }\OperatorTok{=}\NormalTok{ basic}\OperatorTok{.}\FunctionTok{getResponses}\OperatorTok{()[}\DecValTok{0}\OperatorTok{].}\FunctionTok{getNextUpdate}\OperatorTok{();}

\NormalTok{        ocspCache}\OperatorTok{.}\FunctionTok{put}\OperatorTok{(}\NormalTok{serial}\OperatorTok{,}
                \KeywordTok{new} \FunctionTok{CachedOcspResponse}\OperatorTok{(}\NormalTok{respBytes}\OperatorTok{,}\NormalTok{ nextUpdate}\OperatorTok{));}
\NormalTok{        log}\OperatorTok{.}\FunctionTok{debug}\OperatorTok{(}\StringTok{"OCSP cached for serial=\{\}, valid until \{\}"}\OperatorTok{,}
\NormalTok{                serial}\OperatorTok{,}\NormalTok{ nextUpdate}\OperatorTok{);}
    \OperatorTok{\}}
    \ControlFlowTok{return}\NormalTok{ respBytes}\OperatorTok{;}
\OperatorTok{\}}
\end{Highlighting}
\end{Shaded}

Typical OCSP response validity windows:

\begin{itemize}
\item
  \textbf{Commercial CAs}: 1-4 hours (short, to limit exposure from
  compromised responder)
\item
  \textbf{Government CAs}: 4-24 hours (longer, to reduce load on
  centralized infrastructure)
\item
  \textbf{Qualified CAs (eIDAS)}: 24 hours maximum (ETSI EN 319 411-2
  requirement)
\end{itemize}

For batch signing of 100 documents with the same certificate, a cached
OCSP response eliminates 99 round trips to the OCSP responder.

\subsection{CRL Caching}\label{_crl_caching}

CRLs are even more cacheable than OCSP responses. Every CRL contains a
\texttt{nextUpdate} date, and some CAs publish CRLs on fixed schedules
(every 6, 12, or 24 hours). Cache the CRL bytes keyed by the CRL
Distribution Point URL and check \texttt{nextUpdate} before re-fetching:

\begin{Shaded}
\begin{Highlighting}[]
\CommentTok{/**}
 \CommentTok{*}\NormalTok{ CRL cache }\CommentTok{—}\NormalTok{ caches full CRLs by distribution point URL}\CommentTok{.}
 \CommentTok{*}\NormalTok{ CRLs can be large }\CommentTok{(}\NormalTok{hundreds of KB}\CommentTok{),}\NormalTok{ but they change infrequently}\CommentTok{.}
 \CommentTok{*/}
\KeywordTok{private} \DataTypeTok{final} \BuiltInTok{Map}\OperatorTok{\textless{}}\BuiltInTok{String}\OperatorTok{,}\NormalTok{ CachedCrl}\OperatorTok{\textgreater{}}\NormalTok{ crlCache}
        \OperatorTok{=} \KeywordTok{new} \BuiltInTok{ConcurrentHashMap}\OperatorTok{\textless{}\textgreater{}();}

\KeywordTok{public} \DataTypeTok{byte}\OperatorTok{[]} \FunctionTok{getCrl}\OperatorTok{(}\BuiltInTok{String}\NormalTok{ crlUrl}\OperatorTok{)} \KeywordTok{throws} \BuiltInTok{Exception} \OperatorTok{\{}
\NormalTok{    CachedCrl cached }\OperatorTok{=}\NormalTok{ crlCache}\OperatorTok{.}\FunctionTok{get}\OperatorTok{(}\NormalTok{crlUrl}\OperatorTok{);}

    \ControlFlowTok{if} \OperatorTok{(}\NormalTok{cached }\OperatorTok{!=} \KeywordTok{null} \OperatorTok{\&\&}\NormalTok{ cached}\OperatorTok{.}\FunctionTok{nextUpdate}\OperatorTok{().}\FunctionTok{after}\OperatorTok{(}\KeywordTok{new} \BuiltInTok{Date}\OperatorTok{()))} \OperatorTok{\{}
\NormalTok{        log}\OperatorTok{.}\FunctionTok{debug}\OperatorTok{(}\StringTok{"CRL cache hit for \{\}"}\OperatorTok{,}\NormalTok{ crlUrl}\OperatorTok{);}
        \ControlFlowTok{return}\NormalTok{ cached}\OperatorTok{.}\FunctionTok{crlBytes}\OperatorTok{();}
    \OperatorTok{\}}

    \DataTypeTok{byte}\OperatorTok{[]}\NormalTok{ crlBytes }\OperatorTok{=} \FunctionTok{downloadCrl}\OperatorTok{(}\NormalTok{crlUrl}\OperatorTok{);}
    \BuiltInTok{CertificateFactory}\NormalTok{ cf }\OperatorTok{=} \BuiltInTok{CertificateFactory}\OperatorTok{.}\FunctionTok{getInstance}\OperatorTok{(}\StringTok{"X.509"}\OperatorTok{);}
    \BuiltInTok{X509CRL}\NormalTok{ crl }\OperatorTok{=} \OperatorTok{(}\BuiltInTok{X509CRL}\OperatorTok{)}\NormalTok{ cf}\OperatorTok{.}\FunctionTok{generateCRL}\OperatorTok{(}
            \KeywordTok{new} \BuiltInTok{ByteArrayInputStream}\OperatorTok{(}\NormalTok{crlBytes}\OperatorTok{));}

\NormalTok{    crlCache}\OperatorTok{.}\FunctionTok{put}\OperatorTok{(}\NormalTok{crlUrl}\OperatorTok{,}
            \KeywordTok{new} \FunctionTok{CachedCrl}\OperatorTok{(}\NormalTok{crlBytes}\OperatorTok{,}\NormalTok{ crl}\OperatorTok{.}\FunctionTok{getNextUpdate}\OperatorTok{()));}
\NormalTok{    log}\OperatorTok{.}\FunctionTok{debug}\OperatorTok{(}\StringTok{"CRL cached for \{\}, valid until \{\}"}\OperatorTok{,}
\NormalTok{            crlUrl}\OperatorTok{,}\NormalTok{ crl}\OperatorTok{.}\FunctionTok{getNextUpdate}\OperatorTok{());}

    \ControlFlowTok{return}\NormalTok{ crlBytes}\OperatorTok{;}
\OperatorTok{\}}
\end{Highlighting}
\end{Shaded}

\textbf{Cache warming for DSS construction.} When building a DSS
dictionary (for PAdES-B-LT or B-LTA), you need OCSP responses and CRLs
for the signer's entire certificate chain \emph{plus} the TSA's
certificate chain. A single signature can require 4-8 revocation data
lookups. Pre-fetch and cache this data before the signing flow begins
--- the DSS construction should be an offline operation that reads from
cache, not a live network operation.

\section{Signature Timestamps vs. Document
Timestamps}\label{_signature_timestamps_vs_document_timestamps}

PDF supports two kinds of timestamps, and they serve different purposes:

\subsection{Signature Timestamps}\label{_signature_timestamps}

A \textbf{signature timestamp} is embedded \emph{inside} the CMS
\texttt{SignedData} as an unsigned attribute of the \texttt{SignerInfo}.
It timestamps the \emph{signature value} --- proving when the signature
was created.

\begin{itemize}
\item
  \textbf{Where it lives}: CMS
  \texttt{SignerInfo.unsignedAttrs.id-aa-signatureTimeStampToken}
\item
  \textbf{What it timestamps}: The bytes of the
  \texttt{SignerInfo.signature} field
\item
  \textbf{PAdES level}: Required for PAdES-B-T and above
\item
  \textbf{Who creates it}: The signing application, immediately after
  computing the CMS signature
\end{itemize}

\subsection{Document Timestamps}\label{_document_timestamps}

A \textbf{document timestamp} is a separate signature in the PDF --- a
\texttt{/Sig} dictionary with \texttt{/SubFilter\ /ETSI.RFC3161} ---
that timestamps the \emph{entire document}.

\begin{itemize}
\item
  \textbf{Where it lives}: A separate signature dictionary in the PDF,
  applied as an incremental update
\item
  \textbf{What it timestamps}: The entire PDF file (via byte range),
  including all previous signatures and the DSS
\item
  \textbf{PAdES level}: Required for PAdES-B-LTA (the archival
  timestamp)
\item
  \textbf{Who creates it}: The document management system, as a
  maintenance operation
\end{itemize}

The key difference: a signature timestamp proves when a \emph{specific
signature} was created. A document timestamp proves that the
\emph{entire document} (including all signatures and validation data)
existed at a specific time.

\begin{verbatim}
// Signature timestamp (inside CMS)
/Sig <<
  /SubFilter /ETSI.CAdES.detached
  /Contents <CMS SignedData:
    SignerInfo:
      signedAttrs: [contentType, messageDigest, ESS-signing-certificate-v2]
      signature: <RSA/ECDSA signature bytes>
      unsignedAttrs: [
        signatureTimeStampToken: <TimeStampToken>  ← THIS
      ]
  >
>>

// Document timestamp (separate signature)
/Sig <<
  /SubFilter /ETSI.RFC3161             ← THIS
  /ByteRange [0 ... ... ...]          ← Covers entire file
  /Contents <TimeStampToken as CMS>
>>
\end{verbatim}

\section{The Document Security Store
(DSS)}\label{_the_document_security_store_dss}

\subsection{Purpose}\label{_purpose}

The DSS is a dictionary in the document catalog that stores all the data
needed to validate every signature in the document \emph{offline} ---
without contacting any external servers.

Without a DSS, validating a signature requires:

\begin{itemize}
\item
  Fetching the signer's certificate chain from the CMS and/or AIA URLs
\item
  Contacting OCSP responders to check revocation status
\item
  Downloading CRLs from distribution points
\end{itemize}

All of these require network access and functioning servers. Years from
now, these servers may not exist. The DSS solves this by embedding
everything locally.

\subsection{Building the DSS}\label{_building_the_dss}

Building a DSS requires collecting all the validation data and embedding
it as an incremental update:

\begin{Shaded}
\begin{Highlighting}[]
\CommentTok{/**}
 \CommentTok{*}\NormalTok{ Adds a DSS }\CommentTok{(}\NormalTok{Document Security Store}\CommentTok{)}\NormalTok{ dictionary to the PDF}\CommentTok{,}
 \CommentTok{*}\NormalTok{ embedding certificates}\CommentTok{,}\NormalTok{ OCSP responses}\CommentTok{,}\NormalTok{ and CRLs for all signatures}\CommentTok{.}
 \CommentTok{*/}
\KeywordTok{public} \DataTypeTok{void} \FunctionTok{addDSS}\OperatorTok{(}\NormalTok{PDDocument doc}\OperatorTok{,} \BuiltInTok{List}\OperatorTok{\textless{}}\NormalTok{ValidationData}\OperatorTok{\textgreater{}}\NormalTok{ validationData}\OperatorTok{)}
        \KeywordTok{throws} \BuiltInTok{IOException} \OperatorTok{\{}

\NormalTok{    COSDictionary dssDictionary }\OperatorTok{=} \KeywordTok{new} \FunctionTok{COSDictionary}\OperatorTok{();}
\NormalTok{    COSArray certsArray }\OperatorTok{=} \KeywordTok{new} \FunctionTok{COSArray}\OperatorTok{();}
\NormalTok{    COSArray ocspsArray }\OperatorTok{=} \KeywordTok{new} \FunctionTok{COSArray}\OperatorTok{();}
\NormalTok{    COSArray crlsArray }\OperatorTok{=} \KeywordTok{new} \FunctionTok{COSArray}\OperatorTok{();}

    \ControlFlowTok{for} \OperatorTok{(}\NormalTok{ValidationData vd }\OperatorTok{:}\NormalTok{ validationData}\OperatorTok{)} \OperatorTok{\{}
        \CommentTok{// Embed certificates}
        \ControlFlowTok{for} \OperatorTok{(}\DataTypeTok{byte}\OperatorTok{[]}\NormalTok{ certBytes }\OperatorTok{:}\NormalTok{ vd}\OperatorTok{.}\FunctionTok{getCertificateChain}\OperatorTok{())} \OperatorTok{\{}
\NormalTok{            PDStream certStream }\OperatorTok{=} \KeywordTok{new} \FunctionTok{PDStream}\OperatorTok{(}\NormalTok{doc}\OperatorTok{);}
            \ControlFlowTok{try} \OperatorTok{(}\BuiltInTok{OutputStream}\NormalTok{ os }\OperatorTok{=}\NormalTok{ certStream}\OperatorTok{.}\FunctionTok{createOutputStream}\OperatorTok{())} \OperatorTok{\{}
\NormalTok{                os}\OperatorTok{.}\FunctionTok{write}\OperatorTok{(}\NormalTok{certBytes}\OperatorTok{);}
            \OperatorTok{\}}
\NormalTok{            certsArray}\OperatorTok{.}\FunctionTok{add}\OperatorTok{(}\NormalTok{certStream}\OperatorTok{.}\FunctionTok{getCOSObject}\OperatorTok{());}
        \OperatorTok{\}}

        \CommentTok{// Embed OCSP responses}
        \ControlFlowTok{for} \OperatorTok{(}\DataTypeTok{byte}\OperatorTok{[]}\NormalTok{ ocspBytes }\OperatorTok{:}\NormalTok{ vd}\OperatorTok{.}\FunctionTok{getOcspResponses}\OperatorTok{())} \OperatorTok{\{}
\NormalTok{            PDStream ocspStream }\OperatorTok{=} \KeywordTok{new} \FunctionTok{PDStream}\OperatorTok{(}\NormalTok{doc}\OperatorTok{);}
            \ControlFlowTok{try} \OperatorTok{(}\BuiltInTok{OutputStream}\NormalTok{ os }\OperatorTok{=}\NormalTok{ ocspStream}\OperatorTok{.}\FunctionTok{createOutputStream}\OperatorTok{())} \OperatorTok{\{}
\NormalTok{                os}\OperatorTok{.}\FunctionTok{write}\OperatorTok{(}\NormalTok{ocspBytes}\OperatorTok{);}
            \OperatorTok{\}}
\NormalTok{            ocspsArray}\OperatorTok{.}\FunctionTok{add}\OperatorTok{(}\NormalTok{ocspStream}\OperatorTok{.}\FunctionTok{getCOSObject}\OperatorTok{());}
        \OperatorTok{\}}

        \CommentTok{// Embed CRLs}
        \ControlFlowTok{for} \OperatorTok{(}\DataTypeTok{byte}\OperatorTok{[]}\NormalTok{ crlBytes }\OperatorTok{:}\NormalTok{ vd}\OperatorTok{.}\FunctionTok{getCrls}\OperatorTok{())} \OperatorTok{\{}
\NormalTok{            PDStream crlStream }\OperatorTok{=} \KeywordTok{new} \FunctionTok{PDStream}\OperatorTok{(}\NormalTok{doc}\OperatorTok{);}
            \ControlFlowTok{try} \OperatorTok{(}\BuiltInTok{OutputStream}\NormalTok{ os }\OperatorTok{=}\NormalTok{ crlStream}\OperatorTok{.}\FunctionTok{createOutputStream}\OperatorTok{())} \OperatorTok{\{}
\NormalTok{                os}\OperatorTok{.}\FunctionTok{write}\OperatorTok{(}\NormalTok{crlBytes}\OperatorTok{);}
            \OperatorTok{\}}
\NormalTok{            crlsArray}\OperatorTok{.}\FunctionTok{add}\OperatorTok{(}\NormalTok{crlStream}\OperatorTok{.}\FunctionTok{getCOSObject}\OperatorTok{());}
        \OperatorTok{\}}
    \OperatorTok{\}}

    \ControlFlowTok{if} \OperatorTok{(}\NormalTok{certsArray}\OperatorTok{.}\FunctionTok{size}\OperatorTok{()} \OperatorTok{\textgreater{}} \DecValTok{0}\OperatorTok{)}\NormalTok{ dssDictionary}\OperatorTok{.}\FunctionTok{setItem}\OperatorTok{(}\NormalTok{COSName}\OperatorTok{.}\FunctionTok{getPDFName}\OperatorTok{(}\StringTok{"Certs"}\OperatorTok{),}\NormalTok{ certsArray}\OperatorTok{);}
    \ControlFlowTok{if} \OperatorTok{(}\NormalTok{ocspsArray}\OperatorTok{.}\FunctionTok{size}\OperatorTok{()} \OperatorTok{\textgreater{}} \DecValTok{0}\OperatorTok{)}\NormalTok{ dssDictionary}\OperatorTok{.}\FunctionTok{setItem}\OperatorTok{(}\NormalTok{COSName}\OperatorTok{.}\FunctionTok{getPDFName}\OperatorTok{(}\StringTok{"OCSPs"}\OperatorTok{),}\NormalTok{ ocspsArray}\OperatorTok{);}
    \ControlFlowTok{if} \OperatorTok{(}\NormalTok{crlsArray}\OperatorTok{.}\FunctionTok{size}\OperatorTok{()} \OperatorTok{\textgreater{}} \DecValTok{0}\OperatorTok{)}\NormalTok{ dssDictionary}\OperatorTok{.}\FunctionTok{setItem}\OperatorTok{(}\NormalTok{COSName}\OperatorTok{.}\FunctionTok{getPDFName}\OperatorTok{(}\StringTok{"CRLs"}\OperatorTok{),}\NormalTok{ crlsArray}\OperatorTok{);}

    \CommentTok{// Add DSS to document catalog}
\NormalTok{    doc}\OperatorTok{.}\FunctionTok{getDocumentCatalog}\OperatorTok{().}\FunctionTok{getCOSObject}\OperatorTok{()}
            \OperatorTok{.}\FunctionTok{setItem}\OperatorTok{(}\NormalTok{COSName}\OperatorTok{.}\FunctionTok{getPDFName}\OperatorTok{(}\StringTok{"DSS"}\OperatorTok{),}\NormalTok{ dssDictionary}\OperatorTok{);}
\OperatorTok{\}}
\end{Highlighting}
\end{Shaded}

\subsection{VRI: Validation Related
Information}\label{_vri_validation_related_information}

The DSS can optionally include a \texttt{/VRI} sub-dictionary that maps
each signature to its specific validation data. The VRI keys are the
uppercase hex SHA-1 hash of the signature value (the bytes in
\texttt{/Contents}):

\begin{verbatim}
/DSS <<
  /Certs [...]
  /OCSPs [...]
  /VRI <<
    /A1B2C3D4E5F6... <<        % SHA-1 hash of signature 1's /Contents
      /Cert [30 0 R 31 0 R]    % Certs for this specific signature
      /OCSP [33 0 R]           % OCSP for this specific signature
      /TU (D:20260308120000Z)  % Time of validation
    >>
    /F6E5D4C3B2A1... <<        % SHA-1 hash of signature 2's /Contents
      /Cert [30 0 R 32 0 R]
      /OCSP [34 0 R]
      /TU (D:20260309140000Z)
    >>
  >>
>>
\end{verbatim}

VRI is useful for multi-signature documents where each signature has
different validation data. Without VRI, a validator must try all
certificates and OCSP responses in the DSS against each signature to
find the right match. With VRI, the mapping is explicit.

The VRI key is a SHA-1 hash --- yes, SHA-1. This is a historical
artifact from when the DSS was first specified. SHA-1 is used here only
as a key derivation function (to generate a unique dictionary key from
the signature value), not for cryptographic security. The security of
the DSS does not depend on SHA-1 collision resistance.

\section{Fetching Revocation Data}\label{_fetching_revocation_data}

\subsection{OCSP (Online Certificate Status
Protocol)}\label{_ocsp_online_certificate_status_protocol}

OCSP is the preferred mechanism for checking certificate revocation. You
send a request to the CA's OCSP responder with the certificate's serial
number, and it responds with "good," "revoked," or "unknown."

\begin{Shaded}
\begin{Highlighting}[]
\CommentTok{/**}
 \CommentTok{*}\NormalTok{ Fetches an OCSP response for the given certificate}\CommentTok{.}
 \CommentTok{*}\NormalTok{ Returns the DER}\CommentTok{{-}}\NormalTok{encoded OCSP response bytes}\CommentTok{.}
 \CommentTok{*/}
\KeywordTok{public} \DataTypeTok{byte}\OperatorTok{[]} \FunctionTok{fetchOCSPResponse}\OperatorTok{(}\BuiltInTok{X509Certificate}\NormalTok{ cert}\OperatorTok{,}
        \BuiltInTok{X509Certificate}\NormalTok{ issuer}\OperatorTok{)} \KeywordTok{throws} \BuiltInTok{Exception} \OperatorTok{\{}

    \CommentTok{// 1. Find the OCSP responder URL from the certificate\textquotesingle{}s AIA extension}
    \BuiltInTok{String}\NormalTok{ ocspUrl }\OperatorTok{=} \FunctionTok{extractOcspUrl}\OperatorTok{(}\NormalTok{cert}\OperatorTok{);}
    \ControlFlowTok{if} \OperatorTok{(}\NormalTok{ocspUrl }\OperatorTok{==} \KeywordTok{null}\OperatorTok{)} \OperatorTok{\{}
\NormalTok{        log}\OperatorTok{.}\FunctionTok{warn}\OperatorTok{(}\StringTok{"No OCSP URL in certificate AIA for: \{\}"}\OperatorTok{,}
\NormalTok{                cert}\OperatorTok{.}\FunctionTok{getSubjectX500Principal}\OperatorTok{());}
        \ControlFlowTok{return} \KeywordTok{null}\OperatorTok{;}
    \OperatorTok{\}}

    \CommentTok{// 2. Build the OCSP request}
\NormalTok{    CertificateID certId }\OperatorTok{=} \KeywordTok{new} \FunctionTok{CertificateID}\OperatorTok{(}
            \KeywordTok{new} \FunctionTok{JcaDigestCalculatorProviderBuilder}\OperatorTok{().}\FunctionTok{build}\OperatorTok{()}
                    \OperatorTok{.}\FunctionTok{get}\OperatorTok{(}\NormalTok{CertificateID}\OperatorTok{.}\FunctionTok{HASH\_SHA1}\OperatorTok{),}
            \KeywordTok{new} \FunctionTok{JcaX509CertificateHolder}\OperatorTok{(}\NormalTok{issuer}\OperatorTok{),}
\NormalTok{            cert}\OperatorTok{.}\FunctionTok{getSerialNumber}\OperatorTok{());}

\NormalTok{    OCSPReqBuilder reqBuilder }\OperatorTok{=} \KeywordTok{new} \FunctionTok{OCSPReqBuilder}\OperatorTok{();}
\NormalTok{    reqBuilder}\OperatorTok{.}\FunctionTok{addRequest}\OperatorTok{(}\NormalTok{certId}\OperatorTok{);}
\NormalTok{    OCSPReq ocspReq }\OperatorTok{=}\NormalTok{ reqBuilder}\OperatorTok{.}\FunctionTok{build}\OperatorTok{();}

    \CommentTok{// 3. Send via HTTP}
    \BuiltInTok{HttpURLConnection}\NormalTok{ conn }\OperatorTok{=} \OperatorTok{(}\BuiltInTok{HttpURLConnection}\OperatorTok{)}
            \KeywordTok{new} \BuiltInTok{URL}\OperatorTok{(}\NormalTok{ocspUrl}\OperatorTok{).}\FunctionTok{openConnection}\OperatorTok{();}
\NormalTok{    conn}\OperatorTok{.}\FunctionTok{setDoOutput}\OperatorTok{(}\KeywordTok{true}\OperatorTok{);}
\NormalTok{    conn}\OperatorTok{.}\FunctionTok{setRequestMethod}\OperatorTok{(}\StringTok{"POST"}\OperatorTok{);}
\NormalTok{    conn}\OperatorTok{.}\FunctionTok{setRequestProperty}\OperatorTok{(}\StringTok{"Content{-}Type"}\OperatorTok{,} \StringTok{"application/ocsp{-}request"}\OperatorTok{);}
\NormalTok{    conn}\OperatorTok{.}\FunctionTok{getOutputStream}\OperatorTok{().}\FunctionTok{write}\OperatorTok{(}\NormalTok{ocspReq}\OperatorTok{.}\FunctionTok{getEncoded}\OperatorTok{());}

    \DataTypeTok{byte}\OperatorTok{[]}\NormalTok{ respBytes }\OperatorTok{=}\NormalTok{ conn}\OperatorTok{.}\FunctionTok{getInputStream}\OperatorTok{().}\FunctionTok{readAllBytes}\OperatorTok{();}

    \CommentTok{// 4. Validate the response}
\NormalTok{    OCSPResp ocspResp }\OperatorTok{=} \KeywordTok{new} \FunctionTok{OCSPResp}\OperatorTok{(}\NormalTok{respBytes}\OperatorTok{);}
    \ControlFlowTok{if} \OperatorTok{(}\NormalTok{ocspResp}\OperatorTok{.}\FunctionTok{getStatus}\OperatorTok{()} \OperatorTok{!=}\NormalTok{ OCSPResp}\OperatorTok{.}\FunctionTok{SUCCESSFUL}\OperatorTok{)} \OperatorTok{\{}
\NormalTok{        log}\OperatorTok{.}\FunctionTok{error}\OperatorTok{(}\StringTok{"OCSP response status: \{\}"}\OperatorTok{,}\NormalTok{ ocspResp}\OperatorTok{.}\FunctionTok{getStatus}\OperatorTok{());}
        \ControlFlowTok{return} \KeywordTok{null}\OperatorTok{;}
    \OperatorTok{\}}

\NormalTok{    BasicOCSPResp basicResp }\OperatorTok{=} \OperatorTok{(}\NormalTok{BasicOCSPResp}\OperatorTok{)}\NormalTok{ ocspResp}\OperatorTok{.}\FunctionTok{getResponseObject}\OperatorTok{();}
\NormalTok{    SingleResp singleResp }\OperatorTok{=}\NormalTok{ basicResp}\OperatorTok{.}\FunctionTok{getResponses}\OperatorTok{()[}\DecValTok{0}\OperatorTok{];}

    \ControlFlowTok{if} \OperatorTok{(}\NormalTok{singleResp}\OperatorTok{.}\FunctionTok{getCertStatus}\OperatorTok{()} \OperatorTok{==}\NormalTok{ CertificateStatus}\OperatorTok{.}\FunctionTok{GOOD}\OperatorTok{)} \OperatorTok{\{}
\NormalTok{        log}\OperatorTok{.}\FunctionTok{info}\OperatorTok{(}\StringTok{"OCSP: certificate is GOOD (produced: \{\})"}\OperatorTok{,}
\NormalTok{                basicResp}\OperatorTok{.}\FunctionTok{getProducedAt}\OperatorTok{());}
        \ControlFlowTok{return}\NormalTok{ respBytes}\OperatorTok{;}
    \OperatorTok{\}} \ControlFlowTok{else} \OperatorTok{\{}
\NormalTok{        log}\OperatorTok{.}\FunctionTok{error}\OperatorTok{(}\StringTok{"OCSP: certificate status is NOT GOOD: \{\}"}\OperatorTok{,}
\NormalTok{                singleResp}\OperatorTok{.}\FunctionTok{getCertStatus}\OperatorTok{());}
        \ControlFlowTok{return} \KeywordTok{null}\OperatorTok{;}  \CommentTok{// Certificate is revoked or unknown}
    \OperatorTok{\}}
\OperatorTok{\}}
\end{Highlighting}
\end{Shaded}

\subsection{CRL (Certificate Revocation
List)}\label{_crl_certificate_revocation_list}

CRLs are the fallback when OCSP is not available. A CRL is a signed list
of all revoked certificates issued by a CA.

\begin{Shaded}
\begin{Highlighting}[]
\CommentTok{/**}
 \CommentTok{*}\NormalTok{ Downloads a CRL from the certificate}\CommentTok{\textquotesingle{}}\NormalTok{s CRL Distribution Point}\CommentTok{.}
 \CommentTok{*}\NormalTok{ Returns the DER}\CommentTok{{-}}\NormalTok{encoded CRL bytes}\CommentTok{.}
 \CommentTok{*/}
\KeywordTok{public} \DataTypeTok{byte}\OperatorTok{[]} \FunctionTok{fetchCRL}\OperatorTok{(}\BuiltInTok{X509Certificate}\NormalTok{ cert}\OperatorTok{)} \KeywordTok{throws} \BuiltInTok{Exception} \OperatorTok{\{}
    \BuiltInTok{String}\NormalTok{ crlUrl }\OperatorTok{=} \FunctionTok{extractCrlUrl}\OperatorTok{(}\NormalTok{cert}\OperatorTok{);}
    \ControlFlowTok{if} \OperatorTok{(}\NormalTok{crlUrl }\OperatorTok{==} \KeywordTok{null}\OperatorTok{)} \OperatorTok{\{}
\NormalTok{        log}\OperatorTok{.}\FunctionTok{warn}\OperatorTok{(}\StringTok{"No CRL URL in certificate for: \{\}"}\OperatorTok{,}
\NormalTok{                cert}\OperatorTok{.}\FunctionTok{getSubjectX500Principal}\OperatorTok{());}
        \ControlFlowTok{return} \KeywordTok{null}\OperatorTok{;}
    \OperatorTok{\}}

    \BuiltInTok{HttpURLConnection}\NormalTok{ conn }\OperatorTok{=} \OperatorTok{(}\BuiltInTok{HttpURLConnection}\OperatorTok{)} \KeywordTok{new} \BuiltInTok{URL}\OperatorTok{(}\NormalTok{crlUrl}\OperatorTok{).}\FunctionTok{openConnection}\OperatorTok{();}
\NormalTok{    conn}\OperatorTok{.}\FunctionTok{setConnectTimeout}\OperatorTok{(}\DecValTok{10000}\OperatorTok{);}
\NormalTok{    conn}\OperatorTok{.}\FunctionTok{setReadTimeout}\OperatorTok{(}\DecValTok{30000}\OperatorTok{);}

    \DataTypeTok{byte}\OperatorTok{[]}\NormalTok{ crlBytes }\OperatorTok{=}\NormalTok{ conn}\OperatorTok{.}\FunctionTok{getInputStream}\OperatorTok{().}\FunctionTok{readAllBytes}\OperatorTok{();}

    \CommentTok{// Validate the CRL}
    \BuiltInTok{CertificateFactory}\NormalTok{ cf }\OperatorTok{=} \BuiltInTok{CertificateFactory}\OperatorTok{.}\FunctionTok{getInstance}\OperatorTok{(}\StringTok{"X.509"}\OperatorTok{);}
    \BuiltInTok{X509CRL}\NormalTok{ crl }\OperatorTok{=} \OperatorTok{(}\BuiltInTok{X509CRL}\OperatorTok{)}\NormalTok{ cf}\OperatorTok{.}\FunctionTok{generateCRL}\OperatorTok{(}\KeywordTok{new} \BuiltInTok{ByteArrayInputStream}\OperatorTok{(}\NormalTok{crlBytes}\OperatorTok{));}

\NormalTok{    log}\OperatorTok{.}\FunctionTok{info}\OperatorTok{(}\StringTok{"CRL downloaded: issuer=\{\}, entries=\{\}, nextUpdate=\{\}"}\OperatorTok{,}
\NormalTok{            crl}\OperatorTok{.}\FunctionTok{getIssuerX500Principal}\OperatorTok{(),}
\NormalTok{            crl}\OperatorTok{.}\FunctionTok{getRevokedCertificates}\OperatorTok{()} \OperatorTok{!=} \KeywordTok{null} \OperatorTok{?}
\NormalTok{                    crl}\OperatorTok{.}\FunctionTok{getRevokedCertificates}\OperatorTok{().}\FunctionTok{size}\OperatorTok{()} \OperatorTok{:} \DecValTok{0}\OperatorTok{,}
\NormalTok{            crl}\OperatorTok{.}\FunctionTok{getNextUpdate}\OperatorTok{());}

    \ControlFlowTok{return}\NormalTok{ crlBytes}\OperatorTok{;}
\OperatorTok{\}}
\end{Highlighting}
\end{Shaded}

CRLs can be \textbf{very large} --- some CAs have CRLs of 500 KB or
more. Embedding a large CRL in the DSS significantly increases the PDF
file size. Prefer OCSP when available --- OCSP responses are typically
1-3 KB.

Also be aware of \textbf{delta CRLs} --- some CAs publish only the
changes since the last full CRL. If the certificate's CRL Distribution
Point specifies a delta CRL, you may need both the full CRL and the
delta CRL for complete validation.

\section{The Complete B-LTA
Workflow}\label{_the_complete_b_lta_workflow}

Putting it all together, here's the end-to-end workflow for producing a
PAdES-B-LTA signature:

\begin{verbatim}
1. Create PAdES-B-B signature
   └── CMS with ESS-signing-certificate-v2, SHA-256, SubFilter ETSI.CAdES.detached

2. Add signature timestamp → PAdES-B-T
   └── Request RFC 3161 timestamp over the signature value
   └── Embed as unsigned CMS attribute

3. Fetch validation data
   └── Build certificate chain (signer → intermediates → root)
   └── Fetch OCSP response for each certificate
   └── Fetch CRL as fallback where OCSP unavailable
   └── Fetch TSA's certificate chain and revocation data too

4. Create DSS dictionary → PAdES-B-LT
   └── Embed all certificates as streams
   └── Embed all OCSP responses as streams
   └── Embed all CRLs as streams
   └── Optionally create VRI entries
   └── Add DSS to document catalog as incremental update

5. Add document timestamp → PAdES-B-LTA
   └── Request RFC 3161 timestamp over the entire file (including DSS)
   └── Apply as separate /Sig dictionary with SubFilter ETSI.RFC3161
   └── Save as incremental update
\end{verbatim}

Each step is an incremental update to the PDF. The final file contains
the original document, the signature, the DSS, and the document
timestamp --- everything needed for indefinite validation.

\section{Key Takeaways}\label{_key_takeaways_6}

\begin{itemize}
\item
  \textbf{RFC 3161 timestamps} are CMS-signed proofs that data existed
  at a specific time --- issued by Trusted Timestamp Authorities
\item
  \textbf{Signature timestamps} (in CMS unsigned attributes) prove when
  a specific signature was created
\item
  \textbf{Document timestamps} (separate \texttt{/Sig} with SubFilter
  \texttt{ETSI.RFC3161}) prove when the entire document existed --- used
  for archival
\item
  The \textbf{DSS dictionary} embeds certificates, OCSP responses, and
  CRLs directly in the PDF for offline validation
\item
  \textbf{VRI entries} map specific signatures to their validation data
  using SHA-1 hash keys
\item
  \textbf{OCSP} is preferred over CRLs for revocation checking (smaller,
  more specific)
\item
  \textbf{TSA selection} matters: trust store compatibility (AATL/EUTL),
  availability, cost, and response time
\item
  \textbf{PAdES-B-LTA} is the complete package: signature + timestamp +
  validation data + archival timestamp
\item
  \textbf{Re-timestamping} extends validity indefinitely --- a scheduled
  maintenance operation for archival documents
\item
  Always implement \textbf{TSA failover} --- a TSA outage should not
  prevent signing
\item
  \textbf{Probe TSAs} in the background at startup and periodically ---
  validate Content-Type headers, not just HTTP 200
\item
  \textbf{Cache OCSP responses} aggressively --- they have explicit
  validity windows (\texttt{thisUpdate}/\texttt{nextUpdate}) and
  eliminate redundant network calls in batch scenarios
\item
  \textbf{Cache CRLs} by distribution point URL --- CRLs change
  infrequently and can be very large; avoid re-downloading within the
  validity window
\item
  \textbf{TSA responses cannot be cached} --- each is unique by design
  --- but connection pooling and pre-warming significantly reduce
  latency
\end{itemize}

\chapter{eIDAS \& Qualified
Signatures}\label{_eidas_qualified_signatures}

\chapter{The eIDAS Regulation}\label{_the_eidas_regulation}

In 2014, the European Union did something that no other jurisdiction had
managed: it created a legal framework where a digital signature --- a
sequence of bytes produced by a cryptographic algorithm --- carries the
same legal weight as ink on paper, in all 27 member states
simultaneously. That framework is called eIDAS, and whether you love it
or hate it, if your software produces or validates digital signatures
that touch Europe, you need to understand it.

The Regulation (EU) No 910/2014 (eIDAS) --- its full title is
"Regulation (EU) No 910/2014 of the European Parliament and of the
Council of 23 July 2014 on electronic identification and trust services
for electronic transactions in the internal market" --- is a 38-page
regulation that created the legal and institutional infrastructure for
electronic signatures, seals, timestamps, and related trust services
across the entire European Union. It replaced a failed directive that
had been limping along for 15 years, and it did so by making one
critical architectural decision: it chose to be a \emph{regulation}, not
a \emph{directive}.

That distinction matters enormously, and understanding why requires a
brief detour through European law.

\section{A Brief History: The Electronic Signatures
Directive}\label{_a_brief_history_the_electronic_signatures_directive}

\subsection{Directive 1999/93/EC}\label{_directive_199993ec}

Before eIDAS, there was the \textbf{Electronic Signatures Directive}
(Directive 1999/93/EC), adopted on December 13, 1999. The directive was
born in the heady optimism of the dot-com era, when legislators across
the world were scrambling to give electronic commerce a legal
foundation. The United States had passed the E-Sign Act. UNCITRAL had
published a Model Law on Electronic Signatures. Europe wanted its own
framework, and it wanted one that would work across all (then 15) member
states.

The directive's core idea was sound: define what an "advanced electronic
signature" is, define what a "qualified certificate" is, establish that
qualified electronic signatures should not be denied legal effect, and
require member states to set up supervisory bodies for certification
service providers.

\subsection{Why Directives Fail}\label{_why_directives_fail}

Here is where European institutional mechanics become relevant to your
code.

In EU law, a \textbf{directive} tells member states \emph{what result to
achieve} but lets each country decide \emph{how} to achieve it through
national legislation. A \textbf{regulation}, by contrast, is directly
applicable in all member states without any national transposition ---
it becomes law everywhere on the same day, with the same text.

The Electronic Signatures Directive was a directive. This meant that
each member state had to transpose it into national law, and each one
did so differently. Germany created the \emph{Signaturgesetz} (Signature
Law). France created \emph{le decret n. 2001-272}. Italy created the
\emph{Codice dell'Amministrazione Digitale}. Each country interpreted
the directive's requirements through its own legal tradition, its own
bureaucratic structures, and its own technical preferences.

The result was exactly what you would predict: 15 (later 28) different
national implementations of the "same" legal framework, with subtly
different definitions of what counted as an "advanced" signature,
different requirements for certification service providers, different
supervisory structures, and --- critically for developers --- different
technical standards that qualified certificates had to meet.

The Electronic Signatures Directive used the term "Certification Service
Provider" (CSP). eIDAS renamed these to "Trust Service Provider" (TSP)
and, for the supervised ones, "Qualified Trust Service Provider" (QTSP).
The terminology change was deliberate: eIDAS covers more than just
certificates --- it covers timestamps, seals, website authentication,
and electronic delivery services. If you encounter "CSP" in older
documentation, read it as "TSP."

\subsection{Cross-Border Failure}\label{_cross_border_failure}

The most damaging consequence of the directive's fragmentation was the
failure of cross-border recognition. In theory, a qualified electronic
signature from Germany should have been accepted in France, and vice
versa. In practice, it often was not. Each country's government
procurement portals, tax systems, and regulatory filing platforms had
been built to recognize only their own national qualified certificates.
A German doctor trying to sign a cross-border prescription
electronically for a French patient would find that the French system
did not recognize the German qualified certificate.

A 2012 study by the European Commission found that cross-border
recognition of electronic signatures was "marginal at best." Of 13
member states surveyed, only 5 had any mechanism at all for recognizing
qualified certificates from other EU countries. The directive had been
in force for 13 years, and the internal market for electronic signatures
was essentially non-existent.

This was the political context in which eIDAS was born.

\subsection{The Lessons Learned}\label{_the_lessons_learned}

The Commission identified several root causes for the directive's
failure:

\begin{enumerate}
\def\labelenumi{\arabic{enumi}.}
\item
  \textbf{Directive vs. Regulation}: national transposition created
  divergent implementations
\item
  \textbf{No mandatory trust infrastructure}: there was no EU-wide list
  of qualified providers that all member states were required to accept
\item
  \textbf{Vague technical requirements}: the directive referenced
  "secure signature creation devices" without specifying what that meant
  technically
\item
  \textbf{No supervision framework}: how supervisory bodies should
  operate, what they should check, and what consequences non-compliance
  should carry were all left to national discretion
\item
  \textbf{Certificates only}: the directive covered only electronic
  signatures, ignoring related services like timestamps, seals, and
  electronic delivery
\end{enumerate}

eIDAS addressed every single one of these failures.

\section{eIDAS: Architecture of the
Regulation}\label{_eidas_architecture_of_the_regulation}

\subsection{Regulation, Not Directive}\label{_regulation_not_directive}

eIDAS is a regulation. Article 288 of the Treaty on the Functioning of
the European Union says: "A regulation shall have general application.
It shall be binding in its entirety and directly applicable in all
Member States."

This means: - The text is the same in all 27 member states (plus EEA
countries through incorporation) - No national transposition is needed
or permitted --- the regulation \emph{is} the law - National courts
apply the regulation directly; citizens and businesses can invoke it
directly - Divergent national interpretations can be challenged at the
Court of Justice of the European Union

For developers, this has a profound practical consequence: the rules are
the same everywhere. When eIDAS says that a qualified electronic
signature has the legal effect equivalent to a handwritten signature, it
means exactly that in Germany, France, Italy, Estonia, and all other
member states. You do not need to check 27 different national laws to
know what your signature is worth.

\subsection{Scope and Structure}\label{_scope_and_structure}

The eIDAS regulation has two main parts:

\textbf{Chapter II: Electronic Identification} (Articles 6-12) ---
establishes a framework for mutual recognition of national electronic
identification schemes (eID). This is the part that deals with
government-issued digital identities (like the German
\emph{Personalausweis} or the Estonian \emph{ID-kaart}) and how they
should be recognized across borders. We will not cover this in detail
because it relates to identity management rather than document signing.

\textbf{Chapter III: Trust Services} (Articles 13-45) --- this is where
the electronic signature framework lives, and it is the focus of this
chapter and the next four chapters.

The trust services chapter covers: - Electronic signatures (Articles
25-27) - Electronic seals (Articles 35-40) - Electronic timestamps
(Articles 41-42) - Electronic delivery services (Articles 43-44) -
Website authentication certificates (Article 45) - General provisions on
trust service providers (Articles 13-24)

eIDAS does not mandate \emph{how} electronic signatures must be
implemented technically. It defines \emph{legal categories} and
\emph{legal effects}, then defers to ETSI and CEN for the technical
standards. The regulation deliberately separates legal requirements from
technical implementation --- a wise decision that has allowed the
technical standards (PAdES, CAdES, XAdES) to evolve independently of the
legal framework.

When eIDAS says "advanced electronic signature," it is defining a legal
concept with four properties. When ETSI EN 319 142 says "PAdES-B-T," it
is defining a technical implementation that satisfies those properties.
The mapping between legal categories and technical standards is
documented in Commission Implementing Decision (EU) 2015/1506.

\section{The Three Tiers of Electronic
Signatures}\label{_the_three_tiers_of_electronic_signatures}

The heart of eIDAS for developers is its three-tier classification of
electronic signatures. Each tier has progressively stricter requirements
and progressively stronger legal effects.

\subsection{Tier 1: Electronic Signature (Article
3(10))}\label{_tier_1_electronic_signature_article_310}

The broadest tier. Article 3(10) defines an electronic signature as:

\begin{quote}
"electronic signature" means data in electronic form which is attached
to or logically associated with other data in electronic form and which
is used by the signatory to sign

---  Regulation (EU) No 910/2014 (eIDAS), Article 3(10)
\end{quote}

That is it. An electronic signature under eIDAS can be: - A scanned
image of a handwritten signature pasted into a PDF - A typed name at the
bottom of an email - A checkbox on a web form saying "I agree" - A click
of a "Sign" button that records user intent

There is no cryptographic requirement. There is no certificate
requirement. There is no requirement for any specific technology at all.

The legal effect of a basic electronic signature is defined in Article
25(1): it "shall not be denied legal effect and admissibility as
evidence in legal proceedings solely on the grounds that it is in an
electronic form." This is a non-discrimination principle --- courts
cannot reject evidence \emph{just} because it is electronic. But it says
nothing about the \emph{weight} of that evidence. A typed name has much
less evidentiary value than a cryptographic signature with a qualified
certificate.

For developers building signing workflows: many business processes do
not require more than a basic electronic signature. Click-to-sign,
draw-to-sign, and type-to-sign all count as electronic signatures under
eIDAS. If your use case is internal corporate approvals, non-binding
agreements, or informal consent, a basic electronic signature may be
legally sufficient. Check with your legal team, but do not assume you
always need cryptography.

\subsection{Tier 2: Advanced Electronic Signature (Article 3(11) and
Article
26)}\label{_tier_2_advanced_electronic_signature_article_311_and_article_26}

The middle tier. Article 3(11) defines an advanced electronic signature
as an electronic signature that meets the requirements set out in
Article 26.

Article 26 specifies four requirements:

\begin{quote}
An advanced electronic signature shall meet the following requirements:

(a) it is uniquely linked to the signatory;

(b) it is capable of identifying the signatory;

(c) it is created using electronic signature creation data that the
signatory can, with a high level of confidence, use under his sole
control; and

(d) it is linked to the data signed therewith in such a way that any
subsequent change in the data is detectable.

---  Regulation (EU) No 910/2014 (eIDAS), Article 26
\end{quote}

Let us examine each requirement from a technical perspective.

\subsubsection{(a) Uniquely linked to the
signatory}\label{_a_uniquely_linked_to_the_signatory}

The signature must be produced using data (typically a private key) that
is uniquely associated with one specific natural person. In practice,
this means: a private key that belongs to exactly one person, used to
produce a cryptographic signature that can only have been created with
that key.

A shared private key used by multiple employees would fail this
requirement. A department-level signing key would fail this requirement.
The key must be tied to a specific individual.

\subsubsection{(b) Capable of identifying the
signatory}\label{_b_capable_of_identifying_the_signatory}

The signature (together with its associated data, such as a certificate)
must make it possible to determine \emph{who} signed. A raw RSA or ECDSA
signature over a hash does not, by itself, identify anyone --- it just
proves that whoever held the private key produced the signature.
Combined with an X.509 certificate that binds the public key to an
identity (name, organization, email), the requirement is satisfied.

\subsubsection{(c) Sole control}\label{_c_sole_control}

The signatory must be the only person who can use the signing key. "With
a high level of confidence" is the qualifier --- absolute sole control
is impossible to guarantee (keys can be stolen), but the system must
make unauthorized use difficult.

Sole control mechanisms include: - \textbf{Password-protected PKCS\#12
keystore}: the signer must know the password - \textbf{Smart card with
PIN}: the signer must possess the card and know the PIN - \textbf{HSM
with authentication}: the signer must authenticate (username + password
+ OTP) before the HSM will sign - \textbf{Biometric activation}:
fingerprint or face recognition to authorize signing

"Sole control" is where many implementations fail the advanced signature
test. If a signing key is stored in a server-side keystore and an
application signs on behalf of users without meaningful per-signature
authentication, the resulting signatures are arguably \emph{not}
advanced electronic signatures --- because the user does not have sole
control over the key.

The fact that the user logged into the application is not sufficient.
The authentication must be linked to the \emph{signing action}, not just
the \emph{session}. This is the distinction that the SCAL (Sole Control
Assurance Level) framework addresses --- more on this in Chapter 22.

\subsubsection{(d) Linked to the data}\label{_d_linked_to_the_data}

Any change to the signed data must be detectable. This is the integrity
property that every digital signature scheme provides: if you change
even one bit of the signed data, the signature will no longer verify.
For PDF signatures, the byte range mechanism (Chapter 13) ensures this
linkage.

\subsection{Tier 3: Qualified Electronic Signature (Article
3(12))}\label{_tier_3_qualified_electronic_signature_article_312}

The highest tier, and the most consequential. Article 3(12) defines a
qualified electronic signature as:

\begin{quote}
"qualified electronic signature" means an advanced electronic signature
that is created by a qualified electronic signature creation device, and
which is based on a qualified certificate for electronic signatures

---  Regulation (EU) No 910/2014 (eIDAS), Article 3(12)
\end{quote}

A QES is an advanced electronic signature (all four Article 26
properties) that \emph{additionally}: 1. Uses a \textbf{Qualified
Signature Creation Device (QSCD)} --- certified hardware that protects
the signing key 2. Is based on a \textbf{Qualified Certificate} ---
issued by a Qualified Trust Service Provider (QTSP) listed in the EU
Trusted List

The legal effect of a QES is defined in Article 25(2): "A qualified
electronic signature shall have the equivalent legal effect of a
handwritten signature." This is not a non-discrimination principle ---
it is an \emph{equivalence} rule. A QES is legally equal to pen on
paper. No court in any EU member state can treat it differently.

Furthermore, Article 25(3) establishes mutual recognition: "A qualified
electronic signature based on a qualified certificate issued in one
Member State shall be recognised as a qualified electronic signature in
all other Member States." This is the cross-border recognition that the
1999 directive failed to achieve.

\section{Key Articles in Detail}\label{_key_articles_in_detail}

\subsection{Article 25: Legal Effects of Electronic
Signatures}\label{_article_25_legal_effects_of_electronic_signatures}

Article 25 is the most cited provision in eIDAS, and it deserves careful
reading.

\begin{longtable}[]{@{}
  >{\raggedright\arraybackslash}p{(\linewidth - 4\tabcolsep) * \real{0.1429}}
  >{\raggedright\arraybackslash}p{(\linewidth - 4\tabcolsep) * \real{0.4286}}
  >{\raggedright\arraybackslash}p{(\linewidth - 4\tabcolsep) * \real{0.4286}}@{}}
\caption{Article 25 --- Legal effects summarized}\tabularnewline
\toprule\noalign{}
\begin{minipage}[b]{\linewidth}\raggedright
Paragraph
\end{minipage} & \begin{minipage}[b]{\linewidth}\raggedright
Rule
\end{minipage} & \begin{minipage}[b]{\linewidth}\raggedright
Practical Meaning
\end{minipage} \\
\midrule\noalign{}
\endfirsthead
\toprule\noalign{}
\begin{minipage}[b]{\linewidth}\raggedright
Paragraph
\end{minipage} & \begin{minipage}[b]{\linewidth}\raggedright
Rule
\end{minipage} & \begin{minipage}[b]{\linewidth}\raggedright
Practical Meaning
\end{minipage} \\
\midrule\noalign{}
\endhead
\bottomrule\noalign{}
\endlastfoot
25(1) & Electronic signatures shall not be denied legal effect solely
because they are electronic & Non-discrimination: courts cannot reject
electronic evidence just because it is not paper \\
25(2) & QES shall have the equivalent legal effect of a handwritten
signature & Legal equivalence: a QES \emph{is} a handwritten signature
for all legal purposes \\
25(3) & QES based on a qualified certificate from one member state shall
be recognized in all member states & Cross-border recognition: no
country can reject another country's QES \\
\end{longtable}

Note what Article 25 does \emph{not} say: it does not say that advanced
electronic signatures (AdES) have the legal effect of handwritten
signatures. An AdES benefits only from the non-discrimination rule of
25(1) --- it cannot be rejected \emph{solely} because it is electronic,
but a court is free to assign it less evidentiary weight than a
handwritten signature based on the circumstances.

This creates a clear hierarchy: - \textbf{QES}: automatic legal
equivalence to handwritten (Article 25(2)) - \textbf{AdES}: cannot be
dismissed as evidence just for being electronic (Article 25(1)), but the
court decides its weight - \textbf{Basic electronic signature}: same as
AdES --- cannot be dismissed for being electronic, but its evidentiary
weight is at the court's discretion

Many developers (and their managers) misunderstand Article 25(1) as
meaning "all electronic signatures are legally valid." That is not what
it says. It says they cannot be denied legal effect \emph{solely}
because they are electronic. A typed name at the bottom of an email is
an electronic signature under eIDAS, but it provides very little
evidence about who actually sent the email, whether they intended to be
bound by its contents, or whether the email was tampered with. A court
will weigh this evidence accordingly.

Only QES gets automatic equivalence. Everything else is a matter of
evidence and judicial discretion.

\subsection{Article 28: Qualified Certificates for Electronic
Signatures}\label{_article_28_qualified_certificates_for_electronic_signatures}

Article 28 requires that qualified certificates contain specific
information, listed in Annex I of the regulation. A qualified
certificate must contain:

\begin{enumerate}
\def\labelenumi{\arabic{enumi}.}
\item
  An indication (at least in machine-processable form) that the
  certificate is a qualified certificate
\item
  A set of data representing the QTSP, including the member state where
  it is established and its name
\item
  The signatory's name (or a pseudonym, clearly identified as such)
\item
  Electronic signature validation data (the public key) corresponding to
  the electronic signature creation data (the private key)
\item
  The period of validity
\item
  The certificate serial number
\item
  The advanced electronic signature or seal of the issuing QTSP
\item
  The location where the certificate supporting the QTSP's
  signature/seal is available free of charge
\item
  The location of the QTSP's certificate revocation status services
\end{enumerate}

In X.509 terms, this translates to specific extensions and naming
conventions: - The certificate policies extension must contain the OID
for qualified certificates (0.4.0.194112.1.2 for QES) - The subject DN
must contain the signatory's name - The issuer DN identifies the QTSP -
CRL Distribution Points and/or OCSP responder URLs must be present - The
QCStatements extension (OID 0.4.0.1862.1) carries machine-readable
qualification information

\begin{Shaded}
\begin{Highlighting}[]
\CommentTok{{-}{-} QCStatements extension (ETSI EN 319 412{-}5)}
\NormalTok{QCStatements ::= }\DataTypeTok{SEQUENCE} \DataTypeTok{OF}\NormalTok{ QCStatement}

\NormalTok{QCStatement ::= }\DataTypeTok{SEQUENCE}\NormalTok{ \{}
\NormalTok{    statementId   OBJECT IDENTIFIER,}
\NormalTok{    statementInfo ANY DEFINED BY statementId }\KeywordTok{OPTIONAL}
\NormalTok{\}}

\CommentTok{{-}{-} Key OIDs:}
\CommentTok{{-}{-} id{-}etsi{-}qcs{-}QcCompliance  (0.4.0.1862.1.1)  = "this is a qualified cert"}
\CommentTok{{-}{-} id{-}etsi{-}qcs{-}QcSSCD        (0.4.0.1862.1.4)  = "key is on a QSCD"}
\CommentTok{{-}{-} id{-}etsi{-}qcs{-}QcType         (0.4.0.1862.1.6)  = type (esign, eseal, web)}
\end{Highlighting}
\end{Shaded}

\subsection{Article 29: Requirements for Qualified Signature Creation
Devices}\label{_article_29_requirements_for_qualified_signature_creation_devices}

Article 29 states that qualified electronic signature creation data (the
private key) must be managed by a QSCD that meets the requirements of
Annex II.

Annex II requires that QSCDs must:

\begin{enumerate}
\def\labelenumi{\arabic{enumi}.}
\item
  Ensure the confidentiality of the electronic signature creation data
\item
  Ensure that the electronic signature creation data used for signature
  creation can occur only once (practical uniqueness)
\item
  Ensure that the electronic signature creation data cannot be derived
  with reasonable assurance, and the signature is reliably protected
  against forgery using currently available technology
\item
  Ensure that the electronic signature creation data can be reliably
  protected by the legitimate signatory against use by others
\end{enumerate}

In practice, this means: the private key must be generated on, and never
leave, a certified hardware device. The certification is typically
Common Criteria EAL4+ (or equivalent) for the device, plus specific
Protection Profiles defined by CEN and ETSI.

Examples of certified QSCDs: - Smart cards (e.g., IDPrime MD from
Thales, CardOS from Atos) - USB tokens (e.g., SafeNet eToken 5110) -
Hardware Security Modules (e.g., Thales Luna Network HSM, Utimaco
SecurityServer) - Remote QSCDs operated by a QTSP (the key lives in an
HSM in a data center, accessed via API)

The rise of remote QSCDs has been the biggest practical change since
eIDAS entered into force. The original assumption was that QES required
a physical device in the signer's possession (a smart card in their
wallet). But eIDAS Annex II does not require \emph{physical possession}
--- it requires sole control. A private key in an HSM operated by a
QTSP, accessible only through strong multi-factor authentication
specific to the signer and the signing operation, satisfies the
requirement. This opened the door to remote signing (Chapter 22), which
is now the dominant model for QES in practice.

\subsection{Article 35-40: Electronic
Seals}\label{_article_35_40_electronic_seals}

Electronic seals are the organizational counterpart to electronic
signatures. While a signature is created by a natural person (a human),
a seal is created by a legal person (a company, government agency,
association).

The seal articles mirror the signature articles almost exactly: -
Article 35(1): non-discrimination for electronic seals - Article 35(2):
qualified electronic seal enjoys a presumption of integrity and origin -
Article 36: AdESeal requirements (same four properties as AdES, adapted
for legal persons) - Article 38: qualified certificates for electronic
seals (same as Art. 28, but for organizations)

The key difference: Article 35(2) says a qualified seal enjoys a
"presumption of integrity of the data and of correctness of the origin
of that data." Note the careful wording --- it says nothing about the
legal \emph{intent} behind the seal. A signature expresses consent or
approval. A seal proves origin and integrity. Chapter 21 explores this
distinction in depth.

\section{Mutual Recognition}\label{_mutual_recognition}

\subsection{Within the EU}\label{_within_the_eu}

Article 25(3) establishes that a QES based on a qualified certificate
from one member state is recognized as a QES in all member states.
Article 35(3) does the same for qualified electronic seals.

This means: - An Austrian QES is valid in Portugal - A Lithuanian
qualified seal is valid in Spain - An Estonian qualified timestamp is
valid in Germany

Member states \textbf{cannot} impose additional requirements for
cross-border recognition. They cannot require that QES certificates be
issued by domestic QTSPs. They cannot require additional technical
features beyond what eIDAS specifies.

The mechanism that makes this work is the \textbf{EU Trusted List
(EUTL)} --- a machine-readable, cryptographically signed registry of all
QTSPs in the EU. If a certificate chains to a QTSP on the EUTL, it
qualifies for cross-border recognition. Chapter 20 covers the EUTL in
detail.

In practice, cross-border recognition is still imperfect. While the
\emph{legal} framework is clear, many national IT systems have not been
updated to consume the EUTL and accept certificates from other member
states. A 2021 European Commission survey found that while most member
states had updated their legal frameworks, actual system-level
interoperability lagged significantly.

If your software validates signatures against a fixed list of trusted
certificates (rather than dynamically checking the EUTL), you are likely
rejecting valid QES from other member states. This is a compliance
failure under eIDAS --- Article 25(3) does not have exceptions for "we
haven't updated our trust store yet."

\subsection{Non-EU Recognition}\label{_non_eu_recognition}

eIDAS Article 14 allows the European Commission to adopt implementing
acts recognizing electronic identification schemes from third countries,
provided they meet equivalent requirements.

In practice, non-EU recognition is sparse: - \textbf{Switzerland} has a
bilateral agreement with the EU that provides for mutual recognition of
qualified trust services under the ZertES (Swiss Electronic Signatures
Act) - \textbf{United Kingdom} post-Brexit: the UK retained eIDAS as
domestic law (the UK eIDAS regulation) but mutual recognition lapsed; UK
QES are not automatically recognized in the EU and vice versa -
\textbf{United States}: no formal equivalence arrangement; US digital
signatures have no special status under eIDAS (and eIDAS signatures have
no special status under US law)

For developers serving global markets, this means: do not assume that a
qualified certificate from a non-EU QTSP will be recognized as
"qualified" in the EU, even if the QTSP provides technically identical
services. Qualification is a \emph{legal status} granted by an EU member
state's supervisory body, not a technical property of the certificate.

\section{eIDAS 2.0: Regulation (EU)
2024/1183}\label{_eidas_2_0_regulation_eu_20241183}

On April 11, 2024, the revised eIDAS regulation --- commonly called
\textbf{eIDAS 2.0} or Regulation (EU) 2024/1183 (eIDAS 2.0) --- entered
into force. The revision keeps the trust services framework largely
intact but adds a transformative new component: the \textbf{European
Digital Identity Wallet (EUDI Wallet)}.

\subsection{What Changed}\label{_what_changed}

The key changes in eIDAS 2.0:

\begin{enumerate}
\def\labelenumi{\arabic{enumi}.}
\item
  \textbf{EUDI Wallet}: Every member state must provide a digital
  identity wallet to its citizens and residents by 2026. The wallet will
  carry government-issued Person Identification Data (PID) and Qualified
  Electronic Attestations of Attributes (QEAAs). Crucially, the wallet
  will be able to create QES --- making qualified signatures accessible
  to every EU citizen from their smartphone.
\item
  \textbf{Qualified Electronic Attestations of Attributes (QEAAs)}: A
  new trust service type. QEAAs are digitally signed attestations of
  specific attributes (age, professional qualifications, university
  degrees) that can be presented through the EUDI wallet. Think of them
  as digital equivalents of physical certificates and credentials.
\item
  \textbf{Remote QSCD in the wallet}: The EUDI wallet can use a remote
  QSCD operated by a QTSP for creating qualified signatures. The private
  key lives in an HSM in the QTSP's data center, but the wallet provides
  the authentication and sole control mechanisms.
\item
  \textbf{Free QES for citizens}: Article 45e requires that natural
  persons be offered QES creation through the wallet free of charge for
  non-professional use. This is a dramatic change --- until now, QES
  required purchasing a certificate from a commercial QTSP, which
  typically costs 50-300 EUR per year.
\item
  \textbf{Qualified Electronic Ledgers}: A new trust service type for
  distributed ledger technology (blockchain). ETSI is currently
  developing the technical standards.
\end{enumerate}

\subsection{Impact on Developers}\label{_impact_on_developers}

For PDF signature developers, eIDAS 2.0 means:

\begin{itemize}
\item
  \textbf{Expect a surge in QES volume}: Free QES through the wallet
  will dramatically increase the number of qualified signatures in
  circulation. Your validation code needs to handle EUTL lookups
  efficiently.
\item
  \textbf{Wallet-initiated signing flows}: The EUDI wallet defines a new
  interaction pattern where the wallet app initiates a signing flow with
  a remote QSCD. If your application accepts signed documents, you need
  to support the signature formats the wallet produces.
\item
  \textbf{QEAA validation}: You may need to validate not just the
  signature on a document, but also QEAAs presented alongside it (e.g.,
  a doctor signing a prescription presents their medical qualification
  attestation).
\end{itemize}

The EUDI wallet's Architecture Reference Framework (ARF) is publicly
available and under active development. As of early 2026, large-scale
pilots are running in multiple sectors across the EU. The technical
specifications are not yet final, but the direction is clear: QES will
move from a niche enterprise technology to a consumer feature on
smartphones. Plan your architecture accordingly.

\section{The Gap Between Regulation and
Reality}\label{_the_gap_between_regulation_and_reality}

\subsection{What the Regulation Says}\label{_what_the_regulation_says}

eIDAS paints an elegant picture: three clearly defined tiers, qualified
trust services supervised by competent national bodies, mutual
recognition across borders, machine-readable trusted lists that software
can consume programmatically.

\subsection{What Actually Happens}\label{_what_actually_happens}

The reality is more complicated.

\textbf{National implementation varies despite the regulation's direct
applicability.} While member states cannot transpose the regulation
differently (it is directly applicable), they \emph{can} adopt
supplementary national measures. Germany's
\emph{Vertrauensdienstegesetz} (Trust Services Act), France's
\emph{decret relatif a la signature electronique}, and Italy's
\emph{Codice dell'Amministrazione Digitale} all contain national
provisions that fill gaps the regulation left open. These national
provisions sometimes create practical differences in how QES is
implemented and accepted.

\textbf{EUTL consumption is not universal.} Many government and
enterprise systems still validate signatures against hardcoded trust
stores rather than dynamically checking the EUTL. This means they accept
domestic QES but reject equally valid QES from other member states.

\textbf{QSCD certification is slow and expensive.} Getting a signature
creation device certified as a QSCD takes 12-24 months and costs
hundreds of thousands of euros. This limits the number of available
QSCDs and concentrates the market among a few hardware vendors.

\textbf{The "qualified" premium.} Qualified certificates cost
significantly more than standard certificates --- typically 50-300 EUR
per year for individuals, and more for organizations. This price barrier
has limited adoption, particularly among small businesses and
individuals. eIDAS 2.0's free QES through the wallet aims to address
this, but the wallet is not yet widely deployed.

\textbf{Interoperability testing is ongoing.} The European Commission
funds regular interoperability testing events (called "plugtests")
through ETSI, where QTSPs and software vendors test their
implementations against each other. These events consistently reveal
interoperability issues --- signatures that validate in one tool but
fail in another, timestamps that are accepted by some validators but
rejected by others, EUTL entries that are ambiguous about the services a
QTSP provides.

\section{Common Misconceptions}\label{_common_misconceptions}

\subsection{"eIDAS makes all digital signatures
legal"}\label{_eidas_makes_all_digital_signatures_legal}

No. eIDAS creates a \emph{framework} with three tiers. Only QES gets
automatic legal equivalence to handwritten signatures. Basic electronic
signatures and advanced electronic signatures benefit only from the
non-discrimination principle --- they cannot be rejected \emph{solely}
for being electronic, but their legal weight is a matter for the court
to decide.

\subsection{"If my signature uses a certificate, it's
advanced"}\label{_if_my_signature_uses_a_certificate_its_advanced}

Not necessarily. An advanced electronic signature must satisfy all four
requirements of Article 26: uniquely linked to the signatory, capable of
identifying them, under sole control, and linked to the data. A
certificate-based signature where the private key is stored in a shared
server-side keystore, accessible to multiple system administrators,
arguably fails the "sole control" requirement.

\subsection{"PAdES = eIDAS compliance"}\label{_pades_eidas_compliance}

Not automatically. PAdES defines technical profiles for PDF signatures.
Using PAdES-B-LTA with a qualified certificate on a QSCD produces a QES.
Using PAdES-B-LTA with a self-signed certificate produces a technically
well-structured signature with zero eIDAS qualification. The eIDAS tier
depends on \emph{who issued the certificate} and \emph{where the key is
stored}, not on the PDF signature format alone.

\subsection{"My country's signatures are recognized everywhere in the
EU"}\label{_my_countrys_signatures_are_recognized_everywhere_in_the_eu}

Only if they are \emph{qualified}. A non-qualified signature from
Germany has no special status in France --- it is just evidence, to be
evaluated by the French court on its merits. Only QES enjoys automatic
cross-border recognition under Article 25(3).

\subsection{"eIDAS applies outside the
EU"}\label{_eidas_applies_outside_the_eu}

eIDAS is EU law. It applies within the EU, the EEA (through the EEA
Agreement), and countries with specific bilateral arrangements
(Switzerland). It does not apply in the United States, China, Japan, or
anywhere else, unless those jurisdictions choose to recognize
eIDAS-qualified signatures under their own legal frameworks.

However, eIDAS has become a de facto global reference model. Countries
designing their own electronic signature frameworks often look to eIDAS
for inspiration. The three-tier model (basic, advanced, qualified) has
been adopted in various forms by multiple non-EU jurisdictions.

\section{What Your Code Needs to Do}\label{_what_your_code_needs_to_do}

For developers building PDF signature creation or validation software,
eIDAS imposes specific technical requirements depending on the tier you
are targeting:

\subsection{For Basic Electronic
Signatures}\label{_for_basic_electronic_signatures}

Nothing beyond what PDF already requires. Any valid PDF signature ---
even a simple \texttt{/SubFilter\ /adbe.pkcs7.detached} with a
self-signed certificate --- is an electronic signature under Article
3(10).

\subsection{For Advanced Electronic
Signatures}\label{_for_advanced_electronic_signatures}

Your signing workflow must ensure: 1. \textbf{Unique linkage}: one
private key per signer (never share keys) 2. \textbf{Identification}:
the certificate must contain the signer's identity (name, at minimum) 3.
\textbf{Sole control}: meaningful per-signature authentication (not just
session login) 4. \textbf{Data linkage}: the signature must cover the
entire document (standard byte range coverage)

Technically, this means: PAdES-B-B or higher, with a certificate issued
to a specific individual, and a signing flow that authenticates the
signer for each signature.

\subsection{For Qualified Electronic
Signatures}\label{_for_qualified_electronic_signatures}

All of the above, plus: 1. \textbf{Qualified certificate}: issued by a
QTSP listed on the EUTL 2. \textbf{QSCD}: the private key must be on a
certified qualified signature creation device 3. \textbf{Recommendation:
PAdES-B-LTA}: while not strictly required by eIDAS, PAdES-B-LTA ensures
long-term verifiability

\subsection{For Validation}\label{_for_validation}

If your software validates signatures and needs to determine the eIDAS
tier:

\begin{Shaded}
\begin{Highlighting}[]
\CommentTok{/**}
 \CommentTok{*}\NormalTok{ Determines the eIDAS tier of a PDF signature}\CommentTok{.}
 \CommentTok{*}\NormalTok{ This is a simplified illustration }\CommentTok{{-}}\NormalTok{ production code should use EU DSS}
 \CommentTok{*}\NormalTok{ which handles the full complexity of EUTL parsing and chain validation}\CommentTok{.}
 \CommentTok{*/}
\KeywordTok{public}\NormalTok{ EidasTier }\FunctionTok{determineEidasTier}\OperatorTok{(}\BuiltInTok{X509Certificate}\NormalTok{ signerCert}\OperatorTok{,}
                                     \DataTypeTok{boolean}\NormalTok{ keyOnQscd}\OperatorTok{)} \OperatorTok{\{}
    \CommentTok{// Check if the certificate is qualified (issued by a QTSP)}
    \DataTypeTok{boolean}\NormalTok{ isQualifiedCert }\OperatorTok{=} \FunctionTok{isListedOnEutl}\OperatorTok{(}\NormalTok{signerCert}\OperatorTok{);}

    \CommentTok{// Check for QCStatements extension}
    \DataTypeTok{boolean}\NormalTok{ hasQcCompliance }\OperatorTok{=} \FunctionTok{hasQcStatement}\OperatorTok{(}\NormalTok{signerCert}\OperatorTok{,}
            \StringTok{"0.4.0.1862.1.1"}\OperatorTok{);}  \CommentTok{// id{-}etsi{-}qcs{-}QcCompliance}
    \DataTypeTok{boolean}\NormalTok{ hasQscdStatement }\OperatorTok{=} \FunctionTok{hasQcStatement}\OperatorTok{(}\NormalTok{signerCert}\OperatorTok{,}
            \StringTok{"0.4.0.1862.1.4"}\OperatorTok{);}  \CommentTok{// id{-}etsi{-}qcs{-}QcSSCD}

    \ControlFlowTok{if} \OperatorTok{(}\NormalTok{isQualifiedCert }\OperatorTok{\&\&}\NormalTok{ hasQcCompliance }\OperatorTok{\&\&} \OperatorTok{(}\NormalTok{keyOnQscd }\OperatorTok{||}\NormalTok{ hasQscdStatement}\OperatorTok{))} \OperatorTok{\{}
        \ControlFlowTok{return}\NormalTok{ EidasTier}\OperatorTok{.}\FunctionTok{QUALIFIED}\OperatorTok{;}     \CommentTok{// QES {-} Art. 25(2) applies}
    \OperatorTok{\}}

    \ControlFlowTok{if} \OperatorTok{(}\FunctionTok{meetsArticle26Requirements}\OperatorTok{(}\NormalTok{signerCert}\OperatorTok{))} \OperatorTok{\{}
        \ControlFlowTok{return}\NormalTok{ EidasTier}\OperatorTok{.}\FunctionTok{ADVANCED}\OperatorTok{;}      \CommentTok{// AdES {-} Art. 25(1) applies}
    \OperatorTok{\}}

    \ControlFlowTok{return}\NormalTok{ EidasTier}\OperatorTok{.}\FunctionTok{BASIC}\OperatorTok{;}            \CommentTok{// Basic {-} Art. 25(1) applies}
\OperatorTok{\}}

\KeywordTok{private} \DataTypeTok{boolean} \FunctionTok{hasQcStatement}\OperatorTok{(}\BuiltInTok{X509Certificate}\NormalTok{ cert}\OperatorTok{,} \BuiltInTok{String}\NormalTok{ oid}\OperatorTok{)} \OperatorTok{\{}
    \DataTypeTok{byte}\OperatorTok{[]}\NormalTok{ extValue }\OperatorTok{=}\NormalTok{ cert}\OperatorTok{.}\FunctionTok{getExtensionValue}\OperatorTok{(}\StringTok{"1.3.6.1.5.5.7.1.3"}\OperatorTok{);}
    \ControlFlowTok{if} \OperatorTok{(}\NormalTok{extValue }\OperatorTok{==} \KeywordTok{null}\OperatorTok{)} \OperatorTok{\{}
        \ControlFlowTok{return} \KeywordTok{false}\OperatorTok{;}
    \OperatorTok{\}}
    \CommentTok{// Parse QCStatements ASN.1 structure}
\NormalTok{    ASN1Sequence qcStatements }\OperatorTok{=}\NormalTok{ ASN1Sequence}\OperatorTok{.}\FunctionTok{getInstance}\OperatorTok{(}
\NormalTok{            ASN1OctetString}\OperatorTok{.}\FunctionTok{getInstance}\OperatorTok{(}\NormalTok{extValue}\OperatorTok{).}\FunctionTok{getOctets}\OperatorTok{());}
    \ControlFlowTok{for} \OperatorTok{(}\NormalTok{ASN1Encodable statement }\OperatorTok{:}\NormalTok{ qcStatements}\OperatorTok{)} \OperatorTok{\{}
\NormalTok{        ASN1Sequence seq }\OperatorTok{=}\NormalTok{ ASN1Sequence}\OperatorTok{.}\FunctionTok{getInstance}\OperatorTok{(}\NormalTok{statement}\OperatorTok{);}
\NormalTok{        ASN1ObjectIdentifier statementId }\OperatorTok{=}
\NormalTok{                ASN1ObjectIdentifier}\OperatorTok{.}\FunctionTok{getInstance}\OperatorTok{(}\NormalTok{seq}\OperatorTok{.}\FunctionTok{getObjectAt}\OperatorTok{(}\DecValTok{0}\OperatorTok{));}
        \ControlFlowTok{if} \OperatorTok{(}\NormalTok{statementId}\OperatorTok{.}\FunctionTok{getId}\OperatorTok{().}\FunctionTok{equals}\OperatorTok{(}\NormalTok{oid}\OperatorTok{))} \OperatorTok{\{}
            \ControlFlowTok{return} \KeywordTok{true}\OperatorTok{;}
        \OperatorTok{\}}
    \OperatorTok{\}}
    \ControlFlowTok{return} \KeywordTok{false}\OperatorTok{;}
\OperatorTok{\}}
\end{Highlighting}
\end{Shaded}

The code above is a simplified illustration. Production eIDAS tier
determination is significantly more complex because:

\begin{itemize}
\item
  EUTL lookup requires parsing the trusted list XML, validating its
  signature, and traversing the LOT hierarchy
\item
  Certificate chain validation must check revocation status for every
  certificate in the chain
\item
  QSCD status may need to be inferred from the certificate's
  QCStatements rather than verified directly
\item
  Some QTSPs issue certificates that are qualified for some services but
  not others
\end{itemize}

For production use, the \textbf{EU DSS framework} (Chapter 23) provides
a complete, tested implementation of eIDAS tier determination. Do not
implement this from scratch unless you have very specific requirements
that EU DSS cannot meet.

\section{Key Takeaways}\label{_key_takeaways_7}

\begin{itemize}
\item
  eIDAS (Regulation (EU) No 910/2014 (eIDAS)) is an EU regulation ---
  directly applicable in all member states, no national transposition
  required
\item
  It replaced the failed Electronic Signatures Directive (1999/93/EC),
  which suffered from fragmented national implementations
\item
  Three tiers: basic electronic signature (any electronic indication of
  signing intent), advanced electronic signature (four properties:
  unique linkage, identification, sole control, data linkage), qualified
  electronic signature (advanced + QSCD + qualified certificate)
\item
  Only QES has automatic legal equivalence to a handwritten signature
  (Article 25(2))
\item
  Cross-border recognition is mandatory for QES (Article 25(3)) ---
  enforced through the EU Trusted List
\item
  eIDAS 2.0 (Regulation (EU) 2024/1183 (eIDAS 2.0)) adds the EUDI
  wallet, making QES accessible to all EU citizens via smartphone
\item
  The gap between regulation and implementation is real --- code that
  hardcodes trust stores instead of checking the EUTL violates the
  mutual recognition principle
\item
  For production eIDAS validation, use the EU DSS framework (Chapter 23)
  rather than implementing from scratch
\end{itemize}

\chapter{Qualified Certificates \& Trust
Services}\label{_qualified_certificates_trust_services}

A certificate is just a data structure --- some ASN.1 fields, a public
key, and a signature from someone who claims to be a Certificate
Authority. What makes a certificate \emph{qualified} is not a technical
property you can detect by parsing bytes alone. It is a legal status,
conferred by a government supervisory body, published in a
machine-readable trust infrastructure, and backed by an institutional
apparatus that most developers have never heard of, let alone queried
programmatically.

This chapter takes you through that entire apparatus: the Qualified
Trust Service Providers that issue qualified certificates, the EU
Trusted List that publishes who they are, the technical requirements
that separate a qualified certificate from an ordinary one, the QSCD
hardware that protects the keys, and the Java code you need to actually
parse and verify all of this. If Chapter 19 was the legal framework,
this chapter is the plumbing.

\section{Qualified Trust Service Providers
(QTSPs)}\label{_qualified_trust_service_providers_qtsps}

\subsection{What a QTSP Is}\label{_what_a_qtsp_is}

A \textbf{Qualified Trust Service Provider} is an organization that has
been audited, approved, and listed by a national supervisory body within
the EU as meeting the requirements of Regulation (EU) No 910/2014
(eIDAS) Articles 20-24. The word "qualified" is doing serious work
here --- it is not a marketing term, not a self-declaration, and not
something you obtain by installing software. It is a legal status,
analogous to a bank receiving a banking license: you cannot call
yourself a QTSP unless a government body says you are one.

A QTSP can provide one or more qualified trust services:

\begin{itemize}
\item
  \textbf{Qualified certificates for electronic signatures} (natural
  persons)
\item
  \textbf{Qualified certificates for electronic seals} (legal persons)
\item
  \textbf{Qualified timestamps} (RFC 3161 tokens from a qualified TSA)
\item
  \textbf{Qualified electronic delivery services} (certified email
  equivalents)
\item
  \textbf{Qualified website authentication certificates} (the
  controversial Article 45, rewritten in eIDAS 2.0)
\item
  \textbf{Qualified electronic attestations of attributes} (new in
  Regulation (EU) 2024/1183 (eIDAS 2.0))
\end{itemize}

The qualification is \emph{per-service}, not per-organization. A
provider might be qualified for timestamps but not for certificates, or
for seals but not for signatures. This matters when you are parsing the
EU Trusted List --- you need to check not just whether the provider is
listed, but \emph{which services} are qualified.

\subsection{How to Become a QTSP}\label{_how_to_become_a_qtsp}

The path to becoming a QTSP is neither quick nor cheap. Here is the
process, simplified but accurate:

\begin{enumerate}
\def\labelenumi{\arabic{enumi}.}
\item
  \textbf{Establish the infrastructure.} You need HSMs (certified to
  Common Criteria EAL4+ or equivalent), a secure data center (typically
  ISO 27001 certified), redundant systems, trained personnel, and
  documented processes for every aspect of certificate lifecycle
  management --- issuance, suspension, revocation, key generation, key
  destruction.
\item
  \textbf{Write a Certificate Practice Statement (CPS).} This is the
  document that describes, in exhaustive detail, how you operate:
  identity verification procedures, key management, certificate
  profiles, revocation handling, disaster recovery, incident response. A
  CPS for a QTSP is typically 80-150 pages long.
\item
  \textbf{Get a conformity assessment.} An accredited \textbf{Conformity
  Assessment Body (CAB)} evaluates your infrastructure, processes, and
  CPS against Regulation (EU) No 910/2014 (eIDAS) and the applicable
  ETSI standards (primarily ETSI EN 319 401 for general TSP requirements
  and ETSI EN 319 411-1/411-2 for certificate issuance). The assessment
  produces a \textbf{Conformity Assessment Report (CAR)}.
\item
  \textbf{Submit to the national supervisory body.} Each EU member state
  designates a supervisory body for QTSPs: in Germany, the
  \textbf{Bundesnetzagentur} (Federal Network Agency); in France,
  \textbf{ANSSI}; in Italy, \textbf{AgID} (Agenzia per l'Italia
  Digitale). You submit the CAR, your CPS, and supporting documentation.
\item
  \textbf{Supervisory body review and approval.} The supervisory body
  reviews everything, may request additional information or on-site
  inspections, and eventually grants or denies qualified status.
\item
  \textbf{Listing on the national Trusted List.} Once approved, the
  supervisory body adds your service(s) to the national \textbf{Trusted
  Service Status List (TSL)} --- the country's portion of the EU Trusted
  List. Your qualified status takes effect when the updated TSL is
  published.
\item
  \textbf{Ongoing supervision.} Qualification is not a one-time event.
  Regulation (EU) No 910/2014 (eIDAS) Article 20(1) mandates audits at
  least every 24 months. Regulation (EU) No 910/2014 (eIDAS) Article
  19(2) requires QTSPs to report security incidents to the supervisory
  body within 24 hours.
\end{enumerate}

The entire process from initial infrastructure setup to EUTL listing
typically takes \textbf{18-36 months} and costs \textbf{several hundred
thousand to several million euros}, depending on the scope of services
and the jurisdiction. This is by design: the barrier to entry ensures
that only well-resourced, professionally operated organizations can
issue qualified certificates.

If someone offers you "qualified certificates" from a provider that is
not on the EUTL, they are either confused about what "qualified" means
or deliberately misleading you. Always verify against the EUTL.

\subsection{The Supervision Model}\label{_the_supervision_model}

The eIDAS supervision model is a two-layer structure:

\textbf{National layer:} Each member state's supervisory body is
responsible for QTSPs established in its territory. The supervisory body
grants and revokes qualified status, maintains the national TSL,
conducts or commissions audits, handles incident reports, and imposes
sanctions for non-compliance.

\textbf{EU layer:} The European Commission maintains the \textbf{List of
Trusted Lists (LOTL)} --- a master document that points to every
national TSL. The Commission does not supervise individual QTSPs; it
supervises the supervisory bodies indirectly, through the mutual
evaluation mechanism and through the eIDAS Expert Group.

\begin{longtable}[]{@{}
  >{\raggedright\arraybackslash}p{(\linewidth - 4\tabcolsep) * \real{0.1429}}
  >{\raggedright\arraybackslash}p{(\linewidth - 4\tabcolsep) * \real{0.4286}}
  >{\raggedright\arraybackslash}p{(\linewidth - 4\tabcolsep) * \real{0.4286}}@{}}
\caption{Supervision hierarchy}\tabularnewline
\toprule\noalign{}
\begin{minipage}[b]{\linewidth}\raggedright
Level
\end{minipage} & \begin{minipage}[b]{\linewidth}\raggedright
Entity
\end{minipage} & \begin{minipage}[b]{\linewidth}\raggedright
Responsibility
\end{minipage} \\
\midrule\noalign{}
\endfirsthead
\toprule\noalign{}
\begin{minipage}[b]{\linewidth}\raggedright
Level
\end{minipage} & \begin{minipage}[b]{\linewidth}\raggedright
Entity
\end{minipage} & \begin{minipage}[b]{\linewidth}\raggedright
Responsibility
\end{minipage} \\
\midrule\noalign{}
\endhead
\bottomrule\noalign{}
\endlastfoot
EU & European Commission & Maintains LOTL, sets technical standards via
implementing acts, coordinates mutual evaluations \\
National & Supervisory Body (Bundesnetzagentur, ANSSI, AgID, etc.) &
Grants/revokes qualified status, maintains national TSL, audits QTSPs,
handles incidents \\
Provider & QTSP (DigiCert, D-Trust, InfoCert, etc.) & Operates trust
services, issues certificates, complies with ETSI standards, reports
incidents \\
\end{longtable}

\section{The EU Trusted List (EUTL)}\label{_the_eu_trusted_list_eutl}

\subsection{Architecture: The LOT
Hierarchy}\label{_architecture_the_lot_hierarchy}

The EUTL is not a single list --- it is a \textbf{hierarchy of
cryptographically signed XML documents}:

\begin{enumerate}
\def\labelenumi{\arabic{enumi}.}
\item
  The \textbf{LOTL (List of Trusted Lists)} is the root. Maintained by
  the European Commission, published at a well-known URL. The LOTL
  contains pointers (URLs) to every national TSL, along with the
  certificates used to verify each national TSL's signature.
\item
  Each \textbf{national TSL (Trusted Service Status List)} lists the
  trust service providers in that country, their services, the status of
  each service (granted, withdrawn, deprecated), and the CA certificates
  associated with each service.
\item
  Each TSL is \textbf{digitally signed} using an XML Advanced Electronic
  Signature (XAdES) --- the same trust-no-one verification model applied
  to the trust infrastructure itself.
\end{enumerate}

The LOTL URL is:

\begin{verbatim}
https://ec.europa.eu/tools/lotl/eu-lotl.xml
\end{verbatim}

This URL has been stable for years, but production systems should make
it configurable. The official reference for the LOTL location is
Commission Implementing Decision (EU) 2015/1505. The LOTL itself
contains a \texttt{NextUpdate} field indicating when the next version
will be published. National TSLs are updated on varying
schedules --- some monthly, some quarterly, some whenever a change
occurs.

\subsection{XML Schema}\label{_xml_schema}

The TSL XML schema is defined in ETSI TS 119 612 (replacing the older
ETSI TS 102 231). The schema is verbose, deeply nested, and frankly
exhausting to read --- but it is well-structured and consistent.

\begin{Shaded}
\begin{Highlighting}[]
\CommentTok{\textless{}!{-}{-} Simplified LOTL structure {-}{-}\textgreater{}}
\NormalTok{\textless{}}\KeywordTok{TrustServiceStatusList} \OtherTok{xmlns=}\StringTok{"http://uri.etsi.org/02231/v2\#"}\NormalTok{\textgreater{}}
\NormalTok{  \textless{}}\KeywordTok{SchemeInformation}\NormalTok{\textgreater{}}
\NormalTok{    \textless{}}\KeywordTok{TSLVersionIdentifier}\NormalTok{\textgreater{}5\textless{}/}\KeywordTok{TSLVersionIdentifier}\NormalTok{\textgreater{}}
\NormalTok{    \textless{}}\KeywordTok{TSLSequenceNumber}\NormalTok{\textgreater{}321\textless{}/}\KeywordTok{TSLSequenceNumber}\NormalTok{\textgreater{}}
\NormalTok{    \textless{}}\KeywordTok{TSLType}\NormalTok{\textgreater{}}
\NormalTok{      http://uri.etsi.org/TrstSvc/TrustedList/TSLType/EUlistofthelists}
\NormalTok{    \textless{}/}\KeywordTok{TSLType}\NormalTok{\textgreater{}}
\NormalTok{    \textless{}}\KeywordTok{SchemeOperatorName}\NormalTok{\textgreater{}}
\NormalTok{      \textless{}}\KeywordTok{Name} \OtherTok{xml:lang=}\StringTok{"en"}\NormalTok{\textgreater{}European Commission\textless{}/}\KeywordTok{Name}\NormalTok{\textgreater{}}
\NormalTok{    \textless{}/}\KeywordTok{SchemeOperatorName}\NormalTok{\textgreater{}}
\NormalTok{    \textless{}}\KeywordTok{ListIssueDateTime}\NormalTok{\textgreater{}2026{-}02{-}15T00:00:00Z\textless{}/}\KeywordTok{ListIssueDateTime}\NormalTok{\textgreater{}}
\NormalTok{    \textless{}}\KeywordTok{NextUpdate}\NormalTok{\textgreater{}}
\NormalTok{      \textless{}}\KeywordTok{dateTime}\NormalTok{\textgreater{}2026{-}08{-}15T00:00:00Z\textless{}/}\KeywordTok{dateTime}\NormalTok{\textgreater{}}
\NormalTok{    \textless{}/}\KeywordTok{NextUpdate}\NormalTok{\textgreater{}}
\NormalTok{  \textless{}/}\KeywordTok{SchemeInformation}\NormalTok{\textgreater{}}

\NormalTok{  \textless{}}\KeywordTok{TrustServiceProviderList}\NormalTok{\textgreater{}}
\NormalTok{    \textless{}}\KeywordTok{TrustServiceProvider}\NormalTok{\textgreater{}}
\NormalTok{      \textless{}}\KeywordTok{TSPInformation}\NormalTok{\textgreater{}}
\NormalTok{        \textless{}}\KeywordTok{TSPName}\NormalTok{\textgreater{}\textless{}}\KeywordTok{Name} \OtherTok{xml:lang=}\StringTok{"en"}\NormalTok{\textgreater{}Germany\textless{}/}\KeywordTok{Name}\NormalTok{\textgreater{}\textless{}/}\KeywordTok{TSPName}\NormalTok{\textgreater{}}
\NormalTok{      \textless{}/}\KeywordTok{TSPInformation}\NormalTok{\textgreater{}}
\NormalTok{      \textless{}}\KeywordTok{TSPServices}\NormalTok{\textgreater{}}
\NormalTok{        \textless{}}\KeywordTok{TSPService}\NormalTok{\textgreater{}}
\NormalTok{          \textless{}}\KeywordTok{ServiceInformation}\NormalTok{\textgreater{}}
\NormalTok{            \textless{}}\KeywordTok{ServiceTypeIdentifier}\NormalTok{\textgreater{}}
\NormalTok{              http://uri.etsi.org/TrstSvc/Svctype/TLIssuer}
\NormalTok{            \textless{}/}\KeywordTok{ServiceTypeIdentifier}\NormalTok{\textgreater{}}
\NormalTok{            \textless{}}\KeywordTok{ServiceSupplyPoints}\NormalTok{\textgreater{}}
\NormalTok{              \textless{}}\KeywordTok{ServiceSupplyPoint}\NormalTok{\textgreater{}}
\NormalTok{                https://www.bundesnetzagentur.de/.../TrustedList.xml}
\NormalTok{              \textless{}/}\KeywordTok{ServiceSupplyPoint}\NormalTok{\textgreater{}}
\NormalTok{            \textless{}/}\KeywordTok{ServiceSupplyPoints}\NormalTok{\textgreater{}}
\NormalTok{          \textless{}/}\KeywordTok{ServiceInformation}\NormalTok{\textgreater{}}
\NormalTok{        \textless{}/}\KeywordTok{TSPService}\NormalTok{\textgreater{}}
\NormalTok{      \textless{}/}\KeywordTok{TSPServices}\NormalTok{\textgreater{}}
\NormalTok{    \textless{}/}\KeywordTok{TrustServiceProvider}\NormalTok{\textgreater{}}
    \CommentTok{\textless{}!{-}{-} ... 26 more countries ... {-}{-}\textgreater{}}
\NormalTok{  \textless{}/}\KeywordTok{TrustServiceProviderList}\NormalTok{\textgreater{}}
\NormalTok{\textless{}/}\KeywordTok{TrustServiceStatusList}\NormalTok{\textgreater{}}
\end{Highlighting}
\end{Shaded}

Each \textbf{national TSL} then contains the actual QTSP entries with
their service certificates:

\begin{Shaded}
\begin{Highlighting}[]
\CommentTok{\textless{}!{-}{-} Simplified national TSL structure (Germany) {-}{-}\textgreater{}}
\NormalTok{\textless{}}\KeywordTok{TrustServiceStatusList} \OtherTok{xmlns=}\StringTok{"http://uri.etsi.org/02231/v2\#"}\NormalTok{\textgreater{}}
\NormalTok{  \textless{}}\KeywordTok{TrustServiceProviderList}\NormalTok{\textgreater{}}
\NormalTok{    \textless{}}\KeywordTok{TrustServiceProvider}\NormalTok{\textgreater{}}
\NormalTok{      \textless{}}\KeywordTok{TSPInformation}\NormalTok{\textgreater{}}
\NormalTok{        \textless{}}\KeywordTok{TSPName}\NormalTok{\textgreater{}}
\NormalTok{          \textless{}}\KeywordTok{Name} \OtherTok{xml:lang=}\StringTok{"en"}\NormalTok{\textgreater{}D{-}Trust GmbH\textless{}/}\KeywordTok{Name}\NormalTok{\textgreater{}}
\NormalTok{        \textless{}/}\KeywordTok{TSPName}\NormalTok{\textgreater{}}
\NormalTok{      \textless{}/}\KeywordTok{TSPInformation}\NormalTok{\textgreater{}}
\NormalTok{      \textless{}}\KeywordTok{TSPServices}\NormalTok{\textgreater{}}
\NormalTok{        \textless{}}\KeywordTok{TSPService}\NormalTok{\textgreater{}}
\NormalTok{          \textless{}}\KeywordTok{ServiceInformation}\NormalTok{\textgreater{}}
\NormalTok{            \textless{}}\KeywordTok{ServiceTypeIdentifier}\NormalTok{\textgreater{}}
\NormalTok{              http://uri.etsi.org/TrstSvc/Svctype/CA/QC}
\NormalTok{            \textless{}/}\KeywordTok{ServiceTypeIdentifier}\NormalTok{\textgreater{}}
\NormalTok{            \textless{}}\KeywordTok{ServiceStatus}\NormalTok{\textgreater{}}
\NormalTok{              http://uri.etsi.org/TrstSvc/TrustedList/Svcstatus/granted}
\NormalTok{            \textless{}/}\KeywordTok{ServiceStatus}\NormalTok{\textgreater{}}
\NormalTok{            \textless{}}\KeywordTok{StatusStartingTime}\NormalTok{\textgreater{}2018{-}07{-}01T00:00:00Z\textless{}/}\KeywordTok{StatusStartingTime}\NormalTok{\textgreater{}}
\NormalTok{            \textless{}}\KeywordTok{ServiceDigitalIdentity}\NormalTok{\textgreater{}}
\NormalTok{              \textless{}}\KeywordTok{DigitalId}\NormalTok{\textgreater{}}
\NormalTok{                \textless{}}\KeywordTok{X509Certificate}\NormalTok{\textgreater{}MIIE...base64...\textless{}/}\KeywordTok{X509Certificate}\NormalTok{\textgreater{}}
\NormalTok{              \textless{}/}\KeywordTok{DigitalId}\NormalTok{\textgreater{}}
\NormalTok{            \textless{}/}\KeywordTok{ServiceDigitalIdentity}\NormalTok{\textgreater{}}
\NormalTok{            \textless{}}\KeywordTok{ServiceInformationExtensions}\NormalTok{\textgreater{}}
\NormalTok{              \textless{}}\KeywordTok{Extension} \OtherTok{Critical=}\StringTok{"true"}\NormalTok{\textgreater{}}
\NormalTok{                \textless{}}\KeywordTok{Qualifications}\NormalTok{\textgreater{}}
\NormalTok{                  \textless{}}\KeywordTok{QualificationElement}\NormalTok{\textgreater{}}
\NormalTok{                    \textless{}}\KeywordTok{Qualifiers}\NormalTok{\textgreater{}}
\NormalTok{                      \textless{}}\KeywordTok{Qualifier} \OtherTok{uri=}\StringTok{"http://uri.etsi.org/TrstSvc/}
\StringTok{                        TrustedList/SvcInfoExt/QCWithQSCD"}\NormalTok{/\textgreater{}}
\NormalTok{                    \textless{}/}\KeywordTok{Qualifiers}\NormalTok{\textgreater{}}
\NormalTok{                  \textless{}/}\KeywordTok{QualificationElement}\NormalTok{\textgreater{}}
\NormalTok{                \textless{}/}\KeywordTok{Qualifications}\NormalTok{\textgreater{}}
\NormalTok{              \textless{}/}\KeywordTok{Extension}\NormalTok{\textgreater{}}
\NormalTok{            \textless{}/}\KeywordTok{ServiceInformationExtensions}\NormalTok{\textgreater{}}
\NormalTok{          \textless{}/}\KeywordTok{ServiceInformation}\NormalTok{\textgreater{}}
\NormalTok{        \textless{}/}\KeywordTok{TSPService}\NormalTok{\textgreater{}}
\NormalTok{      \textless{}/}\KeywordTok{TSPServices}\NormalTok{\textgreater{}}
\NormalTok{    \textless{}/}\KeywordTok{TrustServiceProvider}\NormalTok{\textgreater{}}
\NormalTok{  \textless{}/}\KeywordTok{TrustServiceProviderList}\NormalTok{\textgreater{}}
\NormalTok{\textless{}/}\KeywordTok{TrustServiceStatusList}\NormalTok{\textgreater{}}
\end{Highlighting}
\end{Shaded}

\subsection{Key Service Types and
Statuses}\label{_key_service_types_and_statuses}

\begin{longtable}[]{@{}
  >{\raggedright\arraybackslash}p{(\linewidth - 2\tabcolsep) * \real{0.4286}}
  >{\raggedright\arraybackslash}p{(\linewidth - 2\tabcolsep) * \real{0.5714}}@{}}
\caption{EUTL service type identifiers for PDF-relevant
services}\tabularnewline
\toprule\noalign{}
\begin{minipage}[b]{\linewidth}\raggedright
Service Type URI
\end{minipage} & \begin{minipage}[b]{\linewidth}\raggedright
Meaning
\end{minipage} \\
\midrule\noalign{}
\endfirsthead
\toprule\noalign{}
\begin{minipage}[b]{\linewidth}\raggedright
Service Type URI
\end{minipage} & \begin{minipage}[b]{\linewidth}\raggedright
Meaning
\end{minipage} \\
\midrule\noalign{}
\endhead
\bottomrule\noalign{}
\endlastfoot
\texttt{…\hspace{0pt}/Svctype/CA/QC} & CA issuing \textbf{qualified
certificates} (check qualifiers for esign/eseal/web) \\
\texttt{…\hspace{0pt}/Svctype/TSA/QTST} & \textbf{Qualified timestamp}
authority \\
\texttt{…\hspace{0pt}/Svctype/Certstatus/OCSP/QC} & \textbf{Qualified
OCSP} responder \\
\texttt{…\hspace{0pt}/Svctype/Certstatus/CRL/QC} & \textbf{Qualified
CRL} issuer \\
\end{longtable}

\begin{longtable}[]{@{}
  >{\raggedright\arraybackslash}p{(\linewidth - 2\tabcolsep) * \real{0.4286}}
  >{\raggedright\arraybackslash}p{(\linewidth - 2\tabcolsep) * \real{0.5714}}@{}}
\caption{EUTL service status values}\tabularnewline
\toprule\noalign{}
\begin{minipage}[b]{\linewidth}\raggedright
Status URI
\end{minipage} & \begin{minipage}[b]{\linewidth}\raggedright
Meaning
\end{minipage} \\
\midrule\noalign{}
\endfirsthead
\toprule\noalign{}
\begin{minipage}[b]{\linewidth}\raggedright
Status URI
\end{minipage} & \begin{minipage}[b]{\linewidth}\raggedright
Meaning
\end{minipage} \\
\midrule\noalign{}
\endhead
\bottomrule\noalign{}
\endlastfoot
\texttt{…\hspace{0pt}/Svcstatus/granted} & Currently qualified --- this
is what you are looking for \\
\texttt{…\hspace{0pt}/Svcstatus/withdrawn} & Qualification has been
withdrawn \\
\texttt{…\hspace{0pt}/Svcstatus/recognisedatnationallevel} & Not
EU-qualified, but recognized under national law \\
\texttt{…\hspace{0pt}/Svcstatus/deprecated} & Service deprecated by the
supervisory body \\
\end{longtable}

A common mistake in EUTL parsing is to check only the current status.
You also need to check the \textbf{service history}. A service that was
\texttt{granted} from 2019-01-01 to 2023-06-30 and then
\texttt{withdrawn} was qualified during that window. A certificate
issued during the \texttt{granted} period is still a qualified
certificate, even though the service is now withdrawn.

The TSL includes a \texttt{ServiceHistory} element with
\texttt{ServiceHistoryInstance} entries recording every status change.
Your validation logic must check whether the certificate was issued (and
the signature created) during a period when the service status was
\texttt{granted}.

\subsection{Cryptographic Signing of the
TSL}\label{_cryptographic_signing_of_the_tsl}

Every TSL (including the LOTL) is signed with an XAdES signature. To
trust the contents, you must:

\begin{enumerate}
\def\labelenumi{\arabic{enumi}.}
\item
  \textbf{Verify the XML signature} on the TSL document
\item
  \textbf{Verify the signing certificate} against the LOTL (for national
  TSLs) or against the Commission's well-known certificate (for the LOTL
  itself)
\item
  \textbf{Check freshness} --- the TSL must not be past its
  \texttt{NextUpdate} date (or not past it by more than a configurable
  grace period)
\end{enumerate}

The LOTL signing certificate is published by the European Commission at
a well-known URL. When the signing certificate rotates, the new
certificate is announced in advance in the LOTL itself --- this is the
\textbf{pivot mechanism}, and it is one of the trickiest parts of EUTL
parsing to implement correctly.

For production systems, implement a scheduled refresh --- daily or
weekly --- that re-downloads, re-verifies, and re-parses the entire EUTL
hierarchy. The EU DSS library (Chapter 23) provides a
\texttt{TrustedListsCertificateSource} that handles all of this
automatically, including caching, refresh scheduling, signature
verification, and the pivot mechanism.

\section{Qualified Certificates: What Makes Them
Qualified}\label{_qualified_certificates_what_makes_them_qualified}

A certificate is "qualified" when three conditions are simultaneously
true:

\begin{enumerate}
\def\labelenumi{\arabic{enumi}.}
\item
  It is issued by a CA whose certificate issuance service has
  \texttt{granted} status in the EUTL
\item
  The certificate contains specific technical indicators (the
  QCStatements extension)
\item
  The CA followed the procedures defined in its CPS, which was approved
  during the conformity assessment
\end{enumerate}

Condition 3 is not something you can verify programmatically --- it is
an institutional guarantee, backed by the supervisory body's ongoing
oversight. Conditions 1 and 2 are what your code checks.

\subsection{The QCStatements
Extension}\label{_the_qcstatements_extension}

The technical indicator of a qualified certificate is the
\textbf{QCStatements extension} (OID \texttt{1.3.6.1.5.5.7.1.3}),
defined in ETSI EN 319 412-5. This extension contains a sequence of
statements, each identified by an OID:

\begin{longtable}[]{@{}
  >{\raggedright\arraybackslash}p{(\linewidth - 4\tabcolsep) * \real{0.3333}}
  >{\raggedright\arraybackslash}p{(\linewidth - 4\tabcolsep) * \real{0.2222}}
  >{\raggedright\arraybackslash}p{(\linewidth - 4\tabcolsep) * \real{0.4444}}@{}}
\caption{QCStatements OIDs relevant to PDF signatures}\tabularnewline
\toprule\noalign{}
\begin{minipage}[b]{\linewidth}\raggedright
OID
\end{minipage} & \begin{minipage}[b]{\linewidth}\raggedright
Short Name
\end{minipage} & \begin{minipage}[b]{\linewidth}\raggedright
Meaning
\end{minipage} \\
\midrule\noalign{}
\endfirsthead
\toprule\noalign{}
\begin{minipage}[b]{\linewidth}\raggedright
OID
\end{minipage} & \begin{minipage}[b]{\linewidth}\raggedright
Short Name
\end{minipage} & \begin{minipage}[b]{\linewidth}\raggedright
Meaning
\end{minipage} \\
\midrule\noalign{}
\endhead
\bottomrule\noalign{}
\endlastfoot
\texttt{0.4.0.1862.1.1} & id-etsi-qcs-QcCompliance & This certificate is
issued as a \textbf{qualified certificate} under Regulation (EU) No
910/2014 (eIDAS) \\
\texttt{0.4.0.1862.1.4} & id-etsi-qcs-QcSSCD & The private key resides
on a \textbf{Qualified Signature/Seal Creation Device} \\
\texttt{0.4.0.1862.1.6} & id-etsi-qcs-QcType & The type: esign
(\texttt{…\hspace{0pt}6.1} = signature), eseal
(\texttt{…\hspace{0pt}6.2} = seal), web (\texttt{…\hspace{0pt}6.3}) \\
\texttt{0.4.0.1862.1.3} & id-etsi-qcs-QcRetentionPeriod & How long the
CA retains registration information after certificate expiry \\
\texttt{0.4.0.1862.1.5} & id-etsi-qcs-QcPDS & URL of the PKI Disclosure
Statement (human-readable CPS summary) \\
\texttt{0.4.0.1862.1.2} & id-etsi-qcs-QcLimitValue & Monetary
transaction limit (rarely used) \\
\end{longtable}

Do \textbf{not} trust the QCStatements extension alone. Anyone can put
\texttt{id-etsi-qcs-QcCompliance} in a self-signed certificate. The
QCStatements assertion is only meaningful when the issuing CA is listed
in the EUTL with \texttt{granted} status for a \texttt{CA/QC} service.

The correct validation logic is:

\begin{enumerate}
\def\labelenumi{\arabic{enumi}.}
\item
  Check that the issuing CA's certificate matches one in the EUTL with
  \texttt{ServiceTypeIdentifier} of \texttt{CA/QC} and status
  \texttt{granted}
\item
  Check the period: the certificate's \texttt{notBefore} date must fall
  within the service's \texttt{granted} status period
\item
  \emph{Then} check QCStatements for the specific type (esign, eseal,
  web)
\item
  \emph{Then} check for QcSSCD if you need to verify QSCD usage
\end{enumerate}

Steps 1-2 without steps 3-4: you know the CA is qualified, but not
necessarily this specific certificate. Steps 3-4 without steps 1-2: you
have a self-asserted claim with no institutional backing. You need all
four.

\subsection{Certificate Policy OIDs}\label{_certificate_policy_oids}

Qualified certificates also carry specific \textbf{Certificate Policy}
OIDs in the \texttt{certificatePolicies} extension (OID
\texttt{2.5.29.32}), defined in ETSI EN 319 411-2:

\begin{longtable}[]{@{}
  >{\raggedright\arraybackslash}p{(\linewidth - 2\tabcolsep) * \real{0.4286}}
  >{\raggedright\arraybackslash}p{(\linewidth - 2\tabcolsep) * \real{0.5714}}@{}}
\caption{Certificate Policy OIDs for qualified
certificates}\tabularnewline
\toprule\noalign{}
\begin{minipage}[b]{\linewidth}\raggedright
OID
\end{minipage} & \begin{minipage}[b]{\linewidth}\raggedright
Meaning
\end{minipage} \\
\midrule\noalign{}
\endfirsthead
\toprule\noalign{}
\begin{minipage}[b]{\linewidth}\raggedright
OID
\end{minipage} & \begin{minipage}[b]{\linewidth}\raggedright
Meaning
\end{minipage} \\
\midrule\noalign{}
\endhead
\bottomrule\noalign{}
\endlastfoot
\texttt{0.4.0.194112.1.2} & QCP-n: Qualified Certificate Policy for
\textbf{natural persons} (signatures) \\
\texttt{0.4.0.194112.1.3} & QCP-l: Qualified Certificate Policy for
\textbf{legal persons} (seals) \\
\texttt{0.4.0.194112.1.1} & QCP-n-qscd: Natural person certificates with
QSCD \\
\texttt{0.4.0.194112.1.4} & QCP-l-qscd: Legal person certificates with
QSCD \\
\texttt{0.4.0.194112.1.0} & QCP-w: Website authentication \\
\end{longtable}

The \texttt{-qscd} variants indicate that the CA has verified that the
subscriber's private key is on a QSCD. This is a stronger assertion than
the QcSSCD statement alone --- the policy OID says "the CA has verified
the key is on a QSCD according to the procedures defined in ETSI EN 319
411-2."

\subsection{QSCD Requirements}\label{_qscd_requirements}

A Qualified Signature Creation Device (or Qualified Seal Creation
Device --- QSCD covers both) must meet the requirements of Regulation
(EU) No 910/2014 (eIDAS) Annex II:

\begin{enumerate}
\def\labelenumi{\arabic{enumi}.}
\item
  The private key must be generated on the device and never leave it
\item
  The signing operation must occur on the device itself
\item
  Access to the private key must require authentication of the
  legitimate holder (PIN, biometric, etc.)
\item
  The private key cannot be derived from signatures or the public key
  with reasonable assurance
\end{enumerate}

\textbf{Certification path:} Common Criteria evaluation (ISO/IEC 15408)
at EAL4+ against Protection Profiles defined in EN 419 211 (for smart
cards and tokens) or EN 419 221 (for HSMs and remote signing). A
national certification body (BSI in Germany, ANSSI in France) performs
the evaluation.

\textbf{Hardware categories in practice:}

\begin{itemize}
\item
  \textbf{Smart cards} (Thales IDPrime MD, Atos CardOS): private key
  generated on chip, never exported, signing performed on-card
\item
  \textbf{USB tokens} (SafeNet eToken 5110, YubiKey 5 FIPS): same
  security model as smart cards, USB form factor
\item
  \textbf{HSMs} (Thales Luna, Utimaco SecurityServer): rack-mounted
  devices for server-side signing; basis for remote signing (Chapter 22)
\end{itemize}

Not every HSM is a certified QSCD. An HSM may be FIPS 140-2 Level 3 or
Common Criteria EAL4+ certified for \emph{general cryptographic
operations} without being certified as a QSCD under the specific EN 419
211/221 Protection Profiles. The QSCD certification must reference the
eIDAS-specific Protection Profiles.

\section{Real-World QTSPs}\label{_real_world_qtsps}

When you are debugging why a particular qualified certificate does or
does not validate, knowing who issued it and where they are listed is
invaluable. Here are the QTSPs you are most likely to encounter:

\textbf{D-Trust GmbH} (Germany) --- a subsidiary of the Bundesdruckerei
(the organization that prints German passports). Listed in the German
TSL, supervised by the Bundesnetzagentur. Dominant in German government
and enterprise contexts.

\textbf{DigiCert / QuoVadis} (Netherlands, Belgium) --- DigiCert
acquired QuoVadis in 2019, inheriting QTSP status in multiple EU
countries. One of the few providers with both AATL and eIDAS qualified
status --- useful when documents must validate in both Adobe Acrobat and
European government systems.

\textbf{GlobalSign} (Belgium) --- listed as a QTSP in Belgium,
supervised by the FPS Economy. Strong enterprise presence. AATL
membership plus EUTL listing gives broad compatibility.

\textbf{Swisscom Trust Services} (Austria, Switzerland) --- the Swiss
telecoms giant operates a QTSP listed in the \emph{Austrian} TSL (not
Swiss --- because Switzerland is not in the EU) plus ZertES
qualification in Switzerland. Their remote signing service is popular
for high-volume enterprise signing.

\textbf{InfoCert} (Italy) --- one of the largest European QTSPs by
volume. Listed in the Italian TSL, supervised by AgID. Part of the
Tinexta Group. Deeply embedded in Italian public administration.

\textbf{Asseco Data Systems / Certum} (Poland) --- the dominant QTSP in
Central/Eastern Europe. If you encounter qualified certificates from
Polish organizations, Certum is the likely issuer.

This list is not exhaustive --- there are over 200 qualified trust
services listed across the EUTL. The EU maintains a browsable version at
\texttt{https://eidas.ec.europa.eu/efda/tl-browser/} where you can
filter by country, service type, and status.

\section{Java Code: Parsing the EUTL}\label{_java_code_parsing_the_eutl}

\subsection{Downloading and Parsing the
LOTL}\label{_downloading_and_parsing_the_lotl}

\begin{Shaded}
\begin{Highlighting}[]
\KeywordTok{import} \ImportTok{java}\OperatorTok{.}\ImportTok{io}\OperatorTok{.}\ImportTok{InputStream}\OperatorTok{;}
\KeywordTok{import} \ImportTok{java}\OperatorTok{.}\ImportTok{net}\OperatorTok{.}\ImportTok{URI}\OperatorTok{;}
\KeywordTok{import} \ImportTok{java}\OperatorTok{.}\ImportTok{util}\OperatorTok{.}\ImportTok{ArrayList}\OperatorTok{;}
\KeywordTok{import} \ImportTok{java}\OperatorTok{.}\ImportTok{util}\OperatorTok{.}\ImportTok{List}\OperatorTok{;}

\KeywordTok{import} \ImportTok{javax}\OperatorTok{.}\ImportTok{xml}\OperatorTok{.}\ImportTok{parsers}\OperatorTok{.}\ImportTok{DocumentBuilderFactory}\OperatorTok{;}

\KeywordTok{import} \ImportTok{org}\OperatorTok{.}\ImportTok{w3c}\OperatorTok{.}\ImportTok{dom}\OperatorTok{.}\ImportTok{Document}\OperatorTok{;}
\KeywordTok{import} \ImportTok{org}\OperatorTok{.}\ImportTok{w3c}\OperatorTok{.}\ImportTok{dom}\OperatorTok{.}\ImportTok{Element}\OperatorTok{;}
\KeywordTok{import} \ImportTok{org}\OperatorTok{.}\ImportTok{w3c}\OperatorTok{.}\ImportTok{dom}\OperatorTok{.}\ImportTok{NodeList}\OperatorTok{;}

\KeywordTok{import} \ImportTok{lombok}\OperatorTok{.}\ImportTok{Data}\OperatorTok{;}
\KeywordTok{import} \ImportTok{lombok}\OperatorTok{.}\ImportTok{extern}\OperatorTok{.}\ImportTok{slf4j}\OperatorTok{.}\ImportTok{Slf4j}\OperatorTok{;}

\CommentTok{/**}
 \CommentTok{*}\NormalTok{ Parses the EU List of Trusted Lists }\CommentTok{(}\NormalTok{LOTL}\CommentTok{)}\NormalTok{ to extract}
 \CommentTok{*}\NormalTok{ pointers to each national Trusted Service Status List}\CommentTok{.}
 \CommentTok{*/}
\AttributeTok{@Slf4j}
\KeywordTok{public} \KeywordTok{class}\NormalTok{ EutlLotlParser }\OperatorTok{\{}

    \KeywordTok{private} \DataTypeTok{static} \DataTypeTok{final} \BuiltInTok{String}\NormalTok{ LOTL\_URL }\OperatorTok{=}
            \StringTok{"https://ec.europa.eu/tools/lotl/eu{-}lotl.xml"}\OperatorTok{;}
    \KeywordTok{private} \DataTypeTok{static} \DataTypeTok{final} \BuiltInTok{String}\NormalTok{ TSL\_NS }\OperatorTok{=}
            \StringTok{"http://uri.etsi.org/02231/v2\#"}\OperatorTok{;}

    \AttributeTok{@Data}
    \KeywordTok{public} \DataTypeTok{static} \KeywordTok{class}\NormalTok{ TslPointer }\OperatorTok{\{}
        \KeywordTok{private} \BuiltInTok{String}\NormalTok{ countryCode}\OperatorTok{;}
        \KeywordTok{private} \BuiltInTok{String}\NormalTok{ tslUrl}\OperatorTok{;}
    \OperatorTok{\}}

    \KeywordTok{public} \BuiltInTok{List}\OperatorTok{\textless{}}\NormalTok{TslPointer}\OperatorTok{\textgreater{}} \FunctionTok{parseLotl}\OperatorTok{()} \KeywordTok{throws} \BuiltInTok{Exception} \OperatorTok{\{}
\NormalTok{        log}\OperatorTok{.}\FunctionTok{info}\OperatorTok{(}\StringTok{"Downloading LOTL from: \{\}"}\OperatorTok{,}\NormalTok{ LOTL\_URL}\OperatorTok{);}

        \BuiltInTok{DocumentBuilderFactory}\NormalTok{ factory }\OperatorTok{=}
                \BuiltInTok{DocumentBuilderFactory}\OperatorTok{.}\FunctionTok{newInstance}\OperatorTok{();}
\NormalTok{        factory}\OperatorTok{.}\FunctionTok{setNamespaceAware}\OperatorTok{(}\KeywordTok{true}\OperatorTok{);}
        \CommentTok{// Prevent XXE attacks}
\NormalTok{        factory}\OperatorTok{.}\FunctionTok{setFeature}\OperatorTok{(}
                \StringTok{"http://apache.org/xml/features/"}
                \OperatorTok{+} \StringTok{"disallow{-}doctype{-}decl"}\OperatorTok{,} \KeywordTok{true}\OperatorTok{);}

        \BuiltInTok{Document}\NormalTok{ doc}\OperatorTok{;}
        \ControlFlowTok{try} \OperatorTok{(}\BuiltInTok{InputStream}\NormalTok{ in }\OperatorTok{=} \BuiltInTok{URI}\OperatorTok{.}\FunctionTok{create}\OperatorTok{(}\NormalTok{LOTL\_URL}\OperatorTok{).}\FunctionTok{toURL}\OperatorTok{()}
                \OperatorTok{.}\FunctionTok{openStream}\OperatorTok{())} \OperatorTok{\{}
\NormalTok{            doc }\OperatorTok{=}\NormalTok{ factory}\OperatorTok{.}\FunctionTok{newDocumentBuilder}\OperatorTok{().}\FunctionTok{parse}\OperatorTok{(}\NormalTok{in}\OperatorTok{);}
        \OperatorTok{\}}

        \BuiltInTok{List}\OperatorTok{\textless{}}\NormalTok{TslPointer}\OperatorTok{\textgreater{}}\NormalTok{ pointers }\OperatorTok{=} \KeywordTok{new} \BuiltInTok{ArrayList}\OperatorTok{\textless{}\textgreater{}();}
        \BuiltInTok{NodeList}\NormalTok{ nodes }\OperatorTok{=}\NormalTok{ doc}\OperatorTok{.}\FunctionTok{getElementsByTagNameNS}\OperatorTok{(}
\NormalTok{                TSL\_NS}\OperatorTok{,} \StringTok{"OtherTSLPointer"}\OperatorTok{);}

        \ControlFlowTok{for} \OperatorTok{(}\DataTypeTok{int}\NormalTok{ i }\OperatorTok{=} \DecValTok{0}\OperatorTok{;}\NormalTok{ i }\OperatorTok{\textless{}}\NormalTok{ nodes}\OperatorTok{.}\FunctionTok{getLength}\OperatorTok{();}\NormalTok{ i}\OperatorTok{++)} \OperatorTok{\{}
            \BuiltInTok{Element}\NormalTok{ pointer }\OperatorTok{=} \OperatorTok{(}\BuiltInTok{Element}\OperatorTok{)}\NormalTok{ nodes}\OperatorTok{.}\FunctionTok{item}\OperatorTok{(}\NormalTok{i}\OperatorTok{);}

            \BuiltInTok{NodeList}\NormalTok{ locations }\OperatorTok{=}\NormalTok{ pointer}\OperatorTok{.}\FunctionTok{getElementsByTagNameNS}\OperatorTok{(}
\NormalTok{                    TSL\_NS}\OperatorTok{,} \StringTok{"TSLLocation"}\OperatorTok{);}
            \ControlFlowTok{if} \OperatorTok{(}\NormalTok{locations}\OperatorTok{.}\FunctionTok{getLength}\OperatorTok{()} \OperatorTok{==} \DecValTok{0}\OperatorTok{)} \OperatorTok{\{}
                \ControlFlowTok{continue}\OperatorTok{;}
            \OperatorTok{\}}
            \BuiltInTok{String}\NormalTok{ url }\OperatorTok{=}\NormalTok{ locations}\OperatorTok{.}\FunctionTok{item}\OperatorTok{(}\DecValTok{0}\OperatorTok{).}\FunctionTok{getTextContent}\OperatorTok{().}\FunctionTok{trim}\OperatorTok{();}

            \BuiltInTok{NodeList}\NormalTok{ territories }\OperatorTok{=}\NormalTok{ pointer}\OperatorTok{.}\FunctionTok{getElementsByTagNameNS}\OperatorTok{(}
\NormalTok{                    TSL\_NS}\OperatorTok{,} \StringTok{"SchemeTerritory"}\OperatorTok{);}
            \BuiltInTok{String}\NormalTok{ country }\OperatorTok{=}\NormalTok{ territories}\OperatorTok{.}\FunctionTok{getLength}\OperatorTok{()} \OperatorTok{\textgreater{}} \DecValTok{0}
                    \OperatorTok{?}\NormalTok{ territories}\OperatorTok{.}\FunctionTok{item}\OperatorTok{(}\DecValTok{0}\OperatorTok{).}\FunctionTok{getTextContent}\OperatorTok{().}\FunctionTok{trim}\OperatorTok{()}
                    \OperatorTok{:} \StringTok{"UNKNOWN"}\OperatorTok{;}

\NormalTok{            TslPointer p }\OperatorTok{=} \KeywordTok{new} \FunctionTok{TslPointer}\OperatorTok{();}
\NormalTok{            p}\OperatorTok{.}\FunctionTok{setCountryCode}\OperatorTok{(}\NormalTok{country}\OperatorTok{);}
\NormalTok{            p}\OperatorTok{.}\FunctionTok{setTslUrl}\OperatorTok{(}\NormalTok{url}\OperatorTok{);}
\NormalTok{            pointers}\OperatorTok{.}\FunctionTok{add}\OperatorTok{(}\NormalTok{p}\OperatorTok{);}

\NormalTok{            log}\OperatorTok{.}\FunctionTok{debug}\OperatorTok{(}\StringTok{"Found TSL: country=\{\}, url=\{\}"}\OperatorTok{,}\NormalTok{ country}\OperatorTok{,}\NormalTok{ url}\OperatorTok{);}
        \OperatorTok{\}}

\NormalTok{        log}\OperatorTok{.}\FunctionTok{info}\OperatorTok{(}\StringTok{"Parsed LOTL: \{\} national TSL pointers"}\OperatorTok{,}
\NormalTok{                pointers}\OperatorTok{.}\FunctionTok{size}\OperatorTok{());}
        \ControlFlowTok{return}\NormalTok{ pointers}\OperatorTok{;}
    \OperatorTok{\}}
\OperatorTok{\}}
\end{Highlighting}
\end{Shaded}

\subsection{Matching Certificates Against the
EUTL}\label{_matching_certificates_against_the_eutl}

\begin{Shaded}
\begin{Highlighting}[]
\KeywordTok{import} \ImportTok{java}\OperatorTok{.}\ImportTok{io}\OperatorTok{.}\ImportTok{ByteArrayInputStream}\OperatorTok{;}
\KeywordTok{import} \ImportTok{java}\OperatorTok{.}\ImportTok{io}\OperatorTok{.}\ImportTok{InputStream}\OperatorTok{;}
\KeywordTok{import} \ImportTok{java}\OperatorTok{.}\ImportTok{net}\OperatorTok{.}\ImportTok{URI}\OperatorTok{;}
\KeywordTok{import} \ImportTok{java}\OperatorTok{.}\ImportTok{security}\OperatorTok{.}\ImportTok{cert}\OperatorTok{.}\ImportTok{CertificateFactory}\OperatorTok{;}
\KeywordTok{import} \ImportTok{java}\OperatorTok{.}\ImportTok{security}\OperatorTok{.}\ImportTok{cert}\OperatorTok{.}\ImportTok{X509Certificate}\OperatorTok{;}
\KeywordTok{import} \ImportTok{java}\OperatorTok{.}\ImportTok{util}\OperatorTok{.}\ImportTok{ArrayList}\OperatorTok{;}
\KeywordTok{import} \ImportTok{java}\OperatorTok{.}\ImportTok{util}\OperatorTok{.}\ImportTok{Arrays}\OperatorTok{;}
\KeywordTok{import} \ImportTok{java}\OperatorTok{.}\ImportTok{util}\OperatorTok{.}\ImportTok{Base64}\OperatorTok{;}
\KeywordTok{import} \ImportTok{java}\OperatorTok{.}\ImportTok{util}\OperatorTok{.}\ImportTok{List}\OperatorTok{;}

\KeywordTok{import} \ImportTok{javax}\OperatorTok{.}\ImportTok{xml}\OperatorTok{.}\ImportTok{parsers}\OperatorTok{.}\ImportTok{DocumentBuilderFactory}\OperatorTok{;}

\KeywordTok{import} \ImportTok{org}\OperatorTok{.}\ImportTok{w3c}\OperatorTok{.}\ImportTok{dom}\OperatorTok{.}\ImportTok{Document}\OperatorTok{;}
\KeywordTok{import} \ImportTok{org}\OperatorTok{.}\ImportTok{w3c}\OperatorTok{.}\ImportTok{dom}\OperatorTok{.}\ImportTok{Element}\OperatorTok{;}
\KeywordTok{import} \ImportTok{org}\OperatorTok{.}\ImportTok{w3c}\OperatorTok{.}\ImportTok{dom}\OperatorTok{.}\ImportTok{NodeList}\OperatorTok{;}

\KeywordTok{import} \ImportTok{lombok}\OperatorTok{.}\ImportTok{Data}\OperatorTok{;}
\KeywordTok{import} \ImportTok{lombok}\OperatorTok{.}\ImportTok{extern}\OperatorTok{.}\ImportTok{slf4j}\OperatorTok{.}\ImportTok{Slf4j}\OperatorTok{;}

\CommentTok{/**}
 \CommentTok{*}\NormalTok{ Parses a national TSL and matches signer certificates}
 \CommentTok{*}\NormalTok{ against qualified CA services listed in the EUTL}\CommentTok{.}
 \CommentTok{*/}
\AttributeTok{@Slf4j}
\KeywordTok{public} \KeywordTok{class}\NormalTok{ EutlCertificateMatcher }\OperatorTok{\{}

    \KeywordTok{private} \DataTypeTok{static} \DataTypeTok{final} \BuiltInTok{String}\NormalTok{ TSL\_NS }\OperatorTok{=}
            \StringTok{"http://uri.etsi.org/02231/v2\#"}\OperatorTok{;}

    \AttributeTok{@Data}
    \KeywordTok{public} \DataTypeTok{static} \KeywordTok{class}\NormalTok{ QualifiedCaEntry }\OperatorTok{\{}
        \KeywordTok{private} \BuiltInTok{String}\NormalTok{ providerName}\OperatorTok{;}
        \KeywordTok{private} \BuiltInTok{String}\NormalTok{ status}\OperatorTok{;}
        \KeywordTok{private} \BuiltInTok{X509Certificate}\NormalTok{ caCertificate}\OperatorTok{;}
    \OperatorTok{\}}

    \CommentTok{/**}
     \CommentTok{*}\NormalTok{ Downloads a national TSL and extracts all CA}\CommentTok{/}\NormalTok{QC services}
     \CommentTok{*}\NormalTok{ with }\CommentTok{"}\NormalTok{granted}\CommentTok{"}\NormalTok{ status}\CommentTok{.}
     \CommentTok{*/}
    \KeywordTok{public} \BuiltInTok{List}\OperatorTok{\textless{}}\NormalTok{QualifiedCaEntry}\OperatorTok{\textgreater{}} \FunctionTok{parseNationalTsl}\OperatorTok{(}\BuiltInTok{String}\NormalTok{ tslUrl}\OperatorTok{)}
            \KeywordTok{throws} \BuiltInTok{Exception} \OperatorTok{\{}
\NormalTok{        log}\OperatorTok{.}\FunctionTok{info}\OperatorTok{(}\StringTok{"Parsing national TSL: \{\}"}\OperatorTok{,}\NormalTok{ tslUrl}\OperatorTok{);}

        \BuiltInTok{DocumentBuilderFactory}\NormalTok{ factory }\OperatorTok{=}
                \BuiltInTok{DocumentBuilderFactory}\OperatorTok{.}\FunctionTok{newInstance}\OperatorTok{();}
\NormalTok{        factory}\OperatorTok{.}\FunctionTok{setNamespaceAware}\OperatorTok{(}\KeywordTok{true}\OperatorTok{);}
\NormalTok{        factory}\OperatorTok{.}\FunctionTok{setFeature}\OperatorTok{(}
                \StringTok{"http://apache.org/xml/features/"}
                \OperatorTok{+} \StringTok{"disallow{-}doctype{-}decl"}\OperatorTok{,} \KeywordTok{true}\OperatorTok{);}

        \BuiltInTok{Document}\NormalTok{ doc}\OperatorTok{;}
        \ControlFlowTok{try} \OperatorTok{(}\BuiltInTok{InputStream}\NormalTok{ in }\OperatorTok{=} \BuiltInTok{URI}\OperatorTok{.}\FunctionTok{create}\OperatorTok{(}\NormalTok{tslUrl}\OperatorTok{).}\FunctionTok{toURL}\OperatorTok{()}
                \OperatorTok{.}\FunctionTok{openStream}\OperatorTok{())} \OperatorTok{\{}
\NormalTok{            doc }\OperatorTok{=}\NormalTok{ factory}\OperatorTok{.}\FunctionTok{newDocumentBuilder}\OperatorTok{().}\FunctionTok{parse}\OperatorTok{(}\NormalTok{in}\OperatorTok{);}
        \OperatorTok{\}}

        \BuiltInTok{List}\OperatorTok{\textless{}}\NormalTok{QualifiedCaEntry}\OperatorTok{\textgreater{}}\NormalTok{ entries }\OperatorTok{=} \KeywordTok{new} \BuiltInTok{ArrayList}\OperatorTok{\textless{}\textgreater{}();}
        \BuiltInTok{NodeList}\NormalTok{ services }\OperatorTok{=}\NormalTok{ doc}\OperatorTok{.}\FunctionTok{getElementsByTagNameNS}\OperatorTok{(}
\NormalTok{                TSL\_NS}\OperatorTok{,} \StringTok{"TSPService"}\OperatorTok{);}
        \BuiltInTok{CertificateFactory}\NormalTok{ cf }\OperatorTok{=}
                \BuiltInTok{CertificateFactory}\OperatorTok{.}\FunctionTok{getInstance}\OperatorTok{(}\StringTok{"X.509"}\OperatorTok{);}

        \ControlFlowTok{for} \OperatorTok{(}\DataTypeTok{int}\NormalTok{ i }\OperatorTok{=} \DecValTok{0}\OperatorTok{;}\NormalTok{ i }\OperatorTok{\textless{}}\NormalTok{ services}\OperatorTok{.}\FunctionTok{getLength}\OperatorTok{();}\NormalTok{ i}\OperatorTok{++)} \OperatorTok{\{}
            \BuiltInTok{Element}\NormalTok{ svc }\OperatorTok{=} \OperatorTok{(}\BuiltInTok{Element}\OperatorTok{)}\NormalTok{ services}\OperatorTok{.}\FunctionTok{item}\OperatorTok{(}\NormalTok{i}\OperatorTok{);}

            \CommentTok{// Check service type}
            \BuiltInTok{NodeList}\NormalTok{ typeNodes }\OperatorTok{=}\NormalTok{ svc}\OperatorTok{.}\FunctionTok{getElementsByTagNameNS}\OperatorTok{(}
\NormalTok{                    TSL\_NS}\OperatorTok{,} \StringTok{"ServiceTypeIdentifier"}\OperatorTok{);}
            \ControlFlowTok{if} \OperatorTok{(}\NormalTok{typeNodes}\OperatorTok{.}\FunctionTok{getLength}\OperatorTok{()} \OperatorTok{==} \DecValTok{0}\OperatorTok{)} \OperatorTok{\{}
                \ControlFlowTok{continue}\OperatorTok{;}
            \OperatorTok{\}}
            \BuiltInTok{String}\NormalTok{ type }\OperatorTok{=}\NormalTok{ typeNodes}\OperatorTok{.}\FunctionTok{item}\OperatorTok{(}\DecValTok{0}\OperatorTok{)}
                    \OperatorTok{.}\FunctionTok{getTextContent}\OperatorTok{().}\FunctionTok{trim}\OperatorTok{();}
            \ControlFlowTok{if} \OperatorTok{(!}\NormalTok{type}\OperatorTok{.}\FunctionTok{contains}\OperatorTok{(}\StringTok{"CA/QC"}\OperatorTok{))} \OperatorTok{\{}
                \ControlFlowTok{continue}\OperatorTok{;}
            \OperatorTok{\}}

            \CommentTok{// Check status}
            \BuiltInTok{NodeList}\NormalTok{ statusNodes }\OperatorTok{=}\NormalTok{ svc}\OperatorTok{.}\FunctionTok{getElementsByTagNameNS}\OperatorTok{(}
\NormalTok{                    TSL\_NS}\OperatorTok{,} \StringTok{"ServiceStatus"}\OperatorTok{);}
            \ControlFlowTok{if} \OperatorTok{(}\NormalTok{statusNodes}\OperatorTok{.}\FunctionTok{getLength}\OperatorTok{()} \OperatorTok{==} \DecValTok{0}\OperatorTok{)} \OperatorTok{\{}
                \ControlFlowTok{continue}\OperatorTok{;}
            \OperatorTok{\}}
            \BuiltInTok{String}\NormalTok{ status }\OperatorTok{=}\NormalTok{ statusNodes}\OperatorTok{.}\FunctionTok{item}\OperatorTok{(}\DecValTok{0}\OperatorTok{)}
                    \OperatorTok{.}\FunctionTok{getTextContent}\OperatorTok{().}\FunctionTok{trim}\OperatorTok{();}
            \ControlFlowTok{if} \OperatorTok{(!}\NormalTok{status}\OperatorTok{.}\FunctionTok{contains}\OperatorTok{(}\StringTok{"granted"}\OperatorTok{))} \OperatorTok{\{}
                \ControlFlowTok{continue}\OperatorTok{;}
            \OperatorTok{\}}

            \CommentTok{// Extract CA certificate}
            \BuiltInTok{NodeList}\NormalTok{ certNodes }\OperatorTok{=}\NormalTok{ svc}\OperatorTok{.}\FunctionTok{getElementsByTagNameNS}\OperatorTok{(}
\NormalTok{                    TSL\_NS}\OperatorTok{,} \StringTok{"X509Certificate"}\OperatorTok{);}
            \ControlFlowTok{if} \OperatorTok{(}\NormalTok{certNodes}\OperatorTok{.}\FunctionTok{getLength}\OperatorTok{()} \OperatorTok{==} \DecValTok{0}\OperatorTok{)} \OperatorTok{\{}
                \ControlFlowTok{continue}\OperatorTok{;}
            \OperatorTok{\}}
            \BuiltInTok{String}\NormalTok{ b64 }\OperatorTok{=}\NormalTok{ certNodes}\OperatorTok{.}\FunctionTok{item}\OperatorTok{(}\DecValTok{0}\OperatorTok{).}\FunctionTok{getTextContent}\OperatorTok{()}
                    \OperatorTok{.}\FunctionTok{replaceAll}\OperatorTok{(}\StringTok{"}\SpecialCharTok{\textbackslash{}\textbackslash{}}\StringTok{s"}\OperatorTok{,} \StringTok{""}\OperatorTok{);}
            \DataTypeTok{byte}\OperatorTok{[]}\NormalTok{ certBytes }\OperatorTok{=}\NormalTok{ Base64}\OperatorTok{.}\FunctionTok{getDecoder}\OperatorTok{().}\FunctionTok{decode}\OperatorTok{(}\NormalTok{b64}\OperatorTok{);}
            \BuiltInTok{X509Certificate}\NormalTok{ cert }\OperatorTok{=} \OperatorTok{(}\BuiltInTok{X509Certificate}\OperatorTok{)}
\NormalTok{                    cf}\OperatorTok{.}\FunctionTok{generateCertificate}\OperatorTok{(}
                        \KeywordTok{new} \BuiltInTok{ByteArrayInputStream}\OperatorTok{(}\NormalTok{certBytes}\OperatorTok{));}

\NormalTok{            QualifiedCaEntry entry }\OperatorTok{=} \KeywordTok{new} \FunctionTok{QualifiedCaEntry}\OperatorTok{();}
\NormalTok{            entry}\OperatorTok{.}\FunctionTok{setStatus}\OperatorTok{(}\NormalTok{status}\OperatorTok{);}
\NormalTok{            entry}\OperatorTok{.}\FunctionTok{setCaCertificate}\OperatorTok{(}\NormalTok{cert}\OperatorTok{);}
\NormalTok{            entries}\OperatorTok{.}\FunctionTok{add}\OperatorTok{(}\NormalTok{entry}\OperatorTok{);}

\NormalTok{            log}\OperatorTok{.}\FunctionTok{info}\OperatorTok{(}\StringTok{"Found qualified CA: \{\}"}\OperatorTok{,}
\NormalTok{                    cert}\OperatorTok{.}\FunctionTok{getSubjectX500Principal}\OperatorTok{());}
        \OperatorTok{\}}

        \ControlFlowTok{return}\NormalTok{ entries}\OperatorTok{;}
    \OperatorTok{\}}

    \CommentTok{/**}
     \CommentTok{*}\NormalTok{ Checks whether the given signer certificate was issued}
     \CommentTok{*}\NormalTok{ by a CA listed in the EUTL with }\CommentTok{"}\NormalTok{granted}\CommentTok{"}\NormalTok{ status}\CommentTok{.}
     \CommentTok{*/}
    \KeywordTok{public} \DataTypeTok{boolean} \FunctionTok{isIssuedByQualifiedCa}\OperatorTok{(}
            \BuiltInTok{X509Certificate}\NormalTok{ signerCert}\OperatorTok{,}
            \BuiltInTok{X509Certificate}\NormalTok{ issuerCert}\OperatorTok{,}
            \BuiltInTok{List}\OperatorTok{\textless{}}\NormalTok{QualifiedCaEntry}\OperatorTok{\textgreater{}}\NormalTok{ eutlEntries}\OperatorTok{)} \OperatorTok{\{}
        \DataTypeTok{byte}\OperatorTok{[]}\NormalTok{ issuerEncoded}\OperatorTok{;}
        \ControlFlowTok{try} \OperatorTok{\{}
\NormalTok{            issuerEncoded }\OperatorTok{=}\NormalTok{ issuerCert}\OperatorTok{.}\FunctionTok{getEncoded}\OperatorTok{();}
        \OperatorTok{\}} \ControlFlowTok{catch} \OperatorTok{(}\BuiltInTok{Exception}\NormalTok{ e}\OperatorTok{)} \OperatorTok{\{}
\NormalTok{            log}\OperatorTok{.}\FunctionTok{error}\OperatorTok{(}\StringTok{"Failed to encode issuer cert"}\OperatorTok{,}\NormalTok{ e}\OperatorTok{);}
            \ControlFlowTok{return} \KeywordTok{false}\OperatorTok{;}
        \OperatorTok{\}}

        \ControlFlowTok{for} \OperatorTok{(}\NormalTok{QualifiedCaEntry entry }\OperatorTok{:}\NormalTok{ eutlEntries}\OperatorTok{)} \OperatorTok{\{}
            \ControlFlowTok{try} \OperatorTok{\{}
                \DataTypeTok{byte}\OperatorTok{[]}\NormalTok{ eutlEncoded }\OperatorTok{=}
\NormalTok{                        entry}\OperatorTok{.}\FunctionTok{getCaCertificate}\OperatorTok{().}\FunctionTok{getEncoded}\OperatorTok{();}
                \ControlFlowTok{if} \OperatorTok{(}\BuiltInTok{Arrays}\OperatorTok{.}\FunctionTok{equals}\OperatorTok{(}\NormalTok{issuerEncoded}\OperatorTok{,}\NormalTok{ eutlEncoded}\OperatorTok{))} \OperatorTok{\{}
\NormalTok{                    log}\OperatorTok{.}\FunctionTok{info}\OperatorTok{(}\StringTok{"Issuer matched EUTL entry: \{\}"}\OperatorTok{,}
\NormalTok{                            entry}\OperatorTok{.}\FunctionTok{getCaCertificate}\OperatorTok{()}
                                \OperatorTok{.}\FunctionTok{getSubjectX500Principal}\OperatorTok{());}
                    \ControlFlowTok{return} \KeywordTok{true}\OperatorTok{;}
                \OperatorTok{\}}
            \OperatorTok{\}} \ControlFlowTok{catch} \OperatorTok{(}\BuiltInTok{Exception}\NormalTok{ e}\OperatorTok{)} \OperatorTok{\{}
\NormalTok{                log}\OperatorTok{.}\FunctionTok{warn}\OperatorTok{(}\StringTok{"Failed to encode EUTL cert"}\OperatorTok{,}\NormalTok{ e}\OperatorTok{);}
            \OperatorTok{\}}
        \OperatorTok{\}}
        \ControlFlowTok{return} \KeywordTok{false}\OperatorTok{;}
    \OperatorTok{\}}
\OperatorTok{\}}
\end{Highlighting}
\end{Shaded}

The code above is educational --- it demonstrates the mechanics of EUTL
parsing and certificate matching. A production implementation must also:

\begin{itemize}
\item
  \textbf{Verify the XML signatures} on the LOTL and every national TSL
\item
  \textbf{Handle the LOTL pivot mechanism} when the signing certificate
  rotates
\item
  \textbf{Parse service history} to support time-based qualification
  checks
\item
  \textbf{Cache and refresh} the hierarchy on a schedule (not on every
  validation)
\item
  \textbf{Parse QCStatements} from the signer's certificate to determine
  the qualification type
\end{itemize}

For production use, the \textbf{EU DSS library} (Chapter 23) provides
\texttt{TrustedListsCertificateSource}, which handles all of this
correctly. Do not reimplement this from scratch unless you have a very
specific reason.

\section{Key Takeaways}\label{_key_takeaways_8}

\begin{itemize}
\item
  A \textbf{QTSP} is not a self-declared status --- it requires
  conformity assessment, supervisory body approval, and listing in the
  EUTL; the process takes 18-36 months and significant investment
\item
  The \textbf{EUTL} is a hierarchy: the LOTL (Commission) points to
  national TSLs (supervisory bodies); every document is
  cryptographically signed with XAdES
\item
  A certificate is qualified when: its issuer is in the EUTL with
  \texttt{granted} status for \texttt{CA/QC}, \emph{and} the certificate
  contains the \texttt{QcCompliance} QCStatement
\item
  \textbf{QCStatements} carry machine-readable qualification data:
  QcCompliance, QcSSCD, QcType (esign/eseal/web)
\item
  \textbf{Certificate Policy OIDs} from ETSI EN 319 411-2 encode the
  qualification policy: QCP-n, QCP-l, with \texttt{-qscd} variants
\item
  \textbf{QSCDs} require Common Criteria EAL4+ certification against EN
  419 211/221 Protection Profiles --- not every HSM qualifies
\item
  Check \textbf{service history} when validating: a certificate issued
  during a \texttt{granted} period is qualified even if the service is
  now \texttt{withdrawn}
\item
  Real-world QTSPs include D-Trust (DE), DigiCert/QuoVadis (NL/BE),
  GlobalSign (BE), Swisscom (AT/CH), InfoCert (IT), Certum (PL), and
  over 200 others
\item
  For production EUTL validation, use the EU DSS framework's
  \texttt{TrustedListsCertificateSource} --- do not reimplement the
  parsing, pivot handling, and signature verification from scratch
\end{itemize}

\chapter{Qualified Signatures vs
Seals}\label{_qualified_signatures_vs_seals}

Here is a question that trips up almost every developer the first time
they encounter it: if an electronic signature and an electronic seal use
exactly the same cryptographic mechanism --- the same RSA or ECDSA
algorithm, the same CMS SignedData structure, the same PAdES profile,
the same qualified certificate from the same QTSP, stored on the same
model of QSCD --- what, exactly, is the difference?

The answer is not technical. It is legal. A signature is created by a
\emph{natural person} (a human being) and expresses that person's
intent. A seal is created by a \emph{legal person} (a company,
government agency, association) and proves the origin and integrity of
data. Same bytes, same math, different legal meaning --- and getting it
wrong can mean the difference between a legally binding contract and an
authenticated invoice.

\section{The Fundamental
Distinction}\label{_the_fundamental_distinction}

\subsection{Signatures: Natural
Persons}\label{_signatures_natural_persons}

An electronic signature under Regulation (EU) No 910/2014 (eIDAS) is
always attributed to a \textbf{natural person}. Article 3(10) defines an
electronic signature as "data in electronic form which is attached to or
logically associated with other data in electronic form and which is
used by the signatory to sign." Article 3(9) defines "signatory" as "a
natural person who creates an electronic signature."

A company cannot sign. A department cannot sign. A server cannot sign.
Only a human can sign, and a signature expresses that human's
\textbf{intent} --- consent, approval, acknowledgment, or whatever the
signing context requires.

\subsection{Seals: Legal Persons}\label{_seals_legal_persons}

An electronic seal is attributed to a \textbf{legal person} --- a
company, government agency, or any entity with legal personality.
Regulation (EU) No 910/2014 (eIDAS) Article 3(25) defines an electronic
seal as "data in electronic form, which is attached to or logically
associated with other data in electronic form to ensure the latter's
origin and integrity."

Note the careful wording. A seal does not express \emph{intent}. It
proves \textbf{origin} (the data came from this organization) and
\textbf{integrity} (the data has not been tampered with). When Acme Corp
seals an invoice, the seal says: "This invoice was produced by Acme
Corp, and it has not been modified since." It says nothing about which
employee generated it or whether anyone reviewed it.

\section{eIDAS Articles 25 vs 35}\label{_eidas_articles_25_vs_35}

\subsection{Article 25: Signatures}\label{_article_25_signatures}

\begin{itemize}
\item
  \textbf{25(1)}: Non-discrimination --- an electronic signature shall
  not be denied legal effect solely because it is electronic
\item
  \textbf{25(2)}: A qualified electronic signature has the
  \textbf{equivalent legal effect of a handwritten signature}
\item
  \textbf{25(3)}: A QES from one member state is recognized in all
  member states
\end{itemize}

The critical phrase in 25(2) is "equivalent legal effect." This is the
strongest possible legal recognition --- a QES \emph{is} a handwritten
signature for all legal purposes.

\subsection{Article 35: Seals}\label{_article_35_seals}

\begin{itemize}
\item
  \textbf{35(1)}: Non-discrimination (same principle as 25(1))
\item
  \textbf{35(2)}: A qualified electronic seal enjoys a
  \textbf{presumption of integrity of the data and of correctness of the
  origin}
\item
  \textbf{35(3)}: Mutual recognition across member states
\item
  \textbf{Article 36}: Advanced seal requirements (same four properties
  as Art. 26, adapted for legal persons)
\item
  \textbf{Article 38}: Qualified certificates for seals must meet Annex
  III requirements
\item
  \textbf{Article 40}: Member states may define additional legal effects
  beyond the presumption
\end{itemize}

\begin{longtable}[]{@{}
  >{\raggedright\arraybackslash}p{(\linewidth - 4\tabcolsep) * \real{0.2500}}
  >{\raggedright\arraybackslash}p{(\linewidth - 4\tabcolsep) * \real{0.3750}}
  >{\raggedright\arraybackslash}p{(\linewidth - 4\tabcolsep) * \real{0.3750}}@{}}
\caption{Legal effects comparison}\tabularnewline
\toprule\noalign{}
\begin{minipage}[b]{\linewidth}\raggedright
Aspect
\end{minipage} & \begin{minipage}[b]{\linewidth}\raggedright
Qualified Signature (QES)
\end{minipage} & \begin{minipage}[b]{\linewidth}\raggedright
Qualified Seal (QESeal)
\end{minipage} \\
\midrule\noalign{}
\endfirsthead
\toprule\noalign{}
\begin{minipage}[b]{\linewidth}\raggedright
Aspect
\end{minipage} & \begin{minipage}[b]{\linewidth}\raggedright
Qualified Signature (QES)
\end{minipage} & \begin{minipage}[b]{\linewidth}\raggedright
Qualified Seal (QESeal)
\end{minipage} \\
\midrule\noalign{}
\endhead
\bottomrule\noalign{}
\endlastfoot
Created by & Natural person (human) & Legal person (organization) \\
Legal effect & Equivalent to handwritten signature (Art. 25(2)) &
Presumption of integrity and origin (Art. 35(2)) \\
Expresses intent? & Yes --- the signer intends to be bound &
No --- proves origin and integrity only \\
Cross-border & Mandatory recognition (Art. 25(3)) & Mandatory
recognition (Art. 35(3)) \\
\end{longtable}

The difference between "legal equivalence" and "presumption of integrity
and origin" is not academic. A QES provides \textbf{non-repudiation of
commitment} --- the signer personally committed to the content. A QESeal
provides \textbf{non-repudiation of origin} --- the data came from the
stated organization.

If a contract requires a signature, a seal will not suffice --- even a
qualified seal. A seal on a contract proves only that the document came
from the company, not that any authorized representative agreed to its
terms. Conversely, for automated invoices, a seal is the \emph{correct}
choice --- attributing personal intent to a batch process would be
legally inappropriate.

\section{Same Code, Different
Certificate}\label{_same_code_different_certificate}

\subsection{What Differs in the
Certificate}\label{_what_differs_in_the_certificate}

From a cryptographic perspective, there is zero difference. Both are
X.509 certificates, both use the same algorithms, both can be issued by
the same QTSP. The differences are in the metadata:

\begin{longtable}[]{@{}
  >{\raggedright\arraybackslash}p{(\linewidth - 4\tabcolsep) * \real{0.2500}}
  >{\raggedright\arraybackslash}p{(\linewidth - 4\tabcolsep) * \real{0.3750}}
  >{\raggedright\arraybackslash}p{(\linewidth - 4\tabcolsep) * \real{0.3750}}@{}}
\caption{Certificate differences: signature vs seal}\tabularnewline
\toprule\noalign{}
\begin{minipage}[b]{\linewidth}\raggedright
Field
\end{minipage} & \begin{minipage}[b]{\linewidth}\raggedright
Signature Certificate
\end{minipage} & \begin{minipage}[b]{\linewidth}\raggedright
Seal Certificate
\end{minipage} \\
\midrule\noalign{}
\endfirsthead
\toprule\noalign{}
\begin{minipage}[b]{\linewidth}\raggedright
Field
\end{minipage} & \begin{minipage}[b]{\linewidth}\raggedright
Signature Certificate
\end{minipage} & \begin{minipage}[b]{\linewidth}\raggedright
Seal Certificate
\end{minipage} \\
\midrule\noalign{}
\endhead
\bottomrule\noalign{}
\endlastfoot
Subject CN & Natural person's name ("Maria Schmidt") & Organization name
("Acme GmbH") \\
Certificate Policy OID & \texttt{0.4.0.194112.1.2} (QCP-n) or
\texttt{…\hspace{0pt}1.1} (QCP-n-qscd) & \texttt{0.4.0.194112.1.3}
(QCP-l) or \texttt{…\hspace{0pt}1.4} (QCP-l-qscd) \\
QCStatements: QcType & \texttt{0.4.0.1862.1.6.1} (esign) &
\texttt{0.4.0.1862.1.6.2} (eseal) \\
eIDAS Annex & Annex I & Annex III \\
\end{longtable}

\subsection{Creating a Seal vs a Signature in
Java}\label{_creating_a_seal_vs_a_signature_in_java}

The code is \textbf{identical}. The only difference is which certificate
and private key you load.

\begin{Shaded}
\begin{Highlighting}[]
\KeywordTok{import} \ImportTok{java}\OperatorTok{.}\ImportTok{io}\OperatorTok{.}\ImportTok{ByteArrayOutputStream}\OperatorTok{;}
\KeywordTok{import} \ImportTok{java}\OperatorTok{.}\ImportTok{security}\OperatorTok{.}\ImportTok{PrivateKey}\OperatorTok{;}
\KeywordTok{import} \ImportTok{java}\OperatorTok{.}\ImportTok{security}\OperatorTok{.}\ImportTok{cert}\OperatorTok{.}\ImportTok{X509Certificate}\OperatorTok{;}

\KeywordTok{import} \ImportTok{org}\OperatorTok{.}\ImportTok{apache}\OperatorTok{.}\ImportTok{pdfbox}\OperatorTok{.}\ImportTok{Loader}\OperatorTok{;}
\KeywordTok{import} \ImportTok{org}\OperatorTok{.}\ImportTok{apache}\OperatorTok{.}\ImportTok{pdfbox}\OperatorTok{.}\ImportTok{pdmodel}\OperatorTok{.}\ImportTok{PDDocument}\OperatorTok{;}
\KeywordTok{import} \ImportTok{org}\OperatorTok{.}\ImportTok{apache}\OperatorTok{.}\ImportTok{pdfbox}\OperatorTok{.}\ImportTok{pdmodel}\OperatorTok{.}\ImportTok{interactive}\OperatorTok{.}\ImportTok{digitalsignature}\OperatorTok{.}\ImportTok{PDSignature}\OperatorTok{;}

\KeywordTok{import} \ImportTok{lombok}\OperatorTok{.}\ImportTok{RequiredArgsConstructor}\OperatorTok{;}
\KeywordTok{import} \ImportTok{lombok}\OperatorTok{.}\ImportTok{extern}\OperatorTok{.}\ImportTok{slf4j}\OperatorTok{.}\ImportTok{Slf4j}\OperatorTok{;}

\CommentTok{/**}
 \CommentTok{*}\NormalTok{ Creates either a qualified signature or a qualified seal on a}
 \CommentTok{*}\NormalTok{ PDF document}\CommentTok{. }\NormalTok{The cryptographic operation is identical }\CommentTok{{-}{-}}\NormalTok{ the}
 \CommentTok{*}\NormalTok{ distinction is entirely in the certificate used}\CommentTok{.}
 \CommentTok{*}
 \CommentTok{*}\NormalTok{ Signature}\CommentTok{:}\NormalTok{ cert has QcType esign }\CommentTok{(0.4.0.1862.1.6.1),}
 \CommentTok{*}\NormalTok{            subject }\CommentTok{=} \CommentTok{"}\NormalTok{CN}\CommentTok{=}\NormalTok{Alice Smith}\CommentTok{"}
 \CommentTok{*}\NormalTok{ Seal}\CommentTok{:}\NormalTok{      cert has QcType eseal }\CommentTok{(0.4.0.1862.1.6.2),}
 \CommentTok{*}\NormalTok{            subject }\CommentTok{=} \CommentTok{"}\NormalTok{O}\CommentTok{=}\NormalTok{Acme Corp}\CommentTok{,}\NormalTok{ C}\CommentTok{=}\NormalTok{DE}\CommentTok{"}
 \CommentTok{*/}
\AttributeTok{@Slf4j}
\AttributeTok{@RequiredArgsConstructor}
\KeywordTok{public} \KeywordTok{class}\NormalTok{ PdfSignOrSealService }\OperatorTok{\{}

    \KeywordTok{private} \DataTypeTok{final} \BuiltInTok{PrivateKey}\NormalTok{ privateKey}\OperatorTok{;}
    \KeywordTok{private} \DataTypeTok{final} \BuiltInTok{X509Certificate}\OperatorTok{[]}\NormalTok{ certificateChain}\OperatorTok{;}

    \KeywordTok{public} \DataTypeTok{byte}\OperatorTok{[]} \FunctionTok{signOrSeal}\OperatorTok{(}\DataTypeTok{byte}\OperatorTok{[]}\NormalTok{ pdfBytes}\OperatorTok{,} \BuiltInTok{String}\NormalTok{ reason}\OperatorTok{,}
            \BuiltInTok{String}\NormalTok{ location}\OperatorTok{)} \KeywordTok{throws} \BuiltInTok{Exception} \OperatorTok{\{}

\NormalTok{        PDDocument doc }\OperatorTok{=}\NormalTok{ Loader}\OperatorTok{.}\FunctionTok{loadPDF}\OperatorTok{(}\NormalTok{pdfBytes}\OperatorTok{);}

\NormalTok{        PDSignature sig }\OperatorTok{=} \KeywordTok{new} \FunctionTok{PDSignature}\OperatorTok{();}
\NormalTok{        sig}\OperatorTok{.}\FunctionTok{setFilter}\OperatorTok{(}\NormalTok{PDSignature}\OperatorTok{.}\FunctionTok{FILTER\_ADOBE\_PPKLITE}\OperatorTok{);}
\NormalTok{        sig}\OperatorTok{.}\FunctionTok{setSubFilter}\OperatorTok{(}
\NormalTok{                PDSignature}\OperatorTok{.}\FunctionTok{SUBFILTER\_ETSI\_CADES\_DETACHED}\OperatorTok{);}
\NormalTok{        sig}\OperatorTok{.}\FunctionTok{setReason}\OperatorTok{(}\NormalTok{reason}\OperatorTok{);}    
\NormalTok{        sig}\OperatorTok{.}\FunctionTok{setLocation}\OperatorTok{(}\NormalTok{location}\OperatorTok{);}

\NormalTok{        CmsSignatureCreator signer }\OperatorTok{=} \KeywordTok{new} \FunctionTok{CmsSignatureCreator}\OperatorTok{(}
\NormalTok{                privateKey}\OperatorTok{,}\NormalTok{ certificateChain}\OperatorTok{);}

\NormalTok{        doc}\OperatorTok{.}\FunctionTok{addSignature}\OperatorTok{(}\NormalTok{sig}\OperatorTok{,}\NormalTok{ signer}\OperatorTok{);}

        \BuiltInTok{ByteArrayOutputStream}\NormalTok{ baos }\OperatorTok{=} \KeywordTok{new} \BuiltInTok{ByteArrayOutputStream}\OperatorTok{();}
\NormalTok{        doc}\OperatorTok{.}\FunctionTok{saveIncremental}\OperatorTok{(}\NormalTok{baos}\OperatorTok{);}
\NormalTok{        doc}\OperatorTok{.}\FunctionTok{close}\OperatorTok{();}

\NormalTok{        log}\OperatorTok{.}\FunctionTok{info}\OperatorTok{(}\StringTok{"PDF signed/sealed: cert=\{\}, reason=\{\}"}\OperatorTok{,}
\NormalTok{                certificateChain}\OperatorTok{[}\DecValTok{0}\OperatorTok{].}\FunctionTok{getSubjectX500Principal}\OperatorTok{(),}
\NormalTok{                reason}\OperatorTok{);}
        \ControlFlowTok{return}\NormalTok{ baos}\OperatorTok{.}\FunctionTok{toByteArray}\OperatorTok{();}
    \OperatorTok{\}}
\OperatorTok{\}}
\end{Highlighting}
\end{Shaded}

\begin{itemize}
\item
  The reason field typically differs: "Contract approval by Alice Smith"
  for a signature vs "Authenticated by Acme Corp" for a seal
\end{itemize}

\subsection{Detecting Signature vs Seal During
Validation}\label{_detecting_signature_vs_seal_during_validation}

\begin{Shaded}
\begin{Highlighting}[]
\KeywordTok{import} \ImportTok{java}\OperatorTok{.}\ImportTok{security}\OperatorTok{.}\ImportTok{cert}\OperatorTok{.}\ImportTok{X509Certificate}\OperatorTok{;}

\KeywordTok{import} \ImportTok{org}\OperatorTok{.}\ImportTok{bouncycastle}\OperatorTok{.}\ImportTok{asn1}\OperatorTok{.}\ImportTok{ASN1Encodable}\OperatorTok{;}
\KeywordTok{import} \ImportTok{org}\OperatorTok{.}\ImportTok{bouncycastle}\OperatorTok{.}\ImportTok{asn1}\OperatorTok{.}\ImportTok{ASN1ObjectIdentifier}\OperatorTok{;}
\KeywordTok{import} \ImportTok{org}\OperatorTok{.}\ImportTok{bouncycastle}\OperatorTok{.}\ImportTok{asn1}\OperatorTok{.}\ImportTok{ASN1OctetString}\OperatorTok{;}
\KeywordTok{import} \ImportTok{org}\OperatorTok{.}\ImportTok{bouncycastle}\OperatorTok{.}\ImportTok{asn1}\OperatorTok{.}\ImportTok{ASN1Sequence}\OperatorTok{;}

\KeywordTok{import} \ImportTok{lombok}\OperatorTok{.}\ImportTok{extern}\OperatorTok{.}\ImportTok{slf4j}\OperatorTok{.}\ImportTok{Slf4j}\OperatorTok{;}

\CommentTok{/**}
 \CommentTok{*}\NormalTok{ Determines whether a certificate is for electronic signatures}
 \CommentTok{*} \CommentTok{(}\NormalTok{natural person}\CommentTok{)}\NormalTok{ or electronic seals }\CommentTok{(}\NormalTok{legal person}\CommentTok{)}\NormalTok{ by}
 \CommentTok{*}\NormalTok{ examining the QcType statement in QCStatements}\CommentTok{.}
 \CommentTok{*/}
\AttributeTok{@Slf4j}
\KeywordTok{public} \KeywordTok{class}\NormalTok{ SignatureOrSealDetector }\OperatorTok{\{}

    \KeywordTok{private} \DataTypeTok{static} \DataTypeTok{final} \BuiltInTok{String}\NormalTok{ QC\_STATEMENTS\_OID }\OperatorTok{=}
            \StringTok{"1.3.6.1.5.5.7.1.3"}\OperatorTok{;}
    \KeywordTok{private} \DataTypeTok{static} \DataTypeTok{final} \BuiltInTok{String}\NormalTok{ QC\_TYPE\_OID }\OperatorTok{=} \StringTok{"0.4.0.1862.1.6"}\OperatorTok{;}
    \KeywordTok{private} \DataTypeTok{static} \DataTypeTok{final} \BuiltInTok{String}\NormalTok{ ESIGN\_OID }\OperatorTok{=} \StringTok{"0.4.0.1862.1.6.1"}\OperatorTok{;}
    \KeywordTok{private} \DataTypeTok{static} \DataTypeTok{final} \BuiltInTok{String}\NormalTok{ ESEAL\_OID }\OperatorTok{=} \StringTok{"0.4.0.1862.1.6.2"}\OperatorTok{;}

    \KeywordTok{public} \KeywordTok{enum}\NormalTok{ CertType }\OperatorTok{\{}\NormalTok{ SIGNATURE}\OperatorTok{,}\NormalTok{ SEAL}\OperatorTok{,} \BuiltInTok{UNKNOWN} \OperatorTok{\}}

    \KeywordTok{public}\NormalTok{ CertType }\FunctionTok{detect}\OperatorTok{(}\BuiltInTok{X509Certificate}\NormalTok{ cert}\OperatorTok{)} \OperatorTok{\{}
        \DataTypeTok{byte}\OperatorTok{[]}\NormalTok{ extValue }\OperatorTok{=}\NormalTok{ cert}\OperatorTok{.}\FunctionTok{getExtensionValue}\OperatorTok{(}\NormalTok{QC\_STATEMENTS\_OID}\OperatorTok{);}
        \ControlFlowTok{if} \OperatorTok{(}\NormalTok{extValue }\OperatorTok{==} \KeywordTok{null}\OperatorTok{)} \OperatorTok{\{}
            \ControlFlowTok{return}\NormalTok{ CertType}\OperatorTok{.}\FunctionTok{UNKNOWN}\OperatorTok{;}
        \OperatorTok{\}}

        \ControlFlowTok{try} \OperatorTok{\{}
            \DataTypeTok{byte}\OperatorTok{[]}\NormalTok{ octets }\OperatorTok{=}\NormalTok{ ASN1OctetString}\OperatorTok{.}\FunctionTok{getInstance}\OperatorTok{(}\NormalTok{extValue}\OperatorTok{)}
                    \OperatorTok{.}\FunctionTok{getOctets}\OperatorTok{();}
\NormalTok{            ASN1Sequence statements }\OperatorTok{=}
\NormalTok{                    ASN1Sequence}\OperatorTok{.}\FunctionTok{getInstance}\OperatorTok{(}\NormalTok{octets}\OperatorTok{);}

            \ControlFlowTok{for} \OperatorTok{(}\NormalTok{ASN1Encodable statement }\OperatorTok{:}\NormalTok{ statements}\OperatorTok{)} \OperatorTok{\{}
\NormalTok{                ASN1Sequence seq }\OperatorTok{=}
\NormalTok{                        ASN1Sequence}\OperatorTok{.}\FunctionTok{getInstance}\OperatorTok{(}\NormalTok{statement}\OperatorTok{);}
                \BuiltInTok{String}\NormalTok{ oid }\OperatorTok{=}\NormalTok{ ASN1ObjectIdentifier}\OperatorTok{.}\FunctionTok{getInstance}\OperatorTok{(}
\NormalTok{                        seq}\OperatorTok{.}\FunctionTok{getObjectAt}\OperatorTok{(}\DecValTok{0}\OperatorTok{)).}\FunctionTok{getId}\OperatorTok{();}

                \ControlFlowTok{if} \OperatorTok{(!}\NormalTok{QC\_TYPE\_OID}\OperatorTok{.}\FunctionTok{equals}\OperatorTok{(}\NormalTok{oid}\OperatorTok{)} \OperatorTok{||}\NormalTok{ seq}\OperatorTok{.}\FunctionTok{size}\OperatorTok{()} \OperatorTok{\textless{}} \DecValTok{2}\OperatorTok{)} \OperatorTok{\{}
                    \ControlFlowTok{continue}\OperatorTok{;}
                \OperatorTok{\}}

\NormalTok{                ASN1Sequence typeSeq }\OperatorTok{=}\NormalTok{ ASN1Sequence}\OperatorTok{.}\FunctionTok{getInstance}\OperatorTok{(}
\NormalTok{                        seq}\OperatorTok{.}\FunctionTok{getObjectAt}\OperatorTok{(}\DecValTok{1}\OperatorTok{));}
                \ControlFlowTok{for} \OperatorTok{(}\NormalTok{ASN1Encodable entry }\OperatorTok{:}\NormalTok{ typeSeq}\OperatorTok{)} \OperatorTok{\{}
                    \BuiltInTok{String}\NormalTok{ typeOid }\OperatorTok{=}\NormalTok{ ASN1ObjectIdentifier}
                            \OperatorTok{.}\FunctionTok{getInstance}\OperatorTok{(}\NormalTok{entry}\OperatorTok{).}\FunctionTok{getId}\OperatorTok{();}
                    \ControlFlowTok{if} \OperatorTok{(}\NormalTok{ESIGN\_OID}\OperatorTok{.}\FunctionTok{equals}\OperatorTok{(}\NormalTok{typeOid}\OperatorTok{))} \OperatorTok{\{}
\NormalTok{                        log}\OperatorTok{.}\FunctionTok{info}\OperatorTok{(}\StringTok{"Certificate type: SIGNATURE"}\OperatorTok{);}
                        \ControlFlowTok{return}\NormalTok{ CertType}\OperatorTok{.}\FunctionTok{SIGNATURE}\OperatorTok{;}
                    \OperatorTok{\}}
                    \ControlFlowTok{if} \OperatorTok{(}\NormalTok{ESEAL\_OID}\OperatorTok{.}\FunctionTok{equals}\OperatorTok{(}\NormalTok{typeOid}\OperatorTok{))} \OperatorTok{\{}
\NormalTok{                        log}\OperatorTok{.}\FunctionTok{info}\OperatorTok{(}\StringTok{"Certificate type: SEAL"}\OperatorTok{);}
                        \ControlFlowTok{return}\NormalTok{ CertType}\OperatorTok{.}\FunctionTok{SEAL}\OperatorTok{;}
                    \OperatorTok{\}}
                \OperatorTok{\}}
            \OperatorTok{\}}
        \OperatorTok{\}} \ControlFlowTok{catch} \OperatorTok{(}\BuiltInTok{Exception}\NormalTok{ e}\OperatorTok{)} \OperatorTok{\{}
\NormalTok{            log}\OperatorTok{.}\FunctionTok{warn}\OperatorTok{(}\StringTok{"Failed to parse QCStatements"}\OperatorTok{,}\NormalTok{ e}\OperatorTok{);}
        \OperatorTok{\}}
        \ControlFlowTok{return}\NormalTok{ CertType}\OperatorTok{.}\FunctionTok{UNKNOWN}\OperatorTok{;}
    \OperatorTok{\}}
\OperatorTok{\}}
\end{Highlighting}
\end{Shaded}

If the QcType statement is absent (some older certificates predate ETSI
EN 319 412-5), fall back to the certificate policy OID:
\texttt{0.4.0.194112.1.2} / \texttt{…\hspace{0pt}1.1} indicate signature
certificates, \texttt{0.4.0.194112.1.3} / \texttt{…\hspace{0pt}1.4}
indicate seal certificates. As a last resort, check the subject DN --- a
certificate with only an organization name (no personal CN) is likely a
seal, though this heuristic is not definitive.

\section{Use Cases: When to Sign, When to
Seal}\label{_use_cases_when_to_sign_when_to_seal}

\subsection{Contracts and Agreements:
Signature}\label{_contracts_and_agreements_signature}

Any document requiring personal commitment --- contracts, NDAs, medical
prescriptions, regulatory filings, court submissions --- needs a
\textbf{signature}. The signer is the specific human being who agrees.
Regulation (EU) No 910/2014 (eIDAS) Article 27 requires member states
that mandate handwritten signatures for public services to accept QES.

\subsection{Invoices and Automated Documents:
Seal}\label{_invoices_and_automated_documents_seal}

Documents needing organizational authentication without personal
approval --- invoices, bank statements, payslips, automated
reports --- are candidates for \textbf{seals}. EU VAT Directive
2006/112/EC (Article 233) recognizes electronic seals as meeting invoice
authenticity and integrity requirements.

\subsection{Hybrid Workflows}\label{_hybrid_workflows}

Many real-world workflows combine both:

\begin{enumerate}
\def\labelenumi{\arabic{enumi}.}
\item
  An organization \textbf{seals} a document template (e.g., a bank seals
  a loan agreement)
\item
  A natural person \textbf{signs} to express personal agreement (e.g.,
  the customer signs the loan)
\item
  The organization's representative \textbf{signs} to accept (a natural
  person acting as agent)
\end{enumerate}

The seal and signatures coexist as separate signature dictionaries in
the same PDF, each with its own byte range and CMS structure. Validators
should display both, clearly distinguishing "Sealed by Acme Corp" from
"Signed by Alice Smith."

\section{Server-Side Seal Creation and Batch
Sealing}\label{_server_side_seal_creation_and_batch_sealing}

Because seals represent organizations rather than individuals, they are
ideal for \textbf{automated, server-side creation}. No human needs to be
present at the moment of sealing.

\subsection{The Architecture}\label{_the_architecture}

\begin{enumerate}
\def\labelenumi{\arabic{enumi}.}
\item
  The organization holds a \textbf{qualified seal certificate} with its
  private key on an \textbf{HSM} (certified as a Qualified Seal Creation
  Device)
\item
  An application server generates documents (invoices, reports,
  certificates)
\item
  The server calls the HSM to sign each document --- no human
  interaction per document
\item
  The HSM authenticates the application server (mutual TLS, API key),
  not a person
\end{enumerate}

This is the critical difference from signatures: for QES, the signatory
must authenticate \emph{per signing operation} (sole control). For
seals, the organization's security policies ensure only authorized
systems can trigger sealing. This makes batch processing feasible.

At high volumes, watch for two bottlenecks:

\begin{enumerate}
\def\labelenumi{\arabic{enumi}.}
\item
  \textbf{HSM throughput}: a single signing operation takes 1-10ms
  (RSA-2048 is slower than ECDSA P-256). At 50,000 documents per minute,
  you may need multiple HSM partitions or load balancing.
\item
  \textbf{TSA throughput}: if adding qualified timestamps (PAdES-B-T),
  the TSA becomes the constraint. Most commercial TSAs handle hundreds
  of requests per second, but verify against your volume requirements.
\end{enumerate}

\subsection{Cost Implications}\label{_cost_implications}

\textbf{Signature certificates} are per-person, per-year: 50-300 EUR per
natural person. For 500 employees, the annual cost is substantial.

\textbf{Seal certificates} are per-organization, per-year. A single seal
certificate can produce millions of seals at a fixed cost. The
per-document cost approaches zero at scale.

This asymmetry drives the architectural decision: seals for volume,
signatures for commitment.

\section{Common Pitfalls}\label{_common_pitfalls}

\subsection{Displaying Seals as
Signatures}\label{_displaying_seals_as_signatures}

If your validation UI says "Signed by Acme Corp" when it encounters a
seal, you are misleading the user. The correct display is "Sealed by" or
"Authenticated by" --- language that communicates origin without
implying personal intent.

\subsection{Using Signature Certificates for
Automation}\label{_using_signature_certificates_for_automation}

If your batch pipeline signs 10,000 documents per hour with Alice's
personal qualified certificate, you have a problem. Alice is not
exercising sole control over each operation, and may not even be aware
of most of them. Use a seal certificate for automation and reserve
signature certificates for workflows requiring individual human
approval.

\subsection{Ignoring the Certificate Type in
Code}\label{_ignoring_the_certificate_type_in_code}

Because the signing code is identical, it is easy to accidentally use a
seal certificate for a signature workflow or vice versa. Your
application must validate the certificate \emph{before} signing: check
QcType, reject mismatches, and never rely on the user to "use the right
certificate."

Some older qualified certificates lack a QcType statement entirely. For
these, you must fall back to the EUTL service entry's qualifiers or the
certificate policy OID. This is yet another reason to use the EU DSS
library (Chapter 23), which handles these edge cases.

\section{Key Takeaways}\label{_key_takeaways_9}

\begin{itemize}
\item
  \textbf{Signatures} are for natural persons and express
  \textbf{intent} --- Article 25(2) gives QES the legal equivalence of a
  handwritten signature
\item
  \textbf{Seals} are for legal persons and prove \textbf{origin and
  integrity} --- Article 35(2) gives qualified seals a presumption of
  integrity and correct origin, not equivalence to a signature
\item
  The cryptographic mechanism is \textbf{identical} --- same CMS, same
  PAdES, same QSCD; the only difference is the certificate (QcType esign
  vs eseal, certificate policy QCP-n vs QCP-l)
\item
  Use \textbf{signatures} for contracts, agreements,
  approvals --- anything requiring personal human commitment
\item
  Use \textbf{seals} for invoices, automated reports, official
  publications --- organizational authentication without personal
  approval
\item
  \textbf{Server-side batch sealing} is the dominant seal use case:
  HSM-backed, no per-operation human auth, millions of documents at
  fixed cost
\item
  \textbf{Enforce the certificate type in code} --- do not allow seal
  certs for signature workflows or vice versa
\item
  Display seals correctly: "Sealed by" not "Signed by" --- the legal
  distinction matters
\item
  Regulation (EU) 2024/1183 (eIDAS 2.0) focuses on personal QES via the
  EUDI wallet; seals continue on HSM infrastructure, with rising volume
  from QEAAs
\end{itemize}

\chapter{Remote Signing}\label{_remote_signing}

For the first two decades of digital signatures, the assumption was
simple: the private key lives with the signer. A smart card in a wallet,
a PKCS\#12 file on a laptop, a USB token plugged into a desktop. The
signer has physical custody of their key material, they authenticate
locally (enter a PIN, type a password), and the signing happens on their
device.

That model is dead --- or at least, it is dying.

In the remote signing model, the private key never leaves a server-side
Hardware Security Module (HSM). The signer authenticates remotely ---
via OAuth 2.0, via SMS OTP, via biometric verification on their phone
--- and the server performs the cryptographic signing operation on their
behalf. The signer never touches the key, never downloads it, never even
knows the key's technical parameters. They authorize a signing
operation, and the server produces the signature.

This chapter covers the architecture of remote signing: why it exists,
how the CSC API standardizes it, how HSMs and Signature Activation
Modules work together, what the SCAL levels mean for sole control, and
how to integrate with remote signing services in practice. If your
application needs to produce qualified electronic signatures (QES) at
scale --- for thousands of users who do not have smart cards and do not
want USB tokens --- remote signing is how you get there.

\section{Why Remote Signing Won}\label{_why_remote_signing_won}

\subsection{The Smart Card Problem}\label{_the_smart_card_problem}

The traditional QES model required every signer to have:

\begin{enumerate}
\def\labelenumi{\arabic{enumi}.}
\item
  A qualified certificate issued by a QTSP
\item
  A Qualified Signature Creation Device (QSCD) --- a certified smart
  card or USB token
\item
  A card reader connected to their computer
\item
  Driver software for the card reader and the QSCD
\item
  Middleware that bridges the QSCD to the signing application (typically
  PKCS\#11 or Microsoft CNG)
\end{enumerate}

This worked in tightly controlled enterprise environments --- banks,
government agencies, regulated industries where IT departments could
provision hardware to every employee. It failed completely for consumer
scenarios, mobile scenarios, and any context where "plug in a USB device
and install a driver" is not feasible.

The practical failure modes were numerous:

\begin{itemize}
\item
  \textbf{Driver hell}: Smart card drivers are notoriously
  platform-dependent. A card that works on Windows 10 may not work on
  Windows 11 without a driver update. macOS support is inconsistent.
  Linux support depends on the specific combination of card, reader, and
  PKCS\#11 module.
\item
  \textbf{Mobile impossibility}: Smartphones have no USB-A ports, no
  smart card readers, and limited support for external cryptographic
  hardware. The NFC-based smart card reading that some phones support is
  slow and unreliable.
\item
  \textbf{Certificate lifecycle}: Qualified certificates typically
  expire after 1-3 years. Renewal requires physically visiting a
  Registration Authority (for identity verification) and provisioning a
  new certificate onto the QSCD. Users forget, certificates expire,
  signing fails.
\item
  \textbf{Cost}: A qualified certificate costs 50-300 EUR per year. A
  smart card costs 20-50 EUR. A card reader costs 15-40 EUR. For an
  organization with 10,000 employees, the hardware cost alone is
  significant --- before you account for help desk tickets from users
  who cannot get their cards to work.
\end{itemize}

\subsection{The Remote Alternative}\label{_the_remote_alternative}

Remote signing eliminates every one of these problems:

\begin{itemize}
\item
  \textbf{No hardware}: The signer needs only a web browser or a mobile
  app. No drivers, no readers, no tokens.
\item
  \textbf{Mobile-native}: Authentication happens through the phone
  itself --- push notifications, biometric confirmation, SMS OTP. The
  signing experience on a smartphone is identical to the desktop
  experience.
\item
  \textbf{Managed lifecycle}: The QTSP manages the key lifecycle
  entirely. Certificate issuance, renewal, and revocation happen
  server-side. The signer never handles key material.
\item
  \textbf{Lower per-user cost}: While the QTSP charges per signature or
  per subscription, the elimination of hardware provisioning and support
  often makes remote signing cheaper at scale.
\end{itemize}

Remote signing does not lower the \emph{security} bar --- it moves the
security boundary. Instead of trusting the signer's device to protect
the key (which, for a password-protected PKCS\#12 on a laptop, is a
dubious proposition), you trust the QTSP's HSM infrastructure and
authentication system. For most threat models, a professionally managed
HSM with multi-factor authentication provides stronger key protection
than a smart card with a four-digit PIN that the user wrote on a sticky
note.

\subsection{eIDAS and Remote QSCDs}\label{_eidas_and_remote_qscds}

The Regulation (EU) No 910/2014 (eIDAS) does not require physical
possession of a QSCD --- it requires sole control (Article 26(c)) and
that the device meet the requirements of Annex II. An HSM operated by a
QTSP, where the signer's key partition is accessible only through strong
multi-factor authentication tied to the specific signing operation,
satisfies both requirements.

The European Committee for Standardization (CEN) formalized this in
\textbf{EN 419 241}, which defines requirements for trustworthy systems
supporting server signing:

\begin{itemize}
\item
  \textbf{EN 419 241-1}: General requirements for server signing
\item
  \textbf{EN 419 241-2}: Protection profile for QSCDs for server signing
\end{itemize}

These standards, combined with the SCAL framework (discussed below),
provide the certification path for remote QSCDs. By 2024, remote signing
had become the dominant model for QES in practice --- more qualified
signatures are produced via remote HSMs than via local smart cards.

\section{The Hash-Signing Model}\label{_the_hash_signing_model}

\subsection{Why Send the Hash, Not the
Document}\label{_why_send_the_hash_not_the_document}

Remote signing uses a \textbf{hash-signing model}: the client computes
the hash of the data to be signed and sends only the hash to the server.
The server applies the private key to the hash (RSA-PSS, ECDSA, etc.)
and returns the raw signature value. The client then embeds the
signature value into the CMS structure and the PDF.

This design has three critical advantages:

\begin{enumerate}
\def\labelenumi{\arabic{enumi}.}
\item
  \textbf{Privacy}: The document content never leaves the client. The
  server sees only a 32-byte (SHA-256) or 64-byte (SHA-512) hash. It
  cannot reconstruct the document from the hash. For sensitive documents
  --- medical records, legal contracts, financial statements --- this is
  essential.
\item
  \textbf{Bandwidth}: Sending a 32-byte hash is dramatically faster than
  sending a 50 MB PDF over a network connection. For mobile clients on
  cellular networks, the difference matters.
\item
  \textbf{Separation of concerns}: The client handles PDF parsing, byte
  range computation, CMS structure building, and timestamp embedding.
  The server handles key management, access control, and the raw
  cryptographic operation. Neither needs to understand the other's
  domain.
\end{enumerate}

\subsection{The Flow Step by Step}\label{_the_flow_step_by_step}

\begin{figure}
\centering
\pandocbounded{\includegraphics[keepaspectratio,alt={seq remote signing}]{diagrams/pdf/seq-remote-signing.pdf}}
\caption{Remote signing flow --- client and server responsibilities}
\end{figure}

Steps 1-3 and 11-14 happen entirely on the client. Only steps 4, 6, and
the server-side steps 5-10 involve the remote server. The document never
crosses the network.

The "hash" sent to the server in step 6 is not the hash of the document
bytes. It is the hash of the \textbf{DER-encoded signed attributes} of
the CMS SignerInfo. The signed attributes include the
\texttt{messageDigest} attribute, which itself contains the hash of the
document bytes. This two-level hashing is standard CMS behavior (defined
in RFC 5652 --- Cryptographic Message Syntax (CMS) Section 5.4) and is
easy to get wrong.

If you hash the document bytes and send that to the remote server, the
resulting signature will not verify --- because CMS validators expect
the signature to be over the signed attributes, not the raw data.

\section{The CSC API (Cloud Signature
Consortium)}\label{_the_csc_api_cloud_signature_consortium}

\subsection{What the CSC Is}\label{_what_the_csc_is}

The \textbf{Cloud Signature Consortium (CSC)} is an industry body that
developed an open standard API for remote signing. Before the CSC API,
every remote signing provider had its own proprietary API --- different
authentication flows, different request/response formats, different
error codes. Integrating with Swisscom AIS required completely different
code than integrating with DigiCert, which required completely different
code than integrating with InfoCert.

The CSC API v2.0 (the current version, published in 2019 and updated
since) standardizes:

\begin{itemize}
\item
  \textbf{Authentication}: OAuth 2.0 with specific grant types
\item
  \textbf{Credential discovery}: listing available signing credentials
  for a user
\item
  \textbf{Signature activation}: the mechanism by which the signer
  authorizes a specific signing operation
\item
  \textbf{Hash signing}: the actual cryptographic operation
\end{itemize}

The API is RESTful, uses JSON for request and response bodies, and is
designed to be consumed by web and mobile clients.

\subsection{API Endpoints}\label{_api_endpoints}

The CSC API defines the following core endpoints:

\begin{longtable}[]{@{}
  >{\raggedright\arraybackslash}p{(\linewidth - 4\tabcolsep) * \real{0.3000}}
  >{\raggedright\arraybackslash}p{(\linewidth - 4\tabcolsep) * \real{0.2000}}
  >{\raggedright\arraybackslash}p{(\linewidth - 4\tabcolsep) * \real{0.5000}}@{}}
\caption{CSC API v2.0 Core Endpoints}\tabularnewline
\toprule\noalign{}
\begin{minipage}[b]{\linewidth}\raggedright
Endpoint
\end{minipage} & \begin{minipage}[b]{\linewidth}\raggedright
Method
\end{minipage} & \begin{minipage}[b]{\linewidth}\raggedright
Purpose
\end{minipage} \\
\midrule\noalign{}
\endfirsthead
\toprule\noalign{}
\begin{minipage}[b]{\linewidth}\raggedright
Endpoint
\end{minipage} & \begin{minipage}[b]{\linewidth}\raggedright
Method
\end{minipage} & \begin{minipage}[b]{\linewidth}\raggedright
Purpose
\end{minipage} \\
\midrule\noalign{}
\endhead
\bottomrule\noalign{}
\endlastfoot
\texttt{/auth/login} & POST & Authenticate the user with the remote
signing service. Returns an OAuth 2.0 access token. \\
\texttt{/auth/revoke} & POST & Revoke an active access token. \\
\texttt{/credentials/list} & POST & List the signing credentials
available to the authenticated user. Returns credential IDs, certificate
chains, key algorithms, and SCAL levels. \\
\texttt{/credentials/info} & POST & Get detailed information about a
specific credential, including the certificate chain and supported
algorithms. \\
\texttt{/credentials/authorize} & POST & Authorize a specific credential
for signing. For SCAL2 credentials, this triggers explicit user
authentication (OTP, biometric) and returns a Signature Activation Data
(SAD) token. \\
\texttt{/credentials/sendOTP} & POST & Send a one-time password to the
user (SMS, email) as part of SCAL2 authorization. \\
\texttt{/signatures/signHash} & POST & Sign one or more hashes using a
previously authorized credential. Returns the raw signature values. \\
\texttt{/signatures/timestamp} & POST & Request an RFC 3161 timestamp.
Some providers expose timestamping through the CSC API as well. \\
\end{longtable}

\subsection{Authentication: OAuth 2.0}\label{_authentication_oauth_2_0}

The CSC API uses OAuth 2.0 for authentication. The specific grant type
depends on the deployment scenario:

\begin{itemize}
\item
  \textbf{Authorization Code Grant}: For web applications where the user
  authenticates in a browser. The signing application redirects the user
  to the QTSP's authorization server, the user logs in, and the
  application receives an authorization code that it exchanges for an
  access token.
\item
  \textbf{Client Credentials Grant}: For server-to-server scenarios
  where a system (not a human) needs to sign. The application
  authenticates with its own client ID and secret.
\item
  \textbf{Resource Owner Password Grant}: For scenarios where the
  application collects the user's credentials directly. Discouraged by
  OAuth 2.0 best practices but still supported by some providers.
\end{itemize}

\begin{Shaded}
\begin{Highlighting}[]
\CommentTok{/**}
 \CommentTok{*}\NormalTok{ Authenticates with a CSC}\CommentTok{{-}}\NormalTok{compliant remote signing service}\CommentTok{.}
 \CommentTok{*}\NormalTok{ Returns an OAuth }\CommentTok{2.0}\NormalTok{ access token for subsequent API calls}\CommentTok{.}
 \CommentTok{*/}
\KeywordTok{public} \BuiltInTok{String} \FunctionTok{authenticate}\OperatorTok{(}\BuiltInTok{String}\NormalTok{ cscBaseUrl}\OperatorTok{,} \BuiltInTok{String}\NormalTok{ username}\OperatorTok{,}
        \BuiltInTok{String}\NormalTok{ password}\OperatorTok{)} \KeywordTok{throws} \BuiltInTok{IOException} \OperatorTok{\{}

    \BuiltInTok{String}\NormalTok{ loginUrl }\OperatorTok{=}\NormalTok{ cscBaseUrl }\OperatorTok{+} \StringTok{"/auth/login"}\OperatorTok{;}
    \BuiltInTok{String}\NormalTok{ requestBody }\OperatorTok{=} \StringTok{"""}
\StringTok{            \{}
\StringTok{                "remoteUser": "\%s",}
\StringTok{                "remoteSecret": "\%s"}
\StringTok{            \}}
\StringTok{            """}\OperatorTok{.}\FunctionTok{formatted}\OperatorTok{(}\NormalTok{username}\OperatorTok{,}\NormalTok{ password}\OperatorTok{);}

    \BuiltInTok{HttpURLConnection}\NormalTok{ conn }\OperatorTok{=} \OperatorTok{(}\BuiltInTok{HttpURLConnection}\OperatorTok{)}
            \KeywordTok{new} \BuiltInTok{URL}\OperatorTok{(}\NormalTok{loginUrl}\OperatorTok{).}\FunctionTok{openConnection}\OperatorTok{();}
\NormalTok{    conn}\OperatorTok{.}\FunctionTok{setDoOutput}\OperatorTok{(}\KeywordTok{true}\OperatorTok{);}
\NormalTok{    conn}\OperatorTok{.}\FunctionTok{setRequestMethod}\OperatorTok{(}\StringTok{"POST"}\OperatorTok{);}
\NormalTok{    conn}\OperatorTok{.}\FunctionTok{setRequestProperty}\OperatorTok{(}\StringTok{"Content{-}Type"}\OperatorTok{,} \StringTok{"application/json"}\OperatorTok{);}
\NormalTok{    conn}\OperatorTok{.}\FunctionTok{setConnectTimeout}\OperatorTok{(}\DecValTok{10000}\OperatorTok{);}
\NormalTok{    conn}\OperatorTok{.}\FunctionTok{setReadTimeout}\OperatorTok{(}\DecValTok{30000}\OperatorTok{);}

    \ControlFlowTok{try} \OperatorTok{(}\BuiltInTok{OutputStream}\NormalTok{ out }\OperatorTok{=}\NormalTok{ conn}\OperatorTok{.}\FunctionTok{getOutputStream}\OperatorTok{())} \OperatorTok{\{}
\NormalTok{        out}\OperatorTok{.}\FunctionTok{write}\OperatorTok{(}\NormalTok{requestBody}\OperatorTok{.}\FunctionTok{getBytes}\OperatorTok{(}\NormalTok{StandardCharsets}\OperatorTok{.}\FunctionTok{UTF\_8}\OperatorTok{));}
    \OperatorTok{\}}

    \ControlFlowTok{if} \OperatorTok{(}\NormalTok{conn}\OperatorTok{.}\FunctionTok{getResponseCode}\OperatorTok{()} \OperatorTok{!=} \DecValTok{200}\OperatorTok{)} \OperatorTok{\{}
        \ControlFlowTok{throw} \KeywordTok{new} \BuiltInTok{IOException}\OperatorTok{(}\StringTok{"CSC auth/login failed: HTTP "} \OperatorTok{+}
\NormalTok{                conn}\OperatorTok{.}\FunctionTok{getResponseCode}\OperatorTok{());}
    \OperatorTok{\}}

    \CommentTok{// Parse JSON response — extract access\_token}
    \BuiltInTok{String}\NormalTok{ responseJson}\OperatorTok{;}
    \ControlFlowTok{try} \OperatorTok{(}\BuiltInTok{InputStream}\NormalTok{ in }\OperatorTok{=}\NormalTok{ conn}\OperatorTok{.}\FunctionTok{getInputStream}\OperatorTok{())} \OperatorTok{\{}
\NormalTok{        responseJson }\OperatorTok{=} \KeywordTok{new} \BuiltInTok{String}\OperatorTok{(}\NormalTok{in}\OperatorTok{.}\FunctionTok{readAllBytes}\OperatorTok{(),}\NormalTok{ StandardCharsets}\OperatorTok{.}\FunctionTok{UTF\_8}\OperatorTok{);}
    \OperatorTok{\}}

    \CommentTok{// In production, use a proper JSON parser (Jackson, Gson)}
    \CommentTok{// This is simplified for illustration}
\NormalTok{    log}\OperatorTok{.}\FunctionTok{info}\OperatorTok{(}\StringTok{"CSC authentication successful for user: \{\}"}\OperatorTok{,}\NormalTok{ username}\OperatorTok{);}
    \ControlFlowTok{return} \FunctionTok{extractJsonField}\OperatorTok{(}\NormalTok{responseJson}\OperatorTok{,} \StringTok{"access\_token"}\OperatorTok{);}
\OperatorTok{\}}
\end{Highlighting}
\end{Shaded}

\subsection{Credential Discovery}\label{_credential_discovery}

Once authenticated, the client can discover what signing credentials are
available:

\begin{Shaded}
\begin{Highlighting}[]
\CommentTok{/**}
 \CommentTok{*}\NormalTok{ Lists the signing credentials available to the authenticated user}\CommentTok{.}
 \CommentTok{*}\NormalTok{ Returns credential IDs that can be used for signing}\CommentTok{.}
 \CommentTok{*/}
\KeywordTok{public} \BuiltInTok{List}\OperatorTok{\textless{}}\NormalTok{CscCredential}\OperatorTok{\textgreater{}} \FunctionTok{listCredentials}\OperatorTok{(}\BuiltInTok{String}\NormalTok{ cscBaseUrl}\OperatorTok{,}
        \BuiltInTok{String}\NormalTok{ accessToken}\OperatorTok{)} \KeywordTok{throws} \BuiltInTok{IOException} \OperatorTok{\{}

    \BuiltInTok{String}\NormalTok{ url }\OperatorTok{=}\NormalTok{ cscBaseUrl }\OperatorTok{+} \StringTok{"/credentials/list"}\OperatorTok{;}
    \BuiltInTok{String}\NormalTok{ requestBody }\OperatorTok{=} \StringTok{"""}
\StringTok{            \{}
\StringTok{                "maxResults": 100}
\StringTok{            \}}
\StringTok{            """}\OperatorTok{;}

    \BuiltInTok{HttpURLConnection}\NormalTok{ conn }\OperatorTok{=} \OperatorTok{(}\BuiltInTok{HttpURLConnection}\OperatorTok{)}
            \KeywordTok{new} \BuiltInTok{URL}\OperatorTok{(}\NormalTok{url}\OperatorTok{).}\FunctionTok{openConnection}\OperatorTok{();}
\NormalTok{    conn}\OperatorTok{.}\FunctionTok{setDoOutput}\OperatorTok{(}\KeywordTok{true}\OperatorTok{);}
\NormalTok{    conn}\OperatorTok{.}\FunctionTok{setRequestMethod}\OperatorTok{(}\StringTok{"POST"}\OperatorTok{);}
\NormalTok{    conn}\OperatorTok{.}\FunctionTok{setRequestProperty}\OperatorTok{(}\StringTok{"Content{-}Type"}\OperatorTok{,} \StringTok{"application/json"}\OperatorTok{);}
\NormalTok{    conn}\OperatorTok{.}\FunctionTok{setRequestProperty}\OperatorTok{(}\StringTok{"Authorization"}\OperatorTok{,} \StringTok{"Bearer "} \OperatorTok{+}\NormalTok{ accessToken}\OperatorTok{);}

    \ControlFlowTok{try} \OperatorTok{(}\BuiltInTok{OutputStream}\NormalTok{ out }\OperatorTok{=}\NormalTok{ conn}\OperatorTok{.}\FunctionTok{getOutputStream}\OperatorTok{())} \OperatorTok{\{}
\NormalTok{        out}\OperatorTok{.}\FunctionTok{write}\OperatorTok{(}\NormalTok{requestBody}\OperatorTok{.}\FunctionTok{getBytes}\OperatorTok{(}\NormalTok{StandardCharsets}\OperatorTok{.}\FunctionTok{UTF\_8}\OperatorTok{));}
    \OperatorTok{\}}

    \BuiltInTok{String}\NormalTok{ responseJson}\OperatorTok{;}
    \ControlFlowTok{try} \OperatorTok{(}\BuiltInTok{InputStream}\NormalTok{ in }\OperatorTok{=}\NormalTok{ conn}\OperatorTok{.}\FunctionTok{getInputStream}\OperatorTok{())} \OperatorTok{\{}
\NormalTok{        responseJson }\OperatorTok{=} \KeywordTok{new} \BuiltInTok{String}\OperatorTok{(}\NormalTok{in}\OperatorTok{.}\FunctionTok{readAllBytes}\OperatorTok{(),}\NormalTok{ StandardCharsets}\OperatorTok{.}\FunctionTok{UTF\_8}\OperatorTok{);}
    \OperatorTok{\}}

    \CommentTok{// Response contains:}
    \CommentTok{// \{}
    \CommentTok{//   "credentialIDs": ["cred{-}001", "cred{-}002"],}
    \CommentTok{//   "credentialInfos": [}
    \CommentTok{//     \{}
    \CommentTok{//       "credentialID": "cred{-}001",}
    \CommentTok{//       "key": \{ "algo": ["1.2.840.113549.1.1.11"], "len": 2048 \},}
    \CommentTok{//       "cert": \{ "certificates": ["MIIC...", "MIID..."] \},}
    \CommentTok{//       "SCAL": "2",}
    \CommentTok{//       "description": "Qualified certificate for John Doe"}
    \CommentTok{//     \}}
    \CommentTok{//   ]}
    \CommentTok{// \}}

\NormalTok{    log}\OperatorTok{.}\FunctionTok{info}\OperatorTok{(}\StringTok{"CSC credentials/list returned credentials for user"}\OperatorTok{);}
    \ControlFlowTok{return} \FunctionTok{parseCscCredentials}\OperatorTok{(}\NormalTok{responseJson}\OperatorTok{);}
\OperatorTok{\}}
\end{Highlighting}
\end{Shaded}

The credential response includes the full certificate chain
(Base64-encoded), which the client needs to build the CMS structure
\emph{before} requesting the hash signature. This is a critical detail:
the client must know the signer's certificate at CMS construction time
(for the ESS-signing-certificate-v2 attribute), not just at signing
time.

\subsection{Credential Authorization and
SAD}\label{_credential_authorization_and_sad}

For SCAL2 credentials (see below), the client must explicitly authorize
the credential before signing. This triggers an additional
authentication step --- the QTSP sends an OTP to the signer, or the
signer confirms via a mobile app, or biometric verification occurs.

The result of successful authorization is a \textbf{Signature Activation
Data (SAD)} token --- an opaque string that the client presents with the
signing request to prove that the signer authorized this specific
operation.

\begin{Shaded}
\begin{Highlighting}[]
\CommentTok{/**}
 \CommentTok{*}\NormalTok{ Authorizes a credential for signing }\CommentTok{(}\NormalTok{SCAL2}\CommentTok{).}
 \CommentTok{*}\NormalTok{ Triggers user authentication }\CommentTok{(}\NormalTok{OTP}\CommentTok{,}\NormalTok{ biometric}\CommentTok{,}\NormalTok{ etc}\CommentTok{.).}
 \CommentTok{*}\NormalTok{ Returns a SAD token valid for a limited time and number of signatures}\CommentTok{.}
 \CommentTok{*/}
\KeywordTok{public} \BuiltInTok{String} \FunctionTok{authorizeCredential}\OperatorTok{(}\BuiltInTok{String}\NormalTok{ cscBaseUrl}\OperatorTok{,} \BuiltInTok{String}\NormalTok{ accessToken}\OperatorTok{,}
        \BuiltInTok{String}\NormalTok{ credentialId}\OperatorTok{,} \DataTypeTok{int}\NormalTok{ numSignatures}\OperatorTok{)} \KeywordTok{throws} \BuiltInTok{IOException} \OperatorTok{\{}

    \CommentTok{// Step 1: Trigger OTP delivery}
    \BuiltInTok{String}\NormalTok{ sendOtpUrl }\OperatorTok{=}\NormalTok{ cscBaseUrl }\OperatorTok{+} \StringTok{"/credentials/sendOTP"}\OperatorTok{;}
    \BuiltInTok{String}\NormalTok{ otpRequestBody }\OperatorTok{=} \StringTok{"""}
\StringTok{            \{}
\StringTok{                "credentialID": "\%s"}
\StringTok{            \}}
\StringTok{            """}\OperatorTok{.}\FunctionTok{formatted}\OperatorTok{(}\NormalTok{credentialId}\OperatorTok{);}

    \BuiltInTok{HttpURLConnection}\NormalTok{ otpConn }\OperatorTok{=} \OperatorTok{(}\BuiltInTok{HttpURLConnection}\OperatorTok{)}
            \KeywordTok{new} \BuiltInTok{URL}\OperatorTok{(}\NormalTok{sendOtpUrl}\OperatorTok{).}\FunctionTok{openConnection}\OperatorTok{();}
\NormalTok{    otpConn}\OperatorTok{.}\FunctionTok{setDoOutput}\OperatorTok{(}\KeywordTok{true}\OperatorTok{);}
\NormalTok{    otpConn}\OperatorTok{.}\FunctionTok{setRequestMethod}\OperatorTok{(}\StringTok{"POST"}\OperatorTok{);}
\NormalTok{    otpConn}\OperatorTok{.}\FunctionTok{setRequestProperty}\OperatorTok{(}\StringTok{"Content{-}Type"}\OperatorTok{,} \StringTok{"application/json"}\OperatorTok{);}
\NormalTok{    otpConn}\OperatorTok{.}\FunctionTok{setRequestProperty}\OperatorTok{(}\StringTok{"Authorization"}\OperatorTok{,} \StringTok{"Bearer "} \OperatorTok{+}\NormalTok{ accessToken}\OperatorTok{);}
    \ControlFlowTok{try} \OperatorTok{(}\BuiltInTok{OutputStream}\NormalTok{ out }\OperatorTok{=}\NormalTok{ otpConn}\OperatorTok{.}\FunctionTok{getOutputStream}\OperatorTok{())} \OperatorTok{\{}
\NormalTok{        out}\OperatorTok{.}\FunctionTok{write}\OperatorTok{(}\NormalTok{otpRequestBody}\OperatorTok{.}\FunctionTok{getBytes}\OperatorTok{(}\NormalTok{StandardCharsets}\OperatorTok{.}\FunctionTok{UTF\_8}\OperatorTok{));}
    \OperatorTok{\}}

\NormalTok{    log}\OperatorTok{.}\FunctionTok{info}\OperatorTok{(}\StringTok{"OTP sent for credential: \{\}"}\OperatorTok{,}\NormalTok{ credentialId}\OperatorTok{);}

    \CommentTok{// Step 2: Collect OTP from user (your UI must handle this)}
    \BuiltInTok{String}\NormalTok{ otp }\OperatorTok{=} \FunctionTok{collectOtpFromUser}\OperatorTok{();}

    \CommentTok{// Step 3: Authorize the credential with the OTP}
    \BuiltInTok{String}\NormalTok{ authUrl }\OperatorTok{=}\NormalTok{ cscBaseUrl }\OperatorTok{+} \StringTok{"/credentials/authorize"}\OperatorTok{;}
    \BuiltInTok{String}\NormalTok{ authRequestBody }\OperatorTok{=} \StringTok{"""}
\StringTok{            \{}
\StringTok{                "credentialID": "\%s",}
\StringTok{                "numSignatures": \%d,}
\StringTok{                "OTP": "\%s"}
\StringTok{            \}}
\StringTok{            """}\OperatorTok{.}\FunctionTok{formatted}\OperatorTok{(}\NormalTok{credentialId}\OperatorTok{,}\NormalTok{ numSignatures}\OperatorTok{,}\NormalTok{ otp}\OperatorTok{);}

    \BuiltInTok{HttpURLConnection}\NormalTok{ authConn }\OperatorTok{=} \OperatorTok{(}\BuiltInTok{HttpURLConnection}\OperatorTok{)}
            \KeywordTok{new} \BuiltInTok{URL}\OperatorTok{(}\NormalTok{authUrl}\OperatorTok{).}\FunctionTok{openConnection}\OperatorTok{();}
\NormalTok{    authConn}\OperatorTok{.}\FunctionTok{setDoOutput}\OperatorTok{(}\KeywordTok{true}\OperatorTok{);}
\NormalTok{    authConn}\OperatorTok{.}\FunctionTok{setRequestMethod}\OperatorTok{(}\StringTok{"POST"}\OperatorTok{);}
\NormalTok{    authConn}\OperatorTok{.}\FunctionTok{setRequestProperty}\OperatorTok{(}\StringTok{"Content{-}Type"}\OperatorTok{,} \StringTok{"application/json"}\OperatorTok{);}
\NormalTok{    authConn}\OperatorTok{.}\FunctionTok{setRequestProperty}\OperatorTok{(}\StringTok{"Authorization"}\OperatorTok{,} \StringTok{"Bearer "} \OperatorTok{+}\NormalTok{ accessToken}\OperatorTok{);}
    \ControlFlowTok{try} \OperatorTok{(}\BuiltInTok{OutputStream}\NormalTok{ out }\OperatorTok{=}\NormalTok{ authConn}\OperatorTok{.}\FunctionTok{getOutputStream}\OperatorTok{())} \OperatorTok{\{}
\NormalTok{        out}\OperatorTok{.}\FunctionTok{write}\OperatorTok{(}\NormalTok{authRequestBody}\OperatorTok{.}\FunctionTok{getBytes}\OperatorTok{(}\NormalTok{StandardCharsets}\OperatorTok{.}\FunctionTok{UTF\_8}\OperatorTok{));}
    \OperatorTok{\}}

    \BuiltInTok{String}\NormalTok{ responseJson}\OperatorTok{;}
    \ControlFlowTok{try} \OperatorTok{(}\BuiltInTok{InputStream}\NormalTok{ in }\OperatorTok{=}\NormalTok{ authConn}\OperatorTok{.}\FunctionTok{getInputStream}\OperatorTok{())} \OperatorTok{\{}
\NormalTok{        responseJson }\OperatorTok{=} \KeywordTok{new} \BuiltInTok{String}\OperatorTok{(}\NormalTok{in}\OperatorTok{.}\FunctionTok{readAllBytes}\OperatorTok{(),}\NormalTok{ StandardCharsets}\OperatorTok{.}\FunctionTok{UTF\_8}\OperatorTok{);}
    \OperatorTok{\}}

    \CommentTok{// Response: \{ "SAD": "eyJhbGciOiJSUzI1NiIs..." \}}
    \BuiltInTok{String}\NormalTok{ sad }\OperatorTok{=} \FunctionTok{extractJsonField}\OperatorTok{(}\NormalTok{responseJson}\OperatorTok{,} \StringTok{"SAD"}\OperatorTok{);}
\NormalTok{    log}\OperatorTok{.}\FunctionTok{info}\OperatorTok{(}\StringTok{"Credential authorized, SAD obtained for \{\} signatures"}\OperatorTok{,}
\NormalTok{            numSignatures}\OperatorTok{);}
    \ControlFlowTok{return}\NormalTok{ sad}\OperatorTok{;}
\OperatorTok{\}}
\end{Highlighting}
\end{Shaded}

The SAD token has a limited lifetime (typically 5-10 minutes) and a
limited number of uses (specified in \texttt{numSignatures}). If you
request authorization for 1 signature, the SAD can only be used once. If
you need to sign multiple documents in a batch, request a higher
\texttt{numSignatures} value --- but be aware that some QTSPs limit the
maximum.

If the SAD expires before you use it, you must re-authorize. Design your
signing flow so that SAD acquisition happens immediately before hash
signing, not minutes before.

\subsection{Hash Signing}\label{_hash_signing}

The core operation --- sending a hash and receiving a signature:

\begin{Shaded}
\begin{Highlighting}[]
\CommentTok{/**}
 \CommentTok{*}\NormalTok{ Signs one or more hashes using a CSC}\CommentTok{{-}}\NormalTok{compliant remote signing service}\CommentTok{.}
 \CommentTok{*}\NormalTok{ Returns raw signature values }\CommentTok{(}\NormalTok{one per hash}\CommentTok{).}
 \CommentTok{*/}
\KeywordTok{public} \BuiltInTok{List}\OperatorTok{\textless{}}\DataTypeTok{byte}\OperatorTok{[]\textgreater{}} \FunctionTok{signHashes}\OperatorTok{(}\BuiltInTok{String}\NormalTok{ cscBaseUrl}\OperatorTok{,} \BuiltInTok{String}\NormalTok{ accessToken}\OperatorTok{,}
        \BuiltInTok{String}\NormalTok{ credentialId}\OperatorTok{,} \BuiltInTok{String}\NormalTok{ sad}\OperatorTok{,}
        \BuiltInTok{List}\OperatorTok{\textless{}}\DataTypeTok{byte}\OperatorTok{[]\textgreater{}}\NormalTok{ hashes}\OperatorTok{)} \KeywordTok{throws} \BuiltInTok{IOException} \OperatorTok{\{}

    \BuiltInTok{String}\NormalTok{ signUrl }\OperatorTok{=}\NormalTok{ cscBaseUrl }\OperatorTok{+} \StringTok{"/signatures/signHash"}\OperatorTok{;}

    \CommentTok{// Convert hashes to Base64 for JSON transport}
    \BuiltInTok{List}\OperatorTok{\textless{}}\BuiltInTok{String}\OperatorTok{\textgreater{}}\NormalTok{ hashesBase64 }\OperatorTok{=} \KeywordTok{new} \BuiltInTok{ArrayList}\OperatorTok{\textless{}\textgreater{}();}
    \ControlFlowTok{for} \OperatorTok{(}\DataTypeTok{byte}\OperatorTok{[]}\NormalTok{ hash }\OperatorTok{:}\NormalTok{ hashes}\OperatorTok{)} \OperatorTok{\{}
\NormalTok{        hashesBase64}\OperatorTok{.}\FunctionTok{add}\OperatorTok{(}\NormalTok{Base64}\OperatorTok{.}\FunctionTok{getEncoder}\OperatorTok{().}\FunctionTok{encodeToString}\OperatorTok{(}\NormalTok{hash}\OperatorTok{));}
    \OperatorTok{\}}

    \CommentTok{// Build JSON array of hashes}
    \BuiltInTok{StringBuilder}\NormalTok{ hashArray }\OperatorTok{=} \KeywordTok{new} \BuiltInTok{StringBuilder}\OperatorTok{(}\StringTok{"["}\OperatorTok{);}
    \ControlFlowTok{for} \OperatorTok{(}\DataTypeTok{int}\NormalTok{ i }\OperatorTok{=} \DecValTok{0}\OperatorTok{;}\NormalTok{ i }\OperatorTok{\textless{}}\NormalTok{ hashesBase64}\OperatorTok{.}\FunctionTok{size}\OperatorTok{();}\NormalTok{ i}\OperatorTok{++)} \OperatorTok{\{}
        \ControlFlowTok{if} \OperatorTok{(}\NormalTok{i }\OperatorTok{\textgreater{}} \DecValTok{0}\OperatorTok{)}\NormalTok{ hashArray}\OperatorTok{.}\FunctionTok{append}\OperatorTok{(}\StringTok{","}\OperatorTok{);}
\NormalTok{        hashArray}\OperatorTok{.}\FunctionTok{append}\OperatorTok{(}\StringTok{"}\SpecialCharTok{\textbackslash{}"}\StringTok{"}\OperatorTok{).}\FunctionTok{append}\OperatorTok{(}\NormalTok{hashesBase64}\OperatorTok{.}\FunctionTok{get}\OperatorTok{(}\NormalTok{i}\OperatorTok{)).}\FunctionTok{append}\OperatorTok{(}\StringTok{"}\SpecialCharTok{\textbackslash{}"}\StringTok{"}\OperatorTok{);}
    \OperatorTok{\}}
\NormalTok{    hashArray}\OperatorTok{.}\FunctionTok{append}\OperatorTok{(}\StringTok{"]"}\OperatorTok{);}

    \BuiltInTok{String}\NormalTok{ requestBody }\OperatorTok{=} \StringTok{"""}
\StringTok{            \{}
\StringTok{                "credentialID": "\%s",}
\StringTok{                "SAD": "\%s",}
\StringTok{                "hashes": \%s,}
\StringTok{                "hashAlgo": "2.16.840.1.101.3.4.2.1",}
\StringTok{                "signAlgo": "1.2.840.113549.1.1.11"}
\StringTok{            \}}
\StringTok{            """}\OperatorTok{.}\FunctionTok{formatted}\OperatorTok{(}\NormalTok{credentialId}\OperatorTok{,}\NormalTok{ sad}\OperatorTok{,}\NormalTok{ hashArray}\OperatorTok{.}\FunctionTok{toString}\OperatorTok{());}

    \BuiltInTok{HttpURLConnection}\NormalTok{ conn }\OperatorTok{=} \OperatorTok{(}\BuiltInTok{HttpURLConnection}\OperatorTok{)}
            \KeywordTok{new} \BuiltInTok{URL}\OperatorTok{(}\NormalTok{signUrl}\OperatorTok{).}\FunctionTok{openConnection}\OperatorTok{();}
\NormalTok{    conn}\OperatorTok{.}\FunctionTok{setDoOutput}\OperatorTok{(}\KeywordTok{true}\OperatorTok{);}
\NormalTok{    conn}\OperatorTok{.}\FunctionTok{setRequestMethod}\OperatorTok{(}\StringTok{"POST"}\OperatorTok{);}
\NormalTok{    conn}\OperatorTok{.}\FunctionTok{setRequestProperty}\OperatorTok{(}\StringTok{"Content{-}Type"}\OperatorTok{,} \StringTok{"application/json"}\OperatorTok{);}
\NormalTok{    conn}\OperatorTok{.}\FunctionTok{setRequestProperty}\OperatorTok{(}\StringTok{"Authorization"}\OperatorTok{,} \StringTok{"Bearer "} \OperatorTok{+}\NormalTok{ accessToken}\OperatorTok{);}

    \ControlFlowTok{try} \OperatorTok{(}\BuiltInTok{OutputStream}\NormalTok{ out }\OperatorTok{=}\NormalTok{ conn}\OperatorTok{.}\FunctionTok{getOutputStream}\OperatorTok{())} \OperatorTok{\{}
\NormalTok{        out}\OperatorTok{.}\FunctionTok{write}\OperatorTok{(}\NormalTok{requestBody}\OperatorTok{.}\FunctionTok{getBytes}\OperatorTok{(}\NormalTok{StandardCharsets}\OperatorTok{.}\FunctionTok{UTF\_8}\OperatorTok{));}
    \OperatorTok{\}}

    \BuiltInTok{String}\NormalTok{ responseJson}\OperatorTok{;}
    \ControlFlowTok{try} \OperatorTok{(}\BuiltInTok{InputStream}\NormalTok{ in }\OperatorTok{=}\NormalTok{ conn}\OperatorTok{.}\FunctionTok{getInputStream}\OperatorTok{())} \OperatorTok{\{}
\NormalTok{        responseJson }\OperatorTok{=} \KeywordTok{new} \BuiltInTok{String}\OperatorTok{(}\NormalTok{in}\OperatorTok{.}\FunctionTok{readAllBytes}\OperatorTok{(),}\NormalTok{ StandardCharsets}\OperatorTok{.}\FunctionTok{UTF\_8}\OperatorTok{);}
    \OperatorTok{\}}

    \CommentTok{// Response: \{ "signatures": ["MEUCIQD...", "MEQCIBD..."] \}}
    \BuiltInTok{List}\OperatorTok{\textless{}}\DataTypeTok{byte}\OperatorTok{[]\textgreater{}}\NormalTok{ signatures }\OperatorTok{=} \FunctionTok{parseSignatureValues}\OperatorTok{(}\NormalTok{responseJson}\OperatorTok{);}
\NormalTok{    log}\OperatorTok{.}\FunctionTok{info}\OperatorTok{(}\StringTok{"CSC signHash returned \{\} signature values"}\OperatorTok{,}\NormalTok{ signatures}\OperatorTok{.}\FunctionTok{size}\OperatorTok{());}
    \ControlFlowTok{return}\NormalTok{ signatures}\OperatorTok{;}
\OperatorTok{\}}
\end{Highlighting}
\end{Shaded}

The OIDs in the request body specify the algorithms:

\begin{longtable}[]{@{}
  >{\raggedright\arraybackslash}p{(\linewidth - 4\tabcolsep) * \real{0.3333}}
  >{\raggedright\arraybackslash}p{(\linewidth - 4\tabcolsep) * \real{0.3333}}
  >{\raggedright\arraybackslash}p{(\linewidth - 4\tabcolsep) * \real{0.3333}}@{}}
\caption{Common algorithm OIDs for CSC signHash}\tabularnewline
\toprule\noalign{}
\begin{minipage}[b]{\linewidth}\raggedright
OID
\end{minipage} & \begin{minipage}[b]{\linewidth}\raggedright
Algorithm
\end{minipage} & \begin{minipage}[b]{\linewidth}\raggedright
Notes
\end{minipage} \\
\midrule\noalign{}
\endfirsthead
\toprule\noalign{}
\begin{minipage}[b]{\linewidth}\raggedright
OID
\end{minipage} & \begin{minipage}[b]{\linewidth}\raggedright
Algorithm
\end{minipage} & \begin{minipage}[b]{\linewidth}\raggedright
Notes
\end{minipage} \\
\midrule\noalign{}
\endhead
\bottomrule\noalign{}
\endlastfoot
\texttt{2.16.840.1.101.3.4.2.1} & SHA-256 & Hash algorithm. The hash you
send must use this algorithm. \\
\texttt{2.16.840.1.101.3.4.2.2} & SHA-384 & Stronger hash
alternative. \\
\texttt{1.2.840.113549.1.1.11} & SHA256withRSA (PKCS\#1 v1.5) & RSA
signature with SHA-256. The server applies PKCS\#1 v1.5 padding. \\
\texttt{1.2.840.113549.1.1.10} & RSASSA-PSS & RSA-PSS signature.
Preferred over PKCS\#1 v1.5 for new deployments. \\
\texttt{1.2.840.10045.4.3.2} & SHA256withECDSA & ECDSA signature with
SHA-256. Shorter signatures, equivalent security. \\
\end{longtable}

\section{SCAL: Sole Control Assurance
Levels}\label{_scal_sole_control_assurance_levels}

\subsection{The Problem SCAL Solves}\label{_the_problem_scal_solves}

The Regulation (EU) No 910/2014 (eIDAS) Article 26(c) requires that an
advanced electronic signature be "created using electronic signature
creation data that the signatory can, with a high level of confidence,
use under his sole control."

For local signing with a smart card, sole control is straightforward:
the signer physically possesses the card and knows the PIN. For remote
signing, sole control is more nuanced. The key lives in a server-side
HSM. Multiple users\textquotesingle{} keys may reside in the same
physical HSM. The QTSP's system administrators have physical access to
the HSM. How do you prove that the signer --- and only the signer ---
authorized a specific signing operation?

The \textbf{SCAL} (Sole Control Assurance Level) framework, defined in
CEN EN 419 241-1, provides two levels of assurance:

\subsection{SCAL1: Server-Side Sole
Control}\label{_scal1_server_side_sole_control}

In SCAL1, the remote signing server ensures sole control through its own
authentication and authorization mechanisms, but there is no explicit
per-signature authorization from the user beyond the initial session
authentication.

Characteristics: - The signer authenticates to the service (username +
password, possibly with MFA) - Once authenticated, the service can sign
on the signer's behalf for the duration of the session - The service's
access control mechanisms prevent unauthorized use - There is no
additional user interaction for each signing operation

SCAL1 is suitable for: - \textbf{Advanced electronic signatures (AdES)}:
The sole control requirement of Article 26(c) can be met if the
authentication is sufficiently strong - \textbf{Automated signing
workflows}: Where a human authorizes a batch of signatures upfront -
\textbf{Internal enterprise signing}: Where the trust model accepts
session-level authorization

SCAL1 is \emph{not} suitable for QES, because the sole control assurance
is not strong enough --- the server could theoretically sign without the
user's knowledge during an active session.

\subsection{SCAL2: Explicit Per-Signature
Authorization}\label{_scal2_explicit_per_signature_authorization}

In SCAL2, every signing operation requires explicit authorization from
the signer, independent of session authentication. The authorization is
linked to the specific data being signed.

Characteristics: - The signer authenticates to the service
(session-level) - For each signing operation, the signer must provide
additional authorization: OTP, biometric confirmation, mobile app push
approval - The authorization produces a SAD (Signature Activation Data)
token that is bound to the specific credential and limited in scope
(time, number of signatures) - The HSM will not perform the signing
operation without a valid SAD

SCAL2 is required for: - \textbf{Qualified electronic signatures (QES)}:
The per-signature authorization satisfies Article 26(c) with a high
level of confidence - \textbf{Any scenario where you need to prove that
the signer actively authorized each specific signature}

\begin{longtable}[]{@{}
  >{\raggedright\arraybackslash}p{(\linewidth - 4\tabcolsep) * \real{0.3333}}
  >{\raggedright\arraybackslash}p{(\linewidth - 4\tabcolsep) * \real{0.3333}}
  >{\raggedright\arraybackslash}p{(\linewidth - 4\tabcolsep) * \real{0.3333}}@{}}
\caption{SCAL1 vs. SCAL2 Comparison}\tabularnewline
\toprule\noalign{}
\begin{minipage}[b]{\linewidth}\raggedright
Aspect
\end{minipage} & \begin{minipage}[b]{\linewidth}\raggedright
SCAL1
\end{minipage} & \begin{minipage}[b]{\linewidth}\raggedright
SCAL2
\end{minipage} \\
\midrule\noalign{}
\endfirsthead
\toprule\noalign{}
\begin{minipage}[b]{\linewidth}\raggedright
Aspect
\end{minipage} & \begin{minipage}[b]{\linewidth}\raggedright
SCAL1
\end{minipage} & \begin{minipage}[b]{\linewidth}\raggedright
SCAL2
\end{minipage} \\
\midrule\noalign{}
\endhead
\bottomrule\noalign{}
\endlastfoot
Authentication & Session-level (login once) & Session-level +
per-signature authorization \\
User interaction per signature & None after login & Required (OTP,
biometric, push) \\
SAD token & Not required & Required for each signing operation \\
Suitable for QES & No & Yes \\
Suitable for AdES & Yes & Yes \\
Batch signing & Easy (sign N documents in a loop) & Possible but
requires batch SAD (N signatures authorized at once) \\
eIDAS Article 26(c) assurance & Moderate & High \\
\end{longtable}

The distinction between SCAL1 and SCAL2 is not merely theoretical --- it
determines whether the resulting signature can be a QES. If your
application needs to produce qualified electronic signatures via remote
signing, you \emph{must} implement a SCAL2 flow. A SCAL1 flow, no matter
how strong the session authentication, will not produce a QES because it
lacks per-signature authorization.

This is the single most common mistake in remote signing integrations:
assuming that a strong login (username + password + TOTP) is sufficient
for QES. It is not. The authentication must be linked to the signing
operation, not just the session.

\section{HSM and SAM Architecture}\label{_hsm_and_sam_architecture}

\subsection{The HSM (Hardware Security
Module)}\label{_the_hsm_hardware_security_module}

The HSM is the physical device that stores private keys and performs
cryptographic operations. In a remote signing deployment, the HSM is
typically:

\begin{itemize}
\item
  \textbf{FIPS 140-2 Level 3} or \textbf{Level 4} certified (US
  standard)
\item
  \textbf{Common Criteria EAL4+} certified (international standard)
\item
  Deployed in a physically secured data center with access logging,
  video surveillance, and multi-person access controls
\item
  Connected to the signing server via a high-speed internal network
  (never directly exposed to the internet)
\end{itemize}

Popular HSMs for remote signing: - \textbf{Thales Luna Network HSM}:
Widely used in European QTSPs. Supports PKCS\#11, JCE (Java Cryptography
Extension), and Microsoft CNG. Available in various throughput tiers. -
\textbf{Utimaco SecurityServer}: German manufacturer, popular with
European banks and government agencies. Strong PKCS\#11 support. -
\textbf{Entrust nShield Connect}: Used by large enterprises and some
QTSPs. Offers a proprietary "Security World" key management architecture
alongside standard PKCS\#11. - \textbf{AWS CloudHSM / Azure Dedicated
HSM}: Cloud-hosted HSMs. Convenient for cloud-native architectures, but
may not meet all QSCD certification requirements (check the specific
certification status).

\subsection{The SAM (Signature Activation
Module)}\label{_the_sam_signature_activation_module}

The SAM sits between the remote signing API and the HSM. Its job is to
enforce sole control:

\begin{enumerate}
\def\labelenumi{\arabic{enumi}.}
\item
  \textbf{Receive the signing request} with the signer's credentials and
  SAD
\item
  \textbf{Verify the SAD} --- check that it is valid, not expired, bound
  to the correct credential, and that the number of allowed signatures
  has not been exhausted
\item
  \textbf{Authorize the HSM operation} --- unlock the signer's key
  partition in the HSM for exactly one (or N) signing operations
\item
  \textbf{Return the signature value} to the API layer
\end{enumerate}

The SAM is the component that makes SCAL2 possible. Without a SAM, the
signing server could theoretically use any key in the HSM at any time.
The SAM ensures that keys are only activated when the legitimate signer
has provided explicit, verifiable authorization.

\begin{verbatim}
Architecture of a Remote Signing Service

                         ┌─────────────────────┐
    Internet             │   Load Balancer /    │
    ─────────────────────│   API Gateway        │
                         └──────────┬───────────┘
                                    │
                         ┌──────────┴───────────┐
                         │   CSC API Server      │
                         │   (REST endpoints)    │
                         │   OAuth 2.0 token     │
                         │   validation          │
                         └──────────┬───────────┘
                                    │
                         ┌──────────┴───────────┐
                         │   Signature           │
                         │   Activation Module   │
                         │   (SAM)               │
                         │   • SAD validation    │
                         │   • SCAL enforcement  │
                         │   • Audit logging     │
                         └──────────┬───────────┘
                                    │ PKCS#11
                         ┌──────────┴───────────┐
                         │   Hardware Security   │
                         │   Module (HSM)        │
                         │   • Key storage       │
                         │   • Crypto operations │
                         │   • Tamper-evident    │
                         └──────────────────────┘
\end{verbatim}

\subsection{PKCS\#11: The HSM
Interface}\label{_pkcs11_the_hsm_interface}

PKCS\#11 (also called Cryptoki) is the standard API for communicating
with HSMs and smart cards. It is a C-based API with Java bindings
available through the JCE \texttt{SunPKCS11} provider or through
dedicated libraries.

In a remote signing server, the PKCS\#11 flow for signing looks like
this:

\begin{Shaded}
\begin{Highlighting}[]
\CommentTok{/**}
 \CommentTok{*}\NormalTok{ Signs a hash using an HSM via PKCS}\CommentTok{\#11.}
 \CommentTok{*}\NormalTok{ This runs on the remote signing server}\CommentTok{,}\NormalTok{ NOT on the client}\CommentTok{.}
 \CommentTok{*/}
\KeywordTok{public} \DataTypeTok{byte}\OperatorTok{[]} \FunctionTok{signWithHsm}\OperatorTok{(}\DataTypeTok{byte}\OperatorTok{[]}\NormalTok{ hash}\OperatorTok{,} \BuiltInTok{String}\NormalTok{ keyAlias}\OperatorTok{,}
        \DataTypeTok{char}\OperatorTok{[]}\NormalTok{ keyPin}\OperatorTok{)} \KeywordTok{throws} \BuiltInTok{Exception} \OperatorTok{\{}

    \CommentTok{// 1. Load the PKCS\#11 provider}
    \CommentTok{//    (configured once at server startup, not per request)}
    \BuiltInTok{Provider}\NormalTok{ pkcs11Provider }\OperatorTok{=} \BuiltInTok{Security}\OperatorTok{.}\FunctionTok{getProvider}\OperatorTok{(}\StringTok{"SunPKCS11"}\OperatorTok{);}

    \CommentTok{// 2. Load the keystore from the HSM}
    \BuiltInTok{KeyStore}\NormalTok{ ks }\OperatorTok{=} \BuiltInTok{KeyStore}\OperatorTok{.}\FunctionTok{getInstance}\OperatorTok{(}\StringTok{"PKCS11"}\OperatorTok{,}\NormalTok{ pkcs11Provider}\OperatorTok{);}
\NormalTok{    ks}\OperatorTok{.}\FunctionTok{load}\OperatorTok{(}\KeywordTok{null}\OperatorTok{,}\NormalTok{ keyPin}\OperatorTok{);}

    \CommentTok{// 3. Retrieve the private key}
    \BuiltInTok{PrivateKey}\NormalTok{ privateKey }\OperatorTok{=} \OperatorTok{(}\BuiltInTok{PrivateKey}\OperatorTok{)}\NormalTok{ ks}\OperatorTok{.}\FunctionTok{getKey}\OperatorTok{(}\NormalTok{keyAlias}\OperatorTok{,}\NormalTok{ keyPin}\OperatorTok{);}
    \ControlFlowTok{if} \OperatorTok{(}\NormalTok{privateKey }\OperatorTok{==} \KeywordTok{null}\OperatorTok{)} \OperatorTok{\{}
        \ControlFlowTok{throw} \KeywordTok{new} \BuiltInTok{IllegalStateException}\OperatorTok{(}
                \StringTok{"Key not found in HSM: "} \OperatorTok{+}\NormalTok{ keyAlias}\OperatorTok{);}
    \OperatorTok{\}}

    \CommentTok{// 4. Sign the hash}
    \CommentTok{//    Note: for RSA, Signature.sign() expects the raw data, not a}
    \CommentTok{//    pre{-}computed hash. Use "NONEwithRSA" to sign a raw hash, or}
    \CommentTok{//    use "SHA256withRSA" and pass the original data.}
    \BuiltInTok{Signature}\NormalTok{ sig }\OperatorTok{=} \BuiltInTok{Signature}\OperatorTok{.}\FunctionTok{getInstance}\OperatorTok{(}\StringTok{"SHA256withRSA"}\OperatorTok{,}\NormalTok{ pkcs11Provider}\OperatorTok{);}
\NormalTok{    sig}\OperatorTok{.}\FunctionTok{initSign}\OperatorTok{(}\NormalTok{privateKey}\OperatorTok{);}
\NormalTok{    sig}\OperatorTok{.}\FunctionTok{update}\OperatorTok{(}\NormalTok{hash}\OperatorTok{);}
    \DataTypeTok{byte}\OperatorTok{[]}\NormalTok{ signatureValue }\OperatorTok{=}\NormalTok{ sig}\OperatorTok{.}\FunctionTok{sign}\OperatorTok{();}

\NormalTok{    log}\OperatorTok{.}\FunctionTok{info}\OperatorTok{(}\StringTok{"HSM signing completed: key=\{\}, sigLen=\{\}"}\OperatorTok{,}
\NormalTok{            keyAlias}\OperatorTok{,}\NormalTok{ signatureValue}\OperatorTok{.}\FunctionTok{length}\OperatorTok{);}
    \ControlFlowTok{return}\NormalTok{ signatureValue}\OperatorTok{;}
\OperatorTok{\}}
\end{Highlighting}
\end{Shaded}

The JCE \texttt{SunPKCS11} provider requires a configuration that
specifies the PKCS\#11 library path, slot number, and other HSM-specific
parameters. This configuration is typically loaded from a file or a
dynamic \texttt{InputStream}:

\begin{verbatim}
# pkcs11.cfg
name = MyHSM
library = /opt/hsm/lib/libCryptoki2_64.so
slot = 0
\end{verbatim}

Different HSM vendors provide different PKCS\#11 shared libraries
(\texttt{.so} on Linux, \texttt{.dll} on Windows). The library path is
the most common source of "HSM not found" errors in deployments.

\subsection{Key Lifecycle in an HSM}\label{_key_lifecycle_in_an_hsm}

Private keys in an HSM follow a strict lifecycle:

\begin{enumerate}
\def\labelenumi{\arabic{enumi}.}
\item
  \textbf{Generation}: The key pair is generated \emph{inside} the HSM.
  The private key never exists outside the HSM's tamper-resistant
  boundary. This is a fundamental requirement for QSCD certification.
\item
  \textbf{Activation}: The key is associated with a certificate (issued
  by the QTSP's CA) and made available for signing. For SCAL2,
  activation for each signing operation requires a valid SAD.
\item
  \textbf{Use}: The HSM performs signing operations using the key. Each
  use is logged for audit purposes.
\item
  \textbf{Suspension}: The key can be temporarily suspended (e.g., if
  the signer reports a compromised authentication factor). The
  certificate is not revoked, but signing is disabled.
\item
  \textbf{Revocation}: The key is permanently disabled, and the
  associated certificate is revoked. The key material may be destroyed
  or archived depending on the QTSP's policy.
\end{enumerate}

For QSCD certification under EN 419 241-2, the HSM must guarantee that
private keys are generated internally and can never be extracted --- not
even by the QTSP's own administrators. This is what "the key never
leaves the HSM" means in practice. The HSM must be designed so that
there is no API, no debug mode, no firmware backdoor that could extract
key material. Physical attacks (opening the HSM casing) must trigger key
destruction.

\section{Vendor Landscape}\label{_vendor_landscape}

\subsection{Major Remote Signing
Providers}\label{_major_remote_signing_providers}

The remote signing market is dominated by a handful of European QTSPs,
each offering CSC-compliant APIs with varying degrees of completeness:

\begin{longtable}[]{@{}
  >{\raggedright\arraybackslash}p{(\linewidth - 8\tabcolsep) * \real{0.2000}}
  >{\raggedright\arraybackslash}p{(\linewidth - 8\tabcolsep) * \real{0.2000}}
  >{\raggedright\arraybackslash}p{(\linewidth - 8\tabcolsep) * \real{0.2000}}
  >{\raggedright\arraybackslash}p{(\linewidth - 8\tabcolsep) * \real{0.2000}}
  >{\raggedright\arraybackslash}p{(\linewidth - 8\tabcolsep) * \real{0.2000}}@{}}
\caption{Remote Signing Provider Comparison}\tabularnewline
\toprule\noalign{}
\begin{minipage}[b]{\linewidth}\raggedright
Provider
\end{minipage} & \begin{minipage}[b]{\linewidth}\raggedright
Headquarters
\end{minipage} & \begin{minipage}[b]{\linewidth}\raggedright
CSC API
\end{minipage} & \begin{minipage}[b]{\linewidth}\raggedright
SCAL Levels
\end{minipage} & \begin{minipage}[b]{\linewidth}\raggedright
Notable Features
\end{minipage} \\
\midrule\noalign{}
\endfirsthead
\toprule\noalign{}
\begin{minipage}[b]{\linewidth}\raggedright
Provider
\end{minipage} & \begin{minipage}[b]{\linewidth}\raggedright
Headquarters
\end{minipage} & \begin{minipage}[b]{\linewidth}\raggedright
CSC API
\end{minipage} & \begin{minipage}[b]{\linewidth}\raggedright
SCAL Levels
\end{minipage} & \begin{minipage}[b]{\linewidth}\raggedright
Notable Features
\end{minipage} \\
\midrule\noalign{}
\endhead
\bottomrule\noalign{}
\endlastfoot
Swisscom AIS & Switzerland & v2.0 & SCAL1, SCAL2 & Integrated with Swiss
law (ZertES). Strong mobile authentication via Mobile ID. High
throughput for batch signing. \\
DigiCert + QuoVadis & US / Switzerland & v2.0 & SCAL1, SCAL2 & Global
presence. QuoVadis (acquired by DigiCert) provides the EU-qualified
certificates. AATL-listed roots. \\
InfoCert & Italy & v2.0 & SCAL1, SCAL2 & Largest QTSP in Europe by
volume. SPID integration for Italian digital identity. Strong in Italian
government sector. \\
Namirial & Italy & v2.0 & SCAL1, SCAL2 & Focus on document workflow
integration. Offers both remote signing and a full document management
platform. \\
ZealID & Lithuania & v2.0 & SCAL2 & Mobile-first approach. Signing via
smartphone app with biometric authentication. Popular in Baltic
states. \\
D-Trust (Bundesdruckerei) & Germany & v2.0 & SCAL1, SCAL2 & German
Federal Printing Office subsidiary. Strong in German government and
healthcare sectors. \\
\end{longtable}

\subsection{The Ascertia/ADSS Pattern}\label{_the_ascertiaadss_pattern}

A notable infrastructure pattern in the CSC ecosystem is the widespread
use of \textbf{Ascertia ADSS Server} as the underlying platform. Many
QTSPs --- especially regional and government-affiliated providers --- do
not build their own CSC API implementation. Instead, they deploy
Ascertia's ADSS (Automated Digital Signature Server), which exposes the
CSC API through a consistent base path:

\begin{verbatim}
https://<provider-domain>/adss/service/ras/csc/v1/
\end{verbatim}

This creates a recognizable fingerprint. If you encounter a CSC endpoint
with \texttt{/adss/service/ras/} in the URL, you are almost certainly
talking to an Ascertia deployment, regardless of the provider's
branding. The practical consequence: the CSC behavior, error formats,
and quirks will be identical across all Ascertia-based providers, even
though they are operated by different organizations in different
countries.

Common variations include:

\begin{itemize}
\item
  \texttt{https://signinghub.ascertia.com/adss/service/ras/csc/v1/}
  (Ascertia's own hosted service)
\item
  \texttt{https://ras.\textless{}provider\textgreater{}.com/adss/service/ras/csc/v1/}
  (white-labeled deployments)
\item
  \texttt{https://sign.\textless{}provider\textgreater{}.com/adss/service/ras/csc/v1/}
  (government integrations)
\end{itemize}

This pattern matters for integration testing: if you have tested against
one Ascertia-based provider, you have effectively tested against all of
them. Conversely, if you encounter an Ascertia-specific bug, expect it
to appear at every provider using the same platform.

\subsection{Sandbox and Test
Endpoints}\label{_sandbox_and_test_endpoints}

Most CSC providers offer sandbox environments for development and
integration testing. These sandboxes are invaluable --- they let you
develop and test the entire signing flow without purchasing production
credentials or risking real certificates.

\begin{longtable}[]{@{}
  >{\raggedright\arraybackslash}p{(\linewidth - 4\tabcolsep) * \real{0.2857}}
  >{\raggedright\arraybackslash}p{(\linewidth - 4\tabcolsep) * \real{0.4286}}
  >{\raggedright\arraybackslash}p{(\linewidth - 4\tabcolsep) * \real{0.2857}}@{}}
\caption{CSC Provider Sandbox and Test Endpoints}\tabularnewline
\toprule\noalign{}
\begin{minipage}[b]{\linewidth}\raggedright
Provider
\end{minipage} & \begin{minipage}[b]{\linewidth}\raggedright
Sandbox URL
\end{minipage} & \begin{minipage}[b]{\linewidth}\raggedright
Notes
\end{minipage} \\
\midrule\noalign{}
\endfirsthead
\toprule\noalign{}
\begin{minipage}[b]{\linewidth}\raggedright
Provider
\end{minipage} & \begin{minipage}[b]{\linewidth}\raggedright
Sandbox URL
\end{minipage} & \begin{minipage}[b]{\linewidth}\raggedright
Notes
\end{minipage} \\
\midrule\noalign{}
\endhead
\bottomrule\noalign{}
\endlastfoot
SSL.com eSigner & \texttt{https://cs-try.ssl.com/csc/v0} & Free sandbox
with demo credentials. Publicly documented test accounts available. Good
starting point for CSC integration. \\
DigiCert DTM &
\texttt{https://demo.one.digicert.com/documentmanager/csc/v1} & Demo
environment for Document Trust Manager. Requires DigiCert partnership
account. \\
ZealiD & \texttt{https://hermes-dev.zealid.com/api/v3.0/csc/v1} &
Development endpoint. Mobile-first --- requires ZealiD test app. \\
itsme & \texttt{https://sign.e2e.itsme.services/csc/v1} & End-to-end
test environment. Belgian Mobile ID integration testing. \\
\end{longtable}

\textbf{Start with SSL.com's sandbox.} It is the most accessible CSC
test environment --- no partnership agreements, no special hardware,
publicly documented demo credentials. Build your integration against
their sandbox first, then adapt for your production provider. The CSC
API surface is standardized enough that switching providers later
requires only URL and credential changes, not code changes.

\subsection{National eID and CSC
Integration}\label{_national_eid_and_csc_integration}

Several European countries have integrated their national electronic
identity systems with CSC-compatible remote signing. This means citizens
can use their existing government-issued digital identity to create
qualified electronic signatures without obtaining separate signing
credentials.

\begin{longtable}[]{@{}
  >{\raggedright\arraybackslash}p{(\linewidth - 4\tabcolsep) * \real{0.2500}}
  >{\raggedright\arraybackslash}p{(\linewidth - 4\tabcolsep) * \real{0.3750}}
  >{\raggedright\arraybackslash}p{(\linewidth - 4\tabcolsep) * \real{0.3750}}@{}}
\caption{National eID systems with CSC-compatible remote
signing}\tabularnewline
\toprule\noalign{}
\begin{minipage}[b]{\linewidth}\raggedright
Country
\end{minipage} & \begin{minipage}[b]{\linewidth}\raggedright
eID System
\end{minipage} & \begin{minipage}[b]{\linewidth}\raggedright
CSC Integration
\end{minipage} \\
\midrule\noalign{}
\endfirsthead
\toprule\noalign{}
\begin{minipage}[b]{\linewidth}\raggedright
Country
\end{minipage} & \begin{minipage}[b]{\linewidth}\raggedright
eID System
\end{minipage} & \begin{minipage}[b]{\linewidth}\raggedright
CSC Integration
\end{minipage} \\
\midrule\noalign{}
\endhead
\bottomrule\noalign{}
\endlastfoot
Italy & SPID (Public Digital Identity System) & InfoCert and Namirial
both offer SPID-authenticated remote signing. Citizens use their
existing SPID identity to access signing credentials hosted by the
QTSP. \\
Norway & BankID & BankID provides remote signing through bank-operated
infrastructure. High adoption rate (\textasciitilde4.5 million active
users). \\
Austria & Handy-Signatur / ID Austria & PrimeSign (Cryptas)
infrastructure. Mobile phone-based signing via the national eID
system. \\
Estonia & eID Easy & Integrates with the Estonian national ID card and
Mobile-ID systems for remote signing. \\
Switzerland & SwissID / Mobile ID & Swisscom Trust Services. Qualified
under Swiss law (ZertES) rather than eIDAS, but cross-recognized. \\
Spain & DNIe / FNMT & FNMT (Fábrica Nacional de Moneda y Timbre)
provides remote signing tied to the national certificate
infrastructure. \\
\end{longtable}

For integrators, the key implication is authentication: when the CSC
provider is backed by a national eID system, the OAuth 2.0 flow
typically redirects through the national identity provider (OpenID
Connect federation). Your application's authentication UI must support
this redirect --- it is not a simple username/password flow.

\subsection{CSC Endpoint Discovery}\label{_csc_endpoint_discovery}

Because the CSC specification requires a public \texttt{/info} endpoint
that responds without authentication, CSC servers can be discovered by
probing for known URL patterns:

\begin{verbatim}
GET /csc/v2/info   → Most CSC v2 providers
GET /csc/v1/info   → Legacy CSC v1 providers
GET /adss/service/ras/csc/v1/info   → Ascertia ADSS deployments
\end{verbatim}

A successful response returns the service capabilities:

\begin{Shaded}
\begin{Highlighting}[]
\FunctionTok{\{}
  \DataTypeTok{"specs"}\FunctionTok{:} \StringTok{"CSC API v2"}\FunctionTok{,}
  \DataTypeTok{"name"}\FunctionTok{:} \StringTok{"Remote Signing Service"}\FunctionTok{,}
  \DataTypeTok{"authType"}\FunctionTok{:} \OtherTok{[}\StringTok{"oauth2"}\OtherTok{]}\FunctionTok{,}
  \DataTypeTok{"methods"}\FunctionTok{:} \OtherTok{[}\StringTok{"credentials/list"}\OtherTok{,} \StringTok{"credentials/info"}\OtherTok{,}
              \StringTok{"credentials/authorize"}\OtherTok{,} \StringTok{"signatures/signHash"}\OtherTok{]}
\FunctionTok{\}}
\end{Highlighting}
\end{Shaded}

This discovery mechanism is useful for integration testing (verifying
that a provider's endpoint is reachable and responsive) and for building
provider registries. The \texttt{/info} endpoint also reveals which CSC
API methods the provider supports --- not all providers implement the
full specification.

\subsection{Choosing a Provider}\label{_choosing_a_provider}

When selecting a remote signing provider, consider:

\begin{enumerate}
\def\labelenumi{\arabic{enumi}.}
\item
  \textbf{Trusted list presence}: The provider must be a QTSP listed in
  the EUTL of the member state where they are supervised. Verify this at
  the EU's official Trusted List Browser.
\item
  \textbf{AATL listing}: If your signed PDFs will be opened in Adobe
  Acrobat, the provider's root certificate should be in the Adobe
  Approved Trust List (AATL). Without AATL listing, Acrobat will show a
  warning even for EU-qualified signatures.
\item
  \textbf{Geographic coverage}: Some providers issue certificates
  qualified in one member state; others have qualified status in
  multiple states. For cross-border use, mutual recognition under eIDAS
  Article 25(3) applies, but AATL listing is a separate concern.
\item
  \textbf{API completeness}: Not all providers implement the full CSC
  API. Some support only SCAL1, some have proprietary extensions, some
  are still on CSC v1.x. Request the provider's API documentation and
  test against their sandbox environment.
\item
  \textbf{Authentication options}: SCAL2 requires per-signature
  authorization. Providers offer different mechanisms: SMS OTP, mobile
  app push, biometric (fingerprint/face), hardware token. Choose based
  on your users\textquotesingle{} capabilities.
\item
  \textbf{Pricing model}: Per-signature pricing (typically 0.10-2.00 EUR
  per signature), per-certificate pricing (annual fee per signer), or
  volume-based tiers. Calculate your total cost based on expected
  signing volume.
\item
  \textbf{SLA and availability}: For production signing, you need a
  provider with a published SLA. Remote signing is a synchronous
  operation in most workflows --- if the provider is down, your users
  cannot sign.
\end{enumerate}

Do not assume that all CSC API implementations are interchangeable. The
CSC specification leaves room for provider-specific behavior in areas
like:

\begin{itemize}
\item
  OTP delivery mechanisms (some providers trigger OTP automatically,
  others require an explicit \texttt{/credentials/sendOTP} call)
\item
  SAD token format and lifetime
\item
  Batch signing limits (how many hashes can be signed with a single SAD)
\item
  Error response formats and error codes
\end{itemize}

Always develop and test against the specific provider's sandbox before
going to production. The CSC specification is a standard, but
implementations have edges.

\section{Complete Integration
Example}\label{_complete_integration_example}

The following example shows a complete remote signing flow --- from PDF
preparation through CSC API calls to final PDF assembly:

\begin{Shaded}
\begin{Highlighting}[]
\CommentTok{/**}
 \CommentTok{*}\NormalTok{ Signs a PDF using a CSC}\CommentTok{{-}}\NormalTok{compliant remote signing service}\CommentTok{.}
 \CommentTok{*}\NormalTok{ Produces a PAdES}\CommentTok{{-}}\NormalTok{B}\CommentTok{{-}}\NormalTok{T compliant signature}\CommentTok{.}
 \CommentTok{*/}
\KeywordTok{public} \DataTypeTok{byte}\OperatorTok{[]} \FunctionTok{signPdfRemotely}\OperatorTok{(}\DataTypeTok{byte}\OperatorTok{[]}\NormalTok{ pdfBytes}\OperatorTok{,} \BuiltInTok{String}\NormalTok{ cscBaseUrl}\OperatorTok{,}
        \BuiltInTok{String}\NormalTok{ accessToken}\OperatorTok{,} \BuiltInTok{String}\NormalTok{ credentialId}\OperatorTok{,}
        \BuiltInTok{String}\NormalTok{ tsaUrl}\OperatorTok{)} \KeywordTok{throws} \BuiltInTok{Exception} \OperatorTok{\{}

    \CommentTok{// {-}{-}{-} Phase 1: Prepare the PDF and CMS structure {-}{-}{-}}

    \CommentTok{// 1a. Get the signer\textquotesingle{}s certificate from the CSC service}
\NormalTok{    CscCredentialInfo credInfo }\OperatorTok{=} \FunctionTok{getCredentialInfo}\OperatorTok{(}
\NormalTok{            cscBaseUrl}\OperatorTok{,}\NormalTok{ accessToken}\OperatorTok{,}\NormalTok{ credentialId}\OperatorTok{);}
    \BuiltInTok{X509Certificate}\NormalTok{ signerCert }\OperatorTok{=}\NormalTok{ credInfo}\OperatorTok{.}\FunctionTok{getCertificateChain}\OperatorTok{().}\FunctionTok{get}\OperatorTok{(}\DecValTok{0}\OperatorTok{);}
    \BuiltInTok{List}\OperatorTok{\textless{}}\BuiltInTok{X509Certificate}\OperatorTok{\textgreater{}}\NormalTok{ certChain }\OperatorTok{=}\NormalTok{ credInfo}\OperatorTok{.}\FunctionTok{getCertificateChain}\OperatorTok{();}

    \CommentTok{// 1b. Load the PDF and prepare the signature dictionary}
\NormalTok{    PDDocument doc }\OperatorTok{=}\NormalTok{ Loader}\OperatorTok{.}\FunctionTok{loadPDF}\OperatorTok{(}\NormalTok{pdfBytes}\OperatorTok{);}
\NormalTok{    PDSignature pdSignature }\OperatorTok{=} \KeywordTok{new} \FunctionTok{PDSignature}\OperatorTok{();}
\NormalTok{    pdSignature}\OperatorTok{.}\FunctionTok{setFilter}\OperatorTok{(}\NormalTok{PDSignature}\OperatorTok{.}\FunctionTok{FILTER\_ADOBE\_PPKLITE}\OperatorTok{);}
\NormalTok{    pdSignature}\OperatorTok{.}\FunctionTok{setSubFilter}\OperatorTok{(}
\NormalTok{            COSName}\OperatorTok{.}\FunctionTok{getPDFName}\OperatorTok{(}\StringTok{"ETSI.CAdES.detached"}\OperatorTok{));}
\NormalTok{    pdSignature}\OperatorTok{.}\FunctionTok{setSignDate}\OperatorTok{(}\BuiltInTok{Calendar}\OperatorTok{.}\FunctionTok{getInstance}\OperatorTok{());}

    \CommentTok{// 1c. Register our SignatureInterface}
    \CommentTok{// (PDFBox calls this to get the CMS bytes)}
\NormalTok{    RemoteSignatureInterface sigInterface }\OperatorTok{=}
            \KeywordTok{new} \FunctionTok{RemoteSignatureInterface}\OperatorTok{(}
\NormalTok{                    cscBaseUrl}\OperatorTok{,}\NormalTok{ accessToken}\OperatorTok{,}\NormalTok{ credentialId}\OperatorTok{,}
\NormalTok{                    signerCert}\OperatorTok{,}\NormalTok{ certChain}\OperatorTok{,}\NormalTok{ tsaUrl}\OperatorTok{);}

\NormalTok{    doc}\OperatorTok{.}\FunctionTok{addSignature}\OperatorTok{(}\NormalTok{pdSignature}\OperatorTok{,}\NormalTok{ sigInterface}\OperatorTok{);}

    \CommentTok{// 1d. Save incrementally — this triggers sigInterface.sign()}
    \BuiltInTok{ByteArrayOutputStream}\NormalTok{ baos }\OperatorTok{=} \KeywordTok{new} \BuiltInTok{ByteArrayOutputStream}\OperatorTok{();}
\NormalTok{    doc}\OperatorTok{.}\FunctionTok{saveIncremental}\OperatorTok{(}\NormalTok{baos}\OperatorTok{);}
\NormalTok{    doc}\OperatorTok{.}\FunctionTok{close}\OperatorTok{();}

\NormalTok{    log}\OperatorTok{.}\FunctionTok{info}\OperatorTok{(}\StringTok{"PDF signed remotely: resultSize=\{\} bytes"}\OperatorTok{,}\NormalTok{ baos}\OperatorTok{.}\FunctionTok{size}\OperatorTok{());}
    \ControlFlowTok{return}\NormalTok{ baos}\OperatorTok{.}\FunctionTok{toByteArray}\OperatorTok{();}
\OperatorTok{\}}

\CommentTok{/**}
 \CommentTok{*}\NormalTok{ PDFBox SignatureInterface implementation for remote signing}\CommentTok{.}
 \CommentTok{*}\NormalTok{ Called by PDFBox during saveIncremental}\CommentTok{()}\NormalTok{ with the byte range content}\CommentTok{.}
 \CommentTok{*/}
\KeywordTok{public} \KeywordTok{class}\NormalTok{ RemoteSignatureInterface}
        \KeywordTok{implements}\NormalTok{ SignatureInterface }\OperatorTok{\{}

    \KeywordTok{private} \DataTypeTok{final} \BuiltInTok{String}\NormalTok{ cscBaseUrl}\OperatorTok{;}
    \KeywordTok{private} \DataTypeTok{final} \BuiltInTok{String}\NormalTok{ accessToken}\OperatorTok{;}
    \KeywordTok{private} \DataTypeTok{final} \BuiltInTok{String}\NormalTok{ credentialId}\OperatorTok{;}
    \KeywordTok{private} \DataTypeTok{final} \BuiltInTok{X509Certificate}\NormalTok{ signerCert}\OperatorTok{;}
    \KeywordTok{private} \DataTypeTok{final} \BuiltInTok{List}\OperatorTok{\textless{}}\BuiltInTok{X509Certificate}\OperatorTok{\textgreater{}}\NormalTok{ certChain}\OperatorTok{;}
    \KeywordTok{private} \DataTypeTok{final} \BuiltInTok{String}\NormalTok{ tsaUrl}\OperatorTok{;}

    \CommentTok{// Constructor omitted for brevity}

    \AttributeTok{@Override}
    \KeywordTok{public} \DataTypeTok{byte}\OperatorTok{[]} \FunctionTok{sign}\OperatorTok{(}\BuiltInTok{InputStream}\NormalTok{ content}\OperatorTok{)} \KeywordTok{throws} \BuiltInTok{IOException} \OperatorTok{\{}
        \ControlFlowTok{try} \OperatorTok{\{}
            \CommentTok{// 2a. Read the byte range content (the data to be signed)}
            \DataTypeTok{byte}\OperatorTok{[]}\NormalTok{ contentBytes }\OperatorTok{=}\NormalTok{ content}\OperatorTok{.}\FunctionTok{readAllBytes}\OperatorTok{();}

            \CommentTok{// 2b. Build the CMS signed attributes}
\NormalTok{            CMSSignedDataGenerator gen }\OperatorTok{=} \KeywordTok{new} \FunctionTok{CMSSignedDataGenerator}\OperatorTok{();}

            \CommentTok{// Add certificates to the CMS}
            \BuiltInTok{List}\OperatorTok{\textless{}}\NormalTok{X509CertificateHolder}\OperatorTok{\textgreater{}}\NormalTok{ certHolders }\OperatorTok{=} \KeywordTok{new} \BuiltInTok{ArrayList}\OperatorTok{\textless{}\textgreater{}();}
            \ControlFlowTok{for} \OperatorTok{(}\BuiltInTok{X509Certificate}\NormalTok{ cert }\OperatorTok{:}\NormalTok{ certChain}\OperatorTok{)} \OperatorTok{\{}
\NormalTok{                certHolders}\OperatorTok{.}\FunctionTok{add}\OperatorTok{(}\KeywordTok{new} \FunctionTok{JcaX509CertificateHolder}\OperatorTok{(}\NormalTok{cert}\OperatorTok{));}
            \OperatorTok{\}}
\NormalTok{            gen}\OperatorTok{.}\FunctionTok{addCertificates}\OperatorTok{(}
                    \KeywordTok{new}\NormalTok{ CollectionStore}\OperatorTok{\textless{}\textgreater{}(}\NormalTok{certHolders}\OperatorTok{));}

            \CommentTok{// 2c. Compute the message digest of the content}
            \BuiltInTok{MessageDigest}\NormalTok{ md }\OperatorTok{=} \BuiltInTok{MessageDigest}\OperatorTok{.}\FunctionTok{getInstance}\OperatorTok{(}\StringTok{"SHA{-}256"}\OperatorTok{);}
            \DataTypeTok{byte}\OperatorTok{[]}\NormalTok{ contentHash }\OperatorTok{=}\NormalTok{ md}\OperatorTok{.}\FunctionTok{digest}\OperatorTok{(}\NormalTok{contentBytes}\OperatorTok{);}

            \CommentTok{// 2d. Build signed attributes manually}
\NormalTok{            ASN1EncodableVector signedAttrs }\OperatorTok{=}
                    \KeywordTok{new} \FunctionTok{ASN1EncodableVector}\OperatorTok{();}

            \CommentTok{// content{-}type attribute}
\NormalTok{            signedAttrs}\OperatorTok{.}\FunctionTok{add}\OperatorTok{(}\KeywordTok{new} \BuiltInTok{Attribute}\OperatorTok{(}
\NormalTok{                    CMSAttributes}\OperatorTok{.}\FunctionTok{contentType}\OperatorTok{,}
                    \KeywordTok{new} \FunctionTok{DERSet}\OperatorTok{(}\NormalTok{CMSObjectIdentifiers}\OperatorTok{.}\FunctionTok{data}\OperatorTok{)));}

            \CommentTok{// message{-}digest attribute}
\NormalTok{            signedAttrs}\OperatorTok{.}\FunctionTok{add}\OperatorTok{(}\KeywordTok{new} \BuiltInTok{Attribute}\OperatorTok{(}
\NormalTok{                    CMSAttributes}\OperatorTok{.}\FunctionTok{messageDigest}\OperatorTok{,}
                    \KeywordTok{new} \FunctionTok{DERSet}\OperatorTok{(}\KeywordTok{new} \FunctionTok{DEROctetString}\OperatorTok{(}\NormalTok{contentHash}\OperatorTok{))));}

            \CommentTok{// ESS{-}signing{-}certificate{-}v2 attribute}
            \DataTypeTok{byte}\OperatorTok{[]}\NormalTok{ certHash }\OperatorTok{=} \BuiltInTok{MessageDigest}\OperatorTok{.}\FunctionTok{getInstance}\OperatorTok{(}\StringTok{"SHA{-}256"}\OperatorTok{)}
                    \OperatorTok{.}\FunctionTok{digest}\OperatorTok{(}\NormalTok{signerCert}\OperatorTok{.}\FunctionTok{getEncoded}\OperatorTok{());}
\NormalTok{            ESSCertIDv2 essCertId }\OperatorTok{=} \KeywordTok{new} \FunctionTok{ESSCertIDv2}\OperatorTok{(}
                    \KeywordTok{new} \FunctionTok{AlgorithmIdentifier}\OperatorTok{(}
\NormalTok{                            NISTObjectIdentifiers}\OperatorTok{.}\FunctionTok{id\_sha256}\OperatorTok{),}
\NormalTok{                    certHash}\OperatorTok{);}
\NormalTok{            SigningCertificateV2 sigCertV2 }\OperatorTok{=}
                    \KeywordTok{new} \FunctionTok{SigningCertificateV2}\OperatorTok{(}
                            \KeywordTok{new}\NormalTok{ ESSCertIDv2}\OperatorTok{[]\{}\NormalTok{essCertId}\OperatorTok{\});}
\NormalTok{            signedAttrs}\OperatorTok{.}\FunctionTok{add}\OperatorTok{(}\KeywordTok{new} \BuiltInTok{Attribute}\OperatorTok{(}
\NormalTok{                    PKCSObjectIdentifiers}
                            \OperatorTok{.}\FunctionTok{id\_aa\_signingCertificateV2}\OperatorTok{,}
                    \KeywordTok{new} \FunctionTok{DERSet}\OperatorTok{(}\NormalTok{sigCertV2}\OperatorTok{)));}

            \CommentTok{// 2e. DER{-}encode the signed attributes}
\NormalTok{            DERSet signedAttrSet }\OperatorTok{=} \KeywordTok{new} \FunctionTok{DERSet}\OperatorTok{(}\NormalTok{signedAttrs}\OperatorTok{);}
            \DataTypeTok{byte}\OperatorTok{[]}\NormalTok{ signedAttrBytes }\OperatorTok{=}\NormalTok{ signedAttrSet}\OperatorTok{.}\FunctionTok{getEncoded}\OperatorTok{();}

            \CommentTok{// 2f. Hash the signed attributes — THIS is what we}
            \CommentTok{//     send to the remote signing server}
            \DataTypeTok{byte}\OperatorTok{[]}\NormalTok{ signedAttrHash }\OperatorTok{=} \BuiltInTok{MessageDigest}
                    \OperatorTok{.}\FunctionTok{getInstance}\OperatorTok{(}\StringTok{"SHA{-}256"}\OperatorTok{)}
                    \OperatorTok{.}\FunctionTok{digest}\OperatorTok{(}\NormalTok{signedAttrBytes}\OperatorTok{);}

            \CommentTok{// {-}{-}{-} Phase 2: Remote signing via CSC API {-}{-}{-}}

            \CommentTok{// 3a. Authorize the credential (SCAL2)}
            \BuiltInTok{String}\NormalTok{ sad }\OperatorTok{=} \FunctionTok{authorizeCredential}\OperatorTok{(}
\NormalTok{                    cscBaseUrl}\OperatorTok{,}\NormalTok{ accessToken}\OperatorTok{,}\NormalTok{ credentialId}\OperatorTok{,} \DecValTok{1}\OperatorTok{);}

            \CommentTok{// 3b. Sign the hash remotely}
            \BuiltInTok{List}\OperatorTok{\textless{}}\DataTypeTok{byte}\OperatorTok{[]\textgreater{}}\NormalTok{ signatures }\OperatorTok{=} \FunctionTok{signHashes}\OperatorTok{(}
\NormalTok{                    cscBaseUrl}\OperatorTok{,}\NormalTok{ accessToken}\OperatorTok{,}\NormalTok{ credentialId}\OperatorTok{,}\NormalTok{ sad}\OperatorTok{,}
                    \BuiltInTok{List}\OperatorTok{.}\FunctionTok{of}\OperatorTok{(}\NormalTok{signedAttrHash}\OperatorTok{));}
            \DataTypeTok{byte}\OperatorTok{[]}\NormalTok{ signatureValue }\OperatorTok{=}\NormalTok{ signatures}\OperatorTok{.}\FunctionTok{get}\OperatorTok{(}\DecValTok{0}\OperatorTok{);}

            \CommentTok{// {-}{-}{-} Phase 3: Assemble the CMS and add timestamp {-}{-}{-}}

            \CommentTok{// 4a. Build the complete CMS SignedData manually}
            \CommentTok{//     (embedding the remote signature value)}
            \DataTypeTok{byte}\OperatorTok{[]}\NormalTok{ cmsBytes }\OperatorTok{=} \FunctionTok{assembleCmsSignedData}\OperatorTok{(}
\NormalTok{                    signedAttrSet}\OperatorTok{,}\NormalTok{ signatureValue}\OperatorTok{,}
\NormalTok{                    certChain}\OperatorTok{,}\NormalTok{ signerCert}\OperatorTok{);}

            \CommentTok{// 4b. Request a timestamp (for PAdES{-}B{-}T)}
            \ControlFlowTok{if} \OperatorTok{(}\NormalTok{tsaUrl }\OperatorTok{!=} \KeywordTok{null}\OperatorTok{)} \OperatorTok{\{}
\NormalTok{                cmsBytes }\OperatorTok{=} \FunctionTok{addSignatureTimestamp}\OperatorTok{(}
\NormalTok{                        cmsBytes}\OperatorTok{,}\NormalTok{ signatureValue}\OperatorTok{,}\NormalTok{ tsaUrl}\OperatorTok{);}
            \OperatorTok{\}}

            \ControlFlowTok{return}\NormalTok{ cmsBytes}\OperatorTok{;}

        \OperatorTok{\}} \ControlFlowTok{catch} \OperatorTok{(}\BuiltInTok{Exception}\NormalTok{ e}\OperatorTok{)} \OperatorTok{\{}
            \ControlFlowTok{throw} \KeywordTok{new} \BuiltInTok{IOException}\OperatorTok{(}\StringTok{"Remote signing failed"}\OperatorTok{,}\NormalTok{ e}\OperatorTok{);}
        \OperatorTok{\}}
    \OperatorTok{\}}
\OperatorTok{\}}
\end{Highlighting}
\end{Shaded}

The code above is a simplified illustration of the remote signing flow.
A production implementation must also handle:

\begin{itemize}
\item
  \textbf{Error recovery}: What happens if the CSC API call fails after
  SAD authorization? The SAD may be consumed. You may need to
  re-authorize.
\item
  \textbf{Timeout handling}: Network calls to the CSC API and TSA can
  fail. Implement timeouts and retries.
\item
  \textbf{Content size estimation}: PDFBox's \texttt{addSignature}
  requires you to pre-estimate the size of the \texttt{/Contents} value.
  If you underestimate, the CMS will not fit. Set
  \texttt{SignatureOptions.setPreferredSignatureSize()} generously
  (e.g., 16384 bytes).
\item
  \textbf{Thread safety}: If multiple users sign concurrently, your CSC
  API client must handle concurrent requests correctly. Use separate
  HTTP connections per request.
\end{itemize}

For production remote signing with 3.0.0, consider using the EU DSS
framework (Chapter 23), which provides a complete, tested implementation
of the remote signing flow with CSC API support.

\section{Security Considerations}\label{_security_considerations}

\subsection{Transport Security}\label{_transport_security}

All communication with the CSC API must use TLS 1.2 or later. The access
token is a bearer token --- anyone who intercepts it can use it to sign.
In production:

\begin{itemize}
\item
  Pin the CSC API server's TLS certificate (or at least its root CA) to
  prevent man-in-the-middle attacks
\item
  Use short-lived access tokens (minutes, not hours)
\item
  Transmit tokens only in the \texttt{Authorization} header, never in
  query parameters or request bodies
\item
  Log all API calls for audit purposes, but never log access tokens or
  SAD values
\end{itemize}

\subsection{Hash Manipulation}\label{_hash_manipulation}

The remote signing server signs whatever hash you send it. It does not
--- and cannot --- verify that the hash corresponds to any particular
document. This means: if an attacker compromises the client application,
they can send arbitrary hashes to the server and obtain valid signatures
over arbitrary data.

Mitigations: - \textbf{SCAL2 authorization}: The signer must actively
authorize each signing operation, limiting the window for automated
abuse - \textbf{Audit logging}: Every signing operation is logged on the
server, creating a forensic trail - \textbf{Rate limiting}: Unusual
signing patterns (hundreds of signatures per minute from a single user)
should trigger alerts - \textbf{Document description in SAD}: Some CSC
implementations allow including a document description in the
authorization request, which the signer can verify before approving

\subsection{Key Compromise}\label{_key_compromise}

If an HSM is compromised, all keys stored in it are compromised. This is
the catastrophic failure mode of remote signing --- a single breach
affects all users.

HSM-level mitigations: - Physical security (data center access controls,
tamper detection) - Multi-person key ceremony for HSM initialization -
Regular firmware updates - Separate HSMs for different security levels
(production vs. test) - Geographic redundancy with key replication

Operational mitigations: - Certificate revocation within minutes of
detection - Automatic alerts for anomalous signing patterns - Regular
penetration testing of the signing infrastructure - Insurance against
breach liability (some QTSPs carry specific cyber insurance)

\section{Key Takeaways}\label{_key_takeaways_10}

\begin{itemize}
\item
  \textbf{Remote signing} moves the private key from the signer's device
  to a server-side HSM operated by a QTSP --- the signer authenticates
  remotely, the server signs
\item
  The \textbf{hash-signing model} preserves document confidentiality:
  only a 32-byte hash crosses the network, never the document content
\item
  The \textbf{CSC API v2.0} standardizes remote signing: OAuth 2.0
  authentication, credential discovery, signature activation, hash
  signing
\item
  \textbf{SCAL1} provides session-level sole control (suitable for
  AdES); \textbf{SCAL2} provides per-signature authorization (required
  for QES)
\item
  The \textbf{SAD token} is the proof that the signer authorized a
  specific signing operation --- it has limited lifetime and use count
\item
  \textbf{HSMs} store keys and perform cryptographic operations;
  \textbf{SAMs} enforce sole control and verify SAD tokens
\item
  PKCS\#11 is the standard interface between the signing server and the
  HSM --- vendor-specific libraries are the most common deployment issue
\item
  Major providers (Swisscom AIS, DigiCert/QuoVadis, InfoCert, Namirial,
  ZealID, D-Trust) all offer CSC-compliant APIs with varying degrees of
  completeness
\item
  \textbf{Always test against the specific provider's sandbox} --- CSC
  API implementations have provider-specific behaviors
\item
  Many providers run \textbf{Ascertia ADSS} as their CSC platform --- if
  you've tested one ADSS deployment, the behavior will be identical
  across all of them
\item
  \textbf{National eID integration} (SPID, BankID, SwissID) routes
  authentication through the national identity provider via OpenID
  Connect --- not a simple username/password flow
\item
  The CSC \texttt{/info} endpoint enables \textbf{endpoint discovery}
  --- it responds without authentication and reveals the provider's
  capabilities and supported methods
\item
  For production PDF signing via remote services, consider the EU DSS
  framework (Chapter 23), which provides tested CSC API integration
\end{itemize}

\chapter{The EU DSS Framework}\label{_the_eu_dss_framework}

Throughout this book, we have been building PDF signing and validation
code from the ground up --- constructing CMS structures by hand with
BouncyCastle, computing byte ranges manually, assembling DSS
dictionaries one stream at a time. This is the right approach for
understanding what is happening, but it is the wrong approach for
production software.

The \textbf{European Digital Signature Services (EU DSS)} framework is
the European Commission's open-source reference implementation for
digital signature creation, extension, and validation. It is maintained
by the European Commission's ISA2 programme (now the Digital Europe
Programme), developed primarily by the Luxembourg company Nowina
Solutions, and used by dozens of EU member states in their national
validation services. It handles PAdES, CAdES, XAdES, and JAdES. It
validates against the EU Trusted List. It implements every PAdES profile
from B-B through B-LTA. It supports remote signing via the CSC API. And
it is free --- released under the LGPL 2.1 license.

If you are building production software that creates or validates
digital signatures in a European context, EU DSS is not just an option
--- it is the benchmark against which your implementation will be
measured. This chapter covers everything you need to use it:
architecture, core modules, the signing flow, validation, Spring Boot
integration, and the honest trade-offs between using EU DSS and rolling
your own.

\section{History and Governance}\label{_history_and_governance}

\subsection{Origins}\label{_origins}

The EU DSS project began in 2011 as part of the European Commission's
\textbf{ISA programme} (Interoperability Solutions for European Public
Administrations). The goal was to provide member states with a reusable
software library that correctly implemented the ETSI signature standards
--- so that each country would not have to write its own implementation
from scratch and risk incompatibilities.

The project was motivated by the same cross-border recognition failures
that led to eIDAS itself (Chapter 19). Even when countries agreed on the
legal framework, their software implementations often produced or
validated signatures differently. A common reference implementation was
the practical answer to "why doesn't my signature validate in your
system?"

\subsection{Development}\label{_development}

The framework has been developed primarily by \textbf{Nowina Solutions}
(formerly Nowina Technologies), a Luxembourg-based company specializing
in electronic signatures and trust services. Nowina won the European
Commission's procurement for DSS development and has been the principal
technical contributor since 2015.

The source code is hosted on the European Commission's \textbf{GitHub
organization} at \texttt{github.com/esig/dss} and on the Commission's
own code repository at
\texttt{ec.europa.eu/digital-building-blocks/code/projects/ESIG}.

\subsection{Version History}\label{_version_history}

\begin{longtable}[]{@{}
  >{\raggedright\arraybackslash}p{(\linewidth - 4\tabcolsep) * \real{0.2000}}
  >{\raggedright\arraybackslash}p{(\linewidth - 4\tabcolsep) * \real{0.2000}}
  >{\raggedright\arraybackslash}p{(\linewidth - 4\tabcolsep) * \real{0.6000}}@{}}
\caption{EU DSS Major Version History}\tabularnewline
\toprule\noalign{}
\begin{minipage}[b]{\linewidth}\raggedright
Version
\end{minipage} & \begin{minipage}[b]{\linewidth}\raggedright
Date
\end{minipage} & \begin{minipage}[b]{\linewidth}\raggedright
Key Changes
\end{minipage} \\
\midrule\noalign{}
\endfirsthead
\toprule\noalign{}
\begin{minipage}[b]{\linewidth}\raggedright
Version
\end{minipage} & \begin{minipage}[b]{\linewidth}\raggedright
Date
\end{minipage} & \begin{minipage}[b]{\linewidth}\raggedright
Key Changes
\end{minipage} \\
\midrule\noalign{}
\endhead
\bottomrule\noalign{}
\endlastfoot
4.x & 2015-2017 & First widely adopted version. PAdES, CAdES, XAdES
support. Java 7/8. \\
5.0 & 2018 & Major rewrite. Modular architecture (separate Maven
modules). Java 8 minimum. Introduced pluggable token architecture. \\
5.6 & 2020 & JAdES (JSON Advanced Electronic Signatures) support. ETSI
EN 319 102-1 validation policy. Improved EU Trusted List parsing. \\
5.10 & 2022 & Java 11 minimum. CSC API support for remote signing.
Updated BouncyCastle dependency (1.70+). \\
5.13 & 2023 & Improved PKCS\#11 token support. PSD2 (Payment Services
Directive) certificate validation. Updated CRL/OCSP handling. \\
6.0 & 2024 & Java 17 minimum. Module system (JPMS) support. Restructured
Maven coordinates (groupId changed from
\texttt{eu.europa.ec.joinup.sd-dss} to \texttt{eu.europa.dss}). Updated
1.80 dependency. \\
6.2 & 2025 & Java 17+. Enhanced EUDI wallet integration. Improved TSL
(Trusted List) refresh mechanisms. This is the current stable release as
of early 2026. \\
\end{longtable}

The groupId change in version 6.0 is a common source of confusion. If
your project depends on \texttt{eu.europa.ec.joinup.sd-dss:dss-pades},
that is DSS 5.x. If it depends on \texttt{eu.europa.dss:dss-pades}, that
is DSS 6.x. The two groupIds are not interchangeable --- do not mix
modules from different major versions.

\section{Architecture}\label{_architecture_2}

\subsection{Design Principles}\label{_design_principles}

EU DSS is built around three core architectural principles:

\begin{enumerate}
\def\labelenumi{\arabic{enumi}.}
\item
  \textbf{Format independence}: The same signing and validation workflow
  works for PAdES (PDF), CAdES (CMS), XAdES (XML), and JAdES (JSON). The
  format-specific logic is encapsulated in format modules; the core
  signing flow is shared.
\item
  \textbf{Pluggable token architecture}: The private key can come from
  anywhere --- a PKCS\#12 file, a PKCS\#11 smart card, a Microsoft CNG
  keystore, a remote CSC API service, or a custom
  \texttt{SignatureTokenConnection} implementation. The signing flow
  does not need to know where the key is.
\item
  \textbf{Policy-based validation}: Validation is driven by a
  configurable policy (an XML file) that defines which checks to perform
  and what their outcomes mean. You can use the default ETSI validation
  policy, or you can define your own --- accepting certain
  non-conformities for your use case while rejecting others.
\end{enumerate}

\subsection{Module Overview}\label{_module_overview}

EU DSS is organized into Maven modules, each with a specific
responsibility. You do not need all of them --- include only what you
use.

\begin{longtable}[]{@{}
  >{\raggedright\arraybackslash}p{(\linewidth - 2\tabcolsep) * \real{0.3333}}
  >{\raggedright\arraybackslash}p{(\linewidth - 2\tabcolsep) * \real{0.6667}}@{}}
\caption{Core EU DSS Modules}\tabularnewline
\toprule\noalign{}
\begin{minipage}[b]{\linewidth}\raggedright
Module
\end{minipage} & \begin{minipage}[b]{\linewidth}\raggedright
Purpose
\end{minipage} \\
\midrule\noalign{}
\endfirsthead
\toprule\noalign{}
\begin{minipage}[b]{\linewidth}\raggedright
Module
\end{minipage} & \begin{minipage}[b]{\linewidth}\raggedright
Purpose
\end{minipage} \\
\midrule\noalign{}
\endhead
\bottomrule\noalign{}
\endlastfoot
\texttt{dss-spi} & Service Provider Interface. Core abstractions:
\texttt{DSSDocument}, \texttt{SignatureTokenConnection},
\texttt{CertificateToken}, \texttt{DSSPrivateKeyEntry}. All other
modules depend on this. \\
\texttt{dss-model} & Data model classes: \texttt{SignatureLevel},
\texttt{DigestAlgorithm}, \texttt{EncryptionAlgorithm},
\texttt{TimestampType}, etc. Enumerations and DTOs used across the
framework. \\
\texttt{dss-token} & Token implementations:
\texttt{Pkcs12SignatureToken}, \texttt{Pkcs11SignatureToken},
\texttt{MSCAPISignatureToken}, \texttt{KeyStoreSignatureToken}. These
connect to key material. \\
\texttt{dss-pades} & PAdES (PDF) signing and validation. Depends on
Apache PDFBox internally. Produces and validates PAdES-B-B through
B-LTA. \\
\texttt{dss-cades} & CAdES (CMS) signing and validation. Used internally
by \texttt{dss-pades} for the CMS structures inside PDF signatures. \\
\texttt{dss-xades} & XAdES (XML) signing and validation. Not needed for
PDF work unless you also handle XML signatures. \\
\texttt{dss-jades} & JAdES (JSON) signing and validation. The newest
format support. \\
\texttt{dss-tsl-validation} & EU Trusted List (TSL/EUTL) parsing,
validation, and caching. Downloads the trusted lists, verifies their
signatures, and builds a trust store that other modules use for
validation. \\
\texttt{dss-crl-parser} & CRL downloading and parsing. \\
\texttt{dss-ocsp-parser} & OCSP response fetching and parsing. \\
\texttt{dss-service} & Higher-level service classes:
\texttt{OnlineOCSPSource}, \texttt{OnlineCRLSource},
\texttt{OnlineTSPSource}, \texttt{CommonsDataLoader} (HTTP client).
Implements the data-fetching operations that validation requires. \\
\texttt{dss-document} & Document-level abstractions:
\texttt{SignatureParameters}, \texttt{ToBeSigned},
\texttt{SignatureValue}, \texttt{SignedDocumentValidator}. The signing
flow API lives here. \\
\texttt{dss-diagnostic-data-jaxb} & JAXB model for the diagnostic data
XML that validation produces. Useful for detailed validation reports. \\
\texttt{dss-detailed-report-jaxb} & JAXB model for the detailed
validation report. Human-readable validation results. \\
\texttt{dss-simple-report-jaxb} & JAXB model for the simple validation
report. Summary-level pass/fail results. \\
\end{longtable}

\subsection{Maven Coordinates}\label{_maven_coordinates}

For a typical PAdES signing and validation project with EU DSS 6.2, you
need:

\begin{Shaded}
\begin{Highlighting}[]
\NormalTok{\textless{}}\KeywordTok{properties}\NormalTok{\textgreater{}}
\NormalTok{    \textless{}}\KeywordTok{dss.version}\NormalTok{\textgreater{}6.2\textless{}/}\KeywordTok{dss.version}\NormalTok{\textgreater{}}
\NormalTok{\textless{}/}\KeywordTok{properties}\NormalTok{\textgreater{}}

\NormalTok{\textless{}}\KeywordTok{dependencies}\NormalTok{\textgreater{}}
    \CommentTok{\textless{}!{-}{-} Core signing and validation for PAdES {-}{-}\textgreater{}}
\NormalTok{    \textless{}}\KeywordTok{dependency}\NormalTok{\textgreater{}}
\NormalTok{        \textless{}}\KeywordTok{groupId}\NormalTok{\textgreater{}eu.europa.dss\textless{}/}\KeywordTok{groupId}\NormalTok{\textgreater{}}
\NormalTok{        \textless{}}\KeywordTok{artifactId}\NormalTok{\textgreater{}dss{-}pades\textless{}/}\KeywordTok{artifactId}\NormalTok{\textgreater{}}
\NormalTok{        \textless{}}\KeywordTok{version}\NormalTok{\textgreater{}$\{dss.version\}\textless{}/}\KeywordTok{version}\NormalTok{\textgreater{}}
\NormalTok{    \textless{}/}\KeywordTok{dependency}\NormalTok{\textgreater{}}

    \CommentTok{\textless{}!{-}{-} Token support (PKCS\#12, PKCS\#11) {-}{-}\textgreater{}}
\NormalTok{    \textless{}}\KeywordTok{dependency}\NormalTok{\textgreater{}}
\NormalTok{        \textless{}}\KeywordTok{groupId}\NormalTok{\textgreater{}eu.europa.dss\textless{}/}\KeywordTok{groupId}\NormalTok{\textgreater{}}
\NormalTok{        \textless{}}\KeywordTok{artifactId}\NormalTok{\textgreater{}dss{-}token\textless{}/}\KeywordTok{artifactId}\NormalTok{\textgreater{}}
\NormalTok{        \textless{}}\KeywordTok{version}\NormalTok{\textgreater{}$\{dss.version\}\textless{}/}\KeywordTok{version}\NormalTok{\textgreater{}}
\NormalTok{    \textless{}/}\KeywordTok{dependency}\NormalTok{\textgreater{}}

    \CommentTok{\textless{}!{-}{-} Online sources for OCSP, CRL, TSA {-}{-}\textgreater{}}
\NormalTok{    \textless{}}\KeywordTok{dependency}\NormalTok{\textgreater{}}
\NormalTok{        \textless{}}\KeywordTok{groupId}\NormalTok{\textgreater{}eu.europa.dss\textless{}/}\KeywordTok{groupId}\NormalTok{\textgreater{}}
\NormalTok{        \textless{}}\KeywordTok{artifactId}\NormalTok{\textgreater{}dss{-}service\textless{}/}\KeywordTok{artifactId}\NormalTok{\textgreater{}}
\NormalTok{        \textless{}}\KeywordTok{version}\NormalTok{\textgreater{}$\{dss.version\}\textless{}/}\KeywordTok{version}\NormalTok{\textgreater{}}
\NormalTok{    \textless{}/}\KeywordTok{dependency}\NormalTok{\textgreater{}}

    \CommentTok{\textless{}!{-}{-} EU Trusted List validation {-}{-}\textgreater{}}
\NormalTok{    \textless{}}\KeywordTok{dependency}\NormalTok{\textgreater{}}
\NormalTok{        \textless{}}\KeywordTok{groupId}\NormalTok{\textgreater{}eu.europa.dss\textless{}/}\KeywordTok{groupId}\NormalTok{\textgreater{}}
\NormalTok{        \textless{}}\KeywordTok{artifactId}\NormalTok{\textgreater{}dss{-}tsl{-}validation\textless{}/}\KeywordTok{artifactId}\NormalTok{\textgreater{}}
\NormalTok{        \textless{}}\KeywordTok{version}\NormalTok{\textgreater{}$\{dss.version\}\textless{}/}\KeywordTok{version}\NormalTok{\textgreater{}}
\NormalTok{    \textless{}/}\KeywordTok{dependency}\NormalTok{\textgreater{}}
\NormalTok{\textless{}/}\KeywordTok{dependencies}\NormalTok{\textgreater{}}
\end{Highlighting}
\end{Shaded}

EU DSS depends on 3.0.0 internally for PDF operations. If your project
also uses PDFBox directly, ensure version alignment --- having two
different PDFBox versions on the classpath causes unpredictable
behavior. Check the DSS BOM (Bill of Materials) for the exact PDFBox
version it expects.

\section{The Signing Flow: ToBeSigned /
SignatureValue}\label{_the_signing_flow_tobesigned_signaturevalue}

\subsection{The Two-Phase Signing
Model}\label{_the_two_phase_signing_model}

EU DSS uses a two-phase signing model that maps directly to the
hash-signing model described in Chapter 22. This separation is what
makes remote signing possible --- the private key operation can happen
anywhere, including on a remote HSM via the CSC API.

\textbf{Phase 1: Compute what needs to be signed}

\begin{Shaded}
\begin{Highlighting}[]
\CommentTok{// 1. Create the PAdES service}
\NormalTok{PAdESService service }\OperatorTok{=} \KeywordTok{new} \FunctionTok{PAdESService}\OperatorTok{(}\NormalTok{certificateVerifier}\OperatorTok{);}

\CommentTok{// 2. Set the TSA source (for B{-}T and above)}
\NormalTok{service}\OperatorTok{.}\FunctionTok{setTspSource}\OperatorTok{(}\NormalTok{onlineTSPSource}\OperatorTok{);}

\CommentTok{// 3. Configure signature parameters}
\NormalTok{PAdESSignatureParameters parameters }\OperatorTok{=} \KeywordTok{new} \FunctionTok{PAdESSignatureParameters}\OperatorTok{();}
\NormalTok{parameters}\OperatorTok{.}\FunctionTok{setSignatureLevel}\OperatorTok{(}\NormalTok{SignatureLevel}\OperatorTok{.}\FunctionTok{PAdES\_BASELINE\_T}\OperatorTok{);}
\NormalTok{parameters}\OperatorTok{.}\FunctionTok{setDigestAlgorithm}\OperatorTok{(}\NormalTok{DigestAlgorithm}\OperatorTok{.}\FunctionTok{SHA256}\OperatorTok{);}
\NormalTok{parameters}\OperatorTok{.}\FunctionTok{setSigningCertificate}\OperatorTok{(}\KeywordTok{new} \FunctionTok{CertificateToken}\OperatorTok{(}\NormalTok{signingCert}\OperatorTok{));}
\NormalTok{parameters}\OperatorTok{.}\FunctionTok{setCertificateChain}\OperatorTok{(}\NormalTok{certificateChainTokens}\OperatorTok{);}

\CommentTok{// 4. Get the data to be signed (the hash)}
\NormalTok{DSSDocument toSignDocument }\OperatorTok{=} \KeywordTok{new} \FunctionTok{InMemoryDocument}\OperatorTok{(}\NormalTok{pdfBytes}\OperatorTok{);}
\NormalTok{ToBeSigned dataToSign }\OperatorTok{=}\NormalTok{ service}\OperatorTok{.}\FunctionTok{getDataToSign}\OperatorTok{(}
\NormalTok{        toSignDocument}\OperatorTok{,}\NormalTok{ parameters}\OperatorTok{);}
\end{Highlighting}
\end{Shaded}

The \texttt{ToBeSigned} object contains the bytes that must be signed
--- specifically, the DER-encoded signed attributes of the CMS
SignerInfo, exactly as described in Chapter 22's hash-signing flow. EU
DSS has already computed the message digest over the PDF byte range,
built the ESS-signing-certificate-v2 attribute, and assembled the signed
attributes.

\textbf{Phase 2: Sign the hash and finalize}

\begin{Shaded}
\begin{Highlighting}[]
\CommentTok{// 5. Sign the hash (locally or remotely)}
\CommentTok{//    This is where the private key is used}
\NormalTok{SignatureValue signatureValue }\OperatorTok{=}\NormalTok{ signingToken}\OperatorTok{.}\FunctionTok{sign}\OperatorTok{(}
\NormalTok{        dataToSign}\OperatorTok{,}\NormalTok{ parameters}\OperatorTok{.}\FunctionTok{getDigestAlgorithm}\OperatorTok{(),}
\NormalTok{        privateKeyEntry}\OperatorTok{);}

\CommentTok{// 6. Finalize the signature — embed in PDF}
\NormalTok{DSSDocument signedDocument }\OperatorTok{=}\NormalTok{ service}\OperatorTok{.}\FunctionTok{signDocument}\OperatorTok{(}
\NormalTok{        toSignDocument}\OperatorTok{,}\NormalTok{ parameters}\OperatorTok{,}\NormalTok{ signatureValue}\OperatorTok{);}

\CommentTok{// 7. Write the result}
\NormalTok{signedDocument}\OperatorTok{.}\FunctionTok{writeTo}\OperatorTok{(}\KeywordTok{new} \BuiltInTok{FileOutputStream}\OperatorTok{(}\StringTok{"signed.pdf"}\OperatorTok{));}
\end{Highlighting}
\end{Shaded}

The beauty of this design: step 5 is the \emph{only} step that touches
the private key. Everything else --- PDF parsing, CMS construction,
timestamp embedding, incremental save --- is handled by EU DSS. If you
want to use a remote signing service instead of a local key, you replace
step 5 with a CSC API call and everything else stays the same.

\subsection{SignatureTokenConnection
Implementations}\label{_signaturetokenconnection_implementations}

The \texttt{SignatureTokenConnection} interface abstracts access to the
private key. EU DSS ships with several implementations:

\begin{Shaded}
\begin{Highlighting}[]
\CommentTok{// PKCS\#12 file (local keystore)}
\NormalTok{Pkcs12SignatureToken token }\OperatorTok{=} \KeywordTok{new} \FunctionTok{Pkcs12SignatureToken}\OperatorTok{(}
        \KeywordTok{new} \BuiltInTok{FileInputStream}\OperatorTok{(}\StringTok{"keystore.p12"}\OperatorTok{),}
        \KeywordTok{new} \BuiltInTok{KeyStore}\OperatorTok{.}\FunctionTok{PasswordProtection}\OperatorTok{(}\StringTok{"password"}\OperatorTok{.}\FunctionTok{toCharArray}\OperatorTok{()));}
\BuiltInTok{List}\OperatorTok{\textless{}}\NormalTok{DSSPrivateKeyEntry}\OperatorTok{\textgreater{}}\NormalTok{ keys }\OperatorTok{=}\NormalTok{ token}\OperatorTok{.}\FunctionTok{getKeys}\OperatorTok{();}
\NormalTok{DSSPrivateKeyEntry privateKey }\OperatorTok{=}\NormalTok{ keys}\OperatorTok{.}\FunctionTok{get}\OperatorTok{(}\DecValTok{0}\OperatorTok{);}

\CommentTok{// PKCS\#11 smart card / HSM}
\NormalTok{Pkcs11SignatureToken pkcs11Token }\OperatorTok{=} \KeywordTok{new} \FunctionTok{Pkcs11SignatureToken}\OperatorTok{(}
        \StringTok{"/usr/lib/opensc{-}pkcs11.so"}\OperatorTok{,}  \CommentTok{// PKCS\#11 library path}
        \KeywordTok{new} \BuiltInTok{KeyStore}\OperatorTok{.}\FunctionTok{PasswordProtection}\OperatorTok{(}\StringTok{"pin"}\OperatorTok{.}\FunctionTok{toCharArray}\OperatorTok{()),}
        \DecValTok{0}\OperatorTok{);}  \CommentTok{// slot index}
\BuiltInTok{List}\OperatorTok{\textless{}}\NormalTok{DSSPrivateKeyEntry}\OperatorTok{\textgreater{}}\NormalTok{ smartCardKeys }\OperatorTok{=}\NormalTok{ pkcs11Token}\OperatorTok{.}\FunctionTok{getKeys}\OperatorTok{();}

\CommentTok{// Microsoft CNG (Windows certificate store)}
\NormalTok{MSCAPISignatureToken mscapiToken }\OperatorTok{=} \KeywordTok{new} \FunctionTok{MSCAPISignatureToken}\OperatorTok{();}
\BuiltInTok{List}\OperatorTok{\textless{}}\NormalTok{DSSPrivateKeyEntry}\OperatorTok{\textgreater{}}\NormalTok{ windowsKeys }\OperatorTok{=}\NormalTok{ mscapiToken}\OperatorTok{.}\FunctionTok{getKeys}\OperatorTok{();}
\end{Highlighting}
\end{Shaded}

For remote signing via CSC API, you implement
\texttt{SignatureTokenConnection} with HTTP calls to the CSC endpoints.
The \texttt{sign()} method sends the hash to
\texttt{/signatures/signHash} and returns the signature value.

\section{Complete PAdES-B-T Signing
Example}\label{_complete_pades_b_t_signing_example}

Here is a complete, working example that signs a PDF with PAdES-B-T
using a PKCS\#12 keystore and a TSA:

\begin{Shaded}
\begin{Highlighting}[]
\KeywordTok{import} \ImportTok{eu}\OperatorTok{.}\ImportTok{europa}\OperatorTok{.}\ImportTok{dss}\OperatorTok{.}\ImportTok{cades}\OperatorTok{.}\ImportTok{signature}\OperatorTok{.}\ImportTok{CAdESService}\OperatorTok{;}
\KeywordTok{import} \ImportTok{eu}\OperatorTok{.}\ImportTok{europa}\OperatorTok{.}\ImportTok{dss}\OperatorTok{.}\ImportTok{enumerations}\OperatorTok{.}\ImportTok{DigestAlgorithm}\OperatorTok{;}
\KeywordTok{import} \ImportTok{eu}\OperatorTok{.}\ImportTok{europa}\OperatorTok{.}\ImportTok{dss}\OperatorTok{.}\ImportTok{enumerations}\OperatorTok{.}\ImportTok{SignatureLevel}\OperatorTok{;}
\KeywordTok{import} \ImportTok{eu}\OperatorTok{.}\ImportTok{europa}\OperatorTok{.}\ImportTok{dss}\OperatorTok{.}\ImportTok{model}\OperatorTok{.}\ImportTok{DSSDocument}\OperatorTok{;}
\KeywordTok{import} \ImportTok{eu}\OperatorTok{.}\ImportTok{europa}\OperatorTok{.}\ImportTok{dss}\OperatorTok{.}\ImportTok{model}\OperatorTok{.}\ImportTok{InMemoryDocument}\OperatorTok{;}
\KeywordTok{import} \ImportTok{eu}\OperatorTok{.}\ImportTok{europa}\OperatorTok{.}\ImportTok{dss}\OperatorTok{.}\ImportTok{model}\OperatorTok{.}\ImportTok{SignatureValue}\OperatorTok{;}
\KeywordTok{import} \ImportTok{eu}\OperatorTok{.}\ImportTok{europa}\OperatorTok{.}\ImportTok{dss}\OperatorTok{.}\ImportTok{model}\OperatorTok{.}\ImportTok{ToBeSigned}\OperatorTok{;}
\KeywordTok{import} \ImportTok{eu}\OperatorTok{.}\ImportTok{europa}\OperatorTok{.}\ImportTok{dss}\OperatorTok{.}\ImportTok{pades}\OperatorTok{.}\ImportTok{PAdESSignatureParameters}\OperatorTok{;}
\KeywordTok{import} \ImportTok{eu}\OperatorTok{.}\ImportTok{europa}\OperatorTok{.}\ImportTok{dss}\OperatorTok{.}\ImportTok{pades}\OperatorTok{.}\ImportTok{signature}\OperatorTok{.}\ImportTok{PAdESService}\OperatorTok{;}
\KeywordTok{import} \ImportTok{eu}\OperatorTok{.}\ImportTok{europa}\OperatorTok{.}\ImportTok{dss}\OperatorTok{.}\ImportTok{service}\OperatorTok{.}\ImportTok{http}\OperatorTok{.}\ImportTok{commons}\OperatorTok{.}\ImportTok{CommonsDataLoader}\OperatorTok{;}
\KeywordTok{import} \ImportTok{eu}\OperatorTok{.}\ImportTok{europa}\OperatorTok{.}\ImportTok{dss}\OperatorTok{.}\ImportTok{service}\OperatorTok{.}\ImportTok{http}\OperatorTok{.}\ImportTok{commons}\OperatorTok{.}\ImportTok{OCSPDataLoader}\OperatorTok{;}
\KeywordTok{import} \ImportTok{eu}\OperatorTok{.}\ImportTok{europa}\OperatorTok{.}\ImportTok{dss}\OperatorTok{.}\ImportTok{service}\OperatorTok{.}\ImportTok{ocsp}\OperatorTok{.}\ImportTok{OnlineOCSPSource}\OperatorTok{;}
\KeywordTok{import} \ImportTok{eu}\OperatorTok{.}\ImportTok{europa}\OperatorTok{.}\ImportTok{dss}\OperatorTok{.}\ImportTok{service}\OperatorTok{.}\ImportTok{crl}\OperatorTok{.}\ImportTok{OnlineCRLSource}\OperatorTok{;}
\KeywordTok{import} \ImportTok{eu}\OperatorTok{.}\ImportTok{europa}\OperatorTok{.}\ImportTok{dss}\OperatorTok{.}\ImportTok{service}\OperatorTok{.}\ImportTok{tsp}\OperatorTok{.}\ImportTok{OnlineTSPSource}\OperatorTok{;}
\KeywordTok{import} \ImportTok{eu}\OperatorTok{.}\ImportTok{europa}\OperatorTok{.}\ImportTok{dss}\OperatorTok{.}\ImportTok{spi}\OperatorTok{.}\ImportTok{x509}\OperatorTok{.}\ImportTok{CertificateSource}\OperatorTok{;}
\KeywordTok{import} \ImportTok{eu}\OperatorTok{.}\ImportTok{europa}\OperatorTok{.}\ImportTok{dss}\OperatorTok{.}\ImportTok{spi}\OperatorTok{.}\ImportTok{x509}\OperatorTok{.}\ImportTok{CommonTrustedCertificateSource}\OperatorTok{;}
\KeywordTok{import} \ImportTok{eu}\OperatorTok{.}\ImportTok{europa}\OperatorTok{.}\ImportTok{dss}\OperatorTok{.}\ImportTok{token}\OperatorTok{.}\ImportTok{DSSPrivateKeyEntry}\OperatorTok{;}
\KeywordTok{import} \ImportTok{eu}\OperatorTok{.}\ImportTok{europa}\OperatorTok{.}\ImportTok{dss}\OperatorTok{.}\ImportTok{token}\OperatorTok{.}\ImportTok{Pkcs12SignatureToken}\OperatorTok{;}
\KeywordTok{import} \ImportTok{eu}\OperatorTok{.}\ImportTok{europa}\OperatorTok{.}\ImportTok{dss}\OperatorTok{.}\ImportTok{validation}\OperatorTok{.}\ImportTok{CommonCertificateVerifier}\OperatorTok{;}

\CommentTok{/**}
 \CommentTok{*}\NormalTok{ Complete PAdES}\CommentTok{{-}}\NormalTok{B}\CommentTok{{-}}\NormalTok{T signing with EU DSS}\CommentTok{.}
 \CommentTok{*}\NormalTok{ This example uses a PKCS}\CommentTok{\#12}\NormalTok{ keystore and a public TSA}\CommentTok{.}
 \CommentTok{*/}
\KeywordTok{public} \DataTypeTok{byte}\OperatorTok{[]} \FunctionTok{signWithEuDss}\OperatorTok{(}\DataTypeTok{byte}\OperatorTok{[]}\NormalTok{ pdfBytes}\OperatorTok{,} \BuiltInTok{String}\NormalTok{ keystorePath}\OperatorTok{,}
        \BuiltInTok{String}\NormalTok{ keystorePassword}\OperatorTok{,} \BuiltInTok{String}\NormalTok{ tsaUrl}\OperatorTok{)} \KeywordTok{throws} \BuiltInTok{Exception} \OperatorTok{\{}

    \CommentTok{// 1. Open the PKCS\#12 token}
\NormalTok{    Pkcs12SignatureToken signingToken }\OperatorTok{=} \KeywordTok{new} \FunctionTok{Pkcs12SignatureToken}\OperatorTok{(}
            \KeywordTok{new} \BuiltInTok{FileInputStream}\OperatorTok{(}\NormalTok{keystorePath}\OperatorTok{),}
            \KeywordTok{new} \BuiltInTok{KeyStore}\OperatorTok{.}\FunctionTok{PasswordProtection}\OperatorTok{(}
\NormalTok{                    keystorePassword}\OperatorTok{.}\FunctionTok{toCharArray}\OperatorTok{()));}

    \CommentTok{// 2. Get the private key entry (first key in the keystore)}
    \BuiltInTok{List}\OperatorTok{\textless{}}\NormalTok{DSSPrivateKeyEntry}\OperatorTok{\textgreater{}}\NormalTok{ keys }\OperatorTok{=}\NormalTok{ signingToken}\OperatorTok{.}\FunctionTok{getKeys}\OperatorTok{();}
\NormalTok{    DSSPrivateKeyEntry privateKey }\OperatorTok{=}\NormalTok{ keys}\OperatorTok{.}\FunctionTok{get}\OperatorTok{(}\DecValTok{0}\OperatorTok{);}

    \CommentTok{// 3. Configure the certificate verifier}
    \CommentTok{//    (used during signing to validate the chain)}
\NormalTok{    CommonCertificateVerifier certificateVerifier }\OperatorTok{=}
            \KeywordTok{new} \FunctionTok{CommonCertificateVerifier}\OperatorTok{();}

    \CommentTok{// OCSP source for online revocation checking}
\NormalTok{    OnlineOCSPSource ocspSource }\OperatorTok{=} \KeywordTok{new} \FunctionTok{OnlineOCSPSource}\OperatorTok{();}
\NormalTok{    ocspSource}\OperatorTok{.}\FunctionTok{setDataLoader}\OperatorTok{(}\KeywordTok{new} \FunctionTok{OCSPDataLoader}\OperatorTok{());}
\NormalTok{    certificateVerifier}\OperatorTok{.}\FunctionTok{setOcspSource}\OperatorTok{(}\NormalTok{ocspSource}\OperatorTok{);}

    \CommentTok{// CRL source as fallback}
\NormalTok{    OnlineCRLSource crlSource }\OperatorTok{=} \KeywordTok{new} \FunctionTok{OnlineCRLSource}\OperatorTok{();}
\NormalTok{    crlSource}\OperatorTok{.}\FunctionTok{setDataLoader}\OperatorTok{(}\KeywordTok{new} \FunctionTok{CommonsDataLoader}\OperatorTok{());}
\NormalTok{    certificateVerifier}\OperatorTok{.}\FunctionTok{setCrlSource}\OperatorTok{(}\NormalTok{crlSource}\OperatorTok{);}

    \CommentTok{// 4. Configure the TSA source}
\NormalTok{    OnlineTSPSource tspSource }\OperatorTok{=} \KeywordTok{new} \FunctionTok{OnlineTSPSource}\OperatorTok{(}\NormalTok{tsaUrl}\OperatorTok{);}
\NormalTok{    tspSource}\OperatorTok{.}\FunctionTok{setDataLoader}\OperatorTok{(}\KeywordTok{new} \FunctionTok{CommonsDataLoader}\OperatorTok{());}

    \CommentTok{// 5. Create the PAdES service}
\NormalTok{    PAdESService padesService }\OperatorTok{=} \KeywordTok{new} \FunctionTok{PAdESService}\OperatorTok{(}\NormalTok{certificateVerifier}\OperatorTok{);}
\NormalTok{    padesService}\OperatorTok{.}\FunctionTok{setTspSource}\OperatorTok{(}\NormalTok{tspSource}\OperatorTok{);}

    \CommentTok{// 6. Configure signature parameters}
\NormalTok{    PAdESSignatureParameters parameters }\OperatorTok{=}
            \KeywordTok{new} \FunctionTok{PAdESSignatureParameters}\OperatorTok{();}
\NormalTok{    parameters}\OperatorTok{.}\FunctionTok{setSignatureLevel}\OperatorTok{(}\NormalTok{SignatureLevel}\OperatorTok{.}\FunctionTok{PAdES\_BASELINE\_T}\OperatorTok{);}
\NormalTok{    parameters}\OperatorTok{.}\FunctionTok{setDigestAlgorithm}\OperatorTok{(}\NormalTok{DigestAlgorithm}\OperatorTok{.}\FunctionTok{SHA256}\OperatorTok{);}
\NormalTok{    parameters}\OperatorTok{.}\FunctionTok{setSigningCertificate}\OperatorTok{(}
\NormalTok{            privateKey}\OperatorTok{.}\FunctionTok{getCertificate}\OperatorTok{());}
\NormalTok{    parameters}\OperatorTok{.}\FunctionTok{setCertificateChain}\OperatorTok{(}
\NormalTok{            privateKey}\OperatorTok{.}\FunctionTok{getCertificateChain}\OperatorTok{());}

    \CommentTok{// Optional: set visual signature parameters}
    \CommentTok{// parameters.setSignatureImageParameters(...);}

    \CommentTok{// 7. Phase 1: Get the data to be signed}
\NormalTok{    DSSDocument toSignDocument }\OperatorTok{=} \KeywordTok{new} \FunctionTok{InMemoryDocument}\OperatorTok{(}\NormalTok{pdfBytes}\OperatorTok{);}
\NormalTok{    ToBeSigned dataToSign }\OperatorTok{=}\NormalTok{ padesService}\OperatorTok{.}\FunctionTok{getDataToSign}\OperatorTok{(}
\NormalTok{            toSignDocument}\OperatorTok{,}\NormalTok{ parameters}\OperatorTok{);}

    \CommentTok{// 8. Phase 2: Sign the hash}
\NormalTok{    SignatureValue signatureValue }\OperatorTok{=}\NormalTok{ signingToken}\OperatorTok{.}\FunctionTok{sign}\OperatorTok{(}
\NormalTok{            dataToSign}\OperatorTok{,}\NormalTok{ parameters}\OperatorTok{.}\FunctionTok{getDigestAlgorithm}\OperatorTok{(),}
\NormalTok{            privateKey}\OperatorTok{);}

    \CommentTok{// 9. Finalize: embed signature in PDF}
\NormalTok{    DSSDocument signedDocument }\OperatorTok{=}\NormalTok{ padesService}\OperatorTok{.}\FunctionTok{signDocument}\OperatorTok{(}
\NormalTok{            toSignDocument}\OperatorTok{,}\NormalTok{ parameters}\OperatorTok{,}\NormalTok{ signatureValue}\OperatorTok{);}

    \CommentTok{// 10. Extract the result}
    \BuiltInTok{ByteArrayOutputStream}\NormalTok{ result }\OperatorTok{=} \KeywordTok{new} \BuiltInTok{ByteArrayOutputStream}\OperatorTok{();}
\NormalTok{    signedDocument}\OperatorTok{.}\FunctionTok{writeTo}\OperatorTok{(}\NormalTok{result}\OperatorTok{);}

\NormalTok{    log}\OperatorTok{.}\FunctionTok{info}\OperatorTok{(}\StringTok{"PAdES{-}B{-}T signature created: \{\} bytes"}\OperatorTok{,}
\NormalTok{            result}\OperatorTok{.}\FunctionTok{size}\OperatorTok{());}
    \ControlFlowTok{return}\NormalTok{ result}\OperatorTok{.}\FunctionTok{toByteArray}\OperatorTok{();}
\OperatorTok{\}}
\end{Highlighting}
\end{Shaded}

That is roughly 60 lines of code --- including logging, comments, and
error handling setup --- to produce a fully compliant PAdES-B-T
signature. Compare this to the hundreds of lines required to build the
CMS structure manually with BouncyCastle (Chapter 12), and the value
proposition of EU DSS becomes clear.

\section{Extending Signatures}\label{_extending_signatures}

One of EU DSS's most powerful features is \textbf{signature extension}:
upgrading a signature's PAdES profile without re-signing.

\begin{Shaded}
\begin{Highlighting}[]
\CommentTok{/**}
 \CommentTok{*}\NormalTok{ Extends a PAdES}\CommentTok{{-}}\NormalTok{B}\CommentTok{{-}}\NormalTok{B or B}\CommentTok{{-}}\NormalTok{T signature to PAdES}\CommentTok{{-}}\NormalTok{B}\CommentTok{{-}}\NormalTok{LTA}\CommentTok{.}
 \CommentTok{*}\NormalTok{ Fetches revocation data}\CommentTok{,}\NormalTok{ builds the DSS}\CommentTok{,}\NormalTok{ adds an archival timestamp}\CommentTok{.}
 \CommentTok{*/}
\KeywordTok{public} \DataTypeTok{byte}\OperatorTok{[]} \FunctionTok{extendToLta}\OperatorTok{(}\DataTypeTok{byte}\OperatorTok{[]}\NormalTok{ signedPdfBytes}\OperatorTok{,}
        \BuiltInTok{String}\NormalTok{ tsaUrl}\OperatorTok{)} \KeywordTok{throws} \BuiltInTok{Exception} \OperatorTok{\{}

    \CommentTok{// 1. Set up the certificate verifier with online sources}
\NormalTok{    CommonCertificateVerifier verifier }\OperatorTok{=}
            \KeywordTok{new} \FunctionTok{CommonCertificateVerifier}\OperatorTok{();}
\NormalTok{    verifier}\OperatorTok{.}\FunctionTok{setOcspSource}\OperatorTok{(}\KeywordTok{new} \FunctionTok{OnlineOCSPSource}\OperatorTok{());}
\NormalTok{    verifier}\OperatorTok{.}\FunctionTok{setCrlSource}\OperatorTok{(}\KeywordTok{new} \FunctionTok{OnlineCRLSource}\OperatorTok{());}

    \CommentTok{// 2. Add the EU Trusted List as a trust source}
    \CommentTok{//    (required for validating qualified certificates)}
    \CommentTok{// TrustedListsCertificateSource is populated by the}
    \CommentTok{// TSL loading job — see Spring Boot integration below}

    \CommentTok{// 3. Create the PAdES service with a TSA}
\NormalTok{    PAdESService service }\OperatorTok{=} \KeywordTok{new} \FunctionTok{PAdESService}\OperatorTok{(}\NormalTok{verifier}\OperatorTok{);}
\NormalTok{    OnlineTSPSource tspSource }\OperatorTok{=} \KeywordTok{new} \FunctionTok{OnlineTSPSource}\OperatorTok{(}\NormalTok{tsaUrl}\OperatorTok{);}
\NormalTok{    service}\OperatorTok{.}\FunctionTok{setTspSource}\OperatorTok{(}\NormalTok{tspSource}\OperatorTok{);}

    \CommentTok{// 4. Configure extension parameters}
\NormalTok{    PAdESSignatureParameters extensionParams }\OperatorTok{=}
            \KeywordTok{new} \FunctionTok{PAdESSignatureParameters}\OperatorTok{();}
\NormalTok{    extensionParams}\OperatorTok{.}\FunctionTok{setSignatureLevel}\OperatorTok{(}
\NormalTok{            SignatureLevel}\OperatorTok{.}\FunctionTok{PAdES\_BASELINE\_LTA}\OperatorTok{);}

    \CommentTok{// 5. Extend the document}
\NormalTok{    DSSDocument signedDoc }\OperatorTok{=}
            \KeywordTok{new} \FunctionTok{InMemoryDocument}\OperatorTok{(}\NormalTok{signedPdfBytes}\OperatorTok{);}
\NormalTok{    DSSDocument extendedDoc }\OperatorTok{=}
\NormalTok{            service}\OperatorTok{.}\FunctionTok{extendDocument}\OperatorTok{(}\NormalTok{signedDoc}\OperatorTok{,}\NormalTok{ extensionParams}\OperatorTok{);}

    \BuiltInTok{ByteArrayOutputStream}\NormalTok{ result }\OperatorTok{=} \KeywordTok{new} \BuiltInTok{ByteArrayOutputStream}\OperatorTok{();}
\NormalTok{    extendedDoc}\OperatorTok{.}\FunctionTok{writeTo}\OperatorTok{(}\NormalTok{result}\OperatorTok{);}

\NormalTok{    log}\OperatorTok{.}\FunctionTok{info}\OperatorTok{(}\StringTok{"Signature extended to PAdES{-}B{-}LTA: \{\} bytes"}\OperatorTok{,}
\NormalTok{            result}\OperatorTok{.}\FunctionTok{size}\OperatorTok{());}
    \ControlFlowTok{return}\NormalTok{ result}\OperatorTok{.}\FunctionTok{toByteArray}\OperatorTok{();}
\OperatorTok{\}}
\end{Highlighting}
\end{Shaded}

The \texttt{extendDocument} method handles all the complexity described
in Chapter 17 (PAdES Profiles) and Chapter 18 (Timestamps and LTV):

\begin{enumerate}
\def\labelenumi{\arabic{enumi}.}
\item
  Parses the existing signature to determine its current level
\item
  Fetches OCSP responses and CRLs for every certificate in the chain
\item
  Builds the DSS dictionary with all validation data
\item
  Requests an archival timestamp from the TSA
\item
  Saves everything as an incremental update
\end{enumerate}

This is the operation you would run as a scheduled job to keep archived
documents verifiable indefinitely.

\section{Validation with EU DSS}\label{_validation_with_eu_dss}

\subsection{Basic Validation}\label{_basic_validation}

Validation is where EU DSS provides the most value. Building a correct
signature validator from scratch requires implementing the entire ETSI
EN 319 102-1 validation algorithm --- certificate chain building,
revocation checking at the correct point in time, timestamp validation,
DSS consumption, policy evaluation. EU DSS does all of this:

\begin{Shaded}
\begin{Highlighting}[]
\CommentTok{/**}
 \CommentTok{*}\NormalTok{ Validates all signatures in a signed PDF}\CommentTok{.}
 \CommentTok{*}\NormalTok{ Returns a structured validation report}\CommentTok{.}
 \CommentTok{*/}
\KeywordTok{public}\NormalTok{ Reports }\FunctionTok{validateSignatures}\OperatorTok{(}\DataTypeTok{byte}\OperatorTok{[]}\NormalTok{ signedPdfBytes}\OperatorTok{)}
        \KeywordTok{throws} \BuiltInTok{Exception} \OperatorTok{\{}

    \CommentTok{// 1. Set up the certificate verifier}
\NormalTok{    CommonCertificateVerifier verifier }\OperatorTok{=}
            \KeywordTok{new} \FunctionTok{CommonCertificateVerifier}\OperatorTok{();}
\NormalTok{    verifier}\OperatorTok{.}\FunctionTok{setOcspSource}\OperatorTok{(}\KeywordTok{new} \FunctionTok{OnlineOCSPSource}\OperatorTok{());}
\NormalTok{    verifier}\OperatorTok{.}\FunctionTok{setCrlSource}\OperatorTok{(}\KeywordTok{new} \FunctionTok{OnlineCRLSource}\OperatorTok{());}

    \CommentTok{// For qualified certificate validation, add the EUTL:}
    \CommentTok{// verifier.addTrustedCertSources(trustedListSource);}

    \CommentTok{// 2. Create a validator for the signed document}
\NormalTok{    DSSDocument document }\OperatorTok{=} \KeywordTok{new} \FunctionTok{InMemoryDocument}\OperatorTok{(}\NormalTok{signedPdfBytes}\OperatorTok{);}
\NormalTok{    SignedDocumentValidator validator }\OperatorTok{=}
\NormalTok{            SignedDocumentValidator}\OperatorTok{.}\FunctionTok{fromDocument}\OperatorTok{(}\NormalTok{document}\OperatorTok{);}
\NormalTok{    validator}\OperatorTok{.}\FunctionTok{setCertificateVerifier}\OperatorTok{(}\NormalTok{verifier}\OperatorTok{);}

    \CommentTok{// 3. Validate — this runs the full ETSI EN 319 102{-}1}
    \CommentTok{//    validation algorithm}
\NormalTok{    Reports reports }\OperatorTok{=}\NormalTok{ validator}\OperatorTok{.}\FunctionTok{validateDocument}\OperatorTok{();}

    \CommentTok{// 4. Access the results}
\NormalTok{    SimpleReport simpleReport }\OperatorTok{=}\NormalTok{ reports}\OperatorTok{.}\FunctionTok{getSimpleReport}\OperatorTok{();}
    \ControlFlowTok{for} \OperatorTok{(}\BuiltInTok{String}\NormalTok{ sigId }\OperatorTok{:}\NormalTok{ simpleReport}\OperatorTok{.}\FunctionTok{getSignatureIdList}\OperatorTok{())} \OperatorTok{\{}
\NormalTok{        log}\OperatorTok{.}\FunctionTok{info}\OperatorTok{(}\StringTok{"Signature \{\}: indication=\{\}, subIndication=\{\}"}\OperatorTok{,}
\NormalTok{                sigId}\OperatorTok{,}
\NormalTok{                simpleReport}\OperatorTok{.}\FunctionTok{getIndication}\OperatorTok{(}\NormalTok{sigId}\OperatorTok{),}
\NormalTok{                simpleReport}\OperatorTok{.}\FunctionTok{getSubIndication}\OperatorTok{(}\NormalTok{sigId}\OperatorTok{));}
\NormalTok{        log}\OperatorTok{.}\FunctionTok{info}\OperatorTok{(}\StringTok{"  Signed by: \{\}"}\OperatorTok{,}
\NormalTok{                simpleReport}\OperatorTok{.}\FunctionTok{getSignedBy}\OperatorTok{(}\NormalTok{sigId}\OperatorTok{));}
\NormalTok{        log}\OperatorTok{.}\FunctionTok{info}\OperatorTok{(}\StringTok{"  Signing time: \{\}"}\OperatorTok{,}
\NormalTok{                simpleReport}\OperatorTok{.}\FunctionTok{getSigningTime}\OperatorTok{(}\NormalTok{sigId}\OperatorTok{));}
\NormalTok{        log}\OperatorTok{.}\FunctionTok{info}\OperatorTok{(}\StringTok{"  Signature level: \{\}"}\OperatorTok{,}
\NormalTok{                simpleReport}\OperatorTok{.}\FunctionTok{getSignatureFormat}\OperatorTok{(}\NormalTok{sigId}\OperatorTok{));}
    \OperatorTok{\}}

    \ControlFlowTok{return}\NormalTok{ reports}\OperatorTok{;}
\OperatorTok{\}}
\end{Highlighting}
\end{Shaded}

\subsection{Validation Reports}\label{_validation_reports}

EU DSS produces three levels of validation reports:

\begin{enumerate}
\def\labelenumi{\arabic{enumi}.}
\item
  \textbf{Simple Report}: A summary of each signature's validity ---
  pass/fail, signer name, signing time, signature level. Suitable for
  end-user display.
\item
  \textbf{Detailed Report}: A structured breakdown of every validation
  check --- certificate chain validation, revocation checking, timestamp
  validation, policy conformance. Each check has an indication (PASSED,
  FAILED, INDETERMINATE) and a sub-indication explaining why.
\item
  \textbf{Diagnostic Data}: The raw data collected during validation ---
  all certificates encountered, all revocation data fetched, all
  timestamps found. This is the lowest level and is useful for debugging
  validation failures.
\end{enumerate}

\begin{Shaded}
\begin{Highlighting}[]
\CommentTok{// Accessing the detailed report}
\NormalTok{DetailedReport detailedReport }\OperatorTok{=}\NormalTok{ reports}\OperatorTok{.}\FunctionTok{getDetailedReport}\OperatorTok{();}

\CommentTok{// Get the basic building blocks conclusion}
\CommentTok{//  (certificate chain and revocation status)}
\ControlFlowTok{for} \OperatorTok{(}\BuiltInTok{String}\NormalTok{ sigId }\OperatorTok{:}\NormalTok{ detailedReport}\OperatorTok{.}\FunctionTok{getSignatureIds}\OperatorTok{())} \OperatorTok{\{}
\NormalTok{    Indication bbsIndication }\OperatorTok{=}
\NormalTok{            detailedReport}\OperatorTok{.}\FunctionTok{getBasicBuildingBlocksIndication}\OperatorTok{(}\NormalTok{sigId}\OperatorTok{);}
\NormalTok{    SubIndication bbsSubIndication }\OperatorTok{=}
\NormalTok{            detailedReport}\OperatorTok{.}\FunctionTok{getBasicBuildingBlocksSubIndication}\OperatorTok{(}\NormalTok{sigId}\OperatorTok{);}

\NormalTok{    log}\OperatorTok{.}\FunctionTok{info}\OperatorTok{(}\StringTok{"Signature \{\} basic building blocks: \{\} / \{\}"}\OperatorTok{,}
\NormalTok{            sigId}\OperatorTok{,}\NormalTok{ bbsIndication}\OperatorTok{,}\NormalTok{ bbsSubIndication}\OperatorTok{);}

    \CommentTok{// Get timestamp validation results}
    \BuiltInTok{List}\OperatorTok{\textless{}}\BuiltInTok{String}\OperatorTok{\textgreater{}}\NormalTok{ tsIds }\OperatorTok{=}
\NormalTok{            detailedReport}\OperatorTok{.}\FunctionTok{getTimestampIds}\OperatorTok{(}\NormalTok{sigId}\OperatorTok{);}
    \ControlFlowTok{for} \OperatorTok{(}\BuiltInTok{String}\NormalTok{ tsId }\OperatorTok{:}\NormalTok{ tsIds}\OperatorTok{)} \OperatorTok{\{}
\NormalTok{        log}\OperatorTok{.}\FunctionTok{info}\OperatorTok{(}\StringTok{"  Timestamp \{\}: \{\}"}\OperatorTok{,}
\NormalTok{                tsId}\OperatorTok{,}
\NormalTok{                detailedReport}\OperatorTok{.}\FunctionTok{getTimestampIndication}\OperatorTok{(}\NormalTok{tsId}\OperatorTok{));}
    \OperatorTok{\}}
\OperatorTok{\}}
\end{Highlighting}
\end{Shaded}

\subsection{EU Trusted List
Integration}\label{_eu_trusted_list_integration}

For validating qualified certificates, EU DSS can load and parse the EU
Trusted List (EUTL):

\begin{Shaded}
\begin{Highlighting}[]
\CommentTok{/**}
 \CommentTok{*}\NormalTok{ Configures EU DSS to validate against the EU Trusted List}\CommentTok{.}
 \CommentTok{*}\NormalTok{ The EUTL is downloaded}\CommentTok{,}\NormalTok{ its signature verified}\CommentTok{,}\NormalTok{ and the qualified}
 \CommentTok{*}\NormalTok{ trust service providers are extracted as trusted certificate sources}\CommentTok{.}
 \CommentTok{*/}
\KeywordTok{public}\NormalTok{ TrustedListsCertificateSource }\FunctionTok{loadEuTrustedList}\OperatorTok{()}
        \KeywordTok{throws} \BuiltInTok{Exception} \OperatorTok{\{}

\NormalTok{    TrustedListsCertificateSource trustedListSource }\OperatorTok{=}
            \KeywordTok{new} \FunctionTok{TrustedListsCertificateSource}\OperatorTok{();}

    \CommentTok{// The LOTL (List of the Lists) is the entry point}
\NormalTok{    LOTLSource lotlSource }\OperatorTok{=} \KeywordTok{new} \FunctionTok{LOTLSource}\OperatorTok{();}
\NormalTok{    lotlSource}\OperatorTok{.}\FunctionTok{setUrl}\OperatorTok{(}
            \StringTok{"https://ec.europa.eu/tools/lotl/eu{-}lotl.xml"}\OperatorTok{);}

    \CommentTok{// The LOTL is signed — we need the signing certificate}
    \CommentTok{// to verify it (bootstrapping problem — this cert is}
    \CommentTok{// distributed as part of the EU DSS distribution)}
\NormalTok{    lotlSource}\OperatorTok{.}\FunctionTok{setCertificateSource}\OperatorTok{(}
            \FunctionTok{officialJournalCertificateSource}\OperatorTok{());}

    \CommentTok{// Configure the TSL validation job}
\NormalTok{    TLValidationJob tlValidationJob }\OperatorTok{=} \KeywordTok{new} \FunctionTok{TLValidationJob}\OperatorTok{();}
\NormalTok{    tlValidationJob}\OperatorTok{.}\FunctionTok{setListOfTrustedListSources}\OperatorTok{(}\NormalTok{lotlSource}\OperatorTok{);}
\NormalTok{    tlValidationJob}\OperatorTok{.}\FunctionTok{setTrustedListCertificateSource}\OperatorTok{(}
\NormalTok{            trustedListSource}\OperatorTok{);}

    \CommentTok{// Use a caching data loader to avoid re{-}downloading}
    \CommentTok{// the entire EUTL on every validation}
\NormalTok{    FileCacheDataLoader cacheDataLoader }\OperatorTok{=}
            \KeywordTok{new} \FunctionTok{FileCacheDataLoader}\OperatorTok{();}
\NormalTok{    cacheDataLoader}\OperatorTok{.}\FunctionTok{setCacheExpirationTime}\OperatorTok{(}
            \DecValTok{24} \OperatorTok{*} \DecValTok{60} \OperatorTok{*} \DecValTok{60} \OperatorTok{*} \DecValTok{1000L}\OperatorTok{);} \CommentTok{// 24 hours}
\NormalTok{    cacheDataLoader}\OperatorTok{.}\FunctionTok{setFileCacheDirectory}\OperatorTok{(}
            \KeywordTok{new} \BuiltInTok{File}\OperatorTok{(}\StringTok{"/tmp/dss{-}tl{-}cache"}\OperatorTok{));}
\NormalTok{    tlValidationJob}\OperatorTok{.}\FunctionTok{setOnlineDataLoader}\OperatorTok{(}\NormalTok{cacheDataLoader}\OperatorTok{);}

    \CommentTok{// Run the initial load}
\NormalTok{    tlValidationJob}\OperatorTok{.}\FunctionTok{onlineRefresh}\OperatorTok{();}

\NormalTok{    log}\OperatorTok{.}\FunctionTok{info}\OperatorTok{(}\StringTok{"EU Trusted List loaded: \{\} trusted certificates"}\OperatorTok{,}
\NormalTok{            trustedListSource}\OperatorTok{.}\FunctionTok{getNumberOfTrustedCertificates}\OperatorTok{());}

    \ControlFlowTok{return}\NormalTok{ trustedListSource}\OperatorTok{;}
\OperatorTok{\}}
\end{Highlighting}
\end{Shaded}

The EU Trusted List is large --- the LOTL (List of the Lists) references
27+ national trusted lists, each containing dozens of trust service
providers and their certificates. The initial download and parsing can
take 30-60 seconds.

In production: - Load the EUTL at application startup, not on the first
validation request - Cache the downloaded lists on disk (use
\texttt{FileCacheDataLoader}) - Refresh periodically (every 24 hours is
typical) rather than on every request - If the EUTL is unreachable, fall
back to the cached version

Never skip EUTL validation --- if you validate qualified signatures
without checking the EUTL, you are accepting certificates that may not
be qualified.

\section{Spring Boot Integration}\label{_spring_boot_integration}

EU DSS provides a Spring Boot starter module that auto-configures most
of the framework:

\begin{Shaded}
\begin{Highlighting}[]
\NormalTok{\textless{}}\KeywordTok{dependency}\NormalTok{\textgreater{}}
\NormalTok{    \textless{}}\KeywordTok{groupId}\NormalTok{\textgreater{}eu.europa.dss\textless{}/}\KeywordTok{groupId}\NormalTok{\textgreater{}}
\NormalTok{    \textless{}}\KeywordTok{artifactId}\NormalTok{\textgreater{}dss{-}spring{-}boot{-}starter\textless{}/}\KeywordTok{artifactId}\NormalTok{\textgreater{}}
\NormalTok{    \textless{}}\KeywordTok{version}\NormalTok{\textgreater{}$\{dss.version\}\textless{}/}\KeywordTok{version}\NormalTok{\textgreater{}}
\NormalTok{\textless{}/}\KeywordTok{dependency}\NormalTok{\textgreater{}}
\end{Highlighting}
\end{Shaded}

\subsection{Configuration Properties}\label{_configuration_properties}

The starter exposes configuration through
\texttt{application.properties}:

\begin{Shaded}
\begin{Highlighting}[]
\NormalTok{\# TSA configuration}
\NormalTok{dss.tsp.url=http://timestamp.digicert.com}
\NormalTok{dss.tsp.digest{-}algorithm=SHA256}

\NormalTok{\# EU Trusted List}
\NormalTok{dss.lotl.url=https://ec.europa.eu/tools/lotl/eu{-}lotl.xml}
\NormalTok{dss.lotl.cache{-}directory=/var/cache/dss/tl}
\NormalTok{dss.lotl.cache{-}expiration=86400000}

\NormalTok{\# OCSP/CRL timeouts}
\NormalTok{dss.ocsp.connect{-}timeout=5000}
\NormalTok{dss.ocsp.read{-}timeout=10000}
\NormalTok{dss.crl.connect{-}timeout=5000}
\NormalTok{dss.crl.read{-}timeout=30000}

\NormalTok{\# Proxy (if behind a corporate proxy)}
\NormalTok{dss.proxy.https.host=proxy.example.com}
\NormalTok{dss.proxy.https.port=8443}
\end{Highlighting}
\end{Shaded}

\subsection{Auto-Configured Beans}\label{_auto_configured_beans}

The starter auto-configures the following Spring beans:

\begin{itemize}
\item
  \texttt{CommonCertificateVerifier} --- pre-configured with OCSP and
  CRL sources
\item
  \texttt{OnlineTSPSource} --- configured from \texttt{dss.tsp.url}
\item
  \texttt{OnlineOCSPSource} and \texttt{OnlineCRLSource} --- with
  timeout configuration
\item
  \texttt{TrustedListsCertificateSource} --- loaded from the EUTL on
  startup
\item
  \texttt{TLValidationJob} --- scheduled to refresh the EUTL
  periodically
\item
  \texttt{CommonsDataLoader} --- HTTP client with proxy and timeout
  support
\end{itemize}

You can inject these directly into your services:

\begin{Shaded}
\begin{Highlighting}[]
\AttributeTok{@Service}
\AttributeTok{@Slf4j}
\KeywordTok{public} \KeywordTok{class}\NormalTok{ SignatureService }\OperatorTok{\{}

    \KeywordTok{private} \DataTypeTok{final}\NormalTok{ CommonCertificateVerifier certificateVerifier}\OperatorTok{;}
    \KeywordTok{private} \DataTypeTok{final}\NormalTok{ OnlineTSPSource tspSource}\OperatorTok{;}

    \KeywordTok{public} \FunctionTok{SignatureService}\OperatorTok{(}
\NormalTok{            CommonCertificateVerifier certificateVerifier}\OperatorTok{,}
\NormalTok{            OnlineTSPSource tspSource}\OperatorTok{)} \OperatorTok{\{}
        \KeywordTok{this}\OperatorTok{.}\FunctionTok{certificateVerifier} \OperatorTok{=}\NormalTok{ certificateVerifier}\OperatorTok{;}
        \KeywordTok{this}\OperatorTok{.}\FunctionTok{tspSource} \OperatorTok{=}\NormalTok{ tspSource}\OperatorTok{;}
    \OperatorTok{\}}

    \KeywordTok{public} \DataTypeTok{byte}\OperatorTok{[]} \FunctionTok{signPdf}\OperatorTok{(}\DataTypeTok{byte}\OperatorTok{[]}\NormalTok{ pdfBytes}\OperatorTok{,}
\NormalTok{            SignatureTokenConnection token}\OperatorTok{,}
\NormalTok{            DSSPrivateKeyEntry keyEntry}\OperatorTok{)} \KeywordTok{throws} \BuiltInTok{Exception} \OperatorTok{\{}

\NormalTok{        PAdESService service }\OperatorTok{=}
                \KeywordTok{new} \FunctionTok{PAdESService}\OperatorTok{(}\NormalTok{certificateVerifier}\OperatorTok{);}
\NormalTok{        service}\OperatorTok{.}\FunctionTok{setTspSource}\OperatorTok{(}\NormalTok{tspSource}\OperatorTok{);}

\NormalTok{        PAdESSignatureParameters params }\OperatorTok{=}
                \KeywordTok{new} \FunctionTok{PAdESSignatureParameters}\OperatorTok{();}
\NormalTok{        params}\OperatorTok{.}\FunctionTok{setSignatureLevel}\OperatorTok{(}
\NormalTok{                SignatureLevel}\OperatorTok{.}\FunctionTok{PAdES\_BASELINE\_T}\OperatorTok{);}
\NormalTok{        params}\OperatorTok{.}\FunctionTok{setDigestAlgorithm}\OperatorTok{(}\NormalTok{DigestAlgorithm}\OperatorTok{.}\FunctionTok{SHA256}\OperatorTok{);}
\NormalTok{        params}\OperatorTok{.}\FunctionTok{setSigningCertificate}\OperatorTok{(}
\NormalTok{                keyEntry}\OperatorTok{.}\FunctionTok{getCertificate}\OperatorTok{());}
\NormalTok{        params}\OperatorTok{.}\FunctionTok{setCertificateChain}\OperatorTok{(}
\NormalTok{                keyEntry}\OperatorTok{.}\FunctionTok{getCertificateChain}\OperatorTok{());}

\NormalTok{        DSSDocument doc }\OperatorTok{=} \KeywordTok{new} \FunctionTok{InMemoryDocument}\OperatorTok{(}\NormalTok{pdfBytes}\OperatorTok{);}
\NormalTok{        ToBeSigned tbs }\OperatorTok{=}\NormalTok{ service}\OperatorTok{.}\FunctionTok{getDataToSign}\OperatorTok{(}\NormalTok{doc}\OperatorTok{,}\NormalTok{ params}\OperatorTok{);}
\NormalTok{        SignatureValue sv }\OperatorTok{=}\NormalTok{ token}\OperatorTok{.}\FunctionTok{sign}\OperatorTok{(}
\NormalTok{                tbs}\OperatorTok{,}\NormalTok{ params}\OperatorTok{.}\FunctionTok{getDigestAlgorithm}\OperatorTok{(),}\NormalTok{ keyEntry}\OperatorTok{);}
\NormalTok{        DSSDocument signed }\OperatorTok{=}
\NormalTok{                service}\OperatorTok{.}\FunctionTok{signDocument}\OperatorTok{(}\NormalTok{doc}\OperatorTok{,}\NormalTok{ params}\OperatorTok{,}\NormalTok{ sv}\OperatorTok{);}

        \BuiltInTok{ByteArrayOutputStream}\NormalTok{ out }\OperatorTok{=} \KeywordTok{new} \BuiltInTok{ByteArrayOutputStream}\OperatorTok{();}
\NormalTok{        signed}\OperatorTok{.}\FunctionTok{writeTo}\OperatorTok{(}\NormalTok{out}\OperatorTok{);}
        \ControlFlowTok{return}\NormalTok{ out}\OperatorTok{.}\FunctionTok{toByteArray}\OperatorTok{();}
    \OperatorTok{\}}
\OperatorTok{\}}
\end{Highlighting}
\end{Shaded}

\section{EU DSS vs. Rolling Your Own}\label{_eu_dss_vs_rolling_your_own}

\subsection{When to Use EU DSS}\label{_when_to_use_eu_dss}

Use EU DSS when:

\begin{itemize}
\item
  \textbf{You need eIDAS compliance}: EU DSS is the reference
  implementation. If your validation disagrees with EU DSS, regulators
  will assume your code is wrong, not EU DSS.
\item
  \textbf{You validate qualified signatures}: EUTL parsing, QCStatements
  interpretation, SCAL assessment --- EU DSS handles all of this.
  Implementing it correctly from scratch is months of work.
\item
  \textbf{You need multiple signature formats}: If your application
  handles PAdES, CAdES, and XAdES, EU DSS's format-independent
  architecture saves enormous effort.
\item
  \textbf{You need long-term maintenance}: EU DSS is actively maintained
  with a clear release cadence. ETSI standards evolve, new algorithms
  are added, old ones deprecated. EU DSS tracks these changes.
\end{itemize}

\subsection{When to Consider
Alternatives}\label{_when_to_consider_alternatives}

EU DSS may not be the right choice when:

\begin{itemize}
\item
  \textbf{You need total control over CMS construction}: EU DSS
  abstracts away the CMS building process. If your use case requires
  non-standard CMS attributes, custom signature dictionary entries, or
  unusual byte range handling, EU DSS may fight you. In these cases,
  using BouncyCastle directly (as shown throughout this book) gives you
  full control.
\item
  \textbf{You have strict dependency constraints}: EU DSS brings a
  significant dependency tree --- PDFBox, BouncyCastle, Apache
  HttpClient, JAXB, SLF4J, and dozens of transitive dependencies. If
  your project has strict dependency policies, the EU DSS footprint may
  be too large.
\item
  \textbf{You only create signatures, never validate}: If you only sign
  PDFs and never validate them, EU DSS brings a lot of validation
  machinery you do not need. A leaner approach with PDFBox +
  BouncyCastle may be more appropriate.
\item
  \textbf{You need to support PDF features beyond signing}: EU DSS is
  focused on digital signatures. It does not help with form filling,
  content editing, font embedding, or other PDF operations. You will
  still need PDFBox or another PDF library for those.
\end{itemize}

\subsection{The Hybrid Approach}\label{_the_hybrid_approach}

In practice, many projects use a hybrid approach:

\begin{itemize}
\item
  \textbf{EU DSS for validation}: Use \texttt{SignedDocumentValidator}
  to validate incoming signed PDFs. This gives you correct ETSI EN 319
  102-1 validation with EUTL support.
\item
  \textbf{PDFBox + BouncyCastle for signing}: Use PDFBox for PDF
  manipulation and BouncyCastle for CMS construction. This gives you
  full control over the signing process while still producing
  standard-compliant signatures.
\item
  \textbf{EU DSS for extension}: Use
  \texttt{PAdESService.extendDocument()} to upgrade signatures from
  B-B/B-T to B-LTA. This is complex to implement manually (DSS
  dictionary construction, archival timestamp embedding) and EU DSS does
  it reliably.
\end{itemize}

This hybrid model gives you the best of both worlds: full control where
you need it, and battle-tested reference code where correctness is
paramount.

If you adopt the hybrid approach, be mindful of PDFBox version
alignment. EU DSS uses a specific PDFBox version internally. If your
direct PDFBox dependency uses a different version, classloader conflicts
can cause subtle, hard-to-diagnose bugs --- wrong method signatures,
missing classes, or (worst case) silent data corruption. Always align
your PDFBox version with the one in EU DSS's dependency tree.

\section{Common Pitfalls}\label{_common_pitfalls_2}

\subsection{1. Forgetting to Set the Certificate
Chain}\label{_1_forgetting_to_set_the_certificate_chain}

\begin{Shaded}
\begin{Highlighting}[]
\CommentTok{// WRONG — only sets the signing certificate, not the chain}
\NormalTok{parameters}\OperatorTok{.}\FunctionTok{setSigningCertificate}\OperatorTok{(}\NormalTok{keyEntry}\OperatorTok{.}\FunctionTok{getCertificate}\OperatorTok{());}
\CommentTok{// Missing: parameters.setCertificateChain(...)}

\CommentTok{// RIGHT — set both}
\NormalTok{parameters}\OperatorTok{.}\FunctionTok{setSigningCertificate}\OperatorTok{(}\NormalTok{keyEntry}\OperatorTok{.}\FunctionTok{getCertificate}\OperatorTok{());}
\NormalTok{parameters}\OperatorTok{.}\FunctionTok{setCertificateChain}\OperatorTok{(}\NormalTok{keyEntry}\OperatorTok{.}\FunctionTok{getCertificateChain}\OperatorTok{());}
\end{Highlighting}
\end{Shaded}

Without the certificate chain, the CMS \texttt{certificates} collection
will contain only the signer's certificate. Validators that do not have
access to the intermediate certificates (from AIA or a local store) will
fail to build the chain.

\subsection{2. Not Configuring a Certificate
Verifier}\label{_2_not_configuring_a_certificate_verifier}

\begin{Shaded}
\begin{Highlighting}[]
\CommentTok{// WRONG — empty verifier, cannot check revocation}
\NormalTok{CommonCertificateVerifier verifier }\OperatorTok{=}
        \KeywordTok{new} \FunctionTok{CommonCertificateVerifier}\OperatorTok{();}
\CommentTok{// No OCSP source, no CRL source set}

\CommentTok{// RIGHT — configure revocation sources}
\NormalTok{CommonCertificateVerifier verifier }\OperatorTok{=}
        \KeywordTok{new} \FunctionTok{CommonCertificateVerifier}\OperatorTok{();}
\NormalTok{verifier}\OperatorTok{.}\FunctionTok{setOcspSource}\OperatorTok{(}\KeywordTok{new} \FunctionTok{OnlineOCSPSource}\OperatorTok{());}
\NormalTok{verifier}\OperatorTok{.}\FunctionTok{setCrlSource}\OperatorTok{(}\KeywordTok{new} \FunctionTok{OnlineCRLSource}\OperatorTok{());}
\end{Highlighting}
\end{Shaded}

An unconfigured verifier cannot check certificate revocation. EU DSS
will still produce a signature, but validation will report INDETERMINATE
for revocation checking.

\subsection{3. Underestimating /Contents
Size}\label{_3_underestimating_contents_size}

When EU DSS constructs the PDF signature, it must pre-allocate space for
the \texttt{/Contents} hex string. If the allocated space is too small,
the CMS signature (with certificates, timestamps, etc.) will not fit. EU
DSS handles this internally with reasonable defaults, but if you are
building very large certificate chains or using multiple timestamps, you
may need to increase the allocation:

\begin{Shaded}
\begin{Highlighting}[]
\NormalTok{PAdESSignatureParameters params }\OperatorTok{=}
        \KeywordTok{new} \FunctionTok{PAdESSignatureParameters}\OperatorTok{();}
\NormalTok{params}\OperatorTok{.}\FunctionTok{setContentSize}\OperatorTok{(}\DecValTok{32768}\OperatorTok{);}  \CommentTok{// 32 KB for /Contents}
\CommentTok{// Default is typically 16384 (16 KB)}
\end{Highlighting}
\end{Shaded}

\subsection{4. Blocking the EUTL Load on First
Request}\label{_4_blocking_the_eutl_load_on_first_request}

The EUTL download takes 30-60 seconds. If you trigger it lazily (on the
first validation request), your first user will experience a 30-second
delay. Always load the EUTL eagerly at application startup:

\begin{Shaded}
\begin{Highlighting}[]
\AttributeTok{@PostConstruct}
\KeywordTok{public} \DataTypeTok{void} \FunctionTok{initTrustedLists}\OperatorTok{()} \OperatorTok{\{}
\NormalTok{    log}\OperatorTok{.}\FunctionTok{info}\OperatorTok{(}\StringTok{"Loading EU Trusted List — this may take 30{-}60 seconds"}\OperatorTok{);}
\NormalTok{    tlValidationJob}\OperatorTok{.}\FunctionTok{onlineRefresh}\OperatorTok{();}
\NormalTok{    log}\OperatorTok{.}\FunctionTok{info}\OperatorTok{(}\StringTok{"EU Trusted List loaded: \{\} certificates"}\OperatorTok{,}
\NormalTok{            trustedListSource}\OperatorTok{.}\FunctionTok{getNumberOfTrustedCertificates}\OperatorTok{());}
\OperatorTok{\}}
\end{Highlighting}
\end{Shaded}

\subsection{5. Ignoring Validation
Sub-Indications}\label{_5_ignoring_validation_sub_indications}

A common mistake is checking only the top-level \texttt{Indication}
(TOTAL\_PASSED, TOTAL\_FAILED, INDETERMINATE) and ignoring the
\texttt{SubIndication}. The sub-indication tells you \emph{why} the
validation reached that conclusion:

\begin{itemize}
\item
  \texttt{NO\_SIGNING\_CERTIFICATE\_FOUND} --- the signer's certificate
  is missing
\item
  \texttt{REVOKED} --- a certificate in the chain has been revoked
\item
  \texttt{EXPIRED} --- a certificate has expired and there is no proof
  that signing occurred before expiry
\item
  \texttt{TRY\_LATER} --- revocation data could not be fetched (OCSP/CRL
  unreachable)
\item
  \texttt{CRYPTO\_CONSTRAINTS\_FAILURE} --- weak algorithms detected
  (e.g., SHA-1)
\item
  \texttt{HASH\_FAILURE} --- the message digest does not match (document
  has been modified)
\end{itemize}

Always log the sub-indication alongside the indication for debugging.

\section{Key Takeaways}\label{_key_takeaways_11}

\begin{itemize}
\item
  \textbf{EU DSS} is the European Commission's open-source reference
  implementation for digital signature creation, extension, and
  validation --- LGPL 2.1 licensed, actively maintained
\item
  The \textbf{two-phase signing model} (ToBeSigned / SignatureValue)
  cleanly separates CMS construction from the private key operation,
  enabling local and remote signing through the same API
\item
  \textbf{SignatureTokenConnection} abstracts key access: PKCS\#12,
  PKCS\#11, MS CAPI, or custom (CSC API) --- swap the token, keep the
  signing flow
\item
  EU DSS handles \textbf{PAdES-B-B through B-LTA} creation and the full
  ETSI EN 319 102-1 \textbf{validation algorithm}, including EUTL lookup
  for qualified certificate determination
\item
  \textbf{Signature extension} (\texttt{extendDocument}) upgrades a
  signature's PAdES profile without re-signing --- the most practical
  feature for archival workflows
\item
  The \textbf{Spring Boot starter} auto-configures certificate
  verifiers, TSA sources, and EUTL loading --- production-ready with
  minimal configuration
\item
  \textbf{Validation reports} come in three levels: Simple (pass/fail),
  Detailed (per-check breakdown), Diagnostic (raw collected data)
\item
  The \textbf{EUTL} must be loaded at startup and cached on disk ---
  lazy loading causes unacceptable latency, missing EUTL data causes
  incorrect qualified status determination
\item
  \textbf{Hybrid approach} is common in practice: EU DSS for validation
  and extension, PDFBox + BouncyCastle for signing with full CMS control
\item
  Always align \textbf{PDFBox versions} between your direct dependency
  and EU DSS's transitive dependency --- version mismatches cause
  classloader conflicts
\end{itemize}

\chapter{Security, Incidents \&
Outlook}\label{_security_incidents_outlook}

\section{PDF Signature Attacks: A Complete Technical
Analysis}\label{_pdf_signature_attacks_a_complete_technical_analysis}

\begin{quote}
We found 21 out of 22 desktop viewer applications and 5 out of 7 online
validation services to be vulnerable to at least one of our attacks.

---  Vladislav Mladenov et al. 1 Trillion Dollar Refund
\end{quote}

For nearly two decades, the PDF signature ecosystem operated on an
assumption that turned out to be wrong: that viewers actually verified
what they claimed to verify. Between 2019 and 2021, a research team at
Ruhr University Bochum systematically dismantled this assumption. They
showed that an attacker could take a validly signed PDF, modify it in
specific ways, and present the result to any major PDF viewer --- Adobe
Acrobat, Foxit Reader, Nitro, LibreOffice, Okular, dozens of others ---
and the viewer would display a green checkmark and the text "Signature
is valid."

The signed content had been altered. The viewer said it hadn't. Every
viewer was wrong.

This chapter is a complete technical analysis of every major attack
against PDF digital signatures --- from the Bochum research program to
earlier CVEs, from the theoretical mechanics to working code that
demonstrates each attack class. The purpose is not to enable attacks but
to make them impossible to misunderstand. If you build or maintain a PDF
signature validation system, you need to understand exactly how these
attacks work so that your code cannot be fooled by them. If you are a
security researcher, this chapter gives you the foundation to find the
next generation of PDF signature vulnerabilities --- because the Bochum
team found the obvious ones, but the attack surface is far from fully
explored.

\subsection{The Researchers}\label{_the_researchers}

The PDF signature attacks of 2019--2021 were the work of the
\textbf{Chair of Network and Data Security} (Lehrstuhl für Netz- und
Datensicherheit) at Ruhr University Bochum, Germany, led by
\textbf{Prof. Dr. Jörg Schwenk}.

The core research team:

\begin{itemize}
\item
  \textbf{Vladislav Mladenov} --- PhD researcher, co-first author of the
  CCS 2019 paper, led the implementation and testing of the USF, ISA,
  and SWA attack variants
\item
  \textbf{Christian Mainka} --- post-doctoral researcher, co-first
  author of both the CCS 2019 and NDSS 2021 papers, designed the Shadow
  Attack taxonomy and coordinated responsible disclosure
\item
  \textbf{Karsten Meyer zu Selhausen} --- researcher and co-founder of
  \textbf{Hackmanit GmbH}, the security consultancy that spun out of the
  Bochum group; contributed to attack implementation and vendor
  coordination
\item
  \textbf{Martin Grothe} --- researcher, contributed to the CCS 2019
  paper's systematic evaluation across 22 viewers
\item
  \textbf{Simon Rohlmann} --- researcher, co-author of the NDSS 2021
  Shadow Attacks paper, developed the Hide-and-Replace variant
\item
  \textbf{Prof. Dr. Jörg Schwenk} --- senior author on both papers,
  provided the research direction and the formal security model for PDF
  signatures
\end{itemize}

The team maintained a coordinated disclosure website at
\textbf{pdf-insecurity.org}, where they published proof-of-concept PDFs,
detailed attack descriptions, and a viewer-by-viewer vulnerability
matrix. Hackmanit GmbH later developed \textbf{PDF Detector}, a tool for
checking PDFs against known attack patterns.

The two landmark papers:

\begin{enumerate}
\def\labelenumi{\arabic{enumi}.}
\item
  \emph{"1 Trillion Dollar Refund: How To Spoof PDF Signatures"} ---
  presented at \textbf{ACM Conference on Computer and Communications
  Security (CCS)}, November 2019, London. This paper introduced three
  attack classes (USF, ISA, SWA) and tested them against 22 desktop
  viewers and 7 online validation services.
\item
  \emph{"Shadow Attacks: Hiding and Replacing Content in Signed PDFs"}
  --- presented at \textbf{Network and Distributed Systems Security
  Symposium (NDSS)}, February 2021 (virtual). This paper introduced a
  fundamentally different attack model where malicious content is
  prepared \emph{before} the document is signed.
\end{enumerate}

Both conferences are top-tier academic venues in computer security. The
CCS 2019 paper won a \textbf{Distinguished Paper Award} --- a
recognition given to only a handful of papers at each conference.

\subsection{Before Bochum: The State of PDF Signature
Validation}\label{_before_bochum_the_state_of_pdf_signature_validation}

To understand why the Bochum attacks worked, you need to understand what
PDF viewers \emph{were} doing when they claimed to "validate" a
signature.

A correct signature validation must verify at least the following:

\begin{enumerate}
\def\labelenumi{\arabic{enumi}.}
\item
  \textbf{Structural integrity}: The byte range specified in the
  \texttt{/ByteRange} array must cover the entire file except the gap
  reserved for the \texttt{/Contents} value. There must be no
  unaccounted bytes before, between, or after the byte ranges.
\item
  \textbf{Cryptographic verification}: The CMS \texttt{SignedData}
  structure in \texttt{/Contents} must contain a valid signature over
  the hash of the byte range data. The digest algorithm, signing
  algorithm, and signed attributes must all be consistent.
\item
  \textbf{Certificate chain validation}: The signer's certificate must
  chain to a trusted root. The chain must be valid at the time of
  signing (or at the time of a trusted timestamp). Revocation status
  must be checked.
\item
  \textbf{Modification detection}: If additional data has been appended
  to the file after the signature (an incremental update), the viewer
  must determine whether those changes are permissible under the
  signature's \texttt{/TransformMethod} (e.g., DocMDP permission level)
  or whether they indicate tampering.
\end{enumerate}

What the Bochum team discovered was that most viewers implemented step 3
reasonably well (they checked certificate chains) but failed
catastrophically at steps 1, 2, and 4. Some viewers didn't check the
cryptographic signature value at all. Some didn't verify that the byte
range covered the entire file. Some allowed arbitrary changes after
signing without any warning.

The result: a perfectly executed certificate chain validation was
protecting\ldots\hspace{0pt} nothing. The viewers were carefully
verifying \emph{who} signed the document, while completely ignoring
\emph{whether the document had been changed since they signed it}.

\subsection{Attack Class 1: Universal Signature Forgery
(USF)}\label{_attack_class_1_universal_signature_forgery_usf}

\subsubsection{The Concept}\label{_the_concept}

Universal Signature Forgery is the simplest and most devastating of the
three original attack classes. The premise is almost embarrassingly
straightforward: what happens if the CMS signature value in
\texttt{/Contents} is wrong?

A correct implementation should reject the signature immediately --- if
the cryptographic value doesn't verify, nothing else matters. What the
Bochum team found was that some viewers never checked.

In the most extreme cases, the \texttt{/Contents} value could be empty
(all zeros), filled with random bytes, or contain a valid CMS structure
that signed completely different data --- and the viewer would still
display "Signature is valid" with a green checkmark.

\subsubsection{How It Works}\label{_how_it_works}

The attack requires no special tools and no knowledge of the signer's
private key. The attacker takes a legitimately signed PDF, modifies the
document content (changes text, replaces pages, adds content), and
re-calculates any internal structures (cross-reference table, object
offsets) needed to keep the file syntactically valid. The
\texttt{/Contents} value --- the CMS \texttt{SignedData} blob --- is
left unchanged, even though it now contains a signature over the
\emph{original} document bytes, not the modified ones.

In a correct implementation, the viewer would:

\begin{enumerate}
\def\labelenumi{\arabic{enumi}.}
\item
  Read the \texttt{/ByteRange} to determine which bytes are signed
\item
  Hash those bytes using the digest algorithm specified in the CMS
  structure
\item
  Extract the signature value from the CMS \texttt{SignerInfo}
\item
  Verify the signature over the hash using the signer's public key
\item
  If verification fails → reject the signature
\end{enumerate}

Vulnerable viewers skipped step 4, step 3, or in some cases both. They
parsed the CMS structure to extract the signer's certificate (for
display purposes), checked that the certificate chained to a trust
anchor, and then displayed the result --- without ever verifying that
the signature value was mathematically valid for the current document
content.

\subsubsection{The Variants}\label{_the_variants}

The Bochum team identified three USF variants with increasing severity:

\textbf{USF Variant 1: Manipulated SignerInfo}

The \texttt{SignerInfo.signature} value (the actual RSASSA-PKCS1-v1\_5
or ECDSA bytes) is replaced with random data. A correct implementation
rejects this immediately during RSA/ECDSA verification. Vulnerable
viewers never reached that step.

\textbf{USF Variant 2: Manipulated ContentInfo}

The outer CMS \texttt{ContentInfo} wrapping is modified so the
\texttt{SignedData} structure is structurally invalid --- wrong OIDs,
missing fields, or inconsistent version numbers. A correct
implementation rejects this during CMS parsing. Vulnerable viewers used
lenient parsing that silently ignored structural errors.

\textbf{USF Variant 3: Manipulated \texttt{/Contents} Value}

The entire \texttt{/Contents} hex string is replaced --- either with
zeros, with a truncated CMS structure, or with a CMS structure copied
from a different document. This tests whether the viewer checks the
relationship between the CMS content and the actual document bytes at
all.

\subsubsection{Demonstration: Detecting USF
Vulnerability}\label{_demonstration_detecting_usf_vulnerability}

The following code demonstrates how to detect whether a signed PDF's CMS
signature actually validates cryptographically --- the check that
USF-vulnerable viewers failed to perform. This is the \emph{defense},
not the attack: if you are building a validator, this is the minimum you
must implement.

\begin{Shaded}
\begin{Highlighting}[]
\KeywordTok{import} \ImportTok{org}\OperatorTok{.}\ImportTok{apache}\OperatorTok{.}\ImportTok{pdfbox}\OperatorTok{.}\ImportTok{Loader}\OperatorTok{;}
\KeywordTok{import} \ImportTok{org}\OperatorTok{.}\ImportTok{apache}\OperatorTok{.}\ImportTok{pdfbox}\OperatorTok{.}\ImportTok{pdmodel}\OperatorTok{.}\ImportTok{PDDocument}\OperatorTok{;}
\KeywordTok{import} \ImportTok{org}\OperatorTok{.}\ImportTok{apache}\OperatorTok{.}\ImportTok{pdfbox}\OperatorTok{.}\ImportTok{pdmodel}\OperatorTok{.}\ImportTok{interactive}\OperatorTok{.}\ImportTok{digitalsignature}\OperatorTok{.}\ImportTok{PDSignature}\OperatorTok{;}
\KeywordTok{import} \ImportTok{org}\OperatorTok{.}\ImportTok{bouncycastle}\OperatorTok{.}\ImportTok{cert}\OperatorTok{.}\ImportTok{X509CertificateHolder}\OperatorTok{;}
\KeywordTok{import} \ImportTok{org}\OperatorTok{.}\ImportTok{bouncycastle}\OperatorTok{.}\ImportTok{cert}\OperatorTok{.}\ImportTok{jcajce}\OperatorTok{.}\ImportTok{JcaX509CertificateConverter}\OperatorTok{;}
\KeywordTok{import} \ImportTok{org}\OperatorTok{.}\ImportTok{bouncycastle}\OperatorTok{.}\ImportTok{cms}\OperatorTok{.}\ImportTok{CMSProcessableByteArray}\OperatorTok{;}
\KeywordTok{import} \ImportTok{org}\OperatorTok{.}\ImportTok{bouncycastle}\OperatorTok{.}\ImportTok{cms}\OperatorTok{.}\ImportTok{CMSSignedData}\OperatorTok{;}
\KeywordTok{import} \ImportTok{org}\OperatorTok{.}\ImportTok{bouncycastle}\OperatorTok{.}\ImportTok{cms}\OperatorTok{.}\ImportTok{SignerInformation}\OperatorTok{;}
\KeywordTok{import} \ImportTok{org}\OperatorTok{.}\ImportTok{bouncycastle}\OperatorTok{.}\ImportTok{cms}\OperatorTok{.}\ImportTok{jcajce}\OperatorTok{.}\ImportTok{JcaSimpleSignerInfoVerifierBuilder}\OperatorTok{;}
\KeywordTok{import} \ImportTok{org}\OperatorTok{.}\ImportTok{bouncycastle}\OperatorTok{.}\ImportTok{jce}\OperatorTok{.}\ImportTok{provider}\OperatorTok{.}\ImportTok{BouncyCastleProvider}\OperatorTok{;}
\KeywordTok{import} \ImportTok{org}\OperatorTok{.}\ImportTok{bouncycastle}\OperatorTok{.}\ImportTok{util}\OperatorTok{.}\ImportTok{Store}\OperatorTok{;}

\KeywordTok{import} \ImportTok{java}\OperatorTok{.}\ImportTok{io}\OperatorTok{.}\ImportTok{FileInputStream}\OperatorTok{;}
\KeywordTok{import} \ImportTok{java}\OperatorTok{.}\ImportTok{io}\OperatorTok{.}\ImportTok{InputStream}\OperatorTok{;}
\KeywordTok{import} \ImportTok{java}\OperatorTok{.}\ImportTok{security}\OperatorTok{.}\ImportTok{Security}\OperatorTok{;}
\KeywordTok{import} \ImportTok{java}\OperatorTok{.}\ImportTok{security}\OperatorTok{.}\ImportTok{cert}\OperatorTok{.}\ImportTok{X509Certificate}\OperatorTok{;}
\KeywordTok{import} \ImportTok{java}\OperatorTok{.}\ImportTok{util}\OperatorTok{.}\ImportTok{Collection}\OperatorTok{;}

\CommentTok{/**}
 \CommentTok{*}\NormalTok{ Demonstrates the cryptographic signature verification that USF}\CommentTok{{-}}\NormalTok{vulnerable}
 \CommentTok{*}\NormalTok{ viewers failed to perform}\CommentTok{. }\NormalTok{This is the MINIMUM check for signature integrity}\CommentTok{.}
 \CommentTok{*/}
\KeywordTok{public} \KeywordTok{class}\NormalTok{ UsfDefenseCheck }\OperatorTok{\{}

    \DataTypeTok{static} \OperatorTok{\{}
        \BuiltInTok{Security}\OperatorTok{.}\FunctionTok{addProvider}\OperatorTok{(}\KeywordTok{new} \FunctionTok{BouncyCastleProvider}\OperatorTok{());}
    \OperatorTok{\}}

    \KeywordTok{public} \DataTypeTok{static} \DataTypeTok{void} \FunctionTok{verifySignature}\OperatorTok{(}\DataTypeTok{byte}\OperatorTok{[]}\NormalTok{ pdfBytes}\OperatorTok{)} \KeywordTok{throws} \BuiltInTok{Exception} \OperatorTok{\{}
\NormalTok{        PDDocument doc }\OperatorTok{=}\NormalTok{ Loader}\OperatorTok{.}\FunctionTok{loadPDF}\OperatorTok{(}\NormalTok{pdfBytes}\OperatorTok{);}

        \ControlFlowTok{for} \OperatorTok{(}\NormalTok{PDSignature sig }\OperatorTok{:}\NormalTok{ doc}\OperatorTok{.}\FunctionTok{getSignatureDictionaries}\OperatorTok{())} \OperatorTok{\{}
            \CommentTok{// Step 1: Extract the signed byte ranges}
            \DataTypeTok{byte}\OperatorTok{[]}\NormalTok{ signedContent }\OperatorTok{=}\NormalTok{ sig}\OperatorTok{.}\FunctionTok{getSignedContent}\OperatorTok{(}
                    \KeywordTok{new} \BuiltInTok{ByteArrayInputStream}\OperatorTok{(}\NormalTok{pdfBytes}\OperatorTok{));} 

            \CommentTok{// Step 2: Extract the CMS SignedData from /Contents}
            \DataTypeTok{byte}\OperatorTok{[]}\NormalTok{ cmsBytes }\OperatorTok{=}\NormalTok{ sig}\OperatorTok{.}\FunctionTok{getContents}\OperatorTok{(}
                    \KeywordTok{new} \BuiltInTok{ByteArrayInputStream}\OperatorTok{(}\NormalTok{pdfBytes}\OperatorTok{));} 

            \CommentTok{// Step 3: Parse the CMS structure}
\NormalTok{            CMSSignedData cmsSignedData }\OperatorTok{=} \KeywordTok{new} \FunctionTok{CMSSignedData}\OperatorTok{(}
                    \KeywordTok{new} \FunctionTok{CMSProcessableByteArray}\OperatorTok{(}\NormalTok{signedContent}\OperatorTok{),}
\NormalTok{                    cmsBytes}\OperatorTok{);} 

            \CommentTok{// Step 4: Verify each signer\textquotesingle{}s cryptographic signature}
\NormalTok{            Store}\OperatorTok{\textless{}}\NormalTok{X509CertificateHolder}\OperatorTok{\textgreater{}}\NormalTok{ certStore }\OperatorTok{=}
\NormalTok{                    cmsSignedData}\OperatorTok{.}\FunctionTok{getCertificates}\OperatorTok{();}

            \ControlFlowTok{for} \OperatorTok{(}\NormalTok{SignerInformation signer }\OperatorTok{:}
\NormalTok{                    cmsSignedData}\OperatorTok{.}\FunctionTok{getSignerInfos}\OperatorTok{().}\FunctionTok{getSigners}\OperatorTok{())} \OperatorTok{\{}

                \BuiltInTok{Collection}\OperatorTok{\textless{}}\NormalTok{X509CertificateHolder}\OperatorTok{\textgreater{}}\NormalTok{ matches }\OperatorTok{=}
\NormalTok{                        certStore}\OperatorTok{.}\FunctionTok{getMatches}\OperatorTok{(}\NormalTok{signer}\OperatorTok{.}\FunctionTok{getSID}\OperatorTok{());}

                \ControlFlowTok{if} \OperatorTok{(}\NormalTok{matches}\OperatorTok{.}\FunctionTok{isEmpty}\OperatorTok{())} \OperatorTok{\{}
                    \BuiltInTok{System}\OperatorTok{.}\FunctionTok{err}\OperatorTok{.}\FunctionTok{println}\OperatorTok{(}
                            \StringTok{"FAIL: No certificate found for signer"}\OperatorTok{);} 
                    \ControlFlowTok{continue}\OperatorTok{;}
                \OperatorTok{\}}

\NormalTok{                X509CertificateHolder certHolder }\OperatorTok{=}\NormalTok{ matches}\OperatorTok{.}\FunctionTok{iterator}\OperatorTok{().}\FunctionTok{next}\OperatorTok{();}
                \BuiltInTok{X509Certificate}\NormalTok{ cert }\OperatorTok{=} \KeywordTok{new} \FunctionTok{JcaX509CertificateConverter}\OperatorTok{()}
                        \OperatorTok{.}\FunctionTok{setProvider}\OperatorTok{(}\StringTok{"BC"}\OperatorTok{)}
                        \OperatorTok{.}\FunctionTok{getCertificate}\OperatorTok{(}\NormalTok{certHolder}\OperatorTok{);}

                \CommentTok{// THIS IS THE CHECK THAT USF{-}VULNERABLE VIEWERS SKIPPED}
                \DataTypeTok{boolean}\NormalTok{ cryptoValid }\OperatorTok{=}\NormalTok{ signer}\OperatorTok{.}\FunctionTok{verify}\OperatorTok{(}
                        \KeywordTok{new} \FunctionTok{JcaSimpleSignerInfoVerifierBuilder}\OperatorTok{()}
                                \OperatorTok{.}\FunctionTok{setProvider}\OperatorTok{(}\StringTok{"BC"}\OperatorTok{)}
                                \OperatorTok{.}\FunctionTok{build}\OperatorTok{(}\NormalTok{cert}\OperatorTok{));} 

                \ControlFlowTok{if} \OperatorTok{(!}\NormalTok{cryptoValid}\OperatorTok{)} \OperatorTok{\{}
                    \BuiltInTok{System}\OperatorTok{.}\FunctionTok{err}\OperatorTok{.}\FunctionTok{println}\OperatorTok{(}
                            \StringTok{"FAIL: Cryptographic signature verification "}
                            \OperatorTok{+} \StringTok{"FAILED for signer: "} \OperatorTok{+}\NormalTok{ cert}\OperatorTok{.}\FunctionTok{getSubjectX500Principal}\OperatorTok{());}
                    \BuiltInTok{System}\OperatorTok{.}\FunctionTok{err}\OperatorTok{.}\FunctionTok{println}\OperatorTok{(}
                            \StringTok{"This document may have been tampered with (USF attack)"}\OperatorTok{);}
                \OperatorTok{\}} \ControlFlowTok{else} \OperatorTok{\{}
                    \BuiltInTok{System}\OperatorTok{.}\FunctionTok{out}\OperatorTok{.}\FunctionTok{println}\OperatorTok{(}
                            \StringTok{"PASS: Cryptographic signature valid for: "}
                            \OperatorTok{+}\NormalTok{ cert}\OperatorTok{.}\FunctionTok{getSubjectX500Principal}\OperatorTok{());}
                \OperatorTok{\}}
            \OperatorTok{\}}
        \OperatorTok{\}}

\NormalTok{        doc}\OperatorTok{.}\FunctionTok{close}\OperatorTok{();}
    \OperatorTok{\}}
\OperatorTok{\}}
\end{Highlighting}
\end{Shaded}

\begin{itemize}
\item
  Extracts the bytes covered by \texttt{/ByteRange} --- this is the data
  that was supposedly signed
\item
  Extracts the raw CMS \texttt{SignedData} blob from the
  \texttt{/Contents} hex string
\item
  Creates a CMS structure with the signed content as the detached data
  --- BouncyCastle will use this for digest comparison
\item
  If no certificate matches the \texttt{SignerInfo.sid}, the CMS
  structure is malformed --- reject immediately
\item
  \textbf{This is the critical line.} \texttt{signer.verify()} performs
  the actual RSA/ECDSA verification: it computes the digest over the
  signed content, verifies the signed attributes, and checks the
  signature value against the signer's public key. USF-vulnerable
  viewers never called the equivalent of this method.
\end{itemize}

If your signature validation code does not include a step equivalent to
\texttt{signer.verify()} --- an actual cryptographic verification of the
signature value against the document bytes --- your code is vulnerable
to Universal Signature Forgery. This is non-negotiable. It must happen
before you display any "valid" indicator to the user.

\subsubsection{Constructing a USF Attack (Proof of
Concept)}\label{_constructing_a_usf_attack_proof_of_concept}

To understand USF fully, here is how an attacker constructs one. This is
\emph{not} production exploit code --- it is a minimal demonstration of
the mechanics that every PDF signature validator must defend against.

\begin{Shaded}
\begin{Highlighting}[]
\KeywordTok{import} \ImportTok{org}\OperatorTok{.}\ImportTok{apache}\OperatorTok{.}\ImportTok{pdfbox}\OperatorTok{.}\ImportTok{Loader}\OperatorTok{;}
\KeywordTok{import} \ImportTok{org}\OperatorTok{.}\ImportTok{apache}\OperatorTok{.}\ImportTok{pdfbox}\OperatorTok{.}\ImportTok{pdmodel}\OperatorTok{.}\ImportTok{PDDocument}\OperatorTok{;}
\KeywordTok{import} \ImportTok{org}\OperatorTok{.}\ImportTok{apache}\OperatorTok{.}\ImportTok{pdfbox}\OperatorTok{.}\ImportTok{pdmodel}\OperatorTok{.}\ImportTok{PDPage}\OperatorTok{;}
\KeywordTok{import} \ImportTok{org}\OperatorTok{.}\ImportTok{apache}\OperatorTok{.}\ImportTok{pdfbox}\OperatorTok{.}\ImportTok{pdmodel}\OperatorTok{.}\ImportTok{PDPageContentStream}\OperatorTok{;}
\KeywordTok{import} \ImportTok{org}\OperatorTok{.}\ImportTok{apache}\OperatorTok{.}\ImportTok{pdfbox}\OperatorTok{.}\ImportTok{pdmodel}\OperatorTok{.}\ImportTok{font}\OperatorTok{.}\ImportTok{PDType1Font}\OperatorTok{;}
\KeywordTok{import} \ImportTok{org}\OperatorTok{.}\ImportTok{apache}\OperatorTok{.}\ImportTok{pdfbox}\OperatorTok{.}\ImportTok{pdmodel}\OperatorTok{.}\ImportTok{font}\OperatorTok{.}\ImportTok{Standard14Fonts}\OperatorTok{;}

\KeywordTok{import} \ImportTok{java}\OperatorTok{.}\ImportTok{io}\OperatorTok{.}\ImportTok{ByteArrayOutputStream}\OperatorTok{;}
\KeywordTok{import} \ImportTok{java}\OperatorTok{.}\ImportTok{io}\OperatorTok{.}\ImportTok{File}\OperatorTok{;}
\KeywordTok{import} \ImportTok{java}\OperatorTok{.}\ImportTok{nio}\OperatorTok{.}\ImportTok{file}\OperatorTok{.}\ImportTok{Files}\OperatorTok{;}

\CommentTok{/**}
 \CommentTok{*}\NormalTok{ Demonstrates USF attack mechanics}\CommentTok{. }\NormalTok{Takes a signed PDF}\CommentTok{,}\NormalTok{ modifies the visible}
 \CommentTok{*}\NormalTok{ content}\CommentTok{,}\NormalTok{ and writes the result}\CommentTok{.}\NormalTok{ The CMS signature value is left unchanged }\CommentTok{{-}{-}}
 \CommentTok{*}\NormalTok{ it now signs the ORIGINAL content}\CommentTok{,}\NormalTok{ not the modified content}\CommentTok{.}
 \CommentTok{*}
 \CommentTok{*}\NormalTok{ A correct validator will reject this because the digest over the byte range}
 \CommentTok{*}\NormalTok{ no longer matches the digest in the CMS SignerInfo}\CommentTok{.}
 \CommentTok{*}\NormalTok{ A USF}\CommentTok{{-}}\NormalTok{vulnerable validator will show }\CommentTok{"}\NormalTok{Signature is valid}\CommentTok{."}
 \CommentTok{*/}
\KeywordTok{public} \KeywordTok{class}\NormalTok{ UsfAttackDemo }\OperatorTok{\{}

    \KeywordTok{public} \DataTypeTok{static} \DataTypeTok{void} \FunctionTok{main}\OperatorTok{(}\BuiltInTok{String}\OperatorTok{[]}\NormalTok{ args}\OperatorTok{)} \KeywordTok{throws} \BuiltInTok{Exception} \OperatorTok{\{}
        \DataTypeTok{byte}\OperatorTok{[]}\NormalTok{ originalPdf }\OperatorTok{=}\NormalTok{ Files}\OperatorTok{.}\FunctionTok{readAllBytes}\OperatorTok{(}
                \KeywordTok{new} \BuiltInTok{File}\OperatorTok{(}\StringTok{"signed{-}document.pdf"}\OperatorTok{).}\FunctionTok{toPath}\OperatorTok{());}

        \CommentTok{// Load the signed PDF}
\NormalTok{        PDDocument doc }\OperatorTok{=}\NormalTok{ Loader}\OperatorTok{.}\FunctionTok{loadPDF}\OperatorTok{(}\NormalTok{originalPdf}\OperatorTok{);}

        \CommentTok{// Add a new page with attacker{-}controlled content}
\NormalTok{        PDPage attackPage }\OperatorTok{=} \KeywordTok{new} \FunctionTok{PDPage}\OperatorTok{();}
\NormalTok{        doc}\OperatorTok{.}\FunctionTok{addPage}\OperatorTok{(}\NormalTok{attackPage}\OperatorTok{);} 

\NormalTok{        PDPageContentStream cs }\OperatorTok{=} \KeywordTok{new} \FunctionTok{PDPageContentStream}\OperatorTok{(}\NormalTok{doc}\OperatorTok{,}\NormalTok{ attackPage}\OperatorTok{);}
\NormalTok{        cs}\OperatorTok{.}\FunctionTok{beginText}\OperatorTok{();}
\NormalTok{        cs}\OperatorTok{.}\FunctionTok{setFont}\OperatorTok{(}
                \KeywordTok{new} \FunctionTok{PDType1Font}\OperatorTok{(}\NormalTok{Standard14Fonts}\OperatorTok{.}\FunctionTok{FontName}\OperatorTok{.}\FunctionTok{HELVETICA\_BOLD}\OperatorTok{),} \DecValTok{24}\OperatorTok{);}
\NormalTok{        cs}\OperatorTok{.}\FunctionTok{newLineAtOffset}\OperatorTok{(}\DecValTok{50}\OperatorTok{,} \DecValTok{700}\OperatorTok{);}
\NormalTok{        cs}\OperatorTok{.}\FunctionTok{showText}\OperatorTok{(}\StringTok{"APPROVED {-} Payment of $1,000,000 authorized"}\OperatorTok{);} 
\NormalTok{        cs}\OperatorTok{.}\FunctionTok{endText}\OperatorTok{();}
\NormalTok{        cs}\OperatorTok{.}\FunctionTok{close}\OperatorTok{();}

        \CommentTok{// Save with full rewrite — this changes object offsets,}
        \CommentTok{// xref table, and file length. The /ByteRange in the original}
        \CommentTok{// signature now covers different bytes, but the /Contents}
        \CommentTok{// CMS value still contains the hash of the ORIGINAL bytes.}
        \BuiltInTok{ByteArrayOutputStream}\NormalTok{ baos }\OperatorTok{=} \KeywordTok{new} \BuiltInTok{ByteArrayOutputStream}\OperatorTok{();}
\NormalTok{        doc}\OperatorTok{.}\FunctionTok{save}\OperatorTok{(}\NormalTok{baos}\OperatorTok{);} 
\NormalTok{        doc}\OperatorTok{.}\FunctionTok{close}\OperatorTok{();}

\NormalTok{        Files}\OperatorTok{.}\FunctionTok{write}\OperatorTok{(}
                \KeywordTok{new} \BuiltInTok{File}\OperatorTok{(}\StringTok{"usf{-}attack{-}result.pdf"}\OperatorTok{).}\FunctionTok{toPath}\OperatorTok{(),}\NormalTok{ baos}\OperatorTok{.}\FunctionTok{toByteArray}\OperatorTok{());}

        \CommentTok{// A correct validator: REJECTS (digest mismatch)}
        \CommentTok{// A USF{-}vulnerable validator: shows green checkmark}
    \OperatorTok{\}}
\OperatorTok{\}}
\end{Highlighting}
\end{Shaded}

\begin{itemize}
\item
  Modifying the document structure after it was signed --- adding a
  page, changing text, anything
\item
  The attacker can insert any content they want
\item
  A full \texttt{save()} (not \texttt{saveIncremental()}) rewrites the
  entire file, changing byte offsets everywhere. The original
  \texttt{/ByteRange} and \texttt{/Contents} are preserved structurally,
  but the bytes they reference have changed. The CMS signature now signs
  data that no longer exists at those byte positions.
\end{itemize}

This code produces a PDF that any \emph{correct} validator will reject.
The point is that in 2019, most validators were not correct. The digest
computed over the new byte range content will not match the digest
inside the CMS \texttt{SignerInfo.digestAlgorithm} --- the very first
check in proper CMS verification detects this. USF worked because
viewers displayed the signer's certificate information \emph{without
performing this check}.

\subsection{Attack Class 2: Incremental Saving Attack
(ISA)}\label{_attack_class_2_incremental_saving_attack_isa}

\subsubsection{The Concept}\label{_the_concept_2}

The Incremental Saving Attack is more subtle than USF and exploits a
feature that is fundamental to how PDF signatures work:
\textbf{incremental updates}.

When a PDF is signed, the signing process:

\begin{enumerate}
\def\labelenumi{\arabic{enumi}.}
\item
  Appends the signature dictionary and any new objects to the end of the
  file
\item
  Appends a new cross-reference section pointing to the new/modified
  objects
\item
  Appends a new trailer
\item
  The \texttt{/ByteRange} in the signature covers everything from the
  beginning of the file through the end of this new revision --- with a
  gap for the \texttt{/Contents} value
\end{enumerate}

The critical insight is that PDF is designed to allow \emph{additional}
incremental updates after signing. This is intentional --- you need to
be able to add a second signature to an already-signed document, or fill
in form fields on a signed form. ISO 32000-2:2020 §12.8.1 explicitly
allows certain modifications after signing, depending on the signature's
\texttt{/TransformMethod} and permission level.

The question is: does the viewer correctly distinguish between
\emph{permissible} incremental updates (adding another signature) and
\emph{impermissible} ones (changing the page content)?

The Bochum team found that most viewers did not.

\subsubsection{How It Works}\label{_how_it_works_2}

The attacker takes a legitimately signed PDF and appends a new
incremental update --- exactly the way PDFBox's
\texttt{saveIncremental()} method does. But instead of adding a second
signature or filling in a form field, the attacker appends objects that
change the visible content: modified page content streams, replaced
images, rewritten text, or entirely new pages.

The original signature remains intact --- its byte range still covers
the original revision, and its CMS signature is still mathematically
valid over those bytes. The appended revision contains the
modifications, complete with a new cross-reference section and trailer
that redirect the viewer to the modified objects.

A correct validator must:

\begin{enumerate}
\def\labelenumi{\arabic{enumi}.}
\item
  Verify the signature is cryptographically valid (steps from USF
  defense)
\item
  Determine the document's revision at the time of signing (the byte
  range endpoint)
\item
  Compare the current document state to the signed revision
\item
  If there are changes beyond the signed revision, evaluate whether
  those changes are permissible under the signature's DocMDP permissions
\item
  If impermissible changes are detected, display a clear warning: "The
  document has been modified after signing"
\end{enumerate}

Vulnerable viewers either:

\begin{itemize}
\item
  Did not compare the current document state to the signed revision at
  all
\item
  Compared but accepted all incremental updates as permissible
\item
  Only checked whether the appended data was "syntactically valid PDF"
  rather than whether it changed the visible content
\end{itemize}

\subsubsection{The Variants}\label{_the_variants_2}

\textbf{ISA Variant 1: Content Replacement via Incremental Update}

The attacker appends a new revision that redefines existing page content
stream objects. The cross-reference in the new revision points to new
objects at the end of the file that replace the original content. The
viewer renders the new content; the signature covers the old content.

\textbf{ISA Variant 2: Page Insertion/Deletion}

The attacker appends a new revision that modifies the page tree ---
adding new pages, removing existing pages, or reordering them. The
\texttt{/Pages} dictionary in the new revision points to a different set
of page objects.

\textbf{ISA Variant 3: Annotation and Form Field Manipulation}

The attacker appends annotations or modifies form field values. This
variant is particularly insidious because some annotation changes
\emph{are} legitimately allowed after signing (e.g., adding signature
fields for subsequent signers). The boundary between "allowed" and "not
allowed" is defined by the DocMDP permission level, and most viewers did
not enforce it correctly.

\subsubsection{Demonstration: ISA Attack
Mechanics}\label{_demonstration_isa_attack_mechanics}

\begin{Shaded}
\begin{Highlighting}[]
\KeywordTok{import} \ImportTok{org}\OperatorTok{.}\ImportTok{apache}\OperatorTok{.}\ImportTok{pdfbox}\OperatorTok{.}\ImportTok{Loader}\OperatorTok{;}
\KeywordTok{import} \ImportTok{org}\OperatorTok{.}\ImportTok{apache}\OperatorTok{.}\ImportTok{pdfbox}\OperatorTok{.}\ImportTok{cos}\OperatorTok{.}\ImportTok{COSDictionary}\OperatorTok{;}
\KeywordTok{import} \ImportTok{org}\OperatorTok{.}\ImportTok{apache}\OperatorTok{.}\ImportTok{pdfbox}\OperatorTok{.}\ImportTok{cos}\OperatorTok{.}\ImportTok{COSName}\OperatorTok{;}
\KeywordTok{import} \ImportTok{org}\OperatorTok{.}\ImportTok{apache}\OperatorTok{.}\ImportTok{pdfbox}\OperatorTok{.}\ImportTok{cos}\OperatorTok{.}\ImportTok{COSArray}\OperatorTok{;}
\KeywordTok{import} \ImportTok{org}\OperatorTok{.}\ImportTok{apache}\OperatorTok{.}\ImportTok{pdfbox}\OperatorTok{.}\ImportTok{pdmodel}\OperatorTok{.}\ImportTok{PDDocument}\OperatorTok{;}
\KeywordTok{import} \ImportTok{org}\OperatorTok{.}\ImportTok{apache}\OperatorTok{.}\ImportTok{pdfbox}\OperatorTok{.}\ImportTok{pdmodel}\OperatorTok{.}\ImportTok{PDPage}\OperatorTok{;}
\KeywordTok{import} \ImportTok{org}\OperatorTok{.}\ImportTok{apache}\OperatorTok{.}\ImportTok{pdfbox}\OperatorTok{.}\ImportTok{pdmodel}\OperatorTok{.}\ImportTok{PDPageContentStream}\OperatorTok{;}
\KeywordTok{import} \ImportTok{org}\OperatorTok{.}\ImportTok{apache}\OperatorTok{.}\ImportTok{pdfbox}\OperatorTok{.}\ImportTok{pdmodel}\OperatorTok{.}\ImportTok{common}\OperatorTok{.}\ImportTok{PDRectangle}\OperatorTok{;}
\KeywordTok{import} \ImportTok{org}\OperatorTok{.}\ImportTok{apache}\OperatorTok{.}\ImportTok{pdfbox}\OperatorTok{.}\ImportTok{pdmodel}\OperatorTok{.}\ImportTok{font}\OperatorTok{.}\ImportTok{PDType1Font}\OperatorTok{;}
\KeywordTok{import} \ImportTok{org}\OperatorTok{.}\ImportTok{apache}\OperatorTok{.}\ImportTok{pdfbox}\OperatorTok{.}\ImportTok{pdmodel}\OperatorTok{.}\ImportTok{font}\OperatorTok{.}\ImportTok{Standard14Fonts}\OperatorTok{;}

\KeywordTok{import} \ImportTok{java}\OperatorTok{.}\ImportTok{io}\OperatorTok{.}\ImportTok{ByteArrayOutputStream}\OperatorTok{;}
\KeywordTok{import} \ImportTok{java}\OperatorTok{.}\ImportTok{io}\OperatorTok{.}\ImportTok{File}\OperatorTok{;}
\KeywordTok{import} \ImportTok{java}\OperatorTok{.}\ImportTok{nio}\OperatorTok{.}\ImportTok{file}\OperatorTok{.}\ImportTok{Files}\OperatorTok{;}

\CommentTok{/**}
 \CommentTok{*}\NormalTok{ Demonstrates ISA mechanics}\CommentTok{. }\NormalTok{Takes a signed PDF and appends an incremental}
 \CommentTok{*}\NormalTok{ update that replaces the visible content of page }\CommentTok{1.}
 \CommentTok{*}
 \CommentTok{*}\NormalTok{ The original signature remains cryptographically valid over the original}
 \CommentTok{*}\NormalTok{ byte range}\CommentTok{.}\NormalTok{ The appended revision changes what the user sees}\CommentTok{.}
 \CommentTok{*}
 \CommentTok{*}\NormalTok{ A correct validator detects the impermissible modification}\CommentTok{.}
 \CommentTok{*}\NormalTok{ An ISA}\CommentTok{{-}}\NormalTok{vulnerable validator shows }\CommentTok{"}\NormalTok{Signature is valid}\CommentTok{"}\NormalTok{ on the modified content}\CommentTok{.}
 \CommentTok{*/}
\KeywordTok{public} \KeywordTok{class}\NormalTok{ IsaAttackDemo }\OperatorTok{\{}

    \KeywordTok{public} \DataTypeTok{static} \DataTypeTok{void} \FunctionTok{main}\OperatorTok{(}\BuiltInTok{String}\OperatorTok{[]}\NormalTok{ args}\OperatorTok{)} \KeywordTok{throws} \BuiltInTok{Exception} \OperatorTok{\{}
        \DataTypeTok{byte}\OperatorTok{[]}\NormalTok{ signedPdf }\OperatorTok{=}\NormalTok{ Files}\OperatorTok{.}\FunctionTok{readAllBytes}\OperatorTok{(}
                \KeywordTok{new} \BuiltInTok{File}\OperatorTok{(}\StringTok{"signed{-}document.pdf"}\OperatorTok{).}\FunctionTok{toPath}\OperatorTok{());}

        \CommentTok{// Load for incremental save — preserves the original bytes}
\NormalTok{        PDDocument doc }\OperatorTok{=}\NormalTok{ Loader}\OperatorTok{.}\FunctionTok{loadPDF}\OperatorTok{(}\NormalTok{signedPdf}\OperatorTok{);}

        \CommentTok{// Get the first page and replace its content}
\NormalTok{        PDPage page }\OperatorTok{=}\NormalTok{ doc}\OperatorTok{.}\FunctionTok{getPage}\OperatorTok{(}\DecValTok{0}\OperatorTok{);}

        \CommentTok{// Create a new content stream that overwrites existing content}
\NormalTok{        PDPageContentStream cs }\OperatorTok{=} \KeywordTok{new} \FunctionTok{PDPageContentStream}\OperatorTok{(}
\NormalTok{                doc}\OperatorTok{,}\NormalTok{ page}\OperatorTok{,}
\NormalTok{                PDPageContentStream}\OperatorTok{.}\FunctionTok{AppendMode}\OperatorTok{.}\FunctionTok{OVERWRITE}\OperatorTok{,}
                \KeywordTok{true}\OperatorTok{);} 

\NormalTok{        cs}\OperatorTok{.}\FunctionTok{beginText}\OperatorTok{();}
\NormalTok{        cs}\OperatorTok{.}\FunctionTok{setFont}\OperatorTok{(}
                \KeywordTok{new} \FunctionTok{PDType1Font}\OperatorTok{(}\NormalTok{Standard14Fonts}\OperatorTok{.}\FunctionTok{FontName}\OperatorTok{.}\FunctionTok{HELVETICA\_BOLD}\OperatorTok{),} \DecValTok{36}\OperatorTok{);}
\NormalTok{        cs}\OperatorTok{.}\FunctionTok{newLineAtOffset}\OperatorTok{(}\DecValTok{50}\OperatorTok{,} \DecValTok{700}\OperatorTok{);}
\NormalTok{        cs}\OperatorTok{.}\FunctionTok{showText}\OperatorTok{(}\StringTok{"CONTRACT VOIDED"}\OperatorTok{);} 
\NormalTok{        cs}\OperatorTok{.}\FunctionTok{endText}\OperatorTok{();}
\NormalTok{        cs}\OperatorTok{.}\FunctionTok{close}\OperatorTok{();}

        \CommentTok{// Save incrementally — appends new objects + xref + trailer}
        \CommentTok{// The original signature\textquotesingle{}s byte range is UNTOUCHED}
        \BuiltInTok{ByteArrayOutputStream}\NormalTok{ baos }\OperatorTok{=} \KeywordTok{new} \BuiltInTok{ByteArrayOutputStream}\OperatorTok{();}
\NormalTok{        doc}\OperatorTok{.}\FunctionTok{saveIncremental}\OperatorTok{(}\NormalTok{baos}\OperatorTok{);} 
\NormalTok{        doc}\OperatorTok{.}\FunctionTok{close}\OperatorTok{();}

\NormalTok{        Files}\OperatorTok{.}\FunctionTok{write}\OperatorTok{(}
                \KeywordTok{new} \BuiltInTok{File}\OperatorTok{(}\StringTok{"isa{-}attack{-}result.pdf"}\OperatorTok{).}\FunctionTok{toPath}\OperatorTok{(),}\NormalTok{ baos}\OperatorTok{.}\FunctionTok{toByteArray}\OperatorTok{());}

        \CommentTok{// The file now looks like this:}
        \CommentTok{// [original PDF bytes ... signature ... \%\%EOF]  \textless{}{-}{-} signed revision}
        \CommentTok{// [new objects ... new xref ... new trailer ... \%\%EOF]  \textless{}{-}{-} attacker\textquotesingle{}s appendage}
        \CommentTok{//}
        \CommentTok{// The signature covers only the first section.}
        \CommentTok{// The viewer renders the second section.}
        \CommentTok{// A correct validator: warns "Document modified after signing"}
        \CommentTok{// An ISA{-}vulnerable validator: shows green checkmark}
    \OperatorTok{\}}
\OperatorTok{\}}
\end{Highlighting}
\end{Shaded}

\begin{itemize}
\item
  \texttt{AppendMode.OVERWRITE} replaces the existing content stream.
  PDFBox tracks this as a modified object, which will be written in the
  incremental update
\item
  The attacker-controlled content --- can be anything
\item
  \texttt{saveIncremental()} appends to the existing file bytes without
  modifying them. The original signature's \texttt{/ByteRange} and
  \texttt{/Contents} remain at their original byte positions. The new
  revision simply adds new objects and a new cross-reference section
  after the original \texttt{\%\%EOF}
\end{itemize}

\subsubsection{Detecting ISA: Revision
Comparison}\label{_detecting_isa_revision_comparison}

The defense against ISA requires comparing the document as rendered from
the signed revision to the document as rendered from the current
(potentially modified) revision.

\begin{Shaded}
\begin{Highlighting}[]
\KeywordTok{import} \ImportTok{org}\OperatorTok{.}\ImportTok{apache}\OperatorTok{.}\ImportTok{pdfbox}\OperatorTok{.}\ImportTok{Loader}\OperatorTok{;}
\KeywordTok{import} \ImportTok{org}\OperatorTok{.}\ImportTok{apache}\OperatorTok{.}\ImportTok{pdfbox}\OperatorTok{.}\ImportTok{pdmodel}\OperatorTok{.}\ImportTok{PDDocument}\OperatorTok{;}
\KeywordTok{import} \ImportTok{org}\OperatorTok{.}\ImportTok{apache}\OperatorTok{.}\ImportTok{pdfbox}\OperatorTok{.}\ImportTok{pdmodel}\OperatorTok{.}\ImportTok{interactive}\OperatorTok{.}\ImportTok{digitalsignature}\OperatorTok{.}\ImportTok{PDSignature}\OperatorTok{;}

\KeywordTok{import} \ImportTok{java}\OperatorTok{.}\ImportTok{io}\OperatorTok{.}\ImportTok{ByteArrayInputStream}\OperatorTok{;}
\KeywordTok{import} \ImportTok{java}\OperatorTok{.}\ImportTok{util}\OperatorTok{.}\ImportTok{Arrays}\OperatorTok{;}

\CommentTok{/**}
 \CommentTok{*}\NormalTok{ Demonstrates how to extract the signed revision of a PDF and compare it}
 \CommentTok{*}\NormalTok{ to the full file }\CommentTok{—}\NormalTok{ the core defense against ISA}\CommentTok{.}
 \CommentTok{*/}
\KeywordTok{public} \KeywordTok{class}\NormalTok{ IsaDefenseCheck }\OperatorTok{\{}

    \KeywordTok{public} \DataTypeTok{static} \DataTypeTok{void} \FunctionTok{checkForAppendedContent}\OperatorTok{(}\DataTypeTok{byte}\OperatorTok{[]}\NormalTok{ pdfBytes}\OperatorTok{)}
            \KeywordTok{throws} \BuiltInTok{Exception} \OperatorTok{\{}

\NormalTok{        PDDocument doc }\OperatorTok{=}\NormalTok{ Loader}\OperatorTok{.}\FunctionTok{loadPDF}\OperatorTok{(}\NormalTok{pdfBytes}\OperatorTok{);}

        \ControlFlowTok{for} \OperatorTok{(}\NormalTok{PDSignature sig }\OperatorTok{:}\NormalTok{ doc}\OperatorTok{.}\FunctionTok{getSignatureDictionaries}\OperatorTok{())} \OperatorTok{\{}
            \DataTypeTok{int}\OperatorTok{[]}\NormalTok{ byteRange }\OperatorTok{=}\NormalTok{ sig}\OperatorTok{.}\FunctionTok{getByteRange}\OperatorTok{();} 

            \ControlFlowTok{if} \OperatorTok{(}\NormalTok{byteRange }\OperatorTok{==} \KeywordTok{null} \OperatorTok{||}\NormalTok{ byteRange}\OperatorTok{.}\FunctionTok{length} \OperatorTok{!=} \DecValTok{4}\OperatorTok{)} \OperatorTok{\{}
                \BuiltInTok{System}\OperatorTok{.}\FunctionTok{err}\OperatorTok{.}\FunctionTok{println}\OperatorTok{(}\StringTok{"FAIL: Invalid or missing /ByteRange"}\OperatorTok{);}
                \ControlFlowTok{continue}\OperatorTok{;}
            \OperatorTok{\}}

            \CommentTok{// The signed region ends at byteRange[2] + byteRange[3]}
            \DataTypeTok{int}\NormalTok{ signedEnd }\OperatorTok{=}\NormalTok{ byteRange}\OperatorTok{[}\DecValTok{2}\OperatorTok{]} \OperatorTok{+}\NormalTok{ byteRange}\OperatorTok{[}\DecValTok{3}\OperatorTok{];} 

            \CommentTok{// Check: does the file extend beyond the signed region?}
            \ControlFlowTok{if} \OperatorTok{(}\NormalTok{pdfBytes}\OperatorTok{.}\FunctionTok{length} \OperatorTok{\textgreater{}}\NormalTok{ signedEnd}\OperatorTok{)} \OperatorTok{\{}
                \DataTypeTok{int}\NormalTok{ extraBytes }\OperatorTok{=}\NormalTok{ pdfBytes}\OperatorTok{.}\FunctionTok{length} \OperatorTok{{-}}\NormalTok{ signedEnd}\OperatorTok{;}
                \BuiltInTok{System}\OperatorTok{.}\FunctionTok{err}\OperatorTok{.}\FunctionTok{println}\OperatorTok{(}
                        \StringTok{"WARNING: File contains "} \OperatorTok{+}\NormalTok{ extraBytes}
                        \OperatorTok{+} \StringTok{" bytes AFTER the signed byte range"}\OperatorTok{);} 
                \BuiltInTok{System}\OperatorTok{.}\FunctionTok{err}\OperatorTok{.}\FunctionTok{println}\OperatorTok{(}
                        \StringTok{"This indicates an incremental update after signing."}\OperatorTok{);}
                \BuiltInTok{System}\OperatorTok{.}\FunctionTok{err}\OperatorTok{.}\FunctionTok{println}\OperatorTok{(}
                        \StringTok{"These bytes are NOT covered by the signature."}\OperatorTok{);}
                \BuiltInTok{System}\OperatorTok{.}\FunctionTok{err}\OperatorTok{.}\FunctionTok{println}\OperatorTok{(}
                        \StringTok{"Investigate whether the changes are permissible "}
                        \OperatorTok{+} \StringTok{"(e.g., additional signatures, form fills) "}
                        \OperatorTok{+} \StringTok{"or malicious (content replacement)."}\OperatorTok{);}

                \CommentTok{// Extract the signed revision for comparison}
                \DataTypeTok{byte}\OperatorTok{[]}\NormalTok{ signedRevision }\OperatorTok{=} \BuiltInTok{Arrays}\OperatorTok{.}\FunctionTok{copyOfRange}\OperatorTok{(}
\NormalTok{                        pdfBytes}\OperatorTok{,} \DecValTok{0}\OperatorTok{,}\NormalTok{ signedEnd}\OperatorTok{);} 

                \CommentTok{// You can now load signedRevision as a separate PDDocument}
                \CommentTok{// and compare its rendered pages to the full document}
                \CommentTok{// to determine exactly what changed}
            \OperatorTok{\}} \ControlFlowTok{else} \ControlFlowTok{if} \OperatorTok{(}\NormalTok{pdfBytes}\OperatorTok{.}\FunctionTok{length} \OperatorTok{==}\NormalTok{ signedEnd}\OperatorTok{)} \OperatorTok{\{}
                \BuiltInTok{System}\OperatorTok{.}\FunctionTok{out}\OperatorTok{.}\FunctionTok{println}\OperatorTok{(}
                        \StringTok{"PASS: File ends exactly at the signed byte range. "}
                        \OperatorTok{+} \StringTok{"No incremental updates after signing."}\OperatorTok{);}
            \OperatorTok{\}} \ControlFlowTok{else} \OperatorTok{\{}
                \BuiltInTok{System}\OperatorTok{.}\FunctionTok{err}\OperatorTok{.}\FunctionTok{println}\OperatorTok{(}
                        \StringTok{"FAIL: File is SHORTER than the signed byte range. "}
                        \OperatorTok{+} \StringTok{"File is truncated or corrupted."}\OperatorTok{);}
            \OperatorTok{\}}
        \OperatorTok{\}}

\NormalTok{        doc}\OperatorTok{.}\FunctionTok{close}\OperatorTok{();}
    \OperatorTok{\}}
\OperatorTok{\}}
\end{Highlighting}
\end{Shaded}

\begin{itemize}
\item
  \texttt{/ByteRange} is an array of four integers:
  \texttt{{[}offset1,\ length1,\ offset2,\ length2{]}}. The two ranges
  cover everything before and after the \texttt{/Contents} gap
\item
  The signed content ends at the start of range 2 plus its length ---
  this is the last byte the signature covers
\item
  If the file extends beyond this point, there is appended content that
  the signature does not cover. This is the ISA detection point
\item
  Extracting the signed revision as a standalone PDF allows you to
  render it separately and compare to the current document state
\end{itemize}

A naive "does the file extend beyond the byte range?" check catches
obvious ISA attacks, but it is not sufficient. A legitimate second
signature also extends the file. The real check requires parsing the
appended incremental update and determining \emph{what} changed --- was
it another signature (allowed), a form fill (allowed under DocMDP P=3),
or a content modification (never allowed)? This is the hard part of ISA
defense, and it is where most viewers failed in 2019.

\subsection{Attack Class 3: Signature Wrapping Attack
(SWA)}\label{_attack_class_3_signature_wrapping_attack_swa}

\subsubsection{The Concept}\label{_the_concept_3}

The Signature Wrapping Attack is the most technically sophisticated of
the three original attack classes. Where USF exploits missing
cryptographic verification and ISA exploits missing modification
detection, SWA exploits the way viewers \emph{locate} the signature
within the file.

The core insight: the \texttt{/ByteRange} array tells the viewer which
bytes are signed, but \emph{nothing prevents an attacker from moving the
signed bytes to a different location in the file} --- as long as the
byte values at the \texttt{/ByteRange} offsets still produce the same
hash.

Actually, SWA is even simpler than that. The attacker doesn't need the
same hash. The attacker modifies the \texttt{/ByteRange} array itself to
point to a \emph{different} region of the file --- one that contains
attacker-controlled content --- while keeping the original CMS signature
that was computed over the original content.

Wait --- if the \texttt{/ByteRange} is modified, wouldn't the signature
be invalid? Yes, \emph{if the viewer reads the \texttt{/ByteRange} from
the signed revision}. But many viewers read the \texttt{/ByteRange} from
the \emph{current} cross-reference table, which may have been modified
by an incremental update. The attacker appends a new revision that:

\begin{enumerate}
\def\labelenumi{\arabic{enumi}.}
\item
  Redefines the signature dictionary object to point to a new
  \texttt{/ByteRange}
\item
  Places attacker-controlled content at the byte positions specified by
  the new \texttt{/ByteRange}
\item
  Ensures the CMS signature (in the original \texttt{/Contents}) is
  structurally present
\end{enumerate}

The viewer follows the new cross-reference to the redefined signature
dictionary, reads the attacker's \texttt{/ByteRange}, hashes the
attacker's content, and --- if it doesn't actually verify the hash
against the CMS signature (shades of USF) --- displays "valid."

In some variants, the viewer \emph{does} check the hash, but the
attacker has been clever about which bytes are at the new
\texttt{/ByteRange} positions, potentially including a copy of the
original signed content or exploiting hash-algorithm weaknesses.

\subsubsection{How It Works}\label{_how_it_works_3}

\begin{verbatim}
Original file layout:
  [header ... objects ... page content ... ]
  [signature dict: /ByteRange [0 840 3000 1200] /Contents <CMS...>]
  [xref ... trailer ... %%EOF]

  Bytes 0-839 and 3000-4199 are signed.
  Bytes 840-2999 are the /Contents gap (CMS signature value).
  The CMS signature is valid over hash(bytes[0..839] + bytes[3000..4199]).

Attacker appends an incremental update:
  [... %%EOF]                              <-- original file ends here
  [new object: modified signature dict     <-- attacker redefines the sig dict
      /ByteRange [0 840 5000 1200]         <-- points to NEW byte positions
      /Contents <original CMS bytes>]      <-- keeps the original CMS value
  [attacker content at byte offset 5000]   <-- whatever the attacker wants
  [new xref pointing to new sig dict]
  [new trailer ... %%EOF]
\end{verbatim}

The viewer follows the new xref, finds the redefined signature
dictionary, reads \texttt{/ByteRange\ {[}0\ 840\ 5000\ 1200{]}}, and
hashes bytes at those positions. If it checks the CMS signature against
this hash, it will fail (because the CMS signed the \emph{original} byte
range content). But if the viewer is also vulnerable to USF (skips the
hash verification), the attack succeeds.

SWA is often combined with elements of ISA and USF. The Bochum team
found that many real-world exploits used a combination of attack
techniques.

\subsubsection{Demonstration: Detecting
SWA}\label{_demonstration_detecting_swa}

\begin{Shaded}
\begin{Highlighting}[]
\KeywordTok{import} \ImportTok{org}\OperatorTok{.}\ImportTok{apache}\OperatorTok{.}\ImportTok{pdfbox}\OperatorTok{.}\ImportTok{Loader}\OperatorTok{;}
\KeywordTok{import} \ImportTok{org}\OperatorTok{.}\ImportTok{apache}\OperatorTok{.}\ImportTok{pdfbox}\OperatorTok{.}\ImportTok{pdmodel}\OperatorTok{.}\ImportTok{PDDocument}\OperatorTok{;}
\KeywordTok{import} \ImportTok{org}\OperatorTok{.}\ImportTok{apache}\OperatorTok{.}\ImportTok{pdfbox}\OperatorTok{.}\ImportTok{pdmodel}\OperatorTok{.}\ImportTok{interactive}\OperatorTok{.}\ImportTok{digitalsignature}\OperatorTok{.}\ImportTok{PDSignature}\OperatorTok{;}

\CommentTok{/**}
 \CommentTok{*}\NormalTok{ Demonstrates SWA detection}\CommentTok{:}\NormalTok{ verify that the }\CommentTok{/}\NormalTok{ByteRange covers the entire}
 \CommentTok{*}\NormalTok{ file without gaps }\CommentTok{(}\NormalTok{except the }\CommentTok{/}\NormalTok{Contents gap}\CommentTok{)}\NormalTok{ and without overlaps}\CommentTok{.}
 \CommentTok{*/}
\KeywordTok{public} \KeywordTok{class}\NormalTok{ SwaDefenseCheck }\OperatorTok{\{}

    \KeywordTok{public} \DataTypeTok{static} \DataTypeTok{void} \FunctionTok{checkByteRangeIntegrity}\OperatorTok{(}\DataTypeTok{byte}\OperatorTok{[]}\NormalTok{ pdfBytes}\OperatorTok{)}
            \KeywordTok{throws} \BuiltInTok{Exception} \OperatorTok{\{}

\NormalTok{        PDDocument doc }\OperatorTok{=}\NormalTok{ Loader}\OperatorTok{.}\FunctionTok{loadPDF}\OperatorTok{(}\NormalTok{pdfBytes}\OperatorTok{);}

        \ControlFlowTok{for} \OperatorTok{(}\NormalTok{PDSignature sig }\OperatorTok{:}\NormalTok{ doc}\OperatorTok{.}\FunctionTok{getSignatureDictionaries}\OperatorTok{())} \OperatorTok{\{}
            \DataTypeTok{int}\OperatorTok{[]}\NormalTok{ byteRange }\OperatorTok{=}\NormalTok{ sig}\OperatorTok{.}\FunctionTok{getByteRange}\OperatorTok{();}

            \ControlFlowTok{if} \OperatorTok{(}\NormalTok{byteRange }\OperatorTok{==} \KeywordTok{null} \OperatorTok{||}\NormalTok{ byteRange}\OperatorTok{.}\FunctionTok{length} \OperatorTok{!=} \DecValTok{4}\OperatorTok{)} \OperatorTok{\{}
                \BuiltInTok{System}\OperatorTok{.}\FunctionTok{err}\OperatorTok{.}\FunctionTok{println}\OperatorTok{(}\StringTok{"FAIL: Missing or malformed /ByteRange"}\OperatorTok{);}
                \ControlFlowTok{continue}\OperatorTok{;}
            \OperatorTok{\}}

            \DataTypeTok{int}\NormalTok{ offset1 }\OperatorTok{=}\NormalTok{ byteRange}\OperatorTok{[}\DecValTok{0}\OperatorTok{];} \CommentTok{// Should be 0}
            \DataTypeTok{int}\NormalTok{ length1 }\OperatorTok{=}\NormalTok{ byteRange}\OperatorTok{[}\DecValTok{1}\OperatorTok{];}
            \DataTypeTok{int}\NormalTok{ offset2 }\OperatorTok{=}\NormalTok{ byteRange}\OperatorTok{[}\DecValTok{2}\OperatorTok{];}
            \DataTypeTok{int}\NormalTok{ length2 }\OperatorTok{=}\NormalTok{ byteRange}\OperatorTok{[}\DecValTok{3}\OperatorTok{];}
            \DataTypeTok{int}\NormalTok{ signedEnd }\OperatorTok{=}\NormalTok{ offset2 }\OperatorTok{+}\NormalTok{ length2}\OperatorTok{;}

            \CommentTok{// Check 1: First range must start at byte 0}
            \ControlFlowTok{if} \OperatorTok{(}\NormalTok{offset1 }\OperatorTok{!=} \DecValTok{0}\OperatorTok{)} \OperatorTok{\{}
                \BuiltInTok{System}\OperatorTok{.}\FunctionTok{err}\OperatorTok{.}\FunctionTok{println}\OperatorTok{(}
                        \StringTok{"FAIL: /ByteRange does not start at byte 0. "}
                        \OperatorTok{+} \StringTok{"Starts at: "} \OperatorTok{+}\NormalTok{ offset1}\OperatorTok{);} 
            \OperatorTok{\}}

            \CommentTok{// Check 2: The gap between ranges is the /Contents value}
            \DataTypeTok{int}\NormalTok{ gapStart }\OperatorTok{=}\NormalTok{ offset1 }\OperatorTok{+}\NormalTok{ length1}\OperatorTok{;}
            \DataTypeTok{int}\NormalTok{ gapEnd }\OperatorTok{=}\NormalTok{ offset2}\OperatorTok{;}
            \DataTypeTok{int}\NormalTok{ gapSize }\OperatorTok{=}\NormalTok{ gapEnd }\OperatorTok{{-}}\NormalTok{ gapStart}\OperatorTok{;}

            \ControlFlowTok{if} \OperatorTok{(}\NormalTok{gapSize }\OperatorTok{\textless{}=} \DecValTok{0}\OperatorTok{)} \OperatorTok{\{}
                \BuiltInTok{System}\OperatorTok{.}\FunctionTok{err}\OperatorTok{.}\FunctionTok{println}\OperatorTok{(}
                        \StringTok{"FAIL: /ByteRange gap is zero or negative. "}
                        \OperatorTok{+} \StringTok{"Possible SWA indicator"}\OperatorTok{);} 
            \OperatorTok{\}}

            \CommentTok{// Check 3: Verify /Contents hex string is at the gap position}
            \CommentTok{// The gap should contain exactly the hex{-}encoded CMS signature}
            \CommentTok{// surrounded by angle brackets: \textless{}hex...\textgreater{}}
            \ControlFlowTok{if} \OperatorTok{(}\NormalTok{pdfBytes}\OperatorTok{[}\NormalTok{gapStart}\OperatorTok{]} \OperatorTok{!=} \CharTok{\textquotesingle{}\textless{}\textquotesingle{}}\OperatorTok{)} \OperatorTok{\{}
                \BuiltInTok{System}\OperatorTok{.}\FunctionTok{err}\OperatorTok{.}\FunctionTok{println}\OperatorTok{(}
                        \StringTok{"FAIL: Expected \textquotesingle{}\textless{}\textquotesingle{} at gap start (offset "}
                        \OperatorTok{+}\NormalTok{ gapStart }\OperatorTok{+} \StringTok{"), found: 0x"}
                        \OperatorTok{+} \BuiltInTok{Integer}\OperatorTok{.}\FunctionTok{toHexString}\OperatorTok{(}
\NormalTok{                                pdfBytes}\OperatorTok{[}\NormalTok{gapStart}\OperatorTok{]} \OperatorTok{\&} \BaseNTok{0xFF}\OperatorTok{));} 
            \OperatorTok{\}}

            \CommentTok{// Check 4: The signed region should cover the entire file}
            \CommentTok{// (for a single{-}signature document with no subsequent updates)}
            \ControlFlowTok{if} \OperatorTok{(}\NormalTok{signedEnd }\OperatorTok{\textless{}}\NormalTok{ pdfBytes}\OperatorTok{.}\FunctionTok{length}\OperatorTok{)} \OperatorTok{\{}
                \BuiltInTok{System}\OperatorTok{.}\FunctionTok{err}\OperatorTok{.}\FunctionTok{println}\OperatorTok{(}
                        \StringTok{"WARNING: /ByteRange ends at byte "} \OperatorTok{+}\NormalTok{ signedEnd}
                        \OperatorTok{+} \StringTok{" but file is "} \OperatorTok{+}\NormalTok{ pdfBytes}\OperatorTok{.}\FunctionTok{length}
                        \OperatorTok{+} \StringTok{" bytes. Unsigned trailing content detected"}\OperatorTok{);}
            \OperatorTok{\}}

            \CommentTok{// Check 5: Cross{-}validate that the /ByteRange in the file}
            \CommentTok{// matches what we\textquotesingle{}d expect from the signed revision}
            \CommentTok{// This requires parsing the original revision\textquotesingle{}s xref}
            \BuiltInTok{System}\OperatorTok{.}\FunctionTok{out}\OperatorTok{.}\FunctionTok{println}\OperatorTok{(}
                    \StringTok{"ByteRange analysis: "}
                    \OperatorTok{+} \StringTok{"range1=["} \OperatorTok{+}\NormalTok{ offset1 }\OperatorTok{+} \StringTok{","} \OperatorTok{+} \OperatorTok{(}\NormalTok{offset1 }\OperatorTok{+}\NormalTok{ length1}\OperatorTok{)} \OperatorTok{+} \StringTok{") "}
                    \OperatorTok{+} \StringTok{"gap=["} \OperatorTok{+}\NormalTok{ gapStart }\OperatorTok{+} \StringTok{","} \OperatorTok{+}\NormalTok{ gapEnd }\OperatorTok{+} \StringTok{") "}
                    \OperatorTok{+} \StringTok{"range2=["} \OperatorTok{+}\NormalTok{ offset2 }\OperatorTok{+} \StringTok{","} \OperatorTok{+}\NormalTok{ signedEnd }\OperatorTok{+} \StringTok{") "}
                    \OperatorTok{+} \StringTok{"fileSize="} \OperatorTok{+}\NormalTok{ pdfBytes}\OperatorTok{.}\FunctionTok{length}\OperatorTok{);}
        \OperatorTok{\}}

\NormalTok{        doc}\OperatorTok{.}\FunctionTok{close}\OperatorTok{();}
    \OperatorTok{\}}
\OperatorTok{\}}
\end{Highlighting}
\end{Shaded}

\begin{itemize}
\item
  The first byte range must start at offset 0. If it doesn't, there are
  unsigned bytes at the beginning of the file --- a clear SWA indicator
\item
  The gap between the two ranges must be positive and should correspond
  to the \texttt{/Contents} hex string. A zero or negative gap indicates
  structural manipulation
\item
  The gap should contain the hex-encoded CMS value. Checking that it
  starts with \texttt{\textless{}} and ends with \texttt{\textgreater{}}
  is a basic sanity check
\end{itemize}

SWA defense requires more than just checking \texttt{/ByteRange}
arithmetic. A sophisticated attacker can construct a valid-looking byte
range that passes arithmetic checks but points to carefully crafted
content. The complete defense is: (a) verify \texttt{/ByteRange} covers
the file correctly, (b) verify the CMS signature cryptographically (USF
defense), and (c) verify that the \texttt{/ByteRange} values come from
the signed revision's cross-reference, not from an appended revision.

\subsection{Attack Class 4: Shadow Attacks
(2021)}\label{_attack_class_4_shadow_attacks_2021}

\subsubsection{A Fundamentally Different Threat
Model}\label{_a_fundamentally_different_threat_model}

The three attacks from CCS 2019 (USF, ISA, SWA) all shared a common
assumption: the attacker modifies a PDF \emph{after} it has been
legitimately signed. The signed document is genuine; the attacker's
manipulation comes afterward.

Shadow Attacks, published at NDSS 2021, inverted the model. The attacker
prepares the PDF \emph{before} it is signed, embedding hidden malicious
content that the signer cannot see during the signing process. The
signer reviews the document, sees nothing suspicious, and signs it. The
signed PDF is then sent to the recipient, who opens it --- and sees
different content from what the signer saw.

The signature is cryptographically valid. The signer genuinely signed
those exact bytes. The bytes just don't mean what the signer thought
they meant.

This is not a cryptographic failure. It is a semantic failure: the gap
between what a human \emph{sees} when they review a document and what
the \emph{bytes} actually encode. And it is far harder to defend against
than USF, ISA, or SWA.

\subsubsection{Shadow Attack Variant 1:
Hide}\label{_shadow_attack_variant_1_hide}

The attacker creates a PDF with two layers of content --- one visible,
one hidden. The hidden content is placed \emph{behind} the visible
content using PDF features like:

\begin{itemize}
\item
  \textbf{Overlapping objects}: A white rectangle drawn on top of text,
  hiding it. The text is in the content stream; the rectangle covers it.
  After signing, the attacker removes the rectangle via an incremental
  update, revealing the hidden text.
\item
  \textbf{Optional Content Groups (OCG)}: PDF's layer mechanism. Content
  can be tagged with \texttt{/OC} (Optional Content) markers and made
  visible or invisible via the \texttt{/OCProperties} dictionary. The
  attacker sets a layer to invisible before signing, then sets it to
  visible after signing.
\item
  \textbf{Referenced XObjects}: Content referenced as
  \texttt{Form\ XObject} but not actually rendered on any page. After
  signing, an incremental update adds the XObject to a page's content
  stream.
\end{itemize}

The key constraint: \emph{all of the hidden content must be present in
the file at the time of signing}, so that it falls within the signed
byte range. The attacker only changes the \emph{visibility} via an
incremental update after signing --- which some viewers allowed without
warning.

\subsubsection{Shadow Attack Variant 2:
Replace}\label{_shadow_attack_variant_2_replace}

The attacker creates a PDF where the visible content can be
\emph{replaced} by different visible content without modifying any
objects in the signed byte range. This exploits \textbf{interactive form
features}:

\begin{itemize}
\item
  A text field with a default value ("Contract for \$1,000") and a
  different actual value ("Contract for \$1,000,000") --- the default is
  displayed initially, but changing which value is rendered requires
  only an annotation flag change
\item
  Font references that can be swapped --- the same Unicode code points
  render as different glyphs depending on which font is active
\item
  Images referenced by name from a page's Resources dictionary ---
  changing which image a name resolves to requires only modifying the
  Resources dictionary
\end{itemize}

The Replace variant is particularly dangerous because it exploits
features that are \emph{designed} to change document appearance (form
filling, annotation updates) --- the viewer must distinguish between a
legitimate form fill and an attacker swapping content.

\subsubsection{Shadow Attack Variant 3:
Hide-and-Replace}\label{_shadow_attack_variant_3_hide_and_replace}

The most powerful variant combines both techniques. The attacker creates
a PDF that:

\begin{enumerate}
\def\labelenumi{\arabic{enumi}.}
\item
  Contains the content the signer will see (layer A)
\item
  Contains the content the recipient will see (layer B), hidden behind
  layer A
\item
  Is structured so that a single incremental update (changing a few
  bytes in an annotation dictionary or OCG configuration) switches from
  layer A to layer B
\end{enumerate}

After signing, the attacker appends the minimal incremental update
needed to switch layers. If the viewer allows annotation or OCG changes
after signing (as many did, because these are often considered
"permissible" changes), the attack succeeds.

\subsubsection{Demonstration: Shadow Attack (Hide Variant)
Preparation}\label{_demonstration_shadow_attack_hide_variant_preparation}

\begin{Shaded}
\begin{Highlighting}[]
\KeywordTok{import} \ImportTok{org}\OperatorTok{.}\ImportTok{apache}\OperatorTok{.}\ImportTok{pdfbox}\OperatorTok{.}\ImportTok{Loader}\OperatorTok{;}
\KeywordTok{import} \ImportTok{org}\OperatorTok{.}\ImportTok{apache}\OperatorTok{.}\ImportTok{pdfbox}\OperatorTok{.}\ImportTok{pdmodel}\OperatorTok{.}\ImportTok{PDDocument}\OperatorTok{;}
\KeywordTok{import} \ImportTok{org}\OperatorTok{.}\ImportTok{apache}\OperatorTok{.}\ImportTok{pdfbox}\OperatorTok{.}\ImportTok{pdmodel}\OperatorTok{.}\ImportTok{PDPage}\OperatorTok{;}
\KeywordTok{import} \ImportTok{org}\OperatorTok{.}\ImportTok{apache}\OperatorTok{.}\ImportTok{pdfbox}\OperatorTok{.}\ImportTok{pdmodel}\OperatorTok{.}\ImportTok{PDPageContentStream}\OperatorTok{;}
\KeywordTok{import} \ImportTok{org}\OperatorTok{.}\ImportTok{apache}\OperatorTok{.}\ImportTok{pdfbox}\OperatorTok{.}\ImportTok{pdmodel}\OperatorTok{.}\ImportTok{common}\OperatorTok{.}\ImportTok{PDRectangle}\OperatorTok{;}
\KeywordTok{import} \ImportTok{org}\OperatorTok{.}\ImportTok{apache}\OperatorTok{.}\ImportTok{pdfbox}\OperatorTok{.}\ImportTok{pdmodel}\OperatorTok{.}\ImportTok{font}\OperatorTok{.}\ImportTok{PDType1Font}\OperatorTok{;}
\KeywordTok{import} \ImportTok{org}\OperatorTok{.}\ImportTok{apache}\OperatorTok{.}\ImportTok{pdfbox}\OperatorTok{.}\ImportTok{pdmodel}\OperatorTok{.}\ImportTok{font}\OperatorTok{.}\ImportTok{Standard14Fonts}\OperatorTok{;}

\KeywordTok{import} \ImportTok{java}\OperatorTok{.}\ImportTok{io}\OperatorTok{.}\ImportTok{ByteArrayOutputStream}\OperatorTok{;}
\KeywordTok{import} \ImportTok{java}\OperatorTok{.}\ImportTok{io}\OperatorTok{.}\ImportTok{File}\OperatorTok{;}
\KeywordTok{import} \ImportTok{java}\OperatorTok{.}\ImportTok{nio}\OperatorTok{.}\ImportTok{file}\OperatorTok{.}\ImportTok{Files}\OperatorTok{;}

\CommentTok{/**}
 \CommentTok{*}\NormalTok{ Demonstrates the preparation phase of a Shadow Attack }\CommentTok{(}\NormalTok{Hide variant}\CommentTok{).}
 \CommentTok{*}\NormalTok{ The attacker creates a PDF with hidden text behind a white overlay}\CommentTok{.}
 \CommentTok{*}
 \CommentTok{*}\NormalTok{ BEFORE signing}\CommentTok{:}\NormalTok{ Signer sees }\CommentTok{"}\NormalTok{Payment}\CommentTok{:} \CommentTok{$1,000"}
 \CommentTok{*}\NormalTok{ AFTER signing }\CommentTok{+}\NormalTok{ reveal}\CommentTok{:}\NormalTok{ Recipient sees }\CommentTok{"}\NormalTok{Payment}\CommentTok{:} \CommentTok{$1,000,000"}
 \CommentTok{*}
 \CommentTok{*}\NormalTok{ The hidden text }\CommentTok{("}\NormalTok{Payment}\CommentTok{:} \CommentTok{$1,000,000")}\NormalTok{ is present in the signed bytes}\CommentTok{.}
 \CommentTok{*}\NormalTok{ The overlay }\CommentTok{(}\NormalTok{white rectangle}\CommentTok{)}\NormalTok{ hides it}\CommentTok{.}\NormalTok{ An incremental update that removes}
 \CommentTok{*}\NormalTok{ the overlay reveals the hidden text }\CommentTok{—}\NormalTok{ and the signature remains valid}\CommentTok{.}
 \CommentTok{*/}
\KeywordTok{public} \KeywordTok{class}\NormalTok{ ShadowAttackHideDemo }\OperatorTok{\{}

    \KeywordTok{public} \DataTypeTok{static} \DataTypeTok{void} \FunctionTok{main}\OperatorTok{(}\BuiltInTok{String}\OperatorTok{[]}\NormalTok{ args}\OperatorTok{)} \KeywordTok{throws} \BuiltInTok{Exception} \OperatorTok{\{}
\NormalTok{        PDDocument doc }\OperatorTok{=} \KeywordTok{new} \FunctionTok{PDDocument}\OperatorTok{();}
\NormalTok{        PDPage page }\OperatorTok{=} \KeywordTok{new} \FunctionTok{PDPage}\OperatorTok{(}\NormalTok{PDRectangle}\OperatorTok{.}\FunctionTok{A4}\OperatorTok{);}
\NormalTok{        doc}\OperatorTok{.}\FunctionTok{addPage}\OperatorTok{(}\NormalTok{page}\OperatorTok{);}

\NormalTok{        PDType1Font font }\OperatorTok{=} \KeywordTok{new} \FunctionTok{PDType1Font}\OperatorTok{(}
\NormalTok{                Standard14Fonts}\OperatorTok{.}\FunctionTok{FontName}\OperatorTok{.}\FunctionTok{HELVETICA}\OperatorTok{);}
\NormalTok{        PDType1Font fontBold }\OperatorTok{=} \KeywordTok{new} \FunctionTok{PDType1Font}\OperatorTok{(}
\NormalTok{                Standard14Fonts}\OperatorTok{.}\FunctionTok{FontName}\OperatorTok{.}\FunctionTok{HELVETICA\_BOLD}\OperatorTok{);}

\NormalTok{        PDPageContentStream cs }\OperatorTok{=} \KeywordTok{new} \FunctionTok{PDPageContentStream}\OperatorTok{(}\NormalTok{doc}\OperatorTok{,}\NormalTok{ page}\OperatorTok{);}

        \CommentTok{// Layer 1: The hidden content (written FIRST, so it\textquotesingle{}s behind)}
        \CommentTok{// This text will be present in the signed bytes}
\NormalTok{        cs}\OperatorTok{.}\FunctionTok{beginText}\OperatorTok{();}
\NormalTok{        cs}\OperatorTok{.}\FunctionTok{setFont}\OperatorTok{(}\NormalTok{fontBold}\OperatorTok{,} \DecValTok{16}\OperatorTok{);}
\NormalTok{        cs}\OperatorTok{.}\FunctionTok{newLineAtOffset}\OperatorTok{(}\DecValTok{50}\OperatorTok{,} \DecValTok{700}\OperatorTok{);}
\NormalTok{        cs}\OperatorTok{.}\FunctionTok{showText}\OperatorTok{(}\StringTok{"Payment: $1,000,000"}\OperatorTok{);} 
\NormalTok{        cs}\OperatorTok{.}\FunctionTok{endText}\OperatorTok{();}

        \CommentTok{// Layer 2: The overlay that hides it}
        \CommentTok{// A white rectangle drawn ON TOP of the hidden text}
\NormalTok{        cs}\OperatorTok{.}\FunctionTok{setNonStrokingColor}\OperatorTok{(}\FloatTok{1.0f}\OperatorTok{,} \FloatTok{1.0f}\OperatorTok{,} \FloatTok{1.0f}\OperatorTok{);} \CommentTok{// White}
\NormalTok{        cs}\OperatorTok{.}\FunctionTok{addRect}\OperatorTok{(}\DecValTok{45}\OperatorTok{,} \DecValTok{695}\OperatorTok{,} \DecValTok{250}\OperatorTok{,} \DecValTok{25}\OperatorTok{);} \CommentTok{// Covers the text above}
\NormalTok{        cs}\OperatorTok{.}\FunctionTok{fill}\OperatorTok{();} 

        \CommentTok{// Layer 3: The visible content (what the signer sees)}
\NormalTok{        cs}\OperatorTok{.}\FunctionTok{setNonStrokingColor}\OperatorTok{(}\FloatTok{0.0f}\OperatorTok{,} \FloatTok{0.0f}\OperatorTok{,} \FloatTok{0.0f}\OperatorTok{);} \CommentTok{// Black}
\NormalTok{        cs}\OperatorTok{.}\FunctionTok{beginText}\OperatorTok{();}
\NormalTok{        cs}\OperatorTok{.}\FunctionTok{setFont}\OperatorTok{(}\NormalTok{fontBold}\OperatorTok{,} \DecValTok{16}\OperatorTok{);}
\NormalTok{        cs}\OperatorTok{.}\FunctionTok{newLineAtOffset}\OperatorTok{(}\DecValTok{50}\OperatorTok{,} \DecValTok{700}\OperatorTok{);}
\NormalTok{        cs}\OperatorTok{.}\FunctionTok{showText}\OperatorTok{(}\StringTok{"Payment: $1,000"}\OperatorTok{);} 
\NormalTok{        cs}\OperatorTok{.}\FunctionTok{endText}\OperatorTok{();}

        \CommentTok{// Legitimate{-}looking document body}
\NormalTok{        cs}\OperatorTok{.}\FunctionTok{beginText}\OperatorTok{();}
\NormalTok{        cs}\OperatorTok{.}\FunctionTok{setFont}\OperatorTok{(}\NormalTok{font}\OperatorTok{,} \DecValTok{12}\OperatorTok{);}
\NormalTok{        cs}\OperatorTok{.}\FunctionTok{newLineAtOffset}\OperatorTok{(}\DecValTok{50}\OperatorTok{,} \DecValTok{650}\OperatorTok{);}
\NormalTok{        cs}\OperatorTok{.}\FunctionTok{showText}\OperatorTok{(}\StringTok{"This agreement confirms the payment amount above."}\OperatorTok{);}
\NormalTok{        cs}\OperatorTok{.}\FunctionTok{endText}\OperatorTok{();}

\NormalTok{        cs}\OperatorTok{.}\FunctionTok{close}\OperatorTok{();}

        \CommentTok{// Save — this PDF is now ready to be sent to the signer}
        \BuiltInTok{ByteArrayOutputStream}\NormalTok{ baos }\OperatorTok{=} \KeywordTok{new} \BuiltInTok{ByteArrayOutputStream}\OperatorTok{();}
\NormalTok{        doc}\OperatorTok{.}\FunctionTok{save}\OperatorTok{(}\NormalTok{baos}\OperatorTok{);}
\NormalTok{        doc}\OperatorTok{.}\FunctionTok{close}\OperatorTok{();}

\NormalTok{        Files}\OperatorTok{.}\FunctionTok{write}\OperatorTok{(}
                \KeywordTok{new} \BuiltInTok{File}\OperatorTok{(}\StringTok{"shadow{-}prepared.pdf"}\OperatorTok{).}\FunctionTok{toPath}\OperatorTok{(),}\NormalTok{ baos}\OperatorTok{.}\FunctionTok{toByteArray}\OperatorTok{());}

        \CommentTok{// The signer opens this PDF, sees "Payment: $1,000", and signs it.}
        \CommentTok{// The signed byte range includes ALL content: both the hidden}
        \CommentTok{// "$1,000,000" text and the white rectangle that covers it.}
        \CommentTok{//}
        \CommentTok{// After signing, the attacker creates an incremental update that}
        \CommentTok{// removes or makes transparent the white rectangle, revealing}
        \CommentTok{// "Payment: $1,000,000" — and the signature validates.}
    \OperatorTok{\}}
\OperatorTok{\}}
\end{Highlighting}
\end{Shaded}

\begin{itemize}
\item
  The hidden text is written to the content stream first. In PDF
  rendering order, later operators paint \emph{over} earlier ones on the
  same area
\item
  A white filled rectangle covers the hidden text completely. The signer
  sees only white space at this position
\item
  The visible text is drawn on top of the white rectangle. This is what
  the signer reads when reviewing the document before signing
\end{itemize}

Shadow Attacks are fundamentally harder to prevent than USF, ISA, or SWA
because the malicious content is \emph{part of the signed data}. The
signature is cryptographically valid over the exact bytes the signer
intended to sign. The defense must go beyond byte-level verification
into \emph{semantic} analysis: does the document contain hidden content?
Are there overlapping objects? Are there unused Form XObjects? Are there
Optional Content Groups with suspicious visibility settings? This is an
open research problem with no complete solution.

\subsubsection{The Scope of Shadow Attack
Vulnerability}\label{_the_scope_of_shadow_attack_vulnerability}

The Bochum team tested 29 PDF viewers against the Shadow Attack
variants. The results were comprehensive:

\begin{longtable}[]{@{}
  >{\raggedright\arraybackslash}p{(\linewidth - 6\tabcolsep) * \real{0.5000}}
  >{\raggedright\arraybackslash}p{(\linewidth - 6\tabcolsep) * \real{0.1667}}
  >{\raggedright\arraybackslash}p{(\linewidth - 6\tabcolsep) * \real{0.1667}}
  >{\raggedright\arraybackslash}p{(\linewidth - 6\tabcolsep) * \real{0.1667}}@{}}
\caption{Shadow Attack test results across PDF viewers (NDSS
2021)}\tabularnewline
\toprule\noalign{}
\begin{minipage}[b]{\linewidth}\raggedright
Viewer
\end{minipage} & \begin{minipage}[b]{\linewidth}\raggedright
Hide
\end{minipage} & \begin{minipage}[b]{\linewidth}\raggedright
Replace
\end{minipage} & \begin{minipage}[b]{\linewidth}\raggedright
Hide-and-Replace
\end{minipage} \\
\midrule\noalign{}
\endfirsthead
\toprule\noalign{}
\begin{minipage}[b]{\linewidth}\raggedright
Viewer
\end{minipage} & \begin{minipage}[b]{\linewidth}\raggedright
Hide
\end{minipage} & \begin{minipage}[b]{\linewidth}\raggedright
Replace
\end{minipage} & \begin{minipage}[b]{\linewidth}\raggedright
Hide-and-Replace
\end{minipage} \\
\midrule\noalign{}
\endhead
\bottomrule\noalign{}
\endlastfoot
Adobe Acrobat Reader DC & Vulnerable & Vulnerable & Vulnerable \\
Foxit Reader & Vulnerable & Vulnerable & Vulnerable \\
LibreOffice Draw & Vulnerable & Partially vulnerable & Vulnerable \\
Nitro Pro & Vulnerable & Vulnerable & Vulnerable \\
PDF-XChange & Vulnerable & Not vulnerable & Vulnerable \\
Okular & Vulnerable & Vulnerable & Vulnerable \\
Evince & Vulnerable & Vulnerable & Vulnerable \\
Chrome PDF Viewer & No signature validation & --- & --- \\
macOS Preview & Vulnerable & Vulnerable & Vulnerable \\
\end{longtable}

This table is a representative subset. The full results covering 29
viewers are available in the NDSS 2021 paper and at pdf-insecurity.org.
Chrome's built-in PDF viewer does not validate signatures at all (it
displays them but does not indicate validity), so Shadow Attacks are not
applicable in the traditional sense --- but the user also receives no
warning that the document has been signed and potentially tampered with.

\subsection{Pre-Bochum Vulnerabilities: Earlier PDF Signature
Issues}\label{_pre_bochum_vulnerabilities_earlier_pdf_signature_issues}

The Bochum research was the most systematic analysis of PDF signature
security, but it was not the first time vulnerabilities were discovered.
Several earlier issues laid the groundwork --- and reveal that the
\emph{category} of bug (failing to verify what you claim to verify) has
been a persistent pattern in PDF signature implementations.

\subsubsection{CVE-2018-16042 --- Adobe Acrobat Signature Bypass
(2018)}\label{_cve_2018_16042_adobe_acrobat_signature_bypass_2018}

In 2018, Adobe patched \textbf{CVE-2018-16042} in Acrobat DC and Acrobat
Reader DC (versions before 2019.008.20081). The vulnerability allowed a
crafted PDF to display a valid signature indicator even when the
document had been modified.

The technical details were not fully disclosed (Adobe's advisory was
characteristically brief), but the vulnerability class --- signature
validation bypass via document manipulation --- is the same category as
the ISA attacks that Bochum would systematize a year later. This CVE was
one of several signals that the PDF signature validation code in major
viewers had not been subject to rigorous adversarial testing.

\textbf{Affected versions:} Acrobat DC/Reader DC before 2019.008.20081\\
\textbf{Impact:} Signature validation bypass --- modified documents
displayed as validly signed\\
\textbf{Fix:} Updated signature validation logic in the January 2019
security update

\subsubsection{CVE-2018-4993 --- Adobe Acrobat NTLM Credential Theft
(2018)}\label{_cve_2018_4993_adobe_acrobat_ntlm_credential_theft_2018}

While not strictly a signature bypass, \textbf{CVE-2018-4993}
demonstrated a related risk: a crafted PDF could trigger NTLM
authentication requests to an attacker-controlled server when opened in
Adobe Acrobat. Combined with a forged signature (via any of the bypass
techniques), this creates a compound attack: the victim opens a "validly
signed" PDF, and their Windows credentials are silently exfiltrated.

\textbf{Affected versions:} Acrobat DC/Reader DC before 2018.011.20040\\
\textbf{Impact:} NTLM hash theft via SMB relay when opening a crafted
PDF\\
\textbf{Fix:} Restricted external resource loading in the May 2018
security update

\subsubsection{SHAttered and the SHA-1 Deprecation
(2017)}\label{_shattered_and_the_sha_1_deprecation_2017}

In February 2017, researchers from Google and CWI Amsterdam published
\textbf{SHAttered} --- the first practical collision attack against
SHA-1. They produced two different PDF documents with the same SHA-1
hash, proving that SHA-1 could no longer be relied upon for collision
resistance.

For PDF signatures, this had immediate implications. The
\texttt{adbe.pkcs7.sha1} SubFilter --- which embeds a SHA-1 hash of the
entire document as the CMS content data --- was now theoretically
vulnerable to a collision attack: an attacker could prepare two
documents with the same SHA-1 hash, have the victim sign one, and then
substitute the other.

The practical impact was limited (the collision required approximately
2\textsuperscript{63} computations and the two colliding documents had
to be carefully constructed), but the theoretical break was sufficient
to trigger industry-wide deprecation:

\begin{itemize}
\item
  ISO 32000-2:2020 deprecated \texttt{adbe.pkcs7.sha1} entirely
\item
  Adobe Acrobat began showing warnings for SHA-1 signed documents
\item
  The CA/Browser Forum required SHA-256 as the minimum for all new
  certificates
\item
  PAdES profiles already mandated SHA-256 as the minimum digest
  algorithm
\end{itemize}

\begin{longtable}[]{@{}
  >{\raggedright\arraybackslash}p{(\linewidth - 2\tabcolsep) * \real{0.2000}}
  >{\raggedright\arraybackslash}p{(\linewidth - 2\tabcolsep) * \real{0.8000}}@{}}
\caption{SHA-1 timeline for PDF signatures}\tabularnewline
\toprule\noalign{}
\begin{minipage}[b]{\linewidth}\raggedright
Year
\end{minipage} & \begin{minipage}[b]{\linewidth}\raggedright
Event
\end{minipage} \\
\midrule\noalign{}
\endfirsthead
\toprule\noalign{}
\begin{minipage}[b]{\linewidth}\raggedright
Year
\end{minipage} & \begin{minipage}[b]{\linewidth}\raggedright
Event
\end{minipage} \\
\midrule\noalign{}
\endhead
\bottomrule\noalign{}
\endlastfoot
2005 & Theoretical weakness in SHA-1 published (Wang et al.) \\
2011 & NIST deprecated SHA-1 for digital signatures in federal use \\
2013 & Microsoft announced SHA-1 certificate sunset for 2016 \\
2017 & SHAttered: first practical SHA-1 collision (Google/CWI). Two PDFs
with identical SHA-1 hashes publicly demonstrated \\
2017 & ISO 32000-2 deprecated \texttt{adbe.pkcs7.sha1} SubFilter \\
2020 & SHA-1 collisions reduced to 2\textsuperscript{63} operations
(Leurent and Peyrin, "SHA-1 is a Shambles"), making chosen-prefix
collisions feasible \\
\end{longtable}

If your code still generates signatures using SHA-1 (either via
\texttt{adbe.pkcs7.sha1} or using SHA-1 as the \texttt{digestAlgorithm}
in the CMS \texttt{SignerInfo}), stop immediately. SHA-1 is broken. Use
SHA-256 or SHA-384. There is no legitimate reason to use SHA-1 for any
new signature created after 2017.

\subsubsection{The Comodo Registration Authority Breach
(2011)}\label{_the_comodo_registration_authority_breach_2011}

In March 2011 --- the same month as DigiNotar, though discovered
slightly earlier --- an attacker compromised a \textbf{registration
authority (RA)} affiliated with Comodo, one of the world's largest
certificate authorities at the time.

The attacker used the compromised RA account to issue nine fraudulent
SSL certificates for high-profile domains: \texttt{login.live.com},
\texttt{mail.google.com}, \texttt{www.google.com},
\texttt{login.yahoo.com}, \texttt{login.skype.com},
\texttt{addons.mozilla.org}, and three others. The certificates were
issued by Comodo's legitimate intermediate CA, meaning they chained
correctly to a trusted root.

The breach was discovered quickly because the attacker tested one of the
certificates against a live server, and certificate pinning in Google
Chrome detected the mismatch. Without Chrome's certificate pinning, the
fraudulent certificates could have been used for undetectable
man-in-the-middle attacks.

\textbf{For PDF signatures:} The Comodo breach demonstrated that even
without compromising a root CA's private key, an attacker could obtain
legitimate-looking certificates through a \emph{registration authority}
compromise --- a weaker link in the chain. A document signed with one of
these fraudulent certificates would have validated correctly in any PDF
viewer that trusted Comodo's root. The signer would appear to be
"Google" or "Microsoft" or "Yahoo" --- entities that obviously did not
sign the document.

The incident led to:

\begin{itemize}
\item
  Mandatory two-factor authentication for all RA accounts at major CAs
\item
  Domain control validation (DCV) requirements tightened by the
  CA/Browser Forum
\item
  Increased scrutiny of the RA model, eventually leading several CAs to
  eliminate external RAs entirely
\end{itemize}

\subsubsection{The Symantec Distrust
(2017--2018)}\label{_the_symantec_distrust_20172018}

Between 2015 and 2017, Google's Chrome security team and the CA/Browser
Forum discovered that \textbf{Symantec} (which at that time owned the
Thawte, GeoTrust, VeriSign, and RapidSSL certificate brands) had
mis-issued thousands of certificates --- including certificates for
domains whose owners had not authorized them.

The problems were systemic:

\begin{itemize}
\item
  Symantec had delegated certificate issuance to third-party partners
  (RAs) without adequate oversight
\item
  Test certificates for real domains were issued to unauthorized parties
\item
  Audit logs were incomplete or missing
\item
  Symantec initially denied the scope of the problem, then acknowledged
  it in stages
\end{itemize}

Google announced a graduated distrust of all Symantec-issued
certificates:

\begin{itemize}
\item
  Chrome 66 (April 2018): distrusted certificates issued before June 1,
  2016
\item
  Chrome 70 (October 2018): distrusted all remaining Symantec
  certificates
\end{itemize}

Mozilla Firefox and Apple Safari followed with their own distrust
timelines. Symantec sold its CA business to DigiCert, which re-issued
certificates for affected customers under DigiCert's infrastructure.

\textbf{For PDF signatures:} Symantec was a major CA in the AATL (Adobe
Approved Trust List). Documents signed with Symantec-issued certificates
before the distrust cutoff dates would validate correctly when signed,
but might show warnings or failures in PDF viewers that updated their
trust stores to reflect the browser distrust. This created a
time-dependent validation problem: a PDF signed with a Symantec
certificate in 2015 was "valid" when signed, "suspect" in 2018, and ---
depending on the viewer's trust configuration --- might show as
"invalid" or "valid with warnings" today.

This is exactly the problem that PAdES-B-LTA (long-term validation with
archival timestamps) was designed to solve: by embedding the full
certificate chain and revocation data at the time of signing, and then
protecting it with periodic archival timestamps, the signature's
validity at the time of signing can be verified even after the issuing
CA is no longer trusted.

\subsection{Exploits in the Wild}\label{_exploits_in_the_wild}

A natural question: were any of these attacks exploited before or after
the Bochum disclosures?

The honest answer is: \textbf{we don't know for certain}.

PDF signatures are used in financial contracts, government documents,
legal filings, tax returns, insurance claims, and medical records. If an
attacker successfully forged a signature on a financial contract, the
victim might not discover it until a dispute arises --- and even then,
the forensic analysis might not identify the specific attack technique.

What we do know:

\subsubsection{Indicators of Real-World
Exploitation}\label{_indicators_of_real_world_exploitation}

\begin{enumerate}
\def\labelenumi{\arabic{enumi}.}
\item
  \textbf{The "1 Trillion Dollar Refund" title was not hyperbole.} The
  paper's title referenced the fact that tax authorities in multiple
  countries accept signed PDFs as official submissions. The German tax
  system (ELSTER), for example, processes millions of electronically
  signed documents per year. If USF or ISA could forge signatures on tax
  submissions, the financial exposure was enormous. The Bochum team
  coordinated with CERT-Bund (Germany's national CERT) before
  publication.
\item
  \textbf{Adobe's response was unusually fast.} Adobe issued
  CVE-2020-9592 and CVE-2020-9596 and released patches within weeks of
  the responsible disclosure. For a company that typically operates on
  monthly security update cycles, the accelerated timeline suggests
  Adobe assessed the risk as immediate and severe --- consistent with
  either active exploitation or intelligence indicating exploitation was
  imminent.
\item
  \textbf{The EU Digital Signature Service (DSS) was vulnerable.} The
  EU's own reference implementation for signature validation --- the
  tool member states use to validate eIDAS-compliant signatures --- was
  vulnerable to ISA attacks. If the EU's own validator could be fooled,
  signatures submitted to government services across 27 member states
  could potentially have been forged without detection.
\item
  \textbf{Commercial signing services had no detection capability.}
  Before the Bochum research, no commercial PDF signing or validation
  service tested for these attack classes. There were no automated tools
  to detect ISA, SWA, or Shadow Attack patterns in submitted PDFs. Any
  exploitation that occurred before 2019 would have been invisible to
  the existing security infrastructure.
\end{enumerate}

\subsubsection{Post-Disclosure Exploitation
Risk}\label{_post_disclosure_exploitation_risk}

After the papers were published and patches released, the risk shifted:

\begin{itemize}
\item
  \textbf{Unpatched viewers remain vulnerable.} Not all users update
  their PDF viewers. Enterprise environments, government agencies with
  slow update cycles, and users in regions with limited internet access
  may still be running vulnerable versions of Acrobat, Foxit, or other
  viewers.
\item
  \textbf{The attack techniques are public.} The Bochum team published
  detailed descriptions and proof-of-concept PDFs. Any attacker with
  basic PDF knowledge can reproduce the attacks.
\item
  \textbf{Shadow Attacks remain partially unresolved.} While most
  viewers now detect the Hide and Replace variants, the fundamental
  problem --- that PDF semantics allow content to be visually hidden ---
  is inherent to the format. New variants of Shadow Attacks are likely
  still discoverable.
\end{itemize}

\subsection{The Patch Response: What
Changed}\label{_the_patch_response_what_changed}

\subsubsection{Adobe's Response}\label{_adobes_response}

Adobe's patches for CVE-2020-9592 and CVE-2020-9596 (released in the May
2020 Acrobat security update, version 2020.006.20042) implemented:

\begin{enumerate}
\def\labelenumi{\arabic{enumi}.}
\item
  \textbf{Cryptographic signature verification} --- closing the USF
  vulnerability by ensuring the CMS signature value is always verified
  against the document hash
\item
  \textbf{Byte range coverage validation} --- checking that the
  \texttt{/ByteRange} covers the entire file from offset 0 to the file's
  end (minus the \texttt{/Contents} gap), closing the SWA vulnerability
\item
  \textbf{Incremental update analysis} --- warning users when an
  incremental update modifies content beyond what the signature's
  permission level allows, closing the ISA vulnerability
\item
  \textbf{Shadow content detection} --- partial detection of Hidden
  content patterns (overlapping objects, OCGs with suspicious visibility
  settings)
\end{enumerate}

Acrobat's signature panel now displays more granular information:

\begin{itemize}
\item
  "Signed and all signatures are valid" --- no modifications detected
\item
  "Signed but the document has been modified since signing" ---
  incremental updates detected, with details about what changed
\item
  "At least one signature has problems" --- cryptographic or structural
  verification failed
\item
  "The document has been altered or corrupted since the signature was
  applied" --- clear indication of tampering
\end{itemize}

\subsubsection{Foxit's Response}\label{_foxits_response}

Foxit Reader patched its signature validation in version 10.1 (released
October 2020), implementing similar checks: CMS verification, byte range
validation, and incremental update analysis. Foxit was notably slower
than Adobe in responding --- roughly five months after Adobe's patch.

\subsubsection{LibreOffice's Response}\label{_libreoffices_response}

LibreOffice (which uses the NSS library for cryptographic operations)
patched its PDF signature validation in version 7.0. The fix was more
complex because LibreOffice's PDF import/export pipeline differs
significantly from dedicated PDF viewers --- it converts PDF to an
internal format for editing, which complicates byte-level signature
verification.

\subsubsection{The Broader Industry}\label{_the_broader_industry}

By mid-2021, most actively maintained PDF viewers had released patches
addressing the three original attack classes (USF, ISA, SWA). Shadow
Attack mitigations were slower to arrive and remain incomplete in some
viewers as of 2026.

The Bochum research also influenced the standards bodies:

\begin{itemize}
\item
  \textbf{ISO TC 171} (the committee responsible for ISO 32000-2:2020)
  incorporated the attack findings into the errata and clarifications
  for the second edition of PDF 2.0
\item
  \textbf{ETSI} updated the PAdES standard (EN 319 142-1 V1.2.1,
  published 2022) to include explicit requirements for byte range
  coverage verification and incremental update validation
\item
  \textbf{The EU Digital Signature Service (DSS)} --- the open-source
  reference validator maintained by the European Commission --- was
  patched and now includes specific test cases for each attack class
\end{itemize}

\subsection{Building a Robust Validator: Lessons from the
Attacks}\label{_building_a_robust_validator_lessons_from_the_attacks}

Every attack in this chapter exploits a specific gap between what a
viewer \emph{should} check and what it \emph{actually} checks. Here is
the complete checklist, informed by every vulnerability discussed above.

\begin{longtable}[]{@{}
  >{\raggedright\arraybackslash}p{(\linewidth - 4\tabcolsep) * \real{0.1667}}
  >{\raggedright\arraybackslash}p{(\linewidth - 4\tabcolsep) * \real{0.5000}}
  >{\raggedright\arraybackslash}p{(\linewidth - 4\tabcolsep) * \real{0.3333}}@{}}
\caption{PDF Signature Validation Checklist
(post-Bochum)}\tabularnewline
\toprule\noalign{}
\begin{minipage}[b]{\linewidth}\raggedright
\#
\end{minipage} & \begin{minipage}[b]{\linewidth}\raggedright
Check
\end{minipage} & \begin{minipage}[b]{\linewidth}\raggedright
Defends Against
\end{minipage} \\
\midrule\noalign{}
\endfirsthead
\toprule\noalign{}
\begin{minipage}[b]{\linewidth}\raggedright
\#
\end{minipage} & \begin{minipage}[b]{\linewidth}\raggedright
Check
\end{minipage} & \begin{minipage}[b]{\linewidth}\raggedright
Defends Against
\end{minipage} \\
\midrule\noalign{}
\endhead
\bottomrule\noalign{}
\endlastfoot
1 & Parse the CMS \texttt{SignedData} structure and verify it is
well-formed (correct OIDs, consistent version numbers, all required
fields present) & USF Variant 2 \\
2 & Verify the \texttt{SignerInfo.signature} value cryptographically
against the signer's public key and the computed digest & USF Variant
1 \\
3 & Verify the message digest in the signed attributes matches the hash
of the byte range content & USF Variant 3, SWA \\
4 & Verify \texttt{/ByteRange} starts at offset 0 & SWA \\
5 & Verify \texttt{/ByteRange} gap contains the \texttt{/Contents} hex
string & SWA \\
6 & Verify \texttt{/ByteRange} covers the entire file (for the signed
revision) --- no unaccounted bytes before, between, or after the ranges
& SWA, ISA \\
7 & If the file extends beyond the signed byte range, parse the
incremental update and determine what changed & ISA \\
8 & Verify that any changes in the incremental update are permissible
under the signature's DocMDP permission level (P=1: no changes; P=2:
form fills + signatures; P=3: form fills + signatures + annotations) &
ISA Variant 3 \\
9 & Verify the certificate chain from the signer's certificate to a
trusted root, checking revocation (OCSP/CRL) and validity dates &
Certificate forgery, CA breaches \\
10 & If a timestamp is present, verify the TSA certificate chain
independently and check that the timestamp covers the correct data &
Backdated signatures \\
11 & Check for hidden content patterns: overlapping objects, Optional
Content Groups with invisible layers, unreferenced Form XObjects, form
fields with discrepant default/actual values & Shadow Attacks (Hide,
Replace) \\
12 & Ensure the digest algorithm is SHA-256 or stronger --- reject SHA-1
& SHAttered, hash collision attacks \\
\end{longtable}

Checks 1--6 are non-negotiable --- any validator that skips them is
vulnerable to at least one of USF, ISA, or SWA. Checks 7--10 are
necessary for real-world security but require more complex
implementation. Checks 11--12 are defense-in-depth against attacks that
remain partially unresolved.

No checklist is complete. The Bochum team found these attack classes by
thinking adversarially about the gap between what the spec
\emph{requires} and what implementations \emph{do}. There will be more
attacks. The defense is not just implementing these twelve checks --- it
is building a culture of adversarial testing into your validation code.

\subsection{A Note on Responsible
Disclosure}\label{_a_note_on_responsible_disclosure}

The Bochum team followed responsible disclosure for both papers. They
notified affected vendors months before publication and provided
proof-of-concept PDFs for testing.

For the CCS 2019 paper, the disclosure timeline was approximately:

\begin{itemize}
\item
  \textbf{March--April 2019}: Initial vendor notifications
\item
  \textbf{May--August 2019}: Back-and-forth with vendors; some requested
  extended timelines; some did not respond
\item
  \textbf{November 2019}: Paper presented at CCS; public disclosure
\item
  \textbf{May 2020}: Adobe patches released (CVE-2020-9592,
  CVE-2020-9596)
\item
  \textbf{October 2020}: Foxit patches released
\end{itemize}

The six-month gap between public disclosure and Adobe's patch --- and
the eleven-month gap before Foxit's patch --- is notable. The Bochum
team publicly criticized the slow response, particularly from viewers
that took more than six months to patch after detailed proof-of-concept
material was shared privately.

For the NDSS 2021 paper (Shadow Attacks), the disclosure process was
similar but vendors responded faster, having already been through the
CCS 2019 cycle.

The research materials --- including proof-of-concept PDFs, the
vulnerability matrix, and detailed attack descriptions --- remain
available at \textbf{pdf-insecurity.org}. The Hackmanit GmbH \textbf{PDF
Detector} tool (available as a web service and standalone library) can
check PDFs for known attack patterns.

\subsection{Summary}\label{_summary_10}

PDF signature security before 2019 was, to put it bluntly, an illusion.
The cryptographic primitives were sound --- RSA, ECDSA, SHA-256, CMS are
all well-designed and well-analyzed. The failure was entirely in how PDF
viewers \emph{used} those primitives. They verified certificates
meticulously and signatures not at all.

The Bochum research (CCS 2019, NDSS 2021) exposed four attack classes:

\begin{itemize}
\item
  \textbf{USF}: The viewer doesn't check the cryptographic signature
  value --- forge anything
\item
  \textbf{ISA}: The viewer doesn't check for modifications after signing
  --- append anything
\item
  \textbf{SWA}: The viewer doesn't validate the byte range --- redirect
  it to anything
\item
  \textbf{Shadow Attacks}: The viewer doesn't analyze what the signed
  content \emph{means} --- hide anything
\end{itemize}

Every major PDF viewer was vulnerable. The patches took months to years.
Some Shadow Attack variants remain partially unresolved.

Earlier events --- the DigiNotar breach (2011), the Comodo RA compromise
(2011), the Symantec distrust (2017--2018), and the SHAttered SHA-1
collision (2017) --- had already weakened the trust infrastructure
underneath PDF signatures. The Bochum attacks showed that even if you
trusted the infrastructure, the validators themselves couldn't be
trusted.

The lesson for implementers is clear: \textbf{verify everything, trust
nothing, and test adversarially}. The validation checklist in this
chapter is the minimum. The next attack class is already being
researched.

\chapter{Certificate Authority
Breaches}\label{_certificate_authority_breaches}

Every PDF digital signature carries an implicit assertion: "I trust the
certificate authority that issued this certificate." That assertion is
not a mathematical fact --- it is an organizational, operational, and
political judgment. And like all judgments, it can turn out to be
catastrophically wrong.

This chapter chronicles the most consequential certificate authority
failures of the past two decades. Not because they are historical
curiosities, but because each one contains lessons that directly affect
how you should build, validate, and maintain PDF digital signatures
today. A CA breach does not just compromise web traffic. It
retroactively undermines every document signed with certificates issued
by that CA --- contracts, tax filings, court submissions, medical
records. The signatures do not become invalid in a cryptographic sense;
the math still checks out. They become \emph{untrustworthy}, which in
the real world is worse.

\section{Why CA Breaches Matter for PDF
Signatures}\label{_why_ca_breaches_matter_for_pdf_signatures}

When you sign a PDF, the recipient validates your signature by chaining
your signing certificate back to a trusted root. That chain passes
through one or more intermediate CAs, each vouching for the next, until
it reaches a root certificate that the validator explicitly trusts ---
either because it is in Adobe's AATL, in the EU Trusted List, in the
operating system's trust store, or in a custom trust configuration.

The entire chain is only as strong as its weakest link. If any CA in the
chain is compromised, an attacker can issue fraudulent certificates ---
certificates that look legitimate but are controlled by the attacker.
Documents signed with those fraudulent certificates will pass validation
in any viewer that trusts the compromised CA.

Worse, when a compromised CA is discovered and revoked, the collateral
damage extends to \emph{every} legitimate certificate that CA ever
issued. If you signed a contract last year with a certificate from CA-X,
and CA-X is discovered to be compromised this year, your signature's
trustworthiness is now in question --- not because your key was
compromised, but because the entity vouching for your identity can no
longer be trusted.

This is the fundamental fragility of the PKI trust model: trust is
transitive, delegated, and revocable. You cannot sign a document and
walk away confident that the signature will still be trusted in five
years, because you cannot predict whether the CA that issued your
certificate will still be trusted in five years.

The only mitigation for CA compromise is \textbf{embedding revocation
evidence at signing time}. A PAdES-B-LTA signature includes OCSP
responses and CRLs captured when the signature was created. If the CA is
later revoked, the embedded evidence proves that the certificate was
valid \emph{at the time of signing} --- before the breach was
discovered. Without LTV data, a validator checking the signature after
the CA revocation has no way to distinguish between a legitimate
pre-breach signature and a fraudulent one.

\section{DigiNotar (2011): The Breach That Broke a
Country}\label{_diginotar_2011_the_breach_that_broke_a_country}

\subsection{Background: A Small CA With a Big
Responsibility}\label{_background_a_small_ca_with_a_big_responsibility}

DigiNotar was a Dutch certificate authority founded in 1998,
headquartered in Beverwijk, a small town near Amsterdam. The company was
originally an offshoot of the Dutch notarial profession --- the name is
a portmanteau of "Digital Notary" --- and it carved out a niche issuing
certificates for Dutch government services, notaries, and businesses.

In 2010, VASCO Data Security International (now OneSpan) acquired
DigiNotar for approximately 10 million euros. At the time, DigiNotar was
one of the root CAs trusted by every major browser and operating system.
Its certificates were embedded in the trust stores of Windows, macOS,
Firefox, Chrome, and by extension, in the trust stores used by PDF
viewers.

DigiNotar also operated \textbf{PKIoverheid} (PKI for Government) ---
the Dutch government's public key infrastructure. PKIoverheid
certificates were used across Dutch government digital services,
including the DigiD electronic identity system, tax filing, and ---
critically for our purposes --- official government document signing.

\subsection{The Attack}\label{_the_attack}

In June 2011, an attacker gained access to DigiNotar's certificate
issuance infrastructure. The attack was later attributed to a
state-sponsored actor operating from Iran, though the attacker also
posted messages claiming to be a lone Iranian hacker named
"Comodohacker" --- the same handle used in the Comodo breach earlier
that year.

The attacker exploited multiple vulnerabilities in DigiNotar's
infrastructure:

\begin{itemize}
\item
  The CA's servers ran outdated, unpatched software
\item
  Network segmentation between the public-facing servers and the
  certificate issuance systems was inadequate
\item
  The HSMs (Hardware Security Modules) that should have required
  physical authorization for certificate issuance were configured to
  allow automated remote issuance
\item
  Audit logs were incomplete and inconsistent --- DigiNotar did not have
  reliable records of which certificates had been issued
\end{itemize}

Once inside, the attacker issued over \textbf{531 fraudulent
certificates} for high-value domains. The confirmed targets included:

\begin{longtable}[]{@{}
  >{\raggedright\arraybackslash}p{(\linewidth - 2\tabcolsep) * \real{0.5000}}
  >{\raggedright\arraybackslash}p{(\linewidth - 2\tabcolsep) * \real{0.5000}}@{}}
\caption{Known fraudulently issued DigiNotar certificates (partial
list)}\tabularnewline
\toprule\noalign{}
\begin{minipage}[b]{\linewidth}\raggedright
Domain
\end{minipage} & \begin{minipage}[b]{\linewidth}\raggedright
Significance
\end{minipage} \\
\midrule\noalign{}
\endfirsthead
\toprule\noalign{}
\begin{minipage}[b]{\linewidth}\raggedright
Domain
\end{minipage} & \begin{minipage}[b]{\linewidth}\raggedright
Significance
\end{minipage} \\
\midrule\noalign{}
\endhead
\bottomrule\noalign{}
\endlastfoot
google.com, www.google.com & Gmail, Google Search --- primary
communication tool for millions \\
login.yahoo.com & Yahoo Mail \\
login.live.com & Microsoft Hotmail/Outlook \\
addons.mozilla.org & Firefox extension repository \\
login.skype.com & Skype communications \\
global.sts.microsoft.com & Microsoft federation services \\
www.sis.gov.uk & UK Secret Intelligence Service (MI6) \\
www.cia.gov & US Central Intelligence Agency \\
mossad.gov.il & Israeli intelligence service \\
azadegi.com & Iranian opposition news website \\
Several wildcard certs & Covering entire TLDs and subdomains \\
\end{longtable}

\subsection{The Use: Surveillance of Iranian
Citizens}\label{_the_use_surveillance_of_iranian_citizens}

The fraudulent certificates were not issued for financial fraud. They
were used for \textbf{mass surveillance}.

In July and August 2011, the fraudulent \texttt{google.com} certificate
was deployed in man-in-the-middle (MITM) attacks targeting Iranian
internet users. The Iranian government (or an actor with access to
Iranian internet infrastructure) intercepted HTTPS connections to Google
services --- primarily Gmail --- by presenting the fraudulent DigiNotar
certificate. Because DigiNotar was a trusted root CA, browsers accepted
the certificate without warning.

The victims were Iranian citizens communicating via Gmail during a
period of significant political tension. The Arab Spring protests had
swept through Tunisia, Egypt, Libya, and Syria. Iran's Green Movement,
which had erupted after the disputed 2009 presidential election, was
still active. The surveillance targeted people who thought their
communications were private.

The scale of the interception was significant. Fox-IT, the Dutch
forensics firm that later investigated the breach, estimated that
approximately \textbf{300,000 unique IP addresses} --- almost all
Iranian --- had been subjected to MITM attacks using the fraudulent
Google certificate.

\subsection{Discovery: How Certificate Pinning Saved the
Day}\label{_discovery_how_certificate_pinning_saved_the_day}

The attack was discovered on August 28, 2011, by an Iranian user of
Google Chrome.

Chrome has a feature called \textbf{certificate pinning} (also known as
HPKP, HTTP Public Key Pinning, or in Chrome's case, a hardcoded pin
set). For a small number of high-value domains --- including google.com
--- Chrome maintains a built-in list of expected certificate public
keys. If the presented certificate's public key does not match any of
the pinned keys, Chrome refuses the connection and displays a prominent
error, even if the certificate otherwise validates correctly against the
trust store.

When the Iranian Chrome user attempted to connect to Google, Chrome
detected that the presented certificate (signed by DigiNotar) did not
match Google's pinned public keys. Chrome displayed an error. The user
reported the error to Google.

Certificate pinning was not designed as a comprehensive defense against
CA compromise. It only worked for the handful of domains that Google had
pre-configured pins for. If the attacker had only issued certificates
for domains that Chrome did not pin --- the intelligence agency
websites, the Iranian opposition news site, the Yahoo and Skype domains
--- the attack might have continued undetected for months.

The lesson is humbling: the discovery of one of the most significant CA
compromises in internet history was essentially an accident --- it
depended on one user, in one country, using one specific browser,
connecting to one specific domain that happened to have certificate
pinning enabled.

\subsection{The Response: Death of a
CA}\label{_the_response_death_of_a_ca}

Once Google confirmed the fraudulent certificate, events moved rapidly:

\begin{itemize}
\item
  \textbf{August 29, 2011}: Google published a Chrome update blocking
  the specific DigiNotar intermediate certificate used to sign the
  fraudulent Google certificate
\item
  \textbf{August 30, 2011}: Mozilla pushed a Firefox update removing
  DigiNotar's root from the trust store entirely
\item
  \textbf{September 2, 2011}: The Dutch government's GovCERT (Government
  Computer Emergency Response Team) issued an advisory acknowledging the
  breach
\item
  \textbf{September 3, 2011}: Google, Microsoft, and Apple all pushed
  updates removing DigiNotar from their trust stores
\item
  \textbf{September 5, 2011}: The Dutch government formally revoked
  DigiNotar's authorization to issue PKIoverheid certificates
\item
  \textbf{September 20, 2011}: DigiNotar's parent company VASCO filed
  for DigiNotar's bankruptcy in Dutch court
\end{itemize}

The entire collapse --- from discovery to bankruptcy --- took less than
four weeks.

Fox-IT's forensic investigation, published in their report "DigiNotar
Certificate Authority breach --- \textquotesingle Operation Black
Tulip\textquotesingle" (named with Dutch irony), found that the breach
was far more extensive than initially reported. The attacker had been
inside DigiNotar's systems since at least June 17, 2011, had compromised
multiple CA servers, and had access to issuing infrastructure for both
DigiNotar's commercial certificates and the PKIoverheid government
certificates.

\subsection{Impact on PDF Signatures}\label{_impact_on_pdf_signatures}

The immediate impact on PDF signatures was severe and irreversible.

Any PDF document signed with a DigiNotar-issued certificate became
\textbf{unverifiable} the moment browsers and operating systems removed
DigiNotar from their trust stores. Adobe Acrobat, which uses the Windows
certificate store (on Windows) or its own AATL (on macOS), stopped
trusting DigiNotar certificates. A signed PDF that showed a green
checkmark on September 2 showed a warning on September 4.

The documents themselves had not changed. The signatures had not been
tampered with. The signing keys had not been compromised (in most cases
--- the legitimate certificate holders were not the ones attacked). But
the trust anchor for those certificates no longer existed.

For Dutch government agencies, this was particularly painful.
PKIoverheid certificates were used to sign official government
documents, notarial deeds, and legal filings. The Dutch government had
to issue emergency replacement certificates through other CAs and
re-sign critical documents --- a process that took months.

This is the nightmare scenario for long-lived signed documents. A
contract signed in 2010 with a DigiNotar certificate and stored in an
archive had its trust chain severed in 2011. If the signature did not
include LTV data (OCSP responses and CRLs embedded at signing time, as
in PAdES-B-LTA), there was \emph{no way} to prove after the fact that
the certificate was valid when the signature was created. The document
was, for all practical purposes, unsigned.

This is why PAdES-B-LTA exists. This is why you embed revocation data.
This is not a theoretical concern --- it happened.

\subsection{Lessons from DigiNotar}\label{_lessons_from_diginotar}

The DigiNotar breach taught the PKI community several painful lessons:

\begin{enumerate}
\def\labelenumi{\arabic{enumi}.}
\item
  \textbf{A single CA compromise can retroactively invalidate thousands
  of signatures.} Trust is not permanent. It is not even stable. It can
  be revoked at any time, with retrospective effect.
\item
  \textbf{CA security is only as strong as the weakest operational
  practice.} DigiNotar had HSMs, but they were misconfigured. They had
  audit logs, but they were incomplete. They had network segmentation,
  but it was insufficient. The technology existed to prevent the breach;
  the organization failed to deploy it correctly.
\item
  \textbf{Certificate pinning is a band-aid, not a solution.} It helped
  detect \emph{this} breach, but it does not scale. You cannot pin every
  domain's certificate in every application.
\item
  \textbf{Government PKI is not immune.} DigiNotar operated critical
  Dutch government infrastructure. Government CAs are not inherently
  more trustworthy than commercial ones --- they are subject to the same
  operational failures.
\item
  \textbf{Bankruptcy is the final revocation.} When a CA goes bankrupt,
  its certificates do not gracefully expire. They are abruptly
  distrusted. There is no migration path, no deprecation period, no
  gentle sunset. One day the certificates work; the next day they do
  not.
\end{enumerate}

\section{Comodo (2011): The Parallel
Breach}\label{_comodo_2011_the_parallel_breach}

\subsection{The Attack}\label{_the_attack_2}

Just five months before DigiNotar's collapse, in March 2011, a similar
attack targeted Comodo --- at the time, one of the world's largest
commercial CAs.

The attacker compromised \textbf{registration authority (RA) accounts}
at Comodo's reseller partners --- specifically, the Italian firm
GlobalTrust.it and the Argentine firm InstantSSL.it. Registration
authorities are entities authorized to perform identity verification on
behalf of a CA; they can request certificate issuance without having
direct access to the CA's root key.

Using the compromised RA credentials, the attacker obtained \textbf{nine
fraudulent certificates} for high-value domains:

\begin{itemize}
\item
  login.live.com (Microsoft)
\item
  mail.google.com (Google)
\item
  www.google.com (Google)
\item
  login.yahoo.com (Yahoo, three separate certificates)
\item
  login.skype.com (Skype)
\item
  addons.mozilla.org (Mozilla)
\item
  An additional "global trustee" certificate
\end{itemize}

The target list was remarkably similar to the DigiNotar targets ---
because the same attacker was behind both breaches.

\subsection{The Response: How Comodo Got It
Right}\label{_the_response_how_comodo_got_it_right}

Comodo's response was dramatically faster and more effective than
DigiNotar's would be five months later.

Comodo detected the fraudulent issuance \textbf{within hours} --- their
monitoring systems flagged the unusual certificate requests from the
compromised RA accounts. The nine certificates were revoked within the
same day. Comodo published a detailed incident report, contacted browser
vendors, and cooperated fully with the investigation.

The key differences between Comodo's response and DigiNotar's later
failure:

\begin{longtable}[]{@{}
  >{\raggedright\arraybackslash}p{(\linewidth - 4\tabcolsep) * \real{0.3333}}
  >{\raggedright\arraybackslash}p{(\linewidth - 4\tabcolsep) * \real{0.3333}}
  >{\raggedright\arraybackslash}p{(\linewidth - 4\tabcolsep) * \real{0.3333}}@{}}
\caption{Comodo vs DigiNotar incident response
comparison}\tabularnewline
\toprule\noalign{}
\begin{minipage}[b]{\linewidth}\raggedright
Factor
\end{minipage} & \begin{minipage}[b]{\linewidth}\raggedright
Comodo
\end{minipage} & \begin{minipage}[b]{\linewidth}\raggedright
DigiNotar
\end{minipage} \\
\midrule\noalign{}
\endfirsthead
\toprule\noalign{}
\begin{minipage}[b]{\linewidth}\raggedright
Factor
\end{minipage} & \begin{minipage}[b]{\linewidth}\raggedright
Comodo
\end{minipage} & \begin{minipage}[b]{\linewidth}\raggedright
DigiNotar
\end{minipage} \\
\midrule\noalign{}
\endhead
\bottomrule\noalign{}
\endlastfoot
Detection time & Hours & Weeks (possibly months) \\
Number of fraudulent certs & 9 & 531+ \\
Disclosure & Immediate and transparent & Delayed and incomplete \\
Certificate revocation & Same day & Gradual, over weeks \\
Root trust impact & Root preserved (reseller revoked) & Entire root
removed \\
Business outcome & Continued operating & Bankruptcy \\
CA audit logs & Complete and useful & Incomplete and unreliable \\
\end{longtable}

Comodo survived because it detected the breach quickly, responded
transparently, and demonstrated that the compromise was limited to
reseller accounts --- not the root CA infrastructure. The attackers got
nine certificates; at DigiNotar, they got the keys to the kingdom.

\subsection{The Attacker's Manifesto}\label{_the_attackers_manifesto}

The attacker posted a manifesto online claiming to be a 21-year-old
Iranian, acting alone, motivated by anger at Western
governments\textquotesingle{} treatment of Iran. The manifesto, written
in broken English, contained enough technical detail about both the
Comodo and DigiNotar breaches to be credible. Whether the attacker was
truly a lone wolf or a front for state-sponsored activity remains
debated. The operational sophistication of the MITM infrastructure
deployed against Iranian citizens (requiring cooperation with or access
to Iranian ISP infrastructure) suggests state involvement, regardless of
who wrote the code.

From a PDF signing perspective, the Comodo incident reinforces a
critical point: \textbf{your signing certificate's trustworthiness
depends not only on your CA, but on every reseller, registration
authority, and partner in the CA's issuance chain}. When selecting a CA
for document signing certificates, ask about their RA model. How many
resellers can request certificate issuance? What monitoring exists for
anomalous issuance patterns? How quickly can fraudulent certificates be
detected and revoked?

\section{Symantec Distrust (2017-2018): The Slow-Motion
Collapse}\label{_symantec_distrust_2017_2018_the_slow_motion_collapse}

\subsection{Background: The World's Largest
CA}\label{_background_the_worlds_largest_ca}

If DigiNotar was a small CA that died quickly, Symantec was a giant that
died slowly.

Through a series of acquisitions --- \textbf{VeriSign} (2010),
\textbf{Thawte} (via VeriSign), \textbf{GeoTrust} (via VeriSign), and
\textbf{RapidSSL} (via GeoTrust) --- Symantec had become the world's
largest certificate authority by 2015. At its peak, Symantec and its
subsidiary brands had issued approximately \textbf{30\% of all TLS
certificates on the public internet}.

That market share made Symantec's CA operation systemically important.
It also made it a tempting target for scrutiny.

\subsection{The Mis-Issuance Scandal}\label{_the_mis_issuance_scandal}

In September 2015, Google's Certificate Transparency (CT) monitoring
detected that Symantec had issued several test certificates for
google.com domains --- without Google's authorization. Symantec claimed
these were test certificates issued by employees during internal
testing, and that they had been revoked quickly.

Google investigated further and discovered that the problem was much
larger:

\begin{itemize}
\item
  Symantec had issued \textbf{at least 164 certificates} for domains
  without proper authorization
\item
  Some were issued by Symantec employees, others by RA partners
  (including CrossCert, a Korean RA)
\item
  The unauthorized issuance spanned multiple years
\item
  Symantec's internal controls had failed to detect the mis-issuances
\item
  When confronted, Symantec's initial disclosures underreported the
  scope of the problem
\end{itemize}

In March 2017, Google engineers Ryan Sleevi and Devon O'Brien published
a detailed analysis on the Chromium bug tracker, proposing a phased
distrust of all Symantec-issued certificates. The proposal was
unprecedented: no one had ever suggested distrusting certificates from
the world's largest CA.

\subsection{The Distrust Timeline}\label{_the_distrust_timeline}

What followed was a two-year negotiation between Google, Mozilla,
Symantec, and the broader CA/Browser Forum. The timeline:

\begin{itemize}
\item
  \textbf{March 2017}: Google proposes phased Symantec distrust in
  Chrome
\item
  \textbf{April 2017}: Symantec contests the proposal, arguing the
  issues were limited
\item
  \textbf{July 2017}: Google, Mozilla, and Symantec reach a compromise
  plan: Symantec must transition to a new infrastructure operated by a
  different CA
\item
  \textbf{August 2017}: DigiCert announces acquisition of Symantec's CA
  business for approximately \$950 million
\item
  \textbf{December 2017}: DigiCert completes the acquisition and begins
  re-issuing certificates under DigiCert roots
\item
  \textbf{April 2018}: Chrome 66 distrusts Symantec certificates issued
  before June 1, 2016
\item
  \textbf{October 2018}: Chrome 70 distrusts ALL remaining Symantec
  certificates
\item
  \textbf{October 2018}: Firefox follows with similar distrust in
  Firefox 63
\end{itemize}

The distrust was progressive: older certificates were distrusted first,
giving website operators time to obtain replacement certificates. But
the endpoint was absolute: after October 2018, no Symantec-issued
certificate was trusted by Chrome or Firefox.

\subsection{Impact on PDF Signatures}\label{_impact_on_pdf_signatures_2}

The Symantec distrust had a less dramatic but more insidious impact on
PDF signatures than the DigiNotar breach.

DigiNotar's collapse was sudden --- trust was revoked in days, and the
relatively small number of affected certificates limited the blast
radius. Symantec's distrust was gradual, but the scale was enormous.

\begin{itemize}
\item
  \textbf{Thousands of document signing certificates} had been issued
  under Symantec, VeriSign, Thawte, and GeoTrust brands
\item
  Many of these certificates were in Adobe's AATL (Adobe Approved Trust
  List)
\item
  As Symantec certificates were distrusted, PDF signatures chaining to
  those certificates showed warnings in validators that followed
  Chrome/Firefox trust decisions
\item
  Adobe's trust store (AATL) moved independently from browsers --- Adobe
  did not automatically distrust Symantec certificates just because
  Chrome did
\item
  This created a confusing situation where a signed PDF might validate
  in Acrobat (using AATL) but fail in a web-based validator (using the
  browser trust store)
\end{itemize}

The Symantec case illustrates an important distinction: \textbf{PDF
trust stores and web browser trust stores are not the same thing}. Adobe
Acrobat maintains its own trust list (AATL), which is updated on its own
schedule, with its own criteria. A certificate can be distrusted by
Chrome but still trusted by Acrobat --- or vice versa. This is both a
feature (PDF signatures are not automatically affected by browser trust
decisions) and a bug (trust status becomes inconsistent across different
validation contexts).

The transition to DigiCert mitigated some of the damage --- DigiCert
re-issued certificates under its own roots, and existing Symantec
customers could obtain replacement certificates. But for documents
already signed with Symantec-chain certificates, the question of
long-term trust remained unresolved.

\begin{Shaded}
\begin{Highlighting}[]
\CommentTok{/**}
 \CommentTok{*}\NormalTok{ Checks whether a signing certificate chains to a currently trusted root}\CommentTok{.}
 \CommentTok{*}\NormalTok{ This is the kind of check you need to perform when validating old signatures}
 \CommentTok{*}\NormalTok{ whose CAs may have been distrusted since signing time}\CommentTok{.}
 \CommentTok{*/}
\KeywordTok{public}\NormalTok{ CertificateChainStatus }\FunctionTok{validateChainTrust}\OperatorTok{(}
        \BuiltInTok{X509Certificate}\NormalTok{ signerCert}\OperatorTok{,}
        \BuiltInTok{List}\OperatorTok{\textless{}}\BuiltInTok{X509Certificate}\OperatorTok{\textgreater{}}\NormalTok{ chain}\OperatorTok{,}
        \BuiltInTok{Set}\OperatorTok{\textless{}}\BuiltInTok{TrustAnchor}\OperatorTok{\textgreater{}}\NormalTok{ currentTrustAnchors}\OperatorTok{,}
        \BuiltInTok{Date}\NormalTok{ signingTime}\OperatorTok{)} \KeywordTok{throws} \BuiltInTok{Exception} \OperatorTok{\{}

    \CommentTok{// Build the certificate path}
    \BuiltInTok{CertificateFactory}\NormalTok{ cf }\OperatorTok{=} \BuiltInTok{CertificateFactory}\OperatorTok{.}\FunctionTok{getInstance}\OperatorTok{(}\StringTok{"X.509"}\OperatorTok{);}
    \BuiltInTok{CertPath}\NormalTok{ certPath }\OperatorTok{=}\NormalTok{ cf}\OperatorTok{.}\FunctionTok{generateCertPath}\OperatorTok{(}\NormalTok{chain}\OperatorTok{);}

    \CommentTok{// Configure the path validator with the signing time}
    \BuiltInTok{PKIXParameters}\NormalTok{ params }\OperatorTok{=} \KeywordTok{new} \BuiltInTok{PKIXParameters}\OperatorTok{(}\NormalTok{currentTrustAnchors}\OperatorTok{);}
\NormalTok{    params}\OperatorTok{.}\FunctionTok{setDate}\OperatorTok{(}\NormalTok{signingTime}\OperatorTok{);}  \CommentTok{// Validate as of signing time, not now}
\NormalTok{    params}\OperatorTok{.}\FunctionTok{setRevocationEnabled}\OperatorTok{(}\KeywordTok{false}\OperatorTok{);}  \CommentTok{// We check revocation separately}

    \ControlFlowTok{try} \OperatorTok{\{}
        \BuiltInTok{CertPathValidator}\NormalTok{ validator }\OperatorTok{=} \BuiltInTok{CertPathValidator}\OperatorTok{.}\FunctionTok{getInstance}\OperatorTok{(}\StringTok{"PKIX"}\OperatorTok{);}
        \BuiltInTok{PKIXCertPathValidatorResult}\NormalTok{ result }\OperatorTok{=}
                \OperatorTok{(}\BuiltInTok{PKIXCertPathValidatorResult}\OperatorTok{)}\NormalTok{ validator}\OperatorTok{.}\FunctionTok{validate}\OperatorTok{(}\NormalTok{certPath}\OperatorTok{,}\NormalTok{ params}\OperatorTok{);}

        \BuiltInTok{TrustAnchor}\NormalTok{ anchor }\OperatorTok{=}\NormalTok{ result}\OperatorTok{.}\FunctionTok{getTrustAnchor}\OperatorTok{();}
\NormalTok{        log}\OperatorTok{.}\FunctionTok{info}\OperatorTok{(}\StringTok{"Certificate chain valid at signing time. Trust anchor: \{\}"}\OperatorTok{,}
\NormalTok{                anchor}\OperatorTok{.}\FunctionTok{getTrustedCert}\OperatorTok{().}\FunctionTok{getSubjectX500Principal}\OperatorTok{());}

        \ControlFlowTok{return}\NormalTok{ CertificateChainStatus}\OperatorTok{.}\FunctionTok{builder}\OperatorTok{()}
                \OperatorTok{.}\FunctionTok{valid}\OperatorTok{(}\KeywordTok{true}\OperatorTok{)}
                \OperatorTok{.}\FunctionTok{trustAnchor}\OperatorTok{(}\NormalTok{anchor}\OperatorTok{.}\FunctionTok{getTrustedCert}\OperatorTok{().}\FunctionTok{getSubjectX500Principal}\OperatorTok{().}\FunctionTok{getName}\OperatorTok{())}
                \OperatorTok{.}\FunctionTok{validationTime}\OperatorTok{(}\NormalTok{signingTime}\OperatorTok{)}
                \OperatorTok{.}\FunctionTok{build}\OperatorTok{();}

    \OperatorTok{\}} \ControlFlowTok{catch} \OperatorTok{(}\BuiltInTok{CertPathValidatorException}\NormalTok{ e}\OperatorTok{)} \OperatorTok{\{}
\NormalTok{        log}\OperatorTok{.}\FunctionTok{warn}\OperatorTok{(}\StringTok{"Certificate chain validation failed at index \{\}: \{\}"}\OperatorTok{,}
\NormalTok{                e}\OperatorTok{.}\FunctionTok{getIndex}\OperatorTok{(),}\NormalTok{ e}\OperatorTok{.}\FunctionTok{getMessage}\OperatorTok{());}

        \ControlFlowTok{return}\NormalTok{ CertificateChainStatus}\OperatorTok{.}\FunctionTok{builder}\OperatorTok{()}
                \OperatorTok{.}\FunctionTok{valid}\OperatorTok{(}\KeywordTok{false}\OperatorTok{)}
                \OperatorTok{.}\FunctionTok{failureReason}\OperatorTok{(}\NormalTok{e}\OperatorTok{.}\FunctionTok{getMessage}\OperatorTok{())}
                \OperatorTok{.}\FunctionTok{failedCertificateIndex}\OperatorTok{(}\NormalTok{e}\OperatorTok{.}\FunctionTok{getIndex}\OperatorTok{())}
                \OperatorTok{.}\FunctionTok{validationTime}\OperatorTok{(}\NormalTok{signingTime}\OperatorTok{)}
                \OperatorTok{.}\FunctionTok{build}\OperatorTok{();}
    \OperatorTok{\}}
\OperatorTok{\}}
\end{Highlighting}
\end{Shaded}

\section{Let's Encrypt: The Free CA's Growing
Pains}\label{_lets_encrypt_the_free_cas_growing_pains}

\subsection{The 2020 Mass Revocation}\label{_the_2020_mass_revocation}

Let's Encrypt, the free, automated CA that has issued over three billion
certificates since its launch in 2015, does not issue document signing
certificates. Its certificates are exclusively for TLS (web server) use.
So why does a Let's Encrypt incident appear in a book about PDF
signatures?

Because the 2020 mass revocation event demonstrated something important
about the fragility of CA infrastructure at scale --- a fragility that
applies equally to every CA in the PDF signing ecosystem.

On March 3, 2020, Let's Encrypt announced that it had discovered a bug
in its \textbf{CAA (Certificate Authority Authorization) checking code}.
CAA is a DNS-based mechanism that allows domain owners to specify which
CAs are authorized to issue certificates for their domain. Before
issuing a certificate, a CA must check the domain's CAA records and
refuse issuance if the domain specifies a different CA.

The bug: Let's Encrypt's validation system checked CAA records when a
domain was first validated, but did not re-check them when issuing a
certificate if the validation result was cached. This meant that if a
domain's CAA record changed after the initial validation (for example,
to explicitly prohibit Let's Encrypt from issuing certificates), Let's
Encrypt would still issue the certificate based on the stale cached
result.

This violated the CA/Browser Forum's Baseline Requirements, which
mandate that CAA records be checked within 8 hours of certificate
issuance.

\subsection{The Chaos}\label{_the_chaos}

Let's Encrypt identified approximately \textbf{3,048,289 certificates}
that had been issued without a compliant CAA check at issuance time. The
Baseline Requirements required these certificates to be revoked within 5
days.

Let's Encrypt gave certificate holders \textbf{24 hours} to renew their
certificates before revocation. For the ACME-based automated renewal
ecosystem that Let's Encrypt had built, this should have been manageable
--- the whole point of Let's Encrypt is automated certificate
management.

In practice, it was chaos. Not every certificate holder had functional
automated renewal. Not every server was properly configured for ACME.
Some certificates were on devices that could not easily be updated. The
announcement came on a Tuesday evening (US time), giving operators one
business day to respond.

\subsection{Relevance to PDF Signing}\label{_relevance_to_pdf_signing}

The Let's Encrypt incident did not directly affect any PDF signatures.
But it demonstrated three principles that every PDF signing architect
should internalize:

\begin{enumerate}
\def\labelenumi{\arabic{enumi}.}
\item
  \textbf{Scale amplifies errors.} A single bug in a CA's issuance
  pipeline can affect millions of certificates. The larger the CA, the
  larger the blast radius. When choosing a CA for document signing,
  understand that the CA's operational quality affects every certificate
  they issue --- including yours.
\item
  \textbf{Revocation at scale is operationally brutal.} Even with 24
  hours notice and automated tooling, millions of certificates could not
  be renewed in time. Now imagine the same scenario with document
  signing certificates, where "renewal" means re-signing every document
  that was signed with the revoked certificate. There is no automated
  tool for that.
\item
  \textbf{Compliance requirements drive revocation timelines.} The
  Baseline Requirements gave Let's Encrypt no discretion --- the
  certificates \emph{had} to be revoked within 5 days, regardless of the
  impact. EU regulations impose similar requirements on Qualified Trust
  Service Providers. If a QTSP discovers a compliance issue, it may be
  required to revoke certificates on a timeline that does not allow for
  graceful migration.
\end{enumerate}

If your QTSP discovers a compliance issue and is required to revoke your
signing certificate, every PDF you signed with that certificate becomes
suspect --- unless you embedded LTV data at signing time. This is not a
far-fetched scenario. It has happened to web certificates. It will
happen to document signing certificates. Plan accordingly.

\section{Certificate Transparency: The Response to CA
Failures}\label{_certificate_transparency_the_response_to_ca_failures}

\subsection{What Certificate Transparency
Is}\label{_what_certificate_transparency_is}

Certificate Transparency (CT) was invented by Google engineers
\textbf{Ben Laurie} and \textbf{Adam Langley} in 2012 (published as RFC
5280 --- Internet X.509 PKI Certificate and CRL Profile companion
document RFC 6962, and later updated in RFC 9162), directly in response
to the DigiNotar and Comodo breaches.

The core insight: if CAs cannot be trusted to only issue authorized
certificates, we need a public audit mechanism that makes \emph{every}
issued certificate visible to domain owners and the broader community.

CT works as follows:

\begin{enumerate}
\def\labelenumi{\arabic{enumi}.}
\item
  When a CA issues a certificate, it submits the certificate to one or
  more \textbf{CT logs} --- append-only, cryptographically verifiable
  public databases
\item
  The CT log returns a \textbf{Signed Certificate Timestamp (SCT)} --- a
  promise that the certificate will be included in the log within a
  specified time (the "Maximum Merge Delay," typically 24 hours)
\item
  The SCT is embedded in the certificate (as an X.509 extension), in the
  TLS handshake (as a TLS extension), or in the OCSP response (as an
  OCSP extension)
\item
  Browsers verify that certificates come with valid SCTs --- if a
  certificate has no SCT, Chrome treats it as untrusted (since April
  2018)
\item
  Domain owners (and anyone else) can monitor CT logs for certificates
  issued for their domains --- if an unauthorized certificate appears,
  they can demand revocation
\end{enumerate}

\subsection{How CT Detects Rogue
Certificates}\label{_how_ct_detects_rogue_certificates}

CT does not \emph{prevent} a compromised CA from issuing rogue
certificates. It makes rogue certificates \emph{detectable}.

If the DigiNotar attack had happened in a CT world, the fraudulent
\texttt{google.com} certificate would have appeared in public CT logs
within 24 hours of issuance. Google's CT monitoring infrastructure would
have detected it immediately --- not because a user in Iran happened to
be running Chrome, but because Google continuously monitors every CT log
for certificates mentioning google.com.

The detection would have been systematic, automatic, and fast.

\subsection{CT and PDF: Limited but
Growing}\label{_ct_and_pdf_limited_but_growing}

Certificate Transparency was designed for the web PKI ecosystem. Its
adoption in the document signing world has been slower:

\begin{itemize}
\item
  \textbf{Document signing certificates are not required to have SCTs.}
  The CA/Browser Forum Baseline Requirements mandate CT for TLS
  certificates, but no equivalent mandate exists for document signing
  certificates.
\item
  \textbf{Adobe's AATL does not require CT.} CAs in the AATL are not
  required to log their document signing certificates in CT logs.
\item
  \textbf{The EU Trusted List does not reference CT.} The eIDAS
  regulatory framework does not mention Certificate Transparency.
\end{itemize}

However, CT is increasingly relevant for PDF signatures:

\begin{itemize}
\item
  Some CAs voluntarily log their document signing certificates in CT
  logs
\item
  Organizations building their own signature validation can incorporate
  CT log monitoring as an additional trust signal
\item
  The EUDI Wallet infrastructure (Chapter 27) may incorporate CT-like
  transparency mechanisms for wallet-issued certificates
\end{itemize}

\begin{Shaded}
\begin{Highlighting}[]
\CommentTok{/**}
 \CommentTok{*}\NormalTok{ Fetches and embeds revocation data }\CommentTok{(}\NormalTok{OCSP and CRL}\CommentTok{)}\NormalTok{ at signing time}\CommentTok{,}
 \CommentTok{*}\NormalTok{ creating the foundation for long}\CommentTok{{-}}\NormalTok{term validation that survives CA revocation}\CommentTok{.}
 \CommentTok{*}\NormalTok{ This is the essential defense against CA breach impact}\CommentTok{.}
 \CommentTok{*/}
\KeywordTok{public}\NormalTok{ LtvData }\FunctionTok{collectRevocationData}\OperatorTok{(}\BuiltInTok{List}\OperatorTok{\textless{}}\BuiltInTok{X509Certificate}\OperatorTok{\textgreater{}}\NormalTok{ chain}\OperatorTok{)}
        \KeywordTok{throws} \BuiltInTok{Exception} \OperatorTok{\{}

    \BuiltInTok{List}\OperatorTok{\textless{}}\DataTypeTok{byte}\OperatorTok{[]\textgreater{}}\NormalTok{ ocspResponses }\OperatorTok{=} \KeywordTok{new} \BuiltInTok{ArrayList}\OperatorTok{\textless{}\textgreater{}();}
    \BuiltInTok{List}\OperatorTok{\textless{}}\DataTypeTok{byte}\OperatorTok{[]\textgreater{}}\NormalTok{ crls }\OperatorTok{=} \KeywordTok{new} \BuiltInTok{ArrayList}\OperatorTok{\textless{}\textgreater{}();}

    \CommentTok{// For each certificate in the chain (except the root),}
    \CommentTok{// fetch OCSP response and/or CRL}
    \ControlFlowTok{for} \OperatorTok{(}\DataTypeTok{int}\NormalTok{ i }\OperatorTok{=} \DecValTok{0}\OperatorTok{;}\NormalTok{ i }\OperatorTok{\textless{}}\NormalTok{ chain}\OperatorTok{.}\FunctionTok{size}\OperatorTok{()} \OperatorTok{{-}} \DecValTok{1}\OperatorTok{;}\NormalTok{ i}\OperatorTok{++)} \OperatorTok{\{}
        \BuiltInTok{X509Certificate}\NormalTok{ cert }\OperatorTok{=}\NormalTok{ chain}\OperatorTok{.}\FunctionTok{get}\OperatorTok{(}\NormalTok{i}\OperatorTok{);}
        \BuiltInTok{X509Certificate}\NormalTok{ issuer }\OperatorTok{=}\NormalTok{ chain}\OperatorTok{.}\FunctionTok{get}\OperatorTok{(}\NormalTok{i }\OperatorTok{+} \DecValTok{1}\OperatorTok{);}

        \CommentTok{// Try OCSP first (preferred — smaller, more specific)}
        \DataTypeTok{byte}\OperatorTok{[]}\NormalTok{ ocspResponse }\OperatorTok{=} \FunctionTok{fetchOcspResponse}\OperatorTok{(}\NormalTok{cert}\OperatorTok{,}\NormalTok{ issuer}\OperatorTok{);}
        \ControlFlowTok{if} \OperatorTok{(}\NormalTok{ocspResponse }\OperatorTok{!=} \KeywordTok{null}\OperatorTok{)} \OperatorTok{\{}
\NormalTok{            log}\OperatorTok{.}\FunctionTok{info}\OperatorTok{(}\StringTok{"OCSP response obtained for: \{\}"}\OperatorTok{,}
\NormalTok{                    cert}\OperatorTok{.}\FunctionTok{getSubjectX500Principal}\OperatorTok{());}
\NormalTok{            ocspResponses}\OperatorTok{.}\FunctionTok{add}\OperatorTok{(}\NormalTok{ocspResponse}\OperatorTok{);}
        \OperatorTok{\}}

        \CommentTok{// Also fetch CRL as fallback}
        \DataTypeTok{byte}\OperatorTok{[]}\NormalTok{ crl }\OperatorTok{=} \FunctionTok{fetchCrl}\OperatorTok{(}\NormalTok{cert}\OperatorTok{);}
        \ControlFlowTok{if} \OperatorTok{(}\NormalTok{crl }\OperatorTok{!=} \KeywordTok{null}\OperatorTok{)} \OperatorTok{\{}
\NormalTok{            log}\OperatorTok{.}\FunctionTok{info}\OperatorTok{(}\StringTok{"CRL obtained for: \{\} (\{\} bytes)"}\OperatorTok{,}
\NormalTok{                    cert}\OperatorTok{.}\FunctionTok{getSubjectX500Principal}\OperatorTok{(),}\NormalTok{ crl}\OperatorTok{.}\FunctionTok{length}\OperatorTok{);}
\NormalTok{            crls}\OperatorTok{.}\FunctionTok{add}\OperatorTok{(}\NormalTok{crl}\OperatorTok{);}
        \OperatorTok{\}}

        \ControlFlowTok{if} \OperatorTok{(}\NormalTok{ocspResponse }\OperatorTok{==} \KeywordTok{null} \OperatorTok{\&\&}\NormalTok{ crl }\OperatorTok{==} \KeywordTok{null}\OperatorTok{)} \OperatorTok{\{}
\NormalTok{            log}\OperatorTok{.}\FunctionTok{warn}\OperatorTok{(}\StringTok{"No revocation data available for: \{\}. "} \OperatorTok{+}
                    \StringTok{"LTV validation may fail for this certificate."}\OperatorTok{,}
\NormalTok{                    cert}\OperatorTok{.}\FunctionTok{getSubjectX500Principal}\OperatorTok{());}
        \OperatorTok{\}}
    \OperatorTok{\}}

    \ControlFlowTok{return}\NormalTok{ LtvData}\OperatorTok{.}\FunctionTok{builder}\OperatorTok{()}
            \OperatorTok{.}\FunctionTok{certificates}\OperatorTok{(}\NormalTok{chain}\OperatorTok{.}\FunctionTok{stream}\OperatorTok{()}
                    \OperatorTok{.}\FunctionTok{map}\OperatorTok{(}\NormalTok{c }\OperatorTok{{-}\textgreater{}} \OperatorTok{\{}
                        \ControlFlowTok{try} \OperatorTok{\{} \ControlFlowTok{return}\NormalTok{ c}\OperatorTok{.}\FunctionTok{getEncoded}\OperatorTok{();} \OperatorTok{\}}
                        \ControlFlowTok{catch} \OperatorTok{(}\BuiltInTok{Exception}\NormalTok{ e}\OperatorTok{)} \OperatorTok{\{} \ControlFlowTok{throw} \KeywordTok{new} \BuiltInTok{RuntimeException}\OperatorTok{(}\NormalTok{e}\OperatorTok{);} \OperatorTok{\}}
                    \OperatorTok{\})}
                    \OperatorTok{.}\FunctionTok{collect}\OperatorTok{(}\NormalTok{Collectors}\OperatorTok{.}\FunctionTok{toList}\OperatorTok{()))}
            \OperatorTok{.}\FunctionTok{ocspResponses}\OperatorTok{(}\NormalTok{ocspResponses}\OperatorTok{)}
            \OperatorTok{.}\FunctionTok{crls}\OperatorTok{(}\NormalTok{crls}\OperatorTok{)}
            \OperatorTok{.}\FunctionTok{collectionTime}\OperatorTok{(}\KeywordTok{new} \BuiltInTok{Date}\OperatorTok{())}
            \OperatorTok{.}\FunctionTok{build}\OperatorTok{();}
\OperatorTok{\}}
\end{Highlighting}
\end{Shaded}

\section{What These Incidents Taught
Us}\label{_what_these_incidents_taught_us}

\subsection{Trust Is Fragile and
Revocable}\label{_trust_is_fragile_and_revocable}

The PKI model assumes that root CAs are trustworthy. DigiNotar proved
that assumption can fail catastrophically. Symantec proved it can fail
gradually. Comodo proved it can fail and be recovered from --- but only
if the CA responds competently.

For PDF signature architects, the lesson is: \textbf{design for CA
failure}. Assume that any CA in your signing certificate's chain might
be distrusted at some point in the future. Your architecture should
handle that scenario gracefully.

\subsection{Long-Term Validity Requires Embedded
Evidence}\label{_long_term_validity_requires_embedded_evidence}

The single most important defense against CA breach impact is
\textbf{embedding revocation data at signing time}.

A PAdES-B-LTA signature includes:

\begin{itemize}
\item
  The complete certificate chain (signer, intermediates, root)
\item
  OCSP responses for every certificate in the chain, obtained at signing
  time
\item
  CRLs as fallback revocation evidence
\item
  A signature timestamp proving when the signature was created
\item
  A document timestamp proving when the validation data was embedded
\end{itemize}

With all this data embedded, a validator can determine that the signing
certificate was valid \emph{at the time of signing} --- even if the CA
is later revoked. The embedded OCSP response from 2024 proves the
certificate was not revoked in 2024, regardless of what happens to the
CA in 2025.

Without this data, the validator is left with an impossible question:
was this certificate valid when the signature was created? If the CA is
gone, there is no one to ask.

\subsection{You Cannot Predict CA
Longevity}\label{_you_cannot_predict_ca_longevity}

DigiNotar: bankrupt in 4 weeks. Symantec CA: sold and sunset in 18
months. Let's Encrypt: a compliance bug affecting 3 million
certificates.

When you select a CA for your document signing certificates, you are
making a bet on that CA's continued trustworthiness. History suggests
this bet is riskier than most organizations realize.

Practical mitigations:

\begin{enumerate}
\def\labelenumi{\arabic{enumi}.}
\item
  \textbf{Use PAdES-B-LTA for anything that matters.} The 5 seconds of
  extra processing time to fetch and embed revocation data is negligible
  compared to the cost of re-signing thousands of documents after a CA
  revocation.
\item
  \textbf{Have a certificate migration plan.} Know how you will
  transition to a new CA if your current CA is compromised, acquired, or
  distrusted. This includes having backup certificates from a second CA.
\item
  \textbf{Monitor your CA's compliance status.} Subscribe to CA/Browser
  Forum mailing lists, Google's security blog, and Mozilla's CA program
  communications. You want to know about problems with your CA before
  your customers do.
\item
  \textbf{Re-timestamp archived documents periodically.} A PAdES-B-LTA
  archival timestamp can be refreshed by adding a new document timestamp
  --- extending the signature's verifiability beyond the original TSA's
  certificate lifetime.
\end{enumerate}

\section{Timeline: Major CA Incidents
(2001-2024)}\label{_timeline_major_ca_incidents_2001_2024}

\begin{longtable}[]{@{}
  >{\raggedright\arraybackslash}p{(\linewidth - 4\tabcolsep) * \real{0.1429}}
  >{\raggedright\arraybackslash}p{(\linewidth - 4\tabcolsep) * \real{0.2857}}
  >{\raggedright\arraybackslash}p{(\linewidth - 4\tabcolsep) * \real{0.5714}}@{}}
\caption{Major CA security incidents and their impact on digital
signatures}\tabularnewline
\toprule\noalign{}
\begin{minipage}[b]{\linewidth}\raggedright
Year
\end{minipage} & \begin{minipage}[b]{\linewidth}\raggedright
Incident
\end{minipage} & \begin{minipage}[b]{\linewidth}\raggedright
Impact and Outcome
\end{minipage} \\
\midrule\noalign{}
\endfirsthead
\toprule\noalign{}
\begin{minipage}[b]{\linewidth}\raggedright
Year
\end{minipage} & \begin{minipage}[b]{\linewidth}\raggedright
Incident
\end{minipage} & \begin{minipage}[b]{\linewidth}\raggedright
Impact and Outcome
\end{minipage} \\
\midrule\noalign{}
\endhead
\bottomrule\noalign{}
\endlastfoot
2001 & VeriSign issues fraudulent Microsoft code-signing certs & Two
certificates issued to an impersonator claiming to be Microsoft. Revoked
quickly, but demonstrated that even the largest CA could be socially
engineered. Led to improved identity verification procedures. \\
2008 & Debian weak key vulnerability (CVE-2008-0166) & A Debian-specific
OpenSSL bug caused all keys generated on Debian/Ubuntu systems for 18
months to be predictable. Affected SSL certificates, SSH keys, and
potentially signing keys. Not a CA breach per se, but CAs had to revoke
certificates whose keys were generated on affected systems. \\
2011 (Mar) & Comodo reseller compromise & 9 fraudulent certificates
issued via compromised RA accounts. Revoked within hours. CA survived.
Attacker claimed Iranian motivation. \\
2011 (Jul-Sep) & DigiNotar breach & 531+ fraudulent certificates. Used
for mass MITM surveillance of Iranian citizens. CA bankrupt within 4
weeks. Dutch government PKI severely disrupted. All DigiNotar-chained
PDF signatures became unverifiable. \\
2012 & Trustwave subordinate CA for DLP & Trustwave issued a subordinate
CA certificate to a private company for use in DLP (Data Loss
Prevention) traffic inspection. Technically allowed the company to issue
certificates for any domain. Caused industry debate about acceptable CA
practices. \\
2013 & ANSSI/Google France intermediate cert & A French government
intermediate CA (operated by ANSSI) issued unauthorized certificates for
Google domains. Google detected via CT/pinning. Limited scope but raised
concerns about government CAs. \\
2014 & India CCA/NIC breach & India's National Informatics Centre (NIC),
operating under the Indian Controller of Certifying Authorities (CCA),
issued unauthorized certificates for Google and Yahoo domains. Google
detected via CT. Indian CCA root was constrained in trust stores. \\
2015 & CNNIC/MCS Holdings & China Internet Network Information Center
(CNNIC) delegated certificate issuance to MCS Holdings (an Egyptian
company), which issued unauthorized certificates for Google domains via
a MITM proxy. Google and Mozilla distrusted CNNIC's root
certificates. \\
2015-2018 & Symantec mis-issuance and distrust & Thousands of
certificates issued without proper authorization. Progressive distrust
by Chrome and Firefox over 2 years. Symantec CA sold to DigiCert for
\$950M. Massive impact on document signing certificates issued under
Symantec, VeriSign, Thawte, and GeoTrust brands. \\
2017 & WoSign/StartCom distrust & Chinese CA WoSign acquired StartCom
(Israeli CA) without disclosure, backdated SHA-1 certificates to avoid
browser deadlines, and made false statements to browser vendors. Both
CAs distrusted by all major browsers. \\
2020 & Let's Encrypt CAA bug revocation & 3 million certificates revoked
due to CAA checking non-compliance. 24-hour notice to certificate
holders. Demonstrated operational fragility of CA infrastructure at
scale. \\
2022 & TrustCor allegations & Investigative reporting linked TrustCor (a
Panamanian CA) to a company involved in distributing spyware. Mozilla,
Google, and Microsoft removed TrustCor from their trust stores. No
evidence of fraudulent certificate issuance, but trust was revoked on
the basis of the CA operator's associations. \\
2023 & e-Tugra suspended from CCADB & Turkish CA e-Tugra suspended from
the Common CA Database (CCADB) after security researchers discovered
critical vulnerabilities in its infrastructure, including exposed admin
panels and improperly secured HSMs. \\
\end{longtable}

This timeline only includes publicly known incidents. The actual number
of CA compromises is certainly higher --- not every breach is detected,
and not every detected breach is publicly disclosed. The incidents we
know about are the ones where detection mechanisms worked. The ones we
should worry about are the ones where they did not.

\section{Impact on PDF Signature Validation: What Happens in
Practice}\label{_impact_on_pdf_signature_validation_what_happens_in_practice}

When a CA is revoked, the practical impact on PDF signature validation
depends on several factors:

\subsection{The Validator's Trust
Store}\label{_the_validators_trust_store}

Different validators use different trust stores:

\begin{itemize}
\item
  \textbf{Adobe Acrobat (Windows)}: Uses AATL plus optionally the
  Windows certificate store
\item
  \textbf{Adobe Acrobat (macOS)}: Uses AATL plus optionally the macOS
  Keychain
\item
  \textbf{Foxit Reader}: Uses the operating system's trust store
\item
  \textbf{LibreOffice}: Uses the operating system's trust store (via NSS
  on Linux)
\item
  \textbf{EU DSS validation service}: Uses the EU Trusted List (EUTL)
\item
  \textbf{Custom Java applications}: Use the JVM's \texttt{cacerts}
  trust store or a custom \texttt{TrustStore}
\end{itemize}

When a CA is revoked, the timing and scope of the revocation varies
across these stores. AATL updates are published by Adobe on their own
schedule. OS trust store updates come via system updates. The EUTL is
updated when a member state's supervisory body modifies its trusted
list.

This means a signed PDF can show different validation results in
different viewers during the transition period --- "valid" in one that
has not yet received the revocation update, "invalid" in one that has.

\subsection{Whether LTV Data Is
Present}\label{_whether_ltv_data_is_present}

If the signature includes embedded revocation data (PAdES-B-LT or
B-LTA), a competent validator can perform \textbf{historical
validation}: it checks the embedded OCSP response or CRL to determine
the certificate's status at the time of signing, rather than its status
now.

If the signature does not include LTV data, the validator must check the
certificate's current status --- and if the CA has been revoked, the
chain cannot be validated at all.

\subsection{The Certificate's Position in the
Chain}\label{_the_certificates_position_in_the_chain}

If only an intermediate CA is revoked (as in the Comodo case), the
impact is limited to certificates issued through that specific
intermediate. Other intermediates under the same root may be unaffected.

If the root CA is revoked (as in DigiNotar's case), \emph{every}
certificate in the chain is affected --- there is no way to validate the
chain without a trusted root.

\section{Key Takeaways}\label{_key_takeaways_12}

\begin{itemize}
\item
  \textbf{CA breaches retroactively undermine PDF signatures} --- the
  math does not change, but the trust does
\item
  \textbf{DigiNotar (2011)} showed that a small CA's failure can have
  outsized impact --- 531 rogue certs, mass surveillance, bankruptcy in
  4 weeks
\item
  \textbf{Comodo (2011)} showed that fast, transparent response can save
  a CA --- 9 rogue certs, revoked in hours, CA survived
\item
  \textbf{Symantec (2015-2018)} showed that even the largest CA is not
  too big to distrust --- 30\% of TLS certificates, distrusted over 2
  years
\item
  \textbf{Let's Encrypt (2020)} showed that compliance bugs at scale
  force brutal revocation timelines --- 3 million certs, 24-hour notice
\item
  \textbf{Certificate Transparency} makes rogue certificates detectable
  but does not prevent them
\item
  \textbf{PAdES-B-LTA is the only reliable defense} --- embedded
  revocation data at signing time proves certificate validity regardless
  of what happens to the CA later
\item
  \textbf{Always have a CA migration plan} --- your CA's trustworthiness
  is not guaranteed for the lifetime of your signed documents
\item
  \textbf{Monitor your CA's status} --- you want to know about problems
  before your validators do
\item
  \textbf{Re-timestamp archived documents} periodically to extend their
  verifiability beyond any single CA's lifetime
\end{itemize}

\chapter{Oddities \& Annoyances}\label{_oddities_annoyances}

Every technology has its dark corners --- the design decisions that made
sense in 1999 but haunt us today, the interoperability failures that no
spec can fix, the political arrangements that masquerade as technical
standards. PDF digital signatures have more dark corners than most.

This chapter is a tour of the things that will make you mutter
unprintable words during implementation. Not bugs, exactly --- bugs get
fixed. These are \emph{structural annoyances}: the problems baked into
the specifications, the ecosystem, and the economics of PDF signing.
Each one has wasted thousands of engineering hours across the industry.
Understanding them will not make them go away, but it will stop you from
wasting your own hours rediscovering what everyone else already knows
and nobody bothered to write down.

\section{The Adobe Approved Trust List
Monopoly}\label{_the_adobe_approved_trust_list_monopoly}

\subsection{What the AATL Is}\label{_what_the_aatl_is}

Adobe Acrobat is the world's dominant PDF viewer. When a user opens a
signed PDF in Acrobat, the first thing they see is a blue or green
banner at the top of the document. Green means "Signed and all
signatures are valid." Blue means "At least one signature has problems."
Yellow means "Signed, but the signature cannot be verified."

The difference between green and yellow, in the vast majority of cases,
is not whether the signature is cryptographically correct. It is whether
the signing certificate chains to a root in the \textbf{Adobe Approved
Trust List (AATL)}.

The AATL is a curated list of root certificates that Adobe trusts for
document signing. It is published at
\texttt{https://helpx.adobe.com/acrobat/kb/approved-trust-list2.html}
and ships embedded in Acrobat. If your signing certificate chains to an
AATL root, you get the green checkmark. If it does not --- even if the
certificate is perfectly valid, issued by a reputable CA, chains to a
trust anchor in the operating system's trust store, and was created with
impeccable cryptographic rigor --- you get yellow.

Yellow, in the real world, means "the customer calls support."

\subsection{The Economics of Green}\label{_the_economics_of_green}

Getting your root certificate into the AATL requires meeting Adobe's
technical requirements (documented in the "AATL Technical Requirements"
PDF, which is itself signed with an AATL certificate --- recursion
intended) and paying a fee. The requirements include:

\begin{itemize}
\item
  The root must be at least 4096-bit RSA or equivalent
\item
  The CA must carry an annual WebTrust for CAs audit or equivalent (ETSI
  EN 319 411)
\item
  The CA must use HSMs for root key protection
\item
  Document signing certificates must include specific key usage
  extensions
\item
  The CA must have a published CPS (Certification Practice Statement)
\end{itemize}

These are reasonable requirements for a public trust anchor. The problem
is not the technical bar --- it is the gatekeeping.

Adobe is a private company. The AATL is Adobe's private trust list.
Adobe decides who gets in, Adobe decides who stays in, and Adobe decides
the criteria. There is no appeals process, no regulatory oversight, and
no obligation to accept any particular CA.

For CAs, inclusion in the AATL is not optional --- it is a commercial
necessity. Enterprise customers demand the green checkmark. Government
procurement specifications require it. A CA that is not in the AATL
effectively cannot sell document signing certificates at competitive
prices, because the certificates will produce yellow warnings in the
software that 80\% of users use to view PDFs.

The AATL creates a two-tier market for document signing certificates.
AATL-listed certificates cost significantly more than non-AATL
certificates --- typically \$200-\$500 per year for an individual
signing certificate, versus \$30-\$50 for a certificate from a non-AATL
CA. The price premium has nothing to do with cryptographic quality and
everything to do with the color of a banner in one company's software.

If you are building a signing service, the single most important
commercial decision you will make is whether your use case requires the
green checkmark. If your recipients will open signed PDFs in Acrobat
(and they almost certainly will), you need AATL. If your recipients use
a custom validation workflow that checks against the EU Trusted List or
your own trust store, you do not.

\subsection{The EUTL Alternative}\label{_the_eutl_alternative}

The EU Trusted List (Regulation (EU) No 910/2014 (eIDAS), Article 22) is
the European counterpart to the AATL. Every Qualified Trust Service
Provider in the EU is listed on the EUTL, and a certificate chaining to
a QTSP on the EUTL qualifies for the highest eIDAS tier regardless of
whether Adobe has heard of the CA.

But here is the catch: Acrobat does not validate against the EUTL by
default. Acrobat uses the AATL. There is an option to enable EUTL
validation in Acrobat (Edit \textgreater{} Preferences \textgreater{}
Signatures \textgreater{} Verification \textgreater{} select "European
Union Trusted Lists"), but it is off by default, and the vast majority
of users never touch it.

This means a signature with a qualified certificate from a QTSP that is
\emph{not} in the AATL will show a yellow warning in Acrobat --- even
though it is legally a qualified electronic signature with the same
standing as a handwritten signature under EU law.

The legal framework says: this is a QES, as valid as ink. The software
says: "The validity of the document certification is UNKNOWN."

The user believes the software.

If you are distributing signed PDFs to recipients who use Acrobat, and
your QTSP is on the EUTL but not the AATL, include instructions for
enabling EUTL validation in Acrobat. Better yet, embed the full
certificate chain and LTV data so that validation depends on the trust
store configuration rather than online lookups. Best of all: choose a
QTSP that is on \emph{both} the AATL and the EUTL. Most major European
QTSPs have pursued AATL membership precisely because of this problem.

\subsection{The Political Dimension}\label{_the_political_dimension}

The AATL gives Adobe significant influence over the digital signing
ecosystem. Adobe can, at its discretion, adjust requirements, add fees,
or remove CAs. This is not a theoretical concern --- Adobe has adjusted
AATL requirements multiple times, each adjustment requiring CAs to
re-certify and sometimes change their infrastructure.

Some in the industry argue that this is appropriate: someone needs to
curate trust, and Adobe has a legitimate interest in ensuring that
signatures validated by its software are trustworthy. Others argue that
a single company should not be the gatekeeper for a technology used in
legal, governmental, and financial transactions across the globe.

The emergence of the EUDI wallet (Chapter 27) may shift this dynamic. If
QES becomes a consumer technology available to every EU citizen, the
volume of qualified signatures will dwarf the current market --- and the
EUTL, not the AATL, will be the authoritative trust anchor for those
signatures. Whether Adobe adjusts its defaults accordingly remains to be
seen.

\section{Broken Validation Across
Viewers}\label{_broken_validation_across_viewers}

\subsection{The Problem}\label{_the_problem_3}

Take a signed PDF --- any signed PDF --- and open it in five different
viewers. You will get five different answers about whether the signature
is valid.

This is not an exaggeration. It is the daily reality of PDF signature
interoperability.

The fundamental problem is that ISO 32000-2:2020 defines the
\emph{structure} of a PDF signature but leaves enormous latitude in how
viewers should \emph{validate} that structure. The specification says
what a \texttt{/Sig} dictionary must contain, what the
\texttt{/ByteRange} must cover, and how the CMS \texttt{SignedData} must
be embedded. It does not say, with sufficient precision, what a viewer
should do when the trust chain is incomplete, when the timestamp is
missing, when the revocation data is stale, or when the certificate uses
an algorithm the viewer does not support.

ETSI EN 319 142 (PAdES) adds more precision, but it constrains
\emph{signing} more than it constrains \emph{validation}. A PAdES-B-LTA
signature is well-defined in terms of what it must contain. What a
validator must do with each piece of that content --- and what it should
report to the user when something is missing or ambiguous --- is still
underspecified.

\subsection{The Comparison}\label{_the_comparison}

The following table shows how five common validators handle the same
edge cases. These are real behaviors observed in testing, not
theoretical predictions.

\begin{longtable}[]{@{}
  >{\raggedright\arraybackslash}p{(\linewidth - 10\tabcolsep) * \real{0.2308}}
  >{\raggedright\arraybackslash}p{(\linewidth - 10\tabcolsep) * \real{0.1538}}
  >{\raggedright\arraybackslash}p{(\linewidth - 10\tabcolsep) * \real{0.1538}}
  >{\raggedright\arraybackslash}p{(\linewidth - 10\tabcolsep) * \real{0.1538}}
  >{\raggedright\arraybackslash}p{(\linewidth - 10\tabcolsep) * \real{0.1538}}
  >{\raggedright\arraybackslash}p{(\linewidth - 10\tabcolsep) * \real{0.1538}}@{}}
\caption{Signature validation behavior across viewers (observed, not
specified)}\tabularnewline
\toprule\noalign{}
\begin{minipage}[b]{\linewidth}\raggedright
Scenario
\end{minipage} & \begin{minipage}[b]{\linewidth}\raggedright
Acrobat DC
\end{minipage} & \begin{minipage}[b]{\linewidth}\raggedright
Foxit Reader
\end{minipage} & \begin{minipage}[b]{\linewidth}\raggedright
LibreOffice
\end{minipage} & \begin{minipage}[b]{\linewidth}\raggedright
macOS Preview
\end{minipage} & \begin{minipage}[b]{\linewidth}\raggedright
PDFBox
\end{minipage} \\
\midrule\noalign{}
\endfirsthead
\toprule\noalign{}
\begin{minipage}[b]{\linewidth}\raggedright
Scenario
\end{minipage} & \begin{minipage}[b]{\linewidth}\raggedright
Acrobat DC
\end{minipage} & \begin{minipage}[b]{\linewidth}\raggedright
Foxit Reader
\end{minipage} & \begin{minipage}[b]{\linewidth}\raggedright
LibreOffice
\end{minipage} & \begin{minipage}[b]{\linewidth}\raggedright
macOS Preview
\end{minipage} & \begin{minipage}[b]{\linewidth}\raggedright
PDFBox
\end{minipage} \\
\midrule\noalign{}
\endhead
\bottomrule\noalign{}
\endlastfoot
Valid PAdES-B-LTA, AATL cert & Green checkmark & Valid & Valid & Valid &
Valid (programmatic) \\
Valid CMS, self-signed cert & Yellow warning & Unknown signer & Valid
(if cert trusted) & Invalid & Valid (if cert in trust store) \\
Valid CMS, expired cert, no timestamp & Invalid (expired) & Invalid &
Valid at signing time & Invalid & Depends on validation date \\
Valid CMS, EUTL cert, not in AATL & Yellow (unknown signer) & Valid (if
EUTL enabled) & Valid & Unknown & Depends on trust store \\
SHA-1 digest algorithm & Warning (weak hash) & Valid (no warning) &
Valid (no warning) & Valid & Valid (no warning) \\
Missing OCSP/CRL data & Warning (revocation unknown) & Valid (ignores) &
Valid (ignores) & Valid (ignores) & No revocation check by default \\
Incremental update after signing & Modified after signing & Modified &
May not detect & Does not detect & Detects (if coded) \\
SubFilter \texttt{ETSI.CAdES.detached} & Full validation & Full
validation & Full validation & Partial & Full (if coded) \\
SubFilter \texttt{adbe.pkcs7.detached} & Full validation & Full
validation & Full validation & Full validation & Full (if coded) \\
SubFilter \texttt{adbe.x509.rsa\_sha1} & Legacy support & May fail & May
fail & Fails & Not supported \\
\end{longtable}

The most dangerous row in this table is "Incremental update after
signing." This is exactly the scenario that the Bochum Incremental
Saving Attack (Chapter 24) exploits. If a viewer does not detect that
the file was modified after signing, it will show a valid signature on a
document whose content has been changed. As of 2026, most major viewers
have patched this --- but "most" is not "all," and not every installed
version has been updated.

\subsection{Why Viewers Disagree}\label{_why_viewers_disagree}

The disagreements stem from several root causes:

\textbf{Different trust stores.} Acrobat uses AATL. Foxit uses the OS
trust store. LibreOffice uses NSS (on Linux) or the OS store (on
Windows/macOS). macOS Preview uses the Keychain. A certificate trusted
by one store may not be trusted by another.

\textbf{Different algorithm support.} Some viewers still accept SHA-1;
others warn or reject. Some support ECDSA; others are RSA-only. Some
parse CAdES extensions; others only understand basic CMS.

\textbf{Different revocation policies.} Acrobat actively checks
revocation status (OCSP and CRL) if the information is available. Foxit
and LibreOffice are more lenient --- they may skip revocation checks if
the data is not embedded. macOS Preview performs minimal revocation
checking.

\textbf{Different byte range verification.} Acrobat carefully verifies
that the byte range covers the entire file except the \texttt{/Contents}
gap. Some other viewers are less rigorous --- they check that the CMS
signature verifies over the byte range data, but do not verify that the
byte range actually covers the whole file.

\textbf{Different handling of incremental updates.} This is the critical
one. A proper validator must compare the signed byte range to the total
file size and, if there are bytes after the signed range, analyze the
incremental update to determine whether it contains permitted changes
(adding signatures, annotations) or impermissible modifications
(changing page content). Not all viewers perform this analysis with
equal rigor.

\subsection{What You Can Do About It}\label{_what_you_can_do_about_it}

If you are a signature \emph{creator}, your goal is to produce
signatures that validate consistently across all major viewers:

\begin{enumerate}
\def\labelenumi{\arabic{enumi}.}
\item
  \textbf{Use PAdES-B-LTA.} It is the most broadly supported profile.
\item
  \textbf{Use SubFilter \texttt{ETSI.CAdES.detached}.} All modern
  viewers support it, and it enables PAdES validation.
\item
  \textbf{Use SHA-256 or SHA-384.} Never SHA-1. Avoid SHA-512 only
  because some older viewers have edge-case issues with it.
\item
  \textbf{Embed the full certificate chain.} Do not rely on the viewer
  fetching missing intermediates --- some viewers will not.
\item
  \textbf{Embed OCSP responses and CRLs.} If the revocation data is in
  the DSS, viewers that check revocation will find it; viewers that do
  not check will simply ignore it.
\item
  \textbf{Allocate sufficient space for \texttt{/Contents}.} We discuss
  this in detail later in this chapter.
\item
  \textbf{Test in at least three viewers.} Acrobat, Foxit/LibreOffice,
  and one programmatic validator (PDFBox or EU DSS). If it validates in
  all three, it will validate almost everywhere.
\end{enumerate}

If you are a signature \emph{validator}, your goal is to implement the
most rigorous checks and report the most precise results:

\begin{enumerate}
\def\labelenumi{\arabic{enumi}.}
\item
  \textbf{Verify the byte range covers the entire file minus the
  \texttt{/Contents} gap.} This catches the Signature Wrapping Attack
  (Chapter 24).
\item
  \textbf{Detect incremental updates after signing.} Compare the signed
  byte range to the file size. If there is trailing data, parse it and
  determine whether it represents permitted or impermissible changes.
\item
  \textbf{Check revocation status.} Use embedded OCSP/CRL data if
  available; fall back to online checks if not; report clearly when
  revocation data is unavailable.
\item
  \textbf{Report the PAdES profile level.} Tell the user whether the
  signature is B-B, B-T, B-LT, or B-LTA --- this is much more
  informative than a binary "valid/invalid."
\end{enumerate}

\section{SubFilter Confusion}\label{_subfilter_confusion}

\subsection{The Five SubFilter Values}\label{_the_five_subfilter_values}

The \texttt{/SubFilter} key in a signature dictionary tells the viewer
what format the signature value (in \texttt{/Contents}) uses. Over the
history of PDF, five SubFilter values have been defined:

\begin{longtable}[]{@{}
  >{\raggedright\arraybackslash}p{(\linewidth - 4\tabcolsep) * \real{0.3333}}
  >{\raggedright\arraybackslash}p{(\linewidth - 4\tabcolsep) * \real{0.2222}}
  >{\raggedright\arraybackslash}p{(\linewidth - 4\tabcolsep) * \real{0.4444}}@{}}
\caption{SubFilter values defined across PDF versions and
PAdES}\tabularnewline
\toprule\noalign{}
\begin{minipage}[b]{\linewidth}\raggedright
SubFilter
\end{minipage} & \begin{minipage}[b]{\linewidth}\raggedright
Defined In
\end{minipage} & \begin{minipage}[b]{\linewidth}\raggedright
Status and Usage
\end{minipage} \\
\midrule\noalign{}
\endfirsthead
\toprule\noalign{}
\begin{minipage}[b]{\linewidth}\raggedright
SubFilter
\end{minipage} & \begin{minipage}[b]{\linewidth}\raggedright
Defined In
\end{minipage} & \begin{minipage}[b]{\linewidth}\raggedright
Status and Usage
\end{minipage} \\
\midrule\noalign{}
\endhead
\bottomrule\noalign{}
\endlastfoot
\texttt{adbe.pkcs7.detached} & PDF 1.3 (1999) & The original. CMS
\texttt{SignedData} where the signed content is the byte range data (not
included in the CMS structure). Still widely used. Supported by every
viewer. Not PAdES-compliant on its own. \\
\texttt{adbe.pkcs7.sha1} & PDF 1.4 (2001) & The SHA-1 hash of the entire
document is embedded as the CMS \texttt{EncapsulatedContentInfo}.
Deprecated in ISO 32000-2:2020. SHA-1 is broken for collision resistance
since the SHAttered attack (2017). Do not use. \\
\texttt{adbe.x509.rsa\_sha1} & PDF 1.3 (1999) & Raw RSA signature over
SHA-1 hash. The certificate is in a separate \texttt{/Cert} key, not
embedded in CMS. Extremely legacy. Most modern viewers cannot validate
it. Do not use. \\
\texttt{ETSI.CAdES.detached} & PAdES (ETSI TS 102 778, 2009) & The PAdES
SubFilter. CMS \texttt{SignedData} with CAdES extensions
(ESS-signing-certificate-v2, optionally timestamp attributes). Required
for all PAdES profiles. The recommended choice for any new
implementation. \\
\texttt{ETSI.RFC3161} & ISO 32000-2:2020 (2017) & Document-level
timestamp. The \texttt{/Contents} contains an RFC 3161
\texttt{TimeStampToken} instead of a regular CMS signature. Used for
archive timestamps in PAdES-B-LTA. Not a "signature" in the user-facing
sense --- it proves the document existed at a specific time. \\
\end{longtable}

\subsection{Which One Should You Use?}\label{_which_one_should_you_use}

If you are implementing a new signing feature in 2026, the answer is
simple: \textbf{\texttt{ETSI.CAdES.detached}} for signatures and
\textbf{\texttt{ETSI.RFC3161}} for document-level timestamps.

There is no reason to use \texttt{adbe.pkcs7.detached} in new
implementations unless you have a specific compatibility requirement
with a legacy system that cannot parse CAdES extensions. The practical
difference between \texttt{adbe.pkcs7.detached} and
\texttt{ETSI.CAdES.detached} is the presence of the
\texttt{ESS-signing-certificate-v2} signed attribute (Chapter 17) ---
adding it makes the signature PAdES-B-B compliant at zero additional
cost.

\texttt{adbe.pkcs7.sha1} is deprecated and cryptographically broken. Do
not use it. If you encounter it in a legacy document, validate it with
the understanding that SHA-1 collision resistance has been publicly
demonstrated to be insufficient since February 2017 (the SHAttered
attack by Stevens et al.).

\texttt{adbe.x509.rsa\_sha1} is a relic of PDF 1.3. It separates the
certificate from the signature in a way that no modern library handles
well. If you encounter it, you are looking at a document from the early
2000s.

\subsection{What Happens with Unrecognized SubFilter
Values}\label{_what_happens_with_unrecognized_subfilter_values}

ISO 32000-2:2020 §12.8.1 states that if a viewer encounters a SubFilter
value it does not recognize, it should not attempt to validate the
signature. In practice, viewers handle this inconsistently:

\begin{itemize}
\item
  \textbf{Acrobat}: displays "The signature format is not supported" ---
  a clear and honest message
\item
  \textbf{Foxit}: may attempt validation anyway and produce confusing
  results
\item
  \textbf{LibreOffice}: typically shows "Signature could not be
  verified" --- accurate but unhelpful
\item
  \textbf{macOS Preview}: ignores the signature entirely --- the
  document appears unsigned
\end{itemize}

If you are implementing a validator and encounter an unknown SubFilter,
the correct behavior is to report that the signature format is not
recognized. Do not attempt to guess the format. Do not silently skip the
signature. Report it.

A common source of SubFilter confusion in the wild: some signing tools
produce a SubFilter of \texttt{adbe.pkcs7.detached} but include CAdES
extensions (ESS-signing-certificate-v2) in the CMS signed attributes.
Technically, this is a valid \texttt{adbe.pkcs7.detached} signature that
\emph{also} happens to be PAdES-compliant --- because PAdES compliance
is about the content of the CMS, not just the SubFilter label. But some
validators treat the SubFilter as a strict indicator: if it says
\texttt{adbe.pkcs7.detached}, they skip CAdES-specific checks. The
result: a signature that is PAdES-compliant by content but not
recognized as such by label.

For maximum compatibility, use the SubFilter that matches your intent.
If you are producing PAdES signatures, use \texttt{ETSI.CAdES.detached}.

\section{PAdES vs ISO 32000-2:2020: The
Divergence}\label{_pades_vs_iso_32000_22020_the_divergence}

\subsection{Two Standards, One
Problem}\label{_two_standards_one_problem}

ISO 32000-2:2020 defines how to embed a digital signature in a PDF. ETSI
EN 319 142 (PAdES) defines additional requirements for that signature to
be legally meaningful in Europe.

These two standards \emph{should} be perfectly aligned. They are not.

The divergence is not dramatic --- PAdES is a strict subset of ISO
32000-2:2020, so a PAdES-compliant signature is always a valid ISO
32000-2:2020 signature. But the two standards were developed by
different organizations (ISO 32000-2:2020 by ISO TC 171, PAdES by ETSI
ESI), on different timelines, with different constituencies. The result
is a set of subtle mismatches that catch developers in the gap.

\subsection{Where They Disagree}\label{_where_they_disagree}

\textbf{Signed attributes.} ISO 32000-2:2020 requires the CMS
\texttt{SignedData} to include \texttt{content-type} and
\texttt{message-digest} signed attributes (RFC 5652 --- Cryptographic
Message Syntax (CMS) §11.1). PAdES \emph{additionally} requires
\texttt{ESS-signing-certificate-v2} (ETSI EN 319 122-1 §5.7.3). A
signature that satisfies ISO 32000-2:2020 but omits
\texttt{ESS-signing-certificate-v2} is a valid PDF signature but
\emph{not} a valid PAdES signature. This is the single most common PAdES
compliance failure.

\textbf{SubFilter semantics.} ISO 32000-2:2020 treats SubFilter as
informational --- it tells the viewer what format to expect, but the
viewer is expected to parse the CMS and determine the format from the
content. PAdES treats \texttt{ETSI.CAdES.detached} as a
\emph{requirement} --- a signature claiming PAdES compliance must use
this SubFilter. A technically identical CMS blob wrapped in
\texttt{adbe.pkcs7.detached} is not PAdES-compliant by label, even if
the CMS content is identical.

\textbf{Timestamp handling.} ISO 32000-2:2020 defines document-level
timestamps (SubFilter \texttt{ETSI.RFC3161}) and signature timestamps
(unsigned attribute \texttt{id-aa-signatureTimeStampToken} in the CMS).
PAdES defines the same mechanisms but adds requirements about timestamp
token format, the relationship between signature timestamps and document
timestamps, and the order in which they must be embedded in PAdES-B-LTA.
The PAdES requirements are stricter --- a timestamp arrangement that is
valid under ISO 32000-2:2020 may not satisfy PAdES-B-LTA's requirements
for archive timestamps.

\textbf{DSS dictionary.} The Document Security Store (DSS) dictionary
--- where certificates, OCSP responses, and CRLs are stored for LTV
validation --- is defined in ISO 32000-2:2020 §12.8.4.3 and also
referenced in PAdES. The specifications agree on the structure but
differ slightly on when DSS data is mandatory. ISO 32000-2:2020 treats
DSS as optional enrichment. PAdES-B-LT and B-LTA \emph{require} DSS data
--- without it, the signature cannot claim the LT or LTA profile level.

\subsection{The Practical Consequence}\label{_the_practical_consequence}

If you are building a signing service that needs to produce signatures
valid in both the ISO 32000-2:2020 sense and the PAdES sense (and you
almost certainly are), always target PAdES. PAdES is the stricter
standard. A PAdES-B-LTA signature satisfies both standards
simultaneously. A signature that targets only ISO 32000-2:2020 without
PAdES awareness will fail PAdES validation --- and PAdES validation is
what matters for EU compliance, for AATL-listed CAs, and for most
enterprise validation workflows.

The EU DSS validation framework (Chapter 23) validates signatures
against PAdES profiles, not against raw ISO 32000-2:2020. If your
signatures validate in EU DSS, they satisfy both standards. If you are
building your own validator, validate against PAdES --- it is the
superset of requirements that covers both.

\section{The Timestamp Chicken-and-Egg
Problem}\label{_the_timestamp_chicken_and_egg_problem}

\subsection{The Problem}\label{_the_problem_4}

A PAdES-B-T signature requires a signature timestamp --- an RFC 3161
\texttt{TimeStampToken} embedded as an unsigned CMS attribute. The
timestamp proves when the signature was created.

To obtain the timestamp, you must:

\begin{enumerate}
\def\labelenumi{\arabic{enumi}.}
\item
  Create the CMS \texttt{SignedData} (including the signed attributes
  and the signature value)
\item
  Hash the signature value
\item
  Send the hash to a Time Stamp Authority (TSA)
\item
  Receive the \texttt{TimeStampToken} from the TSA
\item
  Embed the token as an unsigned attribute in the CMS
  \texttt{SignedData}
\end{enumerate}

Here is the problem: the CMS \texttt{SignedData} is the
\texttt{/Contents} value in the signature dictionary. The
\texttt{/Contents} value is a hex string with a \emph{fixed} length ---
the space for it was allocated when the signature dictionary was written
into the PDF, before the CMS was created.

So the size of \texttt{/Contents} must be determined \emph{before} you
know how large the CMS will be. And the size of the CMS depends on
whether you include a timestamp --- and how large that timestamp is. And
you cannot obtain the timestamp until you have the CMS signature value.
And you cannot have the signature value until you have the byte range.
And the byte range depends on the size of \texttt{/Contents}.

This is the chicken-and-egg.

\subsection{The Standard Solution}\label{_the_standard_solution}

The standard solution, used by every major signing library, is to
\textbf{pre-allocate a fixed-size \texttt{/Contents} field}, create the
CMS signature, obtain the timestamp, embed the timestamp in the CMS, and
then pad the CMS to fit the pre-allocated space.

\begin{Shaded}
\begin{Highlighting}[]
\CommentTok{// Pre{-}allocate /Contents with estimated size}
\CommentTok{// The hex encoding doubles the byte count, so allocate 2x the expected CMS size}
\DataTypeTok{int}\NormalTok{ estimatedCmsSize }\OperatorTok{=} \DecValTok{16384}\OperatorTok{;}  \CommentTok{// 16 KB — usually sufficient for B{-}T}
\NormalTok{SignatureOptions sigOptions }\OperatorTok{=} \KeywordTok{new} \FunctionTok{SignatureOptions}\OperatorTok{();}
\NormalTok{sigOptions}\OperatorTok{.}\FunctionTok{setPreferredSignatureSize}\OperatorTok{(}\NormalTok{estimatedCmsSize}\OperatorTok{);}

\CommentTok{// The PDF is written with a /Contents placeholder of 2 * estimatedCmsSize hex chars}
\CommentTok{// After signing, the actual CMS (with timestamp) is padded with zeros to fill the space}
\end{Highlighting}
\end{Shaded}

The padding is simple: after the CMS \texttt{SignedData} DER encoding,
the remaining bytes in \texttt{/Contents} are filled with \texttt{0x00}.
The hex encoding of these zero bytes produces \texttt{00} characters.
The viewer knows to ignore trailing zeros after the DER-encoded CMS ---
the DER encoding is self-delimiting, so the parser can determine where
the CMS ends and the padding begins.

\subsection{The 32 KB Problem}\label{_the_32_kb_problem}

How much space should you pre-allocate? This is a more consequential
question than it appears.

The CMS \texttt{SignedData} contains:

\begin{itemize}
\item
  The DER-encoded \texttt{SignedData} structure itself
  (\textasciitilde100-300 bytes of overhead)
\item
  The signer's certificate (\textasciitilde1-2 KB)
\item
  Any intermediate certificates (\textasciitilde1-2 KB each, typically
  1-3 intermediates)
\item
  The signature value (\textasciitilde256-512 bytes for RSA-2048/4096,
  \textasciitilde64-72 bytes for ECDSA)
\item
  Signed attributes (\textasciitilde200-500 bytes)
\item
  The signature timestamp token (\textasciitilde3-8 KB, depending on the
  TSA)
\end{itemize}

A typical PAdES-B-T signature with one intermediate certificate and a
timestamp requires approximately 8-12 KB. A PAdES-B-LTA signature with
embedded OCSP responses and CRLs in the CMS (rather than in the DSS) can
require 15-30 KB.

Most implementations pre-allocate \textbf{32 KB} (32,768 bytes,
producing 65,536 hex characters in \texttt{/Contents}). This is
sufficient for the vast majority of signatures.

But not all.

\begin{longtable}[]{@{}
  >{\raggedright\arraybackslash}p{(\linewidth - 2\tabcolsep) * \real{0.3750}}
  >{\raggedright\arraybackslash}p{(\linewidth - 2\tabcolsep) * \real{0.6250}}@{}}
\caption{Scenarios that exceed 32 KB}\tabularnewline
\toprule\noalign{}
\begin{minipage}[b]{\linewidth}\raggedright
Scenario
\end{minipage} & \begin{minipage}[b]{\linewidth}\raggedright
Why It Is Large
\end{minipage} \\
\midrule\noalign{}
\endfirsthead
\toprule\noalign{}
\begin{minipage}[b]{\linewidth}\raggedright
Scenario
\end{minipage} & \begin{minipage}[b]{\linewidth}\raggedright
Why It Is Large
\end{minipage} \\
\midrule\noalign{}
\endhead
\bottomrule\noalign{}
\endlastfoot
Long certificate chain & 4-5 intermediates, each 2 KB. Happens with some
government CAs. \\
Large timestamp token & Some TSAs include their own certificate chain in
the TimeStampToken. If the TSA has 3 intermediates, the token alone can
be 10-15 KB. \\
Embedded CRLs & If a CRL is embedded directly in the CMS (as an unsigned
attribute), it can be 50-200 KB. This is unusual for new implementations
(CRLs go in the DSS instead), but some legacy signing tools do this. \\
Cross-certified chains & Government PKIs that cross-certify with
multiple roots can produce chains with redundant certificates, inflating
the CMS. \\
Multiple timestamps & Some workflows add both a signature timestamp and
a content timestamp, doubling the timestamp overhead. \\
\end{longtable}

If the CMS exceeds the pre-allocated \texttt{/Contents} size, the
signing operation fails. The failure mode depends on the library:

\begin{itemize}
\item
  \textbf{PDFBox}: throws an \texttt{IOException} with a message about
  the signature being too large for the placeholder
\item
  \textbf{iText}: throws a similar exception
\item
  \textbf{OpenPDF}: throws an exception
\end{itemize}

The fix is simple: increase the pre-allocated size. But the fix means
you allocated \emph{more} space that is filled with zero padding ---
wasted bytes in every PDF you sign. There is a tension between
allocating enough space to handle edge cases and not bloating every PDF
with unused padding.

A practical compromise: pre-allocate 32 KB for PAdES-B-T signatures and
64 KB for PAdES-B-LTA signatures. If you know your signing chain is
unusually long (government PKIs, cross-certified CAs), increase to 96 KB
or 128 KB. The wasted space is insignificant compared to the risk of a
signing failure in production.

If you want to be precise, you can estimate the CMS size before
allocating:

\begin{enumerate}
\def\labelenumi{\arabic{enumi}.}
\item
  Build a "dry run" CMS with a dummy signature value and a dummy
  timestamp of the expected size
\item
  Measure its DER-encoded length
\item
  Add a 50\% safety margin
\item
  Use that as your pre-allocation size
\end{enumerate}

This is more complex but avoids both undersizing and excessive
oversizing.

\section{Silent Signing Failures}\label{_silent_signing_failures}

\subsection{The Taxonomy of Silence}\label{_the_taxonomy_of_silence}

PDF signing can fail in ways that produce a valid-looking PDF with an
invalid signature. The document opens. The pages render. But the
signature is broken --- and the failure may not be obvious until someone
tries to validate it.

These silent failures are particularly dangerous because they often are
not detected at signing time. The signing code runs without exceptions.
The PDF is saved. It is emailed, archived, or submitted to a government
portal. And then, days or weeks later, someone discovers that the
signature is invalid.

\subsection{\texorpdfstring{Empty
\texttt{/Contents}}{Empty /Contents}}\label{_empty_contents}

The most egregious silent failure: the \texttt{/Contents} value is
present but contains only zero bytes. No CMS structure at all.

This happens when the signing pipeline fails to inject the CMS into the
\texttt{/Contents} placeholder after the byte range is computed and the
CMS is generated. The PDF has the correct structure --- the signature
dictionary exists, the byte range is set, the \texttt{/Contents} key is
there --- but the value is all zeros.

Some signing libraries have had bugs where exceptions during CMS
generation were swallowed (caught and not re-thrown), leaving the
\texttt{/Contents} empty. The resulting PDF looks signed --- it has a
signature field, possibly even a visual appearance --- but the signature
is cryptographically void.

After every signing operation, \emph{always} verify the result. Open the
output PDF, extract the \texttt{/Contents} value, and confirm it
contains a valid DER-encoded CMS \texttt{SignedData}. This adds a few
milliseconds to your signing pipeline and catches the most catastrophic
failure mode.

\begin{Shaded}
\begin{Highlighting}[]
\CommentTok{// Post{-}signing verification — always do this}
\DataTypeTok{byte}\OperatorTok{[]}\NormalTok{ contents }\OperatorTok{=} \FunctionTok{extractContentsFromSignedPdf}\OperatorTok{(}\NormalTok{signedPdfBytes}\OperatorTok{);}
\NormalTok{CMSSignedData cms }\OperatorTok{=} \KeywordTok{new} \FunctionTok{CMSSignedData}\OperatorTok{(}\NormalTok{contents}\OperatorTok{);}
\ControlFlowTok{if} \OperatorTok{(}\NormalTok{cms}\OperatorTok{.}\FunctionTok{getSignerInfos}\OperatorTok{().}\FunctionTok{getSigners}\OperatorTok{().}\FunctionTok{isEmpty}\OperatorTok{())} \OperatorTok{\{}
    \ControlFlowTok{throw} \KeywordTok{new} \BuiltInTok{IllegalStateException}\OperatorTok{(}
            \StringTok{"Signed PDF has no signers in CMS — signing failed silently"}\OperatorTok{);}
\OperatorTok{\}}
\NormalTok{log}\OperatorTok{.}\FunctionTok{info}\OperatorTok{(}\StringTok{"Post{-}signing verification passed: \{\} signer(s) in CMS"}\OperatorTok{,}
\NormalTok{        cms}\OperatorTok{.}\FunctionTok{getSignerInfos}\OperatorTok{().}\FunctionTok{getSigners}\OperatorTok{().}\FunctionTok{size}\OperatorTok{());}
\end{Highlighting}
\end{Shaded}

\subsection{Bad Byte Range}\label{_bad_byte_range}

The \texttt{/ByteRange} array must specify two ranges that together
cover the entire file except the \texttt{/Contents} hex string. A bad
byte range --- one that does not start at 0, does not end at the file's
last byte, or has gaps other than the \texttt{/Contents} region ---
produces a signature that covers only part of the file.

This can happen when:

\begin{itemize}
\item
  The signing library miscalculates the offset of the \texttt{/Contents}
  value in the file
\item
  The file is modified (even trivially --- whitespace changes, metadata
  updates) between byte range calculation and signature injection
\item
  A streaming writer flushes partial content, causing offset shifts
\end{itemize}

The PDF opens fine. The signature exists. But a proper validator will
detect that the byte range does not cover the entire file and report the
signature as invalid or the document as modified.

\subsection{Missing Certificate Chain}\label{_missing_certificate_chain}

A CMS \texttt{SignedData} should include the signer's certificate and
all intermediate certificates needed to build a chain to a trusted root.
If the intermediates are missing, the validator must fetch them --- and
may not be able to.

Missing intermediates are common in testing (developers use self-signed
certificates and forget to include the chain) but also happen in
production when:

\begin{itemize}
\item
  The signing code includes only the signer's certificate, not the
  intermediates
\item
  The PKCS\#12 keystore contains only the signer's certificate (the
  intermediates were not imported)
\item
  A library upgrade changed the default behavior for chain inclusion
\end{itemize}

The signature validates if the validator can fetch or already has the
missing intermediates. It fails if it cannot. This is a Heisenbug ---
the signature works on the developer's machine (where the intermediates
are in the trust store) but fails on the customer's machine (where they
are not).

\subsection{Expired Timestamp}\label{_expired_timestamp}

A signature timestamp proves when the signature was created. But the
timestamp itself is signed by the TSA's certificate --- and that
certificate can expire.

If the TSA certificate expires and the timestamp token does not include
sufficient chain data for historical validation, validators may reject
the timestamp. The signature itself may still verify, but the timestamp
--- which was the proof of when the signature was created --- becomes
unverifiable.

For PAdES-B-LTA, this is handled by the archive timestamp mechanism: a
document-level timestamp covers the entire file (including the DSS with
revocation data for the signature timestamp). As long as the archive
timestamp chain is valid, the signature timestamp can be verified
historically.

For PAdES-B-T signatures without LTV data, an expired TSA certificate
means the timestamp can no longer be verified. The signature degrades
from B-T to effectively B-B --- you lose the proof of when it was
created.

Monitor the expiration dates of the TSAs you use. TSA certificates
typically have 10-15 year validity periods, so this is not an urgent
problem --- but it \emph{is} a problem that will arrive eventually. When
it does, any PAdES-B-T signature without LTV data that used the expired
TSA will lose its timestamp verifiability.

This is another reason to prefer PAdES-B-LTA: the archive timestamp
mechanism provides a renewable chain that survives individual
certificate expirations.

\section{Reader-Enabled PDFs vs Signing
Conflicts}\label{_reader_enabled_pdfs_vs_signing_conflicts}

\subsection{What Reader-Enabling Is}\label{_what_reader_enabling_is}

Adobe Reader (the free version) historically could not save filled-in
forms. To allow Reader users to fill in and save a form, the PDF author
had to "enable" the PDF for Reader using Adobe Acrobat Pro. This was
done by applying a special \textbf{usage rights signature} --- a
certification-level signature with a specific set of usage rights
extensions that told Reader which features to unlock.

The usage rights signature is stored in the document's permissions
dictionary (\texttt{/Perms}) and uses a special Adobe-issued
certificate. When Reader opens a document with a valid usage rights
signature, it enables the specified features (fill-in, save, comment,
etc.).

\subsection{The Conflict with Document
Signing}\label{_the_conflict_with_document_signing}

The problem arises when a Reader-enabled PDF needs to be signed.

The usage rights signature is a certification signature --- it was the
first signature on the document, and it claims \texttt{/DocMDP}
permissions. Adding an approval signature after the certification
signature requires an incremental update, which is permitted under the
DocMDP permissions level set by the usage rights signature.

But here is the catch: some PDF editors, when they save the document
after adding a signature, perform operations that invalidate the usage
rights signature. This happens because:

\begin{enumerate}
\def\labelenumi{\arabic{enumi}.}
\item
  The editor modifies objects that the usage rights signature's byte
  range covers
\item
  The editor performs a full save instead of an incremental save,
  destroying the incremental structure that the usage rights signature
  depends on
\item
  The editor adds or modifies AcroForm entries in a way that the usage
  rights signature's \texttt{/FieldMDP} transform does not permit
\end{enumerate}

When the usage rights signature is invalidated, Reader reverts to its
restricted mode --- the user can no longer save the form. And if the
document's approval signature depends on the AcroForm structure (which
it does, since signature fields \emph{are} form fields), the interplay
between the invalidated usage rights and the new signature can produce
validation errors in Acrobat.

\subsection{The Historical Context}\label{_the_historical_context}

Reader-enabling was a significant source of revenue for Adobe: Acrobat
Pro was the only tool that could create the usage rights signature
(because it used an Adobe-controlled certificate). Third-party libraries
could not Reader-enable PDFs.

ISO 32000-2:2020 deprecated the \texttt{/Perms} usage rights mechanism.
PDF 2.0 assumes that all viewers can save forms without special
enablement. But the real world is not PDF 2.0 --- millions of
Reader-enabled PDF 1.7 documents are still in circulation, and
enterprises still generate new ones because their workflows depend on
the Reader experience.

If your signing service processes documents that may be Reader-enabled,
test the signing workflow with Reader-enabled PDFs specifically. Ensure
your signing code performs an incremental save (not a full save) and
does not modify objects covered by the usage rights signature. PDFBox's
\texttt{PDDocument.saveIncremental()} is designed for this --- but you
must be careful not to load and re-save objects that the usage rights
signature covers.

\subsection{The Modern Alternative}\label{_the_modern_alternative}

The practical alternative to Reader-enabling is to stop caring about it.
Adobe Acrobat Reader DC (the current free version) can save filled-in
forms without usage rights. The limitation existed in older versions of
Reader (pre-DC) and was removed as Adobe shifted to a subscription model
for Acrobat.

If your user base has upgraded to Reader DC (and by 2026, most have),
Reader-enabling is no longer necessary. If you must support legacy
Reader versions, the only safe approach is to Reader-enable the PDF
\emph{after} all signatures have been applied --- which is usually not
possible, since the usage rights signature must be the first signature.

The real answer is: test your specific workflow, document the behavior,
and warn your users if they are operating in a scenario where
Reader-enabling and signing conflict.

\section{\texorpdfstring{Trailing Bytes and Garbage After
\texttt{\%\%EOF}}{Trailing Bytes and Garbage After \%\%EOF}}\label{_trailing_bytes_and_garbage_after_eof}

\subsection{The Spec Says}\label{_the_spec_says}

ISO 32000-2:2020 §7.5.5 states that a PDF file must end with
\texttt{\%\%EOF}. Anything after the last \texttt{\%\%EOF} marker is not
part of the PDF.

In an incrementally saved PDF (which every signed PDF is --- signing
requires an incremental save), each revision ends with its own
\texttt{\%\%EOF}. The last \texttt{\%\%EOF} in the file marks the end of
the last revision.

\subsection{The Reality}\label{_the_reality}

In practice, many tools append data after \texttt{\%\%EOF}:

\begin{itemize}
\item
  Some PDF generators append a newline or carriage return after
  \texttt{\%\%EOF} (harmless but technically non-conformant)
\item
  Some content management systems append metadata or tracking
  information
\item
  Some signing tools, due to buffer sizing bugs, write a few extra bytes
  after the CMS injection
\item
  Some email systems append MIME boundary markers when the PDF is
  transmitted as an attachment
\end{itemize}

These trailing bytes create a problem for signature validation: the
signed byte range (as specified by \texttt{/ByteRange}) covers the file
up to and including the last byte \emph{as calculated at signing time}.
If additional bytes appear after that, the file is longer than the
signed byte range. A strict validator will flag this as "the document
was modified after signing" --- because, technically, it was, even if
the modification was just a trailing newline.

When your validator detects bytes beyond the signed byte range, examine
what those bytes actually are before reporting "modified after signing."
A few bytes of whitespace or a line break after \texttt{\%\%EOF} is
cosmetically messy but not a security concern. A complete incremental
update after the signed range is a genuine modification that needs
careful analysis.

Distinguish between "there are trailing bytes that are not an
incremental update" (cosmetic) and "there is an incremental update that
was not covered by any signature" (potentially malicious).

\section{Key Takeaways}\label{_key_takeaways_13}

\begin{itemize}
\item
  \textbf{The AATL is a private trust list} controlled by Adobe that
  determines whether Acrobat shows a green or yellow banner --- this
  gatekeeping creates a two-tier market for signing certificates
\item
  \textbf{Validation results differ across viewers} --- the same PDF can
  show valid, invalid, or unknown depending on trust stores, algorithm
  support, revocation policies, and byte range verification rigor
\item
  \textbf{Use \texttt{ETSI.CAdES.detached}} for all new signatures ---
  it is PAdES-compliant and universally supported
\item
  \textbf{PAdES is stricter than ISO 32000-2:2020} --- always target
  PAdES, because a PAdES signature satisfies both standards
\item
  \textbf{Pre-allocate 32-64 KB for \texttt{/Contents}} --- undersizing
  causes signing failures; oversizing wastes a few KB of padding
\item
  \textbf{Always verify the result after signing} --- check that
  \texttt{/Contents} contains valid CMS, the byte range is correct, and
  the chain is complete
\item
  \textbf{Monitor TSA certificate expiration} --- expired TSA certs
  degrade B-T signatures to effectively B-B
\item
  \textbf{Reader-enabled PDFs conflict with signing} --- test
  specifically if your workflow involves Reader-enabled documents
\item
  \textbf{Trailing bytes after \texttt{\%\%EOF}} are common and usually
  harmless, but your validator should distinguish cosmetic trailing data
  from actual post-signature modifications
\end{itemize}

\chapter{EUDI Wallet \& Digital
Identity}\label{_eudi_wallet_digital_identity}

For three decades, qualified electronic signatures have been a niche
technology --- expensive, complicated, and reserved for enterprises,
notaries, and government agencies. In 2024, the European Union decided
to change that. The revised eIDAS regulation puts a qualified signature
creation capability in the pocket of every EU citizen, on a device they
already own: their smartphone.

This is not incremental improvement. This is a category shift. If it
works --- and the "if" is substantial --- QES moves from a tool that a
few thousand specialized users understand to a feature that 450 million
Europeans can use as casually as they use contactless payment.

This chapter covers the regulatory, architectural, and technical
dimensions of the EUDI wallet and its impact on PDF digital signatures.
We start with the political context (because regulation drives
everything in this space), move through the architecture, and end with
what you --- the developer building signing or validation software ---
need to prepare.

\section{eIDAS 2.0: What Changed and
Why}\label{_eidas_2_0_what_changed_and_why}

\subsection{The Regulatory History}\label{_the_regulatory_history}

The original eIDAS regulation (Regulation (EU) No 910/2014 (eIDAS),
Chapter 19) entered into force in 2014 and became fully effective in
2016. It was a success in the narrow sense: it created a functioning
framework for qualified trust services, established the EU Trusted List,
and enabled cross-border recognition of QES.

It was a failure in the broad sense: almost nobody used it.

The adoption numbers tell the story. By 2022, eight years after eIDAS
entered into force, fewer than 14\% of EU citizens had ever used a
qualified electronic signature. The number who used QES regularly was a
fraction of that. In some member states, the number of active QES users
was in the low thousands --- in countries with populations of tens of
millions.

The reasons for low adoption were structural:

\begin{enumerate}
\def\labelenumi{\arabic{enumi}.}
\item
  \textbf{Cost.} A qualified signing certificate costs 50-300 EUR per
  year. Most citizens have no reason to pay that.
\item
  \textbf{Complexity.} Creating a QES required a smart card or USB
  token, driver software, and middleware that worked reliably on exactly
  zero consumer devices.
\item
  \textbf{No killer application.} There was no consumer use case
  compelling enough to justify the cost and complexity. Tax filing, the
  most common candidate, was solved by other means in most member
  states.
\item
  \textbf{eID fragmentation.} Each member state had its own electronic
  identity scheme, with its own technology stack, its own enrollment
  process, and its own level of adoption. Cross-border recognition of
  eIDs was technically possible under eIDAS but practically rare.
\end{enumerate}

The European Commission recognized by 2020 that voluntary adoption was
not going to achieve the digital single market it wanted. Something more
forceful was needed.

\subsection{Regulation (EU) 2024/1183}\label{_regulation_eu_20241183}

On April 11, 2024, the revised eIDAS regulation --- Regulation (EU)
2024/1183 (eIDAS 2.0), commonly called eIDAS 2.0 --- entered into force.
The revision keeps the trust services framework from eIDAS 1.0 largely
intact (the three-tier signature model, QTSPs, the EU Trusted List ---
all unchanged) but adds a transformative new component.

The headline change: \textbf{every EU member state must provide a
European Digital Identity Wallet to its citizens and residents}.

The timeline:

\begin{longtable}[]{@{}
  >{\raggedright\arraybackslash}p{(\linewidth - 2\tabcolsep) * \real{0.2857}}
  >{\raggedright\arraybackslash}p{(\linewidth - 2\tabcolsep) * \real{0.7143}}@{}}
\caption{eIDAS 2.0 key dates}\tabularnewline
\toprule\noalign{}
\begin{minipage}[b]{\linewidth}\raggedright
Date
\end{minipage} & \begin{minipage}[b]{\linewidth}\raggedright
Milestone
\end{minipage} \\
\midrule\noalign{}
\endfirsthead
\toprule\noalign{}
\begin{minipage}[b]{\linewidth}\raggedright
Date
\end{minipage} & \begin{minipage}[b]{\linewidth}\raggedright
Milestone
\end{minipage} \\
\midrule\noalign{}
\endhead
\bottomrule\noalign{}
\endlastfoot
April 11, 2024 & eIDAS 2.0 enters into force \\
By November 2024 & European Commission publishes the first set of
implementing acts defining technical specifications \\
2024-2025 & Large-Scale Pilots (LSPs) across the EU test wallet
implementations in real-world scenarios \\
By end of 2026 & Member states must make the EUDI wallet available to
all citizens and residents \\
2027 onwards & Mandatory acceptance of the wallet by public services
and, for certain use cases, private sector reliance parties \\
\end{longtable}

\subsection{The Political Context}\label{_the_political_context}

eIDAS 2.0 did not emerge from a vacuum. It was shaped by several
converging political forces:

\textbf{COVID-19 and the Digital COVID Certificate.} The EU Digital
COVID Certificate (2021) demonstrated that a pan-European digital
credential system could be built and deployed at scale in months. Over
1.8 billion certificates were issued across 46 countries. The technical
infrastructure --- QR codes, digital signatures, trust frameworks ---
proved that citizens would adopt digital credentials if the use case was
compelling enough. The political lesson was clear: mandates work,
voluntary adoption does not.

\textbf{Strategic autonomy.} The European Commission increasingly frames
digital identity as a matter of strategic autonomy. European citizens
authenticate to online services using Google, Apple, and Facebook
identity providers --- American companies, subject to American law,
storing European identity data on American infrastructure. The EUDI
wallet is, in part, an assertion that European identity should be
managed by European institutions.

\textbf{The digital single market.} The EU's single market depends on
cross-border transactions. Those transactions require identity
verification. If every member state has its own identity system,
cross-border identity verification remains fragmented --- exactly the
problem eIDAS 1.0 was supposed to solve and did not.

\textbf{Article 45 and the browser wars.} The original eIDAS 2.0
proposal included a controversial provision (Article 45) that would have
required web browsers to trust Qualified Website Authentication
Certificates (QWACs) issued by QTSPs, without the ability to apply
additional security requirements. Browser vendors --- particularly
Mozilla and Google --- objected strenuously, arguing that this would
undermine browser security by forcing them to trust certificates that
did not meet their own security standards. The final text was amended to
address some of these concerns, but the episode illustrates the tension
between regulatory mandates and technical security communities.

\section{The EUDI Wallet Concept}\label{_the_eudi_wallet_concept}

\subsection{What the Wallet Is}\label{_what_the_wallet_is}

The EUDI wallet is a smartphone application --- provided by or on behalf
of a member state --- that allows citizens and residents to:

\begin{enumerate}
\def\labelenumi{\arabic{enumi}.}
\item
  \textbf{Store government-issued identity data (PID):} name, date of
  birth, nationality, and a unique persistent identifier --- the digital
  equivalent of a national ID card or passport
\item
  \textbf{Store qualified electronic attestations of attributes
  (QEAAs):} digitally signed attestations of specific attributes ---
  university degrees, professional qualifications, driving licenses, age
  verification, right to work
\item
  \textbf{Present credentials selectively:} share only the specific
  attributes required for a transaction (e.g., "I am over 18" without
  revealing date of birth or name)
\item
  \textbf{Create qualified electronic signatures:} sign documents with
  QES, using a remote QSCD accessed through the wallet
\end{enumerate}

The wallet is \emph{not} a single application developed by the European
Commission. Each member state is responsible for providing its own
wallet implementation. The Commission defines the technical
specifications and ensures interoperability; member states build or
procure the implementations.

\subsection{What the Wallet Is Not}\label{_what_the_wallet_is_not}

The EUDI wallet is not a blockchain. It is not decentralized in the
cryptocurrency sense. It is not self-sovereign identity (SSI), although
it borrows some concepts from that space.

The wallet is a \emph{government-issued} credential system. The identity
data comes from authoritative government sources (civil registries,
passport offices). The trust model is top-down: you trust the wallet
because you trust the member state that issued it, and other member
states trust it because the EU framework mandates mutual recognition.

The wallet is also not mandatory for citizens. eIDAS 2.0 requires member
states to \emph{offer} the wallet; it does not require citizens to
\emph{use} it. However, Article 5b requires that a wide range of public
and private sector entities accept the wallet --- which creates strong
incentives for adoption even without an explicit mandate.

The voluntary nature of the wallet is politically important. A mandatory
government identity app on every citizen's phone would face significant
political resistance, particularly in countries with strong privacy
traditions (Germany, Austria, the Netherlands). By making the wallet
optional for citizens but mandatory for reliance parties, eIDAS 2.0
creates adoption through demand-side pressure rather than supply-side
compulsion.

\section{Architecture Reference Framework
(ARF)}\label{_architecture_reference_framework_arf}

\subsection{Overview}\label{_overview}

The Architecture Reference Framework (ARF) is the technical blueprint
for the EUDI wallet ecosystem. Published and maintained by the European
Commission, the ARF defines the architecture, protocols, data models,
and trust framework for the wallet.

As of early 2026, the ARF is at version 1.4 and still evolving. The core
architecture is stable, but specific protocol choices and data formats
are being refined based on pilot feedback.

The ARF defines four main roles:

\begin{longtable}[]{@{}
  >{\raggedright\arraybackslash}p{(\linewidth - 2\tabcolsep) * \real{0.2857}}
  >{\raggedright\arraybackslash}p{(\linewidth - 2\tabcolsep) * \real{0.7143}}@{}}
\caption{EUDI wallet ecosystem roles}\tabularnewline
\toprule\noalign{}
\begin{minipage}[b]{\linewidth}\raggedright
Role
\end{minipage} & \begin{minipage}[b]{\linewidth}\raggedright
Description
\end{minipage} \\
\midrule\noalign{}
\endfirsthead
\toprule\noalign{}
\begin{minipage}[b]{\linewidth}\raggedright
Role
\end{minipage} & \begin{minipage}[b]{\linewidth}\raggedright
Description
\end{minipage} \\
\midrule\noalign{}
\endhead
\bottomrule\noalign{}
\endlastfoot
\textbf{Wallet User} & The citizen or resident who holds the wallet on
their device \\
\textbf{Wallet Provider} & The member state (or entity acting on its
behalf) that issues and manages the wallet application \\
\textbf{PID Provider / Attestation Provider} & The authoritative source
that issues identity data (PID) or attribute attestations (QEAAs) into
the wallet \\
\textbf{Relying Party} & The entity (public or private sector) that
requests and validates credentials from the wallet \\
\end{longtable}

\subsection{The Trust Model}\label{_the_trust_model}

The EUDI wallet trust model has multiple layers:

\begin{enumerate}
\def\labelenumi{\arabic{enumi}.}
\item
  \textbf{Device trust.} The wallet runs on the user's smartphone. The
  device's secure element or Trusted Execution Environment (TEE)
  provides hardware-backed key storage. The wallet checks that the
  device meets minimum security requirements (OS version, no
  root/jailbreak).
\item
  \textbf{Wallet trust.} Each wallet implementation is certified by the
  member state that issues it. The certification confirms that the
  wallet meets the ARF's security and functional requirements. Wallet
  instances are registered with the member state's infrastructure.
\item
  \textbf{PID/QEAA trust.} PID and QEAAs are digitally signed by the
  issuing authority (the PID Provider for identity data, the Attestation
  Provider for attributes). The signature uses a key that chains to the
  member state's entry in the EU Trusted List. Relying parties validate
  the attestation signatures against the EUTL, just as PDF signature
  validators check the signer's certificate against a trust store.
\item
  \textbf{Relying party trust.} Relying parties that request credentials
  from the wallet must be registered and authenticated. The wallet
  verifies the relying party's identity before releasing any data ---
  the user knows \emph{who} is requesting their data, and the wallet
  ensures the request is legitimate.
\end{enumerate}

The trust model for the EUDI wallet is fundamentally the same trust
model as PDF digital signatures: X.509 certificates, certificate chains,
trusted lists, and CMS/COSE signatures. If you understand how PDF
signature validation works (Chapters 16-17), you understand the core of
the wallet's trust model. The wallet adds selective disclosure and user
consent on top, but the cryptographic foundations are identical.

\section{PID and QEAA}\label{_pid_and_qeaa}

\subsection{Person Identification Data
(PID)}\label{_person_identification_data_pid}

The PID is the wallet's core credential --- the digital equivalent of a
national ID card. It is issued by the member state's PID Provider
(typically the civil registry or passport authority) and contains a
minimum set of identity attributes.

\begin{longtable}[]{@{}
  >{\raggedright\arraybackslash}p{(\linewidth - 2\tabcolsep) * \real{0.4286}}
  >{\raggedright\arraybackslash}p{(\linewidth - 2\tabcolsep) * \real{0.5714}}@{}}
\caption{Mandatory PID attributes (ARF v1.4)}\tabularnewline
\toprule\noalign{}
\begin{minipage}[b]{\linewidth}\raggedright
Attribute
\end{minipage} & \begin{minipage}[b]{\linewidth}\raggedright
Description
\end{minipage} \\
\midrule\noalign{}
\endfirsthead
\toprule\noalign{}
\begin{minipage}[b]{\linewidth}\raggedright
Attribute
\end{minipage} & \begin{minipage}[b]{\linewidth}\raggedright
Description
\end{minipage} \\
\midrule\noalign{}
\endhead
\bottomrule\noalign{}
\endlastfoot
\texttt{family\_name} & Current family name(s) \\
\texttt{given\_name} & Current given name(s) \\
\texttt{birth\_date} & Date of birth \\
\texttt{unique\_id} & A persistent, unique identifier assigned by the
member state \\
\texttt{age\_over\_18} & Boolean: is the holder over 18? (Derived
attribute for age verification without revealing birth date) \\
\texttt{issuance\_date} & When the PID was issued \\
\texttt{expiry\_date} & When the PID expires \\
\texttt{issuing\_authority} & The authority that issued the PID \\
\texttt{issuing\_country} & The member state that issued the PID (ISO
3166-1 alpha-2) \\
\end{longtable}

The PID is signed by the PID Provider using a key that chains to the
member state's trust infrastructure. The signature format is currently
specified as either \textbf{ISO/IEC 18013-5 mdoc} (the same format used
for mobile driving licenses) or \textbf{SD-JWT} (Selective Disclosure
JSON Web Token, IETF draft). The ARF supports both formats, with ongoing
discussion about whether to mandate one or allow both.

\subsection{Qualified Electronic Attestations of Attributes
(QEAA)}\label{_qualified_electronic_attestations_of_attributes_qeaa}

QEAAs are the wallet's extensibility mechanism. While the PID provides
core identity, QEAAs attest to specific attributes beyond identity:

\begin{itemize}
\item
  \textbf{Educational credentials:} university degrees, professional
  certifications
\item
  \textbf{Professional qualifications:} medical licenses, legal bar
  admissions, engineering certifications
\item
  \textbf{Organizational identity:} employment status, company role,
  powers of representation
\item
  \textbf{Financial attributes:} bank account verification,
  creditworthiness attestations
\item
  \textbf{Travel documents:} digital driving licenses, digital travel
  credentials
\end{itemize}

A QEAA is issued by a Qualified Trust Service Provider (QTSP) or a
government authority. When issued by a QTSP, it carries the same legal
weight as the paper credential it replaces --- a QEAA attesting a
university degree, issued by a QTSP after verification with the
university, is legally equivalent to the paper diploma.

The "qualified" designation is critical. An Electronic Attestation of
Attributes (EAA) without the "qualified" prefix does not enjoy the same
legal presumptions. The distinction mirrors the QES/AdES distinction:
qualified attestations carry a legal presumption of correctness;
non-qualified attestations are evidence to be evaluated on their merits.

\section{Selective Disclosure and
SD-JWT}\label{_selective_disclosure_and_sd_jwt}

\subsection{The Privacy Problem}\label{_the_privacy_problem}

Traditional digital credentials have a privacy problem: they are
all-or-nothing. If your credential contains your name, date of birth,
address, and nationality, presenting it to verify your age requires
revealing all four attributes --- even though the relying party only
needs to know that you are over 18.

The EUDI wallet solves this through \textbf{selective disclosure} ---
the ability to present only a subset of the attributes in a credential,
without invalidating the cryptographic signature.

\subsection{How SD-JWT Works}\label{_how_sd_jwt_works}

The ARF specifies \textbf{SD-JWT (Selective Disclosure for JWTs)} as one
of the credential formats for selective disclosure. SD-JWT is an IETF
specification (currently in the RFC track) that extends JSON Web Tokens
with selective disclosure capability.

The mechanism works as follows:

\begin{enumerate}
\def\labelenumi{\arabic{enumi}.}
\item
  \textbf{Issuance.} The issuer creates a JWT containing all attributes,
  but each disclosable attribute is replaced with a hash of the
  attribute value combined with a random salt. The original attribute
  values and their salts are provided to the wallet holder as
  "disclosures" --- separate, individually presentable data items.
\item
  \textbf{Presentation.} When the wallet presents the credential to a
  relying party, it includes the JWT (with the hashed placeholders) plus
  only the disclosures for the attributes the relying party needs. The
  relying party can verify that each disclosed attribute matches its
  hash in the JWT, confirming it was part of the original signed
  credential --- but attributes whose disclosures are withheld remain
  hidden.
\item
  \textbf{Verification.} The relying party verifies the JWT signature
  (confirming it was issued by a trusted authority), then verifies each
  disclosed attribute against its hash in the JWT. Attributes that were
  not disclosed cannot be reconstructed from the hashes (the random
  salts ensure this).
\end{enumerate}

\begin{Shaded}
\begin{Highlighting}[]
\ErrorTok{//} \ErrorTok{Simplified} \ErrorTok{SD{-}JWT} \ErrorTok{structure} \ErrorTok{(conceptual,} \ErrorTok{not} \ErrorTok{a} \ErrorTok{real} \ErrorTok{token)}
\FunctionTok{\{}
  \DataTypeTok{"iss"}\FunctionTok{:} \StringTok{"https://pid.example{-}country.eu"}\FunctionTok{,}
  \DataTypeTok{"iat"}\FunctionTok{:} \DecValTok{1709913600}\FunctionTok{,}
  \DataTypeTok{"exp"}\FunctionTok{:} \DecValTok{1741449600}\FunctionTok{,}
  \DataTypeTok{"cnf"}\FunctionTok{:} \FunctionTok{\{} \DataTypeTok{"jwk"}\FunctionTok{:} \FunctionTok{\{} \ErrorTok{...} \FunctionTok{\}} \FunctionTok{\},}  \ErrorTok{//} \ErrorTok{Holder} \ErrorTok{binding} \ErrorTok{key}

  \ErrorTok{//} \ErrorTok{Selectively} \ErrorTok{disclosable} \ErrorTok{claims} \ErrorTok{(hashed} \ErrorTok{placeholders)}
  \DataTypeTok{"\_sd"}\FunctionTok{:} \OtherTok{[}
    \StringTok{"H3JQXF...base64url\_hash\_of\_family\_name\_disclosure"}\OtherTok{,}
    \StringTok{"K9PLWT...base64url\_hash\_of\_given\_name\_disclosure"}\OtherTok{,}
    \StringTok{"M2VRZN...base64url\_hash\_of\_birth\_date\_disclosure"}
  \OtherTok{]}\FunctionTok{,}

  \ErrorTok{//} \ErrorTok{Non{-}disclosable} \ErrorTok{claims} \ErrorTok{(always} \ErrorTok{visible)}
  \DataTypeTok{"age\_over\_18"}\FunctionTok{:} \KeywordTok{true}\FunctionTok{,}
  \DataTypeTok{"issuing\_country"}\FunctionTok{:} \StringTok{"DE"}
\FunctionTok{\}}
\end{Highlighting}
\end{Shaded}

\subsection{The ISO/IEC 18013-5 mdoc
Alternative}\label{_the_isoiec_18013_5_mdoc_alternative}

The ARF also supports the \textbf{mdoc} format defined in ISO/IEC
18013-5 (originally for mobile driving licenses). mdoc uses CBOR
(Concise Binary Object Representation) instead of JSON and COSE (CBOR
Object Signing and Encryption) instead of JOSE.

Both formats support selective disclosure. The choice between them is
still under discussion in the ARF, with mdoc favored for proximity
(NFC/BLE) use cases and SD-JWT favored for remote (online) use cases.

For PDF signing purposes, the credential format does not directly matter
--- what matters is that the wallet can authenticate the user and
authorize a signing operation. The credential format affects how the
wallet proves the user's identity to the signing service, not how the
PDF signature itself is structured.

If you are building a service that accepts signed PDFs and also needs to
verify the signer's attributes (e.g., a medical prescription service
that needs to confirm the signer is a licensed doctor), you will need to
handle QEAA validation \emph{in addition to} PDF signature validation.
The QEAA is presented separately from the PDF --- it is not embedded in
the signature's CMS. Your validation workflow will have two parallel
tracks: verify the PDF signature, and verify the QEAA.

\section{Signing with the Wallet}\label{_signing_with_the_wallet}

\subsection{The Signing Flow}\label{_the_signing_flow}

The most consequential feature of the EUDI wallet for this book's
audience is its ability to create qualified electronic signatures. eIDAS
2.0 Article 5a(4)(d) requires that the wallet support QES creation.
Article 45e goes further: member states must ensure that QES creation
through the wallet is offered to natural persons free of charge for
non-professional purposes.

Free QES for every EU citizen. That sentence alone justifies this
chapter.

The signing flow works as follows:

\begin{enumerate}
\def\labelenumi{\arabic{enumi}.}
\item
  \textbf{The user initiates a signing request.} This could be from
  within the wallet app itself, or from an external application (a web
  portal, a desktop application, a government service) that redirects to
  the wallet via a protocol defined in the ARF (OpenID4VP --- OpenID for
  Verifiable Presentations).
\item
  \textbf{The wallet authenticates the user.} The user confirms their
  identity using the wallet's authentication mechanism --- typically
  biometric (fingerprint or face recognition) or device PIN. This is the
  "sole control" requirement from eIDAS Article 26(c) and Article 29.
\item
  \textbf{The wallet presents the document hash to a remote QSCD.} The
  wallet does \emph{not} contain the private signing key on the device
  itself. Instead, it communicates with a \textbf{remote QSCD} operated
  by a QTSP. The wallet sends a hash of the document to be signed (not
  the document itself --- the hash is small and privacy-preserving).
\item
  \textbf{The remote QSCD creates the signature.} The QSCD --- an HSM in
  the QTSP's data center, certified as a Qualified Signature Creation
  Device --- signs the hash with the user's private key. The private key
  was generated on and never leaves the QSCD.
\item
  \textbf{The signature is returned to the requesting application.} The
  CMS \texttt{SignedData} (or the raw signature value, depending on the
  integration pattern) is returned to the application that initiated the
  signing request. The application embeds the signature in the PDF.
\end{enumerate}

\begin{longtable}[]{@{}
  >{\raggedright\arraybackslash}p{(\linewidth - 2\tabcolsep) * \real{0.1667}}
  >{\raggedright\arraybackslash}p{(\linewidth - 2\tabcolsep) * \real{0.8333}}@{}}
\caption{EUDI wallet signing flow}\tabularnewline
\toprule\noalign{}
\begin{minipage}[b]{\linewidth}\raggedright
Step
\end{minipage} & \begin{minipage}[b]{\linewidth}\raggedright
Description
\end{minipage} \\
\midrule\noalign{}
\endfirsthead
\toprule\noalign{}
\begin{minipage}[b]{\linewidth}\raggedright
Step
\end{minipage} & \begin{minipage}[b]{\linewidth}\raggedright
Description
\end{minipage} \\
\midrule\noalign{}
\endhead
\bottomrule\noalign{}
\endlastfoot
1 & Application presents document to user and requests signature (via
OpenID4VP) \\
2 & Wallet displays document summary and requests user authentication
(biometric/PIN) \\
3 & Wallet sends document hash to remote QSCD at QTSP (over mTLS,
authenticated by wallet instance credentials) \\
4 & QSCD signs the hash with user's private key; returns CMS
\texttt{SignedData} with qualified certificate and signature
timestamp \\
5 & Application receives CMS, embeds it in PDF \texttt{/Contents}, adds
DSS data for LTV, and saves the signed PDF \\
\end{longtable}

\subsection{Remote QSCD: Why the Key Is Not on the
Phone}\label{_remote_qscd_why_the_key_is_not_on_the_phone}

A natural question: why not put the signing key on the phone itself?
Modern smartphones have secure elements and TEEs that can generate and
store cryptographic keys.

The answer is QSCD certification.

A QSCD must be certified under Common Criteria EAL4+ (or equivalent)
against specific Protection Profiles (EN 419 211 or EN 419 221).
Smartphone secure elements are certified for payment (EMVCo) and some
for FIDO authentication, but as of 2026, no smartphone secure element
has been certified as a QSCD under the eIDAS Protection Profiles.

The certification gap is not primarily technical --- modern secure
elements are arguably as secure as certified smart cards. The gap is
procedural: QSCD certification requires extensive evaluation against
specific attack models, documentation of the manufacturing process, and
ongoing compliance monitoring. Smartphone vendors have not pursued this
certification because, until eIDAS 2.0, there was no market demand for
it on consumer devices.

This may change. The ARF explicitly contemplates the possibility of
"local QSCD" on a device --- a future scenario where smartphone secure
elements achieve QSCD certification and the signing key can reside on
the phone. But for the initial wallet deployments, the architecture is
remote QSCD.

Remote QSCD is not a compromise --- it is the dominant model for QES
today. Even before the EUDI wallet, most QES in the EU were created
using remote QSCDs. Services like DocuSign (through its European
qualified signing infrastructure), Swisscom, InfoCert, and D-Trust all
operate remote QSCDs. The wallet extends this existing model to a new
client (the smartphone) but does not fundamentally change the
architecture.

The Cloud Signature Consortium (CSC) API --- the industry-standard REST
API for remote signing --- is being adapted for wallet integration. If
your application already supports CSC-based remote signing (Chapter 22),
adapting it for wallet-initiated signing is a relatively small change.

\subsection{The Sole Control Assurance
Problem}\label{_the_sole_control_assurance_problem}

eIDAS requires that the signing key be under the signatory's "sole
control" (Article 26(c)). When the key is on a smart card in the
signer's wallet, sole control is straightforward: the signer possesses
the card and knows the PIN.

When the key is in a remote QSCD, sole control is more complex. The
signer does not physically possess the key. The key is in an HSM in a
data center, managed by a QTSP. Sole control must be established through
authentication mechanisms that ensure only the authorized signer can
trigger a signing operation with that key.

The SCAL (Sole Control Assurance Level) framework, defined in CEN EN 419
241-1, distinguishes two levels:

\begin{itemize}
\item
  \textbf{SCAL1}: The QTSP relies on an external authentication
  mechanism (e.g., the wallet app's biometric authentication) to
  authorize signing. The QTSP trusts the authentication assertion from
  the wallet.
\item
  \textbf{SCAL2}: The QTSP independently verifies the signer's
  authorization --- the signer must authenticate directly to the QTSP's
  infrastructure (e.g., through a separate OTP sent by the QTSP), in
  addition to any wallet-level authentication.
\end{itemize}

SCAL2 is required for QES. This means the remote QSCD must have its own
authentication of the signer, not just trust the wallet's claim that the
user authenticated.

In practice, this is implemented through a combination of: -
Wallet-level biometric authentication (proves the user is holding the
phone) - A cryptographic challenge from the QSCD that the wallet's
secure element signs (proves the specific wallet instance is authorized)
- An additional PIN or biometric confirmation presented to the QSCD's
authentication service (proves the user explicitly authorized \emph{this
specific signing operation})

The user experience goal is to make this feel like a single
authentication step (touch fingerprint, confirm signing), while the
underlying protocol satisfies SCAL2's multi-factor requirement.

\section{Impact on PDF Signing}\label{_impact_on_pdf_signing}

\subsection{What Changes}\label{_what_changes}

The EUDI wallet does not change the PDF signature format. A PDF signed
through the wallet contains the same structures as any other PAdES
signature: a \texttt{/Sig} dictionary with \texttt{/SubFilter}
\texttt{ETSI.CAdES.detached}, a \texttt{/ByteRange} covering the file, a
\texttt{/Contents} value containing CMS \texttt{SignedData}, and (for
PAdES-B-LTA) a DSS dictionary with certificates and revocation data.

What changes is the \emph{volume}, the \emph{source}, and the
\emph{validation path}.

\textbf{Volume.} Free QES for 450 million EU citizens will produce
orders of magnitude more qualified signatures than the current market.
Your validation code must handle this volume. If your validation
involves synchronous EUTL lookups for every signature, you need caching.
If your trust store is updated manually, you need automation.

\textbf{Source.} Wallet-generated signatures will come from certificates
issued by QTSPs participating in the wallet ecosystem. These QTSPs will
be on the EUTL. They may or may not be on Adobe's AATL. If they are not
on the AATL, Acrobat users will see yellow warnings on legitimate QES
--- which will create pressure on Adobe to either auto-consume the EUTL
or expand the AATL.

\textbf{Validation path.} Wallet-generated signatures chain through a
trust path that includes: the signer's qualified certificate (issued by
a QTSP) → the QTSP's intermediate/root certificate → the EUTL entry for
the QTSP. This is the same chain as any other QES. The novelty is that
the signer is a consumer on a smartphone, not an enterprise user with a
signing portal.

\subsection{PAdES-B-LTA from Mobile}\label{_pades_b_lta_from_mobile}

The wallet ecosystem targets PAdES-B-LTA as the default profile for
wallet-generated signatures. This is the correct engineering decision:
wallet-generated signatures will be used for contracts, government
submissions, and official documents that require long-term
verifiability.

Producing a PAdES-B-LTA signature from a mobile device requires:

\begin{enumerate}
\def\labelenumi{\arabic{enumi}.}
\item
  \textbf{Signature creation} (CMS \texttt{SignedData} with
  \texttt{ESS-signing-certificate-v2}, via remote QSCD)
\item
  \textbf{Signature timestamp} (RFC 3161 \texttt{TimeStampToken} from a
  TSA, embedded as unsigned CMS attribute)
\item
  \textbf{DSS population} (certificates, OCSP responses, CRLs for the
  entire chain)
\item
  \textbf{Document timestamp} (a separate signature dictionary with
  SubFilter \texttt{ETSI.RFC3161} covering the entire file including
  DSS)
\end{enumerate}

Steps 1-2 are performed by the QSCD/QTSP infrastructure. Steps 3-4 are
performed by the application that embeds the signature in the PDF ---
which may be a server-side service, not the mobile device itself.

The typical architecture:

\begin{enumerate}
\def\labelenumi{\arabic{enumi}.}
\item
  The mobile app sends the document hash to the QTSP via the wallet
\item
  The QTSP signs the hash on the QSCD and obtains a signature timestamp
  from its TSA
\item
  The CMS (with timestamp) is returned to the application server
\item
  The server embeds the CMS in the PDF, populates the DSS, and adds a
  document timestamp
\item
  The signed PDF is returned to the user or submitted to the relying
  party
\end{enumerate}

The mobile device never needs to perform the PDF manipulation itself ---
that is handled by the server. This is a practical necessity: PDF
signing with LTV data is computationally modest but requires PDF parsing
libraries that are not available on all mobile platforms.

\begin{Shaded}
\begin{Highlighting}[]
\CommentTok{/**}
 \CommentTok{*}\NormalTok{ Server}\CommentTok{{-}}\NormalTok{side method that completes a wallet}\CommentTok{{-}}\NormalTok{initiated signing operation}\CommentTok{.}
 \CommentTok{*}\NormalTok{ The wallet has obtained a CMS SignedData from the remote QSCD}\CommentTok{;}
 \CommentTok{*}\NormalTok{ this method embeds it in the PDF and adds LTV data}\CommentTok{.}
 \CommentTok{*/}
\KeywordTok{public} \DataTypeTok{byte}\OperatorTok{[]} \FunctionTok{completeWalletSignature}\OperatorTok{(}\DataTypeTok{byte}\OperatorTok{[]}\NormalTok{ pdfBytes}\OperatorTok{,}
                                       \DataTypeTok{byte}\OperatorTok{[]}\NormalTok{ cmsSignedData}\OperatorTok{,}
                                       \BuiltInTok{List}\OperatorTok{\textless{}}\BuiltInTok{X509Certificate}\OperatorTok{\textgreater{}}\NormalTok{ chain}\OperatorTok{)} \KeywordTok{throws} \BuiltInTok{Exception} \OperatorTok{\{}

\NormalTok{    PDDocument doc }\OperatorTok{=}\NormalTok{ Loader}\OperatorTok{.}\FunctionTok{loadPDF}\OperatorTok{(}\NormalTok{pdfBytes}\OperatorTok{);}

    \CommentTok{// Create signature and set CMS from wallet/QSCD}
\NormalTok{    PDSignature signature }\OperatorTok{=} \KeywordTok{new} \FunctionTok{PDSignature}\OperatorTok{();}
\NormalTok{    signature}\OperatorTok{.}\FunctionTok{setFilter}\OperatorTok{(}\NormalTok{PDSignature}\OperatorTok{.}\FunctionTok{FILTER\_ADOBE\_PPKLITE}\OperatorTok{);}
\NormalTok{    signature}\OperatorTok{.}\FunctionTok{setSubFilter}\OperatorTok{(}\NormalTok{PDSignature}\OperatorTok{.}\FunctionTok{SUBFILTER\_ETSI\_CADES\_DETACHED}\OperatorTok{);}
\NormalTok{    signature}\OperatorTok{.}\FunctionTok{setSignDate}\OperatorTok{(}\BuiltInTok{Calendar}\OperatorTok{.}\FunctionTok{getInstance}\OperatorTok{());}

    \CommentTok{// Add signature with pre{-}built CMS}
\NormalTok{    doc}\OperatorTok{.}\FunctionTok{addSignature}\OperatorTok{(}\NormalTok{signature}\OperatorTok{,}\NormalTok{ content }\OperatorTok{{-}\textgreater{}}\NormalTok{ cmsSignedData}\OperatorTok{);}

    \CommentTok{// Save incrementally}
    \BuiltInTok{ByteArrayOutputStream}\NormalTok{ baos }\OperatorTok{=} \KeywordTok{new} \BuiltInTok{ByteArrayOutputStream}\OperatorTok{();}
\NormalTok{    doc}\OperatorTok{.}\FunctionTok{saveIncremental}\OperatorTok{(}\NormalTok{baos}\OperatorTok{);}
    \DataTypeTok{byte}\OperatorTok{[]}\NormalTok{ signedPdf }\OperatorTok{=}\NormalTok{ baos}\OperatorTok{.}\FunctionTok{toByteArray}\OperatorTok{();}
\NormalTok{    doc}\OperatorTok{.}\FunctionTok{close}\OperatorTok{();}

    \CommentTok{// Add DSS with LTV data (separate incremental save)}
\NormalTok{    signedPdf }\OperatorTok{=} \FunctionTok{addDssWithLtvData}\OperatorTok{(}\NormalTok{signedPdf}\OperatorTok{,}\NormalTok{ chain}\OperatorTok{);}

    \CommentTok{// Add document timestamp (another incremental save)}
\NormalTok{    signedPdf }\OperatorTok{=} \FunctionTok{addDocumentTimestamp}\OperatorTok{(}\NormalTok{signedPdf}\OperatorTok{);}

\NormalTok{    log}\OperatorTok{.}\FunctionTok{info}\OperatorTok{(}\StringTok{"Wallet signature completed: PAdES{-}B{-}LTA, \{\} bytes"}\OperatorTok{,}\NormalTok{ signedPdf}\OperatorTok{.}\FunctionTok{length}\OperatorTok{);}
    \ControlFlowTok{return}\NormalTok{ signedPdf}\OperatorTok{;}
\OperatorTok{\}}
\end{Highlighting}
\end{Shaded}

The code above is simplified for illustration. A production
implementation must handle: pre-allocating the \texttt{/Contents}
placeholder, coordinating the byte range computation with the CMS
injection, error handling for TSA timeouts, retry logic for remote QSCD
communication failures, and verification of the completed signature.
Chapter 29 covers the full PDFBox signing pipeline.

\section{Pilot Programs}\label{_pilot_programs}

\subsection{The Large-Scale Pilots}\label{_the_large_scale_pilots}

The European Commission is funding four Large-Scale Pilots (LSPs) that
test the EUDI wallet in real-world scenarios across multiple member
states. Each pilot covers different use cases and involves different
member states.

\begin{longtable}[]{@{}
  >{\raggedright\arraybackslash}p{(\linewidth - 4\tabcolsep) * \real{0.2222}}
  >{\raggedright\arraybackslash}p{(\linewidth - 4\tabcolsep) * \real{0.3333}}
  >{\raggedright\arraybackslash}p{(\linewidth - 4\tabcolsep) * \real{0.4444}}@{}}
\caption{EUDI wallet Large-Scale Pilots (as of early
2026)}\tabularnewline
\toprule\noalign{}
\begin{minipage}[b]{\linewidth}\raggedright
Pilot
\end{minipage} & \begin{minipage}[b]{\linewidth}\raggedright
Lead Countries
\end{minipage} & \begin{minipage}[b]{\linewidth}\raggedright
Primary Use Cases
\end{minipage} \\
\midrule\noalign{}
\endfirsthead
\toprule\noalign{}
\begin{minipage}[b]{\linewidth}\raggedright
Pilot
\end{minipage} & \begin{minipage}[b]{\linewidth}\raggedright
Lead Countries
\end{minipage} & \begin{minipage}[b]{\linewidth}\raggedright
Primary Use Cases
\end{minipage} \\
\midrule\noalign{}
\endhead
\bottomrule\noalign{}
\endlastfoot
\textbf{POTENTIAL} & France, Germany, Italy, Spain & Government services
(eID authentication), banking (customer onboarding), mobile driving
license, SIM registration, QES signing \\
\textbf{NOBID} & Norway, Denmark, Sweden, Finland, Iceland, Latvia,
Italy & Cross-border payment account opening, PID issuance and
verification, age verification, organizational identity \\
\textbf{EWC (EU Digital Identity Wallet Consortium)} & Sweden, Finland,
France, Germany, Czech Republic, Croatia, Italy, Austria, Greece,
Romania & Travel (digital travel credentials), education (diploma
attestation), digital payments, organizational wallet use cases \\
\textbf{DC4EU (Digital Credentials for Europe)} & Spain, Austria,
Belgium, France, Germany, Italy, Luxembourg, Poland, Romania, Sweden &
Education (European Digital Credentials for Learning), social security
(Professional Qualifications, Portable Document A1) \\
\end{longtable}

\subsection{What the Pilots Are
Testing}\label{_what_the_pilots_are_testing}

The signing-relevant findings from the pilots (as of early 2026, with
pilots still running):

\textbf{QES from mobile works.} The core signing flow --- user
authentication, hash submission to remote QSCD, CMS return, PDF
embedding --- has been demonstrated successfully across multiple
implementations. The user experience typically involves a single
biometric confirmation (fingerprint or face recognition) with a clear
"sign this document" prompt.

\textbf{Cross-border validation is the challenge.} Creating the
signature is (relatively) straightforward. Validating it across borders
--- ensuring that a signature created by a French citizen's wallet
validates correctly in a German government portal --- requires all
parties to consume the EUTL, support the same PAdES profiles, and handle
certificate chains from unfamiliar QTSPs. This is exactly the same
interoperability problem discussed earlier in this chapter, amplified by
scale.

\textbf{TSA integration needs standardization.} Different wallet
implementations use different TSAs, with different timestamp token
sizes, different certificate chains, and different response times. The
pilots have exposed edge cases where timestamp tokens from one TSA cause
validation issues in validators that are not prepared for the TSA's
specific certificate chain length.

\textbf{User comprehension is limited.} In pilot user testing, most
participants understood "sign this document" but did not understand what
a qualified electronic signature \emph{is}, how it differs from a basic
electronic signature, or why it has legal equivalence to a handwritten
signature. This is not a technical problem, but it has technical
consequences: if users do not understand the significance of QES, they
may not treat the signing authentication with appropriate care.

\section{What Developers Need to
Prepare}\label{_what_developers_need_to_prepare}

\subsection{If You Build Signing
Software}\label{_if_you_build_signing_software}

\begin{enumerate}
\def\labelenumi{\arabic{enumi}.}
\item
  \textbf{Support wallet-initiated signing flows.} The wallet
  communicates with your application via OpenID4VP (OpenID for
  Verifiable Presentations). You will need to implement the RP (Relying
  Party) side of the OpenID4VP protocol to receive signing requests from
  the wallet.
\item
  \textbf{Handle remote QSCD signatures.} Your application will receive
  a CMS \texttt{SignedData} blob from the QTSP, not generate it locally.
  Your signing pipeline must support "external CMS" injection ---
  computing the byte range, sending the hash for external signing,
  receiving the CMS, and embedding it in the PDF.
\item
  \textbf{Pre-allocate generously.} Wallet-generated CMS may include
  longer certificate chains than you are accustomed to (government PKI
  chains are often deeper than commercial CA chains). Pre-allocate at
  least 64 KB for PAdES-B-LTA.
\item
  \textbf{Add DSS and document timestamps server-side.} The mobile
  device will not populate the DSS or add document timestamps. Your
  server must do this, which means your server needs access to OCSP
  responders and CRL distribution points for the signer's certificate
  chain, plus a TSA for the document timestamp.
\end{enumerate}

\subsection{If You Build Validation
Software}\label{_if_you_build_validation_software}

\begin{enumerate}
\def\labelenumi{\arabic{enumi}.}
\item
  \textbf{Consume the EUTL programmatically.} Wallet-generated
  signatures chain to QTSPs on the EUTL. If your validator uses a static
  trust store, it will reject these signatures. Implement automated EUTL
  consumption --- fetch the LOT (List of Trusted Lists), parse each
  member state's trusted list, extract the trust anchors, and update
  your trust store on a regular schedule (daily is typical).
\item
  \textbf{Cache the EUTL.} With millions of wallet-generated signatures
  to validate, you cannot fetch the EUTL for every validation. Cache the
  parsed trust anchors with a TTL (time-to-live) and refresh
  periodically.
\item
  \textbf{Support deep certificate chains.} Government PKI chains can be
  4-5 certificates deep (signer → intermediate 1 → intermediate 2 →
  root). Ensure your chain-building code handles this without
  performance degradation.
\item
  \textbf{Report the PAdES profile level.} Wallet-generated signatures
  should be PAdES-B-LTA. If a wallet-generated signature is missing LTV
  data (B-T or B-B instead of B-LTA), this indicates a problem in the
  signing pipeline. Report the profile level to help the signer or their
  application diagnose the issue.
\item
  \textbf{Prepare for QEAA validation.} If your application needs to
  verify not just \emph{who signed} but \emph{what qualifications the
  signer has} (e.g., a medical prescription signed by a doctor), you
  will need to validate QEAAs in addition to the PDF signature. This is
  a separate validation path: the QEAA is a signed credential (SD-JWT or
  mdoc) that you verify against the issuing QTSP's certificate on the
  EUTL.
\end{enumerate}

\subsection{If You Build Enterprise
Systems}\label{_if_you_build_enterprise_systems}

\begin{enumerate}
\def\labelenumi{\arabic{enumi}.}
\item
  \textbf{Update your trust store management.} If your enterprise trust
  store is managed manually (a Java \texttt{cacerts} file, a Windows
  certificate store), automate it. Wallet-generated signatures will come
  from QTSPs that may not be in your current trust store.
\item
  \textbf{Budget for EUTL integration.} Consuming the EUTL is not
  trivial: the LOT and individual trusted lists are XML documents signed
  with specific certificates, using a specific schema (ETSI TS 119 612).
  The EU DSS framework (Chapter 23) provides a reference implementation,
  but integrating it into your architecture requires engineering effort.
\item
  \textbf{Plan for volume.} If your system currently validates hundreds
  of qualified signatures per day, prepare for thousands or tens of
  thousands. The wallet will democratize QES, and your systems will see
  the increase.
\item
  \textbf{Test with pilot certificates.} The EUDI wallet pilots issue
  test certificates that can be used for integration testing. Engage
  with the pilot programs relevant to your member state and obtain test
  credentials for your validation pipeline.
\end{enumerate}

\section{The Open Questions}\label{_the_open_questions}

The EUDI wallet is ambitious, unprecedented, and not yet fully defined.
Several significant questions remain open as of early 2026:

\textbf{Will users actually adopt it?} Mandating reliance parties to
accept the wallet creates demand, but the wallet still needs a
compelling user experience. If the signing flow feels cumbersome ---
multiple authentication steps, slow QSCD response times, confusing
consent screens --- users will avoid it.

\textbf{Will Apple and Google cooperate?} The wallet runs on
smartphones, which means it needs integration with iOS and Android. Both
platforms control access to NFC, secure elements, and biometric APIs.
The EU Digital Markets Act (DMA) requires gatekeepers to allow
third-party wallet apps to access device hardware, but the
implementation details are being negotiated, and the mobile platforms
have commercial interests in their own wallet ecosystems (Apple Pay,
Google Pay).

\textbf{What happens with AATL?} If hundreds of millions of qualified
signatures are created through the wallet but the signing certificates
are not in Adobe's AATL, Acrobat users will see yellow warnings on
legitimate QES. This creates political pressure on Adobe to either
accept the EUTL as a trust anchor by default or face regulatory action.
The outcome will shape the PDF signature ecosystem for years.

\textbf{Will the credential format converge?} The current ARF supports
both SD-JWT and mdoc. Supporting two formats doubles the implementation
effort for relying parties. The industry is watching to see if one
format emerges as dominant or if the dual-format approach persists.

\textbf{How will cross-border QES work in practice?} The technical
framework is defined. The legal framework mandates mutual recognition.
But whether a German government portal will actually accept a Romanian
wallet-generated QES without friction --- that depends on whether all
the implementation details are correct across two different member state
implementations, two different QTSP certificate chains, and two
different validation stacks.

The answers to these questions will determine whether the EUDI wallet
becomes the most significant development in digital signing since eIDAS
itself --- or another ambitious European digital project that works in
theory and disappoints in practice.

The EUDI wallet is a moving target. The ARF is being updated,
implementing acts are being published, and pilots are generating
findings that feed back into the specifications. If you are reading this
in 2027 or later, some of the specifics in this chapter may have
changed. The architectural principles --- remote QSCD, selective
disclosure, EUTL-based trust, PAdES-B-LTA as the target profile --- are
stable. The protocol details and credential format choices may have
evolved.

\section{Key Takeaways}\label{_key_takeaways_14}

\begin{itemize}
\item
  \textbf{eIDAS 2.0 (Regulation (EU) 2024/1183 (eIDAS 2.0))} mandates
  that every EU member state offer a digital identity wallet to citizens
  by end of 2026, with QES creation capability included
\item
  \textbf{The EUDI wallet} is a smartphone application that stores
  government-issued identity data (PID) and qualified attribute
  attestations (QEAAs) and can create qualified electronic signatures
\item
  \textbf{Remote QSCD} is the signing architecture: the private key
  lives in a certified HSM operated by a QTSP, not on the phone itself
  --- the wallet provides authentication and sole control
\item
  \textbf{Selective disclosure} (via SD-JWT or mdoc) allows users to
  share only the attributes a relying party needs, without revealing
  their full identity
\item
  \textbf{Free QES for citizens} (Article 45e) will dramatically
  increase the volume of qualified signatures --- plan your
  infrastructure for orders-of-magnitude growth
\item
  \textbf{PAdES-B-LTA is the target profile} for wallet-generated
  signatures --- server-side DSS population and document timestamping
  are required
\item
  \textbf{EUTL consumption is essential} --- wallet-generated signatures
  chain to QTSPs that may not be in your current trust store; automate
  EUTL updates
\item
  \textbf{Four Large-Scale Pilots} (POTENTIAL, NOBID, EWC, DC4EU) are
  testing the wallet across sectors and member states; engage with them
  for test credentials
\item
  \textbf{Open questions remain} about user adoption, platform
  cooperation (Apple/Google), AATL integration, credential format
  convergence, and real-world cross-border validation
\item
  \textbf{The architectural principles are PDF-native} --- if you
  understand PAdES, certificate chains, and the EUTL, you understand the
  cryptographic foundation of the EUDI wallet
\end{itemize}

\chapter{Code, Tools \&
Implementation}\label{_code_tools_implementation}

\chapter{The Java PDF Library
Landscape}\label{_the_java_pdf_library_landscape}

Before we write a single line of signing code, we need to talk about the
tools we're using --- where they came from, why they exist, and what
baggage they carry. If you've ever wondered why there are three
different "iText" libraries on Maven Central, why PDFBox's API changed
so dramatically between version 2 and version 3, or why a library from
the year 2000 with known bugs is still running in Fortune 500 production
systems --- this chapter is for you.

The Java PDF ecosystem is not just a collection of libraries. It is a
story of open-source idealism, licensing betrayal (depending on who you
ask), corporate acquisitions, community forks, and the kind of technical
debt that accumulates when enterprises build critical infrastructure on
a dependency and then the rules change. Understanding this history is
not optional background reading --- it directly affects the code you
write, the signatures you produce, and the compatibility problems you
will encounter.

\section{The iText Saga: From PhD Project to Corporate Licensing
Drama}\label{_the_itext_saga_from_phd_project_to_corporate_licensing_drama}

\subsection{Bruno Lowagie and the Birth of iText
(2000--2009)}\label{_bruno_lowagie_and_the_birth_of_itext_20002009}

The story of Java PDF processing begins, for practical purposes, with
Bruno Lowagie. In 2000, Lowagie --- a Belgian mathematician working at a
printing company --- needed to generate PDF documents from Java. The
options at the time were essentially nonexistent. Adobe had published
the PDF specification, but the only tools that could create PDFs were
Adobe's own products (Acrobat, Distiller) and a handful of expensive
commercial libraries. There was nothing usable in the open-source Java
world.

So Lowagie wrote his own library. He called it \textbf{iText}, released
it under the \textbf{Mozilla Public License (MPL)} and later the
\textbf{GNU Lesser General Public License (LGPL)}, and published it on
SourceForge --- the GitHub of its era. The LGPL license was critically
important: it allowed companies to use iText in proprietary software
without open-sourcing their own code, as long as they didn't modify the
iText library itself. This is the license model that made iText
ubiquitous.

The timing was perfect. Java was exploding in enterprise environments.
J2EE application servers were the dominant architecture. Companies
needed to generate invoices, reports, contracts, and forms as PDFs ---
millions of them per day, in automated workflows. iText was free, it
worked, and its API was straightforward:

\begin{Shaded}
\begin{Highlighting}[]
\CommentTok{// iText 2.x — the classic API that a generation of Java developers learned}
\BuiltInTok{Document}\NormalTok{ document }\OperatorTok{=} \KeywordTok{new} \BuiltInTok{Document}\OperatorTok{();}
\NormalTok{PdfWriter}\OperatorTok{.}\FunctionTok{getInstance}\OperatorTok{(}\NormalTok{document}\OperatorTok{,} \KeywordTok{new} \BuiltInTok{FileOutputStream}\OperatorTok{(}\StringTok{"output.pdf"}\OperatorTok{));}
\NormalTok{document}\OperatorTok{.}\FunctionTok{open}\OperatorTok{();}
\NormalTok{document}\OperatorTok{.}\FunctionTok{add}\OperatorTok{(}\KeywordTok{new} \FunctionTok{Paragraph}\OperatorTok{(}\StringTok{"Hello, PDF"}\OperatorTok{));}
\NormalTok{document}\OperatorTok{.}\FunctionTok{close}\OperatorTok{();}
\end{Highlighting}
\end{Shaded}

This simplicity was revolutionary. You didn't need to understand PDF
internals. You didn't need to manipulate dictionaries or worry about
cross-reference tables. You told iText what you wanted on the page, and
it produced a valid PDF. The library handled fonts, images, tables,
forms, encryption, and --- crucially for our purposes ---
\textbf{digital signatures}.

By the mid-2000s, iText was everywhere. It was embedded in JBoss,
Alfresco, Jasper Reports, BIRT, and hundreds of other Java projects.
Large banks, insurance companies, government agencies, and healthcare
organizations had built their document generation pipelines on iText.
The Maven coordinates \texttt{com.lowagie:itext} were among the most
downloaded artifacts in the Java ecosystem.

\subsection{The LGPL-to-AGPL Switch: The Day the Rules Changed
(2009)}\label{_the_lgpl_to_agpl_switch_the_day_the_rules_changed_2009}

On December 1, 2009, Bruno Lowagie changed the license of iText from
\textbf{LGPL} to \textbf{AGPL} (GNU Affero General Public License),
starting with version \textbf{5.0.0}.

This was, depending on your perspective, either a reasonable business
decision or an act of betrayal against the community that had made iText
successful. Both perspectives have merit, and understanding the
controversy matters because its consequences are still playing out in
production systems today.

\textbf{What AGPL means in practice:} The AGPL is like the GPL, but with
one additional requirement: if you use an AGPL-licensed library in
software that users interact with over a network (which includes web
applications, SaaS products, and API services), you must release the
\emph{entire source code} of your application under the AGPL. Not just
your modifications to iText --- your entire application.

For the thousands of companies using iText in their proprietary Java
backends --- generating invoices, processing insurance claims, signing
contracts --- this was devastating. They had exactly three options:

\begin{enumerate}
\def\labelenumi{\arabic{enumi}.}
\item
  \textbf{Open-source their entire application} under AGPL ---
  unthinkable for most enterprises
\item
  \textbf{Buy a commercial license} from iText Group NV (the company
  Lowagie founded to monetize iText) --- expensive, typically
  €5,000--€50,000+ per year depending on deployment scale
\item
  \textbf{Stay on iText 2.1.7} --- the last LGPL-licensed version,
  released in 2009 --- forever
\end{enumerate}

Most chose option 3.

\textbf{Why Lowagie made the switch:} Lowagie's stated reasons were
straightforward: he had spent nearly a decade maintaining a complex
library for free, while companies worth billions used it to generate
revenue without contributing back. He wanted to build a business. The
AGPL created a forcing function: use iText for free in open-source
projects, or pay for a commercial license. This is a perfectly
legitimate business model --- MongoDB, Elasticsearch, and Redis would
later make similar moves (though with different licenses).

\textbf{Why the community was angry:} The anger wasn't about the
business model itself --- it was about the trust violation. Enterprise
architects had chosen iText \emph{because} of its LGPL license. They had
written it into procurement documents, gotten legal approval, and built
systems that depended on it. The license change felt like a
bait-and-switch: attract an enormous user base with a permissive
license, then change the terms once everyone is locked in.

The fact that existing versions remained LGPL was cold comfort. iText
2.1.7 was already showing its age in 2009. By 2012, it had known bugs in
font handling, incomplete PDF 1.7 support, and no support for the newer
CMS structures that PAdES was beginning to require. But upgrading to
iText 5.x meant either paying or going AGPL.

If you encounter a PDF signed by a Java application that uses
\texttt{com.lowagie.text.pdf.PdfSignatureAppearance} in its imports ---
that's iText 2.x. The class was renamed to
\texttt{com.itextpdf.text.pdf.PdfSignatureAppearance} in iText 5.x (note
the package change from \texttt{com.lowagie} to \texttt{com.itextpdf}).
The package name tells you which era of iText produced the signature.

\subsection{The CMS Structure Problem: Why Old iText Signatures Look
Different}\label{_the_cms_structure_problem_why_old_itext_signatures_look_different}

Here's where the licensing drama directly impacts signature
compatibility.

iText 2.x (the LGPL version that much of the world is still using)
produces CMS \texttt{SignedData} structures that are technically valid
per RFC 5652 --- Cryptographic Message Syntax (CMS), but differ from
what modern validators expect in several ways:

\textbf{Encapsulated vs. Detached content:} Early iText versions could
produce CMS signatures with \emph{encapsulated} content --- meaning the
signed hash was embedded inside the \texttt{SignedData} structure's
\texttt{encapContentInfo} field rather than being referenced externally
via the PDF byte range. This was not wrong per the CMS spec, but it
created problems:

\begin{itemize}
\item
  Modern PAdES validators expect detached signatures (SubFilter
  \texttt{adbe.pkcs7.detached} or \texttt{ETSI.CAdES.detached}), where
  the signed content is the byte range in the PDF itself --- not a copy
  inside the CMS
\item
  Some validators would see encapsulated content and assume the
  signature was standalone (not bound to the PDF), failing validation
  even though the cryptographic math was correct
\item
  The double-hash problem: if the signed bytes in the CMS didn't match
  what a byte-range-aware validator computed from the PDF file, the
  signature appeared invalid
\end{itemize}

\textbf{Missing signed attributes:} iText 2.x did not always include the
\texttt{ESS-signing-certificate-v2} attribute that PAdES requires.
Without this attribute, the signature is valid per ISO 32000-2 but
\emph{not} PAdES-compliant --- which means it fails validation in any
system that enforces PAdES profiles (including EU government validation
services and Adobe Acrobat in its stricter validation modes).

\textbf{SHA-1 default:} iText 2.1.7 defaulted to SHA-1 for the message
digest. SHA-1 was already theoretically weakened by 2009 (Wang and Yu's
collision attacks were published in 2005), and it was definitively
broken in 2017 when the SHAttered attack produced actual colliding
files. PDFs signed with iText 2.x before SHA-256 became the norm carry
SHA-1 digests, which modern validators increasingly flag or reject.

If your organization has a corpus of PDFs signed between 2005 and 2015
by Java applications, there is a significant chance they were signed
with iText 2.x using SHA-1 and potentially using encapsulated (not
detached) CMS content. These signatures may fail validation in modern
tools.

To verify: extract the CMS SignedData from the \texttt{/Contents} entry
and check:

\begin{enumerate}
\def\labelenumi{\arabic{enumi}.}
\item
  Is \texttt{encapContentInfo.eContentType} set with actual content?
  (encapsulated --- old style)
\item
  Is the digest algorithm SHA-1 (\texttt{OID\ 1.3.14.3.2.26})?
  (deprecated)
\item
  Is the \texttt{ESS-signing-certificate-v2} attribute present in
  \texttt{signedAttrs}? (missing = not PAdES-compliant)
\end{enumerate}

Chapter 31 provides code to perform all three checks.

\subsection{iText 5.x, 7.x, and 8.x: The Commercial
Evolution}\label{_itext_5_x_7_x_and_8_x_the_commercial_evolution}

After the AGPL switch, iText evolved rapidly --- but behind a paywall
for commercial use.

\textbf{iText 5.x (2009--2016):} Maintained the original API design but
added PDF 2.0 features, improved CMS handling, and better PAdES support.
The package name changed from \texttt{com.lowagie} to
\texttt{com.itextpdf}. If you see both package prefixes in the same
project's classpath --- and this happens more often than you'd think ---
you're looking at a project that tried to use both the old LGPL version
and the new AGPL version simultaneously, which is a recipe for
classloading nightmares.

\textbf{iText 7.x (2016--2023):} A complete rewrite with a modular
architecture. The API was entirely different from iText 5.x ---
migration was not trivial. iText 7 introduced separate modules for core
PDF operations, forms, signing, and other features. The signing module
(\texttt{itext7-sign}) included comprehensive PAdES support, including
B-LTA profiles with embedded timestamps and revocation data. The
commercial license became the primary offering, with AGPL remaining for
open-source projects.

\textbf{iText 8.x (2023--present):} The current generation, continuing
the modular architecture with added support for PDF 2.0 features,
improved accessibility tagging, and enhanced signing capabilities. iText
Group NV (now iText Software, headquartered in Ghent, Belgium, with
offices in the US and Singapore) is a substantial commercial entity with
paying enterprise customers worldwide.

The commercial iText is, by all accounts, a high-quality library. Its
signing implementation is mature, well-tested, and PAdES-compliant. But
at commercial license prices that can reach tens of thousands of euros
per year for enterprise deployments, it is not an option for every
project --- which is why the open-source alternatives exist.

\subsection{The IBM Connection}\label{_the_ibm_connection}

IBM's involvement with iText is a lesser-known but important part of the
story. In the early 2000s, IBM was one of the largest users of iText ---
embedding it in several of their enterprise products, including
\textbf{IBM Content Manager}, \textbf{IBM Filenet}, and various
WebSphere-based document processing systems.

IBM contributed patches and improvements back to the project during the
LGPL era, which was entirely within the LGPL's collaborative spirit.
However, when the license changed to AGPL, IBM --- like most large
enterprises --- was unwilling to either open-source their proprietary
products or pay iText's commercial licensing fees for a dependency that
had previously been free.

This created an awkward situation: IBM's products contained iText 2.x
code, their customers expected continued PDF functionality, and the
upgrade path required either legal procurement (commercial iText
license) or engineering investment (migrate to a different library).
IBM's response, like most enterprises in this position, was a
combination of staying on old versions and gradually migrating new
development to Apache PDFBox --- which was, by then, becoming mature
enough for production use.

The IBM story is representative of what happened across the industry:
hundreds of enterprises had iText deeply embedded in their systems, and
the license change created a slow-motion migration that, in many cases,
is still incomplete in 2026.

\section{OpenPDF: The Community Fork That Refuses to
Die}\label{_openpdf_the_community_fork_that_refuses_to_die}

\subsection{The Fork (2016)}\label{_the_fork_2016}

When it became clear that iText 5.x was permanently AGPL for commercial
use, and that iText 2.1.7 was accruing unfixed bugs and missing modern
features, a group of developers forked iText 2.1.7 under the name
\textbf{OpenPDF} (initially called \textbf{LibrePDF}).

The fork was hosted on GitHub under the \texttt{LibrePDF/OpenPDF}
repository, and it was explicitly positioned as the LGPL-licensed
continuation of the original iText codebase. The initial stated goals
were modest: fix bugs, add basic modern features (SHA-256 support,
better font handling), and keep the library usable for projects that
needed LGPL licensing.

\subsection{What OpenPDF Is (and
Isn't)}\label{_what_openpdf_is_and_isnt}

OpenPDF today (version 2.0.x, as of 2026) is a capable PDF generation
library that maintains API compatibility with iText 2.x while adding
features that the original 2.1.7 never had:

\begin{itemize}
\item
  \textbf{SHA-256/384/512 message digests} for signing (critical ---
  SHA-1 is no longer acceptable)
\item
  \textbf{AES-256 encryption} support
\item
  \textbf{Improved font handling}, including better CJK support and
  TrueType subsetting
\item
  \textbf{Bug fixes} for dozens of issues that were never patched in
  iText 2.1.7
\end{itemize}

However, OpenPDF also carries significant limitations:

\begin{itemize}
\item
  It is built on the iText 2.x architecture, which was designed for PDF
  1.7 and earlier
\item
  \textbf{No PDF 2.0 support} --- cross-reference streams, tagged PDF
  improvements, and other 2.0 features are absent
\item
  \textbf{Limited PAdES support} --- basic PAdES-B-B signatures are
  possible, but B-LT and B-LTA profiles (which require DSS dictionary
  construction and archival timestamps) are not natively supported
\item
  \textbf{No built-in LTV (Long-Term Validation)} --- you'll need to
  construct the DSS dictionary yourself
\item
  The contributor community is small --- a handful of maintainers, not
  the large team that iText Group NV employs
\item
  \textbf{Documentation is sparse} --- the iText 2.x book (Lowagie's
  "iText in Action") is the closest thing to comprehensive
  documentation, and it describes an API from 2008
\end{itemize}

OpenPDF lives at \texttt{com.github.librepdf:openpdf} on Maven Central.
Do not confuse it with \texttt{com.lowagie:itext} (the original iText
2.1.7) or \texttt{com.itextpdf:itextpdf} (the AGPL iText 5.x). All three
exist on Maven Central. All three have different licenses, different
capabilities, and different bugs. Getting the wrong one in your
dependency tree is a common source of confusion --- and potential legal
liability.

\subsection{Should You Use OpenPDF?}\label{_should_you_use_openpdf}

For \textbf{PDF generation} (creating invoices, reports, forms) where
you need an LGPL-licensed library and don't require PDF 2.0 features:
OpenPDF is a reasonable choice.

For \textbf{digital signatures}: OpenPDF can produce basic CMS
signatures, but for PAdES compliance, long-term validation, and the kind
of robust signing infrastructure this book teaches, \textbf{PDFBox is
the better choice}. PDFBox's signing API is more actively maintained,
better tested against the Bochum attack scenarios, and has a larger
community contributing security fixes.

That said --- if you have an existing codebase built on iText 2.x that
you cannot easily migrate, OpenPDF is the path of least resistance for
getting off the unmaintained 2.1.7 version while keeping your existing
API calls intact.

\section{Apache PDFBox: The Open-Source
Standard}\label{_apache_pdfbox_the_open_source_standard}

\subsection{Origin Story: Ben Litchfield and the Apache
Incubation}\label{_origin_story_ben_litchfield_and_the_apache_incubation}

Apache PDFBox started as \textbf{PDFBox} --- a standalone project
created by \textbf{Ben Litchfield} in 2002. Litchfield, working as a
Java developer, needed to \emph{read} existing PDF documents --- not
just create them (which was iText's primary strength at the time). He
needed to extract text, parse structures, and manipulate the low-level
PDF object model. iText was designed as a PDF \emph{writer}; PDFBox was
designed as a PDF \emph{reader} that could also write.

This fundamental design difference --- reader-first vs. writer-first ---
is still visible in the APIs today:

\begin{itemize}
\item
  \textbf{iText/OpenPDF} excels at high-level document construction:
  "add a paragraph here, put a table there, generate the PDF"
\item
  \textbf{PDFBox} excels at low-level document manipulation: "give me
  the COS dictionary for object 42, let me modify the annotation
  rectangle, show me the raw font encoding"
\end{itemize}

For digital signature work, PDFBox's reader-first architecture is a
significant advantage. Signing a PDF requires understanding the existing
document structure, computing byte ranges, parsing existing signatures,
and manipulating specific dictionaries --- all tasks where low-level
access is essential.

In 2008, PDFBox entered the \textbf{Apache Software Foundation}
incubator process, and in 2009 it graduated as a top-level Apache
project under the \textbf{Apache License 2.0} --- one of the most
permissive open-source licenses available. Unlike iText's LGPL→AGPL
journey, PDFBox has been Apache-licensed since inception and will remain
so. This licensing certainty is one of the reasons enterprises migrated
to PDFBox.

\subsection{PDFBox 2.x: The Workhorse
(2016--2023)}\label{_pdfbox_2_x_the_workhorse_20162023}

PDFBox 2.0 was the version that brought PDFBox into mainstream
production use. It featured:

\begin{itemize}
\item
  A complete COS (Carousel Object System) model for low-level PDF
  manipulation
\item
  Text extraction with position information
\item
  PDF rendering (via \texttt{PDFRenderer})
\item
  AcroForm support (creating, filling, and flattening forms)
\item
  Digital signature support via \texttt{PDDocument.addSignature()}
\item
  Font handling (Type1, TrueType, CID fonts, subsetting)
\end{itemize}

The signing API in PDFBox 2.x was functional but had rough edges. The
\texttt{SignatureInterface} required you to provide a
\texttt{sign(InputStream)} method that received the bytes to sign and
returned the CMS \texttt{SignedData} as a byte array. The actual CMS
construction was your responsibility --- typically using BouncyCastle
(we'll discuss this later). PDFBox handled the PDF-specific parts:
creating the signature dictionary, computing the byte range, reserving
space for \texttt{/Contents}, and performing the incremental save.

\begin{Shaded}
\begin{Highlighting}[]
\CommentTok{// PDFBox 2.x signing pattern}
\NormalTok{PDDocument doc }\OperatorTok{=}\NormalTok{ PDDocument}\OperatorTok{.}\FunctionTok{load}\OperatorTok{(}\NormalTok{inputFile}\OperatorTok{);}
\NormalTok{PDSignature signature }\OperatorTok{=} \KeywordTok{new} \FunctionTok{PDSignature}\OperatorTok{();}
\NormalTok{signature}\OperatorTok{.}\FunctionTok{setFilter}\OperatorTok{(}\NormalTok{PDSignature}\OperatorTok{.}\FunctionTok{FILTER\_ADOBE\_PPKLITE}\OperatorTok{);}
\NormalTok{signature}\OperatorTok{.}\FunctionTok{setSubFilter}\OperatorTok{(}\NormalTok{PDSignature}\OperatorTok{.}\FunctionTok{SUBFILTER\_ADBE\_PKCS7\_DETACHED}\OperatorTok{);}
\NormalTok{signature}\OperatorTok{.}\FunctionTok{setName}\OperatorTok{(}\StringTok{"Signer Name"}\OperatorTok{);}
\NormalTok{signature}\OperatorTok{.}\FunctionTok{setSignDate}\OperatorTok{(}\BuiltInTok{Calendar}\OperatorTok{.}\FunctionTok{getInstance}\OperatorTok{());}

\NormalTok{doc}\OperatorTok{.}\FunctionTok{addSignature}\OperatorTok{(}\NormalTok{signature}\OperatorTok{,} \KeywordTok{new} \FunctionTok{SignatureInterface}\OperatorTok{()} \OperatorTok{\{}
    \AttributeTok{@Override}
    \KeywordTok{public} \DataTypeTok{byte}\OperatorTok{[]} \FunctionTok{sign}\OperatorTok{(}\BuiltInTok{InputStream}\NormalTok{ content}\OperatorTok{)} \KeywordTok{throws} \BuiltInTok{IOException} \OperatorTok{\{}
        \CommentTok{// Your BouncyCastle CMS construction goes here}
        \ControlFlowTok{return} \FunctionTok{createCMSSignature}\OperatorTok{(}\NormalTok{content}\OperatorTok{,}\NormalTok{ certificate}\OperatorTok{,}\NormalTok{ privateKey}\OperatorTok{);}
    \OperatorTok{\}}
\OperatorTok{\});}

\NormalTok{doc}\OperatorTok{.}\FunctionTok{saveIncremental}\OperatorTok{(}\NormalTok{outputStream}\OperatorTok{);}
\NormalTok{doc}\OperatorTok{.}\FunctionTok{close}\OperatorTok{();}
\end{Highlighting}
\end{Shaded}

This separation of concerns --- PDFBox handles the PDF, you handle the
cryptography --- is elegant in theory but demanding in practice. Getting
the CMS structure right for PAdES compliance requires deep knowledge of
signed attributes, certificate embedding, and timestamp handling.
Chapters 29 and 30 provide the complete implementation.

\subsection{PDFBox 3.x: The Breaking Change
(2024)}\label{_pdfbox_3_x_the_breaking_change_2024}

PDFBox 3.0 was released in 2024 as a major version bump with
\textbf{breaking API changes}. The migration was significant enough to
catch many projects off guard:

PDFBox 3.x is \textbf{not} a drop-in replacement for 2.x. If you're
upgrading, budget significant development time. The changes are
pervasive and affect every area of the API.

The most impactful changes for signing:

\textbf{Loading documents:}

\begin{Shaded}
\begin{Highlighting}[]
\CommentTok{// PDFBox 2.x}
\NormalTok{PDDocument doc }\OperatorTok{=}\NormalTok{ PDDocument}\OperatorTok{.}\FunctionTok{load}\OperatorTok{(}\KeywordTok{new} \BuiltInTok{File}\OperatorTok{(}\StringTok{"document.pdf"}\OperatorTok{));}
\NormalTok{PDDocument doc }\OperatorTok{=}\NormalTok{ PDDocument}\OperatorTok{.}\FunctionTok{load}\OperatorTok{(}\NormalTok{byteArray}\OperatorTok{);}

\CommentTok{// PDFBox 3.x — PDDocument.load() is GONE}
\NormalTok{PDDocument doc }\OperatorTok{=}\NormalTok{ Loader}\OperatorTok{.}\FunctionTok{loadPDF}\OperatorTok{(}\KeywordTok{new} \BuiltInTok{File}\OperatorTok{(}\StringTok{"document.pdf"}\OperatorTok{));}
\NormalTok{PDDocument doc }\OperatorTok{=}\NormalTok{ Loader}\OperatorTok{.}\FunctionTok{loadPDF}\OperatorTok{(}\NormalTok{byteArray}\OperatorTok{);}
\end{Highlighting}
\end{Shaded}

\textbf{Font construction:}

\begin{Shaded}
\begin{Highlighting}[]
\CommentTok{// PDFBox 2.x}
\NormalTok{PDType1Font font }\OperatorTok{=}\NormalTok{ PDType1Font}\OperatorTok{.}\FunctionTok{HELVETICA}\OperatorTok{;}

\CommentTok{// PDFBox 3.x — static font constants are GONE}
\NormalTok{PDType1Font font }\OperatorTok{=} \KeywordTok{new} \FunctionTok{PDType1Font}\OperatorTok{(}\NormalTok{Standard14Fonts}\OperatorTok{.}\FunctionTok{FontName}\OperatorTok{.}\FunctionTok{HELVETICA}\OperatorTok{);}
\end{Highlighting}
\end{Shaded}

\textbf{AcroForm API:}

\begin{Shaded}
\begin{Highlighting}[]
\CommentTok{// PDFBox 2.x}
\DataTypeTok{boolean}\NormalTok{ needsAppearances }\OperatorTok{=}\NormalTok{ acroForm}\OperatorTok{.}\FunctionTok{isNeedAppearances}\OperatorTok{();}

\CommentTok{// PDFBox 3.x — method renamed}
\DataTypeTok{boolean}\NormalTok{ needsAppearances }\OperatorTok{=}\NormalTok{ acroForm}\OperatorTok{.}\FunctionTok{getNeedAppearances}\OperatorTok{();}
\end{Highlighting}
\end{Shaded}

\textbf{Font encoding access:}

\begin{Shaded}
\begin{Highlighting}[]
\CommentTok{// PDFBox 2.x}
\BuiltInTok{Encoding}\NormalTok{ encoding }\OperatorTok{=}\NormalTok{ font}\OperatorTok{.}\FunctionTok{getEncoding}\OperatorTok{();} \CommentTok{// Worked for most font types}

\CommentTok{// PDFBox 3.x — getEncoding() removed from PDFont interface}
\CommentTok{// Must access encoding through COS dictionary directly}
\NormalTok{COSDictionary fontDict }\OperatorTok{=}\NormalTok{ font}\OperatorTok{.}\FunctionTok{getCOSObject}\OperatorTok{();}
\NormalTok{COSBase encoding }\OperatorTok{=}\NormalTok{ fontDict}\OperatorTok{.}\FunctionTok{getDictionaryObject}\OperatorTok{(}\NormalTok{COSName}\OperatorTok{.}\FunctionTok{ENCODING}\OperatorTok{);}
\end{Highlighting}
\end{Shaded}

\textbf{Content streams:}

\begin{Shaded}
\begin{Highlighting}[]
\CommentTok{// PDFBox 2.x — could write to form XObjects with PDPageContentStream}
\NormalTok{PDPageContentStream cs }\OperatorTok{=} \KeywordTok{new} \FunctionTok{PDPageContentStream}\OperatorTok{(}\NormalTok{doc}\OperatorTok{,}\NormalTok{ formXObject}\OperatorTok{,} \KeywordTok{...}\OperatorTok{);}

\CommentTok{// PDFBox 3.x — PDPageContentStream is for pages ONLY}
\CommentTok{// For form XObjects, write raw PDF operators via the stream\textquotesingle{}s output}
\BuiltInTok{OutputStream}\NormalTok{ os }\OperatorTok{=}\NormalTok{ formXObject}\OperatorTok{.}\FunctionTok{getStream}\OperatorTok{().}\FunctionTok{createOutputStream}\OperatorTok{();}
\end{Highlighting}
\end{Shaded}

The signing API itself (\texttt{PDDocument.addSignature()},
\texttt{SignatureInterface}, \texttt{PDSignature}) remained structurally
similar, but the surrounding changes --- especially in document loading,
font handling, and content stream creation --- mean that any signing
code that also creates visual signature appearances needs substantial
rework.

\textbf{Migration strategy:} Don't try to upgrade PDFBox in-place.
Create a new branch, update the dependency, let the compiler tell you
every break, and fix them systematically. The PDFBox team maintains a
\href{https://pdfbox.apache.org/3.0/migration.html}{migration guide}
that covers the major changes, but the edge cases (especially around
font encoding and content stream APIs) will require reading the 3.x
Javadoc carefully.

\subsection{The PDFBox Signing Infrastructure: What's Under the
Hood}\label{_the_pdfbox_signing_infrastructure_whats_under_the_hood}

PDFBox's signing implementation is one of the most impressive pieces of
open-source PDF engineering. It's worth understanding the high-level
architecture, because when something goes wrong during signing --- and
it will --- you need to know where to look.

The signing process in PDFBox involves these internal steps:

\begin{enumerate}
\def\labelenumi{\arabic{enumi}.}
\item
  \textbf{Signature dictionary creation}: PDFBox creates the
  \texttt{/Sig} dictionary with the properties you set
  (\texttt{/Filter}, \texttt{/SubFilter}, \texttt{/Name}, \texttt{/M},
  etc.)
\item
  \textbf{Space estimation}: PDFBox needs to know how much space to
  reserve for the \texttt{/Contents} hex string \emph{before} the
  signature is computed --- because the signature covers the bytes
  outside the \texttt{/Contents} gap, and the file must be written
  before the signature can be calculated. The default is 16KB (32,768
  hex characters), but PAdES-B-LTA signatures with embedded timestamps
  and revocation data can easily exceed this.
\item
  \textbf{Byte range computation}: PDFBox calculates the
  \texttt{/ByteRange} array --- the two regions of the file that the
  signature covers (everything before and after the \texttt{/Contents}
  value)
\item
  \textbf{Incremental save}: The modified PDF is written as an
  incremental update --- appended to the original file, preserving
  previous revisions
\item
  \textbf{Signing callback}: PDFBox passes the byte range data (as an
  \texttt{InputStream}) to your \texttt{SignatureInterface.sign()}
  implementation
\item
  \textbf{Signature injection}: The returned CMS bytes are hex-encoded
  and written into the reserved \texttt{/Contents} space
\end{enumerate}

If the CMS signature is larger than the reserved space, the signing
fails. If it's smaller (which is normal), the remaining space is filled
with zeros --- the zero-padding that you see in the hex dump of any
signed PDF's \texttt{/Contents} entry.

This architecture means that \textbf{PDFBox does not construct the CMS
signature itself}. It provides the bytes to sign, and you provide the
signature. This is both a strength (maximum flexibility --- you can use
any CMS library, any algorithm, any certificate) and a burden (you must
get the CMS right yourself). That's why the next chapter is dedicated
entirely to BouncyCastle CMS construction.

\section{The Legion of the Bouncy
Castle}\label{_the_legion_of_the_bouncy_castle}

\subsection{Why BouncyCastle Exists}\label{_why_bouncycastle_exists_2}

Java has a built-in cryptographic provider --- the \texttt{SunJSSE} and
\texttt{SunJCE} providers that ship with every JDK. So why does almost
every Java PDF signing project use BouncyCastle instead?

Three reasons:

\begin{enumerate}
\def\labelenumi{\arabic{enumi}.}
\item
  \textbf{CMS support}: The JDK's built-in providers do not include a
  CMS (\texttt{SignedData}) implementation. They provide the raw
  cryptographic primitives (RSA, ECDSA, SHA-256, etc.), but not the
  higher-level CMS container format that PDF signatures require.
  BouncyCastle's \texttt{bcpkix} module provides
  \texttt{CMSSignedDataGenerator}, \texttt{CMSSignedData}, and all the
  ASN.1 infrastructure needed to construct and parse CMS structures.
\item
  \textbf{Algorithm breadth}: BouncyCastle supports algorithms that the
  JDK providers don't --- including EdDSA (Ed25519, Ed448), GOST, SM2,
  and various other national or emerging algorithms. While PDF
  signatures currently use RSA and ECDSA almost exclusively, the ability
  to work with less common algorithms is valuable for government and
  specialized applications.
\item
  \textbf{ASN.1 library}: BouncyCastle's ASN.1 library (\texttt{bcutil}
  and \texttt{bcprov}) is the de facto standard for Java ASN.1
  encoding/decoding. When you need to parse a certificate, construct a
  signed attribute, or decode an OCSP response, BouncyCastle's ASN.1
  classes are what you use.
\end{enumerate}

\subsection{The BouncyCastle
Ecosystem}\label{_the_bouncycastle_ecosystem}

BouncyCastle is not a single JAR --- it's a family of related libraries:

\begin{longtable}[]{@{}
  >{\raggedright\arraybackslash}p{(\linewidth - 4\tabcolsep) * \real{0.2500}}
  >{\raggedright\arraybackslash}p{(\linewidth - 4\tabcolsep) * \real{0.1250}}
  >{\raggedright\arraybackslash}p{(\linewidth - 4\tabcolsep) * \real{0.6250}}@{}}
\caption{BouncyCastle modules relevant to PDF signing}\tabularnewline
\toprule\noalign{}
\begin{minipage}[b]{\linewidth}\raggedright
Module
\end{minipage} & \begin{minipage}[b]{\linewidth}\raggedright
Maven Artifact
\end{minipage} & \begin{minipage}[b]{\linewidth}\raggedright
Purpose
\end{minipage} \\
\midrule\noalign{}
\endfirsthead
\toprule\noalign{}
\begin{minipage}[b]{\linewidth}\raggedright
Module
\end{minipage} & \begin{minipage}[b]{\linewidth}\raggedright
Maven Artifact
\end{minipage} & \begin{minipage}[b]{\linewidth}\raggedright
Purpose
\end{minipage} \\
\midrule\noalign{}
\endhead
\bottomrule\noalign{}
\endlastfoot
\textbf{Provider} & \texttt{bcprov-jdk18on} & Core cryptographic
provider: symmetric/asymmetric ciphers, hashes, key generation, random
number generation. Register via
\texttt{Security.addProvider(new\ BouncyCastleProvider())} \\
\textbf{PKIX/CMS/EAC} & \texttt{bcpkix-jdk18on} & CMS
\texttt{SignedData} construction/parsing, PKCS\#10 certificate requests,
OCSP, timestamp protocol (RFC 3161), certificate path validation.
\textbf{This is the module you need for PDF signing.} \\
\textbf{Utilities} & \texttt{bcutil-jdk18on} & PEM reader/writer, ASN.1
dump utilities, base64 encoding. Pulled in as a transitive
dependency. \\
\end{longtable}

You must explicitly register the BouncyCastle provider before using it.
If you forget this step, you'll get cryptic
\texttt{NoSuchAlgorithmException} errors at runtime.

\begin{Shaded}
\begin{Highlighting}[]
\KeywordTok{import} \ImportTok{org}\OperatorTok{.}\ImportTok{bouncycastle}\OperatorTok{.}\ImportTok{jce}\OperatorTok{.}\ImportTok{provider}\OperatorTok{.}\ImportTok{BouncyCastleProvider}\OperatorTok{;}
\KeywordTok{import} \ImportTok{java}\OperatorTok{.}\ImportTok{security}\OperatorTok{.}\ImportTok{Security}\OperatorTok{;}

\CommentTok{// Do this once at application startup}
\ControlFlowTok{if} \OperatorTok{(}\BuiltInTok{Security}\OperatorTok{.}\FunctionTok{getProvider}\OperatorTok{(}\StringTok{"BC"}\OperatorTok{)} \OperatorTok{==} \KeywordTok{null}\OperatorTok{)} \OperatorTok{\{}
    \BuiltInTok{Security}\OperatorTok{.}\FunctionTok{addProvider}\OperatorTok{(}\KeywordTok{new} \FunctionTok{BouncyCastleProvider}\OperatorTok{());}
\OperatorTok{\}}
\end{Highlighting}
\end{Shaded}

In Spring Boot applications, put this in a \texttt{@Configuration} class
with a \texttt{@PostConstruct} method. In standalone applications, put
it in your \texttt{main()} method. Forgetting this registration is the
\#1 cause of "it works in my IDE but fails in production" signing
issues.

\subsection{\texorpdfstring{BouncyCastle Version History and the
\texttt{-jdk18on}
Suffix}{BouncyCastle Version History and the -jdk18on Suffix}}\label{_bouncycastle_version_history_and_the_jdk18on_suffix}

BouncyCastle's artifact naming has changed over the years, which causes
dependency confusion:

\begin{itemize}
\item
  \texttt{bcprov-jdk15} --- for JDK 1.5 (ancient, do not use)
\item
  \texttt{bcprov-jdk15on} --- for JDK 1.5 and later (the long-lived
  artifact ID)
\item
  \texttt{bcprov-jdk18on} --- for JDK 1.8 (Java 8) and later (current,
  use this)
\end{itemize}

The \texttt{jdk18on} variants (introduced with BouncyCastle 1.71)
dropped support for Java versions before 8 and added Java module system
(\texttt{module-info.java}) support. If you're on Java 11 or later (and
for PDF signing with PDFBox 3.x, you should be on Java 17+), use the
\texttt{jdk18on} artifacts exclusively.

Do not mix \texttt{jdk15on} and \texttt{jdk18on} artifacts in the same
classpath. BouncyCastle uses a provider registration model, and having
two versions of the provider registered simultaneously causes
unpredictable failures --- the wrong version may handle a given
algorithm, producing incompatible ASN.1 structures.

\section{EU DSS: The European Commission's Signing
Framework}\label{_eu_dss_the_european_commissions_signing_framework}

\subsection{What EU DSS Is}\label{_what_eu_dss_is}

There is one more library that deserves its own section --- not because
it replaces PDFBox or BouncyCastle, but because it \emph{builds on top
of them} and adds a layer that would take months to implement yourself.

\textbf{EU DSS} (Digital Signature Services) is an open-source Java
framework developed and maintained by \textbf{Nowina Solutions} under
contract from the \textbf{European Commission}. It is the reference
implementation for creating and validating digital signatures that
comply with the European eIDAS regulation --- which means it implements
the full PAdES, CAdES, XAdES, and JAdES signature profiles, including
the complex long-term validation (LTV) workflows that individual
libraries like PDFBox cannot do alone.

The project lives at \texttt{github.com/esig/dss} and is released under
the \textbf{LGPL 2.1} license. The Maven coordinates follow the pattern
\texttt{eu.europa.ec.joinup.sd-dss:dss-*}.

\subsection{Why EU DSS Matters for PDF
Signatures}\label{_why_eu_dss_matters_for_pdf_signatures}

If your application needs to produce \textbf{PAdES-B-LTA} signatures ---
the gold standard for long-term archival signatures that remain
verifiable for decades --- you have two choices:

\begin{enumerate}
\def\labelenumi{\arabic{enumi}.}
\item
  \textbf{Build it yourself}: Use PDFBox for the PDF layer, BouncyCastle
  for CMS construction, write your own timestamp client (RFC 3161),
  implement OCSP and CRL fetching, construct the DSS (Document Security
  Store) dictionary, embed the VRI (Validation Related Information), and
  manage archival timestamp chains. This is what Chapters 29--31 of this
  book teach. It works. It is also hundreds of hours of development and
  testing.
\item
  \textbf{Use EU DSS}: Import the \texttt{dss-pades} module, configure a
  \texttt{PAdESSignatureParameters} object, provide a
  \texttt{SignatureTokenConnection} (which wraps your certificate and
  private key), and let the framework handle everything --- CMS
  construction, signed attribute inclusion, timestamp requests,
  revocation data embedding, DSS dictionary creation, and incremental
  save.
\end{enumerate}

Here's what a PAdES-B-LTA signature looks like with EU DSS:

\begin{Shaded}
\begin{Highlighting}[]
\CommentTok{// EU DSS — PAdES{-}B{-}LTA signing (simplified)}
\NormalTok{PAdESSignatureParameters params }\OperatorTok{=} \KeywordTok{new} \FunctionTok{PAdESSignatureParameters}\OperatorTok{();}
\NormalTok{params}\OperatorTok{.}\FunctionTok{setSignatureLevel}\OperatorTok{(}\NormalTok{SignatureLevel}\OperatorTok{.}\FunctionTok{PAdES\_BASELINE\_LTA}\OperatorTok{);}
\NormalTok{params}\OperatorTok{.}\FunctionTok{setDigestAlgorithm}\OperatorTok{(}\NormalTok{DigestAlgorithm}\OperatorTok{.}\FunctionTok{SHA256}\OperatorTok{);}
\NormalTok{params}\OperatorTok{.}\FunctionTok{setSigningCertificate}\OperatorTok{(}\KeywordTok{new} \FunctionTok{CertificateToken}\OperatorTok{(}\NormalTok{certificate}\OperatorTok{));}
\NormalTok{params}\OperatorTok{.}\FunctionTok{setCertificateChain}\OperatorTok{(}\BuiltInTok{Arrays}\OperatorTok{.}\FunctionTok{asList}\OperatorTok{(}\NormalTok{intermediateCerts}\OperatorTok{));}

\CommentTok{// TSA (Timestamp Authority) for the signature timestamp}
\NormalTok{OnlineTSPSource tspSource }\OperatorTok{=} \KeywordTok{new} \FunctionTok{OnlineTSPSource}\OperatorTok{(}\StringTok{"http://tsa.example.com/timestamp"}\OperatorTok{);}
\NormalTok{params}\OperatorTok{.}\FunctionTok{setContentTimestampParameters}\OperatorTok{(}\KeywordTok{new} \FunctionTok{TimestampParameters}\OperatorTok{(}\NormalTok{DigestAlgorithm}\OperatorTok{.}\FunctionTok{SHA256}\OperatorTok{));}

\NormalTok{PAdESService service }\OperatorTok{=} \KeywordTok{new} \FunctionTok{PAdESService}\OperatorTok{(}\NormalTok{commonCertificateVerifier}\OperatorTok{);}
\NormalTok{service}\OperatorTok{.}\FunctionTok{setTspSource}\OperatorTok{(}\NormalTok{tspSource}\OperatorTok{);}

\CommentTok{// Two{-}step signing: digest, then sign}
\NormalTok{ToBeSigned dataToSign }\OperatorTok{=}\NormalTok{ service}\OperatorTok{.}\FunctionTok{getDataToSign}\OperatorTok{(}\NormalTok{documentToSign}\OperatorTok{,}\NormalTok{ params}\OperatorTok{);}
\NormalTok{SignatureValue signatureValue }\OperatorTok{=}\NormalTok{ signingToken}\OperatorTok{.}\FunctionTok{sign}\OperatorTok{(}\NormalTok{dataToSign}\OperatorTok{,}\NormalTok{ params}\OperatorTok{.}\FunctionTok{getDigestAlgorithm}\OperatorTok{(),}\NormalTok{ privateKey}\OperatorTok{);}
\NormalTok{DSSDocument signedDocument }\OperatorTok{=}\NormalTok{ service}\OperatorTok{.}\FunctionTok{signDocument}\OperatorTok{(}\NormalTok{documentToSign}\OperatorTok{,}\NormalTok{ params}\OperatorTok{,}\NormalTok{ signatureValue}\OperatorTok{);}
\end{Highlighting}
\end{Shaded}

Compare this to implementing the same B-LTA profile manually with PDFBox
+ BouncyCastle --- the EU DSS version handles the timestamp embedding,
OCSP/CRL fetching, DSS dictionary construction, and VRI population
automatically.

\subsection{EU DSS Architecture}\label{_eu_dss_architecture}

EU DSS is modular. You don't need to import the entire framework ---
pick the modules relevant to your signature format:

\begin{longtable}[]{@{}
  >{\raggedright\arraybackslash}p{(\linewidth - 2\tabcolsep) * \real{0.2857}}
  >{\raggedright\arraybackslash}p{(\linewidth - 2\tabcolsep) * \real{0.7143}}@{}}
\caption{EU DSS modules relevant to PDF signing}\tabularnewline
\toprule\noalign{}
\begin{minipage}[b]{\linewidth}\raggedright
Module
\end{minipage} & \begin{minipage}[b]{\linewidth}\raggedright
Purpose
\end{minipage} \\
\midrule\noalign{}
\endfirsthead
\toprule\noalign{}
\begin{minipage}[b]{\linewidth}\raggedright
Module
\end{minipage} & \begin{minipage}[b]{\linewidth}\raggedright
Purpose
\end{minipage} \\
\midrule\noalign{}
\endhead
\bottomrule\noalign{}
\endlastfoot
\texttt{dss-pades} & PAdES signature creation and extension (B-B → B-T →
B-LT → B-LTA). Uses PDFBox internally for PDF manipulation. \\
\texttt{dss-cades} & CAdES signature creation. The PAdES module
delegates CMS construction to this. \\
\texttt{dss-document} & Core document model (\texttt{DSSDocument},
\texttt{InMemoryDocument}, \texttt{FileDocument}) and service
interfaces. \\
\texttt{dss-service} & Online data loaders --- HTTP clients for TSA,
OCSP, and CRL fetching. \\
\texttt{dss-tsl-validation} & EU Trusted List (EUTL) parsing and
validation. Essential for determining whether a signing certificate is
qualified under eIDAS. \\
\texttt{dss-validation-report} & Generates detailed validation reports
following ETSI TS 119 102-2 --- the same validation report format that
EU government validation services produce. \\
\texttt{dss-token} & \texttt{SignatureTokenConnection} implementations:
PKCS\#12 file, PKCS\#11 (HSM), JKS keystore, MS-CAPI (Windows
certificate store). \\
\end{longtable}

\subsection{EU DSS vs. Rolling Your
Own}\label{_eu_dss_vs_rolling_your_own_2}

The honest comparison:

\textbf{Advantages of EU DSS:}

\begin{itemize}
\item
  \textbf{Compliance out of the box}: EU DSS is the European
  Commission's reference implementation. If you use it correctly, your
  signatures will pass validation on the EU's official validation
  service (DSS demonstration webapp) and on member-state validation
  services. This is not a guarantee you get from hand-rolling with
  PDFBox + BouncyCastle.
\item
  \textbf{LTV complexity hidden}: Long-term validation requires fetching
  OCSP responses, downloading CRLs, embedding them in the PDF's DSS
  dictionary, creating VRI entries, and adding archival timestamps. EU
  DSS does all of this. Implementing it yourself is a significant
  effort.
\item
  \textbf{Trusted List integration}: The \texttt{dss-tsl-validation}
  module can load and parse the EU Trusted List, letting you
  programmatically determine whether a signing certificate is qualified.
  Building this yourself means implementing ETSI TS 119 612 XML parsing
  and certificate matching --- nontrivial.
\item
  \textbf{Validation reports}: EU DSS generates ETSI-compliant
  validation reports that are accepted by auditors and compliance
  officers.
\end{itemize}

\textbf{Advantages of rolling your own (PDFBox + BouncyCastle):}

\begin{itemize}
\item
  \textbf{Full control}: You understand every byte. When something
  fails, you know exactly where to look because you wrote every layer.
\item
  \textbf{No framework weight}: EU DSS pulls in a substantial dependency
  tree. If you only need PAdES-B-B (baseline signature without
  timestamps or LTV), you can do it with PDFBox + BouncyCastle in about
  200 lines of code --- EU DSS adds a megabytes of framework for
  something achievable with direct library calls.
\item
  \textbf{Learning}: Building it yourself teaches you the spec. Using EU
  DSS teaches you the framework's API. For a book like this one,
  understanding comes from building.
\item
  \textbf{Customization}: EU DSS makes assumptions about how signatures
  should be constructed. If your use case deviates (unusual SubFilter
  values, custom signed attributes, non-standard timestamp
  configurations), you may fight the framework rather than benefit from
  it.
\end{itemize}

\textbf{Our recommendation}: Learn to build signatures with PDFBox +
BouncyCastle first (Chapters 29--31). Then, when you need PAdES-B-LTA or
EU Trusted List integration in production, use EU DSS --- you'll
understand what it's doing under the hood, and you'll be able to debug
it when it behaves unexpectedly.

Chapter 23 provides a complete walkthrough of EU DSS for PAdES signing
and validation.

\subsection{EU DSS Version History}\label{_eu_dss_version_history}

EU DSS follows a rapid release cycle, with major versions roughly
annually:

\begin{itemize}
\item
  \textbf{DSS 5.x} (2017--2022): The long-lived stable series. Most
  documentation and blog posts you'll find online reference 5.x APIs.
\item
  \textbf{DSS 6.x} (2023--present): The current generation, with Java
  17+ as the minimum, updated dependencies (including PDFBox 3.x for the
  PAdES module), and improved validation reporting. The API changed
  significantly between 5.x and 6.x --- if you're following tutorials,
  make sure they match your version.
\end{itemize}

EU DSS 6.x uses PDFBox 3.x internally. If your project also uses PDFBox
directly (for example, for document manipulation before signing), you
must ensure version alignment. Having DSS pull in PDFBox 3.0.0 while
your code uses PDFBox 3.0.1 can cause subtle \texttt{NoSuchMethodError}
exceptions at runtime. Use Maven's
\texttt{\textless{}dependencyManagement\textgreater{}} to lock the
PDFBox version across your entire project.

\section{Analysis and Debugging
Tools}\label{_analysis_and_debugging_tools}

A signing implementation is only as good as your ability to inspect what
it produces. The following tools are indispensable for PDF signature
development.

\subsection{RUPS: The PDF Object
Explorer}\label{_rups_the_pdf_object_explorer}

\textbf{RUPS} (Reading and Updating PDF Syntax) is a Java Swing
application that lets you browse the internal object structure of a PDF
file. It was originally developed by Bruno Lowagie as a companion to
iText, and it shows every COS object, dictionary key, stream, and
cross-reference entry in the file.

For signature work, RUPS is invaluable because it lets you:

\begin{itemize}
\item
  Inspect the \texttt{/Sig} dictionary and all its keys
  (\texttt{/ByteRange}, \texttt{/Contents}, \texttt{/SubFilter},
  \texttt{/Filter})
\item
  Navigate the field tree to see how signature fields are structured
\item
  View the \texttt{/AcroForm} dictionary and its \texttt{/SigFlags}
\item
  Examine \texttt{/Lock} and \texttt{/SV} (Seed Value) dictionaries
\item
  Browse incremental updates to see what changed between revisions
\item
  Inspect the raw byte content of the \texttt{/Contents} hex string
\end{itemize}

RUPS is available on GitHub at \texttt{LibrePDF/OpenPDF} (it's bundled
with the OpenPDF project). It's also available as a standalone download.
The interface is spartan --- it's a developer tool, not a user-facing
product --- but once you learn to navigate the tree view, it becomes the
fastest way to understand what's actually inside a PDF.

When debugging a signature failure, the first thing you should do is
open the PDF in RUPS and check:

\begin{enumerate}
\def\labelenumi{\arabic{enumi}.}
\item
  Does the \texttt{/Sig} dictionary exist? (Navigate: Root → AcroForm →
  Fields → {[}signature field{]} → V)
\item
  Is the \texttt{/ByteRange} array plausible? (Four integers, the first
  should be 0, the last should be close to the file size)
\item
  Is \texttt{/Contents} non-empty? (Should be a long hex string, not all
  zeros)
\item
  What is the \texttt{/SubFilter}? (Should be
  \texttt{adbe.pkcs7.detached} or \texttt{ETSI.CAdES.detached})
\end{enumerate}

If any of these are wrong, you've found your problem before even looking
at the CMS.

\subsection{veraPDF: The Compliance
Validator}\label{_verapdf_the_compliance_validator}

\textbf{veraPDF} is the industry-standard PDF/A validation tool,
developed by the Open Preservation Foundation and the PDF Association.
While its primary mission is PDF/A compliance checking, it includes PDF
analysis capabilities that are useful for signature work:

\begin{itemize}
\item
  \textbf{Policy-based validation} --- you can define custom validation
  profiles that check signature-specific requirements
\item
  \textbf{Feature extraction} --- enumerate all signatures,
  certificates, and cryptographic properties in a PDF
\item
  \textbf{Structural analysis} --- verify that the PDF structure
  conforms to the specification
\end{itemize}

veraPDF is available as a CLI tool, a Java library (which you can embed
in your validation pipeline), and a GUI application. It's the validator
that national archives and government agencies use to verify PDF/A
compliance --- if your signed PDFs also need to be PDF/A compliant (as
many government document workflows require), veraPDF is essential.

\subsection{QPDF: The Swiss Army
Knife}\label{_qpdf_the_swiss_army_knife}

\textbf{QPDF} is a C++ command-line tool (created by Jay Berkenbilt)
that can inspect, transform, and repair PDF files at the structural
level. It's not Java, but every Java PDF developer should have it
installed.

Key capabilities for signature work:

\begin{itemize}
\item
  \texttt{qpdf\ -\/-show-xref} --- displays the cross-reference table,
  showing where every object lives in the file. Critical for
  understanding byte range issues.
\item
  \texttt{qpdf\ -\/-show-object=N} --- displays a specific object by
  number. Faster than navigating a GUI when you know the object number
  from an error message.
\item
  \texttt{qpdf\ -\/-check} --- validates the PDF's structural integrity.
  If this reports problems, your signing code has a bug.
\item
  \texttt{qpdf\ -\/-qdf} --- produces a "QDF" (QPDF data format) output
  where the PDF is restructured for human readability, with objects in
  order and comments showing object numbers. Invaluable for
  understanding what incremental saves changed.
\item
  \texttt{qpdf\ -\/-linearize} --- linearizes the PDF for fast web
  viewing. \textbf{Do not linearize after signing} --- it destroys the
  signature by reorganizing the file.
\end{itemize}

Several QPDF operations will destroy existing signatures:

\begin{itemize}
\item
  \texttt{-\/-linearize} --- reorganizes the entire file
\item
  \texttt{-\/-normalize-content=y} --- reprocesses content streams
\item
  \texttt{-\/-object-streams=generate} --- restructures object storage
\item
  Any operation that produces a non-incremental output
\end{itemize}

Always use \texttt{-\/-copy-encryption=source.pdf} and check whether the
operation preserves incremental structure before running it on a signed
PDF.

\subsection{MuTool (from MuPDF)}\label{_mutool_from_mupdf}

\textbf{MuTool} is the command-line companion to MuPDF (Artifex's
open-source PDF renderer). Its \texttt{mutool\ show} command provides
quick access to PDF objects:

\begin{Shaded}
\begin{Highlighting}[]
\ExtensionTok{mutool}\NormalTok{ show document.pdf trailer          }\CommentTok{\# Show trailer dictionary}
\ExtensionTok{mutool}\NormalTok{ show document.pdf xref             }\CommentTok{\# Show cross{-}reference table}
\ExtensionTok{mutool}\NormalTok{ show document.pdf 42 0             }\CommentTok{\# Show object 42, generation 0}
\ExtensionTok{mutool}\NormalTok{ show document.pdf outline          }\CommentTok{\# Show document outline/bookmarks}
\end{Highlighting}
\end{Shaded}

MuTool is lightweight and fast --- useful for quick inspections when you
don't want to launch a GUI tool.

\subsection{OpenSSL: CMS Verification Outside the PDF
Context}\label{_openssl_cms_verification_outside_the_pdf_context}

For verifying the CMS signature independently of the PDF wrapper,
OpenSSL is the standard tool:

\begin{Shaded}
\begin{Highlighting}[]
\CommentTok{\# Extract the raw signature bytes from /Contents (after hex{-}decoding)}
\CommentTok{\# and save as signature.der}

\CommentTok{\# Verify the CMS structure}
\ExtensionTok{openssl}\NormalTok{ cms }\AttributeTok{{-}inform}\NormalTok{ DER }\AttributeTok{{-}in}\NormalTok{ signature.der }\AttributeTok{{-}noout} \AttributeTok{{-}cmsout} \AttributeTok{{-}print}

\CommentTok{\# Verify the signature against the signed data}
\ExtensionTok{openssl}\NormalTok{ cms }\AttributeTok{{-}inform}\NormalTok{ DER }\AttributeTok{{-}in}\NormalTok{ signature.der }\AttributeTok{{-}verify} \DataTypeTok{\textbackslash{}}
    \AttributeTok{{-}content}\NormalTok{ signed{-}data.bin }\DataTypeTok{\textbackslash{}}
    \AttributeTok{{-}CAfile}\NormalTok{ ca{-}bundle.pem }\DataTypeTok{\textbackslash{}}
    \AttributeTok{{-}out}\NormalTok{ /dev/null}
\end{Highlighting}
\end{Shaded}

This separates the CMS verification from the PDF byte-range logic --- if
OpenSSL says the CMS is valid but your PDF validator says the signature
is invalid, the problem is in the byte range handling, not the
cryptography.

\subsection{ASN.1 Dump Tools}\label{_asn_1_dump_tools}

When debugging CMS issues, being able to read the raw ASN.1 structure is
essential. Several tools help:

\begin{itemize}
\item
  \textbf{\texttt{openssl\ asn1parse}} --- parses DER-encoded ASN.1,
  showing the tag-length-value tree
\item
  \textbf{BouncyCastle's \texttt{ASN1Dump}} ---
  \texttt{org.bouncycastle.asn1.util.ASN1Dump.dumpAsString(obj)} ---
  prints any ASN.1 object as a readable tree from Java code
\item
  \textbf{\texttt{dumpasn1}} (Peter Gutmann) --- the original ASN.1
  dumper, still the most detailed output
\item
  \textbf{Web-based tools} --- several websites parse ASN.1/DER online;
  useful for quick checks but do not upload production signatures to
  third-party websites
\end{itemize}

\section{Choosing the Right Library for Your
Project}\label{_choosing_the_right_library_for_your_project}

After all this history, here is the practical guidance:

\subsection{Decision Matrix}\label{_decision_matrix}

\begin{longtable}[]{@{}
  >{\raggedright\arraybackslash}p{(\linewidth - 10\tabcolsep) * \real{0.2308}}
  >{\raggedright\arraybackslash}p{(\linewidth - 10\tabcolsep) * \real{0.1538}}
  >{\raggedright\arraybackslash}p{(\linewidth - 10\tabcolsep) * \real{0.1538}}
  >{\raggedright\arraybackslash}p{(\linewidth - 10\tabcolsep) * \real{0.1538}}
  >{\raggedright\arraybackslash}p{(\linewidth - 10\tabcolsep) * \real{0.1538}}
  >{\raggedright\arraybackslash}p{(\linewidth - 10\tabcolsep) * \real{0.1538}}@{}}
\caption{Library selection for PDF signature projects}\tabularnewline
\toprule\noalign{}
\begin{minipage}[b]{\linewidth}\raggedright
Requirement
\end{minipage} & \begin{minipage}[b]{\linewidth}\raggedright
PDFBox + BC
\end{minipage} & \begin{minipage}[b]{\linewidth}\raggedright
EU DSS
\end{minipage} & \begin{minipage}[b]{\linewidth}\raggedright
OpenPDF
\end{minipage} & \begin{minipage}[b]{\linewidth}\raggedright
iText 7/8 (AGPL)
\end{minipage} & \begin{minipage}[b]{\linewidth}\raggedright
iText 7/8 (Commercial)
\end{minipage} \\
\midrule\noalign{}
\endfirsthead
\toprule\noalign{}
\begin{minipage}[b]{\linewidth}\raggedright
Requirement
\end{minipage} & \begin{minipage}[b]{\linewidth}\raggedright
PDFBox + BC
\end{minipage} & \begin{minipage}[b]{\linewidth}\raggedright
EU DSS
\end{minipage} & \begin{minipage}[b]{\linewidth}\raggedright
OpenPDF
\end{minipage} & \begin{minipage}[b]{\linewidth}\raggedright
iText 7/8 (AGPL)
\end{minipage} & \begin{minipage}[b]{\linewidth}\raggedright
iText 7/8 (Commercial)
\end{minipage} \\
\midrule\noalign{}
\endhead
\bottomrule\noalign{}
\endlastfoot
\textbf{License} & Apache 2.0 & LGPL 2.1 & LGPL 2.1 & AGPL 3.0 &
Proprietary \\
\textbf{Cost} & Free & Free & Free & Free (if AGPL-compatible) & €€€€ \\
\textbf{PAdES-B-B} & Yes (manual CMS) & Yes (built-in) & Basic & Yes
(built-in) & Yes (built-in) \\
\textbf{PAdES-B-LTA} & Yes (manual DSS) & Yes (built-in) & No & Yes
(built-in) & Yes (built-in) \\
\textbf{EUTL / eIDAS validation} & No (manual) & Yes (built-in) & No &
Partial & Yes \\
\textbf{PDF 2.0} & Yes & Yes (via PDFBox) & No & Yes & Yes \\
\textbf{Community size} & Large (Apache) & Medium (EU-backed) & Small &
N/A (AGPL) & Vendor-supported \\
\textbf{Signing API maturity} & High & Very high & Moderate & High &
High \\
\textbf{Low-level PDF access} & Excellent & Limited (framework) &
Limited & Good & Good \\
\textbf{Our recommendation} & \textbf{Learn + custom} &
\textbf{Production PAdES} & Legacy migration & Open-source projects &
Budget permits \\
\end{longtable}

\subsection{Our Recommendation}\label{_our_recommendation}

For the projects this book targets --- production signing services,
validation pipelines, signature analysis tools --- \textbf{Apache PDFBox
3.x with BouncyCastle} is the default recommendation:

\begin{enumerate}
\def\labelenumi{\arabic{enumi}.}
\item
  \textbf{Apache 2.0 license}: No surprises, no AGPL obligations, no
  license fees. This will not change.
\item
  \textbf{Low-level access}: Signing requires understanding and
  manipulating the raw PDF structure. PDFBox's COS model gives you
  direct access. iText and OpenPDF abstract this away, which is
  convenient until you need to debug a byte-range computation or inspect
  a malformed signature dictionary.
\item
  \textbf{Active maintenance}: PDFBox has a dedicated Apache community,
  regular releases, and responsive handling of security issues. The
  Bochum attacks were patched quickly.
\item
  \textbf{Separation of concerns}: PDFBox handles the PDF, BouncyCastle
  handles the crypto. You understand both layers because you built both
  layers. When something breaks, you know exactly where to look.
\end{enumerate}

The trade-off is that PDFBox makes you build more yourself ---
especially the CMS layer. That's what Chapters 29--31 are for.

\textbf{The best way to learn is to build it yourself.} Using a library
that does everything for you (like commercial iText) gets you to
production faster, but leaves you helpless when something goes wrong.
With PDFBox + BouncyCastle, you understand every byte of the signature
because you put it there. When a customer reports that their signature
"shows invalid in Adobe but valid in Foxit," you'll know exactly where
to look --- because you wrote the code that constructed the CMS,
computed the byte range, and injected the signature.

\section{The Ecosystem at a Glance}\label{_the_ecosystem_at_a_glance}

To close this chapter, here is a summary of every tool and library we've
discussed, organized by function:

\begin{longtable}[]{@{}
  >{\raggedright\arraybackslash}p{(\linewidth - 6\tabcolsep) * \real{0.2000}}
  >{\raggedright\arraybackslash}p{(\linewidth - 6\tabcolsep) * \real{0.2000}}
  >{\raggedright\arraybackslash}p{(\linewidth - 6\tabcolsep) * \real{0.2000}}
  >{\raggedright\arraybackslash}p{(\linewidth - 6\tabcolsep) * \real{0.4000}}@{}}
\caption{The Java PDF signature ecosystem}\tabularnewline
\toprule\noalign{}
\begin{minipage}[b]{\linewidth}\raggedright
Function
\end{minipage} & \begin{minipage}[b]{\linewidth}\raggedright
Tool/Library
\end{minipage} & \begin{minipage}[b]{\linewidth}\raggedright
License
\end{minipage} & \begin{minipage}[b]{\linewidth}\raggedright
Notes
\end{minipage} \\
\midrule\noalign{}
\endfirsthead
\toprule\noalign{}
\begin{minipage}[b]{\linewidth}\raggedright
Function
\end{minipage} & \begin{minipage}[b]{\linewidth}\raggedright
Tool/Library
\end{minipage} & \begin{minipage}[b]{\linewidth}\raggedright
License
\end{minipage} & \begin{minipage}[b]{\linewidth}\raggedright
Notes
\end{minipage} \\
\midrule\noalign{}
\endhead
\bottomrule\noalign{}
\endlastfoot
PDF manipulation (read/write) & Apache PDFBox 3.x & Apache 2.0 & Default
choice for this book. Reader-first architecture. \\
PDF manipulation (read/write) & OpenPDF 2.x & LGPL 2.1 & Community fork
of iText 2.x. Good for legacy migration. \\
PDF manipulation (read/write) & iText 7/8 & AGPL / Commercial &
Feature-rich but expensive or license-restrictive. \\
CMS / Cryptography & BouncyCastle 1.80+ & MIT (Bouncy Castle License) &
De facto standard for Java CMS. Register provider explicitly. \\
PAdES signing framework & EU DSS 6.x & LGPL 2.1 & European Commission
reference implementation. PAdES-B-B through B-LTA, validation reports,
EU Trusted List integration. \\
PDF structure inspection & RUPS & LGPL & GUI object browser. First tool
to reach for. \\
PDF/A compliance & veraPDF & Open source (MPL 2.0) & Industry-standard
validator. Also useful for general PDF analysis. \\
PDF structural tools & QPDF & Apache 2.0 & CLI Swiss army knife.
\texttt{-\/-show-xref}, \texttt{-\/-qdf}, \texttt{-\/-check}. \\
PDF rendering/inspection & MuTool (MuPDF) & AGPL & Lightweight CLI.
Quick object inspection. \\
CMS verification & OpenSSL & Apache 2.0 & CMS verification independent
of PDF context. \\
ASN.1 inspection & BouncyCastle \texttt{ASN1Dump} & MIT & Java-based
ASN.1 tree dump. \\
ASN.1 inspection & \texttt{dumpasn1} (Gutmann) & Public domain & The
classic ASN.1 dumper. \\
\end{longtable}

Every tool in this table will appear at some point in the remaining
chapters. Install them now --- when you're debugging a signature failure
at 2 AM, you don't want to be downloading tools.

\section{Key Takeaways}\label{_key_takeaways_15}

\begin{itemize}
\item
  \textbf{iText's license change} (LGPL → AGPL in 2009) split the Java
  PDF world in two. The consequences --- old versions in production,
  community forks, enterprise migrations --- are still visible in 2026.
\item
  \textbf{OpenPDF} is the LGPL continuation of iText 2.x. Useful for
  legacy migration but limited for modern PAdES signing.
\item
  \textbf{PDFBox + BouncyCastle} is the recommended stack for new
  signing projects. Apache-licensed, actively maintained, gives you
  low-level control.
\item
  \textbf{PDFBox 3.x} broke API compatibility with 2.x. Budget migration
  time. Key changes: \texttt{Loader.loadPDF()}, new font constructors,
  content stream API changes.
\item
  \textbf{BouncyCastle} provides the CMS layer that PDFBox doesn't.
  Always register the provider explicitly. Use \texttt{jdk18on}
  artifacts.
\item
  \textbf{RUPS, QPDF, veraPDF, OpenSSL} --- know your tools. PDF
  signature debugging is 50\% knowing what to look at and 50\% having
  the right tool to look at it.
\item
  \textbf{The library you choose is a 10-year decision.} Choose based on
  license stability, community health, and technical fit --- not just
  today's feature checklist.
\end{itemize}

\chapter{PDFBox Signing --- Complete
Walkthrough}\label{_pdfbox_signingcomplete_walkthrough}

Chapter 28 gave us the lay of the land. Now we write code. This chapter
is a complete, production-quality walkthrough of PDF signing with Apache
PDFBox 3.0.0 and BouncyCastle 1.80 on Java 21. Every code listing
compiles, every class is self-contained, and every pitfall we describe
is one we have hit in production --- or watched someone else hit at 2 AM
on a Friday.

We start with the core PDFBox signing
API --- \texttt{PDDocument.addSignature()}, \texttt{SignatureInterface},
and \texttt{SignatureOptions} --- then build outward: the
\texttt{PDSignature} dictionary and its fields, visible and invisible
signature appearances, multi-signature workflows, external/HSM signing,
certification signatures with DocMDP, timestamp embedding, and the
migration traps waiting for anyone upgrading from PDFBox 2.x to 3.x.

By the end of this chapter, you will have a complete, production-ready
signing class that produces PAdES-B-B and PAdES-B-T compliant
signatures. Chapter 30 goes deeper into the BouncyCastle CMS
construction that powers the actual cryptography.

\section{The PDFBox Signing API: Three Moving
Parts}\label{_the_pdfbox_signing_api_three_moving_parts}

PDFBox signing revolves around three objects that you wire together
before calling \texttt{saveIncremental()}. Understanding what each one
does --- and what happens when you get the order wrong --- is essential.

\subsection{PDSignature: The Signature
Dictionary}\label{_pdsignature_the_signature_dictionary}

\texttt{PDSignature} is PDFBox's Java representation of the PDF
\texttt{/Sig} dictionary (ISO 32000-2:2020 S12.8.1). When you create a
\texttt{PDSignature} and set its properties, you are building the
dictionary that will be written into the PDF file.

\begin{Shaded}
\begin{Highlighting}[]
\NormalTok{PDSignature signature }\OperatorTok{=} \KeywordTok{new} \FunctionTok{PDSignature}\OperatorTok{();}
\NormalTok{signature}\OperatorTok{.}\FunctionTok{setFilter}\OperatorTok{(}\NormalTok{PDSignature}\OperatorTok{.}\FunctionTok{FILTER\_ADOBE\_PPKLITE}\OperatorTok{);}       \CommentTok{// /Filter}
\NormalTok{signature}\OperatorTok{.}\FunctionTok{setSubFilter}\OperatorTok{(}\NormalTok{COSName}\OperatorTok{.}\FunctionTok{getPDFName}\OperatorTok{(}\StringTok{"ETSI.CAdES.detached"}\OperatorTok{));} \CommentTok{// /SubFilter}
\NormalTok{signature}\OperatorTok{.}\FunctionTok{setSignDate}\OperatorTok{(}\BuiltInTok{Calendar}\OperatorTok{.}\FunctionTok{getInstance}\OperatorTok{());}                \CommentTok{// /M}
\NormalTok{signature}\OperatorTok{.}\FunctionTok{setName}\OperatorTok{(}\StringTok{"Alice Smith"}\OperatorTok{);}                             \CommentTok{// /Name}
\NormalTok{signature}\OperatorTok{.}\FunctionTok{setReason}\OperatorTok{(}\StringTok{"I approve this document"}\OperatorTok{);}               \CommentTok{// /Reason}
\NormalTok{signature}\OperatorTok{.}\FunctionTok{setLocation}\OperatorTok{(}\StringTok{"Berlin, Germany"}\OperatorTok{);}                     \CommentTok{// /Location}
\NormalTok{signature}\OperatorTok{.}\FunctionTok{setContactInfo}\OperatorTok{(}\StringTok{"alice@example.com"}\OperatorTok{);}                \CommentTok{// /ContactInfo}
\end{Highlighting}
\end{Shaded}

Every setter maps directly to a key in the \texttt{/Sig} dictionary. The
following table lists all the fields you can set, what they map to, and
when you should care.

\begin{longtable}[]{@{}
  >{\raggedright\arraybackslash}p{(\linewidth - 6\tabcolsep) * \real{0.2308}}
  >{\raggedright\arraybackslash}p{(\linewidth - 6\tabcolsep) * \real{0.1538}}
  >{\raggedright\arraybackslash}p{(\linewidth - 6\tabcolsep) * \real{0.1538}}
  >{\raggedright\arraybackslash}p{(\linewidth - 6\tabcolsep) * \real{0.4615}}@{}}
\caption{PDSignature Fields --- ISO 32000-2:2020 S12.8.1, Table
255}\tabularnewline
\toprule\noalign{}
\begin{minipage}[b]{\linewidth}\raggedright
Java Method
\end{minipage} & \begin{minipage}[b]{\linewidth}\raggedright
PDF Key
\end{minipage} & \begin{minipage}[b]{\linewidth}\raggedright
Required
\end{minipage} & \begin{minipage}[b]{\linewidth}\raggedright
What It Does
\end{minipage} \\
\midrule\noalign{}
\endfirsthead
\toprule\noalign{}
\begin{minipage}[b]{\linewidth}\raggedright
Java Method
\end{minipage} & \begin{minipage}[b]{\linewidth}\raggedright
PDF Key
\end{minipage} & \begin{minipage}[b]{\linewidth}\raggedright
Required
\end{minipage} & \begin{minipage}[b]{\linewidth}\raggedright
What It Does
\end{minipage} \\
\midrule\noalign{}
\endhead
\bottomrule\noalign{}
\endlastfoot
\texttt{setFilter(COSName)} & \texttt{/Filter} & Yes & The signature
handler. Almost always \texttt{Adobe.PPKLite}. PDFBox provides
\texttt{FILTER\_ADOBE\_PPKLITE} as a constant. Other values
(\texttt{Entrust.PPKEF}, \texttt{CICI.SignIt}, \texttt{VeriSign.PPKVS})
exist in the spec but are rarely seen in the wild. \\
\texttt{setSubFilter(COSName)} & \texttt{/SubFilter} & Yes & The
encoding format of the signature value. For PAdES, use
\texttt{ETSI.CAdES.detached}. For legacy compatibility,
\texttt{adbe.pkcs7.detached}. The deprecated \texttt{adbe.pkcs7.sha1}
embeds the document hash as CMS content --- do not use it in new
code. \\
\texttt{setSignDate(Calendar)} & \texttt{/M} & No & Claimed signing
time. This is \emph{not} the authoritative timestamp --- it is the
signer's local clock. A TSA timestamp (PAdES-B-T) provides cryptographic
proof of time. Set it anyway, because viewers display it when no TSA
timestamp exists. \\
\texttt{setName(String)} & \texttt{/Name} & No & Signer's name as a
display string. PAdES validators typically extract the name from the
signing certificate's CN instead. Still useful for the human reader. \\
\texttt{setReason(String)} & \texttt{/Reason} & No & Why the document
was signed. Free text. Some workflows require specific values
("Approved", "Reviewed"). \\
\texttt{setLocation(String)} & \texttt{/Location} & No & Where the
document was signed. Free text. "Berlin, Germany" or "Server Room
3" --- whatever is meaningful. \\
\texttt{setContactInfo(String)} & \texttt{/ContactInfo} & No & Contact
information for the signer. Rarely used in practice, but mandated by
some government workflows. \\
\texttt{setByteRange(int{[}{]})} & \texttt{/ByteRange} & Auto & The byte
ranges of the signed content. \emph{You do not set this yourself.}
PDFBox computes it during \texttt{saveIncremental()}. If you manually
set it, you will break the signature. \\
\end{longtable}

Do not call \texttt{setByteRange()} manually. PDFBox calculates the byte
range during the incremental save process. If you set it yourself, the
byte range will not match the actual file offsets after writing, and the
signature will be invalid. We have seen production code that does
this --- it always fails.

\subsection{SignatureInterface: The Cryptographic
Callback}\label{_signatureinterface_the_cryptographic_callback}

\texttt{SignatureInterface} is a single-method interface:

\begin{Shaded}
\begin{Highlighting}[]
\KeywordTok{public} \KeywordTok{interface}\NormalTok{ SignatureInterface }\OperatorTok{\{}
    \DataTypeTok{byte}\OperatorTok{[]} \FunctionTok{sign}\OperatorTok{(}\BuiltInTok{InputStream}\NormalTok{ content}\OperatorTok{)} \KeywordTok{throws} \BuiltInTok{IOException}\OperatorTok{;}
\OperatorTok{\}}
\end{Highlighting}
\end{Shaded}

PDFBox calls your \texttt{sign()} method during
\texttt{saveIncremental()}. The \texttt{InputStream} parameter contains
the bytes identified by the \texttt{/ByteRange} --- that is, the entire
PDF file \emph{minus} the gap reserved for the \texttt{/Contents} hex
string. Your job is to read those bytes, produce a CMS
\texttt{SignedData} structure (or whatever your \texttt{/SubFilter}
requires), and return the DER-encoded result.

The returned \texttt{byte{[}{]}} is hex-encoded and written into the
\texttt{/Contents} field of the signature dictionary. If the returned
bytes are larger than the space reserved for \texttt{/Contents}, PDFBox
throws an exception. If they are smaller, the remaining space is
zero-padded.

The \texttt{InputStream} you receive in \texttt{sign()} is the
concatenation of the two byte ranges --- the bytes before the
\texttt{/Contents} gap and the bytes after it. You must read \emph{all}
of them with \texttt{content.readAllBytes()} (or equivalent). If you
only read part of the stream, you are signing incomplete data, and the
signature will fail validation.

Here is the simplest possible implementation --- it signs the content
with RSA/SHA-256 and produces a basic CMS \texttt{SignedData}:

\begin{Shaded}
\begin{Highlighting}[]
\KeywordTok{public} \KeywordTok{class}\NormalTok{ SimpleCmsSignatureInterface }\KeywordTok{implements}\NormalTok{ SignatureInterface }\OperatorTok{\{}

    \KeywordTok{private} \DataTypeTok{final} \BuiltInTok{PrivateKey}\NormalTok{ privateKey}\OperatorTok{;}
    \KeywordTok{private} \DataTypeTok{final} \BuiltInTok{X509Certificate}\NormalTok{ certificate}\OperatorTok{;}
    \KeywordTok{private} \DataTypeTok{final} \BuiltInTok{List}\OperatorTok{\textless{}}\BuiltInTok{X509Certificate}\OperatorTok{\textgreater{}}\NormalTok{ chain}\OperatorTok{;}

    \KeywordTok{public} \FunctionTok{SimpleCmsSignatureInterface}\OperatorTok{(}\BuiltInTok{PrivateKey}\NormalTok{ privateKey}\OperatorTok{,}
                                        \BuiltInTok{X509Certificate}\NormalTok{ certificate}\OperatorTok{,}
                                        \BuiltInTok{List}\OperatorTok{\textless{}}\BuiltInTok{X509Certificate}\OperatorTok{\textgreater{}}\NormalTok{ chain}\OperatorTok{)} \OperatorTok{\{}
        \KeywordTok{this}\OperatorTok{.}\FunctionTok{privateKey} \OperatorTok{=}\NormalTok{ privateKey}\OperatorTok{;}
        \KeywordTok{this}\OperatorTok{.}\FunctionTok{certificate} \OperatorTok{=}\NormalTok{ certificate}\OperatorTok{;}
        \KeywordTok{this}\OperatorTok{.}\FunctionTok{chain} \OperatorTok{=}\NormalTok{ chain}\OperatorTok{;}
    \OperatorTok{\}}

    \AttributeTok{@Override}
    \KeywordTok{public} \DataTypeTok{byte}\OperatorTok{[]} \FunctionTok{sign}\OperatorTok{(}\BuiltInTok{InputStream}\NormalTok{ content}\OperatorTok{)} \KeywordTok{throws} \BuiltInTok{IOException} \OperatorTok{\{}
        \ControlFlowTok{try} \OperatorTok{\{}
            \CommentTok{// Read all byte{-}range content}
            \DataTypeTok{byte}\OperatorTok{[]}\NormalTok{ contentBytes }\OperatorTok{=}\NormalTok{ content}\OperatorTok{.}\FunctionTok{readAllBytes}\OperatorTok{();}

            \CommentTok{// Build the CMS SignedData (Chapter 30 covers this in full detail)}
\NormalTok{            CMSSignedDataGenerator gen }\OperatorTok{=} \KeywordTok{new} \FunctionTok{CMSSignedDataGenerator}\OperatorTok{();}

\NormalTok{            ContentSigner signer }\OperatorTok{=} \KeywordTok{new} \FunctionTok{JcaContentSignerBuilder}\OperatorTok{(}\StringTok{"SHA256withRSA"}\OperatorTok{)}
                    \OperatorTok{.}\FunctionTok{setProvider}\OperatorTok{(}\StringTok{"BC"}\OperatorTok{)}
                    \OperatorTok{.}\FunctionTok{build}\OperatorTok{(}\NormalTok{privateKey}\OperatorTok{);}

\NormalTok{            DigestCalculatorProvider dcProvider }\OperatorTok{=}
                    \KeywordTok{new} \FunctionTok{JcaDigestCalculatorProviderBuilder}\OperatorTok{()}
                            \OperatorTok{.}\FunctionTok{setProvider}\OperatorTok{(}\StringTok{"BC"}\OperatorTok{)}
                            \OperatorTok{.}\FunctionTok{build}\OperatorTok{();}

\NormalTok{            gen}\OperatorTok{.}\FunctionTok{addSignerInfoGenerator}\OperatorTok{(}
                    \KeywordTok{new} \FunctionTok{JcaSignerInfoGeneratorBuilder}\OperatorTok{(}\NormalTok{dcProvider}\OperatorTok{)}
                            \OperatorTok{.}\FunctionTok{build}\OperatorTok{(}\NormalTok{signer}\OperatorTok{,}\NormalTok{ certificate}\OperatorTok{));}

\NormalTok{            gen}\OperatorTok{.}\FunctionTok{addCertificates}\OperatorTok{(}\KeywordTok{new} \FunctionTok{JcaCertStore}\OperatorTok{(}\NormalTok{chain}\OperatorTok{));}

\NormalTok{            CMSProcessableByteArray cmsContent }\OperatorTok{=}
                    \KeywordTok{new} \FunctionTok{CMSProcessableByteArray}\OperatorTok{(}\NormalTok{contentBytes}\OperatorTok{);}
\NormalTok{            CMSSignedData signedData }\OperatorTok{=}\NormalTok{ gen}\OperatorTok{.}\FunctionTok{generate}\OperatorTok{(}\NormalTok{cmsContent}\OperatorTok{,} \KeywordTok{false}\OperatorTok{);} 

            \ControlFlowTok{return}\NormalTok{ signedData}\OperatorTok{.}\FunctionTok{getEncoded}\OperatorTok{();}
        \OperatorTok{\}} \ControlFlowTok{catch} \OperatorTok{(}\BuiltInTok{Exception}\NormalTok{ e}\OperatorTok{)} \OperatorTok{\{}
            \ControlFlowTok{throw} \KeywordTok{new} \BuiltInTok{IOException}\OperatorTok{(}\StringTok{"CMS signing failed: "} \OperatorTok{+}\NormalTok{ e}\OperatorTok{.}\FunctionTok{getMessage}\OperatorTok{(),}\NormalTok{ e}\OperatorTok{);}
        \OperatorTok{\}}
    \OperatorTok{\}}
\OperatorTok{\}}
\end{Highlighting}
\end{Shaded}

\begin{itemize}
\item
  The \texttt{false} parameter means "detached signature" --- the
  content is not embedded inside the CMS structure. This is mandatory
  for PDF signatures: the content lives in the PDF file itself, not
  inside the CMS blob.
\end{itemize}

\subsection{SignatureOptions: Container Size and Visual
Appearance}\label{_signatureoptions_container_size_and_visual_appearance}

\texttt{SignatureOptions} controls two things: the size of the
\texttt{/Contents} container (how many bytes are reserved in the PDF for
the CMS signature) and the visual appearance stream.

\begin{Shaded}
\begin{Highlighting}[]
\NormalTok{SignatureOptions sigOptions }\OperatorTok{=} \KeywordTok{new} \FunctionTok{SignatureOptions}\OperatorTok{();}
\NormalTok{sigOptions}\OperatorTok{.}\FunctionTok{setPreferredSignatureSize}\OperatorTok{(}\DecValTok{32768}\OperatorTok{);} \CommentTok{// 32 KB for the CMS container}
\end{Highlighting}
\end{Shaded}

The \texttt{preferredSignatureSize} is the number of \emph{bytes}
reserved for the hex-encoded signature value in \texttt{/Contents}.
Since the hex encoding doubles the size, a 32768-byte container can hold
a CMS blob of up to about 16 KB.

\begin{longtable}[]{@{}
  >{\raggedright\arraybackslash}p{(\linewidth - 4\tabcolsep) * \real{0.3000}}
  >{\raggedright\arraybackslash}p{(\linewidth - 4\tabcolsep) * \real{0.2000}}
  >{\raggedright\arraybackslash}p{(\linewidth - 4\tabcolsep) * \real{0.5000}}@{}}
\caption{Recommended Container Sizes}\tabularnewline
\toprule\noalign{}
\begin{minipage}[b]{\linewidth}\raggedright
Scenario
\end{minipage} & \begin{minipage}[b]{\linewidth}\raggedright
Size (bytes)
\end{minipage} & \begin{minipage}[b]{\linewidth}\raggedright
Why
\end{minipage} \\
\midrule\noalign{}
\endfirsthead
\toprule\noalign{}
\begin{minipage}[b]{\linewidth}\raggedright
Scenario
\end{minipage} & \begin{minipage}[b]{\linewidth}\raggedright
Size (bytes)
\end{minipage} & \begin{minipage}[b]{\linewidth}\raggedright
Why
\end{minipage} \\
\midrule\noalign{}
\endhead
\bottomrule\noalign{}
\endlastfoot
Basic signature (no timestamp) & 32768 & A CMS \texttt{SignedData} with
one signer, RSA-2048, and a 3-certificate chain typically occupies 4-8
KB. 32 KB gives ample room. \\
Signature with TSA timestamp (PAdES-B-T) & 65536 & The timestamp token
adds another CMS \texttt{SignedData} structure (the TSA's signature, its
certificate chain, and the \texttt{TSTInfo}). 64 KB handles even large
TSA responses. \\
Signature with embedded OCSP/CRL (PAdES-B-LT) & 131072 or more & CRL
responses can be very large (100 KB+). If you embed revocation data in
the CMS structure rather than the DSS, you need a correspondingly large
container. \\
\end{longtable}

If your CMS blob exceeds the reserved container size,
\texttt{saveIncremental()} throws an \texttt{IOException} with a message
like "can't write signature, not enough space." The fix is to increase
\texttt{preferredSignatureSize}. But do not set it to absurd values
(e.g., 10 MB) because the entire container is written as hex zeros into
the PDF --- a 10 MB container adds 20 MB of hex padding to the file.

\section{Complete Basic Signing
Example}\label{_complete_basic_signing_example}

Let us put the three pieces together into a complete, compilable
example. This is the minimum code needed to sign a PDF with PDFBox
3.0.0.

\begin{Shaded}
\begin{Highlighting}[]
\KeywordTok{import} \ImportTok{java}\OperatorTok{.}\ImportTok{io}\OperatorTok{.}\ImportTok{ByteArrayOutputStream}\OperatorTok{;}
\KeywordTok{import} \ImportTok{java}\OperatorTok{.}\ImportTok{io}\OperatorTok{.}\ImportTok{FileInputStream}\OperatorTok{;}
\KeywordTok{import} \ImportTok{java}\OperatorTok{.}\ImportTok{io}\OperatorTok{.}\ImportTok{IOException}\OperatorTok{;}
\KeywordTok{import} \ImportTok{java}\OperatorTok{.}\ImportTok{io}\OperatorTok{.}\ImportTok{InputStream}\OperatorTok{;}
\KeywordTok{import} \ImportTok{java}\OperatorTok{.}\ImportTok{nio}\OperatorTok{.}\ImportTok{file}\OperatorTok{.}\ImportTok{Files}\OperatorTok{;}
\KeywordTok{import} \ImportTok{java}\OperatorTok{.}\ImportTok{nio}\OperatorTok{.}\ImportTok{file}\OperatorTok{.}\ImportTok{Path}\OperatorTok{;}
\KeywordTok{import} \ImportTok{java}\OperatorTok{.}\ImportTok{security}\OperatorTok{.}\ImportTok{KeyStore}\OperatorTok{;}
\KeywordTok{import} \ImportTok{java}\OperatorTok{.}\ImportTok{security}\OperatorTok{.}\ImportTok{PrivateKey}\OperatorTok{;}
\KeywordTok{import} \ImportTok{java}\OperatorTok{.}\ImportTok{security}\OperatorTok{.}\ImportTok{cert}\OperatorTok{.}\ImportTok{X509Certificate}\OperatorTok{;}
\KeywordTok{import} \ImportTok{java}\OperatorTok{.}\ImportTok{util}\OperatorTok{.}\ImportTok{ArrayList}\OperatorTok{;}
\KeywordTok{import} \ImportTok{java}\OperatorTok{.}\ImportTok{util}\OperatorTok{.}\ImportTok{Calendar}\OperatorTok{;}
\KeywordTok{import} \ImportTok{java}\OperatorTok{.}\ImportTok{util}\OperatorTok{.}\ImportTok{Enumeration}\OperatorTok{;}
\KeywordTok{import} \ImportTok{java}\OperatorTok{.}\ImportTok{util}\OperatorTok{.}\ImportTok{List}\OperatorTok{;}

\KeywordTok{import} \ImportTok{org}\OperatorTok{.}\ImportTok{apache}\OperatorTok{.}\ImportTok{pdfbox}\OperatorTok{.}\ImportTok{Loader}\OperatorTok{;}
\KeywordTok{import} \ImportTok{org}\OperatorTok{.}\ImportTok{apache}\OperatorTok{.}\ImportTok{pdfbox}\OperatorTok{.}\ImportTok{cos}\OperatorTok{.}\ImportTok{COSName}\OperatorTok{;}
\KeywordTok{import} \ImportTok{org}\OperatorTok{.}\ImportTok{apache}\OperatorTok{.}\ImportTok{pdfbox}\OperatorTok{.}\ImportTok{pdmodel}\OperatorTok{.}\ImportTok{PDDocument}\OperatorTok{;}
\KeywordTok{import} \ImportTok{org}\OperatorTok{.}\ImportTok{apache}\OperatorTok{.}\ImportTok{pdfbox}\OperatorTok{.}\ImportTok{pdmodel}\OperatorTok{.}\ImportTok{interactive}\OperatorTok{.}\ImportTok{digitalsignature}\OperatorTok{.}\ImportTok{PDSignature}\OperatorTok{;}
\KeywordTok{import} \ImportTok{org}\OperatorTok{.}\ImportTok{apache}\OperatorTok{.}\ImportTok{pdfbox}\OperatorTok{.}\ImportTok{pdmodel}\OperatorTok{.}\ImportTok{interactive}\OperatorTok{.}\ImportTok{digitalsignature}\OperatorTok{.}\ImportTok{SignatureInterface}\OperatorTok{;}
\KeywordTok{import} \ImportTok{org}\OperatorTok{.}\ImportTok{apache}\OperatorTok{.}\ImportTok{pdfbox}\OperatorTok{.}\ImportTok{pdmodel}\OperatorTok{.}\ImportTok{interactive}\OperatorTok{.}\ImportTok{digitalsignature}\OperatorTok{.}\ImportTok{SignatureOptions}\OperatorTok{;}
\KeywordTok{import} \ImportTok{org}\OperatorTok{.}\ImportTok{bouncycastle}\OperatorTok{.}\ImportTok{cert}\OperatorTok{.}\ImportTok{jcajce}\OperatorTok{.}\ImportTok{JcaCertStore}\OperatorTok{;}
\KeywordTok{import} \ImportTok{org}\OperatorTok{.}\ImportTok{bouncycastle}\OperatorTok{.}\ImportTok{cms}\OperatorTok{.}\ImportTok{CMSProcessableByteArray}\OperatorTok{;}
\KeywordTok{import} \ImportTok{org}\OperatorTok{.}\ImportTok{bouncycastle}\OperatorTok{.}\ImportTok{cms}\OperatorTok{.}\ImportTok{CMSSignedData}\OperatorTok{;}
\KeywordTok{import} \ImportTok{org}\OperatorTok{.}\ImportTok{bouncycastle}\OperatorTok{.}\ImportTok{cms}\OperatorTok{.}\ImportTok{CMSSignedDataGenerator}\OperatorTok{;}
\KeywordTok{import} \ImportTok{org}\OperatorTok{.}\ImportTok{bouncycastle}\OperatorTok{.}\ImportTok{cms}\OperatorTok{.}\ImportTok{jcajce}\OperatorTok{.}\ImportTok{JcaSignerInfoGeneratorBuilder}\OperatorTok{;}
\KeywordTok{import} \ImportTok{org}\OperatorTok{.}\ImportTok{bouncycastle}\OperatorTok{.}\ImportTok{jce}\OperatorTok{.}\ImportTok{provider}\OperatorTok{.}\ImportTok{BouncyCastleProvider}\OperatorTok{;}
\KeywordTok{import} \ImportTok{org}\OperatorTok{.}\ImportTok{bouncycastle}\OperatorTok{.}\ImportTok{operator}\OperatorTok{.}\ImportTok{ContentSigner}\OperatorTok{;}
\KeywordTok{import} \ImportTok{org}\OperatorTok{.}\ImportTok{bouncycastle}\OperatorTok{.}\ImportTok{operator}\OperatorTok{.}\ImportTok{jcajce}\OperatorTok{.}\ImportTok{JcaContentSignerBuilder}\OperatorTok{;}
\KeywordTok{import} \ImportTok{org}\OperatorTok{.}\ImportTok{bouncycastle}\OperatorTok{.}\ImportTok{operator}\OperatorTok{.}\ImportTok{jcajce}\OperatorTok{.}\ImportTok{JcaDigestCalculatorProviderBuilder}\OperatorTok{;}

\CommentTok{/**}
 \CommentTok{*}\NormalTok{ Minimal PDF signing example with PDFBox }\CommentTok{3.}\NormalTok{x and BouncyCastle}\CommentTok{.}
 \CommentTok{*}\NormalTok{ Signs a PDF with an invisible signature using a PKCS}\CommentTok{\#12}\NormalTok{ keystore}\CommentTok{.}
 \CommentTok{*/}
\KeywordTok{public} \KeywordTok{class}\NormalTok{ BasicPdfSigner }\OperatorTok{\{}

    \KeywordTok{public} \DataTypeTok{static} \DataTypeTok{void} \FunctionTok{main}\OperatorTok{(}\BuiltInTok{String}\OperatorTok{[]}\NormalTok{ args}\OperatorTok{)} \KeywordTok{throws} \BuiltInTok{Exception} \OperatorTok{\{}
        \CommentTok{// Register BouncyCastle provider {-}{-} must happen before any crypto operation}
\NormalTok{        java}\OperatorTok{.}\FunctionTok{security}\OperatorTok{.}\FunctionTok{Security}\OperatorTok{.}\FunctionTok{addProvider}\OperatorTok{(}\KeywordTok{new} \FunctionTok{BouncyCastleProvider}\OperatorTok{());}

\NormalTok{        Path inputPdf }\OperatorTok{=}\NormalTok{ Path}\OperatorTok{.}\FunctionTok{of}\OperatorTok{(}\StringTok{"document.pdf"}\OperatorTok{);}
\NormalTok{        Path keystorePath }\OperatorTok{=}\NormalTok{ Path}\OperatorTok{.}\FunctionTok{of}\OperatorTok{(}\StringTok{"keystore.p12"}\OperatorTok{);}
        \DataTypeTok{char}\OperatorTok{[]}\NormalTok{ keystorePassword }\OperatorTok{=} \StringTok{"changeit"}\OperatorTok{.}\FunctionTok{toCharArray}\OperatorTok{();}

        \CommentTok{// Load key material from PKCS\#12}
        \BuiltInTok{KeyStore}\NormalTok{ ks }\OperatorTok{=} \BuiltInTok{KeyStore}\OperatorTok{.}\FunctionTok{getInstance}\OperatorTok{(}\StringTok{"PKCS12"}\OperatorTok{);}
        \ControlFlowTok{try} \OperatorTok{(}\BuiltInTok{FileInputStream}\NormalTok{ fis }\OperatorTok{=} \KeywordTok{new} \BuiltInTok{FileInputStream}\OperatorTok{(}\NormalTok{keystorePath}\OperatorTok{.}\FunctionTok{toFile}\OperatorTok{()))} \OperatorTok{\{}
\NormalTok{            ks}\OperatorTok{.}\FunctionTok{load}\OperatorTok{(}\NormalTok{fis}\OperatorTok{,}\NormalTok{ keystorePassword}\OperatorTok{);}
        \OperatorTok{\}}

        \BuiltInTok{Enumeration}\OperatorTok{\textless{}}\BuiltInTok{String}\OperatorTok{\textgreater{}}\NormalTok{ aliases }\OperatorTok{=}\NormalTok{ ks}\OperatorTok{.}\FunctionTok{aliases}\OperatorTok{();}
        \BuiltInTok{String}\NormalTok{ alias }\OperatorTok{=}\NormalTok{ aliases}\OperatorTok{.}\FunctionTok{nextElement}\OperatorTok{();}
        \BuiltInTok{PrivateKey}\NormalTok{ privateKey }\OperatorTok{=} \OperatorTok{(}\BuiltInTok{PrivateKey}\OperatorTok{)}\NormalTok{ ks}\OperatorTok{.}\FunctionTok{getKey}\OperatorTok{(}\NormalTok{alias}\OperatorTok{,}\NormalTok{ keystorePassword}\OperatorTok{);}
        \BuiltInTok{X509Certificate}\NormalTok{ cert }\OperatorTok{=} \OperatorTok{(}\BuiltInTok{X509Certificate}\OperatorTok{)}\NormalTok{ ks}\OperatorTok{.}\FunctionTok{getCertificate}\OperatorTok{(}\NormalTok{alias}\OperatorTok{);}

        \BuiltInTok{List}\OperatorTok{\textless{}}\BuiltInTok{X509Certificate}\OperatorTok{\textgreater{}}\NormalTok{ chain }\OperatorTok{=} \KeywordTok{new} \BuiltInTok{ArrayList}\OperatorTok{\textless{}\textgreater{}();}
        \ControlFlowTok{for} \OperatorTok{(}\NormalTok{java}\OperatorTok{.}\FunctionTok{security}\OperatorTok{.}\FunctionTok{cert}\OperatorTok{.}\FunctionTok{Certificate}\NormalTok{ c }\OperatorTok{:}\NormalTok{ ks}\OperatorTok{.}\FunctionTok{getCertificateChain}\OperatorTok{(}\NormalTok{alias}\OperatorTok{))} \OperatorTok{\{}
\NormalTok{            chain}\OperatorTok{.}\FunctionTok{add}\OperatorTok{((}\BuiltInTok{X509Certificate}\OperatorTok{)}\NormalTok{ c}\OperatorTok{);}
        \OperatorTok{\}}

        \CommentTok{// Load the PDF}
        \DataTypeTok{byte}\OperatorTok{[]}\NormalTok{ pdfBytes }\OperatorTok{=}\NormalTok{ Files}\OperatorTok{.}\FunctionTok{readAllBytes}\OperatorTok{(}\NormalTok{inputPdf}\OperatorTok{);}

        \ControlFlowTok{try} \OperatorTok{(}\NormalTok{PDDocument doc }\OperatorTok{=}\NormalTok{ Loader}\OperatorTok{.}\FunctionTok{loadPDF}\OperatorTok{(}\NormalTok{pdfBytes}\OperatorTok{))} \OperatorTok{\{} 
            \CommentTok{// Create the signature dictionary}
\NormalTok{            PDSignature signature }\OperatorTok{=} \KeywordTok{new} \FunctionTok{PDSignature}\OperatorTok{();}
\NormalTok{            signature}\OperatorTok{.}\FunctionTok{setFilter}\OperatorTok{(}\NormalTok{PDSignature}\OperatorTok{.}\FunctionTok{FILTER\_ADOBE\_PPKLITE}\OperatorTok{);}
\NormalTok{            signature}\OperatorTok{.}\FunctionTok{setSubFilter}\OperatorTok{(}\NormalTok{COSName}\OperatorTok{.}\FunctionTok{getPDFName}\OperatorTok{(}\StringTok{"ETSI.CAdES.detached"}\OperatorTok{));}
\NormalTok{            signature}\OperatorTok{.}\FunctionTok{setSignDate}\OperatorTok{(}\BuiltInTok{Calendar}\OperatorTok{.}\FunctionTok{getInstance}\OperatorTok{());}
\NormalTok{            signature}\OperatorTok{.}\FunctionTok{setName}\OperatorTok{(}\NormalTok{cert}\OperatorTok{.}\FunctionTok{getSubjectX500Principal}\OperatorTok{().}\FunctionTok{getName}\OperatorTok{());}
\NormalTok{            signature}\OperatorTok{.}\FunctionTok{setReason}\OperatorTok{(}\StringTok{"Document approval"}\OperatorTok{);}

            \CommentTok{// Configure container size}
\NormalTok{            SignatureOptions options }\OperatorTok{=} \KeywordTok{new} \FunctionTok{SignatureOptions}\OperatorTok{();}
\NormalTok{            options}\OperatorTok{.}\FunctionTok{setPreferredSignatureSize}\OperatorTok{(}\DecValTok{32768}\OperatorTok{);}

            \CommentTok{// Wire everything together}
\NormalTok{            SignatureInterface sigInterface }\OperatorTok{=} \KeywordTok{new} \FunctionTok{SimpleCmsSignatureInterface}\OperatorTok{(}
\NormalTok{                    privateKey}\OperatorTok{,}\NormalTok{ cert}\OperatorTok{,}\NormalTok{ chain}\OperatorTok{);}
\NormalTok{            doc}\OperatorTok{.}\FunctionTok{addSignature}\OperatorTok{(}\NormalTok{signature}\OperatorTok{,}\NormalTok{ sigInterface}\OperatorTok{,}\NormalTok{ options}\OperatorTok{);} 

            \CommentTok{// Save incrementally {-}{-} this triggers the sign() callback}
            \BuiltInTok{ByteArrayOutputStream}\NormalTok{ baos }\OperatorTok{=} \KeywordTok{new} \BuiltInTok{ByteArrayOutputStream}\OperatorTok{();}
\NormalTok{            doc}\OperatorTok{.}\FunctionTok{saveIncremental}\OperatorTok{(}\NormalTok{baos}\OperatorTok{);} 

            \CommentTok{// Write the signed PDF}
\NormalTok{            Path outputPdf }\OperatorTok{=}\NormalTok{ Path}\OperatorTok{.}\FunctionTok{of}\OperatorTok{(}\StringTok{"document{-}signed.pdf"}\OperatorTok{);}
\NormalTok{            Files}\OperatorTok{.}\FunctionTok{write}\OperatorTok{(}\NormalTok{outputPdf}\OperatorTok{,}\NormalTok{ baos}\OperatorTok{.}\FunctionTok{toByteArray}\OperatorTok{());}

            \BuiltInTok{System}\OperatorTok{.}\FunctionTok{out}\OperatorTok{.}\FunctionTok{println}\OperatorTok{(}\StringTok{"Signed PDF written to "} \OperatorTok{+}\NormalTok{ outputPdf}
                    \OperatorTok{+} \StringTok{" ("} \OperatorTok{+}\NormalTok{ baos}\OperatorTok{.}\FunctionTok{size}\OperatorTok{()} \OperatorTok{+} \StringTok{" bytes)"}\OperatorTok{);}
        \OperatorTok{\}}
    \OperatorTok{\}}
\OperatorTok{\}}
\end{Highlighting}
\end{Shaded}

\begin{itemize}
\item
  PDFBox 3.x uses \texttt{Loader.loadPDF(byte{[}{]})}. The old
  \texttt{PDDocument.load()} is gone. See the migration section later in
  this chapter.
\item
  \texttt{addSignature()} registers the signature dictionary, the
  callback, and the options. Nothing is written yet.
\item
  \texttt{saveIncremental()} writes the PDF as an incremental update,
  computes the byte range, calls your \texttt{sign()} method,
  hex-encodes the result, and patches it into the \texttt{/Contents}
  field. This is where all the magic happens.
\end{itemize}

Always use \texttt{saveIncremental()}, never \texttt{save()}. A regular
\texttt{save()} rewrites the entire PDF, which destroys any existing
signatures. \texttt{saveIncremental()} appends a new revision to the
file, preserving all previous signatures. This is not optional --- it is
the only correct way to save a signed PDF.

\section{The InputStream in sign(): What You Are Actually
Signing}\label{_the_inputstream_in_sign_what_you_are_actually_signing}

The \texttt{InputStream} passed to your \texttt{sign()} method is one of
the most misunderstood parts of the PDFBox signing API. Let us be very
precise about what it contains.

When PDFBox writes the PDF during \texttt{saveIncremental()}, it creates
the file with a placeholder for the \texttt{/Contents} value. The byte
range is computed as two ranges that together cover the entire file
\emph{except} the \texttt{/Contents} hex string:

\begin{verbatim}
/ByteRange [0 offset1 offset2 length2]
\end{verbatim}

Where:

\begin{itemize}
\item
  Range 1: bytes 0 through \texttt{offset1} (everything before the
  \texttt{/Contents} hex string)
\item
  Range 2: bytes \texttt{offset2} through \texttt{offset2\ +\ length2}
  (everything after the \texttt{/Contents} hex string)
\end{itemize}

The \texttt{InputStream} you receive is the concatenation of these two
ranges. You do not see the gap. You do not need to parse the byte range.
PDFBox handles all of that.

Your \texttt{sign()} method receives a continuous stream of bytes, and
you sign all of them.

\begin{Shaded}
\begin{Highlighting}[]
\AttributeTok{@Override}
\KeywordTok{public} \DataTypeTok{byte}\OperatorTok{[]} \FunctionTok{sign}\OperatorTok{(}\BuiltInTok{InputStream}\NormalTok{ content}\OperatorTok{)} \KeywordTok{throws} \BuiltInTok{IOException} \OperatorTok{\{}
    \DataTypeTok{byte}\OperatorTok{[]}\NormalTok{ contentBytes }\OperatorTok{=}\NormalTok{ content}\OperatorTok{.}\FunctionTok{readAllBytes}\OperatorTok{();} 
    \CommentTok{// contentBytes now contains Range1 + Range2}
    \CommentTok{// Sign these bytes and return the CMS SignedData}
\OperatorTok{\}}
\end{Highlighting}
\end{Shaded}

\begin{itemize}
\item
  You \emph{must} read the entire stream. If you read only part of it
  (e.g., using a fixed-size buffer and stopping early), the hash you
  compute will not match the hash a validator computes over the full
  byte range, and the signature will be invalid.
\end{itemize}

A common bug in PDFBox 2.x code (which people copy into 3.x projects) is
using \texttt{IOUtils.toByteArray()} from Apache Commons IO with a size
limit. The byte range content can be large --- for a 100 MB PDF, the
stream is approximately 100 MB. If your byte array buffer is smaller,
you get a truncated hash and a broken signature that validates in your
tests (small PDFs) but fails in production (large PDFs).

\section{Visible Signature
Appearances}\label{_visible_signature_appearances}

An invisible signature is fine for server-side signing workflows, but
many use cases --- contracts, approvals, notarizations --- require a
visible signature block on a specific page. This section shows how to
create visible appearances with PDFBox.

\subsection{The Form XObject Approach}\label{_the_form_xobject_approach}

A signature appearance is a Form XObject --- a self-contained mini-page
with its own coordinate system, resources, and content stream. You build
it, attach it to the signature field's widget annotation, and PDFBox
includes it in the signed content.

\begin{Shaded}
\begin{Highlighting}[]
\KeywordTok{import} \ImportTok{org}\OperatorTok{.}\ImportTok{apache}\OperatorTok{.}\ImportTok{pdfbox}\OperatorTok{.}\ImportTok{pdmodel}\OperatorTok{.}\ImportTok{PDDocument}\OperatorTok{;}
\KeywordTok{import} \ImportTok{org}\OperatorTok{.}\ImportTok{apache}\OperatorTok{.}\ImportTok{pdfbox}\OperatorTok{.}\ImportTok{pdmodel}\OperatorTok{.}\ImportTok{PDResources}\OperatorTok{;}
\KeywordTok{import} \ImportTok{org}\OperatorTok{.}\ImportTok{apache}\OperatorTok{.}\ImportTok{pdfbox}\OperatorTok{.}\ImportTok{pdmodel}\OperatorTok{.}\ImportTok{common}\OperatorTok{.}\ImportTok{PDRectangle}\OperatorTok{;}
\KeywordTok{import} \ImportTok{org}\OperatorTok{.}\ImportTok{apache}\OperatorTok{.}\ImportTok{pdfbox}\OperatorTok{.}\ImportTok{pdmodel}\OperatorTok{.}\ImportTok{font}\OperatorTok{.}\ImportTok{PDType1Font}\OperatorTok{;}
\KeywordTok{import} \ImportTok{org}\OperatorTok{.}\ImportTok{apache}\OperatorTok{.}\ImportTok{pdfbox}\OperatorTok{.}\ImportTok{pdmodel}\OperatorTok{.}\ImportTok{font}\OperatorTok{.}\ImportTok{Standard14Fonts}\OperatorTok{;}
\KeywordTok{import} \ImportTok{org}\OperatorTok{.}\ImportTok{apache}\OperatorTok{.}\ImportTok{pdfbox}\OperatorTok{.}\ImportTok{pdmodel}\OperatorTok{.}\ImportTok{graphics}\OperatorTok{.}\ImportTok{form}\OperatorTok{.}\ImportTok{PDFormXObject}\OperatorTok{;}
\KeywordTok{import} \ImportTok{org}\OperatorTok{.}\ImportTok{apache}\OperatorTok{.}\ImportTok{pdfbox}\OperatorTok{.}\ImportTok{pdmodel}\OperatorTok{.}\ImportTok{interactive}\OperatorTok{.}\ImportTok{form}\OperatorTok{.}\ImportTok{PDSignatureField}\OperatorTok{;}

\KeywordTok{import} \ImportTok{java}\OperatorTok{.}\ImportTok{io}\OperatorTok{.}\ImportTok{OutputStream}\OperatorTok{;}
\KeywordTok{import} \ImportTok{java}\OperatorTok{.}\ImportTok{nio}\OperatorTok{.}\ImportTok{charset}\OperatorTok{.}\ImportTok{StandardCharsets}\OperatorTok{;}

\CommentTok{/**}
 \CommentTok{*}\NormalTok{ Builds a visible text}\CommentTok{{-}}\NormalTok{based signature appearance as a Form XObject}\CommentTok{.}
 \CommentTok{*/}
\KeywordTok{public} \KeywordTok{class}\NormalTok{ SignatureAppearanceBuilder }\OperatorTok{\{}

    \CommentTok{/**}
     \CommentTok{*}\NormalTok{ Creates and attaches a visible text appearance to a signature field}\CommentTok{.}
     \CommentTok{*}
     \CommentTok{*} \AnnotationTok{@param}\CommentTok{ }\AnnotationTok{doc}\CommentTok{      }\NormalTok{the document being signed}
     \CommentTok{*} \AnnotationTok{@param}\CommentTok{ }\AnnotationTok{sigField}\CommentTok{ }\NormalTok{the signature field to attach the appearance to}
     \CommentTok{*} \AnnotationTok{@param}\CommentTok{ }\AnnotationTok{name}\CommentTok{     }\NormalTok{signer name to display}
     \CommentTok{*} \AnnotationTok{@param}\CommentTok{ }\AnnotationTok{reason}\CommentTok{   }\NormalTok{signing reason to display}
     \CommentTok{*} \AnnotationTok{@param}\CommentTok{ }\AnnotationTok{width}\CommentTok{    }\NormalTok{appearance width in points}
     \CommentTok{*} \AnnotationTok{@param}\CommentTok{ }\AnnotationTok{height}\CommentTok{   }\NormalTok{appearance height in points}
     \CommentTok{*/}
    \KeywordTok{public} \DataTypeTok{static} \DataTypeTok{void} \FunctionTok{buildTextAppearance}\OperatorTok{(}\NormalTok{PDDocument doc}\OperatorTok{,}
\NormalTok{                                            PDSignatureField sigField}\OperatorTok{,}
                                            \BuiltInTok{String}\NormalTok{ name}\OperatorTok{,} \BuiltInTok{String}\NormalTok{ reason}\OperatorTok{,}
                                            \DataTypeTok{float}\NormalTok{ width}\OperatorTok{,} \DataTypeTok{float}\NormalTok{ height}\OperatorTok{)}
            \KeywordTok{throws} \BuiltInTok{Exception} \OperatorTok{\{}

\NormalTok{        PDResources resources }\OperatorTok{=} \KeywordTok{new} \FunctionTok{PDResources}\OperatorTok{();}

        \CommentTok{// Use a Standard 14 font {-}{-} no embedding needed}
\NormalTok{        PDType1Font font }\OperatorTok{=} \KeywordTok{new} \FunctionTok{PDType1Font}\OperatorTok{(}
\NormalTok{                Standard14Fonts}\OperatorTok{.}\FunctionTok{FontName}\OperatorTok{.}\FunctionTok{HELVETICA}\OperatorTok{);} 
\NormalTok{        org}\OperatorTok{.}\FunctionTok{apache}\OperatorTok{.}\FunctionTok{pdfbox}\OperatorTok{.}\FunctionTok{cos}\OperatorTok{.}\FunctionTok{COSName}\NormalTok{ fontName }\OperatorTok{=}\NormalTok{ resources}\OperatorTok{.}\FunctionTok{add}\OperatorTok{(}\NormalTok{font}\OperatorTok{);}

        \CommentTok{// Build the content stream}
        \DataTypeTok{float}\NormalTok{ fontSize }\OperatorTok{=} \FloatTok{10f}\OperatorTok{;}
        \DataTypeTok{float}\NormalTok{ lineSpacing }\OperatorTok{=}\NormalTok{ fontSize }\OperatorTok{*} \FloatTok{1.4f}\OperatorTok{;}
        \DataTypeTok{float}\NormalTok{ margin }\OperatorTok{=} \FloatTok{4f}\OperatorTok{;}
        \DataTypeTok{float}\NormalTok{ y }\OperatorTok{=}\NormalTok{ height }\OperatorTok{{-}}\NormalTok{ margin }\OperatorTok{{-}}\NormalTok{ fontSize}\OperatorTok{;}

        \BuiltInTok{StringBuilder}\NormalTok{ ops }\OperatorTok{=} \KeywordTok{new} \BuiltInTok{StringBuilder}\OperatorTok{();}
\NormalTok{        ops}\OperatorTok{.}\FunctionTok{append}\OperatorTok{(}\StringTok{"BT}\SpecialCharTok{\textbackslash{}n}\StringTok{"}\OperatorTok{);}
\NormalTok{        ops}\OperatorTok{.}\FunctionTok{append}\OperatorTok{(}\StringTok{"/"}\OperatorTok{).}\FunctionTok{append}\OperatorTok{(}\NormalTok{fontName}\OperatorTok{.}\FunctionTok{getName}\OperatorTok{())}
                \OperatorTok{.}\FunctionTok{append}\OperatorTok{(}\StringTok{" "}\OperatorTok{).}\FunctionTok{append}\OperatorTok{(}\NormalTok{fontSize}\OperatorTok{).}\FunctionTok{append}\OperatorTok{(}\StringTok{" Tf}\SpecialCharTok{\textbackslash{}n}\StringTok{"}\OperatorTok{);}
\NormalTok{        ops}\OperatorTok{.}\FunctionTok{append}\OperatorTok{(}\StringTok{"0 0 0 rg}\SpecialCharTok{\textbackslash{}n}\StringTok{"}\OperatorTok{);} \CommentTok{// black text}

        \CommentTok{// Line 1: "Digitally signed by:"}
\NormalTok{        ops}\OperatorTok{.}\FunctionTok{append}\OperatorTok{(}\NormalTok{margin}\OperatorTok{).}\FunctionTok{append}\OperatorTok{(}\StringTok{" "}\OperatorTok{).}\FunctionTok{append}\OperatorTok{(}\NormalTok{y}\OperatorTok{).}\FunctionTok{append}\OperatorTok{(}\StringTok{" Td}\SpecialCharTok{\textbackslash{}n}\StringTok{"}\OperatorTok{);}
\NormalTok{        ops}\OperatorTok{.}\FunctionTok{append}\OperatorTok{(}\StringTok{"(Digitally signed by:) Tj}\SpecialCharTok{\textbackslash{}n}\StringTok{"}\OperatorTok{);}

        \CommentTok{// Line 2: signer name}
\NormalTok{        y }\OperatorTok{{-}=}\NormalTok{ lineSpacing}\OperatorTok{;}
\NormalTok{        ops}\OperatorTok{.}\FunctionTok{append}\OperatorTok{(}\NormalTok{margin}\OperatorTok{).}\FunctionTok{append}\OperatorTok{(}\StringTok{" "}\OperatorTok{).}\FunctionTok{append}\OperatorTok{(}\NormalTok{y}\OperatorTok{).}\FunctionTok{append}\OperatorTok{(}\StringTok{" Td}\SpecialCharTok{\textbackslash{}n}\StringTok{"}\OperatorTok{);}
\NormalTok{        ops}\OperatorTok{.}\FunctionTok{append}\OperatorTok{(}\StringTok{"("}\OperatorTok{).}\FunctionTok{append}\OperatorTok{(}\FunctionTok{escapePs}\OperatorTok{(}\NormalTok{name}\OperatorTok{)).}\FunctionTok{append}\OperatorTok{(}\StringTok{") Tj}\SpecialCharTok{\textbackslash{}n}\StringTok{"}\OperatorTok{);}

        \CommentTok{// Line 3: reason}
        \ControlFlowTok{if} \OperatorTok{(}\NormalTok{reason }\OperatorTok{!=} \KeywordTok{null} \OperatorTok{\&\&} \OperatorTok{!}\NormalTok{reason}\OperatorTok{.}\FunctionTok{isEmpty}\OperatorTok{())} \OperatorTok{\{}
\NormalTok{            y }\OperatorTok{{-}=}\NormalTok{ lineSpacing}\OperatorTok{;}
\NormalTok{            ops}\OperatorTok{.}\FunctionTok{append}\OperatorTok{(}\NormalTok{margin}\OperatorTok{).}\FunctionTok{append}\OperatorTok{(}\StringTok{" "}\OperatorTok{).}\FunctionTok{append}\OperatorTok{(}\NormalTok{y}\OperatorTok{).}\FunctionTok{append}\OperatorTok{(}\StringTok{" Td}\SpecialCharTok{\textbackslash{}n}\StringTok{"}\OperatorTok{);}
\NormalTok{            ops}\OperatorTok{.}\FunctionTok{append}\OperatorTok{(}\StringTok{"(Reason: "}\OperatorTok{).}\FunctionTok{append}\OperatorTok{(}\FunctionTok{escapePs}\OperatorTok{(}\NormalTok{reason}\OperatorTok{)).}\FunctionTok{append}\OperatorTok{(}\StringTok{") Tj}\SpecialCharTok{\textbackslash{}n}\StringTok{"}\OperatorTok{);}
        \OperatorTok{\}}

\NormalTok{        ops}\OperatorTok{.}\FunctionTok{append}\OperatorTok{(}\StringTok{"ET}\SpecialCharTok{\textbackslash{}n}\StringTok{"}\OperatorTok{);}

        \DataTypeTok{byte}\OperatorTok{[]}\NormalTok{ streamBytes }\OperatorTok{=}\NormalTok{ ops}\OperatorTok{.}\FunctionTok{toString}\OperatorTok{()}
                \OperatorTok{.}\FunctionTok{getBytes}\OperatorTok{(}\NormalTok{StandardCharsets}\OperatorTok{.}\FunctionTok{US\_ASCII}\OperatorTok{);}

        \CommentTok{// Create the Form XObject}
\NormalTok{        PDFormXObject form }\OperatorTok{=} \KeywordTok{new} \FunctionTok{PDFormXObject}\OperatorTok{(}\NormalTok{doc}\OperatorTok{);}
\NormalTok{        form}\OperatorTok{.}\FunctionTok{setBBox}\OperatorTok{(}\KeywordTok{new} \FunctionTok{PDRectangle}\OperatorTok{(}\NormalTok{width}\OperatorTok{,}\NormalTok{ height}\OperatorTok{));}
\NormalTok{        form}\OperatorTok{.}\FunctionTok{setResources}\OperatorTok{(}\NormalTok{resources}\OperatorTok{);}

        \ControlFlowTok{try} \OperatorTok{(}\BuiltInTok{OutputStream}\NormalTok{ os }\OperatorTok{=}\NormalTok{ form}\OperatorTok{.}\FunctionTok{getStream}\OperatorTok{().}\FunctionTok{createOutputStream}\OperatorTok{())} \OperatorTok{\{} 
\NormalTok{            os}\OperatorTok{.}\FunctionTok{write}\OperatorTok{(}\NormalTok{streamBytes}\OperatorTok{);}
        \OperatorTok{\}}

        \CommentTok{// Attach to the widget annotation}
\NormalTok{        org}\OperatorTok{.}\FunctionTok{apache}\OperatorTok{.}\FunctionTok{pdfbox}\OperatorTok{.}\FunctionTok{cos}\OperatorTok{.}\FunctionTok{COSDictionary}\NormalTok{ apDict }\OperatorTok{=}
                \KeywordTok{new}\NormalTok{ org}\OperatorTok{.}\FunctionTok{apache}\OperatorTok{.}\FunctionTok{pdfbox}\OperatorTok{.}\FunctionTok{cos}\OperatorTok{.}\FunctionTok{COSDictionary}\OperatorTok{();}
\NormalTok{        apDict}\OperatorTok{.}\FunctionTok{setItem}\OperatorTok{(}\NormalTok{org}\OperatorTok{.}\FunctionTok{apache}\OperatorTok{.}\FunctionTok{pdfbox}\OperatorTok{.}\FunctionTok{cos}\OperatorTok{.}\FunctionTok{COSName}\OperatorTok{.}\FunctionTok{N}\OperatorTok{,}\NormalTok{ form}\OperatorTok{);}
\NormalTok{        sigField}\OperatorTok{.}\FunctionTok{getWidgets}\OperatorTok{().}\FunctionTok{get}\OperatorTok{(}\DecValTok{0}\OperatorTok{).}\FunctionTok{getCOSObject}\OperatorTok{()}
                \OperatorTok{.}\FunctionTok{setItem}\OperatorTok{(}\NormalTok{org}\OperatorTok{.}\FunctionTok{apache}\OperatorTok{.}\FunctionTok{pdfbox}\OperatorTok{.}\FunctionTok{cos}\OperatorTok{.}\FunctionTok{COSName}\OperatorTok{.}\FunctionTok{AP}\OperatorTok{,}\NormalTok{ apDict}\OperatorTok{);}
    \OperatorTok{\}}

    \CommentTok{/**}\NormalTok{ Escape parentheses and backslashes for a PDF string literal}\CommentTok{. */}
    \KeywordTok{private} \DataTypeTok{static} \BuiltInTok{String} \FunctionTok{escapePs}\OperatorTok{(}\BuiltInTok{String}\NormalTok{ text}\OperatorTok{)} \OperatorTok{\{}
        \ControlFlowTok{return}\NormalTok{ text}\OperatorTok{.}\FunctionTok{replace}\OperatorTok{(}\StringTok{"}\SpecialCharTok{\textbackslash{}\textbackslash{}}\StringTok{"}\OperatorTok{,} \StringTok{"}\SpecialCharTok{\textbackslash{}\textbackslash{}\textbackslash{}\textbackslash{}}\StringTok{"}\OperatorTok{)}
                   \OperatorTok{.}\FunctionTok{replace}\OperatorTok{(}\StringTok{"("}\OperatorTok{,} \StringTok{"}\SpecialCharTok{\textbackslash{}\textbackslash{}}\StringTok{("}\OperatorTok{)}
                   \OperatorTok{.}\FunctionTok{replace}\OperatorTok{(}\StringTok{")"}\OperatorTok{,} \StringTok{"}\SpecialCharTok{\textbackslash{}\textbackslash{}}\StringTok{)"}\OperatorTok{);}
    \OperatorTok{\}}
\OperatorTok{\}}
\end{Highlighting}
\end{Shaded}

\begin{itemize}
\item
  In PDFBox 3.x, \texttt{PDType1Font.HELVETICA} does not exist. You must
  use \texttt{new\ PDType1Font(Standard14Fonts.FontName.HELVETICA)}.
  This is one of the most common migration errors. See the migration
  section at the end of this chapter.
\item
  \texttt{form.getStream().createOutputStream()} returns an output
  stream that writes directly into the Form XObject's content stream.
  This is the correct way to write appearance content in PDFBox 3.x. You
  cannot use \texttt{PDPageContentStream} with a Form XObject.
\end{itemize}

\subsection{Image-Based Appearances}\label{_image_based_appearances}

For handwritten signature images (e.g., from a drawing canvas or a
scanned signature):

\begin{Shaded}
\begin{Highlighting}[]
\KeywordTok{import} \ImportTok{org}\OperatorTok{.}\ImportTok{apache}\OperatorTok{.}\ImportTok{pdfbox}\OperatorTok{.}\ImportTok{pdmodel}\OperatorTok{.}\ImportTok{graphics}\OperatorTok{.}\ImportTok{image}\OperatorTok{.}\ImportTok{PDImageXObject}\OperatorTok{;}
\KeywordTok{import} \ImportTok{java}\OperatorTok{.}\ImportTok{util}\OperatorTok{.}\ImportTok{Base64}\OperatorTok{;}

\CommentTok{/**}
 \CommentTok{*}\NormalTok{ Creates a signature appearance from a Base64}\CommentTok{{-}}\NormalTok{encoded PNG image}\CommentTok{.}
 \CommentTok{*/}
\KeywordTok{public} \DataTypeTok{static} \DataTypeTok{void} \FunctionTok{buildImageAppearance}\OperatorTok{(}\NormalTok{PDDocument doc}\OperatorTok{,}
\NormalTok{                                         PDSignatureField sigField}\OperatorTok{,}
                                         \BuiltInTok{String}\NormalTok{ base64Png}\OperatorTok{,}
                                         \DataTypeTok{float}\NormalTok{ width}\OperatorTok{,} \DataTypeTok{float}\NormalTok{ height}\OperatorTok{)}
        \KeywordTok{throws} \BuiltInTok{Exception} \OperatorTok{\{}

    \CommentTok{// Decode the image}
    \BuiltInTok{String}\NormalTok{ data }\OperatorTok{=}\NormalTok{ base64Png}\OperatorTok{;}
    \ControlFlowTok{if} \OperatorTok{(}\NormalTok{data}\OperatorTok{.}\FunctionTok{contains}\OperatorTok{(}\StringTok{","}\OperatorTok{))} \OperatorTok{\{}
\NormalTok{        data }\OperatorTok{=}\NormalTok{ data}\OperatorTok{.}\FunctionTok{substring}\OperatorTok{(}\NormalTok{data}\OperatorTok{.}\FunctionTok{indexOf}\OperatorTok{(}\CharTok{\textquotesingle{},\textquotesingle{}}\OperatorTok{)} \OperatorTok{+} \DecValTok{1}\OperatorTok{);}
    \OperatorTok{\}}
    \DataTypeTok{byte}\OperatorTok{[]}\NormalTok{ imageBytes }\OperatorTok{=}\NormalTok{ Base64}\OperatorTok{.}\FunctionTok{getDecoder}\OperatorTok{().}\FunctionTok{decode}\OperatorTok{(}\NormalTok{data}\OperatorTok{);}

\NormalTok{    PDImageXObject pdImage }\OperatorTok{=}\NormalTok{ PDImageXObject}\OperatorTok{.}\FunctionTok{createFromByteArray}\OperatorTok{(}
\NormalTok{            doc}\OperatorTok{,}\NormalTok{ imageBytes}\OperatorTok{,} \StringTok{"signature{-}image"}\OperatorTok{);}

\NormalTok{    PDResources resources }\OperatorTok{=} \KeywordTok{new} \FunctionTok{PDResources}\OperatorTok{();}
\NormalTok{    org}\OperatorTok{.}\FunctionTok{apache}\OperatorTok{.}\FunctionTok{pdfbox}\OperatorTok{.}\FunctionTok{cos}\OperatorTok{.}\FunctionTok{COSName}\NormalTok{ imgName }\OperatorTok{=}\NormalTok{ resources}\OperatorTok{.}\FunctionTok{add}\OperatorTok{(}\NormalTok{pdImage}\OperatorTok{);}

    \CommentTok{// Scale the image to fit the appearance rectangle while maintaining}
    \CommentTok{// aspect ratio}
    \DataTypeTok{float}\NormalTok{ imgW }\OperatorTok{=}\NormalTok{ pdImage}\OperatorTok{.}\FunctionTok{getWidth}\OperatorTok{();}
    \DataTypeTok{float}\NormalTok{ imgH }\OperatorTok{=}\NormalTok{ pdImage}\OperatorTok{.}\FunctionTok{getHeight}\OperatorTok{();}
    \DataTypeTok{float}\NormalTok{ scale }\OperatorTok{=} \BuiltInTok{Math}\OperatorTok{.}\FunctionTok{min}\OperatorTok{(}\NormalTok{width }\OperatorTok{/}\NormalTok{ imgW}\OperatorTok{,}\NormalTok{ height }\OperatorTok{/}\NormalTok{ imgH}\OperatorTok{);}
    \DataTypeTok{float}\NormalTok{ drawW }\OperatorTok{=}\NormalTok{ imgW }\OperatorTok{*}\NormalTok{ scale}\OperatorTok{;}
    \DataTypeTok{float}\NormalTok{ drawH }\OperatorTok{=}\NormalTok{ imgH }\OperatorTok{*}\NormalTok{ scale}\OperatorTok{;}
    \DataTypeTok{float}\NormalTok{ xOffset }\OperatorTok{=} \OperatorTok{(}\NormalTok{width }\OperatorTok{{-}}\NormalTok{ drawW}\OperatorTok{)} \OperatorTok{/} \DecValTok{2}\OperatorTok{;}
    \DataTypeTok{float}\NormalTok{ yOffset }\OperatorTok{=} \OperatorTok{(}\NormalTok{height }\OperatorTok{{-}}\NormalTok{ drawH}\OperatorTok{)} \OperatorTok{/} \DecValTok{2}\OperatorTok{;}

    \CommentTok{// PDF content stream: push graphics state, apply transformation}
    \CommentTok{// matrix, draw image, restore graphics state}
    \BuiltInTok{String}\NormalTok{ ops }\OperatorTok{=} \StringTok{"q}\SpecialCharTok{\textbackslash{}n}\StringTok{"}
            \OperatorTok{+}\NormalTok{ drawW }\OperatorTok{+} \StringTok{" 0 0 "} \OperatorTok{+}\NormalTok{ drawH }\OperatorTok{+} \StringTok{" "}
            \OperatorTok{+}\NormalTok{ xOffset }\OperatorTok{+} \StringTok{" "} \OperatorTok{+}\NormalTok{ yOffset }\OperatorTok{+} \StringTok{" cm}\SpecialCharTok{\textbackslash{}n}\StringTok{"}
            \OperatorTok{+} \StringTok{"/"} \OperatorTok{+}\NormalTok{ imgName}\OperatorTok{.}\FunctionTok{getName}\OperatorTok{()} \OperatorTok{+} \StringTok{" Do}\SpecialCharTok{\textbackslash{}n}\StringTok{"}
            \OperatorTok{+} \StringTok{"Q}\SpecialCharTok{\textbackslash{}n}\StringTok{"}\OperatorTok{;}

\NormalTok{    PDFormXObject form }\OperatorTok{=} \KeywordTok{new} \FunctionTok{PDFormXObject}\OperatorTok{(}\NormalTok{doc}\OperatorTok{);}
\NormalTok{    form}\OperatorTok{.}\FunctionTok{setBBox}\OperatorTok{(}\KeywordTok{new} \FunctionTok{PDRectangle}\OperatorTok{(}\NormalTok{width}\OperatorTok{,}\NormalTok{ height}\OperatorTok{));}
\NormalTok{    form}\OperatorTok{.}\FunctionTok{setResources}\OperatorTok{(}\NormalTok{resources}\OperatorTok{);}

    \ControlFlowTok{try} \OperatorTok{(}\BuiltInTok{OutputStream}\NormalTok{ os }\OperatorTok{=}\NormalTok{ form}\OperatorTok{.}\FunctionTok{getStream}\OperatorTok{().}\FunctionTok{createOutputStream}\OperatorTok{())} \OperatorTok{\{}
\NormalTok{        os}\OperatorTok{.}\FunctionTok{write}\OperatorTok{(}\NormalTok{ops}\OperatorTok{.}\FunctionTok{getBytes}\OperatorTok{(}\NormalTok{java}\OperatorTok{.}\FunctionTok{nio}\OperatorTok{.}\FunctionTok{charset}\OperatorTok{.}\FunctionTok{StandardCharsets}\OperatorTok{.}\FunctionTok{US\_ASCII}\OperatorTok{));}
    \OperatorTok{\}}

\NormalTok{    org}\OperatorTok{.}\FunctionTok{apache}\OperatorTok{.}\FunctionTok{pdfbox}\OperatorTok{.}\FunctionTok{cos}\OperatorTok{.}\FunctionTok{COSDictionary}\NormalTok{ apDict }\OperatorTok{=}
            \KeywordTok{new}\NormalTok{ org}\OperatorTok{.}\FunctionTok{apache}\OperatorTok{.}\FunctionTok{pdfbox}\OperatorTok{.}\FunctionTok{cos}\OperatorTok{.}\FunctionTok{COSDictionary}\OperatorTok{();}
\NormalTok{    apDict}\OperatorTok{.}\FunctionTok{setItem}\OperatorTok{(}\NormalTok{org}\OperatorTok{.}\FunctionTok{apache}\OperatorTok{.}\FunctionTok{pdfbox}\OperatorTok{.}\FunctionTok{cos}\OperatorTok{.}\FunctionTok{COSName}\OperatorTok{.}\FunctionTok{N}\OperatorTok{,}\NormalTok{ form}\OperatorTok{);}
\NormalTok{    sigField}\OperatorTok{.}\FunctionTok{getWidgets}\OperatorTok{().}\FunctionTok{get}\OperatorTok{(}\DecValTok{0}\OperatorTok{).}\FunctionTok{getCOSObject}\OperatorTok{()}
            \OperatorTok{.}\FunctionTok{setItem}\OperatorTok{(}\NormalTok{org}\OperatorTok{.}\FunctionTok{apache}\OperatorTok{.}\FunctionTok{pdfbox}\OperatorTok{.}\FunctionTok{cos}\OperatorTok{.}\FunctionTok{COSName}\OperatorTok{.}\FunctionTok{AP}\OperatorTok{,}\NormalTok{ apDict}\OperatorTok{);}
\OperatorTok{\}}
\end{Highlighting}
\end{Shaded}

The content stream uses the standard PDF imaging operators: \texttt{q}
(save graphics state), \texttt{cm} (concatenate transformation matrix),
\texttt{Do} (paint XObject), \texttt{Q} (restore graphics state). The
transformation matrix
\texttt{{[}drawW\ 0\ 0\ drawH\ xOffset\ yOffset{]}} scales and positions
the image within the appearance rectangle.

\section{Multi-Signature Support}\label{_multi_signature_support}

A PDF can contain multiple signatures. Each signature is an independent
signature field with its own \texttt{/Sig} dictionary and its own byte
range. The key insight: each subsequent signature covers the
\emph{entire} file as it existed after the previous signature, including
all previous signatures.

\subsection{How Multiple Signatures
Work}\label{_how_multiple_signatures_work}

Consider a document signed by Alice, then by Bob:

\begin{enumerate}
\def\labelenumi{\arabic{enumi}.}
\item
  Alice signs the original document. Her byte range covers the entire
  file (minus her \texttt{/Contents} gap). PDFBox writes this as an
  incremental update --- the original file plus a new cross-reference
  section and trailer.
\item
  Bob signs the document \emph{after} Alice. His byte range covers the
  entire file as it now exists --- including Alice's incremental update
  and her signature. Bob's signature is written as a \emph{second}
  incremental update appended after Alice's.
\end{enumerate}

Each signature's byte range covers everything that existed when that
signature was created. Alice's signature does not cover Bob's changes
(because Bob had not signed yet), but Bob's signature covers Alice's.

\subsection{Code: Adding a Second
Signature}\label{_code_adding_a_second_signature}

\begin{Shaded}
\begin{Highlighting}[]
\CommentTok{/**}
 \CommentTok{*}\NormalTok{ Adds a second signature to an already}\CommentTok{{-}}\NormalTok{signed PDF}\CommentTok{.}
 \CommentTok{*}\NormalTok{ The key is that we load the already}\CommentTok{{-}}\NormalTok{signed bytes and use saveIncremental}\CommentTok{()}
 \CommentTok{*}\NormalTok{ again }\CommentTok{{-}{-}}\NormalTok{ PDFBox appends another revision}\CommentTok{.}
 \CommentTok{*/}
\KeywordTok{public} \DataTypeTok{byte}\OperatorTok{[]} \FunctionTok{addSecondSignature}\OperatorTok{(}\DataTypeTok{byte}\OperatorTok{[]}\NormalTok{ alreadySignedPdf}\OperatorTok{,}
                                  \BuiltInTok{PrivateKey}\NormalTok{ privateKey}\OperatorTok{,}
                                  \BuiltInTok{X509Certificate}\NormalTok{ certificate}\OperatorTok{,}
                                  \BuiltInTok{List}\OperatorTok{\textless{}}\BuiltInTok{X509Certificate}\OperatorTok{\textgreater{}}\NormalTok{ chain}\OperatorTok{,}
                                  \BuiltInTok{String}\NormalTok{ fieldName}\OperatorTok{)}
        \KeywordTok{throws} \BuiltInTok{Exception} \OperatorTok{\{}

    \ControlFlowTok{try} \OperatorTok{(}\NormalTok{PDDocument doc }\OperatorTok{=}\NormalTok{ Loader}\OperatorTok{.}\FunctionTok{loadPDF}\OperatorTok{(}\NormalTok{alreadySignedPdf}\OperatorTok{))} \OperatorTok{\{}
        \CommentTok{// Find or create the second signature field}
\NormalTok{        PDAcroForm acroForm }\OperatorTok{=}\NormalTok{ doc}\OperatorTok{.}\FunctionTok{getDocumentCatalog}\OperatorTok{().}\FunctionTok{getAcroForm}\OperatorTok{();}
\NormalTok{        PDSignatureField sigField }\OperatorTok{=} \KeywordTok{null}\OperatorTok{;}

        \CommentTok{// Look for an existing empty signature field}
        \ControlFlowTok{for} \OperatorTok{(}\NormalTok{PDField field }\OperatorTok{:}\NormalTok{ acroForm}\OperatorTok{.}\FunctionTok{getFieldTree}\OperatorTok{())} \OperatorTok{\{}
            \ControlFlowTok{if} \OperatorTok{(}\NormalTok{field }\KeywordTok{instanceof}\NormalTok{ PDSignatureField sf}
                    \OperatorTok{\&\&}\NormalTok{ fieldName}\OperatorTok{.}\FunctionTok{equals}\OperatorTok{(}\NormalTok{sf}\OperatorTok{.}\FunctionTok{getFullyQualifiedName}\OperatorTok{())}
                    \OperatorTok{\&\&}\NormalTok{ sf}\OperatorTok{.}\FunctionTok{getSignature}\OperatorTok{()} \OperatorTok{==} \KeywordTok{null}\OperatorTok{)} \OperatorTok{\{}
\NormalTok{                sigField }\OperatorTok{=}\NormalTok{ sf}\OperatorTok{;}
                \ControlFlowTok{break}\OperatorTok{;}
            \OperatorTok{\}}
        \OperatorTok{\}}

        \ControlFlowTok{if} \OperatorTok{(}\NormalTok{sigField }\OperatorTok{==} \KeywordTok{null}\OperatorTok{)} \OperatorTok{\{}
            \ControlFlowTok{throw} \KeywordTok{new} \BuiltInTok{IllegalStateException}\OperatorTok{(}
                    \StringTok{"No empty signature field found: "} \OperatorTok{+}\NormalTok{ fieldName}\OperatorTok{);}
        \OperatorTok{\}}

        \CommentTok{// Build the second signature}
\NormalTok{        PDSignature signature }\OperatorTok{=} \KeywordTok{new} \FunctionTok{PDSignature}\OperatorTok{();}
\NormalTok{        signature}\OperatorTok{.}\FunctionTok{setFilter}\OperatorTok{(}\NormalTok{PDSignature}\OperatorTok{.}\FunctionTok{FILTER\_ADOBE\_PPKLITE}\OperatorTok{);}
\NormalTok{        signature}\OperatorTok{.}\FunctionTok{setSubFilter}\OperatorTok{(}
\NormalTok{                COSName}\OperatorTok{.}\FunctionTok{getPDFName}\OperatorTok{(}\StringTok{"ETSI.CAdES.detached"}\OperatorTok{));}
\NormalTok{        signature}\OperatorTok{.}\FunctionTok{setSignDate}\OperatorTok{(}\BuiltInTok{Calendar}\OperatorTok{.}\FunctionTok{getInstance}\OperatorTok{());}
\NormalTok{        signature}\OperatorTok{.}\FunctionTok{setName}\OperatorTok{(}\StringTok{"Bob Jones"}\OperatorTok{);}

\NormalTok{        sigField}\OperatorTok{.}\FunctionTok{setValue}\OperatorTok{(}\NormalTok{signature}\OperatorTok{);}

\NormalTok{        SignatureOptions options }\OperatorTok{=} \KeywordTok{new} \FunctionTok{SignatureOptions}\OperatorTok{();}
\NormalTok{        options}\OperatorTok{.}\FunctionTok{setPreferredSignatureSize}\OperatorTok{(}\DecValTok{32768}\OperatorTok{);}

\NormalTok{        SignatureInterface sigInterface }\OperatorTok{=}
                \KeywordTok{new} \FunctionTok{SimpleCmsSignatureInterface}\OperatorTok{(}
\NormalTok{                        privateKey}\OperatorTok{,}\NormalTok{ certificate}\OperatorTok{,}\NormalTok{ chain}\OperatorTok{);}
\NormalTok{        doc}\OperatorTok{.}\FunctionTok{addSignature}\OperatorTok{(}\NormalTok{signature}\OperatorTok{,}\NormalTok{ sigInterface}\OperatorTok{,}\NormalTok{ options}\OperatorTok{);}

        \BuiltInTok{ByteArrayOutputStream}\NormalTok{ baos }\OperatorTok{=} \KeywordTok{new} \BuiltInTok{ByteArrayOutputStream}\OperatorTok{();}
\NormalTok{        doc}\OperatorTok{.}\FunctionTok{saveIncremental}\OperatorTok{(}\NormalTok{baos}\OperatorTok{);} 
        \ControlFlowTok{return}\NormalTok{ baos}\OperatorTok{.}\FunctionTok{toByteArray}\OperatorTok{();}
    \OperatorTok{\}}
\OperatorTok{\}}
\end{Highlighting}
\end{Shaded}

\begin{itemize}
\item
  Each call to \texttt{saveIncremental()} appends a new revision. The
  file grows with each signature. A document with three signatures has
  four revisions (original + three incremental updates).
\end{itemize}

The order in which signatures are applied matters for validation. Adobe
Acrobat validates each signature against its revision --- it checks that
the bytes covered by Alice's byte range have not been modified. If
someone modified Alice's revision after she signed (other than by
appending a new incremental update), Acrobat flags Alice's signature as
invalid, even if Bob's is fine.

\section{External Signing: HSM and Remote Key
Flows}\label{_external_signing_hsm_and_remote_key_flows}

In many production environments, the private key never leaves a Hardware
Security Module (HSM), a cloud KMS, or a remote signing service. The key
cannot be exported, so you cannot pass it to a
\texttt{JcaContentSignerBuilder}. Instead, you need to split the signing
process: compute the hash locally, send the hash to the remote system,
receive the raw signature value, and assemble the CMS structure
yourself.

\subsection{The Hash-Signing Pattern}\label{_the_hash_signing_pattern}

The flow has three steps:

\begin{enumerate}
\def\labelenumi{\arabic{enumi}.}
\item
  \textbf{Prepare}: Create the PDF with a placeholder signature, compute
  the byte range, and extract the bytes to be signed.
\item
  \textbf{Hash}: Compute the SHA-256 digest of the byte range content.
\item
  \textbf{Sign remotely}: Send the hash (not the document) to the
  HSM/cloud service, receive the raw signature bytes, and embed them
  into the CMS structure.
\end{enumerate}

PDFBox supports this through a two-phase API using
\texttt{ExternalSigningSupport}:

\begin{Shaded}
\begin{Highlighting}[]
\KeywordTok{import} \ImportTok{org}\OperatorTok{.}\ImportTok{apache}\OperatorTok{.}\ImportTok{pdfbox}\OperatorTok{.}\ImportTok{pdmodel}\OperatorTok{.}\ImportTok{interactive}\OperatorTok{.}\ImportTok{digitalsignature}\OperatorTok{.}\ImportTok{ExternalSigningSupport}\OperatorTok{;}

\CommentTok{/**}
 \CommentTok{*}\NormalTok{ Signs a PDF using an external key source }\CommentTok{(}\NormalTok{HSM}\CommentTok{,}\NormalTok{ cloud KMS}\CommentTok{,}\NormalTok{ smart card}\CommentTok{).}
 \CommentTok{*}\NormalTok{ The private key never leaves the external system }\CommentTok{{-}{-}}\NormalTok{ only a hash is sent}\CommentTok{.}
 \CommentTok{*/}
\KeywordTok{public} \DataTypeTok{byte}\OperatorTok{[]} \FunctionTok{signWithExternalKey}\OperatorTok{(}\DataTypeTok{byte}\OperatorTok{[]}\NormalTok{ pdfBytes}\OperatorTok{,}
\NormalTok{                                   ExternalKeyService externalKeyService}\OperatorTok{,}
                                   \BuiltInTok{X509Certificate}\NormalTok{ certificate}\OperatorTok{,}
                                   \BuiltInTok{List}\OperatorTok{\textless{}}\BuiltInTok{X509Certificate}\OperatorTok{\textgreater{}}\NormalTok{ chain}\OperatorTok{)}
        \KeywordTok{throws} \BuiltInTok{Exception} \OperatorTok{\{}

    \BuiltInTok{ByteArrayOutputStream}\NormalTok{ result }\OperatorTok{=} \KeywordTok{new} \BuiltInTok{ByteArrayOutputStream}\OperatorTok{();}

    \ControlFlowTok{try} \OperatorTok{(}\NormalTok{PDDocument doc }\OperatorTok{=}\NormalTok{ Loader}\OperatorTok{.}\FunctionTok{loadPDF}\OperatorTok{(}\NormalTok{pdfBytes}\OperatorTok{))} \OperatorTok{\{}
\NormalTok{        PDSignature signature }\OperatorTok{=} \KeywordTok{new} \FunctionTok{PDSignature}\OperatorTok{();}
\NormalTok{        signature}\OperatorTok{.}\FunctionTok{setFilter}\OperatorTok{(}\NormalTok{PDSignature}\OperatorTok{.}\FunctionTok{FILTER\_ADOBE\_PPKLITE}\OperatorTok{);}
\NormalTok{        signature}\OperatorTok{.}\FunctionTok{setSubFilter}\OperatorTok{(}
\NormalTok{                COSName}\OperatorTok{.}\FunctionTok{getPDFName}\OperatorTok{(}\StringTok{"ETSI.CAdES.detached"}\OperatorTok{));}
\NormalTok{        signature}\OperatorTok{.}\FunctionTok{setSignDate}\OperatorTok{(}\BuiltInTok{Calendar}\OperatorTok{.}\FunctionTok{getInstance}\OperatorTok{());}

\NormalTok{        SignatureOptions options }\OperatorTok{=} \KeywordTok{new} \FunctionTok{SignatureOptions}\OperatorTok{();}
\NormalTok{        options}\OperatorTok{.}\FunctionTok{setPreferredSignatureSize}\OperatorTok{(}\DecValTok{32768}\OperatorTok{);}

        \CommentTok{// Register signature WITHOUT a SignatureInterface}
\NormalTok{        doc}\OperatorTok{.}\FunctionTok{addSignature}\OperatorTok{(}\NormalTok{signature}\OperatorTok{,}\NormalTok{ options}\OperatorTok{);}

        \CommentTok{// saveIncrementalForExternalSigning writes the PDF with a}
        \CommentTok{// placeholder and returns an ExternalSigningSupport object}
\NormalTok{        ExternalSigningSupport externalSigning }\OperatorTok{=}
\NormalTok{                doc}\OperatorTok{.}\FunctionTok{saveIncrementalForExternalSigning}\OperatorTok{(}\NormalTok{result}\OperatorTok{);} 

        \CommentTok{// Get the byte range content}
        \BuiltInTok{InputStream}\NormalTok{ content }\OperatorTok{=}\NormalTok{ externalSigning}\OperatorTok{.}\FunctionTok{getContent}\OperatorTok{();} 
        \DataTypeTok{byte}\OperatorTok{[]}\NormalTok{ contentBytes }\OperatorTok{=}\NormalTok{ content}\OperatorTok{.}\FunctionTok{readAllBytes}\OperatorTok{();}

        \CommentTok{// Compute the hash locally}
        \BuiltInTok{MessageDigest}\NormalTok{ digest }\OperatorTok{=} \BuiltInTok{MessageDigest}\OperatorTok{.}\FunctionTok{getInstance}\OperatorTok{(}\StringTok{"SHA{-}256"}\OperatorTok{);}
        \DataTypeTok{byte}\OperatorTok{[]}\NormalTok{ hash }\OperatorTok{=}\NormalTok{ digest}\OperatorTok{.}\FunctionTok{digest}\OperatorTok{(}\NormalTok{contentBytes}\OperatorTok{);}

        \CommentTok{// Send hash to external system, receive raw signature value}
        \DataTypeTok{byte}\OperatorTok{[]}\NormalTok{ rawSignature }\OperatorTok{=}
\NormalTok{                externalKeyService}\OperatorTok{.}\FunctionTok{signHash}\OperatorTok{(}\NormalTok{hash}\OperatorTok{,} \StringTok{"SHA256withRSA"}\OperatorTok{);} 

        \CommentTok{// Build CMS structure around the raw signature}
        \DataTypeTok{byte}\OperatorTok{[]}\NormalTok{ cmsBytes }\OperatorTok{=} \FunctionTok{buildCmsWithExternalSignature}\OperatorTok{(}
\NormalTok{                rawSignature}\OperatorTok{,}\NormalTok{ contentBytes}\OperatorTok{,}\NormalTok{ certificate}\OperatorTok{,}\NormalTok{ chain}\OperatorTok{);}

        \CommentTok{// Write the CMS bytes into the placeholder}
\NormalTok{        externalSigning}\OperatorTok{.}\FunctionTok{setSignature}\OperatorTok{(}\NormalTok{cmsBytes}\OperatorTok{);} 
    \OperatorTok{\}}

    \ControlFlowTok{return}\NormalTok{ result}\OperatorTok{.}\FunctionTok{toByteArray}\OperatorTok{();}
\OperatorTok{\}}
\end{Highlighting}
\end{Shaded}

\begin{itemize}
\item
  \texttt{saveIncrementalForExternalSigning()} writes the PDF to the
  output stream with a placeholder \texttt{/Contents} and returns an
  \texttt{ExternalSigningSupport} handle.
\item
  \texttt{getContent()} returns the byte range content as an
  \texttt{InputStream} --- the same data that would be passed to
  \texttt{sign()} in the callback-based approach.
\item
  This is the only operation that touches the external key. You send a
  32-byte SHA-256 hash and receive a raw signature value (typically 256
  bytes for RSA-2048 or 64 bytes for ECDSA P-256).
\item
  \texttt{setSignature()} hex-encodes the CMS bytes and patches them
  into the \texttt{/Contents} placeholder.
\end{itemize}

Building the CMS structure with a pre-computed raw signature is more
involved than using \texttt{CMSSignedDataGenerator}. You need to
construct the \texttt{SignerInfo} manually with the external signature
value. Chapter 30 shows the complete code for this.

\section{Certification Signatures with
DocMDP}\label{_certification_signatures_with_docmdp}

A certification signature (also called a "certifying signature" or "MDP
signature") is a special first signature that declares permission levels
for future modifications to the document. It uses the \texttt{/DocMDP}
transform method (ISO 32000-2:2020 S12.8.2.2).

\subsection{The Three DocMDP Levels}\label{_the_three_docmdp_levels}

\begin{longtable}[]{@{}
  >{\raggedright\arraybackslash}p{(\linewidth - 4\tabcolsep) * \real{0.1111}}
  >{\raggedright\arraybackslash}p{(\linewidth - 4\tabcolsep) * \real{0.3333}}
  >{\raggedright\arraybackslash}p{(\linewidth - 4\tabcolsep) * \real{0.5556}}@{}}
\caption{DocMDP Permission Levels --- ISO 32000-2:2020 S12.8.2.2, Table
257}\tabularnewline
\toprule\noalign{}
\begin{minipage}[b]{\linewidth}\raggedright
Level
\end{minipage} & \begin{minipage}[b]{\linewidth}\raggedright
Name
\end{minipage} & \begin{minipage}[b]{\linewidth}\raggedright
What Is Allowed After Certification
\end{minipage} \\
\midrule\noalign{}
\endfirsthead
\toprule\noalign{}
\begin{minipage}[b]{\linewidth}\raggedright
Level
\end{minipage} & \begin{minipage}[b]{\linewidth}\raggedright
Name
\end{minipage} & \begin{minipage}[b]{\linewidth}\raggedright
What Is Allowed After Certification
\end{minipage} \\
\midrule\noalign{}
\endhead
\bottomrule\noalign{}
\endlastfoot
1 & No changes & No modifications whatsoever. Even filling in form
fields invalidates the signature. \\
2 & Form fill + sign & Form fields may be filled, signature fields may
be signed, and pages may be instantiated from templates. No other
changes. \\
3 & Form fill + sign + annotate & Everything in level 2, plus
annotations may be added, modified, or deleted. \\
\end{longtable}

Level 2 is by far the most common in practice --- it allows a document
to be certified by the author and then signed by additional parties.

\subsection{Code: Applying a Certification
Signature}\label{_code_applying_a_certification_signature}

\begin{Shaded}
\begin{Highlighting}[]
\KeywordTok{import} \ImportTok{org}\OperatorTok{.}\ImportTok{apache}\OperatorTok{.}\ImportTok{pdfbox}\OperatorTok{.}\ImportTok{cos}\OperatorTok{.}\ImportTok{COSArray}\OperatorTok{;}
\KeywordTok{import} \ImportTok{org}\OperatorTok{.}\ImportTok{apache}\OperatorTok{.}\ImportTok{pdfbox}\OperatorTok{.}\ImportTok{cos}\OperatorTok{.}\ImportTok{COSDictionary}\OperatorTok{;}
\KeywordTok{import} \ImportTok{org}\OperatorTok{.}\ImportTok{apache}\OperatorTok{.}\ImportTok{pdfbox}\OperatorTok{.}\ImportTok{cos}\OperatorTok{.}\ImportTok{COSInteger}\OperatorTok{;}
\KeywordTok{import} \ImportTok{org}\OperatorTok{.}\ImportTok{apache}\OperatorTok{.}\ImportTok{pdfbox}\OperatorTok{.}\ImportTok{cos}\OperatorTok{.}\ImportTok{COSName}\OperatorTok{;}

\CommentTok{/**}
 \CommentTok{*}\NormalTok{ Applies a certification }\CommentTok{(}\NormalTok{DocMDP}\CommentTok{)}\NormalTok{ signature to a PDF}\CommentTok{.}
 \CommentTok{*}\NormalTok{ This must be the FIRST signature in the document}\CommentTok{.}
 \CommentTok{*/}
\KeywordTok{public} \DataTypeTok{void} \FunctionTok{applyCertificationSignature}\OperatorTok{(}\NormalTok{PDDocument doc}\OperatorTok{,}
\NormalTok{                                         PDSignature signature}\OperatorTok{,}
\NormalTok{                                         PDSignatureField sigField}\OperatorTok{,}
                                         \DataTypeTok{int}\NormalTok{ docMdpLevel}\OperatorTok{)} \OperatorTok{\{}

    \CommentTok{// Clamp level to valid range}
    \DataTypeTok{int}\NormalTok{ level }\OperatorTok{=} \OperatorTok{(}\NormalTok{docMdpLevel }\OperatorTok{\textgreater{}=} \DecValTok{1} \OperatorTok{\&\&}\NormalTok{ docMdpLevel }\OperatorTok{\textless{}=} \DecValTok{3}\OperatorTok{)}
            \OperatorTok{?}\NormalTok{ docMdpLevel }\OperatorTok{:} \DecValTok{2}\OperatorTok{;}

    \CommentTok{// Build the /Reference array with a DocMDP transform}
\NormalTok{    COSDictionary sigRef }\OperatorTok{=} \KeywordTok{new} \FunctionTok{COSDictionary}\OperatorTok{();}
\NormalTok{    sigRef}\OperatorTok{.}\FunctionTok{setItem}\OperatorTok{(}\NormalTok{COSName}\OperatorTok{.}\FunctionTok{TYPE}\OperatorTok{,}\NormalTok{ COSName}\OperatorTok{.}\FunctionTok{getPDFName}\OperatorTok{(}\StringTok{"SigRef"}\OperatorTok{));}
\NormalTok{    sigRef}\OperatorTok{.}\FunctionTok{setItem}\OperatorTok{(}\NormalTok{COSName}\OperatorTok{.}\FunctionTok{getPDFName}\OperatorTok{(}\StringTok{"TransformMethod"}\OperatorTok{),}
\NormalTok{            COSName}\OperatorTok{.}\FunctionTok{getPDFName}\OperatorTok{(}\StringTok{"DocMDP"}\OperatorTok{));}

\NormalTok{    COSDictionary transformParams }\OperatorTok{=} \KeywordTok{new} \FunctionTok{COSDictionary}\OperatorTok{();}
\NormalTok{    transformParams}\OperatorTok{.}\FunctionTok{setItem}\OperatorTok{(}\NormalTok{COSName}\OperatorTok{.}\FunctionTok{TYPE}\OperatorTok{,}
\NormalTok{            COSName}\OperatorTok{.}\FunctionTok{getPDFName}\OperatorTok{(}\StringTok{"TransformParams"}\OperatorTok{));}
\NormalTok{    transformParams}\OperatorTok{.}\FunctionTok{setItem}\OperatorTok{(}\NormalTok{COSName}\OperatorTok{.}\FunctionTok{getPDFName}\OperatorTok{(}\StringTok{"P"}\OperatorTok{),}
\NormalTok{            COSInteger}\OperatorTok{.}\FunctionTok{get}\OperatorTok{(}\NormalTok{level}\OperatorTok{));} 
\NormalTok{    transformParams}\OperatorTok{.}\FunctionTok{setItem}\OperatorTok{(}\NormalTok{COSName}\OperatorTok{.}\FunctionTok{getPDFName}\OperatorTok{(}\StringTok{"V"}\OperatorTok{),}
\NormalTok{            COSName}\OperatorTok{.}\FunctionTok{getPDFName}\OperatorTok{(}\StringTok{"1.2"}\OperatorTok{));}
\NormalTok{    sigRef}\OperatorTok{.}\FunctionTok{setItem}\OperatorTok{(}\NormalTok{COSName}\OperatorTok{.}\FunctionTok{getPDFName}\OperatorTok{(}\StringTok{"TransformParams"}\OperatorTok{),}
\NormalTok{            transformParams}\OperatorTok{);}

\NormalTok{    COSArray refArray }\OperatorTok{=} \KeywordTok{new} \FunctionTok{COSArray}\OperatorTok{();}
\NormalTok{    refArray}\OperatorTok{.}\FunctionTok{add}\OperatorTok{(}\NormalTok{sigRef}\OperatorTok{);}
\NormalTok{    signature}\OperatorTok{.}\FunctionTok{getCOSObject}\OperatorTok{()}
            \OperatorTok{.}\FunctionTok{setItem}\OperatorTok{(}\NormalTok{COSName}\OperatorTok{.}\FunctionTok{getPDFName}\OperatorTok{(}\StringTok{"Reference"}\OperatorTok{),}\NormalTok{ refArray}\OperatorTok{);}

    \CommentTok{// Set /Perms /DocMDP in the document catalog {-}{-} this is what}
    \CommentTok{// tells viewers "this document has a certification signature"}
\NormalTok{    COSDictionary perms }\OperatorTok{=} \KeywordTok{new} \FunctionTok{COSDictionary}\OperatorTok{();}
\NormalTok{    perms}\OperatorTok{.}\FunctionTok{setItem}\OperatorTok{(}\NormalTok{COSName}\OperatorTok{.}\FunctionTok{getPDFName}\OperatorTok{(}\StringTok{"DocMDP"}\OperatorTok{),}
\NormalTok{            signature}\OperatorTok{.}\FunctionTok{getCOSObject}\OperatorTok{());} 
\NormalTok{    doc}\OperatorTok{.}\FunctionTok{getDocumentCatalog}\OperatorTok{().}\FunctionTok{getCOSObject}\OperatorTok{()}
            \OperatorTok{.}\FunctionTok{setItem}\OperatorTok{(}\NormalTok{COSName}\OperatorTok{.}\FunctionTok{getPDFName}\OperatorTok{(}\StringTok{"Perms"}\OperatorTok{),}\NormalTok{ perms}\OperatorTok{);}
\OperatorTok{\}}
\end{Highlighting}
\end{Shaded}

\begin{itemize}
\item
  The \texttt{/P} value is the permission level: 1, 2, or 3.
\item
  The \texttt{/Perms} dictionary in the document catalog must reference
  the certification signature's \texttt{/Sig} dictionary. This is how
  viewers discover that the document has been certified.
\end{itemize}

A document can have at most \emph{one} certification signature, and it
must be the first signature applied. If you try to add a certification
signature to a document that already has one, the existing certification
is invalidated. If you apply it as a second signature (after an approval
signature), the result is undefined --- some viewers accept it, some
reject it, and Adobe Acrobat shows a confusing error.

\section{PDFBox 2.x to 3.x
Migration}\label{_pdfbox_2_x_to_3_x_migration}

PDFBox 3.0.0 introduced breaking API changes that affect every piece of
signing code. If you are migrating from PDFBox 2.x (which many projects
still use), these are the changes you need to know.

\subsection{Document Loading}\label{_document_loading}

PDFBox 2.x (removed)

\begin{Shaded}
\begin{Highlighting}[]
\CommentTok{// PDFBox 2.x {-}{-} these methods no longer exist}
\NormalTok{PDDocument doc }\OperatorTok{=}\NormalTok{ PDDocument}\OperatorTok{.}\FunctionTok{load}\OperatorTok{(}\KeywordTok{new} \BuiltInTok{File}\OperatorTok{(}\StringTok{"doc.pdf"}\OperatorTok{));}
\NormalTok{PDDocument doc }\OperatorTok{=}\NormalTok{ PDDocument}\OperatorTok{.}\FunctionTok{load}\OperatorTok{(}\NormalTok{inputStream}\OperatorTok{);}
\NormalTok{PDDocument doc }\OperatorTok{=}\NormalTok{ PDDocument}\OperatorTok{.}\FunctionTok{load}\OperatorTok{(}\NormalTok{byteArray}\OperatorTok{);}
\end{Highlighting}
\end{Shaded}

PDFBox 3.x (current)

\begin{Shaded}
\begin{Highlighting}[]
\CommentTok{// PDFBox 3.x {-}{-} use the Loader class}
\NormalTok{PDDocument doc }\OperatorTok{=}\NormalTok{ Loader}\OperatorTok{.}\FunctionTok{loadPDF}\OperatorTok{(}\KeywordTok{new} \BuiltInTok{File}\OperatorTok{(}\StringTok{"doc.pdf"}\OperatorTok{));}
\NormalTok{PDDocument doc }\OperatorTok{=}\NormalTok{ Loader}\OperatorTok{.}\FunctionTok{loadPDF}\OperatorTok{(}\NormalTok{byteArray}\OperatorTok{);}
\CommentTok{// No direct InputStream loading {-}{-} read to byte[] first}
\DataTypeTok{byte}\OperatorTok{[]}\NormalTok{ bytes }\OperatorTok{=}\NormalTok{ inputStream}\OperatorTok{.}\FunctionTok{readAllBytes}\OperatorTok{();}
\NormalTok{PDDocument doc }\OperatorTok{=}\NormalTok{ Loader}\OperatorTok{.}\FunctionTok{loadPDF}\OperatorTok{(}\NormalTok{bytes}\OperatorTok{);}
\end{Highlighting}
\end{Shaded}

The \texttt{Loader} class is in \texttt{org.apache.pdfbox.Loader}. The
old \texttt{PDDocument.load()} static methods were removed
entirely --- not deprecated, \emph{removed}.

\subsection{Standard 14 Fonts}\label{_standard_14_fonts}

This is the single most common compilation error in PDFBox 2.x-to-3.x
migration.

PDFBox 2.x (removed)

\begin{Shaded}
\begin{Highlighting}[]
\CommentTok{// PDFBox 2.x {-}{-} static font constants}
\NormalTok{PDFont font }\OperatorTok{=}\NormalTok{ PDType1Font}\OperatorTok{.}\FunctionTok{HELVETICA}\OperatorTok{;}
\NormalTok{PDFont fontBold }\OperatorTok{=}\NormalTok{ PDType1Font}\OperatorTok{.}\FunctionTok{HELVETICA\_BOLD}\OperatorTok{;}
\NormalTok{PDFont fontMono }\OperatorTok{=}\NormalTok{ PDType1Font}\OperatorTok{.}\FunctionTok{COURIER}\OperatorTok{;}
\end{Highlighting}
\end{Shaded}

PDFBox 3.x (current)

\begin{Shaded}
\begin{Highlighting}[]
\CommentTok{// PDFBox 3.x {-}{-} constructor with FontName enum}
\NormalTok{PDFont font }\OperatorTok{=} \KeywordTok{new} \FunctionTok{PDType1Font}\OperatorTok{(}\NormalTok{Standard14Fonts}\OperatorTok{.}\FunctionTok{FontName}\OperatorTok{.}\FunctionTok{HELVETICA}\OperatorTok{);}
\NormalTok{PDFont fontBold }\OperatorTok{=} \KeywordTok{new} \FunctionTok{PDType1Font}\OperatorTok{(}\NormalTok{Standard14Fonts}\OperatorTok{.}\FunctionTok{FontName}\OperatorTok{.}\FunctionTok{HELVETICA\_BOLD}\OperatorTok{);}
\NormalTok{PDFont fontMono }\OperatorTok{=} \KeywordTok{new} \FunctionTok{PDType1Font}\OperatorTok{(}\NormalTok{Standard14Fonts}\OperatorTok{.}\FunctionTok{FontName}\OperatorTok{.}\FunctionTok{COURIER}\OperatorTok{);}
\end{Highlighting}
\end{Shaded}

The \texttt{Standard14Fonts.FontName} enum lives in
\texttt{org.apache.pdfbox.pdmodel.font.Standard14Fonts}.

\subsection{AcroForm API Changes}\label{_acroform_api_changes}

PDFBox 2.x

\begin{Shaded}
\begin{Highlighting}[]
\DataTypeTok{boolean}\NormalTok{ needsAppearances }\OperatorTok{=}\NormalTok{ acroForm}\OperatorTok{.}\FunctionTok{isNeedAppearances}\OperatorTok{();} \CommentTok{// gone}
\end{Highlighting}
\end{Shaded}

PDFBox 3.x

\begin{Shaded}
\begin{Highlighting}[]
\DataTypeTok{boolean}\NormalTok{ needsAppearances }\OperatorTok{=}\NormalTok{ acroForm}\OperatorTok{.}\FunctionTok{getNeedAppearances}\OperatorTok{();} 
\end{Highlighting}
\end{Shaded}

\begin{itemize}
\item
  The method was renamed from \texttt{isNeedAppearances()} to
  \texttt{getNeedAppearances()}. A small change, but it breaks at
  compile time.
\end{itemize}

\subsection{Content Stream Writing}\label{_content_stream_writing}

PDFBox 2.x --- could use PDPageContentStream for Form XObjects

\begin{Shaded}
\begin{Highlighting}[]
\CommentTok{// This approach no longer works for Form XObjects in 3.x}
\end{Highlighting}
\end{Shaded}

PDFBox 3.x --- use the stream's output stream directly

\begin{Shaded}
\begin{Highlighting}[]
\NormalTok{PDFormXObject form }\OperatorTok{=} \KeywordTok{new} \FunctionTok{PDFormXObject}\OperatorTok{(}\NormalTok{doc}\OperatorTok{);}
\ControlFlowTok{try} \OperatorTok{(}\BuiltInTok{OutputStream}\NormalTok{ os }\OperatorTok{=}\NormalTok{ form}\OperatorTok{.}\FunctionTok{getStream}\OperatorTok{().}\FunctionTok{createOutputStream}\OperatorTok{())} \OperatorTok{\{}
\NormalTok{    os}\OperatorTok{.}\FunctionTok{write}\OperatorTok{(}\NormalTok{contentStreamBytes}\OperatorTok{);}
\OperatorTok{\}}
\end{Highlighting}
\end{Shaded}

\subsection{Complete Migration
Checklist}\label{_complete_migration_checklist}

\begin{longtable}[]{@{}
  >{\raggedright\arraybackslash}p{(\linewidth - 4\tabcolsep) * \real{0.3636}}
  >{\raggedright\arraybackslash}p{(\linewidth - 4\tabcolsep) * \real{0.3636}}
  >{\raggedright\arraybackslash}p{(\linewidth - 4\tabcolsep) * \real{0.2727}}@{}}
\caption{PDFBox 2.x to 3.x Breaking Changes Relevant to
Signing}\tabularnewline
\toprule\noalign{}
\begin{minipage}[b]{\linewidth}\raggedright
PDFBox 2.x
\end{minipage} & \begin{minipage}[b]{\linewidth}\raggedright
PDFBox 3.x
\end{minipage} & \begin{minipage}[b]{\linewidth}\raggedright
Impact
\end{minipage} \\
\midrule\noalign{}
\endfirsthead
\toprule\noalign{}
\begin{minipage}[b]{\linewidth}\raggedright
PDFBox 2.x
\end{minipage} & \begin{minipage}[b]{\linewidth}\raggedright
PDFBox 3.x
\end{minipage} & \begin{minipage}[b]{\linewidth}\raggedright
Impact
\end{minipage} \\
\midrule\noalign{}
\endhead
\bottomrule\noalign{}
\endlastfoot
\texttt{PDDocument.load(File)} & \texttt{Loader.loadPDF(File)} & Compile
error \\
\texttt{PDDocument.load(byte{[}{]})} &
\texttt{Loader.loadPDF(byte{[}{]})} & Compile error \\
\texttt{PDDocument.load(InputStream)} & \texttt{Loader.loadPDF(bytes)}
(read first) & Compile error \\
\texttt{PDType1Font.HELVETICA} &
\texttt{new\ PDType1Font(Standard14Fonts.FontName.HELVETICA)} & Compile
error \\
\texttt{acroForm.isNeedAppearances()} &
\texttt{acroForm.getNeedAppearances()} & Compile error \\
\texttt{new\ PDPageContentStream(doc,\ formXObject)} &
\texttt{form.getStream().createOutputStream()} & Compile error \\
\texttt{PDFont.getEncoding()} & Not available --- use COS dictionary &
Compile error \\
\texttt{PDFont.getToUnicodeCMap()} & Not visible --- check COS key &
Compile error \\
\end{longtable}

When migrating a large codebase, search for every occurrence of
\texttt{PDDocument.load}, \texttt{PDType1Font.HELVETICA} (and all other
font constants), and \texttt{isNeedAppearances}. These three account for
about 90\% of compilation failures in signing code.

\section{Common Pitfalls}\label{_common_pitfalls_3}

We have collected these from production incidents, Stack Overflow
questions, and our own debugging sessions. Each one has cost someone
hours.

\subsection{Pitfall 1: Not Registering the BouncyCastle
Provider}\label{_pitfall_1_not_registering_the_bouncycastle_provider}

\begin{Shaded}
\begin{Highlighting}[]
\CommentTok{// This line MUST execute before any signing operation}
\NormalTok{java}\OperatorTok{.}\FunctionTok{security}\OperatorTok{.}\FunctionTok{Security}\OperatorTok{.}\FunctionTok{addProvider}\OperatorTok{(}
        \KeywordTok{new}\NormalTok{ org}\OperatorTok{.}\FunctionTok{bouncycastle}\OperatorTok{.}\FunctionTok{jce}\OperatorTok{.}\FunctionTok{provider}\OperatorTok{.}\FunctionTok{BouncyCastleProvider}\OperatorTok{());}
\end{Highlighting}
\end{Shaded}

Without this, you get a \texttt{NoSuchAlgorithmException} or a
\texttt{NoSuchProviderException} at runtime. The error message varies
depending on which BouncyCastle class triggers first, making it
confusing to diagnose. In Spring Boot applications, put this in a
\texttt{@PostConstruct} method of a \texttt{@Configuration} class:

\begin{Shaded}
\begin{Highlighting}[]
\KeywordTok{import} \ImportTok{jakarta}\OperatorTok{.}\ImportTok{annotation}\OperatorTok{.}\ImportTok{PostConstruct}\OperatorTok{;}
\KeywordTok{import} \ImportTok{org}\OperatorTok{.}\ImportTok{bouncycastle}\OperatorTok{.}\ImportTok{jce}\OperatorTok{.}\ImportTok{provider}\OperatorTok{.}\ImportTok{BouncyCastleProvider}\OperatorTok{;}
\KeywordTok{import} \ImportTok{org}\OperatorTok{.}\ImportTok{springframework}\OperatorTok{.}\ImportTok{context}\OperatorTok{.}\ImportTok{annotation}\OperatorTok{.}\ImportTok{Configuration}\OperatorTok{;}

\AttributeTok{@Configuration}
\KeywordTok{public} \KeywordTok{class}\NormalTok{ BouncyCastleConfig }\OperatorTok{\{}

    \AttributeTok{@PostConstruct}
    \KeywordTok{public} \DataTypeTok{void} \FunctionTok{registerProvider}\OperatorTok{()} \OperatorTok{\{}
        \ControlFlowTok{if} \OperatorTok{(}\NormalTok{java}\OperatorTok{.}\FunctionTok{security}\OperatorTok{.}\FunctionTok{Security}\OperatorTok{.}\FunctionTok{getProvider}\OperatorTok{(}\StringTok{"BC"}\OperatorTok{)} \OperatorTok{==} \KeywordTok{null}\OperatorTok{)} \OperatorTok{\{}
\NormalTok{            java}\OperatorTok{.}\FunctionTok{security}\OperatorTok{.}\FunctionTok{Security}\OperatorTok{.}\FunctionTok{addProvider}\OperatorTok{(}
                    \KeywordTok{new} \FunctionTok{BouncyCastleProvider}\OperatorTok{());}
        \OperatorTok{\}}
    \OperatorTok{\}}
\OperatorTok{\}}
\end{Highlighting}
\end{Shaded}

\subsection{Pitfall 2: Using save() Instead of
saveIncremental()}\label{_pitfall_2_using_save_instead_of_saveincremental}

\begin{Shaded}
\begin{Highlighting}[]
\CommentTok{// WRONG {-}{-} destroys all existing signatures}
\NormalTok{doc}\OperatorTok{.}\FunctionTok{save}\OperatorTok{(}\NormalTok{outputStream}\OperatorTok{);}

\CommentTok{// CORRECT {-}{-} preserves existing signatures}
\NormalTok{doc}\OperatorTok{.}\FunctionTok{saveIncremental}\OperatorTok{(}\NormalTok{outputStream}\OperatorTok{);}
\end{Highlighting}
\end{Shaded}

A regular \texttt{save()} rewrites the entire PDF from scratch. This
changes every byte offset in the file, which means every existing
signature's byte range is now wrong. Every signature becomes invalid.
Use \texttt{saveIncremental()}. Always.

\subsection{Pitfall 3: /Contents Container Too
Small}\label{_pitfall_3_contents_container_too_small}

If the CMS blob returned by your \texttt{sign()} method exceeds the
space reserved by \texttt{preferredSignatureSize}, you get:

\begin{verbatim}
java.io.IOException: Can't write signature, not enough space
\end{verbatim}

The fix is straightforward --- increase the container size. But be aware
that the container size is in \emph{bytes of hex encoding}, not bytes of
CMS data. A 32768-byte container holds approximately 16 KB of CMS data
(since each byte becomes two hex characters).

When in doubt, use 65536 (64 KB). The extra space is filled with hex
zeros and adds about 64 KB to the file size --- negligible for most
PDFs.

\subsection{Pitfall 4: Not Closing the
PDDocument}\label{_pitfall_4_not_closing_the_pddocument}

\begin{Shaded}
\begin{Highlighting}[]
\CommentTok{// WRONG {-}{-} resource leak, and on some JVMs the temp files}
\CommentTok{// created by PDFBox are never cleaned up}
\NormalTok{PDDocument doc }\OperatorTok{=}\NormalTok{ Loader}\OperatorTok{.}\FunctionTok{loadPDF}\OperatorTok{(}\NormalTok{bytes}\OperatorTok{);}
\NormalTok{doc}\OperatorTok{.}\FunctionTok{addSignature}\OperatorTok{(}\NormalTok{sig}\OperatorTok{,}\NormalTok{ iface}\OperatorTok{,}\NormalTok{ opts}\OperatorTok{);}
\NormalTok{doc}\OperatorTok{.}\FunctionTok{saveIncremental}\OperatorTok{(}\NormalTok{baos}\OperatorTok{);}
\CommentTok{// doc is never closed!}

\CommentTok{// CORRECT {-}{-} use try{-}with{-}resources}
\ControlFlowTok{try} \OperatorTok{(}\NormalTok{PDDocument doc }\OperatorTok{=}\NormalTok{ Loader}\OperatorTok{.}\FunctionTok{loadPDF}\OperatorTok{(}\NormalTok{bytes}\OperatorTok{))} \OperatorTok{\{}
\NormalTok{    doc}\OperatorTok{.}\FunctionTok{addSignature}\OperatorTok{(}\NormalTok{sig}\OperatorTok{,}\NormalTok{ iface}\OperatorTok{,}\NormalTok{ opts}\OperatorTok{);}
\NormalTok{    doc}\OperatorTok{.}\FunctionTok{saveIncremental}\OperatorTok{(}\NormalTok{baos}\OperatorTok{);}
\OperatorTok{\}}
\end{Highlighting}
\end{Shaded}

\texttt{PDDocument} implements \texttt{Closeable}. If you do not close
it, temporary files and memory-mapped buffers may not be released. In a
server application that signs thousands of documents, this leads to
\texttt{OutOfMemoryError} or file handle exhaustion.

\subsection{Pitfall 5: Signing After Closing the
Document}\label{_pitfall_5_signing_after_closing_the_document}

\begin{Shaded}
\begin{Highlighting}[]
\CommentTok{// WRONG {-}{-} the document is closed before saveIncremental()}
\NormalTok{PDDocument doc }\OperatorTok{=}\NormalTok{ Loader}\OperatorTok{.}\FunctionTok{loadPDF}\OperatorTok{(}\NormalTok{bytes}\OperatorTok{);}
\NormalTok{doc}\OperatorTok{.}\FunctionTok{addSignature}\OperatorTok{(}\NormalTok{sig}\OperatorTok{,}\NormalTok{ iface}\OperatorTok{,}\NormalTok{ opts}\OperatorTok{);}
\NormalTok{doc}\OperatorTok{.}\FunctionTok{close}\OperatorTok{();}
\NormalTok{doc}\OperatorTok{.}\FunctionTok{saveIncremental}\OperatorTok{(}\NormalTok{baos}\OperatorTok{);} \CommentTok{// throws IllegalStateException}
\end{Highlighting}
\end{Shaded}

\texttt{addSignature()} only registers the signature. The actual signing
happens inside \texttt{saveIncremental()}. If you close the document
before calling \texttt{saveIncremental()}, the document's internal
structures are released, and the save fails.

\subsection{Pitfall 6: Reusing a PDSignature
Object}\label{_pitfall_6_reusing_a_pdsignature_object}

\begin{Shaded}
\begin{Highlighting}[]
\CommentTok{// WRONG {-}{-} reusing the same PDSignature for two documents}
\NormalTok{PDSignature sig }\OperatorTok{=} \KeywordTok{new} \FunctionTok{PDSignature}\OperatorTok{();}
\NormalTok{sig}\OperatorTok{.}\FunctionTok{setFilter}\OperatorTok{(}\NormalTok{PDSignature}\OperatorTok{.}\FunctionTok{FILTER\_ADOBE\_PPKLITE}\OperatorTok{);}

\NormalTok{doc1}\OperatorTok{.}\FunctionTok{addSignature}\OperatorTok{(}\NormalTok{sig}\OperatorTok{,}\NormalTok{ iface}\OperatorTok{,}\NormalTok{ opts}\OperatorTok{);}
\NormalTok{doc1}\OperatorTok{.}\FunctionTok{saveIncremental}\OperatorTok{(}\NormalTok{out1}\OperatorTok{);}

\NormalTok{doc2}\OperatorTok{.}\FunctionTok{addSignature}\OperatorTok{(}\NormalTok{sig}\OperatorTok{,}\NormalTok{ iface}\OperatorTok{,}\NormalTok{ opts}\OperatorTok{);} \CommentTok{// \textless{}{-}{-} BROKEN}
\NormalTok{doc2}\OperatorTok{.}\FunctionTok{saveIncremental}\OperatorTok{(}\NormalTok{out2}\OperatorTok{);}
\end{Highlighting}
\end{Shaded}

A \texttt{PDSignature} is a COS dictionary tied to a specific document.
After it is added to \texttt{doc1} and saved, it belongs to
\texttt{doc1}. Create a new \texttt{PDSignature} for each document.

\section{Timestamp Embedding}\label{_timestamp_embedding}

A timestamp proves when a signature was created. Without it, you have a
PAdES-B-B signature. With it, you have PAdES-B-T.

\subsection{Embedding a TSA Timestamp as an Unsigned
Attribute}\label{_embedding_a_tsa_timestamp_as_an_unsigned_attribute}

The timestamp token is an RFC 3161 \texttt{TimeStampToken} embedded as
an unsigned attribute (\texttt{id-aa-signatureTimeStampToken}, OID
1.2.840.113549.1.9.16.2.14) in the CMS \texttt{SignerInfo}.

The flow is:

\begin{enumerate}
\def\labelenumi{\arabic{enumi}.}
\item
  Generate the CMS \texttt{SignedData} (without the timestamp)
\item
  Extract the signer's raw signature value from the \texttt{SignerInfo}
\item
  Hash the signature value with SHA-256
\item
  Send the hash to the TSA, receive a \texttt{TimeStampToken}
\item
  Attach the token as an unsigned attribute on the \texttt{SignerInfo}
\item
  Return the modified CMS \texttt{SignedData}
\end{enumerate}

\begin{Shaded}
\begin{Highlighting}[]
\KeywordTok{import} \ImportTok{org}\OperatorTok{.}\ImportTok{bouncycastle}\OperatorTok{.}\ImportTok{asn1}\OperatorTok{.}\ImportTok{ASN1EncodableVector}\OperatorTok{;}
\KeywordTok{import} \ImportTok{org}\OperatorTok{.}\ImportTok{bouncycastle}\OperatorTok{.}\ImportTok{asn1}\OperatorTok{.}\ImportTok{ASN1Primitive}\OperatorTok{;}
\KeywordTok{import} \ImportTok{org}\OperatorTok{.}\ImportTok{bouncycastle}\OperatorTok{.}\ImportTok{asn1}\OperatorTok{.}\ImportTok{DERSet}\OperatorTok{;}
\KeywordTok{import} \ImportTok{org}\OperatorTok{.}\ImportTok{bouncycastle}\OperatorTok{.}\ImportTok{asn1}\OperatorTok{.}\ImportTok{cms}\OperatorTok{.}\ImportTok{Attribute}\OperatorTok{;}
\KeywordTok{import} \ImportTok{org}\OperatorTok{.}\ImportTok{bouncycastle}\OperatorTok{.}\ImportTok{asn1}\OperatorTok{.}\ImportTok{cms}\OperatorTok{.}\ImportTok{AttributeTable}\OperatorTok{;}
\KeywordTok{import} \ImportTok{org}\OperatorTok{.}\ImportTok{bouncycastle}\OperatorTok{.}\ImportTok{asn1}\OperatorTok{.}\ImportTok{pkcs}\OperatorTok{.}\ImportTok{PKCSObjectIdentifiers}\OperatorTok{;}
\KeywordTok{import} \ImportTok{org}\OperatorTok{.}\ImportTok{bouncycastle}\OperatorTok{.}\ImportTok{cms}\OperatorTok{.}\ImportTok{CMSSignedData}\OperatorTok{;}
\KeywordTok{import} \ImportTok{org}\OperatorTok{.}\ImportTok{bouncycastle}\OperatorTok{.}\ImportTok{cms}\OperatorTok{.}\ImportTok{SignerInformation}\OperatorTok{;}
\KeywordTok{import} \ImportTok{org}\OperatorTok{.}\ImportTok{bouncycastle}\OperatorTok{.}\ImportTok{cms}\OperatorTok{.}\ImportTok{SignerInformationStore}\OperatorTok{;}

\KeywordTok{import} \ImportTok{java}\OperatorTok{.}\ImportTok{util}\OperatorTok{.}\ImportTok{Collections}\OperatorTok{;}

\CommentTok{/**}
 \CommentTok{*}\NormalTok{ Embeds an RFC }\CommentTok{3161}\NormalTok{ timestamp token as an unsigned attribute}
 \CommentTok{*}\NormalTok{ in a CMS SignedData structure}\CommentTok{. }\NormalTok{This upgrades a PAdES}\CommentTok{{-}}\NormalTok{B}\CommentTok{{-}}\NormalTok{B}
 \CommentTok{*}\NormalTok{ signature to PAdES}\CommentTok{{-}}\NormalTok{B}\CommentTok{{-}}\NormalTok{T}\CommentTok{.}
 \CommentTok{*/}
\KeywordTok{public}\NormalTok{ CMSSignedData }\FunctionTok{embedTimestamp}\OperatorTok{(}\NormalTok{CMSSignedData signedData}\OperatorTok{,}
\NormalTok{                                     TsaClient tsaClient}\OperatorTok{,}
                                     \BuiltInTok{String}\NormalTok{ tsaUrl}\OperatorTok{)}
        \KeywordTok{throws} \BuiltInTok{Exception} \OperatorTok{\{}

    \CommentTok{// Get the first (and typically only) signer}
\NormalTok{    SignerInformation signerInfo }\OperatorTok{=}\NormalTok{ signedData}\OperatorTok{.}\FunctionTok{getSignerInfos}\OperatorTok{()}
            \OperatorTok{.}\FunctionTok{getSigners}\OperatorTok{().}\FunctionTok{iterator}\OperatorTok{().}\FunctionTok{next}\OperatorTok{();}

    \CommentTok{// The signature value is what gets timestamped}
    \DataTypeTok{byte}\OperatorTok{[]}\NormalTok{ signatureBytes }\OperatorTok{=}\NormalTok{ signerInfo}\OperatorTok{.}\FunctionTok{getSignature}\OperatorTok{();}

    \CommentTok{// Request timestamp from TSA (see Chapter 18 for TSA details)}
    \DataTypeTok{byte}\OperatorTok{[]}\NormalTok{ tsTokenBytes }\OperatorTok{=}\NormalTok{ tsaClient}\OperatorTok{.}\FunctionTok{getTimestampToken}\OperatorTok{(}
\NormalTok{            tsaUrl}\OperatorTok{,}\NormalTok{ signatureBytes}\OperatorTok{);}

    \CommentTok{// Parse the timestamp token as ASN.1}
\NormalTok{    ASN1Primitive tsToken }\OperatorTok{=}\NormalTok{ ASN1Primitive}\OperatorTok{.}\FunctionTok{fromByteArray}\OperatorTok{(}\NormalTok{tsTokenBytes}\OperatorTok{);}

    \CommentTok{// Build the unsigned attribute}
\NormalTok{    ASN1EncodableVector unsignedAttrsVec }\OperatorTok{=} \KeywordTok{new} \FunctionTok{ASN1EncodableVector}\OperatorTok{();}
\NormalTok{    unsignedAttrsVec}\OperatorTok{.}\FunctionTok{add}\OperatorTok{(}\KeywordTok{new} \BuiltInTok{Attribute}\OperatorTok{(}
\NormalTok{            PKCSObjectIdentifiers}\OperatorTok{.}\FunctionTok{id\_aa\_signatureTimeStampToken}\OperatorTok{,}
            \KeywordTok{new} \FunctionTok{DERSet}\OperatorTok{(}\NormalTok{tsToken}\OperatorTok{)));}
\NormalTok{    AttributeTable unsignedAttrs }\OperatorTok{=} \KeywordTok{new} \FunctionTok{AttributeTable}\OperatorTok{(}\NormalTok{unsignedAttrsVec}\OperatorTok{);}

    \CommentTok{// Replace the signer\textquotesingle{}s unsigned attributes}
\NormalTok{    SignerInformation newSignerInfo }\OperatorTok{=}
\NormalTok{            SignerInformation}\OperatorTok{.}\FunctionTok{replaceUnsignedAttributes}\OperatorTok{(}
\NormalTok{                    signerInfo}\OperatorTok{,}\NormalTok{ unsignedAttrs}\OperatorTok{);}

\NormalTok{    SignerInformationStore newSignerInfos }\OperatorTok{=}
            \KeywordTok{new} \FunctionTok{SignerInformationStore}\OperatorTok{(}
                    \BuiltInTok{Collections}\OperatorTok{.}\FunctionTok{singletonList}\OperatorTok{(}\NormalTok{newSignerInfo}\OperatorTok{));}

    \ControlFlowTok{return}\NormalTok{ CMSSignedData}\OperatorTok{.}\FunctionTok{replaceSigners}\OperatorTok{(}\NormalTok{signedData}\OperatorTok{,}\NormalTok{ newSignerInfos}\OperatorTok{);}
\OperatorTok{\}}
\end{Highlighting}
\end{Shaded}

The timestamp is embedded \emph{after} the CMS signature is generated.
This is important: the timestamp covers the signature value, not the
document content. The document content is covered by the CMS signature
itself. The timestamp proves when that signature was created. This
two-layer structure is fundamental to PAdES-B-T.

\section{Complete Production-Ready Signing
Class}\label{_complete_production_ready_signing_class}

Here is a complete signing class that handles all the scenarios
discussed in this chapter: invisible and visible signatures, PAdES-B-B
and PAdES-B-T profiles, certification signatures, and multi-signature
workflows.

\begin{Shaded}
\begin{Highlighting}[]
\KeywordTok{import} \ImportTok{java}\OperatorTok{.}\ImportTok{io}\OperatorTok{.}\ImportTok{ByteArrayOutputStream}\OperatorTok{;}
\KeywordTok{import} \ImportTok{java}\OperatorTok{.}\ImportTok{io}\OperatorTok{.}\ImportTok{IOException}\OperatorTok{;}
\KeywordTok{import} \ImportTok{java}\OperatorTok{.}\ImportTok{io}\OperatorTok{.}\ImportTok{InputStream}\OperatorTok{;}
\KeywordTok{import} \ImportTok{java}\OperatorTok{.}\ImportTok{io}\OperatorTok{.}\ImportTok{OutputStream}\OperatorTok{;}
\KeywordTok{import} \ImportTok{java}\OperatorTok{.}\ImportTok{nio}\OperatorTok{.}\ImportTok{charset}\OperatorTok{.}\ImportTok{StandardCharsets}\OperatorTok{;}
\KeywordTok{import} \ImportTok{java}\OperatorTok{.}\ImportTok{security}\OperatorTok{.}\ImportTok{MessageDigest}\OperatorTok{;}
\KeywordTok{import} \ImportTok{java}\OperatorTok{.}\ImportTok{security}\OperatorTok{.}\ImportTok{PrivateKey}\OperatorTok{;}
\KeywordTok{import} \ImportTok{java}\OperatorTok{.}\ImportTok{security}\OperatorTok{.}\ImportTok{cert}\OperatorTok{.}\ImportTok{X509Certificate}\OperatorTok{;}
\KeywordTok{import} \ImportTok{java}\OperatorTok{.}\ImportTok{util}\OperatorTok{.}\ImportTok{ArrayList}\OperatorTok{;}
\KeywordTok{import} \ImportTok{java}\OperatorTok{.}\ImportTok{util}\OperatorTok{.}\ImportTok{Calendar}\OperatorTok{;}
\KeywordTok{import} \ImportTok{java}\OperatorTok{.}\ImportTok{util}\OperatorTok{.}\ImportTok{Collections}\OperatorTok{;}
\KeywordTok{import} \ImportTok{java}\OperatorTok{.}\ImportTok{util}\OperatorTok{.}\ImportTok{List}\OperatorTok{;}

\KeywordTok{import} \ImportTok{org}\OperatorTok{.}\ImportTok{apache}\OperatorTok{.}\ImportTok{pdfbox}\OperatorTok{.}\ImportTok{Loader}\OperatorTok{;}
\KeywordTok{import} \ImportTok{org}\OperatorTok{.}\ImportTok{apache}\OperatorTok{.}\ImportTok{pdfbox}\OperatorTok{.}\ImportTok{cos}\OperatorTok{.}\ImportTok{COSArray}\OperatorTok{;}
\KeywordTok{import} \ImportTok{org}\OperatorTok{.}\ImportTok{apache}\OperatorTok{.}\ImportTok{pdfbox}\OperatorTok{.}\ImportTok{cos}\OperatorTok{.}\ImportTok{COSDictionary}\OperatorTok{;}
\KeywordTok{import} \ImportTok{org}\OperatorTok{.}\ImportTok{apache}\OperatorTok{.}\ImportTok{pdfbox}\OperatorTok{.}\ImportTok{cos}\OperatorTok{.}\ImportTok{COSInteger}\OperatorTok{;}
\KeywordTok{import} \ImportTok{org}\OperatorTok{.}\ImportTok{apache}\OperatorTok{.}\ImportTok{pdfbox}\OperatorTok{.}\ImportTok{cos}\OperatorTok{.}\ImportTok{COSName}\OperatorTok{;}
\KeywordTok{import} \ImportTok{org}\OperatorTok{.}\ImportTok{apache}\OperatorTok{.}\ImportTok{pdfbox}\OperatorTok{.}\ImportTok{pdmodel}\OperatorTok{.}\ImportTok{PDDocument}\OperatorTok{;}
\KeywordTok{import} \ImportTok{org}\OperatorTok{.}\ImportTok{apache}\OperatorTok{.}\ImportTok{pdfbox}\OperatorTok{.}\ImportTok{pdmodel}\OperatorTok{.}\ImportTok{PDResources}\OperatorTok{;}
\KeywordTok{import} \ImportTok{org}\OperatorTok{.}\ImportTok{apache}\OperatorTok{.}\ImportTok{pdfbox}\OperatorTok{.}\ImportTok{pdmodel}\OperatorTok{.}\ImportTok{common}\OperatorTok{.}\ImportTok{PDRectangle}\OperatorTok{;}
\KeywordTok{import} \ImportTok{org}\OperatorTok{.}\ImportTok{apache}\OperatorTok{.}\ImportTok{pdfbox}\OperatorTok{.}\ImportTok{pdmodel}\OperatorTok{.}\ImportTok{font}\OperatorTok{.}\ImportTok{PDType1Font}\OperatorTok{;}
\KeywordTok{import} \ImportTok{org}\OperatorTok{.}\ImportTok{apache}\OperatorTok{.}\ImportTok{pdfbox}\OperatorTok{.}\ImportTok{pdmodel}\OperatorTok{.}\ImportTok{font}\OperatorTok{.}\ImportTok{Standard14Fonts}\OperatorTok{;}
\KeywordTok{import} \ImportTok{org}\OperatorTok{.}\ImportTok{apache}\OperatorTok{.}\ImportTok{pdfbox}\OperatorTok{.}\ImportTok{pdmodel}\OperatorTok{.}\ImportTok{graphics}\OperatorTok{.}\ImportTok{form}\OperatorTok{.}\ImportTok{PDFormXObject}\OperatorTok{;}
\KeywordTok{import} \ImportTok{org}\OperatorTok{.}\ImportTok{apache}\OperatorTok{.}\ImportTok{pdfbox}\OperatorTok{.}\ImportTok{pdmodel}\OperatorTok{.}\ImportTok{interactive}\OperatorTok{.}\ImportTok{digitalsignature}\OperatorTok{.}\ImportTok{PDSignature}\OperatorTok{;}
\KeywordTok{import} \ImportTok{org}\OperatorTok{.}\ImportTok{apache}\OperatorTok{.}\ImportTok{pdfbox}\OperatorTok{.}\ImportTok{pdmodel}\OperatorTok{.}\ImportTok{interactive}\OperatorTok{.}\ImportTok{digitalsignature}\OperatorTok{.}\ImportTok{SignatureInterface}\OperatorTok{;}
\KeywordTok{import} \ImportTok{org}\OperatorTok{.}\ImportTok{apache}\OperatorTok{.}\ImportTok{pdfbox}\OperatorTok{.}\ImportTok{pdmodel}\OperatorTok{.}\ImportTok{interactive}\OperatorTok{.}\ImportTok{digitalsignature}\OperatorTok{.}\ImportTok{SignatureOptions}\OperatorTok{;}
\KeywordTok{import} \ImportTok{org}\OperatorTok{.}\ImportTok{apache}\OperatorTok{.}\ImportTok{pdfbox}\OperatorTok{.}\ImportTok{pdmodel}\OperatorTok{.}\ImportTok{interactive}\OperatorTok{.}\ImportTok{form}\OperatorTok{.}\ImportTok{PDAcroForm}\OperatorTok{;}
\KeywordTok{import} \ImportTok{org}\OperatorTok{.}\ImportTok{apache}\OperatorTok{.}\ImportTok{pdfbox}\OperatorTok{.}\ImportTok{pdmodel}\OperatorTok{.}\ImportTok{interactive}\OperatorTok{.}\ImportTok{form}\OperatorTok{.}\ImportTok{PDField}\OperatorTok{;}
\KeywordTok{import} \ImportTok{org}\OperatorTok{.}\ImportTok{apache}\OperatorTok{.}\ImportTok{pdfbox}\OperatorTok{.}\ImportTok{pdmodel}\OperatorTok{.}\ImportTok{interactive}\OperatorTok{.}\ImportTok{form}\OperatorTok{.}\ImportTok{PDSignatureField}\OperatorTok{;}
\KeywordTok{import} \ImportTok{org}\OperatorTok{.}\ImportTok{bouncycastle}\OperatorTok{.}\ImportTok{asn1}\OperatorTok{.}\ImportTok{ASN1EncodableVector}\OperatorTok{;}
\KeywordTok{import} \ImportTok{org}\OperatorTok{.}\ImportTok{bouncycastle}\OperatorTok{.}\ImportTok{asn1}\OperatorTok{.}\ImportTok{ASN1Primitive}\OperatorTok{;}
\KeywordTok{import} \ImportTok{org}\OperatorTok{.}\ImportTok{bouncycastle}\OperatorTok{.}\ImportTok{asn1}\OperatorTok{.}\ImportTok{DERSet}\OperatorTok{;}
\KeywordTok{import} \ImportTok{org}\OperatorTok{.}\ImportTok{bouncycastle}\OperatorTok{.}\ImportTok{asn1}\OperatorTok{.}\ImportTok{cms}\OperatorTok{.}\ImportTok{Attribute}\OperatorTok{;}
\KeywordTok{import} \ImportTok{org}\OperatorTok{.}\ImportTok{bouncycastle}\OperatorTok{.}\ImportTok{asn1}\OperatorTok{.}\ImportTok{cms}\OperatorTok{.}\ImportTok{AttributeTable}\OperatorTok{;}
\KeywordTok{import} \ImportTok{org}\OperatorTok{.}\ImportTok{bouncycastle}\OperatorTok{.}\ImportTok{asn1}\OperatorTok{.}\ImportTok{ess}\OperatorTok{.}\ImportTok{ESSCertIDv2}\OperatorTok{;}
\KeywordTok{import} \ImportTok{org}\OperatorTok{.}\ImportTok{bouncycastle}\OperatorTok{.}\ImportTok{asn1}\OperatorTok{.}\ImportTok{ess}\OperatorTok{.}\ImportTok{SigningCertificateV2}\OperatorTok{;}
\KeywordTok{import} \ImportTok{org}\OperatorTok{.}\ImportTok{bouncycastle}\OperatorTok{.}\ImportTok{asn1}\OperatorTok{.}\ImportTok{nist}\OperatorTok{.}\ImportTok{NISTObjectIdentifiers}\OperatorTok{;}
\KeywordTok{import} \ImportTok{org}\OperatorTok{.}\ImportTok{bouncycastle}\OperatorTok{.}\ImportTok{asn1}\OperatorTok{.}\ImportTok{pkcs}\OperatorTok{.}\ImportTok{PKCSObjectIdentifiers}\OperatorTok{;}
\KeywordTok{import} \ImportTok{org}\OperatorTok{.}\ImportTok{bouncycastle}\OperatorTok{.}\ImportTok{asn1}\OperatorTok{.}\ImportTok{x509}\OperatorTok{.}\ImportTok{AlgorithmIdentifier}\OperatorTok{;}
\KeywordTok{import} \ImportTok{org}\OperatorTok{.}\ImportTok{bouncycastle}\OperatorTok{.}\ImportTok{cert}\OperatorTok{.}\ImportTok{jcajce}\OperatorTok{.}\ImportTok{JcaCertStore}\OperatorTok{;}
\KeywordTok{import} \ImportTok{org}\OperatorTok{.}\ImportTok{bouncycastle}\OperatorTok{.}\ImportTok{cms}\OperatorTok{.}\ImportTok{CMSProcessableByteArray}\OperatorTok{;}
\KeywordTok{import} \ImportTok{org}\OperatorTok{.}\ImportTok{bouncycastle}\OperatorTok{.}\ImportTok{cms}\OperatorTok{.}\ImportTok{CMSSignedData}\OperatorTok{;}
\KeywordTok{import} \ImportTok{org}\OperatorTok{.}\ImportTok{bouncycastle}\OperatorTok{.}\ImportTok{cms}\OperatorTok{.}\ImportTok{CMSSignedDataGenerator}\OperatorTok{;}
\KeywordTok{import} \ImportTok{org}\OperatorTok{.}\ImportTok{bouncycastle}\OperatorTok{.}\ImportTok{cms}\OperatorTok{.}\ImportTok{DefaultSignedAttributeTableGenerator}\OperatorTok{;}
\KeywordTok{import} \ImportTok{org}\OperatorTok{.}\ImportTok{bouncycastle}\OperatorTok{.}\ImportTok{cms}\OperatorTok{.}\ImportTok{SignerInfoGenerator}\OperatorTok{;}
\KeywordTok{import} \ImportTok{org}\OperatorTok{.}\ImportTok{bouncycastle}\OperatorTok{.}\ImportTok{cms}\OperatorTok{.}\ImportTok{SignerInformation}\OperatorTok{;}
\KeywordTok{import} \ImportTok{org}\OperatorTok{.}\ImportTok{bouncycastle}\OperatorTok{.}\ImportTok{cms}\OperatorTok{.}\ImportTok{SignerInformationStore}\OperatorTok{;}
\KeywordTok{import} \ImportTok{org}\OperatorTok{.}\ImportTok{bouncycastle}\OperatorTok{.}\ImportTok{cms}\OperatorTok{.}\ImportTok{jcajce}\OperatorTok{.}\ImportTok{JcaSignerInfoGeneratorBuilder}\OperatorTok{;}
\KeywordTok{import} \ImportTok{org}\OperatorTok{.}\ImportTok{bouncycastle}\OperatorTok{.}\ImportTok{operator}\OperatorTok{.}\ImportTok{ContentSigner}\OperatorTok{;}
\KeywordTok{import} \ImportTok{org}\OperatorTok{.}\ImportTok{bouncycastle}\OperatorTok{.}\ImportTok{operator}\OperatorTok{.}\ImportTok{DigestCalculatorProvider}\OperatorTok{;}
\KeywordTok{import} \ImportTok{org}\OperatorTok{.}\ImportTok{bouncycastle}\OperatorTok{.}\ImportTok{operator}\OperatorTok{.}\ImportTok{jcajce}\OperatorTok{.}\ImportTok{JcaContentSignerBuilder}\OperatorTok{;}
\KeywordTok{import} \ImportTok{org}\OperatorTok{.}\ImportTok{bouncycastle}\OperatorTok{.}\ImportTok{operator}\OperatorTok{.}\ImportTok{jcajce}\OperatorTok{.}\ImportTok{JcaDigestCalculatorProviderBuilder}\OperatorTok{;}

\CommentTok{/**}
 \CommentTok{*}\NormalTok{ Production}\CommentTok{{-}}\NormalTok{ready PDF signing service using PDFBox }\CommentTok{\{}\NormalTok{pdfbox}\CommentTok{{-}}\NormalTok{version}\CommentTok{\}}
 \CommentTok{*}\NormalTok{ and BouncyCastle }\CommentTok{\{}\NormalTok{bc}\CommentTok{{-}}\NormalTok{version}\CommentTok{\}.}
 \CommentTok{*}
 \CommentTok{*}\NormalTok{ Supports}\CommentTok{:}
 \CommentTok{*} \CommentTok{{-}}\NormalTok{ PAdES}\CommentTok{{-}}\NormalTok{B}\CommentTok{{-}}\NormalTok{B }\CommentTok{(}\NormalTok{basic}\CommentTok{)}\NormalTok{ and PAdES}\CommentTok{{-}}\NormalTok{B}\CommentTok{{-}}\NormalTok{T }\CommentTok{(}\NormalTok{timestamped}\CommentTok{)}\NormalTok{ profiles}
 \CommentTok{*} \CommentTok{{-}}\NormalTok{ Visible and invisible signatures}
 \CommentTok{*} \CommentTok{{-}}\NormalTok{ Approval and certification }\CommentTok{(}\NormalTok{DocMDP}\CommentTok{)}\NormalTok{ signatures}
 \CommentTok{*} \CommentTok{{-}}\NormalTok{ Multi}\CommentTok{{-}}\NormalTok{signature workflows}
 \CommentTok{*/}
\KeywordTok{public} \KeywordTok{class}\NormalTok{ ProductionPdfSigner }\OperatorTok{\{}

    \CommentTok{/**}\NormalTok{ Container size for signatures without timestamps}\CommentTok{. */}
    \KeywordTok{private} \DataTypeTok{static} \DataTypeTok{final} \DataTypeTok{int}\NormalTok{ CONTAINER\_SIZE\_BASIC }\OperatorTok{=} \DecValTok{32768}\OperatorTok{;}

    \CommentTok{/**}\NormalTok{ Container size for signatures with TSA timestamps}\CommentTok{. */}
    \KeywordTok{private} \DataTypeTok{static} \DataTypeTok{final} \DataTypeTok{int}\NormalTok{ CONTAINER\_SIZE\_WITH\_TSA }\OperatorTok{=} \DecValTok{65536}\OperatorTok{;}

    \CommentTok{/**}
     \CommentTok{*}\NormalTok{ Signs a PDF document}\CommentTok{.}
     \CommentTok{*}
     \CommentTok{*} \AnnotationTok{@param}\CommentTok{ }\AnnotationTok{pdfBytes}\CommentTok{    }\NormalTok{the original PDF bytes}
     \CommentTok{*} \AnnotationTok{@param}\CommentTok{ }\AnnotationTok{privateKey}\CommentTok{  }\NormalTok{the signer}\CommentTok{\textquotesingle{}}\NormalTok{s private key}
     \CommentTok{*} \AnnotationTok{@param}\CommentTok{ }\AnnotationTok{certificate}\CommentTok{ }\NormalTok{the signer}\CommentTok{\textquotesingle{}}\NormalTok{s X}\CommentTok{.509}\NormalTok{ certificate}
     \CommentTok{*} \AnnotationTok{@param}\CommentTok{ }\AnnotationTok{chain}\CommentTok{       }\NormalTok{the full certificate chain }\CommentTok{(}\NormalTok{signer }\CommentTok{+}\NormalTok{ intermediates}\CommentTok{)}
     \CommentTok{*} \AnnotationTok{@param}\CommentTok{ }\AnnotationTok{config}\CommentTok{      }\NormalTok{signing configuration}
     \CommentTok{*} \AnnotationTok{@return}\NormalTok{ the signed PDF bytes }\CommentTok{(}\NormalTok{incremental save}\CommentTok{)}
     \CommentTok{*/}
    \KeywordTok{public} \DataTypeTok{byte}\OperatorTok{[]} \FunctionTok{sign}\OperatorTok{(}\DataTypeTok{byte}\OperatorTok{[]}\NormalTok{ pdfBytes}\OperatorTok{,}
                       \BuiltInTok{PrivateKey}\NormalTok{ privateKey}\OperatorTok{,}
                       \BuiltInTok{X509Certificate}\NormalTok{ certificate}\OperatorTok{,}
                       \BuiltInTok{List}\OperatorTok{\textless{}}\BuiltInTok{X509Certificate}\OperatorTok{\textgreater{}}\NormalTok{ chain}\OperatorTok{,}
\NormalTok{                       SigningConfig config}\OperatorTok{)} \KeywordTok{throws} \BuiltInTok{Exception} \OperatorTok{\{}

        \ControlFlowTok{try} \OperatorTok{(}\NormalTok{PDDocument doc }\OperatorTok{=}\NormalTok{ Loader}\OperatorTok{.}\FunctionTok{loadPDF}\OperatorTok{(}\NormalTok{pdfBytes}\OperatorTok{))} \OperatorTok{\{}

            \CommentTok{// Resolve the target signature field}
\NormalTok{            PDAcroForm acroForm }\OperatorTok{=}\NormalTok{ doc}\OperatorTok{.}\FunctionTok{getDocumentCatalog}\OperatorTok{().}\FunctionTok{getAcroForm}\OperatorTok{();}
            \ControlFlowTok{if} \OperatorTok{(}\NormalTok{acroForm }\OperatorTok{==} \KeywordTok{null}\OperatorTok{)} \OperatorTok{\{}
                \ControlFlowTok{throw} \KeywordTok{new} \BuiltInTok{IllegalStateException}\OperatorTok{(}\StringTok{"PDF has no AcroForm"}\OperatorTok{);}
            \OperatorTok{\}}

\NormalTok{            PDSignatureField sigField }\OperatorTok{=} \FunctionTok{findEmptySignatureField}\OperatorTok{(}
\NormalTok{                    acroForm}\OperatorTok{,}\NormalTok{ config}\OperatorTok{.}\FunctionTok{fieldName}\OperatorTok{());}
            \ControlFlowTok{if} \OperatorTok{(}\NormalTok{sigField }\OperatorTok{==} \KeywordTok{null}\OperatorTok{)} \OperatorTok{\{}
                \ControlFlowTok{throw} \KeywordTok{new} \BuiltInTok{IllegalStateException}\OperatorTok{(}
                        \StringTok{"No empty signature field: "} \OperatorTok{+}\NormalTok{ config}\OperatorTok{.}\FunctionTok{fieldName}\OperatorTok{());}
            \OperatorTok{\}}

            \CommentTok{// Build the /Sig dictionary}
\NormalTok{            PDSignature signature }\OperatorTok{=} \KeywordTok{new} \FunctionTok{PDSignature}\OperatorTok{();}
\NormalTok{            signature}\OperatorTok{.}\FunctionTok{setFilter}\OperatorTok{(}\NormalTok{PDSignature}\OperatorTok{.}\FunctionTok{FILTER\_ADOBE\_PPKLITE}\OperatorTok{);}
\NormalTok{            signature}\OperatorTok{.}\FunctionTok{setSubFilter}\OperatorTok{(}
\NormalTok{                    COSName}\OperatorTok{.}\FunctionTok{getPDFName}\OperatorTok{(}\StringTok{"ETSI.CAdES.detached"}\OperatorTok{));}
\NormalTok{            signature}\OperatorTok{.}\FunctionTok{setSignDate}\OperatorTok{(}\BuiltInTok{Calendar}\OperatorTok{.}\FunctionTok{getInstance}\OperatorTok{());}
\NormalTok{            signature}\OperatorTok{.}\FunctionTok{setName}\OperatorTok{(}\FunctionTok{extractCN}\OperatorTok{(}\NormalTok{certificate}\OperatorTok{));}

            \ControlFlowTok{if} \OperatorTok{(}\NormalTok{config}\OperatorTok{.}\FunctionTok{reason}\OperatorTok{()} \OperatorTok{!=} \KeywordTok{null}\OperatorTok{)} \OperatorTok{\{}
\NormalTok{                signature}\OperatorTok{.}\FunctionTok{setReason}\OperatorTok{(}\NormalTok{config}\OperatorTok{.}\FunctionTok{reason}\OperatorTok{());}
            \OperatorTok{\}}
            \ControlFlowTok{if} \OperatorTok{(}\NormalTok{config}\OperatorTok{.}\FunctionTok{location}\OperatorTok{()} \OperatorTok{!=} \KeywordTok{null}\OperatorTok{)} \OperatorTok{\{}
\NormalTok{                signature}\OperatorTok{.}\FunctionTok{setLocation}\OperatorTok{(}\NormalTok{config}\OperatorTok{.}\FunctionTok{location}\OperatorTok{());}
            \OperatorTok{\}}

            \CommentTok{// Certification signature (DocMDP)}
            \ControlFlowTok{if} \OperatorTok{(}\NormalTok{config}\OperatorTok{.}\FunctionTok{certification}\OperatorTok{())} \OperatorTok{\{}
                \FunctionTok{applyCertification}\OperatorTok{(}\NormalTok{doc}\OperatorTok{,}\NormalTok{ signature}\OperatorTok{,}\NormalTok{ config}\OperatorTok{.}\FunctionTok{docMdpLevel}\OperatorTok{());}
            \OperatorTok{\}}

\NormalTok{            sigField}\OperatorTok{.}\FunctionTok{setValue}\OperatorTok{(}\NormalTok{signature}\OperatorTok{);}

            \CommentTok{// Determine container size}
            \DataTypeTok{boolean}\NormalTok{ useTsa }\OperatorTok{=}\NormalTok{ config}\OperatorTok{.}\FunctionTok{tsaUrl}\OperatorTok{()} \OperatorTok{!=} \KeywordTok{null}
                    \OperatorTok{\&\&} \OperatorTok{!}\NormalTok{config}\OperatorTok{.}\FunctionTok{tsaUrl}\OperatorTok{().}\FunctionTok{isBlank}\OperatorTok{();}
            \DataTypeTok{int}\NormalTok{ containerSize }\OperatorTok{=}\NormalTok{ useTsa}
                    \OperatorTok{?}\NormalTok{ CONTAINER\_SIZE\_WITH\_TSA }\OperatorTok{:}\NormalTok{ CONTAINER\_SIZE\_BASIC}\OperatorTok{;}

\NormalTok{            SignatureOptions options }\OperatorTok{=} \KeywordTok{new} \FunctionTok{SignatureOptions}\OperatorTok{();}
\NormalTok{            options}\OperatorTok{.}\FunctionTok{setPreferredSignatureSize}\OperatorTok{(}\NormalTok{containerSize}\OperatorTok{);}

            \CommentTok{// Visible appearance}
            \ControlFlowTok{if} \OperatorTok{(}\NormalTok{config}\OperatorTok{.}\FunctionTok{visible}\OperatorTok{())} \OperatorTok{\{}
                \FunctionTok{buildTextAppearance}\OperatorTok{(}\NormalTok{doc}\OperatorTok{,}\NormalTok{ sigField}\OperatorTok{,}
                        \FunctionTok{extractCN}\OperatorTok{(}\NormalTok{certificate}\OperatorTok{),}\NormalTok{ config}\OperatorTok{.}\FunctionTok{reason}\OperatorTok{());}
            \OperatorTok{\}}

            \CommentTok{// CMS signing callback}
\NormalTok{            PadesCmsCallback callback }\OperatorTok{=} \KeywordTok{new} \FunctionTok{PadesCmsCallback}\OperatorTok{(}
\NormalTok{                    privateKey}\OperatorTok{,}\NormalTok{ certificate}\OperatorTok{,}\NormalTok{ chain}\OperatorTok{,}
\NormalTok{                    useTsa }\OperatorTok{?}\NormalTok{ config}\OperatorTok{.}\FunctionTok{tsaUrl}\OperatorTok{()} \OperatorTok{:} \KeywordTok{null}\OperatorTok{);}

\NormalTok{            doc}\OperatorTok{.}\FunctionTok{addSignature}\OperatorTok{(}\NormalTok{signature}\OperatorTok{,}\NormalTok{ callback}\OperatorTok{,}\NormalTok{ options}\OperatorTok{);}

            \BuiltInTok{ByteArrayOutputStream}\NormalTok{ baos }\OperatorTok{=} \KeywordTok{new} \BuiltInTok{ByteArrayOutputStream}\OperatorTok{();}
\NormalTok{            doc}\OperatorTok{.}\FunctionTok{saveIncremental}\OperatorTok{(}\NormalTok{baos}\OperatorTok{);}
            \ControlFlowTok{return}\NormalTok{ baos}\OperatorTok{.}\FunctionTok{toByteArray}\OperatorTok{();}
        \OperatorTok{\}}
    \OperatorTok{\}}

    \CommentTok{// ── Configuration record ─────────────────────────────────────────}

    \CommentTok{/**}
     \CommentTok{*}\NormalTok{ Immutable configuration for a signing operation}\CommentTok{.}
     \CommentTok{*/}
    \KeywordTok{public} \KeywordTok{record} \FunctionTok{SigningConfig}\OperatorTok{(}
            \BuiltInTok{String}\NormalTok{ fieldName}\OperatorTok{,}
            \BuiltInTok{String}\NormalTok{ reason}\OperatorTok{,}
            \BuiltInTok{String}\NormalTok{ location}\OperatorTok{,}
            \DataTypeTok{boolean}\NormalTok{ visible}\OperatorTok{,}
            \DataTypeTok{boolean}\NormalTok{ certification}\OperatorTok{,}
            \DataTypeTok{int}\NormalTok{ docMdpLevel}\OperatorTok{,}
            \BuiltInTok{String}\NormalTok{ tsaUrl}\OperatorTok{)} \OperatorTok{\{}

        \KeywordTok{public}\NormalTok{ SigningConfig }\OperatorTok{\{}
            \ControlFlowTok{if} \OperatorTok{(}\NormalTok{fieldName }\OperatorTok{==} \KeywordTok{null} \OperatorTok{||}\NormalTok{ fieldName}\OperatorTok{.}\FunctionTok{isBlank}\OperatorTok{())} \OperatorTok{\{}
                \ControlFlowTok{throw} \KeywordTok{new} \BuiltInTok{IllegalArgumentException}\OperatorTok{(}
                        \StringTok{"fieldName is required"}\OperatorTok{);}
            \OperatorTok{\}}
        \OperatorTok{\}}
    \OperatorTok{\}}

    \CommentTok{// ── CMS callback ─────────────────────────────────────────────────}

    \KeywordTok{private} \DataTypeTok{static} \KeywordTok{class}\NormalTok{ PadesCmsCallback }\KeywordTok{implements}\NormalTok{ SignatureInterface }\OperatorTok{\{}

        \KeywordTok{private} \DataTypeTok{final} \BuiltInTok{PrivateKey}\NormalTok{ privateKey}\OperatorTok{;}
        \KeywordTok{private} \DataTypeTok{final} \BuiltInTok{X509Certificate}\NormalTok{ certificate}\OperatorTok{;}
        \KeywordTok{private} \DataTypeTok{final} \BuiltInTok{List}\OperatorTok{\textless{}}\BuiltInTok{X509Certificate}\OperatorTok{\textgreater{}}\NormalTok{ chain}\OperatorTok{;}
        \KeywordTok{private} \DataTypeTok{final} \BuiltInTok{String}\NormalTok{ tsaUrl}\OperatorTok{;}

        \FunctionTok{PadesCmsCallback}\OperatorTok{(}\BuiltInTok{PrivateKey}\NormalTok{ privateKey}\OperatorTok{,}
                         \BuiltInTok{X509Certificate}\NormalTok{ certificate}\OperatorTok{,}
                         \BuiltInTok{List}\OperatorTok{\textless{}}\BuiltInTok{X509Certificate}\OperatorTok{\textgreater{}}\NormalTok{ chain}\OperatorTok{,}
                         \BuiltInTok{String}\NormalTok{ tsaUrl}\OperatorTok{)} \OperatorTok{\{}
            \KeywordTok{this}\OperatorTok{.}\FunctionTok{privateKey} \OperatorTok{=}\NormalTok{ privateKey}\OperatorTok{;}
            \KeywordTok{this}\OperatorTok{.}\FunctionTok{certificate} \OperatorTok{=}\NormalTok{ certificate}\OperatorTok{;}
            \KeywordTok{this}\OperatorTok{.}\FunctionTok{chain} \OperatorTok{=}\NormalTok{ chain}\OperatorTok{;}
            \KeywordTok{this}\OperatorTok{.}\FunctionTok{tsaUrl} \OperatorTok{=}\NormalTok{ tsaUrl}\OperatorTok{;}
        \OperatorTok{\}}

        \AttributeTok{@Override}
        \KeywordTok{public} \DataTypeTok{byte}\OperatorTok{[]} \FunctionTok{sign}\OperatorTok{(}\BuiltInTok{InputStream}\NormalTok{ content}\OperatorTok{)} \KeywordTok{throws} \BuiltInTok{IOException} \OperatorTok{\{}
            \ControlFlowTok{try} \OperatorTok{\{}
                \DataTypeTok{byte}\OperatorTok{[]}\NormalTok{ contentBytes }\OperatorTok{=}\NormalTok{ content}\OperatorTok{.}\FunctionTok{readAllBytes}\OperatorTok{();}

                \CommentTok{// Determine signature algorithm from key type}
                \BuiltInTok{String}\NormalTok{ keyAlgo }\OperatorTok{=}\NormalTok{ privateKey}\OperatorTok{.}\FunctionTok{getAlgorithm}\OperatorTok{();}
                \BuiltInTok{String}\NormalTok{ sigAlgo }\OperatorTok{=} \StringTok{"EC"}\OperatorTok{.}\FunctionTok{equals}\OperatorTok{(}\NormalTok{keyAlgo}\OperatorTok{)}
                        \OperatorTok{?} \StringTok{"SHA256withECDSA"} \OperatorTok{:} \StringTok{"SHA256withRSA"}\OperatorTok{;}

                \CommentTok{// Build certificate store}
\NormalTok{                JcaCertStore certStore }\OperatorTok{=} \KeywordTok{new} \FunctionTok{JcaCertStore}\OperatorTok{(}\NormalTok{chain}\OperatorTok{);}

                \CommentTok{// Content signer}
\NormalTok{                ContentSigner signer }\OperatorTok{=} \KeywordTok{new} \FunctionTok{JcaContentSignerBuilder}\OperatorTok{(}\NormalTok{sigAlgo}\OperatorTok{)}
                        \OperatorTok{.}\FunctionTok{setProvider}\OperatorTok{(}\StringTok{"BC"}\OperatorTok{)}
                        \OperatorTok{.}\FunctionTok{build}\OperatorTok{(}\NormalTok{privateKey}\OperatorTok{);}

                \CommentTok{// Digest calculator}
\NormalTok{                DigestCalculatorProvider dcProvider }\OperatorTok{=}
                        \KeywordTok{new} \FunctionTok{JcaDigestCalculatorProviderBuilder}\OperatorTok{()}
                                \OperatorTok{.}\FunctionTok{setProvider}\OperatorTok{(}\StringTok{"BC"}\OperatorTok{).}\FunctionTok{build}\OperatorTok{();}

                \CommentTok{// Build signer info with ESS{-}signing{-}certificate{-}v2}
\NormalTok{                JcaSignerInfoGeneratorBuilder signerBuilder }\OperatorTok{=}
                        \KeywordTok{new} \FunctionTok{JcaSignerInfoGeneratorBuilder}\OperatorTok{(}\NormalTok{dcProvider}\OperatorTok{);}

\NormalTok{                AttributeTable signedAttrs }\OperatorTok{=}
                        \FunctionTok{buildEssSigningCertAttribute}\OperatorTok{(}\NormalTok{certificate}\OperatorTok{);}
\NormalTok{                signerBuilder}\OperatorTok{.}\FunctionTok{setSignedAttributeGenerator}\OperatorTok{(}
                        \KeywordTok{new} \FunctionTok{DefaultSignedAttributeTableGenerator}\OperatorTok{(}
\NormalTok{                                signedAttrs}\OperatorTok{));}

\NormalTok{                SignerInfoGenerator signerInfoGen }\OperatorTok{=}
\NormalTok{                        signerBuilder}\OperatorTok{.}\FunctionTok{build}\OperatorTok{(}\NormalTok{signer}\OperatorTok{,}\NormalTok{ certificate}\OperatorTok{);}

                \CommentTok{// Assemble the CMS SignedData}
\NormalTok{                CMSSignedDataGenerator gen }\OperatorTok{=}
                        \KeywordTok{new} \FunctionTok{CMSSignedDataGenerator}\OperatorTok{();}
\NormalTok{                gen}\OperatorTok{.}\FunctionTok{addSignerInfoGenerator}\OperatorTok{(}\NormalTok{signerInfoGen}\OperatorTok{);}
\NormalTok{                gen}\OperatorTok{.}\FunctionTok{addCertificates}\OperatorTok{(}\NormalTok{certStore}\OperatorTok{);}

\NormalTok{                CMSProcessableByteArray cmsContent }\OperatorTok{=}
                        \KeywordTok{new} \FunctionTok{CMSProcessableByteArray}\OperatorTok{(}\NormalTok{contentBytes}\OperatorTok{);}
\NormalTok{                CMSSignedData signedData }\OperatorTok{=}
\NormalTok{                        gen}\OperatorTok{.}\FunctionTok{generate}\OperatorTok{(}\NormalTok{cmsContent}\OperatorTok{,} \KeywordTok{false}\OperatorTok{);}

                \CommentTok{// PAdES{-}B{-}T: add timestamp}
                \ControlFlowTok{if} \OperatorTok{(}\NormalTok{tsaUrl }\OperatorTok{!=} \KeywordTok{null} \OperatorTok{\&\&} \OperatorTok{!}\NormalTok{tsaUrl}\OperatorTok{.}\FunctionTok{isBlank}\OperatorTok{())} \OperatorTok{\{}
\NormalTok{                    signedData }\OperatorTok{=} \FunctionTok{addTimestamp}\OperatorTok{(}\NormalTok{signedData}\OperatorTok{,}\NormalTok{ tsaUrl}\OperatorTok{);}
                \OperatorTok{\}}

                \ControlFlowTok{return}\NormalTok{ signedData}\OperatorTok{.}\FunctionTok{getEncoded}\OperatorTok{();}

            \OperatorTok{\}} \ControlFlowTok{catch} \OperatorTok{(}\BuiltInTok{Exception}\NormalTok{ e}\OperatorTok{)} \OperatorTok{\{}
                \ControlFlowTok{throw} \KeywordTok{new} \BuiltInTok{IOException}\OperatorTok{(}
                        \StringTok{"CMS signing failed: "} \OperatorTok{+}\NormalTok{ e}\OperatorTok{.}\FunctionTok{getMessage}\OperatorTok{(),}\NormalTok{ e}\OperatorTok{);}
            \OperatorTok{\}}
        \OperatorTok{\}}

        \KeywordTok{private}\NormalTok{ AttributeTable }\FunctionTok{buildEssSigningCertAttribute}\OperatorTok{(}
                \BuiltInTok{X509Certificate}\NormalTok{ cert}\OperatorTok{)} \KeywordTok{throws} \BuiltInTok{Exception} \OperatorTok{\{}
            \BuiltInTok{MessageDigest}\NormalTok{ md }\OperatorTok{=} \BuiltInTok{MessageDigest}\OperatorTok{.}\FunctionTok{getInstance}\OperatorTok{(}\StringTok{"SHA{-}256"}\OperatorTok{);}
            \DataTypeTok{byte}\OperatorTok{[]}\NormalTok{ certHash }\OperatorTok{=}\NormalTok{ md}\OperatorTok{.}\FunctionTok{digest}\OperatorTok{(}\NormalTok{cert}\OperatorTok{.}\FunctionTok{getEncoded}\OperatorTok{());}

\NormalTok{            AlgorithmIdentifier hashAlgo }\OperatorTok{=} \KeywordTok{new} \FunctionTok{AlgorithmIdentifier}\OperatorTok{(}
\NormalTok{                    NISTObjectIdentifiers}\OperatorTok{.}\FunctionTok{id\_sha256}\OperatorTok{);}
\NormalTok{            ESSCertIDv2 essCertId }\OperatorTok{=}
                    \KeywordTok{new} \FunctionTok{ESSCertIDv2}\OperatorTok{(}\NormalTok{hashAlgo}\OperatorTok{,}\NormalTok{ certHash}\OperatorTok{);}
\NormalTok{            SigningCertificateV2 sigCertV2 }\OperatorTok{=}
                    \KeywordTok{new} \FunctionTok{SigningCertificateV2}\OperatorTok{(}
                            \KeywordTok{new}\NormalTok{ ESSCertIDv2}\OperatorTok{[]\{}\NormalTok{essCertId}\OperatorTok{\});}

\NormalTok{            ASN1EncodableVector v }\OperatorTok{=} \KeywordTok{new} \FunctionTok{ASN1EncodableVector}\OperatorTok{();}
\NormalTok{            v}\OperatorTok{.}\FunctionTok{add}\OperatorTok{(}\KeywordTok{new} \BuiltInTok{Attribute}\OperatorTok{(}
\NormalTok{                    PKCSObjectIdentifiers}\OperatorTok{.}\FunctionTok{id\_aa\_signingCertificateV2}\OperatorTok{,}
                    \KeywordTok{new} \FunctionTok{DERSet}\OperatorTok{(}\NormalTok{sigCertV2}\OperatorTok{)));}

            \ControlFlowTok{return} \KeywordTok{new} \FunctionTok{AttributeTable}\OperatorTok{(}\NormalTok{v}\OperatorTok{);}
        \OperatorTok{\}}

        \KeywordTok{private}\NormalTok{ CMSSignedData }\FunctionTok{addTimestamp}\OperatorTok{(}\NormalTok{CMSSignedData signedData}\OperatorTok{,}
                                            \BuiltInTok{String}\NormalTok{ tsaUrl}\OperatorTok{)}
                \KeywordTok{throws} \BuiltInTok{Exception} \OperatorTok{\{}
\NormalTok{            SignerInformation signerInfo }\OperatorTok{=}\NormalTok{ signedData}
                    \OperatorTok{.}\FunctionTok{getSignerInfos}\OperatorTok{().}\FunctionTok{getSigners}\OperatorTok{().}\FunctionTok{iterator}\OperatorTok{().}\FunctionTok{next}\OperatorTok{();}
            \DataTypeTok{byte}\OperatorTok{[]}\NormalTok{ sigBytes }\OperatorTok{=}\NormalTok{ signerInfo}\OperatorTok{.}\FunctionTok{getSignature}\OperatorTok{();}

            \CommentTok{// Request timestamp (see TsaClient class below)}
            \DataTypeTok{byte}\OperatorTok{[]}\NormalTok{ tsToken }\OperatorTok{=}\NormalTok{ TsaClient}\OperatorTok{.}\FunctionTok{getTimestamp}\OperatorTok{(}\NormalTok{tsaUrl}\OperatorTok{,}\NormalTok{ sigBytes}\OperatorTok{);}

\NormalTok{            ASN1Primitive tsAsn1 }\OperatorTok{=}
\NormalTok{                    ASN1Primitive}\OperatorTok{.}\FunctionTok{fromByteArray}\OperatorTok{(}\NormalTok{tsToken}\OperatorTok{);}

\NormalTok{            ASN1EncodableVector vec }\OperatorTok{=} \KeywordTok{new} \FunctionTok{ASN1EncodableVector}\OperatorTok{();}
\NormalTok{            vec}\OperatorTok{.}\FunctionTok{add}\OperatorTok{(}\KeywordTok{new} \BuiltInTok{Attribute}\OperatorTok{(}
\NormalTok{                    PKCSObjectIdentifiers}\OperatorTok{.}\FunctionTok{id\_aa\_signatureTimeStampToken}\OperatorTok{,}
                    \KeywordTok{new} \FunctionTok{DERSet}\OperatorTok{(}\NormalTok{tsAsn1}\OperatorTok{)));}
\NormalTok{            AttributeTable unsignedAttrs }\OperatorTok{=} \KeywordTok{new} \FunctionTok{AttributeTable}\OperatorTok{(}\NormalTok{vec}\OperatorTok{);}

\NormalTok{            SignerInformation newSigner }\OperatorTok{=}
\NormalTok{                    SignerInformation}\OperatorTok{.}\FunctionTok{replaceUnsignedAttributes}\OperatorTok{(}
\NormalTok{                            signerInfo}\OperatorTok{,}\NormalTok{ unsignedAttrs}\OperatorTok{);}

            \ControlFlowTok{return}\NormalTok{ CMSSignedData}\OperatorTok{.}\FunctionTok{replaceSigners}\OperatorTok{(}\NormalTok{signedData}\OperatorTok{,}
                    \KeywordTok{new} \FunctionTok{SignerInformationStore}\OperatorTok{(}
                            \BuiltInTok{Collections}\OperatorTok{.}\FunctionTok{singletonList}\OperatorTok{(}\NormalTok{newSigner}\OperatorTok{)));}
        \OperatorTok{\}}
    \OperatorTok{\}}

    \CommentTok{// ── TSA client ───────────────────────────────────────────────────}

    \CommentTok{/**}
     \CommentTok{*}\NormalTok{ Minimal RFC }\CommentTok{3161}\NormalTok{ TSA client}\CommentTok{.}
     \CommentTok{*}\NormalTok{ See Chapter }\CommentTok{18}\NormalTok{ for a full discussion of the timestamp protocol}\CommentTok{.}
     \CommentTok{*/}
    \DataTypeTok{static} \KeywordTok{class}\NormalTok{ TsaClient }\OperatorTok{\{}

        \DataTypeTok{static} \DataTypeTok{byte}\OperatorTok{[]} \FunctionTok{getTimestamp}\OperatorTok{(}\BuiltInTok{String}\NormalTok{ tsaUrl}\OperatorTok{,} \DataTypeTok{byte}\OperatorTok{[]}\NormalTok{ signatureBytes}\OperatorTok{)}
                \KeywordTok{throws} \BuiltInTok{Exception} \OperatorTok{\{}
            \BuiltInTok{MessageDigest}\NormalTok{ md }\OperatorTok{=} \BuiltInTok{MessageDigest}\OperatorTok{.}\FunctionTok{getInstance}\OperatorTok{(}\StringTok{"SHA{-}256"}\OperatorTok{);}
            \DataTypeTok{byte}\OperatorTok{[]}\NormalTok{ hash }\OperatorTok{=}\NormalTok{ md}\OperatorTok{.}\FunctionTok{digest}\OperatorTok{(}\NormalTok{signatureBytes}\OperatorTok{);}

\NormalTok{            org}\OperatorTok{.}\FunctionTok{bouncycastle}\OperatorTok{.}\FunctionTok{tsp}\OperatorTok{.}\FunctionTok{TimeStampRequestGenerator}\NormalTok{ reqGen }\OperatorTok{=}
                    \KeywordTok{new}\NormalTok{ org}\OperatorTok{.}\FunctionTok{bouncycastle}\OperatorTok{.}\FunctionTok{tsp}\OperatorTok{.}\FunctionTok{TimeStampRequestGenerator}\OperatorTok{();}
\NormalTok{            reqGen}\OperatorTok{.}\FunctionTok{setCertReq}\OperatorTok{(}\KeywordTok{true}\OperatorTok{);}
\NormalTok{            org}\OperatorTok{.}\FunctionTok{bouncycastle}\OperatorTok{.}\FunctionTok{tsp}\OperatorTok{.}\FunctionTok{TimeStampRequest}\NormalTok{ tsReq }\OperatorTok{=}
\NormalTok{                    reqGen}\OperatorTok{.}\FunctionTok{generate}\OperatorTok{(}\NormalTok{NISTObjectIdentifiers}\OperatorTok{.}\FunctionTok{id\_sha256}\OperatorTok{,}\NormalTok{ hash}\OperatorTok{);}

\NormalTok{            java}\OperatorTok{.}\FunctionTok{net}\OperatorTok{.}\FunctionTok{http}\OperatorTok{.}\FunctionTok{HttpClient}\NormalTok{ client }\OperatorTok{=}
\NormalTok{                    java}\OperatorTok{.}\FunctionTok{net}\OperatorTok{.}\FunctionTok{http}\OperatorTok{.}\FunctionTok{HttpClient}\OperatorTok{.}\FunctionTok{newHttpClient}\OperatorTok{();}
\NormalTok{            java}\OperatorTok{.}\FunctionTok{net}\OperatorTok{.}\FunctionTok{http}\OperatorTok{.}\FunctionTok{HttpRequest}\NormalTok{ request }\OperatorTok{=}
\NormalTok{                    java}\OperatorTok{.}\FunctionTok{net}\OperatorTok{.}\FunctionTok{http}\OperatorTok{.}\FunctionTok{HttpRequest}\OperatorTok{.}\FunctionTok{newBuilder}\OperatorTok{()}
                            \OperatorTok{.}\FunctionTok{uri}\OperatorTok{(}\NormalTok{java}\OperatorTok{.}\FunctionTok{net}\OperatorTok{.}\FunctionTok{URI}\OperatorTok{.}\FunctionTok{create}\OperatorTok{(}\NormalTok{tsaUrl}\OperatorTok{))}
                            \OperatorTok{.}\FunctionTok{header}\OperatorTok{(}\StringTok{"Content{-}Type"}\OperatorTok{,}
                                    \StringTok{"application/timestamp{-}query"}\OperatorTok{)}
                            \OperatorTok{.}\FunctionTok{POST}\OperatorTok{(}\NormalTok{java}\OperatorTok{.}\FunctionTok{net}\OperatorTok{.}\FunctionTok{http}\OperatorTok{.}\FunctionTok{HttpRequest}\OperatorTok{.}\FunctionTok{BodyPublishers}
                                    \OperatorTok{.}\FunctionTok{ofByteArray}\OperatorTok{(}\NormalTok{tsReq}\OperatorTok{.}\FunctionTok{getEncoded}\OperatorTok{()))}
                            \OperatorTok{.}\FunctionTok{build}\OperatorTok{();}

\NormalTok{            java}\OperatorTok{.}\FunctionTok{net}\OperatorTok{.}\FunctionTok{http}\OperatorTok{.}\FunctionTok{HttpResponse}\OperatorTok{\textless{}}\DataTypeTok{byte}\OperatorTok{[]\textgreater{}}\NormalTok{ response }\OperatorTok{=}\NormalTok{ client}\OperatorTok{.}\FunctionTok{send}\OperatorTok{(}
\NormalTok{                    request}\OperatorTok{,}
\NormalTok{                    java}\OperatorTok{.}\FunctionTok{net}\OperatorTok{.}\FunctionTok{http}\OperatorTok{.}\FunctionTok{HttpResponse}\OperatorTok{.}\FunctionTok{BodyHandlers}\OperatorTok{.}\FunctionTok{ofByteArray}\OperatorTok{());}

            \ControlFlowTok{if} \OperatorTok{(}\NormalTok{response}\OperatorTok{.}\FunctionTok{statusCode}\OperatorTok{()} \OperatorTok{!=} \DecValTok{200}\OperatorTok{)} \OperatorTok{\{}
                \ControlFlowTok{throw} \KeywordTok{new} \BuiltInTok{IOException}\OperatorTok{(}\StringTok{"TSA returned HTTP "}
                        \OperatorTok{+}\NormalTok{ response}\OperatorTok{.}\FunctionTok{statusCode}\OperatorTok{());}
            \OperatorTok{\}}

\NormalTok{            org}\OperatorTok{.}\FunctionTok{bouncycastle}\OperatorTok{.}\FunctionTok{tsp}\OperatorTok{.}\FunctionTok{TimeStampResponse}\NormalTok{ tsResp }\OperatorTok{=}
                    \KeywordTok{new}\NormalTok{ org}\OperatorTok{.}\FunctionTok{bouncycastle}\OperatorTok{.}\FunctionTok{tsp}\OperatorTok{.}\FunctionTok{TimeStampResponse}\OperatorTok{(}
\NormalTok{                            response}\OperatorTok{.}\FunctionTok{body}\OperatorTok{());}
\NormalTok{            tsResp}\OperatorTok{.}\FunctionTok{validate}\OperatorTok{(}\NormalTok{tsReq}\OperatorTok{);}

            \ControlFlowTok{return}\NormalTok{ tsResp}\OperatorTok{.}\FunctionTok{getTimeStampToken}\OperatorTok{().}\FunctionTok{getEncoded}\OperatorTok{();}
        \OperatorTok{\}}
    \OperatorTok{\}}

    \CommentTok{// ── Appearance builder ───────────────────────────────────────────}

    \KeywordTok{private} \DataTypeTok{void} \FunctionTok{buildTextAppearance}\OperatorTok{(}\NormalTok{PDDocument doc}\OperatorTok{,}
\NormalTok{                                      PDSignatureField sigField}\OperatorTok{,}
                                      \BuiltInTok{String}\NormalTok{ name}\OperatorTok{,} \BuiltInTok{String}\NormalTok{ reason}\OperatorTok{)}
            \KeywordTok{throws} \BuiltInTok{Exception} \OperatorTok{\{}

\NormalTok{        PDRectangle rect }\OperatorTok{=}\NormalTok{ sigField}\OperatorTok{.}\FunctionTok{getWidgets}\OperatorTok{().}\FunctionTok{isEmpty}\OperatorTok{()} \OperatorTok{?} \KeywordTok{null}
                \OperatorTok{:}\NormalTok{ sigField}\OperatorTok{.}\FunctionTok{getWidgets}\OperatorTok{().}\FunctionTok{get}\OperatorTok{(}\DecValTok{0}\OperatorTok{).}\FunctionTok{getRectangle}\OperatorTok{();}
        \ControlFlowTok{if} \OperatorTok{(}\NormalTok{rect }\OperatorTok{==} \KeywordTok{null} \OperatorTok{||}\NormalTok{ rect}\OperatorTok{.}\FunctionTok{getWidth}\OperatorTok{()} \OperatorTok{\textless{}} \DecValTok{1}
                \OperatorTok{||}\NormalTok{ rect}\OperatorTok{.}\FunctionTok{getHeight}\OperatorTok{()} \OperatorTok{\textless{}} \DecValTok{1}\OperatorTok{)} \OperatorTok{\{}
            \ControlFlowTok{return}\OperatorTok{;}
        \OperatorTok{\}}

        \DataTypeTok{float}\NormalTok{ w }\OperatorTok{=}\NormalTok{ rect}\OperatorTok{.}\FunctionTok{getWidth}\OperatorTok{();}
        \DataTypeTok{float}\NormalTok{ h }\OperatorTok{=}\NormalTok{ rect}\OperatorTok{.}\FunctionTok{getHeight}\OperatorTok{();}
\NormalTok{        PDResources resources }\OperatorTok{=} \KeywordTok{new} \FunctionTok{PDResources}\OperatorTok{();}
\NormalTok{        PDType1Font font }\OperatorTok{=} \KeywordTok{new} \FunctionTok{PDType1Font}\OperatorTok{(}
\NormalTok{                Standard14Fonts}\OperatorTok{.}\FunctionTok{FontName}\OperatorTok{.}\FunctionTok{HELVETICA}\OperatorTok{);}
\NormalTok{        COSName fontName }\OperatorTok{=}\NormalTok{ resources}\OperatorTok{.}\FunctionTok{add}\OperatorTok{(}\NormalTok{font}\OperatorTok{);}

        \DataTypeTok{float}\NormalTok{ fontSize }\OperatorTok{=} \FloatTok{9f}\OperatorTok{;}
        \DataTypeTok{float}\NormalTok{ y }\OperatorTok{=}\NormalTok{ h }\OperatorTok{{-}}\NormalTok{ fontSize }\OperatorTok{{-}} \DecValTok{3}\OperatorTok{;}
        \BuiltInTok{StringBuilder}\NormalTok{ ops }\OperatorTok{=} \KeywordTok{new} \BuiltInTok{StringBuilder}\OperatorTok{();}
\NormalTok{        ops}\OperatorTok{.}\FunctionTok{append}\OperatorTok{(}\StringTok{"BT}\SpecialCharTok{\textbackslash{}n}\StringTok{"}\OperatorTok{);}
\NormalTok{        ops}\OperatorTok{.}\FunctionTok{append}\OperatorTok{(}\StringTok{"/"}\OperatorTok{).}\FunctionTok{append}\OperatorTok{(}\NormalTok{fontName}\OperatorTok{.}\FunctionTok{getName}\OperatorTok{())}
                \OperatorTok{.}\FunctionTok{append}\OperatorTok{(}\StringTok{" "}\OperatorTok{).}\FunctionTok{append}\OperatorTok{(}\NormalTok{fontSize}\OperatorTok{).}\FunctionTok{append}\OperatorTok{(}\StringTok{" Tf}\SpecialCharTok{\textbackslash{}n}\StringTok{"}\OperatorTok{);}
\NormalTok{        ops}\OperatorTok{.}\FunctionTok{append}\OperatorTok{(}\StringTok{"0 0 0 rg}\SpecialCharTok{\textbackslash{}n}\StringTok{"}\OperatorTok{);}
\NormalTok{        ops}\OperatorTok{.}\FunctionTok{append}\OperatorTok{(}\StringTok{"3 "}\OperatorTok{).}\FunctionTok{append}\OperatorTok{(}\NormalTok{y}\OperatorTok{).}\FunctionTok{append}\OperatorTok{(}\StringTok{" Td}\SpecialCharTok{\textbackslash{}n}\StringTok{"}\OperatorTok{);}
\NormalTok{        ops}\OperatorTok{.}\FunctionTok{append}\OperatorTok{(}\StringTok{"(Digitally signed by: "}\OperatorTok{)}
                \OperatorTok{.}\FunctionTok{append}\OperatorTok{(}\FunctionTok{escapePs}\OperatorTok{(}\NormalTok{name}\OperatorTok{)).}\FunctionTok{append}\OperatorTok{(}\StringTok{") Tj}\SpecialCharTok{\textbackslash{}n}\StringTok{"}\OperatorTok{);}
        \ControlFlowTok{if} \OperatorTok{(}\NormalTok{reason }\OperatorTok{!=} \KeywordTok{null}\OperatorTok{)} \OperatorTok{\{}
\NormalTok{            ops}\OperatorTok{.}\FunctionTok{append}\OperatorTok{(}\StringTok{"0 {-}"}\OperatorTok{).}\FunctionTok{append}\OperatorTok{(}\NormalTok{fontSize }\OperatorTok{*} \FloatTok{1.3f}\OperatorTok{).}\FunctionTok{append}\OperatorTok{(}\StringTok{" Td}\SpecialCharTok{\textbackslash{}n}\StringTok{"}\OperatorTok{);}
\NormalTok{            ops}\OperatorTok{.}\FunctionTok{append}\OperatorTok{(}\StringTok{"(Reason: "}\OperatorTok{).}\FunctionTok{append}\OperatorTok{(}\FunctionTok{escapePs}\OperatorTok{(}\NormalTok{reason}\OperatorTok{))}
                    \OperatorTok{.}\FunctionTok{append}\OperatorTok{(}\StringTok{") Tj}\SpecialCharTok{\textbackslash{}n}\StringTok{"}\OperatorTok{);}
        \OperatorTok{\}}
\NormalTok{        ops}\OperatorTok{.}\FunctionTok{append}\OperatorTok{(}\StringTok{"ET}\SpecialCharTok{\textbackslash{}n}\StringTok{"}\OperatorTok{);}

\NormalTok{        PDFormXObject form }\OperatorTok{=} \KeywordTok{new} \FunctionTok{PDFormXObject}\OperatorTok{(}\NormalTok{doc}\OperatorTok{);}
\NormalTok{        form}\OperatorTok{.}\FunctionTok{setBBox}\OperatorTok{(}\KeywordTok{new} \FunctionTok{PDRectangle}\OperatorTok{(}\NormalTok{w}\OperatorTok{,}\NormalTok{ h}\OperatorTok{));}
\NormalTok{        form}\OperatorTok{.}\FunctionTok{setResources}\OperatorTok{(}\NormalTok{resources}\OperatorTok{);}
        \ControlFlowTok{try} \OperatorTok{(}\BuiltInTok{OutputStream}\NormalTok{ os }\OperatorTok{=}\NormalTok{ form}\OperatorTok{.}\FunctionTok{getStream}\OperatorTok{().}\FunctionTok{createOutputStream}\OperatorTok{())} \OperatorTok{\{}
\NormalTok{            os}\OperatorTok{.}\FunctionTok{write}\OperatorTok{(}\NormalTok{ops}\OperatorTok{.}\FunctionTok{toString}\OperatorTok{()}
                    \OperatorTok{.}\FunctionTok{getBytes}\OperatorTok{(}\NormalTok{StandardCharsets}\OperatorTok{.}\FunctionTok{US\_ASCII}\OperatorTok{));}
        \OperatorTok{\}}

\NormalTok{        COSDictionary apDict }\OperatorTok{=} \KeywordTok{new} \FunctionTok{COSDictionary}\OperatorTok{();}
\NormalTok{        apDict}\OperatorTok{.}\FunctionTok{setItem}\OperatorTok{(}\NormalTok{COSName}\OperatorTok{.}\FunctionTok{N}\OperatorTok{,}\NormalTok{ form}\OperatorTok{);}
\NormalTok{        sigField}\OperatorTok{.}\FunctionTok{getWidgets}\OperatorTok{().}\FunctionTok{get}\OperatorTok{(}\DecValTok{0}\OperatorTok{).}\FunctionTok{getCOSObject}\OperatorTok{()}
                \OperatorTok{.}\FunctionTok{setItem}\OperatorTok{(}\NormalTok{COSName}\OperatorTok{.}\FunctionTok{AP}\OperatorTok{,}\NormalTok{ apDict}\OperatorTok{);}
    \OperatorTok{\}}

    \CommentTok{// ── Certification ────────────────────────────────────────────────}

    \KeywordTok{private} \DataTypeTok{void} \FunctionTok{applyCertification}\OperatorTok{(}\NormalTok{PDDocument doc}\OperatorTok{,}
\NormalTok{                                     PDSignature signature}\OperatorTok{,}
                                     \DataTypeTok{int}\NormalTok{ level}\OperatorTok{)} \OperatorTok{\{}
        \DataTypeTok{int}\NormalTok{ p }\OperatorTok{=} \OperatorTok{(}\NormalTok{level }\OperatorTok{\textgreater{}=} \DecValTok{1} \OperatorTok{\&\&}\NormalTok{ level }\OperatorTok{\textless{}=} \DecValTok{3}\OperatorTok{)} \OperatorTok{?}\NormalTok{ level }\OperatorTok{:} \DecValTok{2}\OperatorTok{;}

\NormalTok{        COSDictionary sigRef }\OperatorTok{=} \KeywordTok{new} \FunctionTok{COSDictionary}\OperatorTok{();}
\NormalTok{        sigRef}\OperatorTok{.}\FunctionTok{setItem}\OperatorTok{(}\NormalTok{COSName}\OperatorTok{.}\FunctionTok{TYPE}\OperatorTok{,}\NormalTok{ COSName}\OperatorTok{.}\FunctionTok{getPDFName}\OperatorTok{(}\StringTok{"SigRef"}\OperatorTok{));}
\NormalTok{        sigRef}\OperatorTok{.}\FunctionTok{setItem}\OperatorTok{(}\NormalTok{COSName}\OperatorTok{.}\FunctionTok{getPDFName}\OperatorTok{(}\StringTok{"TransformMethod"}\OperatorTok{),}
\NormalTok{                COSName}\OperatorTok{.}\FunctionTok{getPDFName}\OperatorTok{(}\StringTok{"DocMDP"}\OperatorTok{));}

\NormalTok{        COSDictionary params }\OperatorTok{=} \KeywordTok{new} \FunctionTok{COSDictionary}\OperatorTok{();}
\NormalTok{        params}\OperatorTok{.}\FunctionTok{setItem}\OperatorTok{(}\NormalTok{COSName}\OperatorTok{.}\FunctionTok{TYPE}\OperatorTok{,}
\NormalTok{                COSName}\OperatorTok{.}\FunctionTok{getPDFName}\OperatorTok{(}\StringTok{"TransformParams"}\OperatorTok{));}
\NormalTok{        params}\OperatorTok{.}\FunctionTok{setItem}\OperatorTok{(}\NormalTok{COSName}\OperatorTok{.}\FunctionTok{getPDFName}\OperatorTok{(}\StringTok{"P"}\OperatorTok{),}
\NormalTok{                COSInteger}\OperatorTok{.}\FunctionTok{get}\OperatorTok{(}\NormalTok{p}\OperatorTok{));}
\NormalTok{        params}\OperatorTok{.}\FunctionTok{setItem}\OperatorTok{(}\NormalTok{COSName}\OperatorTok{.}\FunctionTok{getPDFName}\OperatorTok{(}\StringTok{"V"}\OperatorTok{),}
\NormalTok{                COSName}\OperatorTok{.}\FunctionTok{getPDFName}\OperatorTok{(}\StringTok{"1.2"}\OperatorTok{));}
\NormalTok{        sigRef}\OperatorTok{.}\FunctionTok{setItem}\OperatorTok{(}\NormalTok{COSName}\OperatorTok{.}\FunctionTok{getPDFName}\OperatorTok{(}\StringTok{"TransformParams"}\OperatorTok{),}\NormalTok{ params}\OperatorTok{);}

\NormalTok{        COSArray refArray }\OperatorTok{=} \KeywordTok{new} \FunctionTok{COSArray}\OperatorTok{();}
\NormalTok{        refArray}\OperatorTok{.}\FunctionTok{add}\OperatorTok{(}\NormalTok{sigRef}\OperatorTok{);}
\NormalTok{        signature}\OperatorTok{.}\FunctionTok{getCOSObject}\OperatorTok{()}
                \OperatorTok{.}\FunctionTok{setItem}\OperatorTok{(}\NormalTok{COSName}\OperatorTok{.}\FunctionTok{getPDFName}\OperatorTok{(}\StringTok{"Reference"}\OperatorTok{),}\NormalTok{ refArray}\OperatorTok{);}

\NormalTok{        COSDictionary perms }\OperatorTok{=} \KeywordTok{new} \FunctionTok{COSDictionary}\OperatorTok{();}
\NormalTok{        perms}\OperatorTok{.}\FunctionTok{setItem}\OperatorTok{(}\NormalTok{COSName}\OperatorTok{.}\FunctionTok{getPDFName}\OperatorTok{(}\StringTok{"DocMDP"}\OperatorTok{),}
\NormalTok{                signature}\OperatorTok{.}\FunctionTok{getCOSObject}\OperatorTok{());}
\NormalTok{        doc}\OperatorTok{.}\FunctionTok{getDocumentCatalog}\OperatorTok{().}\FunctionTok{getCOSObject}\OperatorTok{()}
                \OperatorTok{.}\FunctionTok{setItem}\OperatorTok{(}\NormalTok{COSName}\OperatorTok{.}\FunctionTok{getPDFName}\OperatorTok{(}\StringTok{"Perms"}\OperatorTok{),}\NormalTok{ perms}\OperatorTok{);}
    \OperatorTok{\}}

    \CommentTok{// ── Helpers ──────────────────────────────────────────────────────}

    \KeywordTok{private}\NormalTok{ PDSignatureField }\FunctionTok{findEmptySignatureField}\OperatorTok{(}\NormalTok{PDAcroForm acroForm}\OperatorTok{,}
                                                      \BuiltInTok{String}\NormalTok{ fieldName}\OperatorTok{)} \OperatorTok{\{}
        \ControlFlowTok{for} \OperatorTok{(}\NormalTok{PDField field }\OperatorTok{:}\NormalTok{ acroForm}\OperatorTok{.}\FunctionTok{getFieldTree}\OperatorTok{())} \OperatorTok{\{}
            \ControlFlowTok{if} \OperatorTok{(}\NormalTok{field }\KeywordTok{instanceof}\NormalTok{ PDSignatureField sf}\OperatorTok{)} \OperatorTok{\{}
                \BuiltInTok{String}\NormalTok{ fqn }\OperatorTok{=}\NormalTok{ sf}\OperatorTok{.}\FunctionTok{getFullyQualifiedName}\OperatorTok{();}
                \BuiltInTok{String}\NormalTok{ partial }\OperatorTok{=}\NormalTok{ sf}\OperatorTok{.}\FunctionTok{getPartialName}\OperatorTok{();}
                \ControlFlowTok{if} \OperatorTok{((}\NormalTok{fieldName}\OperatorTok{.}\FunctionTok{equals}\OperatorTok{(}\NormalTok{fqn}\OperatorTok{)} \OperatorTok{||}\NormalTok{ fieldName}\OperatorTok{.}\FunctionTok{equals}\OperatorTok{(}\NormalTok{partial}\OperatorTok{))}
                        \OperatorTok{\&\&}\NormalTok{ sf}\OperatorTok{.}\FunctionTok{getSignature}\OperatorTok{()} \OperatorTok{==} \KeywordTok{null}\OperatorTok{)} \OperatorTok{\{}
                    \ControlFlowTok{return}\NormalTok{ sf}\OperatorTok{;}
                \OperatorTok{\}}
            \OperatorTok{\}}
        \OperatorTok{\}}
        \ControlFlowTok{return} \KeywordTok{null}\OperatorTok{;}
    \OperatorTok{\}}

    \KeywordTok{private} \BuiltInTok{String} \FunctionTok{extractCN}\OperatorTok{(}\BuiltInTok{X509Certificate}\NormalTok{ cert}\OperatorTok{)} \OperatorTok{\{}
        \BuiltInTok{String}\NormalTok{ dn }\OperatorTok{=}\NormalTok{ cert}\OperatorTok{.}\FunctionTok{getSubjectX500Principal}\OperatorTok{().}\FunctionTok{getName}\OperatorTok{();}
\NormalTok{        java}\OperatorTok{.}\FunctionTok{util}\OperatorTok{.}\FunctionTok{regex}\OperatorTok{.}\FunctionTok{Matcher}\NormalTok{ m }\OperatorTok{=}\NormalTok{ java}\OperatorTok{.}\FunctionTok{util}\OperatorTok{.}\FunctionTok{regex}\OperatorTok{.}\FunctionTok{Pattern}
                \OperatorTok{.}\FunctionTok{compile}\OperatorTok{(}\StringTok{"CN=([\^{},]+)"}\OperatorTok{,}\NormalTok{ java}\OperatorTok{.}\FunctionTok{util}\OperatorTok{.}\FunctionTok{regex}\OperatorTok{.}\FunctionTok{Pattern}\OperatorTok{.}\FunctionTok{CASE\_INSENSITIVE}\OperatorTok{)}
                \OperatorTok{.}\FunctionTok{matcher}\OperatorTok{(}\NormalTok{dn}\OperatorTok{);}
        \ControlFlowTok{return}\NormalTok{ m}\OperatorTok{.}\FunctionTok{find}\OperatorTok{()} \OperatorTok{?}\NormalTok{ m}\OperatorTok{.}\FunctionTok{group}\OperatorTok{(}\DecValTok{1}\OperatorTok{).}\FunctionTok{trim}\OperatorTok{()} \OperatorTok{:} \StringTok{"Unknown"}\OperatorTok{;}
    \OperatorTok{\}}

    \KeywordTok{private} \DataTypeTok{static} \BuiltInTok{String} \FunctionTok{escapePs}\OperatorTok{(}\BuiltInTok{String}\NormalTok{ text}\OperatorTok{)} \OperatorTok{\{}
        \ControlFlowTok{return}\NormalTok{ text}\OperatorTok{.}\FunctionTok{replace}\OperatorTok{(}\StringTok{"}\SpecialCharTok{\textbackslash{}\textbackslash{}}\StringTok{"}\OperatorTok{,} \StringTok{"}\SpecialCharTok{\textbackslash{}\textbackslash{}\textbackslash{}\textbackslash{}}\StringTok{"}\OperatorTok{)}
                   \OperatorTok{.}\FunctionTok{replace}\OperatorTok{(}\StringTok{"("}\OperatorTok{,} \StringTok{"}\SpecialCharTok{\textbackslash{}\textbackslash{}}\StringTok{("}\OperatorTok{)}
                   \OperatorTok{.}\FunctionTok{replace}\OperatorTok{(}\StringTok{")"}\OperatorTok{,} \StringTok{"}\SpecialCharTok{\textbackslash{}\textbackslash{}}\StringTok{)"}\OperatorTok{);}
    \OperatorTok{\}}
\OperatorTok{\}}
\end{Highlighting}
\end{Shaded}

This class covers the full production workflow. It produces PAdES-B-B
compliant signatures by default (with
\texttt{ESS-signing-certificate-v2} and \texttt{ETSI.CAdES.detached}),
upgrades to PAdES-B-T when a TSA URL is provided, supports visible and
invisible appearances, and handles certification signatures with DocMDP.

In the next chapter, we go deep into the BouncyCastle CMS construction
that powers the \texttt{PadesCmsCallback} --- every signed attribute,
every ASN.1 structure, every decision point that determines whether your
signature is a basic CMS blob or a fully PAdES-compliant
\texttt{SignedData}.

\chapter{BouncyCastle CMS
Construction}\label{_bouncycastle_cms_construction}

Chapter 29 showed how PDFBox orchestrates the signing
process --- creating the signature dictionary, computing the byte range,
calling your \texttt{sign()} method, and writing the result into
\texttt{/Contents}. This chapter goes inside that \texttt{sign()}
method. We are now in BouncyCastle territory, building the CMS
\texttt{SignedData} structure that is the actual cryptographic heart of
every PDF signature.

If you have ever looked at a CMS signature in an ASN.1 dump and wondered
what each field means, why certain attributes are mandatory for PAdES,
or how a timestamp token gets attached --- this chapter answers all of
those questions with complete, compilable code.

We cover: \texttt{CMSSignedDataGenerator} and its components, provider
registration, building the \texttt{ContentSigner}, the four mandatory
signed attributes for PAdES-B-B, the \texttt{ESS-signing-certificate-v2}
attribute in detail, algorithm negotiation between RSA and ECDSA,
timestamp tokens as unsigned attributes, debugging CMS with ASN1Dump and
openssl, and complete builder classes for PAdES-B-B and PAdES-B-T.

\section{The CMS SignedData Structure: A Quick
Refresher}\label{_the_cms_signeddata_structure_a_quick_refresher}

Before we write code, let us recall what we are building. A CMS
\texttt{SignedData} structure (RFC 5652 --- Cryptographic Message Syntax
(CMS) S5) contains:

\begin{verbatim}
SignedData ::= SEQUENCE {
    version            CMSVersion,
    digestAlgorithms   SET OF AlgorithmIdentifier,
    encapContentInfo   EncapsulatedContentInfo,
    certificates    [0] IMPLICIT CertificateSet OPTIONAL,
    crls            [1] IMPLICIT RevocationInfoChoices OPTIONAL,
    signerInfos        SET OF SignerInfo
}
\end{verbatim}

For a PDF signature:

\begin{itemize}
\item
  \texttt{version} is 1 (one signer, detached content)
\item
  \texttt{digestAlgorithms} contains \texttt{SHA-256} (or SHA-384,
  SHA-512)
\item
  \texttt{encapContentInfo} has \texttt{contentType\ =\ id-data} and
  \emph{no content} (detached signature --- the content lives in the
  PDF)
\item
  \texttt{certificates} contains the signer's certificate and the chain
  (up to but not including the root)
\item
  \texttt{crls} is typically empty (revocation data goes in the DSS for
  PAdES-B-LT)
\item
  \texttt{signerInfos} contains exactly one \texttt{SignerInfo} with
  signed attributes, the signature value, and optionally unsigned
  attributes (the timestamp)
\end{itemize}

The \texttt{SignerInfo} is where most of the action is:

\begin{verbatim}
SignerInfo ::= SEQUENCE {
    version              CMSVersion,
    sid                  SignerIdentifier,
    digestAlgorithm      AlgorithmIdentifier,
    signedAttrs      [0] IMPLICIT SignedAttributes OPTIONAL,
    signatureAlgorithm   AlgorithmIdentifier,
    signature            SignatureValue,
    unsignedAttrs    [1] IMPLICIT UnsignedAttributes OPTIONAL
}
\end{verbatim}

The \texttt{signedAttrs} (signed attributes) are hashed and
signed --- they form the binding between the document content and the
signer's identity. The \texttt{unsignedAttrs} (unsigned attributes) are
\emph{not} covered by the signature --- this is where the timestamp
goes, because the timestamp is created \emph{after} the signature.

\section{BouncyCastle Provider
Registration}\label{_bouncycastle_provider_registration}

Every BouncyCastle crypto operation requires the provider to be
registered with the JCA (Java Cryptography Architecture). Without it,
calls to \texttt{JcaContentSignerBuilder},
\texttt{JcaDigestCalculatorProviderBuilder}, and \texttt{JcaCertStore}
fail with \texttt{NoSuchProviderException}.

\begin{Shaded}
\begin{Highlighting}[]
\KeywordTok{import} \ImportTok{org}\OperatorTok{.}\ImportTok{bouncycastle}\OperatorTok{.}\ImportTok{jce}\OperatorTok{.}\ImportTok{provider}\OperatorTok{.}\ImportTok{BouncyCastleProvider}\OperatorTok{;}

\CommentTok{// Register once, at application startup}
\ControlFlowTok{if} \OperatorTok{(}\NormalTok{java}\OperatorTok{.}\FunctionTok{security}\OperatorTok{.}\FunctionTok{Security}\OperatorTok{.}\FunctionTok{getProvider}\OperatorTok{(}\StringTok{"BC"}\OperatorTok{)} \OperatorTok{==} \KeywordTok{null}\OperatorTok{)} \OperatorTok{\{}
\NormalTok{    java}\OperatorTok{.}\FunctionTok{security}\OperatorTok{.}\FunctionTok{Security}\OperatorTok{.}\FunctionTok{addProvider}\OperatorTok{(}\KeywordTok{new} \FunctionTok{BouncyCastleProvider}\OperatorTok{());}
\OperatorTok{\}}
\end{Highlighting}
\end{Shaded}

In a Spring Boot application, the cleanest approach is a configuration
class:

\begin{Shaded}
\begin{Highlighting}[]
\KeywordTok{import} \ImportTok{jakarta}\OperatorTok{.}\ImportTok{annotation}\OperatorTok{.}\ImportTok{PostConstruct}\OperatorTok{;}
\KeywordTok{import} \ImportTok{org}\OperatorTok{.}\ImportTok{bouncycastle}\OperatorTok{.}\ImportTok{jce}\OperatorTok{.}\ImportTok{provider}\OperatorTok{.}\ImportTok{BouncyCastleProvider}\OperatorTok{;}
\KeywordTok{import} \ImportTok{org}\OperatorTok{.}\ImportTok{springframework}\OperatorTok{.}\ImportTok{context}\OperatorTok{.}\ImportTok{annotation}\OperatorTok{.}\ImportTok{Configuration}\OperatorTok{;}

\AttributeTok{@Configuration}
\KeywordTok{public} \KeywordTok{class}\NormalTok{ BouncyCastleConfig }\OperatorTok{\{}

    \AttributeTok{@PostConstruct}
    \KeywordTok{public} \DataTypeTok{void} \FunctionTok{init}\OperatorTok{()} \OperatorTok{\{}
        \ControlFlowTok{if} \OperatorTok{(}\NormalTok{java}\OperatorTok{.}\FunctionTok{security}\OperatorTok{.}\FunctionTok{Security}\OperatorTok{.}\FunctionTok{getProvider}\OperatorTok{(}\StringTok{"BC"}\OperatorTok{)} \OperatorTok{==} \KeywordTok{null}\OperatorTok{)} \OperatorTok{\{}
\NormalTok{            java}\OperatorTok{.}\FunctionTok{security}\OperatorTok{.}\FunctionTok{Security}\OperatorTok{.}\FunctionTok{addProvider}\OperatorTok{(}\KeywordTok{new} \FunctionTok{BouncyCastleProvider}\OperatorTok{());}
        \OperatorTok{\}}
    \OperatorTok{\}}
\OperatorTok{\}}
\end{Highlighting}
\end{Shaded}

If you register the provider in a \texttt{static} block, it races with
other static initializers. In complex applications with multiple
classloaders (application servers, OSGi), the provider may not be
visible where you need it. The \texttt{@PostConstruct} approach avoids
both problems.

You can verify the registration by calling
\texttt{java.security.Security.getProvider("BC")}. If it returns
\texttt{null}, BouncyCastle is not registered. If it returns a
\texttt{Provider} object, check its version:
\texttt{provider.getVersion()} should return \texttt{1.80} (or whatever
version you have on the classpath).

\section{Building a Basic CMS
Signature}\label{_building_a_basic_cms_signature}

Let us build a CMS \texttt{SignedData} from scratch, step by step. We
will start with the simplest possible version and then add PAdES
requirements.

\subsection{Step 1: The ContentSigner}\label{_step_1_the_contentsigner}

The \texttt{ContentSigner} is BouncyCastle's abstraction over the
private key and signature algorithm. It provides an
\texttt{OutputStream} that you write data into, and when you are done,
it produces the raw signature bytes.

\begin{Shaded}
\begin{Highlighting}[]
\KeywordTok{import} \ImportTok{org}\OperatorTok{.}\ImportTok{bouncycastle}\OperatorTok{.}\ImportTok{operator}\OperatorTok{.}\ImportTok{ContentSigner}\OperatorTok{;}
\KeywordTok{import} \ImportTok{org}\OperatorTok{.}\ImportTok{bouncycastle}\OperatorTok{.}\ImportTok{operator}\OperatorTok{.}\ImportTok{jcajce}\OperatorTok{.}\ImportTok{JcaContentSignerBuilder}\OperatorTok{;}

\CommentTok{// For RSA keys:}
\NormalTok{ContentSigner signer }\OperatorTok{=} \KeywordTok{new} \FunctionTok{JcaContentSignerBuilder}\OperatorTok{(}\StringTok{"SHA256withRSA"}\OperatorTok{)}
        \OperatorTok{.}\FunctionTok{setProvider}\OperatorTok{(}\StringTok{"BC"}\OperatorTok{)}
        \OperatorTok{.}\FunctionTok{build}\OperatorTok{(}\NormalTok{privateKey}\OperatorTok{);}

\CommentTok{// For EC keys:}
\NormalTok{ContentSigner signer }\OperatorTok{=} \KeywordTok{new} \FunctionTok{JcaContentSignerBuilder}\OperatorTok{(}\StringTok{"SHA256withECDSA"}\OperatorTok{)}
        \OperatorTok{.}\FunctionTok{setProvider}\OperatorTok{(}\StringTok{"BC"}\OperatorTok{)}
        \OperatorTok{.}\FunctionTok{build}\OperatorTok{(}\NormalTok{privateKey}\OperatorTok{);}
\end{Highlighting}
\end{Shaded}

The \texttt{JcaContentSignerBuilder} takes a JCA algorithm string and a
private key. The \texttt{.setProvider("BC")} ensures BouncyCastle
handles the actual signing, not the JDK's default provider (which may
not support all algorithms or may behave differently for edge cases like
PSS padding).

\subsection{Step 2: The
DigestCalculatorProvider}\label{_step_2_the_digestcalculatorprovider}

BouncyCastle needs a way to compute message digests. The
\texttt{DigestCalculatorProvider} is the factory for digest calculators:

\begin{Shaded}
\begin{Highlighting}[]
\KeywordTok{import} \ImportTok{org}\OperatorTok{.}\ImportTok{bouncycastle}\OperatorTok{.}\ImportTok{operator}\OperatorTok{.}\ImportTok{DigestCalculatorProvider}\OperatorTok{;}
\KeywordTok{import} \ImportTok{org}\OperatorTok{.}\ImportTok{bouncycastle}\OperatorTok{.}\ImportTok{operator}\OperatorTok{.}\ImportTok{jcajce}\OperatorTok{.}\ImportTok{JcaDigestCalculatorProviderBuilder}\OperatorTok{;}

\NormalTok{DigestCalculatorProvider dcProvider }\OperatorTok{=}
        \KeywordTok{new} \FunctionTok{JcaDigestCalculatorProviderBuilder}\OperatorTok{()}
                \OperatorTok{.}\FunctionTok{setProvider}\OperatorTok{(}\StringTok{"BC"}\OperatorTok{)}
                \OperatorTok{.}\FunctionTok{build}\OperatorTok{();}
\end{Highlighting}
\end{Shaded}

This provider is used internally by the signer info generator to compute
the message digest of the signed content and the signed attributes.

\subsection{Step 3: The
SignerInfoGenerator}\label{_step_3_the_signerinfogenerator}

The \texttt{SignerInfoGenerator} combines the \texttt{ContentSigner},
the \texttt{DigestCalculatorProvider}, and the signer's certificate into
a component that can produce a \texttt{SignerInfo}:

\begin{Shaded}
\begin{Highlighting}[]
\KeywordTok{import} \ImportTok{org}\OperatorTok{.}\ImportTok{bouncycastle}\OperatorTok{.}\ImportTok{cms}\OperatorTok{.}\ImportTok{SignerInfoGenerator}\OperatorTok{;}
\KeywordTok{import} \ImportTok{org}\OperatorTok{.}\ImportTok{bouncycastle}\OperatorTok{.}\ImportTok{cms}\OperatorTok{.}\ImportTok{jcajce}\OperatorTok{.}\ImportTok{JcaSignerInfoGeneratorBuilder}\OperatorTok{;}

\NormalTok{SignerInfoGenerator signerInfoGen }\OperatorTok{=}
        \KeywordTok{new} \FunctionTok{JcaSignerInfoGeneratorBuilder}\OperatorTok{(}\NormalTok{dcProvider}\OperatorTok{)}
                \OperatorTok{.}\FunctionTok{build}\OperatorTok{(}\NormalTok{signer}\OperatorTok{,}\NormalTok{ certificate}\OperatorTok{);}
\end{Highlighting}
\end{Shaded}

At this point, the \texttt{signerInfoGen} knows which key to use, which
algorithm, which certificate identifies the signer, and how to compute
digests. It does \emph{not} yet know about signed attributes --- we will
add those next.

\subsection{Step 4: The
CMSSignedDataGenerator}\label{_step_4_the_cmssigneddatagenerator}

Now we wire everything together:

\begin{Shaded}
\begin{Highlighting}[]
\KeywordTok{import} \ImportTok{org}\OperatorTok{.}\ImportTok{bouncycastle}\OperatorTok{.}\ImportTok{cert}\OperatorTok{.}\ImportTok{jcajce}\OperatorTok{.}\ImportTok{JcaCertStore}\OperatorTok{;}
\KeywordTok{import} \ImportTok{org}\OperatorTok{.}\ImportTok{bouncycastle}\OperatorTok{.}\ImportTok{cms}\OperatorTok{.}\ImportTok{CMSProcessableByteArray}\OperatorTok{;}
\KeywordTok{import} \ImportTok{org}\OperatorTok{.}\ImportTok{bouncycastle}\OperatorTok{.}\ImportTok{cms}\OperatorTok{.}\ImportTok{CMSSignedData}\OperatorTok{;}
\KeywordTok{import} \ImportTok{org}\OperatorTok{.}\ImportTok{bouncycastle}\OperatorTok{.}\ImportTok{cms}\OperatorTok{.}\ImportTok{CMSSignedDataGenerator}\OperatorTok{;}

\CommentTok{// Certificate store {-}{-} signer cert + chain}
\BuiltInTok{List}\OperatorTok{\textless{}}\BuiltInTok{X509Certificate}\OperatorTok{\textgreater{}}\NormalTok{ certList }\OperatorTok{=} \KeywordTok{new} \BuiltInTok{ArrayList}\OperatorTok{\textless{}\textgreater{}(}\NormalTok{chain}\OperatorTok{);}
\NormalTok{JcaCertStore certStore }\OperatorTok{=} \KeywordTok{new} \FunctionTok{JcaCertStore}\OperatorTok{(}\NormalTok{certList}\OperatorTok{);}

\NormalTok{CMSSignedDataGenerator gen }\OperatorTok{=} \KeywordTok{new} \FunctionTok{CMSSignedDataGenerator}\OperatorTok{();}
\NormalTok{gen}\OperatorTok{.}\FunctionTok{addSignerInfoGenerator}\OperatorTok{(}\NormalTok{signerInfoGen}\OperatorTok{);}
\NormalTok{gen}\OperatorTok{.}\FunctionTok{addCertificates}\OperatorTok{(}\NormalTok{certStore}\OperatorTok{);} 

\CommentTok{// The content to be signed {-}{-} this is the byte range content from PDFBox}
\NormalTok{CMSProcessableByteArray cmsContent }\OperatorTok{=}
        \KeywordTok{new} \FunctionTok{CMSProcessableByteArray}\OperatorTok{(}\NormalTok{contentBytes}\OperatorTok{);}

\CommentTok{// Generate the SignedData {-}{-} false = detached (content not embedded)}
\NormalTok{CMSSignedData signedData }\OperatorTok{=}\NormalTok{ gen}\OperatorTok{.}\FunctionTok{generate}\OperatorTok{(}\NormalTok{cmsContent}\OperatorTok{,} \KeywordTok{false}\OperatorTok{);} 

\CommentTok{// Get the DER{-}encoded result}
\DataTypeTok{byte}\OperatorTok{[]}\NormalTok{ cmsBytes }\OperatorTok{=}\NormalTok{ signedData}\OperatorTok{.}\FunctionTok{getEncoded}\OperatorTok{();}
\end{Highlighting}
\end{Shaded}

\begin{itemize}
\item
  The certificates are embedded in the CMS \texttt{certificates} field.
  Include the signer's certificate and all intermediate certificates. Do
  \emph{not} include the root CA certificate --- validators are expected
  to have the root in their trust store already, and including it wastes
  space without adding trust.
\item
  The second parameter \texttt{false} means detached signature. For PDF
  signatures, this must always be \texttt{false}. If you pass
  \texttt{true}, the content is embedded inside the CMS structure, which
  doubles the size and does not match what PDF validators expect for
  \texttt{ETSI.CAdES.detached} or \texttt{adbe.pkcs7.detached}.
\end{itemize}

\section{Signed Attributes: What PAdES
Requires}\label{_signed_attributes_what_pades_requires}

A basic CMS signature (without signed attributes) signs the content
directly --- the hash of the content is what the signature covers. When
signed attributes are present, the signature covers the hash of the
\emph{signed attributes} instead, and the signed attributes contain the
hash of the content. This indirection is how CMS allows metadata
(signing time, content type, certificate binding) to be included in the
signed data.

For PAdES-B-B compliance (ETSI EN 319 142-1 V1.1.1 (2016-04)), four
signed attributes are required.

\begin{longtable}[]{@{}
  >{\raggedright\arraybackslash}p{(\linewidth - 4\tabcolsep) * \real{0.2727}}
  >{\raggedright\arraybackslash}p{(\linewidth - 4\tabcolsep) * \real{0.2727}}
  >{\raggedright\arraybackslash}p{(\linewidth - 4\tabcolsep) * \real{0.4545}}@{}}
\caption{Mandatory Signed Attributes for PAdES-B-B}\tabularnewline
\toprule\noalign{}
\begin{minipage}[b]{\linewidth}\raggedright
Attribute
\end{minipage} & \begin{minipage}[b]{\linewidth}\raggedright
OID
\end{minipage} & \begin{minipage}[b]{\linewidth}\raggedright
Purpose
\end{minipage} \\
\midrule\noalign{}
\endfirsthead
\toprule\noalign{}
\begin{minipage}[b]{\linewidth}\raggedright
Attribute
\end{minipage} & \begin{minipage}[b]{\linewidth}\raggedright
OID
\end{minipage} & \begin{minipage}[b]{\linewidth}\raggedright
Purpose
\end{minipage} \\
\midrule\noalign{}
\endhead
\bottomrule\noalign{}
\endlastfoot
\texttt{content-type} & 1.2.840.113549.1.9.3 & Identifies the type of
content being signed. Must be \texttt{id-data} (OID
1.2.840.113549.1.7.1). BouncyCastle adds this automatically. \\
\texttt{message-digest} & 1.2.840.113549.1.9.4 & The hash of the signed
content (the byte range bytes). BouncyCastle computes and adds this
automatically. \\
\texttt{signing-time} & 1.2.840.113549.1.9.5 & The claimed signing time.
BouncyCastle adds this automatically when using
\texttt{DefaultSignedAttributeTableGenerator}. Note: PAdES validators
prefer the TSA timestamp over this value. \\
\texttt{ESS-signing-certificate-v2} & 1.2.840.113549.1.9.16.2.47 & Binds
the signature to a specific certificate by its hash. This is the one you
must add manually. Without it, you have a valid CMS signature but
\emph{not} a PAdES-B-B signature. \\
\end{longtable}

BouncyCastle's \texttt{DefaultSignedAttributeTableGenerator}
automatically adds \texttt{content-type}, \texttt{message-digest}, and
\texttt{signing-time}. You only need to explicitly construct
\texttt{ESS-signing-certificate-v2} and pass it in. The generator merges
your custom attributes with the automatic ones.

\section{Creating
ESS-signing-certificate-v2}\label{_creating_ess_signing_certificate_v2}

The \texttt{ESS-signing-certificate-v2} attribute (RFC 5652 ---
Cryptographic Message Syntax (CMS) and RFC 5035) prevents certificate
substitution attacks. It contains a hash of the signer's certificate, so
a validator can verify that the certificate embedded in the CMS
\texttt{certificates} field is the one the signer intended to use.

The ASN.1 structure:

\begin{verbatim}
SigningCertificateV2 ::= SEQUENCE {
    certs    SEQUENCE OF ESSCertIDv2,
    policies SEQUENCE OF PolicyInformation OPTIONAL
}

ESSCertIDv2 ::= SEQUENCE {
    hashAlgorithm  AlgorithmIdentifier DEFAULT {sha-256},
    certHash       OCTET STRING,
    issuerSerial   IssuerSerial OPTIONAL
}
\end{verbatim}

Here is the complete code to build this attribute:

\begin{Shaded}
\begin{Highlighting}[]
\KeywordTok{import} \ImportTok{java}\OperatorTok{.}\ImportTok{security}\OperatorTok{.}\ImportTok{MessageDigest}\OperatorTok{;}
\KeywordTok{import} \ImportTok{java}\OperatorTok{.}\ImportTok{security}\OperatorTok{.}\ImportTok{cert}\OperatorTok{.}\ImportTok{X509Certificate}\OperatorTok{;}

\KeywordTok{import} \ImportTok{org}\OperatorTok{.}\ImportTok{bouncycastle}\OperatorTok{.}\ImportTok{asn1}\OperatorTok{.}\ImportTok{ASN1EncodableVector}\OperatorTok{;}
\KeywordTok{import} \ImportTok{org}\OperatorTok{.}\ImportTok{bouncycastle}\OperatorTok{.}\ImportTok{asn1}\OperatorTok{.}\ImportTok{DERSet}\OperatorTok{;}
\KeywordTok{import} \ImportTok{org}\OperatorTok{.}\ImportTok{bouncycastle}\OperatorTok{.}\ImportTok{asn1}\OperatorTok{.}\ImportTok{cms}\OperatorTok{.}\ImportTok{Attribute}\OperatorTok{;}
\KeywordTok{import} \ImportTok{org}\OperatorTok{.}\ImportTok{bouncycastle}\OperatorTok{.}\ImportTok{asn1}\OperatorTok{.}\ImportTok{cms}\OperatorTok{.}\ImportTok{AttributeTable}\OperatorTok{;}
\KeywordTok{import} \ImportTok{org}\OperatorTok{.}\ImportTok{bouncycastle}\OperatorTok{.}\ImportTok{asn1}\OperatorTok{.}\ImportTok{ess}\OperatorTok{.}\ImportTok{ESSCertIDv2}\OperatorTok{;}
\KeywordTok{import} \ImportTok{org}\OperatorTok{.}\ImportTok{bouncycastle}\OperatorTok{.}\ImportTok{asn1}\OperatorTok{.}\ImportTok{ess}\OperatorTok{.}\ImportTok{SigningCertificateV2}\OperatorTok{;}
\KeywordTok{import} \ImportTok{org}\OperatorTok{.}\ImportTok{bouncycastle}\OperatorTok{.}\ImportTok{asn1}\OperatorTok{.}\ImportTok{nist}\OperatorTok{.}\ImportTok{NISTObjectIdentifiers}\OperatorTok{;}
\KeywordTok{import} \ImportTok{org}\OperatorTok{.}\ImportTok{bouncycastle}\OperatorTok{.}\ImportTok{asn1}\OperatorTok{.}\ImportTok{pkcs}\OperatorTok{.}\ImportTok{PKCSObjectIdentifiers}\OperatorTok{;}
\KeywordTok{import} \ImportTok{org}\OperatorTok{.}\ImportTok{bouncycastle}\OperatorTok{.}\ImportTok{asn1}\OperatorTok{.}\ImportTok{x509}\OperatorTok{.}\ImportTok{AlgorithmIdentifier}\OperatorTok{;}

\CommentTok{/**}
 \CommentTok{*}\NormalTok{ Builds the ESS}\CommentTok{{-}}\NormalTok{signing}\CommentTok{{-}}\NormalTok{certificate}\CommentTok{{-}}\NormalTok{v2 signed attribute}
 \CommentTok{*}\NormalTok{ required for PAdES}\CommentTok{{-}}\NormalTok{B}\CommentTok{{-}}\NormalTok{B compliance}\CommentTok{.}
 \CommentTok{*/}
\KeywordTok{public} \KeywordTok{class}\NormalTok{ EssCertAttributeBuilder }\OperatorTok{\{}

    \CommentTok{/**}
     \CommentTok{*}\NormalTok{ Creates an AttributeTable containing the}
     \CommentTok{*}\NormalTok{ ESS}\CommentTok{{-}}\NormalTok{signing}\CommentTok{{-}}\NormalTok{certificate}\CommentTok{{-}}\NormalTok{v2 attribute}\CommentTok{.}
     \CommentTok{*}
     \CommentTok{*} \AnnotationTok{@param}\CommentTok{ }\AnnotationTok{certificate}\CommentTok{ }\NormalTok{the signer}\CommentTok{\textquotesingle{}}\NormalTok{s X}\CommentTok{.509}\NormalTok{ certificate}
     \CommentTok{*} \AnnotationTok{@return}\NormalTok{ an AttributeTable that can be passed to}
     \CommentTok{*}\NormalTok{         DefaultSignedAttributeTableGenerator}
     \CommentTok{*/}
    \KeywordTok{public} \DataTypeTok{static}\NormalTok{ AttributeTable }\FunctionTok{build}\OperatorTok{(}\BuiltInTok{X509Certificate}\NormalTok{ certificate}\OperatorTok{)}
            \KeywordTok{throws} \BuiltInTok{Exception} \OperatorTok{\{}

        \CommentTok{// 1. Hash the DER{-}encoded certificate with SHA{-}256}
        \BuiltInTok{MessageDigest}\NormalTok{ md }\OperatorTok{=} \BuiltInTok{MessageDigest}\OperatorTok{.}\FunctionTok{getInstance}\OperatorTok{(}\StringTok{"SHA{-}256"}\OperatorTok{);}
        \DataTypeTok{byte}\OperatorTok{[]}\NormalTok{ certHash }\OperatorTok{=}\NormalTok{ md}\OperatorTok{.}\FunctionTok{digest}\OperatorTok{(}\NormalTok{certificate}\OperatorTok{.}\FunctionTok{getEncoded}\OperatorTok{());} 

        \CommentTok{// 2. Create the AlgorithmIdentifier for SHA{-}256}
\NormalTok{        AlgorithmIdentifier hashAlgo }\OperatorTok{=} \KeywordTok{new} \FunctionTok{AlgorithmIdentifier}\OperatorTok{(}
\NormalTok{                NISTObjectIdentifiers}\OperatorTok{.}\FunctionTok{id\_sha256}\OperatorTok{);} 

        \CommentTok{// 3. Create the ESSCertIDv2 (hash + algorithm)}
\NormalTok{        ESSCertIDv2 essCertId }\OperatorTok{=} \KeywordTok{new} \FunctionTok{ESSCertIDv2}\OperatorTok{(}\NormalTok{hashAlgo}\OperatorTok{,}\NormalTok{ certHash}\OperatorTok{);}

        \CommentTok{// 4. Wrap in SigningCertificateV2 (array of ESSCertIDv2)}
\NormalTok{        SigningCertificateV2 sigCertV2 }\OperatorTok{=} \KeywordTok{new} \FunctionTok{SigningCertificateV2}\OperatorTok{(}
                \KeywordTok{new}\NormalTok{ ESSCertIDv2}\OperatorTok{[]\{}\NormalTok{essCertId}\OperatorTok{\});} 

        \CommentTok{// 5. Create the CMS Attribute}
\NormalTok{        ASN1EncodableVector v }\OperatorTok{=} \KeywordTok{new} \FunctionTok{ASN1EncodableVector}\OperatorTok{();}
\NormalTok{        v}\OperatorTok{.}\FunctionTok{add}\OperatorTok{(}\KeywordTok{new} \BuiltInTok{Attribute}\OperatorTok{(}
\NormalTok{                PKCSObjectIdentifiers}\OperatorTok{.}\FunctionTok{id\_aa\_signingCertificateV2}\OperatorTok{,} 
                \KeywordTok{new} \FunctionTok{DERSet}\OperatorTok{(}\NormalTok{sigCertV2}\OperatorTok{)));}

        \ControlFlowTok{return} \KeywordTok{new} \FunctionTok{AttributeTable}\OperatorTok{(}\NormalTok{v}\OperatorTok{);}
    \OperatorTok{\}}
\OperatorTok{\}}
\end{Highlighting}
\end{Shaded}

\begin{itemize}
\item
  \texttt{certificate.getEncoded()} returns the DER-encoded X.509
  certificate. We hash this entire byte sequence, not just the public
  key or the subject name. This ensures that any change to the
  certificate (even adding an extension) changes the hash.
\item
  The hash algorithm identifier must match the algorithm used to compute
  the hash. SHA-256 is the standard choice. If you use SHA-384 or
  SHA-512 for the signature algorithm, you may use the matching digest
  here, but SHA-256 is always acceptable.
\item
  The \texttt{SigningCertificateV2} wraps an array of
  \texttt{ESSCertIDv2}. For a single signer, there is one entry. If your
  signing certificate was issued by an intermediate CA, you may
  optionally include the intermediate's hash as well, but this is not
  required by PAdES.
\item
  The OID \texttt{1.2.840.113549.1.9.16.2.47} identifies this as the
  \texttt{ESS-signing-certificate-v2} attribute. Do not confuse it with
  the older \texttt{ESS-signing-certificate} (v1, OID
  1.2.840.113549.1.9.16.2.12), which uses SHA-1. PAdES requires v2.
\end{itemize}

Using \texttt{ESS-signing-certificate} v1 (with SHA-1) instead of v2
(with SHA-256) is a PAdES compliance failure. Some older code examples
on the internet still use v1. SHA-1 is considered cryptographically
broken for collision resistance since the SHAttered attack (2017).
Always use v2 with SHA-256 or stronger.

\subsection{Wiring ESS-signing-certificate-v2 into the
Generator}\label{_wiring_ess_signing_certificate_v2_into_the_generator}

Now we pass the attribute table to the \texttt{SignerInfoGenerator}:

\begin{Shaded}
\begin{Highlighting}[]
\KeywordTok{import} \ImportTok{org}\OperatorTok{.}\ImportTok{bouncycastle}\OperatorTok{.}\ImportTok{cms}\OperatorTok{.}\ImportTok{DefaultSignedAttributeTableGenerator}\OperatorTok{;}
\KeywordTok{import} \ImportTok{org}\OperatorTok{.}\ImportTok{bouncycastle}\OperatorTok{.}\ImportTok{cms}\OperatorTok{.}\ImportTok{jcajce}\OperatorTok{.}\ImportTok{JcaSignerInfoGeneratorBuilder}\OperatorTok{;}

\CommentTok{// Build the ESS attribute table}
\NormalTok{AttributeTable essAttrs }\OperatorTok{=}\NormalTok{ EssCertAttributeBuilder}\OperatorTok{.}\FunctionTok{build}\OperatorTok{(}\NormalTok{certificate}\OperatorTok{);}

\CommentTok{// Create a generator that merges ESS attribute with the automatic ones}
\NormalTok{DefaultSignedAttributeTableGenerator attrGen }\OperatorTok{=}
        \KeywordTok{new} \FunctionTok{DefaultSignedAttributeTableGenerator}\OperatorTok{(}\NormalTok{essAttrs}\OperatorTok{);} 

\NormalTok{JcaSignerInfoGeneratorBuilder signerBuilder }\OperatorTok{=}
        \KeywordTok{new} \FunctionTok{JcaSignerInfoGeneratorBuilder}\OperatorTok{(}\NormalTok{dcProvider}\OperatorTok{);}
\NormalTok{signerBuilder}\OperatorTok{.}\FunctionTok{setSignedAttributeGenerator}\OperatorTok{(}\NormalTok{attrGen}\OperatorTok{);}

\NormalTok{SignerInfoGenerator signerInfoGen }\OperatorTok{=}
\NormalTok{        signerBuilder}\OperatorTok{.}\FunctionTok{build}\OperatorTok{(}\NormalTok{signer}\OperatorTok{,}\NormalTok{ certificate}\OperatorTok{);}
\end{Highlighting}
\end{Shaded}

\begin{itemize}
\item
  \texttt{DefaultSignedAttributeTableGenerator} accepts an initial
  \texttt{AttributeTable} and \emph{merges} it with the attributes it
  generates automatically (\texttt{content-type},
  \texttt{message-digest}, \texttt{signing-time}). Your
  \texttt{ESS-signing-certificate-v2} is included alongside the
  automatic attributes. You do not need to manually construct the other
  three.
\end{itemize}

\section{Algorithm Negotiation: RSA vs.
ECDSA}\label{_algorithm_negotiation_rsa_vs_ecdsa}

The signature algorithm depends on the private key type. You must match
the algorithm string to the key, or the \texttt{ContentSigner}
construction fails.

\begin{Shaded}
\begin{Highlighting}[]
\CommentTok{/**}
 \CommentTok{*}\NormalTok{ Determines the JCA signature algorithm string for a given private key}\CommentTok{.}
 \CommentTok{*}\NormalTok{ Supports RSA}\CommentTok{,}\NormalTok{ RSA}\CommentTok{{-}}\NormalTok{PSS}\CommentTok{,}\NormalTok{ ECDSA}\CommentTok{,}\NormalTok{ Ed25519}\CommentTok{,}\NormalTok{ and Ed448}\CommentTok{.}
 \CommentTok{*/}
\KeywordTok{public} \DataTypeTok{static} \BuiltInTok{String} \FunctionTok{resolveSignatureAlgorithm}\OperatorTok{(}\BuiltInTok{PrivateKey}\NormalTok{ privateKey}\OperatorTok{)} \OperatorTok{\{}
    \BuiltInTok{String}\NormalTok{ keyAlgo }\OperatorTok{=}\NormalTok{ privateKey}\OperatorTok{.}\FunctionTok{getAlgorithm}\OperatorTok{();}

    \ControlFlowTok{return} \ControlFlowTok{switch} \OperatorTok{(}\NormalTok{keyAlgo}\OperatorTok{)} \OperatorTok{\{}
        \ControlFlowTok{case} \StringTok{"RSA"} \OperatorTok{{-}\textgreater{}} \StringTok{"SHA256withRSA"}\OperatorTok{;}
        \ControlFlowTok{case} \StringTok{"RSASSA{-}PSS"} \OperatorTok{{-}\textgreater{}} \StringTok{"SHA256withRSAandMGF1"}\OperatorTok{;} 
        \ControlFlowTok{case} \StringTok{"EC"}\OperatorTok{,} \StringTok{"ECDSA"} \OperatorTok{{-}\textgreater{}} \StringTok{"SHA256withECDSA"}\OperatorTok{;}
        \ControlFlowTok{case} \StringTok{"Ed25519"} \OperatorTok{{-}\textgreater{}} \StringTok{"Ed25519"}\OperatorTok{;} 
        \ControlFlowTok{case} \StringTok{"Ed448"} \OperatorTok{{-}\textgreater{}} \StringTok{"Ed448"}\OperatorTok{;}
        \KeywordTok{default} \OperatorTok{{-}\textgreater{}} \ControlFlowTok{throw} \KeywordTok{new} \BuiltInTok{IllegalArgumentException}\OperatorTok{(}
                \StringTok{"Unsupported key algorithm: "} \OperatorTok{+}\NormalTok{ keyAlgo}\OperatorTok{);}
    \OperatorTok{\};}
\OperatorTok{\}}
\end{Highlighting}
\end{Shaded}

\begin{itemize}
\item
  RSA-PSS (Probabilistic Signature Scheme) is increasingly preferred
  over PKCS\#1 v1.5 padding (\texttt{SHA256withRSA}). Some EU validation
  services flag PKCS\#1 v1.5 as deprecated. If your key was generated as
  RSASSA-PSS, you must use the PSS algorithm string.
\item
  Ed25519 and Ed448 (Edwards curves) use the algorithm name
  directly --- there is no separate hash algorithm because the hash is
  built into the signature scheme. These are not yet widely supported by
  PDF validators (Adobe Acrobat does not validate Ed25519 signatures as
  of 2026), but they are valid per ISO 32000-2:2020 and will gain
  support over time.
\end{itemize}

\begin{longtable}[]{@{}
  >{\raggedright\arraybackslash}p{(\linewidth - 6\tabcolsep) * \real{0.1818}}
  >{\raggedright\arraybackslash}p{(\linewidth - 6\tabcolsep) * \real{0.2727}}
  >{\raggedright\arraybackslash}p{(\linewidth - 6\tabcolsep) * \real{0.2727}}
  >{\raggedright\arraybackslash}p{(\linewidth - 6\tabcolsep) * \real{0.2727}}@{}}
\caption{Signature Algorithm Matrix}\tabularnewline
\toprule\noalign{}
\begin{minipage}[b]{\linewidth}\raggedright
Key Type
\end{minipage} & \begin{minipage}[b]{\linewidth}\raggedright
Algorithm String
\end{minipage} & \begin{minipage}[b]{\linewidth}\raggedright
Hash
\end{minipage} & \begin{minipage}[b]{\linewidth}\raggedright
Notes
\end{minipage} \\
\midrule\noalign{}
\endfirsthead
\toprule\noalign{}
\begin{minipage}[b]{\linewidth}\raggedright
Key Type
\end{minipage} & \begin{minipage}[b]{\linewidth}\raggedright
Algorithm String
\end{minipage} & \begin{minipage}[b]{\linewidth}\raggedright
Hash
\end{minipage} & \begin{minipage}[b]{\linewidth}\raggedright
Notes
\end{minipage} \\
\midrule\noalign{}
\endhead
\bottomrule\noalign{}
\endlastfoot
RSA-2048+ & \texttt{SHA256withRSA} & SHA-256 & Most compatible. Works
everywhere. \\
RSA-2048+ & \texttt{SHA256withRSAandMGF1} & SHA-256 & RSA-PSS. Better
security properties. Some legacy validators do not support it. \\
RSA-3072+ & \texttt{SHA384withRSA} & SHA-384 & For higher security
requirements. \\
EC P-256 & \texttt{SHA256withECDSA} & SHA-256 & Compact signatures (64
bytes vs. 256 for RSA-2048). Widely supported. \\
EC P-384 & \texttt{SHA384withECDSA} & SHA-384 & Higher security
level. \\
EC P-521 & \texttt{SHA512withECDSA} & SHA-512 & Maximum EC security.
Rarely needed. \\
Ed25519 & \texttt{Ed25519} & Built-in & Modern curve. Limited validator
support (as of 2026). \\
Ed448 & \texttt{Ed448} & Built-in & Higher security Ed curve. Very
limited support. \\
\end{longtable}

For maximum compatibility in 2026, use \texttt{SHA256withRSA} (RSA-2048
key) or \texttt{SHA256withECDSA} (P-256 key). These work with every PDF
validator we have tested, including Adobe Acrobat, Foxit Reader, and the
EU DSS validation service.

\section{Timestamp Token as Unsigned
Attribute}\label{_timestamp_token_as_unsigned_attribute}

A PAdES-B-T signature contains a timestamp token as an unsigned
attribute on the \texttt{SignerInfo}. The timestamp proves when the
signature was created.

The attribute is identified by OID \texttt{1.2.840.113549.1.9.16.2.14}
(\texttt{id-aa-signatureTimeStampToken}), and its value is a complete
CMS \texttt{SignedData} structure --- the timestamp is itself a CMS
signature, created by the TSA.

\subsection{The Timestamping Flow}\label{_the_timestamping_flow}

\begin{enumerate}
\def\labelenumi{\arabic{enumi}.}
\item
  Generate the CMS \texttt{SignedData} for the document signature
  (PAdES-B-B)
\item
  Extract the raw signature value (\texttt{SignerInfo.signature})
\item
  Hash the signature value with SHA-256
\item
  Send the hash to the TSA in an RFC 3161 \texttt{TimeStampRequest}
\item
  Receive the \texttt{TimeStampResponse}, extract the
  \texttt{TimeStampToken}
\item
  Encode the token as an ASN.1 structure
\item
  Attach it as an unsigned attribute on the \texttt{SignerInfo}
\item
  Replace the \texttt{SignerInfo} in the \texttt{CMSSignedData}
\item
  Return the modified \texttt{CMSSignedData}
\end{enumerate}

\subsection{Complete Implementation}\label{_complete_implementation}

\begin{Shaded}
\begin{Highlighting}[]
\KeywordTok{import} \ImportTok{java}\OperatorTok{.}\ImportTok{security}\OperatorTok{.}\ImportTok{MessageDigest}\OperatorTok{;}
\KeywordTok{import} \ImportTok{java}\OperatorTok{.}\ImportTok{util}\OperatorTok{.}\ImportTok{Collections}\OperatorTok{;}

\KeywordTok{import} \ImportTok{org}\OperatorTok{.}\ImportTok{bouncycastle}\OperatorTok{.}\ImportTok{asn1}\OperatorTok{.}\ImportTok{ASN1EncodableVector}\OperatorTok{;}
\KeywordTok{import} \ImportTok{org}\OperatorTok{.}\ImportTok{bouncycastle}\OperatorTok{.}\ImportTok{asn1}\OperatorTok{.}\ImportTok{ASN1Primitive}\OperatorTok{;}
\KeywordTok{import} \ImportTok{org}\OperatorTok{.}\ImportTok{bouncycastle}\OperatorTok{.}\ImportTok{asn1}\OperatorTok{.}\ImportTok{DERSet}\OperatorTok{;}
\KeywordTok{import} \ImportTok{org}\OperatorTok{.}\ImportTok{bouncycastle}\OperatorTok{.}\ImportTok{asn1}\OperatorTok{.}\ImportTok{cms}\OperatorTok{.}\ImportTok{Attribute}\OperatorTok{;}
\KeywordTok{import} \ImportTok{org}\OperatorTok{.}\ImportTok{bouncycastle}\OperatorTok{.}\ImportTok{asn1}\OperatorTok{.}\ImportTok{cms}\OperatorTok{.}\ImportTok{AttributeTable}\OperatorTok{;}
\KeywordTok{import} \ImportTok{org}\OperatorTok{.}\ImportTok{bouncycastle}\OperatorTok{.}\ImportTok{asn1}\OperatorTok{.}\ImportTok{nist}\OperatorTok{.}\ImportTok{NISTObjectIdentifiers}\OperatorTok{;}
\KeywordTok{import} \ImportTok{org}\OperatorTok{.}\ImportTok{bouncycastle}\OperatorTok{.}\ImportTok{asn1}\OperatorTok{.}\ImportTok{pkcs}\OperatorTok{.}\ImportTok{PKCSObjectIdentifiers}\OperatorTok{;}
\KeywordTok{import} \ImportTok{org}\OperatorTok{.}\ImportTok{bouncycastle}\OperatorTok{.}\ImportTok{cms}\OperatorTok{.}\ImportTok{CMSSignedData}\OperatorTok{;}
\KeywordTok{import} \ImportTok{org}\OperatorTok{.}\ImportTok{bouncycastle}\OperatorTok{.}\ImportTok{cms}\OperatorTok{.}\ImportTok{SignerInformation}\OperatorTok{;}
\KeywordTok{import} \ImportTok{org}\OperatorTok{.}\ImportTok{bouncycastle}\OperatorTok{.}\ImportTok{cms}\OperatorTok{.}\ImportTok{SignerInformationStore}\OperatorTok{;}
\KeywordTok{import} \ImportTok{org}\OperatorTok{.}\ImportTok{bouncycastle}\OperatorTok{.}\ImportTok{tsp}\OperatorTok{.}\ImportTok{TimeStampRequest}\OperatorTok{;}
\KeywordTok{import} \ImportTok{org}\OperatorTok{.}\ImportTok{bouncycastle}\OperatorTok{.}\ImportTok{tsp}\OperatorTok{.}\ImportTok{TimeStampRequestGenerator}\OperatorTok{;}
\KeywordTok{import} \ImportTok{org}\OperatorTok{.}\ImportTok{bouncycastle}\OperatorTok{.}\ImportTok{tsp}\OperatorTok{.}\ImportTok{TimeStampResponse}\OperatorTok{;}
\KeywordTok{import} \ImportTok{org}\OperatorTok{.}\ImportTok{bouncycastle}\OperatorTok{.}\ImportTok{tsp}\OperatorTok{.}\ImportTok{TimeStampToken}\OperatorTok{;}

\CommentTok{/**}
 \CommentTok{*}\NormalTok{ Adds an RFC }\CommentTok{3161}\NormalTok{ timestamp to a CMS SignedData}\CommentTok{,}\NormalTok{ upgrading it}
 \CommentTok{*}\NormalTok{ from PAdES}\CommentTok{{-}}\NormalTok{B}\CommentTok{{-}}\NormalTok{B to PAdES}\CommentTok{{-}}\NormalTok{B}\CommentTok{{-}}\NormalTok{T}\CommentTok{.}
 \CommentTok{*/}
\KeywordTok{public} \KeywordTok{class}\NormalTok{ TimestampEmbedder }\OperatorTok{\{}

    \CommentTok{/**}
     \CommentTok{*}\NormalTok{ Embeds a timestamp token as an unsigned attribute}\CommentTok{.}
     \CommentTok{*}
     \CommentTok{*} \AnnotationTok{@param}\CommentTok{ }\AnnotationTok{signedData}\CommentTok{ }\NormalTok{the original CMS SignedData }\CommentTok{(}\NormalTok{PAdES}\CommentTok{{-}}\NormalTok{B}\CommentTok{{-}}\NormalTok{B}\CommentTok{)}
     \CommentTok{*} \AnnotationTok{@param}\CommentTok{ }\AnnotationTok{tsaUrl}\CommentTok{     }\NormalTok{the RFC }\CommentTok{3161}\NormalTok{ TSA endpoint URL}
     \CommentTok{*} \AnnotationTok{@return}\NormalTok{ the modified CMS SignedData }\CommentTok{(}\NormalTok{PAdES}\CommentTok{{-}}\NormalTok{B}\CommentTok{{-}}\NormalTok{T}\CommentTok{)}
     \CommentTok{*/}
    \KeywordTok{public} \DataTypeTok{static}\NormalTok{ CMSSignedData }\FunctionTok{addTimestamp}\OperatorTok{(}\NormalTok{CMSSignedData signedData}\OperatorTok{,}
                                              \BuiltInTok{String}\NormalTok{ tsaUrl}\OperatorTok{)}
            \KeywordTok{throws} \BuiltInTok{Exception} \OperatorTok{\{}

        \CommentTok{// Get the first (only) signer info}
\NormalTok{        SignerInformation signerInfo }\OperatorTok{=}\NormalTok{ signedData}\OperatorTok{.}\FunctionTok{getSignerInfos}\OperatorTok{()}
                \OperatorTok{.}\FunctionTok{getSigners}\OperatorTok{().}\FunctionTok{iterator}\OperatorTok{().}\FunctionTok{next}\OperatorTok{();}

        \CommentTok{// The timestamp covers the signature value, not the content}
        \DataTypeTok{byte}\OperatorTok{[]}\NormalTok{ signatureBytes }\OperatorTok{=}\NormalTok{ signerInfo}\OperatorTok{.}\FunctionTok{getSignature}\OperatorTok{();} 

        \CommentTok{// Hash the signature for the TSA request}
        \BuiltInTok{MessageDigest}\NormalTok{ md }\OperatorTok{=} \BuiltInTok{MessageDigest}\OperatorTok{.}\FunctionTok{getInstance}\OperatorTok{(}\StringTok{"SHA{-}256"}\OperatorTok{);}
        \DataTypeTok{byte}\OperatorTok{[]}\NormalTok{ hash }\OperatorTok{=}\NormalTok{ md}\OperatorTok{.}\FunctionTok{digest}\OperatorTok{(}\NormalTok{signatureBytes}\OperatorTok{);}

        \CommentTok{// Build and send the RFC 3161 request}
\NormalTok{        TimeStampRequestGenerator reqGen }\OperatorTok{=}
                \KeywordTok{new} \FunctionTok{TimeStampRequestGenerator}\OperatorTok{();}
\NormalTok{        reqGen}\OperatorTok{.}\FunctionTok{setCertReq}\OperatorTok{(}\KeywordTok{true}\OperatorTok{);} 
\NormalTok{        TimeStampRequest tsReq }\OperatorTok{=}\NormalTok{ reqGen}\OperatorTok{.}\FunctionTok{generate}\OperatorTok{(}
\NormalTok{                NISTObjectIdentifiers}\OperatorTok{.}\FunctionTok{id\_sha256}\OperatorTok{,}\NormalTok{ hash}\OperatorTok{);}

        \DataTypeTok{byte}\OperatorTok{[]}\NormalTok{ tsRespBytes }\OperatorTok{=} \FunctionTok{sendTsaRequest}\OperatorTok{(}\NormalTok{tsaUrl}\OperatorTok{,}\NormalTok{ tsReq}\OperatorTok{.}\FunctionTok{getEncoded}\OperatorTok{());}

        \CommentTok{// Parse and validate the response}
\NormalTok{        TimeStampResponse tsResponse }\OperatorTok{=}
                \KeywordTok{new} \FunctionTok{TimeStampResponse}\OperatorTok{(}\NormalTok{tsRespBytes}\OperatorTok{);}
\NormalTok{        tsResponse}\OperatorTok{.}\FunctionTok{validate}\OperatorTok{(}\NormalTok{tsReq}\OperatorTok{);} 

\NormalTok{        TimeStampToken tsToken }\OperatorTok{=}\NormalTok{ tsResponse}\OperatorTok{.}\FunctionTok{getTimeStampToken}\OperatorTok{();}
        \ControlFlowTok{if} \OperatorTok{(}\NormalTok{tsToken }\OperatorTok{==} \KeywordTok{null}\OperatorTok{)} \OperatorTok{\{}
            \ControlFlowTok{throw} \KeywordTok{new} \BuiltInTok{IllegalStateException}\OperatorTok{(}
                    \StringTok{"TSA response contains no timestamp token"}\OperatorTok{);}
        \OperatorTok{\}}

        \CommentTok{// Convert to ASN.1 for embedding}
\NormalTok{        ASN1Primitive tsAsn1 }\OperatorTok{=}\NormalTok{ ASN1Primitive}\OperatorTok{.}\FunctionTok{fromByteArray}\OperatorTok{(}
\NormalTok{                tsToken}\OperatorTok{.}\FunctionTok{getEncoded}\OperatorTok{());}

        \CommentTok{// Build the unsigned attribute}
\NormalTok{        ASN1EncodableVector unsignedVec }\OperatorTok{=} \KeywordTok{new} \FunctionTok{ASN1EncodableVector}\OperatorTok{();}
\NormalTok{        unsignedVec}\OperatorTok{.}\FunctionTok{add}\OperatorTok{(}\KeywordTok{new} \BuiltInTok{Attribute}\OperatorTok{(}
\NormalTok{                PKCSObjectIdentifiers}\OperatorTok{.}\FunctionTok{id\_aa\_signatureTimeStampToken}\OperatorTok{,}
                \KeywordTok{new} \FunctionTok{DERSet}\OperatorTok{(}\NormalTok{tsAsn1}\OperatorTok{)));} 
\NormalTok{        AttributeTable unsignedAttrs }\OperatorTok{=} \KeywordTok{new} \FunctionTok{AttributeTable}\OperatorTok{(}\NormalTok{unsignedVec}\OperatorTok{);}

        \CommentTok{// Replace the signer\textquotesingle{}s unsigned attributes}
\NormalTok{        SignerInformation newSignerInfo }\OperatorTok{=}
\NormalTok{                SignerInformation}\OperatorTok{.}\FunctionTok{replaceUnsignedAttributes}\OperatorTok{(}
\NormalTok{                        signerInfo}\OperatorTok{,}\NormalTok{ unsignedAttrs}\OperatorTok{);} 

\NormalTok{        SignerInformationStore newSignerInfos }\OperatorTok{=}
                \KeywordTok{new} \FunctionTok{SignerInformationStore}\OperatorTok{(}
                        \BuiltInTok{Collections}\OperatorTok{.}\FunctionTok{singletonList}\OperatorTok{(}\NormalTok{newSignerInfo}\OperatorTok{));}

        \ControlFlowTok{return}\NormalTok{ CMSSignedData}\OperatorTok{.}\FunctionTok{replaceSigners}\OperatorTok{(}
\NormalTok{                signedData}\OperatorTok{,}\NormalTok{ newSignerInfos}\OperatorTok{);}
    \OperatorTok{\}}

    \CommentTok{/**}\NormalTok{ Sends an RFC }\CommentTok{3161}\NormalTok{ request to a TSA endpoint via HTTP POST}\CommentTok{. */}
    \KeywordTok{private} \DataTypeTok{static} \DataTypeTok{byte}\OperatorTok{[]} \FunctionTok{sendTsaRequest}\OperatorTok{(}\BuiltInTok{String}\NormalTok{ tsaUrl}\OperatorTok{,} \DataTypeTok{byte}\OperatorTok{[]}\NormalTok{ reqBytes}\OperatorTok{)}
            \KeywordTok{throws} \BuiltInTok{Exception} \OperatorTok{\{}

\NormalTok{        java}\OperatorTok{.}\FunctionTok{net}\OperatorTok{.}\FunctionTok{http}\OperatorTok{.}\FunctionTok{HttpClient}\NormalTok{ client }\OperatorTok{=}
\NormalTok{                java}\OperatorTok{.}\FunctionTok{net}\OperatorTok{.}\FunctionTok{http}\OperatorTok{.}\FunctionTok{HttpClient}\OperatorTok{.}\FunctionTok{newBuilder}\OperatorTok{()}
                        \OperatorTok{.}\FunctionTok{connectTimeout}\OperatorTok{(}\NormalTok{java}\OperatorTok{.}\FunctionTok{time}\OperatorTok{.}\FunctionTok{Duration}\OperatorTok{.}\FunctionTok{ofSeconds}\OperatorTok{(}\DecValTok{15}\OperatorTok{))}
                        \OperatorTok{.}\FunctionTok{build}\OperatorTok{();}

\NormalTok{        java}\OperatorTok{.}\FunctionTok{net}\OperatorTok{.}\FunctionTok{http}\OperatorTok{.}\FunctionTok{HttpRequest}\NormalTok{ request }\OperatorTok{=}
\NormalTok{                java}\OperatorTok{.}\FunctionTok{net}\OperatorTok{.}\FunctionTok{http}\OperatorTok{.}\FunctionTok{HttpRequest}\OperatorTok{.}\FunctionTok{newBuilder}\OperatorTok{()}
                        \OperatorTok{.}\FunctionTok{uri}\OperatorTok{(}\NormalTok{java}\OperatorTok{.}\FunctionTok{net}\OperatorTok{.}\FunctionTok{URI}\OperatorTok{.}\FunctionTok{create}\OperatorTok{(}\NormalTok{tsaUrl}\OperatorTok{))}
                        \OperatorTok{.}\FunctionTok{header}\OperatorTok{(}\StringTok{"Content{-}Type"}\OperatorTok{,}
                                \StringTok{"application/timestamp{-}query"}\OperatorTok{)}
                        \OperatorTok{.}\FunctionTok{timeout}\OperatorTok{(}\NormalTok{java}\OperatorTok{.}\FunctionTok{time}\OperatorTok{.}\FunctionTok{Duration}\OperatorTok{.}\FunctionTok{ofSeconds}\OperatorTok{(}\DecValTok{30}\OperatorTok{))}
                        \OperatorTok{.}\FunctionTok{POST}\OperatorTok{(}\NormalTok{java}\OperatorTok{.}\FunctionTok{net}\OperatorTok{.}\FunctionTok{http}\OperatorTok{.}\FunctionTok{HttpRequest}\OperatorTok{.}\FunctionTok{BodyPublishers}
                                \OperatorTok{.}\FunctionTok{ofByteArray}\OperatorTok{(}\NormalTok{reqBytes}\OperatorTok{))}
                        \OperatorTok{.}\FunctionTok{build}\OperatorTok{();}

\NormalTok{        java}\OperatorTok{.}\FunctionTok{net}\OperatorTok{.}\FunctionTok{http}\OperatorTok{.}\FunctionTok{HttpResponse}\OperatorTok{\textless{}}\DataTypeTok{byte}\OperatorTok{[]\textgreater{}}\NormalTok{ response }\OperatorTok{=}\NormalTok{ client}\OperatorTok{.}\FunctionTok{send}\OperatorTok{(}
\NormalTok{                request}\OperatorTok{,}
\NormalTok{                java}\OperatorTok{.}\FunctionTok{net}\OperatorTok{.}\FunctionTok{http}\OperatorTok{.}\FunctionTok{HttpResponse}\OperatorTok{.}\FunctionTok{BodyHandlers}\OperatorTok{.}\FunctionTok{ofByteArray}\OperatorTok{());}

        \ControlFlowTok{if} \OperatorTok{(}\NormalTok{response}\OperatorTok{.}\FunctionTok{statusCode}\OperatorTok{()} \OperatorTok{!=} \DecValTok{200}\OperatorTok{)} \OperatorTok{\{}
            \ControlFlowTok{throw} \KeywordTok{new}\NormalTok{ java}\OperatorTok{.}\FunctionTok{io}\OperatorTok{.}\FunctionTok{IOException}\OperatorTok{(}
                    \StringTok{"TSA returned HTTP "} \OperatorTok{+}\NormalTok{ response}\OperatorTok{.}\FunctionTok{statusCode}\OperatorTok{()}
                            \OperatorTok{+} \StringTok{" from "} \OperatorTok{+}\NormalTok{ tsaUrl}\OperatorTok{);}
        \OperatorTok{\}}

        \DataTypeTok{byte}\OperatorTok{[]}\NormalTok{ body }\OperatorTok{=}\NormalTok{ response}\OperatorTok{.}\FunctionTok{body}\OperatorTok{();}
        \ControlFlowTok{if} \OperatorTok{(}\NormalTok{body }\OperatorTok{==} \KeywordTok{null} \OperatorTok{||}\NormalTok{ body}\OperatorTok{.}\FunctionTok{length} \OperatorTok{==} \DecValTok{0}\OperatorTok{)} \OperatorTok{\{}
            \ControlFlowTok{throw} \KeywordTok{new}\NormalTok{ java}\OperatorTok{.}\FunctionTok{io}\OperatorTok{.}\FunctionTok{IOException}\OperatorTok{(}
                    \StringTok{"TSA returned empty response from "} \OperatorTok{+}\NormalTok{ tsaUrl}\OperatorTok{);}
        \OperatorTok{\}}

        \ControlFlowTok{return}\NormalTok{ body}\OperatorTok{;}
    \OperatorTok{\}}
\OperatorTok{\}}
\end{Highlighting}
\end{Shaded}

\begin{itemize}
\item
  \texttt{getSignature()} returns the raw signature bytes from the
  \texttt{SignerInfo} --- this is the value that was produced by the
  \texttt{ContentSigner}. For RSA-2048, this is 256 bytes. For ECDSA
  P-256, it is approximately 72 bytes (DER-encoded).
\item
  \texttt{setCertReq(true)} tells the TSA to include its signing
  certificate in the response. This is important: without the TSA's
  certificate, a validator cannot verify the timestamp without
  separately fetching the TSA's certificate.
\item
  \texttt{validate()} checks that the response corresponds to our
  request (matching nonce, hash algorithm, and hash value). If the TSA
  returned a response for a different request, this throws a
  \texttt{TSPValidationException}.
\item
  The timestamp token is wrapped in a \texttt{DERSet} and assigned OID
  \texttt{1.2.840.113549.1.9.16.2.14}. This is the standard way to embed
  a timestamp in CMS (RFC 5652 --- Cryptographic Message Syntax (CMS)).
\item
  \texttt{replaceUnsignedAttributes()} creates a new
  \texttt{SignerInformation} with the same signed attributes and
  signature value, but with the specified unsigned attributes. The
  original \texttt{SignerInformation} is immutable.
\end{itemize}

\section{Debugging CMS Structures}\label{_debugging_cms_structures}

When a signature fails validation, you need to look inside the CMS blob.
There are two indispensable tools: BouncyCastle's \texttt{ASN1Dump}
(Java) and OpenSSL (command line).

\subsection{ASN1Dump: Java-Side
Inspection}\label{_asn1dump_java_side_inspection}

\begin{Shaded}
\begin{Highlighting}[]
\KeywordTok{import} \ImportTok{org}\OperatorTok{.}\ImportTok{bouncycastle}\OperatorTok{.}\ImportTok{asn1}\OperatorTok{.}\ImportTok{ASN1InputStream}\OperatorTok{;}
\KeywordTok{import} \ImportTok{org}\OperatorTok{.}\ImportTok{bouncycastle}\OperatorTok{.}\ImportTok{asn1}\OperatorTok{.}\ImportTok{ASN1Primitive}\OperatorTok{;}
\KeywordTok{import} \ImportTok{org}\OperatorTok{.}\ImportTok{bouncycastle}\OperatorTok{.}\ImportTok{asn1}\OperatorTok{.}\ImportTok{util}\OperatorTok{.}\ImportTok{ASN1Dump}\OperatorTok{;}

\CommentTok{/**}
 \CommentTok{*}\NormalTok{ Dumps the ASN}\CommentTok{.1}\NormalTok{ structure of a CMS blob to stdout}\CommentTok{.}
 \CommentTok{*}\NormalTok{ Invaluable for debugging signature construction}\CommentTok{.}
 \CommentTok{*/}
\KeywordTok{public} \DataTypeTok{static} \DataTypeTok{void} \FunctionTok{dumpCms}\OperatorTok{(}\DataTypeTok{byte}\OperatorTok{[]}\NormalTok{ cmsBytes}\OperatorTok{)} \KeywordTok{throws} \BuiltInTok{Exception} \OperatorTok{\{}
    \ControlFlowTok{try} \OperatorTok{(}\NormalTok{ASN1InputStream asn1In }\OperatorTok{=} \KeywordTok{new} \FunctionTok{ASN1InputStream}\OperatorTok{(}\NormalTok{cmsBytes}\OperatorTok{))} \OperatorTok{\{}
\NormalTok{        ASN1Primitive obj }\OperatorTok{=}\NormalTok{ asn1In}\OperatorTok{.}\FunctionTok{readObject}\OperatorTok{();}
        \BuiltInTok{String}\NormalTok{ dump }\OperatorTok{=}\NormalTok{ ASN1Dump}\OperatorTok{.}\FunctionTok{dumpAsString}\OperatorTok{(}\NormalTok{obj}\OperatorTok{,} \KeywordTok{true}\OperatorTok{);} 
        \BuiltInTok{System}\OperatorTok{.}\FunctionTok{out}\OperatorTok{.}\FunctionTok{println}\OperatorTok{(}\NormalTok{dump}\OperatorTok{);}
    \OperatorTok{\}}
\OperatorTok{\}}
\end{Highlighting}
\end{Shaded}

\begin{itemize}
\item
  The \texttt{true} parameter enables verbose mode, which shows the full
  content of OCTET STRING and BIT STRING values. Without it, large
  values are truncated.
\end{itemize}

The output looks like this (truncated for readability):

\begin{verbatim}
SEQUENCE
  OBJECT IDENTIFIER 1.2.840.113549.1.7.2    (signedData)
  [0]
    SEQUENCE
      INTEGER 1                               (version)
      SET
        SEQUENCE
          OBJECT IDENTIFIER 2.16.840.1.101.3.4.2.1 (SHA-256)
      SEQUENCE
        OBJECT IDENTIFIER 1.2.840.113549.1.7.1     (data)
      [0]                                     (certificates)
        SEQUENCE
          ...                                 (X.509 certificate)
      SET
        SEQUENCE                              (SignerInfo)
          INTEGER 1                           (version)
          SEQUENCE                            (sid - issuerAndSerialNumber)
            ...
          SEQUENCE                            (digestAlgorithm)
            OBJECT IDENTIFIER 2.16.840.1.101.3.4.2.1 (SHA-256)
          [0]                                 (signedAttrs)
            SEQUENCE
              OBJECT IDENTIFIER 1.2.840.113549.1.9.3  (content-type)
              ...
            SEQUENCE
              OBJECT IDENTIFIER 1.2.840.113549.1.9.4  (message-digest)
              ...
            SEQUENCE
              OBJECT IDENTIFIER 1.2.840.113549.1.9.5  (signing-time)
              ...
            SEQUENCE
              OBJECT IDENTIFIER 1.2.840.113549.1.9.16.2.47 (ESS-signing-cert-v2)
              ...
          SEQUENCE                            (signatureAlgorithm)
            OBJECT IDENTIFIER 1.2.840.113549.1.1.1 (RSA)
          OCTET STRING                        (signature value)
            ...
          [1]                                 (unsignedAttrs)
            SEQUENCE
              OBJECT IDENTIFIER 1.2.840.113549.1.9.16.2.14 (timestamp)
              ...
\end{verbatim}

When debugging, check for:

\begin{enumerate}
\def\labelenumi{\arabic{enumi}.}
\item
  Is \texttt{ESS-signing-certificate-v2} (OID \ldots\hspace{0pt}16.2.47)
  present in the signed attributes? If not, the signature is not
  PAdES-B-B compliant.
\item
  Is \texttt{message-digest} present? If not, the signature structure is
  malformed.
\item
  Is the timestamp (OID \ldots\hspace{0pt}16.2.14) in the unsigned
  attributes? If not, it is not PAdES-B-T.
\item
  Are certificates present in the \texttt{{[}0{]}} IMPLICIT set? If not,
  validators cannot build the chain.
\end{enumerate}

\subsection{OpenSSL: Command-Line
Inspection}\label{_openssl_command_line_inspection}

Extract the CMS blob from the PDF (it is the hex-encoded content of
\texttt{/Contents}), decode it, and use OpenSSL:

\begin{Shaded}
\begin{Highlighting}[]
\CommentTok{\# Decode the hex string from /Contents to binary}
\ExtensionTok{xxd} \AttributeTok{{-}r} \AttributeTok{{-}p}\NormalTok{ signature{-}hex.txt }\OperatorTok{\textgreater{}}\NormalTok{ signature.der}

\CommentTok{\# Parse the CMS structure}
\ExtensionTok{openssl}\NormalTok{ cms }\AttributeTok{{-}inform}\NormalTok{ DER }\AttributeTok{{-}in}\NormalTok{ signature.der }\AttributeTok{{-}cmsout} \AttributeTok{{-}print}

\CommentTok{\# Verify the CMS signature (requires the signed content)}
\ExtensionTok{openssl}\NormalTok{ cms }\AttributeTok{{-}verify} \AttributeTok{{-}inform}\NormalTok{ DER }\AttributeTok{{-}in}\NormalTok{ signature.der }\DataTypeTok{\textbackslash{}}
    \AttributeTok{{-}content}\NormalTok{ signed{-}bytes.bin }\AttributeTok{{-}noverify} \AttributeTok{{-}out}\NormalTok{ /dev/null}

\CommentTok{\# Extract and examine the signing certificate}
\ExtensionTok{openssl}\NormalTok{ cms }\AttributeTok{{-}inform}\NormalTok{ DER }\AttributeTok{{-}in}\NormalTok{ signature.der }\AttributeTok{{-}cmsout} \AttributeTok{{-}print} \DataTypeTok{\textbackslash{}}
    \KeywordTok{|} \FunctionTok{grep} \AttributeTok{{-}A}\NormalTok{ 50 }\StringTok{"cert:"}

\CommentTok{\# View the timestamp token (if present)}
\ExtensionTok{openssl}\NormalTok{ ts }\AttributeTok{{-}reply} \AttributeTok{{-}in}\NormalTok{ timestamp{-}token.der }\AttributeTok{{-}text}
\end{Highlighting}
\end{Shaded}

To extract the \texttt{/Contents} hex string from a PDF, search for
\texttt{/Contents\ \textless{}} in the raw file. The hex string starts
after the \texttt{\textless{}} and ends at \texttt{\textgreater{}}. It
is zero-padded --- strip trailing zeros before decoding. You can also
use PDFBox's \texttt{PDSignature.getContents(byte{[}{]})} to extract the
raw bytes programmatically.

\section{Complete PAdES-B-B Builder}\label{_complete_pades_b_b_builder}

Here is a self-contained class that builds a PAdES-B-B compliant CMS
\texttt{SignedData}:

\begin{Shaded}
\begin{Highlighting}[]
\KeywordTok{import} \ImportTok{java}\OperatorTok{.}\ImportTok{io}\OperatorTok{.}\ImportTok{IOException}\OperatorTok{;}
\KeywordTok{import} \ImportTok{java}\OperatorTok{.}\ImportTok{io}\OperatorTok{.}\ImportTok{InputStream}\OperatorTok{;}
\KeywordTok{import} \ImportTok{java}\OperatorTok{.}\ImportTok{security}\OperatorTok{.}\ImportTok{MessageDigest}\OperatorTok{;}
\KeywordTok{import} \ImportTok{java}\OperatorTok{.}\ImportTok{security}\OperatorTok{.}\ImportTok{PrivateKey}\OperatorTok{;}
\KeywordTok{import} \ImportTok{java}\OperatorTok{.}\ImportTok{security}\OperatorTok{.}\ImportTok{cert}\OperatorTok{.}\ImportTok{X509Certificate}\OperatorTok{;}
\KeywordTok{import} \ImportTok{java}\OperatorTok{.}\ImportTok{util}\OperatorTok{.}\ImportTok{ArrayList}\OperatorTok{;}
\KeywordTok{import} \ImportTok{java}\OperatorTok{.}\ImportTok{util}\OperatorTok{.}\ImportTok{List}\OperatorTok{;}

\KeywordTok{import} \ImportTok{org}\OperatorTok{.}\ImportTok{bouncycastle}\OperatorTok{.}\ImportTok{asn1}\OperatorTok{.}\ImportTok{ASN1EncodableVector}\OperatorTok{;}
\KeywordTok{import} \ImportTok{org}\OperatorTok{.}\ImportTok{bouncycastle}\OperatorTok{.}\ImportTok{asn1}\OperatorTok{.}\ImportTok{DERSet}\OperatorTok{;}
\KeywordTok{import} \ImportTok{org}\OperatorTok{.}\ImportTok{bouncycastle}\OperatorTok{.}\ImportTok{asn1}\OperatorTok{.}\ImportTok{cms}\OperatorTok{.}\ImportTok{Attribute}\OperatorTok{;}
\KeywordTok{import} \ImportTok{org}\OperatorTok{.}\ImportTok{bouncycastle}\OperatorTok{.}\ImportTok{asn1}\OperatorTok{.}\ImportTok{cms}\OperatorTok{.}\ImportTok{AttributeTable}\OperatorTok{;}
\KeywordTok{import} \ImportTok{org}\OperatorTok{.}\ImportTok{bouncycastle}\OperatorTok{.}\ImportTok{asn1}\OperatorTok{.}\ImportTok{ess}\OperatorTok{.}\ImportTok{ESSCertIDv2}\OperatorTok{;}
\KeywordTok{import} \ImportTok{org}\OperatorTok{.}\ImportTok{bouncycastle}\OperatorTok{.}\ImportTok{asn1}\OperatorTok{.}\ImportTok{ess}\OperatorTok{.}\ImportTok{SigningCertificateV2}\OperatorTok{;}
\KeywordTok{import} \ImportTok{org}\OperatorTok{.}\ImportTok{bouncycastle}\OperatorTok{.}\ImportTok{asn1}\OperatorTok{.}\ImportTok{nist}\OperatorTok{.}\ImportTok{NISTObjectIdentifiers}\OperatorTok{;}
\KeywordTok{import} \ImportTok{org}\OperatorTok{.}\ImportTok{bouncycastle}\OperatorTok{.}\ImportTok{asn1}\OperatorTok{.}\ImportTok{pkcs}\OperatorTok{.}\ImportTok{PKCSObjectIdentifiers}\OperatorTok{;}
\KeywordTok{import} \ImportTok{org}\OperatorTok{.}\ImportTok{bouncycastle}\OperatorTok{.}\ImportTok{asn1}\OperatorTok{.}\ImportTok{x509}\OperatorTok{.}\ImportTok{AlgorithmIdentifier}\OperatorTok{;}
\KeywordTok{import} \ImportTok{org}\OperatorTok{.}\ImportTok{bouncycastle}\OperatorTok{.}\ImportTok{cert}\OperatorTok{.}\ImportTok{jcajce}\OperatorTok{.}\ImportTok{JcaCertStore}\OperatorTok{;}
\KeywordTok{import} \ImportTok{org}\OperatorTok{.}\ImportTok{bouncycastle}\OperatorTok{.}\ImportTok{cms}\OperatorTok{.}\ImportTok{CMSProcessableByteArray}\OperatorTok{;}
\KeywordTok{import} \ImportTok{org}\OperatorTok{.}\ImportTok{bouncycastle}\OperatorTok{.}\ImportTok{cms}\OperatorTok{.}\ImportTok{CMSSignedData}\OperatorTok{;}
\KeywordTok{import} \ImportTok{org}\OperatorTok{.}\ImportTok{bouncycastle}\OperatorTok{.}\ImportTok{cms}\OperatorTok{.}\ImportTok{CMSSignedDataGenerator}\OperatorTok{;}
\KeywordTok{import} \ImportTok{org}\OperatorTok{.}\ImportTok{bouncycastle}\OperatorTok{.}\ImportTok{cms}\OperatorTok{.}\ImportTok{DefaultSignedAttributeTableGenerator}\OperatorTok{;}
\KeywordTok{import} \ImportTok{org}\OperatorTok{.}\ImportTok{bouncycastle}\OperatorTok{.}\ImportTok{cms}\OperatorTok{.}\ImportTok{SignerInfoGenerator}\OperatorTok{;}
\KeywordTok{import} \ImportTok{org}\OperatorTok{.}\ImportTok{bouncycastle}\OperatorTok{.}\ImportTok{cms}\OperatorTok{.}\ImportTok{jcajce}\OperatorTok{.}\ImportTok{JcaSignerInfoGeneratorBuilder}\OperatorTok{;}
\KeywordTok{import} \ImportTok{org}\OperatorTok{.}\ImportTok{bouncycastle}\OperatorTok{.}\ImportTok{operator}\OperatorTok{.}\ImportTok{ContentSigner}\OperatorTok{;}
\KeywordTok{import} \ImportTok{org}\OperatorTok{.}\ImportTok{bouncycastle}\OperatorTok{.}\ImportTok{operator}\OperatorTok{.}\ImportTok{DigestCalculatorProvider}\OperatorTok{;}
\KeywordTok{import} \ImportTok{org}\OperatorTok{.}\ImportTok{bouncycastle}\OperatorTok{.}\ImportTok{operator}\OperatorTok{.}\ImportTok{jcajce}\OperatorTok{.}\ImportTok{JcaContentSignerBuilder}\OperatorTok{;}
\KeywordTok{import} \ImportTok{org}\OperatorTok{.}\ImportTok{bouncycastle}\OperatorTok{.}\ImportTok{operator}\OperatorTok{.}\ImportTok{jcajce}\OperatorTok{.}\ImportTok{JcaDigestCalculatorProviderBuilder}\OperatorTok{;}

\CommentTok{/**}
 \CommentTok{*}\NormalTok{ Builds a PAdES}\CommentTok{{-}}\NormalTok{B}\CommentTok{{-}}\NormalTok{B compliant CMS SignedData structure}\CommentTok{.}
 \CommentTok{*}
 \CommentTok{*}\NormalTok{ Implements SignatureInterface for use with PDFBox}\CommentTok{\textquotesingle{}}\NormalTok{s}
 \CommentTok{*}\NormalTok{ PDDocument}\CommentTok{.}\NormalTok{addSignature}\CommentTok{().}
 \CommentTok{*}
 \CommentTok{*}\NormalTok{ Produces}\CommentTok{:}
 \CommentTok{*} \CommentTok{{-}}\NormalTok{ Detached CMS SignedData }\CommentTok{(}\NormalTok{SubFilter}\CommentTok{:}\NormalTok{ ETSI}\CommentTok{.}\NormalTok{CAdES}\CommentTok{.}\NormalTok{detached}\CommentTok{)}
 \CommentTok{*} \CommentTok{{-}}\NormalTok{ Signed attributes}\CommentTok{:}\NormalTok{ content}\CommentTok{{-}}\NormalTok{type}\CommentTok{,}\NormalTok{ message}\CommentTok{{-}}\NormalTok{digest}\CommentTok{,}\NormalTok{ signing}\CommentTok{{-}}\NormalTok{time}\CommentTok{,}
 \CommentTok{*}\NormalTok{   ESS}\CommentTok{{-}}\NormalTok{signing}\CommentTok{{-}}\NormalTok{certificate}\CommentTok{{-}}\NormalTok{v2}
 \CommentTok{*} \CommentTok{{-}}\NormalTok{ Certificate chain embedded in CMS certificates field}
 \CommentTok{*/}
\KeywordTok{public} \KeywordTok{class}\NormalTok{ PadesBBBuilder}
        \KeywordTok{implements}\NormalTok{ org}\OperatorTok{.}\FunctionTok{apache}\OperatorTok{.}\FunctionTok{pdfbox}\OperatorTok{.}\FunctionTok{pdmodel}\OperatorTok{.}\FunctionTok{interactive}
                \OperatorTok{.}\FunctionTok{digitalsignature}\OperatorTok{.}\FunctionTok{SignatureInterface} \OperatorTok{\{}

    \KeywordTok{private} \DataTypeTok{final} \BuiltInTok{PrivateKey}\NormalTok{ privateKey}\OperatorTok{;}
    \KeywordTok{private} \DataTypeTok{final} \BuiltInTok{X509Certificate}\NormalTok{ certificate}\OperatorTok{;}
    \KeywordTok{private} \DataTypeTok{final} \BuiltInTok{List}\OperatorTok{\textless{}}\BuiltInTok{X509Certificate}\OperatorTok{\textgreater{}}\NormalTok{ chain}\OperatorTok{;}

    \KeywordTok{public} \FunctionTok{PadesBBBuilder}\OperatorTok{(}\BuiltInTok{PrivateKey}\NormalTok{ privateKey}\OperatorTok{,}
                           \BuiltInTok{X509Certificate}\NormalTok{ certificate}\OperatorTok{,}
                           \BuiltInTok{List}\OperatorTok{\textless{}}\BuiltInTok{X509Certificate}\OperatorTok{\textgreater{}}\NormalTok{ chain}\OperatorTok{)} \OperatorTok{\{}
        \KeywordTok{this}\OperatorTok{.}\FunctionTok{privateKey} \OperatorTok{=}\NormalTok{ privateKey}\OperatorTok{;}
        \KeywordTok{this}\OperatorTok{.}\FunctionTok{certificate} \OperatorTok{=}\NormalTok{ certificate}\OperatorTok{;}
        \KeywordTok{this}\OperatorTok{.}\FunctionTok{chain} \OperatorTok{=}\NormalTok{ chain}\OperatorTok{;}
    \OperatorTok{\}}

    \AttributeTok{@Override}
    \KeywordTok{public} \DataTypeTok{byte}\OperatorTok{[]} \FunctionTok{sign}\OperatorTok{(}\BuiltInTok{InputStream}\NormalTok{ content}\OperatorTok{)} \KeywordTok{throws} \BuiltInTok{IOException} \OperatorTok{\{}
        \ControlFlowTok{try} \OperatorTok{\{}
            \DataTypeTok{byte}\OperatorTok{[]}\NormalTok{ contentBytes }\OperatorTok{=}\NormalTok{ content}\OperatorTok{.}\FunctionTok{readAllBytes}\OperatorTok{();}
            \ControlFlowTok{return} \FunctionTok{buildSignedData}\OperatorTok{(}\NormalTok{contentBytes}\OperatorTok{);}
        \OperatorTok{\}} \ControlFlowTok{catch} \OperatorTok{(}\BuiltInTok{IOException}\NormalTok{ e}\OperatorTok{)} \OperatorTok{\{}
            \ControlFlowTok{throw}\NormalTok{ e}\OperatorTok{;}
        \OperatorTok{\}} \ControlFlowTok{catch} \OperatorTok{(}\BuiltInTok{Exception}\NormalTok{ e}\OperatorTok{)} \OperatorTok{\{}
            \ControlFlowTok{throw} \KeywordTok{new} \BuiltInTok{IOException}\OperatorTok{(}\StringTok{"PAdES{-}B{-}B signing failed: "}
                    \OperatorTok{+}\NormalTok{ e}\OperatorTok{.}\FunctionTok{getMessage}\OperatorTok{(),}\NormalTok{ e}\OperatorTok{);}
        \OperatorTok{\}}
    \OperatorTok{\}}

    \CommentTok{/**}
     \CommentTok{*}\NormalTok{ Builds the CMS SignedData from the byte range content}\CommentTok{.}
     \CommentTok{*}\NormalTok{ This method can also be called directly for testing}\CommentTok{.}
     \CommentTok{*/}
    \KeywordTok{public} \DataTypeTok{byte}\OperatorTok{[]} \FunctionTok{buildSignedData}\OperatorTok{(}\DataTypeTok{byte}\OperatorTok{[]}\NormalTok{ contentBytes}\OperatorTok{)} \KeywordTok{throws} \BuiltInTok{Exception} \OperatorTok{\{}

        \CommentTok{// 1. Determine signature algorithm}
        \BuiltInTok{String}\NormalTok{ sigAlgo }\OperatorTok{=} \FunctionTok{resolveAlgorithm}\OperatorTok{(}\NormalTok{privateKey}\OperatorTok{);}

        \CommentTok{// 2. Content signer}
\NormalTok{        ContentSigner signer }\OperatorTok{=} \KeywordTok{new} \FunctionTok{JcaContentSignerBuilder}\OperatorTok{(}\NormalTok{sigAlgo}\OperatorTok{)}
                \OperatorTok{.}\FunctionTok{setProvider}\OperatorTok{(}\StringTok{"BC"}\OperatorTok{)}
                \OperatorTok{.}\FunctionTok{build}\OperatorTok{(}\NormalTok{privateKey}\OperatorTok{);}

        \CommentTok{// 3. Digest calculator provider}
\NormalTok{        DigestCalculatorProvider dcProvider }\OperatorTok{=}
                \KeywordTok{new} \FunctionTok{JcaDigestCalculatorProviderBuilder}\OperatorTok{()}
                        \OperatorTok{.}\FunctionTok{setProvider}\OperatorTok{(}\StringTok{"BC"}\OperatorTok{).}\FunctionTok{build}\OperatorTok{();}

        \CommentTok{// 4. ESS{-}signing{-}certificate{-}v2 attribute}
\NormalTok{        AttributeTable essAttrs }\OperatorTok{=} \FunctionTok{buildEssAttribute}\OperatorTok{(}\NormalTok{certificate}\OperatorTok{);}

        \CommentTok{// 5. Signer info generator with custom signed attributes}
\NormalTok{        JcaSignerInfoGeneratorBuilder signerBuilder }\OperatorTok{=}
                \KeywordTok{new} \FunctionTok{JcaSignerInfoGeneratorBuilder}\OperatorTok{(}\NormalTok{dcProvider}\OperatorTok{);}
\NormalTok{        signerBuilder}\OperatorTok{.}\FunctionTok{setSignedAttributeGenerator}\OperatorTok{(}
                \KeywordTok{new} \FunctionTok{DefaultSignedAttributeTableGenerator}\OperatorTok{(}\NormalTok{essAttrs}\OperatorTok{));}

\NormalTok{        SignerInfoGenerator signerInfoGen }\OperatorTok{=}
\NormalTok{                signerBuilder}\OperatorTok{.}\FunctionTok{build}\OperatorTok{(}\NormalTok{signer}\OperatorTok{,}\NormalTok{ certificate}\OperatorTok{);}

        \CommentTok{// 6. Certificate store}
        \BuiltInTok{List}\OperatorTok{\textless{}}\BuiltInTok{X509Certificate}\OperatorTok{\textgreater{}}\NormalTok{ certs }\OperatorTok{=} \KeywordTok{new} \BuiltInTok{ArrayList}\OperatorTok{\textless{}\textgreater{}(}\NormalTok{chain}\OperatorTok{);}
\NormalTok{        JcaCertStore certStore }\OperatorTok{=} \KeywordTok{new} \FunctionTok{JcaCertStore}\OperatorTok{(}\NormalTok{certs}\OperatorTok{);}

        \CommentTok{// 7. Assemble and generate}
\NormalTok{        CMSSignedDataGenerator gen }\OperatorTok{=} \KeywordTok{new} \FunctionTok{CMSSignedDataGenerator}\OperatorTok{();}
\NormalTok{        gen}\OperatorTok{.}\FunctionTok{addSignerInfoGenerator}\OperatorTok{(}\NormalTok{signerInfoGen}\OperatorTok{);}
\NormalTok{        gen}\OperatorTok{.}\FunctionTok{addCertificates}\OperatorTok{(}\NormalTok{certStore}\OperatorTok{);}

\NormalTok{        CMSProcessableByteArray cmsContent }\OperatorTok{=}
                \KeywordTok{new} \FunctionTok{CMSProcessableByteArray}\OperatorTok{(}\NormalTok{contentBytes}\OperatorTok{);}
\NormalTok{        CMSSignedData signedData }\OperatorTok{=}\NormalTok{ gen}\OperatorTok{.}\FunctionTok{generate}\OperatorTok{(}\NormalTok{cmsContent}\OperatorTok{,} \KeywordTok{false}\OperatorTok{);}

        \ControlFlowTok{return}\NormalTok{ signedData}\OperatorTok{.}\FunctionTok{getEncoded}\OperatorTok{();}
    \OperatorTok{\}}

    \KeywordTok{private} \DataTypeTok{static}\NormalTok{ AttributeTable }\FunctionTok{buildEssAttribute}\OperatorTok{(}
            \BuiltInTok{X509Certificate}\NormalTok{ cert}\OperatorTok{)} \KeywordTok{throws} \BuiltInTok{Exception} \OperatorTok{\{}
        \BuiltInTok{MessageDigest}\NormalTok{ md }\OperatorTok{=} \BuiltInTok{MessageDigest}\OperatorTok{.}\FunctionTok{getInstance}\OperatorTok{(}\StringTok{"SHA{-}256"}\OperatorTok{);}
        \DataTypeTok{byte}\OperatorTok{[]}\NormalTok{ certHash }\OperatorTok{=}\NormalTok{ md}\OperatorTok{.}\FunctionTok{digest}\OperatorTok{(}\NormalTok{cert}\OperatorTok{.}\FunctionTok{getEncoded}\OperatorTok{());}

\NormalTok{        AlgorithmIdentifier hashAlgo }\OperatorTok{=} \KeywordTok{new} \FunctionTok{AlgorithmIdentifier}\OperatorTok{(}
\NormalTok{                NISTObjectIdentifiers}\OperatorTok{.}\FunctionTok{id\_sha256}\OperatorTok{);}
\NormalTok{        ESSCertIDv2 essCertId }\OperatorTok{=} \KeywordTok{new} \FunctionTok{ESSCertIDv2}\OperatorTok{(}\NormalTok{hashAlgo}\OperatorTok{,}\NormalTok{ certHash}\OperatorTok{);}
\NormalTok{        SigningCertificateV2 sigCertV2 }\OperatorTok{=} \KeywordTok{new} \FunctionTok{SigningCertificateV2}\OperatorTok{(}
                \KeywordTok{new}\NormalTok{ ESSCertIDv2}\OperatorTok{[]\{}\NormalTok{essCertId}\OperatorTok{\});}

\NormalTok{        ASN1EncodableVector v }\OperatorTok{=} \KeywordTok{new} \FunctionTok{ASN1EncodableVector}\OperatorTok{();}
\NormalTok{        v}\OperatorTok{.}\FunctionTok{add}\OperatorTok{(}\KeywordTok{new} \BuiltInTok{Attribute}\OperatorTok{(}
\NormalTok{                PKCSObjectIdentifiers}\OperatorTok{.}\FunctionTok{id\_aa\_signingCertificateV2}\OperatorTok{,}
                \KeywordTok{new} \FunctionTok{DERSet}\OperatorTok{(}\NormalTok{sigCertV2}\OperatorTok{)));}
        \ControlFlowTok{return} \KeywordTok{new} \FunctionTok{AttributeTable}\OperatorTok{(}\NormalTok{v}\OperatorTok{);}
    \OperatorTok{\}}

    \KeywordTok{private} \DataTypeTok{static} \BuiltInTok{String} \FunctionTok{resolveAlgorithm}\OperatorTok{(}\BuiltInTok{PrivateKey}\NormalTok{ key}\OperatorTok{)} \OperatorTok{\{}
        \BuiltInTok{String}\NormalTok{ algo }\OperatorTok{=}\NormalTok{ key}\OperatorTok{.}\FunctionTok{getAlgorithm}\OperatorTok{();}
        \ControlFlowTok{return} \ControlFlowTok{switch} \OperatorTok{(}\NormalTok{algo}\OperatorTok{)} \OperatorTok{\{}
            \ControlFlowTok{case} \StringTok{"RSA"} \OperatorTok{{-}\textgreater{}} \StringTok{"SHA256withRSA"}\OperatorTok{;}
            \ControlFlowTok{case} \StringTok{"RSASSA{-}PSS"} \OperatorTok{{-}\textgreater{}} \StringTok{"SHA256withRSAandMGF1"}\OperatorTok{;}
            \ControlFlowTok{case} \StringTok{"EC"}\OperatorTok{,} \StringTok{"ECDSA"} \OperatorTok{{-}\textgreater{}} \StringTok{"SHA256withECDSA"}\OperatorTok{;}
            \ControlFlowTok{case} \StringTok{"Ed25519"} \OperatorTok{{-}\textgreater{}} \StringTok{"Ed25519"}\OperatorTok{;}
            \ControlFlowTok{case} \StringTok{"Ed448"} \OperatorTok{{-}\textgreater{}} \StringTok{"Ed448"}\OperatorTok{;}
            \KeywordTok{default} \OperatorTok{{-}\textgreater{}} \ControlFlowTok{throw} \KeywordTok{new} \BuiltInTok{IllegalArgumentException}\OperatorTok{(}
                    \StringTok{"Unsupported key algorithm: "} \OperatorTok{+}\NormalTok{ algo}\OperatorTok{);}
        \OperatorTok{\};}
    \OperatorTok{\}}
\OperatorTok{\}}
\end{Highlighting}
\end{Shaded}

\section{Complete PAdES-B-T Builder}\label{_complete_pades_b_t_builder}

The PAdES-B-T builder extends B-B by adding a timestamp after the
signature is generated:

\begin{Shaded}
\begin{Highlighting}[]
\KeywordTok{import} \ImportTok{java}\OperatorTok{.}\ImportTok{io}\OperatorTok{.}\ImportTok{IOException}\OperatorTok{;}
\KeywordTok{import} \ImportTok{java}\OperatorTok{.}\ImportTok{io}\OperatorTok{.}\ImportTok{InputStream}\OperatorTok{;}
\KeywordTok{import} \ImportTok{java}\OperatorTok{.}\ImportTok{security}\OperatorTok{.}\ImportTok{MessageDigest}\OperatorTok{;}
\KeywordTok{import} \ImportTok{java}\OperatorTok{.}\ImportTok{security}\OperatorTok{.}\ImportTok{PrivateKey}\OperatorTok{;}
\KeywordTok{import} \ImportTok{java}\OperatorTok{.}\ImportTok{security}\OperatorTok{.}\ImportTok{cert}\OperatorTok{.}\ImportTok{X509Certificate}\OperatorTok{;}
\KeywordTok{import} \ImportTok{java}\OperatorTok{.}\ImportTok{util}\OperatorTok{.}\ImportTok{Collections}\OperatorTok{;}
\KeywordTok{import} \ImportTok{java}\OperatorTok{.}\ImportTok{util}\OperatorTok{.}\ImportTok{List}\OperatorTok{;}

\KeywordTok{import} \ImportTok{org}\OperatorTok{.}\ImportTok{bouncycastle}\OperatorTok{.}\ImportTok{asn1}\OperatorTok{.}\ImportTok{ASN1EncodableVector}\OperatorTok{;}
\KeywordTok{import} \ImportTok{org}\OperatorTok{.}\ImportTok{bouncycastle}\OperatorTok{.}\ImportTok{asn1}\OperatorTok{.}\ImportTok{ASN1Primitive}\OperatorTok{;}
\KeywordTok{import} \ImportTok{org}\OperatorTok{.}\ImportTok{bouncycastle}\OperatorTok{.}\ImportTok{asn1}\OperatorTok{.}\ImportTok{DERSet}\OperatorTok{;}
\KeywordTok{import} \ImportTok{org}\OperatorTok{.}\ImportTok{bouncycastle}\OperatorTok{.}\ImportTok{asn1}\OperatorTok{.}\ImportTok{cms}\OperatorTok{.}\ImportTok{Attribute}\OperatorTok{;}
\KeywordTok{import} \ImportTok{org}\OperatorTok{.}\ImportTok{bouncycastle}\OperatorTok{.}\ImportTok{asn1}\OperatorTok{.}\ImportTok{cms}\OperatorTok{.}\ImportTok{AttributeTable}\OperatorTok{;}
\KeywordTok{import} \ImportTok{org}\OperatorTok{.}\ImportTok{bouncycastle}\OperatorTok{.}\ImportTok{asn1}\OperatorTok{.}\ImportTok{nist}\OperatorTok{.}\ImportTok{NISTObjectIdentifiers}\OperatorTok{;}
\KeywordTok{import} \ImportTok{org}\OperatorTok{.}\ImportTok{bouncycastle}\OperatorTok{.}\ImportTok{asn1}\OperatorTok{.}\ImportTok{pkcs}\OperatorTok{.}\ImportTok{PKCSObjectIdentifiers}\OperatorTok{;}
\KeywordTok{import} \ImportTok{org}\OperatorTok{.}\ImportTok{bouncycastle}\OperatorTok{.}\ImportTok{cms}\OperatorTok{.}\ImportTok{CMSSignedData}\OperatorTok{;}
\KeywordTok{import} \ImportTok{org}\OperatorTok{.}\ImportTok{bouncycastle}\OperatorTok{.}\ImportTok{cms}\OperatorTok{.}\ImportTok{SignerInformation}\OperatorTok{;}
\KeywordTok{import} \ImportTok{org}\OperatorTok{.}\ImportTok{bouncycastle}\OperatorTok{.}\ImportTok{cms}\OperatorTok{.}\ImportTok{SignerInformationStore}\OperatorTok{;}
\KeywordTok{import} \ImportTok{org}\OperatorTok{.}\ImportTok{bouncycastle}\OperatorTok{.}\ImportTok{tsp}\OperatorTok{.}\ImportTok{TimeStampRequest}\OperatorTok{;}
\KeywordTok{import} \ImportTok{org}\OperatorTok{.}\ImportTok{bouncycastle}\OperatorTok{.}\ImportTok{tsp}\OperatorTok{.}\ImportTok{TimeStampRequestGenerator}\OperatorTok{;}
\KeywordTok{import} \ImportTok{org}\OperatorTok{.}\ImportTok{bouncycastle}\OperatorTok{.}\ImportTok{tsp}\OperatorTok{.}\ImportTok{TimeStampResponse}\OperatorTok{;}
\KeywordTok{import} \ImportTok{org}\OperatorTok{.}\ImportTok{bouncycastle}\OperatorTok{.}\ImportTok{tsp}\OperatorTok{.}\ImportTok{TimeStampToken}\OperatorTok{;}

\CommentTok{/**}
 \CommentTok{*}\NormalTok{ Builds a PAdES}\CommentTok{{-}}\NormalTok{B}\CommentTok{{-}}\NormalTok{T compliant CMS SignedData structure}\CommentTok{.}
 \CommentTok{*}
 \CommentTok{*}\NormalTok{ Extends PAdES}\CommentTok{{-}}\NormalTok{B}\CommentTok{{-}}\NormalTok{B by adding an RFC }\CommentTok{3161}\NormalTok{ timestamp token}
 \CommentTok{*}\NormalTok{ as an unsigned attribute on the SignerInfo}\CommentTok{.}
 \CommentTok{*}
 \CommentTok{*}\NormalTok{ Usage with PDFBox}\CommentTok{:}
 \CommentTok{*} \DataTypeTok{\textless{}}\KeywordTok{pre}\DataTypeTok{\textgreater{}}
 \CommentTok{*}\NormalTok{ PadesBTBuilder builder }\CommentTok{=}\NormalTok{ new PadesBTBuilder}\CommentTok{(}
 \CommentTok{*}\NormalTok{     privateKey}\CommentTok{,}\NormalTok{ certificate}\CommentTok{,}\NormalTok{ chain}\CommentTok{,} \CommentTok{"}\NormalTok{http}\CommentTok{://}\NormalTok{tsa}\CommentTok{.}\NormalTok{example}\CommentTok{.}\NormalTok{com}\CommentTok{");}
 \CommentTok{*}\NormalTok{ doc}\CommentTok{.}\NormalTok{addSignature}\CommentTok{(}\NormalTok{signature}\CommentTok{,}\NormalTok{ builder}\CommentTok{,}\NormalTok{ options}\CommentTok{);}
 \CommentTok{*}\NormalTok{ doc}\CommentTok{.}\NormalTok{saveIncremental}\CommentTok{(}\NormalTok{outputStream}\CommentTok{);}
 \CommentTok{*} \DataTypeTok{\textless{}/}\KeywordTok{pre}\DataTypeTok{\textgreater{}}
 \CommentTok{*/}
\KeywordTok{public} \KeywordTok{class}\NormalTok{ PadesBTBuilder}
        \KeywordTok{implements}\NormalTok{ org}\OperatorTok{.}\FunctionTok{apache}\OperatorTok{.}\FunctionTok{pdfbox}\OperatorTok{.}\FunctionTok{pdmodel}\OperatorTok{.}\FunctionTok{interactive}
                \OperatorTok{.}\FunctionTok{digitalsignature}\OperatorTok{.}\FunctionTok{SignatureInterface} \OperatorTok{\{}

    \KeywordTok{private} \DataTypeTok{final}\NormalTok{ PadesBBBuilder bbBuilder}\OperatorTok{;}
    \KeywordTok{private} \DataTypeTok{final} \BuiltInTok{String}\NormalTok{ tsaUrl}\OperatorTok{;}

    \KeywordTok{public} \FunctionTok{PadesBTBuilder}\OperatorTok{(}\BuiltInTok{PrivateKey}\NormalTok{ privateKey}\OperatorTok{,}
                           \BuiltInTok{X509Certificate}\NormalTok{ certificate}\OperatorTok{,}
                           \BuiltInTok{List}\OperatorTok{\textless{}}\BuiltInTok{X509Certificate}\OperatorTok{\textgreater{}}\NormalTok{ chain}\OperatorTok{,}
                           \BuiltInTok{String}\NormalTok{ tsaUrl}\OperatorTok{)} \OperatorTok{\{}
        \KeywordTok{this}\OperatorTok{.}\FunctionTok{bbBuilder} \OperatorTok{=} \KeywordTok{new} \FunctionTok{PadesBBBuilder}\OperatorTok{(}
\NormalTok{                privateKey}\OperatorTok{,}\NormalTok{ certificate}\OperatorTok{,}\NormalTok{ chain}\OperatorTok{);}
        \KeywordTok{this}\OperatorTok{.}\FunctionTok{tsaUrl} \OperatorTok{=}\NormalTok{ tsaUrl}\OperatorTok{;}
    \OperatorTok{\}}

    \AttributeTok{@Override}
    \KeywordTok{public} \DataTypeTok{byte}\OperatorTok{[]} \FunctionTok{sign}\OperatorTok{(}\BuiltInTok{InputStream}\NormalTok{ content}\OperatorTok{)} \KeywordTok{throws} \BuiltInTok{IOException} \OperatorTok{\{}
        \ControlFlowTok{try} \OperatorTok{\{}
            \DataTypeTok{byte}\OperatorTok{[]}\NormalTok{ contentBytes }\OperatorTok{=}\NormalTok{ content}\OperatorTok{.}\FunctionTok{readAllBytes}\OperatorTok{();}

            \CommentTok{// Step 1: Build PAdES{-}B{-}B SignedData}
            \DataTypeTok{byte}\OperatorTok{[]}\NormalTok{ bbBytes }\OperatorTok{=}\NormalTok{ bbBuilder}\OperatorTok{.}\FunctionTok{buildSignedData}\OperatorTok{(}\NormalTok{contentBytes}\OperatorTok{);}
\NormalTok{            CMSSignedData signedData }\OperatorTok{=} \KeywordTok{new} \FunctionTok{CMSSignedData}\OperatorTok{(}\NormalTok{bbBytes}\OperatorTok{);}

            \CommentTok{// Step 2: Add timestamp}
\NormalTok{            CMSSignedData timestamped }\OperatorTok{=}
                    \FunctionTok{embedTimestamp}\OperatorTok{(}\NormalTok{signedData}\OperatorTok{,}\NormalTok{ tsaUrl}\OperatorTok{);}

            \ControlFlowTok{return}\NormalTok{ timestamped}\OperatorTok{.}\FunctionTok{getEncoded}\OperatorTok{();}
        \OperatorTok{\}} \ControlFlowTok{catch} \OperatorTok{(}\BuiltInTok{IOException}\NormalTok{ e}\OperatorTok{)} \OperatorTok{\{}
            \ControlFlowTok{throw}\NormalTok{ e}\OperatorTok{;}
        \OperatorTok{\}} \ControlFlowTok{catch} \OperatorTok{(}\BuiltInTok{Exception}\NormalTok{ e}\OperatorTok{)} \OperatorTok{\{}
            \ControlFlowTok{throw} \KeywordTok{new} \BuiltInTok{IOException}\OperatorTok{(}\StringTok{"PAdES{-}B{-}T signing failed: "}
                    \OperatorTok{+}\NormalTok{ e}\OperatorTok{.}\FunctionTok{getMessage}\OperatorTok{(),}\NormalTok{ e}\OperatorTok{);}
        \OperatorTok{\}}
    \OperatorTok{\}}

    \KeywordTok{private} \DataTypeTok{static}\NormalTok{ CMSSignedData }\FunctionTok{embedTimestamp}\OperatorTok{(}
\NormalTok{            CMSSignedData signedData}\OperatorTok{,} \BuiltInTok{String}\NormalTok{ tsaUrl}\OperatorTok{)}
            \KeywordTok{throws} \BuiltInTok{Exception} \OperatorTok{\{}

        \CommentTok{// Extract signature value}
\NormalTok{        SignerInformation signerInfo }\OperatorTok{=}\NormalTok{ signedData}\OperatorTok{.}\FunctionTok{getSignerInfos}\OperatorTok{()}
                \OperatorTok{.}\FunctionTok{getSigners}\OperatorTok{().}\FunctionTok{iterator}\OperatorTok{().}\FunctionTok{next}\OperatorTok{();}
        \DataTypeTok{byte}\OperatorTok{[]}\NormalTok{ sigBytes }\OperatorTok{=}\NormalTok{ signerInfo}\OperatorTok{.}\FunctionTok{getSignature}\OperatorTok{();}

        \CommentTok{// Hash for TSA request}
        \BuiltInTok{MessageDigest}\NormalTok{ md }\OperatorTok{=} \BuiltInTok{MessageDigest}\OperatorTok{.}\FunctionTok{getInstance}\OperatorTok{(}\StringTok{"SHA{-}256"}\OperatorTok{);}
        \DataTypeTok{byte}\OperatorTok{[]}\NormalTok{ hash }\OperatorTok{=}\NormalTok{ md}\OperatorTok{.}\FunctionTok{digest}\OperatorTok{(}\NormalTok{sigBytes}\OperatorTok{);}

        \CommentTok{// Build RFC 3161 request}
\NormalTok{        TimeStampRequestGenerator reqGen }\OperatorTok{=}
                \KeywordTok{new} \FunctionTok{TimeStampRequestGenerator}\OperatorTok{();}
\NormalTok{        reqGen}\OperatorTok{.}\FunctionTok{setCertReq}\OperatorTok{(}\KeywordTok{true}\OperatorTok{);}
\NormalTok{        TimeStampRequest tsReq }\OperatorTok{=}\NormalTok{ reqGen}\OperatorTok{.}\FunctionTok{generate}\OperatorTok{(}
\NormalTok{                NISTObjectIdentifiers}\OperatorTok{.}\FunctionTok{id\_sha256}\OperatorTok{,}\NormalTok{ hash}\OperatorTok{);}

        \CommentTok{// Send to TSA}
        \DataTypeTok{byte}\OperatorTok{[]}\NormalTok{ tsRespBytes }\OperatorTok{=} \FunctionTok{sendTsaRequest}\OperatorTok{(}\NormalTok{tsaUrl}\OperatorTok{,}\NormalTok{ tsReq}\OperatorTok{.}\FunctionTok{getEncoded}\OperatorTok{());}

        \CommentTok{// Parse and validate}
\NormalTok{        TimeStampResponse tsResp }\OperatorTok{=} \KeywordTok{new} \FunctionTok{TimeStampResponse}\OperatorTok{(}\NormalTok{tsRespBytes}\OperatorTok{);}
\NormalTok{        tsResp}\OperatorTok{.}\FunctionTok{validate}\OperatorTok{(}\NormalTok{tsReq}\OperatorTok{);}

        \DataTypeTok{int}\NormalTok{ status }\OperatorTok{=}\NormalTok{ tsResp}\OperatorTok{.}\FunctionTok{getStatus}\OperatorTok{();}
        \ControlFlowTok{if} \OperatorTok{(}\NormalTok{status }\OperatorTok{!=} \DecValTok{0} \OperatorTok{\&\&}\NormalTok{ status }\OperatorTok{!=} \DecValTok{1}\OperatorTok{)} \OperatorTok{\{} 
            \BuiltInTok{String}\NormalTok{ msg }\OperatorTok{=}\NormalTok{ tsResp}\OperatorTok{.}\FunctionTok{getStatusString}\OperatorTok{();}
            \ControlFlowTok{throw} \KeywordTok{new} \BuiltInTok{IOException}\OperatorTok{(}\StringTok{"TSA error: "}
                    \OperatorTok{+} \OperatorTok{(}\NormalTok{msg }\OperatorTok{!=} \KeywordTok{null} \OperatorTok{?}\NormalTok{ msg }\OperatorTok{:} \StringTok{"status="} \OperatorTok{+}\NormalTok{ status}\OperatorTok{));}
        \OperatorTok{\}}

\NormalTok{        TimeStampToken tsToken }\OperatorTok{=}\NormalTok{ tsResp}\OperatorTok{.}\FunctionTok{getTimeStampToken}\OperatorTok{();}
        \ControlFlowTok{if} \OperatorTok{(}\NormalTok{tsToken }\OperatorTok{==} \KeywordTok{null}\OperatorTok{)} \OperatorTok{\{}
            \ControlFlowTok{throw} \KeywordTok{new} \BuiltInTok{IOException}\OperatorTok{(}\StringTok{"TSA returned no timestamp token"}\OperatorTok{);}
        \OperatorTok{\}}

        \CommentTok{// Embed as unsigned attribute}
\NormalTok{        ASN1Primitive tsAsn1 }\OperatorTok{=}\NormalTok{ ASN1Primitive}\OperatorTok{.}\FunctionTok{fromByteArray}\OperatorTok{(}
\NormalTok{                tsToken}\OperatorTok{.}\FunctionTok{getEncoded}\OperatorTok{());}

\NormalTok{        ASN1EncodableVector vec }\OperatorTok{=} \KeywordTok{new} \FunctionTok{ASN1EncodableVector}\OperatorTok{();}
\NormalTok{        vec}\OperatorTok{.}\FunctionTok{add}\OperatorTok{(}\KeywordTok{new} \BuiltInTok{Attribute}\OperatorTok{(}
\NormalTok{                PKCSObjectIdentifiers}\OperatorTok{.}\FunctionTok{id\_aa\_signatureTimeStampToken}\OperatorTok{,}
                \KeywordTok{new} \FunctionTok{DERSet}\OperatorTok{(}\NormalTok{tsAsn1}\OperatorTok{)));}
\NormalTok{        AttributeTable unsignedAttrs }\OperatorTok{=} \KeywordTok{new} \FunctionTok{AttributeTable}\OperatorTok{(}\NormalTok{vec}\OperatorTok{);}

\NormalTok{        SignerInformation newSignerInfo }\OperatorTok{=}
\NormalTok{                SignerInformation}\OperatorTok{.}\FunctionTok{replaceUnsignedAttributes}\OperatorTok{(}
\NormalTok{                        signerInfo}\OperatorTok{,}\NormalTok{ unsignedAttrs}\OperatorTok{);}

        \ControlFlowTok{return}\NormalTok{ CMSSignedData}\OperatorTok{.}\FunctionTok{replaceSigners}\OperatorTok{(}\NormalTok{signedData}\OperatorTok{,}
                \KeywordTok{new} \FunctionTok{SignerInformationStore}\OperatorTok{(}
                        \BuiltInTok{Collections}\OperatorTok{.}\FunctionTok{singletonList}\OperatorTok{(}\NormalTok{newSignerInfo}\OperatorTok{)));}
    \OperatorTok{\}}

    \KeywordTok{private} \DataTypeTok{static} \DataTypeTok{byte}\OperatorTok{[]} \FunctionTok{sendTsaRequest}\OperatorTok{(}\BuiltInTok{String}\NormalTok{ tsaUrl}\OperatorTok{,} \DataTypeTok{byte}\OperatorTok{[]}\NormalTok{ reqBytes}\OperatorTok{)}
            \KeywordTok{throws} \BuiltInTok{Exception} \OperatorTok{\{}

\NormalTok{        java}\OperatorTok{.}\FunctionTok{net}\OperatorTok{.}\FunctionTok{http}\OperatorTok{.}\FunctionTok{HttpClient}\NormalTok{ client }\OperatorTok{=}
\NormalTok{                java}\OperatorTok{.}\FunctionTok{net}\OperatorTok{.}\FunctionTok{http}\OperatorTok{.}\FunctionTok{HttpClient}\OperatorTok{.}\FunctionTok{newBuilder}\OperatorTok{()}
                        \OperatorTok{.}\FunctionTok{connectTimeout}\OperatorTok{(}\NormalTok{java}\OperatorTok{.}\FunctionTok{time}\OperatorTok{.}\FunctionTok{Duration}\OperatorTok{.}\FunctionTok{ofSeconds}\OperatorTok{(}\DecValTok{15}\OperatorTok{))}
                        \OperatorTok{.}\FunctionTok{followRedirects}\OperatorTok{(}
\NormalTok{                                java}\OperatorTok{.}\FunctionTok{net}\OperatorTok{.}\FunctionTok{http}\OperatorTok{.}\FunctionTok{HttpClient}\OperatorTok{.}\FunctionTok{Redirect}\OperatorTok{.}\FunctionTok{NORMAL}\OperatorTok{)}
                        \OperatorTok{.}\FunctionTok{build}\OperatorTok{();}

\NormalTok{        java}\OperatorTok{.}\FunctionTok{net}\OperatorTok{.}\FunctionTok{http}\OperatorTok{.}\FunctionTok{HttpRequest}\NormalTok{ request }\OperatorTok{=}
\NormalTok{                java}\OperatorTok{.}\FunctionTok{net}\OperatorTok{.}\FunctionTok{http}\OperatorTok{.}\FunctionTok{HttpRequest}\OperatorTok{.}\FunctionTok{newBuilder}\OperatorTok{()}
                        \OperatorTok{.}\FunctionTok{uri}\OperatorTok{(}\NormalTok{java}\OperatorTok{.}\FunctionTok{net}\OperatorTok{.}\FunctionTok{URI}\OperatorTok{.}\FunctionTok{create}\OperatorTok{(}\NormalTok{tsaUrl}\OperatorTok{))}
                        \OperatorTok{.}\FunctionTok{header}\OperatorTok{(}\StringTok{"Content{-}Type"}\OperatorTok{,}
                                \StringTok{"application/timestamp{-}query"}\OperatorTok{)}
                        \OperatorTok{.}\FunctionTok{timeout}\OperatorTok{(}\NormalTok{java}\OperatorTok{.}\FunctionTok{time}\OperatorTok{.}\FunctionTok{Duration}\OperatorTok{.}\FunctionTok{ofSeconds}\OperatorTok{(}\DecValTok{30}\OperatorTok{))}
                        \OperatorTok{.}\FunctionTok{POST}\OperatorTok{(}\NormalTok{java}\OperatorTok{.}\FunctionTok{net}\OperatorTok{.}\FunctionTok{http}\OperatorTok{.}\FunctionTok{HttpRequest}\OperatorTok{.}\FunctionTok{BodyPublishers}
                                \OperatorTok{.}\FunctionTok{ofByteArray}\OperatorTok{(}\NormalTok{reqBytes}\OperatorTok{))}
                        \OperatorTok{.}\FunctionTok{build}\OperatorTok{();}

\NormalTok{        java}\OperatorTok{.}\FunctionTok{net}\OperatorTok{.}\FunctionTok{http}\OperatorTok{.}\FunctionTok{HttpResponse}\OperatorTok{\textless{}}\DataTypeTok{byte}\OperatorTok{[]\textgreater{}}\NormalTok{ response }\OperatorTok{=}\NormalTok{ client}\OperatorTok{.}\FunctionTok{send}\OperatorTok{(}
\NormalTok{                request}\OperatorTok{,}
\NormalTok{                java}\OperatorTok{.}\FunctionTok{net}\OperatorTok{.}\FunctionTok{http}\OperatorTok{.}\FunctionTok{HttpResponse}\OperatorTok{.}\FunctionTok{BodyHandlers}\OperatorTok{.}\FunctionTok{ofByteArray}\OperatorTok{());}

        \ControlFlowTok{if} \OperatorTok{(}\NormalTok{response}\OperatorTok{.}\FunctionTok{statusCode}\OperatorTok{()} \OperatorTok{!=} \DecValTok{200}\OperatorTok{)} \OperatorTok{\{}
            \ControlFlowTok{throw} \KeywordTok{new} \BuiltInTok{IOException}\OperatorTok{(}\StringTok{"TSA HTTP "} \OperatorTok{+}\NormalTok{ response}\OperatorTok{.}\FunctionTok{statusCode}\OperatorTok{()}
                    \OperatorTok{+} \StringTok{" from "} \OperatorTok{+}\NormalTok{ tsaUrl}\OperatorTok{);}
        \OperatorTok{\}}

        \DataTypeTok{byte}\OperatorTok{[]}\NormalTok{ body }\OperatorTok{=}\NormalTok{ response}\OperatorTok{.}\FunctionTok{body}\OperatorTok{();}
        \ControlFlowTok{if} \OperatorTok{(}\NormalTok{body }\OperatorTok{==} \KeywordTok{null} \OperatorTok{||}\NormalTok{ body}\OperatorTok{.}\FunctionTok{length} \OperatorTok{==} \DecValTok{0}\OperatorTok{)} \OperatorTok{\{}
            \ControlFlowTok{throw} \KeywordTok{new} \BuiltInTok{IOException}\OperatorTok{(}\StringTok{"TSA empty response from "} \OperatorTok{+}\NormalTok{ tsaUrl}\OperatorTok{);}
        \OperatorTok{\}}

        \ControlFlowTok{return}\NormalTok{ body}\OperatorTok{;}
    \OperatorTok{\}}
\OperatorTok{\}}
\end{Highlighting}
\end{Shaded}

\begin{itemize}
\item
  TSA status 0 = \texttt{granted}, 1 = \texttt{grantedWithMods}. Any
  other status (2 = rejection, 3 = waiting, 4 = revocationWarning, 5 =
  revocationNotification) means the TSA refused to issue a timestamp.
  Handle this gracefully --- some TSAs reject requests during
  maintenance windows.
\end{itemize}

\section{CMS Construction for Remote Signing (CSC
API)}\label{_cms_construction_for_remote_signing_csc_api}

When using a remote signing service via the CSC API (Chapter 22), you
cannot use BouncyCastle's \texttt{CMSSignedDataGenerator} in the normal
way --- because you do not have the private key locally. The generator
expects a \texttt{ContentSigner} that performs the signing operation,
but with CSC, the signing happens on a remote HSM. You must build the
CMS \texttt{SignedData} structure \emph{manually} using BouncyCastle's
low-level ASN.1 classes.

\subsection{The Challenge}\label{_the_challenge}

The remote signing flow splits the CMS construction into three phases:

\begin{enumerate}
\def\labelenumi{\arabic{enumi}.}
\item
  \textbf{Build signed attributes} (local) --- content-type,
  message-digest, ESS-signing-certificate-v2
\item
  \textbf{Hash and send to CSC} (remote) --- the server signs the
  DER-encoded signed attributes hash
\item
  \textbf{Assemble CMS} (local) --- combine the signed attributes, the
  remote signature value, and the certificate chain into a valid
  \texttt{SignedData} structure
\end{enumerate}

Phase 2 returns raw signature bytes --- not a CMS structure, not a
PKCS\#7 blob, just the naked RSA or ECDSA signature value. You must wrap
those bytes into a complete CMS \texttt{SignedData} yourself.

\subsection{Manual CMS Assembly}\label{_manual_cms_assembly}

\begin{Shaded}
\begin{Highlighting}[]
\CommentTok{/**}
 \CommentTok{*}\NormalTok{ Assembles a CMS SignedData structure from pre}\CommentTok{{-}}\NormalTok{computed components}\CommentTok{.}
 \CommentTok{*}\NormalTok{ Used when the signature value comes from a remote signing service}
 \CommentTok{*} \CommentTok{(}\NormalTok{CSC API}\CommentTok{)}\NormalTok{ and the CMS must be constructed locally}\CommentTok{.}
 \CommentTok{*/}
\KeywordTok{public} \DataTypeTok{byte}\OperatorTok{[]} \FunctionTok{assembleCmsForRemoteSigning}\OperatorTok{(}
        \DataTypeTok{byte}\OperatorTok{[]}\NormalTok{ contentHash}\OperatorTok{,}
        \DataTypeTok{byte}\OperatorTok{[]}\NormalTok{ remoteSignatureValue}\OperatorTok{,}
        \BuiltInTok{X509Certificate}\NormalTok{ signerCert}\OperatorTok{,}
        \BuiltInTok{List}\OperatorTok{\textless{}}\BuiltInTok{X509Certificate}\OperatorTok{\textgreater{}}\NormalTok{ certChain}\OperatorTok{,}
        \BuiltInTok{String}\NormalTok{ signAlgoOid}\OperatorTok{)} \KeywordTok{throws} \BuiltInTok{Exception} \OperatorTok{\{}

    \CommentTok{// {-}{-}{-} 1. Build the signed attributes {-}{-}{-}}
\NormalTok{    ASN1EncodableVector signedAttrs }\OperatorTok{=} \KeywordTok{new} \FunctionTok{ASN1EncodableVector}\OperatorTok{();}

    \CommentTok{// content{-}type = id{-}data}
\NormalTok{    signedAttrs}\OperatorTok{.}\FunctionTok{add}\OperatorTok{(}\KeywordTok{new} \BuiltInTok{Attribute}\OperatorTok{(}
\NormalTok{            CMSAttributes}\OperatorTok{.}\FunctionTok{contentType}\OperatorTok{,}
            \KeywordTok{new} \FunctionTok{DERSet}\OperatorTok{(}\NormalTok{CMSObjectIdentifiers}\OperatorTok{.}\FunctionTok{data}\OperatorTok{)));}

    \CommentTok{// message{-}digest = hash of byte range content}
\NormalTok{    signedAttrs}\OperatorTok{.}\FunctionTok{add}\OperatorTok{(}\KeywordTok{new} \BuiltInTok{Attribute}\OperatorTok{(}
\NormalTok{            CMSAttributes}\OperatorTok{.}\FunctionTok{messageDigest}\OperatorTok{,}
            \KeywordTok{new} \FunctionTok{DERSet}\OperatorTok{(}\KeywordTok{new} \FunctionTok{DEROctetString}\OperatorTok{(}\NormalTok{contentHash}\OperatorTok{))));}

    \CommentTok{// signing{-}time}
\NormalTok{    signedAttrs}\OperatorTok{.}\FunctionTok{add}\OperatorTok{(}\KeywordTok{new} \BuiltInTok{Attribute}\OperatorTok{(}
\NormalTok{            CMSAttributes}\OperatorTok{.}\FunctionTok{signingTime}\OperatorTok{,}
            \KeywordTok{new} \FunctionTok{DERSet}\OperatorTok{(}\KeywordTok{new} \FunctionTok{ASN1UTCTime}\OperatorTok{(}\KeywordTok{new} \BuiltInTok{Date}\OperatorTok{()))));}

    \CommentTok{// ESS{-}signing{-}certificate{-}v2 (mandatory for PAdES)}
    \DataTypeTok{byte}\OperatorTok{[]}\NormalTok{ certHash }\OperatorTok{=} \BuiltInTok{MessageDigest}\OperatorTok{.}\FunctionTok{getInstance}\OperatorTok{(}\StringTok{"SHA{-}256"}\OperatorTok{)}
            \OperatorTok{.}\FunctionTok{digest}\OperatorTok{(}\NormalTok{signerCert}\OperatorTok{.}\FunctionTok{getEncoded}\OperatorTok{());}
\NormalTok{    ESSCertIDv2 essCertId }\OperatorTok{=} \KeywordTok{new} \FunctionTok{ESSCertIDv2}\OperatorTok{(}
            \KeywordTok{new} \FunctionTok{AlgorithmIdentifier}\OperatorTok{(}\NormalTok{NISTObjectIdentifiers}\OperatorTok{.}\FunctionTok{id\_sha256}\OperatorTok{),}
\NormalTok{            certHash}\OperatorTok{);}
\NormalTok{    signedAttrs}\OperatorTok{.}\FunctionTok{add}\OperatorTok{(}\KeywordTok{new} \BuiltInTok{Attribute}\OperatorTok{(}
\NormalTok{            PKCSObjectIdentifiers}\OperatorTok{.}\FunctionTok{id\_aa\_signingCertificateV2}\OperatorTok{,}
            \KeywordTok{new} \FunctionTok{DERSet}\OperatorTok{(}\KeywordTok{new} \FunctionTok{SigningCertificateV2}\OperatorTok{(}
                    \KeywordTok{new}\NormalTok{ ESSCertIDv2}\OperatorTok{[]\{}\NormalTok{essCertId}\OperatorTok{\}))));}

\NormalTok{    DERSet signedAttrSet }\OperatorTok{=} \KeywordTok{new} \FunctionTok{DERSet}\OperatorTok{(}\NormalTok{signedAttrs}\OperatorTok{);}

    \CommentTok{// {-}{-}{-} 2. Build the SignerInfo {-}{-}{-}}
\NormalTok{    X509CertificateHolder certHolder }\OperatorTok{=}
            \KeywordTok{new} \FunctionTok{JcaX509CertificateHolder}\OperatorTok{(}\NormalTok{signerCert}\OperatorTok{);}
\NormalTok{    IssuerAndSerialNumber issuerAndSerial }\OperatorTok{=}
            \KeywordTok{new} \FunctionTok{IssuerAndSerialNumber}\OperatorTok{(}
\NormalTok{                    certHolder}\OperatorTok{.}\FunctionTok{getIssuer}\OperatorTok{(),}
\NormalTok{                    certHolder}\OperatorTok{.}\FunctionTok{getSerialNumber}\OperatorTok{());}

    \CommentTok{// Determine digest algorithm from signature algorithm}
\NormalTok{    AlgorithmIdentifier digestAlgId }\OperatorTok{=} \KeywordTok{new} \FunctionTok{AlgorithmIdentifier}\OperatorTok{(}
\NormalTok{            NISTObjectIdentifiers}\OperatorTok{.}\FunctionTok{id\_sha256}\OperatorTok{);}
\NormalTok{    AlgorithmIdentifier sigAlgId }\OperatorTok{=} \KeywordTok{new} \FunctionTok{AlgorithmIdentifier}\OperatorTok{(}
            \KeywordTok{new} \FunctionTok{ASN1ObjectIdentifier}\OperatorTok{(}\NormalTok{signAlgoOid}\OperatorTok{));}

\NormalTok{    SignerInfo signerInfo }\OperatorTok{=} \KeywordTok{new} \FunctionTok{SignerInfo}\OperatorTok{(}
            \KeywordTok{new} \FunctionTok{SignerIdentifier}\OperatorTok{(}\NormalTok{issuerAndSerial}\OperatorTok{),}
\NormalTok{            digestAlgId}\OperatorTok{,}
\NormalTok{            signedAttrSet}\OperatorTok{,}
\NormalTok{            sigAlgId}\OperatorTok{,}
            \KeywordTok{new} \FunctionTok{DEROctetString}\OperatorTok{(}\NormalTok{remoteSignatureValue}\OperatorTok{),}
            \KeywordTok{null}  \CommentTok{// no unsigned attributes yet (added later for B{-}T)}
    \OperatorTok{);}

    \CommentTok{// {-}{-}{-} 3. Build the certificate set {-}{-}{-}}
\NormalTok{    ASN1EncodableVector certVector }\OperatorTok{=} \KeywordTok{new} \FunctionTok{ASN1EncodableVector}\OperatorTok{();}
    \ControlFlowTok{for} \OperatorTok{(}\BuiltInTok{X509Certificate}\NormalTok{ cert }\OperatorTok{:}\NormalTok{ certChain}\OperatorTok{)} \OperatorTok{\{}
\NormalTok{        certVector}\OperatorTok{.}\FunctionTok{add}\OperatorTok{(}
\NormalTok{                org}\OperatorTok{.}\FunctionTok{bouncycastle}\OperatorTok{.}\FunctionTok{asn1}\OperatorTok{.}\FunctionTok{x509}\OperatorTok{.}\FunctionTok{Certificate}\OperatorTok{.}\FunctionTok{getInstance}\OperatorTok{(}
\NormalTok{                        cert}\OperatorTok{.}\FunctionTok{getEncoded}\OperatorTok{()));}
    \OperatorTok{\}}

    \CommentTok{// {-}{-}{-} 4. Assemble the SignedData {-}{-}{-}}
\NormalTok{    ASN1EncodableVector digestAlgIds }\OperatorTok{=} \KeywordTok{new} \FunctionTok{ASN1EncodableVector}\OperatorTok{();}
\NormalTok{    digestAlgIds}\OperatorTok{.}\FunctionTok{add}\OperatorTok{(}\NormalTok{digestAlgId}\OperatorTok{);}

\NormalTok{    SignedData signedData }\OperatorTok{=} \KeywordTok{new} \FunctionTok{SignedData}\OperatorTok{(}
            \KeywordTok{new} \FunctionTok{DERSet}\OperatorTok{(}\NormalTok{digestAlgIds}\OperatorTok{),}
            \KeywordTok{new} \FunctionTok{ContentInfo}\OperatorTok{(}\NormalTok{CMSObjectIdentifiers}\OperatorTok{.}\FunctionTok{data}\OperatorTok{,} \KeywordTok{null}\OperatorTok{),} \CommentTok{// detached}
            \KeywordTok{new} \FunctionTok{BERSet}\OperatorTok{(}\NormalTok{certVector}\OperatorTok{),}
            \KeywordTok{null}\OperatorTok{,}  \CommentTok{// no CRLs in the CMS itself (they go in the DSS)}
            \KeywordTok{new} \FunctionTok{DERSet}\OperatorTok{(}\NormalTok{signerInfo}\OperatorTok{));}

    \CommentTok{// {-}{-}{-} 5. Wrap in ContentInfo {-}{-}{-}}
\NormalTok{    ContentInfo cmsContentInfo }\OperatorTok{=} \KeywordTok{new} \FunctionTok{ContentInfo}\OperatorTok{(}
\NormalTok{            CMSObjectIdentifiers}\OperatorTok{.}\FunctionTok{signedData}\OperatorTok{,}
\NormalTok{            signedData}\OperatorTok{);}

    \ControlFlowTok{return}\NormalTok{ cmsContentInfo}\OperatorTok{.}\FunctionTok{getEncoded}\OperatorTok{(}\StringTok{"DER"}\OperatorTok{);}
\OperatorTok{\}}
\end{Highlighting}
\end{Shaded}

\subsection{Algorithm Detection for CSC
Credentials}\label{_algorithm_detection_for_csc_credentials}

The CSC credential info response tells you the key algorithm (RSA or EC)
but not always the exact signing algorithm OID. Use the key algorithm
OID to determine the correct signing algorithm:

\begin{Shaded}
\begin{Highlighting}[]
\CommentTok{/**}
 \CommentTok{*}\NormalTok{ Determines the signing algorithm OID from the CSC credential}\CommentTok{\textquotesingle{}}\NormalTok{s}
 \CommentTok{*}\NormalTok{ key algorithm}\CommentTok{. }\NormalTok{Defaults to SHA256withRSA for RSA keys and}
 \CommentTok{*}\NormalTok{ SHA256withECDSA for EC keys}\CommentTok{.}
 \CommentTok{*/}
\KeywordTok{public} \BuiltInTok{String} \FunctionTok{resolveSignAlgoOid}\OperatorTok{(}\BuiltInTok{String}\NormalTok{ keyAlgorithmOid}\OperatorTok{)} \OperatorTok{\{}
    \ControlFlowTok{if} \OperatorTok{(}\NormalTok{keyAlgorithmOid }\OperatorTok{!=} \KeywordTok{null}
            \OperatorTok{\&\&}\NormalTok{ keyAlgorithmOid}\OperatorTok{.}\FunctionTok{startsWith}\OperatorTok{(}\StringTok{"1.2.840.10045"}\OperatorTok{))} \OperatorTok{\{}
        \CommentTok{// EC key — use SHA256withECDSA}
        \ControlFlowTok{return} \StringTok{"1.2.840.10045.4.3.2"}\OperatorTok{;}
    \OperatorTok{\}}
    \CommentTok{// RSA key — use SHA256withRSA (PKCS\#1 v1.5)}
    \ControlFlowTok{return} \StringTok{"1.2.840.113549.1.1.11"}\OperatorTok{;}
\OperatorTok{\}}
\end{Highlighting}
\end{Shaded}

\subsection{Adding an Unsigned Timestamp to the Remote
CMS}\label{_adding_an_unsigned_timestamp_to_the_remote_cms}

After assembling the CMS with the remote signature value, you can add an
RFC 3161 timestamp as an unsigned attribute to achieve PAdES-B-T. This
requires parsing the CMS back, adding the unsigned attribute, and
re-encoding:

\begin{Shaded}
\begin{Highlighting}[]
\CommentTok{/**}
 \CommentTok{*}\NormalTok{ Adds an RFC }\CommentTok{3161}\NormalTok{ timestamp token as an unsigned attribute}
 \CommentTok{*}\NormalTok{ to an existing CMS SignedData}\CommentTok{.}
 \CommentTok{*}\NormalTok{ The timestamp is over the signature value }\CommentTok{(}\NormalTok{not the content}\CommentTok{).}
 \CommentTok{*/}
\KeywordTok{public} \DataTypeTok{byte}\OperatorTok{[]} \FunctionTok{addTimestampToCms}\OperatorTok{(}\DataTypeTok{byte}\OperatorTok{[]}\NormalTok{ cmsBytes}\OperatorTok{,}
        \DataTypeTok{byte}\OperatorTok{[]}\NormalTok{ signatureValue}\OperatorTok{,} \BuiltInTok{String}\NormalTok{ tsaUrl}\OperatorTok{)} \KeywordTok{throws} \BuiltInTok{Exception} \OperatorTok{\{}

    \CommentTok{// 1. Get timestamp token over the signature value}
    \DataTypeTok{byte}\OperatorTok{[]}\NormalTok{ tsToken }\OperatorTok{=} \FunctionTok{requestTimestamp}\OperatorTok{(}\NormalTok{signatureValue}\OperatorTok{,}\NormalTok{ tsaUrl}\OperatorTok{);}

    \CommentTok{// 2. Parse the existing CMS}
\NormalTok{    CMSSignedData signedData }\OperatorTok{=} \KeywordTok{new} \FunctionTok{CMSSignedData}\OperatorTok{(}\NormalTok{cmsBytes}\OperatorTok{);}
\NormalTok{    SignerInformationStore signerStore }\OperatorTok{=}\NormalTok{ signedData}\OperatorTok{.}\FunctionTok{getSignerInfos}\OperatorTok{();}
\NormalTok{    SignerInformation signer }\OperatorTok{=}\NormalTok{ signerStore}\OperatorTok{.}\FunctionTok{getSigners}\OperatorTok{()}
            \OperatorTok{.}\FunctionTok{iterator}\OperatorTok{().}\FunctionTok{next}\OperatorTok{();}

    \CommentTok{// 3. Add the timestamp as an unsigned attribute}
    \BuiltInTok{Hashtable}\OperatorTok{\textless{}}\NormalTok{ASN1ObjectIdentifier}\OperatorTok{,} \BuiltInTok{Attribute}\OperatorTok{\textgreater{}}\NormalTok{ unsignedAttrs }\OperatorTok{=}
            \KeywordTok{new} \BuiltInTok{Hashtable}\OperatorTok{\textless{}\textgreater{}();}
\NormalTok{    unsignedAttrs}\OperatorTok{.}\FunctionTok{put}\OperatorTok{(}
\NormalTok{            PKCSObjectIdentifiers}\OperatorTok{.}\FunctionTok{id\_aa\_signatureTimeStampToken}\OperatorTok{,}
            \KeywordTok{new} \BuiltInTok{Attribute}\OperatorTok{(}
\NormalTok{                    PKCSObjectIdentifiers}\OperatorTok{.}\FunctionTok{id\_aa\_signatureTimeStampToken}\OperatorTok{,}
                    \KeywordTok{new} \FunctionTok{DERSet}\OperatorTok{(}\NormalTok{ASN1Primitive}\OperatorTok{.}\FunctionTok{fromByteArray}\OperatorTok{(}\NormalTok{tsToken}\OperatorTok{))));}

\NormalTok{    AttributeTable unsignedTable }\OperatorTok{=} \KeywordTok{new} \FunctionTok{AttributeTable}\OperatorTok{(}\NormalTok{unsignedAttrs}\OperatorTok{);}
\NormalTok{    SignerInformation timestampedSigner }\OperatorTok{=}
\NormalTok{            SignerInformation}\OperatorTok{.}\FunctionTok{replaceUnsignedAttributes}\OperatorTok{(}
\NormalTok{                    signer}\OperatorTok{,}\NormalTok{ unsignedTable}\OperatorTok{);}

    \CommentTok{// 4. Rebuild the SignedData with the updated SignerInfo}
    \BuiltInTok{List}\OperatorTok{\textless{}}\NormalTok{SignerInformation}\OperatorTok{\textgreater{}}\NormalTok{ newSigners }\OperatorTok{=} \KeywordTok{new} \BuiltInTok{ArrayList}\OperatorTok{\textless{}\textgreater{}();}
\NormalTok{    newSigners}\OperatorTok{.}\FunctionTok{add}\OperatorTok{(}\NormalTok{timestampedSigner}\OperatorTok{);}
\NormalTok{    CMSSignedData timestampedCms }\OperatorTok{=}\NormalTok{ CMSSignedData}\OperatorTok{.}\FunctionTok{replaceSigners}\OperatorTok{(}
\NormalTok{            signedData}\OperatorTok{,} \KeywordTok{new} \FunctionTok{SignerInformationStore}\OperatorTok{(}\NormalTok{newSigners}\OperatorTok{));}

    \ControlFlowTok{return}\NormalTok{ timestampedCms}\OperatorTok{.}\FunctionTok{getEncoded}\OperatorTok{();}
\OperatorTok{\}}
\end{Highlighting}
\end{Shaded}

This pattern --- build CMS without unsigned attributes, then parse and
add the timestamp --- is the standard approach for remote signing. You
cannot add the timestamp during initial CMS construction because the
timestamp must cover the signature value, which does not exist until the
CSC API returns it. The two-phase assembly is inherent to the remote
signing model.

For local signing (where you have the private key), BouncyCastle's
\texttt{CMSSignedDataGenerator} can add the timestamp in a single pass.
The manual assembly shown here is only needed when the signing operation
is split between local CMS construction and remote cryptographic
operation.

\section{Putting It All Together}\label{_putting_it_all_together}

Here is how you use the PAdES-B-B and PAdES-B-T builders with PDFBox:

\begin{Shaded}
\begin{Highlighting}[]
\KeywordTok{import} \ImportTok{org}\OperatorTok{.}\ImportTok{apache}\OperatorTok{.}\ImportTok{pdfbox}\OperatorTok{.}\ImportTok{Loader}\OperatorTok{;}
\KeywordTok{import} \ImportTok{org}\OperatorTok{.}\ImportTok{apache}\OperatorTok{.}\ImportTok{pdfbox}\OperatorTok{.}\ImportTok{cos}\OperatorTok{.}\ImportTok{COSName}\OperatorTok{;}
\KeywordTok{import} \ImportTok{org}\OperatorTok{.}\ImportTok{apache}\OperatorTok{.}\ImportTok{pdfbox}\OperatorTok{.}\ImportTok{pdmodel}\OperatorTok{.}\ImportTok{PDDocument}\OperatorTok{;}
\KeywordTok{import} \ImportTok{org}\OperatorTok{.}\ImportTok{apache}\OperatorTok{.}\ImportTok{pdfbox}\OperatorTok{.}\ImportTok{pdmodel}\OperatorTok{.}\ImportTok{interactive}\OperatorTok{.}\ImportTok{digitalsignature}\OperatorTok{.}\ImportTok{PDSignature}\OperatorTok{;}
\KeywordTok{import} \ImportTok{org}\OperatorTok{.}\ImportTok{apache}\OperatorTok{.}\ImportTok{pdfbox}\OperatorTok{.}\ImportTok{pdmodel}\OperatorTok{.}\ImportTok{interactive}\OperatorTok{.}\ImportTok{digitalsignature}\OperatorTok{.}\ImportTok{SignatureInterface}\OperatorTok{;}
\KeywordTok{import} \ImportTok{org}\OperatorTok{.}\ImportTok{apache}\OperatorTok{.}\ImportTok{pdfbox}\OperatorTok{.}\ImportTok{pdmodel}\OperatorTok{.}\ImportTok{interactive}\OperatorTok{.}\ImportTok{digitalsignature}\OperatorTok{.}\ImportTok{SignatureOptions}\OperatorTok{;}

\KeywordTok{import} \ImportTok{java}\OperatorTok{.}\ImportTok{io}\OperatorTok{.}\ImportTok{ByteArrayOutputStream}\OperatorTok{;}
\KeywordTok{import} \ImportTok{java}\OperatorTok{.}\ImportTok{util}\OperatorTok{.}\ImportTok{Calendar}\OperatorTok{;}

\CommentTok{// {-}{-}{-} PAdES{-}B{-}B (no timestamp) {-}{-}{-}}
\NormalTok{SignatureInterface bbSigner }\OperatorTok{=} \KeywordTok{new} \FunctionTok{PadesBBBuilder}\OperatorTok{(}
\NormalTok{        privateKey}\OperatorTok{,}\NormalTok{ certificate}\OperatorTok{,}\NormalTok{ chain}\OperatorTok{);}

\ControlFlowTok{try} \OperatorTok{(}\NormalTok{PDDocument doc }\OperatorTok{=}\NormalTok{ Loader}\OperatorTok{.}\FunctionTok{loadPDF}\OperatorTok{(}\NormalTok{pdfBytes}\OperatorTok{))} \OperatorTok{\{}
\NormalTok{    PDSignature sig }\OperatorTok{=} \KeywordTok{new} \FunctionTok{PDSignature}\OperatorTok{();}
\NormalTok{    sig}\OperatorTok{.}\FunctionTok{setFilter}\OperatorTok{(}\NormalTok{PDSignature}\OperatorTok{.}\FunctionTok{FILTER\_ADOBE\_PPKLITE}\OperatorTok{);}
\NormalTok{    sig}\OperatorTok{.}\FunctionTok{setSubFilter}\OperatorTok{(}\NormalTok{COSName}\OperatorTok{.}\FunctionTok{getPDFName}\OperatorTok{(}\StringTok{"ETSI.CAdES.detached"}\OperatorTok{));}
\NormalTok{    sig}\OperatorTok{.}\FunctionTok{setSignDate}\OperatorTok{(}\BuiltInTok{Calendar}\OperatorTok{.}\FunctionTok{getInstance}\OperatorTok{());}

\NormalTok{    SignatureOptions opts }\OperatorTok{=} \KeywordTok{new} \FunctionTok{SignatureOptions}\OperatorTok{();}
\NormalTok{    opts}\OperatorTok{.}\FunctionTok{setPreferredSignatureSize}\OperatorTok{(}\DecValTok{32768}\OperatorTok{);}

\NormalTok{    doc}\OperatorTok{.}\FunctionTok{addSignature}\OperatorTok{(}\NormalTok{sig}\OperatorTok{,}\NormalTok{ bbSigner}\OperatorTok{,}\NormalTok{ opts}\OperatorTok{);}

    \BuiltInTok{ByteArrayOutputStream}\NormalTok{ out }\OperatorTok{=} \KeywordTok{new} \BuiltInTok{ByteArrayOutputStream}\OperatorTok{();}
\NormalTok{    doc}\OperatorTok{.}\FunctionTok{saveIncremental}\OperatorTok{(}\NormalTok{out}\OperatorTok{);}
    \CommentTok{// out.toByteArray() is the PAdES{-}B{-}B signed PDF}
\OperatorTok{\}}

\CommentTok{// {-}{-}{-} PAdES{-}B{-}T (with timestamp) {-}{-}{-}}
\NormalTok{SignatureInterface btSigner }\OperatorTok{=} \KeywordTok{new} \FunctionTok{PadesBTBuilder}\OperatorTok{(}
\NormalTok{        privateKey}\OperatorTok{,}\NormalTok{ certificate}\OperatorTok{,}\NormalTok{ chain}\OperatorTok{,}
        \StringTok{"http://timestamp.example.com/tsa"}\OperatorTok{);}

\ControlFlowTok{try} \OperatorTok{(}\NormalTok{PDDocument doc }\OperatorTok{=}\NormalTok{ Loader}\OperatorTok{.}\FunctionTok{loadPDF}\OperatorTok{(}\NormalTok{pdfBytes}\OperatorTok{))} \OperatorTok{\{}
\NormalTok{    PDSignature sig }\OperatorTok{=} \KeywordTok{new} \FunctionTok{PDSignature}\OperatorTok{();}
\NormalTok{    sig}\OperatorTok{.}\FunctionTok{setFilter}\OperatorTok{(}\NormalTok{PDSignature}\OperatorTok{.}\FunctionTok{FILTER\_ADOBE\_PPKLITE}\OperatorTok{);}
\NormalTok{    sig}\OperatorTok{.}\FunctionTok{setSubFilter}\OperatorTok{(}\NormalTok{COSName}\OperatorTok{.}\FunctionTok{getPDFName}\OperatorTok{(}\StringTok{"ETSI.CAdES.detached"}\OperatorTok{));}
\NormalTok{    sig}\OperatorTok{.}\FunctionTok{setSignDate}\OperatorTok{(}\BuiltInTok{Calendar}\OperatorTok{.}\FunctionTok{getInstance}\OperatorTok{());}

\NormalTok{    SignatureOptions opts }\OperatorTok{=} \KeywordTok{new} \FunctionTok{SignatureOptions}\OperatorTok{();}
\NormalTok{    opts}\OperatorTok{.}\FunctionTok{setPreferredSignatureSize}\OperatorTok{(}\DecValTok{65536}\OperatorTok{);} \CommentTok{// larger for timestamp}

\NormalTok{    doc}\OperatorTok{.}\FunctionTok{addSignature}\OperatorTok{(}\NormalTok{sig}\OperatorTok{,}\NormalTok{ btSigner}\OperatorTok{,}\NormalTok{ opts}\OperatorTok{);}

    \BuiltInTok{ByteArrayOutputStream}\NormalTok{ out }\OperatorTok{=} \KeywordTok{new} \BuiltInTok{ByteArrayOutputStream}\OperatorTok{();}
\NormalTok{    doc}\OperatorTok{.}\FunctionTok{saveIncremental}\OperatorTok{(}\NormalTok{out}\OperatorTok{);}
    \CommentTok{// out.toByteArray() is the PAdES{-}B{-}T signed PDF}
\OperatorTok{\}}
\end{Highlighting}
\end{Shaded}

The \texttt{PadesBBBuilder} and \texttt{PadesBTBuilder} classes are
production-ready. They handle algorithm negotiation, ESS attribute
construction, certificate embedding, and timestamp integration. Drop
them into your project alongside the BouncyCastle and PDFBox
dependencies from Chapter 28, register the BouncyCastle provider, and
you have a complete signing pipeline.

In Chapter 31, we reverse the process: given a signed PDF, we verify the
CMS signature, validate the certificate chain, check the timestamp, and
extract every piece of information that these structures contain.

\chapter{Verification \& Analysis with
PDFBox}\label{_verification_analysis_with_pdfbox}

Signing a PDF is a matter of following a recipe. Verifying a signed PDF
is where things get interesting --- and where most implementations get
things wrong. PDFBox gives you every primitive you need: signature
extraction, byte range access, CMS parsing, content stream reading. But
it deliberately does not give you a "verify this signature" method. You
build the pipeline yourself, step by step, and you understand every
failure mode because you wrote the code that handles it.

This chapter walks through a complete verification implementation using
Apache PDFBox 3.0.0 and BouncyCastle 1.80. We start with extracting
signatures from a PDF document, proceed through each verification step
from Chapter 16, and end with a production-ready five-step pipeline that
handles the edge cases real-world PDFs throw at you. If Chapter 16
described \emph{what} to verify, this chapter shows you \emph{how} ---
with code you can compile and run.

\section{Extracting Signatures from a PDF
Document}\label{_extracting_signatures_from_a_pdf_document}

\subsection{The Extraction Path: AcroForm to
PDSignature}\label{_the_extraction_path_acroform_to_pdsignature}

Every PDF signature lives inside the AcroForm field tree. The extraction
path follows the object hierarchy we dissected in Chapter 12: the
document's \texttt{/AcroForm} dictionary contains a \texttt{/Fields}
array, which includes \texttt{PDSignatureField} entries, each of which
references a \texttt{PDSignature} dictionary.

PDFBox models this hierarchy directly. The key classes are:

\begin{itemize}
\item
  \texttt{PDAcroForm} --- the form dictionary, accessed via
  \texttt{PDDocument.getDocumentCatalog().getAcroForm()}
\item
  \texttt{PDSignatureField} --- a form field of type \texttt{/Sig},
  extends \texttt{PDTerminalField}
\item
  \texttt{PDSignature} --- the actual \texttt{/Type\ /Sig} dictionary
  containing \texttt{/Contents}, \texttt{/ByteRange},
  \texttt{/SubFilter}, etc.
\end{itemize}

There are two ways to access signatures in PDFBox. The first walks the
field tree:

\begin{Shaded}
\begin{Highlighting}[]
\CommentTok{/**}
 \CommentTok{*}\NormalTok{ Extract all signature fields from the AcroForm field tree}\CommentTok{.}
 \CommentTok{*}\NormalTok{ This approach gives you access to the field }\CommentTok{(}\NormalTok{widget}\CommentTok{,}\NormalTok{ appearance}\CommentTok{)}
 \CommentTok{*}\NormalTok{ in addition to the signature dictionary}\CommentTok{.}
 \CommentTok{*/}
\KeywordTok{public} \BuiltInTok{List}\OperatorTok{\textless{}}\NormalTok{PDSignatureField}\OperatorTok{\textgreater{}} \FunctionTok{extractSignatureFields}\OperatorTok{(}\NormalTok{PDDocument document}\OperatorTok{)} \OperatorTok{\{}
\NormalTok{    PDAcroForm acroForm }\OperatorTok{=}\NormalTok{ document}\OperatorTok{.}\FunctionTok{getDocumentCatalog}\OperatorTok{().}\FunctionTok{getAcroForm}\OperatorTok{();}
    \ControlFlowTok{if} \OperatorTok{(}\NormalTok{acroForm }\OperatorTok{==} \KeywordTok{null}\OperatorTok{)} \OperatorTok{\{}
\NormalTok{        log}\OperatorTok{.}\FunctionTok{info}\OperatorTok{(}\StringTok{"Document has no AcroForm — no signatures possible"}\OperatorTok{);}
        \ControlFlowTok{return} \BuiltInTok{Collections}\OperatorTok{.}\FunctionTok{emptyList}\OperatorTok{();}
    \OperatorTok{\}}

    \BuiltInTok{List}\OperatorTok{\textless{}}\NormalTok{PDSignatureField}\OperatorTok{\textgreater{}}\NormalTok{ signatureFields }\OperatorTok{=} \KeywordTok{new} \BuiltInTok{ArrayList}\OperatorTok{\textless{}\textgreater{}();}
    \ControlFlowTok{for} \OperatorTok{(}\NormalTok{PDField field }\OperatorTok{:}\NormalTok{ acroForm}\OperatorTok{.}\FunctionTok{getFields}\OperatorTok{())} \OperatorTok{\{}
        \FunctionTok{collectSignatureFields}\OperatorTok{(}\NormalTok{field}\OperatorTok{,}\NormalTok{ signatureFields}\OperatorTok{);}
    \OperatorTok{\}}

\NormalTok{    log}\OperatorTok{.}\FunctionTok{info}\OperatorTok{(}\StringTok{"Found \{\} signature field(s) in AcroForm"}\OperatorTok{,}\NormalTok{ signatureFields}\OperatorTok{.}\FunctionTok{size}\OperatorTok{());}
    \ControlFlowTok{return}\NormalTok{ signatureFields}\OperatorTok{;}
\OperatorTok{\}}

\KeywordTok{private} \DataTypeTok{void} \FunctionTok{collectSignatureFields}\OperatorTok{(}\NormalTok{PDField field}\OperatorTok{,}
        \BuiltInTok{List}\OperatorTok{\textless{}}\NormalTok{PDSignatureField}\OperatorTok{\textgreater{}}\NormalTok{ result}\OperatorTok{)} \OperatorTok{\{}
    \ControlFlowTok{if} \OperatorTok{(}\NormalTok{field }\KeywordTok{instanceof}\NormalTok{ PDSignatureField sigField}\OperatorTok{)} \OperatorTok{\{}
\NormalTok{        result}\OperatorTok{.}\FunctionTok{add}\OperatorTok{(}\NormalTok{sigField}\OperatorTok{);}
    \OperatorTok{\}} \ControlFlowTok{else} \ControlFlowTok{if} \OperatorTok{(}\NormalTok{field }\KeywordTok{instanceof}\NormalTok{ PDNonTerminalField nonTerminal}\OperatorTok{)} \OperatorTok{\{}
        \ControlFlowTok{for} \OperatorTok{(}\NormalTok{PDField child }\OperatorTok{:}\NormalTok{ nonTerminal}\OperatorTok{.}\FunctionTok{getChildren}\OperatorTok{())} \OperatorTok{\{}
            \FunctionTok{collectSignatureFields}\OperatorTok{(}\NormalTok{child}\OperatorTok{,}\NormalTok{ result}\OperatorTok{);}
        \OperatorTok{\}}
    \OperatorTok{\}}
\OperatorTok{\}}
\end{Highlighting}
\end{Shaded}

The second approach bypasses the field tree and goes directly to the
signature dictionaries:

\begin{Shaded}
\begin{Highlighting}[]
\CommentTok{/**}
 \CommentTok{*}\NormalTok{ Extract all PDSignature dictionaries directly}\CommentTok{.}
 \CommentTok{*}\NormalTok{ Faster}\CommentTok{,}\NormalTok{ but you lose access to the field and widget objects}\CommentTok{.}
 \CommentTok{*}\NormalTok{ Useful when you only need the cryptographic data}\CommentTok{.}
 \CommentTok{*/}
\KeywordTok{public} \BuiltInTok{List}\OperatorTok{\textless{}}\NormalTok{PDSignature}\OperatorTok{\textgreater{}} \FunctionTok{extractSignatures}\OperatorTok{(}\NormalTok{PDDocument document}\OperatorTok{)} \OperatorTok{\{}
    \BuiltInTok{List}\OperatorTok{\textless{}}\NormalTok{PDSignature}\OperatorTok{\textgreater{}}\NormalTok{ signatures }\OperatorTok{=}\NormalTok{ document}\OperatorTok{.}\FunctionTok{getSignatureDictionaries}\OperatorTok{();}
\NormalTok{    log}\OperatorTok{.}\FunctionTok{info}\OperatorTok{(}\StringTok{"Found \{\} signature dictionaries"}\OperatorTok{,}\NormalTok{ signatures}\OperatorTok{.}\FunctionTok{size}\OperatorTok{());}
    \ControlFlowTok{return}\NormalTok{ signatures}\OperatorTok{;}
\OperatorTok{\}}
\end{Highlighting}
\end{Shaded}

Use the first approach (field tree walk) when you need the full
picture --- field name, widget annotation, visual appearance, page
location. Use the second approach when you only need the cryptographic
payload --- for pure verification, it is faster and simpler.

\subsection{Reading Signature
Metadata}\label{_reading_signature_metadata}

Once you have a \texttt{PDSignature}, you can extract the metadata that
drives the verification pipeline:

\begin{Shaded}
\begin{Highlighting}[]
\KeywordTok{public} \DataTypeTok{void} \FunctionTok{logSignatureMetadata}\OperatorTok{(}\NormalTok{PDSignature sig}\OperatorTok{)} \OperatorTok{\{}
\NormalTok{    log}\OperatorTok{.}\FunctionTok{info}\OperatorTok{(}\StringTok{"SubFilter: \{\}"}\OperatorTok{,}\NormalTok{ sig}\OperatorTok{.}\FunctionTok{getSubFilter}\OperatorTok{());}
\NormalTok{    log}\OperatorTok{.}\FunctionTok{info}\OperatorTok{(}\StringTok{"Filter: \{\}"}\OperatorTok{,}\NormalTok{ sig}\OperatorTok{.}\FunctionTok{getFilter}\OperatorTok{());}
\NormalTok{    log}\OperatorTok{.}\FunctionTok{info}\OperatorTok{(}\StringTok{"Name: \{\}"}\OperatorTok{,}\NormalTok{ sig}\OperatorTok{.}\FunctionTok{getName}\OperatorTok{());}
\NormalTok{    log}\OperatorTok{.}\FunctionTok{info}\OperatorTok{(}\StringTok{"Location: \{\}"}\OperatorTok{,}\NormalTok{ sig}\OperatorTok{.}\FunctionTok{getLocation}\OperatorTok{());}
\NormalTok{    log}\OperatorTok{.}\FunctionTok{info}\OperatorTok{(}\StringTok{"Reason: \{\}"}\OperatorTok{,}\NormalTok{ sig}\OperatorTok{.}\FunctionTok{getReason}\OperatorTok{());}
\NormalTok{    log}\OperatorTok{.}\FunctionTok{info}\OperatorTok{(}\StringTok{"Sign date: \{\}"}\OperatorTok{,}\NormalTok{ sig}\OperatorTok{.}\FunctionTok{getSignDate}\OperatorTok{());}
\NormalTok{    log}\OperatorTok{.}\FunctionTok{info}\OperatorTok{(}\StringTok{"Contact info: \{\}"}\OperatorTok{,}\NormalTok{ sig}\OperatorTok{.}\FunctionTok{getContactInfo}\OperatorTok{());}

    \DataTypeTok{int}\OperatorTok{[]}\NormalTok{ byteRange }\OperatorTok{=}\NormalTok{ sig}\OperatorTok{.}\FunctionTok{getByteRange}\OperatorTok{();}
    \ControlFlowTok{if} \OperatorTok{(}\NormalTok{byteRange }\OperatorTok{!=} \KeywordTok{null} \OperatorTok{\&\&}\NormalTok{ byteRange}\OperatorTok{.}\FunctionTok{length} \OperatorTok{==} \DecValTok{4}\OperatorTok{)} \OperatorTok{\{}
\NormalTok{        log}\OperatorTok{.}\FunctionTok{info}\OperatorTok{(}\StringTok{"ByteRange: [\{\} \{\} \{\} \{\}]"}\OperatorTok{,}
\NormalTok{                byteRange}\OperatorTok{[}\DecValTok{0}\OperatorTok{],}\NormalTok{ byteRange}\OperatorTok{[}\DecValTok{1}\OperatorTok{],}\NormalTok{ byteRange}\OperatorTok{[}\DecValTok{2}\OperatorTok{],}\NormalTok{ byteRange}\OperatorTok{[}\DecValTok{3}\OperatorTok{]);}
\NormalTok{        log}\OperatorTok{.}\FunctionTok{info}\OperatorTok{(}\StringTok{"Signed bytes: \{\} + \{\} = \{\} total"}\OperatorTok{,}
\NormalTok{                byteRange}\OperatorTok{[}\DecValTok{1}\OperatorTok{],}\NormalTok{ byteRange}\OperatorTok{[}\DecValTok{3}\OperatorTok{],}
\NormalTok{                byteRange}\OperatorTok{[}\DecValTok{1}\OperatorTok{]} \OperatorTok{+}\NormalTok{ byteRange}\OperatorTok{[}\DecValTok{3}\OperatorTok{]);}
\NormalTok{        log}\OperatorTok{.}\FunctionTok{info}\OperatorTok{(}\StringTok{"Gap (Contents): offset \{\} to \{\}, size \{\} bytes"}\OperatorTok{,}
\NormalTok{                byteRange}\OperatorTok{[}\DecValTok{1}\OperatorTok{],}\NormalTok{ byteRange}\OperatorTok{[}\DecValTok{2}\OperatorTok{],}
\NormalTok{                byteRange}\OperatorTok{[}\DecValTok{2}\OperatorTok{]} \OperatorTok{{-}}\NormalTok{ byteRange}\OperatorTok{[}\DecValTok{1}\OperatorTok{]);}
    \OperatorTok{\}}
\OperatorTok{\}}
\end{Highlighting}
\end{Shaded}

The \texttt{/SubFilter} value tells you what kind of signature you are
dealing with:

\begin{itemize}
\item
  \texttt{adbe.pkcs7.detached} --- Classic PKCS\#7 (CMS) detached
  signature
\item
  \texttt{ETSI.CAdES.detached} --- PAdES CAdES-based signature (the
  modern standard)
\item
  \texttt{adbe.pkcs7.sha1} --- Legacy: SHA-1 hash embedded in PKCS\#7
  (deprecated, avoid)
\item
  \texttt{adbe.x509.rsa\_sha1} --- Legacy: bare RSA signature with X.509
  cert (deprecated, avoid)
\end{itemize}

For any new code, you should handle \texttt{adbe.pkcs7.detached} and
\texttt{ETSI.CAdES.detached}. The legacy formats exist in old documents,
and you may encounter them, but you should never produce them.

\section{Parsing the CMS SignedData from
/Contents}\label{_parsing_the_cms_signeddata_from_contents}

\subsection{Extracting the Raw CMS
Bytes}\label{_extracting_the_raw_cms_bytes}

The \texttt{/Contents} value in the signature dictionary is a
hex-encoded byte string containing the DER-encoded CMS
\texttt{SignedData} structure. PDFBox provides this as a byte array:

\begin{Shaded}
\begin{Highlighting}[]
\DataTypeTok{byte}\OperatorTok{[]}\NormalTok{ cmsBytes }\OperatorTok{=}\NormalTok{ signature}\OperatorTok{.}\FunctionTok{getContents}\OperatorTok{();}
\end{Highlighting}
\end{Shaded}

This single line hides a meaningful amount of work. PDFBox reads the hex
string from the PDF (which can be 16KB to 64KB or more of hex
characters), decodes it from hex to binary, and strips any trailing zero
padding. The result is the raw DER-encoded CMS \texttt{SignedData} that
BouncyCastle can parse.

\subsection{Parsing with BouncyCastle}\label{_parsing_with_bouncycastle}

To parse the CMS structure, you also need the signed bytes --- the byte
range segments from the PDF file. For a detached signature (which is
what \texttt{adbe.pkcs7.detached} and \texttt{ETSI.CAdES.detached} are),
the signed content is not inside the CMS; it is the external byte range
data.

\begin{Shaded}
\begin{Highlighting}[]
\CommentTok{/**}
 \CommentTok{*}\NormalTok{ Parse a CMS SignedData from the }\CommentTok{/}\NormalTok{Contents of a PDF signature}\CommentTok{.}
 \CommentTok{*}
 \CommentTok{*} \AnnotationTok{@param}\CommentTok{ }\AnnotationTok{pdfBytes}\CommentTok{      }\NormalTok{the complete PDF file as a byte array}
 \CommentTok{*} \AnnotationTok{@param}\CommentTok{ }\AnnotationTok{signature}\CommentTok{     }\NormalTok{the PDSignature whose CMS we want to parse}
 \CommentTok{*} \AnnotationTok{@return}\NormalTok{ the parsed CMSSignedData}\CommentTok{,}\NormalTok{ ready for verification}
 \CommentTok{*/}
\KeywordTok{public}\NormalTok{ CMSSignedData }\FunctionTok{parseCmsSignedData}\OperatorTok{(}\DataTypeTok{byte}\OperatorTok{[]}\NormalTok{ pdfBytes}\OperatorTok{,}
\NormalTok{        PDSignature signature}\OperatorTok{)} \KeywordTok{throws}\NormalTok{ CMSException }\OperatorTok{\{}

    \CommentTok{// Extract the signed byte ranges}
    \DataTypeTok{int}\OperatorTok{[]}\NormalTok{ byteRange }\OperatorTok{=}\NormalTok{ signature}\OperatorTok{.}\FunctionTok{getByteRange}\OperatorTok{();}
    \DataTypeTok{byte}\OperatorTok{[]}\NormalTok{ signedContent }\OperatorTok{=} \KeywordTok{new} \DataTypeTok{byte}\OperatorTok{[}\NormalTok{byteRange}\OperatorTok{[}\DecValTok{1}\OperatorTok{]} \OperatorTok{+}\NormalTok{ byteRange}\OperatorTok{[}\DecValTok{3}\OperatorTok{]];}
    \BuiltInTok{System}\OperatorTok{.}\FunctionTok{arraycopy}\OperatorTok{(}\NormalTok{pdfBytes}\OperatorTok{,}\NormalTok{ byteRange}\OperatorTok{[}\DecValTok{0}\OperatorTok{],}\NormalTok{ signedContent}\OperatorTok{,}
            \DecValTok{0}\OperatorTok{,}\NormalTok{ byteRange}\OperatorTok{[}\DecValTok{1}\OperatorTok{]);}
    \BuiltInTok{System}\OperatorTok{.}\FunctionTok{arraycopy}\OperatorTok{(}\NormalTok{pdfBytes}\OperatorTok{,}\NormalTok{ byteRange}\OperatorTok{[}\DecValTok{2}\OperatorTok{],}\NormalTok{ signedContent}\OperatorTok{,}
\NormalTok{            byteRange}\OperatorTok{[}\DecValTok{1}\OperatorTok{],}\NormalTok{ byteRange}\OperatorTok{[}\DecValTok{3}\OperatorTok{]);}

    \CommentTok{// Parse the CMS with the signed content as the detached data}
\NormalTok{    CMSProcessableByteArray processable }\OperatorTok{=}
            \KeywordTok{new} \FunctionTok{CMSProcessableByteArray}\OperatorTok{(}\NormalTok{signedContent}\OperatorTok{);}
    \ControlFlowTok{return} \KeywordTok{new} \FunctionTok{CMSSignedData}\OperatorTok{(}\NormalTok{processable}\OperatorTok{,}\NormalTok{ signature}\OperatorTok{.}\FunctionTok{getContents}\OperatorTok{());}
\OperatorTok{\}}
\end{Highlighting}
\end{Shaded}

A common mistake is to construct \texttt{CMSSignedData} with only the
\texttt{/Contents} bytes (without the signed content). This will parse
the CMS structure, but any subsequent verification call will fail
because BouncyCastle cannot compute the message digest without the
signed data. Always pass both the signed content and the CMS bytes.

\subsection{Navigating the CMS
Structure}\label{_navigating_the_cms_structure}

Once parsed, the \texttt{CMSSignedData} object exposes the full CMS tree
from RFC 5652 --- Cryptographic Message Syntax (CMS):

\begin{Shaded}
\begin{Highlighting}[]
\KeywordTok{public} \DataTypeTok{void} \FunctionTok{inspectCmsStructure}\OperatorTok{(}\NormalTok{CMSSignedData cmsData}\OperatorTok{)} \OperatorTok{\{}
    \CommentTok{// Signer information (usually exactly one for PDF signatures)}
\NormalTok{    SignerInformationStore signerStore }\OperatorTok{=}\NormalTok{ cmsData}\OperatorTok{.}\FunctionTok{getSignerInfos}\OperatorTok{();}
\NormalTok{    log}\OperatorTok{.}\FunctionTok{info}\OperatorTok{(}\StringTok{"Number of signers: \{\}"}\OperatorTok{,}\NormalTok{ signerStore}\OperatorTok{.}\FunctionTok{size}\OperatorTok{());}

    \ControlFlowTok{for} \OperatorTok{(}\NormalTok{SignerInformation signer }\OperatorTok{:}\NormalTok{ signerStore}\OperatorTok{.}\FunctionTok{getSigners}\OperatorTok{())} \OperatorTok{\{}
\NormalTok{        log}\OperatorTok{.}\FunctionTok{info}\OperatorTok{(}\StringTok{"Digest algorithm OID: \{\}"}\OperatorTok{,}\NormalTok{ signer}\OperatorTok{.}\FunctionTok{getDigestAlgOID}\OperatorTok{());}
\NormalTok{        log}\OperatorTok{.}\FunctionTok{info}\OperatorTok{(}\StringTok{"Encryption algorithm OID: \{\}"}\OperatorTok{,}
\NormalTok{                signer}\OperatorTok{.}\FunctionTok{getEncryptionAlgOID}\OperatorTok{());}

        \CommentTok{// Signed attributes (authenticated)}
\NormalTok{        AttributeTable signedAttrs }\OperatorTok{=}\NormalTok{ signer}\OperatorTok{.}\FunctionTok{getSignedAttributes}\OperatorTok{();}
        \ControlFlowTok{if} \OperatorTok{(}\NormalTok{signedAttrs }\OperatorTok{!=} \KeywordTok{null}\OperatorTok{)} \OperatorTok{\{}
\NormalTok{            log}\OperatorTok{.}\FunctionTok{info}\OperatorTok{(}\StringTok{"Signed attributes count: \{\}"}\OperatorTok{,}
\NormalTok{                    signedAttrs}\OperatorTok{.}\FunctionTok{size}\OperatorTok{());}
            \FunctionTok{logAttribute}\OperatorTok{(}\NormalTok{signedAttrs}\OperatorTok{,}\NormalTok{ CMSAttributes}\OperatorTok{.}\FunctionTok{contentType}\OperatorTok{,}
                    \StringTok{"contentType"}\OperatorTok{);}
            \FunctionTok{logAttribute}\OperatorTok{(}\NormalTok{signedAttrs}\OperatorTok{,}\NormalTok{ CMSAttributes}\OperatorTok{.}\FunctionTok{messageDigest}\OperatorTok{,}
                    \StringTok{"messageDigest"}\OperatorTok{);}
            \FunctionTok{logAttribute}\OperatorTok{(}\NormalTok{signedAttrs}\OperatorTok{,}\NormalTok{ CMSAttributes}\OperatorTok{.}\FunctionTok{signingTime}\OperatorTok{,}
                    \StringTok{"signingTime"}\OperatorTok{);}
            \FunctionTok{logAttribute}\OperatorTok{(}\NormalTok{signedAttrs}\OperatorTok{,}
\NormalTok{                    PKCSObjectIdentifiers}
                            \OperatorTok{.}\FunctionTok{id\_aa\_signingCertificateV2}\OperatorTok{,}
                    \StringTok{"ESS{-}signing{-}certificate{-}v2"}\OperatorTok{);}
        \OperatorTok{\}}

        \CommentTok{// Unsigned attributes (unauthenticated — timestamps live here)}
\NormalTok{        AttributeTable unsignedAttrs }\OperatorTok{=}\NormalTok{ signer}\OperatorTok{.}\FunctionTok{getUnsignedAttributes}\OperatorTok{();}
        \ControlFlowTok{if} \OperatorTok{(}\NormalTok{unsignedAttrs }\OperatorTok{!=} \KeywordTok{null}\OperatorTok{)} \OperatorTok{\{}
\NormalTok{            log}\OperatorTok{.}\FunctionTok{info}\OperatorTok{(}\StringTok{"Unsigned attributes count: \{\}"}\OperatorTok{,}
\NormalTok{                    unsignedAttrs}\OperatorTok{.}\FunctionTok{size}\OperatorTok{());}
            \FunctionTok{logAttribute}\OperatorTok{(}\NormalTok{unsignedAttrs}\OperatorTok{,}
\NormalTok{                    PKCSObjectIdentifiers}
                            \OperatorTok{.}\FunctionTok{id\_aa\_signatureTimeStampToken}\OperatorTok{,}
                    \StringTok{"signatureTimeStampToken"}\OperatorTok{);}
        \OperatorTok{\}}
    \OperatorTok{\}}

    \CommentTok{// Embedded certificates}
\NormalTok{    Store}\OperatorTok{\textless{}}\NormalTok{X509CertificateHolder}\OperatorTok{\textgreater{}}\NormalTok{ certStore }\OperatorTok{=}\NormalTok{ cmsData}\OperatorTok{.}\FunctionTok{getCertificates}\OperatorTok{();}
    \BuiltInTok{Collection}\OperatorTok{\textless{}}\NormalTok{X509CertificateHolder}\OperatorTok{\textgreater{}}\NormalTok{ certs }\OperatorTok{=}
\NormalTok{            certStore}\OperatorTok{.}\FunctionTok{getMatches}\OperatorTok{(}\KeywordTok{null}\OperatorTok{);}
\NormalTok{    log}\OperatorTok{.}\FunctionTok{info}\OperatorTok{(}\StringTok{"Embedded certificates: \{\}"}\OperatorTok{,}\NormalTok{ certs}\OperatorTok{.}\FunctionTok{size}\OperatorTok{());}
    \ControlFlowTok{for} \OperatorTok{(}\NormalTok{X509CertificateHolder cert }\OperatorTok{:}\NormalTok{ certs}\OperatorTok{)} \OperatorTok{\{}
\NormalTok{        log}\OperatorTok{.}\FunctionTok{info}\OperatorTok{(}\StringTok{"  Subject: \{\}"}\OperatorTok{,}\NormalTok{ cert}\OperatorTok{.}\FunctionTok{getSubject}\OperatorTok{());}
\NormalTok{        log}\OperatorTok{.}\FunctionTok{info}\OperatorTok{(}\StringTok{"  Issuer: \{\}"}\OperatorTok{,}\NormalTok{ cert}\OperatorTok{.}\FunctionTok{getIssuer}\OperatorTok{());}
\NormalTok{        log}\OperatorTok{.}\FunctionTok{info}\OperatorTok{(}\StringTok{"  Serial: \{\}"}\OperatorTok{,}\NormalTok{ cert}\OperatorTok{.}\FunctionTok{getSerialNumber}\OperatorTok{());}
\NormalTok{        log}\OperatorTok{.}\FunctionTok{info}\OperatorTok{(}\StringTok{"  Valid: \{\} to \{\}"}\OperatorTok{,}
\NormalTok{                cert}\OperatorTok{.}\FunctionTok{getNotBefore}\OperatorTok{(),}\NormalTok{ cert}\OperatorTok{.}\FunctionTok{getNotAfter}\OperatorTok{());}
    \OperatorTok{\}}
\OperatorTok{\}}

\KeywordTok{private} \DataTypeTok{void} \FunctionTok{logAttribute}\OperatorTok{(}\NormalTok{AttributeTable table}\OperatorTok{,}
\NormalTok{        ASN1ObjectIdentifier oid}\OperatorTok{,} \BuiltInTok{String}\NormalTok{ name}\OperatorTok{)} \OperatorTok{\{}
    \BuiltInTok{Attribute}\NormalTok{ attr }\OperatorTok{=}\NormalTok{ table}\OperatorTok{.}\FunctionTok{get}\OperatorTok{(}\NormalTok{oid}\OperatorTok{);}
    \ControlFlowTok{if} \OperatorTok{(}\NormalTok{attr }\OperatorTok{!=} \KeywordTok{null}\OperatorTok{)} \OperatorTok{\{}
\NormalTok{        log}\OperatorTok{.}\FunctionTok{info}\OperatorTok{(}\StringTok{"  \{\} present, values: \{\}"}\OperatorTok{,}
\NormalTok{                name}\OperatorTok{,}\NormalTok{ attr}\OperatorTok{.}\FunctionTok{getAttrValues}\OperatorTok{().}\FunctionTok{size}\OperatorTok{());}
    \OperatorTok{\}}
\OperatorTok{\}}
\end{Highlighting}
\end{Shaded}

\section{Byte Range Verification in
Code}\label{_byte_range_verification_in_code}

\subsection{The Three Checks}\label{_the_three_checks}

Byte range verification (Step 1 from Chapter 16) requires three checks:

\begin{enumerate}
\def\labelenumi{\arabic{enumi}.}
\item
  The \texttt{/ByteRange} array has exactly four elements
\item
  The first offset is zero (the signed region starts at the beginning of
  the file)
\item
  The second segment extends to the end of the file (or the end of the
  revision for multi-signature documents)
\end{enumerate}

\begin{Shaded}
\begin{Highlighting}[]
\CommentTok{/**}
 \CommentTok{*}\NormalTok{ Verify the structural integrity of the byte range}\CommentTok{.}
 \CommentTok{*}
 \CommentTok{*} \AnnotationTok{@param}\CommentTok{ }\AnnotationTok{byteRange}\CommentTok{   }\NormalTok{the four}\CommentTok{{-}}\NormalTok{element ByteRange array}
 \CommentTok{*} \AnnotationTok{@param}\CommentTok{ }\AnnotationTok{fileSize}\CommentTok{    }\NormalTok{the total size of the PDF file in bytes}
 \CommentTok{*} \AnnotationTok{@return}\NormalTok{ a ValidationResult indicating success or the specific failure}
 \CommentTok{*/}
\KeywordTok{public}\NormalTok{ ValidationResult }\FunctionTok{verifyByteRange}\OperatorTok{(}\DataTypeTok{int}\OperatorTok{[]}\NormalTok{ byteRange}\OperatorTok{,} \DataTypeTok{long}\NormalTok{ fileSize}\OperatorTok{)} \OperatorTok{\{}
    \ControlFlowTok{if} \OperatorTok{(}\NormalTok{byteRange }\OperatorTok{==} \KeywordTok{null} \OperatorTok{||}\NormalTok{ byteRange}\OperatorTok{.}\FunctionTok{length} \OperatorTok{!=} \DecValTok{4}\OperatorTok{)} \OperatorTok{\{}
\NormalTok{        log}\OperatorTok{.}\FunctionTok{error}\OperatorTok{(}\StringTok{"ByteRange is null or does not have exactly 4 elements"}\OperatorTok{);}
        \ControlFlowTok{return}\NormalTok{ ValidationResult}\OperatorTok{.}\FunctionTok{BYTE\_RANGE\_MALFORMED}\OperatorTok{;}
    \OperatorTok{\}}

    \DataTypeTok{int}\NormalTok{ offset1 }\OperatorTok{=}\NormalTok{ byteRange}\OperatorTok{[}\DecValTok{0}\OperatorTok{];}
    \DataTypeTok{int}\NormalTok{ length1 }\OperatorTok{=}\NormalTok{ byteRange}\OperatorTok{[}\DecValTok{1}\OperatorTok{];}
    \DataTypeTok{int}\NormalTok{ offset2 }\OperatorTok{=}\NormalTok{ byteRange}\OperatorTok{[}\DecValTok{2}\OperatorTok{];}
    \DataTypeTok{int}\NormalTok{ length2 }\OperatorTok{=}\NormalTok{ byteRange}\OperatorTok{[}\DecValTok{3}\OperatorTok{];}

    \CommentTok{// Check 1: First segment starts at byte 0}
    \ControlFlowTok{if} \OperatorTok{(}\NormalTok{offset1 }\OperatorTok{!=} \DecValTok{0}\OperatorTok{)} \OperatorTok{\{}
\NormalTok{        log}\OperatorTok{.}\FunctionTok{error}\OperatorTok{(}\StringTok{"ByteRange offset1 is \{\} — expected 0"}\OperatorTok{,}\NormalTok{ offset1}\OperatorTok{);}
        \ControlFlowTok{return}\NormalTok{ ValidationResult}\OperatorTok{.}\FunctionTok{BYTE\_RANGE\_NOT\_AT\_START}\OperatorTok{;}
    \OperatorTok{\}}

    \CommentTok{// Check 2: Segments don\textquotesingle{}t overlap}
    \ControlFlowTok{if} \OperatorTok{(}\NormalTok{offset2 }\OperatorTok{\textless{}}\NormalTok{ length1}\OperatorTok{)} \OperatorTok{\{}
\NormalTok{        log}\OperatorTok{.}\FunctionTok{error}\OperatorTok{(}\StringTok{"ByteRange segments overlap: "} \OperatorTok{+}
                \StringTok{"segment1 ends at \{\}, segment2 starts at \{\}"}\OperatorTok{,}
\NormalTok{                length1}\OperatorTok{,}\NormalTok{ offset2}\OperatorTok{);}
        \ControlFlowTok{return}\NormalTok{ ValidationResult}\OperatorTok{.}\FunctionTok{BYTE\_RANGE\_OVERLAP}\OperatorTok{;}
    \OperatorTok{\}}

    \CommentTok{// Check 3: Signed region extends to end of file}
    \DataTypeTok{long}\NormalTok{ signedEnd }\OperatorTok{=} \OperatorTok{(}\DataTypeTok{long}\OperatorTok{)}\NormalTok{ offset2 }\OperatorTok{+}\NormalTok{ length2}\OperatorTok{;}
    \ControlFlowTok{if} \OperatorTok{(}\NormalTok{signedEnd }\OperatorTok{!=}\NormalTok{ fileSize}\OperatorTok{)} \OperatorTok{\{}
        \DataTypeTok{long}\NormalTok{ trailingBytes }\OperatorTok{=}\NormalTok{ fileSize }\OperatorTok{{-}}\NormalTok{ signedEnd}\OperatorTok{;}
\NormalTok{        log}\OperatorTok{.}\FunctionTok{warn}\OperatorTok{(}\StringTok{"ByteRange covers \{\} bytes but file is \{\} bytes — "} \OperatorTok{+}
                \StringTok{"\{\} trailing unsigned bytes detected"}\OperatorTok{,}
\NormalTok{                signedEnd}\OperatorTok{,}\NormalTok{ fileSize}\OperatorTok{,}\NormalTok{ trailingBytes}\OperatorTok{);}
        \ControlFlowTok{return}\NormalTok{ ValidationResult}\OperatorTok{.}\FunctionTok{BYTE\_RANGE\_TRAILING\_BYTES}\OperatorTok{;}
    \OperatorTok{\}}

\NormalTok{    log}\OperatorTok{.}\FunctionTok{info}\OperatorTok{(}\StringTok{"ByteRange verified: [\{\} \{\} \{\} \{\}], "} \OperatorTok{+}
            \StringTok{"file size \{\}, no trailing bytes"}\OperatorTok{,}
\NormalTok{            offset1}\OperatorTok{,}\NormalTok{ length1}\OperatorTok{,}\NormalTok{ offset2}\OperatorTok{,}\NormalTok{ length2}\OperatorTok{,}\NormalTok{ fileSize}\OperatorTok{);}
    \ControlFlowTok{return}\NormalTok{ ValidationResult}\OperatorTok{.}\FunctionTok{OK}\OperatorTok{;}
\OperatorTok{\}}
\end{Highlighting}
\end{Shaded}

Trailing bytes are not always an attack. In a multi-signature document,
each signature covers only up to the end of its own incremental
revision. Subsequent revisions (containing later signatures) appear as
"trailing bytes" from the perspective of the earlier signature. Chapter
8 covers how incremental updates work, and Chapter 13 covers the byte
range geometry in multi-signature scenarios.

However, trailing bytes on the \emph{last} signature in a document are
suspicious and may indicate a Shadow Attack (Chapter 24). Your
verification code should distinguish between "trailing bytes that are a
legitimate later revision" and "trailing bytes that are unexplained."

\subsection{Validating the Contents
Gap}\label{_validating_the_contents_gap}

The gap between the two segments should contain exactly the
\texttt{/Contents} hex string. You can validate this by checking that
the gap bytes, when decoded from hex, produce valid DER-encoded data:

\begin{Shaded}
\begin{Highlighting}[]
\CommentTok{/**}
 \CommentTok{*}\NormalTok{ Validate that the gap between byte range segments contains}
 \CommentTok{*}\NormalTok{ a properly formatted hex string with valid CMS data}\CommentTok{.}
 \CommentTok{*/}
\KeywordTok{public} \DataTypeTok{boolean} \FunctionTok{validateContentsGap}\OperatorTok{(}\DataTypeTok{byte}\OperatorTok{[]}\NormalTok{ pdfBytes}\OperatorTok{,} \DataTypeTok{int}\OperatorTok{[]}\NormalTok{ byteRange}\OperatorTok{)} \OperatorTok{\{}
    \DataTypeTok{int}\NormalTok{ gapStart }\OperatorTok{=}\NormalTok{ byteRange}\OperatorTok{[}\DecValTok{1}\OperatorTok{];}
    \DataTypeTok{int}\NormalTok{ gapEnd }\OperatorTok{=}\NormalTok{ byteRange}\OperatorTok{[}\DecValTok{2}\OperatorTok{];}
    \DataTypeTok{int}\NormalTok{ gapSize }\OperatorTok{=}\NormalTok{ gapEnd }\OperatorTok{{-}}\NormalTok{ gapStart}\OperatorTok{;}

    \CommentTok{// The gap should start with \textquotesingle{}\textless{}\textquotesingle{} and end with \textquotesingle{}\textgreater{}\textquotesingle{}}
    \ControlFlowTok{if} \OperatorTok{(}\NormalTok{pdfBytes}\OperatorTok{[}\NormalTok{gapStart}\OperatorTok{]} \OperatorTok{!=} \CharTok{\textquotesingle{}\textless{}\textquotesingle{}} \OperatorTok{||}\NormalTok{ pdfBytes}\OperatorTok{[}\NormalTok{gapEnd }\OperatorTok{{-}} \DecValTok{1}\OperatorTok{]} \OperatorTok{!=} \CharTok{\textquotesingle{}\textgreater{}\textquotesingle{}}\OperatorTok{)} \OperatorTok{\{}
\NormalTok{        log}\OperatorTok{.}\FunctionTok{error}\OperatorTok{(}\StringTok{"Gap does not start with \textquotesingle{}\textless{}\textquotesingle{} or end with \textquotesingle{}\textgreater{}\textquotesingle{}"}\OperatorTok{);}
        \ControlFlowTok{return} \KeywordTok{false}\OperatorTok{;}
    \OperatorTok{\}}

    \CommentTok{// Extract hex content (skip the delimiters)}
    \BuiltInTok{String}\NormalTok{ hexContent }\OperatorTok{=} \KeywordTok{new} \BuiltInTok{String}\OperatorTok{(}\NormalTok{pdfBytes}\OperatorTok{,}\NormalTok{ gapStart }\OperatorTok{+} \DecValTok{1}\OperatorTok{,}
\NormalTok{            gapSize }\OperatorTok{{-}} \DecValTok{2}\OperatorTok{,}\NormalTok{ StandardCharsets}\OperatorTok{.}\FunctionTok{US\_ASCII}\OperatorTok{).}\FunctionTok{trim}\OperatorTok{();}

    \CommentTok{// Verify it\textquotesingle{}s valid hex}
    \ControlFlowTok{if} \OperatorTok{(!}\NormalTok{hexContent}\OperatorTok{.}\FunctionTok{matches}\OperatorTok{(}\StringTok{"[0{-}9a{-}fA{-}F]+"}\OperatorTok{))} \OperatorTok{\{}
\NormalTok{        log}\OperatorTok{.}\FunctionTok{error}\OperatorTok{(}\StringTok{"Gap contains non{-}hex characters"}\OperatorTok{);}
        \ControlFlowTok{return} \KeywordTok{false}\OperatorTok{;}
    \OperatorTok{\}}

\NormalTok{    log}\OperatorTok{.}\FunctionTok{info}\OperatorTok{(}\StringTok{"Contents gap validated: \{\} hex chars = \{\} bytes CMS"}\OperatorTok{,}
\NormalTok{            hexContent}\OperatorTok{.}\FunctionTok{length}\OperatorTok{(),}\NormalTok{ hexContent}\OperatorTok{.}\FunctionTok{length}\OperatorTok{()} \OperatorTok{/} \DecValTok{2}\OperatorTok{);}
    \ControlFlowTok{return} \KeywordTok{true}\OperatorTok{;}
\OperatorTok{\}}
\end{Highlighting}
\end{Shaded}

\section{CMS Cryptographic
Verification}\label{_cms_cryptographic_verification}

\subsection{Building the
SignerInformationVerifier}\label{_building_the_signerinformationverifier}

This is Step 3 from Chapter 16 --- verifying that the CMS
\texttt{SignedData} was actually created by the claimed signer's private
key. BouncyCastle's \texttt{SignerInformation.verify()} method performs
this check, but you need to provide the correct
\texttt{SignerInformationVerifier}:

\begin{Shaded}
\begin{Highlighting}[]
\CommentTok{/**}
 \CommentTok{*}\NormalTok{ Verify the CMS signature cryptographically}\CommentTok{.}
 \CommentTok{*}
 \CommentTok{*} \AnnotationTok{@param}\CommentTok{ }\AnnotationTok{cmsData}\CommentTok{  }\NormalTok{the parsed CMS SignedData }\CommentTok{(}\NormalTok{with signed content}\CommentTok{)}
 \CommentTok{*} \AnnotationTok{@return}\NormalTok{ the signer}\CommentTok{\textquotesingle{}}\NormalTok{s certificate if verification succeeds}\CommentTok{,}\NormalTok{ null otherwise}
 \CommentTok{*/}
\KeywordTok{public}\NormalTok{ X509CertificateHolder }\FunctionTok{verifyCmsSignature}\OperatorTok{(}\NormalTok{CMSSignedData cmsData}\OperatorTok{)}
        \KeywordTok{throws} \BuiltInTok{Exception} \OperatorTok{\{}

\NormalTok{    SignerInformationStore signerStore }\OperatorTok{=}\NormalTok{ cmsData}\OperatorTok{.}\FunctionTok{getSignerInfos}\OperatorTok{();}
    \ControlFlowTok{if} \OperatorTok{(}\NormalTok{signerStore}\OperatorTok{.}\FunctionTok{size}\OperatorTok{()} \OperatorTok{==} \DecValTok{0}\OperatorTok{)} \OperatorTok{\{}
\NormalTok{        log}\OperatorTok{.}\FunctionTok{error}\OperatorTok{(}\StringTok{"CMS contains no signers"}\OperatorTok{);}
        \ControlFlowTok{return} \KeywordTok{null}\OperatorTok{;}
    \OperatorTok{\}}

\NormalTok{    SignerInformation signerInfo }\OperatorTok{=}
\NormalTok{            signerStore}\OperatorTok{.}\FunctionTok{getSigners}\OperatorTok{().}\FunctionTok{iterator}\OperatorTok{().}\FunctionTok{next}\OperatorTok{();}

    \CommentTok{// Find the signer\textquotesingle{}s certificate by matching the SignerIdentifier}
\NormalTok{    Store}\OperatorTok{\textless{}}\NormalTok{X509CertificateHolder}\OperatorTok{\textgreater{}}\NormalTok{ certStore }\OperatorTok{=}\NormalTok{ cmsData}\OperatorTok{.}\FunctionTok{getCertificates}\OperatorTok{();}
    \AttributeTok{@SuppressWarnings}\OperatorTok{(}\StringTok{"unchecked"}\OperatorTok{)}
    \BuiltInTok{Collection}\OperatorTok{\textless{}}\NormalTok{X509CertificateHolder}\OperatorTok{\textgreater{}}\NormalTok{ matches }\OperatorTok{=}
\NormalTok{            certStore}\OperatorTok{.}\FunctionTok{getMatches}\OperatorTok{(}\NormalTok{signerInfo}\OperatorTok{.}\FunctionTok{getSID}\OperatorTok{());}

    \ControlFlowTok{if} \OperatorTok{(}\NormalTok{matches}\OperatorTok{.}\FunctionTok{isEmpty}\OperatorTok{())} \OperatorTok{\{}
\NormalTok{        log}\OperatorTok{.}\FunctionTok{error}\OperatorTok{(}\StringTok{"No certificate matches the SignerIdentifier — "} \OperatorTok{+}
                \StringTok{"the signer\textquotesingle{}s cert may not be embedded in the CMS"}\OperatorTok{);}
        \ControlFlowTok{return} \KeywordTok{null}\OperatorTok{;}
    \OperatorTok{\}}

\NormalTok{    X509CertificateHolder signerCert }\OperatorTok{=}\NormalTok{ matches}\OperatorTok{.}\FunctionTok{iterator}\OperatorTok{().}\FunctionTok{next}\OperatorTok{();}
\NormalTok{    log}\OperatorTok{.}\FunctionTok{info}\OperatorTok{(}\StringTok{"Signer certificate: subject=\{\}, issuer=\{\}"}\OperatorTok{,}
\NormalTok{            signerCert}\OperatorTok{.}\FunctionTok{getSubject}\OperatorTok{(),}\NormalTok{ signerCert}\OperatorTok{.}\FunctionTok{getIssuer}\OperatorTok{());}

    \CommentTok{// Build the verifier and verify}
\NormalTok{    SignerInformationVerifier verifier }\OperatorTok{=}
            \KeywordTok{new} \FunctionTok{JcaSimpleSignerInfoVerifierBuilder}\OperatorTok{()}
                    \OperatorTok{.}\FunctionTok{setProvider}\OperatorTok{(}\StringTok{"BC"}\OperatorTok{)}
                    \OperatorTok{.}\FunctionTok{build}\OperatorTok{(}\NormalTok{signerCert}\OperatorTok{);}

    \DataTypeTok{boolean}\NormalTok{ valid }\OperatorTok{=}\NormalTok{ signerInfo}\OperatorTok{.}\FunctionTok{verify}\OperatorTok{(}\NormalTok{verifier}\OperatorTok{);}
    \ControlFlowTok{if} \OperatorTok{(}\NormalTok{valid}\OperatorTok{)} \OperatorTok{\{}
\NormalTok{        log}\OperatorTok{.}\FunctionTok{info}\OperatorTok{(}\StringTok{"CMS signature verification PASSED"}\OperatorTok{);}
        \ControlFlowTok{return}\NormalTok{ signerCert}\OperatorTok{;}
    \OperatorTok{\}} \ControlFlowTok{else} \OperatorTok{\{}
\NormalTok{        log}\OperatorTok{.}\FunctionTok{error}\OperatorTok{(}\StringTok{"CMS signature verification FAILED — "} \OperatorTok{+}
                \StringTok{"signature was not created by the claimed certificate"}\OperatorTok{);}
        \ControlFlowTok{return} \KeywordTok{null}\OperatorTok{;}
    \OperatorTok{\}}
\OperatorTok{\}}
\end{Highlighting}
\end{Shaded}

\subsection{What verify() Actually
Does}\label{_what_verify_actually_does}

When you call \texttt{signerInfo.verify(verifier)}, BouncyCastle
performs the following internally:

\begin{enumerate}
\def\labelenumi{\arabic{enumi}.}
\item
  Computes the digest of the signed content (the byte range data you
  passed when constructing \texttt{CMSSignedData}) using the algorithm
  from \texttt{SignerInfo.digestAlgorithm}
\item
  Checks that this computed digest matches the \texttt{messageDigest}
  attribute in the signed attributes
\item
  DER-encodes the signed attributes
\item
  Verifies the RSA/ECDSA/DSA signature over the DER-encoded signed
  attributes using the public key from the signer's certificate
\end{enumerate}

If signed attributes are absent (which is technically allowed by CMS but
rare in PDF signatures), BouncyCastle instead verifies the signature
directly over the content digest.

The \texttt{verify()} method performs only the cryptographic
verification --- Steps 2 and 3 combined. It does not check certificate
validity, revocation status, or trust. Those are separate steps you must
implement yourself.

\section{Certificate Chain Reconstruction and
Validation}\label{_certificate_chain_reconstruction_and_validation}

\subsection{The Java PKIX CertPath
API}\label{_the_java_pkix_certpath_api}

Java's \texttt{java.security.cert} package provides a complete PKIX
certificate path validation engine via \texttt{CertPathBuilder} and
\texttt{CertPathValidator}. This is the standard way to validate
certificate chains in Java, and it handles the complexity of chain
building, validity checking, basic constraints verification, and key
usage checking.

\begin{Shaded}
\begin{Highlighting}[]
\CommentTok{/**}
 \CommentTok{*}\NormalTok{ Build and validate the certificate chain from the signer}\CommentTok{\textquotesingle{}}\NormalTok{s certificate}
 \CommentTok{*}\NormalTok{ to a trusted root}\CommentTok{.}
 \CommentTok{*}
 \CommentTok{*} \AnnotationTok{@param}\CommentTok{ }\AnnotationTok{signerCert}\CommentTok{    }\NormalTok{the signer}\CommentTok{\textquotesingle{}}\NormalTok{s certificate }\CommentTok{(}\NormalTok{from CMS verification}\CommentTok{)}
 \CommentTok{*} \AnnotationTok{@param}\CommentTok{ }\AnnotationTok{cmsData}\CommentTok{       }\NormalTok{the CMS containing intermediate certificates}
 \CommentTok{*} \AnnotationTok{@param}\CommentTok{ }\AnnotationTok{trustAnchors}\CommentTok{  }\NormalTok{the set of trusted root certificates}
 \CommentTok{*} \AnnotationTok{@param}\CommentTok{ }\AnnotationTok{signingTime}\CommentTok{   }\NormalTok{the claimed signing time }\CommentTok{(}\NormalTok{for validity checks}\CommentTok{)}
 \CommentTok{*} \AnnotationTok{@return}\NormalTok{ the validated CertPath}\CommentTok{,}\NormalTok{ or null if validation fails}
 \CommentTok{*/}
\KeywordTok{public} \BuiltInTok{CertPath} \FunctionTok{validateCertificateChain}\OperatorTok{(}
\NormalTok{        X509CertificateHolder signerCert}\OperatorTok{,}
\NormalTok{        CMSSignedData cmsData}\OperatorTok{,}
        \BuiltInTok{Set}\OperatorTok{\textless{}}\BuiltInTok{TrustAnchor}\OperatorTok{\textgreater{}}\NormalTok{ trustAnchors}\OperatorTok{,}
        \BuiltInTok{Date}\NormalTok{ signingTime}\OperatorTok{)} \KeywordTok{throws} \BuiltInTok{Exception} \OperatorTok{\{}

    \CommentTok{// Convert BouncyCastle cert holders to Java X509Certificates}
    \BuiltInTok{CertificateFactory}\NormalTok{ cf }\OperatorTok{=} \BuiltInTok{CertificateFactory}\OperatorTok{.}\FunctionTok{getInstance}\OperatorTok{(}\StringTok{"X.509"}\OperatorTok{);}
\NormalTok{    JcaX509CertificateConverter converter }\OperatorTok{=}
            \KeywordTok{new} \FunctionTok{JcaX509CertificateConverter}\OperatorTok{().}\FunctionTok{setProvider}\OperatorTok{(}\StringTok{"BC"}\OperatorTok{);}

    \CommentTok{// Collect all certificates from CMS into a CertStore}
    \BuiltInTok{List}\OperatorTok{\textless{}}\BuiltInTok{X509Certificate}\OperatorTok{\textgreater{}}\NormalTok{ allCerts }\OperatorTok{=} \KeywordTok{new} \BuiltInTok{ArrayList}\OperatorTok{\textless{}\textgreater{}();}
\NormalTok{    Store}\OperatorTok{\textless{}}\NormalTok{X509CertificateHolder}\OperatorTok{\textgreater{}}\NormalTok{ bcStore }\OperatorTok{=}\NormalTok{ cmsData}\OperatorTok{.}\FunctionTok{getCertificates}\OperatorTok{();}
    \ControlFlowTok{for} \OperatorTok{(}\NormalTok{X509CertificateHolder holder }\OperatorTok{:}\NormalTok{ bcStore}\OperatorTok{.}\FunctionTok{getMatches}\OperatorTok{(}\KeywordTok{null}\OperatorTok{))} \OperatorTok{\{}
\NormalTok{        allCerts}\OperatorTok{.}\FunctionTok{add}\OperatorTok{(}\NormalTok{converter}\OperatorTok{.}\FunctionTok{getCertificate}\OperatorTok{(}\NormalTok{holder}\OperatorTok{));}
    \OperatorTok{\}}

    \BuiltInTok{CertStore}\NormalTok{ certStore }\OperatorTok{=} \BuiltInTok{CertStore}\OperatorTok{.}\FunctionTok{getInstance}\OperatorTok{(}\StringTok{"Collection"}\OperatorTok{,}
            \KeywordTok{new} \BuiltInTok{CollectionCertStoreParameters}\OperatorTok{(}\NormalTok{allCerts}\OperatorTok{));}

    \CommentTok{// Configure the target certificate selector}
    \BuiltInTok{X509Certificate}\NormalTok{ targetCert }\OperatorTok{=}\NormalTok{ converter}\OperatorTok{.}\FunctionTok{getCertificate}\OperatorTok{(}\NormalTok{signerCert}\OperatorTok{);}
    \BuiltInTok{X509CertSelector}\NormalTok{ targetSelector }\OperatorTok{=} \KeywordTok{new} \BuiltInTok{X509CertSelector}\OperatorTok{();}
\NormalTok{    targetSelector}\OperatorTok{.}\FunctionTok{setCertificate}\OperatorTok{(}\NormalTok{targetCert}\OperatorTok{);}

    \CommentTok{// Build PKIX parameters}
    \BuiltInTok{PKIXBuilderParameters}\NormalTok{ params }\OperatorTok{=}
            \KeywordTok{new} \BuiltInTok{PKIXBuilderParameters}\OperatorTok{(}\NormalTok{trustAnchors}\OperatorTok{,}\NormalTok{ targetSelector}\OperatorTok{);}
\NormalTok{    params}\OperatorTok{.}\FunctionTok{addCertStore}\OperatorTok{(}\NormalTok{certStore}\OperatorTok{);}
\NormalTok{    params}\OperatorTok{.}\FunctionTok{setRevocationEnabled}\OperatorTok{(}\KeywordTok{false}\OperatorTok{);} \CommentTok{// We handle revocation separately}
\NormalTok{    params}\OperatorTok{.}\FunctionTok{setDate}\OperatorTok{(}\NormalTok{signingTime}\OperatorTok{);}        \CommentTok{// Validate at signing time}

    \CommentTok{// Build the chain}
    \BuiltInTok{CertPathBuilder}\NormalTok{ builder }\OperatorTok{=}
            \BuiltInTok{CertPathBuilder}\OperatorTok{.}\FunctionTok{getInstance}\OperatorTok{(}\StringTok{"PKIX"}\OperatorTok{,} \StringTok{"BC"}\OperatorTok{);}
    \BuiltInTok{PKIXCertPathBuilderResult}\NormalTok{ result }\OperatorTok{=}
            \OperatorTok{(}\BuiltInTok{PKIXCertPathBuilderResult}\OperatorTok{)}\NormalTok{ builder}\OperatorTok{.}\FunctionTok{build}\OperatorTok{(}\NormalTok{params}\OperatorTok{);}

    \BuiltInTok{CertPath}\NormalTok{ certPath }\OperatorTok{=}\NormalTok{ result}\OperatorTok{.}\FunctionTok{getCertPath}\OperatorTok{();}
\NormalTok{    log}\OperatorTok{.}\FunctionTok{info}\OperatorTok{(}\StringTok{"Certificate chain built successfully — \{\} certificates"}\OperatorTok{,}
\NormalTok{            certPath}\OperatorTok{.}\FunctionTok{getCertificates}\OperatorTok{().}\FunctionTok{size}\OperatorTok{());}

    \ControlFlowTok{for} \OperatorTok{(}\DataTypeTok{int}\NormalTok{ i }\OperatorTok{=} \DecValTok{0}\OperatorTok{;}\NormalTok{ i }\OperatorTok{\textless{}}\NormalTok{ certPath}\OperatorTok{.}\FunctionTok{getCertificates}\OperatorTok{().}\FunctionTok{size}\OperatorTok{();}\NormalTok{ i}\OperatorTok{++)} \OperatorTok{\{}
        \BuiltInTok{X509Certificate}\NormalTok{ cert }\OperatorTok{=}
                \OperatorTok{(}\BuiltInTok{X509Certificate}\OperatorTok{)}\NormalTok{ certPath}\OperatorTok{.}\FunctionTok{getCertificates}\OperatorTok{().}\FunctionTok{get}\OperatorTok{(}\NormalTok{i}\OperatorTok{);}
\NormalTok{        log}\OperatorTok{.}\FunctionTok{info}\OperatorTok{(}\StringTok{"  [\{\}] Subject: \{\}, Issuer: \{\}"}\OperatorTok{,}
\NormalTok{                i}\OperatorTok{,}\NormalTok{ cert}\OperatorTok{.}\FunctionTok{getSubjectX500Principal}\OperatorTok{(),}
\NormalTok{                cert}\OperatorTok{.}\FunctionTok{getIssuerX500Principal}\OperatorTok{());}
    \OperatorTok{\}}

\NormalTok{    log}\OperatorTok{.}\FunctionTok{info}\OperatorTok{(}\StringTok{"  [root] Trust anchor: \{\}"}\OperatorTok{,}
\NormalTok{            result}\OperatorTok{.}\FunctionTok{getTrustAnchor}\OperatorTok{().}\FunctionTok{getTrustedCert}\OperatorTok{()}
                    \OperatorTok{.}\FunctionTok{getSubjectX500Principal}\OperatorTok{());}

    \ControlFlowTok{return}\NormalTok{ certPath}\OperatorTok{;}
\OperatorTok{\}}
\end{Highlighting}
\end{Shaded}

\subsection{Building the Trust Anchor
Set}\label{_building_the_trust_anchor_set}

The \texttt{TrustAnchor} set represents your trust store --- the root
certificates you are willing to trust. In a real application, this might
come from the Java default trust store (\texttt{cacerts}), a custom
keystore, or the AATL/EUTL (see Chapter 16, Step 5).

\begin{Shaded}
\begin{Highlighting}[]
\CommentTok{/**}
 \CommentTok{*}\NormalTok{ Load trust anchors from a Java KeyStore }\CommentTok{(}\NormalTok{JKS or PKCS12}\CommentTok{).}
 \CommentTok{*/}
\KeywordTok{public} \BuiltInTok{Set}\OperatorTok{\textless{}}\BuiltInTok{TrustAnchor}\OperatorTok{\textgreater{}} \FunctionTok{loadTrustAnchors}\OperatorTok{(}\NormalTok{Path keystorePath}\OperatorTok{,}
        \DataTypeTok{char}\OperatorTok{[]}\NormalTok{ password}\OperatorTok{)} \KeywordTok{throws} \BuiltInTok{Exception} \OperatorTok{\{}

    \BuiltInTok{KeyStore}\NormalTok{ trustStore }\OperatorTok{=} \BuiltInTok{KeyStore}\OperatorTok{.}\FunctionTok{getInstance}\OperatorTok{(}
\NormalTok{            keystorePath}\OperatorTok{.}\FunctionTok{toString}\OperatorTok{().}\FunctionTok{endsWith}\OperatorTok{(}\StringTok{".p12"}\OperatorTok{)}
                    \OperatorTok{?} \StringTok{"PKCS12"} \OperatorTok{:} \StringTok{"JKS"}\OperatorTok{);}
    \ControlFlowTok{try} \OperatorTok{(}\BuiltInTok{InputStream}\NormalTok{ is }\OperatorTok{=}\NormalTok{ Files}\OperatorTok{.}\FunctionTok{newInputStream}\OperatorTok{(}\NormalTok{keystorePath}\OperatorTok{))} \OperatorTok{\{}
\NormalTok{        trustStore}\OperatorTok{.}\FunctionTok{load}\OperatorTok{(}\NormalTok{is}\OperatorTok{,}\NormalTok{ password}\OperatorTok{);}
    \OperatorTok{\}}

    \BuiltInTok{Set}\OperatorTok{\textless{}}\BuiltInTok{TrustAnchor}\OperatorTok{\textgreater{}}\NormalTok{ anchors }\OperatorTok{=} \KeywordTok{new} \BuiltInTok{HashSet}\OperatorTok{\textless{}\textgreater{}();}
    \BuiltInTok{Enumeration}\OperatorTok{\textless{}}\BuiltInTok{String}\OperatorTok{\textgreater{}}\NormalTok{ aliases }\OperatorTok{=}\NormalTok{ trustStore}\OperatorTok{.}\FunctionTok{aliases}\OperatorTok{();}
    \ControlFlowTok{while} \OperatorTok{(}\NormalTok{aliases}\OperatorTok{.}\FunctionTok{hasMoreElements}\OperatorTok{())} \OperatorTok{\{}
        \BuiltInTok{String}\NormalTok{ alias }\OperatorTok{=}\NormalTok{ aliases}\OperatorTok{.}\FunctionTok{nextElement}\OperatorTok{();}
        \ControlFlowTok{if} \OperatorTok{(}\NormalTok{trustStore}\OperatorTok{.}\FunctionTok{isCertificateEntry}\OperatorTok{(}\NormalTok{alias}\OperatorTok{))} \OperatorTok{\{}
            \BuiltInTok{X509Certificate}\NormalTok{ cert }\OperatorTok{=}
                    \OperatorTok{(}\BuiltInTok{X509Certificate}\OperatorTok{)}\NormalTok{ trustStore}\OperatorTok{.}\FunctionTok{getCertificate}\OperatorTok{(}\NormalTok{alias}\OperatorTok{);}
\NormalTok{            anchors}\OperatorTok{.}\FunctionTok{add}\OperatorTok{(}\KeywordTok{new} \BuiltInTok{TrustAnchor}\OperatorTok{(}\NormalTok{cert}\OperatorTok{,} \KeywordTok{null}\OperatorTok{));}
\NormalTok{            log}\OperatorTok{.}\FunctionTok{info}\OperatorTok{(}\StringTok{"Loaded trust anchor: \{\}"}\OperatorTok{,}
\NormalTok{                    cert}\OperatorTok{.}\FunctionTok{getSubjectX500Principal}\OperatorTok{());}
        \OperatorTok{\}}
    \OperatorTok{\}}

\NormalTok{    log}\OperatorTok{.}\FunctionTok{info}\OperatorTok{(}\StringTok{"Loaded \{\} trust anchors from \{\}"}\OperatorTok{,}\NormalTok{ anchors}\OperatorTok{.}\FunctionTok{size}\OperatorTok{(),}
\NormalTok{            keystorePath}\OperatorTok{);}
    \ControlFlowTok{return}\NormalTok{ anchors}\OperatorTok{;}
\OperatorTok{\}}
\end{Highlighting}
\end{Shaded}

\subsection{Key Usage Validation}\label{_key_usage_validation}

The PKIX \texttt{CertPathValidator} checks basic constraints and key
usage automatically. But if you want to verify specific requirements ---
for example, that the signer's certificate has the
\texttt{digitalSignature} or \texttt{nonRepudiation} key usage --- you
can check explicitly:

\begin{Shaded}
\begin{Highlighting}[]
\CommentTok{/**}
 \CommentTok{*}\NormalTok{ Check that the signer}\CommentTok{\textquotesingle{}}\NormalTok{s certificate permits digital signing}\CommentTok{.}
 \CommentTok{*/}
\KeywordTok{public} \DataTypeTok{boolean} \FunctionTok{hasSigningKeyUsage}\OperatorTok{(}\BuiltInTok{X509Certificate}\NormalTok{ cert}\OperatorTok{)} \OperatorTok{\{}
    \DataTypeTok{boolean}\OperatorTok{[]}\NormalTok{ keyUsage }\OperatorTok{=}\NormalTok{ cert}\OperatorTok{.}\FunctionTok{getKeyUsage}\OperatorTok{();}
    \ControlFlowTok{if} \OperatorTok{(}\NormalTok{keyUsage }\OperatorTok{==} \KeywordTok{null}\OperatorTok{)} \OperatorTok{\{}
\NormalTok{        log}\OperatorTok{.}\FunctionTok{warn}\OperatorTok{(}\StringTok{"Certificate has no KeyUsage extension — "} \OperatorTok{+}
                \StringTok{"signing permitted by default"}\OperatorTok{);}
        \ControlFlowTok{return} \KeywordTok{true}\OperatorTok{;}
    \OperatorTok{\}}

    \CommentTok{// bit 0 = digitalSignature, bit 1 = nonRepudiation (contentCommitment)}
    \DataTypeTok{boolean}\NormalTok{ digitalSignature }\OperatorTok{=}\NormalTok{ keyUsage}\OperatorTok{[}\DecValTok{0}\OperatorTok{];}
    \DataTypeTok{boolean}\NormalTok{ nonRepudiation }\OperatorTok{=}\NormalTok{ keyUsage}\OperatorTok{[}\DecValTok{1}\OperatorTok{];}

    \ControlFlowTok{if} \OperatorTok{(!}\NormalTok{digitalSignature }\OperatorTok{\&\&} \OperatorTok{!}\NormalTok{nonRepudiation}\OperatorTok{)} \OperatorTok{\{}
\NormalTok{        log}\OperatorTok{.}\FunctionTok{error}\OperatorTok{(}\StringTok{"Certificate KeyUsage does not include "} \OperatorTok{+}
                \StringTok{"digitalSignature or nonRepudiation"}\OperatorTok{);}
        \ControlFlowTok{return} \KeywordTok{false}\OperatorTok{;}
    \OperatorTok{\}}

\NormalTok{    log}\OperatorTok{.}\FunctionTok{info}\OperatorTok{(}\StringTok{"Certificate KeyUsage: digitalSignature=\{\}, "} \OperatorTok{+}
            \StringTok{"nonRepudiation=\{\}"}\OperatorTok{,}\NormalTok{ digitalSignature}\OperatorTok{,}\NormalTok{ nonRepudiation}\OperatorTok{);}
    \ControlFlowTok{return} \KeywordTok{true}\OperatorTok{;}
\OperatorTok{\}}
\end{Highlighting}
\end{Shaded}

\section{Timestamp Verification}\label{_timestamp_verification_2}

\subsection{Extracting the Timestamp
Token}\label{_extracting_the_timestamp_token}

A signature timestamp is an unsigned CMS attribute with OID
\texttt{1.2.840.113549.1.9.16.2.14}
(\texttt{id-aa-signatureTimeStampToken}). The attribute value is itself
a CMS \texttt{SignedData} structure --- a timestamp is a signature on a
signature:

\begin{Shaded}
\begin{Highlighting}[]
\CommentTok{/**}
 \CommentTok{*}\NormalTok{ Extract and verify the signature timestamp}\CommentTok{,}\NormalTok{ if present}\CommentTok{.}
 \CommentTok{*}
 \CommentTok{*} \AnnotationTok{@param}\CommentTok{ }\AnnotationTok{signerInfo}\CommentTok{  }\NormalTok{the CMS SignerInformation}
 \CommentTok{*} \AnnotationTok{@return}\NormalTok{ the verified signing time}\CommentTok{,}\NormalTok{ or null if no timestamp or invalid}
 \CommentTok{*/}
\KeywordTok{public} \BuiltInTok{Date} \FunctionTok{verifyTimestamp}\OperatorTok{(}\NormalTok{SignerInformation signerInfo}\OperatorTok{)} \KeywordTok{throws} \BuiltInTok{Exception} \OperatorTok{\{}
\NormalTok{    AttributeTable unsignedAttrs }\OperatorTok{=}\NormalTok{ signerInfo}\OperatorTok{.}\FunctionTok{getUnsignedAttributes}\OperatorTok{();}
    \ControlFlowTok{if} \OperatorTok{(}\NormalTok{unsignedAttrs }\OperatorTok{==} \KeywordTok{null}\OperatorTok{)} \OperatorTok{\{}
\NormalTok{        log}\OperatorTok{.}\FunctionTok{info}\OperatorTok{(}\StringTok{"No unsigned attributes — no timestamp present"}\OperatorTok{);}
        \ControlFlowTok{return} \KeywordTok{null}\OperatorTok{;}
    \OperatorTok{\}}

    \BuiltInTok{Attribute}\NormalTok{ tsAttr }\OperatorTok{=}\NormalTok{ unsignedAttrs}\OperatorTok{.}\FunctionTok{get}\OperatorTok{(}
\NormalTok{            PKCSObjectIdentifiers}\OperatorTok{.}\FunctionTok{id\_aa\_signatureTimeStampToken}\OperatorTok{);}
    \ControlFlowTok{if} \OperatorTok{(}\NormalTok{tsAttr }\OperatorTok{==} \KeywordTok{null}\OperatorTok{)} \OperatorTok{\{}
\NormalTok{        log}\OperatorTok{.}\FunctionTok{info}\OperatorTok{(}\StringTok{"No signatureTimeStampToken attribute found"}\OperatorTok{);}
        \ControlFlowTok{return} \KeywordTok{null}\OperatorTok{;}
    \OperatorTok{\}}

    \CommentTok{// Parse the timestamp token}
    \DataTypeTok{byte}\OperatorTok{[]}\NormalTok{ tsBytes }\OperatorTok{=}\NormalTok{ tsAttr}\OperatorTok{.}\FunctionTok{getAttrValues}\OperatorTok{()}
            \OperatorTok{.}\FunctionTok{getObjectAt}\OperatorTok{(}\DecValTok{0}\OperatorTok{).}\FunctionTok{toASN1Primitive}\OperatorTok{().}\FunctionTok{getEncoded}\OperatorTok{();}
\NormalTok{    CMSSignedData tsSignedData }\OperatorTok{=} \KeywordTok{new} \FunctionTok{CMSSignedData}\OperatorTok{(}\NormalTok{tsBytes}\OperatorTok{);}
\NormalTok{    TimeStampToken tsToken }\OperatorTok{=} \KeywordTok{new} \FunctionTok{TimeStampToken}\OperatorTok{(}\NormalTok{tsSignedData}\OperatorTok{);}

    \CommentTok{// Verify the timestamp\textquotesingle{}s CMS signature}
\NormalTok{    SignerInformation tsSigner }\OperatorTok{=}\NormalTok{ tsSignedData}\OperatorTok{.}\FunctionTok{getSignerInfos}\OperatorTok{()}
            \OperatorTok{.}\FunctionTok{getSigners}\OperatorTok{().}\FunctionTok{iterator}\OperatorTok{().}\FunctionTok{next}\OperatorTok{();}
    \AttributeTok{@SuppressWarnings}\OperatorTok{(}\StringTok{"unchecked"}\OperatorTok{)}
    \BuiltInTok{Collection}\OperatorTok{\textless{}}\NormalTok{X509CertificateHolder}\OperatorTok{\textgreater{}}\NormalTok{ tsMatches }\OperatorTok{=}
\NormalTok{            tsSignedData}\OperatorTok{.}\FunctionTok{getCertificates}\OperatorTok{()}
                    \OperatorTok{.}\FunctionTok{getMatches}\OperatorTok{(}\NormalTok{tsSigner}\OperatorTok{.}\FunctionTok{getSID}\OperatorTok{());}

    \ControlFlowTok{if} \OperatorTok{(}\NormalTok{tsMatches}\OperatorTok{.}\FunctionTok{isEmpty}\OperatorTok{())} \OperatorTok{\{}
\NormalTok{        log}\OperatorTok{.}\FunctionTok{error}\OperatorTok{(}\StringTok{"Timestamp token does not contain "} \OperatorTok{+}
                \StringTok{"the TSA\textquotesingle{}s certificate"}\OperatorTok{);}
        \ControlFlowTok{return} \KeywordTok{null}\OperatorTok{;}
    \OperatorTok{\}}

\NormalTok{    X509CertificateHolder tsCert }\OperatorTok{=}\NormalTok{ tsMatches}\OperatorTok{.}\FunctionTok{iterator}\OperatorTok{().}\FunctionTok{next}\OperatorTok{();}

    \CommentTok{// Verify the TSA certificate has timeStamping EKU}
    \DataTypeTok{boolean}\NormalTok{ hasTsEku }\OperatorTok{=} \KeywordTok{false}\OperatorTok{;}
    \ControlFlowTok{if} \OperatorTok{(}\NormalTok{tsCert}\OperatorTok{.}\FunctionTok{getExtensions}\OperatorTok{()} \OperatorTok{!=} \KeywordTok{null}\OperatorTok{)} \OperatorTok{\{}
\NormalTok{        Extension ekuExt }\OperatorTok{=}\NormalTok{ tsCert}\OperatorTok{.}\FunctionTok{getExtensions}\OperatorTok{().}\FunctionTok{getExtension}\OperatorTok{(}
\NormalTok{                Extension}\OperatorTok{.}\FunctionTok{extendedKeyUsage}\OperatorTok{);}
        \ControlFlowTok{if} \OperatorTok{(}\NormalTok{ekuExt }\OperatorTok{!=} \KeywordTok{null}\OperatorTok{)} \OperatorTok{\{}
\NormalTok{            ExtendedKeyUsage eku }\OperatorTok{=}
\NormalTok{                    ExtendedKeyUsage}\OperatorTok{.}\FunctionTok{getInstance}\OperatorTok{(}\NormalTok{ekuExt}\OperatorTok{.}\FunctionTok{getParsedValue}\OperatorTok{());}
\NormalTok{            hasTsEku }\OperatorTok{=}\NormalTok{ eku}\OperatorTok{.}\FunctionTok{hasKeyPurposeId}\OperatorTok{(}
\NormalTok{                    KeyPurposeId}\OperatorTok{.}\FunctionTok{id\_kp\_timeStamping}\OperatorTok{);}
        \OperatorTok{\}}
    \OperatorTok{\}}
    \ControlFlowTok{if} \OperatorTok{(!}\NormalTok{hasTsEku}\OperatorTok{)} \OperatorTok{\{}
\NormalTok{        log}\OperatorTok{.}\FunctionTok{warn}\OperatorTok{(}\StringTok{"TSA certificate does not have "} \OperatorTok{+}
                \StringTok{"id{-}kp{-}timeStamping EKU — "} \OperatorTok{+}
                \StringTok{"timestamp may not be trusted"}\OperatorTok{);}
    \OperatorTok{\}}

    \CommentTok{// Verify the timestamp\textquotesingle{}s CMS cryptographic signature}
\NormalTok{    SignerInformationVerifier tsVerifier }\OperatorTok{=}
            \KeywordTok{new} \FunctionTok{JcaSimpleSignerInfoVerifierBuilder}\OperatorTok{()}
                    \OperatorTok{.}\FunctionTok{setProvider}\OperatorTok{(}\StringTok{"BC"}\OperatorTok{)}
                    \OperatorTok{.}\FunctionTok{build}\OperatorTok{(}\NormalTok{tsCert}\OperatorTok{);}
    \DataTypeTok{boolean}\NormalTok{ tsValid }\OperatorTok{=}\NormalTok{ tsSigner}\OperatorTok{.}\FunctionTok{verify}\OperatorTok{(}\NormalTok{tsVerifier}\OperatorTok{);}

    \ControlFlowTok{if} \OperatorTok{(!}\NormalTok{tsValid}\OperatorTok{)} \OperatorTok{\{}
\NormalTok{        log}\OperatorTok{.}\FunctionTok{error}\OperatorTok{(}\StringTok{"Timestamp CMS signature verification FAILED"}\OperatorTok{);}
        \ControlFlowTok{return} \KeywordTok{null}\OperatorTok{;}
    \OperatorTok{\}}

    \BuiltInTok{Date}\NormalTok{ signingTime }\OperatorTok{=}\NormalTok{ tsToken}\OperatorTok{.}\FunctionTok{getTimeStampInfo}\OperatorTok{().}\FunctionTok{getGenTime}\OperatorTok{();}
\NormalTok{    log}\OperatorTok{.}\FunctionTok{info}\OperatorTok{(}\StringTok{"Timestamp verified — signing time: \{\}"}\OperatorTok{,}\NormalTok{ signingTime}\OperatorTok{);}
\NormalTok{    log}\OperatorTok{.}\FunctionTok{info}\OperatorTok{(}\StringTok{"TSA: \{\}"}\OperatorTok{,}\NormalTok{ tsCert}\OperatorTok{.}\FunctionTok{getSubject}\OperatorTok{());}
\NormalTok{    log}\OperatorTok{.}\FunctionTok{info}\OperatorTok{(}\StringTok{"Timestamp serial: \{\}"}\OperatorTok{,}
\NormalTok{            tsToken}\OperatorTok{.}\FunctionTok{getTimeStampInfo}\OperatorTok{().}\FunctionTok{getSerialNumber}\OperatorTok{());}
\NormalTok{    log}\OperatorTok{.}\FunctionTok{info}\OperatorTok{(}\StringTok{"Timestamp hash algorithm: \{\}"}\OperatorTok{,}
\NormalTok{            tsToken}\OperatorTok{.}\FunctionTok{getTimeStampInfo}\OperatorTok{().}\FunctionTok{getHashAlgorithm}\OperatorTok{()}
                    \OperatorTok{.}\FunctionTok{getAlgorithm}\OperatorTok{().}\FunctionTok{getId}\OperatorTok{());}

    \ControlFlowTok{return}\NormalTok{ signingTime}\OperatorTok{;}
\OperatorTok{\}}
\end{Highlighting}
\end{Shaded}

\subsection{Verifying the Timestamp Covers the
Signature}\label{_verifying_the_timestamp_covers_the_signature}

A valid timestamp is not just a signed time --- it must also prove that
the \emph{signature value} existed at that time. The timestamp's message
imprint digest must match the hash of the signature value from the
\texttt{SignerInfo}:

\begin{Shaded}
\begin{Highlighting}[]
\CommentTok{/**}
 \CommentTok{*}\NormalTok{ Verify that the timestamp}\CommentTok{\textquotesingle{}}\NormalTok{s message imprint matches the signature value}\CommentTok{.}
 \CommentTok{*/}
\KeywordTok{public} \DataTypeTok{boolean} \FunctionTok{verifyTimestampCoversSignature}\OperatorTok{(}
\NormalTok{        TimeStampToken tsToken}\OperatorTok{,}
\NormalTok{        SignerInformation signerInfo}\OperatorTok{)} \KeywordTok{throws} \BuiltInTok{Exception} \OperatorTok{\{}

    \CommentTok{// The timestamp\textquotesingle{}s imprint should be the hash of the signature value}
    \DataTypeTok{byte}\OperatorTok{[]}\NormalTok{ signatureBytes }\OperatorTok{=}\NormalTok{ signerInfo}\OperatorTok{.}\FunctionTok{getSignature}\OperatorTok{();}
    \BuiltInTok{String}\NormalTok{ tsHashAlg }\OperatorTok{=}\NormalTok{ tsToken}\OperatorTok{.}\FunctionTok{getTimeStampInfo}\OperatorTok{()}
            \OperatorTok{.}\FunctionTok{getHashAlgorithm}\OperatorTok{().}\FunctionTok{getAlgorithm}\OperatorTok{().}\FunctionTok{getId}\OperatorTok{();}
    \BuiltInTok{String}\NormalTok{ algName }\OperatorTok{=} \FunctionTok{mapOidToAlgorithm}\OperatorTok{(}\NormalTok{tsHashAlg}\OperatorTok{);}

    \BuiltInTok{MessageDigest}\NormalTok{ digest }\OperatorTok{=} \BuiltInTok{MessageDigest}\OperatorTok{.}\FunctionTok{getInstance}\OperatorTok{(}\NormalTok{algName}\OperatorTok{);}
    \DataTypeTok{byte}\OperatorTok{[]}\NormalTok{ computedImprint }\OperatorTok{=}\NormalTok{ digest}\OperatorTok{.}\FunctionTok{digest}\OperatorTok{(}\NormalTok{signatureBytes}\OperatorTok{);}

    \DataTypeTok{byte}\OperatorTok{[]}\NormalTok{ claimedImprint }\OperatorTok{=}
\NormalTok{            tsToken}\OperatorTok{.}\FunctionTok{getTimeStampInfo}\OperatorTok{().}\FunctionTok{getMessageImprintDigest}\OperatorTok{();}

    \DataTypeTok{boolean}\NormalTok{ match }\OperatorTok{=} \BuiltInTok{MessageDigest}\OperatorTok{.}\FunctionTok{isEqual}\OperatorTok{(}\NormalTok{computedImprint}\OperatorTok{,}\NormalTok{ claimedImprint}\OperatorTok{);}
    \ControlFlowTok{if} \OperatorTok{(}\NormalTok{match}\OperatorTok{)} \OperatorTok{\{}
\NormalTok{        log}\OperatorTok{.}\FunctionTok{info}\OperatorTok{(}\StringTok{"Timestamp message imprint matches signature value"}\OperatorTok{);}
    \OperatorTok{\}} \ControlFlowTok{else} \OperatorTok{\{}
\NormalTok{        log}\OperatorTok{.}\FunctionTok{error}\OperatorTok{(}\StringTok{"Timestamp message imprint does NOT match "} \OperatorTok{+}
                \StringTok{"signature value — timestamp may cover different data"}\OperatorTok{);}
    \OperatorTok{\}}
    \ControlFlowTok{return}\NormalTok{ match}\OperatorTok{;}
\OperatorTok{\}}
\end{Highlighting}
\end{Shaded}

\section{The Complete Five-Step Verification
Pipeline}\label{_the_complete_five_step_verification_pipeline}

\subsection{Putting It All Together}\label{_putting_it_all_together_2}

Here is the complete verification pipeline that combines every step from
this chapter. Each step is independent --- a failure at one step does
not prevent the other steps from running. The pipeline produces a
comprehensive \texttt{VerificationReport} with per-step results:

\begin{Shaded}
\begin{Highlighting}[]
\CommentTok{/**}
 \CommentTok{*}\NormalTok{ Complete signature verification pipeline}\CommentTok{.}
 \CommentTok{*}\NormalTok{ Runs all five verification steps and produces a comprehensive report}\CommentTok{.}
 \CommentTok{*/}
\KeywordTok{public}\NormalTok{ VerificationReport }\FunctionTok{verifySignature}\OperatorTok{(}\DataTypeTok{byte}\OperatorTok{[]}\NormalTok{ pdfBytes}\OperatorTok{,}
\NormalTok{        PDSignature signature}\OperatorTok{,}
        \BuiltInTok{Set}\OperatorTok{\textless{}}\BuiltInTok{TrustAnchor}\OperatorTok{\textgreater{}}\NormalTok{ trustAnchors}\OperatorTok{)} \OperatorTok{\{}

\NormalTok{    VerificationReport report }\OperatorTok{=} \KeywordTok{new} \FunctionTok{VerificationReport}\OperatorTok{();}

    \ControlFlowTok{try} \OperatorTok{\{}
        \CommentTok{// === Step 1: Byte Range Structural Integrity ===}
\NormalTok{        log}\OperatorTok{.}\FunctionTok{info}\OperatorTok{(}\StringTok{"Step 1: Verifying byte range..."}\OperatorTok{);}
        \DataTypeTok{int}\OperatorTok{[]}\NormalTok{ byteRange }\OperatorTok{=}\NormalTok{ signature}\OperatorTok{.}\FunctionTok{getByteRange}\OperatorTok{();}
\NormalTok{        ValidationResult byteRangeResult }\OperatorTok{=}
                \FunctionTok{verifyByteRange}\OperatorTok{(}\NormalTok{byteRange}\OperatorTok{,}\NormalTok{ pdfBytes}\OperatorTok{.}\FunctionTok{length}\OperatorTok{);}
\NormalTok{        report}\OperatorTok{.}\FunctionTok{setByteRangeResult}\OperatorTok{(}\NormalTok{byteRangeResult}\OperatorTok{);}

        \ControlFlowTok{if} \OperatorTok{(}\NormalTok{byteRangeResult }\OperatorTok{!=}\NormalTok{ ValidationResult}\OperatorTok{.}\FunctionTok{OK}
                \OperatorTok{\&\&}\NormalTok{ byteRangeResult}
                    \OperatorTok{!=}\NormalTok{ ValidationResult}\OperatorTok{.}\FunctionTok{BYTE\_RANGE\_TRAILING\_BYTES}\OperatorTok{)} \OperatorTok{\{}
\NormalTok{            log}\OperatorTok{.}\FunctionTok{error}\OperatorTok{(}\StringTok{"Step 1 FAILED: \{\}"}\OperatorTok{,}\NormalTok{ byteRangeResult}\OperatorTok{);}
            \ControlFlowTok{return}\NormalTok{ report}\OperatorTok{;} \CommentTok{// Cannot proceed without valid byte range}
        \OperatorTok{\}}

        \CommentTok{// === Step 2: Hash Verification ===}
\NormalTok{        log}\OperatorTok{.}\FunctionTok{info}\OperatorTok{(}\StringTok{"Step 2: Verifying message digest..."}\OperatorTok{);}
\NormalTok{        CMSSignedData cmsData }\OperatorTok{=} \FunctionTok{parseCmsSignedData}\OperatorTok{(}\NormalTok{pdfBytes}\OperatorTok{,}\NormalTok{ signature}\OperatorTok{);}
\NormalTok{        SignerInformation signerInfo }\OperatorTok{=}\NormalTok{ cmsData}\OperatorTok{.}\FunctionTok{getSignerInfos}\OperatorTok{()}
                \OperatorTok{.}\FunctionTok{getSigners}\OperatorTok{().}\FunctionTok{iterator}\OperatorTok{().}\FunctionTok{next}\OperatorTok{();}

        \CommentTok{// Compute the hash of the signed bytes}
        \BuiltInTok{String}\NormalTok{ digestAlgOid }\OperatorTok{=}\NormalTok{ signerInfo}\OperatorTok{.}\FunctionTok{getDigestAlgOID}\OperatorTok{();}
        \BuiltInTok{String}\NormalTok{ algName }\OperatorTok{=} \FunctionTok{mapOidToAlgorithm}\OperatorTok{(}\NormalTok{digestAlgOid}\OperatorTok{);}
        \BuiltInTok{MessageDigest}\NormalTok{ md }\OperatorTok{=} \BuiltInTok{MessageDigest}\OperatorTok{.}\FunctionTok{getInstance}\OperatorTok{(}\NormalTok{algName}\OperatorTok{);}
\NormalTok{        md}\OperatorTok{.}\FunctionTok{update}\OperatorTok{(}\NormalTok{pdfBytes}\OperatorTok{,}\NormalTok{ byteRange}\OperatorTok{[}\DecValTok{0}\OperatorTok{],}\NormalTok{ byteRange}\OperatorTok{[}\DecValTok{1}\OperatorTok{]);}
\NormalTok{        md}\OperatorTok{.}\FunctionTok{update}\OperatorTok{(}\NormalTok{pdfBytes}\OperatorTok{,}\NormalTok{ byteRange}\OperatorTok{[}\DecValTok{2}\OperatorTok{],}\NormalTok{ byteRange}\OperatorTok{[}\DecValTok{3}\OperatorTok{]);}
        \DataTypeTok{byte}\OperatorTok{[]}\NormalTok{ computedHash }\OperatorTok{=}\NormalTok{ md}\OperatorTok{.}\FunctionTok{digest}\OperatorTok{();}

        \CommentTok{// Compare to the CMS messageDigest attribute}
        \BuiltInTok{Attribute}\NormalTok{ digestAttr }\OperatorTok{=}\NormalTok{ signerInfo}\OperatorTok{.}\FunctionTok{getSignedAttributes}\OperatorTok{()}
                \OperatorTok{.}\FunctionTok{get}\OperatorTok{(}\NormalTok{CMSAttributes}\OperatorTok{.}\FunctionTok{messageDigest}\OperatorTok{);}
        \DataTypeTok{byte}\OperatorTok{[]}\NormalTok{ claimedHash }\OperatorTok{=}\NormalTok{ ASN1OctetString}\OperatorTok{.}\FunctionTok{getInstance}\OperatorTok{(}
\NormalTok{                digestAttr}\OperatorTok{.}\FunctionTok{getAttrValues}\OperatorTok{().}\FunctionTok{getObjectAt}\OperatorTok{(}\DecValTok{0}\OperatorTok{)).}\FunctionTok{getOctets}\OperatorTok{();}

        \DataTypeTok{boolean}\NormalTok{ hashMatch }\OperatorTok{=}
                \BuiltInTok{MessageDigest}\OperatorTok{.}\FunctionTok{isEqual}\OperatorTok{(}\NormalTok{computedHash}\OperatorTok{,}\NormalTok{ claimedHash}\OperatorTok{);}
\NormalTok{        report}\OperatorTok{.}\FunctionTok{setHashResult}\OperatorTok{(}\NormalTok{hashMatch}
                \OperatorTok{?}\NormalTok{ ValidationResult}\OperatorTok{.}\FunctionTok{OK}
                \OperatorTok{:}\NormalTok{ ValidationResult}\OperatorTok{.}\FunctionTok{HASH\_MISMATCH}\OperatorTok{);}

        \ControlFlowTok{if} \OperatorTok{(!}\NormalTok{hashMatch}\OperatorTok{)} \OperatorTok{\{}
\NormalTok{            log}\OperatorTok{.}\FunctionTok{error}\OperatorTok{(}\StringTok{"Step 2 FAILED: computed hash does not match "} \OperatorTok{+}
                    \StringTok{"CMS messageDigest — document was modified"}\OperatorTok{);}
        \OperatorTok{\}} \ControlFlowTok{else} \OperatorTok{\{}
\NormalTok{            log}\OperatorTok{.}\FunctionTok{info}\OperatorTok{(}\StringTok{"Step 2 PASSED: hash matches"}\OperatorTok{);}
        \OperatorTok{\}}

        \CommentTok{// === Step 3: CMS Cryptographic Verification ===}
\NormalTok{        log}\OperatorTok{.}\FunctionTok{info}\OperatorTok{(}\StringTok{"Step 3: Verifying CMS signature..."}\OperatorTok{);}
\NormalTok{        X509CertificateHolder signerCert }\OperatorTok{=}
                \FunctionTok{verifyCmsSignature}\OperatorTok{(}\NormalTok{cmsData}\OperatorTok{);}
\NormalTok{        report}\OperatorTok{.}\FunctionTok{setCmsResult}\OperatorTok{(}\NormalTok{signerCert }\OperatorTok{!=} \KeywordTok{null}
                \OperatorTok{?}\NormalTok{ ValidationResult}\OperatorTok{.}\FunctionTok{OK}
                \OperatorTok{:}\NormalTok{ ValidationResult}\OperatorTok{.}\FunctionTok{CMS\_INVALID}\OperatorTok{);}

        \ControlFlowTok{if} \OperatorTok{(}\NormalTok{signerCert }\OperatorTok{==} \KeywordTok{null}\OperatorTok{)} \OperatorTok{\{}
\NormalTok{            log}\OperatorTok{.}\FunctionTok{error}\OperatorTok{(}\StringTok{"Step 3 FAILED: CMS verification failed"}\OperatorTok{);}
            \ControlFlowTok{return}\NormalTok{ report}\OperatorTok{;}
        \OperatorTok{\}}
\NormalTok{        log}\OperatorTok{.}\FunctionTok{info}\OperatorTok{(}\StringTok{"Step 3 PASSED: CMS signature valid"}\OperatorTok{);}

        \CommentTok{// === Step 4: Certificate Chain Validation ===}
\NormalTok{        log}\OperatorTok{.}\FunctionTok{info}\OperatorTok{(}\StringTok{"Step 4: Validating certificate chain..."}\OperatorTok{);}
        \BuiltInTok{Date}\NormalTok{ validationDate }\OperatorTok{=} \KeywordTok{new} \BuiltInTok{Date}\OperatorTok{();} \CommentTok{// Default: now}
        \ControlFlowTok{try} \OperatorTok{\{}
            \CommentTok{// Check for timestamp first — use timestamp time if available}
            \BuiltInTok{Date}\NormalTok{ tsTime }\OperatorTok{=} \FunctionTok{verifyTimestamp}\OperatorTok{(}\NormalTok{signerInfo}\OperatorTok{);}
            \ControlFlowTok{if} \OperatorTok{(}\NormalTok{tsTime }\OperatorTok{!=} \KeywordTok{null}\OperatorTok{)} \OperatorTok{\{}
\NormalTok{                validationDate }\OperatorTok{=}\NormalTok{ tsTime}\OperatorTok{;}
\NormalTok{                report}\OperatorTok{.}\FunctionTok{setTimestampTime}\OperatorTok{(}\NormalTok{tsTime}\OperatorTok{);}
\NormalTok{                report}\OperatorTok{.}\FunctionTok{setTimestampResult}\OperatorTok{(}\NormalTok{ValidationResult}\OperatorTok{.}\FunctionTok{OK}\OperatorTok{);}
\NormalTok{                log}\OperatorTok{.}\FunctionTok{info}\OperatorTok{(}\StringTok{"Using timestamp time for chain validation: \{\}"}\OperatorTok{,}
\NormalTok{                        tsTime}\OperatorTok{);}
            \OperatorTok{\}}

            \BuiltInTok{CertPath}\NormalTok{ certPath }\OperatorTok{=} \FunctionTok{validateCertificateChain}\OperatorTok{(}
\NormalTok{                    signerCert}\OperatorTok{,}\NormalTok{ cmsData}\OperatorTok{,}\NormalTok{ trustAnchors}\OperatorTok{,}\NormalTok{ validationDate}\OperatorTok{);}
\NormalTok{            report}\OperatorTok{.}\FunctionTok{setChainResult}\OperatorTok{(}\NormalTok{certPath }\OperatorTok{!=} \KeywordTok{null}
                    \OperatorTok{?}\NormalTok{ ValidationResult}\OperatorTok{.}\FunctionTok{OK}
                    \OperatorTok{:}\NormalTok{ ValidationResult}\OperatorTok{.}\FunctionTok{CHAIN\_INVALID}\OperatorTok{);}
\NormalTok{            report}\OperatorTok{.}\FunctionTok{setCertPath}\OperatorTok{(}\NormalTok{certPath}\OperatorTok{);}
        \OperatorTok{\}} \ControlFlowTok{catch} \OperatorTok{(}\BuiltInTok{CertPathBuilderException}\NormalTok{ e}\OperatorTok{)} \OperatorTok{\{}
\NormalTok{            log}\OperatorTok{.}\FunctionTok{error}\OperatorTok{(}\StringTok{"Step 4 FAILED: \{\}"}\OperatorTok{,}\NormalTok{ e}\OperatorTok{.}\FunctionTok{getMessage}\OperatorTok{());}
\NormalTok{            report}\OperatorTok{.}\FunctionTok{setChainResult}\OperatorTok{(}\NormalTok{ValidationResult}\OperatorTok{.}\FunctionTok{CHAIN\_INVALID}\OperatorTok{);}
\NormalTok{            report}\OperatorTok{.}\FunctionTok{setChainError}\OperatorTok{(}\NormalTok{e}\OperatorTok{.}\FunctionTok{getMessage}\OperatorTok{());}
        \OperatorTok{\}}

        \CommentTok{// === Step 5: Trust Evaluation ===}
\NormalTok{        log}\OperatorTok{.}\FunctionTok{info}\OperatorTok{(}\StringTok{"Step 5: Evaluating trust..."}\OperatorTok{);}
        \ControlFlowTok{if} \OperatorTok{(}\NormalTok{report}\OperatorTok{.}\FunctionTok{getChainResult}\OperatorTok{()} \OperatorTok{==}\NormalTok{ ValidationResult}\OperatorTok{.}\FunctionTok{OK}\OperatorTok{)} \OperatorTok{\{}
\NormalTok{            report}\OperatorTok{.}\FunctionTok{setTrustResult}\OperatorTok{(}\NormalTok{ValidationResult}\OperatorTok{.}\FunctionTok{OK}\OperatorTok{);}
\NormalTok{            log}\OperatorTok{.}\FunctionTok{info}\OperatorTok{(}\StringTok{"Step 5 PASSED: certificate chains to "} \OperatorTok{+}
                    \StringTok{"a trusted root"}\OperatorTok{);}
        \OperatorTok{\}} \ControlFlowTok{else} \OperatorTok{\{}
\NormalTok{            report}\OperatorTok{.}\FunctionTok{setTrustResult}\OperatorTok{(}\NormalTok{ValidationResult}\OperatorTok{.}\FunctionTok{NOT\_TRUSTED}\OperatorTok{);}
\NormalTok{            log}\OperatorTok{.}\FunctionTok{warn}\OperatorTok{(}\StringTok{"Step 5: trust could not be established"}\OperatorTok{);}
        \OperatorTok{\}}

    \OperatorTok{\}} \ControlFlowTok{catch} \OperatorTok{(}\BuiltInTok{Exception}\NormalTok{ e}\OperatorTok{)} \OperatorTok{\{}
\NormalTok{        log}\OperatorTok{.}\FunctionTok{error}\OperatorTok{(}\StringTok{"Verification pipeline failed with exception"}\OperatorTok{,}\NormalTok{ e}\OperatorTok{);}
\NormalTok{        report}\OperatorTok{.}\FunctionTok{setException}\OperatorTok{(}\NormalTok{e}\OperatorTok{);}
    \OperatorTok{\}}

\NormalTok{    log}\OperatorTok{.}\FunctionTok{info}\OperatorTok{(}\StringTok{"Verification complete: \{\}"}\OperatorTok{,}\NormalTok{ report}\OperatorTok{.}\FunctionTok{getSummary}\OperatorTok{());}
    \ControlFlowTok{return}\NormalTok{ report}\OperatorTok{;}
\OperatorTok{\}}
\end{Highlighting}
\end{Shaded}

\subsection{The VerificationReport
Model}\label{_the_verificationreport_model}

\begin{Shaded}
\begin{Highlighting}[]
\AttributeTok{@Data}
\AttributeTok{@Builder}
\KeywordTok{public} \KeywordTok{class}\NormalTok{ VerificationReport }\OperatorTok{\{}
    \KeywordTok{private}\NormalTok{ ValidationResult byteRangeResult}\OperatorTok{;}
    \KeywordTok{private}\NormalTok{ ValidationResult hashResult}\OperatorTok{;}
    \KeywordTok{private}\NormalTok{ ValidationResult cmsResult}\OperatorTok{;}
    \KeywordTok{private}\NormalTok{ ValidationResult chainResult}\OperatorTok{;}
    \KeywordTok{private}\NormalTok{ ValidationResult trustResult}\OperatorTok{;}
    \KeywordTok{private}\NormalTok{ ValidationResult timestampResult}\OperatorTok{;}
    \KeywordTok{private} \BuiltInTok{Date}\NormalTok{ timestampTime}\OperatorTok{;}
    \KeywordTok{private} \BuiltInTok{CertPath}\NormalTok{ certPath}\OperatorTok{;}
    \KeywordTok{private} \BuiltInTok{String}\NormalTok{ chainError}\OperatorTok{;}
    \KeywordTok{private} \BuiltInTok{Exception}\NormalTok{ exception}\OperatorTok{;}

    \KeywordTok{public} \BuiltInTok{String} \FunctionTok{getSummary}\OperatorTok{()} \OperatorTok{\{}
        \ControlFlowTok{if} \OperatorTok{(}\NormalTok{exception }\OperatorTok{!=} \KeywordTok{null}\OperatorTok{)} \OperatorTok{\{}
            \ControlFlowTok{return} \StringTok{"ERROR: "} \OperatorTok{+}\NormalTok{ exception}\OperatorTok{.}\FunctionTok{getMessage}\OperatorTok{();}
        \OperatorTok{\}}
        \DataTypeTok{boolean}\NormalTok{ allPassed }\OperatorTok{=}
\NormalTok{                byteRangeResult }\OperatorTok{==}\NormalTok{ ValidationResult}\OperatorTok{.}\FunctionTok{OK}
                \OperatorTok{\&\&}\NormalTok{ hashResult }\OperatorTok{==}\NormalTok{ ValidationResult}\OperatorTok{.}\FunctionTok{OK}
                \OperatorTok{\&\&}\NormalTok{ cmsResult }\OperatorTok{==}\NormalTok{ ValidationResult}\OperatorTok{.}\FunctionTok{OK}
                \OperatorTok{\&\&}\NormalTok{ chainResult }\OperatorTok{==}\NormalTok{ ValidationResult}\OperatorTok{.}\FunctionTok{OK}
                \OperatorTok{\&\&}\NormalTok{ trustResult }\OperatorTok{==}\NormalTok{ ValidationResult}\OperatorTok{.}\FunctionTok{OK}\OperatorTok{;}
        \ControlFlowTok{return}\NormalTok{ allPassed }\OperatorTok{?} \StringTok{"ALL STEPS PASSED"} \OperatorTok{:} \StringTok{"ONE OR MORE STEPS FAILED"}\OperatorTok{;}
    \OperatorTok{\}}

    \KeywordTok{public} \DataTypeTok{boolean} \FunctionTok{isValid}\OperatorTok{()} \OperatorTok{\{}
        \ControlFlowTok{return}\NormalTok{ exception }\OperatorTok{==} \KeywordTok{null}
                \OperatorTok{\&\&}\NormalTok{ byteRangeResult }\OperatorTok{==}\NormalTok{ ValidationResult}\OperatorTok{.}\FunctionTok{OK}
                \OperatorTok{\&\&}\NormalTok{ hashResult }\OperatorTok{==}\NormalTok{ ValidationResult}\OperatorTok{.}\FunctionTok{OK}
                \OperatorTok{\&\&}\NormalTok{ cmsResult }\OperatorTok{==}\NormalTok{ ValidationResult}\OperatorTok{.}\FunctionTok{OK}\OperatorTok{;}
    \OperatorTok{\}}

    \KeywordTok{public} \DataTypeTok{boolean} \FunctionTok{isTrusted}\OperatorTok{()} \OperatorTok{\{}
        \ControlFlowTok{return} \FunctionTok{isValid}\OperatorTok{()}
                \OperatorTok{\&\&}\NormalTok{ chainResult }\OperatorTok{==}\NormalTok{ ValidationResult}\OperatorTok{.}\FunctionTok{OK}
                \OperatorTok{\&\&}\NormalTok{ trustResult }\OperatorTok{==}\NormalTok{ ValidationResult}\OperatorTok{.}\FunctionTok{OK}\OperatorTok{;}
    \OperatorTok{\}}
\OperatorTok{\}}
\end{Highlighting}
\end{Shaded}

\subsection{The ValidationResult Enum}\label{_the_validationresult_enum}

\begin{Shaded}
\begin{Highlighting}[]
\KeywordTok{public} \KeywordTok{enum}\NormalTok{ ValidationResult }\OperatorTok{\{}
    \FunctionTok{OK}\OperatorTok{(}\StringTok{"Passed"}\OperatorTok{),}
    \FunctionTok{BYTE\_RANGE\_MALFORMED}\OperatorTok{(}\StringTok{"ByteRange array is malformed"}\OperatorTok{),}
    \FunctionTok{BYTE\_RANGE\_NOT\_AT\_START}\OperatorTok{(}\StringTok{"ByteRange does not start at offset 0"}\OperatorTok{),}
    \FunctionTok{BYTE\_RANGE\_OVERLAP}\OperatorTok{(}\StringTok{"ByteRange segments overlap"}\OperatorTok{),}
    \FunctionTok{BYTE\_RANGE\_TRAILING\_BYTES}\OperatorTok{(}\StringTok{"Unsigned trailing bytes detected"}\OperatorTok{),}
    \FunctionTok{HASH\_MISMATCH}\OperatorTok{(}\StringTok{"Computed hash does not match CMS messageDigest"}\OperatorTok{),}
    \FunctionTok{CMS\_INVALID}\OperatorTok{(}\StringTok{"CMS cryptographic verification failed"}\OperatorTok{),}
    \FunctionTok{CHAIN\_INVALID}\OperatorTok{(}\StringTok{"Certificate chain validation failed"}\OperatorTok{),}
    \FunctionTok{NOT\_TRUSTED}\OperatorTok{(}\StringTok{"Signer certificate does not chain to a trusted root"}\OperatorTok{),}
    \FunctionTok{TIMESTAMP\_INVALID}\OperatorTok{(}\StringTok{"Timestamp verification failed"}\OperatorTok{),}
    \FunctionTok{REVOKED}\OperatorTok{(}\StringTok{"A certificate in the chain has been revoked"}\OperatorTok{);}

    \KeywordTok{private} \DataTypeTok{final} \BuiltInTok{String}\NormalTok{ description}\OperatorTok{;}

    \FunctionTok{ValidationResult}\OperatorTok{(}\BuiltInTok{String}\NormalTok{ description}\OperatorTok{)} \OperatorTok{\{}
        \KeywordTok{this}\OperatorTok{.}\FunctionTok{description} \OperatorTok{=}\NormalTok{ description}\OperatorTok{;}
    \OperatorTok{\}}

    \KeywordTok{public} \BuiltInTok{String} \FunctionTok{getDescription}\OperatorTok{()} \OperatorTok{\{}
        \ControlFlowTok{return}\NormalTok{ description}\OperatorTok{;}
    \OperatorTok{\}}
\OperatorTok{\}}
\end{Highlighting}
\end{Shaded}

\section{Modification Detection Between
Revisions}\label{_modification_detection_between_revisions}

\subsection{Why Multi-Signature Documents Need Revision
Comparison}\label{_why_multi_signature_documents_need_revision_comparison}

When a PDF has multiple signatures, each signature covers a specific
revision of the document --- a snapshot of the file at the point when
that signature was applied. Later signatures add incremental updates
(Chapter 8) that extend the file. From the perspective of the first
signature, the second signature's bytes are "trailing bytes" beyond its
byte range.

The critical question is: \textbf{did the incremental update between
revisions only add the new signature, or did it modify existing
content?}

If someone added a new page, changed text, or altered form field values
between signing events, the earlier signature's byte range still
verifies (the signed bytes haven't changed), but the document's visible
content has changed. This is the basis of the Shadow Attacks described
in Chapter 24.

\subsection{Extracting Individual
Revisions}\label{_extracting_individual_revisions}

PDFBox allows you to extract the bytes of each revision by using the
byte range information:

\begin{Shaded}
\begin{Highlighting}[]
\CommentTok{/**}
 \CommentTok{*}\NormalTok{ Extract the signed revision for a specific signature}\CommentTok{.}
 \CommentTok{*}\NormalTok{ Returns the PDF bytes up to the end of the signature}\CommentTok{\textquotesingle{}}\NormalTok{s coverage}\CommentTok{.}
 \CommentTok{*/}
\KeywordTok{public} \DataTypeTok{byte}\OperatorTok{[]} \FunctionTok{extractSignedRevision}\OperatorTok{(}\DataTypeTok{byte}\OperatorTok{[]}\NormalTok{ pdfBytes}\OperatorTok{,}
\NormalTok{        PDSignature signature}\OperatorTok{)} \OperatorTok{\{}
    \DataTypeTok{int}\OperatorTok{[]}\NormalTok{ byteRange }\OperatorTok{=}\NormalTok{ signature}\OperatorTok{.}\FunctionTok{getByteRange}\OperatorTok{();}
    \DataTypeTok{int}\NormalTok{ revisionEnd }\OperatorTok{=}\NormalTok{ byteRange}\OperatorTok{[}\DecValTok{2}\OperatorTok{]} \OperatorTok{+}\NormalTok{ byteRange}\OperatorTok{[}\DecValTok{3}\OperatorTok{];}

    \DataTypeTok{byte}\OperatorTok{[]}\NormalTok{ revision }\OperatorTok{=} \KeywordTok{new} \DataTypeTok{byte}\OperatorTok{[}\NormalTok{revisionEnd}\OperatorTok{];}
    \BuiltInTok{System}\OperatorTok{.}\FunctionTok{arraycopy}\OperatorTok{(}\NormalTok{pdfBytes}\OperatorTok{,} \DecValTok{0}\OperatorTok{,}\NormalTok{ revision}\OperatorTok{,} \DecValTok{0}\OperatorTok{,}\NormalTok{ revisionEnd}\OperatorTok{);}

\NormalTok{    log}\OperatorTok{.}\FunctionTok{info}\OperatorTok{(}\StringTok{"Extracted revision: \{\} bytes (full file: \{\} bytes)"}\OperatorTok{,}
\NormalTok{            revisionEnd}\OperatorTok{,}\NormalTok{ pdfBytes}\OperatorTok{.}\FunctionTok{length}\OperatorTok{);}
    \ControlFlowTok{return}\NormalTok{ revision}\OperatorTok{;}
\OperatorTok{\}}

\CommentTok{/**}
 \CommentTok{*}\NormalTok{ Compare two revisions to detect modifications beyond signature additions}\CommentTok{.}
 \CommentTok{*}\NormalTok{ Returns a list of differences found}\CommentTok{.}
 \CommentTok{*/}
\KeywordTok{public} \BuiltInTok{List}\OperatorTok{\textless{}}\NormalTok{RevisionDifference}\OperatorTok{\textgreater{}} \FunctionTok{compareRevisions}\OperatorTok{(}
        \DataTypeTok{byte}\OperatorTok{[]}\NormalTok{ earlierRevision}\OperatorTok{,} \DataTypeTok{byte}\OperatorTok{[]}\NormalTok{ laterRevision}\OperatorTok{)} \KeywordTok{throws} \BuiltInTok{IOException} \OperatorTok{\{}

    \BuiltInTok{List}\OperatorTok{\textless{}}\NormalTok{RevisionDifference}\OperatorTok{\textgreater{}}\NormalTok{ differences }\OperatorTok{=} \KeywordTok{new} \BuiltInTok{ArrayList}\OperatorTok{\textless{}\textgreater{}();}

    \ControlFlowTok{try} \OperatorTok{(}\NormalTok{PDDocument earlier }\OperatorTok{=}\NormalTok{ Loader}\OperatorTok{.}\FunctionTok{loadPDF}\OperatorTok{(}\NormalTok{earlierRevision}\OperatorTok{);}
\NormalTok{         PDDocument later }\OperatorTok{=}\NormalTok{ Loader}\OperatorTok{.}\FunctionTok{loadPDF}\OperatorTok{(}\NormalTok{laterRevision}\OperatorTok{))} \OperatorTok{\{}

        \CommentTok{// Check page count changes}
        \DataTypeTok{int}\NormalTok{ earlyPages }\OperatorTok{=}\NormalTok{ earlier}\OperatorTok{.}\FunctionTok{getNumberOfPages}\OperatorTok{();}
        \DataTypeTok{int}\NormalTok{ laterPages }\OperatorTok{=}\NormalTok{ later}\OperatorTok{.}\FunctionTok{getNumberOfPages}\OperatorTok{();}
        \ControlFlowTok{if} \OperatorTok{(}\NormalTok{earlyPages }\OperatorTok{!=}\NormalTok{ laterPages}\OperatorTok{)} \OperatorTok{\{}
\NormalTok{            differences}\OperatorTok{.}\FunctionTok{add}\OperatorTok{(}\KeywordTok{new} \FunctionTok{RevisionDifference}\OperatorTok{(}
\NormalTok{                    DifferenceType}\OperatorTok{.}\FunctionTok{PAGE\_COUNT\_CHANGED}\OperatorTok{,}
                    \StringTok{"Page count changed from "} \OperatorTok{+}\NormalTok{ earlyPages }\OperatorTok{+}
                    \StringTok{" to "} \OperatorTok{+}\NormalTok{ laterPages}\OperatorTok{));}
        \OperatorTok{\}}

        \CommentTok{// Check AcroForm field value changes}
\NormalTok{        PDAcroForm earlyForm }\OperatorTok{=}
\NormalTok{                earlier}\OperatorTok{.}\FunctionTok{getDocumentCatalog}\OperatorTok{().}\FunctionTok{getAcroForm}\OperatorTok{();}
\NormalTok{        PDAcroForm laterForm }\OperatorTok{=}
\NormalTok{                later}\OperatorTok{.}\FunctionTok{getDocumentCatalog}\OperatorTok{().}\FunctionTok{getAcroForm}\OperatorTok{();}

        \ControlFlowTok{if} \OperatorTok{(}\NormalTok{earlyForm }\OperatorTok{!=} \KeywordTok{null} \OperatorTok{\&\&}\NormalTok{ laterForm }\OperatorTok{!=} \KeywordTok{null}\OperatorTok{)} \OperatorTok{\{}
            \FunctionTok{compareFormFields}\OperatorTok{(}\NormalTok{earlyForm}\OperatorTok{,}\NormalTok{ laterForm}\OperatorTok{,}\NormalTok{ differences}\OperatorTok{);}
        \OperatorTok{\}}

        \CommentTok{// Check for new annotations}
        \ControlFlowTok{for} \OperatorTok{(}\DataTypeTok{int}\NormalTok{ i }\OperatorTok{=} \DecValTok{0}\OperatorTok{;}
\NormalTok{                i }\OperatorTok{\textless{}} \BuiltInTok{Math}\OperatorTok{.}\FunctionTok{min}\OperatorTok{(}\NormalTok{earlyPages}\OperatorTok{,}\NormalTok{ laterPages}\OperatorTok{);}\NormalTok{ i}\OperatorTok{++)} \OperatorTok{\{}
            \DataTypeTok{int}\NormalTok{ earlyAnnots }\OperatorTok{=}\NormalTok{ earlier}\OperatorTok{.}\FunctionTok{getPage}\OperatorTok{(}\NormalTok{i}\OperatorTok{)}
                    \OperatorTok{.}\FunctionTok{getAnnotations}\OperatorTok{().}\FunctionTok{size}\OperatorTok{();}
            \DataTypeTok{int}\NormalTok{ laterAnnots }\OperatorTok{=}\NormalTok{ later}\OperatorTok{.}\FunctionTok{getPage}\OperatorTok{(}\NormalTok{i}\OperatorTok{)}
                    \OperatorTok{.}\FunctionTok{getAnnotations}\OperatorTok{().}\FunctionTok{size}\OperatorTok{();}
            \ControlFlowTok{if} \OperatorTok{(}\NormalTok{earlyAnnots }\OperatorTok{!=}\NormalTok{ laterAnnots}\OperatorTok{)} \OperatorTok{\{}
\NormalTok{                differences}\OperatorTok{.}\FunctionTok{add}\OperatorTok{(}\KeywordTok{new} \FunctionTok{RevisionDifference}\OperatorTok{(}
\NormalTok{                        DifferenceType}\OperatorTok{.}\FunctionTok{ANNOTATIONS\_CHANGED}\OperatorTok{,}
                        \StringTok{"Page "} \OperatorTok{+} \OperatorTok{(}\NormalTok{i }\OperatorTok{+} \DecValTok{1}\OperatorTok{)} \OperatorTok{+} \StringTok{": annotations changed "} \OperatorTok{+}
                        \StringTok{"from "} \OperatorTok{+}\NormalTok{ earlyAnnots }\OperatorTok{+} \StringTok{" to "} \OperatorTok{+}\NormalTok{ laterAnnots}\OperatorTok{));}
            \OperatorTok{\}}
        \OperatorTok{\}}
    \OperatorTok{\}}

\NormalTok{    log}\OperatorTok{.}\FunctionTok{info}\OperatorTok{(}\StringTok{"Revision comparison found \{\} difference(s)"}\OperatorTok{,}
\NormalTok{            differences}\OperatorTok{.}\FunctionTok{size}\OperatorTok{());}
    \ControlFlowTok{return}\NormalTok{ differences}\OperatorTok{;}
\OperatorTok{\}}
\end{Highlighting}
\end{Shaded}

Revision comparison is inherently incomplete. A sophisticated attacker
can modify content streams in ways that are difficult to detect by
comparing PDF objects. The comparison above catches obvious changes
(page count, form fields, annotations), but it is not a security
guarantee. For high-security applications, consider using the EU DSS
framework, which implements the full ETSI modification detection
specification.

\section{Edge Cases: Self-Signed, Expired, and Missing
Chain}\label{_edge_cases_self_signed_expired_and_missing_chain}

\subsection{Self-Signed Certificates}\label{_self_signed_certificates_2}

A self-signed certificate is one where the issuer and subject are the
same entity --- the certificate signs itself. Self-signed signatures are
cryptographically valid (the CMS verifies), but the certificate chain is
trivially short: one certificate, no intermediates, no root CA.

\begin{Shaded}
\begin{Highlighting}[]
\CommentTok{/**}
 \CommentTok{*}\NormalTok{ Check if a certificate is self}\CommentTok{{-}}\NormalTok{signed}\CommentTok{.}
 \CommentTok{*/}
\KeywordTok{public} \DataTypeTok{boolean} \FunctionTok{isSelfSigned}\OperatorTok{(}\NormalTok{X509CertificateHolder cert}\OperatorTok{)} \OperatorTok{\{}
    \DataTypeTok{boolean}\NormalTok{ selfSigned }\OperatorTok{=}\NormalTok{ cert}\OperatorTok{.}\FunctionTok{getSubject}\OperatorTok{().}\FunctionTok{equals}\OperatorTok{(}\NormalTok{cert}\OperatorTok{.}\FunctionTok{getIssuer}\OperatorTok{());}
    \ControlFlowTok{if} \OperatorTok{(}\NormalTok{selfSigned}\OperatorTok{)} \OperatorTok{\{}
\NormalTok{        log}\OperatorTok{.}\FunctionTok{info}\OperatorTok{(}\StringTok{"Certificate is self{-}signed: subject and issuer "} \OperatorTok{+}
                \StringTok{"are both \{\}"}\OperatorTok{,}\NormalTok{ cert}\OperatorTok{.}\FunctionTok{getSubject}\OperatorTok{());}
    \OperatorTok{\}}
    \ControlFlowTok{return}\NormalTok{ selfSigned}\OperatorTok{;}
\OperatorTok{\}}
\end{Highlighting}
\end{Shaded}

For self-signed certificates, the chain validation in Step 4 will fail
unless you explicitly add the self-signed certificate as a trust anchor.
Most applications handle self-signed signatures by reporting: "The
signature is cryptographically valid, but the identity of the signer
could not be verified because the certificate is self-signed."

This is the correct behavior. A self-signed signature proves that
someone \emph{possessing that private key} signed the document. It does
not prove \emph{who} possesses that key, because no trusted third party
(CA) has vouched for the identity.

\subsection{Expired Certificates}\label{_expired_certificates_2}

Certificate expiration is time-dependent --- the same signature can be
"valid" today and "expired" tomorrow. The handling depends on whether a
timestamp is present:

\begin{Shaded}
\begin{Highlighting}[]
\CommentTok{/**}
 \CommentTok{*}\NormalTok{ Determine if a signature}\CommentTok{\textquotesingle{}}\NormalTok{s certificate is expired}\CommentTok{,}
 \CommentTok{*}\NormalTok{ accounting for timestamps}\CommentTok{.}
 \CommentTok{*/}
\KeywordTok{public}\NormalTok{ CertExpirationStatus }\FunctionTok{checkCertExpiration}\OperatorTok{(}
\NormalTok{        X509CertificateHolder signerCert}\OperatorTok{,}
        \BuiltInTok{Date}\NormalTok{ timestampTime}\OperatorTok{)} \OperatorTok{\{}

    \BuiltInTok{Date}\NormalTok{ now }\OperatorTok{=} \KeywordTok{new} \BuiltInTok{Date}\OperatorTok{();}
    \BuiltInTok{Date}\NormalTok{ notBefore }\OperatorTok{=}\NormalTok{ signerCert}\OperatorTok{.}\FunctionTok{getNotBefore}\OperatorTok{();}
    \BuiltInTok{Date}\NormalTok{ notAfter }\OperatorTok{=}\NormalTok{ signerCert}\OperatorTok{.}\FunctionTok{getNotAfter}\OperatorTok{();}

    \ControlFlowTok{if} \OperatorTok{(}\NormalTok{now}\OperatorTok{.}\FunctionTok{before}\OperatorTok{(}\NormalTok{notAfter}\OperatorTok{))} \OperatorTok{\{}
        \ControlFlowTok{return}\NormalTok{ CertExpirationStatus}\OperatorTok{.}\FunctionTok{VALID}\OperatorTok{;}
    \OperatorTok{\}}

    \CommentTok{// Certificate has expired — check if we have a timestamp}
    \ControlFlowTok{if} \OperatorTok{(}\NormalTok{timestampTime }\OperatorTok{!=} \KeywordTok{null}\OperatorTok{)} \OperatorTok{\{}
        \ControlFlowTok{if} \OperatorTok{(}\NormalTok{timestampTime}\OperatorTok{.}\FunctionTok{before}\OperatorTok{(}\NormalTok{notAfter}\OperatorTok{)}
                \OperatorTok{\&\&}\NormalTok{ timestampTime}\OperatorTok{.}\FunctionTok{after}\OperatorTok{(}\NormalTok{notBefore}\OperatorTok{))} \OperatorTok{\{}
\NormalTok{            log}\OperatorTok{.}\FunctionTok{info}\OperatorTok{(}\StringTok{"Certificate expired on \{\}, but timestamp \{\} "} \OperatorTok{+}
                    \StringTok{"proves signature was created during validity period"}\OperatorTok{,}
\NormalTok{                    notAfter}\OperatorTok{,}\NormalTok{ timestampTime}\OperatorTok{);}
            \ControlFlowTok{return}\NormalTok{ CertExpirationStatus}\OperatorTok{.}\FunctionTok{EXPIRED\_BUT\_TIMESTAMPED}\OperatorTok{;}
        \OperatorTok{\}} \ControlFlowTok{else} \OperatorTok{\{}
\NormalTok{            log}\OperatorTok{.}\FunctionTok{error}\OperatorTok{(}\StringTok{"Certificate expired on \{\}, and timestamp \{\} "} \OperatorTok{+}
                    \StringTok{"is outside the validity period [\{\} to \{\}]"}\OperatorTok{,}
\NormalTok{                    notAfter}\OperatorTok{,}\NormalTok{ timestampTime}\OperatorTok{,}\NormalTok{ notBefore}\OperatorTok{,}\NormalTok{ notAfter}\OperatorTok{);}
            \ControlFlowTok{return}\NormalTok{ CertExpirationStatus}\OperatorTok{.}\FunctionTok{EXPIRED\_AND\_TIMESTAMP\_INVALID}\OperatorTok{;}
        \OperatorTok{\}}
    \OperatorTok{\}}

\NormalTok{    log}\OperatorTok{.}\FunctionTok{warn}\OperatorTok{(}\StringTok{"Certificate expired on \{\} and no timestamp is present"}\OperatorTok{,}
\NormalTok{            notAfter}\OperatorTok{);}
    \ControlFlowTok{return}\NormalTok{ CertExpirationStatus}\OperatorTok{.}\FunctionTok{EXPIRED\_NO\_TIMESTAMP}\OperatorTok{;}
\OperatorTok{\}}
\end{Highlighting}
\end{Shaded}

\subsection{Missing Intermediate
Certificates}\label{_missing_intermediate_certificates}

The most common real-world chain validation failure is a missing
intermediate certificate. The CMS \texttt{SignedData} should include all
intermediate certificates between the signer's certificate and the root,
but some signing implementations omit them.

When intermediates are missing, you have two options:

\begin{enumerate}
\def\labelenumi{\arabic{enumi}.}
\item
  \textbf{AIA chasing}: Follow the Authority Information Access (AIA)
  extension in the signer's certificate to download the missing
  intermediate from the CA's server
\item
  \textbf{Local certificate cache}: Maintain a local cache of known
  intermediate certificates
\end{enumerate}

\begin{Shaded}
\begin{Highlighting}[]
\CommentTok{/**}
 \CommentTok{*}\NormalTok{ Attempt to fetch missing intermediate certificates via AIA}\CommentTok{.}
 \CommentTok{*}
 \CommentTok{*} \AnnotationTok{@param}\CommentTok{ }\AnnotationTok{cert}\CommentTok{  }\NormalTok{the certificate whose issuer we need}
 \CommentTok{*} \AnnotationTok{@return}\NormalTok{ the issuer}\CommentTok{\textquotesingle{}}\NormalTok{s certificate}\CommentTok{,}\NormalTok{ or null if not fetchable}
 \CommentTok{*/}
\KeywordTok{public} \BuiltInTok{X509Certificate} \FunctionTok{fetchIssuerViaAia}\OperatorTok{(}\BuiltInTok{X509Certificate}\NormalTok{ cert}\OperatorTok{)}
        \KeywordTok{throws} \BuiltInTok{Exception} \OperatorTok{\{}

    \DataTypeTok{byte}\OperatorTok{[]}\NormalTok{ aiaExtBytes }\OperatorTok{=}\NormalTok{ cert}\OperatorTok{.}\FunctionTok{getExtensionValue}\OperatorTok{(}
\NormalTok{            Extension}\OperatorTok{.}\FunctionTok{authorityInfoAccess}\OperatorTok{.}\FunctionTok{getId}\OperatorTok{());}
    \ControlFlowTok{if} \OperatorTok{(}\NormalTok{aiaExtBytes }\OperatorTok{==} \KeywordTok{null}\OperatorTok{)} \OperatorTok{\{}
\NormalTok{        log}\OperatorTok{.}\FunctionTok{info}\OperatorTok{(}\StringTok{"No AIA extension — cannot fetch issuer"}\OperatorTok{);}
        \ControlFlowTok{return} \KeywordTok{null}\OperatorTok{;}
    \OperatorTok{\}}

\NormalTok{    ASN1OctetString aiaOctet }\OperatorTok{=}\NormalTok{ ASN1OctetString}\OperatorTok{.}\FunctionTok{getInstance}\OperatorTok{(}\NormalTok{aiaExtBytes}\OperatorTok{);}
\NormalTok{    AuthorityInformationAccess aia }\OperatorTok{=}
\NormalTok{            AuthorityInformationAccess}\OperatorTok{.}\FunctionTok{getInstance}\OperatorTok{(}
\NormalTok{                    ASN1Sequence}\OperatorTok{.}\FunctionTok{getInstance}\OperatorTok{(}\NormalTok{aiaOctet}\OperatorTok{.}\FunctionTok{getOctets}\OperatorTok{()));}

    \ControlFlowTok{for} \OperatorTok{(}\NormalTok{AccessDescription desc }\OperatorTok{:}\NormalTok{ aia}\OperatorTok{.}\FunctionTok{getAccessDescriptions}\OperatorTok{())} \OperatorTok{\{}
        \ControlFlowTok{if} \OperatorTok{(}\NormalTok{desc}\OperatorTok{.}\FunctionTok{getAccessMethod}\OperatorTok{().}\FunctionTok{equals}\OperatorTok{(}
\NormalTok{                AccessDescription}\OperatorTok{.}\FunctionTok{id\_ad\_caIssuers}\OperatorTok{))} \OperatorTok{\{}
\NormalTok{            GeneralName gn }\OperatorTok{=}\NormalTok{ desc}\OperatorTok{.}\FunctionTok{getAccessLocation}\OperatorTok{();}
            \ControlFlowTok{if} \OperatorTok{(}\NormalTok{gn}\OperatorTok{.}\FunctionTok{getTagNo}\OperatorTok{()} \OperatorTok{==}\NormalTok{ GeneralName}\OperatorTok{.}\FunctionTok{uniformResourceIdentifier}\OperatorTok{)} \OperatorTok{\{}
                \BuiltInTok{String}\NormalTok{ url }\OperatorTok{=}\NormalTok{ gn}\OperatorTok{.}\FunctionTok{getName}\OperatorTok{().}\FunctionTok{toString}\OperatorTok{();}
\NormalTok{                log}\OperatorTok{.}\FunctionTok{info}\OperatorTok{(}\StringTok{"Fetching issuer certificate from AIA: \{\}"}\OperatorTok{,}\NormalTok{ url}\OperatorTok{);}

                \ControlFlowTok{try} \OperatorTok{(}\BuiltInTok{InputStream}\NormalTok{ is }\OperatorTok{=} \KeywordTok{new} \BuiltInTok{URL}\OperatorTok{(}\NormalTok{url}\OperatorTok{).}\FunctionTok{openStream}\OperatorTok{())} \OperatorTok{\{}
                    \BuiltInTok{CertificateFactory}\NormalTok{ cf }\OperatorTok{=}
                            \BuiltInTok{CertificateFactory}\OperatorTok{.}\FunctionTok{getInstance}\OperatorTok{(}\StringTok{"X.509"}\OperatorTok{);}
                    \ControlFlowTok{return} \OperatorTok{(}\BuiltInTok{X509Certificate}\OperatorTok{)}\NormalTok{ cf}\OperatorTok{.}\FunctionTok{generateCertificate}\OperatorTok{(}\NormalTok{is}\OperatorTok{);}
                \OperatorTok{\}}
            \OperatorTok{\}}
        \OperatorTok{\}}
    \OperatorTok{\}}

\NormalTok{    log}\OperatorTok{.}\FunctionTok{warn}\OperatorTok{(}\StringTok{"AIA extension present but no CA Issuers URL found"}\OperatorTok{);}
    \ControlFlowTok{return} \KeywordTok{null}\OperatorTok{;}
\OperatorTok{\}}
\end{Highlighting}
\end{Shaded}

In production, always cache downloaded certificates. AIA chasing adds
network latency and introduces a dependency on external servers ---
exactly the kind of dependency that LTV (Chapter 18) is designed to
eliminate. If you are building a signing service, embed all intermediate
certificates in the CMS at signing time, so validators never need to
chase AIA.

\section{Algorithm OID Mapping
Utility}\label{_algorithm_oid_mapping_utility}

Throughout this chapter, we reference a \texttt{mapOidToAlgorithm()}
helper. Here is a complete implementation that covers the algorithm OIDs
you will encounter in real PDF signatures:

\begin{Shaded}
\begin{Highlighting}[]
\CommentTok{/**}
 \CommentTok{*}\NormalTok{ Map a digest or signature algorithm OID to a Java algorithm name}\CommentTok{.}
 \CommentTok{*}\NormalTok{ Covers the algorithms used in real}\CommentTok{{-}}\NormalTok{world PDF signatures}\CommentTok{.}
 \CommentTok{*/}
\KeywordTok{public} \DataTypeTok{static} \BuiltInTok{String} \FunctionTok{mapOidToAlgorithm}\OperatorTok{(}\BuiltInTok{String}\NormalTok{ oid}\OperatorTok{)} \OperatorTok{\{}
    \ControlFlowTok{return} \ControlFlowTok{switch} \OperatorTok{(}\NormalTok{oid}\OperatorTok{)} \OperatorTok{\{}
        \CommentTok{// Digest algorithms}
        \ControlFlowTok{case} \StringTok{"1.3.14.3.2.26"}         \OperatorTok{{-}\textgreater{}} \StringTok{"SHA{-}1"}\OperatorTok{;}
        \ControlFlowTok{case} \StringTok{"2.16.840.1.101.3.4.2.1"} \OperatorTok{{-}\textgreater{}} \StringTok{"SHA{-}256"}\OperatorTok{;}
        \ControlFlowTok{case} \StringTok{"2.16.840.1.101.3.4.2.2"} \OperatorTok{{-}\textgreater{}} \StringTok{"SHA{-}384"}\OperatorTok{;}
        \ControlFlowTok{case} \StringTok{"2.16.840.1.101.3.4.2.3"} \OperatorTok{{-}\textgreater{}} \StringTok{"SHA{-}512"}\OperatorTok{;}
        \CommentTok{// RSA signature algorithms}
        \ControlFlowTok{case} \StringTok{"1.2.840.113549.1.1.5"}  \OperatorTok{{-}\textgreater{}} \StringTok{"SHA1withRSA"}\OperatorTok{;}
        \ControlFlowTok{case} \StringTok{"1.2.840.113549.1.1.11"} \OperatorTok{{-}\textgreater{}} \StringTok{"SHA256withRSA"}\OperatorTok{;}
        \ControlFlowTok{case} \StringTok{"1.2.840.113549.1.1.12"} \OperatorTok{{-}\textgreater{}} \StringTok{"SHA384withRSA"}\OperatorTok{;}
        \ControlFlowTok{case} \StringTok{"1.2.840.113549.1.1.13"} \OperatorTok{{-}\textgreater{}} \StringTok{"SHA512withRSA"}\OperatorTok{;}
        \ControlFlowTok{case} \StringTok{"1.2.840.113549.1.1.10"} \OperatorTok{{-}\textgreater{}} \StringTok{"RSASSA{-}PSS"}\OperatorTok{;}
        \CommentTok{// ECDSA signature algorithms}
        \ControlFlowTok{case} \StringTok{"1.2.840.10045.4.3.2"}   \OperatorTok{{-}\textgreater{}} \StringTok{"SHA256withECDSA"}\OperatorTok{;}
        \ControlFlowTok{case} \StringTok{"1.2.840.10045.4.3.3"}   \OperatorTok{{-}\textgreater{}} \StringTok{"SHA384withECDSA"}\OperatorTok{;}
        \ControlFlowTok{case} \StringTok{"1.2.840.10045.4.3.4"}   \OperatorTok{{-}\textgreater{}} \StringTok{"SHA512withECDSA"}\OperatorTok{;}
        \KeywordTok{default} \OperatorTok{{-}\textgreater{}} \ControlFlowTok{throw} \KeywordTok{new} \BuiltInTok{IllegalArgumentException}\OperatorTok{(}
                \StringTok{"Unknown algorithm OID: "} \OperatorTok{+}\NormalTok{ oid}\OperatorTok{);}
    \OperatorTok{\};}
\OperatorTok{\}}
\end{Highlighting}
\end{Shaded}

\section{Key Takeaways}\label{_key_takeaways_16}

\begin{itemize}
\item
  PDFBox provides signature extraction via two paths: the AcroForm field
  tree (full metadata) or \texttt{getSignatureDictionaries()} (fast,
  crypto-only)
\item
  CMS parsing requires both the \texttt{/Contents} bytes \emph{and} the
  signed content (byte range data) --- always pass both to
  \texttt{CMSSignedData}
\item
  The five-step pipeline is: byte range check, hash verification, CMS
  cryptographic verification, certificate chain validation, trust
  evaluation
\item
  \texttt{SignerInformation.verify()} performs Steps 2 and 3 combined
  --- but not chain validation or trust evaluation
\item
  Java's PKIX \texttt{CertPathBuilder} handles chain construction,
  validity checking, and basic constraints automatically
\item
  Timestamps shift the validation time from "now" to "when the signature
  was created" --- critical for expired certificates
\item
  Revision comparison detects modifications between signatures but is
  not a complete security guarantee
\item
  Self-signed signatures are cryptographically valid but untrusted ---
  always report the distinction clearly
\item
  Missing intermediates are the most common chain failure --- implement
  AIA chasing or embed all intermediates at signing time
\item
  The \texttt{VerificationReport} model gives callers granular per-step
  results rather than a binary pass/fail
\end{itemize}

\chapter{Test Cases \& Reference PDFs}\label{_test_cases_reference_pdfs}

Testing PDF digital signatures is one of the hardest testing problems in
software engineering. A signature that is valid today may be invalid
tomorrow because a certificate expired. A test that passes on your
machine may fail on CI because the TSA server is unreachable. A
verification test that ran perfectly for three years suddenly breaks
because a CRL distribution point returns HTTP 404. Time, network, and
third-party infrastructure are all embedded in your test surface --- and
they are all moving targets.

This chapter covers how to build a test infrastructure that is
deterministic, fast, and comprehensive. We generate our own keystores,
mock external services, freeze time, and construct a test matrix that
exercises every meaningful signature scenario. By the end, you will have
a test suite that runs in seconds, needs no network access, and covers
the full signing-to-verification round trip including PAdES compliance.

\section{Why Testing PDF Signatures Is
Hard}\label{_why_testing_pdf_signatures_is_hard}

\subsection{The Three Enemies of Deterministic Signature
Tests}\label{_the_three_enemies_of_deterministic_signature_tests}

\textbf{Enemy 1: Time.} Certificate validation is time-dependent. A
certificate that is valid right now has a \texttt{notBefore} and
\texttt{notAfter} date range. If your test generates a certificate valid
for one year and your test suite runs 13 months later, every test that
validates that certificate will fail. Timestamps make this worse --- a
timestamp proves the signature was created at a specific time, and the
TSA certificate has its own validity period.

\textbf{Enemy 2: Network.} Real signing workflows depend on external
services: TSA servers for timestamps, OCSP responders for revocation
checks, CRL distribution points for certificate revocation lists. Any of
these can be slow, unreachable, or returning unexpected responses. A
test that calls a real TSA is a test that depends on the internet ---
and the internet is not deterministic.

\textbf{Enemy 3: State.} PDF files are stateful. A signed PDF is a
specific sequence of bytes, and any change --- even adding a single
trailing newline --- invalidates the signature. Tests that generate PDFs
must produce identical byte sequences across runs, or the test
assertions become meaningless. Different versions of PDFBox may
serialize objects in different orders, produce different padding sizes,
or use different default values --- any of which changes the output
bytes.

\subsection{The Testing Strategy}\label{_the_testing_strategy}

Our approach addresses all three enemies:

\begin{enumerate}
\def\labelenumi{\arabic{enumi}.}
\item
  \textbf{Generate test keystores programmatically} with explicit,
  long-lived validity periods (100 years --- we are not testing
  certificate expiration unless we want to)
\item
  \textbf{Mock all external services} (TSA, OCSP) using WireMock or
  BouncyCastle's in-process TSA
\item
  \textbf{Freeze time} using a \texttt{Clock} abstraction so certificate
  and timestamp validation happens at a controlled instant
\item
  \textbf{Generate reference PDFs as part of the test setup} rather than
  shipping binary PDF files in the repository
\end{enumerate}

\section{Test Infrastructure: JUnit 5 and Test Keystore
Generation}\label{_test_infrastructure_junit_5_and_test_keystore_generation}

\subsection{The Test Keystore
Generator}\label{_the_test_keystore_generator}

Every test needs cryptographic key material: a signing certificate, an
intermediate CA, a root CA, and a TSA certificate. We generate all of
these programmatically using BouncyCastle:

\begin{Shaded}
\begin{Highlighting}[]
\CommentTok{/**}
 \CommentTok{*}\NormalTok{ Generates a complete PKI hierarchy for testing}\CommentTok{:}
 \CommentTok{*}\NormalTok{ Root CA }\CommentTok{{-}\textgreater{}}\NormalTok{ Intermediate CA }\CommentTok{{-}\textgreater{}}\NormalTok{ Signer certificate}
 \CommentTok{*}\NormalTok{ Plus a separate TSA certificate signed by the Root CA}\CommentTok{.}
 \CommentTok{*/}
\AttributeTok{@Slf4j}
\KeywordTok{public} \KeywordTok{class}\NormalTok{ TestKeystoreGenerator }\OperatorTok{\{}

    \KeywordTok{private} \DataTypeTok{static} \DataTypeTok{final} \DataTypeTok{int}\NormalTok{ KEY\_SIZE }\OperatorTok{=} \DecValTok{2048}\OperatorTok{;}
    \KeywordTok{private} \DataTypeTok{static} \DataTypeTok{final} \BuiltInTok{String}\NormalTok{ SIGNATURE\_ALGORITHM }\OperatorTok{=} \StringTok{"SHA256withRSA"}\OperatorTok{;}

    \CommentTok{/**}
     \CommentTok{*}\NormalTok{ Generate a complete test PKI and return it as an in}\CommentTok{{-}}\NormalTok{memory KeyStore}\CommentTok{.}
     \CommentTok{*/}
    \KeywordTok{public} \DataTypeTok{static} \BuiltInTok{KeyStore} \FunctionTok{generateTestKeystore}\OperatorTok{(}\DataTypeTok{char}\OperatorTok{[]}\NormalTok{ password}\OperatorTok{)}
            \KeywordTok{throws} \BuiltInTok{Exception} \OperatorTok{\{}
        \CommentTok{// Root CA — self{-}signed, 100{-}year validity}
        \BuiltInTok{KeyPair}\NormalTok{ rootKeyPair }\OperatorTok{=} \FunctionTok{generateKeyPair}\OperatorTok{();}
        \BuiltInTok{X509Certificate}\NormalTok{ rootCert }\OperatorTok{=} \FunctionTok{generateRootCaCertificate}\OperatorTok{(}
\NormalTok{                rootKeyPair}\OperatorTok{,}
                \StringTok{"CN=Test Root CA, O=PDFalyzer Test, C=DE"}\OperatorTok{,}
                \DecValTok{100}\OperatorTok{);}

        \CommentTok{// Intermediate CA — signed by Root, 50{-}year validity}
        \BuiltInTok{KeyPair}\NormalTok{ intermediateKeyPair }\OperatorTok{=} \FunctionTok{generateKeyPair}\OperatorTok{();}
        \BuiltInTok{X509Certificate}\NormalTok{ intermediateCert }\OperatorTok{=}
                \FunctionTok{generateIntermediateCaCertificate}\OperatorTok{(}
\NormalTok{                        intermediateKeyPair}\OperatorTok{,}
                        \StringTok{"CN=Test Intermediate CA, O=PDFalyzer Test, C=DE"}\OperatorTok{,}
\NormalTok{                        rootKeyPair}\OperatorTok{.}\FunctionTok{getPrivate}\OperatorTok{(),}
\NormalTok{                        rootCert}\OperatorTok{,}
                        \DecValTok{50}\OperatorTok{);}

        \CommentTok{// Signer — signed by Intermediate, 10{-}year validity}
        \BuiltInTok{KeyPair}\NormalTok{ signerKeyPair }\OperatorTok{=} \FunctionTok{generateKeyPair}\OperatorTok{();}
        \BuiltInTok{X509Certificate}\NormalTok{ signerCert }\OperatorTok{=} \FunctionTok{generateSignerCertificate}\OperatorTok{(}
\NormalTok{                signerKeyPair}\OperatorTok{,}
                \StringTok{"CN=Test Signer, O=PDFalyzer Test, C=DE"}\OperatorTok{,}
\NormalTok{                intermediateKeyPair}\OperatorTok{.}\FunctionTok{getPrivate}\OperatorTok{(),}
\NormalTok{                intermediateCert}\OperatorTok{,}
                \DecValTok{10}\OperatorTok{);}

        \CommentTok{// TSA — signed by Root, 50{-}year validity, with timeStamping EKU}
        \BuiltInTok{KeyPair}\NormalTok{ tsaKeyPair }\OperatorTok{=} \FunctionTok{generateKeyPair}\OperatorTok{();}
        \BuiltInTok{X509Certificate}\NormalTok{ tsaCert }\OperatorTok{=} \FunctionTok{generateTsaCertificate}\OperatorTok{(}
\NormalTok{                tsaKeyPair}\OperatorTok{,}
                \StringTok{"CN=Test TSA, O=PDFalyzer Test, C=DE"}\OperatorTok{,}
\NormalTok{                rootKeyPair}\OperatorTok{.}\FunctionTok{getPrivate}\OperatorTok{(),}
\NormalTok{                rootCert}\OperatorTok{,}
                \DecValTok{50}\OperatorTok{);}

        \CommentTok{// Assemble into a PKCS12 KeyStore}
        \BuiltInTok{KeyStore}\NormalTok{ ks }\OperatorTok{=} \BuiltInTok{KeyStore}\OperatorTok{.}\FunctionTok{getInstance}\OperatorTok{(}\StringTok{"PKCS12"}\OperatorTok{);}
\NormalTok{        ks}\OperatorTok{.}\FunctionTok{load}\OperatorTok{(}\KeywordTok{null}\OperatorTok{,}\NormalTok{ password}\OperatorTok{);}

        \CommentTok{// Store the signer\textquotesingle{}s key + full chain}
        \BuiltInTok{X509Certificate}\OperatorTok{[]}\NormalTok{ signerChain }\OperatorTok{=} \OperatorTok{\{}
\NormalTok{                signerCert}\OperatorTok{,}\NormalTok{ intermediateCert}\OperatorTok{,}\NormalTok{ rootCert}
        \OperatorTok{\};}
\NormalTok{        ks}\OperatorTok{.}\FunctionTok{setKeyEntry}\OperatorTok{(}\StringTok{"signer"}\OperatorTok{,}\NormalTok{ signerKeyPair}\OperatorTok{.}\FunctionTok{getPrivate}\OperatorTok{(),}
\NormalTok{                password}\OperatorTok{,}\NormalTok{ signerChain}\OperatorTok{);}

        \CommentTok{// Store the TSA\textquotesingle{}s key + chain}
        \BuiltInTok{X509Certificate}\OperatorTok{[]}\NormalTok{ tsaChain }\OperatorTok{=} \OperatorTok{\{}\NormalTok{ tsaCert}\OperatorTok{,}\NormalTok{ rootCert }\OperatorTok{\};}
\NormalTok{        ks}\OperatorTok{.}\FunctionTok{setKeyEntry}\OperatorTok{(}\StringTok{"tsa"}\OperatorTok{,}\NormalTok{ tsaKeyPair}\OperatorTok{.}\FunctionTok{getPrivate}\OperatorTok{(),}
\NormalTok{                password}\OperatorTok{,}\NormalTok{ tsaChain}\OperatorTok{);}

        \CommentTok{// Store the root as a trusted cert entry}
\NormalTok{        ks}\OperatorTok{.}\FunctionTok{setCertificateEntry}\OperatorTok{(}\StringTok{"root"}\OperatorTok{,}\NormalTok{ rootCert}\OperatorTok{);}

\NormalTok{        log}\OperatorTok{.}\FunctionTok{info}\OperatorTok{(}\StringTok{"Generated test keystore with signer, TSA, and root"}\OperatorTok{);}
        \ControlFlowTok{return}\NormalTok{ ks}\OperatorTok{;}
    \OperatorTok{\}}

    \KeywordTok{private} \DataTypeTok{static} \BuiltInTok{KeyPair} \FunctionTok{generateKeyPair}\OperatorTok{()} \KeywordTok{throws} \BuiltInTok{Exception} \OperatorTok{\{}
        \BuiltInTok{KeyPairGenerator}\NormalTok{ gen }\OperatorTok{=} \BuiltInTok{KeyPairGenerator}\OperatorTok{.}\FunctionTok{getInstance}\OperatorTok{(}\StringTok{"RSA"}\OperatorTok{);}
\NormalTok{        gen}\OperatorTok{.}\FunctionTok{initialize}\OperatorTok{(}\NormalTok{KEY\_SIZE}\OperatorTok{);}
        \ControlFlowTok{return}\NormalTok{ gen}\OperatorTok{.}\FunctionTok{generateKeyPair}\OperatorTok{();}
    \OperatorTok{\}}

    \KeywordTok{private} \DataTypeTok{static} \BuiltInTok{X509Certificate} \FunctionTok{generateRootCaCertificate}\OperatorTok{(}
            \BuiltInTok{KeyPair}\NormalTok{ keyPair}\OperatorTok{,} \BuiltInTok{String}\NormalTok{ dn}\OperatorTok{,} \DataTypeTok{int}\NormalTok{ validityYears}\OperatorTok{)}
            \KeywordTok{throws} \BuiltInTok{Exception} \OperatorTok{\{}

\NormalTok{        X500Name issuer }\OperatorTok{=} \KeywordTok{new} \FunctionTok{X500Name}\OperatorTok{(}\NormalTok{dn}\OperatorTok{);}
        \BuiltInTok{BigInteger}\NormalTok{ serial }\OperatorTok{=} \BuiltInTok{BigInteger}\OperatorTok{.}\FunctionTok{valueOf}\OperatorTok{(}
                \BuiltInTok{System}\OperatorTok{.}\FunctionTok{currentTimeMillis}\OperatorTok{());}
        \BuiltInTok{Date}\NormalTok{ notBefore }\OperatorTok{=} \KeywordTok{new} \BuiltInTok{Date}\OperatorTok{();}
        \BuiltInTok{Date}\NormalTok{ notAfter }\OperatorTok{=} \BuiltInTok{Date}\OperatorTok{.}\FunctionTok{from}\OperatorTok{(}\NormalTok{Instant}\OperatorTok{.}\FunctionTok{now}\OperatorTok{()}
                \OperatorTok{.}\FunctionTok{plus}\OperatorTok{(}\BuiltInTok{Duration}\OperatorTok{.}\FunctionTok{ofDays}\OperatorTok{(}\DecValTok{365L} \OperatorTok{*}\NormalTok{ validityYears}\OperatorTok{)));}

\NormalTok{        JcaX509v3CertificateBuilder builder }\OperatorTok{=}
                \KeywordTok{new} \FunctionTok{JcaX509v3CertificateBuilder}\OperatorTok{(}
\NormalTok{                        issuer}\OperatorTok{,}\NormalTok{ serial}\OperatorTok{,}\NormalTok{ notBefore}\OperatorTok{,}\NormalTok{ notAfter}\OperatorTok{,}
\NormalTok{                        issuer}\OperatorTok{,}\NormalTok{ keyPair}\OperatorTok{.}\FunctionTok{getPublic}\OperatorTok{());}

        \CommentTok{// CA extensions}
\NormalTok{        builder}\OperatorTok{.}\FunctionTok{addExtension}\OperatorTok{(}\NormalTok{Extension}\OperatorTok{.}\FunctionTok{basicConstraints}\OperatorTok{,} \KeywordTok{true}\OperatorTok{,}
                \KeywordTok{new} \FunctionTok{BasicConstraints}\OperatorTok{(}\KeywordTok{true}\OperatorTok{));}
\NormalTok{        builder}\OperatorTok{.}\FunctionTok{addExtension}\OperatorTok{(}\NormalTok{Extension}\OperatorTok{.}\FunctionTok{keyUsage}\OperatorTok{,} \KeywordTok{true}\OperatorTok{,}
                \KeywordTok{new} \FunctionTok{KeyUsage}\OperatorTok{(}\NormalTok{KeyUsage}\OperatorTok{.}\FunctionTok{keyCertSign} \OperatorTok{|}\NormalTok{ KeyUsage}\OperatorTok{.}\FunctionTok{cRLSign}\OperatorTok{));}

\NormalTok{        ContentSigner signer }\OperatorTok{=} \KeywordTok{new} \FunctionTok{JcaContentSignerBuilder}\OperatorTok{(}
\NormalTok{                SIGNATURE\_ALGORITHM}\OperatorTok{).}\FunctionTok{setProvider}\OperatorTok{(}\StringTok{"BC"}\OperatorTok{)}
                \OperatorTok{.}\FunctionTok{build}\OperatorTok{(}\NormalTok{keyPair}\OperatorTok{.}\FunctionTok{getPrivate}\OperatorTok{());}

        \ControlFlowTok{return} \KeywordTok{new} \FunctionTok{JcaX509CertificateConverter}\OperatorTok{().}\FunctionTok{setProvider}\OperatorTok{(}\StringTok{"BC"}\OperatorTok{)}
                \OperatorTok{.}\FunctionTok{getCertificate}\OperatorTok{(}\NormalTok{builder}\OperatorTok{.}\FunctionTok{build}\OperatorTok{(}\NormalTok{signer}\OperatorTok{));}
    \OperatorTok{\}}

    \CommentTok{// Similar methods for intermediate, signer, and TSA certificates}
    \CommentTok{// omitted for brevity — full listings in Appendix A}
\OperatorTok{\}}
\end{Highlighting}
\end{Shaded}

\subsection{JUnit 5 Test Base Class}\label{_junit_5_test_base_class}

We use a shared base class that sets up the test PKI and provides
utility methods:

\begin{Shaded}
\begin{Highlighting}[]
\AttributeTok{@Slf4j}
\KeywordTok{abstract} \KeywordTok{class}\NormalTok{ SignatureTestBase }\OperatorTok{\{}

    \KeywordTok{protected} \DataTypeTok{static} \BuiltInTok{KeyStore}\NormalTok{ testKeystore}\OperatorTok{;}
    \KeywordTok{protected} \DataTypeTok{static} \BuiltInTok{PrivateKey}\NormalTok{ signerKey}\OperatorTok{;}
    \KeywordTok{protected} \DataTypeTok{static} \BuiltInTok{X509Certificate}\OperatorTok{[]}\NormalTok{ signerChain}\OperatorTok{;}
    \KeywordTok{protected} \DataTypeTok{static} \BuiltInTok{PrivateKey}\NormalTok{ tsaKey}\OperatorTok{;}
    \KeywordTok{protected} \DataTypeTok{static} \BuiltInTok{X509Certificate}\OperatorTok{[]}\NormalTok{ tsaChain}\OperatorTok{;}
    \KeywordTok{protected} \DataTypeTok{static} \BuiltInTok{Set}\OperatorTok{\textless{}}\BuiltInTok{TrustAnchor}\OperatorTok{\textgreater{}}\NormalTok{ trustAnchors}\OperatorTok{;}
    \KeywordTok{protected} \DataTypeTok{static} \DataTypeTok{final} \DataTypeTok{char}\OperatorTok{[]}\NormalTok{ PASSWORD }\OperatorTok{=} \StringTok{"test"}\OperatorTok{.}\FunctionTok{toCharArray}\OperatorTok{();}

    \AttributeTok{@BeforeAll}
    \DataTypeTok{static} \DataTypeTok{void} \FunctionTok{setupPki}\OperatorTok{()} \KeywordTok{throws} \BuiltInTok{Exception} \OperatorTok{\{}
        \BuiltInTok{Security}\OperatorTok{.}\FunctionTok{addProvider}\OperatorTok{(}\KeywordTok{new} \FunctionTok{BouncyCastleProvider}\OperatorTok{());}

\NormalTok{        testKeystore }\OperatorTok{=}\NormalTok{ TestKeystoreGenerator}
                \OperatorTok{.}\FunctionTok{generateTestKeystore}\OperatorTok{(}\NormalTok{PASSWORD}\OperatorTok{);}

\NormalTok{        signerKey }\OperatorTok{=} \OperatorTok{(}\BuiltInTok{PrivateKey}\OperatorTok{)}\NormalTok{ testKeystore}\OperatorTok{.}\FunctionTok{getKey}\OperatorTok{(}\StringTok{"signer"}\OperatorTok{,}\NormalTok{ PASSWORD}\OperatorTok{);}
\NormalTok{        signerChain }\OperatorTok{=} \BuiltInTok{Arrays}\OperatorTok{.}\FunctionTok{stream}\OperatorTok{(}
\NormalTok{                testKeystore}\OperatorTok{.}\FunctionTok{getCertificateChain}\OperatorTok{(}\StringTok{"signer"}\OperatorTok{))}
                \OperatorTok{.}\FunctionTok{map}\OperatorTok{(}\NormalTok{c }\OperatorTok{{-}\textgreater{}} \OperatorTok{(}\BuiltInTok{X509Certificate}\OperatorTok{)}\NormalTok{ c}\OperatorTok{)}
                \OperatorTok{.}\FunctionTok{toArray}\OperatorTok{(}\BuiltInTok{X509Certificate}\OperatorTok{[]::}\KeywordTok{new}\OperatorTok{);}

\NormalTok{        tsaKey }\OperatorTok{=} \OperatorTok{(}\BuiltInTok{PrivateKey}\OperatorTok{)}\NormalTok{ testKeystore}\OperatorTok{.}\FunctionTok{getKey}\OperatorTok{(}\StringTok{"tsa"}\OperatorTok{,}\NormalTok{ PASSWORD}\OperatorTok{);}
\NormalTok{        tsaChain }\OperatorTok{=} \BuiltInTok{Arrays}\OperatorTok{.}\FunctionTok{stream}\OperatorTok{(}
\NormalTok{                testKeystore}\OperatorTok{.}\FunctionTok{getCertificateChain}\OperatorTok{(}\StringTok{"tsa"}\OperatorTok{))}
                \OperatorTok{.}\FunctionTok{map}\OperatorTok{(}\NormalTok{c }\OperatorTok{{-}\textgreater{}} \OperatorTok{(}\BuiltInTok{X509Certificate}\OperatorTok{)}\NormalTok{ c}\OperatorTok{)}
                \OperatorTok{.}\FunctionTok{toArray}\OperatorTok{(}\BuiltInTok{X509Certificate}\OperatorTok{[]::}\KeywordTok{new}\OperatorTok{);}

        \BuiltInTok{X509Certificate}\NormalTok{ rootCert }\OperatorTok{=} \OperatorTok{(}\BuiltInTok{X509Certificate}\OperatorTok{)}
\NormalTok{                testKeystore}\OperatorTok{.}\FunctionTok{getCertificate}\OperatorTok{(}\StringTok{"root"}\OperatorTok{);}
\NormalTok{        trustAnchors }\OperatorTok{=} \BuiltInTok{Set}\OperatorTok{.}\FunctionTok{of}\OperatorTok{(}\KeywordTok{new} \BuiltInTok{TrustAnchor}\OperatorTok{(}\NormalTok{rootCert}\OperatorTok{,} \KeywordTok{null}\OperatorTok{));}

\NormalTok{        log}\OperatorTok{.}\FunctionTok{info}\OperatorTok{(}\StringTok{"Test PKI initialized — signer: \{\}, root: \{\}"}\OperatorTok{,}
\NormalTok{                signerChain}\OperatorTok{[}\DecValTok{0}\OperatorTok{].}\FunctionTok{getSubjectX500Principal}\OperatorTok{(),}
\NormalTok{                rootCert}\OperatorTok{.}\FunctionTok{getSubjectX500Principal}\OperatorTok{());}
    \OperatorTok{\}}

    \CommentTok{/**}
     \CommentTok{*}\NormalTok{ Generate a minimal unsigned PDF for testing}\CommentTok{.}
     \CommentTok{*/}
    \KeywordTok{protected} \DataTypeTok{byte}\OperatorTok{[]} \FunctionTok{generateTestPdf}\OperatorTok{(}\BuiltInTok{String}\NormalTok{ content}\OperatorTok{)} \KeywordTok{throws} \BuiltInTok{IOException} \OperatorTok{\{}
        \ControlFlowTok{try} \OperatorTok{(}\NormalTok{PDDocument doc }\OperatorTok{=} \KeywordTok{new} \FunctionTok{PDDocument}\OperatorTok{())} \OperatorTok{\{}
\NormalTok{            PDPage page }\OperatorTok{=} \KeywordTok{new} \FunctionTok{PDPage}\OperatorTok{();}
\NormalTok{            doc}\OperatorTok{.}\FunctionTok{addPage}\OperatorTok{(}\NormalTok{page}\OperatorTok{);}
            \ControlFlowTok{try} \OperatorTok{(}\NormalTok{PDPageContentStream cs }\OperatorTok{=}
                    \KeywordTok{new} \FunctionTok{PDPageContentStream}\OperatorTok{(}\NormalTok{doc}\OperatorTok{,}\NormalTok{ page}\OperatorTok{))} \OperatorTok{\{}
\NormalTok{                cs}\OperatorTok{.}\FunctionTok{beginText}\OperatorTok{();}
\NormalTok{                cs}\OperatorTok{.}\FunctionTok{setFont}\OperatorTok{(}\KeywordTok{new} \FunctionTok{PDType1Font}\OperatorTok{(}
\NormalTok{                        Standard14Fonts}\OperatorTok{.}\FunctionTok{FontName}\OperatorTok{.}\FunctionTok{HELVETICA}\OperatorTok{),} \DecValTok{12}\OperatorTok{);}
\NormalTok{                cs}\OperatorTok{.}\FunctionTok{newLineAtOffset}\OperatorTok{(}\DecValTok{72}\OperatorTok{,} \DecValTok{700}\OperatorTok{);}
\NormalTok{                cs}\OperatorTok{.}\FunctionTok{showText}\OperatorTok{(}\NormalTok{content}\OperatorTok{);}
\NormalTok{                cs}\OperatorTok{.}\FunctionTok{endText}\OperatorTok{();}
            \OperatorTok{\}}
            \BuiltInTok{ByteArrayOutputStream}\NormalTok{ baos }\OperatorTok{=} \KeywordTok{new} \BuiltInTok{ByteArrayOutputStream}\OperatorTok{();}
\NormalTok{            doc}\OperatorTok{.}\FunctionTok{save}\OperatorTok{(}\NormalTok{baos}\OperatorTok{);}
            \ControlFlowTok{return}\NormalTok{ baos}\OperatorTok{.}\FunctionTok{toByteArray}\OperatorTok{();}
        \OperatorTok{\}}
    \OperatorTok{\}}
\OperatorTok{\}}
\end{Highlighting}
\end{Shaded}

\section{Test PDF Generation
Scenarios}\label{_test_pdf_generation_scenarios}

\subsection{Generating Signed PDFs
Programmatically}\label{_generating_signed_pdfs_programmatically}

Rather than shipping pre-built signed PDFs (which become stale and
opaque), we generate them as part of each test. This ensures every test
controls its own input and can vary specific parameters:

\begin{Shaded}
\begin{Highlighting}[]
\CommentTok{/**}
 \CommentTok{*}\NormalTok{ Sign a PDF with the test keystore and return the signed bytes}\CommentTok{.}
 \CommentTok{*/}
\KeywordTok{protected} \DataTypeTok{byte}\OperatorTok{[]} \FunctionTok{signPdf}\OperatorTok{(}\DataTypeTok{byte}\OperatorTok{[]}\NormalTok{ pdfBytes}\OperatorTok{,} \DataTypeTok{boolean}\NormalTok{ addTimestamp}\OperatorTok{,}
        \BuiltInTok{String}\NormalTok{ subFilter}\OperatorTok{)} \KeywordTok{throws} \BuiltInTok{Exception} \OperatorTok{\{}

    \ControlFlowTok{try} \OperatorTok{(}\NormalTok{PDDocument doc }\OperatorTok{=}\NormalTok{ Loader}\OperatorTok{.}\FunctionTok{loadPDF}\OperatorTok{(}\NormalTok{pdfBytes}\OperatorTok{))} \OperatorTok{\{}
\NormalTok{        PDSignature pdSig }\OperatorTok{=} \KeywordTok{new} \FunctionTok{PDSignature}\OperatorTok{();}
\NormalTok{        pdSig}\OperatorTok{.}\FunctionTok{setFilter}\OperatorTok{(}\NormalTok{PDSignature}\OperatorTok{.}\FunctionTok{FILTER\_ADOBE\_PPKLITE}\OperatorTok{);}
\NormalTok{        pdSig}\OperatorTok{.}\FunctionTok{setSubFilter}\OperatorTok{(}\NormalTok{COSName}\OperatorTok{.}\FunctionTok{getPDFName}\OperatorTok{(}\NormalTok{subFilter}\OperatorTok{));}
\NormalTok{        pdSig}\OperatorTok{.}\FunctionTok{setName}\OperatorTok{(}\StringTok{"Test Signer"}\OperatorTok{);}
\NormalTok{        pdSig}\OperatorTok{.}\FunctionTok{setReason}\OperatorTok{(}\StringTok{"Integration test"}\OperatorTok{);}
\NormalTok{        pdSig}\OperatorTok{.}\FunctionTok{setSignDate}\OperatorTok{(}\BuiltInTok{Calendar}\OperatorTok{.}\FunctionTok{getInstance}\OperatorTok{());}

\NormalTok{        SignatureInterface sigInterface }\OperatorTok{=}
                \KeywordTok{new} \FunctionTok{TestSignatureInterface}\OperatorTok{(}
\NormalTok{                        signerKey}\OperatorTok{,}\NormalTok{ signerChain}\OperatorTok{,}
\NormalTok{                        addTimestamp }\OperatorTok{?}\NormalTok{ tsaKey }\OperatorTok{:} \KeywordTok{null}\OperatorTok{,}
\NormalTok{                        addTimestamp }\OperatorTok{?}\NormalTok{ tsaChain }\OperatorTok{:} \KeywordTok{null}\OperatorTok{);}

\NormalTok{        doc}\OperatorTok{.}\FunctionTok{addSignature}\OperatorTok{(}\NormalTok{pdSig}\OperatorTok{,}\NormalTok{ sigInterface}\OperatorTok{);}

        \BuiltInTok{ByteArrayOutputStream}\NormalTok{ baos }\OperatorTok{=} \KeywordTok{new} \BuiltInTok{ByteArrayOutputStream}\OperatorTok{();}
\NormalTok{        doc}\OperatorTok{.}\FunctionTok{saveIncremental}\OperatorTok{(}\NormalTok{baos}\OperatorTok{);}
        \ControlFlowTok{return}\NormalTok{ baos}\OperatorTok{.}\FunctionTok{toByteArray}\OperatorTok{();}
    \OperatorTok{\}}
\OperatorTok{\}}
\end{Highlighting}
\end{Shaded}

\subsection{Scenario Generators}\label{_scenario_generators}

Each test scenario generates a specific type of signed PDF:

\begin{Shaded}
\begin{Highlighting}[]
\CommentTok{/**}
 \CommentTok{*}\NormalTok{ Generate a PDF signed with different configurations}\CommentTok{.}
 \CommentTok{*/}
\KeywordTok{protected} \DataTypeTok{byte}\OperatorTok{[]} \FunctionTok{generateScenario}\OperatorTok{(}\NormalTok{TestScenario scenario}\OperatorTok{)}
        \KeywordTok{throws} \BuiltInTok{Exception} \OperatorTok{\{}

    \DataTypeTok{byte}\OperatorTok{[]}\NormalTok{ pdf }\OperatorTok{=} \FunctionTok{generateTestPdf}\OperatorTok{(}\StringTok{"Scenario: "} \OperatorTok{+}\NormalTok{ scenario}\OperatorTok{.}\FunctionTok{name}\OperatorTok{());}

    \ControlFlowTok{return} \ControlFlowTok{switch} \OperatorTok{(}\NormalTok{scenario}\OperatorTok{)} \OperatorTok{\{}
        \ControlFlowTok{case}\NormalTok{ BASIC\_PKCS7 }\OperatorTok{{-}\textgreater{}}
            \FunctionTok{signPdf}\OperatorTok{(}\NormalTok{pdf}\OperatorTok{,} \KeywordTok{false}\OperatorTok{,} \StringTok{"adbe.pkcs7.detached"}\OperatorTok{);}
        \ControlFlowTok{case}\NormalTok{ BASIC\_PADES }\OperatorTok{{-}\textgreater{}}
            \FunctionTok{signPdf}\OperatorTok{(}\NormalTok{pdf}\OperatorTok{,} \KeywordTok{false}\OperatorTok{,} \StringTok{"ETSI.CAdES.detached"}\OperatorTok{);}
        \ControlFlowTok{case}\NormalTok{ PADES\_WITH\_TIMESTAMP }\OperatorTok{{-}\textgreater{}}
            \FunctionTok{signPdf}\OperatorTok{(}\NormalTok{pdf}\OperatorTok{,} \KeywordTok{true}\OperatorTok{,} \StringTok{"ETSI.CAdES.detached"}\OperatorTok{);}
        \ControlFlowTok{case}\NormalTok{ SELF\_SIGNED }\OperatorTok{{-}\textgreater{}}
            \FunctionTok{signPdfWithSelfSigned}\OperatorTok{(}\NormalTok{pdf}\OperatorTok{);}
        \ControlFlowTok{case}\NormalTok{ EXPIRED\_CERT }\OperatorTok{{-}\textgreater{}}
            \FunctionTok{signPdfWithExpiredCert}\OperatorTok{(}\NormalTok{pdf}\OperatorTok{);}
        \ControlFlowTok{case}\NormalTok{ MULTI\_SIGNATURE }\OperatorTok{{-}\textgreater{}}
            \FunctionTok{signPdfTwice}\OperatorTok{(}\NormalTok{pdf}\OperatorTok{);}
        \ControlFlowTok{case}\NormalTok{ TAMPERED\_AFTER\_SIGNING }\OperatorTok{{-}\textgreater{}}
            \FunctionTok{tamperSignedPdf}\OperatorTok{(}\FunctionTok{signPdf}\OperatorTok{(}\NormalTok{pdf}\OperatorTok{,} \KeywordTok{false}\OperatorTok{,}
                    \StringTok{"adbe.pkcs7.detached"}\OperatorTok{));}
        \ControlFlowTok{case}\NormalTok{ MISSING\_INTERMEDIATE }\OperatorTok{{-}\textgreater{}}
            \FunctionTok{signPdfWithMissingIntermediate}\OperatorTok{(}\NormalTok{pdf}\OperatorTok{);}
    \OperatorTok{\};}
\OperatorTok{\}}
\end{Highlighting}
\end{Shaded}

\section{Test Matrix: Scenario to Expected
Result}\label{_test_matrix_scenario_to_expected_result}

The following table defines every test scenario, what it tests, and the
expected outcome at each verification step. This is the specification
for your test suite --- every row should be a test method:

\begin{longtable}[]{@{}
  >{\raggedright\arraybackslash}p{(\linewidth - 12\tabcolsep) * \real{0.2727}}
  >{\raggedright\arraybackslash}p{(\linewidth - 12\tabcolsep) * \real{0.0909}}
  >{\raggedright\arraybackslash}p{(\linewidth - 12\tabcolsep) * \real{0.0909}}
  >{\raggedright\arraybackslash}p{(\linewidth - 12\tabcolsep) * \real{0.0909}}
  >{\raggedright\arraybackslash}p{(\linewidth - 12\tabcolsep) * \real{0.0909}}
  >{\raggedright\arraybackslash}p{(\linewidth - 12\tabcolsep) * \real{0.0909}}
  >{\raggedright\arraybackslash}p{(\linewidth - 12\tabcolsep) * \real{0.2727}}@{}}
\caption{Signature test matrix}\tabularnewline
\toprule\noalign{}
\begin{minipage}[b]{\linewidth}\raggedright
Scenario
\end{minipage} & \begin{minipage}[b]{\linewidth}\raggedright
ByteRange
\end{minipage} & \begin{minipage}[b]{\linewidth}\raggedright
Hash
\end{minipage} & \begin{minipage}[b]{\linewidth}\raggedright
CMS
\end{minipage} & \begin{minipage}[b]{\linewidth}\raggedright
Chain
\end{minipage} & \begin{minipage}[b]{\linewidth}\raggedright
Trust
\end{minipage} & \begin{minipage}[b]{\linewidth}\raggedright
Notes
\end{minipage} \\
\midrule\noalign{}
\endfirsthead
\toprule\noalign{}
\begin{minipage}[b]{\linewidth}\raggedright
Scenario
\end{minipage} & \begin{minipage}[b]{\linewidth}\raggedright
ByteRange
\end{minipage} & \begin{minipage}[b]{\linewidth}\raggedright
Hash
\end{minipage} & \begin{minipage}[b]{\linewidth}\raggedright
CMS
\end{minipage} & \begin{minipage}[b]{\linewidth}\raggedright
Chain
\end{minipage} & \begin{minipage}[b]{\linewidth}\raggedright
Trust
\end{minipage} & \begin{minipage}[b]{\linewidth}\raggedright
Notes
\end{minipage} \\
\midrule\noalign{}
\endhead
\bottomrule\noalign{}
\endlastfoot
Basic PKCS\#7 (adbe.pkcs7.detached) & OK & OK & OK & OK & OK & Baseline
--- everything should pass \\
Basic PAdES (ETSI.CAdES.detached) & OK & OK & OK & OK & OK & Same as
PKCS\#7 but with PAdES SubFilter \\
PAdES with timestamp & OK & OK & OK & OK & OK & Timestamp attribute
present and valid \\
Self-signed certificate & OK & OK & OK & FAIL & FAIL & CMS is valid but
chain has no trusted root \\
Expired certificate (no timestamp) & OK & OK & OK & FAIL & FAIL &
CertPathValidator rejects expired cert \\
Expired certificate (with timestamp) & OK & OK & OK & OK & OK &
Timestamp proves signing was within validity \\
Multi-signature document & WARN & OK & OK & OK & OK & First sig has
trailing bytes (second revision) \\
Tampered after signing & OK & FAIL &  ---  &  ---  &  ---  & Hash
mismatch --- document was modified \\
Missing intermediate certificate & OK & OK & OK & FAIL & FAIL & Chain
cannot be built without intermediate \\
Revoked certificate & OK & OK & OK & FAIL & FAIL & Requires OCSP/CRL
mock returning "revoked" \\
Wrong SubFilter for content & OK & OK & OK & OK & OK & SubFilter says
PAdES but ESS attr is missing \\
Empty /Contents (zero-length CMS) & OK & FAIL & FAIL &  ---  &  ---  &
CMSException during parsing \\
Invisible signature (zero-size widget) & OK & OK & OK & OK & OK &
Cryptographically identical to visible sig \\
\end{longtable}

Each row in this matrix translates directly to a test method. The "-\/-"
entries mean that step is not reached because an earlier step failed
fatally.

\section{Mocking TSA and OCSP}\label{_mocking_tsa_and_ocsp}

\subsection{In-Process TSA with
BouncyCastle}\label{_in_process_tsa_with_bouncycastle}

For timestamp tests, we need a TSA that returns deterministic timestamp
tokens without hitting the network. BouncyCastle provides everything you
need to build one in-process:

\begin{Shaded}
\begin{Highlighting}[]
\CommentTok{/**}
 \CommentTok{*}\NormalTok{ In}\CommentTok{{-}}\NormalTok{process TSA for testing}\CommentTok{.}
 \CommentTok{*}\NormalTok{ Produces valid RFC }\CommentTok{3161}\NormalTok{ timestamp tokens using the test TSA key}\CommentTok{.}
 \CommentTok{*/}
\AttributeTok{@Slf4j}
\KeywordTok{public} \KeywordTok{class}\NormalTok{ InProcessTsa }\OperatorTok{\{}

    \KeywordTok{private} \DataTypeTok{final} \BuiltInTok{PrivateKey}\NormalTok{ tsaKey}\OperatorTok{;}
    \KeywordTok{private} \DataTypeTok{final} \BuiltInTok{X509Certificate}\NormalTok{ tsaCert}\OperatorTok{;}
    \KeywordTok{private} \DataTypeTok{final} \BuiltInTok{AtomicLong}\NormalTok{ serialNumber }\OperatorTok{=} \KeywordTok{new} \BuiltInTok{AtomicLong}\OperatorTok{(}\DecValTok{1}\OperatorTok{);}

    \KeywordTok{public} \FunctionTok{InProcessTsa}\OperatorTok{(}\BuiltInTok{PrivateKey}\NormalTok{ tsaKey}\OperatorTok{,} \BuiltInTok{X509Certificate}\NormalTok{ tsaCert}\OperatorTok{)} \OperatorTok{\{}
        \KeywordTok{this}\OperatorTok{.}\FunctionTok{tsaKey} \OperatorTok{=}\NormalTok{ tsaKey}\OperatorTok{;}
        \KeywordTok{this}\OperatorTok{.}\FunctionTok{tsaCert} \OperatorTok{=}\NormalTok{ tsaCert}\OperatorTok{;}
    \OperatorTok{\}}

    \CommentTok{/**}
     \CommentTok{*}\NormalTok{ Generate a timestamp token for the given data hash}\CommentTok{.}
     \CommentTok{*/}
    \KeywordTok{public}\NormalTok{ TimeStampToken }\FunctionTok{timestamp}\OperatorTok{(}\DataTypeTok{byte}\OperatorTok{[]}\NormalTok{ dataToTimestamp}\OperatorTok{)}
            \KeywordTok{throws} \BuiltInTok{Exception} \OperatorTok{\{}

        \BuiltInTok{MessageDigest}\NormalTok{ digest }\OperatorTok{=} \BuiltInTok{MessageDigest}\OperatorTok{.}\FunctionTok{getInstance}\OperatorTok{(}\StringTok{"SHA{-}256"}\OperatorTok{);}
        \DataTypeTok{byte}\OperatorTok{[]}\NormalTok{ hash }\OperatorTok{=}\NormalTok{ digest}\OperatorTok{.}\FunctionTok{digest}\OperatorTok{(}\NormalTok{dataToTimestamp}\OperatorTok{);}

\NormalTok{        TimeStampRequestGenerator reqGen }\OperatorTok{=}
                \KeywordTok{new} \FunctionTok{TimeStampRequestGenerator}\OperatorTok{();}
\NormalTok{        reqGen}\OperatorTok{.}\FunctionTok{setCertReq}\OperatorTok{(}\KeywordTok{true}\OperatorTok{);}
\NormalTok{        TimeStampRequest request }\OperatorTok{=}
\NormalTok{                reqGen}\OperatorTok{.}\FunctionTok{generate}\OperatorTok{(}\NormalTok{TSPAlgorithms}\OperatorTok{.}\FunctionTok{SHA256}\OperatorTok{,}\NormalTok{ hash}\OperatorTok{);}

        \CommentTok{// Build the timestamp token}
\NormalTok{        JcaSignerInfoGeneratorBuilder sigGenBuilder }\OperatorTok{=}
                \KeywordTok{new} \FunctionTok{JcaSignerInfoGeneratorBuilder}\OperatorTok{(}
                        \KeywordTok{new} \FunctionTok{JcaDigestCalculatorProviderBuilder}\OperatorTok{()}
                                \OperatorTok{.}\FunctionTok{setProvider}\OperatorTok{(}\StringTok{"BC"}\OperatorTok{).}\FunctionTok{build}\OperatorTok{());}

\NormalTok{        ContentSigner contentSigner }\OperatorTok{=} \KeywordTok{new} \FunctionTok{JcaContentSignerBuilder}\OperatorTok{(}
                \StringTok{"SHA256withRSA"}\OperatorTok{).}\FunctionTok{setProvider}\OperatorTok{(}\StringTok{"BC"}\OperatorTok{).}\FunctionTok{build}\OperatorTok{(}\NormalTok{tsaKey}\OperatorTok{);}
\NormalTok{        SignerInfoGenerator sigGen }\OperatorTok{=}\NormalTok{ sigGenBuilder}\OperatorTok{.}\FunctionTok{build}\OperatorTok{(}
\NormalTok{                contentSigner}\OperatorTok{,}
                \KeywordTok{new} \FunctionTok{JcaX509CertificateHolder}\OperatorTok{(}\NormalTok{tsaCert}\OperatorTok{));}

\NormalTok{        DigestCalculatorProvider dcProv }\OperatorTok{=}
                \KeywordTok{new} \FunctionTok{JcaDigestCalculatorProviderBuilder}\OperatorTok{()}
                        \OperatorTok{.}\FunctionTok{setProvider}\OperatorTok{(}\StringTok{"BC"}\OperatorTok{).}\FunctionTok{build}\OperatorTok{();}
\NormalTok{        DigestCalculator dgCalc }\OperatorTok{=}
\NormalTok{                dcProv}\OperatorTok{.}\FunctionTok{get}\OperatorTok{(}\KeywordTok{new} \FunctionTok{AlgorithmIdentifier}\OperatorTok{(}\NormalTok{NISTObjectIdentifiers}\OperatorTok{.}\FunctionTok{id\_sha256}\OperatorTok{));}

\NormalTok{        TimeStampTokenGenerator tsTokenGen }\OperatorTok{=}
                \KeywordTok{new} \FunctionTok{TimeStampTokenGenerator}\OperatorTok{(}\NormalTok{sigGen}\OperatorTok{,}\NormalTok{ dgCalc}\OperatorTok{,}
                        \KeywordTok{new} \FunctionTok{ASN1ObjectIdentifier}\OperatorTok{(}\StringTok{"1.2.3.4.5"}\OperatorTok{));} \CommentTok{// Test policy}
\NormalTok{        tsTokenGen}\OperatorTok{.}\FunctionTok{addCertificates}\OperatorTok{(}
                \KeywordTok{new} \FunctionTok{JcaCertStore}\OperatorTok{(}\BuiltInTok{List}\OperatorTok{.}\FunctionTok{of}\OperatorTok{(}\NormalTok{tsaCert}\OperatorTok{)));}

\NormalTok{        TimeStampResponseGenerator tsRespGen }\OperatorTok{=}
                \KeywordTok{new} \FunctionTok{TimeStampResponseGenerator}\OperatorTok{(}\NormalTok{tsTokenGen}\OperatorTok{,}
\NormalTok{                        TSPAlgorithms}\OperatorTok{.}\FunctionTok{ALLOWED}\OperatorTok{);}

\NormalTok{        TimeStampResponse response }\OperatorTok{=}\NormalTok{ tsRespGen}\OperatorTok{.}\FunctionTok{generate}\OperatorTok{(}
\NormalTok{                request}\OperatorTok{,}\NormalTok{ serialNumber}\OperatorTok{.}\FunctionTok{getAndIncrement}\OperatorTok{(),} \KeywordTok{new} \BuiltInTok{Date}\OperatorTok{());}

\NormalTok{        log}\OperatorTok{.}\FunctionTok{info}\OperatorTok{(}\StringTok{"Generated test timestamp — serial: \{\}, time: \{\}"}\OperatorTok{,}
\NormalTok{                serialNumber}\OperatorTok{.}\FunctionTok{get}\OperatorTok{()} \OperatorTok{{-}} \DecValTok{1}\OperatorTok{,} \KeywordTok{new} \BuiltInTok{Date}\OperatorTok{());}
        \ControlFlowTok{return}\NormalTok{ response}\OperatorTok{.}\FunctionTok{getTimeStampToken}\OperatorTok{();}
    \OperatorTok{\}}
\OperatorTok{\}}
\end{Highlighting}
\end{Shaded}

\subsection{Mocking OCSP with
WireMock}\label{_mocking_ocsp_with_wiremock}

For revocation checking tests, we use WireMock to simulate an OCSP
responder. This lets us return "good," "revoked," or "unknown" responses
without any real PKI:

\begin{Shaded}
\begin{Highlighting}[]
\CommentTok{/**}
 \CommentTok{*}\NormalTok{ Configure a WireMock OCSP responder stub}\CommentTok{.}
 \CommentTok{*}\NormalTok{ Returns a pre}\CommentTok{{-}}\NormalTok{built OCSP response for any request}\CommentTok{.}
 \CommentTok{*/}
\KeywordTok{public} \DataTypeTok{static} \DataTypeTok{void} \FunctionTok{stubOcspResponder}\OperatorTok{(}\NormalTok{WireMockServer wireMock}\OperatorTok{,}
\NormalTok{        OCSPResp response}\OperatorTok{)} \KeywordTok{throws} \BuiltInTok{Exception} \OperatorTok{\{}

\NormalTok{    wireMock}\OperatorTok{.}\FunctionTok{stubFor}\OperatorTok{(}\FunctionTok{post}\OperatorTok{(}\FunctionTok{urlEqualTo}\OperatorTok{(}\StringTok{"/ocsp"}\OperatorTok{))}
            \OperatorTok{.}\FunctionTok{withHeader}\OperatorTok{(}\StringTok{"Content{-}Type"}\OperatorTok{,}
                    \FunctionTok{equalTo}\OperatorTok{(}\StringTok{"application/ocsp{-}request"}\OperatorTok{))}
            \OperatorTok{.}\FunctionTok{willReturn}\OperatorTok{(}\FunctionTok{aResponse}\OperatorTok{()}
                    \OperatorTok{.}\FunctionTok{withStatus}\OperatorTok{(}\DecValTok{200}\OperatorTok{)}
                    \OperatorTok{.}\FunctionTok{withHeader}\OperatorTok{(}\StringTok{"Content{-}Type"}\OperatorTok{,}
                            \StringTok{"application/ocsp{-}response"}\OperatorTok{)}
                    \OperatorTok{.}\FunctionTok{withBody}\OperatorTok{(}\NormalTok{response}\OperatorTok{.}\FunctionTok{getEncoded}\OperatorTok{())));}
\OperatorTok{\}}

\CommentTok{/**}
 \CommentTok{*}\NormalTok{ Build a }\CommentTok{"}\NormalTok{good}\CommentTok{"}\NormalTok{ OCSP response for testing}\CommentTok{.}
 \CommentTok{*/}
\KeywordTok{public} \DataTypeTok{static}\NormalTok{ OCSPResp }\FunctionTok{buildGoodOcspResponse}\OperatorTok{(}
        \BuiltInTok{X509Certificate}\NormalTok{ targetCert}\OperatorTok{,}
        \BuiltInTok{X509Certificate}\NormalTok{ issuerCert}\OperatorTok{,}
        \BuiltInTok{PrivateKey}\NormalTok{ responderKey}\OperatorTok{,}
        \BuiltInTok{X509Certificate}\NormalTok{ responderCert}\OperatorTok{)} \KeywordTok{throws} \BuiltInTok{Exception} \OperatorTok{\{}

\NormalTok{    BasicOCSPRespBuilder respBuilder }\OperatorTok{=}
            \KeywordTok{new} \FunctionTok{JcaBasicOCSPRespBuilder}\OperatorTok{(}
\NormalTok{                    responderCert}\OperatorTok{.}\FunctionTok{getPublicKey}\OperatorTok{(),}
                    \KeywordTok{new} \FunctionTok{JcaDigestCalculatorProviderBuilder}\OperatorTok{()}
                            \OperatorTok{.}\FunctionTok{setProvider}\OperatorTok{(}\StringTok{"BC"}\OperatorTok{).}\FunctionTok{build}\OperatorTok{()}
                            \OperatorTok{.}\FunctionTok{get}\OperatorTok{(}\NormalTok{CertificateID}\OperatorTok{.}\FunctionTok{HASH\_SHA1}\OperatorTok{));}

\NormalTok{    CertificateID certId }\OperatorTok{=} \KeywordTok{new} \FunctionTok{JcaCertificateID}\OperatorTok{(}
            \KeywordTok{new} \FunctionTok{JcaDigestCalculatorProviderBuilder}\OperatorTok{()}
                    \OperatorTok{.}\FunctionTok{setProvider}\OperatorTok{(}\StringTok{"BC"}\OperatorTok{).}\FunctionTok{build}\OperatorTok{()}
                    \OperatorTok{.}\FunctionTok{get}\OperatorTok{(}\NormalTok{CertificateID}\OperatorTok{.}\FunctionTok{HASH\_SHA1}\OperatorTok{),}
            \KeywordTok{new} \FunctionTok{JcaX509CertificateHolder}\OperatorTok{(}\NormalTok{issuerCert}\OperatorTok{),}
\NormalTok{            targetCert}\OperatorTok{.}\FunctionTok{getSerialNumber}\OperatorTok{());}

\NormalTok{    respBuilder}\OperatorTok{.}\FunctionTok{addResponse}\OperatorTok{(}\NormalTok{certId}\OperatorTok{,}\NormalTok{ CertificateStatus}\OperatorTok{.}\FunctionTok{GOOD}\OperatorTok{);}

\NormalTok{    ContentSigner signer }\OperatorTok{=} \KeywordTok{new} \FunctionTok{JcaContentSignerBuilder}\OperatorTok{(}\StringTok{"SHA256withRSA"}\OperatorTok{)}
            \OperatorTok{.}\FunctionTok{setProvider}\OperatorTok{(}\StringTok{"BC"}\OperatorTok{).}\FunctionTok{build}\OperatorTok{(}\NormalTok{responderKey}\OperatorTok{);}
\NormalTok{    BasicOCSPResp basicResp }\OperatorTok{=}\NormalTok{ respBuilder}\OperatorTok{.}\FunctionTok{build}\OperatorTok{(}\NormalTok{signer}\OperatorTok{,}
            \KeywordTok{new}\NormalTok{ X509CertificateHolder}\OperatorTok{[]\{} \KeywordTok{new} \FunctionTok{JcaX509CertificateHolder}\OperatorTok{(}
\NormalTok{                    responderCert}\OperatorTok{)} \OperatorTok{\},}
            \KeywordTok{new} \BuiltInTok{Date}\OperatorTok{());}

    \ControlFlowTok{return} \KeywordTok{new} \FunctionTok{OCSPRespBuilder}\OperatorTok{().}\FunctionTok{build}\OperatorTok{(}
\NormalTok{            OCSPRespBuilder}\OperatorTok{.}\FunctionTok{SUCCESSFUL}\OperatorTok{,}\NormalTok{ basicResp}\OperatorTok{);}
\OperatorTok{\}}
\end{Highlighting}
\end{Shaded}

\section{Time Travel in Tests}\label{_time_travel_in_tests}

\subsection{The Clock Abstraction}\label{_the_clock_abstraction}

The most reliable way to handle time-dependent tests is to inject a
\texttt{Clock} into any component that checks certificate validity or
timestamp times:

\begin{Shaded}
\begin{Highlighting}[]
\CommentTok{/**}
 \CommentTok{*}\NormalTok{ Verification service that accepts an injectable Clock}
 \CommentTok{*}\NormalTok{ for deterministic testing}\CommentTok{.}
 \CommentTok{*/}
\AttributeTok{@Slf4j}
\KeywordTok{public} \KeywordTok{class}\NormalTok{ VerificationService }\OperatorTok{\{}

    \KeywordTok{private} \DataTypeTok{final}\NormalTok{ Clock clock}\OperatorTok{;}

    \CommentTok{/**}
     \CommentTok{*}\NormalTok{ Production constructor }\CommentTok{—}\NormalTok{ uses system clock}\CommentTok{.}
     \CommentTok{*/}
    \KeywordTok{public} \FunctionTok{VerificationService}\OperatorTok{()} \OperatorTok{\{}
        \KeywordTok{this}\OperatorTok{.}\FunctionTok{clock} \OperatorTok{=}\NormalTok{ Clock}\OperatorTok{.}\FunctionTok{systemUTC}\OperatorTok{();}
    \OperatorTok{\}}

    \CommentTok{/**}
     \CommentTok{*}\NormalTok{ Test constructor }\CommentTok{—}\NormalTok{ uses injected clock}\CommentTok{.}
     \CommentTok{*/}
    \KeywordTok{public} \FunctionTok{VerificationService}\OperatorTok{(}\NormalTok{Clock clock}\OperatorTok{)} \OperatorTok{\{}
        \KeywordTok{this}\OperatorTok{.}\FunctionTok{clock} \OperatorTok{=}\NormalTok{ clock}\OperatorTok{;}
    \OperatorTok{\}}

    \KeywordTok{public}\NormalTok{ VerificationReport }\FunctionTok{verify}\OperatorTok{(}\DataTypeTok{byte}\OperatorTok{[]}\NormalTok{ pdfBytes}\OperatorTok{,}
\NormalTok{            PDSignature signature}\OperatorTok{,}
            \BuiltInTok{Set}\OperatorTok{\textless{}}\BuiltInTok{TrustAnchor}\OperatorTok{\textgreater{}}\NormalTok{ trustAnchors}\OperatorTok{)} \OperatorTok{\{}
        \CommentTok{// ... all validation uses this.clock instead of new Date()}
        \BuiltInTok{Date}\NormalTok{ validationTime }\OperatorTok{=} \BuiltInTok{Date}\OperatorTok{.}\FunctionTok{from}\OperatorTok{(}\NormalTok{clock}\OperatorTok{.}\FunctionTok{instant}\OperatorTok{());}
        \CommentTok{// ...}
    \OperatorTok{\}}
\OperatorTok{\}}
\end{Highlighting}
\end{Shaded}

\subsection{Testing Expired
Certificates}\label{_testing_expired_certificates}

With the \texttt{Clock} abstraction, testing expired certificate
scenarios becomes straightforward:

\begin{Shaded}
\begin{Highlighting}[]
\AttributeTok{@Test}
\DataTypeTok{void} \FunctionTok{expiredCertificateWithoutTimestampShouldFail}\OperatorTok{()} \KeywordTok{throws} \BuiltInTok{Exception} \OperatorTok{\{}
    \CommentTok{// Generate a certificate valid for only 1 second}
    \BuiltInTok{KeyPair}\NormalTok{ keyPair }\OperatorTok{=} \FunctionTok{generateKeyPair}\OperatorTok{();}
    \BuiltInTok{X509Certificate}\NormalTok{ shortLivedCert }\OperatorTok{=}
            \FunctionTok{generateCertWithValidity}\OperatorTok{(}\NormalTok{keyPair}\OperatorTok{,} \BuiltInTok{Duration}\OperatorTok{.}\FunctionTok{ofSeconds}\OperatorTok{(}\DecValTok{1}\OperatorTok{));}

    \CommentTok{// Sign the PDF}
    \DataTypeTok{byte}\OperatorTok{[]}\NormalTok{ pdf }\OperatorTok{=} \FunctionTok{generateTestPdf}\OperatorTok{(}\StringTok{"Short{-}lived cert test"}\OperatorTok{);}
    \DataTypeTok{byte}\OperatorTok{[]}\NormalTok{ signed }\OperatorTok{=} \FunctionTok{signPdfWith}\OperatorTok{(}\NormalTok{pdf}\OperatorTok{,}\NormalTok{ keyPair}\OperatorTok{.}\FunctionTok{getPrivate}\OperatorTok{(),}
            \KeywordTok{new} \BuiltInTok{X509Certificate}\OperatorTok{[]\{}\NormalTok{ shortLivedCert }\OperatorTok{\});}

    \CommentTok{// Travel forward in time — certificate is now expired}
\NormalTok{    Clock futureClock }\OperatorTok{=}\NormalTok{ Clock}\OperatorTok{.}\FunctionTok{fixed}\OperatorTok{(}
\NormalTok{            Instant}\OperatorTok{.}\FunctionTok{now}\OperatorTok{().}\FunctionTok{plus}\OperatorTok{(}\BuiltInTok{Duration}\OperatorTok{.}\FunctionTok{ofDays}\OperatorTok{(}\DecValTok{1}\OperatorTok{)),}
\NormalTok{            ZoneOffset}\OperatorTok{.}\FunctionTok{UTC}\OperatorTok{);}
\NormalTok{    VerificationService verifier }\OperatorTok{=} \KeywordTok{new} \FunctionTok{VerificationService}\OperatorTok{(}\NormalTok{futureClock}\OperatorTok{);}

\NormalTok{    VerificationReport report }\OperatorTok{=}\NormalTok{ verifier}\OperatorTok{.}\FunctionTok{verify}\OperatorTok{(}
\NormalTok{            signed}\OperatorTok{,} \FunctionTok{extractSignature}\OperatorTok{(}\NormalTok{signed}\OperatorTok{),}\NormalTok{ trustAnchors}\OperatorTok{);}

    \FunctionTok{assertEquals}\OperatorTok{(}\NormalTok{ValidationResult}\OperatorTok{.}\FunctionTok{OK}\OperatorTok{,}\NormalTok{ report}\OperatorTok{.}\FunctionTok{getHashResult}\OperatorTok{());}
    \FunctionTok{assertEquals}\OperatorTok{(}\NormalTok{ValidationResult}\OperatorTok{.}\FunctionTok{OK}\OperatorTok{,}\NormalTok{ report}\OperatorTok{.}\FunctionTok{getCmsResult}\OperatorTok{());}
    \FunctionTok{assertEquals}\OperatorTok{(}\NormalTok{ValidationResult}\OperatorTok{.}\FunctionTok{CHAIN\_INVALID}\OperatorTok{,}
\NormalTok{            report}\OperatorTok{.}\FunctionTok{getChainResult}\OperatorTok{());}
\OperatorTok{\}}

\AttributeTok{@Test}
\DataTypeTok{void} \FunctionTok{expiredCertificateWithTimestampShouldPass}\OperatorTok{()} \KeywordTok{throws} \BuiltInTok{Exception} \OperatorTok{\{}
    \CommentTok{// Same setup, but this time we add a timestamp}
    \DataTypeTok{byte}\OperatorTok{[]}\NormalTok{ signed }\OperatorTok{=} \FunctionTok{signPdfWith}\OperatorTok{(}\NormalTok{pdf}\OperatorTok{,}\NormalTok{ keyPair}\OperatorTok{.}\FunctionTok{getPrivate}\OperatorTok{(),}
            \KeywordTok{new} \BuiltInTok{X509Certificate}\OperatorTok{[]\{}\NormalTok{ shortLivedCert }\OperatorTok{\},}
            \KeywordTok{true} \CommentTok{/* addTimestamp */}\OperatorTok{);}

    \CommentTok{// The timestamp proves the signature was created}
    \CommentTok{// before the cert expired}
\NormalTok{    Clock futureClock }\OperatorTok{=}\NormalTok{ Clock}\OperatorTok{.}\FunctionTok{fixed}\OperatorTok{(}
\NormalTok{            Instant}\OperatorTok{.}\FunctionTok{now}\OperatorTok{().}\FunctionTok{plus}\OperatorTok{(}\BuiltInTok{Duration}\OperatorTok{.}\FunctionTok{ofDays}\OperatorTok{(}\DecValTok{1}\OperatorTok{)),}
\NormalTok{            ZoneOffset}\OperatorTok{.}\FunctionTok{UTC}\OperatorTok{);}
\NormalTok{    VerificationService verifier }\OperatorTok{=} \KeywordTok{new} \FunctionTok{VerificationService}\OperatorTok{(}\NormalTok{futureClock}\OperatorTok{);}

\NormalTok{    VerificationReport report }\OperatorTok{=}\NormalTok{ verifier}\OperatorTok{.}\FunctionTok{verify}\OperatorTok{(}
\NormalTok{            signed}\OperatorTok{,} \FunctionTok{extractSignature}\OperatorTok{(}\NormalTok{signed}\OperatorTok{),}\NormalTok{ trustAnchors}\OperatorTok{);}

    \FunctionTok{assertEquals}\OperatorTok{(}\NormalTok{ValidationResult}\OperatorTok{.}\FunctionTok{OK}\OperatorTok{,}\NormalTok{ report}\OperatorTok{.}\FunctionTok{getChainResult}\OperatorTok{());}
    \FunctionTok{assertNotNull}\OperatorTok{(}\NormalTok{report}\OperatorTok{.}\FunctionTok{getTimestampTime}\OperatorTok{());}
\OperatorTok{\}}
\end{Highlighting}
\end{Shaded}

\section{Complete Signing-to-Verification Round-Trip
Test}\label{_complete_signing_to_verification_round_trip_test}

This is the most important test in the suite --- it exercises the entire
pipeline end-to-end:

\begin{Shaded}
\begin{Highlighting}[]
\AttributeTok{@Test}
\DataTypeTok{void} \FunctionTok{completeRoundTripSignAndVerify}\OperatorTok{()} \KeywordTok{throws} \BuiltInTok{Exception} \OperatorTok{\{}
    \CommentTok{// 1. Generate a test PDF}
    \DataTypeTok{byte}\OperatorTok{[]}\NormalTok{ unsignedPdf }\OperatorTok{=} \FunctionTok{generateTestPdf}\OperatorTok{(}
            \StringTok{"Round{-}trip test — "} \OperatorTok{+}\NormalTok{ Instant}\OperatorTok{.}\FunctionTok{now}\OperatorTok{());}
\NormalTok{    log}\OperatorTok{.}\FunctionTok{info}\OperatorTok{(}\StringTok{"Generated unsigned PDF: \{\} bytes"}\OperatorTok{,}\NormalTok{ unsignedPdf}\OperatorTok{.}\FunctionTok{length}\OperatorTok{);}

    \CommentTok{// 2. Sign with the test keystore (PAdES + timestamp)}
    \DataTypeTok{byte}\OperatorTok{[]}\NormalTok{ signedPdf }\OperatorTok{=} \FunctionTok{signPdf}\OperatorTok{(}\NormalTok{unsignedPdf}\OperatorTok{,} \KeywordTok{true}\OperatorTok{,}
            \StringTok{"ETSI.CAdES.detached"}\OperatorTok{);}
\NormalTok{    log}\OperatorTok{.}\FunctionTok{info}\OperatorTok{(}\StringTok{"Generated signed PDF: \{\} bytes"}\OperatorTok{,}\NormalTok{ signedPdf}\OperatorTok{.}\FunctionTok{length}\OperatorTok{);}

    \CommentTok{// 3. Extract the signature}
    \ControlFlowTok{try} \OperatorTok{(}\NormalTok{PDDocument doc }\OperatorTok{=}\NormalTok{ Loader}\OperatorTok{.}\FunctionTok{loadPDF}\OperatorTok{(}\NormalTok{signedPdf}\OperatorTok{))} \OperatorTok{\{}
        \BuiltInTok{List}\OperatorTok{\textless{}}\NormalTok{PDSignature}\OperatorTok{\textgreater{}}\NormalTok{ sigs }\OperatorTok{=}\NormalTok{ doc}\OperatorTok{.}\FunctionTok{getSignatureDictionaries}\OperatorTok{();}
        \FunctionTok{assertEquals}\OperatorTok{(}\DecValTok{1}\OperatorTok{,}\NormalTok{ sigs}\OperatorTok{.}\FunctionTok{size}\OperatorTok{(),} \StringTok{"Expected exactly one signature"}\OperatorTok{);}

\NormalTok{        PDSignature sig }\OperatorTok{=}\NormalTok{ sigs}\OperatorTok{.}\FunctionTok{get}\OperatorTok{(}\DecValTok{0}\OperatorTok{);}
        \FunctionTok{assertEquals}\OperatorTok{(}\StringTok{"ETSI.CAdES.detached"}\OperatorTok{,}
\NormalTok{                sig}\OperatorTok{.}\FunctionTok{getSubFilter}\OperatorTok{());}

        \CommentTok{// 4. Run the complete verification pipeline}
\NormalTok{        VerificationService verifier }\OperatorTok{=} \KeywordTok{new} \FunctionTok{VerificationService}\OperatorTok{();}
\NormalTok{        VerificationReport report }\OperatorTok{=}
\NormalTok{                verifier}\OperatorTok{.}\FunctionTok{verify}\OperatorTok{(}\NormalTok{signedPdf}\OperatorTok{,}\NormalTok{ sig}\OperatorTok{,}\NormalTok{ trustAnchors}\OperatorTok{);}

        \CommentTok{// 5. Assert every step passed}
        \FunctionTok{assertEquals}\OperatorTok{(}\NormalTok{ValidationResult}\OperatorTok{.}\FunctionTok{OK}\OperatorTok{,}
\NormalTok{                report}\OperatorTok{.}\FunctionTok{getByteRangeResult}\OperatorTok{(),}
                \StringTok{"ByteRange check should pass"}\OperatorTok{);}
        \FunctionTok{assertEquals}\OperatorTok{(}\NormalTok{ValidationResult}\OperatorTok{.}\FunctionTok{OK}\OperatorTok{,}
\NormalTok{                report}\OperatorTok{.}\FunctionTok{getHashResult}\OperatorTok{(),}
                \StringTok{"Hash verification should pass"}\OperatorTok{);}
        \FunctionTok{assertEquals}\OperatorTok{(}\NormalTok{ValidationResult}\OperatorTok{.}\FunctionTok{OK}\OperatorTok{,}
\NormalTok{                report}\OperatorTok{.}\FunctionTok{getCmsResult}\OperatorTok{(),}
                \StringTok{"CMS verification should pass"}\OperatorTok{);}
        \FunctionTok{assertEquals}\OperatorTok{(}\NormalTok{ValidationResult}\OperatorTok{.}\FunctionTok{OK}\OperatorTok{,}
\NormalTok{                report}\OperatorTok{.}\FunctionTok{getChainResult}\OperatorTok{(),}
                \StringTok{"Certificate chain should validate"}\OperatorTok{);}
        \FunctionTok{assertEquals}\OperatorTok{(}\NormalTok{ValidationResult}\OperatorTok{.}\FunctionTok{OK}\OperatorTok{,}
\NormalTok{                report}\OperatorTok{.}\FunctionTok{getTrustResult}\OperatorTok{(),}
                \StringTok{"Trust evaluation should pass"}\OperatorTok{);}

        \CommentTok{// 6. Verify timestamp is present and valid}
        \FunctionTok{assertNotNull}\OperatorTok{(}\NormalTok{report}\OperatorTok{.}\FunctionTok{getTimestampTime}\OperatorTok{(),}
                \StringTok{"Timestamp should be present"}\OperatorTok{);}
        \FunctionTok{assertEquals}\OperatorTok{(}\NormalTok{ValidationResult}\OperatorTok{.}\FunctionTok{OK}\OperatorTok{,}
\NormalTok{                report}\OperatorTok{.}\FunctionTok{getTimestampResult}\OperatorTok{(),}
                \StringTok{"Timestamp should be valid"}\OperatorTok{);}

        \FunctionTok{assertTrue}\OperatorTok{(}\NormalTok{report}\OperatorTok{.}\FunctionTok{isValid}\OperatorTok{(),} \StringTok{"Signature should be valid"}\OperatorTok{);}
        \FunctionTok{assertTrue}\OperatorTok{(}\NormalTok{report}\OperatorTok{.}\FunctionTok{isTrusted}\OperatorTok{(),} \StringTok{"Signature should be trusted"}\OperatorTok{);}

\NormalTok{        log}\OperatorTok{.}\FunctionTok{info}\OperatorTok{(}\StringTok{"Round{-}trip test PASSED — all steps verified"}\OperatorTok{);}
    \OperatorTok{\}}
\OperatorTok{\}}
\end{Highlighting}
\end{Shaded}

\section{PAdES Compliance Test}\label{_pades_compliance_test}

A PAdES-compliant signature must include specific signed attributes
beyond what plain CMS requires. We test for these explicitly:

\begin{Shaded}
\begin{Highlighting}[]
\AttributeTok{@Test}
\DataTypeTok{void} \FunctionTok{padesSignatureShouldHaveRequiredAttributes}\OperatorTok{()} \KeywordTok{throws} \BuiltInTok{Exception} \OperatorTok{\{}
    \DataTypeTok{byte}\OperatorTok{[]}\NormalTok{ pdf }\OperatorTok{=} \FunctionTok{generateTestPdf}\OperatorTok{(}\StringTok{"PAdES compliance test"}\OperatorTok{);}
    \DataTypeTok{byte}\OperatorTok{[]}\NormalTok{ signed }\OperatorTok{=} \FunctionTok{signPdf}\OperatorTok{(}\NormalTok{pdf}\OperatorTok{,} \KeywordTok{true}\OperatorTok{,} \StringTok{"ETSI.CAdES.detached"}\OperatorTok{);}

    \ControlFlowTok{try} \OperatorTok{(}\NormalTok{PDDocument doc }\OperatorTok{=}\NormalTok{ Loader}\OperatorTok{.}\FunctionTok{loadPDF}\OperatorTok{(}\NormalTok{signed}\OperatorTok{))} \OperatorTok{\{}
\NormalTok{        PDSignature sig }\OperatorTok{=}\NormalTok{ doc}\OperatorTok{.}\FunctionTok{getSignatureDictionaries}\OperatorTok{().}\FunctionTok{get}\OperatorTok{(}\DecValTok{0}\OperatorTok{);}
\NormalTok{        CMSSignedData cmsData }\OperatorTok{=} \FunctionTok{parseCmsSignedData}\OperatorTok{(}\NormalTok{signed}\OperatorTok{,}\NormalTok{ sig}\OperatorTok{);}
\NormalTok{        SignerInformation signerInfo }\OperatorTok{=}\NormalTok{ cmsData}\OperatorTok{.}\FunctionTok{getSignerInfos}\OperatorTok{()}
                \OperatorTok{.}\FunctionTok{getSigners}\OperatorTok{().}\FunctionTok{iterator}\OperatorTok{().}\FunctionTok{next}\OperatorTok{();}

\NormalTok{        AttributeTable signedAttrs }\OperatorTok{=}
\NormalTok{                signerInfo}\OperatorTok{.}\FunctionTok{getSignedAttributes}\OperatorTok{();}

        \CommentTok{// PAdES{-}B{-}B: Required signed attributes}
        \FunctionTok{assertNotNull}\OperatorTok{(}
\NormalTok{                signedAttrs}\OperatorTok{.}\FunctionTok{get}\OperatorTok{(}\NormalTok{CMSAttributes}\OperatorTok{.}\FunctionTok{contentType}\OperatorTok{),}
                \StringTok{"contentType attribute is required"}\OperatorTok{);}
        \FunctionTok{assertNotNull}\OperatorTok{(}
\NormalTok{                signedAttrs}\OperatorTok{.}\FunctionTok{get}\OperatorTok{(}\NormalTok{CMSAttributes}\OperatorTok{.}\FunctionTok{messageDigest}\OperatorTok{),}
                \StringTok{"messageDigest attribute is required"}\OperatorTok{);}
        \FunctionTok{assertNotNull}\OperatorTok{(}
\NormalTok{                signedAttrs}\OperatorTok{.}\FunctionTok{get}\OperatorTok{(}
\NormalTok{                    PKCSObjectIdentifiers}\OperatorTok{.}\FunctionTok{id\_aa\_signingCertificateV2}\OperatorTok{),}
                \StringTok{"ESS{-}signing{-}certificate{-}v2 is required for PAdES"}\OperatorTok{);}

        \CommentTok{// PAdES{-}B{-}B: signingTime SHOULD NOT be present}
        \CommentTok{// (PAdES relies on timestamps, not claimed signing time)}
        \BuiltInTok{Attribute}\NormalTok{ signingTime }\OperatorTok{=}
\NormalTok{                signedAttrs}\OperatorTok{.}\FunctionTok{get}\OperatorTok{(}\NormalTok{CMSAttributes}\OperatorTok{.}\FunctionTok{signingTime}\OperatorTok{);}
        \ControlFlowTok{if} \OperatorTok{(}\NormalTok{signingTime }\OperatorTok{!=} \KeywordTok{null}\OperatorTok{)} \OperatorTok{\{}
\NormalTok{            log}\OperatorTok{.}\FunctionTok{warn}\OperatorTok{(}\StringTok{"signingTime attribute is present — "} \OperatorTok{+}
                    \StringTok{"PAdES recommends omitting it in favor "} \OperatorTok{+}
                    \StringTok{"of a timestamp"}\OperatorTok{);}
        \OperatorTok{\}}

        \CommentTok{// PAdES{-}B{-}T: Timestamp unsigned attribute}
\NormalTok{        AttributeTable unsignedAttrs }\OperatorTok{=}
\NormalTok{                signerInfo}\OperatorTok{.}\FunctionTok{getUnsignedAttributes}\OperatorTok{();}
        \FunctionTok{assertNotNull}\OperatorTok{(}\NormalTok{unsignedAttrs}\OperatorTok{,}
                \StringTok{"PAdES{-}B{-}T requires unsigned attributes"}\OperatorTok{);}
        \FunctionTok{assertNotNull}\OperatorTok{(}
\NormalTok{                unsignedAttrs}\OperatorTok{.}\FunctionTok{get}\OperatorTok{(}
\NormalTok{                    PKCSObjectIdentifiers}
                        \OperatorTok{.}\FunctionTok{id\_aa\_signatureTimeStampToken}\OperatorTok{),}
                \StringTok{"PAdES{-}B{-}T requires a signature timestamp"}\OperatorTok{);}

        \CommentTok{// Verify the digest algorithm is SHA{-}256 or stronger}
        \BuiltInTok{String}\NormalTok{ digestAlg }\OperatorTok{=}\NormalTok{ signerInfo}\OperatorTok{.}\FunctionTok{getDigestAlgOID}\OperatorTok{();}
        \FunctionTok{assertNotEquals}\OperatorTok{(}\StringTok{"1.3.14.3.2.26"}\OperatorTok{,}\NormalTok{ digestAlg}\OperatorTok{,}
                \StringTok{"SHA{-}1 is not allowed for PAdES — "} \OperatorTok{+}
                \StringTok{"must use SHA{-}256 or stronger"}\OperatorTok{);}

\NormalTok{        log}\OperatorTok{.}\FunctionTok{info}\OperatorTok{(}\StringTok{"PAdES compliance test PASSED"}\OperatorTok{);}
    \OperatorTok{\}}
\OperatorTok{\}}
\end{Highlighting}
\end{Shaded}

\section{Reference PDF Corpus: Naming
Convention}\label{_reference_pdf_corpus_naming_convention}

For reference PDFs that you ship with your test suite (for regression
testing or interoperability testing), use a consistent naming convention
that encodes the signature properties in the filename:

\begin{verbatim}
Naming pattern:
  <scenario>_<subfilter>_<timestamp>_<cert-type>_<extras>.pdf

Examples:
  basic_pkcs7-detached_no-ts_ca-signed_v1.pdf
  basic_cades-detached_rfc3161_ca-signed_v1.pdf
  pades-b-t_cades-detached_rfc3161_ca-signed_v1.pdf
  pades-b-lt_cades-detached_rfc3161_ca-signed_dss_v1.pdf
  selfsigned_pkcs7-detached_no-ts_self-signed_v1.pdf
  multi-sig_pkcs7-detached_no-ts_ca-signed_2sigs_v1.pdf
  tampered_pkcs7-detached_no-ts_ca-signed_modified_v1.pdf
  invisible_cades-detached_rfc3161_ca-signed_zero-widget_v1.pdf
\end{verbatim}

The \texttt{\_v1} suffix is a version number. When you regenerate
reference PDFs (because you changed the test keystore or signing code),
increment the version number. This ensures old and new reference files
can coexist temporarily during migration.

Store reference PDFs in \texttt{src/test/resources/reference-pdfs/} and
regenerate them with a dedicated test utility class. Never edit
reference PDFs by hand --- they should always be produced by code.
Commit the generated PDFs to your repository so tests can run without a
generation step, but include the generator so anyone can rebuild them.

\section{Key Takeaways}\label{_key_takeaways_17}

\begin{itemize}
\item
  PDF signature testing is hard because of three enemies: time, network,
  and state --- address all three explicitly
\item
  Generate test keystores programmatically with long validity periods
  (100 years) to avoid time-rot
\item
  Use an in-process TSA built from BouncyCastle instead of hitting real
  TSA servers
\item
  Mock OCSP responders with WireMock for deterministic revocation
  testing
\item
  Inject a \texttt{Clock} abstraction into verification services so you
  can freeze or advance time in tests
\item
  Build a test matrix that maps every scenario to expected per-step
  results
\item
  The round-trip test (generate PDF, sign, extract, verify) is the
  single most important test in the suite
\item
  Test PAdES compliance explicitly: check for
  \texttt{ESS-signing-certificate-v2}, proper SubFilter, and SHA-256+
\item
  Use a consistent naming convention for reference PDFs that encodes all
  signature properties in the filename
\item
  Generate reference PDFs from code, not by hand --- commit both the
  generator and the generated files
\end{itemize}

\section{Complete Code Listings}\label{_complete_code_listings}

This appendix collects the complete, compilable source code for the
major classes developed throughout Parts 5 and 6. Each listing is
self-contained --- you can copy it into your project, adjust the package
name, and compile. All code targets Java 21, Apache PDFBox 3.0.0, and
BouncyCastle 1.80.

\section{SigningService: Complete PDF Signing
Implementation}\label{_signingservice_complete_pdf_signing_implementation}

This is the full signing service from Chapters 29 and 30, consolidated
into a single class. It handles both \texttt{adbe.pkcs7.detached} and
\texttt{ETSI.CAdES.detached} signatures, with optional RFC 3161
timestamps.

\begin{Shaded}
\begin{Highlighting}[]
\KeywordTok{package}\ImportTok{ io}\OperatorTok{.}\ImportTok{pdfalyzer}\OperatorTok{.}\ImportTok{signing}\OperatorTok{;}

\KeywordTok{import} \ImportTok{java}\OperatorTok{.}\ImportTok{io}\OperatorTok{.}\ImportTok{ByteArrayOutputStream}\OperatorTok{;}
\KeywordTok{import} \ImportTok{java}\OperatorTok{.}\ImportTok{io}\OperatorTok{.}\ImportTok{InputStream}\OperatorTok{;}
\KeywordTok{import} \ImportTok{java}\OperatorTok{.}\ImportTok{io}\OperatorTok{.}\ImportTok{OutputStream}\OperatorTok{;}
\KeywordTok{import} \ImportTok{java}\OperatorTok{.}\ImportTok{net}\OperatorTok{.}\ImportTok{HttpURLConnection}\OperatorTok{;}
\KeywordTok{import} \ImportTok{java}\OperatorTok{.}\ImportTok{net}\OperatorTok{.}\ImportTok{URL}\OperatorTok{;}
\KeywordTok{import} \ImportTok{java}\OperatorTok{.}\ImportTok{security}\OperatorTok{.}\ImportTok{MessageDigest}\OperatorTok{;}
\KeywordTok{import} \ImportTok{java}\OperatorTok{.}\ImportTok{security}\OperatorTok{.}\ImportTok{PrivateKey}\OperatorTok{;}
\KeywordTok{import} \ImportTok{java}\OperatorTok{.}\ImportTok{security}\OperatorTok{.}\ImportTok{cert}\OperatorTok{.}\ImportTok{Certificate}\OperatorTok{;}
\KeywordTok{import} \ImportTok{java}\OperatorTok{.}\ImportTok{security}\OperatorTok{.}\ImportTok{cert}\OperatorTok{.}\ImportTok{X509Certificate}\OperatorTok{;}
\KeywordTok{import} \ImportTok{java}\OperatorTok{.}\ImportTok{util}\OperatorTok{.}\ImportTok{ArrayList}\OperatorTok{;}
\KeywordTok{import} \ImportTok{java}\OperatorTok{.}\ImportTok{util}\OperatorTok{.}\ImportTok{Calendar}\OperatorTok{;}
\KeywordTok{import} \ImportTok{java}\OperatorTok{.}\ImportTok{util}\OperatorTok{.}\ImportTok{List}\OperatorTok{;}

\KeywordTok{import} \ImportTok{lombok}\OperatorTok{.}\ImportTok{RequiredArgsConstructor}\OperatorTok{;}
\KeywordTok{import} \ImportTok{lombok}\OperatorTok{.}\ImportTok{extern}\OperatorTok{.}\ImportTok{slf4j}\OperatorTok{.}\ImportTok{Slf4j}\OperatorTok{;}
\KeywordTok{import} \ImportTok{org}\OperatorTok{.}\ImportTok{apache}\OperatorTok{.}\ImportTok{pdfbox}\OperatorTok{.}\ImportTok{Loader}\OperatorTok{;}
\KeywordTok{import} \ImportTok{org}\OperatorTok{.}\ImportTok{apache}\OperatorTok{.}\ImportTok{pdfbox}\OperatorTok{.}\ImportTok{cos}\OperatorTok{.}\ImportTok{COSName}\OperatorTok{;}
\KeywordTok{import} \ImportTok{org}\OperatorTok{.}\ImportTok{apache}\OperatorTok{.}\ImportTok{pdfbox}\OperatorTok{.}\ImportTok{pdmodel}\OperatorTok{.}\ImportTok{PDDocument}\OperatorTok{;}
\KeywordTok{import} \ImportTok{org}\OperatorTok{.}\ImportTok{apache}\OperatorTok{.}\ImportTok{pdfbox}\OperatorTok{.}\ImportTok{pdmodel}\OperatorTok{.}\ImportTok{interactive}\OperatorTok{.}\ImportTok{digitalsignature}\OperatorTok{.}\ImportTok{PDSignature}\OperatorTok{;}
\KeywordTok{import} \ImportTok{org}\OperatorTok{.}\ImportTok{apache}\OperatorTok{.}\ImportTok{pdfbox}\OperatorTok{.}\ImportTok{pdmodel}\OperatorTok{.}\ImportTok{interactive}\OperatorTok{.}\ImportTok{digitalsignature}\OperatorTok{.}\ImportTok{SignatureInterface}\OperatorTok{;}
\KeywordTok{import} \ImportTok{org}\OperatorTok{.}\ImportTok{bouncycastle}\OperatorTok{.}\ImportTok{asn1}\OperatorTok{.}\ImportTok{ASN1ObjectIdentifier}\OperatorTok{;}
\KeywordTok{import} \ImportTok{org}\OperatorTok{.}\ImportTok{bouncycastle}\OperatorTok{.}\ImportTok{asn1}\OperatorTok{.}\ImportTok{DERSet}\OperatorTok{;}
\KeywordTok{import} \ImportTok{org}\OperatorTok{.}\ImportTok{bouncycastle}\OperatorTok{.}\ImportTok{asn1}\OperatorTok{.}\ImportTok{cms}\OperatorTok{.}\ImportTok{Attribute}\OperatorTok{;}
\KeywordTok{import} \ImportTok{org}\OperatorTok{.}\ImportTok{bouncycastle}\OperatorTok{.}\ImportTok{asn1}\OperatorTok{.}\ImportTok{cms}\OperatorTok{.}\ImportTok{AttributeTable}\OperatorTok{;}
\KeywordTok{import} \ImportTok{org}\OperatorTok{.}\ImportTok{bouncycastle}\OperatorTok{.}\ImportTok{asn1}\OperatorTok{.}\ImportTok{cms}\OperatorTok{.}\ImportTok{CMSAttributes}\OperatorTok{;}
\KeywordTok{import} \ImportTok{org}\OperatorTok{.}\ImportTok{bouncycastle}\OperatorTok{.}\ImportTok{asn1}\OperatorTok{.}\ImportTok{ess}\OperatorTok{.}\ImportTok{ESSCertIDv2}\OperatorTok{;}
\KeywordTok{import} \ImportTok{org}\OperatorTok{.}\ImportTok{bouncycastle}\OperatorTok{.}\ImportTok{asn1}\OperatorTok{.}\ImportTok{ess}\OperatorTok{.}\ImportTok{SigningCertificateV2}\OperatorTok{;}
\KeywordTok{import} \ImportTok{org}\OperatorTok{.}\ImportTok{bouncycastle}\OperatorTok{.}\ImportTok{asn1}\OperatorTok{.}\ImportTok{nist}\OperatorTok{.}\ImportTok{NISTObjectIdentifiers}\OperatorTok{;}
\KeywordTok{import} \ImportTok{org}\OperatorTok{.}\ImportTok{bouncycastle}\OperatorTok{.}\ImportTok{asn1}\OperatorTok{.}\ImportTok{pkcs}\OperatorTok{.}\ImportTok{PKCSObjectIdentifiers}\OperatorTok{;}
\KeywordTok{import} \ImportTok{org}\OperatorTok{.}\ImportTok{bouncycastle}\OperatorTok{.}\ImportTok{asn1}\OperatorTok{.}\ImportTok{x509}\OperatorTok{.}\ImportTok{AlgorithmIdentifier}\OperatorTok{;}
\KeywordTok{import} \ImportTok{org}\OperatorTok{.}\ImportTok{bouncycastle}\OperatorTok{.}\ImportTok{cert}\OperatorTok{.}\ImportTok{jcajce}\OperatorTok{.}\ImportTok{JcaCertStore}\OperatorTok{;}
\KeywordTok{import} \ImportTok{org}\OperatorTok{.}\ImportTok{bouncycastle}\OperatorTok{.}\ImportTok{cert}\OperatorTok{.}\ImportTok{jcajce}\OperatorTok{.}\ImportTok{JcaX509CertificateHolder}\OperatorTok{;}
\KeywordTok{import} \ImportTok{org}\OperatorTok{.}\ImportTok{bouncycastle}\OperatorTok{.}\ImportTok{cms}\OperatorTok{.}\ImportTok{CMSProcessableByteArray}\OperatorTok{;}
\KeywordTok{import} \ImportTok{org}\OperatorTok{.}\ImportTok{bouncycastle}\OperatorTok{.}\ImportTok{cms}\OperatorTok{.}\ImportTok{CMSSignedData}\OperatorTok{;}
\KeywordTok{import} \ImportTok{org}\OperatorTok{.}\ImportTok{bouncycastle}\OperatorTok{.}\ImportTok{cms}\OperatorTok{.}\ImportTok{CMSSignedDataGenerator}\OperatorTok{;}
\KeywordTok{import} \ImportTok{org}\OperatorTok{.}\ImportTok{bouncycastle}\OperatorTok{.}\ImportTok{cms}\OperatorTok{.}\ImportTok{DefaultSignedAttributeTableGenerator}\OperatorTok{;}
\KeywordTok{import} \ImportTok{org}\OperatorTok{.}\ImportTok{bouncycastle}\OperatorTok{.}\ImportTok{cms}\OperatorTok{.}\ImportTok{SignerInfoGenerator}\OperatorTok{;}
\KeywordTok{import} \ImportTok{org}\OperatorTok{.}\ImportTok{bouncycastle}\OperatorTok{.}\ImportTok{cms}\OperatorTok{.}\ImportTok{SignerInfoGeneratorBuilder}\OperatorTok{;}
\KeywordTok{import} \ImportTok{org}\OperatorTok{.}\ImportTok{bouncycastle}\OperatorTok{.}\ImportTok{cms}\OperatorTok{.}\ImportTok{jcajce}\OperatorTok{.}\ImportTok{JcaSignerInfoGeneratorBuilder}\OperatorTok{;}
\KeywordTok{import} \ImportTok{org}\OperatorTok{.}\ImportTok{bouncycastle}\OperatorTok{.}\ImportTok{operator}\OperatorTok{.}\ImportTok{ContentSigner}\OperatorTok{;}
\KeywordTok{import} \ImportTok{org}\OperatorTok{.}\ImportTok{bouncycastle}\OperatorTok{.}\ImportTok{operator}\OperatorTok{.}\ImportTok{DigestCalculatorProvider}\OperatorTok{;}
\KeywordTok{import} \ImportTok{org}\OperatorTok{.}\ImportTok{bouncycastle}\OperatorTok{.}\ImportTok{operator}\OperatorTok{.}\ImportTok{jcajce}\OperatorTok{.}\ImportTok{JcaContentSignerBuilder}\OperatorTok{;}
\KeywordTok{import} \ImportTok{org}\OperatorTok{.}\ImportTok{bouncycastle}\OperatorTok{.}\ImportTok{operator}\OperatorTok{.}\ImportTok{jcajce}\OperatorTok{.}\ImportTok{JcaDigestCalculatorProviderBuilder}\OperatorTok{;}
\KeywordTok{import} \ImportTok{org}\OperatorTok{.}\ImportTok{bouncycastle}\OperatorTok{.}\ImportTok{tsp}\OperatorTok{.}\ImportTok{TSPAlgorithms}\OperatorTok{;}
\KeywordTok{import} \ImportTok{org}\OperatorTok{.}\ImportTok{bouncycastle}\OperatorTok{.}\ImportTok{tsp}\OperatorTok{.}\ImportTok{TimeStampRequest}\OperatorTok{;}
\KeywordTok{import} \ImportTok{org}\OperatorTok{.}\ImportTok{bouncycastle}\OperatorTok{.}\ImportTok{tsp}\OperatorTok{.}\ImportTok{TimeStampRequestGenerator}\OperatorTok{;}
\KeywordTok{import} \ImportTok{org}\OperatorTok{.}\ImportTok{bouncycastle}\OperatorTok{.}\ImportTok{tsp}\OperatorTok{.}\ImportTok{TimeStampResponse}\OperatorTok{;}
\KeywordTok{import} \ImportTok{org}\OperatorTok{.}\ImportTok{bouncycastle}\OperatorTok{.}\ImportTok{tsp}\OperatorTok{.}\ImportTok{TimeStampToken}\OperatorTok{;}

\CommentTok{/**}
 \CommentTok{*}\NormalTok{ Complete PDF signing service}\CommentTok{.}
 \CommentTok{*}\NormalTok{ Supports PKCS}\CommentTok{\#7}\NormalTok{ detached and PAdES CAdES}\CommentTok{{-}}\NormalTok{detached signatures}
 \CommentTok{*}\NormalTok{ with optional RFC }\CommentTok{3161}\NormalTok{ timestamps}\CommentTok{.}
 \CommentTok{*/}
\AttributeTok{@Slf4j}
\AttributeTok{@RequiredArgsConstructor}
\KeywordTok{public} \KeywordTok{class}\NormalTok{ SigningService }\OperatorTok{\{}

    \KeywordTok{private} \DataTypeTok{final} \BuiltInTok{PrivateKey}\NormalTok{ privateKey}\OperatorTok{;}
    \KeywordTok{private} \DataTypeTok{final} \BuiltInTok{X509Certificate}\OperatorTok{[]}\NormalTok{ certificateChain}\OperatorTok{;}
    \KeywordTok{private} \DataTypeTok{final} \BuiltInTok{String}\NormalTok{ tsaUrl}\OperatorTok{;} \CommentTok{// null if no timestamp desired}

    \CommentTok{/**}
     \CommentTok{*}\NormalTok{ Sign a PDF document and return the signed bytes}\CommentTok{.}
     \CommentTok{*}
     \CommentTok{*} \AnnotationTok{@param}\CommentTok{ }\AnnotationTok{pdfBytes}\CommentTok{    }\NormalTok{the unsigned PDF}
     \CommentTok{*} \AnnotationTok{@param}\CommentTok{ }\AnnotationTok{subFilter}\CommentTok{   "}\NormalTok{adbe}\CommentTok{.}\NormalTok{pkcs7}\CommentTok{.}\NormalTok{detached}\CommentTok{"}\NormalTok{ or }\CommentTok{"}\NormalTok{ETSI}\CommentTok{.}\NormalTok{CAdES}\CommentTok{.}\NormalTok{detached}\CommentTok{"}
     \CommentTok{*} \AnnotationTok{@param}\CommentTok{ }\AnnotationTok{name}\CommentTok{        }\NormalTok{signer name }\CommentTok{(}\NormalTok{optional}\CommentTok{)}
     \CommentTok{*} \AnnotationTok{@param}\CommentTok{ }\AnnotationTok{reason}\CommentTok{      }\NormalTok{signing reason }\CommentTok{(}\NormalTok{optional}\CommentTok{)}
     \CommentTok{*} \AnnotationTok{@param}\CommentTok{ }\AnnotationTok{location}\CommentTok{    }\NormalTok{signing location }\CommentTok{(}\NormalTok{optional}\CommentTok{)}
     \CommentTok{*} \AnnotationTok{@return}\NormalTok{ the signed PDF bytes}
     \CommentTok{*/}
    \KeywordTok{public} \DataTypeTok{byte}\OperatorTok{[]} \FunctionTok{sign}\OperatorTok{(}\DataTypeTok{byte}\OperatorTok{[]}\NormalTok{ pdfBytes}\OperatorTok{,} \BuiltInTok{String}\NormalTok{ subFilter}\OperatorTok{,}
            \BuiltInTok{String}\NormalTok{ name}\OperatorTok{,} \BuiltInTok{String}\NormalTok{ reason}\OperatorTok{,} \BuiltInTok{String}\NormalTok{ location}\OperatorTok{)}
            \KeywordTok{throws} \BuiltInTok{Exception} \OperatorTok{\{}

        \ControlFlowTok{try} \OperatorTok{(}\NormalTok{PDDocument doc }\OperatorTok{=}\NormalTok{ Loader}\OperatorTok{.}\FunctionTok{loadPDF}\OperatorTok{(}\NormalTok{pdfBytes}\OperatorTok{))} \OperatorTok{\{}
\NormalTok{            PDSignature pdSig }\OperatorTok{=} \KeywordTok{new} \FunctionTok{PDSignature}\OperatorTok{();}
\NormalTok{            pdSig}\OperatorTok{.}\FunctionTok{setFilter}\OperatorTok{(}\NormalTok{PDSignature}\OperatorTok{.}\FunctionTok{FILTER\_ADOBE\_PPKLITE}\OperatorTok{);}
\NormalTok{            pdSig}\OperatorTok{.}\FunctionTok{setSubFilter}\OperatorTok{(}\NormalTok{COSName}\OperatorTok{.}\FunctionTok{getPDFName}\OperatorTok{(}\NormalTok{subFilter}\OperatorTok{));}
            \ControlFlowTok{if} \OperatorTok{(}\NormalTok{name }\OperatorTok{!=} \KeywordTok{null}\OperatorTok{)}\NormalTok{ pdSig}\OperatorTok{.}\FunctionTok{setName}\OperatorTok{(}\NormalTok{name}\OperatorTok{);}
            \ControlFlowTok{if} \OperatorTok{(}\NormalTok{reason }\OperatorTok{!=} \KeywordTok{null}\OperatorTok{)}\NormalTok{ pdSig}\OperatorTok{.}\FunctionTok{setReason}\OperatorTok{(}\NormalTok{reason}\OperatorTok{);}
            \ControlFlowTok{if} \OperatorTok{(}\NormalTok{location }\OperatorTok{!=} \KeywordTok{null}\OperatorTok{)}\NormalTok{ pdSig}\OperatorTok{.}\FunctionTok{setLocation}\OperatorTok{(}\NormalTok{location}\OperatorTok{);}
\NormalTok{            pdSig}\OperatorTok{.}\FunctionTok{setSignDate}\OperatorTok{(}\BuiltInTok{Calendar}\OperatorTok{.}\FunctionTok{getInstance}\OperatorTok{());}

\NormalTok{            SignatureInterface sigInterface }\OperatorTok{=}\NormalTok{ content }\OperatorTok{{-}\textgreater{}} \OperatorTok{\{}
                \DataTypeTok{byte}\OperatorTok{[]}\NormalTok{ data }\OperatorTok{=}\NormalTok{ content}\OperatorTok{.}\FunctionTok{readAllBytes}\OperatorTok{();}
                \ControlFlowTok{return} \FunctionTok{buildCmsSignature}\OperatorTok{(}\NormalTok{data}\OperatorTok{,}\NormalTok{ subFilter}\OperatorTok{);}
            \OperatorTok{\};}

\NormalTok{            doc}\OperatorTok{.}\FunctionTok{addSignature}\OperatorTok{(}\NormalTok{pdSig}\OperatorTok{,}\NormalTok{ sigInterface}\OperatorTok{);}

            \BuiltInTok{ByteArrayOutputStream}\NormalTok{ result }\OperatorTok{=} \KeywordTok{new} \BuiltInTok{ByteArrayOutputStream}\OperatorTok{();}
\NormalTok{            doc}\OperatorTok{.}\FunctionTok{saveIncremental}\OperatorTok{(}\NormalTok{result}\OperatorTok{);}

\NormalTok{            log}\OperatorTok{.}\FunctionTok{info}\OperatorTok{(}\StringTok{"PDF signed successfully — \{\} bytes, SubFilter=\{\}"}\OperatorTok{,}
\NormalTok{                    result}\OperatorTok{.}\FunctionTok{size}\OperatorTok{(),}\NormalTok{ subFilter}\OperatorTok{);}
            \ControlFlowTok{return}\NormalTok{ result}\OperatorTok{.}\FunctionTok{toByteArray}\OperatorTok{();}
        \OperatorTok{\}}
    \OperatorTok{\}}

    \CommentTok{/**}
     \CommentTok{*}\NormalTok{ Build the CMS SignedData for the given content bytes}\CommentTok{.}
     \CommentTok{*/}
    \KeywordTok{private} \DataTypeTok{byte}\OperatorTok{[]} \FunctionTok{buildCmsSignature}\OperatorTok{(}\DataTypeTok{byte}\OperatorTok{[]}\NormalTok{ content}\OperatorTok{,} \BuiltInTok{String}\NormalTok{ subFilter}\OperatorTok{)}
            \KeywordTok{throws} \BuiltInTok{Exception} \OperatorTok{\{}

        \CommentTok{// Digest the content}
        \BuiltInTok{MessageDigest}\NormalTok{ digest }\OperatorTok{=} \BuiltInTok{MessageDigest}\OperatorTok{.}\FunctionTok{getInstance}\OperatorTok{(}\StringTok{"SHA{-}256"}\OperatorTok{);}
        \DataTypeTok{byte}\OperatorTok{[]}\NormalTok{ hash }\OperatorTok{=}\NormalTok{ digest}\OperatorTok{.}\FunctionTok{digest}\OperatorTok{(}\NormalTok{content}\OperatorTok{);}

        \CommentTok{// Build the CMS generator}
\NormalTok{        CMSSignedDataGenerator gen }\OperatorTok{=} \KeywordTok{new} \FunctionTok{CMSSignedDataGenerator}\OperatorTok{();}

\NormalTok{        DigestCalculatorProvider dcProv }\OperatorTok{=}
                \KeywordTok{new} \FunctionTok{JcaDigestCalculatorProviderBuilder}\OperatorTok{()}
                        \OperatorTok{.}\FunctionTok{setProvider}\OperatorTok{(}\StringTok{"BC"}\OperatorTok{).}\FunctionTok{build}\OperatorTok{();}

\NormalTok{        ContentSigner signer }\OperatorTok{=} \KeywordTok{new} \FunctionTok{JcaContentSignerBuilder}\OperatorTok{(}
                \StringTok{"SHA256withRSA"}\OperatorTok{).}\FunctionTok{setProvider}\OperatorTok{(}\StringTok{"BC"}\OperatorTok{).}\FunctionTok{build}\OperatorTok{(}\NormalTok{privateKey}\OperatorTok{);}

\NormalTok{        JcaSignerInfoGeneratorBuilder signerBuilder }\OperatorTok{=}
                \KeywordTok{new} \FunctionTok{JcaSignerInfoGeneratorBuilder}\OperatorTok{(}\NormalTok{dcProv}\OperatorTok{);}

        \CommentTok{// For PAdES: add ESS{-}signing{-}certificate{-}v2}
        \ControlFlowTok{if} \OperatorTok{(}\StringTok{"ETSI.CAdES.detached"}\OperatorTok{.}\FunctionTok{equals}\OperatorTok{(}\NormalTok{subFilter}\OperatorTok{))} \OperatorTok{\{}
            \DataTypeTok{byte}\OperatorTok{[]}\NormalTok{ certHash }\OperatorTok{=} \BuiltInTok{MessageDigest}\OperatorTok{.}\FunctionTok{getInstance}\OperatorTok{(}\StringTok{"SHA{-}256"}\OperatorTok{)}
                    \OperatorTok{.}\FunctionTok{digest}\OperatorTok{(}\NormalTok{certificateChain}\OperatorTok{[}\DecValTok{0}\OperatorTok{].}\FunctionTok{getEncoded}\OperatorTok{());}
\NormalTok{            ESSCertIDv2 certId }\OperatorTok{=} \KeywordTok{new} \FunctionTok{ESSCertIDv2}\OperatorTok{(}
                    \KeywordTok{new} \FunctionTok{AlgorithmIdentifier}\OperatorTok{(}
\NormalTok{                            NISTObjectIdentifiers}\OperatorTok{.}\FunctionTok{id\_sha256}\OperatorTok{),}
\NormalTok{                    certHash}\OperatorTok{);}
\NormalTok{            SigningCertificateV2 sigCertV2 }\OperatorTok{=}
                    \KeywordTok{new} \FunctionTok{SigningCertificateV2}\OperatorTok{(}
                            \KeywordTok{new}\NormalTok{ ESSCertIDv2}\OperatorTok{[]\{}\NormalTok{ certId }\OperatorTok{\});}

            \BuiltInTok{Attribute}\NormalTok{ sigCertAttr }\OperatorTok{=} \KeywordTok{new} \BuiltInTok{Attribute}\OperatorTok{(}
\NormalTok{                    PKCSObjectIdentifiers}\OperatorTok{.}\FunctionTok{id\_aa\_signingCertificateV2}\OperatorTok{,}
                    \KeywordTok{new} \FunctionTok{DERSet}\OperatorTok{(}\NormalTok{sigCertV2}\OperatorTok{));}

\NormalTok{            AttributeTable signedAttrs }\OperatorTok{=} \KeywordTok{new} \FunctionTok{AttributeTable}\OperatorTok{(}
                    \KeywordTok{new} \BuiltInTok{Attribute}\OperatorTok{[]\{}\NormalTok{ sigCertAttr }\OperatorTok{\});}
\NormalTok{            signerBuilder}\OperatorTok{.}\FunctionTok{setSignedAttributeGenerator}\OperatorTok{(}
                    \KeywordTok{new} \FunctionTok{DefaultSignedAttributeTableGenerator}\OperatorTok{(}\NormalTok{signedAttrs}\OperatorTok{));}
        \OperatorTok{\}}

\NormalTok{        gen}\OperatorTok{.}\FunctionTok{addSignerInfoGenerator}\OperatorTok{(}
\NormalTok{                signerBuilder}\OperatorTok{.}\FunctionTok{build}\OperatorTok{(}\NormalTok{signer}\OperatorTok{,}\NormalTok{ certificateChain}\OperatorTok{[}\DecValTok{0}\OperatorTok{]));}

        \CommentTok{// Add all certificates in the chain}
        \BuiltInTok{List}\OperatorTok{\textless{}}\BuiltInTok{Certificate}\OperatorTok{\textgreater{}}\NormalTok{ certList }\OperatorTok{=} \KeywordTok{new} \BuiltInTok{ArrayList}\OperatorTok{\textless{}\textgreater{}();}
        \ControlFlowTok{for} \OperatorTok{(}\BuiltInTok{X509Certificate}\NormalTok{ cert }\OperatorTok{:}\NormalTok{ certificateChain}\OperatorTok{)} \OperatorTok{\{}
\NormalTok{            certList}\OperatorTok{.}\FunctionTok{add}\OperatorTok{(}\NormalTok{cert}\OperatorTok{);}
        \OperatorTok{\}}
\NormalTok{        gen}\OperatorTok{.}\FunctionTok{addCertificates}\OperatorTok{(}\KeywordTok{new} \FunctionTok{JcaCertStore}\OperatorTok{(}\NormalTok{certList}\OperatorTok{));}

        \CommentTok{// Generate the CMS SignedData}
\NormalTok{        CMSProcessableByteArray processable }\OperatorTok{=}
                \KeywordTok{new} \FunctionTok{CMSProcessableByteArray}\OperatorTok{(}\NormalTok{content}\OperatorTok{);}
\NormalTok{        CMSSignedData signedData }\OperatorTok{=}\NormalTok{ gen}\OperatorTok{.}\FunctionTok{generate}\OperatorTok{(}\NormalTok{processable}\OperatorTok{,} \KeywordTok{false}\OperatorTok{);}

        \CommentTok{// Add timestamp if TSA URL is configured}
        \ControlFlowTok{if} \OperatorTok{(}\NormalTok{tsaUrl }\OperatorTok{!=} \KeywordTok{null}\OperatorTok{)} \OperatorTok{\{}
\NormalTok{            signedData }\OperatorTok{=} \FunctionTok{addTimestamp}\OperatorTok{(}\NormalTok{signedData}\OperatorTok{);}
        \OperatorTok{\}}

        \DataTypeTok{byte}\OperatorTok{[]}\NormalTok{ encoded }\OperatorTok{=}\NormalTok{ signedData}\OperatorTok{.}\FunctionTok{getEncoded}\OperatorTok{();}
\NormalTok{        log}\OperatorTok{.}\FunctionTok{info}\OperatorTok{(}\StringTok{"CMS signature built: \{\} bytes, timestamp=\{\}"}\OperatorTok{,}
\NormalTok{                encoded}\OperatorTok{.}\FunctionTok{length}\OperatorTok{,}\NormalTok{ tsaUrl }\OperatorTok{!=} \KeywordTok{null}\OperatorTok{);}
        \ControlFlowTok{return}\NormalTok{ encoded}\OperatorTok{;}
    \OperatorTok{\}}

    \CommentTok{/**}
     \CommentTok{*}\NormalTok{ Request a timestamp from the TSA and add it as an unsigned attribute}\CommentTok{.}
     \CommentTok{*/}
    \KeywordTok{private}\NormalTok{ CMSSignedData }\FunctionTok{addTimestamp}\OperatorTok{(}\NormalTok{CMSSignedData signedData}\OperatorTok{)}
            \KeywordTok{throws} \BuiltInTok{Exception} \OperatorTok{\{}

        \CommentTok{// Hash the signature value}
        \DataTypeTok{byte}\OperatorTok{[]}\NormalTok{ signatureBytes }\OperatorTok{=}\NormalTok{ signedData}\OperatorTok{.}\FunctionTok{getSignerInfos}\OperatorTok{()}
                \OperatorTok{.}\FunctionTok{getSigners}\OperatorTok{().}\FunctionTok{iterator}\OperatorTok{().}\FunctionTok{next}\OperatorTok{().}\FunctionTok{getSignature}\OperatorTok{();}
        \BuiltInTok{MessageDigest}\NormalTok{ tsDigest }\OperatorTok{=} \BuiltInTok{MessageDigest}\OperatorTok{.}\FunctionTok{getInstance}\OperatorTok{(}\StringTok{"SHA{-}256"}\OperatorTok{);}
        \DataTypeTok{byte}\OperatorTok{[]}\NormalTok{ tsHash }\OperatorTok{=}\NormalTok{ tsDigest}\OperatorTok{.}\FunctionTok{digest}\OperatorTok{(}\NormalTok{signatureBytes}\OperatorTok{);}

        \CommentTok{// Build and send the timestamp request}
\NormalTok{        TimeStampRequestGenerator reqGen }\OperatorTok{=}
                \KeywordTok{new} \FunctionTok{TimeStampRequestGenerator}\OperatorTok{();}
\NormalTok{        reqGen}\OperatorTok{.}\FunctionTok{setCertReq}\OperatorTok{(}\KeywordTok{true}\OperatorTok{);}
\NormalTok{        TimeStampRequest tsReq }\OperatorTok{=}
\NormalTok{                reqGen}\OperatorTok{.}\FunctionTok{generate}\OperatorTok{(}\NormalTok{TSPAlgorithms}\OperatorTok{.}\FunctionTok{SHA256}\OperatorTok{,}\NormalTok{ tsHash}\OperatorTok{);}

        \BuiltInTok{HttpURLConnection}\NormalTok{ conn }\OperatorTok{=} \OperatorTok{(}\BuiltInTok{HttpURLConnection}\OperatorTok{)}
                \KeywordTok{new} \BuiltInTok{URL}\OperatorTok{(}\NormalTok{tsaUrl}\OperatorTok{).}\FunctionTok{openConnection}\OperatorTok{();}
\NormalTok{        conn}\OperatorTok{.}\FunctionTok{setDoOutput}\OperatorTok{(}\KeywordTok{true}\OperatorTok{);}
\NormalTok{        conn}\OperatorTok{.}\FunctionTok{setRequestMethod}\OperatorTok{(}\StringTok{"POST"}\OperatorTok{);}
\NormalTok{        conn}\OperatorTok{.}\FunctionTok{setRequestProperty}\OperatorTok{(}\StringTok{"Content{-}Type"}\OperatorTok{,}
                \StringTok{"application/timestamp{-}query"}\OperatorTok{);}
\NormalTok{        conn}\OperatorTok{.}\FunctionTok{getOutputStream}\OperatorTok{().}\FunctionTok{write}\OperatorTok{(}\NormalTok{tsReq}\OperatorTok{.}\FunctionTok{getEncoded}\OperatorTok{());}

\NormalTok{        TimeStampResponse tsResp }\OperatorTok{=}
                \KeywordTok{new} \FunctionTok{TimeStampResponse}\OperatorTok{(}\NormalTok{conn}\OperatorTok{.}\FunctionTok{getInputStream}\OperatorTok{());}
\NormalTok{        tsResp}\OperatorTok{.}\FunctionTok{validate}\OperatorTok{(}\NormalTok{tsReq}\OperatorTok{);}
\NormalTok{        TimeStampToken tsToken }\OperatorTok{=}\NormalTok{ tsResp}\OperatorTok{.}\FunctionTok{getTimeStampToken}\OperatorTok{();}

\NormalTok{        log}\OperatorTok{.}\FunctionTok{info}\OperatorTok{(}\StringTok{"Timestamp obtained — TSA time: \{\}"}\OperatorTok{,}
\NormalTok{                tsToken}\OperatorTok{.}\FunctionTok{getTimeStampInfo}\OperatorTok{().}\FunctionTok{getGenTime}\OperatorTok{());}

        \CommentTok{// Embed as unsigned attribute}
        \BuiltInTok{Attribute}\NormalTok{ tsAttr }\OperatorTok{=} \KeywordTok{new} \BuiltInTok{Attribute}\OperatorTok{(}
\NormalTok{                PKCSObjectIdentifiers}\OperatorTok{.}\FunctionTok{id\_aa\_signatureTimeStampToken}\OperatorTok{,}
                \KeywordTok{new} \FunctionTok{DERSet}\OperatorTok{(}\NormalTok{tsToken}\OperatorTok{.}\FunctionTok{toCMSSignedData}\OperatorTok{().}\FunctionTok{toASN1Structure}\OperatorTok{()));}

        \ControlFlowTok{return}\NormalTok{ CMSSignedData}\OperatorTok{.}\FunctionTok{replaceUnsignedAttributes}\OperatorTok{(}
\NormalTok{                signedData}\OperatorTok{,}
                \KeywordTok{new} \FunctionTok{AttributeTable}\OperatorTok{(}\KeywordTok{new} \BuiltInTok{Attribute}\OperatorTok{[]\{}\NormalTok{ tsAttr }\OperatorTok{\}));}
    \OperatorTok{\}}
\OperatorTok{\}}
\end{Highlighting}
\end{Shaded}

\section{VerificationService: Complete Signature
Verification}\label{_verificationservice_complete_signature_verification}

The full verification pipeline from Chapter 31, as a single compilable
class.

\begin{Shaded}
\begin{Highlighting}[]
\KeywordTok{package}\ImportTok{ io}\OperatorTok{.}\ImportTok{pdfalyzer}\OperatorTok{.}\ImportTok{verification}\OperatorTok{;}

\KeywordTok{import} \ImportTok{java}\OperatorTok{.}\ImportTok{security}\OperatorTok{.}\ImportTok{MessageDigest}\OperatorTok{;}
\KeywordTok{import} \ImportTok{java}\OperatorTok{.}\ImportTok{security}\OperatorTok{.}\ImportTok{cert}\OperatorTok{.}\ImportTok{CertPath}\OperatorTok{;}
\KeywordTok{import} \ImportTok{java}\OperatorTok{.}\ImportTok{security}\OperatorTok{.}\ImportTok{cert}\OperatorTok{.}\ImportTok{CertPathBuilder}\OperatorTok{;}
\KeywordTok{import} \ImportTok{java}\OperatorTok{.}\ImportTok{security}\OperatorTok{.}\ImportTok{cert}\OperatorTok{.}\ImportTok{CertStore}\OperatorTok{;}
\KeywordTok{import} \ImportTok{java}\OperatorTok{.}\ImportTok{security}\OperatorTok{.}\ImportTok{cert}\OperatorTok{.}\ImportTok{CertificateFactory}\OperatorTok{;}
\KeywordTok{import} \ImportTok{java}\OperatorTok{.}\ImportTok{security}\OperatorTok{.}\ImportTok{cert}\OperatorTok{.}\ImportTok{CollectionCertStoreParameters}\OperatorTok{;}
\KeywordTok{import} \ImportTok{java}\OperatorTok{.}\ImportTok{security}\OperatorTok{.}\ImportTok{cert}\OperatorTok{.}\ImportTok{PKIXBuilderParameters}\OperatorTok{;}
\KeywordTok{import} \ImportTok{java}\OperatorTok{.}\ImportTok{security}\OperatorTok{.}\ImportTok{cert}\OperatorTok{.}\ImportTok{PKIXCertPathBuilderResult}\OperatorTok{;}
\KeywordTok{import} \ImportTok{java}\OperatorTok{.}\ImportTok{security}\OperatorTok{.}\ImportTok{cert}\OperatorTok{.}\ImportTok{TrustAnchor}\OperatorTok{;}
\KeywordTok{import} \ImportTok{java}\OperatorTok{.}\ImportTok{security}\OperatorTok{.}\ImportTok{cert}\OperatorTok{.}\ImportTok{X509CertSelector}\OperatorTok{;}
\KeywordTok{import} \ImportTok{java}\OperatorTok{.}\ImportTok{security}\OperatorTok{.}\ImportTok{cert}\OperatorTok{.}\ImportTok{X509Certificate}\OperatorTok{;}
\KeywordTok{import} \ImportTok{java}\OperatorTok{.}\ImportTok{time}\OperatorTok{.}\ImportTok{Clock}\OperatorTok{;}
\KeywordTok{import} \ImportTok{java}\OperatorTok{.}\ImportTok{util}\OperatorTok{.}\ImportTok{ArrayList}\OperatorTok{;}
\KeywordTok{import} \ImportTok{java}\OperatorTok{.}\ImportTok{util}\OperatorTok{.}\ImportTok{Collection}\OperatorTok{;}
\KeywordTok{import} \ImportTok{java}\OperatorTok{.}\ImportTok{util}\OperatorTok{.}\ImportTok{Date}\OperatorTok{;}
\KeywordTok{import} \ImportTok{java}\OperatorTok{.}\ImportTok{util}\OperatorTok{.}\ImportTok{List}\OperatorTok{;}
\KeywordTok{import} \ImportTok{java}\OperatorTok{.}\ImportTok{util}\OperatorTok{.}\ImportTok{Set}\OperatorTok{;}

\KeywordTok{import} \ImportTok{lombok}\OperatorTok{.}\ImportTok{extern}\OperatorTok{.}\ImportTok{slf4j}\OperatorTok{.}\ImportTok{Slf4j}\OperatorTok{;}
\KeywordTok{import} \ImportTok{org}\OperatorTok{.}\ImportTok{apache}\OperatorTok{.}\ImportTok{pdfbox}\OperatorTok{.}\ImportTok{pdmodel}\OperatorTok{.}\ImportTok{interactive}\OperatorTok{.}\ImportTok{digitalsignature}\OperatorTok{.}\ImportTok{PDSignature}\OperatorTok{;}
\KeywordTok{import} \ImportTok{org}\OperatorTok{.}\ImportTok{bouncycastle}\OperatorTok{.}\ImportTok{asn1}\OperatorTok{.}\ImportTok{ASN1OctetString}\OperatorTok{;}
\KeywordTok{import} \ImportTok{org}\OperatorTok{.}\ImportTok{bouncycastle}\OperatorTok{.}\ImportTok{asn1}\OperatorTok{.}\ImportTok{cms}\OperatorTok{.}\ImportTok{CMSAttributes}\OperatorTok{;}
\KeywordTok{import} \ImportTok{org}\OperatorTok{.}\ImportTok{bouncycastle}\OperatorTok{.}\ImportTok{asn1}\OperatorTok{.}\ImportTok{cms}\OperatorTok{.}\ImportTok{Attribute}\OperatorTok{;}
\KeywordTok{import} \ImportTok{org}\OperatorTok{.}\ImportTok{bouncycastle}\OperatorTok{.}\ImportTok{asn1}\OperatorTok{.}\ImportTok{cms}\OperatorTok{.}\ImportTok{AttributeTable}\OperatorTok{;}
\KeywordTok{import} \ImportTok{org}\OperatorTok{.}\ImportTok{bouncycastle}\OperatorTok{.}\ImportTok{asn1}\OperatorTok{.}\ImportTok{pkcs}\OperatorTok{.}\ImportTok{PKCSObjectIdentifiers}\OperatorTok{;}
\KeywordTok{import} \ImportTok{org}\OperatorTok{.}\ImportTok{bouncycastle}\OperatorTok{.}\ImportTok{cert}\OperatorTok{.}\ImportTok{X509CertificateHolder}\OperatorTok{;}
\KeywordTok{import} \ImportTok{org}\OperatorTok{.}\ImportTok{bouncycastle}\OperatorTok{.}\ImportTok{cert}\OperatorTok{.}\ImportTok{jcajce}\OperatorTok{.}\ImportTok{JcaX509CertificateConverter}\OperatorTok{;}
\KeywordTok{import} \ImportTok{org}\OperatorTok{.}\ImportTok{bouncycastle}\OperatorTok{.}\ImportTok{cms}\OperatorTok{.}\ImportTok{CMSProcessableByteArray}\OperatorTok{;}
\KeywordTok{import} \ImportTok{org}\OperatorTok{.}\ImportTok{bouncycastle}\OperatorTok{.}\ImportTok{cms}\OperatorTok{.}\ImportTok{CMSSignedData}\OperatorTok{;}
\KeywordTok{import} \ImportTok{org}\OperatorTok{.}\ImportTok{bouncycastle}\OperatorTok{.}\ImportTok{cms}\OperatorTok{.}\ImportTok{SignerInformation}\OperatorTok{;}
\KeywordTok{import} \ImportTok{org}\OperatorTok{.}\ImportTok{bouncycastle}\OperatorTok{.}\ImportTok{cms}\OperatorTok{.}\ImportTok{SignerInformationStore}\OperatorTok{;}
\KeywordTok{import} \ImportTok{org}\OperatorTok{.}\ImportTok{bouncycastle}\OperatorTok{.}\ImportTok{cms}\OperatorTok{.}\ImportTok{jcajce}\OperatorTok{.}\ImportTok{JcaSimpleSignerInfoVerifierBuilder}\OperatorTok{;}
\KeywordTok{import} \ImportTok{org}\OperatorTok{.}\ImportTok{bouncycastle}\OperatorTok{.}\ImportTok{tsp}\OperatorTok{.}\ImportTok{TimeStampToken}\OperatorTok{;}
\KeywordTok{import} \ImportTok{org}\OperatorTok{.}\ImportTok{bouncycastle}\OperatorTok{.}\ImportTok{util}\OperatorTok{.}\ImportTok{Store}\OperatorTok{;}

\CommentTok{/**}
 \CommentTok{*}\NormalTok{ Complete five}\CommentTok{{-}}\NormalTok{step signature verification service}\CommentTok{.}
 \CommentTok{*}\NormalTok{ Accepts an injectable Clock for deterministic testing}\CommentTok{.}
 \CommentTok{*/}
\AttributeTok{@Slf4j}
\KeywordTok{public} \KeywordTok{class}\NormalTok{ VerificationService }\OperatorTok{\{}

    \KeywordTok{private} \DataTypeTok{final}\NormalTok{ Clock clock}\OperatorTok{;}

    \KeywordTok{public} \FunctionTok{VerificationService}\OperatorTok{()} \OperatorTok{\{}
        \KeywordTok{this}\OperatorTok{.}\FunctionTok{clock} \OperatorTok{=}\NormalTok{ Clock}\OperatorTok{.}\FunctionTok{systemUTC}\OperatorTok{();}
    \OperatorTok{\}}

    \KeywordTok{public} \FunctionTok{VerificationService}\OperatorTok{(}\NormalTok{Clock clock}\OperatorTok{)} \OperatorTok{\{}
        \KeywordTok{this}\OperatorTok{.}\FunctionTok{clock} \OperatorTok{=}\NormalTok{ clock}\OperatorTok{;}
    \OperatorTok{\}}

    \CommentTok{/**}
     \CommentTok{*}\NormalTok{ Run the full five}\CommentTok{{-}}\NormalTok{step verification pipeline}\CommentTok{.}
     \CommentTok{*/}
    \KeywordTok{public}\NormalTok{ VerificationReport }\FunctionTok{verify}\OperatorTok{(}\DataTypeTok{byte}\OperatorTok{[]}\NormalTok{ pdfBytes}\OperatorTok{,}
\NormalTok{            PDSignature signature}\OperatorTok{,}
            \BuiltInTok{Set}\OperatorTok{\textless{}}\BuiltInTok{TrustAnchor}\OperatorTok{\textgreater{}}\NormalTok{ trustAnchors}\OperatorTok{)} \OperatorTok{\{}

\NormalTok{        VerificationReport report }\OperatorTok{=} \KeywordTok{new} \FunctionTok{VerificationReport}\OperatorTok{();}

        \ControlFlowTok{try} \OperatorTok{\{}
            \CommentTok{// Step 1: Byte range}
            \DataTypeTok{int}\OperatorTok{[]}\NormalTok{ byteRange }\OperatorTok{=}\NormalTok{ signature}\OperatorTok{.}\FunctionTok{getByteRange}\OperatorTok{();}
\NormalTok{            report}\OperatorTok{.}\FunctionTok{setByteRangeResult}\OperatorTok{(}
                    \FunctionTok{verifyByteRange}\OperatorTok{(}\NormalTok{byteRange}\OperatorTok{,}\NormalTok{ pdfBytes}\OperatorTok{.}\FunctionTok{length}\OperatorTok{));}

            \CommentTok{// Step 2: Hash}
            \DataTypeTok{byte}\OperatorTok{[]}\NormalTok{ signedContent }\OperatorTok{=} \FunctionTok{extractSignedContent}\OperatorTok{(}
\NormalTok{                    pdfBytes}\OperatorTok{,}\NormalTok{ byteRange}\OperatorTok{);}
\NormalTok{            CMSSignedData cmsData }\OperatorTok{=} \KeywordTok{new} \FunctionTok{CMSSignedData}\OperatorTok{(}
                    \KeywordTok{new} \FunctionTok{CMSProcessableByteArray}\OperatorTok{(}\NormalTok{signedContent}\OperatorTok{),}
\NormalTok{                    signature}\OperatorTok{.}\FunctionTok{getContents}\OperatorTok{());}
\NormalTok{            SignerInformation signerInfo }\OperatorTok{=}\NormalTok{ cmsData}\OperatorTok{.}\FunctionTok{getSignerInfos}\OperatorTok{()}
                    \OperatorTok{.}\FunctionTok{getSigners}\OperatorTok{().}\FunctionTok{iterator}\OperatorTok{().}\FunctionTok{next}\OperatorTok{();}

\NormalTok{            report}\OperatorTok{.}\FunctionTok{setHashResult}\OperatorTok{(}\FunctionTok{verifyHash}\OperatorTok{(}
\NormalTok{                    signedContent}\OperatorTok{,}\NormalTok{ signerInfo}\OperatorTok{));}

            \CommentTok{// Step 3: CMS}
\NormalTok{            X509CertificateHolder signerCert }\OperatorTok{=}
                    \FunctionTok{verifyCms}\OperatorTok{(}\NormalTok{cmsData}\OperatorTok{,}\NormalTok{ signerInfo}\OperatorTok{);}
\NormalTok{            report}\OperatorTok{.}\FunctionTok{setCmsResult}\OperatorTok{(}\NormalTok{signerCert }\OperatorTok{!=} \KeywordTok{null}
                    \OperatorTok{?}\NormalTok{ ValidationResult}\OperatorTok{.}\FunctionTok{OK}
                    \OperatorTok{:}\NormalTok{ ValidationResult}\OperatorTok{.}\FunctionTok{CMS\_INVALID}\OperatorTok{);}

            \ControlFlowTok{if} \OperatorTok{(}\NormalTok{signerCert }\OperatorTok{==} \KeywordTok{null}\OperatorTok{)} \ControlFlowTok{return}\NormalTok{ report}\OperatorTok{;}

            \CommentTok{// Check for timestamp}
            \BuiltInTok{Date}\NormalTok{ validationTime }\OperatorTok{=} \BuiltInTok{Date}\OperatorTok{.}\FunctionTok{from}\OperatorTok{(}\NormalTok{clock}\OperatorTok{.}\FunctionTok{instant}\OperatorTok{());}
            \BuiltInTok{Date}\NormalTok{ tsTime }\OperatorTok{=} \FunctionTok{extractTimestampTime}\OperatorTok{(}\NormalTok{signerInfo}\OperatorTok{);}
            \ControlFlowTok{if} \OperatorTok{(}\NormalTok{tsTime }\OperatorTok{!=} \KeywordTok{null}\OperatorTok{)} \OperatorTok{\{}
\NormalTok{                validationTime }\OperatorTok{=}\NormalTok{ tsTime}\OperatorTok{;}
\NormalTok{                report}\OperatorTok{.}\FunctionTok{setTimestampTime}\OperatorTok{(}\NormalTok{tsTime}\OperatorTok{);}
\NormalTok{                report}\OperatorTok{.}\FunctionTok{setTimestampResult}\OperatorTok{(}\NormalTok{ValidationResult}\OperatorTok{.}\FunctionTok{OK}\OperatorTok{);}
            \OperatorTok{\}}

            \CommentTok{// Step 4: Certificate chain}
\NormalTok{            report}\OperatorTok{.}\FunctionTok{setChainResult}\OperatorTok{(}\FunctionTok{validateChain}\OperatorTok{(}
\NormalTok{                    signerCert}\OperatorTok{,}\NormalTok{ cmsData}\OperatorTok{,}\NormalTok{ trustAnchors}\OperatorTok{,}\NormalTok{ validationTime}\OperatorTok{));}

            \CommentTok{// Step 5: Trust}
\NormalTok{            report}\OperatorTok{.}\FunctionTok{setTrustResult}\OperatorTok{(}
\NormalTok{                    report}\OperatorTok{.}\FunctionTok{getChainResult}\OperatorTok{()} \OperatorTok{==}\NormalTok{ ValidationResult}\OperatorTok{.}\FunctionTok{OK}
                            \OperatorTok{?}\NormalTok{ ValidationResult}\OperatorTok{.}\FunctionTok{OK}
                            \OperatorTok{:}\NormalTok{ ValidationResult}\OperatorTok{.}\FunctionTok{NOT\_TRUSTED}\OperatorTok{);}

        \OperatorTok{\}} \ControlFlowTok{catch} \OperatorTok{(}\BuiltInTok{Exception}\NormalTok{ e}\OperatorTok{)} \OperatorTok{\{}
\NormalTok{            log}\OperatorTok{.}\FunctionTok{error}\OperatorTok{(}\StringTok{"Verification failed"}\OperatorTok{,}\NormalTok{ e}\OperatorTok{);}
\NormalTok{            report}\OperatorTok{.}\FunctionTok{setException}\OperatorTok{(}\NormalTok{e}\OperatorTok{);}
        \OperatorTok{\}}

        \ControlFlowTok{return}\NormalTok{ report}\OperatorTok{;}
    \OperatorTok{\}}

    \KeywordTok{private} \DataTypeTok{byte}\OperatorTok{[]} \FunctionTok{extractSignedContent}\OperatorTok{(}\DataTypeTok{byte}\OperatorTok{[]}\NormalTok{ pdfBytes}\OperatorTok{,}
            \DataTypeTok{int}\OperatorTok{[]}\NormalTok{ byteRange}\OperatorTok{)} \OperatorTok{\{}
        \DataTypeTok{byte}\OperatorTok{[]}\NormalTok{ content }\OperatorTok{=} \KeywordTok{new} \DataTypeTok{byte}\OperatorTok{[}\NormalTok{byteRange}\OperatorTok{[}\DecValTok{1}\OperatorTok{]} \OperatorTok{+}\NormalTok{ byteRange}\OperatorTok{[}\DecValTok{3}\OperatorTok{]];}
        \BuiltInTok{System}\OperatorTok{.}\FunctionTok{arraycopy}\OperatorTok{(}\NormalTok{pdfBytes}\OperatorTok{,}\NormalTok{ byteRange}\OperatorTok{[}\DecValTok{0}\OperatorTok{],}\NormalTok{ content}\OperatorTok{,}
                \DecValTok{0}\OperatorTok{,}\NormalTok{ byteRange}\OperatorTok{[}\DecValTok{1}\OperatorTok{]);}
        \BuiltInTok{System}\OperatorTok{.}\FunctionTok{arraycopy}\OperatorTok{(}\NormalTok{pdfBytes}\OperatorTok{,}\NormalTok{ byteRange}\OperatorTok{[}\DecValTok{2}\OperatorTok{],}\NormalTok{ content}\OperatorTok{,}
\NormalTok{                byteRange}\OperatorTok{[}\DecValTok{1}\OperatorTok{],}\NormalTok{ byteRange}\OperatorTok{[}\DecValTok{3}\OperatorTok{]);}
        \ControlFlowTok{return}\NormalTok{ content}\OperatorTok{;}
    \OperatorTok{\}}

    \KeywordTok{private}\NormalTok{ ValidationResult }\FunctionTok{verifyByteRange}\OperatorTok{(}\DataTypeTok{int}\OperatorTok{[]}\NormalTok{ byteRange}\OperatorTok{,}
            \DataTypeTok{long}\NormalTok{ fileSize}\OperatorTok{)} \OperatorTok{\{}
        \ControlFlowTok{if} \OperatorTok{(}\NormalTok{byteRange }\OperatorTok{==} \KeywordTok{null} \OperatorTok{||}\NormalTok{ byteRange}\OperatorTok{.}\FunctionTok{length} \OperatorTok{!=} \DecValTok{4}\OperatorTok{)} \OperatorTok{\{}
            \ControlFlowTok{return}\NormalTok{ ValidationResult}\OperatorTok{.}\FunctionTok{BYTE\_RANGE\_MALFORMED}\OperatorTok{;}
        \OperatorTok{\}}
        \ControlFlowTok{if} \OperatorTok{(}\NormalTok{byteRange}\OperatorTok{[}\DecValTok{0}\OperatorTok{]} \OperatorTok{!=} \DecValTok{0}\OperatorTok{)} \OperatorTok{\{}
            \ControlFlowTok{return}\NormalTok{ ValidationResult}\OperatorTok{.}\FunctionTok{BYTE\_RANGE\_NOT\_AT\_START}\OperatorTok{;}
        \OperatorTok{\}}
        \DataTypeTok{long}\NormalTok{ end }\OperatorTok{=} \OperatorTok{(}\DataTypeTok{long}\OperatorTok{)}\NormalTok{ byteRange}\OperatorTok{[}\DecValTok{2}\OperatorTok{]} \OperatorTok{+}\NormalTok{ byteRange}\OperatorTok{[}\DecValTok{3}\OperatorTok{];}
        \ControlFlowTok{if} \OperatorTok{(}\NormalTok{end }\OperatorTok{!=}\NormalTok{ fileSize}\OperatorTok{)} \OperatorTok{\{}
            \ControlFlowTok{return}\NormalTok{ ValidationResult}\OperatorTok{.}\FunctionTok{BYTE\_RANGE\_TRAILING\_BYTES}\OperatorTok{;}
        \OperatorTok{\}}
        \ControlFlowTok{return}\NormalTok{ ValidationResult}\OperatorTok{.}\FunctionTok{OK}\OperatorTok{;}
    \OperatorTok{\}}

    \KeywordTok{private}\NormalTok{ ValidationResult }\FunctionTok{verifyHash}\OperatorTok{(}\DataTypeTok{byte}\OperatorTok{[]}\NormalTok{ signedContent}\OperatorTok{,}
\NormalTok{            SignerInformation signerInfo}\OperatorTok{)} \KeywordTok{throws} \BuiltInTok{Exception} \OperatorTok{\{}
        \BuiltInTok{String}\NormalTok{ algName }\OperatorTok{=} \FunctionTok{mapOidToAlgorithm}\OperatorTok{(}
\NormalTok{                signerInfo}\OperatorTok{.}\FunctionTok{getDigestAlgOID}\OperatorTok{());}
        \BuiltInTok{MessageDigest}\NormalTok{ md }\OperatorTok{=} \BuiltInTok{MessageDigest}\OperatorTok{.}\FunctionTok{getInstance}\OperatorTok{(}\NormalTok{algName}\OperatorTok{);}
        \DataTypeTok{byte}\OperatorTok{[]}\NormalTok{ computed }\OperatorTok{=}\NormalTok{ md}\OperatorTok{.}\FunctionTok{digest}\OperatorTok{(}\NormalTok{signedContent}\OperatorTok{);}

        \BuiltInTok{Attribute}\NormalTok{ digestAttr }\OperatorTok{=}\NormalTok{ signerInfo}\OperatorTok{.}\FunctionTok{getSignedAttributes}\OperatorTok{()}
                \OperatorTok{.}\FunctionTok{get}\OperatorTok{(}\NormalTok{CMSAttributes}\OperatorTok{.}\FunctionTok{messageDigest}\OperatorTok{);}
        \DataTypeTok{byte}\OperatorTok{[]}\NormalTok{ claimed }\OperatorTok{=}\NormalTok{ ASN1OctetString}\OperatorTok{.}\FunctionTok{getInstance}\OperatorTok{(}
\NormalTok{                digestAttr}\OperatorTok{.}\FunctionTok{getAttrValues}\OperatorTok{().}\FunctionTok{getObjectAt}\OperatorTok{(}\DecValTok{0}\OperatorTok{)).}\FunctionTok{getOctets}\OperatorTok{();}

        \ControlFlowTok{return} \BuiltInTok{MessageDigest}\OperatorTok{.}\FunctionTok{isEqual}\OperatorTok{(}\NormalTok{computed}\OperatorTok{,}\NormalTok{ claimed}\OperatorTok{)}
                \OperatorTok{?}\NormalTok{ ValidationResult}\OperatorTok{.}\FunctionTok{OK}
                \OperatorTok{:}\NormalTok{ ValidationResult}\OperatorTok{.}\FunctionTok{HASH\_MISMATCH}\OperatorTok{;}
    \OperatorTok{\}}

    \AttributeTok{@SuppressWarnings}\OperatorTok{(}\StringTok{"unchecked"}\OperatorTok{)}
    \KeywordTok{private}\NormalTok{ X509CertificateHolder }\FunctionTok{verifyCms}\OperatorTok{(}\NormalTok{CMSSignedData cmsData}\OperatorTok{,}
\NormalTok{            SignerInformation signerInfo}\OperatorTok{)} \KeywordTok{throws} \BuiltInTok{Exception} \OperatorTok{\{}
\NormalTok{        Store}\OperatorTok{\textless{}}\NormalTok{X509CertificateHolder}\OperatorTok{\textgreater{}}\NormalTok{ certStore }\OperatorTok{=}
\NormalTok{                cmsData}\OperatorTok{.}\FunctionTok{getCertificates}\OperatorTok{();}
        \BuiltInTok{Collection}\OperatorTok{\textless{}}\NormalTok{X509CertificateHolder}\OperatorTok{\textgreater{}}\NormalTok{ matches }\OperatorTok{=}
\NormalTok{                certStore}\OperatorTok{.}\FunctionTok{getMatches}\OperatorTok{(}\NormalTok{signerInfo}\OperatorTok{.}\FunctionTok{getSID}\OperatorTok{());}
        \ControlFlowTok{if} \OperatorTok{(}\NormalTok{matches}\OperatorTok{.}\FunctionTok{isEmpty}\OperatorTok{())} \ControlFlowTok{return} \KeywordTok{null}\OperatorTok{;}

\NormalTok{        X509CertificateHolder cert }\OperatorTok{=}\NormalTok{ matches}\OperatorTok{.}\FunctionTok{iterator}\OperatorTok{().}\FunctionTok{next}\OperatorTok{();}
        \DataTypeTok{boolean}\NormalTok{ valid }\OperatorTok{=}\NormalTok{ signerInfo}\OperatorTok{.}\FunctionTok{verify}\OperatorTok{(}
                \KeywordTok{new} \FunctionTok{JcaSimpleSignerInfoVerifierBuilder}\OperatorTok{()}
                        \OperatorTok{.}\FunctionTok{setProvider}\OperatorTok{(}\StringTok{"BC"}\OperatorTok{).}\FunctionTok{build}\OperatorTok{(}\NormalTok{cert}\OperatorTok{));}
        \ControlFlowTok{return}\NormalTok{ valid }\OperatorTok{?}\NormalTok{ cert }\OperatorTok{:} \KeywordTok{null}\OperatorTok{;}
    \OperatorTok{\}}

    \KeywordTok{private} \BuiltInTok{Date} \FunctionTok{extractTimestampTime}\OperatorTok{(}\NormalTok{SignerInformation signerInfo}\OperatorTok{)}
            \KeywordTok{throws} \BuiltInTok{Exception} \OperatorTok{\{}
\NormalTok{        AttributeTable unsigned }\OperatorTok{=}\NormalTok{ signerInfo}\OperatorTok{.}\FunctionTok{getUnsignedAttributes}\OperatorTok{();}
        \ControlFlowTok{if} \OperatorTok{(}\NormalTok{unsigned }\OperatorTok{==} \KeywordTok{null}\OperatorTok{)} \ControlFlowTok{return} \KeywordTok{null}\OperatorTok{;}

        \BuiltInTok{Attribute}\NormalTok{ tsAttr }\OperatorTok{=}\NormalTok{ unsigned}\OperatorTok{.}\FunctionTok{get}\OperatorTok{(}
\NormalTok{                PKCSObjectIdentifiers}\OperatorTok{.}\FunctionTok{id\_aa\_signatureTimeStampToken}\OperatorTok{);}
        \ControlFlowTok{if} \OperatorTok{(}\NormalTok{tsAttr }\OperatorTok{==} \KeywordTok{null}\OperatorTok{)} \ControlFlowTok{return} \KeywordTok{null}\OperatorTok{;}

        \DataTypeTok{byte}\OperatorTok{[]}\NormalTok{ tsBytes }\OperatorTok{=}\NormalTok{ tsAttr}\OperatorTok{.}\FunctionTok{getAttrValues}\OperatorTok{()}
                \OperatorTok{.}\FunctionTok{getObjectAt}\OperatorTok{(}\DecValTok{0}\OperatorTok{).}\FunctionTok{toASN1Primitive}\OperatorTok{().}\FunctionTok{getEncoded}\OperatorTok{();}
\NormalTok{        TimeStampToken tst }\OperatorTok{=}
                \KeywordTok{new} \FunctionTok{TimeStampToken}\OperatorTok{(}\KeywordTok{new} \FunctionTok{CMSSignedData}\OperatorTok{(}\NormalTok{tsBytes}\OperatorTok{));}
        \ControlFlowTok{return}\NormalTok{ tst}\OperatorTok{.}\FunctionTok{getTimeStampInfo}\OperatorTok{().}\FunctionTok{getGenTime}\OperatorTok{();}
    \OperatorTok{\}}

    \KeywordTok{private}\NormalTok{ ValidationResult }\FunctionTok{validateChain}\OperatorTok{(}
\NormalTok{            X509CertificateHolder signerCert}\OperatorTok{,}
\NormalTok{            CMSSignedData cmsData}\OperatorTok{,}
            \BuiltInTok{Set}\OperatorTok{\textless{}}\BuiltInTok{TrustAnchor}\OperatorTok{\textgreater{}}\NormalTok{ trustAnchors}\OperatorTok{,}
            \BuiltInTok{Date}\NormalTok{ validationTime}\OperatorTok{)} \OperatorTok{\{}
        \ControlFlowTok{try} \OperatorTok{\{}
\NormalTok{            JcaX509CertificateConverter converter }\OperatorTok{=}
                    \KeywordTok{new} \FunctionTok{JcaX509CertificateConverter}\OperatorTok{().}\FunctionTok{setProvider}\OperatorTok{(}\StringTok{"BC"}\OperatorTok{);}
            \BuiltInTok{List}\OperatorTok{\textless{}}\BuiltInTok{X509Certificate}\OperatorTok{\textgreater{}}\NormalTok{ allCerts }\OperatorTok{=} \KeywordTok{new} \BuiltInTok{ArrayList}\OperatorTok{\textless{}\textgreater{}();}
            \ControlFlowTok{for} \OperatorTok{(}\NormalTok{X509CertificateHolder h }\OperatorTok{:}
\NormalTok{                    cmsData}\OperatorTok{.}\FunctionTok{getCertificates}\OperatorTok{().}\FunctionTok{getMatches}\OperatorTok{(}\KeywordTok{null}\OperatorTok{))} \OperatorTok{\{}
\NormalTok{                allCerts}\OperatorTok{.}\FunctionTok{add}\OperatorTok{(}\NormalTok{converter}\OperatorTok{.}\FunctionTok{getCertificate}\OperatorTok{(}\NormalTok{h}\OperatorTok{));}
            \OperatorTok{\}}

            \BuiltInTok{X509Certificate}\NormalTok{ target }\OperatorTok{=}
\NormalTok{                    converter}\OperatorTok{.}\FunctionTok{getCertificate}\OperatorTok{(}\NormalTok{signerCert}\OperatorTok{);}
            \BuiltInTok{X509CertSelector}\NormalTok{ selector }\OperatorTok{=} \KeywordTok{new} \BuiltInTok{X509CertSelector}\OperatorTok{();}
\NormalTok{            selector}\OperatorTok{.}\FunctionTok{setCertificate}\OperatorTok{(}\NormalTok{target}\OperatorTok{);}

            \BuiltInTok{PKIXBuilderParameters}\NormalTok{ params }\OperatorTok{=}
                    \KeywordTok{new} \BuiltInTok{PKIXBuilderParameters}\OperatorTok{(}\NormalTok{trustAnchors}\OperatorTok{,}\NormalTok{ selector}\OperatorTok{);}
\NormalTok{            params}\OperatorTok{.}\FunctionTok{addCertStore}\OperatorTok{(}\BuiltInTok{CertStore}\OperatorTok{.}\FunctionTok{getInstance}\OperatorTok{(}\StringTok{"Collection"}\OperatorTok{,}
                    \KeywordTok{new} \BuiltInTok{CollectionCertStoreParameters}\OperatorTok{(}\NormalTok{allCerts}\OperatorTok{)));}
\NormalTok{            params}\OperatorTok{.}\FunctionTok{setRevocationEnabled}\OperatorTok{(}\KeywordTok{false}\OperatorTok{);}
\NormalTok{            params}\OperatorTok{.}\FunctionTok{setDate}\OperatorTok{(}\NormalTok{validationTime}\OperatorTok{);}

            \BuiltInTok{CertPathBuilder}\OperatorTok{.}\FunctionTok{getInstance}\OperatorTok{(}\StringTok{"PKIX"}\OperatorTok{,} \StringTok{"BC"}\OperatorTok{).}\FunctionTok{build}\OperatorTok{(}\NormalTok{params}\OperatorTok{);}
            \ControlFlowTok{return}\NormalTok{ ValidationResult}\OperatorTok{.}\FunctionTok{OK}\OperatorTok{;}
        \OperatorTok{\}} \ControlFlowTok{catch} \OperatorTok{(}\BuiltInTok{Exception}\NormalTok{ e}\OperatorTok{)} \OperatorTok{\{}
\NormalTok{            log}\OperatorTok{.}\FunctionTok{error}\OperatorTok{(}\StringTok{"Chain validation failed: \{\}"}\OperatorTok{,}\NormalTok{ e}\OperatorTok{.}\FunctionTok{getMessage}\OperatorTok{());}
            \ControlFlowTok{return}\NormalTok{ ValidationResult}\OperatorTok{.}\FunctionTok{CHAIN\_INVALID}\OperatorTok{;}
        \OperatorTok{\}}
    \OperatorTok{\}}

    \KeywordTok{private} \DataTypeTok{static} \BuiltInTok{String} \FunctionTok{mapOidToAlgorithm}\OperatorTok{(}\BuiltInTok{String}\NormalTok{ oid}\OperatorTok{)} \OperatorTok{\{}
        \ControlFlowTok{return} \ControlFlowTok{switch} \OperatorTok{(}\NormalTok{oid}\OperatorTok{)} \OperatorTok{\{}
            \ControlFlowTok{case} \StringTok{"1.3.14.3.2.26"}          \OperatorTok{{-}\textgreater{}} \StringTok{"SHA{-}1"}\OperatorTok{;}
            \ControlFlowTok{case} \StringTok{"2.16.840.1.101.3.4.2.1"} \OperatorTok{{-}\textgreater{}} \StringTok{"SHA{-}256"}\OperatorTok{;}
            \ControlFlowTok{case} \StringTok{"2.16.840.1.101.3.4.2.2"} \OperatorTok{{-}\textgreater{}} \StringTok{"SHA{-}384"}\OperatorTok{;}
            \ControlFlowTok{case} \StringTok{"2.16.840.1.101.3.4.2.3"} \OperatorTok{{-}\textgreater{}} \StringTok{"SHA{-}512"}\OperatorTok{;}
            \KeywordTok{default} \OperatorTok{{-}\textgreater{}} \ControlFlowTok{throw} \KeywordTok{new} \BuiltInTok{IllegalArgumentException}\OperatorTok{(}
                    \StringTok{"Unknown digest OID: "} \OperatorTok{+}\NormalTok{ oid}\OperatorTok{);}
        \OperatorTok{\};}
    \OperatorTok{\}}
\OperatorTok{\}}
\end{Highlighting}
\end{Shaded}

\section{Signature Appearance
Builder}\label{_signature_appearance_builder}

A minimal but complete appearance builder for visible signatures.
Generates an appearance stream with signer name, date, and reason inside
the signature widget rectangle.

\begin{Shaded}
\begin{Highlighting}[]
\KeywordTok{package}\ImportTok{ io}\OperatorTok{.}\ImportTok{pdfalyzer}\OperatorTok{.}\ImportTok{signing}\OperatorTok{;}

\KeywordTok{import} \ImportTok{java}\OperatorTok{.}\ImportTok{io}\OperatorTok{.}\ImportTok{IOException}\OperatorTok{;}
\KeywordTok{import} \ImportTok{java}\OperatorTok{.}\ImportTok{text}\OperatorTok{.}\ImportTok{SimpleDateFormat}\OperatorTok{;}
\KeywordTok{import} \ImportTok{java}\OperatorTok{.}\ImportTok{util}\OperatorTok{.}\ImportTok{Calendar}\OperatorTok{;}

\KeywordTok{import} \ImportTok{lombok}\OperatorTok{.}\ImportTok{extern}\OperatorTok{.}\ImportTok{slf4j}\OperatorTok{.}\ImportTok{Slf4j}\OperatorTok{;}
\KeywordTok{import} \ImportTok{org}\OperatorTok{.}\ImportTok{apache}\OperatorTok{.}\ImportTok{pdfbox}\OperatorTok{.}\ImportTok{cos}\OperatorTok{.}\ImportTok{COSName}\OperatorTok{;}
\KeywordTok{import} \ImportTok{org}\OperatorTok{.}\ImportTok{apache}\OperatorTok{.}\ImportTok{pdfbox}\OperatorTok{.}\ImportTok{pdmodel}\OperatorTok{.}\ImportTok{PDDocument}\OperatorTok{;}
\KeywordTok{import} \ImportTok{org}\OperatorTok{.}\ImportTok{apache}\OperatorTok{.}\ImportTok{pdfbox}\OperatorTok{.}\ImportTok{pdmodel}\OperatorTok{.}\ImportTok{PDPage}\OperatorTok{;}
\KeywordTok{import} \ImportTok{org}\OperatorTok{.}\ImportTok{apache}\OperatorTok{.}\ImportTok{pdfbox}\OperatorTok{.}\ImportTok{pdmodel}\OperatorTok{.}\ImportTok{PDResources}\OperatorTok{;}
\KeywordTok{import} \ImportTok{org}\OperatorTok{.}\ImportTok{apache}\OperatorTok{.}\ImportTok{pdfbox}\OperatorTok{.}\ImportTok{pdmodel}\OperatorTok{.}\ImportTok{common}\OperatorTok{.}\ImportTok{PDRectangle}\OperatorTok{;}
\KeywordTok{import} \ImportTok{org}\OperatorTok{.}\ImportTok{apache}\OperatorTok{.}\ImportTok{pdfbox}\OperatorTok{.}\ImportTok{pdmodel}\OperatorTok{.}\ImportTok{font}\OperatorTok{.}\ImportTok{PDType1Font}\OperatorTok{;}
\KeywordTok{import} \ImportTok{org}\OperatorTok{.}\ImportTok{apache}\OperatorTok{.}\ImportTok{pdfbox}\OperatorTok{.}\ImportTok{pdmodel}\OperatorTok{.}\ImportTok{font}\OperatorTok{.}\ImportTok{Standard14Fonts}\OperatorTok{;}
\KeywordTok{import} \ImportTok{org}\OperatorTok{.}\ImportTok{apache}\OperatorTok{.}\ImportTok{pdfbox}\OperatorTok{.}\ImportTok{pdmodel}\OperatorTok{.}\ImportTok{interactive}\OperatorTok{.}\ImportTok{annotation}\OperatorTok{.}\ImportTok{PDAnnotationWidget}\OperatorTok{;}
\KeywordTok{import} \ImportTok{org}\OperatorTok{.}\ImportTok{apache}\OperatorTok{.}\ImportTok{pdfbox}\OperatorTok{.}\ImportTok{pdmodel}\OperatorTok{.}\ImportTok{interactive}\OperatorTok{.}\ImportTok{digitalsignature}\OperatorTok{.}\ImportTok{PDSignature}\OperatorTok{;}
\KeywordTok{import} \ImportTok{org}\OperatorTok{.}\ImportTok{apache}\OperatorTok{.}\ImportTok{pdfbox}\OperatorTok{.}\ImportTok{pdmodel}\OperatorTok{.}\ImportTok{interactive}\OperatorTok{.}\ImportTok{form}\OperatorTok{.}\ImportTok{PDAcroForm}\OperatorTok{;}
\KeywordTok{import} \ImportTok{org}\OperatorTok{.}\ImportTok{apache}\OperatorTok{.}\ImportTok{pdfbox}\OperatorTok{.}\ImportTok{pdmodel}\OperatorTok{.}\ImportTok{interactive}\OperatorTok{.}\ImportTok{form}\OperatorTok{.}\ImportTok{PDSignatureField}\OperatorTok{;}

\CommentTok{/**}
 \CommentTok{*}\NormalTok{ Builds a visible signature appearance with text content}\CommentTok{.}
 \CommentTok{*/}
\AttributeTok{@Slf4j}
\KeywordTok{public} \KeywordTok{class}\NormalTok{ AppearanceBuilder }\OperatorTok{\{}

    \CommentTok{/**}
     \CommentTok{*}\NormalTok{ Create a visible signature field with a text}\CommentTok{{-}}\NormalTok{based appearance}\CommentTok{.}
     \CommentTok{*}
     \CommentTok{*} \AnnotationTok{@param}\CommentTok{ }\AnnotationTok{doc}\CommentTok{       }\NormalTok{the document to add the appearance to}
     \CommentTok{*} \AnnotationTok{@param}\CommentTok{ }\AnnotationTok{sig}\CommentTok{       }\NormalTok{the PDSignature to associate with the field}
     \CommentTok{*} \AnnotationTok{@param}\CommentTok{ }\AnnotationTok{page}\CommentTok{      }\NormalTok{the page number }\CommentTok{(0{-}}\NormalTok{based}\CommentTok{)}
     \CommentTok{*} \AnnotationTok{@param}\CommentTok{ }\AnnotationTok{rect}\CommentTok{      }\NormalTok{the rectangle for the signature widget}
     \CommentTok{*} \AnnotationTok{@param}\CommentTok{ }\AnnotationTok{name}\CommentTok{      }\NormalTok{signer}\CommentTok{\textquotesingle{}}\NormalTok{s name}
     \CommentTok{*} \AnnotationTok{@param}\CommentTok{ }\AnnotationTok{reason}\CommentTok{    }\NormalTok{signing reason}
     \CommentTok{*} \AnnotationTok{@param}\CommentTok{ }\AnnotationTok{signDate}\CommentTok{  }\NormalTok{the signing date}
     \CommentTok{*/}
    \KeywordTok{public} \DataTypeTok{static} \DataTypeTok{void} \FunctionTok{createVisibleAppearance}\OperatorTok{(}
\NormalTok{            PDDocument doc}\OperatorTok{,}\NormalTok{ PDSignature sig}\OperatorTok{,}
            \DataTypeTok{int}\NormalTok{ page}\OperatorTok{,}\NormalTok{ PDRectangle rect}\OperatorTok{,}
            \BuiltInTok{String}\NormalTok{ name}\OperatorTok{,} \BuiltInTok{String}\NormalTok{ reason}\OperatorTok{,} \BuiltInTok{Calendar}\NormalTok{ signDate}\OperatorTok{)}
            \KeywordTok{throws} \BuiltInTok{IOException} \OperatorTok{\{}

\NormalTok{        PDAcroForm acroForm }\OperatorTok{=}\NormalTok{ doc}\OperatorTok{.}\FunctionTok{getDocumentCatalog}\OperatorTok{().}\FunctionTok{getAcroForm}\OperatorTok{();}
        \ControlFlowTok{if} \OperatorTok{(}\NormalTok{acroForm }\OperatorTok{==} \KeywordTok{null}\OperatorTok{)} \OperatorTok{\{}
\NormalTok{            acroForm }\OperatorTok{=} \KeywordTok{new} \FunctionTok{PDAcroForm}\OperatorTok{(}\NormalTok{doc}\OperatorTok{);}
\NormalTok{            doc}\OperatorTok{.}\FunctionTok{getDocumentCatalog}\OperatorTok{().}\FunctionTok{setAcroForm}\OperatorTok{(}\NormalTok{acroForm}\OperatorTok{);}
        \OperatorTok{\}}

\NormalTok{        PDSignatureField sigField }\OperatorTok{=} \KeywordTok{new} \FunctionTok{PDSignatureField}\OperatorTok{(}\NormalTok{acroForm}\OperatorTok{);}
\NormalTok{        sigField}\OperatorTok{.}\FunctionTok{setValue}\OperatorTok{(}\NormalTok{sig}\OperatorTok{);}
\NormalTok{        acroForm}\OperatorTok{.}\FunctionTok{getFields}\OperatorTok{().}\FunctionTok{add}\OperatorTok{(}\NormalTok{sigField}\OperatorTok{);}

\NormalTok{        PDAnnotationWidget widget }\OperatorTok{=}\NormalTok{ sigField}\OperatorTok{.}\FunctionTok{getWidgets}\OperatorTok{().}\FunctionTok{get}\OperatorTok{(}\DecValTok{0}\OperatorTok{);}
\NormalTok{        widget}\OperatorTok{.}\FunctionTok{setRectangle}\OperatorTok{(}\NormalTok{rect}\OperatorTok{);}
\NormalTok{        widget}\OperatorTok{.}\FunctionTok{setPage}\OperatorTok{(}\NormalTok{doc}\OperatorTok{.}\FunctionTok{getPage}\OperatorTok{(}\NormalTok{page}\OperatorTok{));}

        \CommentTok{// Build appearance content}
\NormalTok{        PDType1Font font }\OperatorTok{=} \KeywordTok{new} \FunctionTok{PDType1Font}\OperatorTok{(}
\NormalTok{                Standard14Fonts}\OperatorTok{.}\FunctionTok{FontName}\OperatorTok{.}\FunctionTok{HELVETICA}\OperatorTok{);}
        \DataTypeTok{float}\NormalTok{ fontSize }\OperatorTok{=} \FloatTok{8f}\OperatorTok{;}
        \DataTypeTok{float}\NormalTok{ leading }\OperatorTok{=}\NormalTok{ fontSize }\OperatorTok{*} \FloatTok{1.4f}\OperatorTok{;}
        \DataTypeTok{float}\NormalTok{ margin }\OperatorTok{=} \FloatTok{4f}\OperatorTok{;}

        \BuiltInTok{String}\NormalTok{ dateStr }\OperatorTok{=} \KeywordTok{new} \BuiltInTok{SimpleDateFormat}\OperatorTok{(}\StringTok{"yyyy{-}MM{-}dd HH:mm:ss z"}\OperatorTok{)}
                \OperatorTok{.}\FunctionTok{format}\OperatorTok{(}\NormalTok{signDate}\OperatorTok{.}\FunctionTok{getTime}\OperatorTok{());}

        \BuiltInTok{String}\OperatorTok{[]}\NormalTok{ lines }\OperatorTok{=} \OperatorTok{\{}
                \StringTok{"Digitally signed by: "} \OperatorTok{+}\NormalTok{ name}\OperatorTok{,}
                \StringTok{"Date: "} \OperatorTok{+}\NormalTok{ dateStr}\OperatorTok{,}
                \StringTok{"Reason: "} \OperatorTok{+} \OperatorTok{(}\NormalTok{reason }\OperatorTok{!=} \KeywordTok{null} \OperatorTok{?}\NormalTok{ reason }\OperatorTok{:} \StringTok{"—"}\OperatorTok{)}
        \OperatorTok{\};}

        \CommentTok{// Create the appearance stream using form XObject}
\NormalTok{        PDRectangle formRect }\OperatorTok{=} \KeywordTok{new} \FunctionTok{PDRectangle}\OperatorTok{(}
\NormalTok{                rect}\OperatorTok{.}\FunctionTok{getWidth}\OperatorTok{(),}\NormalTok{ rect}\OperatorTok{.}\FunctionTok{getHeight}\OperatorTok{());}

\NormalTok{        org}\OperatorTok{.}\FunctionTok{apache}\OperatorTok{.}\FunctionTok{pdfbox}\OperatorTok{.}\FunctionTok{pdmodel}\OperatorTok{.}\FunctionTok{graphics}\OperatorTok{.}\FunctionTok{form}\OperatorTok{.}\FunctionTok{PDFormXObject}\NormalTok{ form }\OperatorTok{=}
                \KeywordTok{new}\NormalTok{ org}\OperatorTok{.}\FunctionTok{apache}\OperatorTok{.}\FunctionTok{pdfbox}\OperatorTok{.}\FunctionTok{pdmodel}\OperatorTok{.}\FunctionTok{graphics}\OperatorTok{.}\FunctionTok{form}\OperatorTok{.}\FunctionTok{PDFormXObject}\OperatorTok{(}\NormalTok{doc}\OperatorTok{);}
\NormalTok{        form}\OperatorTok{.}\FunctionTok{setBBox}\OperatorTok{(}\NormalTok{formRect}\OperatorTok{);}
\NormalTok{        form}\OperatorTok{.}\FunctionTok{setResources}\OperatorTok{(}\KeywordTok{new} \FunctionTok{PDResources}\OperatorTok{());}
\NormalTok{        form}\OperatorTok{.}\FunctionTok{getResources}\OperatorTok{().}\FunctionTok{put}\OperatorTok{(}\NormalTok{COSName}\OperatorTok{.}\FunctionTok{getPDFName}\OperatorTok{(}\StringTok{"F1"}\OperatorTok{),}\NormalTok{ font}\OperatorTok{);}

        \CommentTok{// Write content stream operators directly}
        \BuiltInTok{StringBuilder}\NormalTok{ cs }\OperatorTok{=} \KeywordTok{new} \BuiltInTok{StringBuilder}\OperatorTok{();}
\NormalTok{        cs}\OperatorTok{.}\FunctionTok{append}\OperatorTok{(}\StringTok{"BT}\SpecialCharTok{\textbackslash{}n}\StringTok{"}\OperatorTok{);}
\NormalTok{        cs}\OperatorTok{.}\FunctionTok{append}\OperatorTok{(}\StringTok{"/F1 "}\OperatorTok{).}\FunctionTok{append}\OperatorTok{(}\NormalTok{fontSize}\OperatorTok{).}\FunctionTok{append}\OperatorTok{(}\StringTok{" Tf}\SpecialCharTok{\textbackslash{}n}\StringTok{"}\OperatorTok{);}
\NormalTok{        cs}\OperatorTok{.}\FunctionTok{append}\OperatorTok{(}\NormalTok{margin}\OperatorTok{).}\FunctionTok{append}\OperatorTok{(}\StringTok{" "}\OperatorTok{)}
          \OperatorTok{.}\FunctionTok{append}\OperatorTok{(}\NormalTok{rect}\OperatorTok{.}\FunctionTok{getHeight}\OperatorTok{()} \OperatorTok{{-}}\NormalTok{ margin }\OperatorTok{{-}}\NormalTok{ fontSize}\OperatorTok{)}
          \OperatorTok{.}\FunctionTok{append}\OperatorTok{(}\StringTok{" Td}\SpecialCharTok{\textbackslash{}n}\StringTok{"}\OperatorTok{);}
\NormalTok{        cs}\OperatorTok{.}\FunctionTok{append}\OperatorTok{(}\NormalTok{leading}\OperatorTok{).}\FunctionTok{append}\OperatorTok{(}\StringTok{" TL}\SpecialCharTok{\textbackslash{}n}\StringTok{"}\OperatorTok{);}

        \ControlFlowTok{for} \OperatorTok{(}\BuiltInTok{String}\NormalTok{ line }\OperatorTok{:}\NormalTok{ lines}\OperatorTok{)} \OperatorTok{\{}
\NormalTok{            cs}\OperatorTok{.}\FunctionTok{append}\OperatorTok{(}\StringTok{"("}\OperatorTok{).}\FunctionTok{append}\OperatorTok{(}\FunctionTok{escapePdfString}\OperatorTok{(}\NormalTok{line}\OperatorTok{)).}\FunctionTok{append}\OperatorTok{(}\StringTok{") Tj}\SpecialCharTok{\textbackslash{}n}\StringTok{"}\OperatorTok{);}
\NormalTok{            cs}\OperatorTok{.}\FunctionTok{append}\OperatorTok{(}\StringTok{"T*}\SpecialCharTok{\textbackslash{}n}\StringTok{"}\OperatorTok{);}
        \OperatorTok{\}}
\NormalTok{        cs}\OperatorTok{.}\FunctionTok{append}\OperatorTok{(}\StringTok{"ET}\SpecialCharTok{\textbackslash{}n}\StringTok{"}\OperatorTok{);}

\NormalTok{        form}\OperatorTok{.}\FunctionTok{getStream}\OperatorTok{().}\FunctionTok{createOutputStream}\OperatorTok{().}\FunctionTok{write}\OperatorTok{(}
\NormalTok{                cs}\OperatorTok{.}\FunctionTok{toString}\OperatorTok{().}\FunctionTok{getBytes}\OperatorTok{(}\NormalTok{java}\OperatorTok{.}\FunctionTok{nio}\OperatorTok{.}\FunctionTok{charset}\OperatorTok{.}\FunctionTok{StandardCharsets}\OperatorTok{.}\FunctionTok{US\_ASCII}\OperatorTok{));}
\NormalTok{        form}\OperatorTok{.}\FunctionTok{getStream}\OperatorTok{().}\FunctionTok{createOutputStream}\OperatorTok{().}\FunctionTok{close}\OperatorTok{();}

\NormalTok{        org}\OperatorTok{.}\FunctionTok{apache}\OperatorTok{.}\FunctionTok{pdfbox}\OperatorTok{.}\FunctionTok{pdmodel}\OperatorTok{.}\FunctionTok{interactive}\OperatorTok{.}\FunctionTok{annotation}\OperatorTok{.}\FunctionTok{PDAppearanceDictionary}\NormalTok{ appearance }\OperatorTok{=}
                \KeywordTok{new}\NormalTok{ org}\OperatorTok{.}\FunctionTok{apache}\OperatorTok{.}\FunctionTok{pdfbox}\OperatorTok{.}\FunctionTok{pdmodel}\OperatorTok{.}\FunctionTok{interactive}\OperatorTok{.}\FunctionTok{annotation}\OperatorTok{.}\FunctionTok{PDAppearanceDictionary}\OperatorTok{();}
\NormalTok{        appearance}\OperatorTok{.}\FunctionTok{setNormalAppearance}\OperatorTok{(}
                \KeywordTok{new}\NormalTok{ org}\OperatorTok{.}\FunctionTok{apache}\OperatorTok{.}\FunctionTok{pdfbox}\OperatorTok{.}\FunctionTok{pdmodel}\OperatorTok{.}\FunctionTok{interactive}\OperatorTok{.}\FunctionTok{annotation}\OperatorTok{.}\FunctionTok{PDAppearanceEntry}\OperatorTok{(}\NormalTok{form}\OperatorTok{));}
\NormalTok{        widget}\OperatorTok{.}\FunctionTok{setAppearance}\OperatorTok{(}\NormalTok{appearance}\OperatorTok{);}

\NormalTok{        doc}\OperatorTok{.}\FunctionTok{getPage}\OperatorTok{(}\NormalTok{page}\OperatorTok{).}\FunctionTok{getAnnotations}\OperatorTok{().}\FunctionTok{add}\OperatorTok{(}\NormalTok{widget}\OperatorTok{);}

\NormalTok{        log}\OperatorTok{.}\FunctionTok{info}\OperatorTok{(}\StringTok{"Created visible signature appearance: \{\}x\{\} at (\{\},\{\})"}\OperatorTok{,}
\NormalTok{                rect}\OperatorTok{.}\FunctionTok{getWidth}\OperatorTok{(),}\NormalTok{ rect}\OperatorTok{.}\FunctionTok{getHeight}\OperatorTok{(),}
\NormalTok{                rect}\OperatorTok{.}\FunctionTok{getLowerLeftX}\OperatorTok{(),}\NormalTok{ rect}\OperatorTok{.}\FunctionTok{getLowerLeftY}\OperatorTok{());}
    \OperatorTok{\}}

    \KeywordTok{private} \DataTypeTok{static} \BuiltInTok{String} \FunctionTok{escapePdfString}\OperatorTok{(}\BuiltInTok{String}\NormalTok{ input}\OperatorTok{)} \OperatorTok{\{}
        \ControlFlowTok{return}\NormalTok{ input}
                \OperatorTok{.}\FunctionTok{replace}\OperatorTok{(}\StringTok{"}\SpecialCharTok{\textbackslash{}\textbackslash{}}\StringTok{"}\OperatorTok{,} \StringTok{"}\SpecialCharTok{\textbackslash{}\textbackslash{}\textbackslash{}\textbackslash{}}\StringTok{"}\OperatorTok{)}
                \OperatorTok{.}\FunctionTok{replace}\OperatorTok{(}\StringTok{"("}\OperatorTok{,} \StringTok{"}\SpecialCharTok{\textbackslash{}\textbackslash{}}\StringTok{("}\OperatorTok{)}
                \OperatorTok{.}\FunctionTok{replace}\OperatorTok{(}\StringTok{")"}\OperatorTok{,} \StringTok{"}\SpecialCharTok{\textbackslash{}\textbackslash{}}\StringTok{)"}\OperatorTok{);}
    \OperatorTok{\}}
\OperatorTok{\}}
\end{Highlighting}
\end{Shaded}

\section{Test Keystore Generator: Full
Implementation}\label{_test_keystore_generator_full_implementation}

The complete generator from Chapter 32, including all certificate
generation methods.

\begin{Shaded}
\begin{Highlighting}[]
\KeywordTok{package}\ImportTok{ io}\OperatorTok{.}\ImportTok{pdfalyzer}\OperatorTok{.}\ImportTok{test}\OperatorTok{;}

\KeywordTok{import} \ImportTok{java}\OperatorTok{.}\ImportTok{math}\OperatorTok{.}\ImportTok{BigInteger}\OperatorTok{;}
\KeywordTok{import} \ImportTok{java}\OperatorTok{.}\ImportTok{security}\OperatorTok{.}\ImportTok{KeyPair}\OperatorTok{;}
\KeywordTok{import} \ImportTok{java}\OperatorTok{.}\ImportTok{security}\OperatorTok{.}\ImportTok{KeyPairGenerator}\OperatorTok{;}
\KeywordTok{import} \ImportTok{java}\OperatorTok{.}\ImportTok{security}\OperatorTok{.}\ImportTok{KeyStore}\OperatorTok{;}
\KeywordTok{import} \ImportTok{java}\OperatorTok{.}\ImportTok{security}\OperatorTok{.}\ImportTok{cert}\OperatorTok{.}\ImportTok{X509Certificate}\OperatorTok{;}
\KeywordTok{import} \ImportTok{java}\OperatorTok{.}\ImportTok{time}\OperatorTok{.}\ImportTok{Duration}\OperatorTok{;}
\KeywordTok{import} \ImportTok{java}\OperatorTok{.}\ImportTok{time}\OperatorTok{.}\ImportTok{Instant}\OperatorTok{;}
\KeywordTok{import} \ImportTok{java}\OperatorTok{.}\ImportTok{util}\OperatorTok{.}\ImportTok{Arrays}\OperatorTok{;}
\KeywordTok{import} \ImportTok{java}\OperatorTok{.}\ImportTok{util}\OperatorTok{.}\ImportTok{Date}\OperatorTok{;}

\KeywordTok{import} \ImportTok{lombok}\OperatorTok{.}\ImportTok{extern}\OperatorTok{.}\ImportTok{slf4j}\OperatorTok{.}\ImportTok{Slf4j}\OperatorTok{;}
\KeywordTok{import} \ImportTok{org}\OperatorTok{.}\ImportTok{bouncycastle}\OperatorTok{.}\ImportTok{asn1}\OperatorTok{.}\ImportTok{ASN1ObjectIdentifier}\OperatorTok{;}
\KeywordTok{import} \ImportTok{org}\OperatorTok{.}\ImportTok{bouncycastle}\OperatorTok{.}\ImportTok{asn1}\OperatorTok{.}\ImportTok{x500}\OperatorTok{.}\ImportTok{X500Name}\OperatorTok{;}
\KeywordTok{import} \ImportTok{org}\OperatorTok{.}\ImportTok{bouncycastle}\OperatorTok{.}\ImportTok{asn1}\OperatorTok{.}\ImportTok{x509}\OperatorTok{.}\ImportTok{BasicConstraints}\OperatorTok{;}
\KeywordTok{import} \ImportTok{org}\OperatorTok{.}\ImportTok{bouncycastle}\OperatorTok{.}\ImportTok{asn1}\OperatorTok{.}\ImportTok{x509}\OperatorTok{.}\ImportTok{ExtendedKeyUsage}\OperatorTok{;}
\KeywordTok{import} \ImportTok{org}\OperatorTok{.}\ImportTok{bouncycastle}\OperatorTok{.}\ImportTok{asn1}\OperatorTok{.}\ImportTok{x509}\OperatorTok{.}\ImportTok{Extension}\OperatorTok{;}
\KeywordTok{import} \ImportTok{org}\OperatorTok{.}\ImportTok{bouncycastle}\OperatorTok{.}\ImportTok{asn1}\OperatorTok{.}\ImportTok{x509}\OperatorTok{.}\ImportTok{KeyPurposeId}\OperatorTok{;}
\KeywordTok{import} \ImportTok{org}\OperatorTok{.}\ImportTok{bouncycastle}\OperatorTok{.}\ImportTok{asn1}\OperatorTok{.}\ImportTok{x509}\OperatorTok{.}\ImportTok{KeyUsage}\OperatorTok{;}
\KeywordTok{import} \ImportTok{org}\OperatorTok{.}\ImportTok{bouncycastle}\OperatorTok{.}\ImportTok{cert}\OperatorTok{.}\ImportTok{jcajce}\OperatorTok{.}\ImportTok{JcaX509CertificateConverter}\OperatorTok{;}
\KeywordTok{import} \ImportTok{org}\OperatorTok{.}\ImportTok{bouncycastle}\OperatorTok{.}\ImportTok{cert}\OperatorTok{.}\ImportTok{jcajce}\OperatorTok{.}\ImportTok{JcaX509v3CertificateBuilder}\OperatorTok{;}
\KeywordTok{import} \ImportTok{org}\OperatorTok{.}\ImportTok{bouncycastle}\OperatorTok{.}\ImportTok{operator}\OperatorTok{.}\ImportTok{ContentSigner}\OperatorTok{;}
\KeywordTok{import} \ImportTok{org}\OperatorTok{.}\ImportTok{bouncycastle}\OperatorTok{.}\ImportTok{operator}\OperatorTok{.}\ImportTok{jcajce}\OperatorTok{.}\ImportTok{JcaContentSignerBuilder}\OperatorTok{;}

\CommentTok{/**}
 \CommentTok{*}\NormalTok{ Generates a complete PKI hierarchy for testing PDF signatures}\CommentTok{.}
 \CommentTok{*}\NormalTok{ Root CA }\CommentTok{{-}\textgreater{}}\NormalTok{ Intermediate CA }\CommentTok{{-}\textgreater{}}\NormalTok{ Signer}
 \CommentTok{*}\NormalTok{ Root CA }\CommentTok{{-}\textgreater{}}\NormalTok{ TSA}
 \CommentTok{*/}
\AttributeTok{@Slf4j}
\KeywordTok{public} \KeywordTok{class}\NormalTok{ TestKeystoreGenerator }\OperatorTok{\{}

    \KeywordTok{private} \DataTypeTok{static} \DataTypeTok{final} \DataTypeTok{int}\NormalTok{ KEY\_SIZE }\OperatorTok{=} \DecValTok{2048}\OperatorTok{;}
    \KeywordTok{private} \DataTypeTok{static} \DataTypeTok{final} \BuiltInTok{String}\NormalTok{ SIG\_ALG }\OperatorTok{=} \StringTok{"SHA256withRSA"}\OperatorTok{;}
    \KeywordTok{private} \DataTypeTok{static} \DataTypeTok{long}\NormalTok{ serialCounter }\OperatorTok{=} \BuiltInTok{System}\OperatorTok{.}\FunctionTok{currentTimeMillis}\OperatorTok{();}

    \KeywordTok{public} \DataTypeTok{static} \BuiltInTok{KeyStore} \FunctionTok{generateTestKeystore}\OperatorTok{(}\DataTypeTok{char}\OperatorTok{[]}\NormalTok{ password}\OperatorTok{)}
            \KeywordTok{throws} \BuiltInTok{Exception} \OperatorTok{\{}
        \CommentTok{// Root CA}
        \BuiltInTok{KeyPair}\NormalTok{ rootKp }\OperatorTok{=} \FunctionTok{generateKeyPair}\OperatorTok{();}
\NormalTok{        X500Name rootDn }\OperatorTok{=} \KeywordTok{new} \FunctionTok{X500Name}\OperatorTok{(}
                \StringTok{"CN=Test Root CA, O=PDFalyzer Test, C=DE"}\OperatorTok{);}
        \BuiltInTok{X509Certificate}\NormalTok{ rootCert }\OperatorTok{=} \FunctionTok{buildCaCert}\OperatorTok{(}
\NormalTok{                rootKp}\OperatorTok{,}\NormalTok{ rootDn}\OperatorTok{,}\NormalTok{ rootKp}\OperatorTok{,}\NormalTok{ rootDn}\OperatorTok{,} \DecValTok{100}\OperatorTok{,} \KeywordTok{true}\OperatorTok{);}

        \CommentTok{// Intermediate CA}
        \BuiltInTok{KeyPair}\NormalTok{ intKp }\OperatorTok{=} \FunctionTok{generateKeyPair}\OperatorTok{();}
\NormalTok{        X500Name intDn }\OperatorTok{=} \KeywordTok{new} \FunctionTok{X500Name}\OperatorTok{(}
                \StringTok{"CN=Test Intermediate CA, O=PDFalyzer Test, C=DE"}\OperatorTok{);}
        \BuiltInTok{X509Certificate}\NormalTok{ intCert }\OperatorTok{=} \FunctionTok{buildCaCert}\OperatorTok{(}
\NormalTok{                intKp}\OperatorTok{,}\NormalTok{ intDn}\OperatorTok{,}\NormalTok{ rootKp}\OperatorTok{,}\NormalTok{ rootDn}\OperatorTok{,} \DecValTok{50}\OperatorTok{,} \KeywordTok{false}\OperatorTok{);}

        \CommentTok{// Signer}
        \BuiltInTok{KeyPair}\NormalTok{ signerKp }\OperatorTok{=} \FunctionTok{generateKeyPair}\OperatorTok{();}
\NormalTok{        X500Name signerDn }\OperatorTok{=} \KeywordTok{new} \FunctionTok{X500Name}\OperatorTok{(}
                \StringTok{"CN=Test Signer, O=PDFalyzer Test, C=DE"}\OperatorTok{);}
        \BuiltInTok{X509Certificate}\NormalTok{ signerCert }\OperatorTok{=} \FunctionTok{buildSignerCert}\OperatorTok{(}
\NormalTok{                signerKp}\OperatorTok{,}\NormalTok{ signerDn}\OperatorTok{,}\NormalTok{ intKp}\OperatorTok{,}\NormalTok{ intDn}\OperatorTok{,} \DecValTok{10}\OperatorTok{);}

        \CommentTok{// TSA}
        \BuiltInTok{KeyPair}\NormalTok{ tsaKp }\OperatorTok{=} \FunctionTok{generateKeyPair}\OperatorTok{();}
\NormalTok{        X500Name tsaDn }\OperatorTok{=} \KeywordTok{new} \FunctionTok{X500Name}\OperatorTok{(}
                \StringTok{"CN=Test TSA, O=PDFalyzer Test, C=DE"}\OperatorTok{);}
        \BuiltInTok{X509Certificate}\NormalTok{ tsaCert }\OperatorTok{=} \FunctionTok{buildTsaCert}\OperatorTok{(}
\NormalTok{                tsaKp}\OperatorTok{,}\NormalTok{ tsaDn}\OperatorTok{,}\NormalTok{ rootKp}\OperatorTok{,}\NormalTok{ rootDn}\OperatorTok{,} \DecValTok{50}\OperatorTok{);}

        \CommentTok{// Assemble keystore}
        \BuiltInTok{KeyStore}\NormalTok{ ks }\OperatorTok{=} \BuiltInTok{KeyStore}\OperatorTok{.}\FunctionTok{getInstance}\OperatorTok{(}\StringTok{"PKCS12"}\OperatorTok{);}
\NormalTok{        ks}\OperatorTok{.}\FunctionTok{load}\OperatorTok{(}\KeywordTok{null}\OperatorTok{,}\NormalTok{ password}\OperatorTok{);}
\NormalTok{        ks}\OperatorTok{.}\FunctionTok{setKeyEntry}\OperatorTok{(}\StringTok{"signer"}\OperatorTok{,}\NormalTok{ signerKp}\OperatorTok{.}\FunctionTok{getPrivate}\OperatorTok{(),}\NormalTok{ password}\OperatorTok{,}
                \KeywordTok{new} \BuiltInTok{X509Certificate}\OperatorTok{[]\{}\NormalTok{ signerCert}\OperatorTok{,}\NormalTok{ intCert}\OperatorTok{,}\NormalTok{ rootCert }\OperatorTok{\});}
\NormalTok{        ks}\OperatorTok{.}\FunctionTok{setKeyEntry}\OperatorTok{(}\StringTok{"tsa"}\OperatorTok{,}\NormalTok{ tsaKp}\OperatorTok{.}\FunctionTok{getPrivate}\OperatorTok{(),}\NormalTok{ password}\OperatorTok{,}
                \KeywordTok{new} \BuiltInTok{X509Certificate}\OperatorTok{[]\{}\NormalTok{ tsaCert}\OperatorTok{,}\NormalTok{ rootCert }\OperatorTok{\});}
\NormalTok{        ks}\OperatorTok{.}\FunctionTok{setCertificateEntry}\OperatorTok{(}\StringTok{"root"}\OperatorTok{,}\NormalTok{ rootCert}\OperatorTok{);}
\NormalTok{        ks}\OperatorTok{.}\FunctionTok{setCertificateEntry}\OperatorTok{(}\StringTok{"intermediate"}\OperatorTok{,}\NormalTok{ intCert}\OperatorTok{);}

\NormalTok{        log}\OperatorTok{.}\FunctionTok{info}\OperatorTok{(}\StringTok{"Test keystore generated — 4 entries"}\OperatorTok{);}
        \ControlFlowTok{return}\NormalTok{ ks}\OperatorTok{;}
    \OperatorTok{\}}

    \KeywordTok{private} \DataTypeTok{static} \BuiltInTok{KeyPair} \FunctionTok{generateKeyPair}\OperatorTok{()} \KeywordTok{throws} \BuiltInTok{Exception} \OperatorTok{\{}
        \BuiltInTok{KeyPairGenerator}\NormalTok{ gen }\OperatorTok{=} \BuiltInTok{KeyPairGenerator}\OperatorTok{.}\FunctionTok{getInstance}\OperatorTok{(}\StringTok{"RSA"}\OperatorTok{);}
\NormalTok{        gen}\OperatorTok{.}\FunctionTok{initialize}\OperatorTok{(}\NormalTok{KEY\_SIZE}\OperatorTok{);}
        \ControlFlowTok{return}\NormalTok{ gen}\OperatorTok{.}\FunctionTok{generateKeyPair}\OperatorTok{();}
    \OperatorTok{\}}

    \KeywordTok{private} \DataTypeTok{static} \BuiltInTok{X509Certificate} \FunctionTok{buildCaCert}\OperatorTok{(}
            \BuiltInTok{KeyPair}\NormalTok{ subjectKp}\OperatorTok{,}\NormalTok{ X500Name subjectDn}\OperatorTok{,}
            \BuiltInTok{KeyPair}\NormalTok{ issuerKp}\OperatorTok{,}\NormalTok{ X500Name issuerDn}\OperatorTok{,}
            \DataTypeTok{int}\NormalTok{ validityYears}\OperatorTok{,} \DataTypeTok{boolean}\NormalTok{ isRoot}\OperatorTok{)} \KeywordTok{throws} \BuiltInTok{Exception} \OperatorTok{\{}

        \BuiltInTok{Date}\NormalTok{ notBefore }\OperatorTok{=} \KeywordTok{new} \BuiltInTok{Date}\OperatorTok{();}
        \BuiltInTok{Date}\NormalTok{ notAfter }\OperatorTok{=} \BuiltInTok{Date}\OperatorTok{.}\FunctionTok{from}\OperatorTok{(}\NormalTok{Instant}\OperatorTok{.}\FunctionTok{now}\OperatorTok{()}
                \OperatorTok{.}\FunctionTok{plus}\OperatorTok{(}\BuiltInTok{Duration}\OperatorTok{.}\FunctionTok{ofDays}\OperatorTok{(}\DecValTok{365L} \OperatorTok{*}\NormalTok{ validityYears}\OperatorTok{)));}

\NormalTok{        JcaX509v3CertificateBuilder builder }\OperatorTok{=}
                \KeywordTok{new} \FunctionTok{JcaX509v3CertificateBuilder}\OperatorTok{(}
\NormalTok{                        issuerDn}\OperatorTok{,}
                        \BuiltInTok{BigInteger}\OperatorTok{.}\FunctionTok{valueOf}\OperatorTok{(}\NormalTok{serialCounter}\OperatorTok{++),}
\NormalTok{                        notBefore}\OperatorTok{,}\NormalTok{ notAfter}\OperatorTok{,}
\NormalTok{                        subjectDn}\OperatorTok{,}
\NormalTok{                        subjectKp}\OperatorTok{.}\FunctionTok{getPublic}\OperatorTok{());}

\NormalTok{        builder}\OperatorTok{.}\FunctionTok{addExtension}\OperatorTok{(}\NormalTok{Extension}\OperatorTok{.}\FunctionTok{basicConstraints}\OperatorTok{,} \KeywordTok{true}\OperatorTok{,}
                \KeywordTok{new} \FunctionTok{BasicConstraints}\OperatorTok{(}\NormalTok{isRoot }\OperatorTok{?} \BuiltInTok{Integer}\OperatorTok{.}\FunctionTok{MAX\_VALUE} \OperatorTok{:} \DecValTok{0}\OperatorTok{));}
\NormalTok{        builder}\OperatorTok{.}\FunctionTok{addExtension}\OperatorTok{(}\NormalTok{Extension}\OperatorTok{.}\FunctionTok{keyUsage}\OperatorTok{,} \KeywordTok{true}\OperatorTok{,}
                \KeywordTok{new} \FunctionTok{KeyUsage}\OperatorTok{(}
\NormalTok{                        KeyUsage}\OperatorTok{.}\FunctionTok{keyCertSign} \OperatorTok{|}\NormalTok{ KeyUsage}\OperatorTok{.}\FunctionTok{cRLSign}\OperatorTok{));}

\NormalTok{        ContentSigner signer }\OperatorTok{=} \KeywordTok{new} \FunctionTok{JcaContentSignerBuilder}\OperatorTok{(}\NormalTok{SIG\_ALG}\OperatorTok{)}
                \OperatorTok{.}\FunctionTok{setProvider}\OperatorTok{(}\StringTok{"BC"}\OperatorTok{).}\FunctionTok{build}\OperatorTok{(}\NormalTok{issuerKp}\OperatorTok{.}\FunctionTok{getPrivate}\OperatorTok{());}
        \ControlFlowTok{return} \KeywordTok{new} \FunctionTok{JcaX509CertificateConverter}\OperatorTok{().}\FunctionTok{setProvider}\OperatorTok{(}\StringTok{"BC"}\OperatorTok{)}
                \OperatorTok{.}\FunctionTok{getCertificate}\OperatorTok{(}\NormalTok{builder}\OperatorTok{.}\FunctionTok{build}\OperatorTok{(}\NormalTok{signer}\OperatorTok{));}
    \OperatorTok{\}}

    \KeywordTok{private} \DataTypeTok{static} \BuiltInTok{X509Certificate} \FunctionTok{buildSignerCert}\OperatorTok{(}
            \BuiltInTok{KeyPair}\NormalTok{ subjectKp}\OperatorTok{,}\NormalTok{ X500Name subjectDn}\OperatorTok{,}
            \BuiltInTok{KeyPair}\NormalTok{ issuerKp}\OperatorTok{,}\NormalTok{ X500Name issuerDn}\OperatorTok{,}
            \DataTypeTok{int}\NormalTok{ validityYears}\OperatorTok{)} \KeywordTok{throws} \BuiltInTok{Exception} \OperatorTok{\{}

        \BuiltInTok{Date}\NormalTok{ notBefore }\OperatorTok{=} \KeywordTok{new} \BuiltInTok{Date}\OperatorTok{();}
        \BuiltInTok{Date}\NormalTok{ notAfter }\OperatorTok{=} \BuiltInTok{Date}\OperatorTok{.}\FunctionTok{from}\OperatorTok{(}\NormalTok{Instant}\OperatorTok{.}\FunctionTok{now}\OperatorTok{()}
                \OperatorTok{.}\FunctionTok{plus}\OperatorTok{(}\BuiltInTok{Duration}\OperatorTok{.}\FunctionTok{ofDays}\OperatorTok{(}\DecValTok{365L} \OperatorTok{*}\NormalTok{ validityYears}\OperatorTok{)));}

\NormalTok{        JcaX509v3CertificateBuilder builder }\OperatorTok{=}
                \KeywordTok{new} \FunctionTok{JcaX509v3CertificateBuilder}\OperatorTok{(}
\NormalTok{                        issuerDn}\OperatorTok{,}
                        \BuiltInTok{BigInteger}\OperatorTok{.}\FunctionTok{valueOf}\OperatorTok{(}\NormalTok{serialCounter}\OperatorTok{++),}
\NormalTok{                        notBefore}\OperatorTok{,}\NormalTok{ notAfter}\OperatorTok{,}
\NormalTok{                        subjectDn}\OperatorTok{,}
\NormalTok{                        subjectKp}\OperatorTok{.}\FunctionTok{getPublic}\OperatorTok{());}

\NormalTok{        builder}\OperatorTok{.}\FunctionTok{addExtension}\OperatorTok{(}\NormalTok{Extension}\OperatorTok{.}\FunctionTok{basicConstraints}\OperatorTok{,} \KeywordTok{true}\OperatorTok{,}
                \KeywordTok{new} \FunctionTok{BasicConstraints}\OperatorTok{(}\KeywordTok{false}\OperatorTok{));}
\NormalTok{        builder}\OperatorTok{.}\FunctionTok{addExtension}\OperatorTok{(}\NormalTok{Extension}\OperatorTok{.}\FunctionTok{keyUsage}\OperatorTok{,} \KeywordTok{true}\OperatorTok{,}
                \KeywordTok{new} \FunctionTok{KeyUsage}\OperatorTok{(}\NormalTok{KeyUsage}\OperatorTok{.}\FunctionTok{digitalSignature}
                        \OperatorTok{|}\NormalTok{ KeyUsage}\OperatorTok{.}\FunctionTok{nonRepudiation}\OperatorTok{));}

\NormalTok{        ContentSigner signer }\OperatorTok{=} \KeywordTok{new} \FunctionTok{JcaContentSignerBuilder}\OperatorTok{(}\NormalTok{SIG\_ALG}\OperatorTok{)}
                \OperatorTok{.}\FunctionTok{setProvider}\OperatorTok{(}\StringTok{"BC"}\OperatorTok{).}\FunctionTok{build}\OperatorTok{(}\NormalTok{issuerKp}\OperatorTok{.}\FunctionTok{getPrivate}\OperatorTok{());}
        \ControlFlowTok{return} \KeywordTok{new} \FunctionTok{JcaX509CertificateConverter}\OperatorTok{().}\FunctionTok{setProvider}\OperatorTok{(}\StringTok{"BC"}\OperatorTok{)}
                \OperatorTok{.}\FunctionTok{getCertificate}\OperatorTok{(}\NormalTok{builder}\OperatorTok{.}\FunctionTok{build}\OperatorTok{(}\NormalTok{signer}\OperatorTok{));}
    \OperatorTok{\}}

    \KeywordTok{private} \DataTypeTok{static} \BuiltInTok{X509Certificate} \FunctionTok{buildTsaCert}\OperatorTok{(}
            \BuiltInTok{KeyPair}\NormalTok{ subjectKp}\OperatorTok{,}\NormalTok{ X500Name subjectDn}\OperatorTok{,}
            \BuiltInTok{KeyPair}\NormalTok{ issuerKp}\OperatorTok{,}\NormalTok{ X500Name issuerDn}\OperatorTok{,}
            \DataTypeTok{int}\NormalTok{ validityYears}\OperatorTok{)} \KeywordTok{throws} \BuiltInTok{Exception} \OperatorTok{\{}

        \BuiltInTok{Date}\NormalTok{ notBefore }\OperatorTok{=} \KeywordTok{new} \BuiltInTok{Date}\OperatorTok{();}
        \BuiltInTok{Date}\NormalTok{ notAfter }\OperatorTok{=} \BuiltInTok{Date}\OperatorTok{.}\FunctionTok{from}\OperatorTok{(}\NormalTok{Instant}\OperatorTok{.}\FunctionTok{now}\OperatorTok{()}
                \OperatorTok{.}\FunctionTok{plus}\OperatorTok{(}\BuiltInTok{Duration}\OperatorTok{.}\FunctionTok{ofDays}\OperatorTok{(}\DecValTok{365L} \OperatorTok{*}\NormalTok{ validityYears}\OperatorTok{)));}

\NormalTok{        JcaX509v3CertificateBuilder builder }\OperatorTok{=}
                \KeywordTok{new} \FunctionTok{JcaX509v3CertificateBuilder}\OperatorTok{(}
\NormalTok{                        issuerDn}\OperatorTok{,}
                        \BuiltInTok{BigInteger}\OperatorTok{.}\FunctionTok{valueOf}\OperatorTok{(}\NormalTok{serialCounter}\OperatorTok{++),}
\NormalTok{                        notBefore}\OperatorTok{,}\NormalTok{ notAfter}\OperatorTok{,}
\NormalTok{                        subjectDn}\OperatorTok{,}
\NormalTok{                        subjectKp}\OperatorTok{.}\FunctionTok{getPublic}\OperatorTok{());}

\NormalTok{        builder}\OperatorTok{.}\FunctionTok{addExtension}\OperatorTok{(}\NormalTok{Extension}\OperatorTok{.}\FunctionTok{basicConstraints}\OperatorTok{,} \KeywordTok{true}\OperatorTok{,}
                \KeywordTok{new} \FunctionTok{BasicConstraints}\OperatorTok{(}\KeywordTok{false}\OperatorTok{));}
\NormalTok{        builder}\OperatorTok{.}\FunctionTok{addExtension}\OperatorTok{(}\NormalTok{Extension}\OperatorTok{.}\FunctionTok{keyUsage}\OperatorTok{,} \KeywordTok{true}\OperatorTok{,}
                \KeywordTok{new} \FunctionTok{KeyUsage}\OperatorTok{(}\NormalTok{KeyUsage}\OperatorTok{.}\FunctionTok{digitalSignature}\OperatorTok{));}
\NormalTok{        builder}\OperatorTok{.}\FunctionTok{addExtension}\OperatorTok{(}\NormalTok{Extension}\OperatorTok{.}\FunctionTok{extendedKeyUsage}\OperatorTok{,} \KeywordTok{true}\OperatorTok{,}
                \KeywordTok{new} \FunctionTok{ExtendedKeyUsage}\OperatorTok{(}\NormalTok{KeyPurposeId}\OperatorTok{.}\FunctionTok{id\_kp\_timeStamping}\OperatorTok{));}

\NormalTok{        ContentSigner signer }\OperatorTok{=} \KeywordTok{new} \FunctionTok{JcaContentSignerBuilder}\OperatorTok{(}\NormalTok{SIG\_ALG}\OperatorTok{)}
                \OperatorTok{.}\FunctionTok{setProvider}\OperatorTok{(}\StringTok{"BC"}\OperatorTok{).}\FunctionTok{build}\OperatorTok{(}\NormalTok{issuerKp}\OperatorTok{.}\FunctionTok{getPrivate}\OperatorTok{());}
        \ControlFlowTok{return} \KeywordTok{new} \FunctionTok{JcaX509CertificateConverter}\OperatorTok{().}\FunctionTok{setProvider}\OperatorTok{(}\StringTok{"BC"}\OperatorTok{)}
                \OperatorTok{.}\FunctionTok{getCertificate}\OperatorTok{(}\NormalTok{builder}\OperatorTok{.}\FunctionTok{build}\OperatorTok{(}\NormalTok{signer}\OperatorTok{));}
    \OperatorTok{\}}
\OperatorTok{\}}
\end{Highlighting}
\end{Shaded}

\section{ASN.1 Structure Reference}\label{_asn_1_structure_reference}

This appendix is a reference card for the ASN.1 structures you will
encounter when working with PDF digital signatures. Every CMS
\texttt{SignedData} blob inside a \texttt{/Contents} value, every X.509
certificate in the chain, every OCSP response and timestamp token ---
they are all ASN.1 structures encoded in DER. When your BouncyCastle
code throws a \texttt{ParseException} at 2 AM, this is the appendix you
reach for.

We present the OID table first (because OIDs are the keys to
everything), then the ASN.1 grammar for each major structure, with
annotations showing where each field appears in the context of PDF
signatures.

\section{OID Reference Table}\label{_oid_reference_table}

Object Identifiers (OIDs) are the universal labels of the ASN.1 world.
Every algorithm, every attribute, every extension is identified by an
OID. The following table covers the OIDs you will encounter in PDF
signature processing:

\subsection{Digest Algorithms}\label{_digest_algorithms}

\begin{longtable}[]{@{}
  >{\raggedright\arraybackslash}p{(\linewidth - 4\tabcolsep) * \real{0.3750}}
  >{\raggedright\arraybackslash}p{(\linewidth - 4\tabcolsep) * \real{0.3750}}
  >{\raggedright\arraybackslash}p{(\linewidth - 4\tabcolsep) * \real{0.2500}}@{}}
\caption{Digest algorithm OIDs}\tabularnewline
\toprule\noalign{}
\begin{minipage}[b]{\linewidth}\raggedright
OID
\end{minipage} & \begin{minipage}[b]{\linewidth}\raggedright
Algorithm
\end{minipage} & \begin{minipage}[b]{\linewidth}\raggedright
Status
\end{minipage} \\
\midrule\noalign{}
\endfirsthead
\toprule\noalign{}
\begin{minipage}[b]{\linewidth}\raggedright
OID
\end{minipage} & \begin{minipage}[b]{\linewidth}\raggedright
Algorithm
\end{minipage} & \begin{minipage}[b]{\linewidth}\raggedright
Status
\end{minipage} \\
\midrule\noalign{}
\endhead
\bottomrule\noalign{}
\endlastfoot
1.3.14.3.2.26 & SHA-1 & Deprecated --- do not use for new signatures \\
2.16.840.1.101.3.4.2.1 & SHA-256 & Recommended minimum \\
2.16.840.1.101.3.4.2.2 & SHA-384 & Recommended for high security \\
2.16.840.1.101.3.4.2.3 & SHA-512 & Recommended for high security \\
2.16.840.1.101.3.4.2.4 & SHA-224 & Acceptable but uncommon \\
\end{longtable}

\subsection{Signature Algorithms}\label{_signature_algorithms}

\begin{longtable}[]{@{}
  >{\raggedright\arraybackslash}p{(\linewidth - 4\tabcolsep) * \real{0.3750}}
  >{\raggedright\arraybackslash}p{(\linewidth - 4\tabcolsep) * \real{0.3750}}
  >{\raggedright\arraybackslash}p{(\linewidth - 4\tabcolsep) * \real{0.2500}}@{}}
\caption{Signature algorithm OIDs}\tabularnewline
\toprule\noalign{}
\begin{minipage}[b]{\linewidth}\raggedright
OID
\end{minipage} & \begin{minipage}[b]{\linewidth}\raggedright
Algorithm
\end{minipage} & \begin{minipage}[b]{\linewidth}\raggedright
Notes
\end{minipage} \\
\midrule\noalign{}
\endfirsthead
\toprule\noalign{}
\begin{minipage}[b]{\linewidth}\raggedright
OID
\end{minipage} & \begin{minipage}[b]{\linewidth}\raggedright
Algorithm
\end{minipage} & \begin{minipage}[b]{\linewidth}\raggedright
Notes
\end{minipage} \\
\midrule\noalign{}
\endhead
\bottomrule\noalign{}
\endlastfoot
1.2.840.113549.1.1.5 & SHA-1 with RSA (sha1WithRSAEncryption) &
Deprecated \\
1.2.840.113549.1.1.11 & SHA-256 with RSA (sha256WithRSAEncryption) &
Recommended \\
1.2.840.113549.1.1.12 & SHA-384 with RSA (sha384WithRSAEncryption) &
Recommended \\
1.2.840.113549.1.1.13 & SHA-512 with RSA (sha512WithRSAEncryption) &
Recommended \\
1.2.840.113549.1.1.10 & RSASSA-PSS & Modern RSA padding, preferred over
PKCS\#1 v1.5 \\
1.2.840.10045.4.3.2 & SHA-256 with ECDSA (ecdsa-with-SHA256) &
Recommended for ECDSA \\
1.2.840.10045.4.3.3 & SHA-384 with ECDSA (ecdsa-with-SHA384) &
Recommended for ECDSA \\
1.2.840.10045.4.3.4 & SHA-512 with ECDSA (ecdsa-with-SHA512) &
Recommended for ECDSA \\
\end{longtable}

\subsection{CMS Attribute OIDs}\label{_cms_attribute_oids}

\begin{longtable}[]{@{}
  >{\raggedright\arraybackslash}p{(\linewidth - 4\tabcolsep) * \real{0.3333}}
  >{\raggedright\arraybackslash}p{(\linewidth - 4\tabcolsep) * \real{0.4444}}
  >{\raggedright\arraybackslash}p{(\linewidth - 4\tabcolsep) * \real{0.2222}}@{}}
\caption{CMS signed and unsigned attribute OIDs}\tabularnewline
\toprule\noalign{}
\begin{minipage}[b]{\linewidth}\raggedright
OID
\end{minipage} & \begin{minipage}[b]{\linewidth}\raggedright
Attribute Name
\end{minipage} & \begin{minipage}[b]{\linewidth}\raggedright
Context
\end{minipage} \\
\midrule\noalign{}
\endfirsthead
\toprule\noalign{}
\begin{minipage}[b]{\linewidth}\raggedright
OID
\end{minipage} & \begin{minipage}[b]{\linewidth}\raggedright
Attribute Name
\end{minipage} & \begin{minipage}[b]{\linewidth}\raggedright
Context
\end{minipage} \\
\midrule\noalign{}
\endhead
\bottomrule\noalign{}
\endlastfoot
1.2.840.113549.1.9.3 & contentType & Signed attr (required) \\
1.2.840.113549.1.9.4 & messageDigest & Signed attr (required) \\
1.2.840.113549.1.9.5 & signingTime & Signed attr (optional, discouraged
in PAdES) \\
1.2.840.113549.1.9.16.2.47 & id-aa-signingCertificateV2
(ESS-signing-certificate-v2) & Signed attr (required for PAdES) \\
1.2.840.113549.1.9.16.2.14 & id-aa-signatureTimeStampToken & Unsigned
attr (required for PAdES-B-T) \\
1.2.840.113549.1.9.16.2.25 & id-aa-ets-escTimeStamp & Unsigned attr
(used in CAdES-X) \\
1.2.840.113549.1.9.16.2.21 & id-aa-ets-certValues & Unsigned attr (used
in CAdES-C) \\
1.2.840.113549.1.9.16.2.23 & id-aa-ets-revocationValues & Unsigned attr
(used in CAdES-C) \\
\end{longtable}

\subsection{X.509 Extension OIDs}\label{_x_509_extension_oids}

\begin{longtable}[]{@{}
  >{\raggedright\arraybackslash}p{(\linewidth - 4\tabcolsep) * \real{0.3333}}
  >{\raggedright\arraybackslash}p{(\linewidth - 4\tabcolsep) * \real{0.4444}}
  >{\raggedright\arraybackslash}p{(\linewidth - 4\tabcolsep) * \real{0.2222}}@{}}
\caption{X.509 v3 extension OIDs relevant to PDF
signatures}\tabularnewline
\toprule\noalign{}
\begin{minipage}[b]{\linewidth}\raggedright
OID
\end{minipage} & \begin{minipage}[b]{\linewidth}\raggedright
Extension Name
\end{minipage} & \begin{minipage}[b]{\linewidth}\raggedright
Critical
\end{minipage} \\
\midrule\noalign{}
\endfirsthead
\toprule\noalign{}
\begin{minipage}[b]{\linewidth}\raggedright
OID
\end{minipage} & \begin{minipage}[b]{\linewidth}\raggedright
Extension Name
\end{minipage} & \begin{minipage}[b]{\linewidth}\raggedright
Critical
\end{minipage} \\
\midrule\noalign{}
\endhead
\bottomrule\noalign{}
\endlastfoot
2.5.29.15 & keyUsage & Yes \\
2.5.29.37 & extendedKeyUsage & Varies \\
2.5.29.19 & basicConstraints & Yes (for CAs) \\
2.5.29.35 & authorityKeyIdentifier & No \\
2.5.29.14 & subjectKeyIdentifier & No \\
1.3.6.1.5.5.7.1.1 & authorityInfoAccess (AIA) & No \\
2.5.29.31 & cRLDistributionPoints & No \\
1.3.6.1.5.5.7.3.8 & id-kp-timeStamping (EKU value) & N/A (EKU value, not
standalone extension) \\
\end{longtable}

\section{CMS SignedData Structure (RFC 5652 --- Cryptographic Message
Syntax
(CMS))}\label{_cms_signeddata_structure_rfc_5652_cryptographic_message_syntax_cms}

The CMS \texttt{SignedData} is the core structure inside every PDF
\texttt{/Contents} value. This is the complete ASN.1 grammar from RFC
5652 --- Cryptographic Message Syntax (CMS) Section 5.1, annotated with
PDF-specific notes:

\begin{Shaded}
\begin{Highlighting}[]
\CommentTok{{-}{-} RFC 5652 Section 5.1}
\NormalTok{ContentInfo ::= }\DataTypeTok{SEQUENCE}\NormalTok{ \{}
\NormalTok{    contentType     ContentType,        }\CommentTok{{-}{-} OID 1.2.840.113549.1.7.2 (signedData)}
\NormalTok{    content     [}\DecValTok{0}\NormalTok{] EXPLICIT ANY        }\CommentTok{{-}{-} The SignedData below}
\NormalTok{\}}

\NormalTok{SignedData ::= }\DataTypeTok{SEQUENCE}\NormalTok{ \{}
\NormalTok{    version             CMSVersion,         }\CommentTok{{-}{-} 1 for most PDF signatures}
\NormalTok{    digestAlgorithms    DigestAlgorithmIdentifiers,}
                                            \CommentTok{{-}{-} SET OF: usually \{sha{-}256\}}
\NormalTok{    encapContentInfo    EncapsulatedContentInfo,}
                                            \CommentTok{{-}{-} EMPTY for detached signatures (PDF)}
\NormalTok{    certificates    [}\DecValTok{0}\NormalTok{] IMPLICIT CertificateSet }\KeywordTok{OPTIONAL}\NormalTok{,}
                                            \CommentTok{{-}{-} Signer cert + intermediates}
\NormalTok{    crls            [}\DecValTok{1}\NormalTok{] IMPLICIT RevocationInfoChoices }\KeywordTok{OPTIONAL}\NormalTok{,}
                                            \CommentTok{{-}{-} Rarely used in PDF CMS}
\NormalTok{    signerInfos         SignerInfos         }\CommentTok{{-}{-} SET OF: usually exactly one}
\NormalTok{\}}

\CommentTok{{-}{-} In PDF signatures, encapContentInfo.eContent is ABSENT because the}
\CommentTok{{-}{-} signed data (the byte range) is external to the CMS structure.}
\NormalTok{EncapsulatedContentInfo ::= }\DataTypeTok{SEQUENCE}\NormalTok{ \{}
\NormalTok{    eContentType    ContentType,        }\CommentTok{{-}{-} 1.2.840.113549.1.7.1 (id{-}data)}
\NormalTok{    eContent    [}\DecValTok{0}\NormalTok{] EXPLICIT }\DataTypeTok{OCTET} \DataTypeTok{STRING} \KeywordTok{OPTIONAL}
                                        \CommentTok{{-}{-} ABSENT for detached signatures}
\NormalTok{\}}

\NormalTok{SignerInfo ::= }\DataTypeTok{SEQUENCE}\NormalTok{ \{}
\NormalTok{    version                 CMSVersion,     }\CommentTok{{-}{-} 1 (issuerAndSerialNumber)}
\NormalTok{    sid                     SignerIdentifier,}
\NormalTok{    digestAlgorithm         DigestAlgorithmIdentifier,}
\NormalTok{    signedAttrs         [}\DecValTok{0}\NormalTok{] IMPLICIT SignedAttributes }\KeywordTok{OPTIONAL}\NormalTok{,}
                                            \CommentTok{{-}{-} MUST be present for PDF sigs}
\NormalTok{    signatureAlgorithm      SignatureAlgorithmIdentifier,}
\NormalTok{    signature               SignatureValue,}
\NormalTok{    unsignedAttrs       [}\DecValTok{1}\NormalTok{] IMPLICIT UnsignedAttributes }\KeywordTok{OPTIONAL}
                                            \CommentTok{{-}{-} Timestamp lives here}
\NormalTok{\}}

\NormalTok{SignerIdentifier ::= }\DataTypeTok{CHOICE}\NormalTok{ \{}
\NormalTok{    issuerAndSerialNumber   IssuerAndSerialNumber,}
                                            \CommentTok{{-}{-} Most common in PDF}
\NormalTok{    subjectKeyIdentifier [}\DecValTok{0}\NormalTok{] SubjectKeyIdentifier}
\NormalTok{\}}

\NormalTok{IssuerAndSerialNumber ::= }\DataTypeTok{SEQUENCE}\NormalTok{ \{}
\NormalTok{    issuer          Name,}
\NormalTok{    serialNumber    CertificateSerialNumber}
\NormalTok{\}}

\CommentTok{{-}{-} Signed attributes (DER{-}encoded, then signed)}
\NormalTok{SignedAttributes ::= }\DataTypeTok{SET} \KeywordTok{SIZE}\NormalTok{ (}\DecValTok{1}\NormalTok{..MAX) }\DataTypeTok{OF}\NormalTok{ Attribute}

\CommentTok{{-}{-} The Attribute structure used for both signed and unsigned attrs}
\NormalTok{Attribute ::= }\DataTypeTok{SEQUENCE}\NormalTok{ \{}
\NormalTok{    attrType    OBJECT IDENTIFIER,      }\CommentTok{{-}{-} The attribute OID}
\NormalTok{    attrValues  }\DataTypeTok{SET} \DataTypeTok{OF}\NormalTok{ AttributeValue   }\CommentTok{{-}{-} Usually exactly one value}
\NormalTok{\}}
\end{Highlighting}
\end{Shaded}

For PDF detached signatures, the \texttt{eContent} field inside
\texttt{EncapsulatedContentInfo} is absent. The signed content is the
byte range data from the PDF file, which is passed externally to the CMS
verifier. This is why you must provide the byte range bytes when
constructing a \texttt{CMSSignedData} in BouncyCastle --- without them,
the verifier cannot compute the digest.

\section{X.509 v3 Certificate Structure (RFC 5280 --- Internet X.509 PKI
Certificate and CRL
Profile)}\label{_x_509_v3_certificate_structure_rfc_5280_internet_x_509_pki_certificate_and_crl_profile}

Every certificate in the CMS \texttt{certificates} collection, every CA
in the chain, and every TSA certificate follows this structure:

\begin{Shaded}
\begin{Highlighting}[]
\CommentTok{{-}{-} RFC 5280 Section 4.1}
\NormalTok{Certificate ::= }\DataTypeTok{SEQUENCE}\NormalTok{ \{}
\NormalTok{    tbsCertificate      TBSCertificate,     }\CommentTok{{-}{-} The signed payload}
\NormalTok{    signatureAlgorithm  AlgorithmIdentifier, }\CommentTok{{-}{-} Must match inner sig alg}
\NormalTok{    signatureValue      BIT }\DataTypeTok{STRING}           \CommentTok{{-}{-} Issuer\textquotesingle{}s signature}
\NormalTok{\}}

\NormalTok{TBSCertificate ::= }\DataTypeTok{SEQUENCE}\NormalTok{ \{}
\NormalTok{    version         [}\DecValTok{0}\NormalTok{] EXPLICIT }\DataTypeTok{INTEGER} \KeywordTok{DEFAULT}\NormalTok{ v1,    }\CommentTok{{-}{-} v3 (2)}
\NormalTok{    serialNumber        CertificateSerialNumber,}
\NormalTok{    signature           AlgorithmIdentifier,    }\CommentTok{{-}{-} Inner copy of sig alg}
\NormalTok{    issuer              Name,                   }\CommentTok{{-}{-} DN of the issuing CA}
\NormalTok{    validity            Validity,}
\NormalTok{    subject             Name,                   }\CommentTok{{-}{-} DN of the certificate owner}
\NormalTok{    subjectPublicKeyInfo SubjectPublicKeyInfo,}
    \CommentTok{{-}{-} v3 extensions follow:}
\NormalTok{    extensions      [}\DecValTok{3}\NormalTok{] EXPLICIT Extensions }\KeywordTok{OPTIONAL}
\NormalTok{\}}

\NormalTok{Validity ::= }\DataTypeTok{SEQUENCE}\NormalTok{ \{}
\NormalTok{    notBefore   Time,       }\CommentTok{{-}{-} Certificate valid from this date}
\NormalTok{    notAfter    Time        }\CommentTok{{-}{-} Certificate valid until this date}
\NormalTok{\}}

\CommentTok{{-}{-} Key extensions for PDF signature certificates:}
\CommentTok{{-}{-}}
\CommentTok{{-}{-} keyUsage (2.5.29.15):}
\CommentTok{{-}{-}   digitalSignature (bit 0) — required for signing certificates}
\CommentTok{{-}{-}   nonRepudiation (bit 1) — alternative to digitalSignature}
\CommentTok{{-}{-}   keyCertSign (bit 5) — required for CA certificates}
\CommentTok{{-}{-}}
\CommentTok{{-}{-} basicConstraints (2.5.29.19):}
\CommentTok{{-}{-}   cA: BOOLEAN — TRUE for CA certs, FALSE (or absent) for end{-}entity}
\CommentTok{{-}{-}   pathLenConstraint: INTEGER — max chain depth below this CA}
\CommentTok{{-}{-}}
\CommentTok{{-}{-} extendedKeyUsage (2.5.29.37):}
\CommentTok{{-}{-}   id{-}kp{-}timeStamping (1.3.6.1.5.5.7.3.8) — required for TSA certs}
\CommentTok{{-}{-}}
\CommentTok{{-}{-} authorityInfoAccess (1.3.6.1.5.5.7.1.1):}
\CommentTok{{-}{-}   id{-}ad{-}caIssuers — URL to fetch the issuer\textquotesingle{}s certificate}
\CommentTok{{-}{-}   id{-}ad{-}ocsp — URL of the OCSP responder}
\end{Highlighting}
\end{Shaded}

\section{OCSP Response Structure (RFC 6960 --- Online Certificate Status
Protocol
(OCSP))}\label{_ocsp_response_structure_rfc_6960_online_certificate_status_protocol_ocsp}

OCSP responses are embedded in the DSS (Document Security Store) for
PAdES-B-LT and PAdES-B-LTA signatures. The structure you will find in
\texttt{/DSS\ /OCSPs}:

\begin{Shaded}
\begin{Highlighting}[]
\CommentTok{{-}{-} RFC 6960 Section 4.2}
\NormalTok{OCSPResponse ::= }\DataTypeTok{SEQUENCE}\NormalTok{ \{}
\NormalTok{    responseStatus      OCSPResponseStatus,     }\CommentTok{{-}{-} 0 = successful}
\NormalTok{    responseBytes   [}\DecValTok{0}\NormalTok{] EXPLICIT ResponseBytes }\KeywordTok{OPTIONAL}
\NormalTok{\}}

\NormalTok{ResponseBytes ::= }\DataTypeTok{SEQUENCE}\NormalTok{ \{}
\NormalTok{    responseType    OBJECT IDENTIFIER,  }\CommentTok{{-}{-} 1.3.6.1.5.5.7.48.1.1 (basicResponse)}
\NormalTok{    response        }\DataTypeTok{OCTET} \DataTypeTok{STRING}        \CommentTok{{-}{-} DER{-}encoded BasicOCSPResponse}
\NormalTok{\}}

\NormalTok{BasicOCSPResponse ::= }\DataTypeTok{SEQUENCE}\NormalTok{ \{}
\NormalTok{    tbsResponseData     ResponseData,}
\NormalTok{    signatureAlgorithm  AlgorithmIdentifier,}
\NormalTok{    signature           BIT }\DataTypeTok{STRING}\NormalTok{,}
\NormalTok{    certs           [}\DecValTok{0}\NormalTok{] EXPLICIT }\DataTypeTok{SEQUENCE} \DataTypeTok{OF}\NormalTok{ Certificate }\KeywordTok{OPTIONAL}
\NormalTok{\}}

\NormalTok{ResponseData ::= }\DataTypeTok{SEQUENCE}\NormalTok{ \{}
\NormalTok{    version         [}\DecValTok{0}\NormalTok{] EXPLICIT }\DataTypeTok{INTEGER} \KeywordTok{DEFAULT}\NormalTok{ v1,}
\NormalTok{    responderID         ResponderID,}
\NormalTok{    producedAt          }\DataTypeTok{GeneralizedTime}\NormalTok{,}
\NormalTok{    responses           }\DataTypeTok{SEQUENCE} \DataTypeTok{OF}\NormalTok{ SingleResponse,}
\NormalTok{    responseExtensions  [}\DecValTok{1}\NormalTok{] EXPLICIT Extensions }\KeywordTok{OPTIONAL}
\NormalTok{\}}

\NormalTok{SingleResponse ::= }\DataTypeTok{SEQUENCE}\NormalTok{ \{}
\NormalTok{    certID              CertID,}
\NormalTok{    certStatus          CertStatus,         }\CommentTok{{-}{-} good | revoked | unknown}
\NormalTok{    thisUpdate          }\DataTypeTok{GeneralizedTime}\NormalTok{,}
\NormalTok{    nextUpdate      [}\DecValTok{0}\NormalTok{] EXPLICIT }\DataTypeTok{GeneralizedTime} \KeywordTok{OPTIONAL}\NormalTok{,}
\NormalTok{    singleExtensions [}\DecValTok{1}\NormalTok{] EXPLICIT Extensions }\KeywordTok{OPTIONAL}
\NormalTok{\}}

\NormalTok{CertStatus ::= }\DataTypeTok{CHOICE}\NormalTok{ \{}
\NormalTok{    good        [}\DecValTok{0}\NormalTok{] IMPLICIT }\DataTypeTok{NULL}\NormalTok{,          }\CommentTok{{-}{-} Certificate is not revoked}
\NormalTok{    revoked     [}\DecValTok{1}\NormalTok{] IMPLICIT RevokedInfo,   }\CommentTok{{-}{-} Certificate has been revoked}
\NormalTok{    unknown     [}\DecValTok{2}\NormalTok{] IMPLICIT UnknownInfo    }\CommentTok{{-}{-} Responder has no information}
\NormalTok{\}}
\end{Highlighting}
\end{Shaded}

\section{Timestamp Token Structure (RFC 3161 --- Internet X.509 PKI
Time-Stamp Protocol
(TSP))}\label{_timestamp_token_structure_rfc_3161_internet_x_509_pki_time_stamp_protocol_tsp}

The timestamp token (TST) is the structure inside the
\texttt{id-aa-signatureTimeStampToken} unsigned attribute. It is itself
a CMS \texttt{SignedData} containing a \texttt{TSTInfo} as the
encapsulated content:

\begin{Shaded}
\begin{Highlighting}[]
\CommentTok{{-}{-} RFC 3161 Section 2.4.2}
\CommentTok{{-}{-} The TimeStampToken is a CMS SignedData where the}
\CommentTok{{-}{-} encapContentInfo contains a TSTInfo:}

\CommentTok{{-}{-} id{-}ct{-}TSTInfo OBJECT IDENTIFIER ::= \{ 1.2.840.113549.1.9.16.1.4 \}}

\NormalTok{TSTInfo ::= }\DataTypeTok{SEQUENCE}\NormalTok{ \{}
\NormalTok{    version         }\DataTypeTok{INTEGER}\NormalTok{ \{ v1(}\DecValTok{1}\NormalTok{) \},}
\NormalTok{    policy          TSAPolicyId,            }\CommentTok{{-}{-} OID of the TSA\textquotesingle{}s policy}
\NormalTok{    messageImprint  MessageImprint,         }\CommentTok{{-}{-} Hash of the timestamped data}
\NormalTok{    serialNumber    }\DataTypeTok{INTEGER}\NormalTok{,                }\CommentTok{{-}{-} Unique per TSA}
\NormalTok{    genTime         }\DataTypeTok{GeneralizedTime}\NormalTok{,        }\CommentTok{{-}{-} The timestamp time (UTC)}
\NormalTok{    accuracy        Accuracy }\KeywordTok{OPTIONAL}\NormalTok{,      }\CommentTok{{-}{-} Time precision}
\NormalTok{    ordering        }\DataTypeTok{BOOLEAN} \KeywordTok{DEFAULT} \KeywordTok{FALSE}\NormalTok{,}
\NormalTok{    nonce           }\DataTypeTok{INTEGER} \KeywordTok{OPTIONAL}\NormalTok{,       }\CommentTok{{-}{-} Echoed from request}
\NormalTok{    tsa         [}\DecValTok{0}\NormalTok{] GeneralName }\KeywordTok{OPTIONAL}\NormalTok{,   }\CommentTok{{-}{-} TSA identity}
\NormalTok{    extensions  [}\DecValTok{1}\NormalTok{] IMPLICIT Extensions }\KeywordTok{OPTIONAL}
\NormalTok{\}}

\NormalTok{MessageImprint ::= }\DataTypeTok{SEQUENCE}\NormalTok{ \{}
\NormalTok{    hashAlgorithm   AlgorithmIdentifier,    }\CommentTok{{-}{-} e.g., SHA{-}256}
\NormalTok{    hashedMessage   }\DataTypeTok{OCTET} \DataTypeTok{STRING}            \CommentTok{{-}{-} Hash of the signature value}
\NormalTok{\}}

\NormalTok{Accuracy ::= }\DataTypeTok{SEQUENCE}\NormalTok{ \{}
\NormalTok{    seconds     }\DataTypeTok{INTEGER} \KeywordTok{OPTIONAL}\NormalTok{,}
\NormalTok{    millis  [}\DecValTok{0}\NormalTok{] }\DataTypeTok{INTEGER}\NormalTok{ (}\DecValTok{1}\NormalTok{..}\DecValTok{999}\NormalTok{) }\KeywordTok{OPTIONAL}\NormalTok{,}
\NormalTok{    micros  [}\DecValTok{1}\NormalTok{] }\DataTypeTok{INTEGER}\NormalTok{ (}\DecValTok{1}\NormalTok{..}\DecValTok{999}\NormalTok{) }\KeywordTok{OPTIONAL}
\NormalTok{\}}

\CommentTok{{-}{-} The Timestamp Request (sent to the TSA):}
\NormalTok{TimeStampReq ::= }\DataTypeTok{SEQUENCE}\NormalTok{ \{}
\NormalTok{    version         }\DataTypeTok{INTEGER}\NormalTok{ \{ v1(}\DecValTok{1}\NormalTok{) \},}
\NormalTok{    messageImprint  MessageImprint,     }\CommentTok{{-}{-} Hash to timestamp}
\NormalTok{    reqPolicy       TSAPolicyId }\KeywordTok{OPTIONAL}\NormalTok{,}
\NormalTok{    nonce           }\DataTypeTok{INTEGER} \KeywordTok{OPTIONAL}\NormalTok{,   }\CommentTok{{-}{-} Anti{-}replay}
\NormalTok{    certReq         }\DataTypeTok{BOOLEAN} \KeywordTok{DEFAULT} \KeywordTok{FALSE}\NormalTok{,}
                                        \CommentTok{{-}{-} Request TSA cert in response}
\NormalTok{    extensions  [}\DecValTok{0}\NormalTok{] IMPLICIT Extensions }\KeywordTok{OPTIONAL}
\NormalTok{\}}

\CommentTok{{-}{-} The Timestamp Response (returned by the TSA):}
\NormalTok{TimeStampResp ::= }\DataTypeTok{SEQUENCE}\NormalTok{ \{}
\NormalTok{    status          PKIStatusInfo,      }\CommentTok{{-}{-} 0 = granted}
\NormalTok{    timeStampToken  TimeStampToken }\KeywordTok{OPTIONAL}
                                        \CommentTok{{-}{-} CMS SignedData containing TSTInfo}
\NormalTok{\}}
\end{Highlighting}
\end{Shaded}

The \texttt{messageImprint} in the TSTInfo is the hash of the
\emph{signature value} (the raw bytes of the
\texttt{SignerInfo.signature} field), not the hash of the signed
document. This is a critical distinction. The timestamp proves that a
specific signature value existed at a specific time. The signature
value, in turn, covers the document bytes (via the signed attributes and
byte range). The chain of trust is: timestamp proves signature time,
signature proves document integrity.

\section{Algorithm Identifier
Structure}\label{_algorithm_identifier_structure}

The \texttt{AlgorithmIdentifier} appears throughout all the structures
above. For most algorithms, the parameters field is either absent or
\texttt{NULL}:

\begin{Shaded}
\begin{Highlighting}[]
\NormalTok{AlgorithmIdentifier ::= }\DataTypeTok{SEQUENCE}\NormalTok{ \{}
\NormalTok{    algorithm   OBJECT IDENTIFIER,}
\NormalTok{    parameters  ANY DEFINED BY algorithm }\KeywordTok{OPTIONAL}
\NormalTok{\}}

\CommentTok{{-}{-} For SHA{-}256: algorithm = 2.16.840.1.101.3.4.2.1, parameters = NULL (or absent)}
\CommentTok{{-}{-} For RSA:     algorithm = 1.2.840.113549.1.1.11, parameters = NULL}
\CommentTok{{-}{-} For ECDSA:   algorithm = 1.2.840.10045.4.3.2, parameters = absent}
\CommentTok{{-}{-} For RSA{-}PSS: algorithm = 1.2.840.113549.1.1.10,}
\CommentTok{{-}{-}              parameters = RSASSA{-}PSS{-}params (complex — see RFC 4055)}

\NormalTok{RSASSA{-}PSS{-}params ::= }\DataTypeTok{SEQUENCE}\NormalTok{ \{}
\NormalTok{    hashAlgorithm       [}\DecValTok{0}\NormalTok{] HashAlgorithm }\KeywordTok{DEFAULT}\NormalTok{ sha1,}
\NormalTok{    maskGenAlgorithm    [}\DecValTok{1}\NormalTok{] MaskGenAlgorithm }\KeywordTok{DEFAULT}\NormalTok{ mgf1SHA1,}
\NormalTok{    saltLength          [}\DecValTok{2}\NormalTok{] }\DataTypeTok{INTEGER} \KeywordTok{DEFAULT} \DecValTok{20}\NormalTok{,}
\NormalTok{    trailerField        [}\DecValTok{3}\NormalTok{] TrailerField }\KeywordTok{DEFAULT}\NormalTok{ trailerFieldBC}
\NormalTok{\}}
\end{Highlighting}
\end{Shaded}

\section{Glossary \& Spec References}\label{_glossary_spec_references}

This appendix collects every acronym, term, and specification referenced
throughout the book into two lookup tables. When you encounter an
unfamiliar term in a chapter, this is the place to find a concise
definition. When you need the precise specification for a standard we
reference, the spec table gives you the document number, title, and
where to find it.

\section{Glossary}\label{_glossary}

{\def\LTcaptype{none} % do not increment counter
\begin{longtable}[]{@{}
  >{\raggedright\arraybackslash}p{(\linewidth - 2\tabcolsep) * \real{0.2500}}
  >{\raggedright\arraybackslash}p{(\linewidth - 2\tabcolsep) * \real{0.7500}}@{}}
\toprule\noalign{}
\begin{minipage}[b]{\linewidth}\raggedright
Term
\end{minipage} & \begin{minipage}[b]{\linewidth}\raggedright
Definition
\end{minipage} \\
\midrule\noalign{}
\endhead
\bottomrule\noalign{}
\endlastfoot
AATL & \textbf{Adobe Approved Trust List.} Adobe's curated list of root
CA certificates trusted by Acrobat and Reader for signature validation.
CAs must apply and meet Adobe's requirements to be listed. If a signer's
certificate does not chain to an AATL root, Acrobat shows a warning even
if the signature is cryptographically valid. \\
AcroForm & \textbf{Interactive Form Dictionary.} The PDF structure
(\texttt{/AcroForm} in the document catalog) that contains all form
fields, including signature fields. Defined in ISO 32000-2:2020 Section
12.7. \\
AIA & \textbf{Authority Information Access.} An X.509 v3 certificate
extension (RFC 5280 --- Internet X.509 PKI Certificate and CRL Profile
Section 4.2.2.1) that provides URLs for the OCSP responder and the
issuer's certificate. Used to fetch missing intermediate certificates
during chain building. \\
ASN.1 & \textbf{Abstract Syntax Notation One.} A standard for describing
data structures used in X.509 certificates, CMS, OCSP, and TSP.
Structures are encoded using DER (Distinguished Encoding Rules) for
cryptographic applications. \\
ByteRange & The \texttt{/ByteRange} array in a PDF signature dictionary.
Specifies which bytes of the PDF file are covered by the signature.
Always four integers: \texttt{{[}offset1\ length1\ offset2\ length2{]}}.
See ISO 32000-2:2020 Section 12.8.1. \\
CA & \textbf{Certificate Authority.} An entity that issues X.509
certificates. CAs form a hierarchy: root CAs are self-signed and
explicitly trusted; intermediate CAs are signed by a root or another
intermediate. \\
CAdES & \textbf{CMS Advanced Electronic Signatures.} ETSI standard (EN
319 122) defining enhanced CMS signature profiles for long-term validity
and legal compliance. PAdES builds on CAdES for the CMS portion of PDF
signatures. \\
CMS & \textbf{Cryptographic Message Syntax.} The standard (RFC 5652 ---
Cryptographic Message Syntax (CMS)) for signing and encrypting data. PDF
signatures use CMS \texttt{SignedData} as the cryptographic payload
inside the \texttt{/Contents} value. Successor to PKCS\#7. \\
CRL & \textbf{Certificate Revocation List.} A signed list published by a
CA containing the serial numbers of revoked certificates. Defined in RFC
5280 --- Internet X.509 PKI Certificate and CRL Profile Section 5.
Checked during certificate chain validation. \\
DER & \textbf{Distinguished Encoding Rules.} A deterministic binary
encoding of ASN.1 structures. All cryptographic ASN.1 (certificates,
CMS, OCSP, TSP) uses DER encoding. \\
DSS & \textbf{Document Security Store.} A PDF dictionary (\texttt{/DSS}
in the document catalog) defined in ISO 32000-2:2020 Section 12.8.4.3
that stores certificates, CRLs, and OCSP responses needed for offline
(LTV) validation. \\
ECDSA & \textbf{Elliptic Curve Digital Signature Algorithm.} A signature
algorithm based on elliptic curve cryptography. Uses smaller keys than
RSA for equivalent security (e.g., 256-bit ECDSA approximates 3072-bit
RSA). \\
eIDAS & \textbf{Electronic Identification, Authentication, and Trust
Services.} EU Regulation (No 910/2014) that establishes a legal
framework for electronic signatures, seals, and timestamps across EU
member states. Defines three signature levels: simple, advanced, and
qualified. \\
ESS & \textbf{Extended Security Services.} Defined in RFC 2634 and RFC
5035. The \texttt{ESS-signing-certificate-v2} attribute is required in
PAdES signatures to bind the CMS signature to a specific signer
certificate. \\
ETSI & \textbf{European Telecommunications Standards Institute.} The
standards body that publishes PAdES (EN 319 142), CAdES (EN 319 122),
and the validation standards (TS 119 102) referenced by the eIDAS
regulation. \\
EUTL & \textbf{European Union Trusted List.} The official registry of
Qualified Trust Service Providers (QTSPs) maintained by EU member
states. Analogous to Adobe's AATL but for the EU regulatory
framework. \\
HSM & \textbf{Hardware Security Module.} A physical device that stores
private keys and performs cryptographic operations. Used in production
signing environments to protect key material from extraction. \\
Incremental Update & A mechanism defined in ISO 32000-2:2020 Section
7.5.6 that appends new or modified objects to the end of a PDF file
without rewriting the entire file. Every PDF signature is applied via an
incremental update. \\
LTV & \textbf{Long-Term Validation.} The practice of embedding all data
needed for signature validation (certificates, CRLs, OCSP responses)
within the PDF itself, so the signature can be validated offline,
indefinitely, without contacting external servers. \\
MDP & \textbf{Modification Detection and Prevention.} A PDF mechanism
(ISO 32000-2:2020 Section 12.8.2.2) that specifies what changes are
permitted after a certification signature is applied. Three permission
levels exist (P=1, P=2, P=3). \\
OCSP & \textbf{Online Certificate Status Protocol.} A protocol (RFC 6960
--- Online Certificate Status Protocol (OCSP)) for checking certificate
revocation status in real time. An OCSP responder returns \texttt{good},
\texttt{revoked}, or \texttt{unknown} for a given certificate serial
number. \\
OID & \textbf{Object Identifier.} A globally unique dotted-decimal
number (e.g., \texttt{1.2.840.113549.1.7.2}) that identifies algorithms,
attributes, extensions, and other standardized objects in ASN.1
structures. \\
PAdES & \textbf{PDF Advanced Electronic Signatures.} ETSI standard (EN
319 142) that defines constrained profiles for PDF signatures to meet EU
legal requirements. Four baseline profiles: B-B, B-T, B-LT, B-LTA. \\
PKCS\#7 & \textbf{Public-Key Cryptography Standards \#7.} The
predecessor to CMS. PDF signatures historically used the term "PKCS\#7"
(and the SubFilter \texttt{adbe.pkcs7.detached}), though modern
implementations use CMS. The terms are often used interchangeably. \\
PKCS\#12 & \textbf{Public-Key Cryptography Standards \#12.} A file
format for storing private keys and certificates together, typically
with \texttt{.p12} or \texttt{.pfx} extension. Protected by a
password. \\
PKIX & \textbf{Public-Key Infrastructure (X.509).} The set of standards
and protocols for managing X.509 certificates and building certificate
chains. In Java, \texttt{CertPathBuilder} and \texttt{CertPathValidator}
implement the PKIX algorithm. \\
QTSP & \textbf{Qualified Trust Service Provider.} Under eIDAS, a trust
service provider that has been granted qualified status by an EU member
state's supervisory body. Only QTSPs can issue qualified certificates
and qualified timestamps. \\
RSA & \textbf{Rivest-Shamir-Adleman.} The most widely used public-key
algorithm for digital signatures. Common key sizes for PDF signatures
are 2048 and 4096 bits. \\
SHA-256 & \textbf{Secure Hash Algorithm, 256-bit.} The recommended
minimum hash algorithm for new PDF signatures. Produces a 32-byte
(256-bit) digest. Part of the SHA-2 family. \\
SubFilter & The \texttt{/SubFilter} entry in a PDF signature dictionary
that specifies the CMS encoding format. Values:
\texttt{adbe.pkcs7.detached} (classic), \texttt{ETSI.CAdES.detached}
(PAdES), \texttt{adbe.pkcs7.sha1} (deprecated),
\texttt{adbe.x509.rsa\_sha1} (deprecated). \\
TSA & \textbf{Timestamp Authority.} A trusted server that issues RFC
3161 timestamp tokens. The timestamp token proves that specific data
existed at a specific point in time. \\
TSP & \textbf{Time-Stamp Protocol.} The protocol defined in RFC 3161 ---
Internet X.509 PKI Time-Stamp Protocol (TSP) for requesting and
receiving timestamp tokens from a TSA. \\
TST & \textbf{Time-Stamp Token.} The CMS \texttt{SignedData} structure
returned by a TSA containing a \texttt{TSTInfo} with the timestamp time
and the hash of the timestamped data. \\
VRI & \textbf{Validation Related Information.} A sub-dictionary of the
DSS that groups certificates, CRLs, and OCSP responses for a specific
signature. Keyed by the SHA-1 hash of the signature value. \\
X.509 & The ITU-T standard for public-key certificates. Version 3 (X.509
v3) is used universally in PDF signatures. Defined in RFC 5280 ---
Internet X.509 PKI Certificate and CRL Profile. \\
XAdES & \textbf{XML Advanced Electronic Signatures.} ETSI standard (EN
319 132) for XML-based signatures. Not directly relevant to PDF
signatures, but part of the same eIDAS framework as PAdES and CAdES. \\
\end{longtable}
}

\section{Specification Reference
Table}\label{_specification_reference_table}

\begin{longtable}[]{@{}
  >{\raggedright\arraybackslash}p{(\linewidth - 4\tabcolsep) * \real{0.2222}}
  >{\raggedright\arraybackslash}p{(\linewidth - 4\tabcolsep) * \real{0.4444}}
  >{\raggedright\arraybackslash}p{(\linewidth - 4\tabcolsep) * \real{0.3333}}@{}}
\caption{Standards and specifications referenced in this
book}\tabularnewline
\toprule\noalign{}
\begin{minipage}[b]{\linewidth}\raggedright
Identifier
\end{minipage} & \begin{minipage}[b]{\linewidth}\raggedright
Title
\end{minipage} & \begin{minipage}[b]{\linewidth}\raggedright
Access
\end{minipage} \\
\midrule\noalign{}
\endfirsthead
\toprule\noalign{}
\begin{minipage}[b]{\linewidth}\raggedright
Identifier
\end{minipage} & \begin{minipage}[b]{\linewidth}\raggedright
Title
\end{minipage} & \begin{minipage}[b]{\linewidth}\raggedright
Access
\end{minipage} \\
\midrule\noalign{}
\endhead
\bottomrule\noalign{}
\endlastfoot
ISO 32000-2:2020 & Document management --- Portable document
format --- Part 2: PDF 2.0 &
\url{https://www.iso.org/standard/75839.html} (paid) or Adobe's
published PDF Reference \\
RFC 5652 & Cryptographic Message Syntax (CMS) &
\url{https://www.rfc-editor.org/rfc/rfc5652} \\
RFC 5280 & Internet X.509 Public Key Infrastructure Certificate and CRL
Profile & \url{https://www.rfc-editor.org/rfc/rfc5280} \\
RFC 3161 & Internet X.509 Public Key Infrastructure Time-Stamp Protocol
(TSP) & \url{https://www.rfc-editor.org/rfc/rfc3161} \\
RFC 6960 & X.509 Internet Public Key Infrastructure Online Certificate
Status Protocol - OCSP & \url{https://www.rfc-editor.org/rfc/rfc6960} \\
RFC 2634 & Enhanced Security Services (ESS) for S/MIME &
\url{https://www.rfc-editor.org/rfc/rfc2634} \\
RFC 5035 & Enhanced Security Services (ESS) Update: Adding CertID
Algorithm Agility & \url{https://www.rfc-editor.org/rfc/rfc5035} \\
RFC 4055 & Additional Algorithms and Identifiers for RSA Cryptography
with SHA-256, SHA-384, and SHA-512 &
\url{https://www.rfc-editor.org/rfc/rfc4055} \\
ETSI EN 319 142-1 & PAdES: Building blocks and PAdES baseline signatures
& \url{https://www.etsi.org/deliver/etsi_en/319100_319199/31914201/} \\
ETSI EN 319 142-2 & PAdES: Additional PAdES signatures profiles &
\url{https://www.etsi.org/deliver/etsi_en/319100_319199/31914202/} \\
ETSI EN 319 122-1 & CAdES: Building blocks and CAdES baseline signatures
& \url{https://www.etsi.org/deliver/etsi_en/319100_319199/31912201/} \\
ETSI TS 119 102-1 & Electronic Signatures and Infrastructures:
Procedures for Creation and Validation &
\url{https://www.etsi.org/deliver/etsi_ts/119100_119199/11910201/} \\
ETSI TS 119 102-2 & Electronic Signatures and Infrastructures: Signature
Validation Report &
\url{https://www.etsi.org/deliver/etsi_ts/119100_119199/11910202/} \\
EU Regulation 910/2014 & eIDAS: Electronic Identification,
Authentication, and Trust Services &
\url{https://eur-lex.europa.eu/eli/reg/2014/910/oj} \\
EU Regulation 2024/1183 & eIDAS 2.0: Amending Regulation 910/2014
(European Digital Identity Framework) &
\url{https://eur-lex.europa.eu/eli/reg/2024/1183/oj} \\
Adobe AATL & Adobe Approved Trust List --- Technical Requirements &
\url{https://helpx.adobe.com/acrobat/kb/approved-trust-list1.html} \\
\end{longtable}

The RFCs are freely available and are the definitive reference for CMS,
X.509, OCSP, and TSP. The ETSI standards are available free of charge
from ETSI's website after registration. ISO 32000-2 is a paid standard,
but much of the PDF signature specification is reproduced (with
commentary) in this book and in Adobe's freely published PDF Reference
documents.

\backmatter
\end{document}
